\documentclass[10pt,twoside,openany]{book}
\usepackage{fontspec}
\usepackage[french]{babel}
% Format petit in-octavo / carnet (155 x 225 mm), proche du livre de poche et homothétique au carnet autographe (ratio ~0,69)
\usepackage[paperwidth=155mm,paperheight=225mm,top=11mm,bottom=11mm,inner=12mm,outer=10mm,headheight=13.5pt,headsep=3.5mm,footskip=5mm]{geometry}
\usepackage{graphicx}
\usepackage{xcolor}
\usepackage{tikz}
\usetikzlibrary{arrows.meta,decorations.pathmorphing,positioning,calc}
\usepackage{amsmath,amssymb,amsfonts}
\usepackage{array,booktabs,longtable}
\usepackage{fancyhdr}
\usepackage{titlesec}
\usepackage{hyperref}

% Palette sobre sur fond blanc pur (modèle de l'édition des Thesmophories)
\definecolor{darkgray}{HTML}{3A3A3A}
\definecolor{bordeaux}{HTML}{5A1818}

% Répertoire des fac-similés :
\newcommand{\facsimiledir}{../scans_compresses}
% \renewcommand{\facsimiledir}{../scans_haute_qualite}

% Macros d'en-tête courant indiquant le mois et l'année en cours ainsi que le folio :
\newcommand{\currentmonthyear}{Prolégomènes}
\newcommand{\currentfoliolabel}{Prolégomènes}

% Pages blanches de garde sans en-tête ni foliotage lors des sauts vers une belle page (impaire) ou une page paire :
\makeatletter
\def\cleardoublepage{\clearpage\if@twoside \ifodd\c@page\else
  \hbox{}\thispagestyle{empty}\newpage\if@twocolumn\hbox{}\newpage\fi\fi\fi}
\makeatother

% Saut automatique vers une page paire (gauche / verso) pour garantir le vis-à-vis strict :
\newcommand{\cleardoubleevenpage}{%
  \clearpage
  \ifodd\value{page}%
    \hbox{}\thispagestyle{empty}\newpage
  \fi
}

\pagestyle{fancy}
\fancyhf{}
\fancyhead[LE]{\scriptsize\textsc{Hippolyte Combes — Livre-Journal} \hfill \textbf{\currentmonthyear}}
\fancyhead[RO]{\scriptsize\textbf{\currentmonthyear} \hfill \textsc{\currentfoliolabel}}
\fancyfoot[LE,RO]{\scriptsize\thepage}
\renewcommand{\headrulewidth}{0.3pt}

\titleformat{\chapter}[display]
  {\normalfont\normalsize\bfseries}
  {\textsc{\chaptertitlename\ \Roman{chapter}}}
  {0.8mm}
  {\large\color{black}}
\titlespacing*{\chapter}{0pt}{-4mm}{2.5mm}
\titlespacing*{\section}{0pt}{2mm}{1mm}

\hypersetup{
  colorlinks=true,
  linkcolor=black,
  urlcolor=bordeaux,
  pdftitle={Hippolyte Combes (1876-1970) - Livre-Journal de campagne (25 juillet 1916 - 30 avril 1917) - Édition philologique et diplomatique},
  pdfauthor={Louis Hauseux (dir.)}
}

\setlength{\parindent}{1.0em}
\setlength{\parskip}{0.25em}
\setlength{\emergencystretch}{1.5em}

\begin{document}
\pagenumbering{arabic}

% =============================================================================
% 1. PREMIÈRE DE COUVERTURE (PAGE 1 DU PDF)
% =============================================================================
\thispagestyle{empty}
\begin{center}
  \setlength{\fboxsep}{4.5mm}
  \setlength{\fboxrule}{0.7pt}
  \fcolorbox{black}{white}{%
    \parbox[t][188mm][s]{120mm}{%
      \begin{center}
        \vspace{2mm}
        {\footnotesize\textsc{Archives familiales de Millau, Toulouse \& du Front britannique}}\\[1.5mm]
        \rule{86mm}{0.4pt}\\[4mm]
        {\Large\textbf{HIPPOLYTE COMBES}}\\[1.5mm]
        {\footnotesize\textsc{(Millau, 14 août 1876 — Saint-Chély-d'Apcher, 2 septembre 1970)}}\\[5mm]
        {\LARGE\textbf{LIVRE - JOURNAL}}\\[2mm]
        {\large\textbf{DE CAMPAGNE (1916–1917)}}\\[2.5mm]
        {\normalsize\textit{Du 25 juillet 1916 au 30 avril 1917}}\\[2mm]
        {\footnotesize Interprète militaire auprès de l'Armée britannique\\
        (\textit{Army Service Corps}, puis services d’approvisionnement britanniques)}\\[4.5mm]
        \setlength{\fboxsep}{1.2mm}
        \setlength{\fboxrule}{0.4pt}
        \fcolorbox{black}{white}{%
          \includegraphics[width=58mm,height=50mm,keepaspectratio]{illustrations/vignette_couverture.jpg}%
        }\\[4.5mm]
        {\normalsize\textsc{Édition philologique, diplomatique et critique en vis-à-vis}}\\[2mm]
        {\scriptsize Transcription diplomatique intégrale des 93 folios numérotée de 5 en 5 lignes,\\
        apparat critique marginal aligné, apparat des \textit{loci paralleli},\\
        commentaire historique et fac-similés en regard}\par
        \vfill
        \rule{70mm}{0.4pt}\\[2mm]
        {\footnotesize\textsc{Édition établie, introduite et annotée}}\\
        {\scriptsize d'après le manuscrit autographe et les archives familiales Combes-Calmels de Millau et Montauban}\\[1.5mm]
        {\footnotesize\textbf{MMXXVI}}
        \vspace{1mm}
      \end{center}
    }%
  }
\end{center}
\newpage

% Page blanche de garde (Page 2)
\hbox{}\thispagestyle{empty}\newpage

% =============================================================================
% 2. FAUX-TITRE & PAGE DE TITRE INTÉRIEURE (AVEC EMBLÈME FRAKTUR \mathfrak{Urtext})
% =============================================================================
\thispagestyle{empty}
\begin{center}
  \vspace*{10mm}
  {\large\textsc{Marie Lucien Hippolyte Combes}}\\
  {\footnotesize (Millau, 14 août 1876 — Saint-Chély-d'Apcher, 2 septembre 1970)}\\[10mm]
  {\LARGE\textbf{LIVRE-JOURNAL DE GUERRE}}\\[3mm]
  {\normalsize\textit{25 juillet 1916 — 30 avril 1917}}\\[4mm]
  {\footnotesize Des Flandres à l'Artois et au front de la Somme\\
  auprès de l'\textit{Army Service Corps} (18\textsuperscript{e} Division, XIII\textsuperscript{e} \& V\textsuperscript{e} Corps)}\\[14mm]
  {\Huge $\mathfrak{Urtext}$}\\[14mm]
  \rule{45mm}{0.4pt}\\[6mm]
  {\footnotesize\textbf{Édition philologique, diplomatique et critique intégrale}}\\[1.5mm]
  {\scriptsize accompagnée d'arbres généalogiques, de cartes vectorielles d'itinéraires,\\
  des lettres de guerre inédites (1916–1918), de l'état signalétique des services\\
  et d'un double index des personnes et des lieux}\par
  \vfill
  {\footnotesize\textsc{Collection des Sources et Manuscrits Contemporains}}\\
  {\scriptsize Millau — Toulouse — Montauban — Paris}
  \vspace*{6mm}
\end{center}
\newpage

{\renewcommand{\currentmonthyear}{Table des matières}
\renewcommand{\currentfoliolabel}{Table des matières}
\small\tableofcontents}
\cleardoublepage

% =============================================================================
% 3. PROLÉGOMÈNES SCIENTIFIQUES ET HISTORIQUES
% =============================================================================
\renewcommand{\currentmonthyear}{Prolégomènes}
\renewcommand{\currentfoliolabel}{Prolégomènes}
\chapter*{Prolégomènes : Hippolyte Combes (1876–1970) et son \textit{Livre-Journal} de la Grande Guerre}
\addcontentsline{toc}{chapter}{Prolégomènes : Hippolyte Combes (1876–1970) et son \textit{Livre-Journal} de la Grande Guerre}
\renewcommand{\thesection}{\Roman{section}}
\setcounter{section}{0}
{\small

\section{Destin d'un universitaire rouergat et toulousain auprès de l'Armée britannique (1876–1970)}

L'acte de naissance n\textsuperscript{o}~395 du 16 août 1876 situe la naissance de \textbf{Marie Lucien Hippolyte Combes} à \textbf{Millau}, rue des Barris, le \textbf{14 août 1876}, fils de \textbf{François Augustin Hippolyte Combes} (1825–1899), avocat, et de \textbf{Marie Anaïs Alric} (1838–1907). La fiche matricule porte la date du 22 août 1876 : les deux témoignages sont signalés sans harmonisation silencieuse. Les pièces d'état civil et monographies familiales retracent ses études à l'Institution Sainte-Marie-des-Champs (Le Caousou) et à la Faculté des lettres de Toulouse (47, rue du Taur), son séjour de quatre années en Angleterre (Nottingham) et sa carrière de professeur d'anglais aux collèges de Saint-Yrieix et de Nogent-le-Rotrou (17, rue Saint-Hilaire), puis au lycée Ingres de Montauban jusqu'à sa retraite à Millau et son décès le \textbf{2 septembre 1970} à l'\textbf{Hôtel du Centre} à \textbf{Saint-Chély-d'Apcher} (Lozère), avant son inhumation au cimetière de Millau.

\begin{figure}[htbp]
  \centering
  \begin{minipage}[b]{0.46\textwidth}
    \centering
    \includegraphics[height=42mm,keepaspectratio]{illustrations/hippolyte_1897.jpg}
    \par\vspace{1mm}{\scriptsize\textit{Fig. 1.} Hippolyte Combes en 1897, à l'âge de 21 ans.}
  \end{minipage}\hfill
  \begin{minipage}[b]{0.48\textwidth}
    \centering
    \includegraphics[height=42mm,keepaspectratio]{illustrations/hippolyte_interprete_1.jpg}
    \par\vspace{1mm}{\scriptsize\textit{Fig. 2.} Hippolyte Combes en uniforme d'interprète auprès de la B.E.F. (1916).}
  \end{minipage}
\end{figure}

Exempté dans sa jeunesse pour « faiblesse générale » (\textbf{classe 1896, bureau de recrutement de Toulouse, matricule n\textsuperscript{o}~911}), Hippolyte Combes est classé « Service armé » par le conseil de révision d'Eure-et-Loir le \textbf{18 décembre 1914}. Appelé à l'activité le \textbf{20 mars 1915} et incorporé le \textbf{21 mars 1915} à la 39\textsuperscript{e} batterie du \textbf{10\textsuperscript{e} régiment d'artillerie à pied} à Toulon (Saint-Jean-du-Var), il est détaché dès le \textbf{4 mai 1915} comme interprète militaire à la \textbf{Mission militaire française attachée à l'Armée britannique} (administrativement rattaché au \textbf{19\textsuperscript{e} escadron du Train des équipages militaires} le 15 octobre 1915) :

\begin{itemize}
  \item \textbf{Depuis août 1915 : service dans la 152\textsuperscript{e} Compagnie A.S.C. de la 18\textsuperscript{e} Division britannique.} Dans la copie autographe de sa lettre du 21 février 1916 (Annexe~I, doc.~1, p.~\pageref{annexe:doc1}), Hippolyte Combes écrit avoir été rattaché à l'été 1915 à la \textbf{18\textsuperscript{e} Division britannique} (\textit{18th Eastern Division}, commandée par le major-général Sir Ivor Maxse) qui se formait autour de Flesselles, avoir servi quinze jours au secrétariat de la Mission française (sous l'officier interprète divisionnaire M.~Pizzi), puis être affecté depuis le mois d'août 1915 à une compagnie d'\textit{Army Service Corps}. Le carnet précise le numéro de cette unité, la \textbf{152nd Company A.S.C.}, rattachée au \textit{18th Divisional Train} et chargée du ravitaillement de la \textit{54th Brigade} (f\textsuperscript{o}~33~r\textsuperscript{o}, p.~\pageref{fol:66}), relate sa démarche au \textit{Train Office} auprès du \textbf{lieutenant-colonel D.~C.~E. Grose} (f\textsuperscript{o}~45~r\textsuperscript{o}, p.~\pageref{fol:90}), et mentionne le 9 janvier 1917 une visite amicale à la \textbf{150th Company A.S.C., 18th Division}, près d'Hédauville (f\textsuperscript{o}~60~v\textsuperscript{o}, p.~\pageref{fol:121}).
  \item \textbf{Du Train divisionnaire aux services d'achats et de réquisitions (\textit{Foraging} / \textit{Army Purchase Board}, juillet 1916 – avril 1917).} Le carnet suit les déplacements d'Hippolyte Combes en Flandre intérieure (Lynde, Sercus, Renescure, Hazebrouck, Saint-Jans-Cappel, Bailleul) et dans la vallée de la Lys (La Croix-du-Bac, Bac-Saint-Maur, Sailly-sur-la-Lys, Armentières, Fleurbaix, Estaires), puis son transfert ferroviaire du \textbf{25 août 1916} vers Marquay par Saint-Pol-sur-Ternoise (f\textsuperscript{o}~13~v\textsuperscript{o}–14~r\textsuperscript{o}, p.~\pageref{fol:27}–\pageref{fol:28}) et sa remontée à cheval vers la Somme (Oppy/Beaudricourt, Lucheux, Raincheval, Acheux, Forceville). Après sa première permission à Paris et Toulouse (25 septembre – 6 octobre 1916) et son séjour d'octobre autour de Ribeaucourt et d'Albert, Hippolyte Combes est affecté le 2 novembre 1916 au service des fourrages et réquisitions à \textbf{Acheux-en-Amiénois} puis à \textbf{Hérissart}, d'abord avec le lieutenant Curtis, puis avec le capitaine Williams (décembre 1916 – mars 1917). Le 11 mars 1917, il est maintenu à l'\textit{Army Purchase Board} sur la demande expresse du \textbf{major Levy} (f\textsuperscript{o}~83~r\textsuperscript{o}, p.~\pageref{fol:166}) et se lie d'amitié en avril 1917 avec le sous-lieutenant poète \textbf{Gilbert}, originaire de Jersey (f\textsuperscript{o}~89~v\textsuperscript{o}, p.~\pageref{fol:179}).
  \item \textbf{Mariage, fin de campagne et carrière d'après-guerre.} Hippolyte Combes épouse \textbf{Marie Marguerite Calmels} (1892–1974) à Millau le \textbf{25 octobre 1917}. Affecté à l'École militaire préparatoire des Andelys comme professeur d'anglais le 24 novembre 1918 (nommé caporal le 4 décembre 1918), il est démobilisé le 7 février 1919 et nommé en 1920 professeur au lycée Ingres de Montauban.
\end{itemize}

\begin{center}
\begin{tikzpicture}[
  scale=0.92,
  every node/.style={font=\fontsize{6.0pt}{7.0pt}\selectfont},
  city/.style={circle, fill=black, inner sep=1.1pt},
  hub/.style={circle, draw=bordeaux, fill=white, line width=0.7pt, inner sep=1.5pt},
  route/.style={draw=bordeaux, line width=0.65pt, ->, >=stealth},
  perm/.style={draw=black!65, line width=0.55pt, dashed, ->, >=stealth}
]
% Cadre géographique schématique de la France
\draw[black!40, line width=0.35pt] (-58mm,-36mm) rectangle (62mm,36mm);
% Silhouette schématique des côtes et frontières de la France
\draw[black!30, line width=0.5pt, fill=black!2]
  (-8mm,33mm) -- (4mm,31mm) -- (18mm,24mm) -- (28mm,16mm) -- (26mm,2mm)
  -- (30mm,-10mm) -- (28mm,-22mm) -- (14mm,-29mm) -- (-4mm,-33mm) -- (-22mm,-31mm)
  -- (-28mm,-16mm) -- (-34mm,2mm) -- (-44mm,12mm) -- (-26mm,18mm) -- (-16mm,27mm) -- cycle;

% Nord / Pas-de-Calais / Somme
\node[hub] (flandre) at (-4mm, 30mm) {};
\node[above right] at (flandre) {\textbf{1. Flandre \& Lys} (Lynde, Bailleul, Steenwerck, juill.–août 1916)};

\node[hub] (ternois) at (-8mm, 24mm) {};
\node[left] at (ternois) {\textbf{2. Artois / Ternois} (Saint-Pol, Marquay, août–sept. 1916)};

\node[hub] (somme) at (-5mm, 19mm) {};
\node[right] at (somme) {\textbf{3. Front de la Somme} (Albert, Acheux, Hérissart, sept. 1916 – avr. 1917)};

\node[city] (amiens) at (-9mm, 16mm) {};
\node[left] at (amiens) {Amiens};

% Paris / Neuilly / Nogent
\node[hub] (paris) at (-7mm, 8mm) {};
\node[right] at (paris) {\textbf{Paris \& Neuilly-sur-Seine} (permissions)};
\node[city] (nogent) at (-18mm, 5mm) {};
\node[left] at (nogent) {Nogent-le-Rotrou};
\node[city] (andelys) at (-14mm, 12mm) {};
\node[left] at (andelys) {Les Andelys (1918)};

% Massif central
\node[city] (clermont) at (3mm, -8mm) {};
\node[right] at (clermont) {Clermont-Ferrand / Saint-Flour (retour févr. 1917)};

% Midi : Toulouse, Montauban, Millau, Béziers
\node[city] (montauban) at (-12mm, -21mm) {};
\node[left] at (montauban) {Montauban};
\node[hub] (toulouse) at (-11mm, -26mm) {};
\node[left] at (toulouse) {\textbf{Toulouse} (sept.–oct. 1916 \& janv.–févr. 1917)};
\node[hub] (millau) at (2mm, -22mm) {};
\node[right] at (millau) {\textbf{Millau} (berceau familial, févr. \& oct. 1917)};
\node[city] (beziers) at (5mm, -29mm) {};
\node[right] at (beziers) {Carcassonne \& Béziers (3 févr. 1917)};

% Flèches d'itinéraires
\draw[route] (flandre) -- (ternois);
\draw[route] (ternois) -- (somme);
\draw[perm] (somme) -- (amiens) -- (paris);
\draw[perm] (paris) to[bend right=12] (toulouse);
\draw[perm] (toulouse) to[bend right=15] (beziers) to[bend right=12] (millau);
\draw[perm] (millau) -- (clermont) -- (paris);

% Rose des vents
\draw[->, line width=0.5pt] (53mm, 24mm) -- (53mm, 31mm) node[above]{\textbf{N}};
\end{tikzpicture}
\par\vspace{0.8mm}{\scriptsize\textit{Carte générale.} Les théâtres d'opérations (Flandre, Artois, Somme) et les grands déplacements de permission d'Hippolyte Combes (1916–1917).}
\end{center}


\section{Présentation des protagonistes et arbres généalogiques}

Afin d'éclairer d'emblée les nombreuses mentions familiales, mondaines, militaires et universitaires qui parcourent le \textit{Livre-Journal}, les trois schémas généalogiques et prosopographiques ci-dessous présentent successivement : (I)~l'ascendance, la fratrie et la descendance d'Hippolyte Combes ; (II)~la famille de sa fiancée Marguerite Calmels et ses alliances à Millau, Toulouse et Neuilly-sur-Seine (branches Calmels, Cartairade, Crouzet, Vincens, Delcamps, Layrals et Faure) ; (III)~l'organigramme des officiers britanniques et français, collègues interprètes et universitaires côtoyés par Hippolyte Combes de 1915 à 1918.


\begin{center}
\begin{tikzpicture}[
  every node/.style={font=\fontsize{6.2pt}{7.2pt}\selectfont, align=center},
  box/.style={draw=black!80, line width=0.4pt, fill=white, rounded corners=1.2pt, inner sep=2.2pt},
  hero/.style={draw=bordeaux, line width=0.8pt, fill=white, rounded corners=1.2pt, inner sep=2.6pt},
  edge/.style={draw=black!75, line width=0.4pt}
]
% Grands-parents
\node[box, text width=34mm] (gp1) at (-42mm, 34mm) {\textbf{Jean-Pierre Combes}\\(1788–1875)\\\textit{Avoué licencié à Millau}};
\node[box, text width=34mm] (gm1) at (-4mm, 34mm) {\textbf{Marie-Julie-Rose Descuret}\\(1801–1876)\\\textit{Millau}};
\node[box, text width=30mm] (gp2) at (30mm, 34mm) {\textbf{Augustin Alric}\\\textit{Banquier à Millau}};
\node[box, text width=28mm] (gm2) at (62mm, 34mm) {\textbf{Émilie Noël}\\(\dag~1906)\\\textit{Millau}};

% Parents
\node[box, text width=46mm] (pere) at (-23mm, 18mm) {\textbf{François Augustin Hippolyte Combes}\\(1825–1899)\\\textit{Avocat à Millau puis Toulouse (47, rue du Taur)}};
\node[box, text width=38mm] (mere) at (35mm, 18mm) {\textbf{Marie Anaïs Alric}\\(1838–1907)\\\textit{Mariés en 1870}};
\node[box, text width=24mm] (oncle_alric) at (67mm, 18mm) {\textbf{Hippolyte Alric}\\\textit{Ingénieur ECP}};

\draw[edge] (gp1.south) -- ++(0,-2.5mm) -| (pere.north);
\draw[edge] (gm1.south) -- ++(0,-2.5mm) -| (pere.north);
\draw[edge] (gp2.south) -- ++(0,-2.5mm) -| (mere.north);
\draw[edge] (gm2.south) -- ++(0,-2.5mm) -| (mere.north);
\draw[edge] (gp2.south) -- ++(0,-2.5mm) -| (oncle_alric.north);
\draw[edge] (pere.east) -- node[above, font=\fontsize{5.5pt}{6.5pt}\selectfont]{$\times$} (mere.west);

% Fratrie Combes
\coordinate (frat_mid) at (6mm, 11mm);
\draw[edge] (6mm, 18mm) -- (frat_mid);

\node[box, text width=27mm] (s1) at (-46mm, 1mm) {\textbf{Françoise Combes}\\(1871–1942)\\\textit{Sœur Marie de la Croix}\\Carmel de Périgueux};
\node[box, text width=27mm] (s2) at (-16mm, 1mm) {\textbf{Joséphine Combes}\\(1873 – oct. 1914)\\\textit{\dag~à Nogent-le-Rotrou}};
\node[box, text width=26mm] (s3) at (13mm, 1mm) {\textbf{Maria Combes}\\(1875 – 19 nov. 1903)\\\textit{Religieuse, \dag~à Carthage}};
\node[hero, text width=44mm] (hc) at (52mm, 1mm) {\textbf{Marie Lucien Hippolyte Combes}\\(Millau, 14 août 1876 –\\Saint-Chély-d'Apcher, 2 sept. 1970)\\\textit{Professeur d'anglais, interprète B.E.F.}};

\draw[edge] (frat_mid) -| (s1.north);
\draw[edge] (frat_mid) -| (s2.north);
\draw[edge] (frat_mid) -| (s3.north);
\draw[edge] (frat_mid) -| (hc.north);

% Mariage Hippolyte & Marguerite et enfants
\node[hero, text width=42mm] (mc) at (-18mm, -15mm) {\textbf{Marie Marguerite Calmels}\\(Millau, 15 avril 1892 – 19 mars 1974)\\\textit{« Marg\textsuperscript{te} / M\textsuperscript{te} / Gueïte »}};
\draw[edge, bordeaux, line width=0.6pt] (hc.south) -- ++(0,-3mm) -| node[pos=0.25, above, font=\fontsize{5.6pt}{6.5pt}\selectfont]{\textbf{Mariés à Millau le 25 oct. 1917}} (mc.north);

\node[box, text width=92mm] (enfants) at (22mm, -26mm) {\textbf{Enfants d'Hippolyte et Marguerite Combes :} \textbf{Valentine Combes} (née en 1918), \textbf{Geneviève Combes}, \textbf{Jean Combes} et \textbf{François Combes} (père de \textbf{Pierre Combes})};
\draw[edge] (17mm, -11mm) -- (enfants.north);
\end{tikzpicture}
\par\vspace{0.8mm}{\scriptsize\textit{Arbre généalogique I.} La famille Combes–Alric (Millau, Toulouse, Nogent-le-Rotrou et Montauban).}
\end{center}

\medskip

\begin{center}
\begin{tikzpicture}[
  every node/.style={font=\fontsize{6.0pt}{7.0pt}\selectfont, align=center},
  box/.style={draw=black!80, line width=0.4pt, fill=white, rounded corners=1.2pt, inner sep=2.0pt},
  hero/.style={draw=bordeaux, line width=0.75pt, fill=white, rounded corners=1.2pt, inner sep=2.2pt},
  edge/.style={draw=black!75, line width=0.4pt}
]
% Ligne supérieure : Crouzet-Calmels (Toulouse), Calmels-Cartairade (Millau), Layrals-Cartairade (Millau)
\node[box, text width=38mm] (crouzet) at (-43mm, 24mm) {\textbf{Émile Crouzet} (1842–1921)\\\textit{Agent de change à Toulouse}\\$\times$ \textbf{Élisabeth Calmels}\\(\textit{« Oncle Émile \& Tante »})};
\node[box, text width=42mm] (calmels) at (1mm, 24mm) {\textbf{D\textsuperscript{r} Gabriel Calmels} (1852–1921)\\\textit{Médecin à Millau (10, rue A.-Guibert)}\\$\times$ \textbf{Valentine Cartairade} (1858–1943)};
\node[box, text width=38mm] (layrals) at (45mm, 24mm) {\textbf{M\textsuperscript{e} Charles Layrals}\\\textit{Notaire à Millau}\\$\times$ \textbf{Thérèse Cartairade}\\(\textit{« l'oncle \& tante Thérèse »})};

\draw[edge, dashed] (crouzet.east) -- (calmels.west);
\draw[edge, dashed] (calmels.east) -- (layrals.west);

% Cousin Jean Crouzet et cousins Layrals
\node[box, text width=32mm] (jean_cr) at (-46mm, 8mm) {\textbf{Abbé Jean Crouzet}\\\textit{« Jean », prêtre à Toulouse}\\(futur curé de la Dalbade)};
\draw[edge] (crouzet.south) -- (jean_cr.north);

\node[box, text width=32mm] (fils_layrals) at (48mm, 8mm) {\textbf{Louis Layrals} (\dag~1961, notaire)\\\textbf{Jean Layrals} (\dag~1947, magistrat)};
\draw[edge] (layrals.south) -- (fils_layrals.north);

% Enfants Calmels
\coordinate (cal_mid) at (1mm, 15.5mm);
\draw[edge] (calmels.south) -- (cal_mid);

\node[box, text width=32mm] (mad) at (-38mm, -5mm) {\textbf{Madeleine Calmels}\\(1882–1961, \textit{« Mad\textsuperscript{ne} »})\\$\times$ \textbf{D\textsuperscript{r} Fernand Vincens}\\\textit{Enfants :} \textbf{Germaine} \& \textbf{Robert}};
\node[box, text width=28mm] (mimi) at (-5mm, -5mm) {\textbf{Anne-Marie Calmels}\\(1888–1975, \textit{« Mimi »})\\$\times$ \textbf{M\textsuperscript{e} Joseph Delcamps}\\\textit{Avoué à Millau}};
\node[box, text width=20mm] (pierre) at (22mm, -5mm) {\textbf{Pierre Calmels}\\\textit{Frère de}\\Marguerite\\(Millau)};
\node[hero, text width=33mm] (marg) at (51mm, -5mm) {\textbf{Marguerite Calmels}\\(1892–1974)\\\textit{« Marg\textsuperscript{te} / M\textsuperscript{te} »}\\$\times$ \textbf{Hippolyte Combes}};

\draw[edge] (cal_mid) -| (mad.north);
\draw[edge] (cal_mid) -| (mimi.north);
\draw[edge] (cal_mid) -| (pierre.north);
\draw[edge] (cal_mid) -| (marg.north);

% Branche Faure à Neuilly-sur-Seine
\node[box, text width=122mm] (faure) at (1mm, -18.5mm) {\textbf{Branche cousine de Neuilly-sur-Seine (16, rue d'Armenonville) :} \textbf{M\textsuperscript{me} Faure} et ses cinq enfants : \textbf{Marguerite}, \textbf{Germaine}, \textbf{Marcel Faure} (en permission en janv. 1917), \textbf{Henri Faure} (aspirant) et \textbf{André Faure} (candidat à Polytechnique)};
\end{tikzpicture}
\par\vspace{0.8mm}{\scriptsize\textit{Arbre généalogique II.} La famille Calmels–Cartairade et ses alliances à Millau, Toulouse et Neuilly-sur-Seine.}
\end{center}

\medskip

\begin{center}
\begin{tikzpicture}[
  every node/.style={font=\fontsize{5.9pt}{6.9pt}\selectfont, align=center},
  box/.style={draw=black!80, line width=0.4pt, fill=white, rounded corners=1.2pt, inner sep=2.0pt},
  hero/.style={draw=bordeaux, line width=0.75pt, fill=white, rounded corners=1.2pt, inner sep=2.2pt},
  edge/.style={draw=black!75, line width=0.4pt, ->, >=stealth}
]
\node[box, text width=56mm] (bef) at (-32mm, 22mm) {\textbf{18th (Eastern) Division \& Fifth Army (B.E.F.)}\\\textbf{Maj.-Gen. Sir Ivor Maxse} (cmdt la 18\textsuperscript{e} Division)\\\textbf{Lt.-Col. D.~C.~E. Grose} (cmdt le \textit{18th Divisional Train})};
\node[box, text width=56mm] (mmf) at (33mm, 22mm) {\textbf{Mission militaire française (B.E.F.)}\\\textbf{Lt Pizzi} (18\textsuperscript{e} Div.), \textbf{S/Lt Ferran} (liaison),\\\textbf{Cap. Huguet} (XIII\textsuperscript{e} Corps), \textbf{Cap. Bresson} (V\textsuperscript{e} Armée)};

\node[hero, text width=68mm] (hc) at (0mm, 9mm) {\textbf{Interprète Hippolyte Combes (matr. 911, classe 1896)}\\Détaché auprès de l'Armée britannique (4 mai 1915 – 11 nov. 1918)};

\draw[edge] (bef.south) -- (hc.north west);
\draw[edge] (mmf.south) -- (hc.north east);

\node[box, text width=40mm] (asc) at (-43mm, -6mm) {\textbf{152nd \& 150th Cos. A.S.C.}\\\textit{(Juillet – Octobre 1916)}\\\textbf{Cap. E.~M. West}, \textbf{Cap. E.~M. Ford},\\\textbf{Lt Curtis} (\textit{C\textsuperscript{tis}}), \textbf{Cap. Panton},\\\textbf{Lt Walsh}, \textbf{Lt Michel}, \textbf{Lt First},\\\textbf{Maj. Hughes}, \textbf{Lt Morris}};
\node[box, text width=41mm] (forage) at (0mm, -6mm) {\textbf{Foraging \& Purchase Board}\\\textit{(Acheux \& Hérissart, nov. 1916 – avr. 1917)}\\\textbf{Major Levy} (\textit{Purchase Board}),\\\textbf{Cap. Williams} (déc. 1916 – mars 1917),\\\textbf{2nd Lt Gilbert} (de Jersey, avr. 1917),\\\textbf{Lt Croughton}, \textbf{Lt Godard}, \textbf{Lt Wentworth}};
\node[box, text width=38mm] (coll) at (43mm, -6mm) {\textbf{Collègues \& Universitaires}\\\textit{Interprètes :} \textbf{Maurin}, \textbf{Lépin},\\\textbf{Schumacher}, \textbf{Raucoule}, \textbf{Valet}\\(\textit{2nd Austr. Div.}), \textbf{Brammas}\\\textit{Lettres :} \textbf{E.~Zyromski} (Toulouse),\\\textbf{Fortuné \& L.-F. Flutre} (ENS)};

\draw[edge] (hc.south) -- (-43mm, 1.5mm);
\draw[edge] (hc.south) -- (0mm, 1.5mm);
\draw[edge] (hc.south) -- (43mm, 1.5mm);
\end{tikzpicture}
\par\vspace{0.8mm}{\scriptsize\textit{Organigramme III.} Cadre militaire britannique et français et réseau intellectuel d'Hippolyte Combes (1916–1917).}
\end{center}


\section{Histoire matérielle, codicologie et paléographie du \textit{Livre-Journal}}

Le manuscrit autographe d'Hippolyte Combes se compose du plat supérieur de la couverture (revêtu d'un papier fort rouge à motif de losanges gaufrés sur lequel est collée une étiquette rectangulaire de papier bleu portant à l'encre noire l'intitulé autographe en cinq lignes : \textit{\underline{Livre - Journal} / — / Du 25 juillet 1916 / au 30 avril 1917 / —}) et d'une séquence de \textbf{93 folios} (\texttt{Folio 1 recto} à \texttt{Folio 93 recto}).

\begin{enumerate}
  \item \textbf{Ductus, soulignements et système abréviatif d'Hippolyte Combes.} L'écriture d'Hippolyte Combes est une cursive latine penchée, rapide et déliée, riche en abréviations par suspension et par lettres suscrites : \textit{ds} (\textit{dans}), \textit{pr} (\textit{pour}), \textit{s/} (\textit{sur}), \textit{j'à / j'au / j'aux} ou \textit{ju} (\textit{jusqu'à, jusqu'au, jusqu'aux}), \textit{lgt / lgs / longt} (\textit{longtemps}), \textit{tt / tte / ts} (\textit{tout, toute, tous}), \textit{tjs / tjours} (\textit{toujours}), \textit{qq} (\textit{quelques}), \textit{gd / gde / gds} (\textit{grand, grande, grands}), \textit{bp / bcp} (\textit{beaucoup}), \textit{cme} (\textit{comme}), \textit{chz} (\textit{chez}), ainsi que les finales adverbiales ou nominales suscrites (\textit{ext\textsuperscript{t}, merv\textsuperscript{t}, définit\textsuperscript{t}, admirab\textsuperscript{t}, appart\textsuperscript{t}, cantonnem\textsuperscript{t}, changem\textsuperscript{t}, propri\textsuperscript{re}, cath\textsuperscript{le}, c\textsuperscript{tes} p\textsuperscript{les}}). Conformément à l'usage diplomatique, la présente édition restitue fidèlement :
  \begin{itemize}
    \item les \textbf{soulignements autographes} (\texttt{\textbackslash underline\{...\}}) : dates en tête d'entrée (\textit{\underline{25 Ma.}}, \textit{\underline{26 Mer.}}, etc.), rappels du mois et de l'année dans l'angle supérieur droit (\textit{\underline{Juillet 1916}}, parfois légèrement surélevés au-dessus de la première ligne d'écriture comme au f\textsuperscript{o}~1~r\textsuperscript{o}), et sous-titres topographiques des blocs de dépenses (\textit{\underline{A Renescure}}, \textit{\underline{A Hazebrouck}}, \textit{\underline{En caisse}}) ;
    \item les \textbf{traits horizontaux de séparation} tracés par Hippolyte Combes à travers la page, soit entre deux journées successives réunies sur un même folio, soit au-dessus du relevé comptable ou d'une note infrapaginale en pied de page ;
    \item les \textbf{lettres suscrites} (\textit{8\textsuperscript{H}}, \textit{7\textsuperscript{h} 1/2}, \textit{Marg\textsuperscript{te}}, \textit{lieut\textsuperscript{t}}, \textit{R.~Fus\textsuperscript{rs}}, \textit{Div\textsuperscript{on}}) et les crochets ou accolades comptables (\texttt{[ ... ]}).
  \end{itemize}
  \item \textbf{Particularités de la séquence des folios.} Quelques inversions ponctuelles lors de la prise de vue photographique du carnet sont scrupuleusement signalées dans le commentaire critique et l'apparat des \textit{loci paralleli} tout en respectant l'ordre matériel de la séquence : le \textbf{Folio 17 verso} (4–5 septembre 1916, p.~\pageref{fol:35}) précède la fin du compte du 3 septembre inscrite au \textbf{Folio 18 recto} (p.~\pageref{fol:36}) ; le \textbf{Folio 31 verso} (jeudi 5 octobre 1916, p.~\pageref{fol:63}) et le \textbf{Folio 32 recto} (nuit du 4 au 5 octobre 1916, p.~\pageref{fol:64}), ainsi que le \textbf{Folio 56 recto} (22–24 décembre 1916, p.~\pageref{fol:112}) et le \textbf{Folio 56 verso} (19–21 décembre 1916, p.~\pageref{fol:113}), se présentent dans l'ordre inverse de la chronologie ; enfin, le \textbf{Folio 45 verso} (p.~\pageref{fol:91}) et le \textbf{Folio 46 verso} (p.~\pageref{fol:93}) correspondent à deux prises de vue successives du même feuillet original des 2 et 3 novembre 1916.
  \item \textbf{Caractère éditorial et arbitraire de la division en sept chapitres.} Il importe de souligner que \textbf{le manuscrit original d'Hippolyte Combes ne comporte aucune division en chapitres ni aucun titre de section} : le \textit{Livre-Journal} se présente matériellement comme une suite continue de notes quotidiennes courant sans interruption du 25 juillet 1916 au 30 avril 1917. \textbf{La répartition en sept chapitres adoptée dans la présente édition est donc purement éditoriale et arbitraire} : elle a été conçue pour offrir au lecteur des repères chronologiques et géographiques clairs en suivant pas à pas \textbf{les grands déplacements d'Hippolyte Combes} entre les Flandres, l'Artois, le front de la Somme et ses deux permissions à Paris, Toulouse et Millau.
\end{enumerate}

\section{Principes de l'édition philologique en vis-à-vis et des trois registres d'apparat}

Inspirée des principes de l'édition philologique classique (sur le modèle de l'édition savante des \textit{Thesmophories} d'Aristophane), la présente édition adopte un format maniable (\texttt{155 $\times$ 225 mm}, proche du livre de poche et homothétique au carnet original), une pagination arabe continue depuis la première page et met en regard sur chaque double page :
\begin{itemize}
  \item \textbf{En-tête courant :} L'en-tête indique en permanence, outre le folio (\texttt{Folio X recto} ou \texttt{Folio X verso}), \textbf{le mois et l'année en cours} (\textit{Juillet 1916} à \textit{Avril 1917}) afin de permettre au lecteur de se situer immédiatement dans la chronologie de la campagne.
  \item \textbf{Sur la page de gauche (page paire / verso) :}
    \begin{enumerate}
      \item \textbf{La transcription diplomatique ligne à ligne}, respectant la disposition spatiale, les soulignements, les traits horizontaux de séparation, les lettres suscrites, l'orthographe, la ponctuation et les crochets comptables \texttt{[ ... ]} du manuscrit. Les lignes sont numérotées de cinq en cinq dans la marge gauche (\texttt{5, 10, 15, 20, 25}). Les suppressions et ratures autographes sont placées entre doubles crochets droits ($[\![ \dots ]\!]$).
      \item \textbf{L'apparat critique marginal}, composé en petits caractères dans la colonne de droite dans le prolongement horizontal exact de la ligne concernée : il développe les abréviations tachygraphiques, signale les biffures du manuscrit, identifie les emprunts lexicaux anglais ou latins et précise les particularités matérielles.
      \item \textbf{En bas de page, jusqu'à trois registres dépourvus d'intitulés répétitifs mais immédiatement distingués par des filets horizontaux et un décrochement typographique franc :}
        \begin{itemize}
          \item \textit{Premier registre (critique interne, présent uniquement lorsque des précisions codicologiques ou structurelles excèdent le cadre des notes marginales) :} composé au fer à gauche immédiatement sous le filet horizontal de pleine largeur (\texttt{0,5 pt}) et suivi d'un filet court de \texttt{22 mm}. Lorsqu'une page n'appelle aucune remarque critique interne en sus des notes marginales, ce registre est omis.
          \item \textit{Deuxième registre (apparat des \textbf{loci paralleli}, rédigé en latin) :} composé au fer à gauche entre le commentaire critique et le commentaire historique, et suivi d'un filet horizontal de \texttt{38 mm} (\texttt{0,4 pt}). Lorsqu'Hippolyte Combes mentionne une personne, une unité, un lieu ou un événement sur lequel d'autres passages du carnet (ou de la correspondance de l'Annexe~I) apportent un éclairage complémentaire, cet apparat latin renvoie simultanément \textbf{au folio du carnet} (\texttt{f\textsuperscript{o}~X~r\textsuperscript{o}/v\textsuperscript{o}}) \textbf{et à la page exacte du présent fichier PDF} (\texttt{p.~Y}).
          \item \textit{Troisième registre (commentaire historique) :} composé avec un retrait général de bloc de \texttt{5 mm} sur la gauche, il fournit de manière concise et dense l'élucidation historique, militaire, architecturale, prosopographique et littéraire du folio.
        \end{itemize}
    \end{enumerate}
  \item \textbf{Sur la page de droite (page impaire / recto) :} \textbf{le fac-similé photographique complet} du folio manuscrit correspondant en regard immédiat de sa transcription.
\end{itemize}
\par}

\renewcommand{\thesection}{\thechapter.\arabic{section}}

% =============================================================================
% CHAPITRE I (COUVERTURE (PLAT SUPÉRIEUR) A FOLIO 9 RECTO)
% =============================================================================
\cleardoubleevenpage
\renewcommand{\currentmonthyear}{Juillet 1916 – Avril 1917 – Août 1916}
\renewcommand{\currentfoliolabel}{Chapitre I}
\chapter{De la Flandre intérieure à la vallée de la Lys}

\noindent{\fontsize{7.8pt}{9.0pt}\selectfont\textit{Couverture et Folios 1 recto à 9 recto — Renescure, Lynde, Sercus, Hazebrouck, Saint-Jans-Cappel, Bailleul et La Croix-du-Bac (25 juillet – 12 août 1916)}}\par\vspace{1.5mm}
\noindent\rule{\textwidth}{0.5pt}\par\vspace{0.9mm}
\noindent\textbf{\scriptsize\textsc{Résumé analytique du Chapitre I (Couverture (plat supérieur) à Folio 9 recto)}}\par\vspace{0.5mm}
\noindent{\fontsize{6.9pt}{8.1pt}\selectfont Après la couverture autographe qui fixe les bornes chronologiques du carnet (25 juillet 1916 – 30 avril 1917), le premier chapitre suit Hippolyte Combes durant trois semaines en Flandre intérieure et dans la plaine de la Lys (f\textsuperscript{o}~1~r\textsuperscript{o} à f\textsuperscript{o}~9~r\textsuperscript{o}). Cantonné d'abord à Lynde du 25 au 28 juillet, il accompagne à pied le lieutenant Curtis au rassemblement de Sercus, assiste le lieutenant Michel pour loger trente-six officiers des \textit{Royal Fusiliers} et profite d'une matinée libre pour visiter les deux châteaux et l'église-halle de Renescure ainsi que la ville d'Hazebrouck. Le vendredi 28 juillet, une longue étape à cheval en tête du convoi le conduit au pied du Mont-Noir, à Saint-Jans-Cappel, d'où il rayonne jusqu'au 5 août pour les achats de fourrage vert et découvre avec admiration la Grand-Place, le beffroi et les dentellières de Bailleul. Le 5 août, la compagnie descend dans la vallée alluviale de la Lys et s'établit au hameau de La Croix-du-Bac (commune de Steenwerck), à proximité du point de ravitaillement quotidien (\textit{refilling point}) de Bac-Saint-Maur. Entre la chaleur orageuse, la poussière des convois et les promenades le long de la Lys jusqu'à l'écluse de Sailly-sur-la-Lys (où il s'entretient avec l'éclusier au sujet du général Maxse), la guerre se rappelle brutalement le 11 août par le bombardement allemand de la blanchisserie Barbry aux portes d'Armentières et par la vision d'un soldat britannique soumis au supplice disciplinaire de la roue.}\par\vspace{1.0mm}
\noindent\rule{\textwidth}{0.35pt}\par\vspace{1.4mm}

\noindent\textbf{\scriptsize\textsc{Introduction historique et philologique}}\par\vspace{0.6mm}
{\fontsize{7.1pt}{8.3pt}\selectfont Relevée du front de la Somme le 20 juillet 1916 après vingt jours de combats meurtriers au bois des Trônes et à Montauban-de-Picardie, la \textbf{18\textsuperscript{e} Division britannique} (\textit{18th Eastern Division}, commandée par le major-général Sir Ivor Maxse) est transférée par voie ferrée en Flandre intérieure pour se reconstituer avant de relever les troupes australiennes dans le secteur d'Armentières et de Fleurbaix (Capt. G.~H.~F. Nichols, \textit{The 18th Division in the Great War}, Londres, Blackwood, 1922, p.~81–85). Affecté depuis août 1915 à une compagnie du Train divisionnaire (la \textbf{152\textsuperscript{e} Compagnie de l'\textit{Army Service Corps}}, chargée du ravitaillement de la 54\textsuperscript{e} Brigade d'infanterie), Hippolyte Combes assure la liaison quotidienne entre les officiers britanniques d'approvisionnement (notamment le lieutenant Curtis), les municipalités et les exploitants agricoles flamands. Ses notes attentives sur le patrimoine architectural de la Flandre — l'église Saint-Érasme de Sercus, les châteaux de Zuthove et de Philippe de Comines à Renescure, le beffroi et les pignons Renaissance de Bailleul, ou l'église néo-gothique Notre-Dame-des-Sept-Douleurs de La Croix-du-Bac — revêtent une valeur documentaire d'autant plus précieuse que plusieurs de ces édifices furent anéantis au printemps 1918 lors de la bataille de la Lys.\par}
\vfill
\newpage
% --- Page de droite d'ouverture du chapitre : Carte vectorielle TikZ & Photographies d'archives ---
\noindent\rule{\textwidth}{0.5pt}\par\vspace{0.6mm}
\noindent\textbf{\fontsize{6.8pt}{7.8pt}\selectfont CARTE D'ITINÉRAIRE DU CHAPITRE I} \hfill \textbf{\fontsize{6.8pt}{7.8pt}\selectfont JUILLET 1916 – AVRIL 1917 – AOÛT 1916}\par\vspace{0.25mm}
\noindent{\fontsize{6.7pt}{7.8pt}\selectfont\textit{Déplacements, cantonnements et tournées d'Hippolyte Combes (Couverture (plat supérieur) à Folio 9 recto)}}\par\vspace{0.55mm}
\noindent\rule{\textwidth}{0.35pt}\par\vspace{1.2mm}
\begin{center}

\begin{tikzpicture}[
  scale=0.96,
  every node/.style={font=\fontsize{6.0pt}{7.0pt}\selectfont},
  loc/.style={circle, fill=black, inner sep=1.1pt},
  step/.style={circle, draw=bordeaux, fill=white, line width=0.75pt, inner sep=1.5pt},
  move/.style={draw=bordeaux, line width=0.75pt, ->, >=stealth},
  excur/.style={draw=black!65, line width=0.5pt, dashed, ->, >=stealth},
  river/.style={draw=black!35, line width=1.1pt}
]
\draw[black!50, line width=0.4pt] (-60mm,-42mm) rectangle (60mm,42mm);
\fill[black!3] (-60mm,-42mm) rectangle (60mm,42mm);

% Cours de la Lys et frontière belge
\draw[river] (-55mm,-30mm) .. controls (-20mm,-26mm) and (5mm,-22mm) .. (55mm,-14mm)
  node[pos=0.35, below, font=\fontsize{5.8pt}{6.5pt}\selectfont\itshape]{La Lys (canalisée)};
\draw[black!45, dash pattern=on 3pt off 2pt, line width=0.5pt]
  (-10mm,40mm) -- (5mm,22mm) -- (22mm,14mm) -- (45mm,-8mm)
  node[pos=0.25, right, font=\fontsize{5.6pt}{6.5pt}\selectfont]{\textit{Frontière franco-belge}};

% Ligne de front (août 1916)
\draw[bordeaux!80, line width=0.9pt, decorate, decoration={zigzag, segment length=4pt, amplitude=1pt}]
  (48mm,10mm) -- (46mm,-16mm) -- (42mm,-38mm);
\node[rotate=-82, font=\fontsize{5.6pt}{6.5pt}\selectfont, bordeaux] at (51mm,-14mm) {\textbf{Ligne de front (secteur d'Armentières–Fleurbaix)}};

% Localités Flandre intérieure (Ouest)
\node[loc] (renescure) at (-48mm, 14mm) {};
\node[below left] at (renescure) {\textbf{Renescure} (châteaux, 27 juill.)};

\node[step] (lynde) at (-36mm, 18mm) {};
\node[above left] at (lynde) {\textbf{1. Lynde} (25–28 juill.)};

\node[loc] (sercus) at (-32mm, 10mm) {};
\node[below] at (sercus) {\textbf{Sercus} (25 juill.)};

\node[loc] (hazebrouck) at (-22mm, 16mm) {};
\node[above] at (hazebrouck) {\textbf{Hazebrouck} (27–28 juill.)};

% Monts des Flandres
\node[step] (stjans) at (-2mm, 18mm) {};
\node[above left] at (stjans) {\textbf{2. Saint-Jans-Cappel} (28 juill. – 5 août)};

\node[loc] (montnoir) at (6mm, 22mm) {};
\node[right] at (montnoir) {Mont-Noir (29 juill.)};

\node[step] (bailleul) at (0mm, 6mm) {};
\node[left] at (bailleul) {\textbf{Bailleul} (30 juill. – 5 août)};

% Vallée de la Lys
\node[loc] (steenwerck) at (4mm, -8mm) {};
\node[left] at (steenwerck) {Steenwerck};

\node[step] (croixdubac) at (12mm, -15mm) {};
\node[above right] at (croixdubac) {\textbf{3. La Croix-du-Bac} (5–12 août)};

\node[loc] (bacstmaur) at (14mm, -23mm) {};
\node[below] at (bacstmaur) {\textbf{Bac-Saint-Maur} (\textit{refilling point})};

\node[loc] (sailly) at (6mm, -26mm) {};
\node[below left] at (sailly) {\textbf{Sailly-sur-la-Lys} (écluse, 10 août)};

\node[loc] (armentieres) at (34mm, -16mm) {};
\node[above] at (armentieres) {\textbf{Armentières} \& blanchisserie Barbry (11 août)};

% Flèches de déplacement principal et tournées
\draw[excur] (lynde) -- (sercus);
\draw[excur] (lynde) -- (renescure);
\draw[move] (lynde) -- (hazebrouck) -- (stjans);
\draw[excur] (stjans) -- (montnoir);
\draw[move] (stjans) -- (bailleul) -- (steenwerck) -- (croixdubac);
\draw[excur] (croixdubac) -- (bacstmaur);
\draw[excur] (bacstmaur) -- (sailly);
\draw[excur] (croixdubac) -- (armentieres);

% Échelle et Nord
\draw[line width=0.6pt] (-54mm,-36mm) -- (-34mm,-36mm);
\draw (-54mm,-37mm) -- (-54mm,-35mm) (-34mm,-37mm) -- (-34mm,-35mm);
\node[above] at (-44mm,-36mm) {0 \hfill 5 \hfill 10 km};
\draw[->, line width=0.6pt] (52mm, 28mm) -- (52mm, 37mm) node[above]{\textbf{N}};
\end{tikzpicture}

\end{center}
\vfill
\noindent\rule{\textwidth}{0.3pt}\par\vspace{1.2mm}
\noindent\begin{center}
  \begin{minipage}[b]{0.46\textwidth}
    \centering
    \setlength{\fboxsep}{0.7mm}\setlength{\fboxrule}{0.4pt}%
    \fcolorbox{gray!55}{white}{\includegraphics[width=0.94\linewidth,height=28mm,keepaspectratio]{illustrations/hippolyte_interprete_1.jpg}}
    \par\vspace{0.6mm}{\fontsize{5.8pt}{6.7pt}\selectfont Hippolyte Combes en uniforme d'interprète militaire attaché à l'Armée britannique (képi à bandeau et insignes au sphinx). Collection familiale Combes.}
  \end{minipage}\hfill
  \begin{minipage}[b]{0.46\textwidth}
    \centering
    \setlength{\fboxsep}{0.7mm}\setlength{\fboxrule}{0.4pt}%
    \fcolorbox{gray!55}{white}{\includegraphics[width=0.94\linewidth,height=28mm,keepaspectratio]{illustrations/toulon_artillerie_1915.jpg}}
    \par\vspace{0.6mm}{\fontsize{5.8pt}{6.7pt}\selectfont Hippolyte Combes au 10\textsuperscript{e} régiment d'artillerie à pied à Toulon (Saint-Jean-du-Var) au printemps 1915.}
  \end{minipage}
\end{center}
\vspace*{1mm}
\newpage

% =============================================================================
% VIS-A-VIS SPREAD : COUVERTURE (PLAT SUPÉRIEUR) (20251126_092853.jpg)
% =============================================================================
\cleardoubleevenpage
\renewcommand{\currentmonthyear}{Juillet 1916 – Avril 1917}
\renewcommand{\currentfoliolabel}{Couverture (plat supérieur)}
\phantomsection\label{fol:1}
% --- En-tête du folio (Page paire / gauche : Transcription, apparat & registres de bas de page) ---
\noindent\rule{\textwidth}{0.5pt}\par\vspace{0.55mm}
\noindent\textbf{\footnotesize COUVERTURE (PLAT SUPÉRIEUR)} \hfill \textbf{\footnotesize JUILLET 1916 – AVRIL 1917}\par\vspace{0.25mm}
\noindent{\fontsize{6.9pt}{8.0pt}\selectfont Couverture rouge à motif losangé et étiquette bleue autographe : Livre-Journal, du 25 juillet 1916 au 30 avril 1917}\par\vspace{0.2mm}
\noindent{\fontsize{6.9pt}{8.0pt}\selectfont\textit{25 juillet 1916 – 30 avril 1917} \hfill \textit{Somme, Flandres, Artois, Paris, Toulouse, Millau}}\par\vspace{0.5mm}
\noindent\rule{\textwidth}{0.35pt}\par\vspace{0.8mm}
% --- Corps de la transcription diplomatique et apparat marginal en prolongement ---
{\renewcommand{\arraystretch}{0.87}%
\noindent\hspace*{3.5mm}\begin{tabular}{@{}p{87mm}@{\hspace{3.5mm}}p{39mm}@{}}
      \leavevmode\fontsize{6.9pt}{7.9pt}\selectfont\strut \underline{Livre - Journal}\strut &  \\
      \leavevmode\fontsize{6.9pt}{7.9pt}\selectfont\strut —\strut &  \\
      \leavevmode\fontsize{6.9pt}{7.9pt}\selectfont\strut Du 25 juillet 1916\strut &  \\
      \leavevmode\fontsize{6.9pt}{7.9pt}\selectfont\strut au 30 avril 1917\strut &  \\
      \leavevmode\llap{\fontsize{6.4pt}{7.3pt}\selectfont\color{darkgray}5\hspace{1.8mm}}\fontsize{6.9pt}{7.9pt}\selectfont\strut —\strut &  \\
\end{tabular}}%
\par\vspace*{\fill}
% --- Registres de bas de page (1. Critique interne [si pertinent] ; 2. Loci paralleli en latin ; 3. Commentaire historique en retrait) ---
\noindent\rule{\textwidth}{0.5pt}\par\vspace{0.65mm}
{\fontsize{6.6pt}{7.7pt}\selectfont%
\noindent\normalfont \textbf{Codicologie :} Plat supérieur du carnet revêtu d'un papier fort rouge gaufré à motif de losanges ; au centre, étiquette rectangulaire de papier bleu à pans coupés portant l'intitulé autographe en cinq lignes à l'encre noire, avec soulignement du titre et deux tirets séparateurs.\par\vspace{0.5mm}
\noindent\rule{22mm}{0.3pt}\par\vspace{0.45mm}
\noindent\normalfont\color{black!92} \textbf{1} (\textit{Livre - Journal}) cf. infra f\textsuperscript{o}~1~r\textsuperscript{o} (p.~\pageref{fol:2}), f\textsuperscript{o}~93~r\textsuperscript{o} (p.~\pageref{fol:186}). — \textbf{3} (\textit{25 juillet 1916}) cf. infra f\textsuperscript{o}~1~r\textsuperscript{o} (p.~\pageref{fol:2}). — \textbf{4} (\textit{30 avril 1917}) cf. infra f\textsuperscript{o}~93~r\textsuperscript{o} (p.~\pageref{fol:186}).\par\vspace{0.55mm}
\noindent\rule{38mm}{0.4pt}\par\vspace{0.5mm}
\noindent\hspace*{5mm}\parbox{\dimexpr\textwidth-5mm\relax}{\fontsize{6.6pt}{7.7pt}\selectfont\color{black!84}Les bornes chronologiques portées sur l'étiquette couvrent neuf mois de campagne d'\textbf{Hippolyte Combes} (né à Millau le 14 août 1876, décédé le 2 septembre 1970 à l'Hôtel du Centre à Saint-Chély-d'Apcher, Lozère) auprès du corps expéditionnaire britannique (\textit{British Expeditionary Force}). Affecté depuis août 1915 à la \textbf{152\textsuperscript{e} Compagnie de l'\textit{Army Service Corps}} au sein du Train divisionnaire de la \textbf{18\textsuperscript{e} Division britannique} (\textit{18th Eastern Division}), Hippolyte Combes passe en novembre 1916 au service des achats et réquisitions de fourrage (\textit{Foraging} puis \textit{Army Purchase Board}) dans la Somme.}\par}
\newpage
\phantomsection\label{fac:1}
% --- En-tête du fac-similé (Page impaire / droite : Photographie du manuscrit) ---
\noindent\rule{\textwidth}{0.5pt}\par\vspace{0.55mm}
\noindent\textbf{\fontsize{6.8pt}{7.8pt}\selectfont FAC-SIMILÉ EN VIS-À-VIS — COUVERTURE (PLAT SUPÉRIEUR)} \hfill \textbf{\fontsize{6.8pt}{7.8pt}\selectfont JUILLET 1916 – AVRIL 1917}\par\vspace{0.25mm}
\noindent{\fontsize{6.8pt}{8.0pt}\selectfont\textit{Livre-Journal d'Hippolyte Combes} — 25 juillet 1916 – 30 avril 1917 (Somme, Flandres, Artois, Paris, Toulouse, Millau)}\par\vspace{0.55mm}
\noindent\rule{\textwidth}{0.35pt}\par\vspace*{\fill}
{\centering
  \setlength{\fboxsep}{1.2mm}%
  \setlength{\fboxrule}{0.4pt}%
  \fcolorbox{gray!55}{white}{%
    \includegraphics[width=0.95\textwidth,height=0.81\textheight,keepaspectratio]{\facsimiledir/20251126_092853.jpg}%
  }\par}
\vspace*{\fill}
\newpage

% =============================================================================
% VIS-A-VIS SPREAD : FOLIO 1 RECTO (20251126_092932.jpg)
% =============================================================================
\cleardoubleevenpage
\renewcommand{\currentmonthyear}{Juillet 1916}
\renewcommand{\currentfoliolabel}{Folio 1 recto}
\phantomsection\label{fol:2}
% --- En-tête du folio (Page paire / gauche : Transcription, apparat & registres de bas de page) ---
\noindent\rule{\textwidth}{0.5pt}\par\vspace{0.55mm}
\noindent\textbf{\footnotesize FOLIO 1 RECTO} \hfill \textbf{\footnotesize JUILLET 1916}\par\vspace{0.25mm}
\noindent{\fontsize{6.9pt}{8.0pt}\selectfont Mardi 25 juillet 1916 — Lynde et Sercus : rassemblement avec le lieutenant Curtis devant l'église de Sercus et cantonnement de 36 officiers des Royal Fusiliers}\par\vspace{0.2mm}
\noindent{\fontsize{6.9pt}{8.0pt}\selectfont\textit{Mardi 25 juillet 1916} \hfill \textit{Lynde et Sercus (Nord)}}\par\vspace{0.5mm}
\noindent\rule{\textwidth}{0.35pt}\par\vspace{0.8mm}
% --- Corps de la transcription diplomatique et apparat marginal en prolongement ---
{\renewcommand{\arraystretch}{0.87}%
\noindent\hspace*{3.5mm}\begin{tabular}{@{}p{87mm}@{\hspace{3.5mm}}p{39mm}@{}}
      \leavevmode\fontsize{6.9pt}{7.9pt}\selectfont\strut \underline{25 Ma.} Levé dès 8 H m. Bien reposé \hfill{} \raisebox{0.5ex}{\underline{Juillet 1916}}\strut & \smash{\parbox[t]{39mm}{\fontsize{6.2pt}{7.0pt}\selectfont\raggedright\color{black!88}\textbf{1 Ma.} \textit{scil.} Mardi. \textbf{1 m.} \textit{abbr. pro} matin. — \textbf{2 Margte} \textit{scil.} Marguerite (Calmels).}} \\
      \leavevmode\fontsize{6.9pt}{7.9pt}\selectfont\strut Courte nuit, pensé très vivement à Margte —\strut &  \\
      \leavevmode\fontsize{6.9pt}{7.9pt}\selectfont\strut Toujours même ciel gris et bas : ciel de\strut &  \\
      \leavevmode\fontsize{6.9pt}{7.9pt}\selectfont\strut Flandre. Beaux arbres élancés aux branches\strut &  \\
      \leavevmode\llap{\fontsize{6.4pt}{7.3pt}\selectfont\color{darkgray}5\hspace{1.8mm}}\fontsize{6.9pt}{7.9pt}\selectfont\strut retombantes entourent pâturages ou champs\strut &  \\
      \leavevmode\fontsize{6.9pt}{7.9pt}\selectfont\strut d'avoine et de blé — épis très hauts \& serrés\strut &  \\
      \leavevmode\fontsize{6.9pt}{7.9pt}\selectfont\strut Mélancolie éparse à travers atmosphère\strut &  \\
      \leavevmode\fontsize{6.9pt}{7.9pt}\selectfont\strut Matin, venu à pied, avec fantastique\strut &  \\
      \leavevmode\fontsize{6.9pt}{7.9pt}\selectfont\strut lieutt Curtis au village de Sercus, où notre\strut & \smash{\parbox[t]{39mm}{\fontsize{6.2pt}{7.0pt}\selectfont\raggedright\color{black!88}\textbf{9 lieutt} \textit{abbr. pro} lieutenant.}} \\
      \leavevmode\llap{\fontsize{6.4pt}{7.3pt}\selectfont\color{darkgray}10\hspace{1.8mm}}\fontsize{6.9pt}{7.9pt}\selectfont\strut rassemblement se fait devant cimetière qui\strut &  \\
      \leavevmode\fontsize{6.9pt}{7.9pt}\selectfont\strut entoure église très curieuse —\strut &  \\
      \leavevmode\fontsize{6.9pt}{7.9pt}\selectfont\strut Après-midi : bien reposé chambre — quand\strut &  \\
      \leavevmode\fontsize{6.9pt}{7.9pt}\selectfont\strut terminé déménagement sur le derrière, à\strut &  \\
      \leavevmode\fontsize{6.9pt}{7.9pt}\selectfont\strut cause de l'arrivée 36 off. du R. Fusrs —\strut & \smash{\parbox[t]{39mm}{\fontsize{6.2pt}{7.0pt}\selectfont\raggedright\color{black!88}\textbf{14 off.} \textit{abbr. pro} officier(s). \textbf{14 R. Fusrs} \textit{angl.} Royal Fusiliers.}} \\
      \leavevmode\llap{\fontsize{6.4pt}{7.3pt}\selectfont\color{darkgray}15\hspace{1.8mm}}\fontsize{6.9pt}{7.9pt}\selectfont\strut Visité longuement église village : remarqué\strut &  \\
      \leavevmode\fontsize{6.9pt}{7.9pt}\selectfont\strut 2 rangs de statues (en plâtre, en bois) sous les\strut &  \\
      \leavevmode\fontsize{6.9pt}{7.9pt}\selectfont\strut orgues\strut &  \\
      \leavevmode\fontsize{6.9pt}{7.9pt}\selectfont\strut Vers 7 h 1/2 s., visite des dames du chateau\strut & \smash{\parbox[t]{39mm}{\fontsize{6.2pt}{7.0pt}\selectfont\raggedright\color{black!88}\textbf{18 1/2} \textit{scil.} demi / moitié. \textbf{18 s.} \textit{abbr. pro} soir. — \textbf{19 off.} \textit{abbr. pro} officier(s). \textbf{19 mess} \textit{angl. in ms.} — \textbf{20 Mme} \textit{abbr. pro} Madame. \textbf{20 f} \textit{abbr. pro} franc. — \textbf{21 Varia} \textit{lat. in ms.} \textbf{21 c} \textit{abbr. pro} centimes.}} \\
      \leavevmode\fontsize{6.9pt}{7.9pt}\selectfont\strut et les ai conduites au mess de mes off.\strut &  \\
      \cline{1-1}
      \leavevmode\llap{\fontsize{6.4pt}{7.3pt}\selectfont\color{darkgray}20\hspace{1.8mm}}\fontsize{6.9pt}{7.9pt}\selectfont\strut A Mme Backer, pour repas d'hier \& 4 oeufs : 1 f 75\strut &  \\
      \leavevmode\fontsize{6.9pt}{7.9pt}\selectfont\strut Varia : 75 c\strut &  \\
\end{tabular}}%
\par\vspace*{\fill}
% --- Registres de bas de page (1. Critique interne [si pertinent] ; 2. Loci paralleli en latin ; 3. Commentaire historique en retrait) ---
\noindent\rule{\textwidth}{0.5pt}\par\vspace{0.65mm}
{\fontsize{6.6pt}{7.7pt}\selectfont%
\noindent\normalfont \textbf{Disposition de la page :} Première page écrite du carnet, ouverte \textit{in medias res} le mardi 25 juillet 1916. Hippolyte Combes place la date soulignée (\textit{25 Ma.}) en tête de la l.~1 et le mois souligné (\textit{Juillet 1916}) légèrement surélevé dans l'angle supérieur droit ; au bas du feuillet (l.~20--21), un trait horizontal sépare le journal proprement dit du relevé des dépenses.\par\vspace{0.5mm}
\noindent\rule{22mm}{0.3pt}\par\vspace{0.45mm}
\noindent\normalfont\color{black!92} \textbf{2} (\textit{Margte}) cf. infra f\textsuperscript{o}~3~r\textsuperscript{o} (p.~\pageref{fol:6}), f\textsuperscript{o}~9~v\textsuperscript{o} (p.~\pageref{fol:19}), \textit{et passim}, Appendicem I, doc.~5 (p.~\pageref{annexe:doc5}). — \textbf{9} (\textit{Curtis}) cf. infra f\textsuperscript{o}~3~r\textsuperscript{o} (p.~\pageref{fol:6}), f\textsuperscript{o}~6~v\textsuperscript{o} (p.~\pageref{fol:13}), \textit{et passim}. — \textbf{14} (\textit{R. Fusrs}) cf. infra f\textsuperscript{o}~1~v\textsuperscript{o} (p.~\pageref{fol:3}), f\textsuperscript{o}~37~v\textsuperscript{o} (p.~\pageref{fol:75}).\par\vspace{0.55mm}
\noindent\rule{38mm}{0.4pt}\par\vspace{0.5mm}
\noindent\hspace*{5mm}\parbox{\dimexpr\textwidth-5mm\relax}{\fontsize{6.6pt}{7.7pt}\selectfont\color{black!84}Relevée de la Somme le 20 juillet 1916 après les combats du bois des Trônes, la \textbf{18\textsuperscript{e} Division britannique} (\textit{18th Eastern Division}, commandée par le major-général Sir Ivor Maxse) gagne la Flandre intérieure (Capt. G.~H.~F. Nichols, \textit{The 18th Division in the Great War}, Londres, Blackwood, 1922, p.~81–85). L'église entourée de son cimetière à \textbf{Sercus} (Nord) est l'église \textbf{Saint-Érasme} (soubassements romans du XII\textsuperscript{e}~s., transformée en \textit{hallekerke} à trois nefs au XVI\textsuperscript{e}~s. ; Base Mérimée, \texttt{PA00107816}).}\par}
\newpage
\phantomsection\label{fac:2}
% --- En-tête du fac-similé (Page impaire / droite : Photographie du manuscrit) ---
\noindent\rule{\textwidth}{0.5pt}\par\vspace{0.55mm}
\noindent\textbf{\fontsize{6.8pt}{7.8pt}\selectfont FAC-SIMILÉ EN VIS-À-VIS — FOLIO 1 RECTO} \hfill \textbf{\fontsize{6.8pt}{7.8pt}\selectfont JUILLET 1916}\par\vspace{0.25mm}
\noindent{\fontsize{6.8pt}{8.0pt}\selectfont\textit{Livre-Journal d'Hippolyte Combes} — Mardi 25 juillet 1916 (Lynde et Sercus (Nord))}\par\vspace{0.55mm}
\noindent\rule{\textwidth}{0.35pt}\par\vspace*{\fill}
{\centering
  \setlength{\fboxsep}{1.2mm}%
  \setlength{\fboxrule}{0.4pt}%
  \fcolorbox{gray!55}{white}{%
    \includegraphics[width=0.95\textwidth,height=0.81\textheight,keepaspectratio]{\facsimiledir/20251126_092932.jpg}%
  }\par}
\vspace*{\fill}
\newpage

% =============================================================================
% VIS-A-VIS SPREAD : FOLIO 1 VERSO (20251126_093018.jpg)
% =============================================================================
\cleardoubleevenpage
\renewcommand{\currentmonthyear}{Juillet 1916}
\renewcommand{\currentfoliolabel}{Folio 1 verso}
\phantomsection\label{fol:3}
% --- En-tête du folio (Page paire / gauche : Transcription, apparat & registres de bas de page) ---
\noindent\rule{\textwidth}{0.5pt}\par\vspace{0.55mm}
\noindent\textbf{\footnotesize FOLIO 1 VERSO} \hfill \textbf{\footnotesize JUILLET 1916}\par\vspace{0.25mm}
\noindent{\fontsize{6.9pt}{8.0pt}\selectfont Mercredi 26 juillet 1916 — Lynde : préparation de 36 chambres d'officiers, château de la brasserie et annonce d'un nouveau départ}\par\vspace{0.2mm}
\noindent{\fontsize{6.9pt}{8.0pt}\selectfont\textit{Mercredi 26 juillet 1916} \hfill \textit{Lynde (Nord)}}\par\vspace{0.5mm}
\noindent\rule{\textwidth}{0.35pt}\par\vspace{0.8mm}
% --- Corps de la transcription diplomatique et apparat marginal en prolongement ---
{\renewcommand{\arraystretch}{1.00}%
\noindent\hspace*{3.5mm}\begin{tabular}{@{}p{87mm}@{\hspace{3.5mm}}p{39mm}@{}}
      \leavevmode\fontsize{7.7pt}{9.0pt}\selectfont\strut \underline{26 Mer.} \hfill{} \underline{Juillet /16}\strut & \smash{\parbox[t]{39mm}{\fontsize{6.2pt}{7.0pt}\selectfont\raggedright\color{black!88}\textbf{1 Mer.} \textit{scil.} Mercredi.}} \\
      \leavevmode\fontsize{7.7pt}{9.0pt}\selectfont\strut De 9 H m. j'à 2 H 1/2 s., occupé au cantonnement\strut & \smash{\parbox[t]{39mm}{\fontsize{6.2pt}{7.0pt}\selectfont\raggedright\color{black!88}\textbf{2 j'à} \textit{abbr. pro} jusqu'à. \textbf{2 1/2} \textit{scil.} demi / moitié. — \textbf{3 Lt} \textit{abbr. pro} Lieutenant. \textbf{3 billeting} \textit{angl. in ms.} — \textbf{5 ds} \textit{abbr. pro} dans. — \textbf{6} $[\![$[illis.]$]\!]\ $ \textit{del. in ms.} \textbf{6} \textit{lect. incert. in ms.} — \textbf{8 ds} \textit{abbr. pro} dans.}} \\
      \leavevmode\fontsize{7.7pt}{9.0pt}\selectfont\strut des Royal Fusiliers. Accompagné Lt Michel, billeting\strut &  \\
      \leavevmode\fontsize{7.7pt}{9.0pt}\selectfont\strut officer. d’ 1 extrémité du village à l'autre et\strut &  \\
      \leavevmode\llap{\fontsize{6.4pt}{7.3pt}\selectfont\color{darkgray}5\hspace{1.8mm}}\fontsize{7.7pt}{9.0pt}\selectfont\strut même ds les environs -- pour trouver 36 chambres\strut &  \\
      \leavevmode\fontsize{7.7pt}{9.0pt}\selectfont\strut d’off $[\![$[illis.]$]\!]$ – Libre enfin vers 2H, mais harassé\strut &  \\
      \leavevmode\fontsize{7.7pt}{9.0pt}\selectfont\strut de fatigue et mourant de faim –\strut &  \\
      \leavevmode\fontsize{7.7pt}{9.0pt}\selectfont\strut Tout l'après-midi : ds ma chambre et\strut &  \\
      \leavevmode\fontsize{7.7pt}{9.0pt}\selectfont\strut me suis bien reposé –\strut &  \\
      \leavevmode\llap{\fontsize{6.4pt}{7.3pt}\selectfont\color{darkgray}10\hspace{1.8mm}}\fontsize{7.7pt}{9.0pt}\selectfont\strut Diner pris, vers 9 H, resté lgt devant porte\strut & \smash{\parbox[t]{39mm}{\fontsize{6.2pt}{7.0pt}\selectfont\raggedright\color{black!88}\textbf{10 lgt} \textit{abbr. pro} longtemps. — \textbf{11 cottage} \textit{angl. in ms.}}} \\
      \leavevmode\fontsize{7.7pt}{9.0pt}\selectfont\strut cottage, à regarder alentours : église aux\strut &  \\
      \leavevmode\fontsize{7.7pt}{9.0pt}\selectfont\strut briques rouges, grands potagers, à g. moutons\strut &  \\
      \leavevmode\fontsize{7.7pt}{9.0pt}\selectfont\strut — chateau de la brasserie — Ciel très bas et\strut &  \\
      \leavevmode\fontsize{7.7pt}{9.0pt}\selectfont\strut gris. Silence et solitude funèbre –\strut &  \\
      \leavevmode\llap{\fontsize{6.4pt}{7.3pt}\selectfont\color{darkgray}15\hspace{1.8mm}}\fontsize{7.7pt}{9.0pt}\selectfont\strut — appris avec regret que de nouveau changemt\strut & \smash{\parbox[t]{39mm}{\fontsize{6.2pt}{7.0pt}\selectfont\raggedright\color{black!88}\textbf{15 changemt} \textit{abbr. pro} changement. — \textbf{16 ct} \textit{abbr. pro} courant. — \textbf{17 f.} \textit{abbr. pro} franc(s). \textbf{17 c.} \textit{abbr. pro} centimes.}} \\
      \leavevmode\fontsize{7.7pt}{9.0pt}\selectfont\strut de cantonnement, le 28 ct –\strut &  \\
      \cline{1-1}
      \leavevmode\fontsize{7.7pt}{9.0pt}\selectfont\strut Food : 60 c. – – Washing : 1 f. –\strut &  \\
\end{tabular}}%
\par\vspace*{\fill}
% --- Registres de bas de page (1. Critique interne [si pertinent] ; 2. Loci paralleli en latin ; 3. Commentaire historique en retrait) ---
\noindent\rule{\textwidth}{0.5pt}\par\vspace{0.65mm}
{\fontsize{6.6pt}{7.7pt}\selectfont%
\noindent\normalfont \textbf{Disposition :} Un trait horizontal sépare au bas du feuillet le récit quotidien et le relevé comptable bilingue de la l.~17 (\textit{Food} / \textit{Washing}).\par\vspace{0.5mm}
\noindent\rule{22mm}{0.3pt}\par\vspace{0.45mm}
\noindent\normalfont\color{black!92} \textbf{3} (\textit{Royal Fusiliers}) cf. supra f\textsuperscript{o}~1~r\textsuperscript{o} (p.~\pageref{fol:2}), infra f\textsuperscript{o}~37~v\textsuperscript{o} (p.~\pageref{fol:75}), f\textsuperscript{o}~46~r\textsuperscript{o} (p.~\pageref{fol:92}).\par\vspace{0.55mm}
\noindent\rule{38mm}{0.4pt}\par\vspace{0.5mm}
\noindent\hspace*{5mm}\parbox{\dimexpr\textwidth-5mm\relax}{\fontsize{6.6pt}{7.7pt}\selectfont\color{black!84}À \textbf{Lynde} (Nord), Hippolyte Combes assiste le lieutenant Michel (\textit{billeting officer}) pour préparer le cantonnement de 36 officiers des \textbf{\textit{Royal Fusiliers}} (dont le 11\textsuperscript{e} Bataillon appartient à la 54\textsuperscript{e} Brigade de la 18\textsuperscript{e} Division) avant l'ordre de départ pour la frontière belge.}\par}
\newpage
\phantomsection\label{fac:3}
% --- En-tête du fac-similé (Page impaire / droite : Photographie du manuscrit) ---
\noindent\rule{\textwidth}{0.5pt}\par\vspace{0.55mm}
\noindent\textbf{\fontsize{6.8pt}{7.8pt}\selectfont FAC-SIMILÉ EN VIS-À-VIS — FOLIO 1 VERSO} \hfill \textbf{\fontsize{6.8pt}{7.8pt}\selectfont JUILLET 1916}\par\vspace{0.25mm}
\noindent{\fontsize{6.8pt}{8.0pt}\selectfont\textit{Livre-Journal d'Hippolyte Combes} — Mercredi 26 juillet 1916 (Lynde (Nord))}\par\vspace{0.55mm}
\noindent\rule{\textwidth}{0.35pt}\par\vspace*{\fill}
{\centering
  \setlength{\fboxsep}{1.2mm}%
  \setlength{\fboxrule}{0.4pt}%
  \fcolorbox{gray!55}{white}{%
    \includegraphics[width=0.95\textwidth,height=0.81\textheight,keepaspectratio]{\facsimiledir/20251126_093018.jpg}%
  }\par}
\vspace*{\fill}
\newpage

% =============================================================================
% VIS-A-VIS SPREAD : FOLIO 2 RECTO (20251126_093100.jpg)
% =============================================================================
\cleardoubleevenpage
\renewcommand{\currentmonthyear}{Juillet 1916}
\renewcommand{\currentfoliolabel}{Folio 2 recto}
\phantomsection\label{fol:4}
% --- En-tête du folio (Page paire / gauche : Transcription, apparat & registres de bas de page) ---
\noindent\rule{\textwidth}{0.5pt}\par\vspace{0.55mm}
\noindent\textbf{\footnotesize FOLIO 2 RECTO} \hfill \textbf{\footnotesize JUILLET 1916}\par\vspace{0.25mm}
\noindent{\fontsize{6.9pt}{8.0pt}\selectfont Jeudi 27 juillet 1916 — De Lynde à Renescure et Hazebrouck : châteaux du Moyen Âge, jardins potagers et achats de cartes postales d'Ypres et de Renescure}\par\vspace{0.2mm}
\noindent{\fontsize{6.9pt}{8.0pt}\selectfont\textit{Jeudi 27 juillet 1916} \hfill \textit{Lynde, Renescure et Hazebrouck (Nord)}}\par\vspace{0.5mm}
\noindent\rule{\textwidth}{0.35pt}\par\vspace{0.8mm}
% --- Corps de la transcription diplomatique et apparat marginal en prolongement ---
{\renewcommand{\arraystretch}{0.83}%
\noindent\hspace*{3.5mm}\begin{tabular}{@{}p{87mm}@{\hspace{3.5mm}}p{39mm}@{}}
      \leavevmode\fontsize{6.6pt}{7.5pt}\selectfont\strut \underline{27 J.} \hfill{} \underline{Juillet /16}\strut & \smash{\parbox[t]{39mm}{\fontsize{6.2pt}{7.0pt}\selectfont\raggedright\color{black!88}\textbf{1 J.} \textit{scil.} Jeudi.}} \\
      \leavevmode\fontsize{6.6pt}{7.5pt}\selectfont\strut – Aussitôt après breakfast, venu à pied\strut & \smash{\parbox[t]{39mm}{\fontsize{6.2pt}{7.0pt}\selectfont\raggedright\color{black!88}\textbf{2 breakfast} \textit{angl. in ms.}}} \\
      \leavevmode\fontsize{6.6pt}{7.5pt}\selectfont\strut à Renescure. Tjours même temps gris et\strut &  \\
      \leavevmode\fontsize{6.6pt}{7.5pt}\selectfont\strut saturé de mélancolie.\strut &  \\
      \leavevmode\llap{\fontsize{6.4pt}{7.3pt}\selectfont\color{darkgray}5\hspace{1.8mm}}\fontsize{6.6pt}{7.5pt}\selectfont\strut Village très intéressant\strut &  \\
      \leavevmode\fontsize{6.6pt}{7.5pt}\selectfont\strut 2 superbes châteaux M. Age – fossé entourant\strut &  \\
      \leavevmode\fontsize{6.6pt}{7.5pt}\selectfont\strut château – tours à chaque angle du carré –\strut &  \\
      \leavevmode\fontsize{6.6pt}{7.5pt}\selectfont\strut donjon etc. etc.\strut &  \\
      \leavevmode\fontsize{6.6pt}{7.5pt}\selectfont\strut Eglise m’a aussi retenu lgt – Double nef\strut & \smash{\parbox[t]{39mm}{\fontsize{6.2pt}{7.0pt}\selectfont\raggedright\color{black!88}\textbf{9 lgt} \textit{abbr. pro} longtemps.}} \\
      \leavevmode\llap{\fontsize{6.4pt}{7.3pt}\selectfont\color{darkgray}10\hspace{1.8mm}}\fontsize{6.6pt}{7.5pt}\selectfont\strut Après lunch, venu à cheval à Hazebrouck\strut & \smash{\parbox[t]{39mm}{\fontsize{6.2pt}{7.0pt}\selectfont\raggedright\color{black!88}\textbf{10 lunch} \textit{angl. in ms.}}} \\
      \leavevmode\fontsize{6.6pt}{7.5pt}\selectfont\strut Ville parait assez importante\strut &  \\
      \leavevmode\fontsize{6.6pt}{7.5pt}\selectfont\strut grande place –\strut &  \\
      \leavevmode\fontsize{6.6pt}{7.5pt}\selectfont\strut Eglise – (avec Square à g.)\strut &  \\
      \leavevmode\fontsize{6.6pt}{7.5pt}\selectfont\strut Fin journée mélancolique : arpenté lgt\strut & \smash{\parbox[t]{39mm}{\fontsize{6.2pt}{7.0pt}\selectfont\raggedright\color{black!88}\textbf{14 lgt} \textit{abbr. pro} longtemps.}} \\
      \leavevmode\llap{\fontsize{6.4pt}{7.3pt}\selectfont\color{darkgray}15\hspace{1.8mm}}\fontsize{6.6pt}{7.5pt}\selectfont\strut seul petite allée jardins potagers devant mon\strut &  \\
      \leavevmode\fontsize{6.6pt}{7.5pt}\selectfont\strut cottage. Temps de brouillard. Silence et\strut & \smash{\parbox[t]{39mm}{\fontsize{6.2pt}{7.0pt}\selectfont\raggedright\color{black!88}\textbf{16 cottage} \textit{angl. in ms.}}} \\
      \leavevmode\fontsize{6.6pt}{7.5pt}\selectfont\strut désert funèbre.\strut &  \\
      \leavevmode\fontsize{6.6pt}{7.5pt}\selectfont\strut \underline{A Renescure} : boîte de plumes : 1 f 50. —\strut &  \\
      \leavevmode\fontsize{6.6pt}{7.5pt}\selectfont\strut 2 carnets timbres-poste : 2 f + 2 f. —\strut & \smash{\parbox[t]{39mm}{\fontsize{6.2pt}{7.0pt}\selectfont\raggedright\color{black!88}\textbf{19 f.} \textit{abbr. pro} franc(s).}} \\
      \cline{1-1}
      \leavevmode\llap{\fontsize{6.4pt}{7.3pt}\selectfont\color{darkgray}20\hspace{1.8mm}}\fontsize{6.6pt}{7.5pt}\selectfont\strut Cartes ples : d'Ypres : 1 f 20. – C. P. de Renescure : 60 c +\strut & \smash{\parbox[t]{39mm}{\fontsize{6.2pt}{7.0pt}\selectfont\raggedright\color{black!88}\textbf{20 C. P.} \textit{abbr. pro} Cartes postales. \textbf{20 Cartes ples} \textit{abbr. pro} Cartes postales. — \textbf{21 C. ples} \textit{abbr. pro} Cartes postales. \textbf{21 c.} \textit{abbr. pro} centimes. — \textbf{22 fr.} \textit{abbr. pro} franc(s). — \textbf{23 C. ples} \textit{abbr. pro} Cartes postales. — \textbf{24 c.} \textit{abbr. pro} centimes.}} \\
      \leavevmode\fontsize{6.6pt}{7.5pt}\selectfont\strut  \hfill{} C. ples (varia) : 50 c. —\strut &  \\
      \cline{1-1}
      \leavevmode\fontsize{6.6pt}{7.5pt}\selectfont\strut \underline{A Hazebrouck} : Solde : 84 fr. —\strut &  \\
      \leavevmode\fontsize{6.6pt}{7.5pt}\selectfont\strut C. ples : 60 + 10 c – \hfill{} 1 cahier cartonné : 1 f 50\strut &  \\
      \leavevmode\fontsize{6.6pt}{7.5pt}\selectfont\strut Cte d'Etat-major : 90 c. – \hfill{} 3 élastiq. de rechange\strut &  \\
      \leavevmode\llap{\fontsize{6.4pt}{7.3pt}\selectfont\color{darkgray}25\hspace{1.8mm}}\fontsize{6.6pt}{7.5pt}\selectfont\strut  \hfill{} (bretelles) : 65 c —\strut &  \\
\end{tabular}}%
\par\vspace*{\fill}
% --- Registres de bas de page (1. Critique interne [si pertinent] ; 2. Loci paralleli en latin ; 3. Commentaire historique en retrait) ---
\noindent\rule{\textwidth}{0.5pt}\par\vspace{0.65mm}
{\fontsize{6.6pt}{7.7pt}\selectfont%
\noindent\normalfont \textbf{Structure comptable :} Au bas du feuillet (l.~18--24), le relevé des dépenses et de la solde est séparé du récit par un trait horizontal et subdivisé par deux sous-titres topographiques soulignés (\textit{A Renescure}, \textit{A Hazebrouck}).\par\vspace{0.5mm}
\noindent\rule{22mm}{0.3pt}\par\vspace{0.45mm}
\noindent\normalfont\color{black!92} \textbf{10} (\textit{Hazebrouck}) cf. infra f\textsuperscript{o}~2~v\textsuperscript{o} (p.~\pageref{fol:5}), f\textsuperscript{o}~4~r\textsuperscript{o} (p.~\pageref{fol:8}). — \textbf{20} (\textit{Ypres}) cf. infra f\textsuperscript{o}~4~v\textsuperscript{o} (p.~\pageref{fol:9}). — \textbf{22} (\textit{84 f.}) cf. infra f\textsuperscript{o}~10~v\textsuperscript{o} (p.~\pageref{fol:21}), f\textsuperscript{o}~39~v\textsuperscript{o} (p.~\pageref{fol:79}).\par\vspace{0.55mm}
\noindent\rule{38mm}{0.4pt}\par\vspace{0.5mm}
\noindent\hspace*{5mm}\parbox{\dimexpr\textwidth-5mm\relax}{\fontsize{6.6pt}{7.7pt}\selectfont\color{black!84}À \textbf{Renescure} (Nord), les deux châteaux remarqués par Hippolyte Combes correspondent au \textbf{château de Zuthove} (demeure seigneuriale de briques entourée de douves, reconstruite à la fin du XV\textsuperscript{e}~s. ; Base Mérimée, \texttt{PA00107784}) et au \textbf{château dit de Philippe de Comines} (ancienne motte castrale remaniée et transformée en mairie ; \texttt{PA00107783}), voisin de l'église-halle \textbf{Notre-Dame-de-l'Assomption}. À \textbf{Hazebrouck}, il visite l'église \textbf{Saint-Éloi} (XVI\textsuperscript{e}~s.) et perçoit 84~francs de solde et indemnités auprès de la Mission française.}\par}
\newpage
\phantomsection\label{fac:4}
% --- En-tête du fac-similé (Page impaire / droite : Photographie du manuscrit) ---
\noindent\rule{\textwidth}{0.5pt}\par\vspace{0.55mm}
\noindent\textbf{\fontsize{6.8pt}{7.8pt}\selectfont FAC-SIMILÉ EN VIS-À-VIS — FOLIO 2 RECTO} \hfill \textbf{\fontsize{6.8pt}{7.8pt}\selectfont JUILLET 1916}\par\vspace{0.25mm}
\noindent{\fontsize{6.8pt}{8.0pt}\selectfont\textit{Livre-Journal d'Hippolyte Combes} — Jeudi 27 juillet 1916 (Lynde, Renescure et Hazebrouck (Nord))}\par\vspace{0.55mm}
\noindent\rule{\textwidth}{0.35pt}\par\vspace*{\fill}
{\centering
  \setlength{\fboxsep}{1.2mm}%
  \setlength{\fboxrule}{0.4pt}%
  \fcolorbox{gray!55}{white}{%
    \includegraphics[width=0.95\textwidth,height=0.81\textheight,keepaspectratio]{\facsimiledir/20251126_093100.jpg}%
  }\par}
\vspace*{\fill}
\newpage

% =============================================================================
% VIS-A-VIS SPREAD : FOLIO 2 VERSO (20251126_093141.jpg)
% =============================================================================
\cleardoubleevenpage
\renewcommand{\currentmonthyear}{Juillet 1916}
\renewcommand{\currentfoliolabel}{Folio 2 verso}
\phantomsection\label{fol:5}
% --- En-tête du folio (Page paire / gauche : Transcription, apparat & registres de bas de page) ---
\noindent\rule{\textwidth}{0.5pt}\par\vspace{0.55mm}
\noindent\textbf{\footnotesize FOLIO 2 VERSO} \hfill \textbf{\footnotesize JUILLET 1916}\par\vspace{0.25mm}
\noindent{\fontsize{6.9pt}{8.0pt}\selectfont Vendredi 28 juillet 1916 — Départ de Lynde (Mme Donze) et lente chevauchée sous la chaleur par Hazebrouck jusqu'à Saint-Jans-Cappel}\par\vspace{0.2mm}
\noindent{\fontsize{6.9pt}{8.0pt}\selectfont\textit{Vendredi 28 juillet 1916} \hfill \textit{De Lynde vers Saint-Jans-Cappel par Hazebrouck (Nord)}}\par\vspace{0.5mm}
\noindent\rule{\textwidth}{0.35pt}\par\vspace{0.8mm}
% --- Corps de la transcription diplomatique et apparat marginal en prolongement ---
{\renewcommand{\arraystretch}{0.87}%
\noindent\hspace*{3.5mm}\begin{tabular}{@{}p{87mm}@{\hspace{3.5mm}}p{39mm}@{}}
      \leavevmode\fontsize{6.9pt}{7.9pt}\selectfont\strut \underline{28 Ve.} \hfill{} \underline{Juillet / 16}\strut & \smash{\parbox[t]{39mm}{\fontsize{6.2pt}{7.0pt}\selectfont\raggedright\color{black!88}\textbf{1 Ve.} \textit{scil.} Vendredi.}} \\
      \leavevmode\fontsize{6.9pt}{7.9pt}\selectfont\strut Hier soir, avant de me coucher,\strut &  \\
      \leavevmode\fontsize{6.9pt}{7.9pt}\selectfont\strut terminé presque entièrement mes paquets\strut &  \\
      \leavevmode\fontsize{6.9pt}{7.9pt}\selectfont\strut Nuit agitée – énervé par ts ces départs successifs\strut & \smash{\parbox[t]{39mm}{\fontsize{6.2pt}{7.0pt}\selectfont\raggedright\color{black!88}\textbf{4 ts} \textit{abbr. pro} tous.}} \\
      \leavevmode\llap{\fontsize{6.4pt}{7.3pt}\selectfont\color{darkgray}5\hspace{1.8mm}}\fontsize{6.9pt}{7.9pt}\selectfont\strut Levé vers 5h 1/2. Ciel voilé de brume qui\strut & \smash{\parbox[t]{39mm}{\fontsize{6.2pt}{7.0pt}\selectfont\raggedright\color{black!88}\textbf{5 1/2} \textit{scil.} demi / moitié.}} \\
      \leavevmode\fontsize{6.9pt}{7.9pt}\selectfont\strut fond en petites gouttelettes. — Journée de mouvt. —\strut & \smash{\parbox[t]{39mm}{\fontsize{6.2pt}{7.0pt}\selectfont\raggedright\color{black!88}\textbf{6 mouvt} \textit{abbr. pro} mouvement. — \textbf{7 1/2} \textit{scil.} demi / moitié.}} \\
      \leavevmode\fontsize{6.9pt}{7.9pt}\selectfont\strut Partis à 8h 1/2. Fait route à cheval, en tête\strut &  \\
      \leavevmode\fontsize{6.9pt}{7.9pt}\selectfont\strut du convoi. —\strut &  \\
      \leavevmode\fontsize{6.9pt}{7.9pt}\selectfont\strut Lente chevauchée le long de routes sablonneuses\strut &  \\
      \leavevmode\llap{\fontsize{6.4pt}{7.3pt}\selectfont\color{darkgray}10\hspace{1.8mm}}\fontsize{6.9pt}{7.9pt}\selectfont\strut bordées de chaque coté par champs d’avoine ou de\strut &  \\
      \leavevmode\fontsize{6.9pt}{7.9pt}\selectfont\strut blé – aux épis très beaux \& très hauts –\strut &  \\
      \leavevmode\fontsize{6.9pt}{7.9pt}\selectfont\strut Contourné Hazebrouck par les faubourgs\strut &  \\
      \leavevmode\fontsize{6.9pt}{7.9pt}\selectfont\strut Peu à peu, soleil ardent perce nuages. Après-\strut &  \\
      \leavevmode\fontsize{6.9pt}{7.9pt}\selectfont\strut midi et fin journée : chaleur intense –\strut &  \\
      \leavevmode\llap{\fontsize{6.4pt}{7.3pt}\selectfont\color{darkgray}15\hspace{1.8mm}}\fontsize{6.9pt}{7.9pt}\selectfont\strut Vers 4h 1/2, enfin, parvenus au terme voyage\strut & \smash{\parbox[t]{39mm}{\fontsize{6.2pt}{7.0pt}\selectfont\raggedright\color{black!88}\textbf{15 1/2} \textit{scil.} demi / moitié. — \textbf{16 gde} \textit{abbr. pro} grande.}} \\
      \leavevmode\fontsize{6.9pt}{7.9pt}\selectfont\strut gde prairie autour d’ 1 ferme – Impossible trouver\strut &  \\
      \leavevmode\fontsize{6.9pt}{7.9pt}\selectfont\strut chambre pour moi – Dû me contenter d’ 1 grenier\strut &  \\
      \leavevmode\fontsize{6.9pt}{7.9pt}\selectfont\strut Harassé de fatigue, suffoqué de chaleur\strut &  \\
      \leavevmode\fontsize{6.9pt}{7.9pt}\selectfont\strut Besoin de repos extrême\strut &  \\
      \cline{1-1}
      \leavevmode\llap{\fontsize{6.4pt}{7.3pt}\selectfont\color{darkgray}20\hspace{1.8mm}}\fontsize{6.9pt}{7.9pt}\selectfont\strut Laissé à Mme [Donze ?], de Lynde (pour cuisine 2 f\strut & \smash{\parbox[t]{39mm}{\fontsize{6.2pt}{7.0pt}\selectfont\raggedright\color{black!88}\textbf{20} \textit{lect. incert. in ms.} \textbf{20 Mme} \textit{abbr. pro} Madame.}} \\
      \leavevmode\fontsize{6.9pt}{7.9pt}\selectfont\strut et 1 f pour oeufs \& lait —\strut &  \\
      \cline{1-1}
      \leavevmode\fontsize{6.9pt}{7.9pt}\selectfont\strut Food (extra) : 10 c, 1 f, 1 f, 60 c. – \hfill{} Cartes ples : 10 c\strut & \smash{\parbox[t]{39mm}{\fontsize{6.2pt}{7.0pt}\selectfont\raggedright\color{black!88}\textbf{22 Cartes ples} \textit{abbr. pro} Cartes postales. \textbf{22 c.} \textit{abbr. pro} centimes. — \textbf{23 f.} \textit{abbr. pro} franc(s). \textbf{23 groom} \textit{angl. in ms.}}} \\
      \leavevmode\fontsize{6.9pt}{7.9pt}\selectfont\strut Tip au groom pour soigné poney : 1 f.\strut &  \\
\end{tabular}}%
\par\vspace*{\fill}
% --- Registres de bas de page (1. Critique interne [si pertinent] ; 2. Loci paralleli en latin ; 3. Commentaire historique en retrait) ---
\noindent\rule{\textwidth}{0.5pt}\par\vspace{0.65mm}
{\fontsize{6.6pt}{7.7pt}\selectfont%
\noindent\normalfont \textbf{Disposition :} Les quatre dernières lignes (l.~20--23), consacrées au règlement de la logeuse de Lynde et aux frais de route, sont isolées du récit par un trait horizontal continu.\par\vspace{0.5mm}
\noindent\rule{22mm}{0.3pt}\par\vspace{0.45mm}
\noindent\normalfont\color{black!92} \textbf{12} (\textit{Hazebrouck}) cf. supra f\textsuperscript{o}~2~r\textsuperscript{o} (p.~\pageref{fol:4}), infra f\textsuperscript{o}~4~r\textsuperscript{o} (p.~\pageref{fol:8}), f\textsuperscript{o}~8~r\textsuperscript{o} (p.~\pageref{fol:16}). — \textbf{20} (\textit{Lynde}) cf. supra f\textsuperscript{o}~1~r\textsuperscript{o} (p.~\pageref{fol:2}), f\textsuperscript{o}~2~r\textsuperscript{o} (p.~\pageref{fol:4}).\par\vspace{0.55mm}
\noindent\rule{38mm}{0.4pt}\par\vspace{0.5mm}
\noindent\hspace*{5mm}\parbox{\dimexpr\textwidth-5mm\relax}{\fontsize{6.6pt}{7.7pt}\selectfont\color{black!84}Étape routière de huit heures (8~h~30 -- 16~h~30) transférant le train divisionnaire de la 18\textsuperscript{e} Division depuis \textbf{Lynde} (nommé dans le règlement de pension à la l.~20) jusqu'aux abords de \textbf{Saint-Jans-Cappel} (Nord, nommé dès la vue suivante, f\textsuperscript{o}~3~r\textsuperscript{o}), au pied des monts des Flandres, en contournant le nœud ferroviaire d'Hazebrouck par les faubourgs.}\par}
\newpage
\phantomsection\label{fac:5}
% --- En-tête du fac-similé (Page impaire / droite : Photographie du manuscrit) ---
\noindent\rule{\textwidth}{0.5pt}\par\vspace{0.55mm}
\noindent\textbf{\fontsize{6.8pt}{7.8pt}\selectfont FAC-SIMILÉ EN VIS-À-VIS — FOLIO 2 VERSO} \hfill \textbf{\fontsize{6.8pt}{7.8pt}\selectfont JUILLET 1916}\par\vspace{0.25mm}
\noindent{\fontsize{6.8pt}{8.0pt}\selectfont\textit{Livre-Journal d'Hippolyte Combes} — Vendredi 28 juillet 1916 (De Lynde vers Saint-Jans-Cappel par Hazebrouck (Nord))}\par\vspace{0.55mm}
\noindent\rule{\textwidth}{0.35pt}\par\vspace*{\fill}
{\centering
  \setlength{\fboxsep}{1.2mm}%
  \setlength{\fboxrule}{0.4pt}%
  \fcolorbox{gray!55}{white}{%
    \includegraphics[width=0.95\textwidth,height=0.81\textheight,keepaspectratio]{\facsimiledir/20251126_093141.jpg}%
  }\par}
\vspace*{\fill}
\newpage

% =============================================================================
% VIS-A-VIS SPREAD : FOLIO 3 RECTO (20251126_093204.jpg)
% =============================================================================
\cleardoubleevenpage
\renewcommand{\currentmonthyear}{Juillet 1916}
\renewcommand{\currentfoliolabel}{Folio 3 recto}
\phantomsection\label{fol:6}
% --- En-tête du folio (Page paire / gauche : Transcription, apparat & registres de bas de page) ---
\noindent\rule{\textwidth}{0.5pt}\par\vspace{0.55mm}
\noindent\textbf{\footnotesize FOLIO 3 RECTO} \hfill \textbf{\footnotesize JUILLET 1916}\par\vspace{0.25mm}
\noindent{\fontsize{6.9pt}{8.0pt}\selectfont Samedi 29 juillet 1916 — Tournée avec le lieutenant Curtis à Saint-Jans-Cappel et achat de fourrages au Mont-Noir}\par\vspace{0.2mm}
\noindent{\fontsize{6.9pt}{8.0pt}\selectfont\textit{Samedi 29 juillet 1916} \hfill \textit{Saint-Jans-Cappel et le Mont-Noir (Nord)}}\par\vspace{0.5mm}
\noindent\rule{\textwidth}{0.35pt}\par\vspace{0.8mm}
% --- Corps de la transcription diplomatique et apparat marginal en prolongement ---
{\renewcommand{\arraystretch}{0.87}%
\noindent\hspace*{3.5mm}\begin{tabular}{@{}p{87mm}@{\hspace{3.5mm}}p{39mm}@{}}
      \leavevmode\fontsize{6.9pt}{7.9pt}\selectfont\strut \underline{Sam. 29} \hfill{} \underline{Juillet 16}\strut & \smash{\parbox[t]{39mm}{\fontsize{6.2pt}{7.0pt}\selectfont\raggedright\color{black!88}\textbf{1 Sam.} \textit{scil.} Samedi.}} \\
      \leavevmode\fontsize{6.9pt}{7.9pt}\selectfont\strut Sur pied à 7H 1/2 –\strut & \smash{\parbox[t]{39mm}{\fontsize{6.2pt}{7.0pt}\selectfont\raggedright\color{black!88}\textbf{2 1/2} \textit{scil.} demi / moitié.}} \\
      \leavevmode\fontsize{6.9pt}{7.9pt}\selectfont\strut Breakfast à peine terminé, que ce fou de\strut & \smash{\parbox[t]{39mm}{\fontsize{6.2pt}{7.0pt}\selectfont\raggedright\color{black!88}\textbf{3 Breakfast} \textit{angl. in ms.}}} \\
      \leavevmode\fontsize{6.9pt}{7.9pt}\selectfont\strut Curtis venu me prendre — Sommes allés chercher fourrage\strut &  \\
      \leavevmode\llap{\fontsize{6.4pt}{7.3pt}\selectfont\color{darkgray}5\hspace{1.8mm}}\fontsize{6.9pt}{7.9pt}\selectfont\strut à St-Jans Cappel – et ds le hameau de Le Mont-Noir\strut & \smash{\parbox[t]{39mm}{\fontsize{6.2pt}{7.0pt}\selectfont\raggedright\color{black!88}\textbf{5 ds} \textit{abbr. pro} dans. \textbf{5 St} \textit{abbr. pro} Saint. — \textbf{6 s/} \textit{abbr. pro} sur. \textbf{6 longt} \textit{abbr. pro} longtemps. — \textbf{8 Ctis} \textit{abbr. pro} Curtis. — \textbf{9 ds} \textit{abbr. pro} dans.}} \\
      \leavevmode\fontsize{6.9pt}{7.9pt}\selectfont\strut Rencontré curieux vicaire de ce village et causé longt s/\strut &  \\
      \leavevmode\fontsize{6.9pt}{7.9pt}\selectfont\strut la dernière attaque à l. notre division a pris part.\strut &  \\
      \leavevmode\fontsize{6.9pt}{7.9pt}\selectfont\strut Après-midi, ce fou de Ctis, revenu me\strut &  \\
      \leavevmode\fontsize{6.9pt}{7.9pt}\selectfont\strut prendre. De 2 à 5 H, au plein soleil, ds 1 champ\strut &  \\
      \leavevmode\llap{\fontsize{6.4pt}{7.3pt}\selectfont\color{darkgray}10\hspace{1.8mm}}\fontsize{6.9pt}{7.9pt}\selectfont\strut où fait faucher \& peser fourrage –\strut &  \\
      \leavevmode\fontsize{6.9pt}{7.9pt}\selectfont\strut Oté veste \& gilet – allongé s/ l’herbe – $[\![$[illis.]$]\!]$\strut & \smash{\parbox[t]{39mm}{\fontsize{6.2pt}{7.0pt}\selectfont\raggedright\color{black!88}\textbf{11} $[\![$[illis.]$]\!]\ $ \textit{del. in ms.} \textbf{11} \textit{lect. incert. in ms.} — \textbf{12 lgt} \textit{abbr. pro} longtemps. \textbf{12 Margte} \textit{scil.} Marguerite (Calmels).}} \\
      \leavevmode\fontsize{6.9pt}{7.9pt}\selectfont\strut paysage et pensé lgt à Margte.\strut &  \\
      \leavevmode\fontsize{6.9pt}{7.9pt}\selectfont\strut Région de plaines, de nombreuses houblonnières\strut &  \\
      \leavevmode\fontsize{6.9pt}{7.9pt}\selectfont\strut A l’horizon, ballons captifs. De temps à autre\strut &  \\
      \leavevmode\llap{\fontsize{6.4pt}{7.3pt}\selectfont\color{darkgray}15\hspace{1.8mm}}\fontsize{6.9pt}{7.9pt}\selectfont\strut bombardement lointain –\strut &  \\
      \leavevmode\fontsize{6.9pt}{7.9pt}\selectfont\strut Ciel traversé d'avions anglais.\strut &  \\
      \leavevmode\fontsize{6.9pt}{7.9pt}\selectfont\strut Après diner : venu promener au village de\strut &  \\
      \leavevmode\fontsize{6.9pt}{7.9pt}\selectfont\strut St-Jans-Cappel (1)\strut & \smash{\parbox[t]{39mm}{\fontsize{6.2pt}{7.0pt}\selectfont\raggedright\color{black!88}\textbf{18 St} \textit{abbr. pro} Saint.}} \\
      \cline{1-1}
      \leavevmode\fontsize{6.9pt}{7.9pt}\selectfont\strut Payé à fermière où je loge pour 4 oeufs et\strut &  \\
      \leavevmode\llap{\fontsize{6.4pt}{7.3pt}\selectfont\color{darkgray}20\hspace{1.8mm}}\fontsize{6.9pt}{7.9pt}\selectfont\strut 2 litres 1/2 de lait : 1 f 60. —\strut & \smash{\parbox[t]{39mm}{\fontsize{6.2pt}{7.0pt}\selectfont\raggedright\color{black!88}\textbf{20 1/2} \textit{scil.} demi / moitié. — \textbf{21 cartes ples} \textit{abbr. pro} cartes postales. \textbf{21 c.} \textit{abbr. pro} centimes.}} \\
      \leavevmode\fontsize{6.9pt}{7.9pt}\selectfont\strut 5 cartes ples : 50 c. —\strut &  \\
      \cline{1-1}
      \leavevmode\fontsize{6.9pt}{7.9pt}\selectfont\strut (1) maisons très basses – blancs rideaux aux\strut &  \\
      \leavevmode\fontsize{6.9pt}{7.9pt}\selectfont\strut fenêtres –\strut &  \\
      \leavevmode\fontsize{6.9pt}{7.9pt}\selectfont\strut église au clocher effilé\strut &  \\
\end{tabular}}%
\par\vspace*{\fill}
% --- Registres de bas de page (1. Critique interne [si pertinent] ; 2. Loci paralleli en latin ; 3. Commentaire historique en retrait) ---
\noindent\rule{\textwidth}{0.5pt}\par\vspace{0.65mm}
{\fontsize{6.6pt}{7.7pt}\selectfont%
\noindent\normalfont\color{black!92} \textbf{4} (\textit{Curtis}) cf. supra f\textsuperscript{o}~1~r\textsuperscript{o} (p.~\pageref{fol:2}), infra f\textsuperscript{o}~6~v\textsuperscript{o} (p.~\pageref{fol:13}), f\textsuperscript{o}~7~r\textsuperscript{o} (p.~\pageref{fol:14}), \textit{et passim}. — \textbf{4} (\textit{fourrage}) cf. infra f\textsuperscript{o}~3~v\textsuperscript{o} (p.~\pageref{fol:7}), f\textsuperscript{o}~4~v\textsuperscript{o} (p.~\pageref{fol:9}), \textit{et passim}. — \textbf{12} (\textit{Margte}) cf. supra f\textsuperscript{o}~1~r\textsuperscript{o} (p.~\pageref{fol:2}), infra f\textsuperscript{o}~9~v\textsuperscript{o} (p.~\pageref{fol:19}), f\textsuperscript{o}~27~r\textsuperscript{o} (p.~\pageref{fol:54}), \textit{et passim}.\par\vspace{0.55mm}
\noindent\rule{38mm}{0.4pt}\par\vspace{0.5mm}
\noindent\hspace*{5mm}\parbox{\dimexpr\textwidth-5mm\relax}{\fontsize{6.6pt}{7.7pt}\selectfont\color{black!84}Tournée d'achat et de pesée de fourrage avec le lieutenant \textbf{Curtis} à \textbf{Saint-Jans-Cappel} et au hameau du \textbf{Mont-Noir} (Nord, sur la frontière franco-belge), au cœur du bassin houblonnier flamand.}\par}
\newpage
\phantomsection\label{fac:6}
% --- En-tête du fac-similé (Page impaire / droite : Photographie du manuscrit) ---
\noindent\rule{\textwidth}{0.5pt}\par\vspace{0.55mm}
\noindent\textbf{\fontsize{6.8pt}{7.8pt}\selectfont FAC-SIMILÉ EN VIS-À-VIS — FOLIO 3 RECTO} \hfill \textbf{\fontsize{6.8pt}{7.8pt}\selectfont JUILLET 1916}\par\vspace{0.25mm}
\noindent{\fontsize{6.8pt}{8.0pt}\selectfont\textit{Livre-Journal d'Hippolyte Combes} — Samedi 29 juillet 1916 (Saint-Jans-Cappel et le Mont-Noir (Nord))}\par\vspace{0.55mm}
\noindent\rule{\textwidth}{0.35pt}\par\vspace*{\fill}
{\centering
  \setlength{\fboxsep}{1.2mm}%
  \setlength{\fboxrule}{0.4pt}%
  \fcolorbox{gray!55}{white}{%
    \includegraphics[width=0.95\textwidth,height=0.81\textheight,keepaspectratio]{\facsimiledir/20251126_093204.jpg}%
  }\par}
\vspace*{\fill}
\newpage

% =============================================================================
% VIS-A-VIS SPREAD : FOLIO 3 VERSO (20251126_093219.jpg)
% =============================================================================
\cleardoubleevenpage
\renewcommand{\currentmonthyear}{Juillet 1916}
\renewcommand{\currentfoliolabel}{Folio 3 verso}
\phantomsection\label{fol:7}
% --- En-tête du folio (Page paire / gauche : Transcription, apparat & registres de bas de page) ---
\noindent\rule{\textwidth}{0.5pt}\par\vspace{0.55mm}
\noindent\textbf{\footnotesize FOLIO 3 VERSO} \hfill \textbf{\footnotesize JUILLET 1916}\par\vspace{0.25mm}
\noindent{\fontsize{6.9pt}{8.0pt}\selectfont Messe dominicale à Saint-Jans-Cappel et corvée de fourrage vert vers Bailleul}\par\vspace{0.2mm}
\noindent{\fontsize{6.9pt}{8.0pt}\selectfont\textit{Dimanche 30 et lundi 31 juillet 1916} \hfill \textit{Saint-Jans-Cappel et Bailleul (Nord)}}\par\vspace{0.5mm}
\noindent\rule{\textwidth}{0.35pt}\par\vspace{0.8mm}
% --- Corps de la transcription diplomatique et apparat marginal en prolongement ---
{\renewcommand{\arraystretch}{0.83}%
\noindent\hspace*{3.5mm}\begin{tabular}{@{}p{87mm}@{\hspace{3.5mm}}p{39mm}@{}}
      \leavevmode\fontsize{6.6pt}{7.5pt}\selectfont\strut \underline{Dim. 30.} \hfill{} \underline{Juillet 1916}\strut & \smash{\parbox[t]{39mm}{\fontsize{6.2pt}{7.0pt}\selectfont\raggedright\color{black!88}\textbf{1 Dim.} \textit{scil.} Dimanche.}} \\
      \leavevmode\fontsize{6.6pt}{7.5pt}\selectfont\strut Couché s/ une mauvaise paillasse tte\strut & \smash{\parbox[t]{39mm}{\fontsize{6.2pt}{7.0pt}\selectfont\raggedright\color{black!88}\textbf{2 s/} \textit{abbr. pro} sur. \textbf{2 tte} \textit{abbr. pro} toute.}} \\
      \leavevmode\fontsize{6.6pt}{7.5pt}\selectfont\strut crasseuse – Rats, sarabande au-dessus de ma\strut &  \\
      \leavevmode\fontsize{6.6pt}{7.5pt}\selectfont\strut tête et m’ont tenu éveillé partie nuit\strut &  \\
      \leavevmode\llap{\fontsize{6.4pt}{7.3pt}\selectfont\color{darkgray}5\hspace{1.8mm}}\fontsize{6.6pt}{7.5pt}\selectfont\strut Après breakfast, venu à pied à travers les\strut & \smash{\parbox[t]{39mm}{\fontsize{6.2pt}{7.0pt}\selectfont\raggedright\color{black!88}\textbf{5 breakfast} \textit{angl. in ms.}}} \\
      \leavevmode\fontsize{6.6pt}{7.5pt}\selectfont\strut houblonnières au village de St-Jans Cappel.\strut & \smash{\parbox[t]{39mm}{\fontsize{6.2pt}{7.0pt}\selectfont\raggedright\color{black!88}\textbf{6 St} \textit{abbr. pro} Saint.}} \\
      \leavevmode\fontsize{6.6pt}{7.5pt}\selectfont\strut Eglise archi-comble. Hoes à d., fem. à g.\strut &  \\
      \leavevmode\fontsize{6.6pt}{7.5pt}\selectfont\strut Prêtre en chaire. Assis s/ un petit pont ruisseau\strut & \smash{\parbox[t]{39mm}{\fontsize{6.2pt}{7.0pt}\selectfont\raggedright\color{black!88}\textbf{8 s/} \textit{abbr. pro} sur.}} \\
      \leavevmode\fontsize{6.6pt}{7.5pt}\selectfont\strut presque en face et lu, ruminé à mon aise\strut &  \\
      \leavevmode\llap{\fontsize{6.4pt}{7.3pt}\selectfont\color{darkgray}10\hspace{1.8mm}}\fontsize{6.6pt}{7.5pt}\selectfont\strut Resté grenier, presque tt l’après-midi, dormi,\strut & \smash{\parbox[t]{39mm}{\fontsize{6.2pt}{7.0pt}\selectfont\raggedright\color{black!88}\textbf{10 tt} \textit{abbr. pro} tout.}} \\
      \leavevmode\fontsize{6.6pt}{7.5pt}\selectfont\strut lu, et écrit –\strut &  \\
      \leavevmode\fontsize{6.6pt}{7.5pt}\selectfont\strut D’une fenêtre latérale, regardé la campagne\strut &  \\
      \leavevmode\fontsize{6.6pt}{7.5pt}\selectfont\strut – plaines fertiles, entrecoupées de boqueteaux\strut &  \\
      \leavevmode\fontsize{6.6pt}{7.5pt}\selectfont\strut et de houblonnières. Ciel peu à peu\strut &  \\
      \leavevmode\llap{\fontsize{6.4pt}{7.3pt}\selectfont\color{darkgray}15\hspace{1.8mm}}\fontsize{6.6pt}{7.5pt}\selectfont\strut estompé de rose –\strut &  \\
      \leavevmode\fontsize{6.6pt}{7.5pt}\selectfont\strut – Tip à Hancock, headcook : 5 f –\strut & \smash{\parbox[t]{39mm}{\fontsize{6.2pt}{7.0pt}\selectfont\raggedright\color{black!88}\textbf{16 headcook} \textit{angl. in ms.}}} \\
      \leavevmode\fontsize{6.6pt}{7.5pt}\selectfont\strut – A fermière, pour lait \& bière : 75c\strut &  \\
      \cline{1-1}
      \leavevmode\fontsize{6.6pt}{7.5pt}\selectfont\strut \underline{Lundi 31.} Chaleur accablante – Air irrespirable\strut &  \\
      \leavevmode\fontsize{6.6pt}{7.5pt}\selectfont\strut Matin, en auto avec mes off. : au\strut & \smash{\parbox[t]{39mm}{\fontsize{6.2pt}{7.0pt}\selectfont\raggedright\color{black!88}\textbf{19 off.} \textit{abbr. pro} officier(s). — \textbf{20 refilling point} \textit{angl. in ms.} — \textbf{21} \textit{lect. incert. in ms.}}} \\
      \leavevmode\llap{\fontsize{6.4pt}{7.3pt}\selectfont\color{darkgray}20\hspace{1.8mm}}\fontsize{6.6pt}{7.5pt}\selectfont\strut refilling point, puis à Bailleul, enfin\strut &  \\
      \leavevmode\fontsize{6.6pt}{7.5pt}\selectfont\strut chez fermier [Terrien ?] –\strut &  \\
      \leavevmode\fontsize{6.6pt}{7.5pt}\selectfont\strut Après-midi, en plein soleil cuisant ai\strut &  \\
      \leavevmode\fontsize{6.6pt}{7.5pt}\selectfont\strut dû conduire 1 waggon to the fields\strut & \smash{\parbox[t]{39mm}{\fontsize{6.2pt}{7.0pt}\selectfont\raggedright\color{black!88}\textbf{23 waggon} \textit{angl. in ms.}}} \\
      \leavevmode\fontsize{6.6pt}{7.5pt}\selectfont\strut for green forage. Chaleur extrême\strut & \smash{\parbox[t]{39mm}{\fontsize{6.2pt}{7.0pt}\selectfont\raggedright\color{black!88}\textbf{24 green} \textit{angl. in ms.}}} \\
      \leavevmode\llap{\fontsize{6.4pt}{7.3pt}\selectfont\color{darkgray}25\hspace{1.8mm}}\fontsize{6.6pt}{7.5pt}\selectfont\strut Fin jour, promené et assis champs attenant\strut &  \\
      \leavevmode\fontsize{6.6pt}{7.5pt}\selectfont\strut à ma ferme – Sur crête boisée, moulins\strut &  \\
      \leavevmode\fontsize{6.6pt}{7.5pt}\selectfont\strut à vent, tournent sans cesse\strut &  \\
\end{tabular}}%
\par\vspace*{\fill}
% --- Registres de bas de page (1. Critique interne [si pertinent] ; 2. Loci paralleli en latin ; 3. Commentaire historique en retrait) ---
\noindent\rule{\textwidth}{0.5pt}\par\vspace{0.65mm}
{\fontsize{6.6pt}{7.7pt}\selectfont%
\noindent\normalfont\color{black!92} \textbf{20} (\textit{Bailleul}) cf. infra f\textsuperscript{o}~4~r\textsuperscript{o} (p.~\pageref{fol:8}), f\textsuperscript{o}~4~v\textsuperscript{o} (p.~\pageref{fol:9}), \textit{et passim}. — \textbf{24} (\textit{green forage}) cf. supra f\textsuperscript{o}~3~r\textsuperscript{o} (p.~\pageref{fol:6}), infra f\textsuperscript{o}~4~v\textsuperscript{o} (p.~\pageref{fol:9}), f\textsuperscript{o}~6~v\textsuperscript{o} (p.~\pageref{fol:13}), \textit{et passim}.\par\vspace{0.55mm}
\noindent\rule{38mm}{0.4pt}\par\vspace{0.5mm}
\noindent\hspace*{5mm}\parbox{\dimexpr\textwidth-5mm\relax}{\fontsize{6.6pt}{7.7pt}\selectfont\color{black!84}Le dimanche 30 juillet, Hippolyte Combes note la séparation traditionnelle des hommes et des femmes dans la nef de l'église \textbf{Saint-Jean-Baptiste de Saint-Jans-Cappel} (dont le clocher effilé était signalé au f\textsuperscript{o}~3~r\textsuperscript{o}) et remet 5~fr. de pourboire au maître-cuisinier (\textit{headcook}) \textbf{Hancock}.}\par}
\newpage
\phantomsection\label{fac:7}
% --- En-tête du fac-similé (Page impaire / droite : Photographie du manuscrit) ---
\noindent\rule{\textwidth}{0.5pt}\par\vspace{0.55mm}
\noindent\textbf{\fontsize{6.8pt}{7.8pt}\selectfont FAC-SIMILÉ EN VIS-À-VIS — FOLIO 3 VERSO} \hfill \textbf{\fontsize{6.8pt}{7.8pt}\selectfont JUILLET 1916}\par\vspace{0.25mm}
\noindent{\fontsize{6.8pt}{8.0pt}\selectfont\textit{Livre-Journal d'Hippolyte Combes} — Dimanche 30 et lundi 31 juillet 1916 (Saint-Jans-Cappel et Bailleul (Nord))}\par\vspace{0.55mm}
\noindent\rule{\textwidth}{0.35pt}\par\vspace*{\fill}
{\centering
  \setlength{\fboxsep}{1.2mm}%
  \setlength{\fboxrule}{0.4pt}%
  \fcolorbox{gray!55}{white}{%
    \includegraphics[width=0.95\textwidth,height=0.81\textheight,keepaspectratio]{\facsimiledir/20251126_093219.jpg}%
  }\par}
\vspace*{\fill}
\newpage

% =============================================================================
% VIS-A-VIS SPREAD : FOLIO 4 RECTO (20251126_093240.jpg)
% =============================================================================
\cleardoubleevenpage
\renewcommand{\currentmonthyear}{Août 1916}
\renewcommand{\currentfoliolabel}{Folio 4 recto}
\phantomsection\label{fol:8}
% --- En-tête du folio (Page paire / gauche : Transcription, apparat & registres de bas de page) ---
\noindent\rule{\textwidth}{0.5pt}\par\vspace{0.55mm}
\noindent\textbf{\footnotesize FOLIO 4 RECTO} \hfill \textbf{\footnotesize AOÛT 1916}\par\vspace{0.25mm}
\noindent{\fontsize{6.9pt}{8.0pt}\selectfont Première promenade à Bailleul : la Grand-Place et le beffroi}\par\vspace{0.2mm}
\noindent{\fontsize{6.9pt}{8.0pt}\selectfont\textit{Mardi 1er août 1916} \hfill \textit{Saint-Jans-Cappel et Bailleul (Nord)}}\par\vspace{0.5mm}
\noindent\rule{\textwidth}{0.35pt}\par\vspace{0.8mm}
% --- Corps de la transcription diplomatique et apparat marginal en prolongement ---
{\renewcommand{\arraystretch}{0.87}%
\noindent\hspace*{3.5mm}\begin{tabular}{@{}p{87mm}@{\hspace{3.5mm}}p{39mm}@{}}
      \leavevmode\fontsize{6.9pt}{7.9pt}\selectfont\strut \underline{Mar 1er} \hfill{} \underline{Aout 1916 –}\strut & \smash{\parbox[t]{39mm}{\fontsize{6.2pt}{7.0pt}\selectfont\raggedright\color{black!88}\textbf{1 Mar} \textit{scil.} Mardi. \textbf{1 1er} \textit{scil.} premier. — \textbf{2 tempérre} \textit{abbr. pro} température. — \textbf{3 1/2} \textit{scil.} demi / moitié.}} \\
      \leavevmode\fontsize{6.9pt}{7.9pt}\selectfont\strut Tjours même tempérre asphyxiante. Eté\strut &  \\
      \leavevmode\fontsize{6.9pt}{7.9pt}\selectfont\strut torride 1/2 jour, impossible s'aventurer sur\strut &  \\
      \leavevmode\fontsize{6.9pt}{7.9pt}\selectfont\strut les routes\strut &  \\
      \leavevmode\llap{\fontsize{6.4pt}{7.3pt}\selectfont\color{darkgray}5\hspace{1.8mm}}\fontsize{6.9pt}{7.9pt}\selectfont\strut Liberté entière tt le jour. Resté ds mon\strut & \smash{\parbox[t]{39mm}{\fontsize{6.2pt}{7.0pt}\selectfont\raggedright\color{black!88}\textbf{5 ds} \textit{abbr. pro} dans. \textbf{5 tt} \textit{abbr. pro} tout. — \textbf{6 garret} \textit{angl. in ms.}}} \\
      \leavevmode\fontsize{6.9pt}{7.9pt}\selectfont\strut garret, où ménagé installation sommaire.\strut &  \\
      \leavevmode\fontsize{6.9pt}{7.9pt}\selectfont\strut Après-midi, promenade à Bailleul.\strut &  \\
      \leavevmode\fontsize{6.9pt}{7.9pt}\selectfont\strut Parti à 4h. Mis 30 m. pour parcourir distance\strut &  \\
      \leavevmode\fontsize{6.9pt}{7.9pt}\selectfont\strut Petite ville. Place assez semblable à celle d'Hazebrouck\strut &  \\
      \leavevmode\llap{\fontsize{6.4pt}{7.3pt}\selectfont\color{darkgray}10\hspace{1.8mm}}\fontsize{6.9pt}{7.9pt}\selectfont\strut \& de Béthune Parcouru rues principales et\strut &  \\
      \leavevmode\fontsize{6.9pt}{7.9pt}\selectfont\strut regardé longuement église \& hotel de ville\strut &  \\
      \leavevmode\fontsize{6.9pt}{7.9pt}\selectfont\strut Retour vers 6 H. — Chaleur accablante —\strut &  \\
      \cline{1-1}
      \leavevmode\fontsize{6.9pt}{7.9pt}\selectfont\strut Payé à fermière : lait et bière : hier : 20c\strut &  \\
      \leavevmode\fontsize{6.9pt}{7.9pt}\selectfont\strut et aujourd’hui : 20 c. + 10 c – + 10c\strut & \smash{\parbox[t]{39mm}{\fontsize{6.2pt}{7.0pt}\selectfont\raggedright\color{black!88}\textbf{14 c.} \textit{abbr. pro} centimes.}} \\
      \cline{1-1}
      \leavevmode\llap{\fontsize{6.4pt}{7.3pt}\selectfont\color{darkgray}15\hspace{1.8mm}}\fontsize{6.9pt}{7.9pt}\selectfont\strut \underline{A Bailleul} : cartes ples : 50 c., 60 c., 0.10 varia\strut & \smash{\parbox[t]{39mm}{\fontsize{6.2pt}{7.0pt}\selectfont\raggedright\color{black!88}\textbf{15 cartes ples} \textit{abbr. pro} cartes postales. \textbf{15 c.} \textit{abbr. pro} centimes. — \textbf{16} \textit{lect. incert. in ms.}}} \\
      \leavevmode\fontsize{6.9pt}{7.9pt}\selectfont\strut " : 4 cartes cuivre M. [Guétin ?] : 1 f 20 c\strut &  \\
\end{tabular}}%
\par\vspace*{\fill}
% --- Registres de bas de page (1. Critique interne [si pertinent] ; 2. Loci paralleli en latin ; 3. Commentaire historique en retrait) ---
\noindent\rule{\textwidth}{0.5pt}\par\vspace{0.65mm}
{\fontsize{6.6pt}{7.7pt}\selectfont%
\noindent\normalfont\color{black!92} \textbf{7} (\textit{Bailleul}) cf. supra f\textsuperscript{o}~3~v\textsuperscript{o} (p.~\pageref{fol:7}), infra f\textsuperscript{o}~4~v\textsuperscript{o} (p.~\pageref{fol:9}), f\textsuperscript{o}~5~r\textsuperscript{o} (p.~\pageref{fol:10}), \textit{et passim}. — \textbf{9} (\textit{Hazebrouck}) cf. supra f\textsuperscript{o}~2~r\textsuperscript{o} (p.~\pageref{fol:4}), f\textsuperscript{o}~2~v\textsuperscript{o} (p.~\pageref{fol:5}), infra f\textsuperscript{o}~8~r\textsuperscript{o} (p.~\pageref{fol:16}), f\textsuperscript{o}~13~v\textsuperscript{o} (p.~\pageref{fol:27}).\par\vspace{0.55mm}
\noindent\rule{38mm}{0.4pt}\par\vspace{0.5mm}
\noindent\hspace*{5mm}\parbox{\dimexpr\textwidth-5mm\relax}{\fontsize{6.6pt}{7.7pt}\selectfont\color{black!84}Située à 3~km au sud de Saint-Jans-Cappel (trente minutes de marche), la ville de \textbf{Bailleul} (Nord) s'ordonnait autour de sa vaste Grand-Place flamande dominée par l'église \textbf{Saint-Vaast} et l'\textbf{Hôtel de Ville} à beffroi gothique. La comparaison autographe avec \textbf{Béthune} (Pas-de-Calais) recoupe le séjour antérieur d'Hippolyte Combes dans le secteur de Béthune attesté au f\textsuperscript{o}~46~r\textsuperscript{o} (l.~11--12).}\par}
\newpage
\phantomsection\label{fac:8}
% --- En-tête du fac-similé (Page impaire / droite : Photographie du manuscrit) ---
\noindent\rule{\textwidth}{0.5pt}\par\vspace{0.55mm}
\noindent\textbf{\fontsize{6.8pt}{7.8pt}\selectfont FAC-SIMILÉ EN VIS-À-VIS — FOLIO 4 RECTO} \hfill \textbf{\fontsize{6.8pt}{7.8pt}\selectfont AOÛT 1916}\par\vspace{0.25mm}
\noindent{\fontsize{6.8pt}{8.0pt}\selectfont\textit{Livre-Journal d'Hippolyte Combes} — Mardi 1er août 1916 (Saint-Jans-Cappel et Bailleul (Nord))}\par\vspace{0.55mm}
\noindent\rule{\textwidth}{0.35pt}\par\vspace*{\fill}
{\centering
  \setlength{\fboxsep}{1.2mm}%
  \setlength{\fboxrule}{0.4pt}%
  \fcolorbox{gray!55}{white}{%
    \includegraphics[width=0.95\textwidth,height=0.81\textheight,keepaspectratio]{\facsimiledir/20251126_093240.jpg}%
  }\par}
\vspace*{\fill}
\newpage

% =============================================================================
% VIS-A-VIS SPREAD : FOLIO 4 VERSO (20251126_093301.jpg)
% =============================================================================
\cleardoubleevenpage
\renewcommand{\currentmonthyear}{Août 1916}
\renewcommand{\currentfoliolabel}{Folio 4 verso}
\phantomsection\label{fol:9}
% --- En-tête du folio (Page paire / gauche : Transcription, apparat & registres de bas de page) ---
\noindent\rule{\textwidth}{0.5pt}\par\vspace{0.55mm}
\noindent\textbf{\footnotesize FOLIO 4 VERSO} \hfill \textbf{\footnotesize AOÛT 1916}\par\vspace{0.25mm}
\noindent{\fontsize{6.9pt}{8.0pt}\selectfont Mercredi 2 et jeudi 3 août 1916 — Promenades dans Bailleul, maisons flamandes et dentellières au fuseau}\par\vspace{0.2mm}
\noindent{\fontsize{6.9pt}{8.0pt}\selectfont\textit{Mercredi 2 et jeudi 3 août 1916} \hfill \textit{Bailleul (Nord)}}\par\vspace{0.5mm}
\noindent\rule{\textwidth}{0.35pt}\par\vspace{0.8mm}
% --- Corps de la transcription diplomatique et apparat marginal en prolongement ---
{\renewcommand{\arraystretch}{0.87}%
\noindent\hspace*{3.5mm}\begin{tabular}{@{}p{87mm}@{\hspace{3.5mm}}p{39mm}@{}}
      \leavevmode\fontsize{6.9pt}{7.9pt}\selectfont\strut \underline{2 Me} – Encore très b. journée. \hfill{} \raisebox{0.5ex}{\underline{Aout / 16}}\strut &  \\
      \leavevmode\fontsize{6.9pt}{7.9pt}\selectfont\strut Matinée entière libre. Passé près de 2 h. assis\strut &  \\
      \leavevmode\fontsize{6.9pt}{7.9pt}\selectfont\strut sur l'herbe ds le camp\strut & \smash{\parbox[t]{39mm}{\fontsize{6.2pt}{7.0pt}\selectfont\raggedright\color{black!88}\textbf{3 ds} \textit{abbr. pro} dans.}} \\
      \leavevmode\fontsize{6.9pt}{7.9pt}\selectfont\strut Après-midi, dû conduire 1 wagon au\strut &  \\
      \leavevmode\llap{\fontsize{6.4pt}{7.3pt}\selectfont\color{darkgray}5\hspace{1.8mm}}\fontsize{6.9pt}{7.9pt}\selectfont\strut fourrage de là, continué à pied j'à Bailleul.\strut & \smash{\parbox[t]{39mm}{\fontsize{6.2pt}{7.0pt}\selectfont\raggedright\color{black!88}\textbf{5 j'à} \textit{abbr. pro} jusqu'à.}} \\
      \leavevmode\fontsize{6.9pt}{7.9pt}\selectfont\strut Visité de nouveau et à loisir cette petite ville\strut &  \\
      \leavevmode\fontsize{6.9pt}{7.9pt}\selectfont\strut égl. St-Vaast. Place et hotel de ville\strut & \smash{\parbox[t]{39mm}{\fontsize{6.2pt}{7.0pt}\selectfont\raggedright\color{black!88}\textbf{7 égl.} \textit{abbr. pro} église. \textbf{7 St} \textit{abbr. pro} Saint. — \textbf{8 longt} \textit{abbr. pro} longtemps.}} \\
      \leavevmode\fontsize{6.9pt}{7.9pt}\selectfont\strut Au retour, me suis longt arrêté devant dentellières\strut &  \\
      \leavevmode\fontsize{6.9pt}{7.9pt}\selectfont\strut travaillant rue devant maisons. Très intéressé\strut &  \\
      \leavevmode\llap{\fontsize{6.4pt}{7.3pt}\selectfont\color{darkgray}10\hspace{1.8mm}}\fontsize{6.9pt}{7.9pt}\selectfont\strut et émerveillé par leur ouvrage $[\![$[illis.]$]\!]$\strut & \smash{\parbox[t]{39mm}{\fontsize{6.2pt}{7.0pt}\selectfont\raggedright\color{black!88}\textbf{10} $[\![$[illis.]$]\!]\ $ \textit{del. in ms.} \textbf{10} \textit{lect. incert. in ms.} — \textbf{11 ctes ples} \textit{abbr. pro} cartes postales. \textbf{11 c.} \textit{abbr. pro} centimes. — \textbf{12 pr} \textit{abbr. pro} pour. \textbf{12 c.} \textit{abbr. pro} centimes. — \textbf{13 tempér.} \textit{abbr. pro} température. — \textbf{18 ds} \textit{abbr. pro} dans. \textbf{18 ts} \textit{abbr. pro} tous.}} \\
      \leavevmode\fontsize{6.9pt}{7.9pt}\selectfont\strut — A Bailleul : ctes ples : 20 c. — Chemise : 5 f 50\strut &  \\
      \leavevmode\fontsize{6.9pt}{7.9pt}\selectfont\strut 2m 40 de galon noir pr lorgnons : 35c.. Drink : 10, 10c\strut &  \\
      \cline{1-1}
      \leavevmode\fontsize{6.9pt}{7.9pt}\selectfont\strut \underline{3 J} – \hfill{} Même tempér., vraiment estivale\strut &  \\
      \leavevmode\fontsize{6.9pt}{7.9pt}\selectfont\strut Après-midi, venu de nouveau promener\strut &  \\
      \leavevmode\llap{\fontsize{6.4pt}{7.3pt}\selectfont\color{darkgray}15\hspace{1.8mm}}\fontsize{6.9pt}{7.9pt}\selectfont\strut à pied à Bailleul. Arrêts fréquents dans\strut &  \\
      \leavevmode\fontsize{6.9pt}{7.9pt}\selectfont\strut les coins d'ombre le long des champs d'avoine\strut &  \\
      \leavevmode\fontsize{6.9pt}{7.9pt}\selectfont\strut et de blé.\strut &  \\
      \leavevmode\fontsize{6.9pt}{7.9pt}\selectfont\strut Parcouru ds ts les sens ville de B.\strut &  \\
      \leavevmode\fontsize{6.9pt}{7.9pt}\selectfont\strut remarqué belles maisons rue d’Ypres. rue\strut &  \\
      \leavevmode\llap{\fontsize{6.4pt}{7.3pt}\selectfont\color{darkgray}20\hspace{1.8mm}}\fontsize{6.9pt}{7.9pt}\selectfont\strut de Lille. etc\strut &  \\
      \leavevmode\fontsize{6.9pt}{7.9pt}\selectfont\strut Entré pour me reposer à St Vaast\strut & \smash{\parbox[t]{39mm}{\fontsize{6.2pt}{7.0pt}\selectfont\raggedright\color{black!88}\textbf{21 St} \textit{abbr. pro} Saint.}} \\
      \leavevmode\fontsize{6.9pt}{7.9pt}\selectfont\strut Au retour, assis au fond jardin ferme où loge\strut &  \\
      \leavevmode\fontsize{6.9pt}{7.9pt}\selectfont\strut — Ctes ples : 40 c. — Drink : 10 c —\strut & \smash{\parbox[t]{39mm}{\fontsize{6.2pt}{7.0pt}\selectfont\raggedright\color{black!88}\textbf{23 Ctes ples} \textit{abbr. pro} Cartes postales. \textbf{23 c.} \textit{abbr. pro} centimes.}} \\
\end{tabular}}%
\par\vspace*{\fill}
% --- Registres de bas de page (1. Critique interne [si pertinent] ; 2. Loci paralleli en latin ; 3. Commentaire historique en retrait) ---
\noindent\rule{\textwidth}{0.5pt}\par\vspace{0.65mm}
{\fontsize{6.6pt}{7.7pt}\selectfont%
\noindent\normalfont \textbf{Unité de page :} Ce feuillet réunit deux journées consécutives (mercredi 2 août, l.~1--11, et jeudi 3 août, l.~12--19) séparées par un trait horizontal à mi-page, suivies d'un relevé de dépenses au bas du folio.\par\vspace{0.5mm}
\noindent\rule{22mm}{0.3pt}\par\vspace{0.45mm}
\noindent\normalfont\color{black!92} \textbf{5} (\textit{Bailleul}) cf. supra f\textsuperscript{o}~3~v\textsuperscript{o} (p.~\pageref{fol:7}), f\textsuperscript{o}~4~r\textsuperscript{o} (p.~\pageref{fol:8}), infra f\textsuperscript{o}~5~r\textsuperscript{o} (p.~\pageref{fol:10}), f\textsuperscript{o}~6~r\textsuperscript{o} (p.~\pageref{fol:12}), \textit{et passim}. — \textbf{5} (\textit{fourrage}) cf. supra f\textsuperscript{o}~3~r\textsuperscript{o} (p.~\pageref{fol:6}), f\textsuperscript{o}~3~v\textsuperscript{o} (p.~\pageref{fol:7}), infra f\textsuperscript{o}~6~v\textsuperscript{o} (p.~\pageref{fol:13}), f\textsuperscript{o}~11~v\textsuperscript{o} (p.~\pageref{fol:23}), \textit{et passim}. — \textbf{19} (\textit{Ypres}) cf. supra f\textsuperscript{o}~2~r\textsuperscript{o} (p.~\pageref{fol:4}).\par\vspace{0.55mm}
\noindent\rule{38mm}{0.4pt}\par\vspace{0.5mm}
\noindent\hspace*{5mm}\parbox{\dimexpr\textwidth-5mm\relax}{\fontsize{6.6pt}{7.7pt}\selectfont\color{black!84}Nouvelles promenades dans \textbf{Bailleul} (rues d'Ypres et de Lille, église Saint-Vaast). L'observation des \textbf{dentellières} travaillant en plein air devant leurs portes témoigne de l'industrie dentellière traditionnelle de Bailleul (école et fabriques de dentelle aux fuseaux actives depuis le XVII\textsuperscript{e}~s. jusqu'à la destruction de la ville au printemps 1918).}\par}
\newpage
\phantomsection\label{fac:9}
% --- En-tête du fac-similé (Page impaire / droite : Photographie du manuscrit) ---
\noindent\rule{\textwidth}{0.5pt}\par\vspace{0.55mm}
\noindent\textbf{\fontsize{6.8pt}{7.8pt}\selectfont FAC-SIMILÉ EN VIS-À-VIS — FOLIO 4 VERSO} \hfill \textbf{\fontsize{6.8pt}{7.8pt}\selectfont AOÛT 1916}\par\vspace{0.25mm}
\noindent{\fontsize{6.8pt}{8.0pt}\selectfont\textit{Livre-Journal d'Hippolyte Combes} — Mercredi 2 et jeudi 3 août 1916 (Bailleul (Nord))}\par\vspace{0.55mm}
\noindent\rule{\textwidth}{0.35pt}\par\vspace*{\fill}
{\centering
  \setlength{\fboxsep}{1.2mm}%
  \setlength{\fboxrule}{0.4pt}%
  \fcolorbox{gray!55}{white}{%
    \includegraphics[width=0.95\textwidth,height=0.81\textheight,keepaspectratio]{\facsimiledir/20251126_093301.jpg}%
  }\par}
\vspace*{\fill}
\newpage

% =============================================================================
% VIS-A-VIS SPREAD : FOLIO 5 RECTO (20251126_093316.jpg)
% =============================================================================
\cleardoubleevenpage
\renewcommand{\currentmonthyear}{Août 1916}
\renewcommand{\currentfoliolabel}{Folio 5 recto}
\phantomsection\label{fol:10}
% --- En-tête du folio (Page paire / gauche : Transcription, apparat & registres de bas de page) ---
\noindent\rule{\textwidth}{0.5pt}\par\vspace{0.55mm}
\noindent\textbf{\footnotesize FOLIO 5 RECTO} \hfill \textbf{\footnotesize AOÛT 1916}\par\vspace{0.25mm}
\noindent{\fontsize{6.9pt}{8.0pt}\selectfont Chevauchée à La Croix-du-Bac par Steenwerck et repos dans le « garret »}\par\vspace{0.2mm}
\noindent{\fontsize{6.9pt}{8.0pt}\selectfont\textit{Vendredi 4 août 1916 (matin et après-midi)} \hfill \textit{Saint-Jans-Cappel, Bailleul, Steenwerck, La Croix-du-Bac (Nord)}}\par\vspace{0.5mm}
\noindent\rule{\textwidth}{0.35pt}\par\vspace{0.8mm}
% --- Corps de la transcription diplomatique et apparat marginal en prolongement ---
{\renewcommand{\arraystretch}{0.87}%
\noindent\hspace*{3.5mm}\begin{tabular}{@{}p{87mm}@{\hspace{3.5mm}}p{39mm}@{}}
      \leavevmode\fontsize{6.9pt}{7.9pt}\selectfont\strut \underline{4 Ve.} Après qq hésitation, décidé \hfill{} \raisebox{0.5ex}{\underline{Août 1916}}\strut & \smash{\parbox[t]{39mm}{\fontsize{6.2pt}{7.0pt}\selectfont\raggedright\color{black!88}\textbf{1 Ve.} \textit{scil.} Vendredi. \textbf{1 qq} \textit{abbr. pro} quelque(s).}} \\
      \leavevmode\fontsize{6.9pt}{7.9pt}\selectfont\strut me rendre ce matin à La Croix-du-Bac, où bureau\strut &  \\
      \leavevmode\fontsize{6.9pt}{7.9pt}\selectfont\strut mission fse. Parti à cheval vers 10 H. Mis près\strut & \smash{\parbox[t]{39mm}{\fontsize{6.2pt}{7.0pt}\selectfont\raggedright\color{black!88}\textbf{3 fse} \textit{abbr. pro} française.}} \\
      \leavevmode\fontsize{6.9pt}{7.9pt}\selectfont\strut d’ 1 h. 1/2 pour parcourir distance 15 Kils\strut & \smash{\parbox[t]{39mm}{\fontsize{6.2pt}{7.0pt}\selectfont\raggedright\color{black!88}\textbf{4 Kils} \textit{abbr. pro} Kilomètres. \textbf{4 1/2} \textit{scil.} demi / moitié.}} \\
      \leavevmode\llap{\fontsize{6.4pt}{7.3pt}\selectfont\color{darkgray}5\hspace{1.8mm}}\fontsize{6.9pt}{7.9pt}\selectfont\strut environ. Passé par Bailleul et Steenwerck.\strut &  \\
      \leavevmode\fontsize{6.9pt}{7.9pt}\selectfont\strut Cette dernière localité, gros village — autour\strut &  \\
      \leavevmode\fontsize{6.9pt}{7.9pt}\selectfont\strut église très gde —\strut & \smash{\parbox[t]{39mm}{\fontsize{6.2pt}{7.0pt}\selectfont\raggedright\color{black!88}\textbf{7 gde} \textit{abbr. pro} grande.}} \\
      \leavevmode\fontsize{6.9pt}{7.9pt}\selectfont\strut Ciel d’abord menaçant. qq gouttes de pluie\strut & \smash{\parbox[t]{39mm}{\fontsize{6.2pt}{7.0pt}\selectfont\raggedright\color{black!88}\textbf{8 qq} \textit{abbr. pro} quelque(s).}} \\
      \leavevmode\fontsize{6.9pt}{7.9pt}\selectfont\strut au retour. Puis vent dissipe nuages et soleil\strut &  \\
      \leavevmode\llap{\fontsize{6.4pt}{7.3pt}\selectfont\color{darkgray}10\hspace{1.8mm}}\fontsize{6.9pt}{7.9pt}\selectfont\strut ardent des jours précédents\strut &  \\
      \leavevmode\fontsize{6.9pt}{7.9pt}\selectfont\strut Tout l’après-midi et fin journée, reposé garret.\strut & \smash{\parbox[t]{39mm}{\fontsize{6.2pt}{7.0pt}\selectfont\raggedright\color{black!88}\textbf{11 garret} \textit{angl. in ms.}}} \\
      \cline{1-1}
      \leavevmode\fontsize{6.9pt}{7.9pt}\selectfont\strut [ — Touché au bureau mission, à\strut & \smash{\parbox[t]{39mm}{\fontsize{6.2pt}{7.0pt}\selectfont\raggedright\color{black!88}\textbf{12} \textit{in marg. cum unco} (comptes).}} \\
      \leavevmode\fontsize{6.9pt}{7.9pt}\selectfont\strut La Croix-du-Bac : 80 f 30 c (y compris 8 f\strut &  \\
      \leavevmode\fontsize{6.9pt}{7.9pt}\selectfont\strut supplément 1er juillet) —\strut & \smash{\parbox[t]{39mm}{\fontsize{6.2pt}{7.0pt}\selectfont\raggedright\color{black!88}\textbf{14 1er} \textit{scil.} premier.}} \\
      \leavevmode\llap{\fontsize{6.4pt}{7.3pt}\selectfont\color{darkgray}15\hspace{1.8mm}}\fontsize{6.9pt}{7.9pt}\selectfont\strut Ferme autour de l. notre coy est cantonnée,\strut &  \\
      \leavevmode\fontsize{6.9pt}{7.9pt}\selectfont\strut isolée en rase campagne. Y gd calme et\strut & \smash{\parbox[t]{39mm}{\fontsize{6.2pt}{7.0pt}\selectfont\raggedright\color{black!88}\textbf{16 gd} \textit{abbr. pro} grand.}} \\
      \leavevmode\fontsize{6.9pt}{7.9pt}\selectfont\strut silence —\strut &  \\
\end{tabular}}%
\par\vspace*{\fill}
% --- Registres de bas de page (1. Critique interne [si pertinent] ; 2. Loci paralleli en latin ; 3. Commentaire historique en retrait) ---
\noindent\rule{\textwidth}{0.5pt}\par\vspace{0.65mm}
{\fontsize{6.6pt}{7.7pt}\selectfont%
\noindent\normalfont\color{black!92} \textbf{2} (\textit{La Croix-du-Bac}) cf. infra f\textsuperscript{o}~6~r\textsuperscript{o} (p.~\pageref{fol:12}), f\textsuperscript{o}~6~v\textsuperscript{o} (p.~\pageref{fol:13}). — \textbf{5} (\textit{Bailleul}) cf. supra f\textsuperscript{o}~3~v\textsuperscript{o} (p.~\pageref{fol:7}), f\textsuperscript{o}~4~v\textsuperscript{o} (p.~\pageref{fol:9}), infra f\textsuperscript{o}~6~r\textsuperscript{o} (p.~\pageref{fol:12}), f\textsuperscript{o}~13~v\textsuperscript{o} (p.~\pageref{fol:27}), \textit{et passim}. — \textbf{5} (\textit{Steenwerck}) cf. infra f\textsuperscript{o}~6~r\textsuperscript{o} (p.~\pageref{fol:12}), f\textsuperscript{o}~7~r\textsuperscript{o} (p.~\pageref{fol:14}), \textit{et passim}.\par\vspace{0.55mm}
\noindent\rule{38mm}{0.4pt}\par\vspace{0.5mm}
\noindent\hspace*{5mm}\parbox{\dimexpr\textwidth-5mm\relax}{\fontsize{6.6pt}{7.7pt}\selectfont\color{black!84}Chevauchée de 15~km par Bailleul et \textbf{Steenwerck} (dont la grande église paroissiale Saint-Jean-Baptiste est signalée aux l.~6--7) jusqu'au hameau de \textbf{La Croix-du-Bac} (commune de Steenwerck, Nord), siège du bureau de la Mission militaire française auprès de la 18\textsuperscript{e} Division, où Hippolyte Combes perçoit sa solde et un rappel de supplément à compter du 1\textsuperscript{er} juillet 1916.}\par}
\newpage
\phantomsection\label{fac:10}
% --- En-tête du fac-similé (Page impaire / droite : Photographie du manuscrit) ---
\noindent\rule{\textwidth}{0.5pt}\par\vspace{0.55mm}
\noindent\textbf{\fontsize{6.8pt}{7.8pt}\selectfont FAC-SIMILÉ EN VIS-À-VIS — FOLIO 5 RECTO} \hfill \textbf{\fontsize{6.8pt}{7.8pt}\selectfont AOÛT 1916}\par\vspace{0.25mm}
\noindent{\fontsize{6.8pt}{8.0pt}\selectfont\textit{Livre-Journal d'Hippolyte Combes} — Vendredi 4 août 1916 (matin et après-midi) (Saint-Jans-Cappel, Bailleul, Steenwerck, La Croix-du-Bac (Nord))}\par\vspace{0.55mm}
\noindent\rule{\textwidth}{0.35pt}\par\vspace*{\fill}
{\centering
  \setlength{\fboxsep}{1.2mm}%
  \setlength{\fboxrule}{0.4pt}%
  \fcolorbox{gray!55}{white}{%
    \includegraphics[width=0.95\textwidth,height=0.81\textheight,keepaspectratio]{\facsimiledir/20251126_093316.jpg}%
  }\par}
\vspace*{\fill}
\newpage

% =============================================================================
% VIS-A-VIS SPREAD : FOLIO 5 VERSO (20251126_093332.jpg)
% =============================================================================
\cleardoubleevenpage
\renewcommand{\currentmonthyear}{Août 1916}
\renewcommand{\currentfoliolabel}{Folio 5 verso}
\phantomsection\label{fol:11}
% --- En-tête du folio (Page paire / gauche : Transcription, apparat & registres de bas de page) ---
\noindent\rule{\textwidth}{0.5pt}\par\vspace{0.55mm}
\noindent\textbf{\footnotesize FOLIO 5 VERSO} \hfill \textbf{\footnotesize AOÛT 1916}\par\vspace{0.25mm}
\noindent{\fontsize{6.9pt}{8.0pt}\selectfont Soirée solitaire à la ferme et prière du soir de la famille flamande}\par\vspace{0.2mm}
\noindent{\fontsize{6.9pt}{8.0pt}\selectfont\textit{Vendredi 4 août 1916 (soir)} \hfill \textit{Saint-Jans-Cappel (Nord)}}\par\vspace{0.5mm}
\noindent\rule{\textwidth}{0.35pt}\par\vspace{0.8mm}
% --- Corps de la transcription diplomatique et apparat marginal en prolongement ---
{\renewcommand{\arraystretch}{1.00}%
\noindent\hspace*{3.5mm}\begin{tabular}{@{}p{87mm}@{\hspace{3.5mm}}p{39mm}@{}}
      \leavevmode\fontsize{7.7pt}{9.0pt}\selectfont\strut \underline{Ve 4} — Avant \& après diner, assis \&\strut &  \\
      \leavevmode\fontsize{7.7pt}{9.0pt}\selectfont\strut promené sur pâtures ferme où je loge\strut &  \\
      \leavevmode\fontsize{7.7pt}{9.0pt}\selectfont\strut Ciel nuageux. Vent bruissement à travers\strut &  \\
      \leavevmode\fontsize{7.7pt}{9.0pt}\selectfont\strut feuillages gds arbres \& me plaisais à l’écouter\strut & \smash{\parbox[t]{39mm}{\fontsize{6.2pt}{7.0pt}\selectfont\raggedright\color{black!88}\textbf{4 gds} \textit{abbr. pro} grands.}} \\
      \leavevmode\llap{\fontsize{6.4pt}{7.3pt}\selectfont\color{darkgray}5\hspace{1.8mm}}\fontsize{7.7pt}{9.0pt}\selectfont\strut Impression solitude \& d'abandon. Ns sommes\strut &  \\
      \leavevmode\fontsize{7.7pt}{9.0pt}\selectfont\strut loin de la ville — N'était le vert intense\strut &  \\
      \leavevmode\fontsize{7.7pt}{9.0pt}\selectfont\strut \& l’épaisseur des frondaisons, on se\strut &  \\
      \leavevmode\fontsize{7.7pt}{9.0pt}\selectfont\strut croirait déjà en automne\strut &  \\
      \leavevmode\fontsize{7.7pt}{9.0pt}\selectfont\strut De mon garret, (où suis remonté), la\strut & \smash{\parbox[t]{39mm}{\fontsize{6.2pt}{7.0pt}\selectfont\raggedright\color{black!88}\textbf{9 garret} \textit{angl. in ms.}}} \\
      \leavevmode\llap{\fontsize{6.4pt}{7.3pt}\selectfont\color{darkgray}10\hspace{1.8mm}}\fontsize{7.7pt}{9.0pt}\selectfont\strut nuit tombée, j’entends famille ferme\strut &  \\
      \leavevmode\fontsize{7.7pt}{9.0pt}\selectfont\strut en commun prière du soir : voix grave \&\strut &  \\
      \leavevmode\fontsize{7.7pt}{9.0pt}\selectfont\strut monotone du père ; [gémt ?] en [patois flamand ?]\strut & \smash{\parbox[t]{39mm}{\fontsize{6.2pt}{7.0pt}\selectfont\raggedright\color{black!88}\textbf{12} \textit{lect. incert. in ms.}}} \\
\end{tabular}}%
\par\vspace*{\fill}
% --- Registres de bas de page (1. Critique interne [si pertinent] ; 2. Loci paralleli en latin ; 3. Commentaire historique en retrait) ---
\noindent\rule{\textwidth}{0.5pt}\par\vspace{0.65mm}
{\fontsize{6.6pt}{7.7pt}\selectfont%
\noindent\normalfont\color{black!92} \textbf{1} (\textit{de eodem commorationis loco}) cf. supra f\textsuperscript{o}~5~r\textsuperscript{o} (p.~\pageref{fol:10}).\par\vspace{0.55mm}
\noindent\rule{38mm}{0.4pt}\par\vspace{0.5mm}
\noindent\hspace*{5mm}\parbox{\dimexpr\textwidth-5mm\relax}{\fontsize{6.6pt}{7.7pt}\selectfont\color{black!84}Dernière soirée à la ferme isolée de \textbf{Saint-Jans-Cappel} avant le déplacement du lendemain vers la vallée de la Lys: depuis sa mansarde (\textit{garret}), Hippolyte Combes écoute la prière vespérale récitée en commun par la famille d'agriculteurs flamands.}\par}
\newpage
\phantomsection\label{fac:11}
% --- En-tête du fac-similé (Page impaire / droite : Photographie du manuscrit) ---
\noindent\rule{\textwidth}{0.5pt}\par\vspace{0.55mm}
\noindent\textbf{\fontsize{6.8pt}{7.8pt}\selectfont FAC-SIMILÉ EN VIS-À-VIS — FOLIO 5 VERSO} \hfill \textbf{\fontsize{6.8pt}{7.8pt}\selectfont AOÛT 1916}\par\vspace{0.25mm}
\noindent{\fontsize{6.8pt}{8.0pt}\selectfont\textit{Livre-Journal d'Hippolyte Combes} — Vendredi 4 août 1916 (soir) (Saint-Jans-Cappel (Nord))}\par\vspace{0.55mm}
\noindent\rule{\textwidth}{0.35pt}\par\vspace*{\fill}
{\centering
  \setlength{\fboxsep}{1.2mm}%
  \setlength{\fboxrule}{0.4pt}%
  \fcolorbox{gray!55}{white}{%
    \includegraphics[width=0.95\textwidth,height=0.81\textheight,keepaspectratio]{\facsimiledir/20251126_093332.jpg}%
  }\par}
\vspace*{\fill}
\newpage

% =============================================================================
% VIS-A-VIS SPREAD : FOLIO 6 RECTO (20251126_093347.jpg)
% =============================================================================
\cleardoubleevenpage
\renewcommand{\currentmonthyear}{Août 1916}
\renewcommand{\currentfoliolabel}{Folio 6 recto}
\phantomsection\label{fol:12}
% --- En-tête du folio (Page paire / gauche : Transcription, apparat & registres de bas de page) ---
\noindent\rule{\textwidth}{0.5pt}\par\vspace{0.55mm}
\noindent\textbf{\footnotesize FOLIO 6 RECTO} \hfill \textbf{\footnotesize AOÛT 1916}\par\vspace{0.25mm}
\noindent{\fontsize{6.9pt}{8.0pt}\selectfont Départ du convoi avec les A.S.C. men et étape à La Croix-du-Bac}\par\vspace{0.2mm}
\noindent{\fontsize{6.9pt}{8.0pt}\selectfont\textit{Samedi 5 août 1916} \hfill \textit{De Saint-Jans-Cappel à La Croix-du-Bac / Steenwerck (Nord)}}\par\vspace{0.5mm}
\noindent\rule{\textwidth}{0.35pt}\par\vspace{0.8mm}
% --- Corps de la transcription diplomatique et apparat marginal en prolongement ---
{\renewcommand{\arraystretch}{0.87}%
\noindent\hspace*{3.5mm}\begin{tabular}{@{}p{87mm}@{\hspace{3.5mm}}p{39mm}@{}}
      \leavevmode\fontsize{6.9pt}{7.9pt}\selectfont\strut \underline{Sam. 5} \hfill{} \underline{Août 1916}\strut & \smash{\parbox[t]{39mm}{\fontsize{6.2pt}{7.0pt}\selectfont\raggedright\color{black!88}\textbf{1 Sam.} \textit{scil.} Samedi.}} \\
      \leavevmode\fontsize{6.9pt}{7.9pt}\selectfont\strut Levé vers 6H\strut &  \\
      \leavevmode\fontsize{6.9pt}{7.9pt}\selectfont\strut Bientôt, ts préparatifs terminés.\strut & \smash{\parbox[t]{39mm}{\fontsize{6.2pt}{7.0pt}\selectfont\raggedright\color{black!88}\textbf{3 ts} \textit{abbr. pro} tous.}} \\
      \leavevmode\fontsize{6.9pt}{7.9pt}\selectfont\strut Porté moi-m. en 3 voy., mes sacs à l'office\strut &  \\
      \leavevmode\llap{\fontsize{6.4pt}{7.3pt}\selectfont\color{darkgray}5\hspace{1.8mm}}\fontsize{6.9pt}{7.9pt}\selectfont\strut et les vu charger s/ wagon no 5 —\strut & \smash{\parbox[t]{39mm}{\fontsize{6.2pt}{7.0pt}\selectfont\raggedright\color{black!88}\textbf{5 s/} \textit{abbr. pro} sur.}} \\
      \leavevmode\fontsize{6.9pt}{7.9pt}\selectfont\strut Ciel tjs couvert de gros nuages noirs.\strut & \smash{\parbox[t]{39mm}{\fontsize{6.2pt}{7.0pt}\selectfont\raggedright\color{black!88}\textbf{6 tjs} \textit{abbr. pro} toujours.}} \\
      \leavevmode\fontsize{6.9pt}{7.9pt}\selectfont\strut Souffle vent à peine perceptible\strut &  \\
      \leavevmode\fontsize{6.9pt}{7.9pt}\selectfont\strut Ds le champ tt proche, les A.S.C. men\strut & \smash{\parbox[t]{39mm}{\fontsize{6.2pt}{7.0pt}\selectfont\raggedright\color{black!88}\textbf{8 Ds} \textit{abbr. pro} Dans. \textbf{8 tt} \textit{abbr. pro} tout.}} \\
      \leavevmode\fontsize{6.9pt}{7.9pt}\selectfont\strut se préparent lentement\strut &  \\
      \leavevmode\llap{\fontsize{6.4pt}{7.3pt}\selectfont\color{darkgray}10\hspace{1.8mm}}\fontsize{6.9pt}{7.9pt}\selectfont\strut Partis à 2H. Soleil devenu brulant\strut &  \\
      \leavevmode\fontsize{6.9pt}{7.9pt}\selectfont\strut Route poussiéreuse\strut &  \\
      \leavevmode\fontsize{6.9pt}{7.9pt}\selectfont\strut par Bailleul, parvenus à La Croix\strut &  \\
      \leavevmode\fontsize{6.9pt}{7.9pt}\selectfont\strut du Bac – vers 4H 1/2 –\strut & \smash{\parbox[t]{39mm}{\fontsize{6.2pt}{7.0pt}\selectfont\raggedright\color{black!88}\textbf{13 1/2} \textit{scil.} demi / moitié. — \textbf{14 Horse-lines} \textit{angl. in ms.}}} \\
      \leavevmode\fontsize{6.9pt}{7.9pt}\selectfont\strut Les Horse-lines, à l’entrée village, du\strut &  \\
      \leavevmode\llap{\fontsize{6.4pt}{7.3pt}\selectfont\color{darkgray}15\hspace{1.8mm}}\fontsize{6.9pt}{7.9pt}\selectfont\strut côté de Steenwerck.\strut &  \\
      \leavevmode\fontsize{6.9pt}{7.9pt}\selectfont\strut Cherché chambres pr. offs. Réussi à me\strut & \smash{\parbox[t]{39mm}{\fontsize{6.2pt}{7.0pt}\selectfont\raggedright\color{black!88}\textbf{16 pr} \textit{abbr. pro} pour. \textbf{16 offs} \textit{abbr. pro} officiers. — \textbf{17 Coy} \textit{angl.} Company.}} \\
      \leavevmode\fontsize{6.9pt}{7.9pt}\selectfont\strut loger maison en face de ma Coy —\strut &  \\
      \leavevmode\fontsize{6.9pt}{7.9pt}\selectfont\strut Fin journée vraiment superbe\strut &  \\
      \leavevmode\fontsize{6.9pt}{7.9pt}\selectfont\strut Lumière vive éclaire paysage plat\strut &  \\
      \leavevmode\llap{\fontsize{6.4pt}{7.3pt}\selectfont\color{darkgray}20\hspace{1.8mm}}\fontsize{6.9pt}{7.9pt}\selectfont\strut A fermière de St Jans-Cappel (pour oeufs,\strut & \smash{\parbox[t]{39mm}{\fontsize{6.2pt}{7.0pt}\selectfont\raggedright\color{black!88}\textbf{20 St} \textit{abbr. pro} Saint.}} \\
      \leavevmode\fontsize{6.9pt}{7.9pt}\selectfont\strut lait \& bière) : 1 f 60. et for washing : 3 f —\strut & \smash{\parbox[t]{39mm}{\fontsize{6.2pt}{7.0pt}\selectfont\raggedright\color{black!88}\textbf{21 washing} \textit{angl. in ms.}}} \\
      \leavevmode\fontsize{6.9pt}{7.9pt}\selectfont\strut Consommations : 60 c (bière)\strut &  \\
      \leavevmode\fontsize{6.9pt}{7.9pt}\selectfont\strut A Connelly, mon groom : 1f –\strut & \smash{\parbox[t]{39mm}{\fontsize{6.2pt}{7.0pt}\selectfont\raggedright\color{black!88}\textbf{23 groom} \textit{angl. in ms.}}} \\
\end{tabular}}%
\par\vspace*{\fill}
% --- Registres de bas de page (1. Critique interne [si pertinent] ; 2. Loci paralleli en latin ; 3. Commentaire historique en retrait) ---
\noindent\rule{\textwidth}{0.5pt}\par\vspace{0.65mm}
{\fontsize{6.6pt}{7.7pt}\selectfont%
\noindent\normalfont\color{black!92} \textbf{12} (\textit{Bailleul}) cf. supra f\textsuperscript{o}~3~v\textsuperscript{o} (p.~\pageref{fol:7}), f\textsuperscript{o}~5~r\textsuperscript{o} (p.~\pageref{fol:10}), infra f\textsuperscript{o}~13~v\textsuperscript{o} (p.~\pageref{fol:27}), f\textsuperscript{o}~14~r\textsuperscript{o} (p.~\pageref{fol:28}), \textit{et passim}. — \textbf{15} (\textit{Steenwerck}) cf. supra f\textsuperscript{o}~5~r\textsuperscript{o} (p.~\pageref{fol:10}), infra f\textsuperscript{o}~7~r\textsuperscript{o} (p.~\pageref{fol:14}), f\textsuperscript{o}~7~v\textsuperscript{o} (p.~\pageref{fol:15}), \textit{et passim}. — \textbf{23} (\textit{Connelly}) cf. infra f\textsuperscript{o}~19~v\textsuperscript{o} (p.~\pageref{fol:39}), f\textsuperscript{o}~20~v\textsuperscript{o} (p.~\pageref{fol:41}).\par\vspace{0.55mm}
\noindent\rule{38mm}{0.4pt}\par\vspace{0.5mm}
\noindent\hspace*{5mm}\parbox{\dimexpr\textwidth-5mm\relax}{\fontsize{6.6pt}{7.7pt}\selectfont\color{black!84}Le samedi 5 août 1916, la compagnie de l'\textit{Army Service Corps} (\textit{A.S.C. men}) quitte Saint-Jans-Cappel par Bailleul pour descendre dans la plaine de la Lys et établir ses lignes de chevaux (\textit{Horse-lines}) à l'entrée de \textbf{La Croix-du-Bac} (commune de Steenwerck), en arrière du secteur de Fleurbaix--Armentières que la 18\textsuperscript{e} Division tient alors en ligne.}\par}
\newpage
\phantomsection\label{fac:12}
% --- En-tête du fac-similé (Page impaire / droite : Photographie du manuscrit) ---
\noindent\rule{\textwidth}{0.5pt}\par\vspace{0.55mm}
\noindent\textbf{\fontsize{6.8pt}{7.8pt}\selectfont FAC-SIMILÉ EN VIS-À-VIS — FOLIO 6 RECTO} \hfill \textbf{\fontsize{6.8pt}{7.8pt}\selectfont AOÛT 1916}\par\vspace{0.25mm}
\noindent{\fontsize{6.8pt}{8.0pt}\selectfont\textit{Livre-Journal d'Hippolyte Combes} — Samedi 5 août 1916 (De Saint-Jans-Cappel à La Croix-du-Bac / Steenwerck (Nord))}\par\vspace{0.55mm}
\noindent\rule{\textwidth}{0.35pt}\par\vspace*{\fill}
{\centering
  \setlength{\fboxsep}{1.2mm}%
  \setlength{\fboxrule}{0.4pt}%
  \fcolorbox{gray!55}{white}{%
    \includegraphics[width=0.95\textwidth,height=0.81\textheight,keepaspectratio]{\facsimiledir/20251126_093347.jpg}%
  }\par}
\vspace*{\fill}
\newpage

% =============================================================================
% VIS-A-VIS SPREAD : FOLIO 6 VERSO (20251126_093412.jpg)
% =============================================================================
\cleardoubleevenpage
\renewcommand{\currentmonthyear}{Août 1916}
\renewcommand{\currentfoliolabel}{Folio 6 verso}
\phantomsection\label{fol:13}
% --- En-tête du folio (Page paire / gauche : Transcription, apparat & registres de bas de page) ---
\noindent\rule{\textwidth}{0.5pt}\par\vspace{0.55mm}
\noindent\textbf{\footnotesize FOLIO 6 VERSO} \hfill \textbf{\footnotesize AOÛT 1916}\par\vspace{0.25mm}
\noindent{\fontsize{6.9pt}{8.0pt}\selectfont Dimanche 6 août 1916 — Au refilling point sur les bords de la Lys avec le lieutenant Curtis, rencontre de Maurin et concert militaire}\par\vspace{0.2mm}
\noindent{\fontsize{6.9pt}{8.0pt}\selectfont\textit{Dimanche 6 août 1916} \hfill \textit{La Croix-du-Bac et Bac-Saint-Maur (Nord)}}\par\vspace{0.5mm}
\noindent\rule{\textwidth}{0.35pt}\par\vspace{0.8mm}
% --- Corps de la transcription diplomatique et apparat marginal en prolongement ---
{\renewcommand{\arraystretch}{0.87}%
\noindent\hspace*{3.5mm}\begin{tabular}{@{}p{87mm}@{\hspace{3.5mm}}p{39mm}@{}}
      \leavevmode\fontsize{6.9pt}{7.9pt}\selectfont\strut \underline{Dim. 6.} \hfill{} \underline{Aout 1916}\strut & \smash{\parbox[t]{39mm}{\fontsize{6.2pt}{7.0pt}\selectfont\raggedright\color{black!88}\textbf{1 Dim.} \textit{scil.} Dimanche.}} \\
      \leavevmode\fontsize{6.9pt}{7.9pt}\selectfont\strut Levé dès 7H 1/2 –\strut &  \\
      \leavevmode\fontsize{6.9pt}{7.9pt}\selectfont\strut Temps magnifique\strut &  \\
      \leavevmode\fontsize{6.9pt}{7.9pt}\selectfont\strut Breakfast terminé à peine, que Lt Curtis\strut & \smash{\parbox[t]{39mm}{\fontsize{6.2pt}{7.0pt}\selectfont\raggedright\color{black!88}\textbf{4 Lt} \textit{abbr. pro} Lieutenant. \textbf{4 Breakfast} \textit{angl. in ms.}}} \\
      \leavevmode\llap{\fontsize{6.4pt}{7.3pt}\selectfont\color{darkgray}5\hspace{1.8mm}}\fontsize{6.9pt}{7.9pt}\selectfont\strut Venu me prendre – Ensemble allés à pied\strut &  \\
      \leavevmode\fontsize{6.9pt}{7.9pt}\selectfont\strut au refilling point, sur rive canal (près\strut & \smash{\parbox[t]{39mm}{\fontsize{6.2pt}{7.0pt}\selectfont\raggedright\color{black!88}\textbf{6 refilling point} \textit{angl. in ms.}}} \\
      \leavevmode\fontsize{6.9pt}{7.9pt}\selectfont\strut d'écluse — Parcouru campagne à la\strut &  \\
      \leavevmode\fontsize{6.9pt}{7.9pt}\selectfont\strut recherche fourrage\strut &  \\
      \leavevmode\fontsize{6.9pt}{7.9pt}\selectfont\strut Seul, continué j.au village Bac St-Maur\strut & \smash{\parbox[t]{39mm}{\fontsize{6.2pt}{7.0pt}\selectfont\raggedright\color{black!88}\textbf{9 St} \textit{abbr. pro} Saint.}} \\
      \leavevmode\llap{\fontsize{6.4pt}{7.3pt}\selectfont\color{darkgray}10\hspace{1.8mm}}\fontsize{6.9pt}{7.9pt}\selectfont\strut \& visité église, à côté usine\strut &  \\
      \leavevmode\fontsize{6.9pt}{7.9pt}\selectfont\strut Après-midi, retrouvé Maurin.\strut &  \\
      \leavevmode\fontsize{6.9pt}{7.9pt}\selectfont\strut Avec lui, assis sur l’herbe, écouté\strut &  \\
      \leavevmode\fontsize{6.9pt}{7.9pt}\selectfont\strut musique militre anglse –\strut & \smash{\parbox[t]{39mm}{\fontsize{6.2pt}{7.0pt}\selectfont\raggedright\color{black!88}\textbf{13 militre} \textit{abbr. pro} militaire.}} \\
      \leavevmode\fontsize{6.9pt}{7.9pt}\selectfont\strut Installation provisoire où je suis\strut &  \\
      \leavevmode\llap{\fontsize{6.4pt}{7.3pt}\selectfont\color{darkgray}15\hspace{1.8mm}}\fontsize{6.9pt}{7.9pt}\selectfont\strut — Payé à ma propre : oeufs, lait \&\strut & \smash{\parbox[t]{39mm}{\fontsize{6.2pt}{7.0pt}\selectfont\raggedright\color{black!88}\textbf{15 propre} \textit{abbr. pro} propriétaire. — \textbf{16 c.} \textit{abbr. pro} centimes. — \textbf{17 Ctes ples} \textit{abbr. pro} Cartes postales.}} \\
      \leavevmode\fontsize{6.9pt}{7.9pt}\selectfont\strut bière : 90 c.. — 1 litre lait : 25 + 0.05 —\strut &  \\
      \leavevmode\fontsize{6.9pt}{7.9pt}\selectfont\strut — Ctes ples : 10 c\strut &  \\
\end{tabular}}%
\par\vspace*{\fill}
% --- Registres de bas de page (1. Critique interne [si pertinent] ; 2. Loci paralleli en latin ; 3. Commentaire historique en retrait) ---
\noindent\rule{\textwidth}{0.5pt}\par\vspace{0.65mm}
{\fontsize{6.6pt}{7.7pt}\selectfont%
\noindent\normalfont \textbf{Disposition :} Au bas du feuillet (l.~20--23), une grande accolade marginale gauche réunit les quatre lignes du compte de \textit{mess} et d'achats.\par\vspace{0.5mm}
\noindent\rule{22mm}{0.3pt}\par\vspace{0.45mm}
\noindent\normalfont\color{black!92} \textbf{4} (\textit{Curtis}) cf. supra f\textsuperscript{o}~1~r\textsuperscript{o} (p.~\pageref{fol:2}), f\textsuperscript{o}~3~r\textsuperscript{o} (p.~\pageref{fol:6}), infra f\textsuperscript{o}~7~r\textsuperscript{o} (p.~\pageref{fol:14}), f\textsuperscript{o}~9~v\textsuperscript{o} (p.~\pageref{fol:19}), \textit{et passim}. — \textbf{8} (\textit{fourrage}) cf. supra f\textsuperscript{o}~3~r\textsuperscript{o} (p.~\pageref{fol:6}), f\textsuperscript{o}~4~v\textsuperscript{o} (p.~\pageref{fol:9}), infra f\textsuperscript{o}~11~v\textsuperscript{o} (p.~\pageref{fol:23}), f\textsuperscript{o}~15~r\textsuperscript{o} (p.~\pageref{fol:30}), \textit{et passim}. — \textbf{11} (\textit{Maurin}) cf. infra f\textsuperscript{o}~7~r\textsuperscript{o} (p.~\pageref{fol:14}), f\textsuperscript{o}~9~v\textsuperscript{o} (p.~\pageref{fol:19}), \textit{et passim}.\par\vspace{0.55mm}
\noindent\rule{38mm}{0.4pt}\par\vspace{0.5mm}
\noindent\hspace*{5mm}\parbox{\dimexpr\textwidth-5mm\relax}{\fontsize{6.6pt}{7.7pt}\selectfont\color{black!84}Le point de ravitaillement quotidien (\textit{refilling point}) est établi sur la Lys canalisée près de l'écluse de \textbf{Bac-Saint-Maur} (Sailly-sur-la-Lys). Au hameau de \textbf{La Croix-du-Bac} (Steenwerck), l'église de briques rouges décrite par Hippolyte Combes est l'église néo-gothique \textbf{Notre-Dame-des-Sept-Douleurs}, bâtie en 1869–1870 par l'architecte lillois Charles Leroy (Base Mérimée, \texttt{IA59000348}).}\par}
\newpage
\phantomsection\label{fac:13}
% --- En-tête du fac-similé (Page impaire / droite : Photographie du manuscrit) ---
\noindent\rule{\textwidth}{0.5pt}\par\vspace{0.55mm}
\noindent\textbf{\fontsize{6.8pt}{7.8pt}\selectfont FAC-SIMILÉ EN VIS-À-VIS — FOLIO 6 VERSO} \hfill \textbf{\fontsize{6.8pt}{7.8pt}\selectfont AOÛT 1916}\par\vspace{0.25mm}
\noindent{\fontsize{6.8pt}{8.0pt}\selectfont\textit{Livre-Journal d'Hippolyte Combes} — Dimanche 6 août 1916 (La Croix-du-Bac et Bac-Saint-Maur (Nord))}\par\vspace{0.55mm}
\noindent\rule{\textwidth}{0.35pt}\par\vspace*{\fill}
{\centering
  \setlength{\fboxsep}{1.2mm}%
  \setlength{\fboxrule}{0.4pt}%
  \fcolorbox{gray!55}{white}{%
    \includegraphics[width=0.95\textwidth,height=0.81\textheight,keepaspectratio]{\facsimiledir/20251126_093412.jpg}%
  }\par}
\vspace*{\fill}
\newpage

% =============================================================================
% VIS-A-VIS SPREAD : FOLIO 7 RECTO (20251126_093426.jpg)
% =============================================================================
\cleardoubleevenpage
\renewcommand{\currentmonthyear}{Août 1916}
\renewcommand{\currentfoliolabel}{Folio 7 recto}
\phantomsection\label{fol:14}
% --- En-tête du folio (Page paire / gauche : Transcription, apparat & registres de bas de page) ---
\noindent\rule{\textwidth}{0.5pt}\par\vspace{0.55mm}
\noindent\textbf{\footnotesize FOLIO 7 RECTO} \hfill \textbf{\footnotesize AOÛT 1916}\par\vspace{0.25mm}
\noindent{\fontsize{6.9pt}{8.0pt}\selectfont Quête de cantonnements avec le lieutenant Curtis, départ inopiné de la compagnie et nouvelle chambre}\par\vspace{0.2mm}
\noindent{\fontsize{6.9pt}{8.0pt}\selectfont\textit{Lundi 7 août 1916} \hfill \textit{Steenwerck et route de Bac-Saint-Maur (Nord)}}\par\vspace{0.5mm}
\noindent\rule{\textwidth}{0.35pt}\par\vspace{0.8mm}
% --- Corps de la transcription diplomatique et apparat marginal en prolongement ---
{\renewcommand{\arraystretch}{0.87}%
\noindent\hspace*{3.5mm}\begin{tabular}{@{}p{87mm}@{\hspace{3.5mm}}p{39mm}@{}}
      \leavevmode\fontsize{6.9pt}{7.9pt}\selectfont\strut \underline{Lundi 7} – \hfill{} \underline{Chaleur accablante Aout 1916}\strut &  \\
      \leavevmode\fontsize{6.9pt}{7.9pt}\selectfont\strut Toute matinée, avec cet idiot de Ctis, cherché\strut & \smash{\parbox[t]{39mm}{\fontsize{6.2pt}{7.0pt}\selectfont\raggedright\color{black!88}\textbf{2 Ctis} \textit{abbr. pro} Curtis.}} \\
      \leavevmode\fontsize{6.9pt}{7.9pt}\selectfont\strut nouveaux billets et messroom pour mes off.\strut & \smash{\parbox[t]{39mm}{\fontsize{6.2pt}{7.0pt}\selectfont\raggedright\color{black!88}\textbf{3 off.} \textit{abbr. pro} officier(s). \textbf{3 billets} \textit{angl. in ms.}}} \\
      \leavevmode\fontsize{6.9pt}{7.9pt}\selectfont\strut Libre, avant midi – mais alors harassé\strut &  \\
      \leavevmode\llap{\fontsize{6.4pt}{7.3pt}\selectfont\color{darkgray}5\hspace{1.8mm}}\fontsize{6.9pt}{7.9pt}\selectfont\strut de fatigue –\strut &  \\
      \leavevmode\fontsize{6.9pt}{7.9pt}\selectfont\strut Quand venu prendre mes rations : trouvé\strut &  \\
      \leavevmode\fontsize{6.9pt}{7.9pt}\selectfont\strut Coy partie. [Barning ?] m'a appris que tte la Coy\strut & \smash{\parbox[t]{39mm}{\fontsize{6.2pt}{7.0pt}\selectfont\raggedright\color{black!88}\textbf{7} \textit{lect. incert. in ms.} \textbf{7 tte} \textit{abbr. pro} toute. — \textbf{8 shifté} \textit{angl. in ms.} — \textbf{9} \textit{lect. incert. in ms.}}} \\
      \leavevmode\fontsize{6.9pt}{7.9pt}\selectfont\strut a shifté ce matin, à 9 H — Damn!\strut &  \\
      \leavevmode\fontsize{6.9pt}{7.9pt}\selectfont\strut [illis.] à l'autre extrémité village —\strut &  \\
      \leavevmode\llap{\fontsize{6.4pt}{7.3pt}\selectfont\color{darkgray}10\hspace{1.8mm}}\fontsize{6.9pt}{7.9pt}\selectfont\strut Dû aussi changer de quartier.\strut &  \\
      \leavevmode\fontsize{6.9pt}{7.9pt}\selectfont\strut Suis installé maison au-dessus Post-office,\strut &  \\
      \leavevmode\fontsize{6.9pt}{7.9pt}\selectfont\strut route de Bac St Maur, avec gde chambre\strut & \smash{\parbox[t]{39mm}{\fontsize{6.2pt}{7.0pt}\selectfont\raggedright\color{black!88}\textbf{12 gde} \textit{abbr. pro} grande. \textbf{12 St} \textit{abbr. pro} Saint. — \textbf{13 1er} \textit{scil.} premier.}} \\
      \leavevmode\fontsize{6.9pt}{7.9pt}\selectfont\strut au 1er étage. Y couché dès ce soir.\strut &  \\
      \leavevmode\fontsize{6.9pt}{7.9pt}\selectfont\strut Terminé journée avec Maurin\strut &  \\
      \cline{1-1}
      \leavevmode\llap{\fontsize{6.4pt}{7.3pt}\selectfont\color{darkgray}15\hspace{1.8mm}}\fontsize{6.9pt}{7.9pt}\selectfont\strut Payé à ma propre : cooking : 1 f\strut & \smash{\parbox[t]{39mm}{\fontsize{6.2pt}{7.0pt}\selectfont\raggedright\color{black!88}\textbf{15 propre} \textit{abbr. pro} propriétaire. \textbf{15 cooking} \textit{angl. in ms.}}} \\
      \leavevmode\fontsize{6.9pt}{7.9pt}\selectfont\strut Oeufs \& lait : 75 c (2 jours)\strut &  \\
      \leavevmode\fontsize{6.9pt}{7.9pt}\selectfont\strut 1 litre bière : 30 c\strut &  \\
\end{tabular}}%
\par\vspace*{\fill}
% --- Registres de bas de page (1. Critique interne [si pertinent] ; 2. Loci paralleli en latin ; 3. Commentaire historique en retrait) ---
\noindent\rule{\textwidth}{0.5pt}\par\vspace{0.65mm}
{\fontsize{6.6pt}{7.7pt}\selectfont%
\noindent\normalfont\color{black!92} \textbf{2} (\textit{Ctis}) cf. supra f\textsuperscript{o}~1~r\textsuperscript{o} (p.~\pageref{fol:2}), f\textsuperscript{o}~6~v\textsuperscript{o} (p.~\pageref{fol:13}), infra f\textsuperscript{o}~9~v\textsuperscript{o} (p.~\pageref{fol:19}), f\textsuperscript{o}~10~v\textsuperscript{o} (p.~\pageref{fol:21}), \textit{et passim}. — \textbf{14} (\textit{Maurin}) cf. supra f\textsuperscript{o}~6~v\textsuperscript{o} (p.~\pageref{fol:13}), infra f\textsuperscript{o}~9~v\textsuperscript{o} (p.~\pageref{fol:19}), f\textsuperscript{o}~10~v\textsuperscript{o} (p.~\pageref{fol:21}), \textit{et passim}.\par\vspace{0.55mm}
\noindent\rule{38mm}{0.4pt}\par\vspace{0.5mm}
\noindent\hspace*{5mm}\parbox{\dimexpr\textwidth-5mm\relax}{\fontsize{6.6pt}{7.7pt}\selectfont\color{black!84}Tandis qu'Hippolyte Combes et le lieutenant \textbf{Curtis} (\textit{C\textsuperscript{tis}}) cherchent des logements (\textit{billets}) et une salle de popote (\textit{messroom}), la compagnie se déplace (\textit{a shifté}) dès 9~h à l'autre extrémité du village; l'interprète s'installe au premier étage d'une maison située au-dessus du bureau de poste (\textit{Post-office}) sur la route de Bac-Saint-Maur.}\par}
\newpage
\phantomsection\label{fac:14}
% --- En-tête du fac-similé (Page impaire / droite : Photographie du manuscrit) ---
\noindent\rule{\textwidth}{0.5pt}\par\vspace{0.55mm}
\noindent\textbf{\fontsize{6.8pt}{7.8pt}\selectfont FAC-SIMILÉ EN VIS-À-VIS — FOLIO 7 RECTO} \hfill \textbf{\fontsize{6.8pt}{7.8pt}\selectfont AOÛT 1916}\par\vspace{0.25mm}
\noindent{\fontsize{6.8pt}{8.0pt}\selectfont\textit{Livre-Journal d'Hippolyte Combes} — Lundi 7 août 1916 (Steenwerck et route de Bac-Saint-Maur (Nord))}\par\vspace{0.55mm}
\noindent\rule{\textwidth}{0.35pt}\par\vspace*{\fill}
{\centering
  \setlength{\fboxsep}{1.2mm}%
  \setlength{\fboxrule}{0.4pt}%
  \fcolorbox{gray!55}{white}{%
    \includegraphics[width=0.95\textwidth,height=0.81\textheight,keepaspectratio]{\facsimiledir/20251126_093426.jpg}%
  }\par}
\vspace*{\fill}
\newpage

% =============================================================================
% VIS-A-VIS SPREAD : FOLIO 7 VERSO (20251126_093455.jpg)
% =============================================================================
\cleardoubleevenpage
\renewcommand{\currentmonthyear}{Août 1916}
\renewcommand{\currentfoliolabel}{Folio 7 verso}
\phantomsection\label{fol:15}
% --- En-tête du folio (Page paire / gauche : Transcription, apparat & registres de bas de page) ---
\noindent\rule{\textwidth}{0.5pt}\par\vspace{0.55mm}
\noindent\textbf{\footnotesize FOLIO 7 VERSO} \hfill \textbf{\footnotesize AOÛT 1916}\par\vspace{0.25mm}
\noindent{\fontsize{6.9pt}{8.0pt}\selectfont Deux journées de canicule, de poussière et de moustiques dans la vallée de la Lys}\par\vspace{0.2mm}
\noindent{\fontsize{6.9pt}{8.0pt}\selectfont\textit{Mardi 8 et mercredi 9 août 1916} \hfill \textit{Route de Bac-Saint-Maur, Steenwerck (Nord)}}\par\vspace{0.5mm}
\noindent\rule{\textwidth}{0.35pt}\par\vspace{0.8mm}
% --- Corps de la transcription diplomatique et apparat marginal en prolongement ---
{\renewcommand{\arraystretch}{1.00}%
\noindent\hspace*{3.5mm}\begin{tabular}{@{}p{87mm}@{\hspace{3.5mm}}p{39mm}@{}}
      \leavevmode\fontsize{7.7pt}{9.0pt}\selectfont\strut \underline{Mardi 8.} \hfill{} \underline{Aout 1916 –}\strut &  \\
      \leavevmode\fontsize{7.7pt}{9.0pt}\selectfont\strut Tjours tempérre extt chaude —\strut & \smash{\parbox[t]{39mm}{\fontsize{6.2pt}{7.0pt}\selectfont\raggedright\color{black!88}\textbf{2 tempérre} \textit{abbr. pro} température. \textbf{2 extt} \textit{abbr. pro} extrêmement.}} \\
      \leavevmode\fontsize{7.7pt}{9.0pt}\selectfont\strut Bien reposé nouvelle chambre\strut &  \\
      \leavevmode\fontsize{7.7pt}{9.0pt}\selectfont\strut Achevé de m'installer et le mieux possible\strut &  \\
      \leavevmode\llap{\fontsize{6.4pt}{7.3pt}\selectfont\color{darkgray}5\hspace{1.8mm}}\fontsize{7.7pt}{9.0pt}\selectfont\strut Journée entière de liberté –\strut &  \\
      \leavevmode\fontsize{7.7pt}{9.0pt}\selectfont\strut Resté chambre à me reposer –\strut &  \\
      \leavevmode\fontsize{7.7pt}{9.0pt}\selectfont\strut Food : eggs : 50 c\strut & \smash{\parbox[t]{39mm}{\fontsize{6.2pt}{7.0pt}\selectfont\raggedright\color{black!88}\textbf{7 Food} \textit{angl. in ms.}}} \\
      \leavevmode\fontsize{7.7pt}{9.0pt}\selectfont\strut mercredi 9.\strut &  \\
      \leavevmode\fontsize{7.7pt}{9.0pt}\selectfont\strut Encore chalr accablante\strut & \smash{\parbox[t]{39mm}{\fontsize{6.2pt}{7.0pt}\selectfont\raggedright\color{black!88}\textbf{9 chalr} \textit{abbr. pro} chaleur.}} \\
      \leavevmode\llap{\fontsize{6.4pt}{7.3pt}\selectfont\color{darkgray}10\hspace{1.8mm}}\fontsize{7.7pt}{9.0pt}\selectfont\strut Resté le + possible ds ma chambre\strut & \smash{\parbox[t]{39mm}{\fontsize{6.2pt}{7.0pt}\selectfont\raggedright\color{black!88}\textbf{10 ds} \textit{abbr. pro} dans.}} \\
      \leavevmode\fontsize{7.7pt}{9.0pt}\selectfont\strut Routes si poussiéreuses que impossible\strut &  \\
      \leavevmode\fontsize{7.7pt}{9.0pt}\selectfont\strut tenir ma fenêtre ouverte —\strut &  \\
      \leavevmode\fontsize{7.7pt}{9.0pt}\selectfont\strut Nuit, dévoré par les moustiques\strut &  \\
\end{tabular}}%
\par\vspace*{\fill}
% --- Registres de bas de page (1. Critique interne [si pertinent] ; 2. Loci paralleli en latin ; 3. Commentaire historique en retrait) ---
\noindent\rule{\textwidth}{0.5pt}\par\vspace{0.65mm}
{\fontsize{6.6pt}{7.7pt}\selectfont%
\noindent\normalfont\color{black!92} \textbf{1} (\textit{de eodem commorationis loco}) cf. supra f\textsuperscript{o}~7~r\textsuperscript{o} (p.~\pageref{fol:14}).\par\vspace{0.55mm}
\noindent\rule{38mm}{0.4pt}\par\vspace{0.5mm}
\noindent\hspace*{5mm}\parbox{\dimexpr\textwidth-5mm\relax}{\fontsize{6.6pt}{7.7pt}\selectfont\color{black!84}Observation des conditions sanitaires et matérielles de l'été 1916 dans la basse plaine alluviale de la Lys (route de Bac-Saint-Maur à Steenwerck), sillonnée de fossés humides (\textit{watergangs}) et de convois hippomobiles et automobiles soulevant la poussière.}\par}
\newpage
\phantomsection\label{fac:15}
% --- En-tête du fac-similé (Page impaire / droite : Photographie du manuscrit) ---
\noindent\rule{\textwidth}{0.5pt}\par\vspace{0.55mm}
\noindent\textbf{\fontsize{6.8pt}{7.8pt}\selectfont FAC-SIMILÉ EN VIS-À-VIS — FOLIO 7 VERSO} \hfill \textbf{\fontsize{6.8pt}{7.8pt}\selectfont AOÛT 1916}\par\vspace{0.25mm}
\noindent{\fontsize{6.8pt}{8.0pt}\selectfont\textit{Livre-Journal d'Hippolyte Combes} — Mardi 8 et mercredi 9 août 1916 (Route de Bac-Saint-Maur, Steenwerck (Nord))}\par\vspace{0.55mm}
\noindent\rule{\textwidth}{0.35pt}\par\vspace*{\fill}
{\centering
  \setlength{\fboxsep}{1.2mm}%
  \setlength{\fboxrule}{0.4pt}%
  \fcolorbox{gray!55}{white}{%
    \includegraphics[width=0.95\textwidth,height=0.81\textheight,keepaspectratio]{\facsimiledir/20251126_093455.jpg}%
  }\par}
\vspace*{\fill}
\newpage

% =============================================================================
% VIS-A-VIS SPREAD : FOLIO 8 RECTO (20251126_093506.jpg)
% =============================================================================
\cleardoubleevenpage
\renewcommand{\currentmonthyear}{Août 1916}
\renewcommand{\currentfoliolabel}{Folio 8 recto}
\phantomsection\label{fol:16}
% --- En-tête du folio (Page paire / gauche : Transcription, apparat & registres de bas de page) ---
\noindent\rule{\textwidth}{0.5pt}\par\vspace{0.55mm}
\noindent\textbf{\footnotesize FOLIO 8 RECTO} \hfill \textbf{\footnotesize AOÛT 1916}\par\vspace{0.25mm}
\noindent{\fontsize{6.9pt}{8.0pt}\selectfont Promenade à Sailly-sur-la-Lys, l'église bombardée et conversation avec l'éclusier sur le général Maxse}\par\vspace{0.2mm}
\noindent{\fontsize{6.9pt}{8.0pt}\selectfont\textit{Jeudi 10 août 1916} \hfill \textit{Bac-Saint-Maur et Sailly-sur-la-Lys (Nord / Pas-de-Calais)}}\par\vspace{0.5mm}
\noindent\rule{\textwidth}{0.35pt}\par\vspace{0.8mm}
% --- Corps de la transcription diplomatique et apparat marginal en prolongement ---
{\renewcommand{\arraystretch}{0.87}%
\noindent\hspace*{3.5mm}\begin{tabular}{@{}p{87mm}@{\hspace{3.5mm}}p{39mm}@{}}
      \leavevmode\fontsize{6.9pt}{7.9pt}\selectfont\strut \underline{10 Jeudi.} \hfill{} \underline{Aout / 16 –}\strut &  \\
      \leavevmode\fontsize{6.9pt}{7.9pt}\selectfont\strut Enfin, matinée et partie soirée\strut &  \\
      \leavevmode\fontsize{6.9pt}{7.9pt}\selectfont\strut pluie fine \& intermittente — Très\strut &  \\
      \leavevmode\fontsize{6.9pt}{7.9pt}\selectfont\strut agréable de marcher le long routes\strut &  \\
      \leavevmode\llap{\fontsize{6.4pt}{7.3pt}\selectfont\color{darkgray}5\hspace{1.8mm}}\fontsize{6.9pt}{7.9pt}\selectfont\strut légèrt arrosées —\strut & \smash{\parbox[t]{39mm}{\fontsize{6.2pt}{7.0pt}\selectfont\raggedright\color{black!88}\textbf{5 légèrt} \textit{abbr. pro} légèrement.}} \\
      \leavevmode\fontsize{6.9pt}{7.9pt}\selectfont\strut Après-midi venu promener à\strut &  \\
      \leavevmode\fontsize{6.9pt}{7.9pt}\selectfont\strut pied au village de $[\![$Maur$]\!]$ Sailly s/ la Lys\strut & \smash{\parbox[t]{39mm}{\fontsize{6.2pt}{7.0pt}\selectfont\raggedright\color{black!88}\textbf{7} $[\![$Maur$]\!]\ $ \textit{del. in ms.} \textbf{7 s/} \textit{abbr. pro} sur.}} \\
      \leavevmode\fontsize{6.9pt}{7.9pt}\selectfont\strut à 3 ou 4 kilm. — Ne pas dépassé pont\strut &  \\
      \leavevmode\fontsize{6.9pt}{7.9pt}\selectfont\strut Remarqué église brûlée par les Boches\strut &  \\
      \leavevmode\llap{\fontsize{6.4pt}{7.3pt}\selectfont\color{darkgray}10\hspace{1.8mm}}\fontsize{6.9pt}{7.9pt}\selectfont\strut Il ne reste que murs et haut clocher\strut &  \\
      \leavevmode\fontsize{6.9pt}{7.9pt}\selectfont\strut Retour en longeant berges\strut &  \\
      \leavevmode\fontsize{6.9pt}{7.9pt}\selectfont\strut de la rivière\strut &  \\
      \leavevmode\fontsize{6.9pt}{7.9pt}\selectfont\strut Arrêté à écluse et causé\strut &  \\
      \leavevmode\fontsize{6.9pt}{7.9pt}\selectfont\strut lgt av. l'éclusier de notre gal\strut & \smash{\parbox[t]{39mm}{\fontsize{6.2pt}{7.0pt}\selectfont\raggedright\color{black!88}\textbf{14 lgt} \textit{abbr. pro} longtemps. \textbf{14 gal} \textit{abbr. pro} général.}} \\
      \leavevmode\llap{\fontsize{6.4pt}{7.3pt}\selectfont\color{darkgray}15\hspace{1.8mm}}\fontsize{6.9pt}{7.9pt}\selectfont\strut en chef, Maxse etc —\strut &  \\
      \leavevmode\fontsize{6.9pt}{7.9pt}\selectfont\strut Très beau coucher de\strut &  \\
      \leavevmode\fontsize{6.9pt}{7.9pt}\selectfont\strut soleil contemplé de ma fenêtre\strut &  \\
      \leavevmode\fontsize{6.9pt}{7.9pt}\selectfont\strut vers Estaires \& Hazebrouck\strut &  \\
      \leavevmode\fontsize{6.9pt}{7.9pt}\selectfont\strut Les Australiens\strut &  \\
      \cline{1-1}
      \leavevmode\llap{\fontsize{6.4pt}{7.3pt}\selectfont\color{darkgray}20\hspace{1.8mm}}\fontsize{6.9pt}{7.9pt}\selectfont\strut Ctes ples : 1 f. — Lait : 50 c\strut & \smash{\parbox[t]{39mm}{\fontsize{6.2pt}{7.0pt}\selectfont\raggedright\color{black!88}\textbf{20 Ctes ples} \textit{abbr. pro} Cartes postales. \textbf{20 f.} \textit{abbr. pro} franc(s).}} \\
\end{tabular}}%
\par\vspace*{\fill}
% --- Registres de bas de page (1. Critique interne [si pertinent] ; 2. Loci paralleli en latin ; 3. Commentaire historique en retrait) ---
\noindent\rule{\textwidth}{0.5pt}\par\vspace{0.65mm}
{\fontsize{6.6pt}{7.7pt}\selectfont%
\noindent\normalfont \textbf{Unité de page :} La page regroupe le récit du jeudi 10 août (l.~1--18), un relevé de dépenses (l.~19) et une brève notation datée du mercredi 9 août (l.~20--22) ajoutée en bas de page.\par\vspace{0.5mm}
\noindent\rule{22mm}{0.3pt}\par\vspace{0.45mm}
\noindent\normalfont\color{black!92} \textbf{7} (\textit{Lys}) cf. supra f\textsuperscript{o}~6~v\textsuperscript{o} (p.~\pageref{fol:13}), f\textsuperscript{o}~7~v\textsuperscript{o} (p.~\pageref{fol:15}), infra f\textsuperscript{o}~8~v\textsuperscript{o} (p.~\pageref{fol:17}), f\textsuperscript{o}~9~r\textsuperscript{o} (p.~\pageref{fol:18}), \textit{et passim}. — \textbf{15} (\textit{Maxse}) cf. . — \textbf{18} (\textit{Hazebrouck}) cf. supra f\textsuperscript{o}~2~r\textsuperscript{o} (p.~\pageref{fol:4}), f\textsuperscript{o}~4~r\textsuperscript{o} (p.~\pageref{fol:8}), infra f\textsuperscript{o}~13~v\textsuperscript{o} (p.~\pageref{fol:27}).\par\vspace{0.55mm}
\noindent\rule{38mm}{0.4pt}\par\vspace{0.5mm}
\noindent\hspace*{5mm}\parbox{\dimexpr\textwidth-5mm\relax}{\fontsize{6.6pt}{7.7pt}\selectfont\color{black!84}Promenade le long de la Lys jusqu'au pont de \textbf{Sailly-sur-la-Lys} (Pas-de-Calais), dont l'église Saint-Vaast avait été incendiée lors des combats d'octobre 1914. À l'écluse, Hippolyte Combes s'entretient avec l'éclusier au sujet du \textbf{major-général Sir Ivor Maxse} (1862–1958), commandant la \textbf{18\textsuperscript{e} Division britannique} de 1914 à janvier 1917 ; la mention des \textit{Australiens} fait écho à l'engagement des divisions de l'I~ et du II~ANZAC Corps dans le secteur voisin de Fromelles et d'Armentières en juillet–août 1916.}\par}
\newpage
\phantomsection\label{fac:16}
% --- En-tête du fac-similé (Page impaire / droite : Photographie du manuscrit) ---
\noindent\rule{\textwidth}{0.5pt}\par\vspace{0.55mm}
\noindent\textbf{\fontsize{6.8pt}{7.8pt}\selectfont FAC-SIMILÉ EN VIS-À-VIS — FOLIO 8 RECTO} \hfill \textbf{\fontsize{6.8pt}{7.8pt}\selectfont AOÛT 1916}\par\vspace{0.25mm}
\noindent{\fontsize{6.8pt}{8.0pt}\selectfont\textit{Livre-Journal d'Hippolyte Combes} — Jeudi 10 août 1916 (Bac-Saint-Maur et Sailly-sur-la-Lys (Nord / Pas-de-Calais))}\par\vspace{0.55mm}
\noindent\rule{\textwidth}{0.35pt}\par\vspace*{\fill}
{\centering
  \setlength{\fboxsep}{1.2mm}%
  \setlength{\fboxrule}{0.4pt}%
  \fcolorbox{gray!55}{white}{%
    \includegraphics[width=0.95\textwidth,height=0.81\textheight,keepaspectratio]{\facsimiledir/20251126_093506.jpg}%
  }\par}
\vspace*{\fill}
\newpage

% =============================================================================
% VIS-A-VIS SPREAD : FOLIO 8 VERSO (20251126_093529.jpg)
% =============================================================================
\cleardoubleevenpage
\renewcommand{\currentmonthyear}{Août 1916}
\renewcommand{\currentfoliolabel}{Folio 8 verso}
\phantomsection\label{fol:17}
% --- En-tête du folio (Page paire / gauche : Transcription, apparat & registres de bas de page) ---
\noindent\rule{\textwidth}{0.5pt}\par\vspace{0.55mm}
\noindent\textbf{\footnotesize FOLIO 8 VERSO} \hfill \textbf{\footnotesize AOÛT 1916}\par\vspace{0.25mm}
\noindent{\fontsize{6.9pt}{8.0pt}\selectfont Vendredi 11 août 1916 — Bombardement meurtrier autour de la grande blanchisserie Barbry à Steenwerck}\par\vspace{0.2mm}
\noindent{\fontsize{6.9pt}{8.0pt}\selectfont\textit{Vendredi 11 août 1916} \hfill \textit{Steenwerck / Erquinghem-Lys (Nord)}}\par\vspace{0.5mm}
\noindent\rule{\textwidth}{0.35pt}\par\vspace{0.8mm}
% --- Corps de la transcription diplomatique et apparat marginal en prolongement ---
{\renewcommand{\arraystretch}{0.87}%
\noindent\hspace*{3.5mm}\begin{tabular}{@{}p{87mm}@{\hspace{3.5mm}}p{39mm}@{}}
      \leavevmode\fontsize{6.9pt}{7.9pt}\selectfont\strut \underline{11 Ven.} \hfill{} \underline{Aout / 16}\strut & \smash{\parbox[t]{39mm}{\fontsize{6.2pt}{7.0pt}\selectfont\raggedright\color{black!88}\textbf{1 Ven.} \textit{scil.} Vendredi.}} \\
      \leavevmode\fontsize{6.9pt}{7.9pt}\selectfont\strut De nouveau tempérre caniculaire.\strut & \smash{\parbox[t]{39mm}{\fontsize{6.2pt}{7.0pt}\selectfont\raggedright\color{black!88}\textbf{2 tempérre} \textit{abbr. pro} température.}} \\
      \leavevmode\fontsize{6.9pt}{7.9pt}\selectfont\strut Après-midi, de 2 à 3, 40 ou 50 obus\strut &  \\
      \leavevmode\fontsize{6.9pt}{7.9pt}\selectfont\strut tombés autour de nos parages\strut &  \\
      \leavevmode\llap{\fontsize{6.4pt}{7.3pt}\selectfont\color{darkgray}5\hspace{1.8mm}}\fontsize{6.9pt}{7.9pt}\selectfont\strut Maison du coin, à 50 mètres : 1\strut &  \\
      \leavevmode\fontsize{6.9pt}{7.9pt}\selectfont\strut soldat, $[\![$décapité [?]$]\!]$. — $[\![$bras emportés [?]$]\!]$\strut & \smash{\parbox[t]{39mm}{\fontsize{6.2pt}{7.0pt}\selectfont\raggedright\color{black!88}\textbf{6} $[\![$décapité [?]$]\!]\ $ \textit{del. in ms.} \textbf{6} $[\![$bras emportés [?]$]\!]\ $ \textit{del. in ms.} — \textbf{7 gde} \textit{abbr. pro} grande. — \textbf{8 bcp} \textit{abbr. pro} beaucoup. — \textbf{9 s/} \textit{abbr. pro} sur. — \textbf{11 bombardt} \textit{abbr. pro} bombardement. — \textbf{12} $[\![$perpétuelle [?]$]\!]\ $ \textit{del. in ms.}}} \\
      \leavevmode\fontsize{6.9pt}{7.9pt}\selectfont\strut A la gde blanchisserie\strut &  \\
      \leavevmode\fontsize{6.9pt}{7.9pt}\selectfont\strut Barbry, bcp de dégâts. Obus\strut &  \\
      \leavevmode\fontsize{6.9pt}{7.9pt}\selectfont\strut non éclatés, encore s/ l'herbe —\strut &  \\
      \leavevmode\llap{\fontsize{6.4pt}{7.3pt}\selectfont\color{darkgray}10\hspace{1.8mm}}\fontsize{6.9pt}{7.9pt}\selectfont\strut Descendus cave pendant le\strut &  \\
      \leavevmode\fontsize{6.9pt}{7.9pt}\selectfont\strut bombardt. Énervement \&\strut &  \\
      \leavevmode\fontsize{6.9pt}{7.9pt}\selectfont\strut malaise — Insécurité $[\![$perpétuelle [?]$]\!]$ continuelle\strut &  \\
      \leavevmode\fontsize{6.9pt}{7.9pt}\selectfont\strut Oeufs : 50 + 0.05. — —\strut &  \\
      \leavevmode\fontsize{6.9pt}{7.9pt}\selectfont\strut Question nouveau départ dimanche\strut &  \\
      \leavevmode\llap{\fontsize{6.4pt}{7.3pt}\selectfont\color{darkgray}15\hspace{1.8mm}}\fontsize{6.9pt}{7.9pt}\selectfont\strut gde chambre, haute fenêtre, belle\strut & \smash{\parbox[t]{39mm}{\fontsize{6.2pt}{7.0pt}\selectfont\raggedright\color{black!88}\textbf{15 gde} \textit{abbr. pro} grande.}} \\
      \leavevmode\fontsize{6.9pt}{7.9pt}\selectfont\strut vue campagne très étendue\strut &  \\
      \leavevmode\fontsize{6.9pt}{7.9pt}\selectfont\strut ouest d'Armentières.\strut &  \\
\end{tabular}}%
\par\vspace*{\fill}
% --- Registres de bas de page (1. Critique interne [si pertinent] ; 2. Loci paralleli en latin ; 3. Commentaire historique en retrait) ---
\noindent\rule{\textwidth}{0.5pt}\par\vspace{0.65mm}
{\fontsize{6.6pt}{7.7pt}\selectfont%
\noindent\normalfont \textbf{Autocensure autographe (l.~20--23) :} Aux quatre dernières lignes du récit, les mots décrivant le supplice infligé au soldat britannique (\textit{attaché à la roue}, \textit{les mains liées}, \textit{les pieds liés}) ont été biffés après coup de traits obliques serrés par Hippolyte Combes, mais demeurent déchiffrables sous la rature.\par\vspace{0.5mm}
\noindent\rule{22mm}{0.3pt}\par\vspace{0.45mm}
\noindent\normalfont\color{black!92} \textbf{8} (\textit{Barbry}) cf. infra f\textsuperscript{o}~11~r\textsuperscript{o} (p.~\pageref{fol:22}), f\textsuperscript{o}~11~v\textsuperscript{o} (p.~\pageref{fol:23}). — \textbf{17} (\textit{Armentières}) cf. infra f\textsuperscript{o}~9~v\textsuperscript{o} (p.~\pageref{fol:19}), f\textsuperscript{o}~10~v\textsuperscript{o} (p.~\pageref{fol:21}).\par\vspace{0.55mm}
\noindent\rule{38mm}{0.4pt}\par\vspace{0.5mm}
\noindent\hspace*{5mm}\parbox{\dimexpr\textwidth-5mm\relax}{\fontsize{6.6pt}{7.7pt}\selectfont\color{black!84}Le bombardement allemand du 11 août 1916 sur les abords d'\textbf{Armentières} touche un soldat britannique et endommage la \textbf{blanchisserie industrielle de fils et toiles Barbry}, établie en bordure de la Lys près du pont de Nieppe et dirigée par \textbf{Émile Barbry} (Base Mérimée, \texttt{IA59001501}). Aux l.~20–23, le passage raturé décrit la peine disciplinaire réglementaire britannique dite \textit{Field Punishment No.~1} (\textit{crucifixion}), consistant à attacher le condamné par les poignets et les chevilles à une roue de caisson.}\par}
\newpage
\phantomsection\label{fac:17}
% --- En-tête du fac-similé (Page impaire / droite : Photographie du manuscrit) ---
\noindent\rule{\textwidth}{0.5pt}\par\vspace{0.55mm}
\noindent\textbf{\fontsize{6.8pt}{7.8pt}\selectfont FAC-SIMILÉ EN VIS-À-VIS — FOLIO 8 VERSO} \hfill \textbf{\fontsize{6.8pt}{7.8pt}\selectfont AOÛT 1916}\par\vspace{0.25mm}
\noindent{\fontsize{6.8pt}{8.0pt}\selectfont\textit{Livre-Journal d'Hippolyte Combes} — Vendredi 11 août 1916 (Steenwerck / Erquinghem-Lys (Nord))}\par\vspace{0.55mm}
\noindent\rule{\textwidth}{0.35pt}\par\vspace*{\fill}
{\centering
  \setlength{\fboxsep}{1.2mm}%
  \setlength{\fboxrule}{0.4pt}%
  \fcolorbox{gray!55}{white}{%
    \includegraphics[width=0.95\textwidth,height=0.81\textheight,keepaspectratio]{\facsimiledir/20251126_093529.jpg}%
  }\par}
\vspace*{\fill}
\newpage

% =============================================================================
% VIS-A-VIS SPREAD : FOLIO 9 RECTO (20251126_093544.jpg)
% =============================================================================
\cleardoubleevenpage
\renewcommand{\currentmonthyear}{Août 1916}
\renewcommand{\currentfoliolabel}{Folio 9 recto}
\phantomsection\label{fol:18}
% --- En-tête du folio (Page paire / gauche : Transcription, apparat & registres de bas de page) ---
\noindent\rule{\textwidth}{0.5pt}\par\vspace{0.55mm}
\noindent\textbf{\footnotesize FOLIO 9 RECTO} \hfill \textbf{\footnotesize AOÛT 1916}\par\vspace{0.25mm}
\noindent{\fontsize{6.9pt}{8.0pt}\selectfont Sieste au bord de la Lys dans un champ de luzerne face aux blés mûrs}\par\vspace{0.2mm}
\noindent{\fontsize{6.9pt}{8.0pt}\selectfont\textit{Samedi 12 août 1916} \hfill \textit{Rives de la Lys, Bac-Saint-Maur (Nord)}}\par\vspace{0.5mm}
\noindent\rule{\textwidth}{0.35pt}\par\vspace{0.8mm}
% --- Corps de la transcription diplomatique et apparat marginal en prolongement ---
{\renewcommand{\arraystretch}{1.00}%
\noindent\hspace*{3.5mm}\begin{tabular}{@{}p{87mm}@{\hspace{3.5mm}}p{39mm}@{}}
      \leavevmode\fontsize{7.7pt}{9.0pt}\selectfont\strut \underline{12 Sam.} \hfill{} \underline{Aout 1916}\strut & \smash{\parbox[t]{39mm}{\fontsize{6.2pt}{7.0pt}\selectfont\raggedright\color{black!88}\textbf{1 Sam.} \textit{scil.} Samedi.}} \\
      \leavevmode\fontsize{7.7pt}{9.0pt}\selectfont\strut Encore journée de chaleur caniculaire\strut &  \\
      \leavevmode\fontsize{7.7pt}{9.0pt}\selectfont\strut Fin après-midi, sorti pour me\strut &  \\
      \leavevmode\fontsize{7.7pt}{9.0pt}\selectfont\strut distraire –\strut &  \\
      \leavevmode\llap{\fontsize{6.4pt}{7.3pt}\selectfont\color{darkgray}5\hspace{1.8mm}}\fontsize{7.7pt}{9.0pt}\selectfont\strut Venu promener le long rives\strut &  \\
      \leavevmode\fontsize{7.7pt}{9.0pt}\selectfont\strut la Lys. Causé lgs avec éclusier\strut & \smash{\parbox[t]{39mm}{\fontsize{6.2pt}{7.0pt}\selectfont\raggedright\color{black!88}\textbf{6 lgs} \textit{abbr. pro} longtemps.}} \\
      \leavevmode\fontsize{7.7pt}{9.0pt}\selectfont\strut puis continué route sur berge\strut &  \\
      \leavevmode\fontsize{7.7pt}{9.0pt}\selectfont\strut gazonnée — Trouvé oasis d'ombre\strut &  \\
      \leavevmode\fontsize{7.7pt}{9.0pt}\selectfont\strut champ de luzerne, $[\![$devant$]\!]$ en face de champs\strut & \smash{\parbox[t]{39mm}{\fontsize{6.2pt}{7.0pt}\selectfont\raggedright\color{black!88}\textbf{9} $[\![$devant$]\!]\ $ \textit{del. in ms.}}} \\
      \leavevmode\llap{\fontsize{6.4pt}{7.3pt}\selectfont\color{darkgray}10\hspace{1.8mm}}\fontsize{7.7pt}{9.0pt}\selectfont\strut d'épis de blé tout jaunis par soleil\strut &  \\
      \leavevmode\fontsize{7.7pt}{9.0pt}\selectfont\strut Allongé s/ le dos, gouté repos\strut & \smash{\parbox[t]{39mm}{\fontsize{6.2pt}{7.0pt}\selectfont\raggedright\color{black!88}\textbf{11 s/} \textit{abbr. pro} sur.}} \\
      \leavevmode\fontsize{7.7pt}{9.0pt}\selectfont\strut délicieux – admirant nuages floconneux\strut &  \\
      \leavevmode\fontsize{7.7pt}{9.0pt}\selectfont\strut collés s/ ciel bleu clair –\strut & \smash{\parbox[t]{39mm}{\fontsize{6.2pt}{7.0pt}\selectfont\raggedright\color{black!88}\textbf{13 s/} \textit{abbr. pro} sur.}} \\
      \leavevmode\fontsize{7.7pt}{9.0pt}\selectfont\strut 1/2 litre lait : 15 c\strut &  \\
      \leavevmode\llap{\fontsize{6.4pt}{7.3pt}\selectfont\color{darkgray}15\hspace{1.8mm}}\fontsize{7.7pt}{9.0pt}\selectfont\strut Après lunch, resté chambre à\strut & \smash{\parbox[t]{39mm}{\fontsize{6.2pt}{7.0pt}\selectfont\raggedright\color{black!88}\textbf{15 lunch} \textit{angl. in ms.}}} \\
      \leavevmode\fontsize{7.7pt}{9.0pt}\selectfont\strut lire – impossible s'éloigner\strut &  \\
      \leavevmode\fontsize{7.7pt}{9.0pt}\selectfont\strut air brulant, soleil torride\strut &  \\
\end{tabular}}%
\par\vspace*{\fill}
% --- Registres de bas de page (1. Critique interne [si pertinent] ; 2. Loci paralleli en latin ; 3. Commentaire historique en retrait) ---
\noindent\rule{\textwidth}{0.5pt}\par\vspace{0.65mm}
{\fontsize{6.6pt}{7.7pt}\selectfont%
\noindent\normalfont\color{black!92} \textbf{6} (\textit{Lys}) cf. supra f\textsuperscript{o}~6~v\textsuperscript{o} (p.~\pageref{fol:13}), f\textsuperscript{o}~8~v\textsuperscript{o} (p.~\pageref{fol:17}), infra f\textsuperscript{o}~9~v\textsuperscript{o} (p.~\pageref{fol:19}), f\textsuperscript{o}~11~r\textsuperscript{o} (p.~\pageref{fol:22}), \textit{et passim}.\par\vspace{0.55mm}
\noindent\rule{38mm}{0.4pt}\par\vspace{0.5mm}
\noindent\hspace*{5mm}\parbox{\dimexpr\textwidth-5mm\relax}{\fontsize{6.6pt}{7.7pt}\selectfont\color{black!84}Après le bombardement de la veille, journée d'accalmie sur les berges gazonnées de la Lys canalisée près de l'écluse de Bac-Saint-Maur, au milieu des champs de luzerne et de blé mûr.}\par}
\newpage
\phantomsection\label{fac:18}
% --- En-tête du fac-similé (Page impaire / droite : Photographie du manuscrit) ---
\noindent\rule{\textwidth}{0.5pt}\par\vspace{0.55mm}
\noindent\textbf{\fontsize{6.8pt}{7.8pt}\selectfont FAC-SIMILÉ EN VIS-À-VIS — FOLIO 9 RECTO} \hfill \textbf{\fontsize{6.8pt}{7.8pt}\selectfont AOÛT 1916}\par\vspace{0.25mm}
\noindent{\fontsize{6.8pt}{8.0pt}\selectfont\textit{Livre-Journal d'Hippolyte Combes} — Samedi 12 août 1916 (Rives de la Lys, Bac-Saint-Maur (Nord))}\par\vspace{0.55mm}
\noindent\rule{\textwidth}{0.35pt}\par\vspace*{\fill}
{\centering
  \setlength{\fboxsep}{1.2mm}%
  \setlength{\fboxrule}{0.4pt}%
  \fcolorbox{gray!55}{white}{%
    \includegraphics[width=0.95\textwidth,height=0.81\textheight,keepaspectratio]{\facsimiledir/20251126_093544.jpg}%
  }\par}
\vspace*{\fill}
\newpage

% =============================================================================
% CHAPITRE II (FOLIO 9 VERSO A FOLIO 24 VERSO)
% =============================================================================
\cleardoubleevenpage
\renewcommand{\currentmonthyear}{Août 1916 – Septembre 1916}
\renewcommand{\currentfoliolabel}{Chapitre II}
\chapter{De la vallée de la Lys à l’Artois et à la Somme}

\noindent{\fontsize{7.8pt}{9.0pt}\selectfont\textit{Folios 9 verso à 24 verso — Armentières, Estaires, Marquay, Saint-Pol, Lucheux, Raincheval et Forceville (13 août – 24 septembre 1916)}}\par\vspace{1.5mm}
\noindent\rule{\textwidth}{0.5pt}\par\vspace{0.9mm}
\noindent\textbf{\scriptsize\textsc{Résumé analytique du Chapitre II (Folio 9 verso à Folio 24 verso)}}\par\vspace{0.5mm}
\noindent{\fontsize{6.9pt}{8.1pt}\selectfont Ce deuxième chapitre (f\textsuperscript{o}~9~v\textsuperscript{o} à f\textsuperscript{o}~24~v\textsuperscript{o}) retrace la transition de la vallée de la Lys vers les plateaux d'Artois puis vers l'arrière-front immédiat de la bataille de la Somme. Du 13 au 24 août, Hippolyte Combes célèbre son quarantième anniversaire (14 août 1916), effectue avec le lieutenant Curtis une reconnaissance à cheval sous les averses jusqu'aux villages ruinés de Fleurbaix et d'Armentières, partage le café chez l'industriel Émile Barbry avec ses collègues interprètes Lépin et Schumacher, expédie un sac de bagages vers sa maison de Nogent-le-Rotrou et visite l'église et l'hôtel de ville d'Estaires. Le vendredi 25 août, embarquée en gare de Bailleul, la compagnie gagne en train Saint-Pol-sur-Ternoise et cantonne durant deux semaines à Marquay (25 août – 9 septembre) : entre les corvées de fourrage vert (\textit{green forage}) pour les fourgons des \textit{Bedfords}, Hippolyte Combes note l'entrée en guerre de la Roumanie, explore les églises de Ligny-Saint-Flochel et de Saint-Michel-sur-Ternoise et flâne à plusieurs reprises dans Saint-Pol. Du 9 au 11 septembre, une marche par étapes à cheval achemine la division vers le sud par le hameau d'Oppy (Beaudricourt) et la cité médiévale de Lucheux jusqu'à Raincheval (11–22 septembre), où le capitaine West l'admet à la table du \textit{mess} des officiers, avant le mouvement final vers Forceville (22–24 septembre) à la veille de l'assaut britannique contre Thiepval.}\par\vspace{1.0mm}
\noindent\rule{\textwidth}{0.35pt}\par\vspace{1.4mm}

\noindent\textbf{\scriptsize\textsc{Introduction historique et philologique}}\par\vspace{0.6mm}
{\fontsize{7.1pt}{8.3pt}\selectfont Le déplacement de la \textbf{18\textsuperscript{e} Division britannique} de la Lys vers le Ternois puis vers le secteur de l'Ancre s'inscrit dans la préparation de la grande offensive de septembre 1916 : rattachée au \textit{II~Corps} (lieutenant-général Sir Claud Jacob) au sein de la \textit{Reserve Army} (future \textit{Fifth Army}) du général Sir Hubert Gough, la division s'entraîne sur les plateaux autour de Saint-Pol avant de monter en ligne devant Thiepval, qu'elle prendra d'assaut le 26 septembre 1916 (Nichols, \textit{op.~cit.}, p.~86–92). Sur le plan linguistique, l'écriture quotidienne d'Hippolyte Combes témoigne d'une imprégnation croissante du vocabulaire militaire et logistique anglais (\textit{mess}, \textit{billet}, \textit{refilling point}, \textit{green forage}, \textit{horse-lines}, \textit{orderly N.C.O.}, \textit{on duty}, \textit{groom}, \textit{saddle soap}), tandis que ses méditations du 25 août dans le train de Saint-Pol font directement écho au pessimisme lucide de sa lettre confidentielle du 16 juin 1916 à sa sœur carmélite Françoise (Annexe~I, doc.~3).\par}
\vfill
\newpage
% --- Page de droite d'ouverture du chapitre : Carte vectorielle TikZ & Photographies d'archives ---
\noindent\rule{\textwidth}{0.5pt}\par\vspace{0.6mm}
\noindent\textbf{\fontsize{6.8pt}{7.8pt}\selectfont CARTE D'ITINÉRAIRE DU CHAPITRE II} \hfill \textbf{\fontsize{6.8pt}{7.8pt}\selectfont AOÛT 1916 – SEPTEMBRE 1916}\par\vspace{0.25mm}
\noindent{\fontsize{6.7pt}{7.8pt}\selectfont\textit{Déplacements, cantonnements et tournées d'Hippolyte Combes (Folio 9 verso à Folio 24 verso)}}\par\vspace{0.55mm}
\noindent\rule{\textwidth}{0.35pt}\par\vspace{1.2mm}
\begin{center}

\begin{tikzpicture}[
  scale=0.96,
  every node/.style={font=\fontsize{6.0pt}{7.0pt}\selectfont},
  loc/.style={circle, fill=black, inner sep=1.1pt},
  step/.style={circle, draw=bordeaux, fill=white, line width=0.75pt, inner sep=1.5pt},
  rail/.style={draw=black!80, line width=0.75pt, double distance=0.8pt, ->, >=stealth},
  move/.style={draw=bordeaux, line width=0.75pt, ->, >=stealth},
  excur/.style={draw=black!65, line width=0.5pt, dashed, ->, >=stealth},
  river/.style={draw=black!35, line width=1.0pt}
]
\draw[black!50, line width=0.4pt] (-60mm,-42mm) rectangle (60mm,42mm);
\fill[black!3] (-60mm,-42mm) rectangle (60mm,42mm);

% Rivières : la Lys au nord, la Ternoise au centre, l'Authie et l'Ancre au sud
\draw[river] (-20mm,31mm) -- (35mm,34mm) node[pos=0.8, below]{\textit{La Lys}};
\draw[river] (-45mm,10mm) .. controls (-20mm,6mm) and (-5mm,4mm) .. (10mm,2mm) node[right]{\textit{La Ternoise}};
\draw[river] (-40mm,-14mm) .. controls (-10mm,-16mm) and (10mm,-18mm) .. (25mm,-20mm) node[right]{\textit{L'Authie}};

% Ligne de front (Armentières - Arras - Thiepval)
\draw[bordeaux!80, line width=0.85pt, dash pattern=on 4pt off 2pt]
  (38mm,38mm) -- (32mm,14mm) -- (36mm,-4mm) -- (34mm,-36mm);
\node[rotate=-86, font=\fontsize{5.5pt}{6.5pt}\selectfont, bordeaux] at (40mm,2mm) {\textbf{Ligne de front (Armentières – Arras – Thiepval)}};

% Secteur 1 : Vallée de la Lys (13-25 août)
\node[step] (bailleul) at (-4mm, 36mm) {};
\node[above] at (bailleul) {\textbf{1. Bailleul} (embarquement 25 août)};
\node[loc] (croix) at (12mm, 33mm) {};
\node[below right] at (croix) {\textbf{La Croix-du-Bac}, Fleurbaix, Armentières (13–24 août)};
\node[loc] (estaires) at (4mm, 28mm) {};
\node[below left] at (estaires) {Estaires (23 août)};
\node[loc] (hazebrouck) at (-22mm, 34mm) {};
\node[above] at (hazebrouck) {Hazebrouck};

% Secteur 2 : Ternois / Artois (25 août - 9 sept.)
\node[step] (stpol) at (-14mm, 8mm) {};
\node[above left] at (stpol) {\textbf{Saint-Pol-sur-Ternoise} (25 août – 7 sept.)};
\node[step] (marquay) at (-2mm, 6mm) {};
\node[above right] at (marquay) {\textbf{2. Marquay} (cantonnement, 25 août – 9 sept.)};
\node[loc] (ligny) at (-4mm, 0mm) {};
\node[right] at (ligny) {Ligny-Saint-Flochel, Roëllecourt, Bailleul-aux-Cornailles};

% Secteur 3 : Marche vers la Somme (9-24 sept.)
\node[step] (oppy) at (4mm, -6mm) {};
\node[right] at (oppy) {\textbf{3. Oppy / Beaudricourt} (9–10 sept.)};
\node[step] (lucheux) at (-4mm, -15mm) {};
\node[left] at (lucheux) {\textbf{4. Lucheux} (11 sept.)};
\node[loc] (thievres) at (4mm, -22mm) {};
\node[right] at (thievres) {Thièvres};
\node[step] (raincheval) at (-2mm, -27mm) {};
\node[left] at (raincheval) {\textbf{5. Raincheval} (11–22 sept.)};
\node[loc] (acheux) at (10mm, -29mm) {};
\node[right] at (acheux) {\textbf{Acheux-en-Amiénois}};
\node[step] (forceville) at (14mm, -35mm) {};
\node[below] at (forceville) {\textbf{6. Forceville} \& Hédauville (22–24 sept.)};

% Trajets
\draw[excur] (croix) -- (estaires);
\draw[move] (croix) -- (bailleul);
\draw[rail] (bailleul) -- (hazebrouck) to[bend right=18] node[pos=0.5, left, font=\fontsize{5.6pt}{6.5pt}\selectfont]{\textit{Train militaire (25 août)}} (stpol);
\draw[move] (stpol) -- (marquay);
\draw[move] (marquay) -- (oppy) -- (lucheux) -- (thievres) -- (raincheval) -- (acheux) -- (forceville);

% Échelle et Nord
\draw[line width=0.6pt] (-54mm,-36mm) -- (-32mm,-36mm);
\node[above] at (-43mm,-36mm) {0 \hfill 10 \hfill 20 km};
\draw[->, line width=0.6pt] (52mm, 28mm) -- (52mm, 37mm) node[above]{\textbf{N}};
\end{tikzpicture}

\end{center}
\vfill
\noindent\rule{\textwidth}{0.3pt}\par\vspace{1.2mm}
\noindent\begin{center}
  \begin{minipage}[b]{0.42\textwidth}
    \centering
    \setlength{\fboxsep}{0.7mm}\setlength{\fboxrule}{0.4pt}%
    \fcolorbox{gray!55}{white}{\includegraphics[width=0.94\linewidth,height=28mm,keepaspectratio]{illustrations/hippolyte_interprete_2.jpg}}
    \par\vspace{0.6mm}{\fontsize{5.8pt}{6.7pt}\selectfont Portrait en buste d'Hippolyte Combes sur le front britannique (1916–1917), portant les insignes métalliques d'interprète détaché auprès de la B.E.F.}
  \end{minipage}\hfill
  \begin{minipage}[b]{0.42\textwidth}
    \centering
    \setlength{\fboxsep}{0.7mm}\setlength{\fboxrule}{0.4pt}%
    \fcolorbox{gray!55}{white}{\includegraphics[width=0.94\linewidth,height=28mm,keepaspectratio]{illustrations/hippolyte_1897.jpg}}
    \par\vspace{0.6mm}{\fontsize{5.8pt}{6.7pt}\selectfont Hippolyte Combes jeune homme en 1897 (à l'âge de vingt et un ans), étudiant à Toulouse avant son séjour de quatre ans à Nottingham.}
  \end{minipage}
\end{center}
\vspace*{1mm}
\newpage

% =============================================================================
% VIS-A-VIS SPREAD : FOLIO 9 VERSO (20251126_093618.jpg)
% =============================================================================
\cleardoubleevenpage
\renewcommand{\currentmonthyear}{Août 1916}
\renewcommand{\currentfoliolabel}{Folio 9 verso}
\phantomsection\label{fol:19}
% --- En-tête du folio (Page paire / gauche : Transcription, apparat & registres de bas de page) ---
\noindent\rule{\textwidth}{0.5pt}\par\vspace{0.55mm}
\noindent\textbf{\footnotesize FOLIO 9 VERSO} \hfill \textbf{\footnotesize AOÛT 1916}\par\vspace{0.25mm}
\noindent{\fontsize{6.9pt}{8.0pt}\selectfont Attente vaine d'une lettre de Marguerite, « pot-au-noir » et orage sur la Lys}\par\vspace{0.2mm}
\noindent{\fontsize{6.9pt}{8.0pt}\selectfont\textit{Dimanche 13 et lundi 14 août 1916} \hfill \textit{Bac-Saint-Maur et direction d'Armentières (Nord)}}\par\vspace{0.5mm}
\noindent\rule{\textwidth}{0.35pt}\par\vspace{0.8mm}
% --- Corps de la transcription diplomatique et apparat marginal en prolongement ---
{\renewcommand{\arraystretch}{0.87}%
\noindent\hspace*{3.5mm}\begin{tabular}{@{}p{87mm}@{\hspace{3.5mm}}p{39mm}@{}}
      \leavevmode\fontsize{6.9pt}{7.9pt}\selectfont\strut \underline{13 D} – \hfill{} \underline{Aout 1916}\strut &  \\
      \leavevmode\fontsize{6.9pt}{7.9pt}\selectfont\strut Soleil ardent\strut &  \\
      \leavevmode\fontsize{6.9pt}{7.9pt}\selectfont\strut Déception pénible : attendu en vain un\strut &  \\
      \leavevmode\fontsize{6.9pt}{7.9pt}\selectfont\strut mot de Margte –\strut & \smash{\parbox[t]{39mm}{\fontsize{6.2pt}{7.0pt}\selectfont\raggedright\color{black!88}\textbf{4 Margte} \textit{scil.} Marguerite (Calmels). — \textbf{5 Tt} \textit{abbr. pro} Tout.}} \\
      \leavevmode\llap{\fontsize{6.4pt}{7.3pt}\selectfont\color{darkgray}5\hspace{1.8mm}}\fontsize{6.9pt}{7.9pt}\selectfont\strut Tt le jour, pot-au-noir –\strut &  \\
      \leavevmode\fontsize{6.9pt}{7.9pt}\selectfont\strut Après-midi, venu retrouver Maurin\strut &  \\
      \leavevmode\fontsize{6.9pt}{7.9pt}\selectfont\strut et avec lui, promené campagne\strut &  \\
      \leavevmode\fontsize{6.9pt}{7.9pt}\selectfont\strut le long rivière de la Lys –\strut &  \\
      \cline{1-1}
      \leavevmode\fontsize{6.9pt}{7.9pt}\selectfont\strut \underline{14 L} – Ciel couvert nuages –\strut &  \\
      \leavevmode\llap{\fontsize{6.4pt}{7.3pt}\selectfont\color{darkgray}10\hspace{1.8mm}}\fontsize{6.9pt}{7.9pt}\selectfont\strut Enfin, vers 1 H, pluie très forte\strut &  \\
      \leavevmode\fontsize{6.9pt}{7.9pt}\selectfont\strut rafraichi l'air — et mouille (colle\strut &  \\
      \leavevmode\fontsize{6.9pt}{7.9pt}\selectfont\strut poussière routes.\strut &  \\
      \leavevmode\fontsize{6.9pt}{7.9pt}\selectfont\strut Pris plaisir à regarder averse\strut &  \\
      \leavevmode\fontsize{6.9pt}{7.9pt}\selectfont\strut zébrer air de raies obliques, à\strut &  \\
      \leavevmode\llap{\fontsize{6.4pt}{7.3pt}\selectfont\color{darkgray}15\hspace{1.8mm}}\fontsize{6.9pt}{7.9pt}\selectfont\strut écouter la pluie crépiter s/ vitre\strut & \smash{\parbox[t]{39mm}{\fontsize{6.2pt}{7.0pt}\selectfont\raggedright\color{black!88}\textbf{15 s/} \textit{abbr. pro} sur.}} \\
      \leavevmode\fontsize{6.9pt}{7.9pt}\selectfont\strut de ma fenêtre\strut &  \\
      \leavevmode\fontsize{6.9pt}{7.9pt}\selectfont\strut Avec Lt Curtis, course à cheval\strut & \smash{\parbox[t]{39mm}{\fontsize{6.2pt}{7.0pt}\selectfont\raggedright\color{black!88}\textbf{17 Lt} \textit{abbr. pro} Lieutenant. — \textbf{18 ds} \textit{abbr. pro} dans.}} \\
      \leavevmode\fontsize{6.9pt}{7.9pt}\selectfont\strut ds la direction d'Armentières —\strut &  \\
      \leavevmode\fontsize{6.9pt}{7.9pt}\selectfont\strut Note de blanchissage : 50c\strut &  \\
      \leavevmode\llap{\fontsize{6.4pt}{7.3pt}\selectfont\color{darkgray}20\hspace{1.8mm}}\fontsize{6.9pt}{7.9pt}\selectfont\strut Lait : 15 c + 0.05\strut &  \\
\end{tabular}}%
\par\vspace*{\fill}
% --- Registres de bas de page (1. Critique interne [si pertinent] ; 2. Loci paralleli en latin ; 3. Commentaire historique en retrait) ---
\noindent\rule{\textwidth}{0.5pt}\par\vspace{0.65mm}
{\fontsize{6.6pt}{7.7pt}\selectfont%
\noindent\normalfont\color{black!92} \textbf{4} (\textit{Margte}) cf. supra f\textsuperscript{o}~1~r\textsuperscript{o} (p.~\pageref{fol:2}), f\textsuperscript{o}~3~r\textsuperscript{o} (p.~\pageref{fol:6}), infra f\textsuperscript{o}~27~r\textsuperscript{o} (p.~\pageref{fol:54}), f\textsuperscript{o}~27~v\textsuperscript{o} (p.~\pageref{fol:55}), \textit{et passim}, Appendicem I, doc.~5 (p.~\pageref{annexe:doc5}). — \textbf{6} (\textit{Maurin}) cf. supra f\textsuperscript{o}~6~v\textsuperscript{o} (p.~\pageref{fol:13}), f\textsuperscript{o}~7~r\textsuperscript{o} (p.~\pageref{fol:14}), infra f\textsuperscript{o}~10~v\textsuperscript{o} (p.~\pageref{fol:21}), f\textsuperscript{o}~12~r\textsuperscript{o} (p.~\pageref{fol:24}), \textit{et passim}. — \textbf{8} (\textit{Lys}) cf. supra f\textsuperscript{o}~6~v\textsuperscript{o} (p.~\pageref{fol:13}), f\textsuperscript{o}~9~r\textsuperscript{o} (p.~\pageref{fol:18}), infra f\textsuperscript{o}~11~r\textsuperscript{o} (p.~\pageref{fol:22}), f\textsuperscript{o}~11~v\textsuperscript{o} (p.~\pageref{fol:23}), \textit{et passim}.\par\vspace{0.55mm}
\noindent\rule{38mm}{0.4pt}\par\vspace{0.5mm}
\noindent\hspace*{5mm}\parbox{\dimexpr\textwidth-5mm\relax}{\fontsize{6.6pt}{7.7pt}\selectfont\color{black!84}Le dimanche 13 août 1916, l'absence de courrier de sa fiancée \textbf{Marguerite Calmels} (1892--1974) assombrit Hippolyte Combes; le lundi \textbf{14 août 1916} — jour du quarantième anniversaire d'Hippolyte Combes, né à Millau le 14 août 1876 — est marqué par une forte averse d'orage suivie d'une reconnaissance à cheval avec le lieutenant Curtis vers Armentières.}\par}
\newpage
\phantomsection\label{fac:19}
% --- En-tête du fac-similé (Page impaire / droite : Photographie du manuscrit) ---
\noindent\rule{\textwidth}{0.5pt}\par\vspace{0.55mm}
\noindent\textbf{\fontsize{6.8pt}{7.8pt}\selectfont FAC-SIMILÉ EN VIS-À-VIS — FOLIO 9 VERSO} \hfill \textbf{\fontsize{6.8pt}{7.8pt}\selectfont AOÛT 1916}\par\vspace{0.25mm}
\noindent{\fontsize{6.8pt}{8.0pt}\selectfont\textit{Livre-Journal d'Hippolyte Combes} — Dimanche 13 et lundi 14 août 1916 (Bac-Saint-Maur et direction d'Armentières (Nord))}\par\vspace{0.55mm}
\noindent\rule{\textwidth}{0.35pt}\par\vspace*{\fill}
{\centering
  \setlength{\fboxsep}{1.2mm}%
  \setlength{\fboxrule}{0.4pt}%
  \fcolorbox{gray!55}{white}{%
    \includegraphics[width=0.95\textwidth,height=0.81\textheight,keepaspectratio]{\facsimiledir/20251126_093618.jpg}%
  }\par}
\vspace*{\fill}
\newpage

% =============================================================================
% VIS-A-VIS SPREAD : FOLIO 10 RECTO (20251126_093746.jpg)
% =============================================================================
\cleardoubleevenpage
\renewcommand{\currentmonthyear}{Août 1916}
\renewcommand{\currentfoliolabel}{Folio 10 recto}
\phantomsection\label{fol:20}
% --- En-tête du folio (Page paire / gauche : Transcription, apparat & registres de bas de page) ---
\noindent\rule{\textwidth}{0.5pt}\par\vspace{0.55mm}
\noindent\textbf{\footnotesize FOLIO 10 RECTO} \hfill \textbf{\footnotesize AOÛT 1916}\par\vspace{0.25mm}
\noindent{\fontsize{6.9pt}{8.0pt}\selectfont Journée pluvieuse du 15 août dans la mansarde et comptes de blanchissage}\par\vspace{0.2mm}
\noindent{\fontsize{6.9pt}{8.0pt}\selectfont\textit{Mardi 15 août 1916 (Assomption)} \hfill \textit{Bac-Saint-Maur / Steenwerck (Nord)}}\par\vspace{0.5mm}
\noindent\rule{\textwidth}{0.35pt}\par\vspace{0.8mm}
% --- Corps de la transcription diplomatique et apparat marginal en prolongement ---
{\renewcommand{\arraystretch}{1.00}%
\noindent\hspace*{3.5mm}\begin{tabular}{@{}p{87mm}@{\hspace{3.5mm}}p{39mm}@{}}
      \leavevmode\fontsize{7.7pt}{9.0pt}\selectfont\strut \underline{15 Ma – Assompt.} \hfill{} \raisebox{0.5ex}{\underline{Aout / 16}}\strut &  \\
      \leavevmode\fontsize{7.7pt}{9.0pt}\selectfont\strut Toute matinée et 1 peu après-midi :\strut &  \\
      \leavevmode\fontsize{7.7pt}{9.0pt}\selectfont\strut fortes ondées. Routes détrempées. Aucun\strut &  \\
      \leavevmode\fontsize{7.7pt}{9.0pt}\selectfont\strut envie aller promener. Resté partie journée\strut &  \\
      \leavevmode\llap{\fontsize{6.4pt}{7.3pt}\selectfont\color{darkgray}5\hspace{1.8mm}}\fontsize{7.7pt}{9.0pt}\selectfont\strut ds ma chambre. Large pièce : très\strut & \smash{\parbox[t]{39mm}{\fontsize{6.2pt}{7.0pt}\selectfont\raggedright\color{black!88}\textbf{5 ds} \textit{abbr. pro} dans.}} \\
      \leavevmode\fontsize{7.7pt}{9.0pt}\selectfont\strut haute de plafond. Belle vue de la\strut &  \\
      \leavevmode\fontsize{7.7pt}{9.0pt}\selectfont\strut fenêtre, où me tiens sans cesse\strut &  \\
      \leavevmode\fontsize{7.7pt}{9.0pt}\selectfont\strut — Presque ts les matins, à l'aube,\strut & \smash{\parbox[t]{39mm}{\fontsize{6.2pt}{7.0pt}\selectfont\raggedright\color{black!88}\textbf{8 ts} \textit{abbr. pro} tous.}} \\
      \leavevmode\fontsize{7.7pt}{9.0pt}\selectfont\strut épais brouillards\strut &  \\
      \leavevmode\llap{\fontsize{6.4pt}{7.3pt}\selectfont\color{darkgray}10\hspace{1.8mm}}\fontsize{7.7pt}{9.0pt}\selectfont\strut — Donné à ma propre, pour\strut & \smash{\parbox[t]{39mm}{\fontsize{6.2pt}{7.0pt}\selectfont\raggedright\color{black!88}\textbf{10 propre} \textit{abbr. pro} propriétaire. — \textbf{11 cooking} \textit{angl. in ms.}}} \\
      \leavevmode\fontsize{7.7pt}{9.0pt}\selectfont\strut cooking semaine dernière et 2 l. lait : 5 f —\strut &  \\
\end{tabular}}%
\par\vspace*{\fill}
% --- Registres de bas de page (1. Critique interne [si pertinent] ; 2. Loci paralleli en latin ; 3. Commentaire historique en retrait) ---
\noindent\rule{\textwidth}{0.5pt}\par\vspace{0.65mm}
{\fontsize{6.6pt}{7.7pt}\selectfont%
\noindent\normalfont\color{black!92} \textbf{1} (\textit{de eodem commorationis loco}) cf. supra f\textsuperscript{o}~9~v\textsuperscript{o} (p.~\pageref{fol:19}).\par\vspace{0.55mm}
\noindent\rule{38mm}{0.4pt}\par\vspace{0.5mm}
\noindent\hspace*{5mm}\parbox{\dimexpr\textwidth-5mm\relax}{\fontsize{6.6pt}{7.7pt}\selectfont\color{black!84}Journée pluvieuse de la fête de l'Assomption (15 août 1916) passée dans sa chambre de la route de Bac-Saint-Maur (Steenwerck); le tarif hebdomadaire versé à sa logeuse pour la préparation des repas (\textit{cooking}) et deux litres de lait s'élève à 5~francs.}\par}
\newpage
\phantomsection\label{fac:20}
% --- En-tête du fac-similé (Page impaire / droite : Photographie du manuscrit) ---
\noindent\rule{\textwidth}{0.5pt}\par\vspace{0.55mm}
\noindent\textbf{\fontsize{6.8pt}{7.8pt}\selectfont FAC-SIMILÉ EN VIS-À-VIS — FOLIO 10 RECTO} \hfill \textbf{\fontsize{6.8pt}{7.8pt}\selectfont AOÛT 1916}\par\vspace{0.25mm}
\noindent{\fontsize{6.8pt}{8.0pt}\selectfont\textit{Livre-Journal d'Hippolyte Combes} — Mardi 15 août 1916 (Assomption) (Bac-Saint-Maur / Steenwerck (Nord))}\par\vspace{0.55mm}
\noindent\rule{\textwidth}{0.35pt}\par\vspace*{\fill}
{\centering
  \setlength{\fboxsep}{1.2mm}%
  \setlength{\fboxrule}{0.4pt}%
  \fcolorbox{gray!55}{white}{%
    \includegraphics[width=0.95\textwidth,height=0.81\textheight,keepaspectratio]{\facsimiledir/20251126_093746.jpg}%
  }\par}
\vspace*{\fill}
\newpage

% =============================================================================
% VIS-A-VIS SPREAD : FOLIO 10 VERSO (20251126_093843.jpg)
% =============================================================================
\cleardoubleevenpage
\renewcommand{\currentmonthyear}{Août 1916}
\renewcommand{\currentfoliolabel}{Folio 10 verso}
\phantomsection\label{fol:21}
% --- En-tête du folio (Page paire / gauche : Transcription, apparat & registres de bas de page) ---
\noindent\rule{\textwidth}{0.5pt}\par\vspace{0.55mm}
\noindent\textbf{\footnotesize FOLIO 10 VERSO} \hfill \textbf{\footnotesize AOÛT 1916}\par\vspace{0.25mm}
\noindent{\fontsize{6.9pt}{8.0pt}\selectfont Vendredi 18 et Samedi 19 août 1916 — Giboulées d'août, perception de la solde et chevauchée sous l'orage avec le lieutenant Curtis}\par\vspace{0.2mm}
\noindent{\fontsize{6.9pt}{8.0pt}\selectfont\textit{Vendredi 18 et samedi 19 août 1916} \hfill \textit{Steenwerck, Fleurbaix et secteur d'Armentières (Nord)}}\par\vspace{0.5mm}
\noindent\rule{\textwidth}{0.35pt}\par\vspace{0.8mm}
% --- Corps de la transcription diplomatique et apparat marginal en prolongement ---
{\renewcommand{\arraystretch}{0.87}%
\noindent\hspace*{3.5mm}\begin{tabular}{@{}p{87mm}@{\hspace{3.5mm}}p{39mm}@{}}
      \leavevmode\fontsize{6.9pt}{7.9pt}\selectfont\strut \underline{18 Ve.} \hfill{} \underline{– Aout / 16 –}\strut & \smash{\parbox[t]{39mm}{\fontsize{6.2pt}{7.0pt}\selectfont\raggedright\color{black!88}\textbf{1 Ve.} \textit{scil.} Vendredi.}} \\
      \leavevmode\fontsize{6.9pt}{7.9pt}\selectfont\strut Depuis 2 ou 3 jours, vraies giboulées,\strut &  \\
      \leavevmode\fontsize{6.9pt}{7.9pt}\selectfont\strut comme en mars – Averses prolongées. Puis éclaircies\strut &  \\
      \leavevmode\fontsize{6.9pt}{7.9pt}\selectfont\strut et soleil reparaît\strut &  \\
      \leavevmode\llap{\fontsize{6.4pt}{7.3pt}\selectfont\color{darkgray}5\hspace{1.8mm}}\fontsize{6.9pt}{7.9pt}\selectfont\strut Hier après-midi, au refilling point et ds 1\strut & \smash{\parbox[t]{39mm}{\fontsize{6.2pt}{7.0pt}\selectfont\raggedright\color{black!88}\textbf{5 ds} \textit{abbr. pro} dans. \textbf{5 refilling point} \textit{angl. in ms.}}} \\
      \leavevmode\fontsize{6.9pt}{7.9pt}\selectfont\strut champ de trèfle, surpris par la pluie et\strut &  \\
      \leavevmode\fontsize{6.9pt}{7.9pt}\selectfont\strut mouillé j. à la peau –\strut & \smash{\parbox[t]{39mm}{\fontsize{6.2pt}{7.0pt}\selectfont\raggedright\color{black!88}\textbf{7 j. à} \textit{abbr. pro} jusqu'à.}} \\
      \leavevmode\fontsize{6.9pt}{7.9pt}\selectfont\strut Fin journée, assez belle. Venu retrouver Maurin,\strut &  \\
      \leavevmode\fontsize{6.9pt}{7.9pt}\selectfont\strut promené ensemble et assis, pour causer\strut &  \\
      \leavevmode\llap{\fontsize{6.4pt}{7.3pt}\selectfont\color{darkgray}10\hspace{1.8mm}}\fontsize{6.9pt}{7.9pt}\selectfont\strut sur l'herbe sèche, près des lignes de ma\strut &  \\
      \leavevmode\fontsize{6.9pt}{7.9pt}\selectfont\strut Coy. \hfill{} Vers 7 H 1/2 orage\strut & \smash{\parbox[t]{39mm}{\fontsize{6.2pt}{7.0pt}\selectfont\raggedright\color{black!88}\textbf{11 1/2} \textit{scil.} demi / moitié. — \textbf{12} \textit{in marg. cum unco} (comptes). \textbf{12 fr.} \textit{abbr. pro} franc(s). — \textbf{13 tte} \textit{abbr. pro} toute. — \textbf{15 délict} \textit{abbr. pro} délicieusement. — \textbf{16 commt} \textit{abbr. pro} commencement.}} \\
      \leavevmode\fontsize{6.9pt}{7.9pt}\selectfont\strut [ – Solde \& indemnité de fonctions : 84 fr.\strut &  \\
      \cline{1-1}
      \leavevmode\fontsize{6.9pt}{7.9pt}\selectfont\strut \underline{19 S.} — Plu tte nuit —\strut &  \\
      \leavevmode\fontsize{6.9pt}{7.9pt}\selectfont\strut Dès 8 H matin, pluie torrentielle. Routes\strut &  \\
      \leavevmode\llap{\fontsize{6.4pt}{7.3pt}\selectfont\color{darkgray}15\hspace{1.8mm}}\fontsize{6.9pt}{7.9pt}\selectfont\strut détrempées. Température rafraîchie délict.\strut &  \\
      \leavevmode\fontsize{6.9pt}{7.9pt}\selectfont\strut Toute matinée et commt après-midi :\strut &  \\
      \leavevmode\fontsize{6.9pt}{7.9pt}\selectfont\strut averses très fortes — comme grêle.\strut &  \\
      \leavevmode\fontsize{6.9pt}{7.9pt}\selectfont\strut Vers 2 H 30, ce crétin de Curtis venu\strut &  \\
      \leavevmode\fontsize{6.9pt}{7.9pt}\selectfont\strut me prendre : dû l'accompagner\strut &  \\
      \leavevmode\llap{\fontsize{6.4pt}{7.3pt}\selectfont\color{darkgray}20\hspace{1.8mm}}\fontsize{6.9pt}{7.9pt}\selectfont\strut à cheval, – \& cette fois sur un bon poney.\strut & \smash{\parbox[t]{39mm}{\fontsize{6.2pt}{7.0pt}\selectfont\raggedright\color{black!88}\textbf{20 poney} \textit{angl. in ms.}}} \\
      \leavevmode\fontsize{6.9pt}{7.9pt}\selectfont\strut Longue chevauchée zone\strut &  \\
      \leavevmode\fontsize{6.9pt}{7.9pt}\selectfont\strut exposée, derrière ligne de feu\strut &  \\
\end{tabular}}%
\par\vspace*{\fill}
% --- Registres de bas de page (1. Critique interne [si pertinent] ; 2. Loci paralleli en latin ; 3. Commentaire historique en retrait) ---
\noindent\rule{\textwidth}{0.5pt}\par\vspace{0.65mm}
{\fontsize{6.6pt}{7.7pt}\selectfont%
\noindent\normalfont \textbf{Unité de page :} Deux entrées successives (vendredi 18 août, l.~1--7, et samedi 19 août, l.~8--22) séparées par un trait horizontal ; la phrase de la l.~22 se poursuit sans interruption en tête du f\textsuperscript{o}~11~r\textsuperscript{o}.\par\vspace{0.5mm}
\noindent\rule{22mm}{0.3pt}\par\vspace{0.45mm}
\noindent\normalfont\color{black!92} \textbf{8} (\textit{Maurin}) cf. supra f\textsuperscript{o}~6~v\textsuperscript{o} (p.~\pageref{fol:13}), f\textsuperscript{o}~9~v\textsuperscript{o} (p.~\pageref{fol:19}), infra f\textsuperscript{o}~12~r\textsuperscript{o} (p.~\pageref{fol:24}), f\textsuperscript{o}~13~r\textsuperscript{o} (p.~\pageref{fol:26}), \textit{et passim}. — \textbf{12} (\textit{84 f.}) cf. supra f\textsuperscript{o}~2~r\textsuperscript{o} (p.~\pageref{fol:4}), infra f\textsuperscript{o}~39~v\textsuperscript{o} (p.~\pageref{fol:79}), f\textsuperscript{o}~56~v\textsuperscript{o} (p.~\pageref{fol:113}). — \textbf{18} (\textit{Curtis}) cf. supra f\textsuperscript{o}~1~r\textsuperscript{o} (p.~\pageref{fol:2}), f\textsuperscript{o}~9~v\textsuperscript{o} (p.~\pageref{fol:19}), infra f\textsuperscript{o}~11~v\textsuperscript{o} (p.~\pageref{fol:23}), f\textsuperscript{o}~15~r\textsuperscript{o} (p.~\pageref{fol:30}), \textit{et passim}.\par\vspace{0.55mm}
\noindent\rule{38mm}{0.4pt}\par\vspace{0.5mm}
\noindent\hspace*{5mm}\parbox{\dimexpr\textwidth-5mm\relax}{\fontsize{6.6pt}{7.7pt}\selectfont\color{black!84}Après avoir perçu le 18 août sa solde et son indemnité de fonctions (84~fr.), Hippolyte Combes accompagne à cheval le lieutenant Curtis le samedi 19 août à 14~h~30, sous des averses glaciales, dans la zone exposée immédiatement en arrière des tranchées de première ligne du secteur de Fleurbaix et d'Armentières (dont l'itinéraire détaillé se poursuit au sommet du f\textsuperscript{o}~11~r\textsuperscript{o}).}\par}
\newpage
\phantomsection\label{fac:21}
% --- En-tête du fac-similé (Page impaire / droite : Photographie du manuscrit) ---
\noindent\rule{\textwidth}{0.5pt}\par\vspace{0.55mm}
\noindent\textbf{\fontsize{6.8pt}{7.8pt}\selectfont FAC-SIMILÉ EN VIS-À-VIS — FOLIO 10 VERSO} \hfill \textbf{\fontsize{6.8pt}{7.8pt}\selectfont AOÛT 1916}\par\vspace{0.25mm}
\noindent{\fontsize{6.8pt}{8.0pt}\selectfont\textit{Livre-Journal d'Hippolyte Combes} — Vendredi 18 et samedi 19 août 1916 (Steenwerck, Fleurbaix et secteur d'Armentières (Nord))}\par\vspace{0.55mm}
\noindent\rule{\textwidth}{0.35pt}\par\vspace*{\fill}
{\centering
  \setlength{\fboxsep}{1.2mm}%
  \setlength{\fboxrule}{0.4pt}%
  \fcolorbox{gray!55}{white}{%
    \includegraphics[width=0.95\textwidth,height=0.81\textheight,keepaspectratio]{\facsimiledir/20251126_093843.jpg}%
  }\par}
\vspace*{\fill}
\newpage

% =============================================================================
% VIS-A-VIS SPREAD : FOLIO 11 RECTO (20251126_093857.jpg)
% =============================================================================
\cleardoubleevenpage
\renewcommand{\currentmonthyear}{Août 1916}
\renewcommand{\currentfoliolabel}{Folio 11 recto}
\phantomsection\label{fol:22}
% --- En-tête du folio (Page paire / gauche : Transcription, apparat & registres de bas de page) ---
\noindent\rule{\textwidth}{0.5pt}\par\vspace{0.55mm}
\noindent\textbf{\footnotesize FOLIO 11 RECTO} \hfill \textbf{\footnotesize AOÛT 1916}\par\vspace{0.25mm}
\noindent{\fontsize{6.9pt}{8.0pt}\selectfont Traversée de Fleurbaix et d'Armentières dévastées par les obus}\par\vspace{0.2mm}
\noindent{\fontsize{6.9pt}{8.0pt}\selectfont\textit{Samedi 19 août 1916 (suite)} \hfill \textit{Erquinghem-Lys, Fleurbaix et Armentières (Pas-de-Calais / Nord)}}\par\vspace{0.5mm}
\noindent\rule{\textwidth}{0.35pt}\par\vspace{0.8mm}
% --- Corps de la transcription diplomatique et apparat marginal en prolongement ---
{\renewcommand{\arraystretch}{0.87}%
\noindent\hspace*{3.5mm}\begin{tabular}{@{}p{87mm}@{\hspace{3.5mm}}p{39mm}@{}}
      \leavevmode\fontsize{6.9pt}{7.9pt}\selectfont\strut et galop à travers champs \hfill{} \raisebox{0.5ex}{\underline{19 Sam –}}\strut &  \\
      \leavevmode\fontsize{6.9pt}{7.9pt}\selectfont\strut et arrivé j. à la blanchisserie Barbry\strut & \smash{\parbox[t]{39mm}{\fontsize{6.2pt}{7.0pt}\selectfont\raggedright\color{black!88}\textbf{2 j. à} \textit{abbr. pro} jusqu'à.}} \\
      \leavevmode\fontsize{6.9pt}{7.9pt}\selectfont\strut Puis dirigé sur Fleurbaix. Au centre du\strut &  \\
      \leavevmode\fontsize{6.9pt}{7.9pt}\selectfont\strut village (tourné à dr.) et remonté lentement rue\strut &  \\
      \leavevmode\llap{\fontsize{6.4pt}{7.3pt}\selectfont\color{darkgray}5\hspace{1.8mm}}\fontsize{6.9pt}{7.9pt}\selectfont\strut dont presque ttes les maisons ont reçu des\strut & \smash{\parbox[t]{39mm}{\fontsize{6.2pt}{7.0pt}\selectfont\raggedright\color{black!88}\textbf{5 ttes} \textit{abbr. pro} toutes.}} \\
      \leavevmode\fontsize{6.9pt}{7.9pt}\selectfont\strut obus. –\strut &  \\
      \leavevmode\fontsize{6.9pt}{7.9pt}\selectfont\strut Continué ensuite sur Armentières\strut &  \\
      \leavevmode\fontsize{6.9pt}{7.9pt}\selectfont\strut Arrivé par gare (l. est complt déserte)\strut & \smash{\parbox[t]{39mm}{\fontsize{6.2pt}{7.0pt}\selectfont\raggedright\color{black!88}\textbf{8 complt} \textit{abbr. pro} complètement.}} \\
      \leavevmode\fontsize{6.9pt}{7.9pt}\selectfont\strut Parcouru principaux quartiers.\strut &  \\
      \leavevmode\llap{\fontsize{6.4pt}{7.3pt}\selectfont\color{darkgray}10\hspace{1.8mm}}\fontsize{6.9pt}{7.9pt}\selectfont\strut Ville importante, riche — très b. maisons\strut &  \\
      \leavevmode\fontsize{6.9pt}{7.9pt}\selectfont\strut Effets des bombardements : grands\strut &  \\
      \leavevmode\fontsize{6.9pt}{7.9pt}\selectfont\strut dommages. Vu, çà \& là, de toits\strut &  \\
      \leavevmode\fontsize{6.9pt}{7.9pt}\selectfont\strut crevés, de façades trouées, cloisons\strut &  \\
      \leavevmode\fontsize{6.9pt}{7.9pt}\selectfont\strut d'appartement mises à nu –\strut &  \\
      \leavevmode\llap{\fontsize{6.4pt}{7.3pt}\selectfont\color{darkgray}15\hspace{1.8mm}}\fontsize{6.9pt}{7.9pt}\selectfont\strut Église à briques rouges, sur 1 place\strut &  \\
      \leavevmode\fontsize{6.9pt}{7.9pt}\selectfont\strut très atteinte. Clocher découronné\strut &  \\
      \leavevmode\fontsize{6.9pt}{7.9pt}\selectfont\strut Toiture abattue\strut &  \\
      \leavevmode\fontsize{6.9pt}{7.9pt}\selectfont\strut La + des immeubles inoccupés,\strut &  \\
      \leavevmode\fontsize{6.9pt}{7.9pt}\selectfont\strut abandonnés par leurs habitants.\strut &  \\
      \leavevmode\llap{\fontsize{6.4pt}{7.3pt}\selectfont\color{darkgray}20\hspace{1.8mm}}\fontsize{6.9pt}{7.9pt}\selectfont\strut Écossais\strut &  \\
      \leavevmode\fontsize{6.9pt}{7.9pt}\selectfont\strut Rentrés par Gde route – Arrêté à\strut & \smash{\parbox[t]{39mm}{\fontsize{6.2pt}{7.0pt}\selectfont\raggedright\color{black!88}\textbf{21 Gde} \textit{abbr. pro} Grande.}} \\
      \leavevmode\fontsize{6.9pt}{7.9pt}\selectfont\strut Erquinghem et thé avec 1 off des Sussex\strut & \smash{\parbox[t]{39mm}{\fontsize{6.2pt}{7.0pt}\selectfont\raggedright\color{black!88}\textbf{22 off} \textit{abbr. pro} officier(s).}} \\
      \leavevmode\fontsize{6.9pt}{7.9pt}\selectfont\strut Fin journée ensoleillée et\strut &  \\
      \leavevmode\fontsize{6.9pt}{7.9pt}\selectfont\strut très belle.\strut &  \\
\end{tabular}}%
\par\vspace*{\fill}
% --- Registres de bas de page (1. Critique interne [si pertinent] ; 2. Loci paralleli en latin ; 3. Commentaire historique en retrait) ---
\noindent\rule{\textwidth}{0.5pt}\par\vspace{0.65mm}
{\fontsize{6.6pt}{7.7pt}\selectfont%
\noindent\normalfont \textbf{Enchaînement matériel :} La l.~1 poursuit directement la phrase commencée au bas du f\textsuperscript{o}~10~v\textsuperscript{o} ; la mention de date (\textit{19 Sam --}) est rejetée à l'extrémité droite de la première ligne.\par\vspace{0.5mm}
\noindent\rule{22mm}{0.3pt}\par\vspace{0.45mm}
\noindent\normalfont\color{black!92} \textbf{2} (\textit{Barbry}) cf. supra f\textsuperscript{o}~8~v\textsuperscript{o} (p.~\pageref{fol:17}), infra f\textsuperscript{o}~11~v\textsuperscript{o} (p.~\pageref{fol:23}). — \textbf{3} (\textit{Fleurbaix}) cf. supra f\textsuperscript{o}~10~v\textsuperscript{o} (p.~\pageref{fol:21}). — \textbf{7} (\textit{Armentières}) cf. supra f\textsuperscript{o}~8~v\textsuperscript{o} (p.~\pageref{fol:17}), f\textsuperscript{o}~10~v\textsuperscript{o} (p.~\pageref{fol:21}).\par\vspace{0.55mm}
\noindent\rule{38mm}{0.4pt}\par\vspace{0.5mm}
\noindent\hspace*{5mm}\parbox{\dimexpr\textwidth-5mm\relax}{\fontsize{6.6pt}{7.7pt}\selectfont\color{black!84}Suite directe de la chevauchée du 19 août 1916 avec le lieutenant Curtis: par la blanchisserie Barbry et \textbf{Fleurbaix} (Pas-de-Calais, village du front criblé d'obus), ils gagnent \textbf{Armentières} (Nord), grande cité linière évacuée en partie sous les bombardements allemands (dont l'église de briques rouges au clocher découronné et à la toiture abattue correspond à l'église \textbf{Saint-Vaast} ou \textbf{Notre-Dame du Sacré-Cœur} d'Armentières).}\par}
\newpage
\phantomsection\label{fac:22}
% --- En-tête du fac-similé (Page impaire / droite : Photographie du manuscrit) ---
\noindent\rule{\textwidth}{0.5pt}\par\vspace{0.55mm}
\noindent\textbf{\fontsize{6.8pt}{7.8pt}\selectfont FAC-SIMILÉ EN VIS-À-VIS — FOLIO 11 RECTO} \hfill \textbf{\fontsize{6.8pt}{7.8pt}\selectfont AOÛT 1916}\par\vspace{0.25mm}
\noindent{\fontsize{6.8pt}{8.0pt}\selectfont\textit{Livre-Journal d'Hippolyte Combes} — Samedi 19 août 1916 (suite) (Erquinghem-Lys, Fleurbaix et Armentières (Pas-de-Calais / Nord))}\par\vspace{0.55mm}
\noindent\rule{\textwidth}{0.35pt}\par\vspace*{\fill}
{\centering
  \setlength{\fboxsep}{1.2mm}%
  \setlength{\fboxrule}{0.4pt}%
  \fcolorbox{gray!55}{white}{%
    \includegraphics[width=0.95\textwidth,height=0.81\textheight,keepaspectratio]{\facsimiledir/20251126_093857.jpg}%
  }\par}
\vspace*{\fill}
\newpage

% =============================================================================
% VIS-A-VIS SPREAD : FOLIO 11 VERSO (20251126_093922.jpg)
% =============================================================================
\cleardoubleevenpage
\renewcommand{\currentmonthyear}{Août 1916}
\renewcommand{\currentfoliolabel}{Folio 11 verso}
\phantomsection\label{fol:23}
% --- En-tête du folio (Page paire / gauche : Transcription, apparat & registres de bas de page) ---
\noindent\rule{\textwidth}{0.5pt}\par\vspace{0.55mm}
\noindent\textbf{\footnotesize FOLIO 11 VERSO} \hfill \textbf{\footnotesize AOÛT 1916}\par\vspace{0.25mm}
\noindent{\fontsize{6.9pt}{8.0pt}\selectfont Café chez M. Barbry avec les adjudants Lupin et Schumacher et préparatifs de départ}\par\vspace{0.2mm}
\noindent{\fontsize{6.9pt}{8.0pt}\selectfont\textit{Dimanche 20 et lundi 21 août 1916} \hfill \textit{Erquinghem-Lys et Steenwerck (Nord)}}\par\vspace{0.5mm}
\noindent\rule{\textwidth}{0.35pt}\par\vspace{0.8mm}
% --- Corps de la transcription diplomatique et apparat marginal en prolongement ---
{\renewcommand{\arraystretch}{0.87}%
\noindent\hspace*{3.5mm}\begin{tabular}{@{}p{87mm}@{\hspace{3.5mm}}p{39mm}@{}}
      \leavevmode\fontsize{6.9pt}{7.9pt}\selectfont\strut \underline{20 D} – \hfill{} \underline{Aout 1916}\strut &  \\
      \leavevmode\fontsize{6.9pt}{7.9pt}\selectfont\strut Tour à tour averses \& sunshine\strut & \smash{\parbox[t]{39mm}{\fontsize{6.2pt}{7.0pt}\selectfont\raggedright\color{black!88}\textbf{2 sunshine} \textit{angl. in ms.}}} \\
      \leavevmode\fontsize{6.9pt}{7.9pt}\selectfont\strut Partie de l'après-midi : assez beau –\strut &  \\
      \leavevmode\fontsize{6.9pt}{7.9pt}\selectfont\strut Avec cet idiot de Ctis, to the fields for\strut & \smash{\parbox[t]{39mm}{\fontsize{6.2pt}{7.0pt}\selectfont\raggedright\color{black!88}\textbf{4 Ctis} \textit{abbr. pro} Curtis.}} \\
      \leavevmode\llap{\fontsize{6.4pt}{7.3pt}\selectfont\color{darkgray}5\hspace{1.8mm}}\fontsize{6.9pt}{7.9pt}\selectfont\strut green forage, puis, café et rhum chez\strut & \smash{\parbox[t]{39mm}{\fontsize{6.2pt}{7.0pt}\selectfont\raggedright\color{black!88}\textbf{5 green} \textit{angl. in ms.}}} \\
      \leavevmode\fontsize{6.9pt}{7.9pt}\selectfont\strut le propre, M. Barbry — y $[\![$[illis.]$]\!]$ amené\strut & \smash{\parbox[t]{39mm}{\fontsize{6.2pt}{7.0pt}\selectfont\raggedright\color{black!88}\textbf{6} $[\![$[illis.]$]\!]\ $ \textit{del. in ms.} \textbf{6} \textit{lect. incert. in ms.}}} \\
      \leavevmode\fontsize{6.9pt}{7.9pt}\selectfont\strut adjudants Lupin \& Schumacher\strut &  \\
      \leavevmode\fontsize{6.9pt}{7.9pt}\selectfont\strut Bien reposé chambre fin journée\strut &  \\
      \leavevmode\fontsize{6.9pt}{7.9pt}\selectfont\strut Moral devenu excellent. Réagi contre\strut &  \\
      \leavevmode\llap{\fontsize{6.4pt}{7.3pt}\selectfont\color{darkgray}10\hspace{1.8mm}}\fontsize{6.9pt}{7.9pt}\selectfont\strut dépression des jours précédents. Me $[\![$[illis.]$]\!]$\strut & \smash{\parbox[t]{39mm}{\fontsize{6.2pt}{7.0pt}\selectfont\raggedright\color{black!88}\textbf{10} $[\![$[illis.]$]\!]\ $ \textit{del. in ms.} \textbf{10} \textit{lect. incert. in ms.}}} \\
      \leavevmode\fontsize{6.9pt}{7.9pt}\selectfont\strut prépare déjà départ semaine prochaine\strut &  \\
      \leavevmode\fontsize{6.9pt}{7.9pt}\selectfont\strut de ma fenêtre : haute \& large, découvre\strut &  \\
      \leavevmode\fontsize{6.9pt}{7.9pt}\selectfont\strut gd horizon : prairies, champs, boqueteaux\strut & \smash{\parbox[t]{39mm}{\fontsize{6.2pt}{7.0pt}\selectfont\raggedright\color{black!88}\textbf{13 gd} \textit{abbr. pro} grand.}} \\
      \leavevmode\fontsize{6.9pt}{7.9pt}\selectfont\strut Paysage d’ 1 banalité \& platitude extrême\strut &  \\
      \cline{1-1}
      \leavevmode\llap{\fontsize{6.4pt}{7.3pt}\selectfont\color{darkgray}15\hspace{1.8mm}}\fontsize{6.9pt}{7.9pt}\selectfont\strut \underline{21 L.} Journée very belle. Soleil, brillé\strut & \smash{\parbox[t]{39mm}{\fontsize{6.2pt}{7.0pt}\selectfont\raggedright\color{black!88}\textbf{15 L.} \textit{scil.} Lundi. \textbf{15 very} \textit{angl. in ms.}}} \\
      \leavevmode\fontsize{6.9pt}{7.9pt}\selectfont\strut matin \& soir. Ciel couvert pendant 1 c. temps\strut &  \\
      \leavevmode\fontsize{6.9pt}{7.9pt}\selectfont\strut Fin jour très belle mervx sunset, que\strut & \smash{\parbox[t]{39mm}{\fontsize{6.2pt}{7.0pt}\selectfont\raggedright\color{black!88}\textbf{17 mervx} \textit{abbr. pro} merveilleux. \textbf{17 sunset} \textit{angl. in ms.}}} \\
      \leavevmode\fontsize{6.9pt}{7.9pt}\selectfont\strut j'ai admiré de ma fenêtre. Lignes de\strut &  \\
      \leavevmode\fontsize{6.9pt}{7.9pt}\selectfont\strut nuages frangés d'or étincelant\strut &  \\
      \leavevmode\llap{\fontsize{6.4pt}{7.3pt}\selectfont\color{darkgray}20\hspace{1.8mm}}\fontsize{6.9pt}{7.9pt}\selectfont\strut Allégé mes bagages en\strut &  \\
      \leavevmode\fontsize{6.9pt}{7.9pt}\selectfont\strut expédiant 1 sac à Nogent.\strut &  \\
      \leavevmode\fontsize{6.9pt}{7.9pt}\selectfont\strut Matin, venu à cheval à Steenwerk\strut &  \\
      \cline{1-1}
      \leavevmode\fontsize{6.9pt}{7.9pt}\selectfont\strut [ — A Proprie : 5 f (cooking, washing, \& eggs)\strut & \smash{\parbox[t]{39mm}{\fontsize{6.2pt}{7.0pt}\selectfont\raggedright\color{black!88}\textbf{23} \textit{in marg. cum unco} (comptes). \textbf{23 washing} \textit{angl. in ms.} — \textbf{24 groom} \textit{angl. in ms.}}} \\
      \leavevmode\fontsize{6.9pt}{7.9pt}\selectfont\strut A mon groom : 2 f – Varia : 10 c –\strut &  \\
\end{tabular}}%
\par\vspace*{\fill}
% --- Registres de bas de page (1. Critique interne [si pertinent] ; 2. Loci paralleli en latin ; 3. Commentaire historique en retrait) ---
\noindent\rule{\textwidth}{0.5pt}\par\vspace{0.65mm}
{\fontsize{6.6pt}{7.7pt}\selectfont%
\noindent\normalfont \textbf{Unité de page :} Le feuillet réunit les journées du dimanche 20 août (l.~1--15) et du lundi 21 août (l.~16--20), séparées par un trait horizontal.\par\vspace{0.5mm}
\noindent\rule{22mm}{0.3pt}\par\vspace{0.45mm}
\noindent\normalfont\color{black!92} \textbf{4} (\textit{Ctis}) cf. supra f\textsuperscript{o}~1~r\textsuperscript{o} (p.~\pageref{fol:2}), f\textsuperscript{o}~10~v\textsuperscript{o} (p.~\pageref{fol:21}), infra f\textsuperscript{o}~15~r\textsuperscript{o} (p.~\pageref{fol:30}), f\textsuperscript{o}~18~v\textsuperscript{o} (p.~\pageref{fol:37}), \textit{et passim}. — \textbf{5} (\textit{green forage}) cf. supra f\textsuperscript{o}~3~r\textsuperscript{o} (p.~\pageref{fol:6}), f\textsuperscript{o}~6~v\textsuperscript{o} (p.~\pageref{fol:13}), infra f\textsuperscript{o}~15~r\textsuperscript{o} (p.~\pageref{fol:30}), f\textsuperscript{o}~15~v\textsuperscript{o} (p.~\pageref{fol:31}), \textit{et passim}. — \textbf{6} (\textit{Barbry}) cf. supra f\textsuperscript{o}~8~v\textsuperscript{o} (p.~\pageref{fol:17}), f\textsuperscript{o}~11~r\textsuperscript{o} (p.~\pageref{fol:22}).\par\vspace{0.55mm}
\noindent\rule{38mm}{0.4pt}\par\vspace{0.5mm}
\noindent\hspace*{5mm}\parbox{\dimexpr\textwidth-5mm\relax}{\fontsize{6.6pt}{7.7pt}\selectfont\color{black!84}Le dimanche 20 août, après la corvée de fourrage vert avec Curtis, Hippolyte Combes prend le café et le rhum chez l'industriel blanchisseur \textbf{M.~Barbry} (propriétaire de la blanchisserie bombardée le 11 août, f\textsuperscript{o}~8~v\textsuperscript{o}) avec les adjudants interprètes français \textbf{Lupin} (ou \textit{Liepin}/\textit{Lépin}, f\textsuperscript{o}~42~v\textsuperscript{o} et f\textsuperscript{o}~45~r\textsuperscript{o}) et \textbf{Schumacher} — ce dernier déjà mentionné avec ironie dans la lettre d'Hippolyte Combes du 21 février 1916 comme ancien contremaître d'usine métallurgique l'ayant remplacé au secrétariat de la Mission française de la 18\textsuperscript{e} Division (Annexe~I, doc.~1).}\par}
\newpage
\phantomsection\label{fac:23}
% --- En-tête du fac-similé (Page impaire / droite : Photographie du manuscrit) ---
\noindent\rule{\textwidth}{0.5pt}\par\vspace{0.55mm}
\noindent\textbf{\fontsize{6.8pt}{7.8pt}\selectfont FAC-SIMILÉ EN VIS-À-VIS — FOLIO 11 VERSO} \hfill \textbf{\fontsize{6.8pt}{7.8pt}\selectfont AOÛT 1916}\par\vspace{0.25mm}
\noindent{\fontsize{6.8pt}{8.0pt}\selectfont\textit{Livre-Journal d'Hippolyte Combes} — Dimanche 20 et lundi 21 août 1916 (Erquinghem-Lys et Steenwerck (Nord))}\par\vspace{0.55mm}
\noindent\rule{\textwidth}{0.35pt}\par\vspace*{\fill}
{\centering
  \setlength{\fboxsep}{1.2mm}%
  \setlength{\fboxrule}{0.4pt}%
  \fcolorbox{gray!55}{white}{%
    \includegraphics[width=0.95\textwidth,height=0.81\textheight,keepaspectratio]{\facsimiledir/20251126_093922.jpg}%
  }\par}
\vspace*{\fill}
\newpage

% =============================================================================
% VIS-A-VIS SPREAD : FOLIO 12 RECTO (20251126_093941.jpg)
% =============================================================================
\cleardoubleevenpage
\renewcommand{\currentmonthyear}{Août 1916}
\renewcommand{\currentfoliolabel}{Folio 12 recto}
\phantomsection\label{fol:24}
% --- En-tête du folio (Page paire / gauche : Transcription, apparat & registres de bas de page) ---
\noindent\rule{\textwidth}{0.5pt}\par\vspace{0.55mm}
\noindent\textbf{\footnotesize FOLIO 12 RECTO} \hfill \textbf{\footnotesize AOÛT 1916}\par\vspace{0.25mm}
\noindent{\fontsize{6.9pt}{8.0pt}\selectfont Mardi 22 août 1916 — Bière chez un charcutier avec Maurin évoquant les beautés d'Antibes et envoi d’un sac de linge à Nogent}\par\vspace{0.2mm}
\noindent{\fontsize{6.9pt}{8.0pt}\selectfont\textit{Mardi 22 août 1916} \hfill \textit{Steenwerck, Bac-Saint-Maur et Erquinghem-Lys (Nord)}}\par\vspace{0.5mm}
\noindent\rule{\textwidth}{0.35pt}\par\vspace{0.8mm}
% --- Corps de la transcription diplomatique et apparat marginal en prolongement ---
{\renewcommand{\arraystretch}{0.87}%
\noindent\hspace*{3.5mm}\begin{tabular}{@{}p{87mm}@{\hspace{3.5mm}}p{39mm}@{}}
      \leavevmode\fontsize{6.9pt}{7.9pt}\selectfont\strut \underline{22 Ma} \hfill{} \underline{Aout / 16}\strut &  \\
      \leavevmode\fontsize{6.9pt}{7.9pt}\selectfont\strut Ciel nuageux – s'est tt à fait éclairci\strut & \smash{\parbox[t]{39mm}{\fontsize{6.2pt}{7.0pt}\selectfont\raggedright\color{black!88}\textbf{2 tt} \textit{abbr. pro} tout.}} \\
      \leavevmode\fontsize{6.9pt}{7.9pt}\selectfont\strut après-midi et très beau soleil\strut &  \\
      \leavevmode\fontsize{6.9pt}{7.9pt}\selectfont\strut Venu à pied à Steenwerck : 55 m.\strut &  \\
      \leavevmode\llap{\fontsize{6.4pt}{7.3pt}\selectfont\color{darkgray}5\hspace{1.8mm}}\fontsize{6.9pt}{7.9pt}\selectfont\strut Promenade lente, très agréable\strut &  \\
      \leavevmode\fontsize{6.9pt}{7.9pt}\selectfont\strut Vers 5H 1/2, retrouvé Maurin, et allés\strut &  \\
      \leavevmode\fontsize{6.9pt}{7.9pt}\selectfont\strut ensemble à Bac St-Maur. En savourant\strut & \smash{\parbox[t]{39mm}{\fontsize{6.2pt}{7.0pt}\selectfont\raggedright\color{black!88}\textbf{7 St} \textit{abbr. pro} Saint.}} \\
      \leavevmode\fontsize{6.9pt}{7.9pt}\selectfont\strut bonne bière ds coquette salle à manger\strut & \smash{\parbox[t]{39mm}{\fontsize{6.2pt}{7.0pt}\selectfont\raggedright\color{black!88}\textbf{8 ds} \textit{abbr. pro} dans.}} \\
      \leavevmode\fontsize{6.9pt}{7.9pt}\selectfont\strut chz 1 charcutier, causé Provence. Maurin\strut &  \\
      \leavevmode\llap{\fontsize{6.4pt}{7.3pt}\selectfont\color{darkgray}10\hspace{1.8mm}}\fontsize{6.9pt}{7.9pt}\selectfont\strut m'a décrit av. lyrisme beautés d'Antibes.\strut &  \\
      \leavevmode\fontsize{6.9pt}{7.9pt}\selectfont\strut Magnifique sunset\strut & \smash{\parbox[t]{39mm}{\fontsize{6.2pt}{7.0pt}\selectfont\raggedright\color{black!88}\textbf{11 sunset} \textit{angl. in ms.}}} \\
      \leavevmode\fontsize{6.9pt}{7.9pt}\selectfont\strut — Fait expédier par 1 commre d'Erquinghem\strut & \smash{\parbox[t]{39mm}{\fontsize{6.2pt}{7.0pt}\selectfont\raggedright\color{black!88}\textbf{12 commre} \textit{abbr. pro} commissionnaire.}} \\
      \leavevmode\fontsize{6.9pt}{7.9pt}\selectfont\strut 1 sac linge à Nogent, par gare de\strut &  \\
      \leavevmode\fontsize{6.9pt}{7.9pt}\selectfont\strut Steenwerck : 13 kgs; 3f 80; factage à l'arrivée\strut &  \\
      \leavevmode\llap{\fontsize{6.4pt}{7.3pt}\selectfont\color{darkgray}15\hspace{1.8mm}}\fontsize{6.9pt}{7.9pt}\selectfont\strut compris (40c). – Emballage au départ : 50c\strut &  \\
      \cline{1-1}
      \leavevmode\fontsize{6.9pt}{7.9pt}\selectfont\strut Cartes ples à Steenwerck : 30 c\strut & \smash{\parbox[t]{39mm}{\fontsize{6.2pt}{7.0pt}\selectfont\raggedright\color{black!88}\textbf{16 Cartes ples} \textit{abbr. pro} Cartes postales.}} \\
      \leavevmode\fontsize{6.9pt}{7.9pt}\selectfont\strut Varia : 10 c\strut &  \\
      \cline{1-1}
      \leavevmode\fontsize{6.9pt}{7.9pt}\selectfont\strut Bain : 1 f 25\strut &  \\
\end{tabular}}%
\par\vspace*{\fill}
% --- Registres de bas de page (1. Critique interne [si pertinent] ; 2. Loci paralleli en latin ; 3. Commentaire historique en retrait) ---
\noindent\rule{\textwidth}{0.5pt}\par\vspace{0.65mm}
{\fontsize{6.6pt}{7.7pt}\selectfont%
\noindent\normalfont\color{black!92} \textbf{4} (\textit{Steenwerck}) cf. supra f\textsuperscript{o}~5~r\textsuperscript{o} (p.~\pageref{fol:10}), f\textsuperscript{o}~11~v\textsuperscript{o} (p.~\pageref{fol:23}), infra f\textsuperscript{o}~13~r\textsuperscript{o} (p.~\pageref{fol:26}), f\textsuperscript{o}~13~v\textsuperscript{o} (p.~\pageref{fol:27}), \textit{et passim}. — \textbf{6} (\textit{Maurin}) cf. supra f\textsuperscript{o}~6~v\textsuperscript{o} (p.~\pageref{fol:13}), f\textsuperscript{o}~10~v\textsuperscript{o} (p.~\pageref{fol:21}), infra f\textsuperscript{o}~13~r\textsuperscript{o} (p.~\pageref{fol:26}), \textit{et passim}. — \textbf{12} (\textit{Erquinghem}) cf. supra f\textsuperscript{o}~8~v\textsuperscript{o} (p.~\pageref{fol:17}), f\textsuperscript{o}~11~v\textsuperscript{o} (p.~\pageref{fol:23}).\par\vspace{0.55mm}
\noindent\rule{38mm}{0.4pt}\par\vspace{0.5mm}
\noindent\hspace*{5mm}\parbox{\dimexpr\textwidth-5mm\relax}{\fontsize{6.6pt}{7.7pt}\selectfont\color{black!84}Entretien amical chez un charcutier de \textbf{Bac-Saint-Maur} avec son collègue \textbf{Maurin} évoquant la Provence et Antibes, et détail tarifaire de l'expédition ferroviaire en petite vitesse d'un sac de 13~kg de linge depuis la gare de \textbf{Steenwerck} jusqu'à \textbf{Nogent-le-Rotrou} (3~fr.~80 factage compris, plus 0~fr.~50 de timbre de récépissé).}\par}
\newpage
\phantomsection\label{fac:24}
% --- En-tête du fac-similé (Page impaire / droite : Photographie du manuscrit) ---
\noindent\rule{\textwidth}{0.5pt}\par\vspace{0.55mm}
\noindent\textbf{\fontsize{6.8pt}{7.8pt}\selectfont FAC-SIMILÉ EN VIS-À-VIS — FOLIO 12 RECTO} \hfill \textbf{\fontsize{6.8pt}{7.8pt}\selectfont AOÛT 1916}\par\vspace{0.25mm}
\noindent{\fontsize{6.8pt}{8.0pt}\selectfont\textit{Livre-Journal d'Hippolyte Combes} — Mardi 22 août 1916 (Steenwerck, Bac-Saint-Maur et Erquinghem-Lys (Nord))}\par\vspace{0.55mm}
\noindent\rule{\textwidth}{0.35pt}\par\vspace*{\fill}
{\centering
  \setlength{\fboxsep}{1.2mm}%
  \setlength{\fboxrule}{0.4pt}%
  \fcolorbox{gray!55}{white}{%
    \includegraphics[width=0.95\textwidth,height=0.81\textheight,keepaspectratio]{\facsimiledir/20251126_093941.jpg}%
  }\par}
\vspace*{\fill}
\newpage

% =============================================================================
% VIS-A-VIS SPREAD : FOLIO 12 VERSO (20251126_093955.jpg)
% =============================================================================
\cleardoubleevenpage
\renewcommand{\currentmonthyear}{Août 1916}
\renewcommand{\currentfoliolabel}{Folio 12 verso}
\phantomsection\label{fol:25}
% --- En-tête du folio (Page paire / gauche : Transcription, apparat & registres de bas de page) ---
\noindent\rule{\textwidth}{0.5pt}\par\vspace{0.55mm}
\noindent\textbf{\footnotesize FOLIO 12 VERSO} \hfill \textbf{\footnotesize AOÛT 1916}\par\vspace{0.25mm}
\noindent{\fontsize{6.9pt}{8.0pt}\selectfont Chevauchée jusqu'à Estaires : l'église Saint-Vaast, l'Hôtel de Ville et l'hygiène par l'exercice}\par\vspace{0.2mm}
\noindent{\fontsize{6.9pt}{8.0pt}\selectfont\textit{Mercredi 23 août 1916} \hfill \textit{Estaires (Nord)}}\par\vspace{0.5mm}
\noindent\rule{\textwidth}{0.35pt}\par\vspace{0.8mm}
% --- Corps de la transcription diplomatique et apparat marginal en prolongement ---
{\renewcommand{\arraystretch}{0.93}%
\noindent\hspace*{3.5mm}\begin{tabular}{@{}p{87mm}@{\hspace{3.5mm}}p{39mm}@{}}
      \leavevmode\fontsize{7.3pt}{8.4pt}\selectfont\strut \underline{23 Merc} \hfill{} \underline{Aout / 16}\strut & \smash{\parbox[t]{39mm}{\fontsize{6.2pt}{7.0pt}\selectfont\raggedright\color{black!88}\textbf{1 Merc} \textit{scil.} Mercredi.}} \\
      \leavevmode\fontsize{7.3pt}{8.4pt}\selectfont\strut Temps couvert, matinée et surtout\strut &  \\
      \leavevmode\fontsize{7.3pt}{8.4pt}\selectfont\strut fin journée – Dès 6H 1/2, cme 1 légère\strut & \smash{\parbox[t]{39mm}{\fontsize{6.2pt}{7.0pt}\selectfont\raggedright\color{black!88}\textbf{3 cme} \textit{abbr. pro} comme. \textbf{3 1/2} \textit{scil.} demi / moitié. — \textbf{4 ds} \textit{abbr. pro} dans.}} \\
      \leavevmode\fontsize{7.3pt}{8.4pt}\selectfont\strut brume répandue ds l'air. –\strut &  \\
      \leavevmode\llap{\fontsize{6.4pt}{7.3pt}\selectfont\color{darkgray}5\hspace{1.8mm}}\fontsize{7.3pt}{8.4pt}\selectfont\strut Après-midi, promenade à cheval\strut &  \\
      \leavevmode\fontsize{7.3pt}{8.4pt}\selectfont\strut arrivé j-à Estaires. Visité longuement\strut & \smash{\parbox[t]{39mm}{\fontsize{6.2pt}{7.0pt}\selectfont\raggedright\color{black!88}\textbf{6 j-à} \textit{abbr. pro} jusqu'à.}} \\
      \leavevmode\fontsize{7.3pt}{8.4pt}\selectfont\strut cette petite ville, très animée \& x magasins\strut &  \\
      \leavevmode\fontsize{7.3pt}{8.4pt}\selectfont\strut Très belle église (chœur, stalles ogivales\strut &  \\
      \leavevmode\fontsize{7.3pt}{8.4pt}\selectfont\strut et très artistiques statuettes en bois (S. Ambroise etc.)\strut &  \\
      \leavevmode\llap{\fontsize{6.4pt}{7.3pt}\selectfont\color{darkgray}10\hspace{1.8mm}}\fontsize{7.3pt}{8.4pt}\selectfont\strut Remarqué hotel de ville – écusson armoiries\strut &  \\
      \leavevmode\fontsize{7.3pt}{8.4pt}\selectfont\strut peint mur au dessus du perron –\strut &  \\
      \leavevmode\fontsize{7.3pt}{8.4pt}\selectfont\strut De nouveau\strut &  \\
      \leavevmode\fontsize{7.3pt}{8.4pt}\selectfont\strut Poussière routes desséchées.\strut &  \\
      \leavevmode\fontsize{7.3pt}{8.4pt}\selectfont\strut Environs extt plats — \& vulgaires —\strut & \smash{\parbox[t]{39mm}{\fontsize{6.2pt}{7.0pt}\selectfont\raggedright\color{black!88}\textbf{14 extt} \textit{abbr. pro} extrêmement.}} \\
      \leavevmode\llap{\fontsize{6.4pt}{7.3pt}\selectfont\color{darkgray}15\hspace{1.8mm}}\fontsize{7.3pt}{8.4pt}\selectfont\strut aucun intérêt pittoresque.\strut &  \\
      \leavevmode\fontsize{7.3pt}{8.4pt}\selectfont\strut Campagne extt monotone \& laide.\strut & \smash{\parbox[t]{39mm}{\fontsize{6.2pt}{7.0pt}\selectfont\raggedright\color{black!88}\textbf{16 extt} \textit{abbr. pro} extrêmement.}} \\
      \leavevmode\fontsize{7.3pt}{8.4pt}\selectfont\strut Dépense d'activité \& fatigues physiques\strut &  \\
      \leavevmode\fontsize{7.3pt}{8.4pt}\selectfont\strut ft [?] de bien à mon moral —\strut & \smash{\parbox[t]{39mm}{\fontsize{6.2pt}{7.0pt}\selectfont\raggedright\color{black!88}\textbf{18 ft} \textit{lect. incert. in ms.}}} \\
      \leavevmode\fontsize{7.3pt}{8.4pt}\selectfont\strut Sensation d'équilibre \& santé parfaite\strut &  \\
      \cline{1-1}
      \leavevmode\llap{\fontsize{6.4pt}{7.3pt}\selectfont\color{darkgray}20\hspace{1.8mm}}\fontsize{7.3pt}{8.4pt}\selectfont\strut Ctes ples, à Estaires : 35 c\strut & \smash{\parbox[t]{39mm}{\fontsize{6.2pt}{7.0pt}\selectfont\raggedright\color{black!88}\textbf{20 Ctes ples} \textit{abbr. pro} Cartes postales.}} \\
      \leavevmode\fontsize{7.3pt}{8.4pt}\selectfont\strut ¼ livre bonbons chocolat : 75 c –\strut &  \\
\end{tabular}}%
\par\vspace*{\fill}
% --- Registres de bas de page (1. Critique interne [si pertinent] ; 2. Loci paralleli en latin ; 3. Commentaire historique en retrait) ---
\noindent\rule{\textwidth}{0.5pt}\par\vspace{0.65mm}
{\fontsize{6.6pt}{7.7pt}\selectfont%
\noindent\normalfont\color{black!92} \textbf{6} (\textit{Estaires}) cf. supra f\textsuperscript{o}~8~r\textsuperscript{o} (p.~\pageref{fol:16}).\par\vspace{0.55mm}
\noindent\rule{38mm}{0.4pt}\par\vspace{0.5mm}
\noindent\hspace*{5mm}\parbox{\dimexpr\textwidth-5mm\relax}{\fontsize{6.6pt}{7.7pt}\selectfont\color{black!84}Promenade à cheval jusqu'à la petite ville commerçante d'\textbf{Estaires} (Nord, sur la Lys).}\par}
\newpage
\phantomsection\label{fac:25}
% --- En-tête du fac-similé (Page impaire / droite : Photographie du manuscrit) ---
\noindent\rule{\textwidth}{0.5pt}\par\vspace{0.55mm}
\noindent\textbf{\fontsize{6.8pt}{7.8pt}\selectfont FAC-SIMILÉ EN VIS-À-VIS — FOLIO 12 VERSO} \hfill \textbf{\fontsize{6.8pt}{7.8pt}\selectfont AOÛT 1916}\par\vspace{0.25mm}
\noindent{\fontsize{6.8pt}{8.0pt}\selectfont\textit{Livre-Journal d'Hippolyte Combes} — Mercredi 23 août 1916 (Estaires (Nord))}\par\vspace{0.55mm}
\noindent\rule{\textwidth}{0.35pt}\par\vspace*{\fill}
{\centering
  \setlength{\fboxsep}{1.2mm}%
  \setlength{\fboxrule}{0.4pt}%
  \fcolorbox{gray!55}{white}{%
    \includegraphics[width=0.95\textwidth,height=0.81\textheight,keepaspectratio]{\facsimiledir/20251126_093955.jpg}%
  }\par}
\vspace*{\fill}
\newpage

% =============================================================================
% VIS-A-VIS SPREAD : FOLIO 13 RECTO (20251126_094019.jpg)
% =============================================================================
\cleardoubleevenpage
\renewcommand{\currentmonthyear}{Août 1916}
\renewcommand{\currentfoliolabel}{Folio 13 recto}
\phantomsection\label{fol:26}
% --- En-tête du folio (Page paire / gauche : Transcription, apparat & registres de bas de page) ---
\noindent\rule{\textwidth}{0.5pt}\par\vspace{0.55mm}
\noindent\textbf{\footnotesize FOLIO 13 RECTO} \hfill \textbf{\footnotesize AOÛT 1916}\par\vspace{0.25mm}
\noindent{\fontsize{6.9pt}{8.0pt}\selectfont Jeudi 24 août 1916 — Église de Steenwerck, transport des sacs sous la tente de l'office par Fisher et adieux aux Flandres}\par\vspace{0.2mm}
\noindent{\fontsize{6.9pt}{8.0pt}\selectfont\textit{Jeudi 24 août 1916} \hfill \textit{Steenwerck (Nord)}}\par\vspace{0.5mm}
\noindent\rule{\textwidth}{0.35pt}\par\vspace{0.8mm}
% --- Corps de la transcription diplomatique et apparat marginal en prolongement ---
{\renewcommand{\arraystretch}{1.00}%
\noindent\hspace*{3.5mm}\begin{tabular}{@{}p{87mm}@{\hspace{3.5mm}}p{39mm}@{}}
      \leavevmode\fontsize{7.7pt}{9.0pt}\selectfont\strut \underline{Jeudi 24} \hfill{} \underline{Août 1916}\strut &  \\
      \leavevmode\fontsize{7.7pt}{9.0pt}\selectfont\strut Matinée temps nuageux –\strut &  \\
      \leavevmode\fontsize{7.7pt}{9.0pt}\selectfont\strut Après-midi, vent éclairci ciel \& brillé très\strut &  \\
      \leavevmode\fontsize{7.7pt}{9.0pt}\selectfont\strut chaud soleil.\strut &  \\
      \leavevmode\llap{\fontsize{6.4pt}{7.3pt}\selectfont\color{darkgray}5\hspace{1.8mm}}\fontsize{7.7pt}{9.0pt}\selectfont\strut Venu à pied à Steenwerck. Entré église pr me\strut & \smash{\parbox[t]{39mm}{\fontsize{6.2pt}{7.0pt}\selectfont\raggedright\color{black!88}\textbf{5 pr} \textit{abbr. pro} pour.}} \\
      \leavevmode\fontsize{7.7pt}{9.0pt}\selectfont\strut reposer 1 peu \& y découvert relique de St Hippte,\strut & \smash{\parbox[t]{39mm}{\fontsize{6.2pt}{7.0pt}\selectfont\raggedright\color{black!88}\textbf{6 Hippte} \textit{scil.} Hippolyte. \textbf{6 St} \textit{abbr. pro} Saint.}} \\
      \leavevmode\fontsize{7.7pt}{9.0pt}\selectfont\strut mon patron!\strut &  \\
      \leavevmode\fontsize{7.7pt}{9.0pt}\selectfont\strut Revu Maurin \& causé avec lui, longuement\strut &  \\
      \leavevmode\fontsize{7.7pt}{9.0pt}\selectfont\strut Vers 7 H, Fisher venu prendre ts mes sacs\strut & \smash{\parbox[t]{39mm}{\fontsize{6.2pt}{7.0pt}\selectfont\raggedright\color{black!88}\textbf{9 ts} \textit{abbr. pro} tous.}} \\
      \leavevmode\llap{\fontsize{6.4pt}{7.3pt}\selectfont\color{darkgray}10\hspace{1.8mm}}\fontsize{7.7pt}{9.0pt}\selectfont\strut et les a portés sous la tente de l'office\strut &  \\
      \leavevmode\fontsize{7.7pt}{9.0pt}\selectfont\strut ds le champ —\strut & \smash{\parbox[t]{39mm}{\fontsize{6.2pt}{7.0pt}\selectfont\raggedright\color{black!88}\textbf{11 ds} \textit{abbr. pro} dans.}} \\
      \leavevmode\fontsize{7.7pt}{9.0pt}\selectfont\strut Impression de mélancolie, au moment\strut &  \\
      \leavevmode\fontsize{7.7pt}{9.0pt}\selectfont\strut de m'éloigner de ce village pour toujours.\strut &  \\
      \leavevmode\fontsize{7.7pt}{9.0pt}\selectfont\strut Steenwerck : ctes ples : 1 f. — Varia : 1 f\strut & \smash{\parbox[t]{39mm}{\fontsize{6.2pt}{7.0pt}\selectfont\raggedright\color{black!88}\textbf{14 ctes ples} \textit{abbr. pro} cartes postales. \textbf{14 f.} \textit{abbr. pro} franc(s). — \textbf{16 Tip} \textit{angl. in ms.} — \textbf{17 pr} \textit{abbr. pro} pour. \textbf{17 Mme} \textit{abbr. pro} Madame. — \textbf{18 fr.} \textit{abbr. pro} franc(s). \textbf{18 washing} \textit{angl. in ms.}}} \\
      \leavevmode\llap{\fontsize{6.4pt}{7.3pt}\selectfont\color{darkgray}15\hspace{1.8mm}}\fontsize{7.7pt}{9.0pt}\selectfont\strut Matin, réparation à 1 soulier : 40c\strut &  \\
      \cline{1-1}
      \leavevmode\fontsize{7.7pt}{9.0pt}\selectfont\strut Tip à Fisher : 50 c\strut &  \\
      \leavevmode\fontsize{7.7pt}{9.0pt}\selectfont\strut Donné à Mme Cliquennois, pr cooking\strut &  \\
      \leavevmode\fontsize{7.7pt}{9.0pt}\selectfont\strut washing et divers : 4 fr. —\strut &  \\
\end{tabular}}%
\par\vspace*{\fill}
% --- Registres de bas de page (1. Critique interne [si pertinent] ; 2. Loci paralleli en latin ; 3. Commentaire historique en retrait) ---
\noindent\rule{\textwidth}{0.5pt}\par\vspace{0.65mm}
{\fontsize{6.6pt}{7.7pt}\selectfont%
\noindent\normalfont\color{black!92} \textbf{5} (\textit{Steenwerck}) cf. supra f\textsuperscript{o}~5~r\textsuperscript{o} (p.~\pageref{fol:10}), f\textsuperscript{o}~12~r\textsuperscript{o} (p.~\pageref{fol:24}), infra f\textsuperscript{o}~13~v\textsuperscript{o} (p.~\pageref{fol:27}), \textit{et passim}. — \textbf{8} (\textit{Maurin}) cf. supra f\textsuperscript{o}~6~v\textsuperscript{o} (p.~\pageref{fol:13}), f\textsuperscript{o}~12~r\textsuperscript{o} (p.~\pageref{fol:24}), \textit{et passim}.\par\vspace{0.55mm}
\noindent\rule{38mm}{0.4pt}\par\vspace{0.5mm}
\noindent\hspace*{5mm}\parbox{\dimexpr\textwidth-5mm\relax}{\fontsize{6.6pt}{7.7pt}\selectfont\color{black!84}À la veille du départ de la 18\textsuperscript{e} Division vers l'Artois et la Somme, Hippolyte Combes entre dans l'église \textbf{Saint-Jean-Baptiste de Steenwerck} (édifice qui sera détruit lors de l'offensive allemande d'avril 1918) et y remarque un reliquaire de son saint patron, saint Hippolyte Combes.}\par}
\newpage
\phantomsection\label{fac:26}
% --- En-tête du fac-similé (Page impaire / droite : Photographie du manuscrit) ---
\noindent\rule{\textwidth}{0.5pt}\par\vspace{0.55mm}
\noindent\textbf{\fontsize{6.8pt}{7.8pt}\selectfont FAC-SIMILÉ EN VIS-À-VIS — FOLIO 13 RECTO} \hfill \textbf{\fontsize{6.8pt}{7.8pt}\selectfont AOÛT 1916}\par\vspace{0.25mm}
\noindent{\fontsize{6.8pt}{8.0pt}\selectfont\textit{Livre-Journal d'Hippolyte Combes} — Jeudi 24 août 1916 (Steenwerck (Nord))}\par\vspace{0.55mm}
\noindent\rule{\textwidth}{0.35pt}\par\vspace*{\fill}
{\centering
  \setlength{\fboxsep}{1.2mm}%
  \setlength{\fboxrule}{0.4pt}%
  \fcolorbox{gray!55}{white}{%
    \includegraphics[width=0.95\textwidth,height=0.81\textheight,keepaspectratio]{\facsimiledir/20251126_094019.jpg}%
  }\par}
\vspace*{\fill}
\newpage

% =============================================================================
% VIS-A-VIS SPREAD : FOLIO 13 VERSO (20251126_094045.jpg)
% =============================================================================
\cleardoubleevenpage
\renewcommand{\currentmonthyear}{Août 1916}
\renewcommand{\currentfoliolabel}{Folio 13 verso}
\phantomsection\label{fol:27}
% --- En-tête du folio (Page paire / gauche : Transcription, apparat & registres de bas de page) ---
\noindent\rule{\textwidth}{0.5pt}\par\vspace{0.55mm}
\noindent\textbf{\footnotesize FOLIO 13 VERSO} \hfill \textbf{\footnotesize AOÛT 1916}\par\vspace{0.25mm}
\noindent{\fontsize{6.9pt}{8.0pt}\selectfont Embarquement ferroviaire à Bailleul et descente par Hazebrouck et Lillers vers Saint-Pol-sur-Ternoise}\par\vspace{0.2mm}
\noindent{\fontsize{6.9pt}{8.0pt}\selectfont\textit{Vendredi 25 août 1916} \hfill \textit{De Steenwerck et Bailleul à Saint-Pol-sur-Ternoise (Pas-de-Calais)}}\par\vspace{0.5mm}
\noindent\rule{\textwidth}{0.35pt}\par\vspace{0.8mm}
% --- Corps de la transcription diplomatique et apparat marginal en prolongement ---
{\renewcommand{\arraystretch}{0.87}%
\noindent\hspace*{3.5mm}\begin{tabular}{@{}p{87mm}@{\hspace{3.5mm}}p{39mm}@{}}
      \leavevmode\fontsize{6.9pt}{7.9pt}\selectfont\strut \underline{Ven. 25} \hfill{} \underline{Août 1916}\strut & \smash{\parbox[t]{39mm}{\fontsize{6.2pt}{7.0pt}\selectfont\raggedright\color{black!88}\textbf{1 Ven.} \textit{scil.} Vendredi.}} \\
      \leavevmode\fontsize{6.9pt}{7.9pt}\selectfont\strut Partis vers 8H m.. – à cheval, soleil s'est\strut & \smash{\parbox[t]{39mm}{\fontsize{6.2pt}{7.0pt}\selectfont\raggedright\color{black!88}\textbf{2 m.} \textit{abbr. pro} matin.}} \\
      \leavevmode\fontsize{6.9pt}{7.9pt}\selectfont\strut levé peu à peu. devint brulant\strut &  \\
      \leavevmode\fontsize{6.9pt}{7.9pt}\selectfont\strut A 11H m. – à Bailleul.\strut & \smash{\parbox[t]{39mm}{\fontsize{6.2pt}{7.0pt}\selectfont\raggedright\color{black!88}\textbf{4 m.} \textit{abbr. pro} matin.}} \\
      \leavevmode\llap{\fontsize{6.4pt}{7.3pt}\selectfont\color{darkgray}5\hspace{1.8mm}}\fontsize{6.9pt}{7.9pt}\selectfont\strut Embarquement chevaux \& wagons terminé\strut &  \\
      \leavevmode\fontsize{6.9pt}{7.9pt}\selectfont\strut à midi 1/2.\strut &  \\
      \leavevmode\fontsize{6.9pt}{7.9pt}\selectfont\strut Train démarre v. 1H 1/2\strut & \smash{\parbox[t]{39mm}{\fontsize{6.2pt}{7.0pt}\selectfont\raggedright\color{black!88}\textbf{7 1/2} \textit{scil.} demi / moitié.}} \\
      \leavevmode\fontsize{6.9pt}{7.9pt}\selectfont\strut Traversé plusies gares imptes sans arrêt\strut & \smash{\parbox[t]{39mm}{\fontsize{6.2pt}{7.0pt}\selectfont\raggedright\color{black!88}\textbf{8 plusies} \textit{abbr. pro} plusieurs. \textbf{8 imptes} \textit{abbr. pro} importantes. — \textbf{9} \textit{lect. incert. in ms.} — \textbf{10 s.} \textit{abbr. pro} soir.}} \\
      \leavevmode\fontsize{6.9pt}{7.9pt}\selectfont\strut A Lillers, vers 2H [45 ?] \hfill{} Hazebrouck : 2H\strut &  \\
      \leavevmode\llap{\fontsize{6.4pt}{7.3pt}\selectfont\color{darkgray}10\hspace{1.8mm}}\fontsize{6.9pt}{7.9pt}\selectfont\strut A St. Pol : à 4H 30 s. \hfill{} Berguette : 2H30\strut &  \\
      \leavevmode\fontsize{6.9pt}{7.9pt}\selectfont\strut Pris route à 5H – d'abord bruine,\strut &  \\
      \leavevmode\fontsize{6.9pt}{7.9pt}\selectfont\strut puis pluie battante\strut &  \\
      \leavevmode\fontsize{6.9pt}{7.9pt}\selectfont\strut Pays pittoresque. Ce n'est plus la\strut &  \\
      \leavevmode\fontsize{6.9pt}{7.9pt}\selectfont\strut campagne plate et monotone du nord.\strut &  \\
      \leavevmode\llap{\fontsize{6.4pt}{7.3pt}\selectfont\color{darkgray}15\hspace{1.8mm}}\fontsize{6.9pt}{7.9pt}\selectfont\strut Région très boisée. Verdures fraîches\strut &  \\
      \leavevmode\fontsize{6.9pt}{7.9pt}\selectfont\strut comme au printemps. Si champs n'étaient\strut &  \\
      \leavevmode\fontsize{6.9pt}{7.9pt}\selectfont\strut pas moissonnés et les gerbes d'or entassées\strut &  \\
      \leavevmode\fontsize{6.9pt}{7.9pt}\selectfont\strut en meules, on ne croirait pas être\strut &  \\
      \leavevmode\fontsize{6.9pt}{7.9pt}\selectfont\strut à la fin de l'été\strut &  \\
\end{tabular}}%
\par\vspace*{\fill}
% --- Registres de bas de page (1. Critique interne [si pertinent] ; 2. Loci paralleli en latin ; 3. Commentaire historique en retrait) ---
\noindent\rule{\textwidth}{0.5pt}\par\vspace{0.65mm}
{\fontsize{6.6pt}{7.7pt}\selectfont%
\noindent\normalfont\color{black!92} \textbf{4} (\textit{Bailleul}) cf. supra f\textsuperscript{o}~3~v\textsuperscript{o} (p.~\pageref{fol:7}), f\textsuperscript{o}~6~r\textsuperscript{o} (p.~\pageref{fol:12}), infra f\textsuperscript{o}~14~r\textsuperscript{o} (p.~\pageref{fol:28}), \textit{et passim}. — \textbf{9} (\textit{Hazebrouck}) cf. supra f\textsuperscript{o}~2~r\textsuperscript{o} (p.~\pageref{fol:4}), f\textsuperscript{o}~8~r\textsuperscript{o} (p.~\pageref{fol:16}).\par\vspace{0.55mm}
\noindent\rule{38mm}{0.4pt}\par\vspace{0.5mm}
\noindent\hspace*{5mm}\parbox{\dimexpr\textwidth-5mm\relax}{\fontsize{6.6pt}{7.7pt}\selectfont\color{black!84}Le vendredi 25 août 1916, le train divisionnaire de la \textbf{18\textsuperscript{e} Division britannique} quitte la vallée de la Lys: embarquée en gare de \textbf{Bailleul} à midi et demi, l'unité descend par la voie ferrée stratégique d'\textbf{Hazebrouck} (14~h), \textbf{Berguette} (14~h~30) et \textbf{Lillers} jusqu'à \textbf{Saint-Pol-sur-Ternoise} (Pas-de-Calais, 16~h~30), où la division se regroupe pour son entraînement avant de redescendre sur le front de la Somme (Nichols, \textit{The 18th Division in the Great War}, 1922, p.~96).}\par}
\newpage
\phantomsection\label{fac:27}
% --- En-tête du fac-similé (Page impaire / droite : Photographie du manuscrit) ---
\noindent\rule{\textwidth}{0.5pt}\par\vspace{0.55mm}
\noindent\textbf{\fontsize{6.8pt}{7.8pt}\selectfont FAC-SIMILÉ EN VIS-À-VIS — FOLIO 13 VERSO} \hfill \textbf{\fontsize{6.8pt}{7.8pt}\selectfont AOÛT 1916}\par\vspace{0.25mm}
\noindent{\fontsize{6.8pt}{8.0pt}\selectfont\textit{Livre-Journal d'Hippolyte Combes} — Vendredi 25 août 1916 (De Steenwerck et Bailleul à Saint-Pol-sur-Ternoise (Pas-de-Calais))}\par\vspace{0.55mm}
\noindent\rule{\textwidth}{0.35pt}\par\vspace*{\fill}
{\centering
  \setlength{\fboxsep}{1.2mm}%
  \setlength{\fboxrule}{0.4pt}%
  \fcolorbox{gray!55}{white}{%
    \includegraphics[width=0.95\textwidth,height=0.81\textheight,keepaspectratio]{\facsimiledir/20251126_094045.jpg}%
  }\par}
\vspace*{\fill}
\newpage

% =============================================================================
% VIS-A-VIS SPREAD : FOLIO 14 RECTO (20251126_094058.jpg)
% =============================================================================
\cleardoubleevenpage
\renewcommand{\currentmonthyear}{Août 1916}
\renewcommand{\currentfoliolabel}{Folio 14 recto}
\phantomsection\label{fol:28}
% --- En-tête du folio (Page paire / gauche : Transcription, apparat & registres de bas de page) ---
\noindent\rule{\textwidth}{0.5pt}\par\vspace{0.55mm}
\noindent\textbf{\footnotesize FOLIO 14 RECTO} \hfill \textbf{\footnotesize AOÛT 1916}\par\vspace{0.25mm}
\noindent{\fontsize{6.9pt}{8.0pt}\selectfont Mélancolie du voyage en train au son de l'harmonica et arrivée nocturne à Marquay}\par\vspace{0.2mm}
\noindent{\fontsize{6.9pt}{8.0pt}\selectfont\textit{Vendredi 25 août 1916 (suite)} \hfill \textit{De Bailleul à Saint-Pol et Marquay (Pas-de-Calais)}}\par\vspace{0.5mm}
\noindent\rule{\textwidth}{0.35pt}\par\vspace{0.8mm}
% --- Corps de la transcription diplomatique et apparat marginal en prolongement ---
{\renewcommand{\arraystretch}{0.87}%
\noindent\hspace*{3.5mm}\begin{tabular}{@{}p{87mm}@{\hspace{3.5mm}}p{39mm}@{}}
      \leavevmode\fontsize{6.9pt}{7.9pt}\selectfont\strut \underline{Ven. 25 (suite)} \hfill{} \raisebox{0.5ex}{\underline{Août 1916}}\strut & \smash{\parbox[t]{39mm}{\fontsize{6.2pt}{7.0pt}\selectfont\raggedright\color{black!88}\textbf{1 Ven.} \textit{scil.} Vendredi.}} \\
      \leavevmode\fontsize{6.9pt}{7.9pt}\selectfont\strut En gare de Bailleul, vers 1H – ds le\strut & \smash{\parbox[t]{39mm}{\fontsize{6.2pt}{7.0pt}\selectfont\raggedright\color{black!88}\textbf{2 ds} \textit{abbr. pro} dans.}} \\
      \leavevmode\fontsize{6.9pt}{7.9pt}\selectfont\strut train, joué par tommies, musique\strut & \smash{\parbox[t]{39mm}{\fontsize{6.2pt}{7.0pt}\selectfont\raggedright\color{black!88}\textbf{3 tommies} \textit{angl. in ms.}}} \\
      \leavevmode\fontsize{6.9pt}{7.9pt}\selectfont\strut d'harmonica, monotone \& banal\strut &  \\
      \leavevmode\llap{\fontsize{6.4pt}{7.3pt}\selectfont\color{darkgray}5\hspace{1.8mm}}\fontsize{6.9pt}{7.9pt}\selectfont\strut sous le ciel pommelé de nuées grises \& blanchâtres\strut &  \\
      \leavevmode\fontsize{6.9pt}{7.9pt}\selectfont\strut Réflexions amères. L'Etat [illis.] mobiliser nos\strut & \smash{\parbox[t]{39mm}{\fontsize{6.2pt}{7.0pt}\selectfont\raggedright\color{black!88}\textbf{6} \textit{lect. incert. in ms.}}} \\
      \leavevmode\fontsize{6.9pt}{7.9pt}\selectfont\strut corps, [illis.] a encore p. notre intellge, notre\strut & \smash{\parbox[t]{39mm}{\fontsize{6.2pt}{7.0pt}\selectfont\raggedright\color{black!88}\textbf{7} \textit{lect. incert. in ms.} \textbf{7 intellge} \textit{abbr. pro} intelligence.}} \\
      \leavevmode\fontsize{6.9pt}{7.9pt}\selectfont\strut cœur —\strut &  \\
      \leavevmode\fontsize{6.9pt}{7.9pt}\selectfont\strut Ciel de + en + bas \& nuageux\strut &  \\
      \leavevmode\llap{\fontsize{6.4pt}{7.3pt}\selectfont\color{darkgray}10\hspace{1.8mm}}\fontsize{6.9pt}{7.9pt}\selectfont\strut orage menaçant. 1 légère brise souffle\strut &  \\
      \leavevmode\fontsize{6.9pt}{7.9pt}\selectfont\strut qui me caresse visage délicieust. Bruit\strut & \smash{\parbox[t]{39mm}{\fontsize{6.2pt}{7.0pt}\selectfont\raggedright\color{black!88}\textbf{11 délicieust} \textit{abbr. pro} délicieusement.}} \\
      \leavevmode\fontsize{6.9pt}{7.9pt}\selectfont\strut doux de feuillages froissés par la bise de\strut &  \\
      \leavevmode\fontsize{6.9pt}{7.9pt}\selectfont\strut tilleuls en bordure voie ferrée –\strut &  \\
      \leavevmode\fontsize{6.9pt}{7.9pt}\selectfont\strut Pluie nous a trempé j’à la peau.\strut &  \\
      \leavevmode\llap{\fontsize{6.4pt}{7.3pt}\selectfont\color{darkgray}15\hspace{1.8mm}}\fontsize{6.9pt}{7.9pt}\selectfont\strut Vers 7H enfin, arrivé à Marquay\strut &  \\
      \leavevmode\fontsize{6.9pt}{7.9pt}\selectfont\strut petit hameau perdu en haut d’ 1 plateau\strut &  \\
      \leavevmode\fontsize{6.9pt}{7.9pt}\selectfont\strut d'Artois. Avec bp de peine, pu trouver\strut & \smash{\parbox[t]{39mm}{\fontsize{6.2pt}{7.0pt}\selectfont\raggedright\color{black!88}\textbf{17 bp} \textit{abbr. pro} beaucoup.}} \\
      \leavevmode\fontsize{6.9pt}{7.9pt}\selectfont\strut chambrette exiguë où coucherai lit à côté de\strut &  \\
      \leavevmode\fontsize{6.9pt}{7.9pt}\selectfont\strut celui du [illis.] d'office\strut & \smash{\parbox[t]{39mm}{\fontsize{6.2pt}{7.0pt}\selectfont\raggedright\color{black!88}\textbf{19} \textit{lect. incert. in ms.}}} \\
      \leavevmode\llap{\fontsize{6.4pt}{7.3pt}\selectfont\color{darkgray}20\hspace{1.8mm}}\fontsize{6.9pt}{7.9pt}\selectfont\strut Invité à diner par offs mess\strut & \smash{\parbox[t]{39mm}{\fontsize{6.2pt}{7.0pt}\selectfont\raggedright\color{black!88}\textbf{20 offs} \textit{abbr. pro} officiers. \textbf{20 mess} \textit{angl. in ms.} — \textbf{21 groom} \textit{angl. in ms.} — \textbf{22 St} \textit{abbr. pro} Saint. \textbf{22 c.} \textit{abbr. pro} centimes.}} \\
      \cline{1-1}
      \leavevmode\fontsize{6.9pt}{7.9pt}\selectfont\strut Tip au groom : 1 f\strut &  \\
      \leavevmode\fontsize{6.9pt}{7.9pt}\selectfont\strut Food : Bailleul : 75 c. – St-Pol : 1 f\strut &  \\
\end{tabular}}%
\par\vspace*{\fill}
% --- Registres de bas de page (1. Critique interne [si pertinent] ; 2. Loci paralleli en latin ; 3. Commentaire historique en retrait) ---
\noindent\rule{\textwidth}{0.5pt}\par\vspace{0.65mm}
{\fontsize{6.6pt}{7.7pt}\selectfont%
\noindent\normalfont \textbf{Continuité :} Feuillet de suite du vendredi 25 août 1916 (\textit{Ven. 25 (suite)}) ; les réflexions personnelles des l.~5--14 sont encadrées de tirets longs caractéristiques de la ponctuation expressive d'Hippolyte Combes.\par\vspace{0.5mm}
\noindent\rule{22mm}{0.3pt}\par\vspace{0.45mm}
\noindent\normalfont\color{black!92} \textbf{2} (\textit{Bailleul}) cf. supra f\textsuperscript{o}~3~v\textsuperscript{o} (p.~\pageref{fol:7}), f\textsuperscript{o}~13~v\textsuperscript{o} (p.~\pageref{fol:27}), \textit{et passim}. — \textbf{15} (\textit{Marquay}) cf. infra f\textsuperscript{o}~14~v\textsuperscript{o} (p.~\pageref{fol:29}), f\textsuperscript{o}~15~r\textsuperscript{o} (p.~\pageref{fol:30}), \textit{et passim}. — \textbf{22} (\textit{St-Pol}) cf. supra f\textsuperscript{o}~13~v\textsuperscript{o} (p.~\pageref{fol:27}), infra f\textsuperscript{o}~16~r\textsuperscript{o} (p.~\pageref{fol:32}), f\textsuperscript{o}~17~r\textsuperscript{o} (p.~\pageref{fol:34}), \textit{et passim}.\par\vspace{0.55mm}
\noindent\rule{38mm}{0.4pt}\par\vspace{0.5mm}
\noindent\hspace*{5mm}\parbox{\dimexpr\textwidth-5mm\relax}{\fontsize{6.6pt}{7.7pt}\selectfont\color{black!84}Ce feuillet complémentaire du 25 août 1916 livre une réflexion critique sur l'emprise morale de l'État en temps de guerre, en écho direct à la lettre confidentielle d'Hippolyte Combes à sa sœur carmélite Françoise du 16 juin 1916 (Annexe~I, doc.~3). Débarquée à Saint-Pol à 17~h sous une pluie battante, la compagnie gagne à 19~h le village de \textbf{Marquay} (Pas-de-Calais, à 6~km à l'est de Saint-Pol-sur-Ternoise), où elle cantonne jusqu'au 9 septembre.}\par}
\newpage
\phantomsection\label{fac:28}
% --- En-tête du fac-similé (Page impaire / droite : Photographie du manuscrit) ---
\noindent\rule{\textwidth}{0.5pt}\par\vspace{0.55mm}
\noindent\textbf{\fontsize{6.8pt}{7.8pt}\selectfont FAC-SIMILÉ EN VIS-À-VIS — FOLIO 14 RECTO} \hfill \textbf{\fontsize{6.8pt}{7.8pt}\selectfont AOÛT 1916}\par\vspace{0.25mm}
\noindent{\fontsize{6.8pt}{8.0pt}\selectfont\textit{Livre-Journal d'Hippolyte Combes} — Vendredi 25 août 1916 (suite) (De Bailleul à Saint-Pol et Marquay (Pas-de-Calais))}\par\vspace{0.55mm}
\noindent\rule{\textwidth}{0.35pt}\par\vspace*{\fill}
{\centering
  \setlength{\fboxsep}{1.2mm}%
  \setlength{\fboxrule}{0.4pt}%
  \fcolorbox{gray!55}{white}{%
    \includegraphics[width=0.95\textwidth,height=0.81\textheight,keepaspectratio]{\facsimiledir/20251126_094058.jpg}%
  }\par}
\vspace*{\fill}
\newpage

% =============================================================================
% VIS-A-VIS SPREAD : FOLIO 14 VERSO (20251126_094159.jpg)
% =============================================================================
\cleardoubleevenpage
\renewcommand{\currentmonthyear}{Août 1916}
\renewcommand{\currentfoliolabel}{Folio 14 verso}
\phantomsection\label{fol:29}
% --- En-tête du folio (Page paire / gauche : Transcription, apparat & registres de bas de page) ---
\noindent\rule{\textwidth}{0.5pt}\par\vspace{0.55mm}
\noindent\textbf{\footnotesize FOLIO 14 VERSO} \hfill \textbf{\footnotesize AOÛT 1916}\par\vspace{0.25mm}
\noindent{\fontsize{6.9pt}{8.0pt}\selectfont Promenade sur le plateau de Bailleul-aux-Cornailles et portrait au vitriol de la logeuse}\par\vspace{0.2mm}
\noindent{\fontsize{6.9pt}{8.0pt}\selectfont\textit{Samedi 26 août 1916} \hfill \textit{Marquay et Bailleul-aux-Cornailles (Pas-de-Calais)}}\par\vspace{0.5mm}
\noindent\rule{\textwidth}{0.35pt}\par\vspace{0.8mm}
% --- Corps de la transcription diplomatique et apparat marginal en prolongement ---
{\renewcommand{\arraystretch}{0.93}%
\noindent\hspace*{3.5mm}\begin{tabular}{@{}p{87mm}@{\hspace{3.5mm}}p{39mm}@{}}
      \leavevmode\fontsize{7.3pt}{8.4pt}\selectfont\strut \underline{Sam. 26} – \hfill{} \underline{Août / 16}\strut & \smash{\parbox[t]{39mm}{\fontsize{6.2pt}{7.0pt}\selectfont\raggedright\color{black!88}\textbf{1 Sam.} \textit{scil.} Samedi.}} \\
      \leavevmode\fontsize{7.3pt}{8.4pt}\selectfont\strut Bien reposé, quoique dans chambrette étroite\strut &  \\
      \leavevmode\fontsize{7.3pt}{8.4pt}\selectfont\strut aux murs moisis... \hfill{} Levé dès 7H 1/2\strut & \smash{\parbox[t]{39mm}{\fontsize{6.2pt}{7.0pt}\selectfont\raggedright\color{black!88}\textbf{3 1/2} \textit{scil.} demi / moitié.}} \\
      \leavevmode\fontsize{7.3pt}{8.4pt}\selectfont\strut Matinée, de nouveau, très fortes ondées.\strut &  \\
      \leavevmode\llap{\fontsize{6.4pt}{7.3pt}\selectfont\color{darkgray}5\hspace{1.8mm}}\fontsize{7.3pt}{8.4pt}\selectfont\strut Après-midi, enfin, pluie a\strut &  \\
      \leavevmode\fontsize{7.3pt}{8.4pt}\selectfont\strut cessé –\strut &  \\
      \leavevmode\fontsize{7.3pt}{8.4pt}\selectfont\strut Venu me promener au village\strut &  \\
      \leavevmode\fontsize{7.3pt}{8.4pt}\selectfont\strut de Bailleul-aux-Cornailles. à 2 ou 3\strut &  \\
      \leavevmode\fontsize{7.3pt}{8.4pt}\selectfont\strut kilomètres de Marquay\strut &  \\
      \leavevmode\llap{\fontsize{6.4pt}{7.3pt}\selectfont\color{darkgray}10\hspace{1.8mm}}\fontsize{7.3pt}{8.4pt}\selectfont\strut Route à découvert, entre champs de\strut &  \\
      \leavevmode\fontsize{7.3pt}{8.4pt}\selectfont\strut blé déjà moissonnés et champs d'avoine\strut &  \\
      \leavevmode\fontsize{7.3pt}{8.4pt}\selectfont\strut où bruit le vent –\strut &  \\
      \leavevmode\fontsize{7.3pt}{8.4pt}\selectfont\strut Plateau d'où l'on domine horizons\strut &  \\
      \leavevmode\fontsize{7.3pt}{8.4pt}\selectfont\strut de bois; d'où émergent clochers\strut &  \\
      \leavevmode\llap{\fontsize{6.4pt}{7.3pt}\selectfont\color{darkgray}15\hspace{1.8mm}}\fontsize{7.3pt}{8.4pt}\selectfont\strut d'églises\strut &  \\
      \leavevmode\fontsize{7.3pt}{8.4pt}\selectfont\strut Avant diner, que pris très tard,\strut &  \\
      \leavevmode\fontsize{7.3pt}{8.4pt}\selectfont\strut promené sur la route de.. – à\strut &  \\
      \leavevmode\fontsize{7.3pt}{8.4pt}\selectfont\strut l'horizon de l. très beau sunset\strut & \smash{\parbox[t]{39mm}{\fontsize{6.2pt}{7.0pt}\selectfont\raggedright\color{black!88}\textbf{18 sunset} \textit{angl. in ms.}}} \\
      \leavevmode\fontsize{7.3pt}{8.4pt}\selectfont\strut Mauvaise installat° où je suis\strut &  \\
      \leavevmode\llap{\fontsize{6.4pt}{7.3pt}\selectfont\color{darkgray}20\hspace{1.8mm}}\fontsize{7.3pt}{8.4pt}\selectfont\strut gosses, coqueluche — maison très mal\strut &  \\
      \leavevmode\fontsize{7.3pt}{8.4pt}\selectfont\strut tenue, sale, moisie — quelle [torture ?]!\strut & \smash{\parbox[t]{39mm}{\fontsize{6.2pt}{7.0pt}\selectfont\raggedright\color{black!88}\textbf{21} \textit{lect. incert. in ms.}}} \\
\end{tabular}}%
\par\vspace*{\fill}
% --- Registres de bas de page (1. Critique interne [si pertinent] ; 2. Loci paralleli en latin ; 3. Commentaire historique en retrait) ---
\noindent\rule{\textwidth}{0.5pt}\par\vspace{0.65mm}
{\fontsize{6.6pt}{7.7pt}\selectfont%
\noindent\normalfont\color{black!92} \textbf{8} (\textit{Bailleul-aux-Cornailles}) cf. infra f\textsuperscript{o}~15~v\textsuperscript{o} (p.~\pageref{fol:31}), f\textsuperscript{o}~18~v\textsuperscript{o} (p.~\pageref{fol:37}). — \textbf{9} (\textit{Marquay}) cf. supra f\textsuperscript{o}~14~r\textsuperscript{o} (p.~\pageref{fol:28}), infra f\textsuperscript{o}~15~r\textsuperscript{o} (p.~\pageref{fol:30}), f\textsuperscript{o}~15~v\textsuperscript{o} (p.~\pageref{fol:31}), \textit{et passim}.\par\vspace{0.55mm}
\noindent\rule{38mm}{0.4pt}\par\vspace{0.5mm}
\noindent\hspace*{5mm}\parbox{\dimexpr\textwidth-5mm\relax}{\fontsize{6.6pt}{7.7pt}\selectfont\color{black!84}Première promenade sur le plateau crayeux du Ternois jusqu'au village voisin de \textbf{Bailleul-aux-Cornailles} (Pas-de-Calais, à 2,5~km au nord-est de Marquay) et constat de l'insalubrité de sa chambrette de cantonnement à Marquay.}\par}
\newpage
\phantomsection\label{fac:29}
% --- En-tête du fac-similé (Page impaire / droite : Photographie du manuscrit) ---
\noindent\rule{\textwidth}{0.5pt}\par\vspace{0.55mm}
\noindent\textbf{\fontsize{6.8pt}{7.8pt}\selectfont FAC-SIMILÉ EN VIS-À-VIS — FOLIO 14 VERSO} \hfill \textbf{\fontsize{6.8pt}{7.8pt}\selectfont AOÛT 1916}\par\vspace{0.25mm}
\noindent{\fontsize{6.8pt}{8.0pt}\selectfont\textit{Livre-Journal d'Hippolyte Combes} — Samedi 26 août 1916 (Marquay et Bailleul-aux-Cornailles (Pas-de-Calais))}\par\vspace{0.55mm}
\noindent\rule{\textwidth}{0.35pt}\par\vspace*{\fill}
{\centering
  \setlength{\fboxsep}{1.2mm}%
  \setlength{\fboxrule}{0.4pt}%
  \fcolorbox{gray!55}{white}{%
    \includegraphics[width=0.95\textwidth,height=0.81\textheight,keepaspectratio]{\facsimiledir/20251126_094159.jpg}%
  }\par}
\vspace*{\fill}
\newpage

% =============================================================================
% VIS-A-VIS SPREAD : FOLIO 15 RECTO (20251126_094215.jpg)
% =============================================================================
\cleardoubleevenpage
\renewcommand{\currentmonthyear}{Août 1916}
\renewcommand{\currentfoliolabel}{Folio 15 recto}
\phantomsection\label{fol:30}
% --- En-tête du folio (Page paire / gauche : Transcription, apparat & registres de bas de page) ---
\noindent\rule{\textwidth}{0.5pt}\par\vspace{0.55mm}
\noindent\textbf{\footnotesize FOLIO 15 RECTO} \hfill \textbf{\footnotesize AOÛT 1916}\par\vspace{0.25mm}
\noindent{\fontsize{6.9pt}{8.0pt}\selectfont Lundi 28 et Mardi 29 août 1916 — Bourbiers d'Artois, corvées de fourrage avec le lieutenant Curtis et entrée en guerre de la Roumanie}\par\vspace{0.2mm}
\noindent{\fontsize{6.9pt}{8.0pt}\selectfont\textit{Lundi 28 et mardi 29 août 1916} \hfill \textit{Marquay (Pas-de-Calais)}}\par\vspace{0.5mm}
\noindent\rule{\textwidth}{0.35pt}\par\vspace{0.8mm}
% --- Corps de la transcription diplomatique et apparat marginal en prolongement ---
{\renewcommand{\arraystretch}{0.87}%
\noindent\hspace*{3.5mm}\begin{tabular}{@{}p{87mm}@{\hspace{3.5mm}}p{39mm}@{}}
      \leavevmode\fontsize{6.9pt}{7.9pt}\selectfont\strut \underline{28 Lundi.} \hfill{} \underline{Août 1916}\strut &  \\
      \leavevmode\fontsize{6.9pt}{7.9pt}\selectfont\strut Aujourd’hui encore, giboulées \& éclaircies –\strut &  \\
      \leavevmode\fontsize{6.9pt}{7.9pt}\selectfont\strut comme au mois de mars\strut &  \\
      \leavevmode\fontsize{6.9pt}{7.9pt}\selectfont\strut Surtout matinée, très fortes ondées. Routes\strut &  \\
      \leavevmode\llap{\fontsize{6.4pt}{7.3pt}\selectfont\color{darkgray}5\hspace{1.8mm}}\fontsize{6.9pt}{7.9pt}\selectfont\strut détrempées, cloaques de boue. Impossible\strut &  \\
      \leavevmode\fontsize{6.9pt}{7.9pt}\selectfont\strut de faire longues promenades aux environs\strut &  \\
      \leavevmode\fontsize{6.9pt}{7.9pt}\selectfont\strut On ne sort que par nécessité\strut &  \\
      \leavevmode\fontsize{6.9pt}{7.9pt}\selectfont\strut Matin \& soir, dû aller champs pour green\strut & \smash{\parbox[t]{39mm}{\fontsize{6.2pt}{7.0pt}\selectfont\raggedright\color{black!88}\textbf{8 green} \textit{angl. in ms.}}} \\
      \leavevmode\fontsize{6.9pt}{7.9pt}\selectfont\strut fourrage\strut &  \\
      \leavevmode\llap{\fontsize{6.4pt}{7.3pt}\selectfont\color{darkgray}10\hspace{1.8mm}}\fontsize{6.9pt}{7.9pt}\selectfont\strut Fin après-midi, très beau soleil.\strut &  \\
      \leavevmode\fontsize{6.9pt}{7.9pt}\selectfont\strut Aucun confort. Jamais + mal logé\strut &  \\
      \leavevmode\fontsize{6.9pt}{7.9pt}\selectfont\strut Vie végétative déprimante \& abrutissante\strut &  \\
      \leavevmode\fontsize{6.9pt}{7.9pt}\selectfont\strut Dîné tôt. puis promené ds la campagne\strut & \smash{\parbox[t]{39mm}{\fontsize{6.2pt}{7.0pt}\selectfont\raggedright\color{black!88}\textbf{13 ds} \textit{abbr. pro} dans.}} \\
      \leavevmode\fontsize{6.9pt}{7.9pt}\selectfont\strut Regardé paysans travailler champs – machine\strut &  \\
      \leavevmode\llap{\fontsize{6.4pt}{7.3pt}\selectfont\color{darkgray}15\hspace{1.8mm}}\fontsize{6.9pt}{7.9pt}\selectfont\strut moissonneuse-lieuse – j-aux dernières lueurs jour\strut & \smash{\parbox[t]{39mm}{\fontsize{6.2pt}{7.0pt}\selectfont\raggedright\color{black!88}\textbf{15 j-aux} \textit{abbr. pro} jusqu'aux.}} \\
      \cline{1-1}
      \leavevmode\fontsize{6.9pt}{7.9pt}\selectfont\strut \underline{29 Ma} Matin \& soir, dû accompagner\strut &  \\
      \leavevmode\fontsize{6.9pt}{7.9pt}\selectfont\strut cet idiot de Curtis — A travers champs\strut &  \\
      \leavevmode\fontsize{6.9pt}{7.9pt}\selectfont\strut – malgré pluie \& boue – fourrage –\strut &  \\
      \leavevmode\fontsize{6.9pt}{7.9pt}\selectfont\strut Après-midi, très violent orage\strut &  \\
      \leavevmode\llap{\fontsize{6.4pt}{7.3pt}\selectfont\color{darkgray}20\hspace{1.8mm}}\fontsize{6.9pt}{7.9pt}\selectfont\strut Cour de la ferme, transformée en mare\strut &  \\
      \leavevmode\fontsize{6.9pt}{7.9pt}\selectfont\strut Longues heures de mélancolie\strut &  \\
      \leavevmode\fontsize{6.9pt}{7.9pt}\selectfont\strut et d'ennui –\strut &  \\
      \leavevmode\fontsize{6.9pt}{7.9pt}\selectfont\strut Confirmations nouvelles (Roumanie déclare guerre\strut &  \\
      \leavevmode\fontsize{6.9pt}{7.9pt}\selectfont\strut [illis.] grecque [?] etc. [?]\strut & \smash{\parbox[t]{39mm}{\fontsize{6.2pt}{7.0pt}\selectfont\raggedright\color{black!88}\textbf{24 grecque} \textit{lect. incert. in ms.} \textbf{24 etc} \textit{lect. incert. in ms.}}} \\
\end{tabular}}%
\par\vspace*{\fill}
% --- Registres de bas de page (1. Critique interne [si pertinent] ; 2. Loci paralleli en latin ; 3. Commentaire historique en retrait) ---
\noindent\rule{\textwidth}{0.5pt}\par\vspace{0.65mm}
{\fontsize{6.6pt}{7.7pt}\selectfont%
\noindent\normalfont\color{black!92} \textbf{9} (\textit{fourrage}) cf. supra f\textsuperscript{o}~3~r\textsuperscript{o} (p.~\pageref{fol:6}), f\textsuperscript{o}~11~v\textsuperscript{o} (p.~\pageref{fol:23}), infra f\textsuperscript{o}~15~v\textsuperscript{o} (p.~\pageref{fol:31}), f\textsuperscript{o}~19~r\textsuperscript{o} (p.~\pageref{fol:38}), \textit{et passim}. — \textbf{17} (\textit{Curtis}) cf. supra f\textsuperscript{o}~1~r\textsuperscript{o} (p.~\pageref{fol:2}), f\textsuperscript{o}~11~v\textsuperscript{o} (p.~\pageref{fol:23}), infra f\textsuperscript{o}~18~v\textsuperscript{o} (p.~\pageref{fol:37}), f\textsuperscript{o}~44~v\textsuperscript{o} (p.~\pageref{fol:89}), \textit{et passim}.\par\vspace{0.55mm}
\noindent\rule{38mm}{0.4pt}\par\vspace{0.5mm}
\noindent\hspace*{5mm}\parbox{\dimexpr\textwidth-5mm\relax}{\fontsize{6.6pt}{7.7pt}\selectfont\color{black!84}Entre les corvées de fourrage vert dans la boue avec le lieutenant Curtis et l'observation des moissons tardives d'Artois à la moissonneuse-lieuse, Hippolyte Combes consigne le mardi 29 août 1916 la confirmation de l'entrée en guerre du \textbf{royaume de Roumanie} aux côtés de l'Entente (déclaration de guerre à l'Autriche-Hongrie du dimanche \textbf{27 août 1916} au soir, suivie le 28 août de la déclaration de guerre de l'Allemagne à la Roumanie).}\par}
\newpage
\phantomsection\label{fac:30}
% --- En-tête du fac-similé (Page impaire / droite : Photographie du manuscrit) ---
\noindent\rule{\textwidth}{0.5pt}\par\vspace{0.55mm}
\noindent\textbf{\fontsize{6.8pt}{7.8pt}\selectfont FAC-SIMILÉ EN VIS-À-VIS — FOLIO 15 RECTO} \hfill \textbf{\fontsize{6.8pt}{7.8pt}\selectfont AOÛT 1916}\par\vspace{0.25mm}
\noindent{\fontsize{6.8pt}{8.0pt}\selectfont\textit{Livre-Journal d'Hippolyte Combes} — Lundi 28 et mardi 29 août 1916 (Marquay (Pas-de-Calais))}\par\vspace{0.55mm}
\noindent\rule{\textwidth}{0.35pt}\par\vspace*{\fill}
{\centering
  \setlength{\fboxsep}{1.2mm}%
  \setlength{\fboxrule}{0.4pt}%
  \fcolorbox{gray!55}{white}{%
    \includegraphics[width=0.95\textwidth,height=0.81\textheight,keepaspectratio]{\facsimiledir/20251126_094215.jpg}%
  }\par}
\vspace*{\fill}
\newpage

% =============================================================================
% VIS-A-VIS SPREAD : FOLIO 15 VERSO (20251126_094238.jpg)
% =============================================================================
\cleardoubleevenpage
\renewcommand{\currentmonthyear}{Août 1916}
\renewcommand{\currentfoliolabel}{Folio 15 verso}
\phantomsection\label{fol:31}
% --- En-tête du folio (Page paire / gauche : Transcription, apparat & registres de bas de page) ---
\noindent\rule{\textwidth}{0.5pt}\par\vspace{0.55mm}
\noindent\textbf{\footnotesize FOLIO 15 VERSO} \hfill \textbf{\footnotesize AOÛT 1916}\par\vspace{0.25mm}
\noindent{\fontsize{6.9pt}{8.0pt}\selectfont Tempête de vent dans les grands arbres du voisinage et fourrage sous les ondées}\par\vspace{0.2mm}
\noindent{\fontsize{6.9pt}{8.0pt}\selectfont\textit{Mercredi 30 août 1916} \hfill \textit{Marquay et Bailleul-aux-Cornailles (Pas-de-Calais)}}\par\vspace{0.5mm}
\noindent\rule{\textwidth}{0.35pt}\par\vspace{0.8mm}
% --- Corps de la transcription diplomatique et apparat marginal en prolongement ---
{\renewcommand{\arraystretch}{1.00}%
\noindent\hspace*{3.5mm}\begin{tabular}{@{}p{87mm}@{\hspace{3.5mm}}p{39mm}@{}}
      \leavevmode\fontsize{7.7pt}{9.0pt}\selectfont\strut \underline{30 Me} \hfill{} \underline{Août 1916}\strut &  \\
      \leavevmode\fontsize{7.7pt}{9.0pt}\selectfont\strut Tout le jour, pluie. et rafales\strut &  \\
      \leavevmode\fontsize{7.7pt}{9.0pt}\selectfont\strut Vent tourmente grands arbres du\strut &  \\
      \leavevmode\fontsize{7.7pt}{9.0pt}\selectfont\strut voisinage (boqueteau de l'autre côté de\strut &  \\
      \leavevmode\llap{\fontsize{6.4pt}{7.3pt}\selectfont\color{darkgray}5\hspace{1.8mm}}\fontsize{7.7pt}{9.0pt}\selectfont\strut la route et prairies derrière fermes et\strut &  \\
      \leavevmode\fontsize{7.7pt}{9.0pt}\selectfont\strut fait bruit effrayant\strut &  \\
      \leavevmode\fontsize{7.7pt}{9.0pt}\selectfont\strut Matin \& après-midi, en dépit ondées\strut &  \\
      \leavevmode\fontsize{7.7pt}{9.0pt}\selectfont\strut très fortes, dû aller, champs détrempés\strut &  \\
      \leavevmode\fontsize{7.7pt}{9.0pt}\selectfont\strut pour faire prendre fourrage vert\strut &  \\
\end{tabular}}%
\par\vspace*{\fill}
% --- Registres de bas de page (1. Critique interne [si pertinent] ; 2. Loci paralleli en latin ; 3. Commentaire historique en retrait) ---
\noindent\rule{\textwidth}{0.5pt}\par\vspace{0.65mm}
{\fontsize{6.6pt}{7.7pt}\selectfont%
\noindent\normalfont\color{black!92} \textbf{9} (\textit{fourrage}) cf. supra f\textsuperscript{o}~3~r\textsuperscript{o} (p.~\pageref{fol:6}), f\textsuperscript{o}~15~r\textsuperscript{o} (p.~\pageref{fol:30}), infra f\textsuperscript{o}~19~r\textsuperscript{o} (p.~\pageref{fol:38}), f\textsuperscript{o}~21~r\textsuperscript{o} (p.~\pageref{fol:42}), \textit{et passim}.\par\vspace{0.55mm}
\noindent\rule{38mm}{0.4pt}\par\vspace{0.5mm}
\noindent\hspace*{5mm}\parbox{\dimexpr\textwidth-5mm\relax}{\fontsize{6.6pt}{7.7pt}\selectfont\color{black!84}La tempête de la fin août 1916 sur les plateaux d'Artois (Marquay) n'interrompt pas le service quotidien d'achat et de coupe du fourrage vert (\textit{green forage}) pour les chevaux du train divisionnaire.}\par}
\newpage
\phantomsection\label{fac:31}
% --- En-tête du fac-similé (Page impaire / droite : Photographie du manuscrit) ---
\noindent\rule{\textwidth}{0.5pt}\par\vspace{0.55mm}
\noindent\textbf{\fontsize{6.8pt}{7.8pt}\selectfont FAC-SIMILÉ EN VIS-À-VIS — FOLIO 15 VERSO} \hfill \textbf{\fontsize{6.8pt}{7.8pt}\selectfont AOÛT 1916}\par\vspace{0.25mm}
\noindent{\fontsize{6.8pt}{8.0pt}\selectfont\textit{Livre-Journal d'Hippolyte Combes} — Mercredi 30 août 1916 (Marquay et Bailleul-aux-Cornailles (Pas-de-Calais))}\par\vspace{0.55mm}
\noindent\rule{\textwidth}{0.35pt}\par\vspace*{\fill}
{\centering
  \setlength{\fboxsep}{1.2mm}%
  \setlength{\fboxrule}{0.4pt}%
  \fcolorbox{gray!55}{white}{%
    \includegraphics[width=0.95\textwidth,height=0.81\textheight,keepaspectratio]{\facsimiledir/20251126_094238.jpg}%
  }\par}
\vspace*{\fill}
\newpage

% =============================================================================
% VIS-A-VIS SPREAD : FOLIO 16 RECTO (20251126_094250.jpg)
% =============================================================================
\cleardoubleevenpage
\renewcommand{\currentmonthyear}{Août 1916}
\renewcommand{\currentfoliolabel}{Folio 16 recto}
\phantomsection\label{fol:32}
% --- En-tête du folio (Page paire / gauche : Transcription, apparat & registres de bas de page) ---
\noindent\rule{\textwidth}{0.5pt}\par\vspace{0.55mm}
\noindent\textbf{\footnotesize FOLIO 16 RECTO} \hfill \textbf{\footnotesize AOÛT 1916}\par\vspace{0.25mm}
\noindent{\fontsize{6.9pt}{8.0pt}\selectfont Marche à pied par Ramecourt jusqu'à Saint-Pol-sur-Ternoise et souvenir de saint Benoît Labre}\par\vspace{0.2mm}
\noindent{\fontsize{6.9pt}{8.0pt}\selectfont\textit{Jeudi 31 août 1916} \hfill \textit{Marquay, Ramecourt, Saint-Michel, Saint-Pol-sur-Ternoise et Ostreville (Pas-de-Calais)}}\par\vspace{0.5mm}
\noindent\rule{\textwidth}{0.35pt}\par\vspace{0.8mm}
% --- Corps de la transcription diplomatique et apparat marginal en prolongement ---
{\renewcommand{\arraystretch}{0.87}%
\noindent\hspace*{3.5mm}\begin{tabular}{@{}p{87mm}@{\hspace{3.5mm}}p{39mm}@{}}
      \leavevmode\fontsize{6.9pt}{7.9pt}\selectfont\strut \underline{31 J.} – Enfin, belle journée. – \hfill{} \raisebox{0.5ex}{\underline{Aout / 16}}\strut & \smash{\parbox[t]{39mm}{\fontsize{6.2pt}{7.0pt}\selectfont\raggedright\color{black!88}\textbf{1 J.} \textit{scil.} Jeudi.}} \\
      \leavevmode\fontsize{6.9pt}{7.9pt}\selectfont\strut Soleil ardent transformt paysages \& humeurs\strut &  \\
      \leavevmode\fontsize{6.9pt}{7.9pt}\selectfont\strut Après-midi, promenade à pied j.à St-Pol.\strut & \smash{\parbox[t]{39mm}{\fontsize{6.2pt}{7.0pt}\selectfont\raggedright\color{black!88}\textbf{3 j. à} \textit{abbr. pro} jusqu'à. \textbf{3 St} \textit{abbr. pro} Saint.}} \\
      \leavevmode\fontsize{6.9pt}{7.9pt}\selectfont\strut Au départ, coupé à travers champs, en\strut &  \\
      \leavevmode\llap{\fontsize{6.4pt}{7.3pt}\selectfont\color{darkgray}5\hspace{1.8mm}}\fontsize{6.9pt}{7.9pt}\selectfont\strut haut plateau de Ramecourt \& arrivé au\strut &  \\
      \leavevmode\fontsize{6.9pt}{7.9pt}\selectfont\strut village de St. Michel-s-Ternoise. puis\strut &  \\
      \leavevmode\fontsize{6.9pt}{7.9pt}\selectfont\strut suivi gde route Arras. St. Pol.\strut & \smash{\parbox[t]{39mm}{\fontsize{6.2pt}{7.0pt}\selectfont\raggedright\color{black!88}\textbf{7 gde} \textit{abbr. pro} grande.}} \\
      \leavevmode\fontsize{6.9pt}{7.9pt}\selectfont\strut Au fond vallon, s'étend cette petite\strut &  \\
      \leavevmode\fontsize{6.9pt}{7.9pt}\selectfont\strut sous-préf[re ?] de 5000 habts\strut & \smash{\parbox[t]{39mm}{\fontsize{6.2pt}{7.0pt}\selectfont\raggedright\color{black!88}\textbf{9} \textit{lect. incert. in ms.} \textbf{9 habts} \textit{abbr. pro} habitants. — \textbf{10 gde} \textit{abbr. pro} grande.}} \\
      \leavevmode\llap{\fontsize{6.4pt}{7.3pt}\selectfont\color{darkgray}10\hspace{1.8mm}}\fontsize{6.9pt}{7.9pt}\selectfont\strut Rues en pente aboutissant à gde place\strut &  \\
      \leavevmode\fontsize{6.9pt}{7.9pt}\selectfont\strut Eglise briques rouges, style hollandais,\strut &  \\
      \leavevmode\fontsize{6.9pt}{7.9pt}\selectfont\strut (1697). – Remarqué intérieur plaque\strut &  \\
      \leavevmode\fontsize{6.9pt}{7.9pt}\selectfont\strut marbre s/ 1 pilier : rappelant que\strut & \smash{\parbox[t]{39mm}{\fontsize{6.2pt}{7.0pt}\selectfont\raggedright\color{black!88}\textbf{13 s/} \textit{abbr. pro} sur.}} \\
      \leavevmode\fontsize{6.9pt}{7.9pt}\selectfont\strut mars 1767 le bienh Benoît Labre – passa\strut &  \\
      \leavevmode\llap{\fontsize{6.4pt}{7.3pt}\selectfont\color{darkgray}15\hspace{1.8mm}}\fontsize{6.9pt}{7.9pt}\selectfont\strut la nuit en prière ds cette = église\strut & \smash{\parbox[t]{39mm}{\fontsize{6.2pt}{7.0pt}\selectfont\raggedright\color{black!88}\textbf{15 ds} \textit{abbr. pro} dans.}} \\
      \leavevmode\fontsize{6.9pt}{7.9pt}\selectfont\strut Rentré par Ostreville\strut &  \\
      \leavevmode\fontsize{6.9pt}{7.9pt}\selectfont\strut Course m'a fait bp de bien\strut & \smash{\parbox[t]{39mm}{\fontsize{6.2pt}{7.0pt}\selectfont\raggedright\color{black!88}\textbf{17 bp} \textit{abbr. pro} beaucoup.}} \\
      \cline{1-1}
      \leavevmode\fontsize{6.9pt}{7.9pt}\selectfont\strut \underline{A Saint-Pol} : 12 ctes ples : 50 c — Varia : [10/40 ?] c\strut & \smash{\parbox[t]{39mm}{\fontsize{6.2pt}{7.0pt}\selectfont\raggedright\color{black!88}\textbf{18} \textit{lect. incert. in ms.} \textbf{18 ctes ples} \textit{abbr. pro} cartes postales. — \textbf{19} \textit{lect. incert. in ms.} \textbf{19 cartes ples} \textit{abbr. pro} cartes postales. — \textbf{20} \textit{lect. incert. in ms.} \textbf{20 Cartes ples} \textit{abbr. pro} Cartes postales.}} \\
      \leavevmode\fontsize{6.9pt}{7.9pt}\selectfont\strut Carte du départt : 70 c. — Div. cartes ples : 0 f [50 ?]\strut &  \\
      \cline{1-1}
      \leavevmode\llap{\fontsize{6.4pt}{7.3pt}\selectfont\color{darkgray}20\hspace{1.8mm}}\fontsize{6.9pt}{7.9pt}\selectfont\strut Cartes ples militaires : 40 c. — \hfill{} [compte ajouté : 0.80]\strut &  \\
\end{tabular}}%
\par\vspace*{\fill}
% --- Registres de bas de page (1. Critique interne [si pertinent] ; 2. Loci paralleli en latin ; 3. Commentaire historique en retrait) ---
\noindent\rule{\textwidth}{0.5pt}\par\vspace{0.65mm}
{\fontsize{6.6pt}{7.7pt}\selectfont%
\noindent\normalfont\color{black!92} \textbf{3} (\textit{St-Pol}) cf. supra f\textsuperscript{o}~13~v\textsuperscript{o} (p.~\pageref{fol:27}), f\textsuperscript{o}~14~r\textsuperscript{o} (p.~\pageref{fol:28}), infra f\textsuperscript{o}~17~r\textsuperscript{o} (p.~\pageref{fol:34}), f\textsuperscript{o}~18~r\textsuperscript{o} (p.~\pageref{fol:36}), \textit{et passim}, Appendicem I, doc.~5 (p.~\pageref{annexe:doc5}). — \textbf{7} (\textit{Arras}) cf. infra f\textsuperscript{o}~59~r\textsuperscript{o} (p.~\pageref{fol:118}), f\textsuperscript{o}~89~r\textsuperscript{o} (p.~\pageref{fol:178}).\par\vspace{0.55mm}
\noindent\rule{38mm}{0.4pt}\par\vspace{0.5mm}
\noindent\hspace*{5mm}\parbox{\dimexpr\textwidth-5mm\relax}{\fontsize{6.6pt}{7.7pt}\selectfont\color{black!84}Descendu de Marquay à \textbf{Saint-Pol-sur-Ternoise} (Pas-de-Calais) par Ramecourt et Saint-Michel-sur-Ternoise, Hippolyte Combes visite l'ancienne chapelle conventuelle des \textbf{Carmes} (XVII\textsuperscript{e}~s., bâtie en briques rouges en 1697 et devenue église paroissiale Saint-Paul après la destruction de l'ancienne église au XVI\textsuperscript{e}~s.), où une plaque commémore la nuit de prière passée en mars 1767 par saint \textbf{Benoît-Joseph Labre} (1748–1783).}\par}
\newpage
\phantomsection\label{fac:32}
% --- En-tête du fac-similé (Page impaire / droite : Photographie du manuscrit) ---
\noindent\rule{\textwidth}{0.5pt}\par\vspace{0.55mm}
\noindent\textbf{\fontsize{6.8pt}{7.8pt}\selectfont FAC-SIMILÉ EN VIS-À-VIS — FOLIO 16 RECTO} \hfill \textbf{\fontsize{6.8pt}{7.8pt}\selectfont AOÛT 1916}\par\vspace{0.25mm}
\noindent{\fontsize{6.8pt}{8.0pt}\selectfont\textit{Livre-Journal d'Hippolyte Combes} — Jeudi 31 août 1916 (Marquay, Ramecourt, Saint-Michel, Saint-Pol-sur-Ternoise et Ostreville (Pas-de-Calais))}\par\vspace{0.55mm}
\noindent\rule{\textwidth}{0.35pt}\par\vspace*{\fill}
{\centering
  \setlength{\fboxsep}{1.2mm}%
  \setlength{\fboxrule}{0.4pt}%
  \fcolorbox{gray!55}{white}{%
    \includegraphics[width=0.95\textwidth,height=0.81\textheight,keepaspectratio]{\facsimiledir/20251126_094250.jpg}%
  }\par}
\vspace*{\fill}
\newpage

% =============================================================================
% VIS-A-VIS SPREAD : FOLIO 16 VERSO (20251126_094310.jpg)
% =============================================================================
\cleardoubleevenpage
\renewcommand{\currentmonthyear}{Septembre 1916}
\renewcommand{\currentfoliolabel}{Folio 16 verso}
\phantomsection\label{fol:33}
% --- En-tête du folio (Page paire / gauche : Transcription, apparat & registres de bas de page) ---
\noindent\rule{\textwidth}{0.5pt}\par\vspace{0.55mm}
\noindent\textbf{\footnotesize FOLIO 16 VERSO} \hfill \textbf{\footnotesize SEPTEMBRE 1916}\par\vspace{0.25mm}
\noindent{\fontsize{6.9pt}{8.0pt}\selectfont L'église médiévale de Ligny-Saint-Flochel et les « manœuvres spéciales » britanniques}\par\vspace{0.2mm}
\noindent{\fontsize{6.9pt}{8.0pt}\selectfont\textit{Vendredi 1er septembre 1916} \hfill \textit{Marquay et Ligny-Saint-Flochel (Pas-de-Calais)}}\par\vspace{0.5mm}
\noindent\rule{\textwidth}{0.35pt}\par\vspace{0.8mm}
% --- Corps de la transcription diplomatique et apparat marginal en prolongement ---
{\renewcommand{\arraystretch}{0.93}%
\noindent\hspace*{3.5mm}\begin{tabular}{@{}p{87mm}@{\hspace{3.5mm}}p{39mm}@{}}
      \leavevmode\fontsize{7.3pt}{8.4pt}\selectfont\strut \underline{1er V.} \hfill{} \underline{Septbre 1916}\strut & \smash{\parbox[t]{39mm}{\fontsize{6.2pt}{7.0pt}\selectfont\raggedright\color{black!88}\textbf{1 1er} \textit{scil.} premier.}} \\
      \leavevmode\fontsize{7.3pt}{8.4pt}\selectfont\strut Matinée, temps variable\strut &  \\
      \leavevmode\fontsize{7.3pt}{8.4pt}\selectfont\strut Tout l'après-midi, qq rayons de soleil et puis\strut & \smash{\parbox[t]{39mm}{\fontsize{6.2pt}{7.0pt}\selectfont\raggedright\color{black!88}\textbf{3 qq} \textit{abbr. pro} quelque(s).}} \\
      \leavevmode\fontsize{7.3pt}{8.4pt}\selectfont\strut ciel couvert de nuages\strut &  \\
      \leavevmode\llap{\fontsize{6.4pt}{7.3pt}\selectfont\color{darkgray}5\hspace{1.8mm}}\fontsize{7.3pt}{8.4pt}\selectfont\strut Venu promener au village de Ligny-\strut &  \\
      \leavevmode\fontsize{7.3pt}{8.4pt}\selectfont\strut Saint-Flochel, à 2 kil 5 — Route en pente longe\strut &  \\
      \leavevmode\fontsize{7.3pt}{8.4pt}\selectfont\strut bois où régiments anglais faisaient les manœuvres\strut &  \\
      \leavevmode\fontsize{7.3pt}{8.4pt}\selectfont\strut spéciales\strut &  \\
      \leavevmode\fontsize{7.3pt}{8.4pt}\selectfont\strut Visité un instant petite eglise du hameau St Fl.\strut & \smash{\parbox[t]{39mm}{\fontsize{6.2pt}{7.0pt}\selectfont\raggedright\color{black!88}\textbf{9 St} \textit{abbr. pro} Saint.}} \\
      \leavevmode\llap{\fontsize{6.4pt}{7.3pt}\selectfont\color{darkgray}10\hspace{1.8mm}}\fontsize{7.3pt}{8.4pt}\selectfont\strut est du XIIIe s. clocher très curieux, octogonal,\strut &  \\
      \leavevmode\fontsize{7.3pt}{8.4pt}\selectfont\strut avec b. pierres sculptées. Tout autour cimetière\strut &  \\
      \leavevmode\fontsize{7.3pt}{8.4pt}\selectfont\strut village, où gazon \& plantes cachent presque les\strut &  \\
      \leavevmode\fontsize{7.3pt}{8.4pt}\selectfont\strut tombes. — Y éprouvé impression\strut &  \\
      \leavevmode\fontsize{7.3pt}{8.4pt}\selectfont\strut mélancolie pénétrante — Ciel nuageux et\strut &  \\
      \leavevmode\llap{\fontsize{6.4pt}{7.3pt}\selectfont\color{darkgray}15\hspace{1.8mm}}\fontsize{7.3pt}{8.4pt}\selectfont\strut bas. Légère brise agitant ramures. Air doux.\strut &  \\
      \leavevmode\fontsize{7.3pt}{8.4pt}\selectfont\strut [Bon/où ?] s'arrêter \& réfléchir\strut & \smash{\parbox[t]{39mm}{\fontsize{6.2pt}{7.0pt}\selectfont\raggedright\color{black!88}\textbf{16} \textit{lect. incert. in ms.}}} \\
      \cline{1-1}
      \leavevmode\fontsize{7.3pt}{8.4pt}\selectfont\strut [ Tip au headcook : 5 f –\strut & \smash{\parbox[t]{39mm}{\fontsize{6.2pt}{7.0pt}\selectfont\raggedright\color{black!88}\textbf{17} \textit{in marg. cum unco} (comptes). \textbf{17 headcook} \textit{angl. in ms.}}} \\
      \leavevmode\fontsize{7.3pt}{8.4pt}\selectfont\strut Jours se succèdent monotones \&\strut &  \\
      \leavevmode\fontsize{7.3pt}{8.4pt}\selectfont\strut ternes\strut &  \\
      \leavevmode\llap{\fontsize{6.4pt}{7.3pt}\selectfont\color{darkgray}20\hspace{1.8mm}}\fontsize{7.3pt}{8.4pt}\selectfont\strut Espoir départ prochain, lundi\strut &  \\
\end{tabular}}%
\par\vspace*{\fill}
% --- Registres de bas de page (1. Critique interne [si pertinent] ; 2. Loci paralleli en latin ; 3. Commentaire historique en retrait) ---
\noindent\rule{\textwidth}{0.5pt}\par\vspace{0.65mm}
{\fontsize{6.6pt}{7.7pt}\selectfont%
\noindent\normalfont\color{black!92} \textbf{5} (\textit{Ligny}) cf. infra f\textsuperscript{o}~19~v\textsuperscript{o} (p.~\pageref{fol:39}).\par\vspace{0.55mm}
\noindent\rule{38mm}{0.4pt}\par\vspace{0.5mm}
\noindent\hspace*{5mm}\parbox{\dimexpr\textwidth-5mm\relax}{\fontsize{6.6pt}{7.7pt}\selectfont\color{black!84}À \textbf{Ligny-Saint-Flochel} (Pas-de-Calais), les manœuvres observées en lisière des bois s'inscrivent dans l'entraînement tactique de la \textbf{18\textsuperscript{e} Division britannique} avant son retour sur la Somme (Nichols, 1922, p.~86). L'église \textbf{Saint-Flochel} conserve un clocher de pierre des XV\textsuperscript{e}–XVI\textsuperscript{e}~siècles classé monument historique (Base Mérimée, \texttt{PA00108332}).}\par}
\newpage
\phantomsection\label{fac:33}
% --- En-tête du fac-similé (Page impaire / droite : Photographie du manuscrit) ---
\noindent\rule{\textwidth}{0.5pt}\par\vspace{0.55mm}
\noindent\textbf{\fontsize{6.8pt}{7.8pt}\selectfont FAC-SIMILÉ EN VIS-À-VIS — FOLIO 16 VERSO} \hfill \textbf{\fontsize{6.8pt}{7.8pt}\selectfont SEPTEMBRE 1916}\par\vspace{0.25mm}
\noindent{\fontsize{6.8pt}{8.0pt}\selectfont\textit{Livre-Journal d'Hippolyte Combes} — Vendredi 1er septembre 1916 (Marquay et Ligny-Saint-Flochel (Pas-de-Calais))}\par\vspace{0.55mm}
\noindent\rule{\textwidth}{0.35pt}\par\vspace*{\fill}
{\centering
  \setlength{\fboxsep}{1.2mm}%
  \setlength{\fboxrule}{0.4pt}%
  \fcolorbox{gray!55}{white}{%
    \includegraphics[width=0.95\textwidth,height=0.81\textheight,keepaspectratio]{\facsimiledir/20251126_094310.jpg}%
  }\par}
\vspace*{\fill}
\newpage

% =============================================================================
% VIS-A-VIS SPREAD : FOLIO 17 RECTO (20251126_094324.jpg)
% =============================================================================
\cleardoubleevenpage
\renewcommand{\currentmonthyear}{Septembre 1916}
\renewcommand{\currentfoliolabel}{Folio 17 recto}
\phantomsection\label{fol:34}
% --- En-tête du folio (Page paire / gauche : Transcription, apparat & registres de bas de page) ---
\noindent\rule{\textwidth}{0.5pt}\par\vspace{0.55mm}
\noindent\textbf{\footnotesize FOLIO 17 RECTO} \hfill \textbf{\footnotesize SEPTEMBRE 1916}\par\vspace{0.25mm}
\noindent{\fontsize{6.9pt}{8.0pt}\selectfont Rosée matinale dans les terres argileuses, liste de blanchissage en anglais et marche à Saint-Pol}\par\vspace{0.2mm}
\noindent{\fontsize{6.9pt}{8.0pt}\selectfont\textit{Samedi 2 et dimanche 3 septembre 1916} \hfill \textit{Marquay et Saint-Pol-sur-Ternoise (Pas-de-Calais)}}\par\vspace{0.5mm}
\noindent\rule{\textwidth}{0.35pt}\par\vspace{0.8mm}
% --- Corps de la transcription diplomatique et apparat marginal en prolongement ---
{\renewcommand{\arraystretch}{0.87}%
\noindent\hspace*{3.5mm}\begin{tabular}{@{}p{87mm}@{\hspace{3.5mm}}p{39mm}@{}}
      \leavevmode\fontsize{6.9pt}{7.9pt}\selectfont\strut \underline{2 S.} \hfill{} \underline{Septbre 16}\strut &  \\
      \leavevmode\fontsize{6.9pt}{7.9pt}\selectfont\strut Matinée, ds les champs. Herbes\strut & \smash{\parbox[t]{39mm}{\fontsize{6.2pt}{7.0pt}\selectfont\raggedright\color{black!88}\textbf{2 ds} \textit{abbr. pro} dans.}} \\
      \leavevmode\fontsize{6.9pt}{7.9pt}\selectfont\strut couvertes de rosée mouillant souliers et eau\strut &  \\
      \leavevmode\fontsize{6.9pt}{7.9pt}\selectfont\strut arrive j'à la peau. Terre friable ([boueuse ?])\strut & \smash{\parbox[t]{39mm}{\fontsize{6.2pt}{7.0pt}\selectfont\raggedright\color{black!88}\textbf{4} \textit{lect. incert. in ms.} \textbf{4 j'à} \textit{abbr. pro} jusqu'à.}} \\
      \leavevmode\llap{\fontsize{6.4pt}{7.3pt}\selectfont\color{darkgray}5\hspace{1.8mm}}\fontsize{6.9pt}{7.9pt}\selectfont\strut aux semelles en fortes mottes.\strut &  \\
      \leavevmode\fontsize{6.9pt}{7.9pt}\selectfont\strut Ciel tout couvert de gros nuages – A\strut &  \\
      \leavevmode\fontsize{6.9pt}{7.9pt}\selectfont\strut l'horizon, pointent clochers villages voisins\strut &  \\
      \leavevmode\fontsize{6.9pt}{7.9pt}\selectfont\strut au sein futaies \& boqueteaux –\strut &  \\
      \leavevmode\fontsize{6.9pt}{7.9pt}\selectfont\strut Après-midi, entiert libre\strut & \smash{\parbox[t]{39mm}{\fontsize{6.2pt}{7.0pt}\selectfont\raggedright\color{black!88}\textbf{9 entiert} \textit{abbr. pro} entièrement. — \textbf{10} \textit{in marg. cum unco} (comptes). \textbf{10 Washing} \textit{angl. in ms.} — \textbf{12 Dim.} \textit{scil.} Dimanche. — \textbf{13 j-à} \textit{abbr. pro} jusqu'à. \textbf{13 St} \textit{abbr. pro} Saint. — \textbf{14 gd} \textit{abbr. pro} grand.}} \\
      \leavevmode\llap{\fontsize{6.4pt}{7.3pt}\selectfont\color{darkgray}10\hspace{1.8mm}}\fontsize{6.9pt}{7.9pt}\selectfont\strut [ — Washing : 1 f \hfill{} — Varia : 10 c\strut &  \\
      \leavevmode\fontsize{6.9pt}{7.9pt}\selectfont\strut (1 shirt. 1 vest / 1 p. socks. 1 handk) ]\strut &  \\
      \cline{1-1}
      \leavevmode\fontsize{6.9pt}{7.9pt}\selectfont\strut \underline{3 Dim.} Malgré ciel menaçant\strut &  \\
      \leavevmode\fontsize{6.9pt}{7.9pt}\selectfont\strut allé à pied j-à St-Pol.\strut &  \\
      \leavevmode\fontsize{6.9pt}{7.9pt}\selectfont\strut Sensation enivrante de gd air \& d'espace\strut &  \\
      \leavevmode\llap{\fontsize{6.4pt}{7.3pt}\selectfont\color{darkgray}15\hspace{1.8mm}}\fontsize{6.9pt}{7.9pt}\selectfont\strut – quand cheminé sur plateau d'où très\strut &  \\
      \leavevmode\fontsize{6.9pt}{7.9pt}\selectfont\strut belle vue circulaire!\strut &  \\
      \cline{1-1}
      \leavevmode\fontsize{6.9pt}{7.9pt}\selectfont\strut \underline{A St-Pol}, reposé lgt église — avant \& après\strut & \smash{\parbox[t]{39mm}{\fontsize{6.2pt}{7.0pt}\selectfont\raggedright\color{black!88}\textbf{17 lgt} \textit{abbr. pro} longtemps. \textbf{17 St} \textit{abbr. pro} Saint. — \textbf{18 plusies} \textit{abbr. pro} plusieurs. — \textbf{20} \textit{lect. incert. in ms.}}} \\
      \leavevmode\fontsize{6.9pt}{7.9pt}\selectfont\strut Parcouru plusies fois et lentement rues\strut &  \\
      \leavevmode\fontsize{6.9pt}{7.9pt}\selectfont\strut petite ville banale\strut &  \\
      \leavevmode\llap{\fontsize{6.4pt}{7.3pt}\selectfont\color{darkgray}20\hspace{1.8mm}}\fontsize{6.9pt}{7.9pt}\selectfont\strut Au retour, grosses nuées encre ciel — [mot surchargé] orage\strut &  \\
      \leavevmode\fontsize{6.9pt}{7.9pt}\selectfont\strut Rentré avant pluie —\strut &  \\
\end{tabular}}%
\par\vspace*{\fill}
% --- Registres de bas de page (1. Critique interne [si pertinent] ; 2. Loci paralleli en latin ; 3. Commentaire historique en retrait) ---
\noindent\rule{\textwidth}{0.5pt}\par\vspace{0.65mm}
{\fontsize{6.6pt}{7.7pt}\selectfont%
\noindent\normalfont \textbf{Unité de page :} Ce feuillet réunit le samedi 2 septembre (l.~1--8) et le début du dimanche 3 septembre (l.~9--22), dont la suite et les comptes se trouvent au f\textsuperscript{o}~18~r\textsuperscript{o}.\par\vspace{0.5mm}
\noindent\rule{22mm}{0.3pt}\par\vspace{0.45mm}
\noindent\normalfont\color{black!92} \textbf{13} (\textit{St-Pol}) cf. supra f\textsuperscript{o}~13~v\textsuperscript{o} (p.~\pageref{fol:27}), f\textsuperscript{o}~16~r\textsuperscript{o} (p.~\pageref{fol:32}), infra f\textsuperscript{o}~18~r\textsuperscript{o} (p.~\pageref{fol:36}), f\textsuperscript{o}~18~v\textsuperscript{o} (p.~\pageref{fol:37}), \textit{et passim}, Appendicem I, doc.~5 (p.~\pageref{annexe:doc5}).\par\vspace{0.55mm}
\noindent\rule{38mm}{0.4pt}\par\vspace{0.5mm}
\noindent\hspace*{5mm}\parbox{\dimexpr\textwidth-5mm\relax}{\fontsize{6.6pt}{7.7pt}\selectfont\color{black!84}Deux journées au cantonnement de Marquay: corvée matinale du samedi 2 septembre dans les champs détrempés de rosée, et promenade dominicale à pied le 3 septembre jusqu'à \textbf{Saint-Pol-sur-Ternoise} par le plateau dominant la vallée de la Ternoise (les notes descriptives et les comptes de cette sortie du 3 septembre se poursuivent au f\textsuperscript{o}~18~r\textsuperscript{o}).}\par}
\newpage
\phantomsection\label{fac:34}
% --- En-tête du fac-similé (Page impaire / droite : Photographie du manuscrit) ---
\noindent\rule{\textwidth}{0.5pt}\par\vspace{0.55mm}
\noindent\textbf{\fontsize{6.8pt}{7.8pt}\selectfont FAC-SIMILÉ EN VIS-À-VIS — FOLIO 17 RECTO} \hfill \textbf{\fontsize{6.8pt}{7.8pt}\selectfont SEPTEMBRE 1916}\par\vspace{0.25mm}
\noindent{\fontsize{6.8pt}{8.0pt}\selectfont\textit{Livre-Journal d'Hippolyte Combes} — Samedi 2 et dimanche 3 septembre 1916 (Marquay et Saint-Pol-sur-Ternoise (Pas-de-Calais))}\par\vspace{0.55mm}
\noindent\rule{\textwidth}{0.35pt}\par\vspace*{\fill}
{\centering
  \setlength{\fboxsep}{1.2mm}%
  \setlength{\fboxrule}{0.4pt}%
  \fcolorbox{gray!55}{white}{%
    \includegraphics[width=0.95\textwidth,height=0.81\textheight,keepaspectratio]{\facsimiledir/20251126_094324.jpg}%
  }\par}
\vspace*{\fill}
\newpage

% =============================================================================
% VIS-A-VIS SPREAD : FOLIO 17 VERSO (20251126_094345.jpg)
% =============================================================================
\cleardoubleevenpage
\renewcommand{\currentmonthyear}{Septembre 1916}
\renewcommand{\currentfoliolabel}{Folio 17 verso}
\phantomsection\label{fol:35}
% --- En-tête du folio (Page paire / gauche : Transcription, apparat & registres de bas de page) ---
\noindent\rule{\textwidth}{0.5pt}\par\vspace{0.55mm}
\noindent\textbf{\footnotesize FOLIO 17 VERSO} \hfill \textbf{\footnotesize SEPTEMBRE 1916}\par\vspace{0.25mm}
\noindent{\fontsize{6.9pt}{8.0pt}\selectfont Courses à Roëllecourt, préparatifs inutiles et déception du départ retardé}\par\vspace{0.2mm}
\noindent{\fontsize{6.9pt}{8.0pt}\selectfont\textit{Lundi 4 et mardi 5 septembre 1916} \hfill \textit{Marquay et Roëllecourt (Pas-de-Calais)}}\par\vspace{0.5mm}
\noindent\rule{\textwidth}{0.35pt}\par\vspace{0.8mm}
% --- Corps de la transcription diplomatique et apparat marginal en prolongement ---
{\renewcommand{\arraystretch}{0.93}%
\noindent\hspace*{3.5mm}\begin{tabular}{@{}p{87mm}@{\hspace{3.5mm}}p{39mm}@{}}
      \leavevmode\fontsize{7.3pt}{8.4pt}\selectfont\strut \underline{4 L} \hfill{} \underline{Sept 1916}\strut &  \\
      \leavevmode\fontsize{7.3pt}{8.4pt}\selectfont\strut Tout le jour, pluie\strut &  \\
      \leavevmode\fontsize{7.3pt}{8.4pt}\selectfont\strut – tour à tour fine \& averses\strut &  \\
      \leavevmode\fontsize{7.3pt}{8.4pt}\selectfont\strut Matin, venu à Roellecourt, sur 1 wagon\strut &  \\
      \leavevmode\llap{\fontsize{6.4pt}{7.3pt}\selectfont\color{darkgray}5\hspace{1.8mm}}\fontsize{7.3pt}{8.4pt}\selectfont\strut Retourné à pied.\strut &  \\
      \leavevmode\fontsize{7.3pt}{8.4pt}\selectfont\strut Après-midi, fait avec lenteur ts mes\strut & \smash{\parbox[t]{39mm}{\fontsize{6.2pt}{7.0pt}\selectfont\raggedright\color{black!88}\textbf{6 ts} \textit{abbr. pro} tous.}} \\
      \leavevmode\fontsize{7.3pt}{8.4pt}\selectfont\strut paquets : Ts mes préparatifs terminés à 5H1/2\strut & \smash{\parbox[t]{39mm}{\fontsize{6.2pt}{7.0pt}\selectfont\raggedright\color{black!88}\textbf{7 Ts} \textit{abbr. pro} Tous.}} \\
      \leavevmode\fontsize{7.3pt}{8.4pt}\selectfont\strut Me réjouis, à la pensée de quitter cet\strut &  \\
      \leavevmode\fontsize{7.3pt}{8.4pt}\selectfont\strut obscur \& boueux hameau, où ns avons été\strut & \smash{\parbox[t]{39mm}{\fontsize{6.2pt}{7.0pt}\selectfont\raggedright\color{black!88}\textbf{9 ns} \textit{abbr. pro} nous.}} \\
      \leavevmode\llap{\fontsize{6.4pt}{7.3pt}\selectfont\color{darkgray}10\hspace{1.8mm}}\fontsize{7.3pt}{8.4pt}\selectfont\strut si mal.\strut &  \\
      \cline{1-1}
      \leavevmode\fontsize{7.3pt}{8.4pt}\selectfont\strut \underline{5 Ma.} Matin, prairie, appris par 1 orderly\strut & \smash{\parbox[t]{39mm}{\fontsize{6.2pt}{7.0pt}\selectfont\raggedright\color{black!88}\textbf{11 Ma.} \textit{scil.} Mardi.}} \\
      \leavevmode\fontsize{7.3pt}{8.4pt}\selectfont\strut N.C.O. que départ retardé. Sorti pr me renseigner\strut & \smash{\parbox[t]{39mm}{\fontsize{6.2pt}{7.0pt}\selectfont\raggedright\color{black!88}\textbf{12 pr} \textit{abbr. pro} pour. \textbf{12 N.C.O.} \textit{angl.} Non-Commissioned Officer. — \textbf{13 confirmaton} \textit{abbr. pro} confirmation.}} \\
      \leavevmode\fontsize{7.3pt}{8.4pt}\selectfont\strut eu confirmaton mauvaise nouvelle : encore 8\strut &  \\
      \leavevmode\fontsize{7.3pt}{8.4pt}\selectfont\strut jours à moisir ici, hélas!\strut &  \\
      \leavevmode\llap{\fontsize{6.4pt}{7.3pt}\selectfont\color{darkgray}15\hspace{1.8mm}}\fontsize{7.3pt}{8.4pt}\selectfont\strut Et Toujours, même temps affreux –\strut &  \\
      \leavevmode\fontsize{7.3pt}{8.4pt}\selectfont\strut Ciel couvert de nuages – Pluie intermittente\strut &  \\
      \leavevmode\fontsize{7.3pt}{8.4pt}\selectfont\strut Après-midi, malgré averses\strut &  \\
      \leavevmode\fontsize{7.3pt}{8.4pt}\selectfont\strut menaçantes, venu à pied à Roellecourt.\strut &  \\
      \leavevmode\fontsize{7.3pt}{8.4pt}\selectfont\strut Au retour, sous pluie battante\strut &  \\
      \leavevmode\llap{\fontsize{6.4pt}{7.3pt}\selectfont\color{darkgray}20\hspace{1.8mm}}\fontsize{7.3pt}{8.4pt}\selectfont\strut Temps abominable jusqu'à fin jour –\strut &  \\
      \leavevmode\fontsize{7.3pt}{8.4pt}\selectfont\strut Mission Fr.se : solde \& Indemnité :\strut & \smash{\parbox[t]{39mm}{\fontsize{6.2pt}{7.0pt}\selectfont\raggedright\color{black!88}\textbf{21 Fr.se} \textit{abbr. pro} Française.}} \\
      \leavevmode\fontsize{7.3pt}{8.4pt}\selectfont\strut 78 f 75c\strut &  \\
\end{tabular}}%
\par\vspace*{\fill}
% --- Registres de bas de page (1. Critique interne [si pertinent] ; 2. Loci paralleli en latin ; 3. Commentaire historique en retrait) ---
\noindent\rule{\textwidth}{0.5pt}\par\vspace{0.65mm}
{\fontsize{6.6pt}{7.7pt}\selectfont%
\noindent\normalfont \textbf{Ordre des folios :} Le f\textsuperscript{o}~17~v\textsuperscript{o} porte les entrées des lundi 4 et mardi 5 septembre 1916, tandis que le f\textsuperscript{o}~18~r\textsuperscript{o} contient la fin de l'entrée du dimanche 3 septembre : Hippolyte Combes a rédigé les journées des 4 et 5 septembre sur la page de gauche avant ou après avoir complété le compte du 3 septembre sur la page en regard.\par\vspace{0.5mm}
\noindent\rule{22mm}{0.3pt}\par\vspace{0.45mm}
\noindent\normalfont\color{black!92} \textbf{1} (\textit{de eodem commorationis loco}) cf. supra f\textsuperscript{o}~17~r\textsuperscript{o} (p.~\pageref{fol:34}).\par\vspace{0.55mm}
\noindent\rule{38mm}{0.4pt}\par\vspace{0.5mm}
\noindent\hspace*{5mm}\parbox{\dimexpr\textwidth-5mm\relax}{\fontsize{6.6pt}{7.7pt}\selectfont\color{black!84}Après avoir préparé ses bagages le lundi 4 septembre, Hippolyte Combes apprend le mardi 5 septembre par un sous-officier de service (\textit{orderly N.C.O.}) que le départ de la division est différé (il aura finalement lieu le samedi 9 septembre, f\textsuperscript{o}~19~v\textsuperscript{o}).}\par}
\newpage
\phantomsection\label{fac:35}
% --- En-tête du fac-similé (Page impaire / droite : Photographie du manuscrit) ---
\noindent\rule{\textwidth}{0.5pt}\par\vspace{0.55mm}
\noindent\textbf{\fontsize{6.8pt}{7.8pt}\selectfont FAC-SIMILÉ EN VIS-À-VIS — FOLIO 17 VERSO} \hfill \textbf{\fontsize{6.8pt}{7.8pt}\selectfont SEPTEMBRE 1916}\par\vspace{0.25mm}
\noindent{\fontsize{6.8pt}{8.0pt}\selectfont\textit{Livre-Journal d'Hippolyte Combes} — Lundi 4 et mardi 5 septembre 1916 (Marquay et Roëllecourt (Pas-de-Calais))}\par\vspace{0.55mm}
\noindent\rule{\textwidth}{0.35pt}\par\vspace*{\fill}
{\centering
  \setlength{\fboxsep}{1.2mm}%
  \setlength{\fboxrule}{0.4pt}%
  \fcolorbox{gray!55}{white}{%
    \includegraphics[width=0.95\textwidth,height=0.81\textheight,keepaspectratio]{\facsimiledir/20251126_094345.jpg}%
  }\par}
\vspace*{\fill}
\newpage

% =============================================================================
% VIS-A-VIS SPREAD : FOLIO 18 RECTO (20251126_094449.jpg)
% =============================================================================
\cleardoubleevenpage
\renewcommand{\currentmonthyear}{Septembre 1916}
\renewcommand{\currentfoliolabel}{Folio 18 recto}
\phantomsection\label{fol:36}
% --- En-tête du folio (Page paire / gauche : Transcription, apparat & registres de bas de page) ---
\noindent\rule{\textwidth}{0.5pt}\par\vspace{0.55mm}
\noindent\textbf{\footnotesize FOLIO 18 RECTO} \hfill \textbf{\footnotesize SEPTEMBRE 1916}\par\vspace{0.25mm}
\noindent{\fontsize{6.9pt}{8.0pt}\selectfont Champs d'avoine sous la brise et achat de la carte du département de la Somme}\par\vspace{0.2mm}
\noindent{\fontsize{6.9pt}{8.0pt}\selectfont\textit{Dimanche 3 (addition) / mardi 5 septembre 1916} \hfill \textit{Marquay et Saint-Pol-sur-Ternoise (Pas-de-Calais)}}\par\vspace{0.5mm}
\noindent\rule{\textwidth}{0.35pt}\par\vspace{0.8mm}
% --- Corps de la transcription diplomatique et apparat marginal en prolongement ---
{\renewcommand{\arraystretch}{1.00}%
\noindent\hspace*{3.5mm}\begin{tabular}{@{}p{87mm}@{\hspace{3.5mm}}p{39mm}@{}}
      \leavevmode\fontsize{7.7pt}{9.0pt}\selectfont\strut \underline{3 D.} \hfill{} \underline{Septbre /16}\strut & \smash{\parbox[t]{39mm}{\fontsize{6.2pt}{7.0pt}\selectfont\raggedright\color{black!88}\textbf{1 D.} \textit{scil.} Dimanche.}} \\
      \leavevmode\fontsize{7.7pt}{9.0pt}\selectfont\strut Champs d'avoine non encore\strut &  \\
      \leavevmode\fontsize{7.7pt}{9.0pt}\selectfont\strut fauchés, tout jaunes — et qui\strut &  \\
      \leavevmode\fontsize{7.7pt}{9.0pt}\selectfont\strut bruissaient monotone, sous la caresse\strut &  \\
      \leavevmode\llap{\fontsize{6.4pt}{7.3pt}\selectfont\color{darkgray}5\hspace{1.8mm}}\fontsize{7.7pt}{9.0pt}\selectfont\strut de la brise. –\strut &  \\
      \cline{1-1}
      \leavevmode\fontsize{7.7pt}{9.0pt}\selectfont\strut \underline{A Saint-Pol} : ctes ples : 30 c. + [2 f / 25 c ?]\strut & \smash{\parbox[t]{39mm}{\fontsize{6.2pt}{7.0pt}\selectfont\raggedright\color{black!88}\textbf{6} \textit{lect. incert. in ms.} \textbf{6 ctes ples} \textit{abbr. pro} cartes postales. — \textbf{7 c.} \textit{abbr. pro} centimes. — \textbf{8 saddle soap} \textit{angl. in ms.} — \textbf{9 c.} \textit{abbr. pro} centimes.}} \\
      \leavevmode\fontsize{7.7pt}{9.0pt}\selectfont\strut Carte du dépt Somme : 70 c. \hfill{} Géog. du Pas de Calais : 35 c\strut &  \\
      \leavevmode\fontsize{7.7pt}{9.0pt}\selectfont\strut 1 boîte de saddle soap : 1 f 50 \hfill{} Varia : 30 c\strut &  \\
      \leavevmode\fontsize{7.7pt}{9.0pt}\selectfont\strut 1 éponge : 60 c. – – 1/4 livre choc Poulain : 1f\strut &  \\
\end{tabular}}%
\par\vspace*{\fill}
% --- Registres de bas de page (1. Critique interne [si pertinent] ; 2. Loci paralleli en latin ; 3. Commentaire historique en retrait) ---
\noindent\rule{\textwidth}{0.5pt}\par\vspace{0.65mm}
{\fontsize{6.6pt}{7.7pt}\selectfont%
\noindent\normalfont \textbf{Continuité :} Fin du récit et relevé des dépenses de la journée du dimanche 3 septembre 1916, faisant suite au f\textsuperscript{o}~17~r\textsuperscript{o}.\par\vspace{0.5mm}
\noindent\rule{22mm}{0.3pt}\par\vspace{0.45mm}
\noindent\normalfont\color{black!92} \textbf{6} (\textit{Saint-Pol}) cf. supra f\textsuperscript{o}~13~v\textsuperscript{o} (p.~\pageref{fol:27}), f\textsuperscript{o}~17~r\textsuperscript{o} (p.~\pageref{fol:34}), infra f\textsuperscript{o}~18~v\textsuperscript{o} (p.~\pageref{fol:37}), \textit{et passim}, Appendicem I, doc.~5 (p.~\pageref{annexe:doc5}).\par\vspace{0.55mm}
\noindent\rule{38mm}{0.4pt}\par\vspace{0.5mm}
\noindent\hspace*{5mm}\parbox{\dimexpr\textwidth-5mm\relax}{\fontsize{6.6pt}{7.7pt}\selectfont\color{black!84}Suite des notes et comptes du dimanche 3 septembre 1916 à \textbf{Saint-Pol-sur-Ternoise}.}\par}
\newpage
\phantomsection\label{fac:36}
% --- En-tête du fac-similé (Page impaire / droite : Photographie du manuscrit) ---
\noindent\rule{\textwidth}{0.5pt}\par\vspace{0.55mm}
\noindent\textbf{\fontsize{6.8pt}{7.8pt}\selectfont FAC-SIMILÉ EN VIS-À-VIS — FOLIO 18 RECTO} \hfill \textbf{\fontsize{6.8pt}{7.8pt}\selectfont SEPTEMBRE 1916}\par\vspace{0.25mm}
\noindent{\fontsize{6.8pt}{8.0pt}\selectfont\textit{Livre-Journal d'Hippolyte Combes} — Dimanche 3 (addition) / mardi 5 septembre 1916 (Marquay et Saint-Pol-sur-Ternoise (Pas-de-Calais))}\par\vspace{0.55mm}
\noindent\rule{\textwidth}{0.35pt}\par\vspace*{\fill}
{\centering
  \setlength{\fboxsep}{1.2mm}%
  \setlength{\fboxrule}{0.4pt}%
  \fcolorbox{gray!55}{white}{%
    \includegraphics[width=0.95\textwidth,height=0.81\textheight,keepaspectratio]{\facsimiledir/20251126_094449.jpg}%
  }\par}
\vspace*{\fill}
\newpage

% =============================================================================
% VIS-A-VIS SPREAD : FOLIO 18 VERSO (20251126_094511.jpg)
% =============================================================================
\cleardoubleevenpage
\renewcommand{\currentmonthyear}{Septembre 1916}
\renewcommand{\currentfoliolabel}{Folio 18 verso}
\phantomsection\label{fol:37}
% --- En-tête du folio (Page paire / gauche : Transcription, apparat & registres de bas de page) ---
\noindent\rule{\textwidth}{0.5pt}\par\vspace{0.55mm}
\noindent\textbf{\footnotesize FOLIO 18 VERSO} \hfill \textbf{\footnotesize SEPTEMBRE 1916}\par\vspace{0.25mm}
\noindent{\fontsize{6.9pt}{8.0pt}\selectfont Conversation avec Mme Vauchel et halte au cimetière de Saint-Michel-sur-Ternoise}\par\vspace{0.2mm}
\noindent{\fontsize{6.9pt}{8.0pt}\selectfont\textit{Mercredi 6 et jeudi 7 septembre 1916} \hfill \textit{Marquay, Bailleul-aux-Cornailles, Saint-Pol et Saint-Michel-sur-Ternoise (Pas-de-Calais)}}\par\vspace{0.5mm}
\noindent\rule{\textwidth}{0.35pt}\par\vspace{0.8mm}
% --- Corps de la transcription diplomatique et apparat marginal en prolongement ---
{\renewcommand{\arraystretch}{0.83}%
\noindent\hspace*{3.5mm}\begin{tabular}{@{}p{87mm}@{\hspace{3.5mm}}p{39mm}@{}}
      \leavevmode\fontsize{6.6pt}{7.5pt}\selectfont\strut \underline{6 Me} Enfin, beau temps. – \hfill{} \raisebox{0.5ex}{\underline{Septbre 1916}}\strut &  \\
      \leavevmode\fontsize{6.6pt}{7.5pt}\selectfont\strut Après-midi superbe Chaud soleil.\strut &  \\
      \leavevmode\fontsize{6.6pt}{7.5pt}\selectfont\strut Resté 2 h sur pature, assis chaise à\strut &  \\
      \leavevmode\fontsize{6.6pt}{7.5pt}\selectfont\strut l'ombre de vieux arbres. [Bain ?] de gd air.\strut & \smash{\parbox[t]{39mm}{\fontsize{6.2pt}{7.0pt}\selectfont\raggedright\color{black!88}\textbf{4} \textit{lect. incert. in ms.} \textbf{4 gd} \textit{abbr. pro} grand. — \textbf{5 Lt} \textit{abbr. pro} Lieutenant. \textbf{5 Ctis} \textit{abbr. pro} Curtis. — \textbf{7 Mme} \textit{abbr. pro} Madame. \textbf{7 propre} \textit{abbr. pro} propriétaire. — \textbf{8 admirabt} \textit{abbr. pro} admirablement. — \textbf{9 ds} \textit{abbr. pro} dans. \textbf{9 tjs} \textit{abbr. pro} toujours. — \textbf{12 déjeun.} \textit{abbr. pro} déjeuné / déjeuner. — \textbf{13 ds} \textit{abbr. pro} dans. — \textbf{14 ds} \textit{abbr. pro} dans.}} \\
      \leavevmode\llap{\fontsize{6.4pt}{7.3pt}\selectfont\color{darkgray}5\hspace{1.8mm}}\fontsize{6.6pt}{7.5pt}\selectfont\strut A 6 H s., dû accompagner Lt Ctis à\strut &  \\
      \leavevmode\fontsize{6.6pt}{7.5pt}\selectfont\strut ferme voisinage. Y très bien accueillis\strut &  \\
      \leavevmode\fontsize{6.6pt}{7.5pt}\selectfont\strut par propre, Mme Vauchel, femme très\strut &  \\
      \leavevmode\fontsize{6.6pt}{7.5pt}\selectfont\strut intelligente \& raisonnant admirabt\strut &  \\
      \leavevmode\fontsize{6.6pt}{7.5pt}\selectfont\strut Languis tjs ds ce trou perdu, et\strut &  \\
      \leavevmode\llap{\fontsize{6.4pt}{7.3pt}\selectfont\color{darkgray}10\hspace{1.8mm}}\fontsize{6.6pt}{7.5pt}\selectfont\strut sale, où aucune distraction —\strut &  \\
      \cline{1-1}
      \leavevmode\fontsize{6.6pt}{7.5pt}\selectfont\strut \underline{7 J} Encore une b. journée –\strut &  \\
      \leavevmode\fontsize{6.6pt}{7.5pt}\selectfont\strut Sitôt après-déjeun, parti campagne. Marché\strut &  \\
      \leavevmode\fontsize{6.6pt}{7.5pt}\selectfont\strut lentement et fait haltes fréquentes ds le\strut &  \\
      \leavevmode\fontsize{6.6pt}{7.5pt}\selectfont\strut carré gazonné de calvaires, ds 1 bois ou au\strut &  \\
      \leavevmode\llap{\fontsize{6.4pt}{7.3pt}\selectfont\color{darkgray}15\hspace{1.8mm}}\fontsize{6.6pt}{7.5pt}\selectfont\strut bord champs d'avoine\strut &  \\
      \leavevmode\fontsize{6.6pt}{7.5pt}\selectfont\strut Soleil ardent\strut &  \\
      \leavevmode\fontsize{6.6pt}{7.5pt}\selectfont\strut Passé [1/2 ?] longue[s?] heure[s?] ds cette petite ville de\strut & \smash{\parbox[t]{39mm}{\fontsize{6.2pt}{7.0pt}\selectfont\raggedright\color{black!88}\textbf{17} \textit{lect. incert. in ms.} \textbf{17 ds} \textit{abbr. pro} dans. — \textbf{18 St} \textit{abbr. pro} Saint. — \textbf{19 gde} \textit{abbr. pro} grande.}} \\
      \leavevmode\fontsize{6.6pt}{7.5pt}\selectfont\strut St-Pol : aux maisons briques. rues en pente\strut &  \\
      \leavevmode\fontsize{6.6pt}{7.5pt}\selectfont\strut aboutissant à gde place et eglise\strut &  \\
      \leavevmode\llap{\fontsize{6.4pt}{7.3pt}\selectfont\color{darkgray}20\hspace{1.8mm}}\fontsize{6.6pt}{7.5pt}\selectfont\strut Retour par St. Michel-sur-Ternoise : arrêt prolongé\strut &  \\
      \leavevmode\fontsize{6.6pt}{7.5pt}\selectfont\strut ds le petit cimetière ombragé par vieux [ifs/arbres ?] et tt\strut & \smash{\parbox[t]{39mm}{\fontsize{6.2pt}{7.0pt}\selectfont\raggedright\color{black!88}\textbf{21} \textit{lect. incert. in ms.} \textbf{21 ds} \textit{abbr. pro} dans. — \textbf{22 1/2} \textit{scil.} demi / moitié.}} \\
      \leavevmode\fontsize{6.6pt}{7.5pt}\selectfont\strut envahi par hautes herbes — au 1/2 s'élève l'église\strut &  \\
      \leavevmode\fontsize{6.6pt}{7.5pt}\selectfont\strut dont le clocher octogonal en pierre sculptée est\strut &  \\
      \leavevmode\fontsize{6.6pt}{7.5pt}\selectfont\strut intéressant\strut &  \\
      \leavevmode\llap{\fontsize{6.4pt}{7.3pt}\selectfont\color{darkgray}25\hspace{1.8mm}}\fontsize{6.6pt}{7.5pt}\selectfont\strut [groupe biffé illis.] allée de pins \hfill{} [ Ctes ples : 1 f. — Hair cut : 50 c\strut & \smash{\parbox[t]{39mm}{\fontsize{6.2pt}{7.0pt}\selectfont\raggedright\color{black!88}\textbf{25} \textit{lect. incert. in ms.} \textbf{25} \textit{in marg. cum unco} (comptes).}} \\
      \leavevmode\fontsize{6.6pt}{7.5pt}\selectfont\strut  \hfill{} Consommations : 50 c ]\strut &  \\
\end{tabular}}%
\par\vspace*{\fill}
% --- Registres de bas de page (1. Critique interne [si pertinent] ; 2. Loci paralleli en latin ; 3. Commentaire historique en retrait) ---
\noindent\rule{\textwidth}{0.5pt}\par\vspace{0.65mm}
{\fontsize{6.6pt}{7.7pt}\selectfont%
\noindent\normalfont\color{black!92} \textbf{5} (\textit{Ctis}) cf. supra f\textsuperscript{o}~1~r\textsuperscript{o} (p.~\pageref{fol:2}), f\textsuperscript{o}~15~r\textsuperscript{o} (p.~\pageref{fol:30}), infra f\textsuperscript{o}~44~v\textsuperscript{o} (p.~\pageref{fol:89}), f\textsuperscript{o}~45~r\textsuperscript{o} (p.~\pageref{fol:90}), \textit{et passim}. — \textbf{18} (\textit{St-Pol}) cf. supra f\textsuperscript{o}~13~v\textsuperscript{o} (p.~\pageref{fol:27}), f\textsuperscript{o}~18~r\textsuperscript{o} (p.~\pageref{fol:36}), \textit{et passim}.\par\vspace{0.55mm}
\noindent\rule{38mm}{0.4pt}\par\vspace{0.5mm}
\noindent\hspace*{5mm}\parbox{\dimexpr\textwidth-5mm\relax}{\fontsize{6.6pt}{7.7pt}\selectfont\color{black!84}Le mercredi 6 septembre, visite de service avec le lieutenant Curtis chez une exploitante de Marquay, \textbf{Mme Vauchel}. Le jeudi 7 septembre, revenant à pied de Saint-Pol, Hippolyte Combes s'attarde dans le cimetière entourant l'église \textbf{Saint-Michel de Saint-Michel-sur-Ternoise} (Pas-de-Calais), édifice de pierre blanche reconstruit au XVI\textsuperscript{e}~s. (chœur de 1591) et doté d'un clocher de pierre du XVIII\textsuperscript{e}~s..}\par}
\newpage
\phantomsection\label{fac:37}
% --- En-tête du fac-similé (Page impaire / droite : Photographie du manuscrit) ---
\noindent\rule{\textwidth}{0.5pt}\par\vspace{0.55mm}
\noindent\textbf{\fontsize{6.8pt}{7.8pt}\selectfont FAC-SIMILÉ EN VIS-À-VIS — FOLIO 18 VERSO} \hfill \textbf{\fontsize{6.8pt}{7.8pt}\selectfont SEPTEMBRE 1916}\par\vspace{0.25mm}
\noindent{\fontsize{6.8pt}{8.0pt}\selectfont\textit{Livre-Journal d'Hippolyte Combes} — Mercredi 6 et jeudi 7 septembre 1916 (Marquay, Bailleul-aux-Cornailles, Saint-Pol et Saint-Michel-sur-Ternoise (Pas-de-Calais))}\par\vspace{0.55mm}
\noindent\rule{\textwidth}{0.35pt}\par\vspace*{\fill}
{\centering
  \setlength{\fboxsep}{1.2mm}%
  \setlength{\fboxrule}{0.4pt}%
  \fcolorbox{gray!55}{white}{%
    \includegraphics[width=0.95\textwidth,height=0.81\textheight,keepaspectratio]{\facsimiledir/20251126_094511.jpg}%
  }\par}
\vspace*{\fill}
\newpage

% =============================================================================
% VIS-A-VIS SPREAD : FOLIO 19 RECTO (20251126_094524.jpg)
% =============================================================================
\cleardoubleevenpage
\renewcommand{\currentmonthyear}{Septembre 1916}
\renewcommand{\currentfoliolabel}{Folio 19 recto}
\phantomsection\label{fol:38}
% --- En-tête du folio (Page paire / gauche : Transcription, apparat & registres de bas de page) ---
\noindent\rule{\textwidth}{0.5pt}\par\vspace{0.55mm}
\noindent\textbf{\footnotesize FOLIO 19 RECTO} \hfill \textbf{\footnotesize SEPTEMBRE 1916}\par\vspace{0.25mm}
\noindent{\fontsize{6.9pt}{8.0pt}\selectfont Attente des fourgons des « Bedfords » et promenade au clair de lune à Ryonval}\par\vspace{0.2mm}
\noindent{\fontsize{6.9pt}{8.0pt}\selectfont\textit{Vendredi 8 septembre 1916} \hfill \textit{Marquay et ferme de Ryonval (Pas-de-Calais)}}\par\vspace{0.5mm}
\noindent\rule{\textwidth}{0.35pt}\par\vspace{0.8mm}
% --- Corps de la transcription diplomatique et apparat marginal en prolongement ---
{\renewcommand{\arraystretch}{1.00}%
\noindent\hspace*{3.5mm}\begin{tabular}{@{}p{87mm}@{\hspace{3.5mm}}p{39mm}@{}}
      \leavevmode\fontsize{7.7pt}{9.0pt}\selectfont\strut \underline{8 Ve.} \hfill{} \underline{Septbre /16}\strut & \smash{\parbox[t]{39mm}{\fontsize{6.2pt}{7.0pt}\selectfont\raggedright\color{black!88}\textbf{1 Ve.} \textit{scil.} Vendredi.}} \\
      \leavevmode\fontsize{7.7pt}{9.0pt}\selectfont\strut Tout l'après-midi,\strut &  \\
      \leavevmode\fontsize{7.7pt}{9.0pt}\selectfont\strut dû attendre et inutilt 2 wagons de\strut & \smash{\parbox[t]{39mm}{\fontsize{6.2pt}{7.0pt}\selectfont\raggedright\color{black!88}\textbf{3 inutilt} \textit{abbr. pro} inutilement. — \textbf{4 green} \textit{angl. in ms.} — \textbf{5 lgt} \textit{abbr. pro} longtemps. — \textbf{6 ds} \textit{abbr. pro} dans.}} \\
      \leavevmode\fontsize{7.7pt}{9.0pt}\selectfont\strut Bedfords : for green forage\strut &  \\
      \leavevmode\llap{\fontsize{6.4pt}{7.3pt}\selectfont\color{darkgray}5\hspace{1.8mm}}\fontsize{7.7pt}{9.0pt}\selectfont\strut Après diner, soir, promené lgt\strut &  \\
      \leavevmode\fontsize{7.7pt}{9.0pt}\selectfont\strut ds la campagne\strut &  \\
      \leavevmode\fontsize{7.7pt}{9.0pt}\selectfont\strut Clair de lune superbe\strut &  \\
      \leavevmode\fontsize{7.7pt}{9.0pt}\selectfont\strut arbres nettement silhouettés\strut &  \\
      \leavevmode\fontsize{7.7pt}{9.0pt}\selectfont\strut me suis attardé petit chemin de Ryonval\strut &  \\
      \leavevmode\llap{\fontsize{6.4pt}{7.3pt}\selectfont\color{darkgray}10\hspace{1.8mm}}\fontsize{7.7pt}{9.0pt}\selectfont\strut ds la pâture derrière ferme Collet\strut & \smash{\parbox[t]{39mm}{\fontsize{6.2pt}{7.0pt}\selectfont\raggedright\color{black!88}\textbf{10 ds} \textit{abbr. pro} dans.}} \\
      \leavevmode\fontsize{7.7pt}{9.0pt}\selectfont\strut Joui délict du silence nuit.\strut & \smash{\parbox[t]{39mm}{\fontsize{6.2pt}{7.0pt}\selectfont\raggedright\color{black!88}\textbf{11 délict} \textit{abbr. pro} délicieusement. — \textbf{12 tooth powder} \textit{angl. in ms.}}} \\
      \leavevmode\fontsize{7.7pt}{9.0pt}\selectfont\strut — Box of tooth powder : 60 c\strut &  \\
\end{tabular}}%
\par\vspace*{\fill}
% --- Registres de bas de page (1. Critique interne [si pertinent] ; 2. Loci paralleli en latin ; 3. Commentaire historique en retrait) ---
\noindent\rule{\textwidth}{0.5pt}\par\vspace{0.65mm}
{\fontsize{6.6pt}{7.7pt}\selectfont%
\noindent\normalfont\color{black!92} \textbf{4} (\textit{Bedfords}) cf. infra f\textsuperscript{o}~23~v\textsuperscript{o} (p.~\pageref{fol:47}). — \textbf{4} (\textit{green forage}) cf. supra f\textsuperscript{o}~3~r\textsuperscript{o} (p.~\pageref{fol:6}), f\textsuperscript{o}~15~v\textsuperscript{o} (p.~\pageref{fol:31}), infra f\textsuperscript{o}~21~r\textsuperscript{o} (p.~\pageref{fol:42}), f\textsuperscript{o}~22~r\textsuperscript{o} (p.~\pageref{fol:44}), \textit{et passim}.\par\vspace{0.55mm}
\noindent\rule{38mm}{0.4pt}\par\vspace{0.5mm}
\noindent\hspace*{5mm}\parbox{\dimexpr\textwidth-5mm\relax}{\fontsize{6.6pt}{7.7pt}\selectfont\color{black!84}Les deux fourgons (\textit{waggons}) venus charger le fourrage vert derrière la ferme Collet appartiennent au \textbf{\textit{Bedfordshire Regiment}} (\textit{Bedfords}), dont le 7\textsuperscript{e} Bataillon (\textit{Service}) fait partie de la 54\textsuperscript{e} Brigade de la \textbf{18\textsuperscript{e} Division britannique} (Nichols, 1922, p.~86).}\par}
\newpage
\phantomsection\label{fac:38}
% --- En-tête du fac-similé (Page impaire / droite : Photographie du manuscrit) ---
\noindent\rule{\textwidth}{0.5pt}\par\vspace{0.55mm}
\noindent\textbf{\fontsize{6.8pt}{7.8pt}\selectfont FAC-SIMILÉ EN VIS-À-VIS — FOLIO 19 RECTO} \hfill \textbf{\fontsize{6.8pt}{7.8pt}\selectfont SEPTEMBRE 1916}\par\vspace{0.25mm}
\noindent{\fontsize{6.8pt}{8.0pt}\selectfont\textit{Livre-Journal d'Hippolyte Combes} — Vendredi 8 septembre 1916 (Marquay et ferme de Ryonval (Pas-de-Calais))}\par\vspace{0.55mm}
\noindent\rule{\textwidth}{0.35pt}\par\vspace*{\fill}
{\centering
  \setlength{\fboxsep}{1.2mm}%
  \setlength{\fboxrule}{0.4pt}%
  \fcolorbox{gray!55}{white}{%
    \includegraphics[width=0.95\textwidth,height=0.81\textheight,keepaspectratio]{\facsimiledir/20251126_094524.jpg}%
  }\par}
\vspace*{\fill}
\newpage

% =============================================================================
% VIS-A-VIS SPREAD : FOLIO 19 VERSO (20251126_094545.jpg)
% =============================================================================
\cleardoubleevenpage
\renewcommand{\currentmonthyear}{Septembre 1916}
\renewcommand{\currentfoliolabel}{Folio 19 verso}
\phantomsection\label{fol:39}
% --- En-tête du folio (Page paire / gauche : Transcription, apparat & registres de bas de page) ---
\noindent\rule{\textwidth}{0.5pt}\par\vspace{0.55mm}
\noindent\textbf{\footnotesize FOLIO 19 VERSO} \hfill \textbf{\footnotesize SEPTEMBRE 1916}\par\vspace{0.25mm}
\noindent{\fontsize{6.9pt}{8.0pt}\selectfont Départ de Marquay, l'école de mortiers de tranchée et l'ancien monastère de Monts-en-Ternois}\par\vspace{0.2mm}
\noindent{\fontsize{6.9pt}{8.0pt}\selectfont\textit{Samedi 9 septembre 1916} \hfill \textit{De Marquay et Ligny-Saint-Flochel à Monts-en-Ternois (Pas-de-Calais)}}\par\vspace{0.5mm}
\noindent\rule{\textwidth}{0.35pt}\par\vspace{0.8mm}
% --- Corps de la transcription diplomatique et apparat marginal en prolongement ---
{\renewcommand{\arraystretch}{0.87}%
\noindent\hspace*{3.5mm}\begin{tabular}{@{}p{87mm}@{\hspace{3.5mm}}p{39mm}@{}}
      \leavevmode\fontsize{6.9pt}{7.9pt}\selectfont\strut \underline{Sam. 9.} \hfill{} \underline{Sept / 16}\strut & \smash{\parbox[t]{39mm}{\fontsize{6.2pt}{7.0pt}\selectfont\raggedright\color{black!88}\textbf{1 Sam.} \textit{scil.} Samedi.}} \\
      \leavevmode\fontsize{6.9pt}{7.9pt}\selectfont\strut Quitté Marquay vers 9 H 1/2 matin\strut &  \\
      \leavevmode\fontsize{6.9pt}{7.9pt}\selectfont\strut 3 brigades division défilé avant ns\strut & \smash{\parbox[t]{39mm}{\fontsize{6.2pt}{7.0pt}\selectfont\raggedright\color{black!88}\textbf{3 ns} \textit{abbr. pro} nous.}} \\
      \leavevmode\fontsize{6.9pt}{7.9pt}\selectfont\strut Très lente chevauchée; arrêts fréquents\strut &  \\
      \leavevmode\llap{\fontsize{6.4pt}{7.3pt}\selectfont\color{darkgray}5\hspace{1.8mm}}\fontsize{6.9pt}{7.9pt}\selectfont\strut et prolongés.\strut &  \\
      \leavevmode\fontsize{6.9pt}{7.9pt}\selectfont\strut Passé par Ligny - St Flochel\strut & \smash{\parbox[t]{39mm}{\fontsize{6.2pt}{7.0pt}\selectfont\raggedright\color{black!88}\textbf{6 St} \textit{abbr. pro} Saint.}} \\
      \leavevmode\fontsize{6.9pt}{7.9pt}\selectfont\strut puis atteint plateau où trench mortar\strut &  \\
      \leavevmode\fontsize{6.9pt}{7.9pt}\selectfont\strut school, et d'où vue panoramique\strut &  \\
      \leavevmode\fontsize{6.9pt}{7.9pt}\selectfont\strut contrée environnante\strut &  \\
      \leavevmode\llap{\fontsize{6.4pt}{7.3pt}\selectfont\color{darkgray}10\hspace{1.8mm}}\fontsize{6.9pt}{7.9pt}\selectfont\strut Temps tt fait beau – Soleil supportable\strut & \smash{\parbox[t]{39mm}{\fontsize{6.2pt}{7.0pt}\selectfont\raggedright\color{black!88}\textbf{10 tt} \textit{abbr. pro} tout.}} \\
      \leavevmode\fontsize{6.9pt}{7.9pt}\selectfont\strut [Peu/Flots ?] de poussière route[s?]\strut & \smash{\parbox[t]{39mm}{\fontsize{6.2pt}{7.0pt}\selectfont\raggedright\color{black!88}\textbf{11} \textit{lect. incert. in ms.}}} \\
      \leavevmode\fontsize{6.9pt}{7.9pt}\selectfont\strut Vers 1H s., terminé 1ère étape, au\strut & \smash{\parbox[t]{39mm}{\fontsize{6.2pt}{7.0pt}\selectfont\raggedright\color{black!88}\textbf{12 1ère} \textit{scil.} première. \textbf{12 s.} \textit{abbr. pro} soir.}} \\
      \leavevmode\fontsize{6.9pt}{7.9pt}\selectfont\strut hameau de Monts-en-Ternois.\strut &  \\
      \leavevmode\fontsize{6.9pt}{7.9pt}\selectfont\strut Cantonnés ds gde ferme, à l'[entrée ?] du\strut & \smash{\parbox[t]{39mm}{\fontsize{6.2pt}{7.0pt}\selectfont\raggedright\color{black!88}\textbf{14} \textit{lect. incert. in ms.} \textbf{14 ds} \textit{abbr. pro} dans.}} \\
      \leavevmode\llap{\fontsize{6.4pt}{7.3pt}\selectfont\color{darkgray}15\hspace{1.8mm}}\fontsize{6.9pt}{7.9pt}\selectfont\strut village – maison d'habitation, ancien\strut &  \\
      \leavevmode\fontsize{6.9pt}{7.9pt}\selectfont\strut monastère –\strut &  \\
      \leavevmode\fontsize{6.9pt}{7.9pt}\selectfont\strut Impression de calme reposant jardin\strut &  \\
      \leavevmode\fontsize{6.9pt}{7.9pt}\selectfont\strut fleuri — allongé s/ pelouse\strut & \smash{\parbox[t]{39mm}{\fontsize{6.2pt}{7.0pt}\selectfont\raggedright\color{black!88}\textbf{18 s/} \textit{abbr. pro} sur.}} \\
      \leavevmode\fontsize{6.9pt}{7.9pt}\selectfont\strut Invité par Cap. West à messer av. offs\strut & \smash{\parbox[t]{39mm}{\fontsize{6.2pt}{7.0pt}\selectfont\raggedright\color{black!88}\textbf{19 offs} \textit{abbr. pro} officiers. \textbf{19 Cap.} \textit{abbr. pro} Capitaine. — \textbf{21 ds} \textit{abbr. pro} dans.}} \\
      \leavevmode\llap{\fontsize{6.4pt}{7.3pt}\selectfont\color{darkgray}20\hspace{1.8mm}}\fontsize{6.9pt}{7.9pt}\selectfont\strut Et le soir\strut &  \\
      \leavevmode\fontsize{6.9pt}{7.9pt}\selectfont\strut Dormi sur 2 bancs de jardin ds la laiterie\strut &  \\
      \leavevmode\fontsize{6.9pt}{7.9pt}\selectfont\strut excellente nuit.\strut &  \\
      \leavevmode\fontsize{6.9pt}{7.9pt}\selectfont\strut — Donné 3 f à M. Dupuich à Marquay —\strut &  \\
      \leavevmode\fontsize{6.9pt}{7.9pt}\selectfont\strut \& 1 tip au groom [Connelly ?] —\strut & \smash{\parbox[t]{39mm}{\fontsize{6.2pt}{7.0pt}\selectfont\raggedright\color{black!88}\textbf{24} \textit{lect. incert. in ms.} \textbf{24 groom} \textit{angl. in ms.}}} \\
\end{tabular}}%
\par\vspace*{\fill}
% --- Registres de bas de page (1. Critique interne [si pertinent] ; 2. Loci paralleli en latin ; 3. Commentaire historique en retrait) ---
\noindent\rule{\textwidth}{0.5pt}\par\vspace{0.65mm}
{\fontsize{6.6pt}{7.7pt}\selectfont%
\noindent\normalfont\color{black!92} \textbf{2} (\textit{Marquay}) cf. supra f\textsuperscript{o}~14~r\textsuperscript{o} (p.~\pageref{fol:28}), f\textsuperscript{o}~19~r\textsuperscript{o} (p.~\pageref{fol:38}), \textit{et passim}. — \textbf{6} (\textit{Ligny}) cf. supra f\textsuperscript{o}~16~v\textsuperscript{o} (p.~\pageref{fol:33}). — \textbf{19} (\textit{Cap. West}) cf. infra f\textsuperscript{o}~21~r\textsuperscript{o} (p.~\pageref{fol:42}), f\textsuperscript{o}~24~r\textsuperscript{o} (p.~\pageref{fol:48}).\par\vspace{0.55mm}
\noindent\rule{38mm}{0.4pt}\par\vspace{0.5mm}
\noindent\hspace*{5mm}\parbox{\dimexpr\textwidth-5mm\relax}{\fontsize{6.6pt}{7.7pt}\selectfont\color{black!84}Le samedi 9 septembre 1916 débute le mouvement par étapes de la \textbf{18\textsuperscript{e} Division britannique} depuis le Ternois vers le front de l'Ancre: après le défilé des trois brigades d'infanterie de la division (53\textsuperscript{e}, 54\textsuperscript{e} et 55\textsuperscript{e} Brigades) et le passage devant l'école de mortiers de tranchée (\textit{trench mortar school}) sur le plateau de Ligny-Saint-Flochel, la compagnie fait étape à \textbf{Monts-en-Ternois} (Pas-de-Calais, à 6~km au sud de Marquay), où le \textbf{capitaine West} (officier de la \textit{152nd Company A.S.C.}) invite Hippolyte Combes à dîner au mess des officiers.}\par}
\newpage
\phantomsection\label{fac:39}
% --- En-tête du fac-similé (Page impaire / droite : Photographie du manuscrit) ---
\noindent\rule{\textwidth}{0.5pt}\par\vspace{0.55mm}
\noindent\textbf{\fontsize{6.8pt}{7.8pt}\selectfont FAC-SIMILÉ EN VIS-À-VIS — FOLIO 19 VERSO} \hfill \textbf{\fontsize{6.8pt}{7.8pt}\selectfont SEPTEMBRE 1916}\par\vspace{0.25mm}
\noindent{\fontsize{6.8pt}{8.0pt}\selectfont\textit{Livre-Journal d'Hippolyte Combes} — Samedi 9 septembre 1916 (De Marquay et Ligny-Saint-Flochel à Monts-en-Ternois (Pas-de-Calais))}\par\vspace{0.55mm}
\noindent\rule{\textwidth}{0.35pt}\par\vspace*{\fill}
{\centering
  \setlength{\fboxsep}{1.2mm}%
  \setlength{\fboxrule}{0.4pt}%
  \fcolorbox{gray!55}{white}{%
    \includegraphics[width=0.95\textwidth,height=0.81\textheight,keepaspectratio]{\facsimiledir/20251126_094545.jpg}%
  }\par}
\vspace*{\fill}
\newpage

% =============================================================================
% VIS-A-VIS SPREAD : FOLIO 20 RECTO (20251126_094600.jpg)
% =============================================================================
\cleardoubleevenpage
\renewcommand{\currentmonthyear}{Septembre 1916}
\renewcommand{\currentfoliolabel}{Folio 20 recto}
\phantomsection\label{fol:40}
% --- En-tête du folio (Page paire / gauche : Transcription, apparat & registres de bas de page) ---
\noindent\rule{\textwidth}{0.5pt}\par\vspace{0.55mm}
\noindent\textbf{\footnotesize FOLIO 20 RECTO} \hfill \textbf{\footnotesize SEPTEMBRE 1916}\par\vspace{0.25mm}
\noindent{\fontsize{6.9pt}{8.0pt}\selectfont Nuit devant le fournil à Monts-en-Ternois et étape au hameau d'Oppy (Beaudricourt)}\par\vspace{0.2mm}
\noindent{\fontsize{6.9pt}{8.0pt}\selectfont\textit{Dimanche 10 septembre 1916} \hfill \textit{De Monts-en-Ternois à Oppy et Beaudricourt (Pas-de-Calais)}}\par\vspace{0.5mm}
\noindent\rule{\textwidth}{0.35pt}\par\vspace{0.8mm}
% --- Corps de la transcription diplomatique et apparat marginal en prolongement ---
{\renewcommand{\arraystretch}{1.00}%
\noindent\hspace*{3.5mm}\begin{tabular}{@{}p{87mm}@{\hspace{3.5mm}}p{39mm}@{}}
      \leavevmode\fontsize{7.7pt}{9.0pt}\selectfont\strut \underline{Dim. 10.} \hfill{} \underline{Septbre 16}\strut & \smash{\parbox[t]{39mm}{\fontsize{6.2pt}{7.0pt}\selectfont\raggedright\color{black!88}\textbf{1 Dim.} \textit{scil.} Dimanche.}} \\
      \leavevmode\fontsize{7.7pt}{9.0pt}\selectfont\strut Très bonne nuit couché s/ 2 bancs de\strut & \smash{\parbox[t]{39mm}{\fontsize{6.2pt}{7.0pt}\selectfont\raggedright\color{black!88}\textbf{2 s/} \textit{abbr. pro} sur.}} \\
      \leavevmode\fontsize{7.7pt}{9.0pt}\selectfont\strut jardin rapprochés – devant le four –\strut &  \\
      \leavevmode\fontsize{7.7pt}{9.0pt}\selectfont\strut Temps d'abord nuageux\strut &  \\
      \leavevmode\llap{\fontsize{6.4pt}{7.3pt}\selectfont\color{darkgray}5\hspace{1.8mm}}\fontsize{7.7pt}{9.0pt}\selectfont\strut Partis à 11H m\strut & \smash{\parbox[t]{39mm}{\fontsize{6.2pt}{7.0pt}\selectfont\raggedright\color{black!88}\textbf{5 m.} \textit{abbr. pro} matin.}} \\
      \leavevmode\fontsize{7.7pt}{9.0pt}\selectfont\strut Voyage très agréable. – paysages très\strut &  \\
      \leavevmode\fontsize{7.7pt}{9.0pt}\selectfont\strut variés – d'abord plaines. puis coteaux\strut &  \\
      \leavevmode\fontsize{7.7pt}{9.0pt}\selectfont\strut et peu à peu hauteurs boisées.\strut &  \\
      \leavevmode\fontsize{7.7pt}{9.0pt}\selectfont\strut A 1 H [s. ?], arrivés à Oppy\strut & \smash{\parbox[t]{39mm}{\fontsize{6.2pt}{7.0pt}\selectfont\raggedright\color{black!88}\textbf{9} \textit{lect. incert. in ms.}}} \\
      \leavevmode\llap{\fontsize{6.4pt}{7.3pt}\selectfont\color{darkgray}10\hspace{1.8mm}}\fontsize{7.7pt}{9.0pt}\selectfont\strut Campagne verdoyante. régions fertiles.\strut &  \\
      \leavevmode\fontsize{7.7pt}{9.0pt}\selectfont\strut Oppy, situé tt en haut coteau. De ce\strut & \smash{\parbox[t]{39mm}{\fontsize{6.2pt}{7.0pt}\selectfont\raggedright\color{black!88}\textbf{11 tt} \textit{abbr. pro} tout.}} \\
      \leavevmode\fontsize{7.7pt}{9.0pt}\selectfont\strut plateau, très b. vue. Ptit hameau\strut & \smash{\parbox[t]{39mm}{\fontsize{6.2pt}{7.0pt}\selectfont\raggedright\color{black!88}\textbf{12 Ptit} \textit{abbr. pro} Petit.}} \\
      \leavevmode\fontsize{7.7pt}{9.0pt}\selectfont\strut de 6 maisons.\strut &  \\
      \leavevmode\fontsize{7.7pt}{9.0pt}\selectfont\strut Obligé coucher grenier ferme où mess.\strut & \smash{\parbox[t]{39mm}{\fontsize{6.2pt}{7.0pt}\selectfont\raggedright\color{black!88}\textbf{14 mess} \textit{angl. in ms.}}} \\
      \leavevmode\llap{\fontsize{6.4pt}{7.3pt}\selectfont\color{darkgray}15\hspace{1.8mm}}\fontsize{7.7pt}{9.0pt}\selectfont\strut Fin après-midi, venu à pied à\strut &  \\
      \leavevmode\fontsize{7.7pt}{9.0pt}\selectfont\strut Beaudricourt, on duty, chz le maire\strut &  \\
      \leavevmode\fontsize{7.7pt}{9.0pt}\selectfont\strut Tout jour, mangé au mess off\strut & \smash{\parbox[t]{39mm}{\fontsize{6.2pt}{7.0pt}\selectfont\raggedright\color{black!88}\textbf{17 off} \textit{abbr. pro} officier(s). \textbf{17 mess} \textit{angl. in ms.}}} \\
      \leavevmode\fontsize{7.7pt}{9.0pt}\selectfont\strut Couché tôt\strut &  \\
\end{tabular}}%
\par\vspace*{\fill}
% --- Registres de bas de page (1. Critique interne [si pertinent] ; 2. Loci paralleli en latin ; 3. Commentaire historique en retrait) ---
\noindent\rule{\textwidth}{0.5pt}\par\vspace{0.65mm}
{\fontsize{6.6pt}{7.7pt}\selectfont%
\noindent\normalfont\color{black!92} \textbf{9} (\textit{Oppy}) cf. infra f\textsuperscript{o}~20~v\textsuperscript{o} (p.~\pageref{fol:41}). — \textbf{16} (\textit{Beaudricourt}) cf. infra f\textsuperscript{o}~20~v\textsuperscript{o} (p.~\pageref{fol:41}).\par\vspace{0.55mm}
\noindent\rule{38mm}{0.4pt}\par\vspace{0.5mm}
\noindent\hspace*{5mm}\parbox{\dimexpr\textwidth-5mm\relax}{\fontsize{6.6pt}{7.7pt}\selectfont\color{black!84}Le hameau d'\textbf{Oppy} où cantonne la compagnie le 9 septembre 1916 est l'écart d'Oppy rattaché à la commune de \textbf{Beaudricourt} (Pas-de-Calais, canton d'Avesnes-le-Comte, à l'ouest d'Arras), où Hippolyte Combes se rend à pied le 10 septembre auprès du maire pour le service de cantonnement.}\par}
\newpage
\phantomsection\label{fac:40}
% --- En-tête du fac-similé (Page impaire / droite : Photographie du manuscrit) ---
\noindent\rule{\textwidth}{0.5pt}\par\vspace{0.55mm}
\noindent\textbf{\fontsize{6.8pt}{7.8pt}\selectfont FAC-SIMILÉ EN VIS-À-VIS — FOLIO 20 RECTO} \hfill \textbf{\fontsize{6.8pt}{7.8pt}\selectfont SEPTEMBRE 1916}\par\vspace{0.25mm}
\noindent{\fontsize{6.8pt}{8.0pt}\selectfont\textit{Livre-Journal d'Hippolyte Combes} — Dimanche 10 septembre 1916 (De Monts-en-Ternois à Oppy et Beaudricourt (Pas-de-Calais))}\par\vspace{0.55mm}
\noindent\rule{\textwidth}{0.35pt}\par\vspace*{\fill}
{\centering
  \setlength{\fboxsep}{1.2mm}%
  \setlength{\fboxrule}{0.4pt}%
  \fcolorbox{gray!55}{white}{%
    \includegraphics[width=0.95\textwidth,height=0.81\textheight,keepaspectratio]{\facsimiledir/20251126_094600.jpg}%
  }\par}
\vspace*{\fill}
\newpage

% =============================================================================
% VIS-A-VIS SPREAD : FOLIO 20 VERSO (20251126_094621.jpg)
% =============================================================================
\cleardoubleevenpage
\renewcommand{\currentmonthyear}{Septembre 1916}
\renewcommand{\currentfoliolabel}{Folio 20 verso}
\phantomsection\label{fol:41}
% --- En-tête du folio (Page paire / gauche : Transcription, apparat & registres de bas de page) ---
\noindent\rule{\textwidth}{0.5pt}\par\vspace{0.55mm}
\noindent\textbf{\footnotesize FOLIO 20 VERSO} \hfill \textbf{\footnotesize SEPTEMBRE 1916}\par\vspace{0.25mm}
\noindent{\fontsize{6.9pt}{8.0pt}\selectfont Longue chevauchée d'Oppy à Raincheval par Lucheux (l'Édit des postes de Louis XI) et Thièvres}\par\vspace{0.2mm}
\noindent{\fontsize{6.9pt}{8.0pt}\selectfont\textit{Lundi 11 septembre 1916} \hfill \textit{D'Oppy (Beaudricourt) à Lucheux, Thièvres et Raincheval (Pas-de-Calais / Somme)}}\par\vspace{0.5mm}
\noindent\rule{\textwidth}{0.35pt}\par\vspace{0.8mm}
% --- Corps de la transcription diplomatique et apparat marginal en prolongement ---
{\renewcommand{\arraystretch}{0.87}%
\noindent\hspace*{3.5mm}\begin{tabular}{@{}p{87mm}@{\hspace{3.5mm}}p{39mm}@{}}
      \leavevmode\fontsize{6.9pt}{7.9pt}\selectfont\strut \underline{Lundi 11} — Sur pied, dès 5 H — \hfill{} \raisebox{0.5ex}{\underline{Septbre /16}}\strut &  \\
      \leavevmode\fontsize{6.9pt}{7.9pt}\selectfont\strut dérangé sur matin\strut &  \\
      \leavevmode\fontsize{6.9pt}{7.9pt}\selectfont\strut Quitté Oppy 7 H. Ciel couvert\strut &  \\
      \leavevmode\fontsize{6.9pt}{7.9pt}\selectfont\strut Chevauché matinée entière — Après 3 H soir,\strut &  \\
      \leavevmode\llap{\fontsize{6.4pt}{7.3pt}\selectfont\color{darkgray}5\hspace{1.8mm}}\fontsize{6.9pt}{7.9pt}\selectfont\strut seulement, arrivés à Raincheval\strut &  \\
      \leavevmode\fontsize{6.9pt}{7.9pt}\selectfont\strut Aujd’hui souffert poussière.\strut &  \\
      \leavevmode\fontsize{6.9pt}{7.9pt}\selectfont\strut A 11H m. soleil perce nuages – et\strut & \smash{\parbox[t]{39mm}{\fontsize{6.2pt}{7.0pt}\selectfont\raggedright\color{black!88}\textbf{7 m.} \textit{abbr. pro} matin.}} \\
      \leavevmode\fontsize{6.9pt}{7.9pt}\selectfont\strut très ardent\strut &  \\
      \leavevmode\fontsize{6.9pt}{7.9pt}\selectfont\strut Traversé régions pittoresques : coteaux\strut &  \\
      \leavevmode\llap{\fontsize{6.4pt}{7.3pt}\selectfont\color{darkgray}10\hspace{1.8mm}}\fontsize{6.9pt}{7.9pt}\selectfont\strut boisés, petite vallée verdoyante,\strut &  \\
      \leavevmode\fontsize{6.9pt}{7.9pt}\selectfont\strut Arrêt à Lucheux : Edit de Louis XI. 1464\strut &  \\
      \leavevmode\fontsize{6.9pt}{7.9pt}\selectfont\strut des postes\strut &  \\
      \leavevmode\fontsize{6.9pt}{7.9pt}\selectfont\strut Raincheval, très gros village aux\strut &  \\
      \leavevmode\fontsize{6.9pt}{7.9pt}\selectfont\strut maisons mal entretenues – fourmillant\strut &  \\
      \leavevmode\llap{\fontsize{6.4pt}{7.3pt}\selectfont\color{darkgray}15\hspace{1.8mm}}\fontsize{6.9pt}{7.9pt}\selectfont\strut de tommies. Très nombreuses troupes y\strut & \smash{\parbox[t]{39mm}{\fontsize{6.2pt}{7.0pt}\selectfont\raggedright\color{black!88}\textbf{15 tommies} \textit{angl. in ms.}}} \\
      \leavevmode\fontsize{6.9pt}{7.9pt}\selectfont\strut cantonnées.\strut &  \\
      \leavevmode\fontsize{6.9pt}{7.9pt}\selectfont\strut Trouvé chambrette famille de pauvres\strut &  \\
      \leavevmode\fontsize{6.9pt}{7.9pt}\selectfont\strut mais braves gens. Fille aînée, prête à accoucher\strut &  \\
      \leavevmode\fontsize{6.9pt}{7.9pt}\selectfont\strut mangé tjs au mess off\strut & \smash{\parbox[t]{39mm}{\fontsize{6.2pt}{7.0pt}\selectfont\raggedright\color{black!88}\textbf{19 tjs} \textit{abbr. pro} toujours. \textbf{19 off} \textit{abbr. pro} officier(s). — \textbf{20} \textit{in marg. cum unco} (comptes). \textbf{20 groom} \textit{angl. in ms.} — \textbf{21 c.} \textit{abbr. pro} centimes. — \textbf{22 cartes ples} \textit{abbr. pro} cartes postales.}} \\
      \cline{1-1}
      \leavevmode\llap{\fontsize{6.4pt}{7.3pt}\selectfont\color{darkgray}20\hspace{1.8mm}}\fontsize{6.9pt}{7.9pt}\selectfont\strut [ — Tip au groom Connelly : 1 f –\strut &  \\
      \leavevmode\fontsize{6.9pt}{7.9pt}\selectfont\strut Oeufs : 3 à Thièvres : 50 c. – 4 à Raincheval : 70 c\strut &  \\
      \leavevmode\fontsize{6.9pt}{7.9pt}\selectfont\strut 6 cartes ples de Raincheval : 30 c\strut &  \\
      \leavevmode\fontsize{6.9pt}{7.9pt}\selectfont\strut 1 litre bière : 50 + 10 c –\strut &  \\
\end{tabular}}%
\par\vspace*{\fill}
% --- Registres de bas de page (1. Critique interne [si pertinent] ; 2. Loci paralleli en latin ; 3. Commentaire historique en retrait) ---
\noindent\rule{\textwidth}{0.5pt}\par\vspace{0.65mm}
{\fontsize{6.6pt}{7.7pt}\selectfont%
\noindent\normalfont\color{black!92} \textbf{3} (\textit{Oppy}) cf. supra f\textsuperscript{o}~20~r\textsuperscript{o} (p.~\pageref{fol:40}). — \textbf{5} (\textit{Raincheval}) cf. infra f\textsuperscript{o}~21~r\textsuperscript{o} (p.~\pageref{fol:42}), f\textsuperscript{o}~21~v\textsuperscript{o} (p.~\pageref{fol:43}), \textit{et passim}. — \textbf{20} (\textit{Connelly}) cf. supra f\textsuperscript{o}~6~r\textsuperscript{o} (p.~\pageref{fol:12}), f\textsuperscript{o}~19~v\textsuperscript{o} (p.~\pageref{fol:39}), infra f\textsuperscript{o}~22~v\textsuperscript{o} (p.~\pageref{fol:45}), f\textsuperscript{o}~37~r\textsuperscript{o} (p.~\pageref{fol:74}).\par\vspace{0.55mm}
\noindent\rule{38mm}{0.4pt}\par\vspace{0.5mm}
\noindent\hspace*{5mm}\parbox{\dimexpr\textwidth-5mm\relax}{\fontsize{6.6pt}{7.7pt}\selectfont\color{black!84}Longue étape à cheval conduisant la colonne du Pas-de-Calais dans la Somme par \textbf{Lucheux}, \textbf{Thièvres} et \textbf{Raincheval}. À \textbf{Lucheux} (Somme), Hippolyte Combes remarque les ruines du château fort des comtes de Saint-Pol (XII\textsuperscript{e}–XIV\textsuperscript{e}~s. ; Base Mérimée, \texttt{PA00116194}) et le beffroi abritant le chêne historique sous lequel la tradition locale situe la signature par \textbf{Louis~XI} de l'édit du 19 juin 1464 instituant la poste royale.}\par}
\newpage
\phantomsection\label{fac:41}
% --- En-tête du fac-similé (Page impaire / droite : Photographie du manuscrit) ---
\noindent\rule{\textwidth}{0.5pt}\par\vspace{0.55mm}
\noindent\textbf{\fontsize{6.8pt}{7.8pt}\selectfont FAC-SIMILÉ EN VIS-À-VIS — FOLIO 20 VERSO} \hfill \textbf{\fontsize{6.8pt}{7.8pt}\selectfont SEPTEMBRE 1916}\par\vspace{0.25mm}
\noindent{\fontsize{6.8pt}{8.0pt}\selectfont\textit{Livre-Journal d'Hippolyte Combes} — Lundi 11 septembre 1916 (D'Oppy (Beaudricourt) à Lucheux, Thièvres et Raincheval (Pas-de-Calais / Somme))}\par\vspace{0.55mm}
\noindent\rule{\textwidth}{0.35pt}\par\vspace*{\fill}
{\centering
  \setlength{\fboxsep}{1.2mm}%
  \setlength{\fboxrule}{0.4pt}%
  \fcolorbox{gray!55}{white}{%
    \includegraphics[width=0.95\textwidth,height=0.81\textheight,keepaspectratio]{\facsimiledir/20251126_094621.jpg}%
  }\par}
\vspace*{\fill}
\newpage

% =============================================================================
% VIS-A-VIS SPREAD : FOLIO 21 RECTO (20251126_094639.jpg)
% =============================================================================
\cleardoubleevenpage
\renewcommand{\currentmonthyear}{Septembre 1916}
\renewcommand{\currentfoliolabel}{Folio 21 recto}
\phantomsection\label{fol:42}
% --- En-tête du folio (Page paire / gauche : Transcription, apparat & registres de bas de page) ---
\noindent\rule{\textwidth}{0.5pt}\par\vspace{0.55mm}
\noindent\textbf{\footnotesize FOLIO 21 RECTO} \hfill \textbf{\footnotesize SEPTEMBRE 1916}\par\vspace{0.25mm}
\noindent{\fontsize{6.9pt}{8.0pt}\selectfont Admission permanente au mess des officiers du capitaine West et promenade en forêt de Raincheval}\par\vspace{0.2mm}
\noindent{\fontsize{6.9pt}{8.0pt}\selectfont\textit{Mardi 12 et mercredi 13 septembre 1916} \hfill \textit{Raincheval (Somme)}}\par\vspace{0.5mm}
\noindent\rule{\textwidth}{0.35pt}\par\vspace{0.8mm}
% --- Corps de la transcription diplomatique et apparat marginal en prolongement ---
{\renewcommand{\arraystretch}{0.87}%
\noindent\hspace*{3.5mm}\begin{tabular}{@{}p{87mm}@{\hspace{3.5mm}}p{39mm}@{}}
      \leavevmode\fontsize{6.9pt}{7.9pt}\selectfont\strut 12 mardi. \hfill{} \underline{Sept / 16}\strut &  \\
      \leavevmode\fontsize{6.9pt}{7.9pt}\selectfont\strut Après excte nuit, les fatigues oubliées.\strut & \smash{\parbox[t]{39mm}{\fontsize{6.2pt}{7.0pt}\selectfont\raggedright\color{black!88}\textbf{2 excte} \textit{abbr. pro} excellente. — \textbf{3 ds} \textit{abbr. pro} dans.}} \\
      \leavevmode\fontsize{6.9pt}{7.9pt}\selectfont\strut Déjeuné midi ds ma chambre\strut &  \\
      \leavevmode\fontsize{6.9pt}{7.9pt}\selectfont\strut Après-midi, Cap. West, offert de venir à\strut & \smash{\parbox[t]{39mm}{\fontsize{6.2pt}{7.0pt}\selectfont\raggedright\color{black!88}\textbf{4 Cap.} \textit{abbr. pro} Capitaine. — \textbf{5 mess} \textit{angl. in ms.}}} \\
      \leavevmode\llap{\fontsize{6.4pt}{7.3pt}\selectfont\color{darkgray}5\hspace{1.8mm}}\fontsize{6.9pt}{7.9pt}\selectfont\strut leur mess d'1 manière permanente — et,\strut &  \\
      \leavevmode\fontsize{6.9pt}{7.9pt}\selectfont\strut devant son insistance, [ai accepté / dû accepter ?]\strut & \smash{\parbox[t]{39mm}{\fontsize{6.2pt}{7.0pt}\selectfont\raggedright\color{black!88}\textbf{6} \textit{lect. incert. in ms.}}} \\
      \leavevmode\fontsize{6.9pt}{7.9pt}\selectfont\strut Dès le soir, membre du mess des off à\strut & \smash{\parbox[t]{39mm}{\fontsize{6.2pt}{7.0pt}\selectfont\raggedright\color{black!88}\textbf{7 off} \textit{abbr. pro} officier(s). \textbf{7 mess} \textit{angl. in ms.}}} \\
      \leavevmode\fontsize{6.9pt}{7.9pt}\selectfont\strut titre définitif\strut &  \\
      \leavevmode\fontsize{6.9pt}{7.9pt}\selectfont\strut Cuisine quelconque — mais grands avantages.\strut &  \\
      \leavevmode\llap{\fontsize{6.4pt}{7.3pt}\selectfont\color{darkgray}10\hspace{1.8mm}}\fontsize{6.9pt}{7.9pt}\selectfont\strut converser anglais \& sorte de club où me\strut &  \\
      \leavevmode\fontsize{6.9pt}{7.9pt}\selectfont\strut retirer pour lire revues \& journaux.\strut &  \\
      \leavevmode\fontsize{6.9pt}{7.9pt}\selectfont\strut – Temps menaçant – (Matin \& après-midi\strut &  \\
      \leavevmode\fontsize{6.9pt}{7.9pt}\selectfont\strut sur les hauteurs for\strut &  \\
      \leavevmode\fontsize{6.9pt}{7.9pt}\selectfont\strut green forage\strut & \smash{\parbox[t]{39mm}{\fontsize{6.2pt}{7.0pt}\selectfont\raggedright\color{black!88}\textbf{14 green} \textit{angl. in ms.}}} \\
      \cline{1-1}
      \leavevmode\llap{\fontsize{6.4pt}{7.3pt}\selectfont\color{darkgray}15\hspace{1.8mm}}\fontsize{6.9pt}{7.9pt}\selectfont\strut \underline{13 Me} Temps couvert après-midi,\strut &  \\
      \leavevmode\fontsize{6.9pt}{7.9pt}\selectfont\strut averses de courte durée\strut &  \\
      \leavevmode\fontsize{6.9pt}{7.9pt}\selectfont\strut A g. du chateau, forêt magnifique\strut &  \\
      \leavevmode\fontsize{6.9pt}{7.9pt}\selectfont\strut Venu m'y promener après le lunch\strut & \smash{\parbox[t]{39mm}{\fontsize{6.2pt}{7.0pt}\selectfont\raggedright\color{black!88}\textbf{18 lunch} \textit{angl. in ms.}}} \\
      \leavevmode\fontsize{6.9pt}{7.9pt}\selectfont\strut Arbres séculaires – routes \& sentiers à\strut &  \\
      \leavevmode\llap{\fontsize{6.4pt}{7.3pt}\selectfont\color{darkgray}20\hspace{1.8mm}}\fontsize{6.9pt}{7.9pt}\selectfont\strut angle droit. Echappées de ciel à l'extré\strut &  \\
      \leavevmode\fontsize{6.9pt}{7.9pt}\selectfont\strut d'allées profondes – Bruissement des feuilles\strut &  \\
      \leavevmode\fontsize{6.9pt}{7.9pt}\selectfont\strut Remonté plateau pour green forage\strut & \smash{\parbox[t]{39mm}{\fontsize{6.2pt}{7.0pt}\selectfont\raggedright\color{black!88}\textbf{22 green} \textit{angl. in ms.}}} \\
      \leavevmode\fontsize{6.9pt}{7.9pt}\selectfont\strut Tout le jour, pr 1ère fois, mess des offs\strut & \smash{\parbox[t]{39mm}{\fontsize{6.2pt}{7.0pt}\selectfont\raggedright\color{black!88}\textbf{23 pr} \textit{abbr. pro} pour. \textbf{23 offs} \textit{abbr. pro} officiers.}} \\
\end{tabular}}%
\par\vspace*{\fill}
% --- Registres de bas de page (1. Critique interne [si pertinent] ; 2. Loci paralleli en latin ; 3. Commentaire historique en retrait) ---
\noindent\rule{\textwidth}{0.5pt}\par\vspace{0.65mm}
{\fontsize{6.6pt}{7.7pt}\selectfont%
\noindent\normalfont\color{black!92} \textbf{4} (\textit{Cap. West}) cf. supra f\textsuperscript{o}~19~v\textsuperscript{o} (p.~\pageref{fol:39}), infra f\textsuperscript{o}~24~r\textsuperscript{o} (p.~\pageref{fol:48}). — \textbf{14} (\textit{green forage}) cf. supra f\textsuperscript{o}~3~r\textsuperscript{o} (p.~\pageref{fol:6}), f\textsuperscript{o}~19~r\textsuperscript{o} (p.~\pageref{fol:38}), infra f\textsuperscript{o}~22~r\textsuperscript{o} (p.~\pageref{fol:44}), f\textsuperscript{o}~24~v\textsuperscript{o} (p.~\pageref{fol:49}), \textit{et passim}.\par\vspace{0.55mm}
\noindent\rule{38mm}{0.4pt}\par\vspace{0.5mm}
\noindent\hspace*{5mm}\parbox{\dimexpr\textwidth-5mm\relax}{\fontsize{6.6pt}{7.7pt}\selectfont\color{black!84}À \textbf{Raincheval} (Somme), le mardi 12 septembre 1916, le \textbf{capitaine West} (de la \textit{152nd Company A.S.C.}) admet définitivement l'interprète français à la table du mess des officiers britanniques. Le lendemain 13 septembre, entre deux corvées de fourrage vert sur le plateau, Hippolyte Combes parcourt les allées du grand bois jouxtant le château de Raincheval.}\par}
\newpage
\phantomsection\label{fac:42}
% --- En-tête du fac-similé (Page impaire / droite : Photographie du manuscrit) ---
\noindent\rule{\textwidth}{0.5pt}\par\vspace{0.55mm}
\noindent\textbf{\fontsize{6.8pt}{7.8pt}\selectfont FAC-SIMILÉ EN VIS-À-VIS — FOLIO 21 RECTO} \hfill \textbf{\fontsize{6.8pt}{7.8pt}\selectfont SEPTEMBRE 1916}\par\vspace{0.25mm}
\noindent{\fontsize{6.8pt}{8.0pt}\selectfont\textit{Livre-Journal d'Hippolyte Combes} — Mardi 12 et mercredi 13 septembre 1916 (Raincheval (Somme))}\par\vspace{0.55mm}
\noindent\rule{\textwidth}{0.35pt}\par\vspace*{\fill}
{\centering
  \setlength{\fboxsep}{1.2mm}%
  \setlength{\fboxrule}{0.4pt}%
  \fcolorbox{gray!55}{white}{%
    \includegraphics[width=0.95\textwidth,height=0.81\textheight,keepaspectratio]{\facsimiledir/20251126_094639.jpg}%
  }\par}
\vspace*{\fill}
\newpage

% =============================================================================
% VIS-A-VIS SPREAD : FOLIO 21 VERSO (20251126_094654.jpg)
% =============================================================================
\cleardoubleevenpage
\renewcommand{\currentmonthyear}{Septembre 1916}
\renewcommand{\currentfoliolabel}{Folio 21 verso}
\phantomsection\label{fol:43}
% --- En-tête du folio (Page paire / gauche : Transcription, apparat & registres de bas de page) ---
\noindent\rule{\textwidth}{0.5pt}\par\vspace{0.55mm}
\noindent\textbf{\footnotesize FOLIO 21 VERSO} \hfill \textbf{\footnotesize SEPTEMBRE 1916}\par\vspace{0.25mm}
\noindent{\fontsize{6.9pt}{8.0pt}\selectfont Visite à la Mission française au parc du château d'Acheux et exercices de grenades à Raincheval}\par\vspace{0.2mm}
\noindent{\fontsize{6.9pt}{8.0pt}\selectfont\textit{Jeudi 14 et vendredi 15 septembre 1916} \hfill \textit{Raincheval et Acheux-en-Amiénois (Somme)}}\par\vspace{0.5mm}
\noindent\rule{\textwidth}{0.35pt}\par\vspace{0.8mm}
% --- Corps de la transcription diplomatique et apparat marginal en prolongement ---
{\renewcommand{\arraystretch}{0.93}%
\noindent\hspace*{3.5mm}\begin{tabular}{@{}p{87mm}@{\hspace{3.5mm}}p{39mm}@{}}
      \leavevmode\fontsize{7.3pt}{8.4pt}\selectfont\strut \underline{14 J.} matin, ciel bleu \& soleil \hfill{} \raisebox{0.5ex}{\underline{Sept 16}}\strut & \smash{\parbox[t]{39mm}{\fontsize{6.2pt}{7.0pt}\selectfont\raggedright\color{black!88}\textbf{1 J.} \textit{scil.} Jeudi.}} \\
      \leavevmode\fontsize{7.3pt}{8.4pt}\selectfont\strut radieux.\strut &  \\
      \leavevmode\fontsize{7.3pt}{8.4pt}\selectfont\strut Après le lunch, course à cheval. Venu\strut & \smash{\parbox[t]{39mm}{\fontsize{6.2pt}{7.0pt}\selectfont\raggedright\color{black!88}\textbf{3 lunch} \textit{angl. in ms.}}} \\
      \leavevmode\fontsize{7.3pt}{8.4pt}\selectfont\strut à Acheux; bureaux mission fr; dans\strut &  \\
      \leavevmode\llap{\fontsize{6.4pt}{7.3pt}\selectfont\color{darkgray}5\hspace{1.8mm}}\fontsize{7.3pt}{8.4pt}\selectfont\strut les parcs du chateau – maisonnette toile\strut &  \\
      \leavevmode\fontsize{7.3pt}{8.4pt}\selectfont\strut \& planches\strut &  \\
      \leavevmode\fontsize{7.3pt}{8.4pt}\selectfont\strut Fait route avec Lereux [?]\strut & \smash{\parbox[t]{39mm}{\fontsize{6.2pt}{7.0pt}\selectfont\raggedright\color{black!88}\textbf{7 Lereux} \textit{lect. incert. in ms.}}} \\
      \leavevmode\fontsize{7.3pt}{8.4pt}\selectfont\strut Bourg très coquet\strut &  \\
      \leavevmode\fontsize{7.3pt}{8.4pt}\selectfont\strut Fin jour, 1er froid d'automne\strut & \smash{\parbox[t]{39mm}{\fontsize{6.2pt}{7.0pt}\selectfont\raggedright\color{black!88}\textbf{9 1er} \textit{scil.} premier.}} \\
      \leavevmode\llap{\fontsize{6.4pt}{7.3pt}\selectfont\color{darkgray}10\hspace{1.8mm}}\fontsize{7.3pt}{8.4pt}\selectfont\strut [ — Ctes ples, à Acheux : 30 c.\strut & \smash{\parbox[t]{39mm}{\fontsize{6.2pt}{7.0pt}\selectfont\raggedright\color{black!88}\textbf{10} \textit{in marg. cum unco} (comptes). \textbf{10 Ctes ples} \textit{abbr. pro} Cartes postales. — \textbf{11 batman} \textit{angl. in ms.} — \textbf{12 extt} \textit{abbr. pro} extrêmement. — \textbf{14 commt} \textit{abbr. pro} commencement. — \textbf{16 canteen} \textit{angl. in ms.}}} \\
      \leavevmode\fontsize{7.3pt}{8.4pt}\selectfont\strut To a batman for saddling up my poney : 50c\strut &  \\
      \cline{1-1}
      \leavevmode\fontsize{7.3pt}{8.4pt}\selectfont\strut \underline{15 Ve} – nuit froide – matinée extt fraiche\strut &  \\
      \leavevmode\fontsize{7.3pt}{8.4pt}\selectfont\strut très b. journée\strut &  \\
      \leavevmode\fontsize{7.3pt}{8.4pt}\selectfont\strut Impressions de fin été \& commt automne\strut &  \\
      \leavevmode\llap{\fontsize{6.4pt}{7.3pt}\selectfont\color{darkgray}15\hspace{1.8mm}}\fontsize{7.3pt}{8.4pt}\selectfont\strut Après-midi, venu à travers champs à\strut &  \\
      \leavevmode\fontsize{7.3pt}{8.4pt}\selectfont\strut la canteen, station de Raincheval\strut &  \\
      \leavevmode\fontsize{7.3pt}{8.4pt}\selectfont\strut Exerc. de grenades \& mines de champs\strut &  \\
      \leavevmode\fontsize{7.3pt}{8.4pt}\selectfont\strut en face gare\strut &  \\
      \leavevmode\fontsize{7.3pt}{8.4pt}\selectfont\strut — Canteen station de Raincheval :\strut & \smash{\parbox[t]{39mm}{\fontsize{6.2pt}{7.0pt}\selectfont\raggedright\color{black!88}\textbf{19 Canteen} \textit{angl. in ms.}}} \\
      \leavevmode\llap{\fontsize{6.4pt}{7.3pt}\selectfont\color{darkgray}20\hspace{1.8mm}}\fontsize{7.3pt}{8.4pt}\selectfont\strut Canada in Flanders... 1f 25\strut &  \\
\end{tabular}}%
\par\vspace*{\fill}
% --- Registres de bas de page (1. Critique interne [si pertinent] ; 2. Loci paralleli en latin ; 3. Commentaire historique en retrait) ---
\noindent\rule{\textwidth}{0.5pt}\par\vspace{0.65mm}
{\fontsize{6.6pt}{7.7pt}\selectfont%
\noindent\normalfont\color{black!92} \textbf{4} (\textit{Acheux}) cf. infra f\textsuperscript{o}~22~v\textsuperscript{o} (p.~\pageref{fol:45}), f\textsuperscript{o}~23~r\textsuperscript{o} (p.~\pageref{fol:46}), \textit{et passim}. — \textbf{16} (\textit{Raincheval}) cf. supra f\textsuperscript{o}~20~v\textsuperscript{o} (p.~\pageref{fol:41}), f\textsuperscript{o}~21~r\textsuperscript{o} (p.~\pageref{fol:42}), infra f\textsuperscript{o}~22~r\textsuperscript{o} (p.~\pageref{fol:44}), f\textsuperscript{o}~22~v\textsuperscript{o} (p.~\pageref{fol:45}), \textit{et passim}. — \textbf{16} (\textit{canteen}) cf. infra f\textsuperscript{o}~33~v\textsuperscript{o} (p.~\pageref{fol:67}), f\textsuperscript{o}~59~r\textsuperscript{o} (p.~\pageref{fol:118}), \textit{et passim}.\par\vspace{0.55mm}
\noindent\rule{38mm}{0.4pt}\par\vspace{0.5mm}
\noindent\hspace*{5mm}\parbox{\dimexpr\textwidth-5mm\relax}{\fontsize{6.6pt}{7.7pt}\selectfont\color{black!84}Le 14 septembre, première visite à \textbf{Acheux-en-Amiénois} (Somme), où le bureau de la Mission française de la 18\textsuperscript{e} Division occupe une baraque dans le parc du château.}\par}
\newpage
\phantomsection\label{fac:43}
% --- En-tête du fac-similé (Page impaire / droite : Photographie du manuscrit) ---
\noindent\rule{\textwidth}{0.5pt}\par\vspace{0.55mm}
\noindent\textbf{\fontsize{6.8pt}{7.8pt}\selectfont FAC-SIMILÉ EN VIS-À-VIS — FOLIO 21 VERSO} \hfill \textbf{\fontsize{6.8pt}{7.8pt}\selectfont SEPTEMBRE 1916}\par\vspace{0.25mm}
\noindent{\fontsize{6.8pt}{8.0pt}\selectfont\textit{Livre-Journal d'Hippolyte Combes} — Jeudi 14 et vendredi 15 septembre 1916 (Raincheval et Acheux-en-Amiénois (Somme))}\par\vspace{0.55mm}
\noindent\rule{\textwidth}{0.35pt}\par\vspace*{\fill}
{\centering
  \setlength{\fboxsep}{1.2mm}%
  \setlength{\fboxrule}{0.4pt}%
  \fcolorbox{gray!55}{white}{%
    \includegraphics[width=0.95\textwidth,height=0.81\textheight,keepaspectratio]{\facsimiledir/20251126_094654.jpg}%
  }\par}
\vspace*{\fill}
\newpage

% =============================================================================
% VIS-A-VIS SPREAD : FOLIO 22 RECTO (20251126_094706.jpg)
% =============================================================================
\cleardoubleevenpage
\renewcommand{\currentmonthyear}{Septembre 1916}
\renewcommand{\currentfoliolabel}{Folio 22 recto}
\phantomsection\label{fol:44}
% --- En-tête du folio (Page paire / gauche : Transcription, apparat & registres de bas de page) ---
\noindent\rule{\textwidth}{0.5pt}\par\vspace{0.55mm}
\noindent\textbf{\footnotesize FOLIO 22 RECTO} \hfill \textbf{\footnotesize SEPTEMBRE 1916}\par\vspace{0.25mm}
\noindent{\fontsize{6.9pt}{8.0pt}\selectfont Les grands chênes du bois de Raincheval, le vieux collectionneur de faïences et l’office méthodiste (Wesleyans)}\par\vspace{0.2mm}
\noindent{\fontsize{6.9pt}{8.0pt}\selectfont\textit{Dimanche 17 septembre 1916} \hfill \textit{Raincheval (Somme)}}\par\vspace{0.5mm}
\noindent\rule{\textwidth}{0.35pt}\par\vspace{0.8mm}
% --- Corps de la transcription diplomatique et apparat marginal en prolongement ---
{\renewcommand{\arraystretch}{0.93}%
\noindent\hspace*{3.5mm}\begin{tabular}{@{}p{87mm}@{\hspace{3.5mm}}p{39mm}@{}}
      \leavevmode\fontsize{7.3pt}{8.4pt}\selectfont\strut \underline{17 Dim.} \hfill{} \underline{Septre 1916}\strut & \smash{\parbox[t]{39mm}{\fontsize{6.2pt}{7.0pt}\selectfont\raggedright\color{black!88}\textbf{1 Dim.} \textit{scil.} Dimanche.}} \\
      \leavevmode\fontsize{7.3pt}{8.4pt}\selectfont\strut Tjs ds l'attente notre départ\strut & \smash{\parbox[t]{39mm}{\fontsize{6.2pt}{7.0pt}\selectfont\raggedright\color{black!88}\textbf{2 ds} \textit{abbr. pro} dans. \textbf{2 Tjs} \textit{abbr. pro} Toujours. — \textbf{3 green} \textit{angl. in ms.}}} \\
      \leavevmode\fontsize{7.3pt}{8.4pt}\selectfont\strut Comme occupations : green forage. Dois\strut &  \\
      \leavevmode\fontsize{7.3pt}{8.4pt}\selectfont\strut conduire wagon champ de trèfle rouge.\strut &  \\
      \leavevmode\llap{\fontsize{6.4pt}{7.3pt}\selectfont\color{darkgray}5\hspace{1.8mm}}\fontsize{7.3pt}{8.4pt}\selectfont\strut Temps très beau – un peu froid, le soir\strut &  \\
      \leavevmode\fontsize{7.3pt}{8.4pt}\selectfont\strut \& surtout le matin. [Le/1/2 ?] jour, soleil brulant.\strut & \smash{\parbox[t]{39mm}{\fontsize{6.2pt}{7.0pt}\selectfont\raggedright\color{black!88}\textbf{6} \textit{lect. incert. in ms.}}} \\
      \leavevmode\fontsize{7.3pt}{8.4pt}\selectfont\strut Reviens pl. fois par jour me promener\strut & \smash{\parbox[t]{39mm}{\fontsize{6.2pt}{7.0pt}\selectfont\raggedright\color{black!88}\textbf{7 pl.} \textit{abbr. pro} plusieurs.}} \\
      \leavevmode\fontsize{7.3pt}{8.4pt}\selectfont\strut ds le grand bois qui s'étend à coté\strut & \smash{\parbox[t]{39mm}{\fontsize{6.2pt}{7.0pt}\selectfont\raggedright\color{black!88}\textbf{8 ds} \textit{abbr. pro} dans.}} \\
      \leavevmode\fontsize{7.3pt}{8.4pt}\selectfont\strut de nos horse-lines, derrière le mur du\strut & \smash{\parbox[t]{39mm}{\fontsize{6.2pt}{7.0pt}\selectfont\raggedright\color{black!88}\textbf{9 horse-lines} \textit{angl. in ms.}}} \\
      \leavevmode\llap{\fontsize{6.4pt}{7.3pt}\selectfont\color{darkgray}10\hspace{1.8mm}}\fontsize{7.3pt}{8.4pt}\selectfont\strut château. Ombre fraîche \& délicieuse.\strut &  \\
      \leavevmode\fontsize{7.3pt}{8.4pt}\selectfont\strut gds arbres séculaires – Chênes très droits\strut & \smash{\parbox[t]{39mm}{\fontsize{6.2pt}{7.0pt}\selectfont\raggedright\color{black!88}\textbf{11 gds} \textit{abbr. pro} grands.}} \\
      \leavevmode\fontsize{7.3pt}{8.4pt}\selectfont\strut aux troncs énormes\strut &  \\
      \leavevmode\fontsize{7.3pt}{8.4pt}\selectfont\strut Matin, rencontré intte [?] à la sortie église.\strut & \smash{\parbox[t]{39mm}{\fontsize{6.2pt}{7.0pt}\selectfont\raggedright\color{black!88}\textbf{13 intte} \textit{lect. incert. in ms.}}} \\
      \leavevmode\fontsize{7.3pt}{8.4pt}\selectfont\strut et avec lui venu chez le petit vieux, qui\strut &  \\
      \leavevmode\llap{\fontsize{6.4pt}{7.3pt}\selectfont\color{darkgray}15\hspace{1.8mm}}\fontsize{7.3pt}{8.4pt}\selectfont\strut possède des collections très belles de faïences,\strut &  \\
      \leavevmode\fontsize{7.3pt}{8.4pt}\selectfont\strut pierres, médailles etc etc\strut &  \\
      \leavevmode\fontsize{7.3pt}{8.4pt}\selectfont\strut Après-midi, assisté church service (Wesleyans)\strut &  \\
      \leavevmode\fontsize{7.3pt}{8.4pt}\selectfont\strut ds petite salle en face église village\strut & \smash{\parbox[t]{39mm}{\fontsize{6.2pt}{7.0pt}\selectfont\raggedright\color{black!88}\textbf{18 ds} \textit{abbr. pro} dans.}} \\
      \leavevmode\fontsize{7.3pt}{8.4pt}\selectfont\strut — Washing : 1 f. —\strut & \smash{\parbox[t]{39mm}{\fontsize{6.2pt}{7.0pt}\selectfont\raggedright\color{black!88}\textbf{19 f.} \textit{abbr. pro} franc(s). \textbf{19 Washing} \textit{angl. in ms.}}} \\
\end{tabular}}%
\par\vspace*{\fill}
% --- Registres de bas de page (1. Critique interne [si pertinent] ; 2. Loci paralleli en latin ; 3. Commentaire historique en retrait) ---
\noindent\rule{\textwidth}{0.5pt}\par\vspace{0.65mm}
{\fontsize{6.6pt}{7.7pt}\selectfont%
\noindent\normalfont\color{black!92} \textbf{3} (\textit{green forage}) cf. supra f\textsuperscript{o}~3~r\textsuperscript{o} (p.~\pageref{fol:6}), f\textsuperscript{o}~21~r\textsuperscript{o} (p.~\pageref{fol:42}), infra f\textsuperscript{o}~24~v\textsuperscript{o} (p.~\pageref{fol:49}), f\textsuperscript{o}~25~r\textsuperscript{o} (p.~\pageref{fol:50}), \textit{et passim}.\par\vspace{0.55mm}
\noindent\rule{38mm}{0.4pt}\par\vspace{0.5mm}
\noindent\hspace*{5mm}\parbox{\dimexpr\textwidth-5mm\relax}{\fontsize{6.6pt}{7.7pt}\selectfont\color{black!84}Dimanche d'attente à \textbf{Raincheval}: promenades sous les chênes séculaires bordant le parc du château, visite avec un collègue chez un collectionneur âgé du village (faïences, pierres et médailles), et présence l'après-midi à un office religieux méthodiste (\textit{Wesleyans}) célébré pour les soldats britanniques dans une salle faisant face à l'église Saint-Éloi de Raincheval.}\par}
\newpage
\phantomsection\label{fac:44}
% --- En-tête du fac-similé (Page impaire / droite : Photographie du manuscrit) ---
\noindent\rule{\textwidth}{0.5pt}\par\vspace{0.55mm}
\noindent\textbf{\fontsize{6.8pt}{7.8pt}\selectfont FAC-SIMILÉ EN VIS-À-VIS — FOLIO 22 RECTO} \hfill \textbf{\fontsize{6.8pt}{7.8pt}\selectfont SEPTEMBRE 1916}\par\vspace{0.25mm}
\noindent{\fontsize{6.8pt}{8.0pt}\selectfont\textit{Livre-Journal d'Hippolyte Combes} — Dimanche 17 septembre 1916 (Raincheval (Somme))}\par\vspace{0.55mm}
\noindent\rule{\textwidth}{0.35pt}\par\vspace*{\fill}
{\centering
  \setlength{\fboxsep}{1.2mm}%
  \setlength{\fboxrule}{0.4pt}%
  \fcolorbox{gray!55}{white}{%
    \includegraphics[width=0.95\textwidth,height=0.81\textheight,keepaspectratio]{\facsimiledir/20251126_094706.jpg}%
  }\par}
\vspace*{\fill}
\newpage

% =============================================================================
% VIS-A-VIS SPREAD : FOLIO 22 VERSO (20251126_094726.jpg)
% =============================================================================
\cleardoubleevenpage
\renewcommand{\currentmonthyear}{Septembre 1916}
\renewcommand{\currentfoliolabel}{Folio 22 verso}
\phantomsection\label{fol:45}
% --- En-tête du folio (Page paire / gauche : Transcription, apparat & registres de bas de page) ---
\noindent\rule{\textwidth}{0.5pt}\par\vspace{0.55mm}
\noindent\textbf{\footnotesize FOLIO 22 VERSO} \hfill \textbf{\footnotesize SEPTEMBRE 1916}\par\vspace{0.25mm}
\noindent{\fontsize{6.9pt}{8.0pt}\selectfont Chevauchée sous l'averse par Hédauville et Forceville jusqu'à Acheux et malaise gastrique}\par\vspace{0.2mm}
\noindent{\fontsize{6.9pt}{8.0pt}\selectfont\textit{Lundi 18 septembre 1916} \hfill \textit{Raincheval, Hédauville, Forceville et Acheux-en-Amiénois (Somme)}}\par\vspace{0.5mm}
\noindent\rule{\textwidth}{0.35pt}\par\vspace{0.8mm}
% --- Corps de la transcription diplomatique et apparat marginal en prolongement ---
{\renewcommand{\arraystretch}{0.87}%
\noindent\hspace*{3.5mm}\begin{tabular}{@{}p{87mm}@{\hspace{3.5mm}}p{39mm}@{}}
      \leavevmode\fontsize{6.9pt}{7.9pt}\selectfont\strut \underline{18 Lundi.} \hfill{} \underline{Septbr 1916}\strut &  \\
      \leavevmode\fontsize{6.9pt}{7.9pt}\selectfont\strut Tout le jour, pluie très forte\strut &  \\
      \leavevmode\fontsize{6.9pt}{7.9pt}\selectfont\strut Reçu hier soir message : obligé venir mission\strut &  \\
      \leavevmode\fontsize{6.9pt}{7.9pt}\selectfont\strut frse. Malgré averses incessantes, parti\strut & \smash{\parbox[t]{39mm}{\fontsize{6.2pt}{7.0pt}\selectfont\raggedright\color{black!88}\textbf{4 frse} \textit{abbr. pro} française.}} \\
      \leavevmode\llap{\fontsize{6.4pt}{7.3pt}\selectfont\color{darkgray}5\hspace{1.8mm}}\fontsize{6.9pt}{7.9pt}\selectfont\strut à cheval vers 1H 1/2. Toute l'après-midi par\strut &  \\
      \leavevmode\fontsize{6.9pt}{7.9pt}\selectfont\strut monts \& par vaux, sur routes boueuses\strut &  \\
      \leavevmode\fontsize{6.9pt}{7.9pt}\selectfont\strut A Hédauville, appris [groupe abrégé] divon, contrt\strut & \smash{\parbox[t]{39mm}{\fontsize{6.2pt}{7.0pt}\selectfont\raggedright\color{black!88}\textbf{7} \textit{lect. incert. in ms.} \textbf{7 contrt} \textit{abbr. pro} contrairement. — \textbf{8} \textit{lect. incert. in ms.} \textbf{8 divon} \textit{abbr. pro} division. — \textbf{10} \textit{lect. incert. in ms.} — \textbf{11 frse} \textit{abbr. pro} française. \textbf{11 divon} \textit{abbr. pro} division. — \textbf{12 cie} \textit{abbr. pro} compagnie. — \textbf{15 dérangt} \textit{abbr. pro} dérangement. — \textbf{16 qq} \textit{abbr. pro} quelque(s). \textbf{16 changemt} \textit{abbr. pro} changement. — \textbf{17 offs} \textit{abbr. pro} officiers. \textbf{17 mess} \textit{angl. in ms.}}} \\
      \leavevmode\fontsize{6.9pt}{7.9pt}\selectfont\strut [bruit ?], pas quitté Acheux — Le H.Q. 18e divon reste\strut &  \\
      \leavevmode\fontsize{6.9pt}{7.9pt}\selectfont\strut où elle était.\strut &  \\
      \leavevmode\llap{\fontsize{6.4pt}{7.3pt}\selectfont\color{darkgray}10\hspace{1.8mm}}\fontsize{6.9pt}{7.9pt}\selectfont\strut [mot surchargé] par Forceville à Acheux — où\strut &  \\
      \leavevmode\fontsize{6.9pt}{7.9pt}\selectfont\strut enfin retrouvé mission frse de notre divon —\strut &  \\
      \leavevmode\fontsize{6.9pt}{7.9pt}\selectfont\strut Au retour, fait route en cie de Brammas\strut &  \\
      \leavevmode\fontsize{6.9pt}{7.9pt}\selectfont\strut Temps affreux. Cape traversée par la pluie\strut &  \\
      \leavevmode\fontsize{6.9pt}{7.9pt}\selectfont\strut Habits trempés. Souliers remplis d'eau\strut &  \\
      \leavevmode\llap{\fontsize{6.4pt}{7.3pt}\selectfont\color{darkgray}15\hspace{1.8mm}}\fontsize{6.9pt}{7.9pt}\selectfont\strut En outre, dérangt d'entrailles, depuis\strut &  \\
      \leavevmode\fontsize{6.9pt}{7.9pt}\selectfont\strut qq jours, par suite changemt de régime\strut &  \\
      \leavevmode\fontsize{6.9pt}{7.9pt}\selectfont\strut (mess d'offs\strut &  \\
      \leavevmode\fontsize{6.9pt}{7.9pt}\selectfont\strut Rentré chambre, mal à l'aise. Selles\strut &  \\
      \leavevmode\fontsize{6.9pt}{7.9pt}\selectfont\strut fréquentes. Rendu enfin repas de midi\strut &  \\
      \leavevmode\llap{\fontsize{6.4pt}{7.3pt}\selectfont\color{darkgray}20\hspace{1.8mm}}\fontsize{6.9pt}{7.9pt}\selectfont\strut qui me chargeait l'estomac\strut &  \\
      \leavevmode\fontsize{6.9pt}{7.9pt}\selectfont\strut Couché de b. heure\strut &  \\
      \cline{1-1}
      \leavevmode\fontsize{6.9pt}{7.9pt}\selectfont\strut Tip à Connelly mon groom : 2 f –\strut & \smash{\parbox[t]{39mm}{\fontsize{6.2pt}{7.0pt}\selectfont\raggedright\color{black!88}\textbf{22 groom} \textit{angl. in ms.}}} \\
\end{tabular}}%
\par\vspace*{\fill}
% --- Registres de bas de page (1. Critique interne [si pertinent] ; 2. Loci paralleli en latin ; 3. Commentaire historique en retrait) ---
\noindent\rule{\textwidth}{0.5pt}\par\vspace{0.65mm}
{\fontsize{6.6pt}{7.7pt}\selectfont%
\noindent\normalfont\color{black!92} \textbf{7} (\textit{Hédauville}) cf. infra f\textsuperscript{o}~38~v\textsuperscript{o} (p.~\pageref{fol:77}), f\textsuperscript{o}~54~r\textsuperscript{o} (p.~\pageref{fol:108}). — \textbf{8} (\textit{Acheux}) cf. supra f\textsuperscript{o}~21~v\textsuperscript{o} (p.~\pageref{fol:43}), infra f\textsuperscript{o}~23~r\textsuperscript{o} (p.~\pageref{fol:46}), f\textsuperscript{o}~24~r\textsuperscript{o} (p.~\pageref{fol:48}), \textit{et passim}. — \textbf{10} (\textit{Forceville}) cf. infra f\textsuperscript{o}~24~r\textsuperscript{o} (p.~\pageref{fol:48}), f\textsuperscript{o}~24~v\textsuperscript{o} (p.~\pageref{fol:49}).\par\vspace{0.55mm}
\noindent\rule{38mm}{0.4pt}\par\vspace{0.5mm}
\noindent\hspace*{5mm}\parbox{\dimexpr\textwidth-5mm\relax}{\fontsize{6.6pt}{7.7pt}\selectfont\color{black!84}Convoqué sous une pluie battante le 18 septembre 1916, Hippolyte Combes chevauche jusqu'à \textbf{Hédauville} puis par \textbf{Forceville} jusqu'à \textbf{Acheux-en-Amiénois}, où se trouvent toujours l'état-major (\textit{H.Q.}) de la \textbf{18\textsuperscript{e} Division britannique} et sa Mission française (la division prépare alors l'assaut de Thiepval du 26 septembre; Nichols, \textit{op.~cit.}, p.~100).}\par}
\newpage
\phantomsection\label{fac:45}
% --- En-tête du fac-similé (Page impaire / droite : Photographie du manuscrit) ---
\noindent\rule{\textwidth}{0.5pt}\par\vspace{0.55mm}
\noindent\textbf{\fontsize{6.8pt}{7.8pt}\selectfont FAC-SIMILÉ EN VIS-À-VIS — FOLIO 22 VERSO} \hfill \textbf{\fontsize{6.8pt}{7.8pt}\selectfont SEPTEMBRE 1916}\par\vspace{0.25mm}
\noindent{\fontsize{6.8pt}{8.0pt}\selectfont\textit{Livre-Journal d'Hippolyte Combes} — Lundi 18 septembre 1916 (Raincheval, Hédauville, Forceville et Acheux-en-Amiénois (Somme))}\par\vspace{0.55mm}
\noindent\rule{\textwidth}{0.35pt}\par\vspace*{\fill}
{\centering
  \setlength{\fboxsep}{1.2mm}%
  \setlength{\fboxrule}{0.4pt}%
  \fcolorbox{gray!55}{white}{%
    \includegraphics[width=0.95\textwidth,height=0.81\textheight,keepaspectratio]{\facsimiledir/20251126_094726.jpg}%
  }\par}
\vspace*{\fill}
\newpage

% =============================================================================
% VIS-A-VIS SPREAD : FOLIO 23 RECTO (20251126_094738.jpg)
% =============================================================================
\cleardoubleevenpage
\renewcommand{\currentmonthyear}{Septembre 1916}
\renewcommand{\currentfoliolabel}{Folio 23 recto}
\phantomsection\label{fol:46}
% --- En-tête du folio (Page paire / gauche : Transcription, apparat & registres de bas de page) ---
\noindent\rule{\textwidth}{0.5pt}\par\vspace{0.55mm}
\noindent\textbf{\footnotesize FOLIO 23 RECTO} \hfill \textbf{\footnotesize SEPTEMBRE 1916}\par\vspace{0.25mm}
\noindent{\fontsize{6.9pt}{8.0pt}\selectfont Routes transformées en ornières et perception de la solde à la hutte du château d'Acheux}\par\vspace{0.2mm}
\noindent{\fontsize{6.9pt}{8.0pt}\selectfont\textit{Jeudi 21 septembre 1916} \hfill \textit{Raincheval et Acheux-en-Amiénois (Somme)}}\par\vspace{0.5mm}
\noindent\rule{\textwidth}{0.35pt}\par\vspace{0.8mm}
% --- Corps de la transcription diplomatique et apparat marginal en prolongement ---
{\renewcommand{\arraystretch}{1.00}%
\noindent\hspace*{3.5mm}\begin{tabular}{@{}p{87mm}@{\hspace{3.5mm}}p{39mm}@{}}
      \leavevmode\fontsize{7.7pt}{9.0pt}\selectfont\strut \underline{21 J.} \hfill{} \underline{Septbre /16}\strut & \smash{\parbox[t]{39mm}{\fontsize{6.2pt}{7.0pt}\selectfont\raggedright\color{black!88}\textbf{1 J.} \textit{scil.} Jeudi.}} \\
      \leavevmode\fontsize{7.7pt}{9.0pt}\selectfont\strut Ts ces derniers jours, averses continuelles.\strut & \smash{\parbox[t]{39mm}{\fontsize{6.2pt}{7.0pt}\selectfont\raggedright\color{black!88}\textbf{2 Ts} \textit{abbr. pro} Tous.}} \\
      \leavevmode\fontsize{7.7pt}{9.0pt}\selectfont\strut Routes boueuses transformées en ornières\strut &  \\
      \leavevmode\fontsize{7.7pt}{9.0pt}\selectfont\strut Enfin ce matin, malgré temps\strut &  \\
      \leavevmode\llap{\fontsize{6.4pt}{7.3pt}\selectfont\color{darkgray}5\hspace{1.8mm}}\fontsize{7.7pt}{9.0pt}\selectfont\strut menaçant, pluie pas tombée.\strut &  \\
      \leavevmode\fontsize{7.7pt}{9.0pt}\selectfont\strut Après-midi, ciel tjs nuageux — mais\strut & \smash{\parbox[t]{39mm}{\fontsize{6.2pt}{7.0pt}\selectfont\raggedright\color{black!88}\textbf{6 tjs} \textit{abbr. pro} toujours.}} \\
      \leavevmode\fontsize{7.7pt}{9.0pt}\selectfont\strut soleil brille un peu\strut &  \\
      \leavevmode\fontsize{7.7pt}{9.0pt}\selectfont\strut Venu à cheval à Acheux – bureau\strut &  \\
      \leavevmode\fontsize{7.7pt}{9.0pt}\selectfont\strut mission fr, hutte s/ la pelouse chateau\strut & \smash{\parbox[t]{39mm}{\fontsize{6.2pt}{7.0pt}\selectfont\raggedright\color{black!88}\textbf{9 s/} \textit{abbr. pro} sur.}} \\
      \leavevmode\llap{\fontsize{6.4pt}{7.3pt}\selectfont\color{darkgray}10\hspace{1.8mm}}\fontsize{7.7pt}{9.0pt}\selectfont\strut A 4H 1/2, retour – et sans avoir été mouillé\strut &  \\
      \leavevmode\fontsize{7.7pt}{9.0pt}\selectfont\strut Fin journée, assez belle\strut &  \\
      \cline{1-1}
      \leavevmode\fontsize{7.7pt}{9.0pt}\selectfont\strut \underline{A Acheux} : solde : 78 f 75 c\strut &  \\
      \leavevmode\fontsize{7.7pt}{9.0pt}\selectfont\strut — cartes ples : 0.25 c.\strut & \smash{\parbox[t]{39mm}{\fontsize{6.2pt}{7.0pt}\selectfont\raggedright\color{black!88}\textbf{13 cartes ples} \textit{abbr. pro} cartes postales. \textbf{13 c.} \textit{abbr. pro} centimes.}} \\
\end{tabular}}%
\par\vspace*{\fill}
% --- Registres de bas de page (1. Critique interne [si pertinent] ; 2. Loci paralleli en latin ; 3. Commentaire historique en retrait) ---
\noindent\rule{\textwidth}{0.5pt}\par\vspace{0.65mm}
{\fontsize{6.6pt}{7.7pt}\selectfont%
\noindent\normalfont\color{black!92} \textbf{8} (\textit{Acheux}) cf. supra f\textsuperscript{o}~21~v\textsuperscript{o} (p.~\pageref{fol:43}), f\textsuperscript{o}~22~v\textsuperscript{o} (p.~\pageref{fol:45}), infra f\textsuperscript{o}~24~r\textsuperscript{o} (p.~\pageref{fol:48}), f\textsuperscript{o}~24~v\textsuperscript{o} (p.~\pageref{fol:49}), \textit{et passim}.\par\vspace{0.55mm}
\noindent\rule{38mm}{0.4pt}\par\vspace{0.5mm}
\noindent\hspace*{5mm}\parbox{\dimexpr\textwidth-5mm\relax}{\fontsize{6.6pt}{7.7pt}\selectfont\color{black!84}Profitant d'une accalmie après trois journées d'averses qui ont détrempé les chemins de la Somme, Hippolyte Combes se rend à cheval au bureau de la Mission française installé dans une baraque (\textit{hutte}) sur la pelouse du château d'\textbf{Acheux-en-Amiénois} et y perçoit sa solde de quinzaine (78~fr.~75), à quatre jours de son départ en permission.}\par}
\newpage
\phantomsection\label{fac:46}
% --- En-tête du fac-similé (Page impaire / droite : Photographie du manuscrit) ---
\noindent\rule{\textwidth}{0.5pt}\par\vspace{0.55mm}
\noindent\textbf{\fontsize{6.8pt}{7.8pt}\selectfont FAC-SIMILÉ EN VIS-À-VIS — FOLIO 23 RECTO} \hfill \textbf{\fontsize{6.8pt}{7.8pt}\selectfont SEPTEMBRE 1916}\par\vspace{0.25mm}
\noindent{\fontsize{6.8pt}{8.0pt}\selectfont\textit{Livre-Journal d'Hippolyte Combes} — Jeudi 21 septembre 1916 (Raincheval et Acheux-en-Amiénois (Somme))}\par\vspace{0.55mm}
\noindent\rule{\textwidth}{0.35pt}\par\vspace*{\fill}
{\centering
  \setlength{\fboxsep}{1.2mm}%
  \setlength{\fboxrule}{0.4pt}%
  \fcolorbox{gray!55}{white}{%
    \includegraphics[width=0.95\textwidth,height=0.81\textheight,keepaspectratio]{\facsimiledir/20251126_094738.jpg}%
  }\par}
\vspace*{\fill}
\newpage

% =============================================================================
% VIS-A-VIS SPREAD : FOLIO 23 VERSO (20251126_094822.jpg)
% =============================================================================
\cleardoubleevenpage
\renewcommand{\currentmonthyear}{Septembre 1916}
\renewcommand{\currentfoliolabel}{Folio 23 verso}
\phantomsection\label{fol:47}
% --- En-tête du folio (Page paire / gauche : Transcription, apparat & registres de bas de page) ---
\noindent\rule{\textwidth}{0.5pt}\par\vspace{0.55mm}
\noindent\textbf{\footnotesize FOLIO 23 VERSO} \hfill \textbf{\footnotesize SEPTEMBRE 1916}\par\vspace{0.25mm}
\noindent{\fontsize{6.9pt}{8.0pt}\selectfont Match de football des « Bedfs », annonce du départ (« we are moving demain ») et nuit sombre}\par\vspace{0.2mm}
\noindent{\fontsize{6.9pt}{8.0pt}\selectfont\textit{Vendredi 22 septembre 1916} \hfill \textit{Raincheval (Somme)}}\par\vspace{0.5mm}
\noindent\rule{\textwidth}{0.35pt}\par\vspace{0.8mm}
% --- Corps de la transcription diplomatique et apparat marginal en prolongement ---
{\renewcommand{\arraystretch}{0.87}%
\noindent\hspace*{3.5mm}\begin{tabular}{@{}p{87mm}@{\hspace{3.5mm}}p{39mm}@{}}
      \leavevmode\fontsize{6.9pt}{7.9pt}\selectfont\strut \underline{22 Ve} – Enfin, temps splendide Sept / 16\strut &  \\
      \leavevmode\fontsize{6.9pt}{7.9pt}\selectfont\strut ciel d'1 bleu intense. Soleil ardent [t./1/2 ?] jour\strut & \smash{\parbox[t]{39mm}{\fontsize{6.2pt}{7.0pt}\selectfont\raggedright\color{black!88}\textbf{2} \textit{lect. incert. in ms.}}} \\
      \leavevmode\fontsize{6.9pt}{7.9pt}\selectfont\strut Après-midi, promené j'à la station de\strut & \smash{\parbox[t]{39mm}{\fontsize{6.2pt}{7.0pt}\selectfont\raggedright\color{black!88}\textbf{3 j'à} \textit{abbr. pro} jusqu'à.}} \\
      \leavevmode\fontsize{6.9pt}{7.9pt}\selectfont\strut Raincheval et me suis renseigné horaire trains\strut &  \\
      \leavevmode\llap{\fontsize{6.4pt}{7.3pt}\selectfont\color{darkgray}5\hspace{1.8mm}}\fontsize{6.9pt}{7.9pt}\selectfont\strut Au retour, assisté match de football\strut &  \\
      \leavevmode\fontsize{6.9pt}{7.9pt}\selectfont\strut s/ prairie à coté horse-lines the Bedfs\strut & \smash{\parbox[t]{39mm}{\fontsize{6.2pt}{7.0pt}\selectfont\raggedright\color{black!88}\textbf{6 s/} \textit{abbr. pro} sur. \textbf{6 horse-lines} \textit{angl. in ms.}}} \\
      \leavevmode\fontsize{6.9pt}{7.9pt}\selectfont\strut maillots rouge vif, \& blancs — à l'arrière-\strut &  \\
      \leavevmode\fontsize{6.9pt}{7.9pt}\selectfont\strut plan, très beaux arbres —\strut &  \\
      \leavevmode\fontsize{6.9pt}{7.9pt}\selectfont\strut Cme tournais en face chateau, Cap.\strut & \smash{\parbox[t]{39mm}{\fontsize{6.2pt}{7.0pt}\selectfont\raggedright\color{black!88}\textbf{9 Cme} \textit{abbr. pro} Comme. \textbf{9 Cap.} \textit{abbr. pro} Capitaine. — \textbf{10 moving} \textit{angl. in ms.} — \textbf{11 gde} \textit{abbr. pro} grande.}} \\
      \leavevmode\llap{\fontsize{6.4pt}{7.3pt}\selectfont\color{darkgray}10\hspace{1.8mm}}\fontsize{6.9pt}{7.9pt}\selectfont\strut Panton m'a appris que we are moving demain.\strut &  \\
      \leavevmode\fontsize{6.9pt}{7.9pt}\selectfont\strut gde déception — \& mélancolie —\strut &  \\
      \leavevmode\fontsize{6.9pt}{7.9pt}\selectfont\strut Dernière nuit chambre, préoccupé\strut &  \\
      \leavevmode\fontsize{6.9pt}{7.9pt}\selectfont\strut paquets à faire et voyage à Toulouse\strut &  \\
      \leavevmode\fontsize{6.9pt}{7.9pt}\selectfont\strut départ le 25\strut &  \\
      \leavevmode\llap{\fontsize{6.4pt}{7.3pt}\selectfont\color{darkgray}15\hspace{1.8mm}}\fontsize{6.9pt}{7.9pt}\selectfont\strut Nuit très sombre — qq stars au ciel noir.\strut & \smash{\parbox[t]{39mm}{\fontsize{6.2pt}{7.0pt}\selectfont\raggedright\color{black!88}\textbf{15 qq} \textit{abbr. pro} quelque(s). — \textbf{16} \textit{in marg. cum unco} (comptes). \textbf{16 tins} \textit{angl. in ms.} — \textbf{17} \textit{lect. incert. in ms.} — \textbf{18 food} \textit{angl. in ms.}}} \\
      \cline{1-1}
      \leavevmode\fontsize{6.9pt}{7.9pt}\selectfont\strut [ — 2 tins of coffee milk, à la\strut &  \\
      \leavevmode\fontsize{6.9pt}{7.9pt}\selectfont\strut Cantine : 3 f [25 ?]\strut &  \\
      \leavevmode\fontsize{6.9pt}{7.9pt}\selectfont\strut Eggs \& food : 1 f 10 — ]\strut &  \\
\end{tabular}}%
\par\vspace*{\fill}
% --- Registres de bas de page (1. Critique interne [si pertinent] ; 2. Loci paralleli en latin ; 3. Commentaire historique en retrait) ---
\noindent\rule{\textwidth}{0.5pt}\par\vspace{0.65mm}
{\fontsize{6.6pt}{7.7pt}\selectfont%
\noindent\normalfont\color{black!92} \textbf{4} (\textit{Raincheval}) cf. supra f\textsuperscript{o}~20~v\textsuperscript{o} (p.~\pageref{fol:41}), f\textsuperscript{o}~23~r\textsuperscript{o} (p.~\pageref{fol:46}), infra f\textsuperscript{o}~24~r\textsuperscript{o} (p.~\pageref{fol:48}), f\textsuperscript{o}~45~v\textsuperscript{o} (p.~\pageref{fol:91}), \textit{et passim}. — \textbf{6} (\textit{Bedfs}) cf. supra f\textsuperscript{o}~19~r\textsuperscript{o} (p.~\pageref{fol:38}). — \textbf{10} (\textit{Panton}) cf. infra f\textsuperscript{o}~34~r\textsuperscript{o} (p.~\pageref{fol:68}), f\textsuperscript{o}~36~v\textsuperscript{o} (p.~\pageref{fol:73}).\par\vspace{0.55mm}
\noindent\rule{38mm}{0.4pt}\par\vspace{0.5mm}
\noindent\hspace*{5mm}\parbox{\dimexpr\textwidth-5mm\relax}{\fontsize{6.6pt}{7.7pt}\selectfont\color{black!84}Après s'être informé des trains à la gare de Raincheval en vue de son départ en permission pour Toulouse fixé au 25 septembre, Hippolyte Combes assiste à un match de football près des lignes de chevaux des \textit{Bedfords} (quasi-certainement le \textbf{7\textsuperscript{e} Bataillon du \textit{Bedfordshire Regiment}}, 54\textsuperscript{e} Brigade de la 18\textsuperscript{e} Division, déjà mentionné au f\textsuperscript{o}~19~r\textsuperscript{o}).}\par}
\newpage
\phantomsection\label{fac:47}
% --- En-tête du fac-similé (Page impaire / droite : Photographie du manuscrit) ---
\noindent\rule{\textwidth}{0.5pt}\par\vspace{0.55mm}
\noindent\textbf{\fontsize{6.8pt}{7.8pt}\selectfont FAC-SIMILÉ EN VIS-À-VIS — FOLIO 23 VERSO} \hfill \textbf{\fontsize{6.8pt}{7.8pt}\selectfont SEPTEMBRE 1916}\par\vspace{0.25mm}
\noindent{\fontsize{6.8pt}{8.0pt}\selectfont\textit{Livre-Journal d'Hippolyte Combes} — Vendredi 22 septembre 1916 (Raincheval (Somme))}\par\vspace{0.55mm}
\noindent\rule{\textwidth}{0.35pt}\par\vspace*{\fill}
{\centering
  \setlength{\fboxsep}{1.2mm}%
  \setlength{\fboxrule}{0.4pt}%
  \fcolorbox{gray!55}{white}{%
    \includegraphics[width=0.95\textwidth,height=0.81\textheight,keepaspectratio]{\facsimiledir/20251126_094822.jpg}%
  }\par}
\vspace*{\fill}
\newpage

% =============================================================================
% VIS-A-VIS SPREAD : FOLIO 24 RECTO (20251126_094835.jpg)
% =============================================================================
\cleardoubleevenpage
\renewcommand{\currentmonthyear}{Septembre 1916}
\renewcommand{\currentfoliolabel}{Folio 24 recto}
\phantomsection\label{fol:48}
% --- En-tête du folio (Page paire / gauche : Transcription, apparat & registres de bas de page) ---
\noindent\rule{\textwidth}{0.5pt}\par\vspace{0.55mm}
\noindent\textbf{\footnotesize FOLIO 24 RECTO} \hfill \textbf{\footnotesize SEPTEMBRE 1916}\par\vspace{0.25mm}
\noindent{\fontsize{6.9pt}{8.0pt}\selectfont Samedi 23 septembre 1916 — Marche vers Léalvillers, message apporté en auto par le lieutenant Pizzi et bivouac sous la tente à Forceville}\par\vspace{0.2mm}
\noindent{\fontsize{6.9pt}{8.0pt}\selectfont\textit{Samedi 23 septembre 1916} \hfill \textit{Raincheval, Léalvillers, Acheux et Forceville (Somme)}}\par\vspace{0.5mm}
\noindent\rule{\textwidth}{0.35pt}\par\vspace{0.8mm}
% --- Corps de la transcription diplomatique et apparat marginal en prolongement ---
{\renewcommand{\arraystretch}{0.87}%
\noindent\hspace*{3.5mm}\begin{tabular}{@{}p{87mm}@{\hspace{3.5mm}}p{39mm}@{}}
      \leavevmode\fontsize{6.9pt}{7.9pt}\selectfont\strut \underline{Sam 23} \hfill{} \underline{Sept 16}\strut &  \\
      \leavevmode\fontsize{6.9pt}{7.9pt}\selectfont\strut Très belle journée.\strut &  \\
      \leavevmode\fontsize{6.9pt}{7.9pt}\selectfont\strut Levé 6H 1/2 et fait ts préparatifs.\strut & \smash{\parbox[t]{39mm}{\fontsize{6.2pt}{7.0pt}\selectfont\raggedright\color{black!88}\textbf{3 ts} \textit{abbr. pro} tous.}} \\
      \leavevmode\fontsize{6.9pt}{7.9pt}\selectfont\strut Vers 7H 1/2, Cap West venu me dire que départ\strut &  \\
      \leavevmode\llap{\fontsize{6.4pt}{7.3pt}\selectfont\color{darkgray}5\hspace{1.8mm}}\fontsize{6.9pt}{7.9pt}\selectfont\strut fixé à 9H 30.\strut &  \\
      \leavevmode\fontsize{6.9pt}{7.9pt}\selectfont\strut Quitté village heure indiquée\strut &  \\
      \leavevmode\fontsize{6.9pt}{7.9pt}\selectfont\strut Temps superbe Soleil ardent\strut &  \\
      \leavevmode\fontsize{6.9pt}{7.9pt}\selectfont\strut A peine arrivions-ns vers garage de Léalvillers\strut & \smash{\parbox[t]{39mm}{\fontsize{6.2pt}{7.0pt}\selectfont\raggedright\color{black!88}\textbf{8 ns} \textit{abbr. pro} nous.}} \\
      \leavevmode\fontsize{6.9pt}{7.9pt}\selectfont\strut que Pizzi, en auto, à notre rencontre av. message\strut &  \\
      \leavevmode\llap{\fontsize{6.4pt}{7.3pt}\selectfont\color{darkgray}10\hspace{1.8mm}}\fontsize{6.9pt}{7.9pt}\selectfont\strut Dû tourner bride et aller attendre champs à\strut &  \\
      \leavevmode\fontsize{6.9pt}{7.9pt}\selectfont\strut coté refilling point\strut & \smash{\parbox[t]{39mm}{\fontsize{6.2pt}{7.0pt}\selectfont\raggedright\color{black!88}\textbf{11 refilling point} \textit{angl. in ms.} — \textbf{12 j. à} \textit{abbr. pro} jusqu'à. — \textbf{13 Cap.} \textit{abbr. pro} Capitaine. \textbf{13 1/2} \textit{scil.} demi / moitié. — \textbf{14 tte} \textit{abbr. pro} toute. \textbf{14 Coy} \textit{angl.} Company.}} \\
      \leavevmode\fontsize{6.9pt}{7.9pt}\selectfont\strut Repartis à cheval seul j. à gare d'Acheux\strut &  \\
      \leavevmode\fontsize{6.9pt}{7.9pt}\selectfont\strut Cap. West, vers 2 H 1/2, m'a rejoint et aussitôt\strut &  \\
      \leavevmode\fontsize{6.9pt}{7.9pt}\selectfont\strut tte Coy file. –\strut &  \\
      \leavevmode\llap{\fontsize{6.4pt}{7.3pt}\selectfont\color{darkgray}15\hspace{1.8mm}}\fontsize{6.9pt}{7.9pt}\selectfont\strut Revenus route gare j'au tournant route de\strut &  \\
      \leavevmode\fontsize{6.9pt}{7.9pt}\selectfont\strut Forceville\strut &  \\
      \leavevmode\fontsize{6.9pt}{7.9pt}\selectfont\strut Cantonnés in the open – champs au flanc\strut &  \\
      \leavevmode\fontsize{6.9pt}{7.9pt}\selectfont\strut coteau.\strut &  \\
      \leavevmode\fontsize{6.9pt}{7.9pt}\selectfont\strut Très b. fin journée.\strut &  \\
      \leavevmode\llap{\fontsize{6.4pt}{7.3pt}\selectfont\color{darkgray}20\hspace{1.8mm}}\fontsize{6.9pt}{7.9pt}\selectfont\strut Me suis arrangé hutte avec épaisse [bâche ?]\strut & \smash{\parbox[t]{39mm}{\fontsize{6.2pt}{7.0pt}\selectfont\raggedright\color{black!88}\textbf{20} \textit{lect. incert. in ms.}}} \\
      \leavevmode\fontsize{6.9pt}{7.9pt}\selectfont\strut verte cme litière — Couché bonne nuit\strut & \smash{\parbox[t]{39mm}{\fontsize{6.2pt}{7.0pt}\selectfont\raggedright\color{black!88}\textbf{21 cme} \textit{abbr. pro} comme.}} \\
      \cline{1-1}
      \leavevmode\fontsize{6.9pt}{7.9pt}\selectfont\strut A Mme — propre à Acheux : 1 f (cooking\strut & \smash{\parbox[t]{39mm}{\fontsize{6.2pt}{7.0pt}\selectfont\raggedright\color{black!88}\textbf{22 Mme} \textit{abbr. pro} Madame. \textbf{22 propre} \textit{abbr. pro} propriétaire. — \textbf{23} \textit{in marg. cum unco} (comptes).}} \\
      \leavevmode\fontsize{6.9pt}{7.9pt}\selectfont\strut [ Carte du Petit Journal : 60 c —\strut &  \\
\end{tabular}}%
\par\vspace*{\fill}
% --- Registres de bas de page (1. Critique interne [si pertinent] ; 2. Loci paralleli en latin ; 3. Commentaire historique en retrait) ---
\noindent\rule{\textwidth}{0.5pt}\par\vspace{0.65mm}
{\fontsize{6.6pt}{7.7pt}\selectfont%
\noindent\normalfont\color{black!92} \textbf{4} (\textit{West}) cf. supra f\textsuperscript{o}~19~v\textsuperscript{o} (p.~\pageref{fol:39}), f\textsuperscript{o}~21~r\textsuperscript{o} (p.~\pageref{fol:42}). — \textbf{8} (\textit{Léalvillers}) cf. infra f\textsuperscript{o}~47~r\textsuperscript{o} (p.~\pageref{fol:94}), f\textsuperscript{o}~47~v\textsuperscript{o} (p.~\pageref{fol:95}), \textit{et passim}. — \textbf{9} (\textit{Pizzi}) cf. infra f\textsuperscript{o}~42~v\textsuperscript{o} (p.~\pageref{fol:85}), f\textsuperscript{o}~45~r\textsuperscript{o} (p.~\pageref{fol:90}).\par\vspace{0.55mm}
\noindent\rule{38mm}{0.4pt}\par\vspace{0.5mm}
\noindent\hspace*{5mm}\parbox{\dimexpr\textwidth-5mm\relax}{\fontsize{6.6pt}{7.7pt}\selectfont\color{black!84}Le samedi 23 septembre 1916 — à trois jours de l'attaque de la 18\textsuperscript{e} Division contre Thiepval —, la compagnie quitte Raincheval à 9~h~30 sous les ordres du \textbf{capitaine West}.}\par}
\newpage
\phantomsection\label{fac:48}
% --- En-tête du fac-similé (Page impaire / droite : Photographie du manuscrit) ---
\noindent\rule{\textwidth}{0.5pt}\par\vspace{0.55mm}
\noindent\textbf{\fontsize{6.8pt}{7.8pt}\selectfont FAC-SIMILÉ EN VIS-À-VIS — FOLIO 24 RECTO} \hfill \textbf{\fontsize{6.8pt}{7.8pt}\selectfont SEPTEMBRE 1916}\par\vspace{0.25mm}
\noindent{\fontsize{6.8pt}{8.0pt}\selectfont\textit{Livre-Journal d'Hippolyte Combes} — Samedi 23 septembre 1916 (Raincheval, Léalvillers, Acheux et Forceville (Somme))}\par\vspace{0.55mm}
\noindent\rule{\textwidth}{0.35pt}\par\vspace*{\fill}
{\centering
  \setlength{\fboxsep}{1.2mm}%
  \setlength{\fboxrule}{0.4pt}%
  \fcolorbox{gray!55}{white}{%
    \includegraphics[width=0.95\textwidth,height=0.81\textheight,keepaspectratio]{\facsimiledir/20251126_094835.jpg}%
  }\par}
\vspace*{\fill}
\newpage

% =============================================================================
% VIS-A-VIS SPREAD : FOLIO 24 VERSO (20251126_103003.jpg)
% =============================================================================
\cleardoubleevenpage
\renewcommand{\currentmonthyear}{Septembre 1916}
\renewcommand{\currentfoliolabel}{Folio 24 verso}
\phantomsection\label{fol:49}
% --- En-tête du folio (Page paire / gauche : Transcription, apparat & registres de bas de page) ---
\noindent\rule{\textwidth}{0.5pt}\par\vspace{0.55mm}
\noindent\textbf{\footnotesize FOLIO 24 VERSO} \hfill \textbf{\footnotesize SEPTEMBRE 1916}\par\vspace{0.25mm}
\noindent{\fontsize{6.9pt}{8.0pt}\selectfont Dimanche 24 septembre 1916 — Construction d'une hutte derrière la haie d'aubépine à Forceville et préparatifs du départ en permission}\par\vspace{0.2mm}
\noindent{\fontsize{6.9pt}{8.0pt}\selectfont\textit{Dimanche 24 septembre 1916} \hfill \textit{Forceville et Acheux-en-Amiénois (Somme)}}\par\vspace{0.5mm}
\noindent\rule{\textwidth}{0.35pt}\par\vspace{0.8mm}
% --- Corps de la transcription diplomatique et apparat marginal en prolongement ---
{\renewcommand{\arraystretch}{0.87}%
\noindent\hspace*{3.5mm}\begin{tabular}{@{}p{87mm}@{\hspace{3.5mm}}p{39mm}@{}}
      \leavevmode\fontsize{6.9pt}{7.9pt}\selectfont\strut \underline{Dim. 24} Levé 7H 1/2 \hfill{} \raisebox{0.5ex}{\underline{Sept 16}}\strut & \smash{\parbox[t]{39mm}{\fontsize{6.2pt}{7.0pt}\selectfont\raggedright\color{black!88}\textbf{1 Dim.} \textit{scil.} Dimanche.}} \\
      \leavevmode\fontsize{6.9pt}{7.9pt}\selectfont\strut Après 8 H, épais brouillard. A 11 H seulemt,\strut &  \\
      \leavevmode\fontsize{6.9pt}{7.9pt}\selectfont\strut soleil l'a dissipé –\strut &  \\
      \leavevmode\fontsize{6.9pt}{7.9pt}\selectfont\strut Brume \& temps splendide\strut &  \\
      \leavevmode\llap{\fontsize{6.4pt}{7.3pt}\selectfont\color{darkgray}5\hspace{1.8mm}}\fontsize{6.9pt}{7.9pt}\selectfont\strut Achevé de construire ma hutte, [longeant ?]\strut & \smash{\parbox[t]{39mm}{\fontsize{6.2pt}{7.0pt}\selectfont\raggedright\color{black!88}\textbf{5} \textit{lect. incert. in ms.}}} \\
      \leavevmode\fontsize{6.9pt}{7.9pt}\selectfont\strut haie d'aubépine — plus spacieuse et haute\strut &  \\
      \leavevmode\fontsize{6.9pt}{7.9pt}\selectfont\strut Ds la matinée, dû venir à pied à\strut & \smash{\parbox[t]{39mm}{\fontsize{6.2pt}{7.0pt}\selectfont\raggedright\color{black!88}\textbf{7 Ds} \textit{abbr. pro} Dans.}} \\
      \leavevmode\fontsize{6.9pt}{7.9pt}\selectfont\strut Forceville, for green forage. Parcouru\strut & \smash{\parbox[t]{39mm}{\fontsize{6.2pt}{7.0pt}\selectfont\raggedright\color{black!88}\textbf{8 green} \textit{angl. in ms.}}} \\
      \leavevmode\fontsize{6.9pt}{7.9pt}\selectfont\strut village\strut &  \\
      \leavevmode\llap{\fontsize{6.4pt}{7.3pt}\selectfont\color{darkgray}10\hspace{1.8mm}}\fontsize{6.9pt}{7.9pt}\selectfont\strut Après-midi, chaleur intense –\strut &  \\
      \leavevmode\fontsize{6.9pt}{7.9pt}\selectfont\strut Lunch pris (Cap. Newton [?] invité), me suis\strut & \smash{\parbox[t]{39mm}{\fontsize{6.2pt}{7.0pt}\selectfont\raggedright\color{black!88}\textbf{11 Newton} \textit{lect. incert. in ms.} \textbf{11 Cap.} \textit{abbr. pro} Capitaine. — \textbf{12} \textit{lect. incert. in ms.} \textbf{12 ts} \textit{abbr. pro} tous. — \textbf{14 tea} \textit{angl. in ms.} — \textbf{15 milre} \textit{abbr. pro} militaire.}} \\
      \leavevmode\fontsize{6.9pt}{7.9pt}\selectfont\strut retiré sous ma tente et fait ts mes [bagages ?]\strut &  \\
      \leavevmode\fontsize{6.9pt}{7.9pt}\selectfont\strut en 2 h. tout était prêt\strut &  \\
      \leavevmode\fontsize{6.9pt}{7.9pt}\selectfont\strut Après le tea, venu me promener à\strut &  \\
      \leavevmode\llap{\fontsize{6.4pt}{7.3pt}\selectfont\color{darkgray}15\hspace{1.8mm}}\fontsize{6.9pt}{7.9pt}\selectfont\strut Acheux. Musique milre jouait sur la\strut &  \\
      \leavevmode\fontsize{6.9pt}{7.9pt}\selectfont\strut place. à g. de l'église et devant entrée du\strut &  \\
      \leavevmode\fontsize{6.9pt}{7.9pt}\selectfont\strut Chateau\strut &  \\
      \leavevmode\fontsize{6.9pt}{7.9pt}\selectfont\strut Rentré au camp à nuit tombante\strut &  \\
      \leavevmode\fontsize{6.9pt}{7.9pt}\selectfont\strut Mieux dormi ds ma hutte\strut & \smash{\parbox[t]{39mm}{\fontsize{6.2pt}{7.0pt}\selectfont\raggedright\color{black!88}\textbf{19 ds} \textit{abbr. pro} dans.}} \\
      \leavevmode\llap{\fontsize{6.4pt}{7.3pt}\selectfont\color{darkgray}20\hspace{1.8mm}}\fontsize{6.9pt}{7.9pt}\selectfont\strut Magnifique nuit étoilée –\strut &  \\
      \leavevmode\fontsize{6.9pt}{7.9pt}\selectfont\strut Canon gronde tout le long du front\strut &  \\
\end{tabular}}%
\par\vspace*{\fill}
% --- Registres de bas de page (1. Critique interne [si pertinent] ; 2. Loci paralleli en latin ; 3. Commentaire historique en retrait) ---
\noindent\rule{\textwidth}{0.5pt}\par\vspace{0.65mm}
{\fontsize{6.6pt}{7.7pt}\selectfont%
\noindent\normalfont \textbf{Disposition :} Le feuillet réunit le dimanche 24 septembre (l.~1--16) et le relevé récapitulatif des dépenses des 23 et 24 septembre (l.~17--22) encadré d'un crochet comptable.\par\vspace{0.5mm}
\noindent\rule{22mm}{0.3pt}\par\vspace{0.45mm}
\noindent\normalfont\color{black!92} \textbf{8} (\textit{Forceville}) cf. supra f\textsuperscript{o}~22~v\textsuperscript{o} (p.~\pageref{fol:45}), f\textsuperscript{o}~24~r\textsuperscript{o} (p.~\pageref{fol:48}), infra f\textsuperscript{o}~25~r\textsuperscript{o} (p.~\pageref{fol:50}), f\textsuperscript{o}~48~r\textsuperscript{o} (p.~\pageref{fol:96}). — \textbf{8} (\textit{green forage}) cf. supra f\textsuperscript{o}~3~r\textsuperscript{o} (p.~\pageref{fol:6}), f\textsuperscript{o}~22~r\textsuperscript{o} (p.~\pageref{fol:44}), infra f\textsuperscript{o}~25~r\textsuperscript{o} (p.~\pageref{fol:50}), f\textsuperscript{o}~35~v\textsuperscript{o} (p.~\pageref{fol:71}), \textit{et passim}. — \textbf{15} (\textit{Acheux}) cf. supra f\textsuperscript{o}~21~v\textsuperscript{o} (p.~\pageref{fol:43}), f\textsuperscript{o}~24~r\textsuperscript{o} (p.~\pageref{fol:48}), infra f\textsuperscript{o}~25~r\textsuperscript{o} (p.~\pageref{fol:50}), f\textsuperscript{o}~43~v\textsuperscript{o} (p.~\pageref{fol:87}), \textit{et passim}.\par\vspace{0.55mm}
\noindent\rule{38mm}{0.4pt}\par\vspace{0.5mm}
\noindent\hspace*{5mm}\parbox{\dimexpr\textwidth-5mm\relax}{\fontsize{6.6pt}{7.7pt}\selectfont\color{black!84}Au bivouac de \textbf{Forceville} le dimanche 24 septembre 1916, Hippolyte Combes achève sa hutte, assure une corvée matinale de fourrage vert et prépare ses bagages de permission avant d'écouter la musique militaire sur la place d'Acheux.}\par}
\newpage
\phantomsection\label{fac:49}
% --- En-tête du fac-similé (Page impaire / droite : Photographie du manuscrit) ---
\noindent\rule{\textwidth}{0.5pt}\par\vspace{0.55mm}
\noindent\textbf{\fontsize{6.8pt}{7.8pt}\selectfont FAC-SIMILÉ EN VIS-À-VIS — FOLIO 24 VERSO} \hfill \textbf{\fontsize{6.8pt}{7.8pt}\selectfont SEPTEMBRE 1916}\par\vspace{0.25mm}
\noindent{\fontsize{6.8pt}{8.0pt}\selectfont\textit{Livre-Journal d'Hippolyte Combes} — Dimanche 24 septembre 1916 (Forceville et Acheux-en-Amiénois (Somme))}\par\vspace{0.55mm}
\noindent\rule{\textwidth}{0.35pt}\par\vspace*{\fill}
{\centering
  \setlength{\fboxsep}{1.2mm}%
  \setlength{\fboxrule}{0.4pt}%
  \fcolorbox{gray!55}{white}{%
    \includegraphics[width=0.95\textwidth,height=0.81\textheight,keepaspectratio]{\facsimiledir/20251126_103003.jpg}%
  }\par}
\vspace*{\fill}
\newpage

% =============================================================================
% CHAPITRE III (FOLIO 25 RECTO A FOLIO 33 RECTO)
% =============================================================================
\cleardoubleevenpage
\renewcommand{\currentmonthyear}{Septembre 1916 – Octobre 1916}
\renewcommand{\currentfoliolabel}{Chapitre III}
\chapter{La première permission : Paris et Toulouse}

\noindent{\fontsize{7.8pt}{9.0pt}\selectfont\textit{Folios 25 recto à 33 recto — De Forceville et Acheux à Paris, Neuilly-sur-Seine et Toulouse, puis retour vers la Somme (25 septembre – 6 octobre 1916)}}\par\vspace{1.5mm}
\noindent\rule{\textwidth}{0.5pt}\par\vspace{0.9mm}
\noindent\textbf{\scriptsize\textsc{Résumé analytique du Chapitre III (Folio 25 recto à Folio 33 recto)}}\par\vspace{0.5mm}
\noindent{\fontsize{6.9pt}{8.1pt}\selectfont Le troisième chapitre (f\textsuperscript{o}~25~r\textsuperscript{o} à f\textsuperscript{o}~33~r\textsuperscript{o}) est entièrement consacré à la première permission du carnet, du 25 septembre au 6 octobre 1916. Quittant le plateau de Forceville le lundi 25 septembre sur un fourgon de ravitaillement jusqu'à la gare régulatrice (\textit{railhead}) d'Acheux-Varennes, Hippolyte Combes gagne Amiens puis Paris à l'aube du 26 septembre. Après une visite matinale à Neuilly-sur-Seine chez les cousines Faure, un déjeuner au boulevard Saint-Germain et l'achat au Quartier latin d'ouvrages d'Ernest Zyromski et de Maurice Maeterlinck, il prend le soir même à la gare d'Orsay l'express de Toulouse. Du 27 septembre au 4 octobre, il vit huit journées lumineuses à Toulouse auprès de sa fiancée Marguerite Calmels, logée chez son oncle et sa tante Émile et Élisabeth Crouzet : promenades sur les coteaux de Pouvourville et de Pech-David, achat d'un sautoir en or de fiançailles (295~francs) le 29 septembre, flâneries le long du canal du Midi et au Grand-Rond, offices à la cathédrale Saint-Étienne et rangement de sa bibliothèque dans son appartement toulousain. Reparti le 4 octobre au soir par l'express de Paris (où il admire le 5 octobre l'église Saint-Vincent-de-Paul avant de reprendre le train du front à la gare du Nord), il rejoint le 6 octobre par Doullens et le \textit{railhead} de Canaples la 152\textsuperscript{e} Compagnie A.S.C.}\par\vspace{1.0mm}
\noindent\rule{\textwidth}{0.35pt}\par\vspace{1.4mm}

\noindent\textbf{\scriptsize\textsc{Introduction historique et philologique}}\par\vspace{0.6mm}
{\fontsize{7.1pt}{8.3pt}\selectfont Cette parenthèse civile au cœur de la guerre révèle l'ancrage intellectuel et familial d'Hippolyte Combes à Paris et à Toulouse. Les archives familiales de Millau et Montauban éclairent avec précision le cercle toulousain qui entoure sa fiancée \textbf{Marguerite Calmels} (1892–1974) : l'agent de change \textbf{Émile Crouzet} (1842–1921, \textit{M.~C\textsuperscript{zt}}), son épouse \textbf{Élisabeth Calmels} (tante paternelle de Marguerite) et leur fils l'\textbf{abbé Jean Crouzet}, futur curé de Notre-Dame de la Dalbade, ainsi que les cousines \textbf{Faure} installées au 16, rue d'Armenonville à Neuilly-sur-Seine. Les achats en librairie du 26 septembre témoignent de la fidélité d'Hippolyte Combes à son ancien maître de la Faculté des lettres de Toulouse, \textbf{Ernest Zyromski} (\textit{L'Orgueil humain}, 1904), et de sa prédilection constante pour l'œuvre de \textbf{Maurice Maeterlinck} (\textit{Les Débris de la guerre}, 1916). Enfin, le f\textsuperscript{o}~33~r\textsuperscript{o} (6 octobre 1916) fournit la mention autographe décisive qui permet d'identifier formellement l'unité d'Hippolyte Combes au sein du \textit{18th Divisional Train} : la \textbf{152\textsuperscript{e} Compagnie A.S.C.} (\textit{152nd Company A.S.C.}), affectée au ravitaillement de la \textbf{54\textsuperscript{e} Brigade d'infanterie}.\par}
\vfill
\newpage
% --- Page de droite d'ouverture du chapitre : Carte vectorielle TikZ & Photographies d'archives ---
\noindent\rule{\textwidth}{0.5pt}\par\vspace{0.6mm}
\noindent\textbf{\fontsize{6.8pt}{7.8pt}\selectfont CARTE D'ITINÉRAIRE DU CHAPITRE III} \hfill \textbf{\fontsize{6.8pt}{7.8pt}\selectfont SEPTEMBRE 1916 – OCTOBRE 1916}\par\vspace{0.25mm}
\noindent{\fontsize{6.7pt}{7.8pt}\selectfont\textit{Déplacements, cantonnements et tournées d'Hippolyte Combes (Folio 25 recto à Folio 33 recto)}}\par\vspace{0.55mm}
\noindent\rule{\textwidth}{0.35pt}\par\vspace{1.2mm}
\begin{center}

\begin{tikzpicture}[
  scale=0.96,
  every node/.style={font=\fontsize{6.0pt}{7.0pt}\selectfont},
  loc/.style={circle, fill=black, inner sep=1.1pt},
  step/.style={circle, draw=bordeaux, fill=white, line width=0.75pt, inner sep=1.5pt},
  rail/.style={draw=bordeaux, line width=0.75pt, ->, >=stealth}
]
\draw[black!50, line width=0.4pt] (-60mm,-42mm) rectangle (60mm,42mm);
\fill[black!3] (-60mm,-42mm) rectangle (60mm,42mm);

% Partie gauche : Grand axe ferroviaire Somme - Paris - Toulouse (-58mm à -2mm)
\node[font=\fontsize{6.3pt}{7.2pt}\selectfont\bfseries] at (-30mm, 37mm) {A. Trajet ferroviaire (25 sept. – 6 oct. 1916)};
\node[step] (acheux) at (-32mm, 30mm) {};
\node[right] at (acheux) {\textbf{1. Forceville / Acheux} (25 sept.)};
\node[loc] (amiens) at (-36mm, 23mm) {};
\node[left] at (amiens) {Amiens \& Creil};
\node[step] (paris) at (-32mm, 14mm) {};
\node[right] at (paris) {\textbf{2. Paris \& Neuilly} (26 sept. / 5 oct.)};
\node[loc] (limoges) at (-36mm, -2mm) {};
\node[left] at (limoges) {Limoges, Brive, Cahors};
\node[loc] (montauban) at (-33mm, -16mm) {};
\node[left] at (montauban) {Montauban};
\node[step] (toulouse) at (-30mm, -28mm) {};
\node[right] at (toulouse) {\textbf{3. Toulouse} (27 sept. – 4 oct.)};
\node[step] (canaples) at (-18mm, 26mm) {};
\node[right] at (canaples) {\textbf{4. Doullens / Canaples} (6 oct.)};

\draw[rail] (acheux) -- (amiens) -- (paris);
\draw[rail] (paris) -- (limoges) -- (montauban) -- (toulouse);
\draw[rail, dashed] (toulouse) to[bend right=22] (paris);
\draw[rail, dashed] (paris) to[bend right=15] (canaples);

% Filet vertical de séparation
\draw[black!40, line width=0.4pt] (-1mm,-38mm) -- (-1mm,38mm);

% Partie droite : Plan vectoriel de Toulouse (0mm à 58mm)
\node[font=\fontsize{6.3pt}{7.2pt}\selectfont\bfseries] at (29mm, 37mm) {B. Plan des lieux fréquentés à Toulouse};
% La Garonne et le Canal du Midi
\draw[black!35, line width=1.6pt] (6mm,-34mm) .. controls (12mm,-10mm) and (14mm,10mm) .. (6mm,30mm)
  node[pos=0.18, right, font=\fontsize{5.6pt}{6.5pt}\selectfont\itshape]{La Garonne};
\draw[black!45, line width=0.8pt] (28mm,32mm) .. controls (40mm,18mm) and (42mm,-4mm) .. (36mm,-28mm)
  node[pos=0.5, right, font=\fontsize{5.5pt}{6.5pt}\selectfont\itshape]{Canal du Midi};

\node[loc] (matabiau) at (34mm, 26mm) {};
\node[left] at (matabiau) {Gare Matabiau};
\node[step] (taur) at (20mm, 18mm) {};
\node[right] at (taur) {\textbf{Rue du Taur} (n\textsuperscript{o}~47) \& St-Sernin};
\node[loc] (capitole) at (19mm, 11mm) {};
\node[right] at (capitole) {Capitole \& Place Esquirol};
\node[step] (pharaon) at (17mm, 3mm) {};
\node[right] at (pharaon) {\textbf{Rue Pharaon} (famille Crouzet)};
\node[step] (stetienne) at (28mm, 6mm) {};
\node[above right] at (stetienne) {\textbf{Cath. Saint-Étienne}};
\node[loc] (grandrond) at (32mm, -2mm) {};
\node[left] at (grandrond) {Grand-Rond};
\node[loc] (stemarie) at (40mm, -8mm) {};
\node[right] at (stemarie) {\textbf{Ste-Marie-des-Champs}};
\node[loc] (pechdavid) at (18mm, -26mm) {};
\node[right] at (pechdavid) {\textbf{Pouvourville \& Pech-David} (28 sept.)};
\end{tikzpicture}

\end{center}
\vfill
\noindent\rule{\textwidth}{0.3pt}\par\vspace{1.2mm}
\noindent\begin{center}
  \begin{minipage}[b]{0.42\textwidth}
    \centering
    \setlength{\fboxsep}{0.7mm}\setlength{\fboxrule}{0.4pt}%
    \fcolorbox{gray!55}{white}{\includegraphics[width=0.94\linewidth,height=28mm,keepaspectratio]{illustrations/marguerite_1917.jpg}}
    \par\vspace{0.6mm}{\fontsize{5.8pt}{6.7pt}\selectfont Marguerite Calmels (1892–1974), photographiée par le studio Jublin à Millau à l'époque de ses fiançailles avec Hippolyte Combes.}
  \end{minipage}\hfill
  \begin{minipage}[b]{0.46\textwidth}
    \centering
    \setlength{\fboxsep}{0.7mm}\setlength{\fboxrule}{0.4pt}%
    \fcolorbox{gray!55}{white}{\includegraphics[width=0.94\linewidth,height=28mm,keepaspectratio]{illustrations/maison_millau.jpg}}
    \par\vspace{0.6mm}{\fontsize{5.8pt}{6.7pt}\selectfont La demeure familiale de Millau (Aveyron), berceau de la famille Combes.}
  \end{minipage}
\end{center}
\vspace*{1mm}
\newpage

% =============================================================================
% VIS-A-VIS SPREAD : FOLIO 25 RECTO (20251126_103019.jpg)
% =============================================================================
\cleardoubleevenpage
\renewcommand{\currentmonthyear}{Septembre 1916}
\renewcommand{\currentfoliolabel}{Folio 25 recto}
\phantomsection\label{fol:50}
% --- En-tête du folio (Page paire / gauche : Transcription, apparat & registres de bas de page) ---
\noindent\rule{\textwidth}{0.5pt}\par\vspace{0.55mm}
\noindent\textbf{\footnotesize FOLIO 25 RECTO} \hfill \textbf{\footnotesize SEPTEMBRE 1916}\par\vspace{0.25mm}
\noindent{\fontsize{6.9pt}{8.0pt}\selectfont Lundi 25 septembre 1916 — Réveil à 7 h, corvée de green forage à Forceville avec le Lt. Ford et départ en permission à la gare d'Acheux-Varennes}\par\vspace{0.2mm}
\noindent{\fontsize{6.9pt}{8.0pt}\selectfont\textit{Lundi 25 septembre 1916 (matin et après-midi)} \hfill \textit{Forceville et gare d'Acheux-Varennes (Somme)}}\par\vspace{0.5mm}
\noindent\rule{\textwidth}{0.35pt}\par\vspace{0.8mm}
% --- Corps de la transcription diplomatique et apparat marginal en prolongement ---
{\renewcommand{\arraystretch}{0.87}%
\noindent\hspace*{3.5mm}\begin{tabular}{@{}p{87mm}@{\hspace{3.5mm}}p{39mm}@{}}
      \leavevmode\fontsize{6.9pt}{7.9pt}\selectfont\strut \underline{25 L.} \hfill{} \underline{Sept /16}\strut & \smash{\parbox[t]{39mm}{\fontsize{6.2pt}{7.0pt}\selectfont\raggedright\color{black!88}\textbf{1 L.} \textit{scil.} Lundi.}} \\
      \leavevmode\fontsize{6.9pt}{7.9pt}\selectfont\strut Vers 7 H, sur pied.\strut &  \\
      \leavevmode\fontsize{6.9pt}{7.9pt}\selectfont\strut Au firmament bleuissant, léger croissant lunaire\strut &  \\
      \leavevmode\fontsize{6.9pt}{7.9pt}\selectfont\strut et 2 étoiles très brillantes.\strut &  \\
      \leavevmode\llap{\fontsize{6.4pt}{7.3pt}\selectfont\color{darkgray}5\hspace{1.8mm}}\fontsize{6.9pt}{7.9pt}\selectfont\strut Temps superbe.\strut &  \\
      \leavevmode\fontsize{6.9pt}{7.9pt}\selectfont\strut Bientôt, entiert prêt à partir.\strut & \smash{\parbox[t]{39mm}{\fontsize{6.2pt}{7.0pt}\selectfont\raggedright\color{black!88}\textbf{6 entiert} \textit{abbr. pro} entièrement. — \textbf{7 tt} \textit{abbr. pro} tout. — \textbf{8 Lt} \textit{abbr. pro} Lieutenant. \textbf{8 green} \textit{angl. in ms.} — \textbf{9 forage} \textit{angl. in ms.}}} \\
      \leavevmode\fontsize{6.9pt}{7.9pt}\selectfont\strut Promené champs tt en haut plateau.\strut &  \\
      \leavevmode\fontsize{6.9pt}{7.9pt}\selectfont\strut Dû venir à Forceville av. Lt Ford, for green\strut &  \\
      \leavevmode\fontsize{6.9pt}{7.9pt}\selectfont\strut forage. Chaleur très forte.\strut &  \\
      \leavevmode\llap{\fontsize{6.4pt}{7.3pt}\selectfont\color{darkgray}10\hspace{1.8mm}}\fontsize{6.9pt}{7.9pt}\selectfont\strut A midi enfin, monté sur le dernier wagon\strut &  \\
      \leavevmode\fontsize{6.9pt}{7.9pt}\selectfont\strut du convoy allant au railhead, puis continué\strut & \smash{\parbox[t]{39mm}{\fontsize{6.2pt}{7.0pt}\selectfont\raggedright\color{black!88}\textbf{11 railhead} \textit{angl. in ms.}}} \\
      \leavevmode\fontsize{6.9pt}{7.9pt}\selectfont\strut à travers chps j'à station d'Acheux-Varennes\strut & \smash{\parbox[t]{39mm}{\fontsize{6.2pt}{7.0pt}\selectfont\raggedright\color{black!88}\textbf{12 j'à} \textit{abbr. pro} jusqu'à. \textbf{12 chps} \textit{abbr. pro} champs. — \textbf{13 R.T.O.} \textit{angl.} Railway Transport Officer. — \textbf{15 ds} \textit{abbr. pro} dans. \textbf{15 lgt} \textit{abbr. pro} longtemps.}} \\
      \leavevmode\fontsize{6.9pt}{7.9pt}\selectfont\strut Laissé mes colis à l'office de l'R.T.O. \& venu\strut &  \\
      \leavevmode\fontsize{6.9pt}{7.9pt}\selectfont\strut village, déjeuné.\strut &  \\
      \leavevmode\llap{\fontsize{6.4pt}{7.3pt}\selectfont\color{darkgray}15\hspace{1.8mm}}\fontsize{6.9pt}{7.9pt}\selectfont\strut Puis attendu lgt. ds mon compartimt petit\strut &  \\
      \leavevmode\fontsize{6.9pt}{7.9pt}\selectfont\strut train départemental en gare —\strut &  \\
      \leavevmode\fontsize{6.9pt}{7.9pt}\selectfont\strut chaleur intense\strut &  \\
      \cline{1-1}
      \leavevmode\fontsize{6.9pt}{7.9pt}\selectfont\strut [ Tip à mon groom : 2 f — Biscuits : 0.50 c\strut & \smash{\parbox[t]{39mm}{\fontsize{6.2pt}{7.0pt}\selectfont\raggedright\color{black!88}\textbf{18} \textit{in marg. cum unco} (comptes). \textbf{18 groom} \textit{angl. in ms.} — \textbf{19 Lunch} \textit{angl. in ms.}}} \\
      \leavevmode\fontsize{6.9pt}{7.9pt}\selectfont\strut Lunch à Acheux : 1 f 50 ]\strut &  \\
\end{tabular}}%
\par\vspace*{\fill}
% --- Registres de bas de page (1. Critique interne [si pertinent] ; 2. Loci paralleli en latin ; 3. Commentaire historique en retrait) ---
\noindent\rule{\textwidth}{0.5pt}\par\vspace{0.65mm}
{\fontsize{6.6pt}{7.7pt}\selectfont%
\noindent\normalfont\color{black!92} \textbf{8} (\textit{Forceville}) cf. supra f\textsuperscript{o}~22~v\textsuperscript{o} (p.~\pageref{fol:45}), f\textsuperscript{o}~24~v\textsuperscript{o} (p.~\pageref{fol:49}), infra f\textsuperscript{o}~48~r\textsuperscript{o} (p.~\pageref{fol:96}). — \textbf{12} (\textit{Acheux-Varennes}) cf. supra f\textsuperscript{o}~21~v\textsuperscript{o} (p.~\pageref{fol:43}), f\textsuperscript{o}~24~v\textsuperscript{o} (p.~\pageref{fol:49}), infra f\textsuperscript{o}~43~v\textsuperscript{o} (p.~\pageref{fol:87}), f\textsuperscript{o}~44~v\textsuperscript{o} (p.~\pageref{fol:89}), \textit{et passim}. — \textbf{12} (\textit{Varennes}) cf. infra f\textsuperscript{o}~48~r\textsuperscript{o} (p.~\pageref{fol:96}), f\textsuperscript{o}~48~v\textsuperscript{o} (p.~\pageref{fol:97}), \textit{et passim}.\par\vspace{0.55mm}
\noindent\rule{38mm}{0.4pt}\par\vspace{0.5mm}
\noindent\hspace*{5mm}\parbox{\dimexpr\textwidth-5mm\relax}{\fontsize{6.6pt}{7.7pt}\selectfont\color{black!84}Ce feuillet du lundi 25 septembre 1916 (omis dans le premier document de transcription mais conservé à sa place dans la séquence photographique) relate le départ d'Hippolyte Combes pour sa première permission: après une corvée matinale à Forceville avec le \textbf{lieutenant Ford}, il gagne à midi en fourgon la tête de ligne ferroviaire (\textit{railhead}) d'\textbf{Acheux-Varennes}, confie ses bagages au bureau du \textit{Railway Transport Officer} (\textit{R.T.O.}) et prend le train départemental à voie métrique vers Gézaincourt.}\par}
\newpage
\phantomsection\label{fac:50}
% --- En-tête du fac-similé (Page impaire / droite : Photographie du manuscrit) ---
\noindent\rule{\textwidth}{0.5pt}\par\vspace{0.55mm}
\noindent\textbf{\fontsize{6.8pt}{7.8pt}\selectfont FAC-SIMILÉ EN VIS-À-VIS — FOLIO 25 RECTO} \hfill \textbf{\fontsize{6.8pt}{7.8pt}\selectfont SEPTEMBRE 1916}\par\vspace{0.25mm}
\noindent{\fontsize{6.8pt}{8.0pt}\selectfont\textit{Livre-Journal d'Hippolyte Combes} — Lundi 25 septembre 1916 (matin et après-midi) (Forceville et gare d'Acheux-Varennes (Somme))}\par\vspace{0.55mm}
\noindent\rule{\textwidth}{0.35pt}\par\vspace*{\fill}
{\centering
  \setlength{\fboxsep}{1.2mm}%
  \setlength{\fboxrule}{0.4pt}%
  \fcolorbox{gray!55}{white}{%
    \includegraphics[width=0.95\textwidth,height=0.81\textheight,keepaspectratio]{\facsimiledir/20251126_103019.jpg}%
  }\par}
\vspace*{\fill}
\newpage

% =============================================================================
% VIS-A-VIS SPREAD : FOLIO 25 VERSO (20251126_103036.jpg)
% =============================================================================
\cleardoubleevenpage
\renewcommand{\currentmonthyear}{Septembre 1916}
\renewcommand{\currentfoliolabel}{Folio 25 verso}
\phantomsection\label{fol:51}
% --- En-tête du folio (Page paire / gauche : Transcription, apparat & registres de bas de page) ---
\noindent\rule{\textwidth}{0.5pt}\par\vspace{0.55mm}
\noindent\textbf{\footnotesize FOLIO 25 VERSO} \hfill \textbf{\footnotesize SEPTEMBRE 1916}\par\vspace{0.25mm}
\noindent{\fontsize{6.9pt}{8.0pt}\selectfont De Gézaincourt à Amiens et raid nocturne de « Taubes » sur la gare de Longueau}\par\vspace{0.2mm}
\noindent{\fontsize{6.9pt}{8.0pt}\selectfont\textit{Lundi 25 septembre 1916 (soir et nuit)} \hfill \textit{Gézaincourt, Amiens, Longueau et train vers Paris}}\par\vspace{0.5mm}
\noindent\rule{\textwidth}{0.35pt}\par\vspace{0.8mm}
% --- Corps de la transcription diplomatique et apparat marginal en prolongement ---
{\renewcommand{\arraystretch}{0.87}%
\noindent\hspace*{3.5mm}\begin{tabular}{@{}p{87mm}@{\hspace{3.5mm}}p{39mm}@{}}
      \leavevmode\fontsize{6.9pt}{7.9pt}\selectfont\strut \underline{25 Lundi.} (Suite) \hfill{} \raisebox{0.5ex}{\underline{Sept / 16}}\strut &  \\
      \leavevmode\fontsize{6.9pt}{7.9pt}\selectfont\strut au ciel bleu, qq nuages, nettement détachés.\strut & \smash{\parbox[t]{39mm}{\fontsize{6.2pt}{7.0pt}\selectfont\raggedright\color{black!88}\textbf{2 qq} \textit{abbr. pro} quelque(s).}} \\
      \leavevmode\fontsize{6.9pt}{7.9pt}\selectfont\strut Enfin, vers 3H 1/2, départ\strut &  \\
      \leavevmode\fontsize{6.9pt}{7.9pt}\selectfont\strut Doucement bercé par mouvement train. Et caresse\strut &  \\
      \leavevmode\llap{\fontsize{6.4pt}{7.3pt}\selectfont\color{darkgray}5\hspace{1.8mm}}\fontsize{6.9pt}{7.9pt}\selectfont\strut brise légère\strut &  \\
      \leavevmode\fontsize{6.9pt}{7.9pt}\selectfont\strut Verdures étincelantes de lumière. Bois ensoleillés.\strut &  \\
      \leavevmode\fontsize{6.9pt}{7.9pt}\selectfont\strut A Gézaincourt, vers 4H 1/2, changé train.\strut &  \\
      \leavevmode\fontsize{6.9pt}{7.9pt}\selectfont\strut Reparti à 5 H ds le train venant de Doullens\strut & \smash{\parbox[t]{39mm}{\fontsize{6.2pt}{7.0pt}\selectfont\raggedright\color{black!88}\textbf{8 ds} \textit{abbr. pro} dans.}} \\
      \leavevmode\fontsize{6.9pt}{7.9pt}\selectfont\strut \underline{A Amiens}, vers 8H s\strut & \smash{\parbox[t]{39mm}{\fontsize{6.2pt}{7.0pt}\selectfont\raggedright\color{black!88}\textbf{9 s.} \textit{abbr. pro} soir.}} \\
      \leavevmode\llap{\fontsize{6.4pt}{7.3pt}\selectfont\color{darkgray}10\hspace{1.8mm}}\fontsize{6.9pt}{7.9pt}\selectfont\strut Monté train de permissionns pr Paris. Compartimt\strut & \smash{\parbox[t]{39mm}{\fontsize{6.2pt}{7.0pt}\selectfont\raggedright\color{black!88}\textbf{10 pr} \textit{abbr. pro} pour.}} \\
      \leavevmode\fontsize{6.9pt}{7.9pt}\selectfont\strut de 3e Cl., avec grossiers poilus, barbus \& sales\strut &  \\
      \leavevmode\fontsize{6.9pt}{7.9pt}\selectfont\strut Roulé tte nuit\strut & \smash{\parbox[t]{39mm}{\fontsize{6.2pt}{7.0pt}\selectfont\raggedright\color{black!88}\textbf{12 tte} \textit{abbr. pro} toute.}} \\
      \leavevmode\fontsize{6.9pt}{7.9pt}\selectfont\strut Après avoir quitté Amiens, et au croisement\strut &  \\
      \leavevmode\fontsize{6.9pt}{7.9pt}\selectfont\strut de Longueau, survolé par des Taubes qui\strut &  \\
      \leavevmode\llap{\fontsize{6.4pt}{7.3pt}\selectfont\color{darkgray}15\hspace{1.8mm}}\fontsize{6.9pt}{7.9pt}\selectfont\strut jeté bombes — obus ds les airs. Ciel fouillé\strut & \smash{\parbox[t]{39mm}{\fontsize{6.2pt}{7.0pt}\selectfont\raggedright\color{black!88}\textbf{15 ds} \textit{abbr. pro} dans.}} \\
      \leavevmode\fontsize{6.9pt}{7.9pt}\selectfont\strut par les projecteurs. Très beau spectacle\strut &  \\
      \leavevmode\fontsize{6.9pt}{7.9pt}\selectfont\strut Pas descendu à Creil (en retard [?])\strut & \smash{\parbox[t]{39mm}{\fontsize{6.2pt}{7.0pt}\selectfont\raggedright\color{black!88}\textbf{17 retard} \textit{lect. incert. in ms.} — \textbf{18 tte} \textit{abbr. pro} toute.}} \\
      \leavevmode\fontsize{6.9pt}{7.9pt}\selectfont\strut Voyage tte la nuit\strut &  \\
\end{tabular}}%
\par\vspace*{\fill}
% --- Registres de bas de page (1. Critique interne [si pertinent] ; 2. Loci paralleli en latin ; 3. Commentaire historique en retrait) ---
\noindent\rule{\textwidth}{0.5pt}\par\vspace{0.65mm}
{\fontsize{6.6pt}{7.7pt}\selectfont%
\noindent\normalfont\color{black!92} \textbf{8} (\textit{Doullens}) cf. infra f\textsuperscript{o}~32~v\textsuperscript{o} (p.~\pageref{fol:65}), f\textsuperscript{o}~33~r\textsuperscript{o} (p.~\pageref{fol:66}), \textit{et passim}. — \textbf{9} (\textit{Amiens}) cf. infra f\textsuperscript{o}~32~v\textsuperscript{o} (p.~\pageref{fol:65}), f\textsuperscript{o}~44~v\textsuperscript{o} (p.~\pageref{fol:89}), \textit{et passim}. — \textbf{10} (\textit{Paris}) cf. supra in tegumento (p.~\pageref{fol:1}), infra f\textsuperscript{o}~26~r\textsuperscript{o} (p.~\pageref{fol:52}), f\textsuperscript{o}~26~v\textsuperscript{o} (p.~\pageref{fol:53}), \textit{et passim}.\par\vspace{0.55mm}
\noindent\rule{38mm}{0.4pt}\par\vspace{0.5mm}
\noindent\hspace*{5mm}\parbox{\dimexpr\textwidth-5mm\relax}{\fontsize{6.6pt}{7.7pt}\selectfont\color{black!84}Après la correspondance à la gare de bifurcation de \textbf{Gézaincourt} (au sud de Doullens) et le départ d'\textbf{Amiens} à 20~h dans le train de permissionnaires vers Paris, le convoi est pris au triage ferroviaire de \textbf{Longueau} (Somme) sous un raid nocturne d'avions allemands (désignés sous le terme générique de \textit{Taubes}) combattus par les projecteurs et la défense contre avions.}\par}
\newpage
\phantomsection\label{fac:51}
% --- En-tête du fac-similé (Page impaire / droite : Photographie du manuscrit) ---
\noindent\rule{\textwidth}{0.5pt}\par\vspace{0.55mm}
\noindent\textbf{\fontsize{6.8pt}{7.8pt}\selectfont FAC-SIMILÉ EN VIS-À-VIS — FOLIO 25 VERSO} \hfill \textbf{\fontsize{6.8pt}{7.8pt}\selectfont SEPTEMBRE 1916}\par\vspace{0.25mm}
\noindent{\fontsize{6.8pt}{8.0pt}\selectfont\textit{Livre-Journal d'Hippolyte Combes} — Lundi 25 septembre 1916 (soir et nuit) (Gézaincourt, Amiens, Longueau et train vers Paris)}\par\vspace{0.55mm}
\noindent\rule{\textwidth}{0.35pt}\par\vspace*{\fill}
{\centering
  \setlength{\fboxsep}{1.2mm}%
  \setlength{\fboxrule}{0.4pt}%
  \fcolorbox{gray!55}{white}{%
    \includegraphics[width=0.95\textwidth,height=0.81\textheight,keepaspectratio]{\facsimiledir/20251126_103036.jpg}%
  }\par}
\vspace*{\fill}
\newpage

% =============================================================================
% VIS-A-VIS SPREAD : FOLIO 26 RECTO (20251126_103109.jpg)
% =============================================================================
\cleardoubleevenpage
\renewcommand{\currentmonthyear}{Septembre 1916}
\renewcommand{\currentfoliolabel}{Folio 26 recto}
\phantomsection\label{fol:52}
% --- En-tête du folio (Page paire / gauche : Transcription, apparat & registres de bas de page) ---
\noindent\rule{\textwidth}{0.5pt}\par\vspace{0.55mm}
\noindent\textbf{\footnotesize FOLIO 26 RECTO} \hfill \textbf{\footnotesize SEPTEMBRE 1916}\par\vspace{0.25mm}
\noindent{\fontsize{6.9pt}{8.0pt}\selectfont Arrivée matinale au Bourget, traversée de Paris et démarches aux Invalides}\par\vspace{0.2mm}
\noindent{\fontsize{6.9pt}{8.0pt}\selectfont\textit{Mardi 26 septembre 1916 (matin et après-midi)} \hfill \textit{Le Bourget et Paris (Gare du Nord, Tuileries, Orsay, Invalides, rive gauche)}}\par\vspace{0.5mm}
\noindent\rule{\textwidth}{0.35pt}\par\vspace{0.8mm}
% --- Corps de la transcription diplomatique et apparat marginal en prolongement ---
{\renewcommand{\arraystretch}{0.87}%
\noindent\hspace*{3.5mm}\begin{tabular}{@{}p{87mm}@{\hspace{3.5mm}}p{39mm}@{}}
      \leavevmode\fontsize{6.9pt}{7.9pt}\selectfont\strut \underline{Mardi 26} \hfill{} \underline{Sept 16}\strut &  \\
      \leavevmode\fontsize{6.9pt}{7.9pt}\selectfont\strut Vers 6H m., encore nuit, dû descendre gare\strut & \smash{\parbox[t]{39mm}{\fontsize{6.2pt}{7.0pt}\selectfont\raggedright\color{black!88}\textbf{2 m.} \textit{abbr. pro} matin.}} \\
      \leavevmode\fontsize{6.9pt}{7.9pt}\selectfont\strut du Bourget.\strut &  \\
      \leavevmode\fontsize{6.9pt}{7.9pt}\selectfont\strut Suivi voies j'que bureau militre. Là\strut & \smash{\parbox[t]{39mm}{\fontsize{6.2pt}{7.0pt}\selectfont\raggedright\color{black!88}\textbf{4 j'que} \textit{abbr. pro} jusque. \textbf{4 militre} \textit{abbr. pro} militaire. — \textbf{5 autorison} \textit{abbr. pro} autorisation. — \textbf{6 j'à} \textit{abbr. pro} jusqu'à. — \textbf{8 j-à} \textit{abbr. pro} jusqu'à.}} \\
      \leavevmode\llap{\fontsize{6.4pt}{7.3pt}\selectfont\color{darkgray}5\hspace{1.8mm}}\fontsize{6.9pt}{7.9pt}\selectfont\strut obtenu l'autorison de sortir et accompagné\strut &  \\
      \leavevmode\fontsize{6.9pt}{7.9pt}\selectfont\strut par un planton j'à la porte\strut &  \\
      \leavevmode\fontsize{6.9pt}{7.9pt}\selectfont\strut Du Bourget, sortie de la rue de Flandre\strut &  \\
      \leavevmode\fontsize{6.9pt}{7.9pt}\selectfont\strut pris tramway qui m'a conduit j-à gare\strut &  \\
      \leavevmode\fontsize{6.9pt}{7.9pt}\selectfont\strut du Nord. De là, par métro, arrivé\strut &  \\
      \leavevmode\llap{\fontsize{6.4pt}{7.3pt}\selectfont\color{darkgray}10\hspace{1.8mm}}\fontsize{6.9pt}{7.9pt}\selectfont\strut à Station Tuileries. (changé au Chatelet)\strut &  \\
      \leavevmode\fontsize{6.9pt}{7.9pt}\selectfont\strut Traversé jardin et\strut &  \\
      \leavevmode\fontsize{6.9pt}{7.9pt}\selectfont\strut Venu gare d'Orsay\strut &  \\
      \leavevmode\fontsize{6.9pt}{7.9pt}\selectfont\strut Me suis rendu à la Place (Hotel)\strut &  \\
      \leavevmode\fontsize{6.9pt}{7.9pt}\selectfont\strut des Invalides, pr demande autorison de\strut & \smash{\parbox[t]{39mm}{\fontsize{6.2pt}{7.0pt}\selectfont\raggedright\color{black!88}\textbf{14 pr} \textit{abbr. pro} pour. \textbf{14 autorison} \textit{abbr. pro} autorisation.}} \\
      \leavevmode\llap{\fontsize{6.4pt}{7.3pt}\selectfont\color{darkgray}15\hspace{1.8mm}}\fontsize{6.9pt}{7.9pt}\selectfont\strut voyager à mes frais en express. Revenu\strut &  \\
      \leavevmode\fontsize{6.9pt}{7.9pt}\selectfont\strut gare et pu enfin prendre billet de 2e cl.\strut &  \\
      \leavevmode\fontsize{6.9pt}{7.9pt}\selectfont\strut Toulouse\strut &  \\
      \leavevmode\fontsize{6.9pt}{7.9pt}\selectfont\strut Puis descendu rue du Bac\strut &  \\
      \leavevmode\fontsize{6.9pt}{7.9pt}\selectfont\strut Au Bon Marché\strut &  \\
      \leavevmode\llap{\fontsize{6.4pt}{7.3pt}\selectfont\color{darkgray}20\hspace{1.8mm}}\fontsize{6.9pt}{7.9pt}\selectfont\strut Rue de Rennes\strut &  \\
      \leavevmode\fontsize{6.9pt}{7.9pt}\selectfont\strut Jardin du Luxembourg\strut &  \\
      \leavevmode\fontsize{6.9pt}{7.9pt}\selectfont\strut Chez Colin\strut &  \\
\end{tabular}}%
\par\vspace*{\fill}
% --- Registres de bas de page (1. Critique interne [si pertinent] ; 2. Loci paralleli en latin ; 3. Commentaire historique en retrait) ---
\noindent\rule{\textwidth}{0.5pt}\par\vspace{0.65mm}
{\fontsize{6.6pt}{7.7pt}\selectfont%
\noindent\normalfont \textbf{Continuité :} Première page consacrée à la journée parisienne du mardi 26 septembre 1916 (arrivée à la gare du Nord et matinée à Neuilly), poursuivie au f\textsuperscript{o}~26~v\textsuperscript{o}.\par\vspace{0.5mm}
\noindent\rule{22mm}{0.3pt}\par\vspace{0.45mm}
\noindent\normalfont\color{black!92} \textbf{17} (\textit{Toulouse}) cf. supra in tegumento (p.~\pageref{fol:1}), f\textsuperscript{o}~23~v\textsuperscript{o} (p.~\pageref{fol:47}), infra f\textsuperscript{o}~26~v\textsuperscript{o} (p.~\pageref{fol:53}), f\textsuperscript{o}~27~r\textsuperscript{o} (p.~\pageref{fol:54}), \textit{et passim}, Appendicem I, doc.~1 et 3 (p.~\pageref{annexe:doc1}).\par\vspace{0.55mm}
\noindent\rule{38mm}{0.4pt}\par\vspace{0.5mm}
\noindent\hspace*{5mm}\parbox{\dimexpr\textwidth-5mm\relax}{\fontsize{6.6pt}{7.7pt}\selectfont\color{black!84}Débarqué à 6~h du matin à la gare militaire de triage du \textbf{Bourget}, Hippolyte Combes gagne la gare du Nord en tramway puis la \textbf{gare d'Orsay} (tête de ligne de la Compagnie du Paris-Orléans vers Toulouse) et le gouvernement militaire de Paris à l'\textbf{Hôtel des Invalides} pour obtenir l'autorisation de voyager à ses frais en train express (2\textsuperscript{e} classe, au quart de tarif militaire), avant de parcourir la rive gauche (rue du Bac, Bon Marché, rue de Rennes, jardin du Luxembourg et librairie Armand Colin, 103 boulevard Saint-Michel).}\par}
\newpage
\phantomsection\label{fac:52}
% --- En-tête du fac-similé (Page impaire / droite : Photographie du manuscrit) ---
\noindent\rule{\textwidth}{0.5pt}\par\vspace{0.55mm}
\noindent\textbf{\fontsize{6.8pt}{7.8pt}\selectfont FAC-SIMILÉ EN VIS-À-VIS — FOLIO 26 RECTO} \hfill \textbf{\fontsize{6.8pt}{7.8pt}\selectfont SEPTEMBRE 1916}\par\vspace{0.25mm}
\noindent{\fontsize{6.8pt}{8.0pt}\selectfont\textit{Livre-Journal d'Hippolyte Combes} — Mardi 26 septembre 1916 (matin et après-midi) (Le Bourget et Paris (Gare du Nord, Tuileries, Orsay, Invalides, rive gauche))}\par\vspace{0.55mm}
\noindent\rule{\textwidth}{0.35pt}\par\vspace*{\fill}
{\centering
  \setlength{\fboxsep}{1.2mm}%
  \setlength{\fboxrule}{0.4pt}%
  \fcolorbox{gray!55}{white}{%
    \includegraphics[width=0.95\textwidth,height=0.81\textheight,keepaspectratio]{\facsimiledir/20251126_103109.jpg}%
  }\par}
\vspace*{\fill}
\newpage

% =============================================================================
% VIS-A-VIS SPREAD : FOLIO 26 VERSO (20251126_103131.jpg)
% =============================================================================
\cleardoubleevenpage
\renewcommand{\currentmonthyear}{Septembre 1916}
\renewcommand{\currentfoliolabel}{Folio 26 verso}
\phantomsection\label{fol:53}
% --- En-tête du folio (Page paire / gauche : Transcription, apparat & registres de bas de page) ---
\noindent\rule{\textwidth}{0.5pt}\par\vspace{0.55mm}
\noindent\textbf{\footnotesize FOLIO 26 VERSO} \hfill \textbf{\footnotesize SEPTEMBRE 1916}\par\vspace{0.25mm}
\noindent{\fontsize{6.9pt}{8.0pt}\selectfont Mardi 26 septembre 1916 (suite) — Square de la Tour Saint-Jacques, rue de Rivoli, achats d'ouvrages d'Ernest Zyromski et de Maurice Maeterlinck chez Armand Colin et départ à la gare d'Orsay}\par\vspace{0.2mm}
\noindent{\fontsize{6.9pt}{8.0pt}\selectfont\textit{Mardi 26 septembre 1916 (après-midi et soir)} \hfill \textit{Paris (Tour Saint-Jacques, rue de Rivoli, librairie Colin, gare d'Orsay)}}\par\vspace{0.5mm}
\noindent\rule{\textwidth}{0.35pt}\par\vspace{0.8mm}
% --- Corps de la transcription diplomatique et apparat marginal en prolongement ---
{\renewcommand{\arraystretch}{0.87}%
\noindent\hspace*{3.5mm}\begin{tabular}{@{}p{87mm}@{\hspace{3.5mm}}p{39mm}@{}}
      \leavevmode\fontsize{6.9pt}{7.9pt}\selectfont\strut \underline{Mardi 26} \hfill{} \underline{Sept 16}\strut &  \\
      \leavevmode\fontsize{6.9pt}{7.9pt}\selectfont\strut Déjeuné – Bd St-Germain\strut & \smash{\parbox[t]{39mm}{\fontsize{6.2pt}{7.0pt}\selectfont\raggedright\color{black!88}\textbf{2 Bd} \textit{abbr. pro} Boulevard. \textbf{2 St} \textit{abbr. pro} Saint.}} \\
      \leavevmode\fontsize{6.9pt}{7.9pt}\selectfont\strut Reposé – jardin Notre - Dame\strut &  \\
      \leavevmode\fontsize{6.9pt}{7.9pt}\selectfont\strut Magasins de l'Hotel - de - Ville\strut &  \\
      \leavevmode\llap{\fontsize{6.4pt}{7.3pt}\selectfont\color{darkgray}5\hspace{1.8mm}}\fontsize{6.9pt}{7.9pt}\selectfont\strut [Lgt/1/2 h ?] assis banc ds le square de la\strut & \smash{\parbox[t]{39mm}{\fontsize{6.2pt}{7.0pt}\selectfont\raggedright\color{black!88}\textbf{5} \textit{lect. incert. in ms.} \textbf{5 ds} \textit{abbr. pro} dans. — \textbf{6 mouvt} \textit{abbr. pro} mouvement. \textbf{6 St} \textit{abbr. pro} Saint. — \textbf{7 ordinre} \textit{abbr. pro} ordinaire. — \textbf{8 lentemt} \textit{abbr. pro} lentement.}} \\
      \leavevmode\fontsize{6.9pt}{7.9pt}\selectfont\strut Tour St-Jacques — Gros mouvt de rue de\strut &  \\
      \leavevmode\fontsize{6.9pt}{7.9pt}\selectfont\strut Rivoli. Animation ordinre. —\strut &  \\
      \leavevmode\fontsize{6.9pt}{7.9pt}\selectfont\strut Après diner, venu lentemt, par les\strut &  \\
      \leavevmode\fontsize{6.9pt}{7.9pt}\selectfont\strut quais à gare d'Orsay\strut &  \\
      \leavevmode\llap{\fontsize{6.4pt}{7.3pt}\selectfont\color{darkgray}10\hspace{1.8mm}}\fontsize{6.9pt}{7.9pt}\selectfont\strut Départ train vers 8H S\strut &  \\
      \leavevmode\fontsize{6.9pt}{7.9pt}\selectfont\strut Tramway : Bourget - gare du Nord : 50c\strut &  \\
      \leavevmode\fontsize{6.9pt}{7.9pt}\selectfont\strut Métro : 1 f. — Envoi[s?] : 10 c, 1 f, 2 f. ([mot incertain])\strut & \smash{\parbox[t]{39mm}{\fontsize{6.2pt}{7.0pt}\selectfont\raggedright\color{black!88}\textbf{12} \textit{lect. incert. in ms.} \textbf{12 f.} \textit{abbr. pro} franc(s). — \textbf{13 chvx} \textit{abbr. pro} cheveux. — \textbf{14} \textit{lect. incert. in ms.}}} \\
      \leavevmode\fontsize{6.9pt}{7.9pt}\selectfont\strut Billet 2e cl. Toulouse 1/4 place : 15 f 15 \hfill{} Coupe chvx : 1 f\strut &  \\
      \leavevmode\fontsize{6.9pt}{7.9pt}\selectfont\strut Christofle : 12 f 50, remise, une [mot illis.] : 10 f 50\strut &  \\
      \leavevmode\llap{\fontsize{6.4pt}{7.3pt}\selectfont\color{darkgray}15\hspace{1.8mm}}\fontsize{6.9pt}{7.9pt}\selectfont\strut Zyromski : L'orgueil humain : 4 f : net : 3 f 20\strut &  \\
      \leavevmode\fontsize{6.9pt}{7.9pt}\selectfont\strut Maeterlinck : Les Débris : 3 f 50. — Cirage : 25 c\strut &  \\
      \leavevmode\fontsize{6.9pt}{7.9pt}\selectfont\strut Déjeuner : 1 f 50. — Dîner : 1 f 60. — Ctes ples : 15 c\strut & \smash{\parbox[t]{39mm}{\fontsize{6.2pt}{7.0pt}\selectfont\raggedright\color{black!88}\textbf{17 Ctes ples} \textit{abbr. pro} Cartes postales. — \textbf{18 artiflles} \textit{abbr. pro} artificielles. — \textbf{20} \textit{lect. incert. in ms.} — \textbf{21} \textit{lect. incert. in ms.}}} \\
      \leavevmode\fontsize{6.9pt}{7.9pt}\selectfont\strut Fleurs artiflles, Hotel de Ville : 8 f 50 couronne\strut &  \\
      \leavevmode\fontsize{6.9pt}{7.9pt}\selectfont\strut Képi : au Bon Marché : 8 f 75. — Consigne : 15 c\strut &  \\
      \leavevmode\llap{\fontsize{6.4pt}{7.3pt}\selectfont\color{darkgray}20\hspace{1.8mm}}\fontsize{6.9pt}{7.9pt}\selectfont\strut Carte de Verdun : 50 c — [Entretien ?] souliers : 30 c\strut &  \\
      \leavevmode\fontsize{6.9pt}{7.9pt}\selectfont\strut Carte Roumanie : 30 c — [Divers ?] : W.C. : [15 c ?]\strut &  \\
\end{tabular}}%
\par\vspace*{\fill}
% --- Registres de bas de page (1. Critique interne [si pertinent] ; 2. Loci paralleli en latin ; 3. Commentaire historique en retrait) ---
\noindent\rule{\textwidth}{0.5pt}\par\vspace{0.65mm}
{\fontsize{6.6pt}{7.7pt}\selectfont%
\noindent\normalfont \textbf{Structure comptable :} Seconde page du mardi 26 septembre 1916, occupée aux l.~12--23 par le tableau détaillé des achats effectués par Hippolyte Combes à Paris avant son départ pour Toulouse.\par\vspace{0.5mm}
\noindent\rule{22mm}{0.3pt}\par\vspace{0.45mm}
\noindent\normalfont\color{black!92} \textbf{13} (\textit{Toulouse}) cf. supra in tegumento (p.~\pageref{fol:1}), f\textsuperscript{o}~26~r\textsuperscript{o} (p.~\pageref{fol:52}), infra f\textsuperscript{o}~27~r\textsuperscript{o} (p.~\pageref{fol:54}), f\textsuperscript{o}~27~v\textsuperscript{o} (p.~\pageref{fol:55}), \textit{et passim}, Appendicem I, doc.~1 et 3 (p.~\pageref{annexe:doc1}). — \textbf{15} (\textit{Zyromski}) cf. infra f\textsuperscript{o}~67~v\textsuperscript{o} (p.~\pageref{fol:135}). — \textbf{16} (\textit{Maeterlinck}) cf. .\par\vspace{0.55mm}
\noindent\rule{38mm}{0.4pt}\par\vspace{0.5mm}
\noindent\hspace*{5mm}\parbox{\dimexpr\textwidth-5mm\relax}{\fontsize{6.6pt}{7.7pt}\selectfont\color{black!84}À la librairie Armand Colin (boulevard Saint-Michel), Hippolyte Combes acquiert deux ouvrages contemporains : l'essai d'\textbf{Ernest Zyromski} (1862–1933, professeur de littérature française à la Faculté des lettres de Toulouse), \textbf{\textit{L'Orgueil humain}} (Armand Colin, 1904), et le recueil d'essais de guerre de \textbf{Maurice Maeterlinck}, \textbf{\textit{Les Débris de la guerre}} (Fasquelle, Bibliothèque-Charpentier, 1916).}\par}
\newpage
\phantomsection\label{fac:53}
% --- En-tête du fac-similé (Page impaire / droite : Photographie du manuscrit) ---
\noindent\rule{\textwidth}{0.5pt}\par\vspace{0.55mm}
\noindent\textbf{\fontsize{6.8pt}{7.8pt}\selectfont FAC-SIMILÉ EN VIS-À-VIS — FOLIO 26 VERSO} \hfill \textbf{\fontsize{6.8pt}{7.8pt}\selectfont SEPTEMBRE 1916}\par\vspace{0.25mm}
\noindent{\fontsize{6.8pt}{8.0pt}\selectfont\textit{Livre-Journal d'Hippolyte Combes} — Mardi 26 septembre 1916 (après-midi et soir) (Paris (Tour Saint-Jacques, rue de Rivoli, librairie Colin, gare d'Orsay))}\par\vspace{0.55mm}
\noindent\rule{\textwidth}{0.35pt}\par\vspace*{\fill}
{\centering
  \setlength{\fboxsep}{1.2mm}%
  \setlength{\fboxrule}{0.4pt}%
  \fcolorbox{gray!55}{white}{%
    \includegraphics[width=0.95\textwidth,height=0.81\textheight,keepaspectratio]{\facsimiledir/20251126_103131.jpg}%
  }\par}
\vspace*{\fill}
\newpage

% =============================================================================
% VIS-A-VIS SPREAD : FOLIO 27 RECTO (20251126_103146.jpg)
% =============================================================================
\cleardoubleevenpage
\renewcommand{\currentmonthyear}{Septembre 1916}
\renewcommand{\currentfoliolabel}{Folio 27 recto}
\phantomsection\label{fol:54}
% --- En-tête du folio (Page paire / gauche : Transcription, apparat & registres de bas de page) ---
\noindent\rule{\textwidth}{0.5pt}\par\vspace{0.55mm}
\noindent\textbf{\footnotesize FOLIO 27 RECTO} \hfill \textbf{\footnotesize SEPTEMBRE 1916}\par\vspace{0.25mm}
\noindent{\fontsize{6.9pt}{8.0pt}\selectfont Mercredi 27 septembre 1916 — Arrivée à Toulouse, montée à Sainte-Marie-des-Champs et déjeuner intime chez les Crouzet}\par\vspace{0.2mm}
\noindent{\fontsize{6.9pt}{8.0pt}\selectfont\textit{Mercredi 27 septembre 1916} \hfill \textit{Toulouse (Sainte-Marie-des-Champs)}}\par\vspace{0.5mm}
\noindent\rule{\textwidth}{0.35pt}\par\vspace{0.8mm}
% --- Corps de la transcription diplomatique et apparat marginal en prolongement ---
{\renewcommand{\arraystretch}{0.87}%
\noindent\hspace*{3.5mm}\begin{tabular}{@{}p{87mm}@{\hspace{3.5mm}}p{39mm}@{}}
      \leavevmode\fontsize{6.9pt}{7.9pt}\selectfont\strut \underline{27 Mer.} \hfill{} \underline{Sept 1916}\strut & \smash{\parbox[t]{39mm}{\fontsize{6.2pt}{7.0pt}\selectfont\raggedright\color{black!88}\textbf{1 Mer.} \textit{scil.} Mercredi.}} \\
      \leavevmode\fontsize{6.9pt}{7.9pt}\selectfont\strut Nuit passée ds compmt. Dormi 1 peu\strut & \smash{\parbox[t]{39mm}{\fontsize{6.2pt}{7.0pt}\selectfont\raggedright\color{black!88}\textbf{2 ds} \textit{abbr. pro} dans.}} \\
      \leavevmode\fontsize{6.9pt}{7.9pt}\selectfont\strut A 7H 30 m., arrivé à Toulouse\strut & \smash{\parbox[t]{39mm}{\fontsize{6.2pt}{7.0pt}\selectfont\raggedright\color{black!88}\textbf{3 m.} \textit{abbr. pro} matin.}} \\
      \leavevmode\fontsize{6.9pt}{7.9pt}\selectfont\strut Pris tram et vers 8H, ds mon appartement\strut & \smash{\parbox[t]{39mm}{\fontsize{6.2pt}{7.0pt}\selectfont\raggedright\color{black!88}\textbf{4 ds} \textit{abbr. pro} dans.}} \\
      \leavevmode\llap{\fontsize{6.4pt}{7.3pt}\selectfont\color{darkgray}5\hspace{1.8mm}}\fontsize{6.9pt}{7.9pt}\selectfont\strut A 11 H m. — monté à Ste Marie-des-Champs. Les dames\strut & \smash{\parbox[t]{39mm}{\fontsize{6.2pt}{7.0pt}\selectfont\raggedright\color{black!88}\textbf{5 Ste} \textit{abbr. pro} Sainte. \textbf{5 m.} \textit{abbr. pro} matin.}} \\
      \leavevmode\fontsize{6.9pt}{7.9pt}\selectfont\strut Calmels, pas e. arrivées, ne seront là que ce soir —\strut &  \\
      \leavevmode\fontsize{6.9pt}{7.9pt}\selectfont\strut Déjeuné chz les Crouzet, ds la + gde intimité —\strut & \smash{\parbox[t]{39mm}{\fontsize{6.2pt}{7.0pt}\selectfont\raggedright\color{black!88}\textbf{7 ds} \textit{abbr. pro} dans. \textbf{7 gde} \textit{abbr. pro} grande.}} \\
      \leavevmode\fontsize{6.9pt}{7.9pt}\selectfont\strut Redescendu après-midi en ville et fait courses\strut &  \\
      \leavevmode\fontsize{6.9pt}{7.9pt}\selectfont\strut diverses\strut &  \\
      \leavevmode\llap{\fontsize{6.4pt}{7.3pt}\selectfont\color{darkgray}10\hspace{1.8mm}}\fontsize{6.9pt}{7.9pt}\selectfont\strut A 5H, remonté Cote Pavée. Enfin, revu\strut &  \\
      \leavevmode\fontsize{6.9pt}{7.9pt}\selectfont\strut Margte. Charmante de simplicité \& de tendresse\strut & \smash{\parbox[t]{39mm}{\fontsize{6.2pt}{7.0pt}\selectfont\raggedright\color{black!88}\textbf{11 Margte} \textit{scil.} Marguerite (Calmels).}} \\
      \leavevmode\fontsize{6.9pt}{7.9pt}\selectfont\strut confiante.\strut &  \\
      \leavevmode\fontsize{6.9pt}{7.9pt}\selectfont\strut Trams : 10, 10, 20 c..\strut & \smash{\parbox[t]{39mm}{\fontsize{6.2pt}{7.0pt}\selectfont\raggedright\color{black!88}\textbf{13 c.} \textit{abbr. pro} centimes.}} \\
      \cline{1-1}
      \leavevmode\fontsize{6.9pt}{7.9pt}\selectfont\strut Bain, aux Thermes du Languedoc : 90 c + [15 c / 1 f ?]\strut & \smash{\parbox[t]{39mm}{\fontsize{6.2pt}{7.0pt}\selectfont\raggedright\color{black!88}\textbf{14} \textit{lect. incert. in ms.}}} \\
      \leavevmode\llap{\fontsize{6.4pt}{7.3pt}\selectfont\color{darkgray}15\hspace{1.8mm}}\fontsize{6.9pt}{7.9pt}\selectfont\strut Cartes ples :\strut & \smash{\parbox[t]{39mm}{\fontsize{6.2pt}{7.0pt}\selectfont\raggedright\color{black!88}\textbf{15 Cartes ples} \textit{abbr. pro} Cartes postales.}} \\
\end{tabular}}%
\par\vspace*{\fill}
% --- Registres de bas de page (1. Critique interne [si pertinent] ; 2. Loci paralleli en latin ; 3. Commentaire historique en retrait) ---
\noindent\rule{\textwidth}{0.5pt}\par\vspace{0.65mm}
{\fontsize{6.6pt}{7.7pt}\selectfont%
\noindent\normalfont\color{black!92} \textbf{3} (\textit{Toulouse}) cf. supra in tegumento (p.~\pageref{fol:1}), f\textsuperscript{o}~26~v\textsuperscript{o} (p.~\pageref{fol:53}), infra f\textsuperscript{o}~27~v\textsuperscript{o} (p.~\pageref{fol:55}), f\textsuperscript{o}~28~r\textsuperscript{o} (p.~\pageref{fol:56}), \textit{et passim}, Appendicem I, doc.~1 et 3 (p.~\pageref{annexe:doc1}). — \textbf{7} (\textit{Crouzet}) cf. infra f\textsuperscript{o}~27~v\textsuperscript{o} (p.~\pageref{fol:55}), f\textsuperscript{o}~31~r\textsuperscript{o} (p.~\pageref{fol:62}), \textit{et passim}. — \textbf{11} (\textit{Margte}) cf. supra f\textsuperscript{o}~1~r\textsuperscript{o} (p.~\pageref{fol:2}), f\textsuperscript{o}~9~v\textsuperscript{o} (p.~\pageref{fol:19}), infra f\textsuperscript{o}~27~v\textsuperscript{o} (p.~\pageref{fol:55}), f\textsuperscript{o}~28~r\textsuperscript{o} (p.~\pageref{fol:56}), \textit{et passim}.\par\vspace{0.55mm}
\noindent\rule{38mm}{0.4pt}\par\vspace{0.5mm}
\noindent\hspace*{5mm}\parbox{\dimexpr\textwidth-5mm\relax}{\fontsize{6.6pt}{7.7pt}\selectfont\color{black!84}Arrivé à \textbf{Toulouse} pour sa première permission, Hippolyte Combes retrouve sa fiancée \textbf{Marguerite Calmels} (1892–1974) chez l'oncle et la tante de celle-ci, \textbf{Émile Crouzet} (1842–1921, agent de change à Toulouse) et \textbf{Élisabeth Calmels}, et revoit le pensionnat jésuite de \textbf{Sainte-Marie-des-Champs} (Le Caousou) où il avait accompli ses études secondaires.}\par}
\newpage
\phantomsection\label{fac:54}
% --- En-tête du fac-similé (Page impaire / droite : Photographie du manuscrit) ---
\noindent\rule{\textwidth}{0.5pt}\par\vspace{0.55mm}
\noindent\textbf{\fontsize{6.8pt}{7.8pt}\selectfont FAC-SIMILÉ EN VIS-À-VIS — FOLIO 27 RECTO} \hfill \textbf{\fontsize{6.8pt}{7.8pt}\selectfont SEPTEMBRE 1916}\par\vspace{0.25mm}
\noindent{\fontsize{6.8pt}{8.0pt}\selectfont\textit{Livre-Journal d'Hippolyte Combes} — Mercredi 27 septembre 1916 (Toulouse (Sainte-Marie-des-Champs))}\par\vspace{0.55mm}
\noindent\rule{\textwidth}{0.35pt}\par\vspace*{\fill}
{\centering
  \setlength{\fboxsep}{1.2mm}%
  \setlength{\fboxrule}{0.4pt}%
  \fcolorbox{gray!55}{white}{%
    \includegraphics[width=0.95\textwidth,height=0.81\textheight,keepaspectratio]{\facsimiledir/20251126_103146.jpg}%
  }\par}
\vspace*{\fill}
\newpage

% =============================================================================
% VIS-A-VIS SPREAD : FOLIO 27 VERSO (20251126_103215.jpg)
% =============================================================================
\cleardoubleevenpage
\renewcommand{\currentmonthyear}{Septembre 1916}
\renewcommand{\currentfoliolabel}{Folio 27 verso}
\phantomsection\label{fol:55}
% --- En-tête du folio (Page paire / gauche : Transcription, apparat & registres de bas de page) ---
\noindent\rule{\textwidth}{0.5pt}\par\vspace{0.55mm}
\noindent\textbf{\footnotesize FOLIO 27 VERSO} \hfill \textbf{\footnotesize SEPTEMBRE 1916}\par\vspace{0.25mm}
\noindent{\fontsize{6.9pt}{8.0pt}\selectfont Jeudi 28 septembre 1916 — Rangement de l'appartement toulousain, bouquet pour Marguerite et promenade par Pouvourville et la crête de Pech-David}\par\vspace{0.2mm}
\noindent{\fontsize{6.9pt}{8.0pt}\selectfont\textit{Jeudi 28 septembre 1916} \hfill \textit{Toulouse (Sainte-Marie, barrière Saint-Agne, Pouvourville, Pech-David)}}\par\vspace{0.5mm}
\noindent\rule{\textwidth}{0.35pt}\par\vspace{0.8mm}
% --- Corps de la transcription diplomatique et apparat marginal en prolongement ---
{\renewcommand{\arraystretch}{0.87}%
\noindent\hspace*{3.5mm}\begin{tabular}{@{}p{87mm}@{\hspace{3.5mm}}p{39mm}@{}}
      \leavevmode\fontsize{6.9pt}{7.9pt}\selectfont\strut \underline{28 Jeudi.} \hfill{} \underline{Sept 16}\strut &  \\
      \leavevmode\fontsize{6.9pt}{7.9pt}\selectfont\strut Me suis attardé ds mon appartt — Pris un\strut & \smash{\parbox[t]{39mm}{\fontsize{6.2pt}{7.0pt}\selectfont\raggedright\color{black!88}\textbf{2 ds} \textit{abbr. pro} dans. \textbf{2 appartt} \textit{abbr. pro} appartement.}} \\
      \leavevmode\fontsize{6.9pt}{7.9pt}\selectfont\strut plaisir extrême à regarder \& ranger mes\strut &  \\
      \leavevmode\fontsize{6.9pt}{7.9pt}\selectfont\strut affaires, et à laisser peu à peu\strut &  \\
      \leavevmode\llap{\fontsize{6.4pt}{7.3pt}\selectfont\color{darkgray}5\hspace{1.8mm}}\fontsize{6.9pt}{7.9pt}\selectfont\strut s'éveiller p moi des souvenirs du passé.\strut &  \\
      \leavevmode\fontsize{6.9pt}{7.9pt}\selectfont\strut Monté à Ste-Marie, vers 11 H 1/2. Apporté à Mte\strut & \smash{\parbox[t]{39mm}{\fontsize{6.2pt}{7.0pt}\selectfont\raggedright\color{black!88}\textbf{6 Mte} \textit{scil.} Marguerite (Calmels). \textbf{6 Ste} \textit{abbr. pro} Sainte.}} \\
      \leavevmode\fontsize{6.9pt}{7.9pt}\selectfont\strut un très joli bouquet fleurs délicates.\strut &  \\
      \leavevmode\fontsize{6.9pt}{7.9pt}\selectfont\strut Mêmes impressions au sujet de ma fiancée :\strut &  \\
      \leavevmode\fontsize{6.9pt}{7.9pt}\selectfont\strut tjs simple, affectueuse, charmante\strut & \smash{\parbox[t]{39mm}{\fontsize{6.2pt}{7.0pt}\selectfont\raggedright\color{black!88}\textbf{9 tjs} \textit{abbr. pro} toujours.}} \\
      \leavevmode\llap{\fontsize{6.4pt}{7.3pt}\selectfont\color{darkgray}10\hspace{1.8mm}}\fontsize{6.9pt}{7.9pt}\selectfont\strut Après déjeun., promenade en cie Mme C. et\strut & \smash{\parbox[t]{39mm}{\fontsize{6.2pt}{7.0pt}\selectfont\raggedright\color{black!88}\textbf{10 déjeun.} \textit{abbr. pro} déjeuné / déjeuner. \textbf{10 Mme} \textit{abbr. pro} Madame. — \textbf{11 St} \textit{abbr. pro} Saint. \textbf{11 Czt} \textit{abbr. pro} Crouzet. — \textbf{12 j'à} \textit{abbr. pro} jusqu'à. \textbf{12 St} \textit{abbr. pro} Saint. — \textbf{13 j-à} \textit{abbr. pro} jusqu'à. \textbf{13 St} \textit{abbr. pro} Saint. — \textbf{18 St} \textit{abbr. pro} Saint.}} \\
      \leavevmode\fontsize{6.9pt}{7.9pt}\selectfont\strut M. Czt. Descendu à pied côte vers St Joseph et\strut &  \\
      \leavevmode\fontsize{6.9pt}{7.9pt}\selectfont\strut arrivé j'à la barrière St-Agne. De là, pris trams\strut &  \\
      \leavevmode\fontsize{6.9pt}{7.9pt}\selectfont\strut j-à St-Agne. Remonté par route ombragée\strut &  \\
      \leavevmode\fontsize{6.9pt}{7.9pt}\selectfont\strut j'au village de Pouvourville, église —\strut &  \\
      \leavevmode\llap{\fontsize{6.4pt}{7.3pt}\selectfont\color{darkgray}15\hspace{1.8mm}}\fontsize{6.9pt}{7.9pt}\selectfont\strut Puis, sur la hauteur, d'où vue superbe\strut &  \\
      \leavevmode\fontsize{6.9pt}{7.9pt}\selectfont\strut suivi crête de Pech David — j'au carref.\strut &  \\
      \leavevmode\fontsize{6.9pt}{7.9pt}\selectfont\strut Récollets et Caserne\strut &  \\
      \leavevmode\fontsize{6.9pt}{7.9pt}\selectfont\strut Allé prendre trams à la barrière St-Agne\strut &  \\
      \leavevmode\fontsize{6.9pt}{7.9pt}\selectfont\strut Exc.te après-midi.\strut &  \\
      \leavevmode\llap{\fontsize{6.4pt}{7.3pt}\selectfont\color{darkgray}20\hspace{1.8mm}}\fontsize{6.9pt}{7.9pt}\selectfont\strut Vers 9 H 30 s, redescendu à pied Cote Pavée\strut & \smash{\parbox[t]{39mm}{\fontsize{6.2pt}{7.0pt}\selectfont\raggedright\color{black!88}\textbf{20 s.} \textit{abbr. pro} soir.}} \\
      \leavevmode\fontsize{6.9pt}{7.9pt}\selectfont\strut Trams : 20 c, 40, 40, 40 c. —\strut & \smash{\parbox[t]{39mm}{\fontsize{6.2pt}{7.0pt}\selectfont\raggedright\color{black!88}\textbf{21 c.} \textit{abbr. pro} centimes.}} \\
      \leavevmode\fontsize{6.9pt}{7.9pt}\selectfont\strut Fruits (raisins \& grenades) : 1 f — Varia : 10 c.\strut & \smash{\parbox[t]{39mm}{\fontsize{6.2pt}{7.0pt}\selectfont\raggedright\color{black!88}\textbf{22 c.} \textit{abbr. pro} centimes.}} \\
      \leavevmode\fontsize{6.9pt}{7.9pt}\selectfont\strut Bouquet (place Esquirol) : 3 f —\strut &  \\
\end{tabular}}%
\par\vspace*{\fill}
% --- Registres de bas de page (1. Critique interne [si pertinent] ; 2. Loci paralleli en latin ; 3. Commentaire historique en retrait) ---
\noindent\rule{\textwidth}{0.5pt}\par\vspace{0.65mm}
{\fontsize{6.6pt}{7.7pt}\selectfont%
\noindent\normalfont\color{black!92} \textbf{6} (\textit{Mte}) cf. supra f\textsuperscript{o}~1~r\textsuperscript{o} (p.~\pageref{fol:2}), f\textsuperscript{o}~27~r\textsuperscript{o} (p.~\pageref{fol:54}), infra f\textsuperscript{o}~28~r\textsuperscript{o} (p.~\pageref{fol:56}), f\textsuperscript{o}~28~v\textsuperscript{o} (p.~\pageref{fol:57}), \textit{et passim}, Appendicem I, doc.~5 (p.~\pageref{annexe:doc5}). — \textbf{6} (\textit{Ste-Marie}) cf. infra f\textsuperscript{o}~66~r\textsuperscript{o} (p.~\pageref{fol:132}), f\textsuperscript{o}~66~v\textsuperscript{o} (p.~\pageref{fol:133}), \textit{et passim}. — \textbf{11} (\textit{Czt}) cf. supra f\textsuperscript{o}~27~r\textsuperscript{o} (p.~\pageref{fol:54}), infra f\textsuperscript{o}~31~r\textsuperscript{o} (p.~\pageref{fol:62}), f\textsuperscript{o}~32~r\textsuperscript{o} (p.~\pageref{fol:64}), \textit{et passim}.\par\vspace{0.55mm}
\noindent\rule{38mm}{0.4pt}\par\vspace{0.5mm}
\noindent\hspace*{5mm}\parbox{\dimexpr\textwidth-5mm\relax}{\fontsize{6.6pt}{7.7pt}\selectfont\color{black!84}Après avoir offert à Marguerite Calmels un bouquet acheté place Esquirol, Hippolyte Combes effectue avec elle, \textbf{Élisabeth Calmels-Crouzet} et \textbf{Émile Crouzet} (\textit{M.~Czt}) une promenade sur les coteaux dominant la Garonne au sud de Toulouse, par l'église de \textbf{Pouvourville} et la crête de \textbf{Pech-David}.}\par}
\newpage
\phantomsection\label{fac:55}
% --- En-tête du fac-similé (Page impaire / droite : Photographie du manuscrit) ---
\noindent\rule{\textwidth}{0.5pt}\par\vspace{0.55mm}
\noindent\textbf{\fontsize{6.8pt}{7.8pt}\selectfont FAC-SIMILÉ EN VIS-À-VIS — FOLIO 27 VERSO} \hfill \textbf{\fontsize{6.8pt}{7.8pt}\selectfont SEPTEMBRE 1916}\par\vspace{0.25mm}
\noindent{\fontsize{6.8pt}{8.0pt}\selectfont\textit{Livre-Journal d'Hippolyte Combes} — Jeudi 28 septembre 1916 (Toulouse (Sainte-Marie, barrière Saint-Agne, Pouvourville, Pech-David))}\par\vspace{0.55mm}
\noindent\rule{\textwidth}{0.35pt}\par\vspace*{\fill}
{\centering
  \setlength{\fboxsep}{1.2mm}%
  \setlength{\fboxrule}{0.4pt}%
  \fcolorbox{gray!55}{white}{%
    \includegraphics[width=0.95\textwidth,height=0.81\textheight,keepaspectratio]{\facsimiledir/20251126_103215.jpg}%
  }\par}
\vspace*{\fill}
\newpage

% =============================================================================
% VIS-A-VIS SPREAD : FOLIO 28 RECTO (20251126_103227.jpg)
% =============================================================================
\cleardoubleevenpage
\renewcommand{\currentmonthyear}{Septembre 1916}
\renewcommand{\currentfoliolabel}{Folio 28 recto}
\phantomsection\label{fol:56}
% --- En-tête du folio (Page paire / gauche : Transcription, apparat & registres de bas de page) ---
\noindent\rule{\textwidth}{0.5pt}\par\vspace{0.55mm}
\noindent\textbf{\footnotesize FOLIO 28 RECTO} \hfill \textbf{\footnotesize SEPTEMBRE 1916}\par\vspace{0.25mm}
\noindent{\fontsize{6.9pt}{8.0pt}\selectfont Vendredi 29 septembre 1916 — Visite à Mme de Labarthe, promenade avec sa fiancée Marguerite à Blagnac (ruines de l'amphithéâtre romain) et achat d'un sautoir en or chez Pfister (295 fr.)}\par\vspace{0.2mm}
\noindent{\fontsize{6.9pt}{8.0pt}\selectfont\textit{Vendredi 29 septembre 1916} \hfill \textit{Toulouse, Purpan, Blagnac et les Ponts-Jumeaux (Haute-Garonne)}}\par\vspace{0.5mm}
\noindent\rule{\textwidth}{0.35pt}\par\vspace{0.8mm}
% --- Corps de la transcription diplomatique et apparat marginal en prolongement ---
{\renewcommand{\arraystretch}{0.93}%
\noindent\hspace*{3.5mm}\begin{tabular}{@{}p{87mm}@{\hspace{3.5mm}}p{39mm}@{}}
      \leavevmode\fontsize{7.3pt}{8.4pt}\selectfont\strut \underline{29 Ve.} \hfill{} \underline{Sept / 16}\strut & \smash{\parbox[t]{39mm}{\fontsize{6.2pt}{7.0pt}\selectfont\raggedright\color{black!88}\textbf{1 Ve.} \textit{scil.} Vendredi.}} \\
      \leavevmode\fontsize{7.3pt}{8.4pt}\selectfont\strut Tout le jour, ciel nuageux.\strut &  \\
      \leavevmode\fontsize{7.3pt}{8.4pt}\selectfont\strut Fin après-midi, courte averse\strut &  \\
      \leavevmode\fontsize{7.3pt}{8.4pt}\selectfont\strut Matin, venu voir Mme de Labarthe. gde amie\strut & \smash{\parbox[t]{39mm}{\fontsize{6.2pt}{7.0pt}\selectfont\raggedright\color{black!88}\textbf{4 gde} \textit{abbr. pro} grande. \textbf{4 Mme} \textit{abbr. pro} Madame.}} \\
      \leavevmode\llap{\fontsize{6.4pt}{7.3pt}\selectfont\color{darkgray}5\hspace{1.8mm}}\fontsize{7.3pt}{8.4pt}\selectfont\strut affectueuse \& attachée.\strut &  \\
      \leavevmode\fontsize{7.3pt}{8.4pt}\selectfont\strut Marché — déjeuné restaurant du Chge — et retrouvé Mte\strut & \smash{\parbox[t]{39mm}{\fontsize{6.2pt}{7.0pt}\selectfont\raggedright\color{black!88}\textbf{6 Mte} \textit{scil.} Marguerite (Calmels).}} \\
      \leavevmode\fontsize{7.3pt}{8.4pt}\selectfont\strut Affection réciproque de + en + profonde. Intimité\strut &  \\
      \leavevmode\fontsize{7.3pt}{8.4pt}\selectfont\strut tjs + gde —\strut & \smash{\parbox[t]{39mm}{\fontsize{6.2pt}{7.0pt}\selectfont\raggedright\color{black!88}\textbf{8 tjs} \textit{abbr. pro} toujours. \textbf{8 gde} \textit{abbr. pro} grande.}} \\
      \leavevmode\fontsize{7.3pt}{8.4pt}\selectfont\strut Après-midi, course en tram. Esquirol – Purpan\strut &  \\
      \leavevmode\llap{\fontsize{6.4pt}{7.3pt}\selectfont\color{darkgray}10\hspace{1.8mm}}\fontsize{7.3pt}{8.4pt}\selectfont\strut Longue halte aux Ruines. Assis côte à côte\strut &  \\
      \leavevmode\fontsize{7.3pt}{8.4pt}\selectfont\strut sur reste de pilier, doux épanchts \& baisers intimes\strut & \smash{\parbox[t]{39mm}{\fontsize{6.2pt}{7.0pt}\selectfont\raggedright\color{black!88}\textbf{11 épanchts} \textit{abbr. pro} épanchements. \textbf{11 épanchts} \textit{abbr. pro} épanchements. — \textbf{13 j'à} \textit{abbr. pro} jusqu'à. \textbf{13 ds} \textit{abbr. pro} dans.}} \\
      \leavevmode\fontsize{7.3pt}{8.4pt}\selectfont\strut av. ma Femme —\strut &  \\
      \leavevmode\fontsize{7.3pt}{8.4pt}\selectfont\strut Continué à pied j'à Blagnac. Reposé ds\strut &  \\
      \leavevmode\fontsize{7.3pt}{8.4pt}\selectfont\strut la chapelle du monastère.\strut &  \\
      \leavevmode\llap{\fontsize{6.4pt}{7.3pt}\selectfont\color{darkgray}15\hspace{1.8mm}}\fontsize{7.3pt}{8.4pt}\selectfont\strut Allé ensuite prairie au bord de la Garonne\strut &  \\
      \leavevmode\fontsize{7.3pt}{8.4pt}\selectfont\strut Assis sur l'herbe haute, en face petite grève\strut &  \\
      \leavevmode\fontsize{7.3pt}{8.4pt}\selectfont\strut Venus à pied j'aux Ponts Jumeaux, d'où pris\strut &  \\
      \leavevmode\fontsize{7.3pt}{8.4pt}\selectfont\strut tram j'à pl. Alsace. —\strut & \smash{\parbox[t]{39mm}{\fontsize{6.2pt}{7.0pt}\selectfont\raggedright\color{black!88}\textbf{18 j'à} \textit{abbr. pro} jusqu'à.}} \\
      \leavevmode\fontsize{7.3pt}{8.4pt}\selectfont\strut — Trams : 20 c, 40 c, 60 c, 40, 40 c –\strut &  \\
      \leavevmode\llap{\fontsize{6.4pt}{7.3pt}\selectfont\color{darkgray}20\hspace{1.8mm}}\fontsize{7.3pt}{8.4pt}\selectfont\strut Chz Pfister, sautoir en or — : 300 f net 295 f\strut &  \\
      \leavevmode\fontsize{7.3pt}{8.4pt}\selectfont\strut Remboursé à mes nièces [?] (oeufs, beurre etc.) : 2 f 75\strut & \smash{\parbox[t]{39mm}{\fontsize{6.2pt}{7.0pt}\selectfont\raggedright\color{black!88}\textbf{21 nièces} \textit{lect. incert. in ms.} — \textbf{22 c.} \textit{abbr. pro} centimes.}} \\
      \leavevmode\fontsize{7.3pt}{8.4pt}\selectfont\strut Fruits : 2f + 55 c. – Fleurs : 2 f\strut &  \\
\end{tabular}}%
\par\vspace*{\fill}
% --- Registres de bas de page (1. Critique interne [si pertinent] ; 2. Loci paralleli en latin ; 3. Commentaire historique en retrait) ---
\noindent\rule{\textwidth}{0.5pt}\par\vspace{0.65mm}
{\fontsize{6.6pt}{7.7pt}\selectfont%
\noindent\normalfont\color{black!92} \textbf{6} (\textit{Mte}) cf. supra f\textsuperscript{o}~1~r\textsuperscript{o} (p.~\pageref{fol:2}), f\textsuperscript{o}~27~v\textsuperscript{o} (p.~\pageref{fol:55}), infra f\textsuperscript{o}~28~v\textsuperscript{o} (p.~\pageref{fol:57}), f\textsuperscript{o}~29~r\textsuperscript{o} (p.~\pageref{fol:58}), \textit{et passim}, Appendicem I, doc.~5 (p.~\pageref{annexe:doc5}). — \textbf{20} (\textit{Pfister}) cf. infra f\textsuperscript{o}~68~v\textsuperscript{o} (p.~\pageref{fol:137}), f\textsuperscript{o}~69~r\textsuperscript{o} (p.~\pageref{fol:138}).\par\vspace{0.55mm}
\noindent\rule{38mm}{0.4pt}\par\vspace{0.5mm}
\noindent\hspace*{5mm}\parbox{\dimexpr\textwidth-5mm\relax}{\fontsize{6.6pt}{7.7pt}\selectfont\color{black!84}Journée de fiançailles à \textbf{Toulouse} le vendredi 29 septembre 1916 : Hippolyte Combes offre à \textbf{Marguerite Calmels} — qu'il appelle déjà dans l'intimité du carnet \textit{ma Femme}, un an avant leur mariage célébré à Millau le 25 octobre 1917 — un sautoir en or acquis 295~francs.}\par}
\newpage
\phantomsection\label{fac:56}
% --- En-tête du fac-similé (Page impaire / droite : Photographie du manuscrit) ---
\noindent\rule{\textwidth}{0.5pt}\par\vspace{0.55mm}
\noindent\textbf{\fontsize{6.8pt}{7.8pt}\selectfont FAC-SIMILÉ EN VIS-À-VIS — FOLIO 28 RECTO} \hfill \textbf{\fontsize{6.8pt}{7.8pt}\selectfont SEPTEMBRE 1916}\par\vspace{0.25mm}
\noindent{\fontsize{6.8pt}{8.0pt}\selectfont\textit{Livre-Journal d'Hippolyte Combes} — Vendredi 29 septembre 1916 (Toulouse, Purpan, Blagnac et les Ponts-Jumeaux (Haute-Garonne))}\par\vspace{0.55mm}
\noindent\rule{\textwidth}{0.35pt}\par\vspace*{\fill}
{\centering
  \setlength{\fboxsep}{1.2mm}%
  \setlength{\fboxrule}{0.4pt}%
  \fcolorbox{gray!55}{white}{%
    \includegraphics[width=0.95\textwidth,height=0.81\textheight,keepaspectratio]{\facsimiledir/20251126_103227.jpg}%
  }\par}
\vspace*{\fill}
\newpage

% =============================================================================
% VIS-A-VIS SPREAD : FOLIO 28 VERSO (20251126_103247.jpg)
% =============================================================================
\cleardoubleevenpage
\renewcommand{\currentmonthyear}{Septembre 1916}
\renewcommand{\currentfoliolabel}{Folio 28 verso}
\phantomsection\label{fol:57}
% --- En-tête du folio (Page paire / gauche : Transcription, apparat & registres de bas de page) ---
\noindent\rule{\textwidth}{0.5pt}\par\vspace{0.55mm}
\noindent\textbf{\footnotesize FOLIO 28 VERSO} \hfill \textbf{\footnotesize SEPTEMBRE 1916}\par\vspace{0.25mm}
\noindent{\fontsize{6.9pt}{8.0pt}\selectfont Les cinq cloches du clocher de Montaudran et le serment au cimetière de Terre-Cabade}\par\vspace{0.2mm}
\noindent{\fontsize{6.9pt}{8.0pt}\selectfont\textit{Samedi 30 septembre 1916} \hfill \textit{Toulouse (Sainte-Marie, Montaudran et cimetière de Terre-Cabade)}}\par\vspace{0.5mm}
\noindent\rule{\textwidth}{0.35pt}\par\vspace{0.8mm}
% --- Corps de la transcription diplomatique et apparat marginal en prolongement ---
{\renewcommand{\arraystretch}{0.83}%
\noindent\hspace*{3.5mm}\begin{tabular}{@{}p{87mm}@{\hspace{3.5mm}}p{39mm}@{}}
      \leavevmode\fontsize{6.6pt}{7.5pt}\selectfont\strut \underline{30 Samedi.} \hfill{} \underline{Septb /16}\strut &  \\
      \leavevmode\fontsize{6.6pt}{7.5pt}\selectfont\strut Tout le jour, temps pluvieux. Fines ondées qui durent peu.\strut &  \\
      \leavevmode\fontsize{6.6pt}{7.5pt}\selectfont\strut Matin, venu voir Mme de Labarthe, au 12.. Ai\strut & \smash{\parbox[t]{39mm}{\fontsize{6.2pt}{7.0pt}\selectfont\raggedright\color{black!88}\textbf{3 Mme} \textit{abbr. pro} Madame.}} \\
      \leavevmode\fontsize{6.6pt}{7.5pt}\selectfont\strut pour cette amie, attachement profond. –\strut &  \\
      \leavevmode\llap{\fontsize{6.4pt}{7.3pt}\selectfont\color{darkgray}5\hspace{1.8mm}}\fontsize{6.6pt}{7.5pt}\selectfont\strut Monté à Ste Marie vers 11H 1/2 —\strut & \smash{\parbox[t]{39mm}{\fontsize{6.2pt}{7.0pt}\selectfont\raggedright\color{black!88}\textbf{5 Ste} \textit{abbr. pro} Sainte. \textbf{5 1/2} \textit{scil.} demi / moitié. — \textbf{6 qq} \textit{abbr. pro} quelque(s). — \textbf{7 Margte} \textit{scil.} Marguerite (Calmels). — \textbf{9 ds} \textit{abbr. pro} dans.}} \\
      \leavevmode\fontsize{6.6pt}{7.5pt}\selectfont\strut Dans petit salon, qq minutes de tête à tête délicieux\strut &  \\
      \leavevmode\fontsize{6.6pt}{7.5pt}\selectfont\strut avec Margte..\strut &  \\
      \leavevmode\fontsize{6.6pt}{7.5pt}\selectfont\strut Après midi, promenade à pied au village\strut &  \\
      \leavevmode\fontsize{6.6pt}{7.5pt}\selectfont\strut de Montaudran. Reposé ds l'église. Monté à la\strut &  \\
      \leavevmode\llap{\fontsize{6.4pt}{7.3pt}\selectfont\color{darkgray}10\hspace{1.8mm}}\fontsize{6.6pt}{7.5pt}\selectfont\strut tribune puis au clocher, par escalier très\strut &  \\
      \leavevmode\fontsize{6.6pt}{7.5pt}\selectfont\strut raide \& très étroit, presque 1 échelle. Là-haut,\strut &  \\
      \leavevmode\fontsize{6.6pt}{7.5pt}\selectfont\strut très b. vue panoramique de la plaine du Lauraguais\strut &  \\
      \leavevmode\fontsize{6.6pt}{7.5pt}\selectfont\strut Paysage nullement banal. Verdure \& lignes d'arbres\strut &  \\
      \leavevmode\fontsize{6.6pt}{7.5pt}\selectfont\strut clochers de villages s/ l'horizon. – Ciel tout\strut & \smash{\parbox[t]{39mm}{\fontsize{6.2pt}{7.0pt}\selectfont\raggedright\color{black!88}\textbf{14 s/} \textit{abbr. pro} sur.}} \\
      \leavevmode\llap{\fontsize{6.4pt}{7.3pt}\selectfont\color{darkgray}15\hspace{1.8mm}}\fontsize{6.6pt}{7.5pt}\selectfont\strut couvert de nuages noirs.\strut &  \\
      \leavevmode\fontsize{6.6pt}{7.5pt}\selectfont\strut Fait vibrer chacune des 5 cloches d'airain – –\strut &  \\
      \leavevmode\fontsize{6.6pt}{7.5pt}\selectfont\strut Sons différents \& très mélodieux — [ajout marginal : p. Rem ?]\strut & \smash{\parbox[t]{39mm}{\fontsize{6.2pt}{7.0pt}\selectfont\raggedright\color{black!88}\textbf{17} \textit{add. in marg. (lect. incert.).} — \textbf{18 j-au} \textit{abbr. pro} jusqu'au.}} \\
      \leavevmode\fontsize{6.6pt}{7.5pt}\selectfont\strut Continué promenade j-au Cimetière de\strut &  \\
      \leavevmode\fontsize{6.6pt}{7.5pt}\selectfont\strut Terre-Cabade. Admiré tombeau imité d'un\strut &  \\
      \leavevmode\llap{\fontsize{6.4pt}{7.3pt}\selectfont\color{darkgray}20\hspace{1.8mm}}\fontsize{6.6pt}{7.5pt}\selectfont\strut temple grec\strut &  \\
      \leavevmode\fontsize{6.6pt}{7.5pt}\selectfont\strut Venu devant caveau des P.P. Jésuites \& celui du\strut &  \\
      \leavevmode\fontsize{6.6pt}{7.5pt}\selectfont\strut P. Marie-Antoine. Margte \& moi, échangé serments\strut & \smash{\parbox[t]{39mm}{\fontsize{6.2pt}{7.0pt}\selectfont\raggedright\color{black!88}\textbf{22 Margte} \textit{scil.} Marguerite (Calmels). — \textbf{23 gd} \textit{abbr. pro} grand. \textbf{23 profondt} \textit{abbr. pro} profondément.}} \\
      \leavevmode\fontsize{6.6pt}{7.5pt}\selectfont\strut profondt émus au souvenir de notre gd ami \& de Josephe\strut &  \\
      \leavevmode\fontsize{6.6pt}{7.5pt}\selectfont\strut Retour du cimetière à pied par rues \& chemins de\strut &  \\
      \leavevmode\llap{\fontsize{6.4pt}{7.3pt}\selectfont\color{darkgray}25\hspace{1.8mm}}\fontsize{6.6pt}{7.5pt}\selectfont\strut traverse – Pluie cesse bientôt \& soleil couché\strut &  \\
      \leavevmode\fontsize{6.6pt}{7.5pt}\selectfont\strut horizon nuageux\strut &  \\
      \leavevmode\fontsize{6.6pt}{7.5pt}\selectfont\strut Tram : 20 c. — Sweets \& cakes : 3 f — Bière \& limonade : 80 c.\strut & \smash{\parbox[t]{39mm}{\fontsize{6.2pt}{7.0pt}\selectfont\raggedright\color{black!88}\textbf{27 c.} \textit{abbr. pro} centimes. \textbf{27 Sweets} \textit{angl. in ms.}}} \\
\end{tabular}}%
\par\vspace*{\fill}
% --- Registres de bas de page (1. Critique interne [si pertinent] ; 2. Loci paralleli en latin ; 3. Commentaire historique en retrait) ---
\noindent\rule{\textwidth}{0.5pt}\par\vspace{0.65mm}
{\fontsize{6.6pt}{7.7pt}\selectfont%
\noindent\normalfont\color{black!92} \textbf{7} (\textit{Margte}) cf. supra f\textsuperscript{o}~1~r\textsuperscript{o} (p.~\pageref{fol:2}), f\textsuperscript{o}~28~r\textsuperscript{o} (p.~\pageref{fol:56}), infra f\textsuperscript{o}~29~r\textsuperscript{o} (p.~\pageref{fol:58}), f\textsuperscript{o}~29~v\textsuperscript{o} (p.~\pageref{fol:59}), \textit{et passim}, Appendicem I, doc.~5 (p.~\pageref{annexe:doc5}).\par\vspace{0.55mm}
\noindent\rule{38mm}{0.4pt}\par\vspace{0.5mm}
\noindent\hspace*{5mm}\parbox{\dimexpr\textwidth-5mm\relax}{\fontsize{6.6pt}{7.7pt}\selectfont\color{black!84}Promenade toulousaine avec Marguerite Calmels le long des allées et du \textbf{canal du Midi} jusqu'au faubourg Saint-Étienne, rythmée par la sonnerie des cinq cloches de la cathédrale \textbf{Saint-Étienne de Toulouse}.}\par}
\newpage
\phantomsection\label{fac:57}
% --- En-tête du fac-similé (Page impaire / droite : Photographie du manuscrit) ---
\noindent\rule{\textwidth}{0.5pt}\par\vspace{0.55mm}
\noindent\textbf{\fontsize{6.8pt}{7.8pt}\selectfont FAC-SIMILÉ EN VIS-À-VIS — FOLIO 28 VERSO} \hfill \textbf{\fontsize{6.8pt}{7.8pt}\selectfont SEPTEMBRE 1916}\par\vspace{0.25mm}
\noindent{\fontsize{6.8pt}{8.0pt}\selectfont\textit{Livre-Journal d'Hippolyte Combes} — Samedi 30 septembre 1916 (Toulouse (Sainte-Marie, Montaudran et cimetière de Terre-Cabade))}\par\vspace{0.55mm}
\noindent\rule{\textwidth}{0.35pt}\par\vspace*{\fill}
{\centering
  \setlength{\fboxsep}{1.2mm}%
  \setlength{\fboxrule}{0.4pt}%
  \fcolorbox{gray!55}{white}{%
    \includegraphics[width=0.95\textwidth,height=0.81\textheight,keepaspectratio]{\facsimiledir/20251126_103247.jpg}%
  }\par}
\vspace*{\fill}
\newpage

% =============================================================================
% VIS-A-VIS SPREAD : FOLIO 29 RECTO (20251126_103301.jpg)
% =============================================================================
\cleardoubleevenpage
\renewcommand{\currentmonthyear}{Octobre 1916}
\renewcommand{\currentfoliolabel}{Folio 29 recto}
\phantomsection\label{fol:58}
% --- En-tête du folio (Page paire / gauche : Transcription, apparat & registres de bas de page) ---
\noindent\rule{\textwidth}{0.5pt}\par\vspace{0.55mm}
\noindent\textbf{\footnotesize FOLIO 29 RECTO} \hfill \textbf{\footnotesize OCTOBRE 1916}\par\vspace{0.25mm}
\noindent{\fontsize{6.9pt}{8.0pt}\selectfont Messe matinale à la cathédrale Saint-Étienne et promenade au bord du Touch à Lardenne}\par\vspace{0.2mm}
\noindent{\fontsize{6.9pt}{8.0pt}\selectfont\textit{Dimanche 1er octobre 1916} \hfill \textit{Toulouse (Cathédrale Saint-Étienne, Sainte-Marie et Lardenne)}}\par\vspace{0.5mm}
\noindent\rule{\textwidth}{0.35pt}\par\vspace{0.8mm}
% --- Corps de la transcription diplomatique et apparat marginal en prolongement ---
{\renewcommand{\arraystretch}{0.93}%
\noindent\hspace*{3.5mm}\begin{tabular}{@{}p{87mm}@{\hspace{3.5mm}}p{39mm}@{}}
      \leavevmode\fontsize{7.3pt}{8.4pt}\selectfont\strut \underline{1er D.} \hfill{} \underline{Octobre /16}\strut & \smash{\parbox[t]{39mm}{\fontsize{6.2pt}{7.0pt}\selectfont\raggedright\color{black!88}\textbf{1 D.} \textit{scil.} Dimanche. \textbf{1 1er} \textit{scil.} premier. — \textbf{2 St} \textit{abbr. pro} Saint. — \textbf{3} \textit{lect. incert. in ms.} \textbf{3 Mte} \textit{scil.} Marguerite (Calmels). — \textbf{4 ds} \textit{abbr. pro} dans. \textbf{4 appartt} \textit{abbr. pro} appartement. — \textbf{6 déjeun.} \textit{abbr. pro} déjeuné / déjeuner. \textbf{6 Ste} \textit{abbr. pro} Sainte. — \textbf{8 j-au} \textit{abbr. pro} jusqu'au. — \textbf{9 s/} \textit{abbr. pro} sur. — \textbf{10 Tjs} \textit{abbr. pro} Toujours. \textbf{10 Mte} \textit{scil.} Marguerite (Calmels). — \textbf{14 c.} \textit{abbr. pro} centimes. — \textbf{15 Sweets} \textit{angl. in ms.}}} \\
      \leavevmode\fontsize{7.3pt}{8.4pt}\selectfont\strut Matin, à St-Etienne. Entendu\strut &  \\
      \leavevmode\fontsize{7.3pt}{8.4pt}\selectfont\strut messe [curé ?] — 8 H autel paroissial. Mte, à ma droite\strut &  \\
      \leavevmode\fontsize{7.3pt}{8.4pt}\selectfont\strut Matinée, ds mon appartt, à ranger mes livres\strut &  \\
      \leavevmode\llap{\fontsize{6.4pt}{7.3pt}\selectfont\color{darkgray}5\hspace{1.8mm}}\fontsize{7.3pt}{8.4pt}\selectfont\strut Joui du calme \& silence —\strut &  \\
      \leavevmode\fontsize{7.3pt}{8.4pt}\selectfont\strut Monté déjeun à Ste Marie. – Temps tr beau\strut &  \\
      \leavevmode\fontsize{7.3pt}{8.4pt}\selectfont\strut Après repas, pris trams et venu à Lardenne\strut &  \\
      \leavevmode\fontsize{7.3pt}{8.4pt}\selectfont\strut De l'église, continué à pied j-au pont du Touch\strut &  \\
      \leavevmode\fontsize{7.3pt}{8.4pt}\selectfont\strut Assis s/ l'herbe prairie longeant rivière –\strut &  \\
      \leavevmode\llap{\fontsize{6.4pt}{7.3pt}\selectfont\color{darkgray}10\hspace{1.8mm}}\fontsize{7.3pt}{8.4pt}\selectfont\strut \& Longue causerie intime avec Mte. – Tjs\strut &  \\
      \leavevmode\fontsize{7.3pt}{8.4pt}\selectfont\strut charmante, simple \& affectueuse\strut &  \\
      \leavevmode\fontsize{7.3pt}{8.4pt}\selectfont\strut Au retour, trouvé Jean maison.\strut &  \\
      \leavevmode\fontsize{7.3pt}{8.4pt}\selectfont\strut Redescendus à 7 H 30 av. Jean en tram. —\strut &  \\
      \leavevmode\fontsize{7.3pt}{8.4pt}\selectfont\strut Trams : 20 c. 40 c, 80, 30 c\strut &  \\
      \leavevmode\llap{\fontsize{6.4pt}{7.3pt}\selectfont\color{darkgray}15\hspace{1.8mm}}\fontsize{7.3pt}{8.4pt}\selectfont\strut Sweets : 20 c\strut &  \\
\end{tabular}}%
\par\vspace*{\fill}
% --- Registres de bas de page (1. Critique interne [si pertinent] ; 2. Loci paralleli en latin ; 3. Commentaire historique en retrait) ---
\noindent\rule{\textwidth}{0.5pt}\par\vspace{0.65mm}
{\fontsize{6.6pt}{7.7pt}\selectfont%
\noindent\normalfont\color{black!92} \textbf{2} (\textit{St-Etienne}) cf. infra f\textsuperscript{o}~30~r\textsuperscript{o} (p.~\pageref{fol:60}). — \textbf{3} (\textit{Mte}) cf. supra f\textsuperscript{o}~1~r\textsuperscript{o} (p.~\pageref{fol:2}), f\textsuperscript{o}~28~v\textsuperscript{o} (p.~\pageref{fol:57}), infra f\textsuperscript{o}~29~v\textsuperscript{o} (p.~\pageref{fol:59}), f\textsuperscript{o}~30~r\textsuperscript{o} (p.~\pageref{fol:60}), \textit{et passim}. — \textbf{12} (\textit{Jean}) cf. infra f\textsuperscript{o}~29~v\textsuperscript{o} (p.~\pageref{fol:59}), f\textsuperscript{o}~66~v\textsuperscript{o} (p.~\pageref{fol:133}).\par\vspace{0.55mm}
\noindent\rule{38mm}{0.4pt}\par\vspace{0.5mm}
\noindent\hspace*{5mm}\parbox{\dimexpr\textwidth-5mm\relax}{\fontsize{6.6pt}{7.7pt}\selectfont\color{black!84}Dimanche 1\textsuperscript{er} octobre 1916 à \textbf{Toulouse} : après la messe matinale à la cathédrale \textbf{Saint-Étienne}, Hippolyte Combes range sa bibliothèque personnelle dans son appartement toulousain avant une promenade vespérale au Grand-Rond avec Marguerite Calmels.}\par}
\newpage
\phantomsection\label{fac:58}
% --- En-tête du fac-similé (Page impaire / droite : Photographie du manuscrit) ---
\noindent\rule{\textwidth}{0.5pt}\par\vspace{0.55mm}
\noindent\textbf{\fontsize{6.8pt}{7.8pt}\selectfont FAC-SIMILÉ EN VIS-À-VIS — FOLIO 29 RECTO} \hfill \textbf{\fontsize{6.8pt}{7.8pt}\selectfont OCTOBRE 1916}\par\vspace{0.25mm}
\noindent{\fontsize{6.8pt}{8.0pt}\selectfont\textit{Livre-Journal d'Hippolyte Combes} — Dimanche 1er octobre 1916 (Toulouse (Cathédrale Saint-Étienne, Sainte-Marie et Lardenne))}\par\vspace{0.55mm}
\noindent\rule{\textwidth}{0.35pt}\par\vspace*{\fill}
{\centering
  \setlength{\fboxsep}{1.2mm}%
  \setlength{\fboxrule}{0.4pt}%
  \fcolorbox{gray!55}{white}{%
    \includegraphics[width=0.95\textwidth,height=0.81\textheight,keepaspectratio]{\facsimiledir/20251126_103301.jpg}%
  }\par}
\vspace*{\fill}
\newpage

% =============================================================================
% VIS-A-VIS SPREAD : FOLIO 29 VERSO (20251126_103321.jpg)
% =============================================================================
\cleardoubleevenpage
\renewcommand{\currentmonthyear}{Octobre 1916}
\renewcommand{\currentfoliolabel}{Folio 29 verso}
\phantomsection\label{fol:59}
% --- En-tête du folio (Page paire / gauche : Transcription, apparat & registres de bas de page) ---
\noindent\rule{\textwidth}{0.5pt}\par\vspace{0.55mm}
\noindent\textbf{\footnotesize FOLIO 29 VERSO} \hfill \textbf{\footnotesize OCTOBRE 1916}\par\vspace{0.25mm}
\noindent{\fontsize{6.9pt}{8.0pt}\selectfont Méditation solitaire dans l'appartement toulousain et visite de Marguerite et de sa mère}\par\vspace{0.2mm}
\noindent{\fontsize{6.9pt}{8.0pt}\selectfont\textit{Lundi 2 octobre 1916} \hfill \textit{Toulouse (Appartement, Sainte-Marie, chapelle de Nazareth, rue d'Alsace-Lorraine)}}\par\vspace{0.5mm}
\noindent\rule{\textwidth}{0.35pt}\par\vspace{0.8mm}
% --- Corps de la transcription diplomatique et apparat marginal en prolongement ---
{\renewcommand{\arraystretch}{0.87}%
\noindent\hspace*{3.5mm}\begin{tabular}{@{}p{87mm}@{\hspace{3.5mm}}p{39mm}@{}}
      \leavevmode\fontsize{6.9pt}{7.9pt}\selectfont\strut \underline{2 L.} Temps splendide \hfill{} \raisebox{0.5ex}{\underline{Octbre 1916}}\strut & \smash{\parbox[t]{39mm}{\fontsize{6.2pt}{7.0pt}\selectfont\raggedright\color{black!88}\textbf{1 L.} \textit{scil.} Lundi.}} \\
      \leavevmode\fontsize{6.9pt}{7.9pt}\selectfont\strut Matinée, cme [tjs ?] resté appartt \& trouvé\strut & \smash{\parbox[t]{39mm}{\fontsize{6.2pt}{7.0pt}\selectfont\raggedright\color{black!88}\textbf{2} \textit{lect. incert. in ms.} \textbf{2 cme} \textit{abbr. pro} comme.}} \\
      \leavevmode\fontsize{6.9pt}{7.9pt}\selectfont\strut charme mélancolique à laisser souvenirs\strut &  \\
      \leavevmode\fontsize{6.9pt}{7.9pt}\selectfont\strut s'éveiller peu à peu ds mon esprit\strut & \smash{\parbox[t]{39mm}{\fontsize{6.2pt}{7.0pt}\selectfont\raggedright\color{black!88}\textbf{4 ds} \textit{abbr. pro} dans.}} \\
      \leavevmode\llap{\fontsize{6.4pt}{7.3pt}\selectfont\color{darkgray}5\hspace{1.8mm}}\fontsize{6.9pt}{7.9pt}\selectfont\strut Monté à Ste Marie vers 11H 1/2 –\strut & \smash{\parbox[t]{39mm}{\fontsize{6.2pt}{7.0pt}\selectfont\raggedright\color{black!88}\textbf{5 Ste} \textit{abbr. pro} Sainte.}} \\
      \leavevmode\fontsize{6.9pt}{7.9pt}\selectfont\strut Avant et surtout après déjeuner, ds le petit\strut & \smash{\parbox[t]{39mm}{\fontsize{6.2pt}{7.0pt}\selectfont\raggedright\color{black!88}\textbf{6 ds} \textit{abbr. pro} dans.}} \\
      \leavevmode\fontsize{6.9pt}{7.9pt}\selectfont\strut salon, très long \& intime tête à tête avec\strut &  \\
      \leavevmode\fontsize{6.9pt}{7.9pt}\selectfont\strut Mte — Ainsi pu deviner profondeur sentimt\strut & \smash{\parbox[t]{39mm}{\fontsize{6.2pt}{7.0pt}\selectfont\raggedright\color{black!88}\textbf{8 sentimt} \textit{abbr. pro} sentiment. \textbf{8 Mte} \textit{scil.} Marguerite (Calmels). — \textbf{10 s.} \textit{abbr. pro} soir. — \textbf{11 Margte} \textit{scil.} Marguerite (Calmels). — \textbf{12 appartt} \textit{abbr. pro} appartement. \textbf{12 Margte} \textit{scil.} Marguerite (Calmels). — \textbf{14} \textit{lect. incert. in ms.} — \textbf{15 Ts} \textit{abbr. pro} Tous.}} \\
      \leavevmode\fontsize{6.9pt}{7.9pt}\selectfont\strut qu'elle a pour moi...\strut &  \\
      \leavevmode\llap{\fontsize{6.4pt}{7.3pt}\selectfont\color{darkgray}10\hspace{1.8mm}}\fontsize{6.9pt}{7.9pt}\selectfont\strut Descendu seul vers 4 H s. Diverses courses\strut &  \\
      \leavevmode\fontsize{6.9pt}{7.9pt}\selectfont\strut Dès 5 H, chz moi. Visite de Margte \& de sa\strut &  \\
      \leavevmode\fontsize{6.9pt}{7.9pt}\selectfont\strut mère à mon appartt. Margte, charmante\strut &  \\
      \leavevmode\fontsize{6.9pt}{7.9pt}\selectfont\strut de simplicité \& de grâce\strut &  \\
      \leavevmode\fontsize{6.9pt}{7.9pt}\selectfont\strut A Chapelle Nazareth — à 6 H [1/2 / s. ?]\strut &  \\
      \leavevmode\llap{\fontsize{6.4pt}{7.3pt}\selectfont\color{darkgray}15\hspace{1.8mm}}\fontsize{6.9pt}{7.9pt}\selectfont\strut Ts 3, promené rue Alsace, regardé\strut &  \\
      \leavevmode\fontsize{6.9pt}{7.9pt}\selectfont\strut vitrines des magasins \& remonté tram Côte-Pavée\strut &  \\
      \leavevmode\fontsize{6.9pt}{7.9pt}\selectfont\strut A 9 H 30, av. Jean, partis —\strut &  \\
      \leavevmode\fontsize{6.9pt}{7.9pt}\selectfont\strut Très belle nuit étoilée — Air doux \& [parfumé ?]\strut & \smash{\parbox[t]{39mm}{\fontsize{6.2pt}{7.0pt}\selectfont\raggedright\color{black!88}\textbf{18} \textit{lect. incert. in ms.}}} \\
      \leavevmode\fontsize{6.9pt}{7.9pt}\selectfont\strut Trams : 10, 10, 30 – – Chapelet pr. Bne Maman : 1 f 25\strut & \smash{\parbox[t]{39mm}{\fontsize{6.2pt}{7.0pt}\selectfont\raggedright\color{black!88}\textbf{19 pr} \textit{abbr. pro} pour. \textbf{19 Bne} \textit{abbr. pro} Bonne.}} \\
\end{tabular}}%
\par\vspace*{\fill}
% --- Registres de bas de page (1. Critique interne [si pertinent] ; 2. Loci paralleli en latin ; 3. Commentaire historique en retrait) ---
\noindent\rule{\textwidth}{0.5pt}\par\vspace{0.65mm}
{\fontsize{6.6pt}{7.7pt}\selectfont%
\noindent\normalfont\color{black!92} \textbf{8} (\textit{Mte}) cf. supra f\textsuperscript{o}~1~r\textsuperscript{o} (p.~\pageref{fol:2}), f\textsuperscript{o}~29~r\textsuperscript{o} (p.~\pageref{fol:58}), infra f\textsuperscript{o}~30~r\textsuperscript{o} (p.~\pageref{fol:60}), f\textsuperscript{o}~31~r\textsuperscript{o} (p.~\pageref{fol:62}), \textit{et passim}, Appendicem I, doc.~5 (p.~\pageref{annexe:doc5}). — \textbf{17} (\textit{Jean}) cf. supra f\textsuperscript{o}~29~r\textsuperscript{o} (p.~\pageref{fol:58}), infra f\textsuperscript{o}~66~v\textsuperscript{o} (p.~\pageref{fol:133}), f\textsuperscript{o}~69~v\textsuperscript{o} (p.~\pageref{fol:139}).\par\vspace{0.55mm}
\noindent\rule{38mm}{0.4pt}\par\vspace{0.5mm}
\noindent\hspace*{5mm}\parbox{\dimexpr\textwidth-5mm\relax}{\fontsize{6.6pt}{7.7pt}\selectfont\color{black!84}Le lundi 2 octobre à 17~h, Marguerite et sa mère \textbf{Mme Valentine Calmels} rendent visite à Hippolyte Combes dans son appartement toulousain, puis tous trois se rendent à la chapelle \textbf{Notre-Dame-de-Nazareth} (rue Pharaon / quartier des Carmes, sanctuaire familial où Hippolyte Combes déclare au f\textsuperscript{o}~68~v\textsuperscript{o}, l.~10: \textit{A Nazareth, revu tte ma famille}) et flânent rue d'Alsace-Lorraine avant de remonter à la Côte Pavée, d'où Hippolyte Combes repart à 21~h~30 avec l'abbé \textbf{Jean Crouzet}.}\par}
\newpage
\phantomsection\label{fac:59}
% --- En-tête du fac-similé (Page impaire / droite : Photographie du manuscrit) ---
\noindent\rule{\textwidth}{0.5pt}\par\vspace{0.55mm}
\noindent\textbf{\fontsize{6.8pt}{7.8pt}\selectfont FAC-SIMILÉ EN VIS-À-VIS — FOLIO 29 VERSO} \hfill \textbf{\fontsize{6.8pt}{7.8pt}\selectfont OCTOBRE 1916}\par\vspace{0.25mm}
\noindent{\fontsize{6.8pt}{8.0pt}\selectfont\textit{Livre-Journal d'Hippolyte Combes} — Lundi 2 octobre 1916 (Toulouse (Appartement, Sainte-Marie, chapelle de Nazareth, rue d'Alsace-Lorraine))}\par\vspace{0.55mm}
\noindent\rule{\textwidth}{0.35pt}\par\vspace*{\fill}
{\centering
  \setlength{\fboxsep}{1.2mm}%
  \setlength{\fboxrule}{0.4pt}%
  \fcolorbox{gray!55}{white}{%
    \includegraphics[width=0.95\textwidth,height=0.81\textheight,keepaspectratio]{\facsimiledir/20251126_103321.jpg}%
  }\par}
\vspace*{\fill}
\newpage

% =============================================================================
% VIS-A-VIS SPREAD : FOLIO 30 RECTO (20251126_103335.jpg)
% =============================================================================
\cleardoubleevenpage
\renewcommand{\currentmonthyear}{Octobre 1916}
\renewcommand{\currentfoliolabel}{Folio 30 recto}
\phantomsection\label{fol:60}
% --- En-tête du folio (Page paire / gauche : Transcription, apparat & registres de bas de page) ---
\noindent\rule{\textwidth}{0.5pt}\par\vspace{0.55mm}
\noindent\textbf{\footnotesize FOLIO 30 RECTO} \hfill \textbf{\footnotesize OCTOBRE 1916}\par\vspace{0.25mm}
\noindent{\fontsize{6.9pt}{8.0pt}\selectfont Mardi 3 octobre 1916 — Messe à la cathédrale Saint-Étienne, déjeuner à l'Automate, promenade vers la route de Castanet et présentation de sa fiancée à Mme de Labarthe}\par\vspace{0.2mm}
\noindent{\fontsize{6.9pt}{8.0pt}\selectfont\textit{Mardi 3 octobre 1916} \hfill \textit{Toulouse (Cathédrale Saint-Étienne, rue d'Alsace-Lorraine, marché Victor-Hugo, route de Castanet)}}\par\vspace{0.5mm}
\noindent\rule{\textwidth}{0.35pt}\par\vspace{0.8mm}
% --- Corps de la transcription diplomatique et apparat marginal en prolongement ---
{\renewcommand{\arraystretch}{0.87}%
\noindent\hspace*{3.5mm}\begin{tabular}{@{}p{87mm}@{\hspace{3.5mm}}p{39mm}@{}}
      \leavevmode\fontsize{6.9pt}{7.9pt}\selectfont\strut \underline{3 Ma} \hfill{} \underline{Temps mervt beau \textbackslash\{\}hfill Octobre /16}\strut & \smash{\parbox[t]{39mm}{\fontsize{6.2pt}{7.0pt}\selectfont\raggedright\color{black!88}\textbf{1 mervt} \textit{abbr. pro} merveilleusement. — \textbf{2 St} \textit{abbr. pro} Saint. — \textbf{3 Margte} \textit{scil.} Marguerite (Calmels). \textbf{3 Mlle} \textit{abbr. pro} Mademoiselle. — \textbf{4 Margte} \textit{scil.} Marguerite (Calmels). \textbf{4 Mme} \textit{abbr. pro} Madame. — \textbf{8 pr} \textit{abbr. pro} pour. \textbf{8 qq} \textit{abbr. pro} quelque(s). — \textbf{9} \textit{lect. incert. in ms.} — \textbf{10} \textit{lect. incert. in ms.} \textbf{10 St} \textit{abbr. pro} Saint. — \textbf{12 Mme} \textit{abbr. pro} Madame.}} \\
      \leavevmode\fontsize{6.9pt}{7.9pt}\selectfont\strut Matin, à St-Etienne, messe annivre pour\strut &  \\
      \leavevmode\fontsize{6.9pt}{7.9pt}\selectfont\strut Mlle Calmels. A coté de Margte\strut &  \\
      \leavevmode\fontsize{6.9pt}{7.9pt}\selectfont\strut En sortant cathle, parlé av. Mme C. \& Margte\strut &  \\
      \leavevmode\llap{\fontsize{6.4pt}{7.3pt}\selectfont\color{darkgray}5\hspace{1.8mm}}\fontsize{6.9pt}{7.9pt}\selectfont\strut rue Alsace et déjeuné à l'Automate\strut &  \\
      \leavevmode\fontsize{6.9pt}{7.9pt}\selectfont\strut Fait ensemble marché à Victor Hugo\strut &  \\
      \leavevmode\fontsize{6.9pt}{7.9pt}\selectfont\strut Séparés devant magasins Capitole\strut &  \\
      \leavevmode\fontsize{6.9pt}{7.9pt}\selectfont\strut Rentré pr qq visites \& arrangt appartt\strut &  \\
      \leavevmode\fontsize{6.9pt}{7.9pt}\selectfont\strut Après déjeuner, venus [groupe abrégé illis.]\strut &  \\
      \leavevmode\llap{\fontsize{6.4pt}{7.3pt}\selectfont\color{darkgray}10\hspace{1.8mm}}\fontsize{6.9pt}{7.9pt}\selectfont\strut [Castanet/cravate ?]. Pris tram rue de St Joseph\strut &  \\
      \leavevmode\fontsize{6.9pt}{7.9pt}\selectfont\strut et descendus à terminus. Continué route à\strut &  \\
      \leavevmode\fontsize{6.9pt}{7.9pt}\selectfont\strut pied. Mme de Labarthe, charmante.\strut &  \\
      \leavevmode\fontsize{6.9pt}{7.9pt}\selectfont\strut Présenté ma Fiancée —\strut &  \\
      \leavevmode\fontsize{6.9pt}{7.9pt}\selectfont\strut Revenus à pied. Promenade très agréable\strut &  \\
      \leavevmode\llap{\fontsize{6.4pt}{7.3pt}\selectfont\color{darkgray}15\hspace{1.8mm}}\fontsize{6.9pt}{7.9pt}\selectfont\strut Soir, vers 8 H, pris tram \& venu voir\strut &  \\
      \leavevmode\fontsize{6.9pt}{7.9pt}\selectfont\strut Mme [Laudet/Lamolet ?].. —\strut & \smash{\parbox[t]{39mm}{\fontsize{6.2pt}{7.0pt}\selectfont\raggedright\color{black!88}\textbf{16} \textit{lect. incert. in ms.} \textbf{16 Mme} \textit{abbr. pro} Madame.}} \\
      \leavevmode\fontsize{6.9pt}{7.9pt}\selectfont\strut Trams : 20, 30 Marché V. Hugo :\strut &  \\
      \leavevmode\fontsize{6.9pt}{7.9pt}\selectfont\strut Automate : 30 + 20 c \hfill{} fruits : 80 c. Gâteaux : 3 f\strut & \smash{\parbox[t]{39mm}{\fontsize{6.2pt}{7.0pt}\selectfont\raggedright\color{black!88}\textbf{18 c.} \textit{abbr. pro} centimes.}} \\
\end{tabular}}%
\par\vspace*{\fill}
% --- Registres de bas de page (1. Critique interne [si pertinent] ; 2. Loci paralleli en latin ; 3. Commentaire historique en retrait) ---
\noindent\rule{\textwidth}{0.5pt}\par\vspace{0.65mm}
{\fontsize{6.6pt}{7.7pt}\selectfont%
\noindent\normalfont\color{black!92} \textbf{2} (\textit{St-Etienne}) cf. supra f\textsuperscript{o}~29~r\textsuperscript{o} (p.~\pageref{fol:58}). — \textbf{3} (\textit{Margte}) cf. supra f\textsuperscript{o}~1~r\textsuperscript{o} (p.~\pageref{fol:2}), f\textsuperscript{o}~29~v\textsuperscript{o} (p.~\pageref{fol:59}), infra f\textsuperscript{o}~31~r\textsuperscript{o} (p.~\pageref{fol:62}), f\textsuperscript{o}~32~r\textsuperscript{o} (p.~\pageref{fol:64}), \textit{et passim}.\par\vspace{0.55mm}
\noindent\rule{38mm}{0.4pt}\par\vspace{0.5mm}
\noindent\hspace*{5mm}\parbox{\dimexpr\textwidth-5mm\relax}{\fontsize{6.6pt}{7.7pt}\selectfont\color{black!84}Après une messe anniversaire à la cathédrale \textbf{Saint-Étienne} pour une parente de la famille Calmels (\textit{Mlle Calmels}), un passage au bar-restaurant automatique (\textit{l'Automate}) et les courses au marché couvert de la place \textbf{Victor-Hugo} avec Mme Calmels et Marguerite, Hippolyte Combes conduit l'après-midi sa fiancée et sa future belle-mère par le tramway de la rue Saint-Joseph jusqu'à la propriété de sa grande amie \textbf{Mme de Labarthe}, à qui il présente officiellement Marguerite.}\par}
\newpage
\phantomsection\label{fac:60}
% --- En-tête du fac-similé (Page impaire / droite : Photographie du manuscrit) ---
\noindent\rule{\textwidth}{0.5pt}\par\vspace{0.55mm}
\noindent\textbf{\fontsize{6.8pt}{7.8pt}\selectfont FAC-SIMILÉ EN VIS-À-VIS — FOLIO 30 RECTO} \hfill \textbf{\fontsize{6.8pt}{7.8pt}\selectfont OCTOBRE 1916}\par\vspace{0.25mm}
\noindent{\fontsize{6.8pt}{8.0pt}\selectfont\textit{Livre-Journal d'Hippolyte Combes} — Mardi 3 octobre 1916 (Toulouse (Cathédrale Saint-Étienne, rue d'Alsace-Lorraine, marché Victor-Hugo, route de Castanet))}\par\vspace{0.55mm}
\noindent\rule{\textwidth}{0.35pt}\par\vspace*{\fill}
{\centering
  \setlength{\fboxsep}{1.2mm}%
  \setlength{\fboxrule}{0.4pt}%
  \fcolorbox{gray!55}{white}{%
    \includegraphics[width=0.95\textwidth,height=0.81\textheight,keepaspectratio]{\facsimiledir/20251126_103335.jpg}%
  }\par}
\vspace*{\fill}
\newpage

% =============================================================================
% VIS-A-VIS SPREAD : FOLIO 30 VERSO (20251126_103455.jpg)
% =============================================================================
\cleardoubleevenpage
\renewcommand{\currentmonthyear}{Octobre 1916}
\renewcommand{\currentfoliolabel}{Folio 30 verso}
\phantomsection\label{fol:61}
% --- En-tête du folio (Page paire / gauche : Transcription, apparat & registres de bas de page) ---
\noindent\rule{\textwidth}{0.5pt}\par\vspace{0.55mm}
\noindent\textbf{\footnotesize FOLIO 30 VERSO} \hfill \textbf{\footnotesize OCTOBRE 1916}\par\vspace{0.25mm}
\noindent{\fontsize{6.9pt}{8.0pt}\selectfont Règlement du tailleur Guiraud, du loyer semestriel à Mlle Guilhem et thé chez Mme de Labarthe}\par\vspace{0.2mm}
\noindent{\fontsize{6.9pt}{8.0pt}\selectfont\textit{Mardi 3 octobre 1916 (suite)} \hfill \textit{Toulouse}}\par\vspace{0.5mm}
\noindent\rule{\textwidth}{0.35pt}\par\vspace{0.8mm}
% --- Corps de la transcription diplomatique et apparat marginal en prolongement ---
{\renewcommand{\arraystretch}{1.00}%
\noindent\hspace*{3.5mm}\begin{tabular}{@{}p{87mm}@{\hspace{3.5mm}}p{39mm}@{}}
      \leavevmode\fontsize{7.7pt}{9.0pt}\selectfont\strut \underline{3 Ma} \hfill{} \underline{Octbre 1916}\strut &  \\
      \leavevmode\fontsize{7.7pt}{9.0pt}\selectfont\strut – Payé à Guiraud, pour façon\strut &  \\
      \leavevmode\fontsize{7.7pt}{9.0pt}\selectfont\strut \& fournitures tunique : 38 f\strut &  \\
      \leavevmode\fontsize{7.7pt}{9.0pt}\selectfont\strut – Semestre à Mlle Guilhem : 90 f –\strut & \smash{\parbox[t]{39mm}{\fontsize{6.2pt}{7.0pt}\selectfont\raggedright\color{black!88}\textbf{4 Mlle} \textit{abbr. pro} Mademoiselle. — \textbf{5} \textit{lect. incert. in ms.} — \textbf{6 Cie} \textit{lect. incert. in ms.}}} \\
      \leavevmode\llap{\fontsize{6.4pt}{7.3pt}\selectfont\color{darkgray}5\hspace{1.8mm}}\fontsize{7.7pt}{9.0pt}\selectfont\strut (Juillet [1916 ajouté sous la ligne] — Janvier 1917)\strut &  \\
      \leavevmode\fontsize{7.7pt}{9.0pt}\selectfont\strut — Cie [?] Française : 20 cm doublure satin : 60 c\strut &  \\
      \leavevmode\fontsize{7.7pt}{9.0pt}\selectfont\strut – Fleurs : 1f 10\strut &  \\
      \leavevmode\fontsize{7.7pt}{9.0pt}\selectfont\strut Délicieux ds la propriété de Mme de Labarthe\strut & \smash{\parbox[t]{39mm}{\fontsize{6.2pt}{7.0pt}\selectfont\raggedright\color{black!88}\textbf{8 ds} \textit{abbr. pro} dans. \textbf{8 Mme} \textit{abbr. pro} Madame. — \textbf{9 ns} \textit{abbr. pro} nous.}} \\
      \leavevmode\fontsize{7.7pt}{9.0pt}\selectfont\strut Elle ns a fait visiter maison.\strut &  \\
      \leavevmode\llap{\fontsize{6.4pt}{7.3pt}\selectfont\color{darkgray}10\hspace{1.8mm}}\fontsize{7.7pt}{9.0pt}\selectfont\strut Simplicité \& abandon charmants\strut &  \\
\end{tabular}}%
\par\vspace*{\fill}
% --- Registres de bas de page (1. Critique interne [si pertinent] ; 2. Loci paralleli en latin ; 3. Commentaire historique en retrait) ---
\noindent\rule{\textwidth}{0.5pt}\par\vspace{0.65mm}
{\fontsize{6.6pt}{7.7pt}\selectfont%
\noindent\normalfont \textbf{Tableau financier :} Page entièrement consacrée au bilan comptable de la permission toulousaine du 3 octobre 1916, articulé autour d'un solde encadré (\textit{En caisse : 272 f.}) suivi des dépenses détaillées.\par\vspace{0.5mm}
\noindent\rule{22mm}{0.3pt}\par\vspace{0.45mm}
\noindent\normalfont\color{black!92} \textbf{1} (\textit{de eodem commorationis loco}) cf. supra f\textsuperscript{o}~30~r\textsuperscript{o} (p.~\pageref{fol:60}).\par\vspace{0.55mm}
\noindent\rule{38mm}{0.4pt}\par\vspace{0.5mm}
\noindent\hspace*{5mm}\parbox{\dimexpr\textwidth-5mm\relax}{\fontsize{6.6pt}{7.7pt}\selectfont\color{black!84}Avant de quitter Toulouse, Hippolyte Combes règle 38~fr. à son tailleur \textbf{Guiraud} pour la façon et les fournitures de sa tunique militaire (plus 60~c. de doublure de satin) ainsi que 90~fr. à sa propriétaire \textbf{Mlle Guilhem} pour le loyer semestriel (juillet 1916 -- janvier 1917) de son appartement toulousain, où il conserve ses meubles et souvenirs de famille.}\par}
\newpage
\phantomsection\label{fac:61}
% --- En-tête du fac-similé (Page impaire / droite : Photographie du manuscrit) ---
\noindent\rule{\textwidth}{0.5pt}\par\vspace{0.55mm}
\noindent\textbf{\fontsize{6.8pt}{7.8pt}\selectfont FAC-SIMILÉ EN VIS-À-VIS — FOLIO 30 VERSO} \hfill \textbf{\fontsize{6.8pt}{7.8pt}\selectfont OCTOBRE 1916}\par\vspace{0.25mm}
\noindent{\fontsize{6.8pt}{8.0pt}\selectfont\textit{Livre-Journal d'Hippolyte Combes} — Mardi 3 octobre 1916 (suite) (Toulouse)}\par\vspace{0.55mm}
\noindent\rule{\textwidth}{0.35pt}\par\vspace*{\fill}
{\centering
  \setlength{\fboxsep}{1.2mm}%
  \setlength{\fboxrule}{0.4pt}%
  \fcolorbox{gray!55}{white}{%
    \includegraphics[width=0.95\textwidth,height=0.81\textheight,keepaspectratio]{\facsimiledir/20251126_103455.jpg}%
  }\par}
\vspace*{\fill}
\newpage

% =============================================================================
% VIS-A-VIS SPREAD : FOLIO 31 RECTO (20251126_103514.jpg)
% =============================================================================
\cleardoubleevenpage
\renewcommand{\currentmonthyear}{Octobre 1916}
\renewcommand{\currentfoliolabel}{Folio 31 recto}
\phantomsection\label{fol:62}
% --- En-tête du folio (Page paire / gauche : Transcription, apparat & registres de bas de page) ---
\noindent\rule{\textwidth}{0.5pt}\par\vspace{0.55mm}
\noindent\textbf{\footnotesize FOLIO 31 RECTO} \hfill \textbf{\footnotesize OCTOBRE 1916}\par\vspace{0.25mm}
\noindent{\fontsize{6.9pt}{8.0pt}\selectfont Mercredi 4 octobre 1916 — Dernier après-midi à Sainte-Marie et séparation déchirante à la gare Matabiau}\par\vspace{0.2mm}
\noindent{\fontsize{6.9pt}{8.0pt}\selectfont\textit{Mercredi 4 octobre 1916} \hfill \textit{Toulouse (Sainte-Marie et gare Matabiau)}}\par\vspace{0.5mm}
\noindent\rule{\textwidth}{0.35pt}\par\vspace{0.8mm}
% --- Corps de la transcription diplomatique et apparat marginal en prolongement ---
{\renewcommand{\arraystretch}{0.87}%
\noindent\hspace*{3.5mm}\begin{tabular}{@{}p{87mm}@{\hspace{3.5mm}}p{39mm}@{}}
      \leavevmode\fontsize{6.9pt}{7.9pt}\selectfont\strut \underline{Mer. 4} \hfill{} \underline{Octbre 16}\strut & \smash{\parbox[t]{39mm}{\fontsize{6.2pt}{7.0pt}\selectfont\raggedright\color{black!88}\textbf{1 Mer.} \textit{scil.} Mercredi.}} \\
      \leavevmode\fontsize{6.9pt}{7.9pt}\selectfont\strut Sur pied dès 6H. –\strut &  \\
      \leavevmode\fontsize{6.9pt}{7.9pt}\selectfont\strut Lentement préparé départ. Enlevé \& rentré\strut &  \\
      \leavevmode\fontsize{6.9pt}{7.9pt}\selectfont\strut rideaux – mis tt en ordre ds l'appartt\strut & \smash{\parbox[t]{39mm}{\fontsize{6.2pt}{7.0pt}\selectfont\raggedright\color{black!88}\textbf{4 ds} \textit{abbr. pro} dans. \textbf{4 tt} \textit{abbr. pro} tout.}} \\
      \leavevmode\llap{\fontsize{6.4pt}{7.3pt}\selectfont\color{darkgray}5\hspace{1.8mm}}\fontsize{6.9pt}{7.9pt}\selectfont\strut Venu gare. puis aux Thermes du Languedoc\strut &  \\
      \leavevmode\fontsize{6.9pt}{7.9pt}\selectfont\strut A 11 H m., quitté maison. Monté à\strut & \smash{\parbox[t]{39mm}{\fontsize{6.2pt}{7.0pt}\selectfont\raggedright\color{black!88}\textbf{6 m.} \textit{abbr. pro} matin.}} \\
      \leavevmode\fontsize{6.9pt}{7.9pt}\selectfont\strut Ste Marie –\strut & \smash{\parbox[t]{39mm}{\fontsize{6.2pt}{7.0pt}\selectfont\raggedright\color{black!88}\textbf{7 Ste} \textit{abbr. pro} Sainte.}} \\
      \leavevmode\fontsize{6.9pt}{7.9pt}\selectfont\strut Passé tt l'après-midi en cie de Mte\strut & \smash{\parbox[t]{39mm}{\fontsize{6.2pt}{7.0pt}\selectfont\raggedright\color{black!88}\textbf{8 tt} \textit{abbr. pro} tout. \textbf{8 Mte} \textit{scil.} Marguerite (Calmels).}} \\
      \leavevmode\fontsize{6.9pt}{7.9pt}\selectfont\strut Longs tête à tête très fort intimes.\strut &  \\
      \leavevmode\llap{\fontsize{6.4pt}{7.3pt}\selectfont\color{darkgray}10\hspace{1.8mm}}\fontsize{6.9pt}{7.9pt}\selectfont\strut Promené au fond du parc. Joué\strut &  \\
      \leavevmode\fontsize{6.9pt}{7.9pt}\selectfont\strut avec les enfants\strut &  \\
      \leavevmode\fontsize{6.9pt}{7.9pt}\selectfont\strut Puis causé lgt, assis sous gd arbre\strut & \smash{\parbox[t]{39mm}{\fontsize{6.2pt}{7.0pt}\selectfont\raggedright\color{black!88}\textbf{12 lgt} \textit{abbr. pro} longtemps. \textbf{12 gd} \textit{abbr. pro} grand. — \textbf{13} \textit{lect. incert. in ms.} — \textbf{14 absolt} \textit{abbr. pro} absolument. \textbf{14 Mte} \textit{scil.} Marguerite (Calmels). — \textbf{15} \textit{lect. incert. in ms.} — \textbf{16 ds} \textit{abbr. pro} dans. — \textbf{17 Mte} \textit{scil.} Marguerite (Calmels). — \textbf{20 ns} \textit{abbr. pro} nous.}} \\
      \leavevmode\fontsize{6.9pt}{7.9pt}\selectfont\strut Diner à 6 H [1/2 ?]\strut &  \\
      \leavevmode\fontsize{6.9pt}{7.9pt}\selectfont\strut Mte, voulu absolt m'accompagner gare\strut &  \\
      \leavevmode\llap{\fontsize{6.4pt}{7.3pt}\selectfont\color{darkgray}15\hspace{1.8mm}}\fontsize{6.9pt}{7.9pt}\selectfont\strut [groupe surchargé illis.]\strut &  \\
      \leavevmode\fontsize{6.9pt}{7.9pt}\selectfont\strut \& choisi coin ds 1 compartiment\strut &  \\
      \leavevmode\fontsize{6.9pt}{7.9pt}\selectfont\strut Puis ressorti salle d'attente. Avec Mte\strut &  \\
      \leavevmode\fontsize{6.9pt}{7.9pt}\selectfont\strut \& oncle Crouzet. promené au dehors\strut &  \\
      \leavevmode\fontsize{6.9pt}{7.9pt}\selectfont\strut le long trottoirs cour de G.V.\strut &  \\
      \leavevmode\llap{\fontsize{6.4pt}{7.3pt}\selectfont\color{darkgray}20\hspace{1.8mm}}\fontsize{6.9pt}{7.9pt}\selectfont\strut Enfin, dû ns séparer, par\strut &  \\
      \leavevmode\fontsize{6.9pt}{7.9pt}\selectfont\strut ménagement pour l'oncle —\strut &  \\
\end{tabular}}%
\par\vspace*{\fill}
% --- Registres de bas de page (1. Critique interne [si pertinent] ; 2. Loci paralleli en latin ; 3. Commentaire historique en retrait) ---
\noindent\rule{\textwidth}{0.5pt}\par\vspace{0.65mm}
{\fontsize{6.6pt}{7.7pt}\selectfont%
\noindent\normalfont\color{black!92} \textbf{8} (\textit{Mte}) cf. supra f\textsuperscript{o}~1~r\textsuperscript{o} (p.~\pageref{fol:2}), f\textsuperscript{o}~30~r\textsuperscript{o} (p.~\pageref{fol:60}), infra f\textsuperscript{o}~32~r\textsuperscript{o} (p.~\pageref{fol:64}), f\textsuperscript{o}~56~r\textsuperscript{o} (p.~\pageref{fol:112}), \textit{et passim}, Appendicem I, doc.~5 (p.~\pageref{annexe:doc5}). — \textbf{18} (\textit{Crouzet}) cf. supra f\textsuperscript{o}~27~r\textsuperscript{o} (p.~\pageref{fol:54}), f\textsuperscript{o}~27~v\textsuperscript{o} (p.~\pageref{fol:55}), infra f\textsuperscript{o}~32~r\textsuperscript{o} (p.~\pageref{fol:64}), f\textsuperscript{o}~66~r\textsuperscript{o} (p.~\pageref{fol:132}), \textit{et passim}.\par\vspace{0.55mm}
\noindent\rule{38mm}{0.4pt}\par\vspace{0.5mm}
\noindent\hspace*{5mm}\parbox{\dimexpr\textwidth-5mm\relax}{\fontsize{6.6pt}{7.7pt}\selectfont\color{black!84}Dernière journée de la permission d'automne à Toulouse (mercredi 4 octobre 1916): après avoir fermé son appartement et pris un bain aux \textbf{Thermes du Languedoc}, Hippolyte Combes passe l'après-midi auprès de Marguerite dans le parc de Sainte-Marie, puis gagne le soir la \textbf{gare Matabiau} (cour de la Grande Vitesse, \textit{G.V.}) accompagné de sa fiancée et de l'oncle \textbf{Émile Crouzet} (1842--1921).}\par}
\newpage
\phantomsection\label{fac:62}
% --- En-tête du fac-similé (Page impaire / droite : Photographie du manuscrit) ---
\noindent\rule{\textwidth}{0.5pt}\par\vspace{0.55mm}
\noindent\textbf{\fontsize{6.8pt}{7.8pt}\selectfont FAC-SIMILÉ EN VIS-À-VIS — FOLIO 31 RECTO} \hfill \textbf{\fontsize{6.8pt}{7.8pt}\selectfont OCTOBRE 1916}\par\vspace{0.25mm}
\noindent{\fontsize{6.8pt}{8.0pt}\selectfont\textit{Livre-Journal d'Hippolyte Combes} — Mercredi 4 octobre 1916 (Toulouse (Sainte-Marie et gare Matabiau))}\par\vspace{0.55mm}
\noindent\rule{\textwidth}{0.35pt}\par\vspace*{\fill}
{\centering
  \setlength{\fboxsep}{1.2mm}%
  \setlength{\fboxrule}{0.4pt}%
  \fcolorbox{gray!55}{white}{%
    \includegraphics[width=0.95\textwidth,height=0.81\textheight,keepaspectratio]{\facsimiledir/20251126_103514.jpg}%
  }\par}
\vspace*{\fill}
\newpage

% =============================================================================
% VIS-A-VIS SPREAD : FOLIO 31 VERSO (20251126_103538.jpg)
% =============================================================================
\cleardoubleevenpage
\renewcommand{\currentmonthyear}{Octobre 1916}
\renewcommand{\currentfoliolabel}{Folio 31 verso}
\phantomsection\label{fol:63}
% --- En-tête du folio (Page paire / gauche : Transcription, apparat & registres de bas de page) ---
\noindent\rule{\textwidth}{0.5pt}\par\vspace{0.55mm}
\noindent\textbf{\footnotesize FOLIO 31 VERSO} \hfill \textbf{\footnotesize OCTOBRE 1916}\par\vspace{0.25mm}
\noindent{\fontsize{6.9pt}{8.0pt}\selectfont Traversée de Paris sous la pluie : l'église Saint-Vincent-de-Paul et les salons du Louvre}\par\vspace{0.2mm}
\noindent{\fontsize{6.9pt}{8.0pt}\selectfont\textit{Jeudi 5 octobre 1916} \hfill \textit{Paris (Gare d'Orsay, Gare du Nord, église Saint-Vincent-de-Paul, Grands Magasins du Louvre)}}\par\vspace{0.5mm}
\noindent\rule{\textwidth}{0.35pt}\par\vspace{0.8mm}
% --- Corps de la transcription diplomatique et apparat marginal en prolongement ---
{\renewcommand{\arraystretch}{0.93}%
\noindent\hspace*{3.5mm}\begin{tabular}{@{}p{87mm}@{\hspace{3.5mm}}p{39mm}@{}}
      \leavevmode\fontsize{7.3pt}{8.4pt}\selectfont\strut \underline{Jeudi 5} \hfill{} \underline{Octb / 16}\strut &  \\
      \leavevmode\fontsize{7.3pt}{8.4pt}\selectfont\strut Le jour parut peu à peu\strut &  \\
      \leavevmode\fontsize{7.3pt}{8.4pt}\selectfont\strut A 8H m., gare d'Orsay\strut & \smash{\parbox[t]{39mm}{\fontsize{6.2pt}{7.0pt}\selectfont\raggedright\color{black!88}\textbf{3 m.} \textit{abbr. pro} matin.}} \\
      \leavevmode\fontsize{7.3pt}{8.4pt}\selectfont\strut Sous pluie fine \& serrée, venu à pied\strut &  \\
      \leavevmode\llap{\fontsize{6.4pt}{7.3pt}\selectfont\color{darkgray}5\hspace{1.8mm}}\fontsize{7.3pt}{8.4pt}\selectfont\strut station métro Tuileries. Changé au Chatelet pr\strut & \smash{\parbox[t]{39mm}{\fontsize{6.2pt}{7.0pt}\selectfont\raggedright\color{black!88}\textbf{5 pr} \textit{abbr. pro} pour.}} \\
      \leavevmode\fontsize{7.3pt}{8.4pt}\selectfont\strut gare du Nord\strut &  \\
      \leavevmode\fontsize{7.3pt}{8.4pt}\selectfont\strut Laissé paquets consigne gare du Nord\strut &  \\
      \leavevmode\fontsize{7.3pt}{8.4pt}\selectfont\strut Descendu en flânant rue La Fayette\strut &  \\
      \leavevmode\fontsize{7.3pt}{8.4pt}\selectfont\strut Entré et repos Egl St-Vincent-de-Paul.\strut & \smash{\parbox[t]{39mm}{\fontsize{6.2pt}{7.0pt}\selectfont\raggedright\color{black!88}\textbf{9 Egl.} \textit{abbr. pro} Église. \textbf{9 St} \textit{abbr. pro} Saint.}} \\
      \leavevmode\llap{\fontsize{6.4pt}{7.3pt}\selectfont\color{darkgray}10\hspace{1.8mm}}\fontsize{7.3pt}{8.4pt}\selectfont\strut Admiré lentement peintures murales\strut &  \\
      \leavevmode\fontsize{7.3pt}{8.4pt}\selectfont\strut et rangées superposées de belles colonnes\strut &  \\
      \leavevmode\fontsize{7.3pt}{8.4pt}\selectfont\strut de marbre marron\strut &  \\
      \leavevmode\fontsize{7.3pt}{8.4pt}\selectfont\strut Continué promenade j-aux Mag\strut & \smash{\parbox[t]{39mm}{\fontsize{6.2pt}{7.0pt}\selectfont\raggedright\color{black!88}\textbf{13 j-aux} \textit{abbr. pro} jusqu'aux.}} \\
      \leavevmode\fontsize{7.3pt}{8.4pt}\selectfont\strut du Louvre. Monté au salon de\strut &  \\
      \leavevmode\llap{\fontsize{6.4pt}{7.3pt}\selectfont\color{darkgray}15\hspace{1.8mm}}\fontsize{7.3pt}{8.4pt}\selectfont\strut Corresp – et écrit 2 lettres\strut & \smash{\parbox[t]{39mm}{\fontsize{6.2pt}{7.0pt}\selectfont\raggedright\color{black!88}\textbf{15 Corresp} \textit{abbr. pro} Correspondance.}} \\
      \leavevmode\fontsize{7.3pt}{8.4pt}\selectfont\strut Pluie cesse. Soleil commence briller 1 peu\strut &  \\
      \leavevmode\fontsize{7.3pt}{8.4pt}\selectfont\strut Trop long arrêt au Louvre. Pas\strut &  \\
      \leavevmode\fontsize{7.3pt}{8.4pt}\selectfont\strut assez temps pr déjeuner restaurant\strut & \smash{\parbox[t]{39mm}{\fontsize{6.2pt}{7.0pt}\selectfont\raggedright\color{black!88}\textbf{18 pr} \textit{abbr. pro} pour.}} \\
      \leavevmode\fontsize{7.3pt}{8.4pt}\selectfont\strut Dû revenir par métro à gare du Nord\strut &  \\
\end{tabular}}%
\par\vspace*{\fill}
% --- Registres de bas de page (1. Critique interne [si pertinent] ; 2. Loci paralleli en latin ; 3. Commentaire historique en retrait) ---
\noindent\rule{\textwidth}{0.5pt}\par\vspace{0.65mm}
{\fontsize{6.6pt}{7.7pt}\selectfont%
\noindent\normalfont \textbf{Inversion de la séquence photographique :} Le f\textsuperscript{o}~31~v\textsuperscript{o} (journée du jeudi 5 octobre 1916 à Paris et retour vers le front) et le f\textsuperscript{o}~32~r\textsuperscript{o} (soirée du mercredi 4 octobre et nuit du 4 au 5 octobre dans l'express Toulouse–Paris) ont été photographiés dans l'ordre inverse de la chronologie.\par\vspace{0.5mm}
\noindent\rule{22mm}{0.3pt}\par\vspace{0.45mm}
\noindent\normalfont\color{black!92} \textbf{1} (\textit{Jeudi 5}) de ordine chronologico cf. infra f\textsuperscript{o}~32~r\textsuperscript{o} (p.~\pageref{fol:64}).\par\vspace{0.55mm}
\noindent\rule{38mm}{0.4pt}\par\vspace{0.5mm}
\noindent\hspace*{5mm}\parbox{\dimexpr\textwidth-5mm\relax}{\fontsize{6.6pt}{7.7pt}\selectfont\color{black!84}Traversée de Paris le jeudi 5 octobre 1916 entre la gare d'Orsay et la gare du Nord.}\par}
\newpage
\phantomsection\label{fac:63}
% --- En-tête du fac-similé (Page impaire / droite : Photographie du manuscrit) ---
\noindent\rule{\textwidth}{0.5pt}\par\vspace{0.55mm}
\noindent\textbf{\fontsize{6.8pt}{7.8pt}\selectfont FAC-SIMILÉ EN VIS-À-VIS — FOLIO 31 VERSO} \hfill \textbf{\fontsize{6.8pt}{7.8pt}\selectfont OCTOBRE 1916}\par\vspace{0.25mm}
\noindent{\fontsize{6.8pt}{8.0pt}\selectfont\textit{Livre-Journal d'Hippolyte Combes} — Jeudi 5 octobre 1916 (Paris (Gare d'Orsay, Gare du Nord, église Saint-Vincent-de-Paul, Grands Magasins du Louvre))}\par\vspace{0.55mm}
\noindent\rule{\textwidth}{0.35pt}\par\vspace*{\fill}
{\centering
  \setlength{\fboxsep}{1.2mm}%
  \setlength{\fboxrule}{0.4pt}%
  \fcolorbox{gray!55}{white}{%
    \includegraphics[width=0.95\textwidth,height=0.81\textheight,keepaspectratio]{\facsimiledir/20251126_103538.jpg}%
  }\par}
\vspace*{\fill}
\newpage

% =============================================================================
% VIS-A-VIS SPREAD : FOLIO 32 RECTO (20251126_103554.jpg)
% =============================================================================
\cleardoubleevenpage
\renewcommand{\currentmonthyear}{Octobre 1916}
\renewcommand{\currentfoliolabel}{Folio 32 recto}
\phantomsection\label{fol:64}
% --- En-tête du folio (Page paire / gauche : Transcription, apparat & registres de bas de page) ---
\noindent\rule{\textwidth}{0.5pt}\par\vspace{0.55mm}
\noindent\textbf{\footnotesize FOLIO 32 RECTO} \hfill \textbf{\footnotesize OCTOBRE 1916}\par\vspace{0.25mm}
\noindent{\fontsize{6.9pt}{8.0pt}\selectfont Fin des adieux en gare de Toulouse, nuit blanche dans l'express d'Orsay et comptes de départ}\par\vspace{0.2mm}
\noindent{\fontsize{6.9pt}{8.0pt}\selectfont\textit{Mercredi 4 octobre 1916 (soir et nuit)} \hfill \textit{Gare de Toulouse et train de nuit Toulouse – Paris-Orsay}}\par\vspace{0.5mm}
\noindent\rule{\textwidth}{0.35pt}\par\vspace{0.8mm}
% --- Corps de la transcription diplomatique et apparat marginal en prolongement ---
{\renewcommand{\arraystretch}{0.87}%
\noindent\hspace*{3.5mm}\begin{tabular}{@{}p{87mm}@{\hspace{3.5mm}}p{39mm}@{}}
      \leavevmode\fontsize{6.9pt}{7.9pt}\selectfont\strut ... Mte \& moi, très\strut & \smash{\parbox[t]{39mm}{\fontsize{6.2pt}{7.0pt}\selectfont\raggedright\color{black!88}\textbf{1 Mte} \textit{scil.} Marguerite (Calmels). — \textbf{2 Mer.} \textit{scil.} Mercredi. — \textbf{3 s.} \textit{abbr. pro} soir.}} \\
      \leavevmode\fontsize{6.9pt}{7.9pt}\selectfont\strut émus.. – \hfill{} Mer. 4 Oct 16\strut &  \\
      \leavevmode\fontsize{6.9pt}{7.9pt}\selectfont\strut Train, démarré vers 8 H s —\strut &  \\
      \leavevmode\fontsize{6.9pt}{7.9pt}\selectfont\strut Voyage – toute la nuit. – n'ai pu\strut &  \\
      \leavevmode\llap{\fontsize{6.4pt}{7.3pt}\selectfont\color{darkgray}5\hspace{1.8mm}}\fontsize{6.9pt}{7.9pt}\selectfont\strut m'assoupir 1 peu que sur le\strut &  \\
      \leavevmode\fontsize{6.9pt}{7.9pt}\selectfont\strut matin.\strut &  \\
      \leavevmode\fontsize{6.9pt}{7.9pt}\selectfont\strut Hier 2 images de Ste Thérèse — : 40 c\strut & \smash{\parbox[t]{39mm}{\fontsize{6.2pt}{7.0pt}\selectfont\raggedright\color{black!88}\textbf{7 Ste} \textit{abbr. pro} Sainte.}} \\
      \leavevmode\fontsize{6.9pt}{7.9pt}\selectfont\strut \& Fleurs pour Mte : 1f 70\strut & \smash{\parbox[t]{39mm}{\fontsize{6.2pt}{7.0pt}\selectfont\raggedright\color{black!88}\textbf{8 Mte} \textit{scil.} Marguerite (Calmels).}} \\
      \leavevmode\fontsize{6.9pt}{7.9pt}\selectfont\strut – Trams : 20, 10, 30 c\strut &  \\
      \leavevmode\llap{\fontsize{6.4pt}{7.3pt}\selectfont\color{darkgray}10\hspace{1.8mm}}\fontsize{6.9pt}{7.9pt}\selectfont\strut Billet 3e cl. p. Orsay : 11 f 75\strut &  \\
      \cline{1-1}
      \leavevmode\fontsize{6.9pt}{7.9pt}\selectfont\strut Bain, aux Thermes – : 90 + 10c\strut &  \\
      \leavevmode\fontsize{6.9pt}{7.9pt}\selectfont\strut Etrenne à Maria, bonne des\strut &  \\
      \leavevmode\fontsize{6.9pt}{7.9pt}\selectfont\strut Crouzet : 5 f –\strut &  \\
\end{tabular}}%
\par\vspace*{\fill}
% --- Registres de bas de page (1. Critique interne [si pertinent] ; 2. Loci paralleli en latin ; 3. Commentaire historique en retrait) ---
\noindent\rule{\textwidth}{0.5pt}\par\vspace{0.65mm}
{\fontsize{6.6pt}{7.7pt}\selectfont%
\noindent\normalfont \textbf{Chronologie interne :} Ce feuillet (fin de la journée du 4 octobre 1916 au départ de Toulouse et voyage de nuit vers Paris) précède chronologiquement le f\textsuperscript{o}~31~v\textsuperscript{o}.\par\vspace{0.5mm}
\noindent\rule{22mm}{0.3pt}\par\vspace{0.45mm}
\noindent\normalfont\color{black!92} \textbf{1} (\textit{4 au 5 Oct.}) de ordine chronologico cf. supra f\textsuperscript{o}~31~v\textsuperscript{o} (p.~\pageref{fol:63}). — \textbf{1} (\textit{Mte}) cf. supra f\textsuperscript{o}~1~r\textsuperscript{o} (p.~\pageref{fol:2}), f\textsuperscript{o}~31~r\textsuperscript{o} (p.~\pageref{fol:62}), infra f\textsuperscript{o}~56~r\textsuperscript{o} (p.~\pageref{fol:112}), f\textsuperscript{o}~57~r\textsuperscript{o} (p.~\pageref{fol:114}), \textit{et passim}, Appendicem I, doc.~5 (p.~\pageref{annexe:doc5}). — \textbf{13} (\textit{Crouzet}) cf. supra f\textsuperscript{o}~27~r\textsuperscript{o} (p.~\pageref{fol:54}), f\textsuperscript{o}~31~r\textsuperscript{o} (p.~\pageref{fol:62}), infra f\textsuperscript{o}~66~r\textsuperscript{o} (p.~\pageref{fol:132}), f\textsuperscript{o}~66~v\textsuperscript{o} (p.~\pageref{fol:133}), \textit{et passim}.\par\vspace{0.55mm}
\noindent\rule{38mm}{0.4pt}\par\vspace{0.5mm}
\noindent\hspace*{5mm}\parbox{\dimexpr\textwidth-5mm\relax}{\fontsize{6.6pt}{7.7pt}\selectfont\color{black!84}Suite immédiate du f\textsuperscript{o}~31~r\textsuperscript{o} (mercredi 4 octobre 1916 au soir à la gare Matabiau de Toulouse, photographiée après le f\textsuperscript{o}~31~v\textsuperscript{o}): pendant le trajet nocturne de Toulouse à Paris-Orsay (billet militaire de 3\textsuperscript{e} classe à 11~fr.~75), Hippolyte Combes récapitule ses dernières dépenses toulousaines, dont deux images pieuses de sainte Thérèse (40~c.), un bouquet pour Marguerite (1~fr.~70) et 5~fr. d'étrennes à \textbf{Maria}, domestique de la famille Crouzet à Sainte-Marie.}\par}
\newpage
\phantomsection\label{fac:64}
% --- En-tête du fac-similé (Page impaire / droite : Photographie du manuscrit) ---
\noindent\rule{\textwidth}{0.5pt}\par\vspace{0.55mm}
\noindent\textbf{\fontsize{6.8pt}{7.8pt}\selectfont FAC-SIMILÉ EN VIS-À-VIS — FOLIO 32 RECTO} \hfill \textbf{\fontsize{6.8pt}{7.8pt}\selectfont OCTOBRE 1916}\par\vspace{0.25mm}
\noindent{\fontsize{6.8pt}{8.0pt}\selectfont\textit{Livre-Journal d'Hippolyte Combes} — Mercredi 4 octobre 1916 (soir et nuit) (Gare de Toulouse et train de nuit Toulouse – Paris-Orsay)}\par\vspace{0.55mm}
\noindent\rule{\textwidth}{0.35pt}\par\vspace*{\fill}
{\centering
  \setlength{\fboxsep}{1.2mm}%
  \setlength{\fboxrule}{0.4pt}%
  \fcolorbox{gray!55}{white}{%
    \includegraphics[width=0.95\textwidth,height=0.81\textheight,keepaspectratio]{\facsimiledir/20251126_103554.jpg}%
  }\par}
\vspace*{\fill}
\newpage

% =============================================================================
% VIS-A-VIS SPREAD : FOLIO 32 VERSO (20251126_103743.jpg)
% =============================================================================
\cleardoubleevenpage
\renewcommand{\currentmonthyear}{Octobre 1916}
\renewcommand{\currentfoliolabel}{Folio 32 verso}
\phantomsection\label{fol:65}
% --- En-tête du folio (Page paire / gauche : Transcription, apparat & registres de bas de page) ---
\noindent\rule{\textwidth}{0.5pt}\par\vspace{0.55mm}
\noindent\textbf{\footnotesize FOLIO 32 VERSO} \hfill \textbf{\footnotesize OCTOBRE 1916}\par\vspace{0.25mm}
\noindent{\fontsize{6.9pt}{8.0pt}\selectfont Train de l’après-midi de Paris-Nord à Amiens et arrivée à minuit en gare de Doullens}\par\vspace{0.2mm}
\noindent{\fontsize{6.9pt}{8.0pt}\selectfont\textit{Jeudi 5 octobre 1916 (après-midi et soir)} \hfill \textit{Paris (Gare du Nord), Amiens et Doullens (Somme)}}\par\vspace{0.5mm}
\noindent\rule{\textwidth}{0.35pt}\par\vspace{0.8mm}
% --- Corps de la transcription diplomatique et apparat marginal en prolongement ---
{\renewcommand{\arraystretch}{0.87}%
\noindent\hspace*{3.5mm}\begin{tabular}{@{}p{87mm}@{\hspace{3.5mm}}p{39mm}@{}}
      \leavevmode\fontsize{6.9pt}{7.9pt}\selectfont\strut Parti train 1 H 15 s (espéré [groupe abrégé illis.] à 10 H 45 matin)\strut & \smash{\parbox[t]{39mm}{\fontsize{6.2pt}{7.0pt}\selectfont\raggedright\color{black!88}\textbf{1} \textit{lect. incert. in ms.} \textbf{1 s.} \textit{abbr. pro} soir. — \textbf{2 s.} \textit{abbr. pro} soir. — \textbf{3 s.} \textit{abbr. pro} soir.}} \\
      \leavevmode\fontsize{6.9pt}{7.9pt}\selectfont\strut \underline{A Amiens}, à 6 H s —\strut &  \\
      \leavevmode\fontsize{6.9pt}{7.9pt}\selectfont\strut Départ à 6 H 30 s\strut &  \\
      \cline{1-1}
      \leavevmode\fontsize{6.9pt}{7.9pt}\selectfont\strut \underline{A Doullens}, vers minuit\strut &  \\
      \leavevmode\llap{\fontsize{6.4pt}{7.3pt}\selectfont\color{darkgray}5\hspace{1.8mm}}\fontsize{6.9pt}{7.9pt}\selectfont\strut Dû coucher gare —\strut &  \\
      \leavevmode\fontsize{6.9pt}{7.9pt}\selectfont\strut $[\![$Paris$]\!]$ Métro : 15 c, 15 c\strut & \smash{\parbox[t]{39mm}{\fontsize{6.2pt}{7.0pt}\selectfont\raggedright\color{black!88}\textbf{6} $[\![$Paris$]\!]\ $ \textit{del. in ms.} — \textbf{7 c.} \textit{abbr. pro} centimes. \textbf{7 Food} \textit{angl. in ms.}}} \\
      \cline{1-1}
      \leavevmode\fontsize{6.9pt}{7.9pt}\selectfont\strut Food... : 25 c, 25 c, 50 c., Drink : 10 c.\strut &  \\
      \leavevmode\fontsize{6.9pt}{7.9pt}\selectfont\strut Consigne : 10 c\strut &  \\
\end{tabular}}%
\par\vspace*{\fill}
% --- Registres de bas de page (1. Critique interne [si pertinent] ; 2. Loci paralleli en latin ; 3. Commentaire historique en retrait) ---
\noindent\rule{\textwidth}{0.5pt}\par\vspace{0.65mm}
{\fontsize{6.6pt}{7.7pt}\selectfont%
\noindent\normalfont\color{black!92} \textbf{2} (\textit{Amiens}) cf. supra f\textsuperscript{o}~25~v\textsuperscript{o} (p.~\pageref{fol:51}), infra f\textsuperscript{o}~44~v\textsuperscript{o} (p.~\pageref{fol:89}), f\textsuperscript{o}~49~r\textsuperscript{o} (p.~\pageref{fol:98}), \textit{et passim}. — \textbf{4} (\textit{Doullens}) cf. supra f\textsuperscript{o}~25~v\textsuperscript{o} (p.~\pageref{fol:51}), infra f\textsuperscript{o}~33~r\textsuperscript{o} (p.~\pageref{fol:66}), f\textsuperscript{o}~33~v\textsuperscript{o} (p.~\pageref{fol:67}), \textit{et passim}. — \textbf{6} (\textit{Paris}) cf. supra in tegumento (p.~\pageref{fol:1}), f\textsuperscript{o}~32~r\textsuperscript{o} (p.~\pageref{fol:64}), infra f\textsuperscript{o}~49~v\textsuperscript{o} (p.~\pageref{fol:99}), f\textsuperscript{o}~58~r\textsuperscript{o} (p.~\pageref{fol:116}), \textit{et passim}.\par\vspace{0.55mm}
\noindent\rule{38mm}{0.4pt}\par\vspace{0.5mm}
\noindent\hspace*{5mm}\parbox{\dimexpr\textwidth-5mm\relax}{\fontsize{6.6pt}{7.7pt}\selectfont\color{black!84}Suite de la journée du jeudi 5 octobre 1916: parti de Paris-Nord à 13~h~15 pour \textbf{Amiens} (18~h), Hippolyte Combes reprend à 18~h~30 la ligne de Saint-Pol et débarque vers minuit en gare de \textbf{Doullens} (Somme), grand nœud ferroviaire d'arrière-front où il doit passer le reste de la nuit sur un banc de la salle d'attente.}\par}
\newpage
\phantomsection\label{fac:65}
% --- En-tête du fac-similé (Page impaire / droite : Photographie du manuscrit) ---
\noindent\rule{\textwidth}{0.5pt}\par\vspace{0.55mm}
\noindent\textbf{\fontsize{6.8pt}{7.8pt}\selectfont FAC-SIMILÉ EN VIS-À-VIS — FOLIO 32 VERSO} \hfill \textbf{\fontsize{6.8pt}{7.8pt}\selectfont OCTOBRE 1916}\par\vspace{0.25mm}
\noindent{\fontsize{6.8pt}{8.0pt}\selectfont\textit{Livre-Journal d'Hippolyte Combes} — Jeudi 5 octobre 1916 (après-midi et soir) (Paris (Gare du Nord), Amiens et Doullens (Somme))}\par\vspace{0.55mm}
\noindent\rule{\textwidth}{0.35pt}\par\vspace*{\fill}
{\centering
  \setlength{\fboxsep}{1.2mm}%
  \setlength{\fboxrule}{0.4pt}%
  \fcolorbox{gray!55}{white}{%
    \includegraphics[width=0.95\textwidth,height=0.81\textheight,keepaspectratio]{\facsimiledir/20251126_103743.jpg}%
  }\par}
\vspace*{\fill}
\newpage

% =============================================================================
% VIS-A-VIS SPREAD : FOLIO 33 RECTO (20251126_103755.jpg)
% =============================================================================
\cleardoubleevenpage
\renewcommand{\currentmonthyear}{Octobre 1916}
\renewcommand{\currentfoliolabel}{Folio 33 recto}
\phantomsection\label{fol:66}
% --- En-tête du folio (Page paire / gauche : Transcription, apparat & registres de bas de page) ---
\noindent\rule{\textwidth}{0.5pt}\par\vspace{0.55mm}
\noindent\textbf{\footnotesize FOLIO 33 RECTO} \hfill \textbf{\footnotesize OCTOBRE 1916}\par\vspace{0.25mm}
\noindent{\fontsize{6.9pt}{8.0pt}\selectfont Vendredi 6 octobre 1916 — Fin de la nuit à la gare de Doullens, train de 7 h 25 et retour en lorry de la 54e Brigade auprès de la 152nd Coy A.S.C.}\par\vspace{0.2mm}
\noindent{\fontsize{6.9pt}{8.0pt}\selectfont\textit{Vendredi 6 octobre 1916 (matin et midi)} \hfill \textit{Doullens, Canaples et Barlette (Somme)}}\par\vspace{0.5mm}
\noindent\rule{\textwidth}{0.35pt}\par\vspace{0.8mm}
% --- Corps de la transcription diplomatique et apparat marginal en prolongement ---
{\renewcommand{\arraystretch}{0.87}%
\noindent\hspace*{3.5mm}\begin{tabular}{@{}p{87mm}@{\hspace{3.5mm}}p{39mm}@{}}
      \leavevmode\fontsize{6.9pt}{7.9pt}\selectfont\strut \underline{Vendredi 6} \hfill{} \underline{Octb 1916}\strut &  \\
      \leavevmode\fontsize{6.9pt}{7.9pt}\selectfont\strut Achevé nuit salle départ\strut &  \\
      \leavevmode\fontsize{6.9pt}{7.9pt}\selectfont\strut gare de Doullens. Dormi allongé sur 1 banc.\strut &  \\
      \leavevmode\fontsize{6.9pt}{7.9pt}\selectfont\strut Au petit jour, levé —\strut &  \\
      \leavevmode\llap{\fontsize{6.4pt}{7.3pt}\selectfont\color{darkgray}5\hspace{1.8mm}}\fontsize{6.9pt}{7.9pt}\selectfont\strut Venu promener en ville. Rues encore\strut &  \\
      \leavevmode\fontsize{6.9pt}{7.9pt}\selectfont\strut désertes. qq magasins déjà ouverts –\strut & \smash{\parbox[t]{39mm}{\fontsize{6.2pt}{7.0pt}\selectfont\raggedright\color{black!88}\textbf{6 qq} \textit{abbr. pro} quelque(s).}} \\
      \leavevmode\fontsize{6.9pt}{7.9pt}\selectfont\strut Entré église \& y rêvé à mon aise\strut &  \\
      \leavevmode\fontsize{6.9pt}{7.9pt}\selectfont\strut Tjours pluie fine, incessante —\strut &  \\
      \leavevmode\fontsize{6.9pt}{7.9pt}\selectfont\strut Parti de Doullens, train 7 H 25. Appris à\strut &  \\
      \leavevmode\llap{\fontsize{6.4pt}{7.3pt}\selectfont\color{darkgray}10\hspace{1.8mm}}\fontsize{6.9pt}{7.9pt}\selectfont\strut l'office de l'R.T.O. que ma division\strut & \smash{\parbox[t]{39mm}{\fontsize{6.2pt}{7.0pt}\selectfont\raggedright\color{black!88}\textbf{10 R.T.O.} \textit{angl.} Railway Transport Officer. — \textbf{11 shifté} \textit{angl. in ms.}}} \\
      \leavevmode\fontsize{6.9pt}{7.9pt}\selectfont\strut a shifté. Le railhead est Canaples\strut &  \\
      \leavevmode\fontsize{6.9pt}{7.9pt}\selectfont\strut Descendu à cette station, me suis\strut &  \\
      \leavevmode\fontsize{6.9pt}{7.9pt}\selectfont\strut informé – Coïncidence heureuse : les\strut &  \\
      \leavevmode\fontsize{6.9pt}{7.9pt}\selectfont\strut lorries (supply column) de la 54e\strut & \smash{\parbox[t]{39mm}{\fontsize{6.2pt}{7.0pt}\selectfont\raggedright\color{black!88}\textbf{14 lorries} \textit{angl. in ms.}}} \\
      \leavevmode\llap{\fontsize{6.4pt}{7.3pt}\selectfont\color{darkgray}15\hspace{1.8mm}}\fontsize{6.9pt}{7.9pt}\selectfont\strut brigade, encore là. Monté sur 1\strut &  \\
      \leavevmode\fontsize{6.9pt}{7.9pt}\selectfont\strut de ces lorries et ainsi arrivé j'au\strut & \smash{\parbox[t]{39mm}{\fontsize{6.2pt}{7.0pt}\selectfont\raggedright\color{black!88}\textbf{16 lorries} \textit{angl. in ms.}}} \\
      \leavevmode\fontsize{6.9pt}{7.9pt}\selectfont\strut refilling point de ma brigade\strut & \smash{\parbox[t]{39mm}{\fontsize{6.2pt}{7.0pt}\selectfont\raggedright\color{black!88}\textbf{17 refilling point} \textit{angl. in ms.} — \textbf{18 A.S.C.} \textit{angl.} Army Service Corps.}} \\
      \leavevmode\fontsize{6.9pt}{7.9pt}\selectfont\strut Là retrouvé la 152 Coy. A.S.C.\strut &  \\
      \leavevmode\fontsize{6.9pt}{7.9pt}\selectfont\strut A 1H soir –\strut &  \\
\end{tabular}}%
\par\vspace*{\fill}
% --- Registres de bas de page (1. Critique interne [si pertinent] ; 2. Loci paralleli en latin ; 3. Commentaire historique en retrait) ---
\noindent\rule{\textwidth}{0.5pt}\par\vspace{0.65mm}
{\fontsize{6.6pt}{7.7pt}\selectfont%
\noindent\normalfont\color{black!92} \textbf{3} (\textit{Doullens}) cf. supra f\textsuperscript{o}~25~v\textsuperscript{o} (p.~\pageref{fol:51}), f\textsuperscript{o}~32~v\textsuperscript{o} (p.~\pageref{fol:65}), infra f\textsuperscript{o}~33~v\textsuperscript{o} (p.~\pageref{fol:67}), f\textsuperscript{o}~49~r\textsuperscript{o} (p.~\pageref{fol:98}), \textit{et passim}. — \textbf{18} (\textit{152 Coy}) cf. infra f\textsuperscript{o}~45~v\textsuperscript{o} (p.~\pageref{fol:91}), f\textsuperscript{o}~46~r\textsuperscript{o} (p.~\pageref{fol:92}), \textit{et passim}.\par\vspace{0.55mm}
\noindent\rule{38mm}{0.4pt}\par\vspace{0.5mm}
\noindent\hspace*{5mm}\parbox{\dimexpr\textwidth-5mm\relax}{\fontsize{6.6pt}{7.7pt}\selectfont\color{black!84}Première mention explicite dans le carnet du numéro de l'unité britannique d'Hippolyte Combes : la \textbf{152\textsuperscript{e} Compagnie de l'\textit{Army Service Corps}} (\textit{152nd Company A.S.C.}), rattachée au \textit{18th Divisional Train} et chargée du ravitaillement de la \textbf{54\textsuperscript{e} Brigade d'infanterie} de la 18\textsuperscript{e} Division, qui vient de livrer la bataille de la crête de Thiepval (26–30 septembre 1916) et a établi sa tête de ligne ferroviaire (\textit{railhead}) à \textbf{Canaples} (Somme).}\par}
\newpage
\phantomsection\label{fac:66}
% --- En-tête du fac-similé (Page impaire / droite : Photographie du manuscrit) ---
\noindent\rule{\textwidth}{0.5pt}\par\vspace{0.55mm}
\noindent\textbf{\fontsize{6.8pt}{7.8pt}\selectfont FAC-SIMILÉ EN VIS-À-VIS — FOLIO 33 RECTO} \hfill \textbf{\fontsize{6.8pt}{7.8pt}\selectfont OCTOBRE 1916}\par\vspace{0.25mm}
\noindent{\fontsize{6.8pt}{8.0pt}\selectfont\textit{Livre-Journal d'Hippolyte Combes} — Vendredi 6 octobre 1916 (matin et midi) (Doullens, Canaples et Barlette (Somme))}\par\vspace{0.55mm}
\noindent\rule{\textwidth}{0.35pt}\par\vspace*{\fill}
{\centering
  \setlength{\fboxsep}{1.2mm}%
  \setlength{\fboxrule}{0.4pt}%
  \fcolorbox{gray!55}{white}{%
    \includegraphics[width=0.95\textwidth,height=0.81\textheight,keepaspectratio]{\facsimiledir/20251126_103755.jpg}%
  }\par}
\vspace*{\fill}
\newpage

% =============================================================================
% CHAPITRE IV (FOLIO 33 VERSO A FOLIO 52 RECTO)
% =============================================================================
\cleardoubleevenpage
\renewcommand{\currentmonthyear}{Octobre 1916 – Novembre 1916}
\renewcommand{\currentfoliolabel}{Chapitre IV}
\chapter{L’automne dans la Somme : de Ribeaucourt à Albert et Acheux}

\noindent{\fontsize{7.8pt}{9.0pt}\selectfont\textit{Folios 33 verso à 52 recto — Ribeaucourt, Beauval, Warloy-Baillon, Bouzincourt, Albert et Acheux-en-Amiénois (6 octobre – 30 novembre 1916)}}\par\vspace{1.5mm}
\noindent\rule{\textwidth}{0.5pt}\par\vspace{0.9mm}
\noindent\textbf{\scriptsize\textsc{Résumé analytique du Chapitre IV (Folio 33 verso à Folio 52 recto)}}\par\vspace{0.5mm}
\noindent{\fontsize{6.9pt}{8.1pt}\selectfont Le quatrième chapitre (f\textsuperscript{o}~33~v\textsuperscript{o} à f\textsuperscript{o}~52~r\textsuperscript{o}) couvre les deux mois d'octobre et de novembre 1916 sur le front de la Somme. Du 6 au 15 octobre, la 152\textsuperscript{e} Compagnie A.S.C. jouit d'une semaine de repos dans les vergers du Ponthieu, à Ribeaucourt et au hameau de Barlette, où Hippolyte Combes visite le château de Ribeaucourt et le village voisin de Domqueur. Le 15 octobre, la colonne reprend la route du front par Beauval et Warloy-Baillon pour bivouaquer sous la tente dans la boue glaciale de Bouzincourt (17–19 octobre) puis du camp des briqueteries d'Albert (\textit{Brickfields Area}, 19–31 octobre), où Hippolyte Combes trouve une chambre chez la famille Duchâteau, décrit avec émotion les ruines de la basilique Notre-Dame de Brebières et sa Vierge penchée, et rassure son statut d'interprète auprès du lieutenant Pizzi. Au début de novembre s'opère un tournant majeur dans son service : détaché le 2 novembre par le lieutenant-colonel Grose auprès du service des fourrages et achats d'armée (dirigé par le major Levy), il s'installe à Acheux-en-Amiénois chez le tailleur M.~Carluvant (n\textsuperscript{o}~1, route de Léalvillers) et sillonne désormais en automobile ou à cheval les communes de l'arrière-front picard (Varennes, Léalvillers, Arquèves et la ferme de Belle-Église, Louvencourt, Bertrancourt, Bus-lès-Artois, Raincheval, Rubempré, Hérissart et le château de Martinsart).}\par\vspace{1.0mm}
\noindent\rule{\textwidth}{0.35pt}\par\vspace{1.4mm}

\noindent\textbf{\scriptsize\textsc{Introduction historique et philologique}}\par\vspace{0.6mm}
{\fontsize{7.1pt}{8.3pt}\selectfont Durant les mois d'octobre et de novembre 1916, tandis que les brigades de la \textbf{18\textsuperscript{e} Division britannique} livrent de rudes combats d'usure sous la pluie devant la redoute Schwaben, la tranchée Regina (\textit{Regina Trench}) et les hauteurs de l'Ancre (Nichols, \textit{op.~cit.}, chap.~VIII–IX), le \textit{18th Divisional Train}, commandé par le \textbf{lieutenant-colonel Daniel Charles Evans Grose, A.S.C.} (\texttt{Col.~Grose}, f\textsuperscript{o}~45~r\textsuperscript{o}), assure le ravitaillement quotidien des unités depuis ses bases de Warloy-Baillon et d'Albert. Le détachement d'Hippolyte Combes à \textbf{Acheux-en-Amiénois} le 2 novembre 1916 répond aux besoins considérables en paille et en foin de la cavalerie et des trains hippomobiles britanniques : rattaché aux services d'achats et de réquisitions (\textit{Foraging} et \textit{Army Purchase Board}, sous l'autorité du \textbf{major Levy}), il contrôle les chantiers de battage mécanique (notamment à la commanderie de Belle-Église chez la famille belge Houthouan) et négocie directement avec les maires et les cultivateurs de la Somme.\par}
\vfill
\newpage
% --- Page de droite d'ouverture du chapitre : Carte vectorielle TikZ & Photographies d'archives ---
\noindent\rule{\textwidth}{0.5pt}\par\vspace{0.6mm}
\noindent\textbf{\fontsize{6.8pt}{7.8pt}\selectfont CARTE D'ITINÉRAIRE DU CHAPITRE IV} \hfill \textbf{\fontsize{6.8pt}{7.8pt}\selectfont OCTOBRE 1916 – NOVEMBRE 1916}\par\vspace{0.25mm}
\noindent{\fontsize{6.7pt}{7.8pt}\selectfont\textit{Déplacements, cantonnements et tournées d'Hippolyte Combes (Folio 33 verso à Folio 52 recto)}}\par\vspace{0.55mm}
\noindent\rule{\textwidth}{0.35pt}\par\vspace{1.2mm}
\begin{center}

\begin{tikzpicture}[
  scale=0.96,
  every node/.style={font=\fontsize{6.0pt}{7.0pt}\selectfont},
  loc/.style={circle, fill=black, inner sep=1.1pt},
  step/.style={circle, draw=bordeaux, fill=white, line width=0.75pt, inner sep=1.5pt},
  move/.style={draw=bordeaux, line width=0.75pt, ->, >=stealth},
  excur/.style={draw=black!65, line width=0.5pt, dashed, ->, >=stealth},
  river/.style={draw=black!35, line width=1.1pt}
]
\draw[black!50, line width=0.4pt] (-60mm,-42mm) rectangle (60mm,42mm);
\fill[black!3] (-60mm,-42mm) rectangle (60mm,42mm);

% Vallée de l'Ancre
\draw[river] (36mm,26mm) .. controls (30mm,6mm) and (26mm,-12mm) .. (20mm,-36mm)
  node[pos=0.75, right, font=\fontsize{5.6pt}{6.5pt}\selectfont\itshape]{L'Ancre};

% Ligne de front de l'Ancre (Thiepval - Beaumont-Hamel)
\draw[bordeaux!80, line width=0.9pt, decorate, decoration={zigzag, segment length=4pt, amplitude=1pt}]
  (38mm,36mm) -- (42mm,12mm) -- (45mm,-14mm);
\node[rotate=-84, font=\fontsize{5.5pt}{6.5pt}\selectfont, bordeaux] at (47mm,10mm) {\textbf{Front de l'Ancre (Thiepval – Redoute Schwaben)}};

% Cantonnement de repos dans le Ponthieu (6-15 oct.)
\node[step] (ribeaucourt) at (-46mm, 24mm) {};
\node[above right] at (ribeaucourt) {\textbf{1. Ribeaucourt \& Barlette} (6–15 oct.)};
\node[loc] (domqueur) at (-48mm, 14mm) {};
\node[below right] at (domqueur) {Domqueur (9 oct.)};
\node[loc] (canaples) at (-32mm, 2mm) {};
\node[left] at (canaples) {Canaples (\textit{railhead}, 6 oct.)};

% Marche vers Albert (15-17 oct.)
\node[step] (beauval) at (-14mm, 20mm) {};
\node[above] at (beauval) {\textbf{2. Beauval} (15–16 oct.)};
\node[loc] (rubempre) at (-8mm, -8mm) {};
\node[left] at (rubempre) {Rubempré \& Septenville};
\node[step] (warloy) at (6mm, -10mm) {};
\node[below left] at (warloy) {\textbf{3. Warloy-Baillon} (16 oct.)};
\node[loc] (senlis) at (16mm, -6mm) {};
\node[left] at (senlis) {Senlis-le-Sec};

% Secteur d'Albert (17-31 oct.)
\node[step] (bouzincourt) at (22mm, -2mm) {};
\node[above left] at (bouzincourt) {\textbf{4. Bouzincourt} (17–19 oct.)};
\node[step] (albert) at (28mm, -12mm) {};
\node[below right] at (albert) {\textbf{5. Albert (\textit{Brickfields})} (19–31 oct.)};
\node[loc] (martinsart) at (30mm, 4mm) {};
\node[right] at (martinsart) {Martinsart \& Aveluy};

% Secteur d'Acheux (novembre 1916)
\node[step] (acheux) at (12mm, 18mm) {};
\node[above left] at (acheux) {\textbf{6. Acheux-en-Amiénois} (2–30 nov.)};
\node[loc] (varennes) at (16mm, 12mm) {};
\node[right] at (varennes) {Varennes \& Léalvillers};
\node[loc] (raincheval) at (-2mm, 10mm) {};
\node[left] at (raincheval) {Raincheval \& Arquèves};
\node[loc] (louvencourt) at (10mm, 28mm) {};
\node[above] at (louvencourt) {Louvencourt, Bertrancourt, Bus-lès-Artois};
\node[loc] (herissart) at (-6mm, -1mm) {};
\node[left] at (herissart) {Hérissart \& Puchevillers};

% Flèches d'itinéraires
\draw[move] (canaples) -- (ribeaucourt);
\draw[excur] (ribeaucourt) -- (domqueur);
\draw[move] (ribeaucourt) -- (beauval) -- (warloy) -- (bouzincourt) -- (albert);
\draw[move] (albert) to[bend left=18] (acheux);
\draw[excur] (acheux) -- (louvencourt);
\draw[excur] (acheux) -- (raincheval);
\draw[excur] (acheux) -- (varennes);
\draw[excur] (acheux) -- (herissart);

\draw[line width=0.6pt] (-54mm,-36mm) -- (-34mm,-36mm);
\node[above] at (-44mm,-36mm) {0 \hfill 5 \hfill 10 km};
\draw[->, line width=0.6pt] (52mm, 28mm) -- (52mm, 37mm) node[above]{\textbf{N}};
\end{tikzpicture}

\end{center}
\vfill
\noindent\rule{\textwidth}{0.3pt}\par\vspace{1.2mm}
\noindent\begin{center}
  \begin{minipage}[b]{0.76\textwidth}
    \centering
    \setlength{\fboxsep}{0.7mm}\setlength{\fboxrule}{0.4pt}%
    \fcolorbox{gray!55}{white}{\includegraphics[width=0.94\linewidth,height=28mm,keepaspectratio]{illustrations/fiche_matricule.jpg}}
    \par\vspace{0.6mm}{\fontsize{5.8pt}{6.7pt}\selectfont Extrait du registre matricule d'Hippolyte Combes (classe 1896, matricule 911, bureau de recrutement de Toulouse) détaillant ses services auprès de la Mission militaire française attachée à l'Armée britannique (4 mai 1915 – 11 novembre 1918).}
  \end{minipage}
\end{center}
\vspace*{1mm}
\newpage

% =============================================================================
% VIS-A-VIS SPREAD : FOLIO 33 VERSO (20251126_103816.jpg)
% =============================================================================
\cleardoubleevenpage
\renewcommand{\currentmonthyear}{Octobre 1916}
\renewcommand{\currentfoliolabel}{Folio 33 verso}
\phantomsection\label{fol:67}
% --- En-tête du folio (Page paire / gauche : Transcription, apparat & registres de bas de page) ---
\noindent\rule{\textwidth}{0.5pt}\par\vspace{0.55mm}
\noindent\textbf{\footnotesize FOLIO 33 VERSO} \hfill \textbf{\footnotesize OCTOBRE 1916}\par\vspace{0.25mm}
\noindent{\fontsize{6.9pt}{8.0pt}\selectfont Installation au hameau de Barlette (Ribeaucourt) et visa de permission à Bernaville}\par\vspace{0.2mm}
\noindent{\fontsize{6.9pt}{8.0pt}\selectfont\textit{Vendredi 6 octobre 1916 (après-midi et soir)} \hfill \textit{Hameau de Barlette (Ribeaucourt) et Bernaville (Somme)}}\par\vspace{0.5mm}
\noindent\rule{\textwidth}{0.35pt}\par\vspace{0.8mm}
% --- Corps de la transcription diplomatique et apparat marginal en prolongement ---
{\renewcommand{\arraystretch}{0.87}%
\noindent\hspace*{3.5mm}\begin{tabular}{@{}p{87mm}@{\hspace{3.5mm}}p{39mm}@{}}
      \leavevmode\fontsize{6.9pt}{7.9pt}\selectfont\strut \underline{Ven. 6} \hfill{} \underline{Oct / 16}\strut & \smash{\parbox[t]{39mm}{\fontsize{6.2pt}{7.0pt}\selectfont\raggedright\color{black!88}\textbf{1 Ven.} \textit{scil.} Vendredi.}} \\
      \leavevmode\fontsize{6.9pt}{7.9pt}\selectfont\strut Ma Coy, cantonnée au hameau\strut & \smash{\parbox[t]{39mm}{\fontsize{6.2pt}{7.0pt}\selectfont\raggedright\color{black!88}\textbf{2 Coy} \textit{angl.} Company.}} \\
      \leavevmode\fontsize{6.9pt}{7.9pt}\selectfont\strut de Barlette, près de Ribeaucourt\strut &  \\
      \leavevmode\fontsize{6.9pt}{7.9pt}\selectfont\strut qq maisons. Population très symphe\strut & \smash{\parbox[t]{39mm}{\fontsize{6.2pt}{7.0pt}\selectfont\raggedright\color{black!88}\textbf{4 qq} \textit{abbr. pro} quelque(s). \textbf{4 symphe} \textit{abbr. pro} sympathique. — \textbf{5 extt} \textit{abbr. pro} extrêmement. — \textbf{6 lunch} \textit{angl. in ms.} — \textbf{8 Fse} \textit{abbr. pro} Française. — \textbf{9 gde} \textit{abbr. pro} grande. — \textbf{10 Mme} \textit{abbr. pro} Madame.}} \\
      \leavevmode\llap{\fontsize{6.4pt}{7.3pt}\selectfont\color{darkgray}5\hspace{1.8mm}}\fontsize{6.9pt}{7.9pt}\selectfont\strut Région très boisée, extt pittoresque\strut &  \\
      \leavevmode\fontsize{6.9pt}{7.9pt}\selectfont\strut Après lunch, venu sur 1 lorry\strut &  \\
      \leavevmode\fontsize{6.9pt}{7.9pt}\selectfont\strut à Bernaville et visé mon titre de\strut &  \\
      \leavevmode\fontsize{6.9pt}{7.9pt}\selectfont\strut permission au bureau de la Mission Fse\strut &  \\
      \leavevmode\fontsize{6.9pt}{7.9pt}\selectfont\strut Bientôt trouvé 1 billet = gde\strut &  \\
      \leavevmode\llap{\fontsize{6.4pt}{7.3pt}\selectfont\color{darkgray}10\hspace{1.8mm}}\fontsize{6.9pt}{7.9pt}\selectfont\strut chambre, chz Mme Ve —\strut &  \\
      \leavevmode\fontsize{6.9pt}{7.9pt}\selectfont\strut à 900 mètres de Barlette\strut &  \\
      \leavevmode\fontsize{6.9pt}{7.9pt}\selectfont\strut Ameublement très sommaire, mais\strut &  \\
      \leavevmode\fontsize{6.9pt}{7.9pt}\selectfont\strut chambre très belle \& spacieuse\strut &  \\
      \cline{1-1}
      \leavevmode\fontsize{6.9pt}{7.9pt}\selectfont\strut \underline{A Doullens}, Cartes ples : 50 c\strut & \smash{\parbox[t]{39mm}{\fontsize{6.2pt}{7.0pt}\selectfont\raggedright\color{black!88}\textbf{14 Cartes ples} \textit{abbr. pro} Cartes postales. — \textbf{15 canteen} \textit{angl. in ms.} — \textbf{16 cartes ples} \textit{abbr. pro} cartes postales. — \textbf{17 lorry} \textit{angl. in ms.} — \textbf{18 Q.M.S.} \textit{angl.} Quartermaster Sergeant. \textbf{18 fr.} \textit{abbr. pro} franc(s).}} \\
      \leavevmode\llap{\fontsize{6.4pt}{7.3pt}\selectfont\color{darkgray}15\hspace{1.8mm}}\fontsize{6.9pt}{7.9pt}\selectfont\strut A ", canteen : 20, 20, 20, 20 c\strut &  \\
      \leavevmode\fontsize{6.9pt}{7.9pt}\selectfont\strut A Bernaville, cartes ples : 10 c\strut &  \\
      \leavevmode\fontsize{6.9pt}{7.9pt}\selectfont\strut A Barlette : Tip au driver du lorry : 1f\strut &  \\
      \leavevmode\fontsize{6.9pt}{7.9pt}\selectfont\strut Gift au Q.M.S. Tait : 5 fr. –\strut &  \\
\end{tabular}}%
\par\vspace*{\fill}
% --- Registres de bas de page (1. Critique interne [si pertinent] ; 2. Loci paralleli en latin ; 3. Commentaire historique en retrait) ---
\noindent\rule{\textwidth}{0.5pt}\par\vspace{0.65mm}
{\fontsize{6.6pt}{7.7pt}\selectfont%
\noindent\normalfont\color{black!92} \textbf{3} (\textit{Ribeaucourt}) cf. infra f\textsuperscript{o}~34~r\textsuperscript{o} (p.~\pageref{fol:68}), f\textsuperscript{o}~34~v\textsuperscript{o} (p.~\pageref{fol:69}), \textit{et passim}. — \textbf{3} (\textit{Barlette}) cf. supra f\textsuperscript{o}~33~r\textsuperscript{o} (p.~\pageref{fol:66}), infra f\textsuperscript{o}~34~r\textsuperscript{o} (p.~\pageref{fol:68}), f\textsuperscript{o}~34~v\textsuperscript{o} (p.~\pageref{fol:69}), \textit{et passim}. — \textbf{14} (\textit{Doullens}) cf. supra f\textsuperscript{o}~25~v\textsuperscript{o} (p.~\pageref{fol:51}), f\textsuperscript{o}~33~r\textsuperscript{o} (p.~\pageref{fol:66}), infra f\textsuperscript{o}~49~r\textsuperscript{o} (p.~\pageref{fol:98}), f\textsuperscript{o}~57~v\textsuperscript{o} (p.~\pageref{fol:115}), \textit{et passim}.\par\vspace{0.55mm}
\noindent\rule{38mm}{0.4pt}\par\vspace{0.5mm}
\noindent\hspace*{5mm}\parbox{\dimexpr\textwidth-5mm\relax}{\fontsize{6.6pt}{7.7pt}\selectfont\color{black!84}Tandis que les brigades de la 18\textsuperscript{e} Division sont relevées après la prise de Thiepval, la \textbf{152\textsuperscript{e} Compagnie A.S.C.} cantonne au repos du 6 au 15 octobre 1916 au hameau boisé de \textbf{Barlette} (commune de \textbf{Ribeaucourt}, Somme, canton de Domart-en-Ponthieu).}\par}
\newpage
\phantomsection\label{fac:67}
% --- En-tête du fac-similé (Page impaire / droite : Photographie du manuscrit) ---
\noindent\rule{\textwidth}{0.5pt}\par\vspace{0.55mm}
\noindent\textbf{\fontsize{6.8pt}{7.8pt}\selectfont FAC-SIMILÉ EN VIS-À-VIS — FOLIO 33 VERSO} \hfill \textbf{\fontsize{6.8pt}{7.8pt}\selectfont OCTOBRE 1916}\par\vspace{0.25mm}
\noindent{\fontsize{6.8pt}{8.0pt}\selectfont\textit{Livre-Journal d'Hippolyte Combes} — Vendredi 6 octobre 1916 (après-midi et soir) (Hameau de Barlette (Ribeaucourt) et Bernaville (Somme))}\par\vspace{0.55mm}
\noindent\rule{\textwidth}{0.35pt}\par\vspace*{\fill}
{\centering
  \setlength{\fboxsep}{1.2mm}%
  \setlength{\fboxrule}{0.4pt}%
  \fcolorbox{gray!55}{white}{%
    \includegraphics[width=0.95\textwidth,height=0.81\textheight,keepaspectratio]{\facsimiledir/20251126_103816.jpg}%
  }\par}
\vspace*{\fill}
\newpage

% =============================================================================
% VIS-A-VIS SPREAD : FOLIO 34 RECTO (20251126_103828.jpg)
% =============================================================================
\cleardoubleevenpage
\renewcommand{\currentmonthyear}{Octobre 1916}
\renewcommand{\currentfoliolabel}{Folio 34 recto}
\phantomsection\label{fol:68}
% --- En-tête du folio (Page paire / gauche : Transcription, apparat & registres de bas de page) ---
\noindent\rule{\textwidth}{0.5pt}\par\vspace{0.55mm}
\noindent\textbf{\footnotesize FOLIO 34 RECTO} \hfill \textbf{\footnotesize OCTOBRE 1916}\par\vspace{0.25mm}
\noindent{\fontsize{6.9pt}{8.0pt}\selectfont Tempête nocturne sur les vergers de Barlette et règlement des frais de mess au capitaine Panton}\par\vspace{0.2mm}
\noindent{\fontsize{6.9pt}{8.0pt}\selectfont\textit{Samedi 7 octobre 1916} \hfill \textit{Hameau de Barlette, Ribeaucourt (Somme)}}\par\vspace{0.5mm}
\noindent\rule{\textwidth}{0.35pt}\par\vspace{0.8mm}
% --- Corps de la transcription diplomatique et apparat marginal en prolongement ---
{\renewcommand{\arraystretch}{1.00}%
\noindent\hspace*{3.5mm}\begin{tabular}{@{}p{87mm}@{\hspace{3.5mm}}p{39mm}@{}}
      \leavevmode\fontsize{7.7pt}{9.0pt}\selectfont\strut \underline{Sam. 7.} \hfill{} \underline{Octbre 1916}\strut & \smash{\parbox[t]{39mm}{\fontsize{6.2pt}{7.0pt}\selectfont\raggedright\color{black!88}\textbf{1 Sam.} \textit{scil.} Samedi.}} \\
      \leavevmode\fontsize{7.7pt}{9.0pt}\selectfont\strut Bonne nuit de repos s/ lit improvisé.\strut & \smash{\parbox[t]{39mm}{\fontsize{6.2pt}{7.0pt}\selectfont\raggedright\color{black!88}\textbf{2 s/} \textit{abbr. pro} sur.}} \\
      \leavevmode\fontsize{7.7pt}{9.0pt}\selectfont\strut Levé vers 6H 1/2 et commencé à m'installer. Ouvert sacs\strut & \smash{\parbox[t]{39mm}{\fontsize{6.2pt}{7.0pt}\selectfont\raggedright\color{black!88}\textbf{3 1/2} \textit{scil.} demi / moitié.}} \\
      \leavevmode\fontsize{7.7pt}{9.0pt}\selectfont\strut Durant nuit, averses torrentielles et rafales.\strut &  \\
      \leavevmode\llap{\fontsize{6.4pt}{7.3pt}\selectfont\color{darkgray}5\hspace{1.8mm}}\fontsize{7.7pt}{9.0pt}\selectfont\strut Les arbres et les arbustes, en face maison,\strut &  \\
      \leavevmode\fontsize{7.7pt}{9.0pt}\selectfont\strut tourmentés par le vent, bruissaient plaintivt\strut & \smash{\parbox[t]{39mm}{\fontsize{6.2pt}{7.0pt}\selectfont\raggedright\color{black!88}\textbf{6 plaintivt} \textit{abbr. pro} plaintivement.}} \\
      \leavevmode\fontsize{7.7pt}{9.0pt}\selectfont\strut – On se serait cru près de la mer.\strut &  \\
      \leavevmode\fontsize{7.7pt}{9.0pt}\selectfont\strut Tout le jour, pluie et bourrasques.\strut &  \\
      \leavevmode\fontsize{7.7pt}{9.0pt}\selectfont\strut ciel couvert de gros nuages – –\strut &  \\
      \leavevmode\llap{\fontsize{6.4pt}{7.3pt}\selectfont\color{darkgray}10\hspace{1.8mm}}\fontsize{7.7pt}{9.0pt}\selectfont\strut Me plais bp ds mon billet. Ai très\strut & \smash{\parbox[t]{39mm}{\fontsize{6.2pt}{7.0pt}\selectfont\raggedright\color{black!88}\textbf{10 ds} \textit{abbr. pro} dans. \textbf{10 bp} \textit{abbr. pro} beaucoup. — \textbf{11 gde} \textit{abbr. pro} grande. — \textbf{12 gd} \textit{abbr. pro} grand. \textbf{12 extt} \textit{abbr. pro} extrêmement. — \textbf{14 gds} \textit{abbr. pro} grands. — \textbf{15 Cap.} \textit{abbr. pro} Capitaine. — \textbf{16 mess} \textit{angl. in ms.}}} \\
      \leavevmode\fontsize{7.7pt}{9.0pt}\selectfont\strut gde chambre, haute de plafond. C'était le\strut &  \\
      \leavevmode\fontsize{7.7pt}{9.0pt}\selectfont\strut salon — vieille gd-mère, extt sympathique\strut &  \\
      \leavevmode\fontsize{7.7pt}{9.0pt}\selectfont\strut Devant maison, pâture ac pommiers\strut &  \\
      \leavevmode\fontsize{7.7pt}{9.0pt}\selectfont\strut chargés de fruit Au delà, gds bois. –\strut &  \\
      \leavevmode\llap{\fontsize{6.4pt}{7.3pt}\selectfont\color{darkgray}15\hspace{1.8mm}}\fontsize{7.7pt}{9.0pt}\selectfont\strut Au Cap. Panton payé ma part dépenses\strut &  \\
      \leavevmode\fontsize{7.7pt}{9.0pt}\selectfont\strut mess en septbre, pour — — jours : 29 f —\strut &  \\
\end{tabular}}%
\par\vspace*{\fill}
% --- Registres de bas de page (1. Critique interne [si pertinent] ; 2. Loci paralleli en latin ; 3. Commentaire historique en retrait) ---
\noindent\rule{\textwidth}{0.5pt}\par\vspace{0.65mm}
{\fontsize{6.6pt}{7.7pt}\selectfont%
\noindent\normalfont\color{black!92} \textbf{15} (\textit{Cap. Panton}) cf. supra f\textsuperscript{o}~23~v\textsuperscript{o} (p.~\pageref{fol:47}), infra f\textsuperscript{o}~36~v\textsuperscript{o} (p.~\pageref{fol:73}).\par\vspace{0.55mm}
\noindent\rule{38mm}{0.4pt}\par\vspace{0.5mm}
\noindent\hspace*{5mm}\parbox{\dimexpr\textwidth-5mm\relax}{\fontsize{6.6pt}{7.7pt}\selectfont\color{black!84}Installé dans l'ancien salon d'une maison entourée de vergers de pommiers et de grands bois à \textbf{Barlette} (Ribeaucourt), Hippolyte Combes règle au \textbf{capitaine Panton} (officier de la \textit{152nd Company A.S.C.}, gérant de la popote des officiers) sa quote-part des frais de \textit{mess} du mois de septembre (29~fr., soit environ 2~fr.~20 par jour depuis son admission permanente au mess le 12 septembre, f\textsuperscript{o}~21~r\textsuperscript{o}).}\par}
\newpage
\phantomsection\label{fac:68}
% --- En-tête du fac-similé (Page impaire / droite : Photographie du manuscrit) ---
\noindent\rule{\textwidth}{0.5pt}\par\vspace{0.55mm}
\noindent\textbf{\fontsize{6.8pt}{7.8pt}\selectfont FAC-SIMILÉ EN VIS-À-VIS — FOLIO 34 RECTO} \hfill \textbf{\fontsize{6.8pt}{7.8pt}\selectfont OCTOBRE 1916}\par\vspace{0.25mm}
\noindent{\fontsize{6.8pt}{8.0pt}\selectfont\textit{Livre-Journal d'Hippolyte Combes} — Samedi 7 octobre 1916 (Hameau de Barlette, Ribeaucourt (Somme))}\par\vspace{0.55mm}
\noindent\rule{\textwidth}{0.35pt}\par\vspace*{\fill}
{\centering
  \setlength{\fboxsep}{1.2mm}%
  \setlength{\fboxrule}{0.4pt}%
  \fcolorbox{gray!55}{white}{%
    \includegraphics[width=0.95\textwidth,height=0.81\textheight,keepaspectratio]{\facsimiledir/20251126_103828.jpg}%
  }\par}
\vspace*{\fill}
\newpage

% =============================================================================
% VIS-A-VIS SPREAD : FOLIO 34 VERSO (20251126_103845.jpg)
% =============================================================================
\cleardoubleevenpage
\renewcommand{\currentmonthyear}{Octobre 1916}
\renewcommand{\currentfoliolabel}{Folio 34 verso}
\phantomsection\label{fol:69}
% --- En-tête du folio (Page paire / gauche : Transcription, apparat & registres de bas de page) ---
\noindent\rule{\textwidth}{0.5pt}\par\vspace{0.55mm}
\noindent\textbf{\footnotesize FOLIO 34 VERSO} \hfill \textbf{\footnotesize OCTOBRE 1916}\par\vspace{0.25mm}
\noindent{\fontsize{6.9pt}{8.0pt}\selectfont Dimanche de repos à Barlette et perception de la solde à pied à Bernaville sous la rafale}\par\vspace{0.2mm}
\noindent{\fontsize{6.9pt}{8.0pt}\selectfont\textit{Dimanche 8 octobre 1916} \hfill \textit{Barlette (Ribeaucourt) et Bernaville (Somme)}}\par\vspace{0.5mm}
\noindent\rule{\textwidth}{0.35pt}\par\vspace{0.8mm}
% --- Corps de la transcription diplomatique et apparat marginal en prolongement ---
{\renewcommand{\arraystretch}{1.00}%
\noindent\hspace*{3.5mm}\begin{tabular}{@{}p{87mm}@{\hspace{3.5mm}}p{39mm}@{}}
      \leavevmode\fontsize{7.7pt}{9.0pt}\selectfont\strut \underline{Dim. 8.} \hfill{} \underline{Octb 1916}\strut & \smash{\parbox[t]{39mm}{\fontsize{6.2pt}{7.0pt}\selectfont\raggedright\color{black!88}\textbf{1 Dim.} \textit{scil.} Dimanche.}} \\
      \leavevmode\fontsize{7.7pt}{9.0pt}\selectfont\strut Installation terminée...\strut &  \\
      \leavevmode\fontsize{7.7pt}{9.0pt}\selectfont\strut Jouis repos complet village paisible\strut &  \\
      \leavevmode\fontsize{7.7pt}{9.0pt}\selectfont\strut Tout le jour, ciel couvert de nuages\strut &  \\
      \leavevmode\llap{\fontsize{6.4pt}{7.3pt}\selectfont\color{darkgray}5\hspace{1.8mm}}\fontsize{7.7pt}{9.0pt}\selectfont\strut que le vent chassait avec violence\strut &  \\
      \leavevmode\fontsize{7.7pt}{9.0pt}\selectfont\strut Après-midi, venu à Bernaville,\strut &  \\
      \leavevmode\fontsize{7.7pt}{9.0pt}\selectfont\strut au bureau de la mission fse –\strut & \smash{\parbox[t]{39mm}{\fontsize{6.2pt}{7.0pt}\selectfont\raggedright\color{black!88}\textbf{7 fse} \textit{abbr. pro} française.}} \\
      \leavevmode\fontsize{7.7pt}{9.0pt}\selectfont\strut Aller \& retour à pied. Violentes rafales\strut &  \\
      \leavevmode\fontsize{7.7pt}{9.0pt}\selectfont\strut Pays très verdoyant \& boisé.\strut &  \\
      \leavevmode\llap{\fontsize{6.4pt}{7.3pt}\selectfont\color{darkgray}10\hspace{1.8mm}}\fontsize{7.7pt}{9.0pt}\selectfont\strut A Bernaville, solde \& indemnité : 78 f 75\strut &  \\
\end{tabular}}%
\par\vspace*{\fill}
% --- Registres de bas de page (1. Critique interne [si pertinent] ; 2. Loci paralleli en latin ; 3. Commentaire historique en retrait) ---
\noindent\rule{\textwidth}{0.5pt}\par\vspace{0.65mm}
{\fontsize{6.6pt}{7.7pt}\selectfont%
\noindent\normalfont\color{black!92} \textbf{1} (\textit{de eodem commorationis loco}) cf. supra f\textsuperscript{o}~34~r\textsuperscript{o} (p.~\pageref{fol:68}).\par\vspace{0.55mm}
\noindent\rule{38mm}{0.4pt}\par\vspace{0.5mm}
\noindent\hspace*{5mm}\parbox{\dimexpr\textwidth-5mm\relax}{\fontsize{6.6pt}{7.7pt}\selectfont\color{black!84}Marche dominicale aller et retour à travers le bocage du Ponthieu entre Barlette (Ribeaucourt) et \textbf{Bernaville} (Somme, à 5~km à l'est), où siège pendant la période de repos le bureau de la Mission française de la 18\textsuperscript{e} Division, pour percevoir sa solde et son indemnité de quinzaine (78~fr.~75).}\par}
\newpage
\phantomsection\label{fac:69}
% --- En-tête du fac-similé (Page impaire / droite : Photographie du manuscrit) ---
\noindent\rule{\textwidth}{0.5pt}\par\vspace{0.55mm}
\noindent\textbf{\fontsize{6.8pt}{7.8pt}\selectfont FAC-SIMILÉ EN VIS-À-VIS — FOLIO 34 VERSO} \hfill \textbf{\fontsize{6.8pt}{7.8pt}\selectfont OCTOBRE 1916}\par\vspace{0.25mm}
\noindent{\fontsize{6.8pt}{8.0pt}\selectfont\textit{Livre-Journal d'Hippolyte Combes} — Dimanche 8 octobre 1916 (Barlette (Ribeaucourt) et Bernaville (Somme))}\par\vspace{0.55mm}
\noindent\rule{\textwidth}{0.35pt}\par\vspace*{\fill}
{\centering
  \setlength{\fboxsep}{1.2mm}%
  \setlength{\fboxrule}{0.4pt}%
  \fcolorbox{gray!55}{white}{%
    \includegraphics[width=0.95\textwidth,height=0.81\textheight,keepaspectratio]{\facsimiledir/20251126_103845.jpg}%
  }\par}
\vspace*{\fill}
\newpage

% =============================================================================
% VIS-A-VIS SPREAD : FOLIO 35 RECTO (20251126_103857.jpg)
% =============================================================================
\cleardoubleevenpage
\renewcommand{\currentmonthyear}{Octobre 1916}
\renewcommand{\currentfoliolabel}{Folio 35 recto}
\phantomsection\label{fol:70}
% --- En-tête du folio (Page paire / gauche : Transcription, apparat & registres de bas de page) ---
\noindent\rule{\textwidth}{0.5pt}\par\vspace{0.55mm}
\noindent\textbf{\footnotesize FOLIO 35 RECTO} \hfill \textbf{\footnotesize OCTOBRE 1916}\par\vspace{0.25mm}
\noindent{\fontsize{6.9pt}{8.0pt}\selectfont Promenade sous les futaies centenaires de Ribeaucourt et clair de lune d'octobre}\par\vspace{0.2mm}
\noindent{\fontsize{6.9pt}{8.0pt}\selectfont\textit{Lundi 9 octobre 1916} \hfill \textit{Barlette et forêt de Ribeaucourt (Somme)}}\par\vspace{0.5mm}
\noindent\rule{\textwidth}{0.35pt}\par\vspace{0.8mm}
% --- Corps de la transcription diplomatique et apparat marginal en prolongement ---
{\renewcommand{\arraystretch}{0.93}%
\noindent\hspace*{3.5mm}\begin{tabular}{@{}p{87mm}@{\hspace{3.5mm}}p{39mm}@{}}
      \leavevmode\fontsize{7.3pt}{8.4pt}\selectfont\strut \underline{Lundi 9} \hfill{} \underline{Oct – / 16}\strut &  \\
      \leavevmode\fontsize{7.3pt}{8.4pt}\selectfont\strut Tout le jour, ciel assombri gros nuages\strut &  \\
      \leavevmode\fontsize{7.3pt}{8.4pt}\selectfont\strut couleur de plomb\strut &  \\
      \leavevmode\fontsize{7.3pt}{8.4pt}\selectfont\strut Après-midi, pénétré forêt qui s'étend à 1\strut &  \\
      \leavevmode\llap{\fontsize{6.4pt}{7.3pt}\selectfont\color{darkgray}5\hspace{1.8mm}}\fontsize{7.3pt}{8.4pt}\selectfont\strut faible distance de la route et parallelt –\strut & \smash{\parbox[t]{39mm}{\fontsize{6.2pt}{7.0pt}\selectfont\raggedright\color{black!88}\textbf{5 parallelt} \textit{abbr. pro} parallèlement.}} \\
      \leavevmode\fontsize{7.3pt}{8.4pt}\selectfont\strut Sentiers encore boueux et jonchés de feuilles rouillées.\strut &  \\
      \leavevmode\fontsize{7.3pt}{8.4pt}\selectfont\strut Feuillages épais \& verts. Ds les taillis, enchevêtrt\strut & \smash{\parbox[t]{39mm}{\fontsize{6.2pt}{7.0pt}\selectfont\raggedright\color{black!88}\textbf{7 Ds} \textit{abbr. pro} Dans. \textbf{7 enchevêtrt} \textit{abbr. pro} enchevêtrement. — \textbf{9 plusis} \textit{abbr. pro} plusieurs. — \textbf{10 ds} \textit{abbr. pro} dans. \textbf{10 lgt} \textit{abbr. pro} longtemps.}} \\
      \leavevmode\fontsize{7.3pt}{8.4pt}\selectfont\strut d'arbustes \& de lianes. Des troncs énormes\strut &  \\
      \leavevmode\fontsize{7.3pt}{8.4pt}\selectfont\strut s'élèvent très haut, arbres plusis fois\strut &  \\
      \leavevmode\llap{\fontsize{6.4pt}{7.3pt}\selectfont\color{darkgray}10\hspace{1.8mm}}\fontsize{7.3pt}{8.4pt}\selectfont\strut séculaires. Resté lgt ds cet océan de\strut &  \\
      \leavevmode\fontsize{7.3pt}{8.4pt}\selectfont\strut verdure. Ecouté, ravi, le bruissement\strut &  \\
      \leavevmode\fontsize{7.3pt}{8.4pt}\selectfont\strut du vent à travers les branches\strut &  \\
      \leavevmode\fontsize{7.3pt}{8.4pt}\selectfont\strut feuillues\strut &  \\
      \leavevmode\fontsize{7.3pt}{8.4pt}\selectfont\strut Le soir, en revenant du mess,\strut & \smash{\parbox[t]{39mm}{\fontsize{6.2pt}{7.0pt}\selectfont\raggedright\color{black!88}\textbf{14 mess} \textit{angl. in ms.}}} \\
      \leavevmode\llap{\fontsize{6.4pt}{7.3pt}\selectfont\color{darkgray}15\hspace{1.8mm}}\fontsize{7.3pt}{8.4pt}\selectfont\strut admiré ciel [mervt ?] éclairé par lune [ajout interl. : pleine].\strut & \smash{\parbox[t]{39mm}{\fontsize{6.2pt}{7.0pt}\selectfont\raggedright\color{black!88}\textbf{15} \textit{lect. incert. in ms.} \textbf{15 pleine} \textit{add. interl. in ms.} — \textbf{16 qq} \textit{abbr. pro} quelque(s).}} \\
      \leavevmode\fontsize{7.3pt}{8.4pt}\selectfont\strut qq étoiles, çà et là. Campagne curieuse\strut &  \\
      \leavevmode\fontsize{7.3pt}{8.4pt}\selectfont\strut à regarder les moindres details, bien\strut &  \\
      \leavevmode\fontsize{7.3pt}{8.4pt}\selectfont\strut détachés –\strut &  \\
      \leavevmode\fontsize{7.3pt}{8.4pt}\selectfont\strut Me plais infinimt ds ma gde et\strut & \smash{\parbox[t]{39mm}{\fontsize{6.2pt}{7.0pt}\selectfont\raggedright\color{black!88}\textbf{19 ds} \textit{abbr. pro} dans. \textbf{19 gde} \textit{abbr. pro} grande.}} \\
      \leavevmode\llap{\fontsize{6.4pt}{7.3pt}\selectfont\color{darkgray}20\hspace{1.8mm}}\fontsize{7.3pt}{8.4pt}\selectfont\strut belle chambre, 1 vrai salon.\strut &  \\
\end{tabular}}%
\par\vspace*{\fill}
% --- Registres de bas de page (1. Critique interne [si pertinent] ; 2. Loci paralleli en latin ; 3. Commentaire historique en retrait) ---
\noindent\rule{\textwidth}{0.5pt}\par\vspace{0.65mm}
{\fontsize{6.6pt}{7.7pt}\selectfont%
\noindent\normalfont\color{black!92} \textbf{1} (\textit{de eodem commorationis loco}) cf. supra f\textsuperscript{o}~34~v\textsuperscript{o} (p.~\pageref{fol:69}).\par\vspace{0.55mm}
\noindent\rule{38mm}{0.4pt}\par\vspace{0.5mm}
\noindent\hspace*{5mm}\parbox{\dimexpr\textwidth-5mm\relax}{\fontsize{6.6pt}{7.7pt}\selectfont\color{black!84}Journée de repos au cantonnement de Barlette: longue promenade solitaire sous les futaies centenaires du \textbf{bois de Ribeaucourt} (entre Ribeaucourt, Franqueville et Bernaville) et retour nocturne du mess sous la pleine lune d'octobre.}\par}
\newpage
\phantomsection\label{fac:70}
% --- En-tête du fac-similé (Page impaire / droite : Photographie du manuscrit) ---
\noindent\rule{\textwidth}{0.5pt}\par\vspace{0.55mm}
\noindent\textbf{\fontsize{6.8pt}{7.8pt}\selectfont FAC-SIMILÉ EN VIS-À-VIS — FOLIO 35 RECTO} \hfill \textbf{\fontsize{6.8pt}{7.8pt}\selectfont OCTOBRE 1916}\par\vspace{0.25mm}
\noindent{\fontsize{6.8pt}{8.0pt}\selectfont\textit{Livre-Journal d'Hippolyte Combes} — Lundi 9 octobre 1916 (Barlette et forêt de Ribeaucourt (Somme))}\par\vspace{0.55mm}
\noindent\rule{\textwidth}{0.35pt}\par\vspace*{\fill}
{\centering
  \setlength{\fboxsep}{1.2mm}%
  \setlength{\fboxrule}{0.4pt}%
  \fcolorbox{gray!55}{white}{%
    \includegraphics[width=0.95\textwidth,height=0.81\textheight,keepaspectratio]{\facsimiledir/20251126_103857.jpg}%
  }\par}
\vspace*{\fill}
\newpage

% =============================================================================
% VIS-A-VIS SPREAD : FOLIO 35 VERSO (20251126_103913.jpg)
% =============================================================================
\cleardoubleevenpage
\renewcommand{\currentmonthyear}{Octobre 1916}
\renewcommand{\currentfoliolabel}{Folio 35 verso}
\phantomsection\label{fol:71}
% --- En-tête du folio (Page paire / gauche : Transcription, apparat & registres de bas de page) ---
\noindent\rule{\textwidth}{0.5pt}\par\vspace{0.55mm}
\noindent\textbf{\footnotesize FOLIO 35 VERSO} \hfill \textbf{\footnotesize OCTOBRE 1916}\par\vspace{0.25mm}
\noindent{\fontsize{6.9pt}{8.0pt}\selectfont Corvée de fourrage vert en lisière des bois aux couleurs d'automne et nuit de pleine lune}\par\vspace{0.2mm}
\noindent{\fontsize{6.9pt}{8.0pt}\selectfont\textit{Mardi 10 octobre 1916} \hfill \textit{Barlette et bois de Ribeaucourt (Somme)}}\par\vspace{0.5mm}
\noindent\rule{\textwidth}{0.35pt}\par\vspace{0.8mm}
% --- Corps de la transcription diplomatique et apparat marginal en prolongement ---
{\renewcommand{\arraystretch}{1.00}%
\noindent\hspace*{3.5mm}\begin{tabular}{@{}p{87mm}@{\hspace{3.5mm}}p{39mm}@{}}
      \leavevmode\fontsize{7.7pt}{9.0pt}\selectfont\strut \underline{Mardi 10} \hfill{} \underline{Oct 1916}\strut &  \\
      \leavevmode\fontsize{7.7pt}{9.0pt}\selectfont\strut Enfin, très belle journée.\strut &  \\
      \leavevmode\fontsize{7.7pt}{9.0pt}\selectfont\strut Matin \& après-midi, occupé pour green\strut & \smash{\parbox[t]{39mm}{\fontsize{6.2pt}{7.0pt}\selectfont\raggedright\color{black!88}\textbf{3 green} \textit{angl. in ms.}}} \\
      \leavevmode\fontsize{7.7pt}{9.0pt}\selectfont\strut forage. Conduit fourragères ds les\strut & \smash{\parbox[t]{39mm}{\fontsize{6.2pt}{7.0pt}\selectfont\raggedright\color{black!88}\textbf{4 ds} \textit{abbr. pro} dans. \textbf{4 forage} \textit{angl. in ms.}}} \\
      \leavevmode\llap{\fontsize{6.4pt}{7.3pt}\selectfont\color{darkgray}5\hspace{1.8mm}}\fontsize{7.7pt}{9.0pt}\selectfont\strut champs. Tout en surveillant le\strut &  \\
      \leavevmode\fontsize{7.7pt}{9.0pt}\selectfont\strut chargement, assis s/ 1 botte de\strut & \smash{\parbox[t]{39mm}{\fontsize{6.2pt}{7.0pt}\selectfont\raggedright\color{black!88}\textbf{6 s/} \textit{abbr. pro} sur.}} \\
      \leavevmode\fontsize{7.7pt}{9.0pt}\selectfont\strut luzerne – et rempli yeux vue\strut &  \\
      \leavevmode\fontsize{7.7pt}{9.0pt}\selectfont\strut des bois. Devant moi, à qq 40 ou 50\strut & \smash{\parbox[t]{39mm}{\fontsize{6.2pt}{7.0pt}\selectfont\raggedright\color{black!88}\textbf{8 qq} \textit{abbr. pro} quelque(s).}} \\
      \leavevmode\fontsize{7.7pt}{9.0pt}\selectfont\strut mètres, lisière forêt arbres\strut &  \\
      \leavevmode\llap{\fontsize{6.4pt}{7.3pt}\selectfont\color{darkgray}10\hspace{1.8mm}}\fontsize{7.7pt}{9.0pt}\selectfont\strut massés les 1 à côté des autres\strut &  \\
      \leavevmode\fontsize{7.7pt}{9.0pt}\selectfont\strut très belles teintes que le soleil rendait\strut &  \\
      \leavevmode\fontsize{7.7pt}{9.0pt}\selectfont\strut + éclatantes et qui se détachaient\strut &  \\
      \leavevmode\fontsize{7.7pt}{9.0pt}\selectfont\strut sur le ciel bleu –\strut &  \\
      \leavevmode\fontsize{7.7pt}{9.0pt}\selectfont\strut Comme ts les soirs, nuit\strut & \smash{\parbox[t]{39mm}{\fontsize{6.2pt}{7.0pt}\selectfont\raggedright\color{black!88}\textbf{14 ts} \textit{abbr. pro} tous.}} \\
      \leavevmode\llap{\fontsize{6.4pt}{7.3pt}\selectfont\color{darkgray}15\hspace{1.8mm}}\fontsize{7.7pt}{9.0pt}\selectfont\strut merveilleuse. Lune pleine\strut &  \\
      \leavevmode\fontsize{7.7pt}{9.0pt}\selectfont\strut brillant d'1 vive lueur, ds l'immensité\strut & \smash{\parbox[t]{39mm}{\fontsize{6.2pt}{7.0pt}\selectfont\raggedright\color{black!88}\textbf{16 ds} \textit{abbr. pro} dans.}} \\
      \leavevmode\fontsize{7.7pt}{9.0pt}\selectfont\strut firmament. qq rares étoiles\strut & \smash{\parbox[t]{39mm}{\fontsize{6.2pt}{7.0pt}\selectfont\raggedright\color{black!88}\textbf{17 qq} \textit{abbr. pro} quelque(s). — \textbf{18 gde} \textit{abbr. pro} grande.}} \\
      \leavevmode\fontsize{7.7pt}{9.0pt}\selectfont\strut gde paix épandue ——————\strut &  \\
\end{tabular}}%
\par\vspace*{\fill}
% --- Registres de bas de page (1. Critique interne [si pertinent] ; 2. Loci paralleli en latin ; 3. Commentaire historique en retrait) ---
\noindent\rule{\textwidth}{0.5pt}\par\vspace{0.65mm}
{\fontsize{6.6pt}{7.7pt}\selectfont%
\noindent\normalfont\color{black!92} \textbf{1} (\textit{de eodem commorationis loco}) cf. supra f\textsuperscript{o}~35~r\textsuperscript{o} (p.~\pageref{fol:70}).\par\vspace{0.55mm}
\noindent\rule{38mm}{0.4pt}\par\vspace{0.5mm}
\noindent\hspace*{5mm}\parbox{\dimexpr\textwidth-5mm\relax}{\fontsize{6.6pt}{7.7pt}\selectfont\color{black!84}Reprise du service quotidien de ravitaillement en fourrage vert (\textit{green forage}) pour les attelages de la 152\textsuperscript{e} Compagnie A.S.C., en lisière du bois de Ribeaucourt paré des couleurs d'automne.}\par}
\newpage
\phantomsection\label{fac:71}
% --- En-tête du fac-similé (Page impaire / droite : Photographie du manuscrit) ---
\noindent\rule{\textwidth}{0.5pt}\par\vspace{0.55mm}
\noindent\textbf{\fontsize{6.8pt}{7.8pt}\selectfont FAC-SIMILÉ EN VIS-À-VIS — FOLIO 35 VERSO} \hfill \textbf{\fontsize{6.8pt}{7.8pt}\selectfont OCTOBRE 1916}\par\vspace{0.25mm}
\noindent{\fontsize{6.8pt}{8.0pt}\selectfont\textit{Livre-Journal d'Hippolyte Combes} — Mardi 10 octobre 1916 (Barlette et bois de Ribeaucourt (Somme))}\par\vspace{0.55mm}
\noindent\rule{\textwidth}{0.35pt}\par\vspace*{\fill}
{\centering
  \setlength{\fboxsep}{1.2mm}%
  \setlength{\fboxrule}{0.4pt}%
  \fcolorbox{gray!55}{white}{%
    \includegraphics[width=0.95\textwidth,height=0.81\textheight,keepaspectratio]{\facsimiledir/20251126_103913.jpg}%
  }\par}
\vspace*{\fill}
\newpage

% =============================================================================
% VIS-A-VIS SPREAD : FOLIO 36 RECTO (20251126_103948.jpg)
% =============================================================================
\cleardoubleevenpage
\renewcommand{\currentmonthyear}{Octobre 1916}
\renewcommand{\currentfoliolabel}{Folio 36 recto}
\phantomsection\label{fol:72}
% --- En-tête du folio (Page paire / gauche : Transcription, apparat & registres de bas de page) ---
\noindent\rule{\textwidth}{0.5pt}\par\vspace{0.55mm}
\noindent\textbf{\footnotesize FOLIO 36 RECTO} \hfill \textbf{\footnotesize OCTOBRE 1916}\par\vspace{0.25mm}
\noindent{\fontsize{6.9pt}{8.0pt}\selectfont Convoi de luzerne sous une pluie battante derrière le grand bois de Ribeaucourt}\par\vspace{0.2mm}
\noindent{\fontsize{6.9pt}{8.0pt}\selectfont\textit{Mercredi 11 octobre 1916} \hfill \textit{Barlette et Ribeaucourt (Somme)}}\par\vspace{0.5mm}
\noindent\rule{\textwidth}{0.35pt}\par\vspace{0.8mm}
% --- Corps de la transcription diplomatique et apparat marginal en prolongement ---
{\renewcommand{\arraystretch}{1.00}%
\noindent\hspace*{3.5mm}\begin{tabular}{@{}p{87mm}@{\hspace{3.5mm}}p{39mm}@{}}
      \leavevmode\fontsize{7.7pt}{9.0pt}\selectfont\strut \underline{Mer. 11} \hfill{} \underline{Oct / 16}\strut & \smash{\parbox[t]{39mm}{\fontsize{6.2pt}{7.0pt}\selectfont\raggedright\color{black!88}\textbf{1 Mer.} \textit{scil.} Mercredi.}} \\
      \leavevmode\fontsize{7.7pt}{9.0pt}\selectfont\strut Matinée employée de =\strut & \smash{\parbox[t]{39mm}{\fontsize{6.2pt}{7.0pt}\selectfont\raggedright\color{black!88}\textbf{2 de =} \textit{signum aequalitatis pro} de même.}} \\
      \leavevmode\fontsize{7.7pt}{9.0pt}\selectfont\strut manière qu'hier. Sous pluie battante,\strut &  \\
      \leavevmode\fontsize{7.7pt}{9.0pt}\selectfont\strut dû conduire wagons ds 1 champ de luzerne\strut & \smash{\parbox[t]{39mm}{\fontsize{6.2pt}{7.0pt}\selectfont\raggedright\color{black!88}\textbf{4 ds} \textit{abbr. pro} dans.}} \\
      \leavevmode\llap{\fontsize{6.4pt}{7.3pt}\selectfont\color{darkgray}5\hspace{1.8mm}}\fontsize{7.7pt}{9.0pt}\selectfont\strut derrière le grand bois de Ribeaucourt\strut &  \\
      \leavevmode\fontsize{7.7pt}{9.0pt}\selectfont\strut Rentré billet, casque \& cape tt ruisselants\strut & \smash{\parbox[t]{39mm}{\fontsize{6.2pt}{7.0pt}\selectfont\raggedright\color{black!88}\textbf{6 tt} \textit{abbr. pro} tout.}} \\
      \leavevmode\fontsize{7.7pt}{9.0pt}\selectfont\strut Après-midi, liberté entière –\strut &  \\
      \leavevmode\fontsize{7.7pt}{9.0pt}\selectfont\strut Tt le jour, ciel nuageux et temps\strut & \smash{\parbox[t]{39mm}{\fontsize{6.2pt}{7.0pt}\selectfont\raggedright\color{black!88}\textbf{8 Tt} \textit{abbr. pro} Tout.}} \\
      \leavevmode\fontsize{7.7pt}{9.0pt}\selectfont\strut pluvieux. – –\strut &  \\
\end{tabular}}%
\par\vspace*{\fill}
% --- Registres de bas de page (1. Critique interne [si pertinent] ; 2. Loci paralleli en latin ; 3. Commentaire historique en retrait) ---
\noindent\rule{\textwidth}{0.5pt}\par\vspace{0.65mm}
{\fontsize{6.6pt}{7.7pt}\selectfont%
\noindent\normalfont\color{black!92} \textbf{5} (\textit{Ribeaucourt}) cf. supra f\textsuperscript{o}~33~v\textsuperscript{o} (p.~\pageref{fol:67}), f\textsuperscript{o}~35~v\textsuperscript{o} (p.~\pageref{fol:71}), infra f\textsuperscript{o}~36~v\textsuperscript{o} (p.~\pageref{fol:73}), f\textsuperscript{o}~37~r\textsuperscript{o} (p.~\pageref{fol:74}), \textit{et passim}.\par\vspace{0.55mm}
\noindent\rule{38mm}{0.4pt}\par\vspace{0.5mm}
\noindent\hspace*{5mm}\parbox{\dimexpr\textwidth-5mm\relax}{\fontsize{6.6pt}{7.7pt}\selectfont\color{black!84}Corvée matinale de fourrage vert sous les averses d'octobre derrière le grand bois de \textbf{Ribeaucourt} (Somme), suivie d'une après-midi de liberté dans sa chambre de Barlette.}\par}
\newpage
\phantomsection\label{fac:72}
% --- En-tête du fac-similé (Page impaire / droite : Photographie du manuscrit) ---
\noindent\rule{\textwidth}{0.5pt}\par\vspace{0.55mm}
\noindent\textbf{\fontsize{6.8pt}{7.8pt}\selectfont FAC-SIMILÉ EN VIS-À-VIS — FOLIO 36 RECTO} \hfill \textbf{\fontsize{6.8pt}{7.8pt}\selectfont OCTOBRE 1916}\par\vspace{0.25mm}
\noindent{\fontsize{6.8pt}{8.0pt}\selectfont\textit{Livre-Journal d'Hippolyte Combes} — Mercredi 11 octobre 1916 (Barlette et Ribeaucourt (Somme))}\par\vspace{0.55mm}
\noindent\rule{\textwidth}{0.35pt}\par\vspace*{\fill}
{\centering
  \setlength{\fboxsep}{1.2mm}%
  \setlength{\fboxrule}{0.4pt}%
  \fcolorbox{gray!55}{white}{%
    \includegraphics[width=0.95\textwidth,height=0.81\textheight,keepaspectratio]{\facsimiledir/20251126_103948.jpg}%
  }\par}
\vspace*{\fill}
\newpage

% =============================================================================
% VIS-A-VIS SPREAD : FOLIO 36 VERSO (20251126_104011.jpg)
% =============================================================================
\cleardoubleevenpage
\renewcommand{\currentmonthyear}{Octobre 1916}
\renewcommand{\currentfoliolabel}{Folio 36 verso}
\phantomsection\label{fol:73}
% --- En-tête du folio (Page paire / gauche : Transcription, apparat & registres de bas de page) ---
\noindent\rule{\textwidth}{0.5pt}\par\vspace{0.55mm}
\noindent\textbf{\footnotesize FOLIO 36 VERSO} \hfill \textbf{\footnotesize OCTOBRE 1916}\par\vspace{0.25mm}
\noindent{\fontsize{6.9pt}{8.0pt}\selectfont Conduite d'un « limber G.S. », symphonie du vent dans les bois et chevauchée à Prouville avec le capitaine Panton}\par\vspace{0.2mm}
\noindent{\fontsize{6.9pt}{8.0pt}\selectfont\textit{Jeudi 12 et vendredi 13 octobre 1916} \hfill \textit{Barlette, Ribeaucourt et Prouville (Somme)}}\par\vspace{0.5mm}
\noindent\rule{\textwidth}{0.35pt}\par\vspace{0.8mm}
% --- Corps de la transcription diplomatique et apparat marginal en prolongement ---
{\renewcommand{\arraystretch}{1.00}%
\noindent\hspace*{3.5mm}\begin{tabular}{@{}p{87mm}@{\hspace{3.5mm}}p{39mm}@{}}
      \leavevmode\fontsize{7.7pt}{9.0pt}\selectfont\strut \underline{Jeudi 12} \hfill{} \underline{Octb / 16}\strut &  \\
      \leavevmode\fontsize{7.7pt}{9.0pt}\selectfont\strut Temps pluvieux\strut &  \\
      \leavevmode\fontsize{7.7pt}{9.0pt}\selectfont\strut Après-midi, conduit limber \& G.S.\strut & \smash{\parbox[t]{39mm}{\fontsize{6.2pt}{7.0pt}\selectfont\raggedright\color{black!88}\textbf{3 limber} \textit{angl. in ms.}}} \\
      \leavevmode\fontsize{7.7pt}{9.0pt}\selectfont\strut wagons ds 1 gd champ où luzerne\strut & \smash{\parbox[t]{39mm}{\fontsize{6.2pt}{7.0pt}\selectfont\raggedright\color{black!88}\textbf{4 ds} \textit{abbr. pro} dans. \textbf{4 gd} \textit{abbr. pro} grand.}} \\
      \leavevmode\llap{\fontsize{6.4pt}{7.3pt}\selectfont\color{darkgray}5\hspace{1.8mm}}\fontsize{7.7pt}{9.0pt}\selectfont\strut coupée –\strut &  \\
      \leavevmode\fontsize{7.7pt}{9.0pt}\selectfont\strut A l'aller \& au retour, me suis\strut &  \\
      \leavevmode\fontsize{7.7pt}{9.0pt}\selectfont\strut attardé ds bois : vent soufflant\strut & \smash{\parbox[t]{39mm}{\fontsize{6.2pt}{7.0pt}\selectfont\raggedright\color{black!88}\textbf{7 ds} \textit{abbr. pro} dans.}} \\
      \leavevmode\fontsize{7.7pt}{9.0pt}\selectfont\strut en rafales produisait 1 bruit de\strut &  \\
      \leavevmode\fontsize{7.7pt}{9.0pt}\selectfont\strut houle.. – Resté lgs à écouter cette\strut & \smash{\parbox[t]{39mm}{\fontsize{6.2pt}{7.0pt}\selectfont\raggedright\color{black!88}\textbf{9 lgs} \textit{abbr. pro} longtemps.}} \\
      \leavevmode\llap{\fontsize{6.4pt}{7.3pt}\selectfont\color{darkgray}10\hspace{1.8mm}}\fontsize{7.7pt}{9.0pt}\selectfont\strut symphonie très forte\strut &  \\
      \cline{1-1}
      \leavevmode\fontsize{7.7pt}{9.0pt}\selectfont\strut \underline{Ven. 13.} \hfill{} Tjours ciel bas et gris\strut & \smash{\parbox[t]{39mm}{\fontsize{6.2pt}{7.0pt}\selectfont\raggedright\color{black!88}\textbf{11 Ven.} \textit{scil.} Vendredi.}} \\
      \leavevmode\fontsize{7.7pt}{9.0pt}\selectfont\strut de cendre. grosses nuées courent,\strut &  \\
      \leavevmode\fontsize{7.7pt}{9.0pt}\selectfont\strut chassées par le vent\strut &  \\
      \leavevmode\fontsize{7.7pt}{9.0pt}\selectfont\strut Après-midi, course à\strut &  \\
      \leavevmode\llap{\fontsize{6.4pt}{7.3pt}\selectfont\color{darkgray}15\hspace{1.8mm}}\fontsize{7.7pt}{9.0pt}\selectfont\strut cheval av. Cap. Pant. au village de\strut & \smash{\parbox[t]{39mm}{\fontsize{6.2pt}{7.0pt}\selectfont\raggedright\color{black!88}\textbf{15 Cap.} \textit{abbr. pro} Capitaine.}} \\
      \leavevmode\fontsize{7.7pt}{9.0pt}\selectfont\strut Prouville, distant de 4 k..\strut &  \\
      \leavevmode\fontsize{7.7pt}{9.0pt}\selectfont\strut Soir, au diner, appris que ns\strut & \smash{\parbox[t]{39mm}{\fontsize{6.2pt}{7.0pt}\selectfont\raggedright\color{black!88}\textbf{17 ns} \textit{abbr. pro} nous.}} \\
      \leavevmode\fontsize{7.7pt}{9.0pt}\selectfont\strut shifterons dimanche. Regret de partir!\strut & \smash{\parbox[t]{39mm}{\fontsize{6.2pt}{7.0pt}\selectfont\raggedright\color{black!88}\textbf{18 shifterons} \textit{angl. in ms.}}} \\
\end{tabular}}%
\par\vspace*{\fill}
% --- Registres de bas de page (1. Critique interne [si pertinent] ; 2. Loci paralleli en latin ; 3. Commentaire historique en retrait) ---
\noindent\rule{\textwidth}{0.5pt}\par\vspace{0.65mm}
{\fontsize{6.6pt}{7.7pt}\selectfont%
\noindent\normalfont\color{black!92} \textbf{1} (\textit{de eodem commorationis loco}) cf. supra f\textsuperscript{o}~36~r\textsuperscript{o} (p.~\pageref{fol:72}).\par\vspace{0.55mm}
\noindent\rule{38mm}{0.4pt}\par\vspace{0.5mm}
\noindent\hspace*{5mm}\parbox{\dimexpr\textwidth-5mm\relax}{\fontsize{6.6pt}{7.7pt}\selectfont\color{black!84}Le jeudi 12 octobre, Hippolyte Combes conduit aux champs de luzerne un avant-train articulé (\textit{limber}) et des fourgons réglementaires à quatre roues de l'A.S.C. (\textit{General Service wagons}).}\par}
\newpage
\phantomsection\label{fac:73}
% --- En-tête du fac-similé (Page impaire / droite : Photographie du manuscrit) ---
\noindent\rule{\textwidth}{0.5pt}\par\vspace{0.55mm}
\noindent\textbf{\fontsize{6.8pt}{7.8pt}\selectfont FAC-SIMILÉ EN VIS-À-VIS — FOLIO 36 VERSO} \hfill \textbf{\fontsize{6.8pt}{7.8pt}\selectfont OCTOBRE 1916}\par\vspace{0.25mm}
\noindent{\fontsize{6.8pt}{8.0pt}\selectfont\textit{Livre-Journal d'Hippolyte Combes} — Jeudi 12 et vendredi 13 octobre 1916 (Barlette, Ribeaucourt et Prouville (Somme))}\par\vspace{0.55mm}
\noindent\rule{\textwidth}{0.35pt}\par\vspace*{\fill}
{\centering
  \setlength{\fboxsep}{1.2mm}%
  \setlength{\fboxrule}{0.4pt}%
  \fcolorbox{gray!55}{white}{%
    \includegraphics[width=0.95\textwidth,height=0.81\textheight,keepaspectratio]{\facsimiledir/20251126_104011.jpg}%
  }\par}
\vspace*{\fill}
\newpage

% =============================================================================
% VIS-A-VIS SPREAD : FOLIO 37 RECTO (20251126_104027.jpg)
% =============================================================================
\cleardoubleevenpage
\renewcommand{\currentmonthyear}{Octobre 1916}
\renewcommand{\currentfoliolabel}{Folio 37 recto}
\phantomsection\label{fol:74}
% --- En-tête du folio (Page paire / gauche : Transcription, apparat & registres de bas de page) ---
\noindent\rule{\textwidth}{0.5pt}\par\vspace{0.55mm}
\noindent\textbf{\footnotesize FOLIO 37 RECTO} \hfill \textbf{\footnotesize OCTOBRE 1916}\par\vspace{0.25mm}
\noindent{\fontsize{6.9pt}{8.0pt}\selectfont Départ définitif de Barlette sous les averses et arrivée de la colonne à Beauval}\par\vspace{0.2mm}
\noindent{\fontsize{6.9pt}{8.0pt}\selectfont\textit{Dimanche 15 octobre 1916} \hfill \textit{De Barlette (Ribeaucourt) à Beauval (Somme)}}\par\vspace{0.5mm}
\noindent\rule{\textwidth}{0.35pt}\par\vspace{0.8mm}
% --- Corps de la transcription diplomatique et apparat marginal en prolongement ---
{\renewcommand{\arraystretch}{0.87}%
\noindent\hspace*{3.5mm}\begin{tabular}{@{}p{87mm}@{\hspace{3.5mm}}p{39mm}@{}}
      \leavevmode\fontsize{6.9pt}{7.9pt}\selectfont\strut \underline{Dim. 15} \hfill{} \underline{Oct / 16}\strut & \smash{\parbox[t]{39mm}{\fontsize{6.2pt}{7.0pt}\selectfont\raggedright\color{black!88}\textbf{1 Dim.} \textit{scil.} Dimanche.}} \\
      \leavevmode\fontsize{6.9pt}{7.9pt}\selectfont\strut Levé 6 H 1/2. Ts paquets terminés. Quitté maison définitt\strut & \smash{\parbox[t]{39mm}{\fontsize{6.2pt}{7.0pt}\selectfont\raggedright\color{black!88}\textbf{2 Ts} \textit{abbr. pro} Tous. \textbf{2 définitt} \textit{abbr. pro} définitivement. — \textbf{3 lgt} \textit{abbr. pro} longtemps. \textbf{3 mess} \textit{angl. in ms.} — \textbf{6 m.} \textit{abbr. pro} matin.}} \\
      \leavevmode\fontsize{6.9pt}{7.9pt}\selectfont\strut Au mess, lgt avant 9 h.\strut &  \\
      \leavevmode\fontsize{6.9pt}{7.9pt}\selectfont\strut Ciel couvert de gros nuages. Pluie fine\strut &  \\
      \leavevmode\llap{\fontsize{6.4pt}{7.3pt}\selectfont\color{darkgray}5\hspace{1.8mm}}\fontsize{6.9pt}{7.9pt}\selectfont\strut Matinée perdue à attendre moment départ\strut &  \\
      \leavevmode\fontsize{6.9pt}{7.9pt}\selectfont\strut Enfin, en route vers 11 H 15 m.\strut &  \\
      \leavevmode\fontsize{6.9pt}{7.9pt}\selectfont\strut Lente chevauchée à travers région paisible et pittoresque\strut &  \\
      \leavevmode\fontsize{6.9pt}{7.9pt}\selectfont\strut Arrêts fréquents le long routes\strut &  \\
      \leavevmode\fontsize{6.9pt}{7.9pt}\selectfont\strut Au delà de [blanc], averses très fortes\strut & \smash{\parbox[t]{39mm}{\fontsize{6.2pt}{7.0pt}\selectfont\raggedright\color{black!88}\textbf{9} \textit{spatium relictum in ms.}}} \\
      \leavevmode\llap{\fontsize{6.4pt}{7.3pt}\selectfont\color{darkgray}10\hspace{1.8mm}}\fontsize{6.9pt}{7.9pt}\selectfont\strut Enfin à 3 H 1/2 s, arrivés à Beauval. Alors ciel\strut & \smash{\parbox[t]{39mm}{\fontsize{6.2pt}{7.0pt}\selectfont\raggedright\color{black!88}\textbf{10 1/2} \textit{scil.} demi / moitié. \textbf{10 s.} \textit{abbr. pro} soir.}} \\
      \leavevmode\fontsize{6.9pt}{7.9pt}\selectfont\strut éclairci soleil brille\strut &  \\
      \leavevmode\fontsize{6.9pt}{7.9pt}\selectfont\strut Petite ville très coquette — remplie de troupes\strut &  \\
      \leavevmode\fontsize{6.9pt}{7.9pt}\selectfont\strut Trouvé petite chambre en face de notre mess –\strut & \smash{\parbox[t]{39mm}{\fontsize{6.2pt}{7.0pt}\selectfont\raggedright\color{black!88}\textbf{13 mess} \textit{angl. in ms.}}} \\
      \leavevmode\fontsize{6.9pt}{7.9pt}\selectfont\strut A Barlette : washing \& darning : 70 c. — Connelly : 1 f —\strut & \smash{\parbox[t]{39mm}{\fontsize{6.2pt}{7.0pt}\selectfont\raggedright\color{black!88}\textbf{14 c.} \textit{abbr. pro} centimes. \textbf{14 washing} \textit{angl. in ms.} — \textbf{15 Cartes ples} \textit{abbr. pro} Cartes postales.}} \\
      \cline{1-1}
      \leavevmode\llap{\fontsize{6.4pt}{7.3pt}\selectfont\color{darkgray}15\hspace{1.8mm}}\fontsize{6.9pt}{7.9pt}\selectfont\strut Cartes ples à Beauval : 15 c\strut &  \\
\end{tabular}}%
\par\vspace*{\fill}
% --- Registres de bas de page (1. Critique interne [si pertinent] ; 2. Loci paralleli en latin ; 3. Commentaire historique en retrait) ---
\noindent\rule{\textwidth}{0.5pt}\par\vspace{0.65mm}
{\fontsize{6.6pt}{7.7pt}\selectfont%
\noindent\normalfont \textbf{Renvoi infrapaginal autographe :} À la l.~10, Hippolyte Combes insère un appel de note \texttt{(1)} renvoyant au post-scriptum de cinq lignes (l.~17--21) tracé au bas du feuillet sous un trait horizontal de séparation.\par\vspace{0.5mm}
\noindent\rule{22mm}{0.3pt}\par\vspace{0.45mm}
\noindent\normalfont\color{black!92} \textbf{10} (\textit{Beauval}) cf. infra f\textsuperscript{o}~37~v\textsuperscript{o} (p.~\pageref{fol:75}), f\textsuperscript{o}~38~r\textsuperscript{o} (p.~\pageref{fol:76}). — \textbf{14} (\textit{Connelly}) cf. supra f\textsuperscript{o}~6~r\textsuperscript{o} (p.~\pageref{fol:12}), f\textsuperscript{o}~22~v\textsuperscript{o} (p.~\pageref{fol:45}). — \textbf{14} (\textit{Barlette}) cf. supra f\textsuperscript{o}~33~r\textsuperscript{o} (p.~\pageref{fol:66}), f\textsuperscript{o}~36~v\textsuperscript{o} (p.~\pageref{fol:73}), \textit{et passim}.\par\vspace{0.55mm}
\noindent\rule{38mm}{0.4pt}\par\vspace{0.5mm}
\noindent\hspace*{5mm}\parbox{\dimexpr\textwidth-5mm\relax}{\fontsize{6.6pt}{7.7pt}\selectfont\color{black!84}Le dimanche 15 octobre 1916 débute la remontée en ligne du train divisionnaire de la \textbf{18\textsuperscript{e} Division britannique} vers le secteur d'Albert et de la tranchée Regina: partie de Barlette à 11~h~15 sous les averses, la colonne fait étape à 15~h~30 dans le bourg industriel de \textbf{Beauval} (Somme, sur la route de Doullens à Amiens), encombré de troupes britanniques en mouvement.}\par}
\newpage
\phantomsection\label{fac:74}
% --- En-tête du fac-similé (Page impaire / droite : Photographie du manuscrit) ---
\noindent\rule{\textwidth}{0.5pt}\par\vspace{0.55mm}
\noindent\textbf{\fontsize{6.8pt}{7.8pt}\selectfont FAC-SIMILÉ EN VIS-À-VIS — FOLIO 37 RECTO} \hfill \textbf{\fontsize{6.8pt}{7.8pt}\selectfont OCTOBRE 1916}\par\vspace{0.25mm}
\noindent{\fontsize{6.8pt}{8.0pt}\selectfont\textit{Livre-Journal d'Hippolyte Combes} — Dimanche 15 octobre 1916 (De Barlette (Ribeaucourt) à Beauval (Somme))}\par\vspace{0.55mm}
\noindent\rule{\textwidth}{0.35pt}\par\vspace*{\fill}
{\centering
  \setlength{\fboxsep}{1.2mm}%
  \setlength{\fboxrule}{0.4pt}%
  \fcolorbox{gray!55}{white}{%
    \includegraphics[width=0.95\textwidth,height=0.81\textheight,keepaspectratio]{\facsimiledir/20251126_104027.jpg}%
  }\par}
\vspace*{\fill}
\newpage

% =============================================================================
% VIS-A-VIS SPREAD : FOLIO 37 VERSO (20251126_104057.jpg)
% =============================================================================
\cleardoubleevenpage
\renewcommand{\currentmonthyear}{Octobre 1916}
\renewcommand{\currentfoliolabel}{Folio 37 verso}
\phantomsection\label{fol:75}
% --- En-tête du folio (Page paire / gauche : Transcription, apparat & registres de bas de page) ---
\noindent\rule{\textwidth}{0.5pt}\par\vspace{0.55mm}
\noindent\textbf{\footnotesize FOLIO 37 VERSO} \hfill \textbf{\footnotesize OCTOBRE 1916}\par\vspace{0.25mm}
\noindent{\fontsize{6.9pt}{8.0pt}\selectfont Départ de Beauval sous un air glacial et marche vers Warloy-Baillon}\par\vspace{0.2mm}
\noindent{\fontsize{6.9pt}{8.0pt}\selectfont\textit{16 octobre 1916} \hfill \textit{Beauval, Beauquesne, Puchevillers, Warloy-Baillon (Somme)}}\par\vspace{0.5mm}
\noindent\rule{\textwidth}{0.35pt}\par\vspace{0.8mm}
% --- Corps de la transcription diplomatique et apparat marginal en prolongement ---
{\renewcommand{\arraystretch}{0.87}%
\noindent\hspace*{3.5mm}\begin{tabular}{@{}p{87mm}@{\hspace{3.5mm}}p{39mm}@{}}
      \leavevmode\fontsize{6.9pt}{7.9pt}\selectfont\strut \underline{Lundi 16} — | Sur pied à 5h 1/2 — Très bonne nuit. Oct /16\strut & \smash{\parbox[t]{39mm}{\fontsize{6.2pt}{7.0pt}\selectfont\raggedright\color{black!88}\textbf{1 1/2} \textit{scil.} demi / moitié.}} \\
      \leavevmode\fontsize{6.9pt}{7.9pt}\selectfont\strut De ma fenêtre, pu contempler pleine lune — admirable —\strut &  \\
      \leavevmode\fontsize{6.9pt}{7.9pt}\selectfont\strut Ciel très beau —\strut &  \\
      \leavevmode\fontsize{6.9pt}{7.9pt}\selectfont\strut Quitté Beauval à 7h 30 m.\strut & \smash{\parbox[t]{39mm}{\fontsize{6.2pt}{7.0pt}\selectfont\raggedright\color{black!88}\textbf{4 m.} \textit{abbr. pro} matin.}} \\
      \leavevmode\llap{\fontsize{6.4pt}{7.3pt}\selectfont\color{darkgray}5\hspace{1.8mm}}\fontsize{6.9pt}{7.9pt}\selectfont\strut Colonne, convoi, partie de la colline où campé — à côté du\strut &  \\
      \leavevmode\fontsize{6.9pt}{7.9pt}\selectfont\strut cimetière.\strut &  \\
      \leavevmode\fontsize{6.9pt}{7.9pt}\selectfont\strut Passé par Beauquesne\strut &  \\
      \leavevmode\fontsize{6.9pt}{7.9pt}\selectfont\strut monté plateau dominant village \& contrée environnante\strut &  \\
      \leavevmode\fontsize{6.9pt}{7.9pt}\selectfont\strut Fait halte au-delà de Puchevillers\strut &  \\
      \leavevmode\llap{\fontsize{6.4pt}{7.3pt}\selectfont\color{darkgray}10\hspace{1.8mm}}\fontsize{6.9pt}{7.9pt}\selectfont\strut Matinée, air glacial, vent pénétrant\strut &  \\
      \leavevmode\fontsize{6.9pt}{7.9pt}\selectfont\strut Tout le long route, très nombreuses unités en marche,\strut &  \\
      \leavevmode\fontsize{6.9pt}{7.9pt}\selectfont\strut convois d'A.S.C., R.F.A. etc. etc. \& infantry battalions\strut & \smash{\parbox[t]{39mm}{\fontsize{6.2pt}{7.0pt}\selectfont\raggedright\color{black!88}\textbf{12 A.S.C.} \textit{angl.} Army Service Corps. \textbf{12 R.F.A.} \textit{angl.} Royal Field Artillery. — \textbf{13 s.} \textit{abbr. pro} soir. — \textbf{14 tommies} \textit{angl. in ms.} — \textbf{15 ds} \textit{abbr. pro} dans. — \textbf{17 ds} \textit{abbr. pro} dans.}} \\
      \leavevmode\fontsize{6.9pt}{7.9pt}\selectfont\strut Vers 6h s., enfin atteint Warloy-Baillon. Village\strut &  \\
      \leavevmode\fontsize{6.9pt}{7.9pt}\selectfont\strut vraie fourmilière de tommies.\strut &  \\
      \leavevmode\llap{\fontsize{6.4pt}{7.3pt}\selectfont\color{darkgray}15\hspace{1.8mm}}\fontsize{6.9pt}{7.9pt}\selectfont\strut Campés sous tentes, ds 1 champ, en contrebas de la\strut &  \\
      \leavevmode\fontsize{6.9pt}{7.9pt}\selectfont\strut route de Senlis.\strut &  \\
      \leavevmode\fontsize{6.9pt}{7.9pt}\selectfont\strut Visité ds la brasserie et revu amis.\strut &  \\
\end{tabular}}%
\par\vspace*{\fill}
% --- Registres de bas de page (1. Critique interne [si pertinent] ; 2. Loci paralleli en latin ; 3. Commentaire historique en retrait) ---
\noindent\rule{\textwidth}{0.5pt}\par\vspace{0.65mm}
{\fontsize{6.6pt}{7.7pt}\selectfont%
\noindent\normalfont \textbf{Unité de page :} La page s'ouvre par un bloc comptable entre crochets (l.~1--4) soldant les frais de logement à Ribeaucourt, suivi de l'entrée du lundi 16 octobre 1916 (l.~5--21).\par\vspace{0.5mm}
\noindent\rule{22mm}{0.3pt}\par\vspace{0.45mm}
\noindent\normalfont\color{black!92} \textbf{4} (\textit{Beauval}) cf. supra f\textsuperscript{o}~37~r\textsuperscript{o} (p.~\pageref{fol:74}), infra f\textsuperscript{o}~38~r\textsuperscript{o} (p.~\pageref{fol:76}). — \textbf{7} (\textit{Beauquesne}) cf. infra f\textsuperscript{o}~49~r\textsuperscript{o} (p.~\pageref{fol:98}), f\textsuperscript{o}~56~r\textsuperscript{o} (p.~\pageref{fol:112}), \textit{et passim}. — \textbf{9} (\textit{Puchevillers}) cf. infra f\textsuperscript{o}~52~r\textsuperscript{o} (p.~\pageref{fol:104}), f\textsuperscript{o}~55~r\textsuperscript{o} (p.~\pageref{fol:110}), \textit{et passim}.\par\vspace{0.55mm}
\noindent\rule{38mm}{0.4pt}\par\vspace{0.5mm}
\noindent\hspace*{5mm}\parbox{\dimexpr\textwidth-5mm\relax}{\fontsize{6.6pt}{7.7pt}\selectfont\color{black!84}Deuxième étape de concentration vers le front de l'Ancre (lundi 16 octobre 1916): en dix heures et demie de marche par \textbf{Beauquesne} et \textbf{Puchevillers} au milieu d'un flot continu de colonnes d'artillerie de campagne (\textit{R.F.A.}), du train (\textit{A.S.C.}) et de bataillons d'infanterie, la compagnie gagne \textbf{Warloy-Baillon} (Somme, à 7~km à l'ouest d'Albert), où elle bivouaque sous tentes en contrebas de la route de Senlis-le-Sec.}\par}
\newpage
\phantomsection\label{fac:75}
% --- En-tête du fac-similé (Page impaire / droite : Photographie du manuscrit) ---
\noindent\rule{\textwidth}{0.5pt}\par\vspace{0.55mm}
\noindent\textbf{\fontsize{6.8pt}{7.8pt}\selectfont FAC-SIMILÉ EN VIS-À-VIS — FOLIO 37 VERSO} \hfill \textbf{\fontsize{6.8pt}{7.8pt}\selectfont OCTOBRE 1916}\par\vspace{0.25mm}
\noindent{\fontsize{6.8pt}{8.0pt}\selectfont\textit{Livre-Journal d'Hippolyte Combes} — 16 octobre 1916 (Beauval, Beauquesne, Puchevillers, Warloy-Baillon (Somme))}\par\vspace{0.55mm}
\noindent\rule{\textwidth}{0.35pt}\par\vspace*{\fill}
{\centering
  \setlength{\fboxsep}{1.2mm}%
  \setlength{\fboxrule}{0.4pt}%
  \fcolorbox{gray!55}{white}{%
    \includegraphics[width=0.95\textwidth,height=0.81\textheight,keepaspectratio]{\facsimiledir/20251126_104057.jpg}%
  }\par}
\vspace*{\fill}
\newpage

% =============================================================================
% VIS-A-VIS SPREAD : FOLIO 38 RECTO (20251126_104118.jpg)
% =============================================================================
\cleardoubleevenpage
\renewcommand{\currentmonthyear}{Octobre 1916}
\renewcommand{\currentfoliolabel}{Folio 38 recto}
\phantomsection\label{fol:76}
% --- En-tête du folio (Page paire / gauche : Transcription, apparat & registres de bas de page) ---
\noindent\rule{\textwidth}{0.5pt}\par\vspace{0.55mm}
\noindent\textbf{\footnotesize FOLIO 38 RECTO} \hfill \textbf{\footnotesize OCTOBRE 1916}\par\vspace{0.25mm}
\noindent{\fontsize{6.9pt}{8.0pt}\selectfont Nuit glaciale sous la tente à Warloy-Baillon et départ de la colonne}\par\vspace{0.2mm}
\noindent{\fontsize{6.9pt}{8.0pt}\selectfont\textit{16 (comptes) – 17 octobre 1916 (matin)} \hfill \textit{Beauval, Warloy-Baillon (Somme)}}\par\vspace{0.5mm}
\noindent\rule{\textwidth}{0.35pt}\par\vspace{0.8mm}
% --- Corps de la transcription diplomatique et apparat marginal en prolongement ---
{\renewcommand{\arraystretch}{1.00}%
\noindent\hspace*{3.5mm}\begin{tabular}{@{}p{87mm}@{\hspace{3.5mm}}p{39mm}@{}}
      \leavevmode\fontsize{7.7pt}{9.0pt}\selectfont\strut \underline{16 Lundi} \hfill{} \underline{Oct /16}\strut &  \\
      \leavevmode\fontsize{7.7pt}{9.0pt}\selectfont\strut [ A Beauval, room : 1 f — groom : 1 f\strut & \smash{\parbox[t]{39mm}{\fontsize{6.2pt}{7.0pt}\selectfont\raggedright\color{black!88}\textbf{2} \textit{in marg. cum unco} (comptes). \textbf{2 groom} \textit{angl. in ms.} — \textbf{3 ctes ples} \textit{abbr. pro} cartes postales. — \textbf{5 off.} \textit{abbr. pro} officier(s).}} \\
      \leavevmode\fontsize{7.7pt}{9.0pt}\selectfont\strut A Warloy, ctes ples : 25 c — ]\strut &  \\
      \cline{1-1}
      \leavevmode\fontsize{7.7pt}{9.0pt}\selectfont\strut \underline{17 Mardi}\strut &  \\
      \leavevmode\llap{\fontsize{6.4pt}{7.3pt}\selectfont\color{darkgray}5\hspace{1.8mm}}\fontsize{7.7pt}{9.0pt}\selectfont\strut Couché sous tente, à côté 2 off.\strut &  \\
      \leavevmode\fontsize{7.7pt}{9.0pt}\selectfont\strut Nuit très froide.\strut &  \\
      \leavevmode\fontsize{7.7pt}{9.0pt}\selectfont\strut Sensations nouvelles bivouac jour et nuit constant.\strut &  \\
      \leavevmode\fontsize{7.7pt}{9.0pt}\selectfont\strut Levé vers 6h m. Au ciel bleu pâle, 1/2 lune argent\strut & \smash{\parbox[t]{39mm}{\fontsize{6.2pt}{7.0pt}\selectfont\raggedright\color{black!88}\textbf{8 1/2} \textit{scil.} demi / moitié. \textbf{8 m.} \textit{abbr. pro} matin. — \textbf{9 ds} \textit{abbr. pro} dans. \textbf{9 tt} \textit{abbr. pro} tout. — \textbf{10 infrie} \textit{abbr. pro} infanterie. — \textbf{11 m.} \textit{abbr. pro} matin. — \textbf{12 tte} \textit{abbr. pro} toute. \textbf{12 tjrs} \textit{abbr. pro} toujours.}} \\
      \leavevmode\fontsize{7.7pt}{9.0pt}\selectfont\strut Grande animation ds tt le camp\strut &  \\
      \leavevmode\llap{\fontsize{6.4pt}{7.3pt}\selectfont\color{darkgray}10\hspace{1.8mm}}\fontsize{7.7pt}{9.0pt}\selectfont\strut Régiments d'infrie, partis d'abord —\strut &  \\
      \leavevmode\fontsize{7.7pt}{9.0pt}\selectfont\strut En route, à 10h m. —\strut &  \\
      \leavevmode\fontsize{7.7pt}{9.0pt}\selectfont\strut Routes tjrs encombrées et sillonnées convois de tte\strut &  \\
      \leavevmode\fontsize{7.7pt}{9.0pt}\selectfont\strut sorte\strut &  \\
\end{tabular}}%
\par\vspace*{\fill}
% --- Registres de bas de page (1. Critique interne [si pertinent] ; 2. Loci paralleli en latin ; 3. Commentaire historique en retrait) ---
\noindent\rule{\textwidth}{0.5pt}\par\vspace{0.65mm}
{\fontsize{6.6pt}{7.7pt}\selectfont%
\noindent\normalfont\color{black!92} \textbf{2} (\textit{Beauval}) cf. supra f\textsuperscript{o}~37~r\textsuperscript{o} (p.~\pageref{fol:74}), f\textsuperscript{o}~37~v\textsuperscript{o} (p.~\pageref{fol:75}). — \textbf{3} (\textit{Warloy}) cf. supra f\textsuperscript{o}~37~v\textsuperscript{o} (p.~\pageref{fol:75}), infra f\textsuperscript{o}~48~r\textsuperscript{o} (p.~\pageref{fol:96}), f\textsuperscript{o}~84~v\textsuperscript{o} (p.~\pageref{fol:169}), \textit{et passim}.\par\vspace{0.55mm}
\noindent\rule{38mm}{0.4pt}\par\vspace{0.5mm}
\noindent\hspace*{5mm}\parbox{\dimexpr\textwidth-5mm\relax}{\fontsize{6.6pt}{7.7pt}\selectfont\color{black!84}Dans le village de \textbf{Warloy-Baillon} saturé de troupes et d'ambulances (\textit{Field Ambulances}), Hippolyte Combes bivouaque pour la première fois sous la tente avec deux officiers par une nuit glaciale, avant le départ de la colonne le mardi 17 octobre à 10~h derrière les bataillons d'infanterie de la division (la suite de l'étape figure au f\textsuperscript{o}~38~v\textsuperscript{o}).}\par}
\newpage
\phantomsection\label{fac:76}
% --- En-tête du fac-similé (Page impaire / droite : Photographie du manuscrit) ---
\noindent\rule{\textwidth}{0.5pt}\par\vspace{0.55mm}
\noindent\textbf{\fontsize{6.8pt}{7.8pt}\selectfont FAC-SIMILÉ EN VIS-À-VIS — FOLIO 38 RECTO} \hfill \textbf{\fontsize{6.8pt}{7.8pt}\selectfont OCTOBRE 1916}\par\vspace{0.25mm}
\noindent{\fontsize{6.8pt}{8.0pt}\selectfont\textit{Livre-Journal d'Hippolyte Combes} — 16 (comptes) – 17 octobre 1916 (matin) (Beauval, Warloy-Baillon (Somme))}\par\vspace{0.55mm}
\noindent\rule{\textwidth}{0.35pt}\par\vspace*{\fill}
{\centering
  \setlength{\fboxsep}{1.2mm}%
  \setlength{\fboxrule}{0.4pt}%
  \fcolorbox{gray!55}{white}{%
    \includegraphics[width=0.95\textwidth,height=0.81\textheight,keepaspectratio]{\facsimiledir/20251126_104118.jpg}%
  }\par}
\vspace*{\fill}
\newpage

% =============================================================================
% VIS-A-VIS SPREAD : FOLIO 38 VERSO (20251126_104143.jpg)
% =============================================================================
\cleardoubleevenpage
\renewcommand{\currentmonthyear}{Octobre 1916}
\renewcommand{\currentfoliolabel}{Folio 38 verso}
\phantomsection\label{fol:77}
% --- En-tête du folio (Page paire / gauche : Transcription, apparat & registres de bas de page) ---
\noindent\rule{\textwidth}{0.5pt}\par\vspace{0.55mm}
\noindent\textbf{\footnotesize FOLIO 38 VERSO} \hfill \textbf{\footnotesize OCTOBRE 1916}\par\vspace{0.25mm}
\noindent{\fontsize{6.9pt}{8.0pt}\selectfont Bivouac dans un champ dévasté entre Bouzincourt et Hédauville}\par\vspace{0.2mm}
\noindent{\fontsize{6.9pt}{8.0pt}\selectfont\textit{17 octobre 1916 (suite)} \hfill \textit{Senlis-le-Sec, Bouzincourt, Hédauville (Somme)}}\par\vspace{0.5mm}
\noindent\rule{\textwidth}{0.35pt}\par\vspace{0.8mm}
% --- Corps de la transcription diplomatique et apparat marginal en prolongement ---
{\renewcommand{\arraystretch}{0.87}%
\noindent\hspace*{3.5mm}\begin{tabular}{@{}p{87mm}@{\hspace{3.5mm}}p{39mm}@{}}
      \leavevmode\fontsize{6.9pt}{7.9pt}\selectfont\strut Dépassé Senlis. Campé ds 1 champ dévasté, 1/2 route entre\strut & \smash{\parbox[t]{39mm}{\fontsize{6.2pt}{7.0pt}\selectfont\raggedright\color{black!88}\textbf{1 ds} \textit{abbr. pro} dans. \textbf{1 1/2} \textit{scil.} demi / moitié.}} \\
      \leavevmode\fontsize{6.9pt}{7.9pt}\selectfont\strut Bouzincourt \& Hédauville. Lieu laissé par ceux qui\strut &  \\
      \leavevmode\fontsize{6.9pt}{7.9pt}\selectfont\strut séjourné ici avant ns. Dû faire installation sommaire,\strut & \smash{\parbox[t]{39mm}{\fontsize{6.2pt}{7.0pt}\selectfont\raggedright\color{black!88}\textbf{3 ns} \textit{abbr. pro} nous.}} \\
      \leavevmode\fontsize{6.9pt}{7.9pt}\selectfont\strut ménagé sorte de hutte ouverte par devant pour abriter mes\strut &  \\
      \leavevmode\llap{\fontsize{6.4pt}{7.3pt}\selectfont\color{darkgray}5\hspace{1.8mm}}\fontsize{6.9pt}{7.9pt}\selectfont\strut sacs.\strut &  \\
      \leavevmode\fontsize{6.9pt}{7.9pt}\selectfont\strut Couché ds mess-tent, allongé s/ fauteuil-toile\strut & \smash{\parbox[t]{39mm}{\fontsize{6.2pt}{7.0pt}\selectfont\raggedright\color{black!88}\textbf{6 ds} \textit{abbr. pro} dans. \textbf{6 s/} \textit{abbr. pro} sur. — \textbf{7 Tte} \textit{abbr. pro} Toute.}} \\
      \leavevmode\fontsize{6.9pt}{7.9pt}\selectfont\strut Tte la nuit, plu à verse. Froid pénétrant\strut &  \\
      \leavevmode\fontsize{6.9pt}{7.9pt}\selectfont\strut Aucun confort —\strut &  \\
      \cline{1-1}
      \leavevmode\fontsize{6.9pt}{7.9pt}\selectfont\strut \underline{17 Mardi} \hfill{} Oct /16\strut &  \\
\end{tabular}}%
\par\vspace*{\fill}
% --- Registres de bas de page (1. Critique interne [si pertinent] ; 2. Loci paralleli en latin ; 3. Commentaire historique en retrait) ---
\noindent\rule{\textwidth}{0.5pt}\par\vspace{0.65mm}
{\fontsize{6.6pt}{7.7pt}\selectfont%
\noindent\normalfont\color{black!92} \textbf{1} (\textit{Senlis}) cf. supra f\textsuperscript{o}~37~v\textsuperscript{o} (p.~\pageref{fol:75}), infra f\textsuperscript{o}~43~v\textsuperscript{o} (p.~\pageref{fol:87}), f\textsuperscript{o}~48~r\textsuperscript{o} (p.~\pageref{fol:96}), \textit{et passim}. — \textbf{2} (\textit{Hédauville}) cf. supra f\textsuperscript{o}~22~v\textsuperscript{o} (p.~\pageref{fol:45}), infra f\textsuperscript{o}~54~r\textsuperscript{o} (p.~\pageref{fol:108}), f\textsuperscript{o}~60~v\textsuperscript{o} (p.~\pageref{fol:121}). — \textbf{2} (\textit{Bouzincourt}) cf. infra f\textsuperscript{o}~39~r\textsuperscript{o} (p.~\pageref{fol:78}), f\textsuperscript{o}~39~v\textsuperscript{o} (p.~\pageref{fol:79}), \textit{et passim}.\par\vspace{0.55mm}
\noindent\rule{38mm}{0.4pt}\par\vspace{0.5mm}
\noindent\hspace*{5mm}\parbox{\dimexpr\textwidth-5mm\relax}{\fontsize{6.6pt}{7.7pt}\selectfont\color{black!84}Suite de l'étape du mardi 17 octobre 1916: dépassant \textbf{Senlis-le-Sec}, la 152\textsuperscript{e} Compagnie A.S.C. s'établit en bivouac dans un champ défoncé à mi-chemin entre \textbf{Hédauville} et \textbf{Bouzincourt} (Somme, à 4~km au nord-ouest d'Albert), zone de déploiement logistique immédiat du II\textsuperscript{e} Corps britannique face à la crête de Thiepval.}\par}
\newpage
\phantomsection\label{fac:77}
% --- En-tête du fac-similé (Page impaire / droite : Photographie du manuscrit) ---
\noindent\rule{\textwidth}{0.5pt}\par\vspace{0.55mm}
\noindent\textbf{\fontsize{6.8pt}{7.8pt}\selectfont FAC-SIMILÉ EN VIS-À-VIS — FOLIO 38 VERSO} \hfill \textbf{\fontsize{6.8pt}{7.8pt}\selectfont OCTOBRE 1916}\par\vspace{0.25mm}
\noindent{\fontsize{6.8pt}{8.0pt}\selectfont\textit{Livre-Journal d'Hippolyte Combes} — 17 octobre 1916 (suite) (Senlis-le-Sec, Bouzincourt, Hédauville (Somme))}\par\vspace{0.55mm}
\noindent\rule{\textwidth}{0.35pt}\par\vspace*{\fill}
{\centering
  \setlength{\fboxsep}{1.2mm}%
  \setlength{\fboxrule}{0.4pt}%
  \fcolorbox{gray!55}{white}{%
    \includegraphics[width=0.95\textwidth,height=0.81\textheight,keepaspectratio]{\facsimiledir/20251126_104143.jpg}%
  }\par}
\vspace*{\fill}
\newpage

% =============================================================================
% VIS-A-VIS SPREAD : FOLIO 39 RECTO (20251126_104157.jpg)
% =============================================================================
\cleardoubleevenpage
\renewcommand{\currentmonthyear}{Octobre 1916}
\renewcommand{\currentfoliolabel}{Folio 39 recto}
\phantomsection\label{fol:78}
% --- En-tête du folio (Page paire / gauche : Transcription, apparat & registres de bas de page) ---
\noindent\rule{\textwidth}{0.5pt}\par\vspace{0.55mm}
\noindent\textbf{\footnotesize FOLIO 39 RECTO} \hfill \textbf{\footnotesize OCTOBRE 1916}\par\vspace{0.25mm}
\noindent{\fontsize{6.9pt}{8.0pt}\selectfont Visite d'Albert bombardée et de sa basilique en ruines}\par\vspace{0.2mm}
\noindent{\fontsize{6.9pt}{8.0pt}\selectfont\textit{18 octobre 1916} \hfill \textit{Camp de Bouzincourt, Albert (Somme)}}\par\vspace{0.5mm}
\noindent\rule{\textwidth}{0.35pt}\par\vspace{0.8mm}
% --- Corps de la transcription diplomatique et apparat marginal en prolongement ---
{\renewcommand{\arraystretch}{1.00}%
\noindent\hspace*{3.5mm}\begin{tabular}{@{}p{87mm}@{\hspace{3.5mm}}p{39mm}@{}}
      \leavevmode\fontsize{7.7pt}{9.0pt}\selectfont\strut \underline{Merc. 18} \hfill{} \underline{Oct /16}\strut & \smash{\parbox[t]{39mm}{\fontsize{6.2pt}{7.0pt}\selectfont\raggedright\color{black!88}\textbf{1 Merc.} \textit{scil.} Mercredi.}} \\
      \leavevmode\fontsize{7.7pt}{9.0pt}\selectfont\strut Dormi mess-tent — à l'autre extrémité,\strut & \smash{\parbox[t]{39mm}{\fontsize{6.2pt}{7.0pt}\selectfont\raggedright\color{black!88}\textbf{2 mess} \textit{angl. in ms.}}} \\
      \leavevmode\fontsize{7.7pt}{9.0pt}\selectfont\strut les 2 cooks. — Allongé sur 1 fauteuil de jardin!\strut &  \\
      \leavevmode\fontsize{7.7pt}{9.0pt}\selectfont\strut Nuit très froide —\strut &  \\
      \leavevmode\llap{\fontsize{6.4pt}{7.3pt}\selectfont\color{darkgray}5\hspace{1.8mm}}\fontsize{7.7pt}{9.0pt}\selectfont\strut Après breakfast, eu a lift s/ 1 lorry chargé\strut & \smash{\parbox[t]{39mm}{\fontsize{6.2pt}{7.0pt}\selectfont\raggedright\color{black!88}\textbf{5 s/} \textit{abbr. pro} sur. \textbf{5 breakfast} \textit{angl. in ms.} — \textbf{6 frse} \textit{abbr. pro} française. — \textbf{7 lgt} \textit{abbr. pro} longtemps. \textbf{7 passim} \textit{lat.} çà et là. — \textbf{9 qq} \textit{abbr. pro} quelque(s). — \textbf{10 s/} \textit{abbr. pro} sur. \textbf{10 lorry} \textit{angl. in ms.} — \textbf{11 de =} \textit{signum aequalitatis pro} de même. — \textbf{12 Tt} \textit{abbr. pro} Tout.}} \\
      \leavevmode\fontsize{7.7pt}{9.0pt}\selectfont\strut d'obus et venu à Albert. Mission frse.\strut &  \\
      \leavevmode\fontsize{7.7pt}{9.0pt}\selectfont\strut Puis erré passim. Arrêté lgt devant basilique\strut &  \\
      \leavevmode\fontsize{7.7pt}{9.0pt}\selectfont\strut ravagée par les obus\strut &  \\
      \leavevmode\fontsize{7.7pt}{9.0pt}\selectfont\strut Parcouru qq rues : maisons vides et abandonnées.\strut &  \\
      \leavevmode\llap{\fontsize{6.4pt}{7.3pt}\selectfont\color{darkgray}10\hspace{1.8mm}}\fontsize{7.7pt}{9.0pt}\selectfont\strut Après-midi, retourné à Albert s/ 1 lorry et revenu\strut &  \\
      \leavevmode\fontsize{7.7pt}{9.0pt}\selectfont\strut de = à notre camp.\strut &  \\
      \leavevmode\fontsize{7.7pt}{9.0pt}\selectfont\strut Tt le jour, ciel couvert. Boue épaisse.\strut &  \\
      \leavevmode\fontsize{7.7pt}{9.0pt}\selectfont\strut Fin journée, resté en dehors de mon abri. Assis sur\strut &  \\
      \leavevmode\fontsize{7.7pt}{9.0pt}\selectfont\strut 1 caisse, regardé à loisir camp et au-delà route\strut &  \\
\end{tabular}}%
\par\vspace*{\fill}
% --- Registres de bas de page (1. Critique interne [si pertinent] ; 2. Loci paralleli en latin ; 3. Commentaire historique en retrait) ---
\noindent\rule{\textwidth}{0.5pt}\par\vspace{0.65mm}
{\fontsize{6.6pt}{7.7pt}\selectfont%
\noindent\normalfont\color{black!92} \textbf{6} (\textit{Albert}) cf. infra f\textsuperscript{o}~39~v\textsuperscript{o} (p.~\pageref{fol:79}), f\textsuperscript{o}~40~r\textsuperscript{o} (p.~\pageref{fol:80}), \textit{et passim}.\par\vspace{0.55mm}
\noindent\rule{38mm}{0.4pt}\par\vspace{0.5mm}
\noindent\hspace*{5mm}\parbox{\dimexpr\textwidth-5mm\relax}{\fontsize{6.6pt}{7.7pt}\selectfont\color{black!84}Le mercredi 18 octobre 1916, profitant d'un camion de munitions depuis le bivouac de Bouzincourt, Hippolyte Combes découvre \textbf{Albert} (Somme) et s'arrête longuement devant les ruines de la basilique de pèlerinage \textbf{Notre-Dame de Brebières}, qu'il décrira plus en détail le 22 octobre (f\textsuperscript{o}~41~v\textsuperscript{o}).}\par}
\newpage
\phantomsection\label{fac:78}
% --- En-tête du fac-similé (Page impaire / droite : Photographie du manuscrit) ---
\noindent\rule{\textwidth}{0.5pt}\par\vspace{0.55mm}
\noindent\textbf{\fontsize{6.8pt}{7.8pt}\selectfont FAC-SIMILÉ EN VIS-À-VIS — FOLIO 39 RECTO} \hfill \textbf{\fontsize{6.8pt}{7.8pt}\selectfont OCTOBRE 1916}\par\vspace{0.25mm}
\noindent{\fontsize{6.8pt}{8.0pt}\selectfont\textit{Livre-Journal d'Hippolyte Combes} — 18 octobre 1916 (Camp de Bouzincourt, Albert (Somme))}\par\vspace{0.55mm}
\noindent\rule{\textwidth}{0.35pt}\par\vspace*{\fill}
{\centering
  \setlength{\fboxsep}{1.2mm}%
  \setlength{\fboxrule}{0.4pt}%
  \fcolorbox{gray!55}{white}{%
    \includegraphics[width=0.95\textwidth,height=0.81\textheight,keepaspectratio]{\facsimiledir/20251126_104157.jpg}%
  }\par}
\vspace*{\fill}
\newpage

% =============================================================================
% VIS-A-VIS SPREAD : FOLIO 39 VERSO (20251126_104222.jpg)
% =============================================================================
\cleardoubleevenpage
\renewcommand{\currentmonthyear}{Octobre 1916}
\renewcommand{\currentfoliolabel}{Folio 39 verso}
\phantomsection\label{fol:79}
% --- En-tête du folio (Page paire / gauche : Transcription, apparat & registres de bas de page) ---
\noindent\rule{\textwidth}{0.5pt}\par\vspace{0.55mm}
\noindent\textbf{\footnotesize FOLIO 39 VERSO} \hfill \textbf{\footnotesize OCTOBRE 1916}\par\vspace{0.25mm}
\noindent{\fontsize{6.9pt}{8.0pt}\selectfont Le flot ininterrompu des convois vers Albert et la solde à la Mission}\par\vspace{0.2mm}
\noindent{\fontsize{6.9pt}{8.0pt}\selectfont\textit{18 octobre 1916 (suite)} \hfill \textit{Camp de Bouzincourt, Albert (Somme)}}\par\vspace{0.5mm}
\noindent\rule{\textwidth}{0.35pt}\par\vspace{0.8mm}
% --- Corps de la transcription diplomatique et apparat marginal en prolongement ---
{\renewcommand{\arraystretch}{0.87}%
\noindent\hspace*{3.5mm}\begin{tabular}{@{}p{87mm}@{\hspace{3.5mm}}p{39mm}@{}}
      \leavevmode\fontsize{6.9pt}{7.9pt}\selectfont\strut sur l., nuit \& jour, flot ininterrompu de véhicules de\strut &  \\
      \leavevmode\fontsize{6.9pt}{7.9pt}\selectfont\strut ttes sortes : lorries, caissons de munitions etc. Autant de\strut & \smash{\parbox[t]{39mm}{\fontsize{6.2pt}{7.0pt}\selectfont\raggedright\color{black!88}\textbf{2 ttes} \textit{abbr. pro} toutes. \textbf{2 lorries} \textit{angl. in ms.} — \textbf{3 gde} \textit{abbr. pro} grande. — \textbf{4 off.} \textit{abbr. pro} officier(s). \textbf{4 dinner} \textit{angl. in ms.} — \textbf{5 ds} \textit{abbr. pro} dans. \textbf{5 mess} \textit{angl. in ms.} — \textbf{6 Merc.} \textit{scil.} Mercredi. \textbf{6 cartes ples} \textit{abbr. pro} cartes postales. — \textbf{7 frse} \textit{abbr. pro} française. \textbf{7 fr.} \textit{abbr. pro} franc(s).}} \\
      \leavevmode\fontsize{6.9pt}{7.9pt}\selectfont\strut mouvement que le long d'1 gde artère de capitale\strut &  \\
      \leavevmode\fontsize{6.9pt}{7.9pt}\selectfont\strut Quand, après dinner, les off. partis, enfin seul\strut &  \\
      \leavevmode\llap{\fontsize{6.4pt}{7.3pt}\selectfont\color{darkgray}5\hspace{1.8mm}}\fontsize{6.9pt}{7.9pt}\selectfont\strut ds la mess-tent —\strut &  \\
      \cline{1-1}
      \leavevmode\fontsize{6.9pt}{7.9pt}\selectfont\strut \underline{Merc. 18} | Albert : cartes ples : 40 c \hfill{} Oct /16\strut &  \\
      \leavevmode\fontsize{6.9pt}{7.9pt}\selectfont\strut | Albert, mission frse : 84 fr. —\strut &  \\
\end{tabular}}%
\par\vspace*{\fill}
% --- Registres de bas de page (1. Critique interne [si pertinent] ; 2. Loci paralleli en latin ; 3. Commentaire historique en retrait) ---
\noindent\rule{\textwidth}{0.5pt}\par\vspace{0.65mm}
{\fontsize{6.6pt}{7.7pt}\selectfont%
\noindent\normalfont\color{black!92} \textbf{6} (\textit{Albert}) cf. supra f\textsuperscript{o}~39~r\textsuperscript{o} (p.~\pageref{fol:78}), infra f\textsuperscript{o}~40~r\textsuperscript{o} (p.~\pageref{fol:80}), f\textsuperscript{o}~40~v\textsuperscript{o} (p.~\pageref{fol:81}), \textit{et passim}. — \textbf{7} (\textit{84 f.}) cf. supra f\textsuperscript{o}~2~r\textsuperscript{o} (p.~\pageref{fol:4}), f\textsuperscript{o}~10~v\textsuperscript{o} (p.~\pageref{fol:21}), infra f\textsuperscript{o}~56~v\textsuperscript{o} (p.~\pageref{fol:113}), f\textsuperscript{o}~62~v\textsuperscript{o} (p.~\pageref{fol:125}).\par\vspace{0.55mm}
\noindent\rule{38mm}{0.4pt}\par\vspace{0.5mm}
\noindent\hspace*{5mm}\parbox{\dimexpr\textwidth-5mm\relax}{\fontsize{6.6pt}{7.7pt}\selectfont\color{black!84}Suite immédiate du f\textsuperscript{o}~39~r\textsuperscript{o} (18 octobre 1916): observation de l'intensité du trafic logistique britannique sur la chaussée de Bouzincourt à \textbf{Albert} (camions automobiles et caissons de munitions ravitaillant le II\textsuperscript{e} Corps en vue de l'attaque du 21 octobre sur la tranchée Regina), et perception de 84~fr. de solde au bureau de la Mission française d'Albert.}\par}
\newpage
\phantomsection\label{fac:79}
% --- En-tête du fac-similé (Page impaire / droite : Photographie du manuscrit) ---
\noindent\rule{\textwidth}{0.5pt}\par\vspace{0.55mm}
\noindent\textbf{\fontsize{6.8pt}{7.8pt}\selectfont FAC-SIMILÉ EN VIS-À-VIS — FOLIO 39 VERSO} \hfill \textbf{\fontsize{6.8pt}{7.8pt}\selectfont OCTOBRE 1916}\par\vspace{0.25mm}
\noindent{\fontsize{6.8pt}{8.0pt}\selectfont\textit{Livre-Journal d'Hippolyte Combes} — 18 octobre 1916 (suite) (Camp de Bouzincourt, Albert (Somme))}\par\vspace{0.55mm}
\noindent\rule{\textwidth}{0.35pt}\par\vspace*{\fill}
{\centering
  \setlength{\fboxsep}{1.2mm}%
  \setlength{\fboxrule}{0.4pt}%
  \fcolorbox{gray!55}{white}{%
    \includegraphics[width=0.95\textwidth,height=0.81\textheight,keepaspectratio]{\facsimiledir/20251126_104222.jpg}%
  }\par}
\vspace*{\fill}
\newpage

% =============================================================================
% VIS-A-VIS SPREAD : FOLIO 40 RECTO (20251126_104235.jpg)
% =============================================================================
\cleardoubleevenpage
\renewcommand{\currentmonthyear}{Octobre 1916}
\renewcommand{\currentfoliolabel}{Folio 40 recto}
\phantomsection\label{fol:80}
% --- En-tête du folio (Page paire / gauche : Transcription, apparat & registres de bas de page) ---
\noindent\rule{\textwidth}{0.5pt}\par\vspace{0.55mm}
\noindent\textbf{\footnotesize FOLIO 40 RECTO} \hfill \textbf{\footnotesize OCTOBRE 1916}\par\vspace{0.25mm}
\noindent{\fontsize{6.9pt}{8.0pt}\selectfont Pluie torrentielle et installation au Brickfield Area d'Albert}\par\vspace{0.2mm}
\noindent{\fontsize{6.9pt}{8.0pt}\selectfont\textit{19 octobre 1916} \hfill \textit{Camp de Bouzincourt, Albert (Brickfield Area, Somme)}}\par\vspace{0.5mm}
\noindent\rule{\textwidth}{0.35pt}\par\vspace{0.8mm}
% --- Corps de la transcription diplomatique et apparat marginal en prolongement ---
{\renewcommand{\arraystretch}{0.93}%
\noindent\hspace*{3.5mm}\begin{tabular}{@{}p{87mm}@{\hspace{3.5mm}}p{39mm}@{}}
      \leavevmode\fontsize{7.3pt}{8.4pt}\selectfont\strut \underline{Jeudi 19} \hfill{} \underline{Oct /16}\strut &  \\
      \leavevmode\fontsize{7.3pt}{8.4pt}\selectfont\strut Levé 6h\strut &  \\
      \leavevmode\fontsize{7.3pt}{8.4pt}\selectfont\strut Réveillé par orage\strut &  \\
      \leavevmode\fontsize{7.3pt}{8.4pt}\selectfont\strut Tte nuit, pluie torrentielle. Durant tte\strut & \smash{\parbox[t]{39mm}{\fontsize{6.2pt}{7.0pt}\selectfont\raggedright\color{black!88}\textbf{4 tte} \textit{abbr. pro} toute.}} \\
      \leavevmode\llap{\fontsize{6.4pt}{7.3pt}\selectfont\color{darkgray}5\hspace{1.8mm}}\fontsize{7.3pt}{8.4pt}\selectfont\strut matinée, averses n'ont pas cessé. Sol\strut &  \\
      \leavevmode\fontsize{7.3pt}{8.4pt}\selectfont\strut détrempé, impraticable\strut &  \\
      \leavevmode\fontsize{7.3pt}{8.4pt}\selectfont\strut Abrité sous tente verte du mess,\strut & \smash{\parbox[t]{39mm}{\fontsize{6.2pt}{7.0pt}\selectfont\raggedright\color{black!88}\textbf{7 mess} \textit{angl. in ms.}}} \\
      \leavevmode\fontsize{7.3pt}{8.4pt}\selectfont\strut attendu départ\strut &  \\
      \leavevmode\fontsize{7.3pt}{8.4pt}\selectfont\strut Tout ruisselant d'eau.\strut &  \\
      \leavevmode\llap{\fontsize{6.4pt}{7.3pt}\selectfont\color{darkgray}10\hspace{1.8mm}}\fontsize{7.3pt}{8.4pt}\selectfont\strut Quitté camp vers 11h. Monté s/ 1 wagon.\strut & \smash{\parbox[t]{39mm}{\fontsize{6.2pt}{7.0pt}\selectfont\raggedright\color{black!88}\textbf{10 s/} \textit{abbr. pro} sur.}} \\
      \leavevmode\fontsize{7.3pt}{8.4pt}\selectfont\strut Vers 1h, arrivé au brickfield area,\strut &  \\
      \leavevmode\fontsize{7.3pt}{8.4pt}\selectfont\strut où ns devons bivouaquer\strut & \smash{\parbox[t]{39mm}{\fontsize{6.2pt}{7.0pt}\selectfont\raggedright\color{black!88}\textbf{12 ns} \textit{abbr. pro} nous.}} \\
      \leavevmode\fontsize{7.3pt}{8.4pt}\selectfont\strut Tentes et abris, dressés ds la boue et\strut & \smash{\parbox[t]{39mm}{\fontsize{6.2pt}{7.0pt}\selectfont\raggedright\color{black!88}\textbf{13 ds} \textit{abbr. pro} dans.}} \\
      \leavevmode\fontsize{7.3pt}{8.4pt}\selectfont\strut sous pluie battante.\strut &  \\
      \leavevmode\llap{\fontsize{6.4pt}{7.3pt}\selectfont\color{darkgray}15\hspace{1.8mm}}\fontsize{7.3pt}{8.4pt}\selectfont\strut Réussi par trouver petite chambre ds\strut & \smash{\parbox[t]{39mm}{\fontsize{6.2pt}{7.0pt}\selectfont\raggedright\color{black!88}\textbf{15 ds} \textit{abbr. pro} dans.}} \\
      \leavevmode\fontsize{7.3pt}{8.4pt}\selectfont\strut maison rue voisine, à qq 3 ou 400 mètres\strut & \smash{\parbox[t]{39mm}{\fontsize{6.2pt}{7.0pt}\selectfont\raggedright\color{black!88}\textbf{16 qq} \textit{abbr. pro} quelque(s).}} \\
      \leavevmode\fontsize{7.3pt}{8.4pt}\selectfont\strut du camp. Y installé!\strut &  \\
      \leavevmode\fontsize{7.3pt}{8.4pt}\selectfont\strut Toute nuit, bruit ininterrompu\strut &  \\
      \leavevmode\fontsize{7.3pt}{8.4pt}\selectfont\strut très forte canonnade\strut &  \\
      \cline{1-1}
      \leavevmode\llap{\fontsize{6.4pt}{7.3pt}\selectfont\color{darkgray}20\hspace{1.8mm}}\fontsize{7.3pt}{8.4pt}\selectfont\strut [ — Cartes ples d'Albert : 20 c ]\strut & \smash{\parbox[t]{39mm}{\fontsize{6.2pt}{7.0pt}\selectfont\raggedright\color{black!88}\textbf{20} \textit{in marg. cum unco} (comptes). \textbf{20 Cartes ples} \textit{abbr. pro} Cartes postales.}} \\
\end{tabular}}%
\par\vspace*{\fill}
% --- Registres de bas de page (1. Critique interne [si pertinent] ; 2. Loci paralleli en latin ; 3. Commentaire historique en retrait) ---
\noindent\rule{\textwidth}{0.5pt}\par\vspace{0.65mm}
{\fontsize{6.6pt}{7.7pt}\selectfont%
\noindent\normalfont\color{black!92} \textbf{20} (\textit{Albert}) cf. supra f\textsuperscript{o}~39~r\textsuperscript{o} (p.~\pageref{fol:78}), f\textsuperscript{o}~39~v\textsuperscript{o} (p.~\pageref{fol:79}), infra f\textsuperscript{o}~40~v\textsuperscript{o} (p.~\pageref{fol:81}), f\textsuperscript{o}~41~r\textsuperscript{o} (p.~\pageref{fol:82}), \textit{et passim}.\par\vspace{0.55mm}
\noindent\rule{38mm}{0.4pt}\par\vspace{0.5mm}
\noindent\hspace*{5mm}\parbox{\dimexpr\textwidth-5mm\relax}{\fontsize{6.6pt}{7.7pt}\selectfont\color{black!84}Le jeudi 19 octobre 1916, la compagnie quitte le champ de Bouzincourt pour le grand bivouac boueux du \textbf{\textit{Brickfields Area}} (secteur des briqueteries à l'entrée nord-ouest d'\textbf{Albert}), tandis qu'Hippolyte Combes parvient à se loger à 300 ou 400~mètres du camp chez la famille \textbf{Duchâteau} (nommée aux f\textsuperscript{o}~41~r\textsuperscript{o}, f\textsuperscript{o}~43~v\textsuperscript{o} et f\textsuperscript{o}~45~v\textsuperscript{o}), sous le roulement continu de la préparation d'artillerie britannique.}\par}
\newpage
\phantomsection\label{fac:80}
% --- En-tête du fac-similé (Page impaire / droite : Photographie du manuscrit) ---
\noindent\rule{\textwidth}{0.5pt}\par\vspace{0.55mm}
\noindent\textbf{\fontsize{6.8pt}{7.8pt}\selectfont FAC-SIMILÉ EN VIS-À-VIS — FOLIO 40 RECTO} \hfill \textbf{\fontsize{6.8pt}{7.8pt}\selectfont OCTOBRE 1916}\par\vspace{0.25mm}
\noindent{\fontsize{6.8pt}{8.0pt}\selectfont\textit{Livre-Journal d'Hippolyte Combes} — 19 octobre 1916 (Camp de Bouzincourt, Albert (Brickfield Area, Somme))}\par\vspace{0.55mm}
\noindent\rule{\textwidth}{0.35pt}\par\vspace*{\fill}
{\centering
  \setlength{\fboxsep}{1.2mm}%
  \setlength{\fboxrule}{0.4pt}%
  \fcolorbox{gray!55}{white}{%
    \includegraphics[width=0.95\textwidth,height=0.81\textheight,keepaspectratio]{\facsimiledir/20251126_104235.jpg}%
  }\par}
\vspace*{\fill}
\newpage

% =============================================================================
% VIS-A-VIS SPREAD : FOLIO 40 VERSO (20251126_104314.jpg)
% =============================================================================
\cleardoubleevenpage
\renewcommand{\currentmonthyear}{Octobre 1916}
\renewcommand{\currentfoliolabel}{Folio 40 verso}
\phantomsection\label{fol:81}
% --- En-tête du folio (Page paire / gauche : Transcription, apparat & registres de bas de page) ---
\noindent\rule{\textwidth}{0.5pt}\par\vspace{0.55mm}
\noindent\textbf{\footnotesize FOLIO 40 VERSO} \hfill \textbf{\footnotesize OCTOBRE 1916}\par\vspace{0.25mm}
\noindent{\fontsize{6.9pt}{8.0pt}\selectfont Duel d'artillerie formidable et horizon embrasé au-dessus d'Albert}\par\vspace{0.2mm}
\noindent{\fontsize{6.9pt}{8.0pt}\selectfont\textit{20 octobre 1916} \hfill \textit{Albert (Brickfield Area, Somme)}}\par\vspace{0.5mm}
\noindent\rule{\textwidth}{0.35pt}\par\vspace{0.8mm}
% --- Corps de la transcription diplomatique et apparat marginal en prolongement ---
{\renewcommand{\arraystretch}{0.93}%
\noindent\hspace*{3.5mm}\begin{tabular}{@{}p{87mm}@{\hspace{3.5mm}}p{39mm}@{}}
      \leavevmode\fontsize{7.3pt}{8.4pt}\selectfont\strut \underline{Vendredi 20} \hfill{} \underline{Oct /16}\strut &  \\
      \leavevmode\fontsize{7.3pt}{8.4pt}\selectfont\strut Enfin, soleil brille 1 peu, mais bise\strut &  \\
      \leavevmode\fontsize{7.3pt}{8.4pt}\selectfont\strut glaciale. Boue tjrs épaisse et liquide\strut & \smash{\parbox[t]{39mm}{\fontsize{6.2pt}{7.0pt}\selectfont\raggedright\color{black!88}\textbf{3 tjrs} \textit{abbr. pro} toujours.}} \\
      \leavevmode\fontsize{7.3pt}{8.4pt}\selectfont\strut Tt jour, bombardemt terrifiant à qq\strut & \smash{\parbox[t]{39mm}{\fontsize{6.2pt}{7.0pt}\selectfont\raggedright\color{black!88}\textbf{4 Tt} \textit{abbr. pro} Tout. \textbf{4 qq} \textit{abbr. pro} quelque(s).}} \\
      \leavevmode\llap{\fontsize{6.4pt}{7.3pt}\selectfont\color{darkgray}5\hspace{1.8mm}}\fontsize{7.3pt}{8.4pt}\selectfont\strut kilomètres. Duel d'artillerie formidable.\strut &  \\
      \leavevmode\fontsize{7.3pt}{8.4pt}\selectfont\strut Vitres et portes tremblent continuellement.\strut &  \\
      \leavevmode\fontsize{7.3pt}{8.4pt}\selectfont\strut Grondement canons finit par énerver.\strut &  \\
      \leavevmode\fontsize{7.3pt}{8.4pt}\selectfont\strut Sommeil, impossible\strut &  \\
      \leavevmode\fontsize{7.3pt}{8.4pt}\selectfont\strut Monté pour mes repas plateau voisin\strut &  \\
      \leavevmode\llap{\fontsize{6.4pt}{7.3pt}\selectfont\color{darkgray}10\hspace{1.8mm}}\fontsize{7.3pt}{8.4pt}\selectfont\strut où mess-room ds 1 petite cahute faite\strut & \smash{\parbox[t]{39mm}{\fontsize{6.2pt}{7.0pt}\selectfont\raggedright\color{black!88}\textbf{10 ds} \textit{abbr. pro} dans. \textbf{10 mess} \textit{angl. in ms.}}} \\
      \leavevmode\fontsize{7.3pt}{8.4pt}\selectfont\strut avec des caisses vides de munitions.\strut &  \\
      \leavevmode\fontsize{7.3pt}{8.4pt}\selectfont\strut Me suis vite organisé ds ma chambrette.\strut & \smash{\parbox[t]{39mm}{\fontsize{6.2pt}{7.0pt}\selectfont\raggedright\color{black!88}\textbf{12 ds} \textit{abbr. pro} dans.}} \\
      \leavevmode\fontsize{7.3pt}{8.4pt}\selectfont\strut Temps froid pénétrant.\strut &  \\
      \leavevmode\fontsize{7.3pt}{8.4pt}\selectfont\strut Nuit, du haut colline où notre\strut &  \\
      \leavevmode\llap{\fontsize{6.4pt}{7.3pt}\selectfont\color{darkgray}15\hspace{1.8mm}}\fontsize{7.3pt}{8.4pt}\selectfont\strut Coy bivouaque, très impressionnant spectacle\strut & \smash{\parbox[t]{39mm}{\fontsize{6.2pt}{7.0pt}\selectfont\raggedright\color{black!88}\textbf{15 Coy} \textit{angl.} Company.}} \\
      \leavevmode\fontsize{7.3pt}{8.4pt}\selectfont\strut à g. : horizon tout illuminé par\strut &  \\
      \leavevmode\fontsize{7.3pt}{8.4pt}\selectfont\strut les éclairs d'innombrables bouches à feu.\strut &  \\
      \leavevmode\fontsize{7.3pt}{8.4pt}\selectfont\strut Canonnade ne cesse pas et, par\strut &  \\
      \leavevmode\fontsize{7.3pt}{8.4pt}\selectfont\strut moments, d'1 intensité et violence\strut &  \\
      \leavevmode\llap{\fontsize{6.4pt}{7.3pt}\selectfont\color{darkgray}20\hspace{1.8mm}}\fontsize{7.3pt}{8.4pt}\selectfont\strut extrême.\strut &  \\
\end{tabular}}%
\par\vspace*{\fill}
% --- Registres de bas de page (1. Critique interne [si pertinent] ; 2. Loci paralleli en latin ; 3. Commentaire historique en retrait) ---
\noindent\rule{\textwidth}{0.5pt}\par\vspace{0.65mm}
{\fontsize{6.6pt}{7.7pt}\selectfont%
\noindent\normalfont\color{black!92} \textbf{1} (\textit{de eodem commorationis loco}) cf. supra f\textsuperscript{o}~40~r\textsuperscript{o} (p.~\pageref{fol:80}).\par\vspace{0.55mm}
\noindent\rule{38mm}{0.4pt}\par\vspace{0.5mm}
\noindent\hspace*{5mm}\parbox{\dimexpr\textwidth-5mm\relax}{\fontsize{6.6pt}{7.7pt}\selectfont\color{black!84}Ce « bombardement terrifiant » observé le vendredi \textbf{20 octobre 1916} depuis la colline des briqueteries d'Albert correspond à la préparation d'artillerie finale précédant l'attaque générale lancée le lendemain \textbf{21 octobre 1916} à 12~h~06 par le \textbf{II\textsuperscript{e} Corps britannique} — dont fait partie la \textbf{18\textsuperscript{e} Division} aux côtés des 25\textsuperscript{e} et 39\textsuperscript{e} Divisions et de la 4\textsuperscript{e} Division canadienne — contre \textbf{\textit{Regina Trench}} et \textbf{\textit{Stuff Trench}} lors de la bataille des hauteurs de l'Ancre (Nichols, \textit{The 18th Division in the Great War}, 1922, chap.~VIII, p.~124--132).}\par}
\newpage
\phantomsection\label{fac:81}
% --- En-tête du fac-similé (Page impaire / droite : Photographie du manuscrit) ---
\noindent\rule{\textwidth}{0.5pt}\par\vspace{0.55mm}
\noindent\textbf{\fontsize{6.8pt}{7.8pt}\selectfont FAC-SIMILÉ EN VIS-À-VIS — FOLIO 40 VERSO} \hfill \textbf{\fontsize{6.8pt}{7.8pt}\selectfont OCTOBRE 1916}\par\vspace{0.25mm}
\noindent{\fontsize{6.8pt}{8.0pt}\selectfont\textit{Livre-Journal d'Hippolyte Combes} — 20 octobre 1916 (Albert (Brickfield Area, Somme))}\par\vspace{0.55mm}
\noindent\rule{\textwidth}{0.35pt}\par\vspace*{\fill}
{\centering
  \setlength{\fboxsep}{1.2mm}%
  \setlength{\fboxrule}{0.4pt}%
  \fcolorbox{gray!55}{white}{%
    \includegraphics[width=0.95\textwidth,height=0.81\textheight,keepaspectratio]{\facsimiledir/20251126_104314.jpg}%
  }\par}
\vspace*{\fill}
\newpage

% =============================================================================
% VIS-A-VIS SPREAD : FOLIO 41 RECTO (20251126_104331.jpg)
% =============================================================================
\cleardoubleevenpage
\renewcommand{\currentmonthyear}{Octobre 1916}
\renewcommand{\currentfoliolabel}{Folio 41 recto}
\phantomsection\label{fol:82}
% --- En-tête du folio (Page paire / gauche : Transcription, apparat & registres de bas de page) ---
\noindent\rule{\textwidth}{0.5pt}\par\vspace{0.55mm}
\noindent\textbf{\footnotesize FOLIO 41 RECTO} \hfill \textbf{\footnotesize OCTOBRE 1916}\par\vspace{0.25mm}
\noindent{\fontsize{6.9pt}{8.0pt}\selectfont Gelée d'octobre, reconnaissance vers Aveluy et veillée près du poêle}\par\vspace{0.2mm}
\noindent{\fontsize{6.9pt}{8.0pt}\selectfont\textit{21 octobre 1916} \hfill \textit{Albert, Aveluy (Somme)}}\par\vspace{0.5mm}
\noindent\rule{\textwidth}{0.35pt}\par\vspace{0.8mm}
% --- Corps de la transcription diplomatique et apparat marginal en prolongement ---
{\renewcommand{\arraystretch}{1.00}%
\noindent\hspace*{3.5mm}\begin{tabular}{@{}p{87mm}@{\hspace{3.5mm}}p{39mm}@{}}
      \leavevmode\fontsize{7.7pt}{9.0pt}\selectfont\strut \underline{Samedi 21} \hfill{} \underline{Oct /16}\strut &  \\
      \leavevmode\fontsize{7.7pt}{9.0pt}\selectfont\strut Nuit très froide. Gelé fortement.\strut &  \\
      \leavevmode\fontsize{7.7pt}{9.0pt}\selectfont\strut Ds la journée, soleil réchauffe 1 peu température.\strut & \smash{\parbox[t]{39mm}{\fontsize{6.2pt}{7.0pt}\selectfont\raggedright\color{black!88}\textbf{3 Ds} \textit{abbr. pro} Dans.}} \\
      \leavevmode\fontsize{7.7pt}{9.0pt}\selectfont\strut Matinée, promenade de reconnaissance en ville\strut &  \\
      \leavevmode\llap{\fontsize{6.4pt}{7.3pt}\selectfont\color{darkgray}5\hspace{1.8mm}}\fontsize{7.7pt}{9.0pt}\selectfont\strut le long rues dont les maisons atteintes par\strut &  \\
      \leavevmode\fontsize{7.7pt}{9.0pt}\selectfont\strut obus. Venu à l'Ordonnance, Aveluy.\strut &  \\
      \leavevmode\fontsize{7.7pt}{9.0pt}\selectfont\strut Tout l'après-midi, ds ma chambre. Achevé\strut & \smash{\parbox[t]{39mm}{\fontsize{6.2pt}{7.0pt}\selectfont\raggedright\color{black!88}\textbf{7 ds} \textit{abbr. pro} dans.}} \\
      \leavevmode\fontsize{7.7pt}{9.0pt}\selectfont\strut de m'installer pour le mieux.\strut &  \\
      \leavevmode\fontsize{7.7pt}{9.0pt}\selectfont\strut Après lunch, me suis attardé sur le\strut & \smash{\parbox[t]{39mm}{\fontsize{6.2pt}{7.0pt}\selectfont\raggedright\color{black!88}\textbf{9 lunch} \textit{angl. in ms.}}} \\
      \leavevmode\llap{\fontsize{6.4pt}{7.3pt}\selectfont\color{darkgray}10\hspace{1.8mm}}\fontsize{7.7pt}{9.0pt}\selectfont\strut plateau où notre camp dressé : toute la\strut &  \\
      \leavevmode\fontsize{7.7pt}{9.0pt}\selectfont\strut colline, couverte de troupes au bivouac.\strut &  \\
      \leavevmode\fontsize{7.7pt}{9.0pt}\selectfont\strut Très beau soleil, mais froid piquant.\strut &  \\
      \leavevmode\fontsize{7.7pt}{9.0pt}\selectfont\strut Pas ressorti, quand rentré home après\strut & \smash{\parbox[t]{39mm}{\fontsize{6.2pt}{7.0pt}\selectfont\raggedright\color{black!88}\textbf{13 home} \textit{angl. in ms.}}} \\
      \leavevmode\fontsize{7.7pt}{9.0pt}\selectfont\strut le thé. Lu ds ma chambre, puis resté\strut & \smash{\parbox[t]{39mm}{\fontsize{6.2pt}{7.0pt}\selectfont\raggedright\color{black!88}\textbf{14 ds} \textit{abbr. pro} dans.}} \\
      \leavevmode\llap{\fontsize{6.4pt}{7.3pt}\selectfont\color{darkgray}15\hspace{1.8mm}}\fontsize{7.7pt}{9.0pt}\selectfont\strut lgts devant cuisinière, en bas, en cie\strut & \smash{\parbox[t]{39mm}{\fontsize{6.2pt}{7.0pt}\selectfont\raggedright\color{black!88}\textbf{15 lgts} \textit{abbr. pro} longtemps. \textbf{15 cie} \textit{abbr. pro} compagnie. — \textbf{17 tjrs} \textit{abbr. pro} toujours.}} \\
      \leavevmode\fontsize{7.7pt}{9.0pt}\selectfont\strut du + jeune fils maison, Abel...\strut &  \\
      \leavevmode\fontsize{7.7pt}{9.0pt}\selectfont\strut Canonnade tjrs très violente —\strut &  \\
\end{tabular}}%
\par\vspace*{\fill}
% --- Registres de bas de page (1. Critique interne [si pertinent] ; 2. Loci paralleli en latin ; 3. Commentaire historique en retrait) ---
\noindent\rule{\textwidth}{0.5pt}\par\vspace{0.65mm}
{\fontsize{6.6pt}{7.7pt}\selectfont%
\noindent\normalfont\color{black!92} \textbf{6} (\textit{Aveluy}) cf. infra f\textsuperscript{o}~75~r\textsuperscript{o} (p.~\pageref{fol:150}), f\textsuperscript{o}~75~v\textsuperscript{o} (p.~\pageref{fol:151}).\par\vspace{0.55mm}
\noindent\rule{38mm}{0.4pt}\par\vspace{0.5mm}
\noindent\hspace*{5mm}\parbox{\dimexpr\textwidth-5mm\relax}{\fontsize{6.6pt}{7.7pt}\selectfont\color{black!84}Le samedi \textbf{21 octobre 1916} — jour même où les 53\textsuperscript{e} et 55\textsuperscript{e} Brigades de la 18\textsuperscript{e} Division enlèvent d'assaut la partie orientale de \textit{Regina Trench} au nord de Courcelette et de Thiepval —, Hippolyte Combes se rend au dépôt du matériel (\textit{Army Ordnance Corps}) à \textbf{Aveluy} (Somme, à 2~km au nord d'Albert dans la vallée de l'Ancre), puis passe la soirée près du poêle avec le jeune \textbf{Abel Duchâteau}, fils cadet de son logeur d'Albert.}\par}
\newpage
\phantomsection\label{fac:82}
% --- En-tête du fac-similé (Page impaire / droite : Photographie du manuscrit) ---
\noindent\rule{\textwidth}{0.5pt}\par\vspace{0.55mm}
\noindent\textbf{\fontsize{6.8pt}{7.8pt}\selectfont FAC-SIMILÉ EN VIS-À-VIS — FOLIO 41 RECTO} \hfill \textbf{\fontsize{6.8pt}{7.8pt}\selectfont OCTOBRE 1916}\par\vspace{0.25mm}
\noindent{\fontsize{6.8pt}{8.0pt}\selectfont\textit{Livre-Journal d'Hippolyte Combes} — 21 octobre 1916 (Albert, Aveluy (Somme))}\par\vspace{0.55mm}
\noindent\rule{\textwidth}{0.35pt}\par\vspace*{\fill}
{\centering
  \setlength{\fboxsep}{1.2mm}%
  \setlength{\fboxrule}{0.4pt}%
  \fcolorbox{gray!55}{white}{%
    \includegraphics[width=0.95\textwidth,height=0.81\textheight,keepaspectratio]{\facsimiledir/20251126_104331.jpg}%
  }\par}
\vspace*{\fill}
\newpage

% =============================================================================
% VIS-A-VIS SPREAD : FOLIO 41 VERSO (20251126_104346.jpg)
% =============================================================================
\cleardoubleevenpage
\renewcommand{\currentmonthyear}{Octobre 1916}
\renewcommand{\currentfoliolabel}{Folio 41 verso}
\phantomsection\label{fol:83}
% --- En-tête du folio (Page paire / gauche : Transcription, apparat & registres de bas de page) ---
\noindent\rule{\textwidth}{0.5pt}\par\vspace{0.55mm}
\noindent\textbf{\footnotesize FOLIO 41 VERSO} \hfill \textbf{\footnotesize OCTOBRE 1916}\par\vspace{0.25mm}
\noindent{\fontsize{6.9pt}{8.0pt}\selectfont La « Vierge penchée » d'Albert et la rumeur de rappel des jeunes R.A.T.}\par\vspace{0.2mm}
\noindent{\fontsize{6.9pt}{8.0pt}\selectfont\textit{22 octobre 1916} \hfill \textit{Albert (Somme)}}\par\vspace{0.5mm}
\noindent\rule{\textwidth}{0.35pt}\par\vspace{0.8mm}
% --- Corps de la transcription diplomatique et apparat marginal en prolongement ---
{\renewcommand{\arraystretch}{0.87}%
\noindent\hspace*{3.5mm}\begin{tabular}{@{}p{87mm}@{\hspace{3.5mm}}p{39mm}@{}}
      \leavevmode\fontsize{6.9pt}{7.9pt}\selectfont\strut \underline{Dimanche 22.} \hfill{} \underline{Oct /16}\strut &  \\
      \leavevmode\fontsize{6.9pt}{7.9pt}\selectfont\strut Encore cette nuit, gelé très dur.\strut &  \\
      \leavevmode\fontsize{6.9pt}{7.9pt}\selectfont\strut Canonnade, pas cessé un instant.\strut &  \\
      \leavevmode\fontsize{6.9pt}{7.9pt}\selectfont\strut 1/2 jour, très beau soleil. Journée superbe.\strut & \smash{\parbox[t]{39mm}{\fontsize{6.2pt}{7.0pt}\selectfont\raggedright\color{black!88}\textbf{4 1/2} \textit{scil.} demi / moitié.}} \\
      \leavevmode\llap{\fontsize{6.4pt}{7.3pt}\selectfont\color{darkgray}5\hspace{1.8mm}}\fontsize{6.9pt}{7.9pt}\selectfont\strut Après lunch, promené ville : resté lgts arrêté\strut & \smash{\parbox[t]{39mm}{\fontsize{6.2pt}{7.0pt}\selectfont\raggedright\color{black!88}\textbf{5 lgts} \textit{abbr. pro} longtemps. \textbf{5 lunch} \textit{angl. in ms.}}} \\
      \leavevmode\fontsize{6.9pt}{7.9pt}\selectfont\strut devant église, murs de briques rouges, façade\strut &  \\
      \leavevmode\fontsize{6.9pt}{7.9pt}\selectfont\strut si artistique, et haute tour, atteinte par\strut &  \\
      \leavevmode\fontsize{6.9pt}{7.9pt}\selectfont\strut obus bombardement 1914 [?]. Vierge, horizle,\strut & \smash{\parbox[t]{39mm}{\fontsize{6.2pt}{7.0pt}\selectfont\raggedright\color{black!88}\textbf{8 1914} \textit{lect. incert. in ms.} \textbf{8 horizle} \textit{abbr. pro} horizontale. — \textbf{10 qq} \textit{abbr. pro} quelque(s).}} \\
      \leavevmode\fontsize{6.9pt}{7.9pt}\selectfont\strut penche Jésus bras ouverts au-dessus place.\strut &  \\
      \leavevmode\llap{\fontsize{6.4pt}{7.3pt}\selectfont\color{darkgray}10\hspace{1.8mm}}\fontsize{6.9pt}{7.9pt}\selectfont\strut Rencontré qq interprètes. Schumacher m'a\strut &  \\
      \leavevmode\fontsize{6.9pt}{7.9pt}\selectfont\strut appris que les 2/3 des interprètes, jeunes R.A.T.\strut &  \\
      \leavevmode\fontsize{6.9pt}{7.9pt}\selectfont\strut seront renvoyés régiments frs...\strut &  \\
      \leavevmode\fontsize{6.9pt}{7.9pt}\selectfont\strut Canons se sont tus fin jour.\strut &  \\
      \leavevmode\fontsize{6.9pt}{7.9pt}\selectfont\strut Au retour du mess, vers 9h s., lost my way\strut & \smash{\parbox[t]{39mm}{\fontsize{6.2pt}{7.0pt}\selectfont\raggedright\color{black!88}\textbf{14 s.} \textit{abbr. pro} soir. \textbf{14 mess} \textit{angl. in ms.}}} \\
      \leavevmode\llap{\fontsize{6.4pt}{7.3pt}\selectfont\color{darkgray}15\hspace{1.8mm}}\fontsize{6.9pt}{7.9pt}\selectfont\strut sur le plateau, à cause nuit très noire.\strut &  \\
      \leavevmode\fontsize{6.9pt}{7.9pt}\selectfont\strut Dû errer passim avant revenir Albert\strut & \smash{\parbox[t]{39mm}{\fontsize{6.2pt}{7.0pt}\selectfont\raggedright\color{black!88}\textbf{16 passim} \textit{lat.} çà et là.}} \\
      \leavevmode\fontsize{6.9pt}{7.9pt}\selectfont\strut et retrouver mon billet.\strut &  \\
      \leavevmode\fontsize{6.9pt}{7.9pt}\selectfont\strut Ce soir, exceptt, calme absolu.\strut &  \\
\end{tabular}}%
\par\vspace*{\fill}
% --- Registres de bas de page (1. Critique interne [si pertinent] ; 2. Loci paralleli en latin ; 3. Commentaire historique en retrait) ---
\noindent\rule{\textwidth}{0.5pt}\par\vspace{0.65mm}
{\fontsize{6.6pt}{7.7pt}\selectfont%
\noindent\normalfont\color{black!92} \textbf{10} (\textit{Schumacher}) cf. supra f\textsuperscript{o}~11~v\textsuperscript{o} (p.~\pageref{fol:23}), Appendicem I, doc.~1 (p.~\pageref{annexe:doc1}). — \textbf{16} (\textit{Albert}) cf. supra f\textsuperscript{o}~39~r\textsuperscript{o} (p.~\pageref{fol:78}), f\textsuperscript{o}~41~r\textsuperscript{o} (p.~\pageref{fol:82}), infra f\textsuperscript{o}~42~r\textsuperscript{o} (p.~\pageref{fol:84}), f\textsuperscript{o}~42~v\textsuperscript{o} (p.~\pageref{fol:85}), \textit{et passim}.\par\vspace{0.55mm}
\noindent\rule{38mm}{0.4pt}\par\vspace{0.5mm}
\noindent\hspace*{5mm}\parbox{\dimexpr\textwidth-5mm\relax}{\fontsize{6.6pt}{7.7pt}\selectfont\color{black!84}À \textbf{Albert} (Somme), Hippolyte Combes décrit la basilique \textbf{Notre-Dame de Brebières} (édifiée en briques rouges de 1885 à 1895 par l'architecte Edmond Duthoit), dont le dôme de 70~m frappé par les obus allemands en janvier 1915 retenait à l'horizontale la célèbre statue dorée de la \textbf{Vierge penchée} d'Albert Roze. L'adjudant interprète \textbf{Schumacher} lui annonce le projet ministériel de relève des jeunes interprètes de la réserve de l'armée territoriale (\textit{R.A.T.}).}\par}
\newpage
\phantomsection\label{fac:83}
% --- En-tête du fac-similé (Page impaire / droite : Photographie du manuscrit) ---
\noindent\rule{\textwidth}{0.5pt}\par\vspace{0.55mm}
\noindent\textbf{\fontsize{6.8pt}{7.8pt}\selectfont FAC-SIMILÉ EN VIS-À-VIS — FOLIO 41 VERSO} \hfill \textbf{\fontsize{6.8pt}{7.8pt}\selectfont OCTOBRE 1916}\par\vspace{0.25mm}
\noindent{\fontsize{6.8pt}{8.0pt}\selectfont\textit{Livre-Journal d'Hippolyte Combes} — 22 octobre 1916 (Albert (Somme))}\par\vspace{0.55mm}
\noindent\rule{\textwidth}{0.35pt}\par\vspace*{\fill}
{\centering
  \setlength{\fboxsep}{1.2mm}%
  \setlength{\fboxrule}{0.4pt}%
  \fcolorbox{gray!55}{white}{%
    \includegraphics[width=0.95\textwidth,height=0.81\textheight,keepaspectratio]{\facsimiledir/20251126_104346.jpg}%
  }\par}
\vspace*{\fill}
\newpage

% =============================================================================
% VIS-A-VIS SPREAD : FOLIO 42 RECTO (20251126_104400.jpg)
% =============================================================================
\cleardoubleevenpage
\renewcommand{\currentmonthyear}{Octobre 1916}
\renewcommand{\currentfoliolabel}{Folio 42 recto}
\phantomsection\label{fol:84}
% --- En-tête du folio (Page paire / gauche : Transcription, apparat & registres de bas de page) ---
\noindent\rule{\textwidth}{0.5pt}\par\vspace{0.55mm}
\noindent\textbf{\footnotesize FOLIO 42 RECTO} \hfill \textbf{\footnotesize OCTOBRE 1916}\par\vspace{0.25mm}
\noindent{\fontsize{6.9pt}{8.0pt}\selectfont Visite des abris souterrains allemands et du grand cratère de La Boisselle}\par\vspace{0.2mm}
\noindent{\fontsize{6.9pt}{8.0pt}\selectfont\textit{23 octobre 1916} \hfill \textit{Albert, La Boisselle (Somme)}}\par\vspace{0.5mm}
\noindent\rule{\textwidth}{0.35pt}\par\vspace{0.8mm}
% --- Corps de la transcription diplomatique et apparat marginal en prolongement ---
{\renewcommand{\arraystretch}{1.00}%
\noindent\hspace*{3.5mm}\begin{tabular}{@{}p{87mm}@{\hspace{3.5mm}}p{39mm}@{}}
      \leavevmode\fontsize{7.7pt}{9.0pt}\selectfont\strut \underline{Lundi 23} \hfill{} \underline{Oct. /16}\strut & \smash{\parbox[t]{39mm}{\fontsize{6.2pt}{7.0pt}\selectfont\raggedright\color{black!88}\textbf{1 Oct. /16} \textit{abbr. pro} Octobre 1916.}} \\
      \leavevmode\fontsize{7.7pt}{9.0pt}\selectfont\strut Levé très tard.\strut &  \\
      \leavevmode\fontsize{7.7pt}{9.0pt}\selectfont\strut Tout le jour, temps brumeux, mais moins\strut &  \\
      \leavevmode\fontsize{7.7pt}{9.0pt}\selectfont\strut froid.\strut &  \\
      \leavevmode\llap{\fontsize{6.4pt}{7.3pt}\selectfont\color{darkgray}5\hspace{1.8mm}}\fontsize{7.7pt}{9.0pt}\selectfont\strut Après-midi, avec 2 de mes off., venu\strut & \smash{\parbox[t]{39mm}{\fontsize{6.2pt}{7.0pt}\selectfont\raggedright\color{black!88}\textbf{5 off.} \textit{abbr. pro} officier(s). — \textbf{6 s/} \textit{abbr. pro} sur. \textbf{6 lorry} \textit{angl. in ms.}}} \\
      \leavevmode\fontsize{7.7pt}{9.0pt}\selectfont\strut s/ 1 lorry à La Boisselle. Environs d'Albert\strut &  \\
      \leavevmode\fontsize{7.7pt}{9.0pt}\selectfont\strut et jusqu'à l'horizon, troupes au bivouac.\strut &  \\
      \leavevmode\fontsize{7.7pt}{9.0pt}\selectfont\strut Tentes innombrables et dépôts de ttes sortes.\strut & \smash{\parbox[t]{39mm}{\fontsize{6.2pt}{7.0pt}\selectfont\raggedright\color{black!88}\textbf{8 ttes} \textit{abbr. pro} toutes.}} \\
      \leavevmode\fontsize{7.7pt}{9.0pt}\selectfont\strut Visité longuement german dug-outs. Ces\strut & \smash{\parbox[t]{39mm}{\fontsize{6.2pt}{7.0pt}\selectfont\raggedright\color{black!88}\textbf{9 dug-outs} \textit{angl. in ms.}}} \\
      \leavevmode\llap{\fontsize{6.4pt}{7.3pt}\selectfont\color{darkgray}10\hspace{1.8mm}}\fontsize{7.7pt}{9.0pt}\selectfont\strut abris souterrains, très habilement aménagés.\strut &  \\
      \leavevmode\fontsize{7.7pt}{9.0pt}\selectfont\strut Énorme cratère s/ la droite.\strut & \smash{\parbox[t]{39mm}{\fontsize{6.2pt}{7.0pt}\selectfont\raggedright\color{black!88}\textbf{11 s/} \textit{abbr. pro} sur.}} \\
      \leavevmode\fontsize{7.7pt}{9.0pt}\selectfont\strut Impression funèbre et désolation — —\strut &  \\
      \cline{1-1}
      \leavevmode\fontsize{7.7pt}{9.0pt}\selectfont\strut [ — Cartes ples d'Albert : 70 c. ]\strut & \smash{\parbox[t]{39mm}{\fontsize{6.2pt}{7.0pt}\selectfont\raggedright\color{black!88}\textbf{13} \textit{in marg. cum unco} (comptes). \textbf{13 Cartes ples} \textit{abbr. pro} Cartes postales.}} \\
\end{tabular}}%
\par\vspace*{\fill}
% --- Registres de bas de page (1. Critique interne [si pertinent] ; 2. Loci paralleli en latin ; 3. Commentaire historique en retrait) ---
\noindent\rule{\textwidth}{0.5pt}\par\vspace{0.65mm}
{\fontsize{6.6pt}{7.7pt}\selectfont%
\noindent\normalfont\color{black!92} \textbf{6} (\textit{Albert}) cf. supra f\textsuperscript{o}~39~r\textsuperscript{o} (p.~\pageref{fol:78}), f\textsuperscript{o}~41~v\textsuperscript{o} (p.~\pageref{fol:83}), infra f\textsuperscript{o}~42~v\textsuperscript{o} (p.~\pageref{fol:85}), f\textsuperscript{o}~43~r\textsuperscript{o} (p.~\pageref{fol:86}), \textit{et passim}.\par\vspace{0.55mm}
\noindent\rule{38mm}{0.4pt}\par\vspace{0.5mm}
\noindent\hspace*{5mm}\parbox{\dimexpr\textwidth-5mm\relax}{\fontsize{6.6pt}{7.7pt}\selectfont\color{black!84}Excursion en camion avec deux officiers de sa compagnie jusqu'aux ruines de \textbf{La Boisselle} (commune d'Ovillers-la-Boisselle, Somme, à 3,5~km à l'est d'Albert sur la route de Bapaume), conquises en juillet 1916.}\par}
\newpage
\phantomsection\label{fac:84}
% --- En-tête du fac-similé (Page impaire / droite : Photographie du manuscrit) ---
\noindent\rule{\textwidth}{0.5pt}\par\vspace{0.55mm}
\noindent\textbf{\fontsize{6.8pt}{7.8pt}\selectfont FAC-SIMILÉ EN VIS-À-VIS — FOLIO 42 RECTO} \hfill \textbf{\fontsize{6.8pt}{7.8pt}\selectfont OCTOBRE 1916}\par\vspace{0.25mm}
\noindent{\fontsize{6.8pt}{8.0pt}\selectfont\textit{Livre-Journal d'Hippolyte Combes} — 23 octobre 1916 (Albert, La Boisselle (Somme))}\par\vspace{0.55mm}
\noindent\rule{\textwidth}{0.35pt}\par\vspace*{\fill}
{\centering
  \setlength{\fboxsep}{1.2mm}%
  \setlength{\fboxrule}{0.4pt}%
  \fcolorbox{gray!55}{white}{%
    \includegraphics[width=0.95\textwidth,height=0.81\textheight,keepaspectratio]{\facsimiledir/20251126_104400.jpg}%
  }\par}
\vspace*{\fill}
\newpage

% =============================================================================
% VIS-A-VIS SPREAD : FOLIO 42 VERSO (20251126_104418.jpg)
% =============================================================================
\cleardoubleevenpage
\renewcommand{\currentmonthyear}{Octobre 1916}
\renewcommand{\currentfoliolabel}{Folio 42 verso}
\phantomsection\label{fol:85}
% --- En-tête du folio (Page paire / gauche : Transcription, apparat & registres de bas de page) ---
\noindent\rule{\textwidth}{0.5pt}\par\vspace{0.55mm}
\noindent\textbf{\footnotesize FOLIO 42 VERSO} \hfill \textbf{\footnotesize OCTOBRE 1916}\par\vspace{0.25mm}
\noindent{\fontsize{6.9pt}{8.0pt}\selectfont Maintien confirmé à la Mission française et averses dans les fondrières d'Albert}\par\vspace{0.2mm}
\noindent{\fontsize{6.9pt}{8.0pt}\selectfont\textit{24 – 25 octobre 1916} \hfill \textit{Albert (Somme)}}\par\vspace{0.5mm}
\noindent\rule{\textwidth}{0.35pt}\par\vspace{0.8mm}
% --- Corps de la transcription diplomatique et apparat marginal en prolongement ---
{\renewcommand{\arraystretch}{0.87}%
\noindent\hspace*{3.5mm}\begin{tabular}{@{}p{87mm}@{\hspace{3.5mm}}p{39mm}@{}}
      \leavevmode\fontsize{6.9pt}{7.9pt}\selectfont\strut \underline{Mardi 24} \hfill{} \underline{Oct. /16}\strut & \smash{\parbox[t]{39mm}{\fontsize{6.2pt}{7.0pt}\selectfont\raggedright\color{black!88}\textbf{1 Oct. /16} \textit{abbr. pro} Octobre 1916.}} \\
      \leavevmode\fontsize{6.9pt}{7.9pt}\selectfont\strut Plu nuit entière — Temps affreux —\strut &  \\
      \leavevmode\fontsize{6.9pt}{7.9pt}\selectfont\strut Durant matinée, averses fines tombent sans\strut &  \\
      \leavevmode\fontsize{6.9pt}{7.9pt}\selectfont\strut discontinuer. Boue épaisse où l'on\strut &  \\
      \leavevmode\llap{\fontsize{6.4pt}{7.3pt}\selectfont\color{darkgray}5\hspace{1.8mm}}\fontsize{6.9pt}{7.9pt}\selectfont\strut enfonce j'à la cheville.\strut & \smash{\parbox[t]{39mm}{\fontsize{6.2pt}{7.0pt}\selectfont\raggedright\color{black!88}\textbf{5 j'à} \textit{abbr. pro} jusqu'à.}} \\
      \leavevmode\fontsize{6.9pt}{7.9pt}\selectfont\strut Brume s'étend de + en + et cache vue\strut &  \\
      \leavevmode\fontsize{6.9pt}{7.9pt}\selectfont\strut objets à 30 ou 40 mètres —\strut &  \\
      \leavevmode\fontsize{6.9pt}{7.9pt}\selectfont\strut Matin, venu mission frse : adjudant\strut & \smash{\parbox[t]{39mm}{\fontsize{6.2pt}{7.0pt}\selectfont\raggedright\color{black!88}\textbf{8 frse} \textit{abbr. pro} française. — \textbf{9} \textit{lect. incert. in ms.} \textbf{9 Lt} \textit{abbr. pro} Lieutenant. — \textbf{11 s/} \textit{abbr. pro} sur. \textbf{11 qq} \textit{abbr. pro} quelque(s).}} \\
      \leavevmode\fontsize{6.9pt}{7.9pt}\selectfont\strut [Lépin ?] \& Lt Pizzi m'ont appris que ne quitterai\strut &  \\
      \leavevmode\llap{\fontsize{6.4pt}{7.3pt}\selectfont\color{darkgray}10\hspace{1.8mm}}\fontsize{6.9pt}{7.9pt}\selectfont\strut pas mission —\strut &  \\
      \leavevmode\fontsize{6.9pt}{7.9pt}\selectfont\strut Vers 11h m., qq obus s/ prairie, à\strut &  \\
      \leavevmode\fontsize{6.9pt}{7.9pt}\selectfont\strut 400 ou 500 mètres chambre.\strut &  \\
      \leavevmode\fontsize{6.9pt}{7.9pt}\selectfont\strut Nuit calme. Pas de canonnade.\strut &  \\
      \cline{1-1}
      \leavevmode\fontsize{6.9pt}{7.9pt}\selectfont\strut \underline{Merc. 25} | Tjrs même temps, pluvieux,\strut & \smash{\parbox[t]{39mm}{\fontsize{6.2pt}{7.0pt}\selectfont\raggedright\color{black!88}\textbf{14 Merc.} \textit{scil.} Mercredi. \textbf{14 Tjrs} \textit{abbr. pro} Toujours.}} \\
      \leavevmode\llap{\fontsize{6.4pt}{7.3pt}\selectfont\color{darkgray}15\hspace{1.8mm}}\fontsize{6.9pt}{7.9pt}\selectfont\strut mais pas froid.\strut &  \\
      \leavevmode\fontsize{6.9pt}{7.9pt}\selectfont\strut Terrain où ns bivouaquons, vraie\strut & \smash{\parbox[t]{39mm}{\fontsize{6.2pt}{7.0pt}\selectfont\raggedright\color{black!88}\textbf{16 ns} \textit{abbr. pro} nous.}} \\
      \leavevmode\fontsize{6.9pt}{7.9pt}\selectfont\strut fondrière. On patauge ds la boue.\strut & \smash{\parbox[t]{39mm}{\fontsize{6.2pt}{7.0pt}\selectfont\raggedright\color{black!88}\textbf{17 ds} \textit{abbr. pro} dans.}} \\
      \leavevmode\fontsize{6.9pt}{7.9pt}\selectfont\strut Liberté entière. Retiré ds ma\strut & \smash{\parbox[t]{39mm}{\fontsize{6.2pt}{7.0pt}\selectfont\raggedright\color{black!88}\textbf{18 ds} \textit{abbr. pro} dans.}} \\
      \leavevmode\fontsize{6.9pt}{7.9pt}\selectfont\strut chambrette, lis ou écris. Vie\strut &  \\
      \leavevmode\llap{\fontsize{6.4pt}{7.3pt}\selectfont\color{darkgray}20\hspace{1.8mm}}\fontsize{6.9pt}{7.9pt}\selectfont\strut d'1 extrême monotonie —\strut &  \\
\end{tabular}}%
\par\vspace*{\fill}
% --- Registres de bas de page (1. Critique interne [si pertinent] ; 2. Loci paralleli en latin ; 3. Commentaire historique en retrait) ---
\noindent\rule{\textwidth}{0.5pt}\par\vspace{0.65mm}
{\fontsize{6.6pt}{7.7pt}\selectfont%
\noindent\normalfont\color{black!92} \textbf{9} (\textit{Lépin}) cf. supra f\textsuperscript{o}~11~v\textsuperscript{o} (p.~\pageref{fol:23}), infra f\textsuperscript{o}~45~r\textsuperscript{o} (p.~\pageref{fol:90}). — \textbf{9} (\textit{Pizzi}) cf. supra f\textsuperscript{o}~24~r\textsuperscript{o} (p.~\pageref{fol:48}), infra f\textsuperscript{o}~45~r\textsuperscript{o} (p.~\pageref{fol:90}), f\textsuperscript{o}~46~r\textsuperscript{o} (p.~\pageref{fol:92}).\par\vspace{0.55mm}
\noindent\rule{38mm}{0.4pt}\par\vspace{0.5mm}
\noindent\hspace*{5mm}\parbox{\dimexpr\textwidth-5mm\relax}{\fontsize{6.6pt}{7.7pt}\selectfont\color{black!84}Au bureau de la Mission française de la 18\textsuperscript{e} Division à Albert, l'adjudant \textbf{Liepin} (ou \textit{Lupin}/\textit{Lépin}, f\textsuperscript{o}~11~v\textsuperscript{o} et f\textsuperscript{o}~45~r\textsuperscript{o}) et le lieutenant interprète chef de détachement \textbf{Pizzi} rassurent Hippolyte Combes (qui appartient à la classe 1896, âgée de quarante ans): il n'est pas touché par le renvoi des jeunes classes R.A.T. dans les corps de troupe évoqué au f\textsuperscript{o}~41~v\textsuperscript{o}.}\par}
\newpage
\phantomsection\label{fac:85}
% --- En-tête du fac-similé (Page impaire / droite : Photographie du manuscrit) ---
\noindent\rule{\textwidth}{0.5pt}\par\vspace{0.55mm}
\noindent\textbf{\fontsize{6.8pt}{7.8pt}\selectfont FAC-SIMILÉ EN VIS-À-VIS — FOLIO 42 VERSO} \hfill \textbf{\fontsize{6.8pt}{7.8pt}\selectfont OCTOBRE 1916}\par\vspace{0.25mm}
\noindent{\fontsize{6.8pt}{8.0pt}\selectfont\textit{Livre-Journal d'Hippolyte Combes} — 24 – 25 octobre 1916 (Albert (Somme))}\par\vspace{0.55mm}
\noindent\rule{\textwidth}{0.35pt}\par\vspace*{\fill}
{\centering
  \setlength{\fboxsep}{1.2mm}%
  \setlength{\fboxrule}{0.4pt}%
  \fcolorbox{gray!55}{white}{%
    \includegraphics[width=0.95\textwidth,height=0.81\textheight,keepaspectratio]{\facsimiledir/20251126_104418.jpg}%
  }\par}
\vspace*{\fill}
\newpage

% =============================================================================
% VIS-A-VIS SPREAD : FOLIO 43 RECTO (20251126_104440.jpg)
% =============================================================================
\cleardoubleevenpage
\renewcommand{\currentmonthyear}{Octobre 1916}
\renewcommand{\currentfoliolabel}{Folio 43 recto}
\phantomsection\label{fol:86}
% --- En-tête du folio (Page paire / gauche : Transcription, apparat & registres de bas de page) ---
\noindent\rule{\textwidth}{0.5pt}\par\vspace{0.55mm}
\noindent\textbf{\footnotesize FOLIO 43 RECTO} \hfill \textbf{\footnotesize OCTOBRE 1916}\par\vspace{0.25mm}
\noindent{\fontsize{6.9pt}{8.0pt}\selectfont Journées de pluie, d'inaction et d'accalmie dans Albert dévastée}\par\vspace{0.2mm}
\noindent{\fontsize{6.9pt}{8.0pt}\selectfont\textit{26 – 28 octobre 1916} \hfill \textit{Albert (Somme)}}\par\vspace{0.5mm}
\noindent\rule{\textwidth}{0.35pt}\par\vspace{0.8mm}
% --- Corps de la transcription diplomatique et apparat marginal en prolongement ---
{\renewcommand{\arraystretch}{0.93}%
\noindent\hspace*{3.5mm}\begin{tabular}{@{}p{87mm}@{\hspace{3.5mm}}p{39mm}@{}}
      \leavevmode\fontsize{7.3pt}{8.4pt}\selectfont\strut \underline{Jeudi 26} \hfill{} \underline{Oct. /16}\strut & \smash{\parbox[t]{39mm}{\fontsize{6.2pt}{7.0pt}\selectfont\raggedright\color{black!88}\textbf{1 Oct. /16} \textit{abbr. pro} Octobre 1916. — \textbf{2 tjrs} \textit{abbr. pro} toujours. — \textbf{3 ds} \textit{abbr. pro} dans. — \textbf{4 absolt} \textit{abbr. pro} absolument.}} \\
      \leavevmode\fontsize{7.3pt}{8.4pt}\selectfont\strut Même temps. Ciel tjrs couvert de nuages.\strut &  \\
      \leavevmode\fontsize{7.3pt}{8.4pt}\selectfont\strut Averses, ds la journée. Plateau et routes,\strut &  \\
      \leavevmode\fontsize{7.3pt}{8.4pt}\selectfont\strut absolt impraticables. Lacs de boue.\strut &  \\
      \leavevmode\llap{\fontsize{6.4pt}{7.3pt}\selectfont\color{darkgray}5\hspace{1.8mm}}\fontsize{7.3pt}{8.4pt}\selectfont\strut Humidité pénétrante pieds \& jambes.\strut &  \\
      \leavevmode\fontsize{7.3pt}{8.4pt}\selectfont\strut Aucune occupation comme interprète. Libre\strut &  \\
      \leavevmode\fontsize{7.3pt}{8.4pt}\selectfont\strut de ttes mes heures. Me retire ds ma\strut & \smash{\parbox[t]{39mm}{\fontsize{6.2pt}{7.0pt}\selectfont\raggedright\color{black!88}\textbf{7 ds} \textit{abbr. pro} dans. \textbf{7 ttes} \textit{abbr. pro} toutes.}} \\
      \leavevmode\fontsize{7.3pt}{8.4pt}\selectfont\strut chambrette \& y rumine à mon aise.\strut &  \\
      \leavevmode\fontsize{7.3pt}{8.4pt}\selectfont\strut Nuits assez calmes. Canonnade s/ le matin\strut & \smash{\parbox[t]{39mm}{\fontsize{6.2pt}{7.0pt}\selectfont\raggedright\color{black!88}\textbf{9 s/} \textit{abbr. pro} sur.}} \\
      \leavevmode\llap{\fontsize{6.4pt}{7.3pt}\selectfont\color{darkgray}10\hspace{1.8mm}}\fontsize{7.3pt}{8.4pt}\selectfont\strut très violente, mais ne dure pas.\strut &  \\
      \leavevmode\fontsize{7.3pt}{8.4pt}\selectfont\strut [ — Hair cut, par un [R.E. ?] : 50 c ]\strut & \smash{\parbox[t]{39mm}{\fontsize{6.2pt}{7.0pt}\selectfont\raggedright\color{black!88}\textbf{11} \textit{in marg. cum unco} (comptes). \textbf{11 R.E.} \textit{vel} R.F., \textit{lect. incert.} \textbf{11 Hair cut} \textit{angl. in ms.} — \textbf{12 Sam.} \textit{scil.} Samedi. — \textbf{13 Tjrs} \textit{abbr. pro} Toujours. — \textbf{15 auj'd'hui} \textit{abbr. pro} aujourd'hui. — \textbf{17 ts} \textit{abbr. pro} tous. \textbf{17 passim} \textit{lat.} çà et là.}} \\
      \cline{1-1}
      \leavevmode\fontsize{7.3pt}{8.4pt}\selectfont\strut \underline{Sam. 28} — Encore temps pluvieux —\strut &  \\
      \leavevmode\fontsize{7.3pt}{8.4pt}\selectfont\strut Tjrs gros nuages.\strut &  \\
      \leavevmode\fontsize{7.3pt}{8.4pt}\selectfont\strut Boue épaisse où l'on s'enlise.\strut &  \\
      \leavevmode\llap{\fontsize{6.4pt}{7.3pt}\selectfont\color{darkgray}15\hspace{1.8mm}}\fontsize{7.3pt}{8.4pt}\selectfont\strut Matin, hier \& auj'd'hui, descendu\strut &  \\
      \leavevmode\fontsize{7.3pt}{8.4pt}\selectfont\strut en ville, malgré mauvais état routes \&\strut &  \\
      \leavevmode\fontsize{7.3pt}{8.4pt}\selectfont\strut rues. Promené lentement passim. De ts\strut &  \\
      \leavevmode\fontsize{7.3pt}{8.4pt}\selectfont\strut côtés, maisons vides et endommagées par\strut &  \\
      \leavevmode\fontsize{7.3pt}{8.4pt}\selectfont\strut obus. Long arrêt devant église.\strut &  \\
      \leavevmode\llap{\fontsize{6.4pt}{7.3pt}\selectfont\color{darkgray}20\hspace{1.8mm}}\fontsize{7.3pt}{8.4pt}\selectfont\strut Journée \& nuit entt calmes. Aucune\strut &  \\
      \leavevmode\fontsize{7.3pt}{8.4pt}\selectfont\strut canonnade — Silence complet\strut &  \\
      \leavevmode\fontsize{7.3pt}{8.4pt}\selectfont\strut à l'horizon — [ Cartes ples : 75 c ]\strut & \smash{\parbox[t]{39mm}{\fontsize{6.2pt}{7.0pt}\selectfont\raggedright\color{black!88}\textbf{22} \textit{in marg. cum unco} (comptes). \textbf{22 Cartes ples} \textit{abbr. pro} Cartes postales.}} \\
\end{tabular}}%
\par\vspace*{\fill}
% --- Registres de bas de page (1. Critique interne [si pertinent] ; 2. Loci paralleli en latin ; 3. Commentaire historique en retrait) ---
\noindent\rule{\textwidth}{0.5pt}\par\vspace{0.65mm}
{\fontsize{6.6pt}{7.7pt}\selectfont%
\noindent\normalfont\color{black!92} \textbf{1} (\textit{de eodem commorationis loco}) cf. supra f\textsuperscript{o}~42~v\textsuperscript{o} (p.~\pageref{fol:85}).\par\vspace{0.55mm}
\noindent\rule{38mm}{0.4pt}\par\vspace{0.5mm}
\noindent\hspace*{5mm}\parbox{\dimexpr\textwidth-5mm\relax}{\fontsize{6.6pt}{7.7pt}\selectfont\color{black!84}Les pluies persistantes de la dernière décade d'octobre 1916 transforment le plateau crayeux d'Albert en bourbier et imposent des reports successifs aux opérations de la \textit{Reserve Army} (renommée \textit{Fifth Army} le 30 octobre) sur l'Ancre. À la l.~11, Hippolyte Combes se fait couper les cheveux pour 50~c. par un soldat des \textbf{\textit{Royal Fusiliers}} (\textit{R.F.}, 18\textsuperscript{e} Division).}\par}
\newpage
\phantomsection\label{fac:86}
% --- En-tête du fac-similé (Page impaire / droite : Photographie du manuscrit) ---
\noindent\rule{\textwidth}{0.5pt}\par\vspace{0.55mm}
\noindent\textbf{\fontsize{6.8pt}{7.8pt}\selectfont FAC-SIMILÉ EN VIS-À-VIS — FOLIO 43 RECTO} \hfill \textbf{\fontsize{6.8pt}{7.8pt}\selectfont OCTOBRE 1916}\par\vspace{0.25mm}
\noindent{\fontsize{6.8pt}{8.0pt}\selectfont\textit{Livre-Journal d'Hippolyte Combes} — 26 – 28 octobre 1916 (Albert (Somme))}\par\vspace{0.55mm}
\noindent\rule{\textwidth}{0.35pt}\par\vspace*{\fill}
{\centering
  \setlength{\fboxsep}{1.2mm}%
  \setlength{\fboxrule}{0.4pt}%
  \fcolorbox{gray!55}{white}{%
    \includegraphics[width=0.95\textwidth,height=0.81\textheight,keepaspectratio]{\facsimiledir/20251126_104440.jpg}%
  }\par}
\vspace*{\fill}
\newpage

% =============================================================================
% VIS-A-VIS SPREAD : FOLIO 43 VERSO (20251126_104459.jpg)
% =============================================================================
\cleardoubleevenpage
\renewcommand{\currentmonthyear}{Octobre 1916}
\renewcommand{\currentfoliolabel}{Folio 43 verso}
\phantomsection\label{fol:87}
% --- En-tête du folio (Page paire / gauche : Transcription, apparat & registres de bas de page) ---
\noindent\rule{\textwidth}{0.5pt}\par\vspace{0.55mm}
\noindent\textbf{\footnotesize FOLIO 43 VERSO} \hfill \textbf{\footnotesize OCTOBRE 1916}\par\vspace{0.25mm}
\noindent{\fontsize{6.9pt}{8.0pt}\selectfont Enquête sur les pailles à Acheux, rencontre d'un artilleur belge et tempête}\par\vspace{0.2mm}
\noindent{\fontsize{6.9pt}{8.0pt}\selectfont\textit{29 – 30 octobre 1916} \hfill \textit{Albert, Acheux, Senlis-le-Sec (Somme)}}\par\vspace{0.5mm}
\noindent\rule{\textwidth}{0.35pt}\par\vspace{0.8mm}
% --- Corps de la transcription diplomatique et apparat marginal en prolongement ---
{\renewcommand{\arraystretch}{0.83}%
\noindent\hspace*{3.5mm}\begin{tabular}{@{}p{87mm}@{\hspace{3.5mm}}p{39mm}@{}}
      \leavevmode\fontsize{6.6pt}{7.5pt}\selectfont\strut \underline{29 D.} \hfill{} \underline{Octbre /16}\strut & \smash{\parbox[t]{39mm}{\fontsize{6.2pt}{7.0pt}\selectfont\raggedright\color{black!88}\textbf{1 D.} \textit{scil.} Dimanche.}} \\
      \leavevmode\fontsize{6.6pt}{7.5pt}\selectfont\strut Durant tt le jour, au ciel, gros\strut & \smash{\parbox[t]{39mm}{\fontsize{6.2pt}{7.0pt}\selectfont\raggedright\color{black!88}\textbf{2 tt} \textit{abbr. pro} tout.}} \\
      \leavevmode\fontsize{6.6pt}{7.5pt}\selectfont\strut nuages chargés de pluie. Averses intermittentes\strut &  \\
      \leavevmode\fontsize{6.6pt}{7.5pt}\selectfont\strut et tjrs, boue liquide où l'on\strut & \smash{\parbox[t]{39mm}{\fontsize{6.2pt}{7.0pt}\selectfont\raggedright\color{black!88}\textbf{4 tjrs} \textit{abbr. pro} toujours.}} \\
      \leavevmode\llap{\fontsize{6.4pt}{7.3pt}\selectfont\color{darkgray}5\hspace{1.8mm}}\fontsize{6.6pt}{7.5pt}\selectfont\strut enfonce partout —\strut &  \\
      \leavevmode\fontsize{6.6pt}{7.5pt}\selectfont\strut Après-midi, course en auto à Acheux\strut &  \\
      \leavevmode\fontsize{6.6pt}{7.5pt}\selectfont\strut \& Senlis — enquête au sujet de straw \& hay\strut &  \\
      \leavevmode\fontsize{6.6pt}{7.5pt}\selectfont\strut disponibles — —\strut &  \\
      \leavevmode\fontsize{6.6pt}{7.5pt}\selectfont\strut Fils aîné de M. Duchâteau, maral logis\strut & \smash{\parbox[t]{39mm}{\fontsize{6.2pt}{7.0pt}\selectfont\raggedright\color{black!88}\textbf{9 maral} \textit{abbr. pro} maréchal. — \textbf{10 Gd} \textit{abbr. pro} Grand(e).}} \\
      \leavevmode\llap{\fontsize{6.4pt}{7.3pt}\selectfont\color{darkgray}10\hspace{1.8mm}}\fontsize{6.6pt}{7.5pt}\selectfont\strut artrie armée belge, en permission. Gd \&\strut &  \\
      \leavevmode\fontsize{6.6pt}{7.5pt}\selectfont\strut beau j. homme, blond, énergique \& simple,\strut &  \\
      \leavevmode\fontsize{6.6pt}{7.5pt}\selectfont\strut extt sympathique. Écouté avec intérêt récit\strut & \smash{\parbox[t]{39mm}{\fontsize{6.2pt}{7.0pt}\selectfont\raggedright\color{black!88}\textbf{12 extt} \textit{abbr. pro} extrêmement.}} \\
      \leavevmode\fontsize{6.6pt}{7.5pt}\selectfont\strut des divers combats auxquels il a pris part.\strut &  \\
      \leavevmode\fontsize{6.6pt}{7.5pt}\selectfont\strut [ Cartes ples église d'Albert : 50 c\strut & \smash{\parbox[t]{39mm}{\fontsize{6.2pt}{7.0pt}\selectfont\raggedright\color{black!88}\textbf{14} \textit{in marg. cum unco} (comptes). \textbf{14 Cartes ples} \textit{abbr. pro} Cartes postales. — \textbf{16 L.} \textit{scil.} Lundi.}} \\
      \leavevmode\llap{\fontsize{6.4pt}{7.3pt}\selectfont\color{darkgray}15\hspace{1.8mm}}\fontsize{6.6pt}{7.5pt}\selectfont\strut Acheux, 2 cafés : 40 c ]\strut &  \\
      \cline{1-1}
      \leavevmode\fontsize{6.6pt}{7.5pt}\selectfont\strut \underline{30 L.} — Matin, temps couvert \& menaçant.\strut &  \\
      \leavevmode\fontsize{6.6pt}{7.5pt}\selectfont\strut Après-midi, durant 3 h., pluie\strut &  \\
      \leavevmode\fontsize{6.6pt}{7.5pt}\selectfont\strut torrentielle. Dû rester j'après le thé ds\strut & \smash{\parbox[t]{39mm}{\fontsize{6.2pt}{7.0pt}\selectfont\raggedright\color{black!88}\textbf{18 ds} \textit{abbr. pro} dans.}} \\
      \leavevmode\fontsize{6.6pt}{7.5pt}\selectfont\strut la cabane qui ns sert de mess-room.\strut & \smash{\parbox[t]{39mm}{\fontsize{6.2pt}{7.0pt}\selectfont\raggedright\color{black!88}\textbf{19 ns} \textit{abbr. pro} nous. \textbf{19 mess} \textit{angl. in ms.}}} \\
      \leavevmode\llap{\fontsize{6.4pt}{7.3pt}\selectfont\color{darkgray}20\hspace{1.8mm}}\fontsize{6.6pt}{7.5pt}\selectfont\strut Rafales d'1 violence extrême.\strut &  \\
      \leavevmode\fontsize{6.6pt}{7.5pt}\selectfont\strut Coucher de soleil : lueurs très vives, jaunes d'or,\strut &  \\
      \leavevmode\fontsize{6.6pt}{7.5pt}\selectfont\strut à l'horizon.\strut &  \\
      \leavevmode\fontsize{6.6pt}{7.5pt}\selectfont\strut Terrain détrempé, énormes flaques d'eau.\strut &  \\
      \leavevmode\fontsize{6.6pt}{7.5pt}\selectfont\strut Toute la nuit, encore vent de tempête\strut &  \\
      \leavevmode\llap{\fontsize{6.4pt}{7.3pt}\selectfont\color{darkgray}25\hspace{1.8mm}}\fontsize{6.6pt}{7.5pt}\selectfont\strut et très fortes averses.\strut &  \\
\end{tabular}}%
\par\vspace*{\fill}
% --- Registres de bas de page (1. Critique interne [si pertinent] ; 2. Loci paralleli en latin ; 3. Commentaire historique en retrait) ---
\noindent\rule{\textwidth}{0.5pt}\par\vspace{0.65mm}
{\fontsize{6.6pt}{7.7pt}\selectfont%
\noindent\normalfont\color{black!92} \textbf{6} (\textit{Acheux}) cf. supra f\textsuperscript{o}~21~v\textsuperscript{o} (p.~\pageref{fol:43}), f\textsuperscript{o}~25~r\textsuperscript{o} (p.~\pageref{fol:50}), infra f\textsuperscript{o}~44~v\textsuperscript{o} (p.~\pageref{fol:89}), f\textsuperscript{o}~45~r\textsuperscript{o} (p.~\pageref{fol:90}), \textit{et passim}. — \textbf{7} (\textit{Senlis}) cf. supra f\textsuperscript{o}~37~v\textsuperscript{o} (p.~\pageref{fol:75}), f\textsuperscript{o}~38~v\textsuperscript{o} (p.~\pageref{fol:77}), infra f\textsuperscript{o}~48~r\textsuperscript{o} (p.~\pageref{fol:96}), f\textsuperscript{o}~55~r\textsuperscript{o} (p.~\pageref{fol:110}), \textit{et passim}. — \textbf{14} (\textit{Albert}) cf. supra f\textsuperscript{o}~39~r\textsuperscript{o} (p.~\pageref{fol:78}), f\textsuperscript{o}~43~r\textsuperscript{o} (p.~\pageref{fol:86}), infra f\textsuperscript{o}~44~r\textsuperscript{o} (p.~\pageref{fol:88}), f\textsuperscript{o}~44~v\textsuperscript{o} (p.~\pageref{fol:89}), \textit{et passim}.\par\vspace{0.55mm}
\noindent\rule{38mm}{0.4pt}\par\vspace{0.5mm}
\noindent\hspace*{5mm}\parbox{\dimexpr\textwidth-5mm\relax}{\fontsize{6.6pt}{7.7pt}\selectfont\color{black!84}La tournée en automobile du dimanche 29 octobre 1916 à \textbf{Acheux-en-Amiénois} et \textbf{Senlis-le-Sec} pour recenser les pailles et foins (\textit{straw \& hay}) disponibles préfigure le détachement permanent d'Hippolyte Combes et du lieutenant Curtis à Acheux quatre jours plus tard (f\textsuperscript{o}~44~v\textsuperscript{o}--f\textsuperscript{o}~45~r\textsuperscript{o}).}\par}
\newpage
\phantomsection\label{fac:87}
% --- En-tête du fac-similé (Page impaire / droite : Photographie du manuscrit) ---
\noindent\rule{\textwidth}{0.5pt}\par\vspace{0.55mm}
\noindent\textbf{\fontsize{6.8pt}{7.8pt}\selectfont FAC-SIMILÉ EN VIS-À-VIS — FOLIO 43 VERSO} \hfill \textbf{\fontsize{6.8pt}{7.8pt}\selectfont OCTOBRE 1916}\par\vspace{0.25mm}
\noindent{\fontsize{6.8pt}{8.0pt}\selectfont\textit{Livre-Journal d'Hippolyte Combes} — 29 – 30 octobre 1916 (Albert, Acheux, Senlis-le-Sec (Somme))}\par\vspace{0.55mm}
\noindent\rule{\textwidth}{0.35pt}\par\vspace*{\fill}
{\centering
  \setlength{\fboxsep}{1.2mm}%
  \setlength{\fboxrule}{0.4pt}%
  \fcolorbox{gray!55}{white}{%
    \includegraphics[width=0.95\textwidth,height=0.81\textheight,keepaspectratio]{\facsimiledir/20251126_104459.jpg}%
  }\par}
\vspace*{\fill}
\newpage

% =============================================================================
% VIS-A-VIS SPREAD : FOLIO 44 RECTO (20251126_104515.jpg)
% =============================================================================
\cleardoubleevenpage
\renewcommand{\currentmonthyear}{Octobre 1916}
\renewcommand{\currentfoliolabel}{Folio 44 recto}
\phantomsection\label{fol:88}
% --- En-tête du folio (Page paire / gauche : Transcription, apparat & registres de bas de page) ---
\noindent\rule{\textwidth}{0.5pt}\par\vspace{0.55mm}
\noindent\textbf{\footnotesize FOLIO 44 RECTO} \hfill \textbf{\footnotesize OCTOBRE 1916}\par\vspace{0.25mm}
\noindent{\fontsize{6.9pt}{8.0pt}\selectfont Éclaircie de fin octobre et bombardement d'Albert}\par\vspace{0.2mm}
\noindent{\fontsize{6.9pt}{8.0pt}\selectfont\textit{31 octobre 1916} \hfill \textit{Albert (Somme)}}\par\vspace{0.5mm}
\noindent\rule{\textwidth}{0.35pt}\par\vspace{0.8mm}
% --- Corps de la transcription diplomatique et apparat marginal en prolongement ---
{\renewcommand{\arraystretch}{1.00}%
\noindent\hspace*{3.5mm}\begin{tabular}{@{}p{87mm}@{\hspace{3.5mm}}p{39mm}@{}}
      \leavevmode\fontsize{7.7pt}{9.0pt}\selectfont\strut \underline{31 Ma.} \hfill{} \underline{Oct. /16}\strut & \smash{\parbox[t]{39mm}{\fontsize{6.2pt}{7.0pt}\selectfont\raggedright\color{black!88}\textbf{1 Ma.} \textit{scil.} Mardi. \textbf{1 Oct. /16} \textit{abbr. pro} Octobre 1916. — \textbf{2 chassé} \textit{scil.} ont chassé (ellipse de l'auxiliaire).}} \\
      \leavevmode\fontsize{7.7pt}{9.0pt}\selectfont\strut Matin, rafales peu à peu chassé\strut &  \\
      \leavevmode\fontsize{7.7pt}{9.0pt}\selectfont\strut nuages. Tout l'après-midi, enfin, soleil\strut &  \\
      \leavevmode\fontsize{7.7pt}{9.0pt}\selectfont\strut reparu. Belle fin de journée.\strut &  \\
      \leavevmode\llap{\fontsize{6.4pt}{7.3pt}\selectfont\color{darkgray}5\hspace{1.8mm}}\fontsize{7.7pt}{9.0pt}\selectfont\strut Nuit étoilée.\strut &  \\
      \leavevmode\fontsize{7.7pt}{9.0pt}\selectfont\strut Durant jour entier, ville bombardée.\strut &  \\
\end{tabular}}%
\par\vspace*{\fill}
% --- Registres de bas de page (1. Critique interne [si pertinent] ; 2. Loci paralleli en latin ; 3. Commentaire historique en retrait) ---
\noindent\rule{\textwidth}{0.5pt}\par\vspace{0.65mm}
{\fontsize{6.6pt}{7.7pt}\selectfont%
\noindent\normalfont\color{black!92} \textbf{1} (\textit{de eodem commorationis loco}) cf. supra f\textsuperscript{o}~43~v\textsuperscript{o} (p.~\pageref{fol:87}).\par\vspace{0.55mm}
\noindent\rule{38mm}{0.4pt}\par\vspace{0.5mm}
\noindent\hspace*{5mm}\parbox{\dimexpr\textwidth-5mm\relax}{\fontsize{6.6pt}{7.7pt}\selectfont\color{black!84}Profitant de l'éclaircie du mardi 31 octobre 1916 qui rétablit l'observation aérienne et terrestre, l'artillerie allemande soumet toute la journée la ville et le nœud routier d'\textbf{Albert} à des tirs de harcèlement.}\par}
\newpage
\phantomsection\label{fac:88}
% --- En-tête du fac-similé (Page impaire / droite : Photographie du manuscrit) ---
\noindent\rule{\textwidth}{0.5pt}\par\vspace{0.55mm}
\noindent\textbf{\fontsize{6.8pt}{7.8pt}\selectfont FAC-SIMILÉ EN VIS-À-VIS — FOLIO 44 RECTO} \hfill \textbf{\fontsize{6.8pt}{7.8pt}\selectfont OCTOBRE 1916}\par\vspace{0.25mm}
\noindent{\fontsize{6.8pt}{8.0pt}\selectfont\textit{Livre-Journal d'Hippolyte Combes} — 31 octobre 1916 (Albert (Somme))}\par\vspace{0.55mm}
\noindent\rule{\textwidth}{0.35pt}\par\vspace*{\fill}
{\centering
  \setlength{\fboxsep}{1.2mm}%
  \setlength{\fboxrule}{0.4pt}%
  \fcolorbox{gray!55}{white}{%
    \includegraphics[width=0.95\textwidth,height=0.81\textheight,keepaspectratio]{\facsimiledir/20251126_104515.jpg}%
  }\par}
\vspace*{\fill}
\newpage

% =============================================================================
% VIS-A-VIS SPREAD : FOLIO 44 VERSO (20251126_104547.jpg)
% =============================================================================
\cleardoubleevenpage
\renewcommand{\currentmonthyear}{Octobre 1916}
\renewcommand{\currentfoliolabel}{Folio 44 verso}
\phantomsection\label{fol:89}
% --- En-tête du folio (Page paire / gauche : Transcription, apparat & registres de bas de page) ---
\noindent\rule{\textwidth}{0.5pt}\par\vspace{0.55mm}
\noindent\textbf{\footnotesize FOLIO 44 VERSO} \hfill \textbf{\footnotesize OCTOBRE 1916}\par\vspace{0.25mm}
\noindent{\fontsize{6.9pt}{8.0pt}\selectfont Toussaint 1916 au cimetière dévasté d'Albert, prise d'armes écossaise et ordre de départ}\par\vspace{0.2mm}
\noindent{\fontsize{6.9pt}{8.0pt}\selectfont\textit{1er novembre 1916} \hfill \textit{Albert (Somme)}}\par\vspace{0.5mm}
\noindent\rule{\textwidth}{0.35pt}\par\vspace{0.8mm}
% --- Corps de la transcription diplomatique et apparat marginal en prolongement ---
{\renewcommand{\arraystretch}{0.87}%
\noindent\hspace*{3.5mm}\begin{tabular}{@{}p{87mm}@{\hspace{3.5mm}}p{39mm}@{}}
      \leavevmode\fontsize{6.9pt}{7.9pt}\selectfont\strut \underline{1er Me.} Toussaint. — Novembre 1916 —\strut & \smash{\parbox[t]{39mm}{\fontsize{6.2pt}{7.0pt}\selectfont\raggedright\color{black!88}\textbf{1 1er} \textit{scil.} premier.}} \\
      \leavevmode\fontsize{6.9pt}{7.9pt}\selectfont\strut Ciel bas et couleur de cendre.\strut &  \\
      \leavevmode\fontsize{6.9pt}{7.9pt}\selectfont\strut Bise froide. Pluie menaçante —\strut &  \\
      \leavevmode\fontsize{6.9pt}{7.9pt}\selectfont\strut Matin, visite au cimetière d'Albert. Gros\strut &  \\
      \leavevmode\llap{\fontsize{6.4pt}{7.3pt}\selectfont\color{darkgray}5\hspace{1.8mm}}\fontsize{6.9pt}{7.9pt}\selectfont\strut nuages assombrissant l'air. Allées jonchées de\strut &  \\
      \leavevmode\fontsize{6.9pt}{7.9pt}\selectfont\strut feuilles mortes. Tombes dévastées par obus. Pierres\strut &  \\
      \leavevmode\fontsize{6.9pt}{7.9pt}\selectfont\strut brisées. Caveaux béants. Croix renversées au 1/2 de\strut & \smash{\parbox[t]{39mm}{\fontsize{6.2pt}{7.0pt}\selectfont\raggedright\color{black!88}\textbf{7 1/2} \textit{scil.} demi / moitié.}} \\
      \leavevmode\fontsize{6.9pt}{7.9pt}\selectfont\strut pierres effondrées...\strut &  \\
      \leavevmode\fontsize{6.9pt}{7.9pt}\selectfont\strut A côté, champ tt rempli de croix blanches \&\strut & \smash{\parbox[t]{39mm}{\fontsize{6.2pt}{7.0pt}\selectfont\raggedright\color{black!88}\textbf{9 tt} \textit{abbr. pro} tout.}} \\
      \leavevmode\llap{\fontsize{6.4pt}{7.3pt}\selectfont\color{darkgray}10\hspace{1.8mm}}\fontsize{6.9pt}{7.9pt}\selectfont\strut jaunes : anglais killed in action —\strut &  \\
      \leavevmode\fontsize{6.9pt}{7.9pt}\selectfont\strut Après-midi, 3h, place du monument\strut &  \\
      \leavevmode\fontsize{6.9pt}{7.9pt}\selectfont\strut prise d'armes. Devant pl. bataillons rassemblés,\strut &  \\
      \leavevmode\fontsize{6.9pt}{7.9pt}\selectfont\strut génal écossais épingle div. décorations s/ tuniques\strut & \smash{\parbox[t]{39mm}{\fontsize{6.2pt}{7.0pt}\selectfont\raggedright\color{black!88}\textbf{13 s/} \textit{abbr. pro} sur. \textbf{13 div.} \textit{abbr. pro} diverses (devant décorations). — \textbf{14 N.C.O's} \textit{angl.} Non-Commissioned Officers. \textbf{14 tommies} \textit{angl. in ms.} — \textbf{15 s/} \textit{abbr. pro} sur. — \textbf{19 ds} \textit{abbr. pro} dans. — \textbf{20 qq} \textit{abbr. pro} quelque(s).}} \\
      \leavevmode\fontsize{6.9pt}{7.9pt}\selectfont\strut de tommies \& N.C.O's qui se sont distingués offensive\strut &  \\
      \leavevmode\llap{\fontsize{6.4pt}{7.3pt}\selectfont\color{darkgray}15\hspace{1.8mm}}\fontsize{6.9pt}{7.9pt}\selectfont\strut s/ la Somme. Pluie fine. Défilé admirable.\strut &  \\
      \leavevmode\fontsize{6.9pt}{7.9pt}\selectfont\strut Derrière moi, jardin public, non entretenu.\strut &  \\
      \leavevmode\fontsize{6.9pt}{7.9pt}\selectfont\strut Feuilles sèches, arbres roussis par l'automne.\strut &  \\
      \leavevmode\fontsize{6.9pt}{7.9pt}\selectfont\strut Soir, venu à Church meeting, route\strut &  \\
      \leavevmode\fontsize{6.9pt}{7.9pt}\selectfont\strut d'Amiens. Rues ds l'obscurité complète.\strut &  \\
      \leavevmode\llap{\fontsize{6.4pt}{7.3pt}\selectfont\color{darkgray}20\hspace{1.8mm}}\fontsize{6.9pt}{7.9pt}\selectfont\strut Çà et là, qq falots \& lanternes de M.P.'s\strut &  \\
      \leavevmode\fontsize{6.9pt}{7.9pt}\selectfont\strut aux carrefours.\strut &  \\
      \leavevmode\fontsize{6.9pt}{7.9pt}\selectfont\strut A 9h soir, Lt Curtis, venu m'apporter\strut & \smash{\parbox[t]{39mm}{\fontsize{6.2pt}{7.0pt}\selectfont\raggedright\color{black!88}\textbf{22 Lt} \textit{abbr. pro} Lieutenant.}} \\
      \leavevmode\fontsize{6.9pt}{7.9pt}\selectfont\strut ordres. Demain matin, devrai partir avec lui\strut &  \\
      \leavevmode\fontsize{6.9pt}{7.9pt}\selectfont\strut à Acheux. Ns sommes détachés de la Coy Transport —\strut & \smash{\parbox[t]{39mm}{\fontsize{6.2pt}{7.0pt}\selectfont\raggedright\color{black!88}\textbf{24 Coy} \textit{angl.} Company.}} \\
\end{tabular}}%
\par\vspace*{\fill}
% --- Registres de bas de page (1. Critique interne [si pertinent] ; 2. Loci paralleli en latin ; 3. Commentaire historique en retrait) ---
\noindent\rule{\textwidth}{0.5pt}\par\vspace{0.65mm}
{\fontsize{6.6pt}{7.7pt}\selectfont%
\noindent\normalfont\color{black!92} \textbf{4} (\textit{Albert}) cf. supra f\textsuperscript{o}~39~r\textsuperscript{o} (p.~\pageref{fol:78}), f\textsuperscript{o}~44~r\textsuperscript{o} (p.~\pageref{fol:88}), infra f\textsuperscript{o}~45~r\textsuperscript{o} (p.~\pageref{fol:90}), f\textsuperscript{o}~45~v\textsuperscript{o} (p.~\pageref{fol:91}), \textit{et passim}. — \textbf{19} (\textit{Amiens}) cf. supra f\textsuperscript{o}~25~v\textsuperscript{o} (p.~\pageref{fol:51}), f\textsuperscript{o}~32~v\textsuperscript{o} (p.~\pageref{fol:65}), infra f\textsuperscript{o}~49~r\textsuperscript{o} (p.~\pageref{fol:98}), f\textsuperscript{o}~59~v\textsuperscript{o} (p.~\pageref{fol:119}), \textit{et passim}. — \textbf{22} (\textit{Curtis}) cf. supra f\textsuperscript{o}~1~r\textsuperscript{o} (p.~\pageref{fol:2}), f\textsuperscript{o}~18~v\textsuperscript{o} (p.~\pageref{fol:37}), infra f\textsuperscript{o}~45~r\textsuperscript{o} (p.~\pageref{fol:90}), f\textsuperscript{o}~45~v\textsuperscript{o} (p.~\pageref{fol:91}), \textit{et passim}.\par\vspace{0.55mm}
\noindent\rule{38mm}{0.4pt}\par\vspace{0.5mm}
\noindent\hspace*{5mm}\parbox{\dimexpr\textwidth-5mm\relax}{\fontsize{6.6pt}{7.7pt}\selectfont\color{black!84}Le mercredi \textbf{1\textsuperscript{er} novembre 1916} (jour de la Toussaint), Hippolyte Combes décrit le cimetière civil d'Albert ravagé par les obus et le carré militaire britannique contigu (\textit{Albert Communal Cemetery Extension}, ouvert par l'armée britannique en août 1915; notice CWGC), puis assiste à 15~h sur la place d'Albert à une remise de décorations aux combattants de la Somme.}\par}
\newpage
\phantomsection\label{fac:89}
% --- En-tête du fac-similé (Page impaire / droite : Photographie du manuscrit) ---
\noindent\rule{\textwidth}{0.5pt}\par\vspace{0.55mm}
\noindent\textbf{\fontsize{6.8pt}{7.8pt}\selectfont FAC-SIMILÉ EN VIS-À-VIS — FOLIO 44 VERSO} \hfill \textbf{\fontsize{6.8pt}{7.8pt}\selectfont OCTOBRE 1916}\par\vspace{0.25mm}
\noindent{\fontsize{6.8pt}{8.0pt}\selectfont\textit{Livre-Journal d'Hippolyte Combes} — 1er novembre 1916 (Albert (Somme))}\par\vspace{0.55mm}
\noindent\rule{\textwidth}{0.35pt}\par\vspace*{\fill}
{\centering
  \setlength{\fboxsep}{1.2mm}%
  \setlength{\fboxrule}{0.4pt}%
  \fcolorbox{gray!55}{white}{%
    \includegraphics[width=0.95\textwidth,height=0.81\textheight,keepaspectratio]{\facsimiledir/20251126_104547.jpg}%
  }\par}
\vspace*{\fill}
\newpage

% =============================================================================
% VIS-A-VIS SPREAD : FOLIO 45 RECTO (20251126_104601.jpg)
% =============================================================================
\cleardoubleevenpage
\renewcommand{\currentmonthyear}{Novembre 1916}
\renewcommand{\currentfoliolabel}{Folio 45 recto}
\phantomsection\label{fol:90}
% --- En-tête du folio (Page paire / gauche : Transcription, apparat & registres de bas de page) ---
\noindent\rule{\textwidth}{0.5pt}\par\vspace{0.55mm}
\noindent\textbf{\footnotesize FOLIO 45 RECTO} \hfill \textbf{\footnotesize NOVEMBRE 1916}\par\vspace{0.25mm}
\noindent{\fontsize{6.9pt}{8.0pt}\selectfont Ordre de mission du colonel Grose et départ sous la pluie vers Acheux}\par\vspace{0.2mm}
\noindent{\fontsize{6.9pt}{8.0pt}\selectfont\textit{2 novembre 1916} \hfill \textit{Albert, Acheux-en-Amiénois (Somme)}}\par\vspace{0.5mm}
\noindent\rule{\textwidth}{0.35pt}\par\vspace{0.8mm}
% --- Corps de la transcription diplomatique et apparat marginal en prolongement ---
{\renewcommand{\arraystretch}{0.87}%
\noindent\hspace*{3.5mm}\begin{tabular}{@{}p{87mm}@{\hspace{3.5mm}}p{39mm}@{}}
      \leavevmode\fontsize{6.9pt}{7.9pt}\selectfont\strut \underline{2 J.} — \hfill{} \underline{Novbre /16}\strut & \smash{\parbox[t]{39mm}{\fontsize{6.2pt}{7.0pt}\selectfont\raggedright\color{black!88}\textbf{1 J.} \textit{scil.} Jeudi.}} \\
      \leavevmode\fontsize{6.9pt}{7.9pt}\selectfont\strut Levé 6h m. En 1 h.,\strut & \smash{\parbox[t]{39mm}{\fontsize{6.2pt}{7.0pt}\selectfont\raggedright\color{black!88}\textbf{2 m.} \textit{abbr. pro} matin.}} \\
      \leavevmode\fontsize{6.9pt}{7.9pt}\selectfont\strut ts paquets terminés.\strut & \smash{\parbox[t]{39mm}{\fontsize{6.2pt}{7.0pt}\selectfont\raggedright\color{black!88}\textbf{3 ts} \textit{abbr. pro} tous.}} \\
      \leavevmode\fontsize{6.9pt}{7.9pt}\selectfont\strut Monté au mess-hut. Là, au breakfast, eu\strut & \smash{\parbox[t]{39mm}{\fontsize{6.2pt}{7.0pt}\selectfont\raggedright\color{black!88}\textbf{4 breakfast} \textit{angl. in ms.}}} \\
      \leavevmode\llap{\fontsize{6.4pt}{7.3pt}\selectfont\color{darkgray}5\hspace{1.8mm}}\fontsize{6.9pt}{7.9pt}\selectfont\strut l'idée de passer bureau mission avant partir.\strut &  \\
      \leavevmode\fontsize{6.9pt}{7.9pt}\selectfont\strut Cap. Pant., approuvé.\strut & \smash{\parbox[t]{39mm}{\fontsize{6.2pt}{7.0pt}\selectfont\raggedright\color{black!88}\textbf{6 Cap.} \textit{abbr. pro} Capitaine. — \textbf{7} \textit{lect. incert. in ms.}}} \\
      \leavevmode\fontsize{6.9pt}{7.9pt}\selectfont\strut Au bureau, [Lépin ?] et surtout Serpin, m'ont\strut &  \\
      \leavevmode\fontsize{6.9pt}{7.9pt}\selectfont\strut conseillé demander ordre écrit, après avoir vu Lt Pizzi.\strut & \smash{\parbox[t]{39mm}{\fontsize{6.2pt}{7.0pt}\selectfont\raggedright\color{black!88}\textbf{8 Lt} \textit{abbr. pro} Lieutenant.}} \\
      \leavevmode\fontsize{6.9pt}{7.9pt}\selectfont\strut Allé trouver ds sa chambre Lt Pizzi —\strut & \smash{\parbox[t]{39mm}{\fontsize{6.2pt}{7.0pt}\selectfont\raggedright\color{black!88}\textbf{9 ds} \textit{abbr. pro} dans. \textbf{9 Lt} \textit{abbr. pro} Lieutenant.}} \\
      \leavevmode\llap{\fontsize{6.4pt}{7.3pt}\selectfont\color{darkgray}10\hspace{1.8mm}}\fontsize{6.9pt}{7.9pt}\selectfont\strut Me suis rendu au Train Office. Major Hughes\strut &  \\
      \leavevmode\fontsize{6.9pt}{7.9pt}\selectfont\strut tempêté. Col. Grose, signé ordre d'accompagner\strut & \smash{\parbox[t]{39mm}{\fontsize{6.2pt}{7.0pt}\selectfont\raggedright\color{black!88}\textbf{11 Col.} \textit{abbr. pro} Colonel.}} \\
      \leavevmode\fontsize{6.9pt}{7.9pt}\selectfont\strut Lt Curtis, à Acheux.\strut & \smash{\parbox[t]{39mm}{\fontsize{6.2pt}{7.0pt}\selectfont\raggedright\color{black!88}\textbf{12 Lt} \textit{abbr. pro} Lieutenant. — \textbf{13 1/2} \textit{scil.} demi / moitié. — \textbf{14 ts} \textit{abbr. pro} tous. — \textbf{15} \textit{lect. incert. in ms.} \textbf{15 Lt} \textit{abbr. pro} Lieutenant. \textbf{15 kit} \textit{angl. in ms.}}} \\
      \leavevmode\fontsize{6.9pt}{7.9pt}\selectfont\strut A 9h 1/2, enfin, auto venue me prendre\strut &  \\
      \leavevmode\fontsize{6.9pt}{7.9pt}\selectfont\strut à mon billet \& ts mes bagages. Les\strut &  \\
      \leavevmode\llap{\fontsize{6.4pt}{7.3pt}\selectfont\color{darkgray}15\hspace{1.8mm}}\fontsize{6.9pt}{7.9pt}\selectfont\strut [verbe illis.] [au ?] kit de Lt Curtis.\strut &  \\
      \leavevmode\fontsize{6.9pt}{7.9pt}\selectfont\strut Pluie battante. Routes affreuses.\strut &  \\
      \leavevmode\fontsize{6.9pt}{7.9pt}\selectfont\strut Vers 11h —, arrivés à Acheux.\strut &  \\
      \leavevmode\fontsize{6.9pt}{7.9pt}\selectfont\strut Vu Town Major. Recherché billets \& mess-room.\strut & \smash{\parbox[t]{39mm}{\fontsize{6.2pt}{7.0pt}\selectfont\raggedright\color{black!88}\textbf{18 billets} \textit{angl. in ms.}}} \\
      \leavevmode\fontsize{6.9pt}{7.9pt}\selectfont\strut Installations provisoires.\strut &  \\
      \leavevmode\llap{\fontsize{6.4pt}{7.3pt}\selectfont\color{darkgray}20\hspace{1.8mm}}\fontsize{6.9pt}{7.9pt}\selectfont\strut Rues ville état lamentable. Ornières.\strut &  \\
      \leavevmode\fontsize{6.9pt}{7.9pt}\selectfont\strut Tout l'après-midi, perdu temps —\strut &  \\
\end{tabular}}%
\par\vspace*{\fill}
% --- Registres de bas de page (1. Critique interne [si pertinent] ; 2. Loci paralleli en latin ; 3. Commentaire historique en retrait) ---
\noindent\rule{\textwidth}{0.5pt}\par\vspace{0.65mm}
{\fontsize{6.6pt}{7.7pt}\selectfont%
\noindent\normalfont \textbf{Discours direct bilingue :} Aux l.~17--19, Hippolyte Combes rapporte au style direct, en français puis en anglais entre guillemets, sa réplique au major Hughes et l'accueil du colonel Grose.\par\vspace{0.5mm}
\noindent\rule{22mm}{0.3pt}\par\vspace{0.45mm}
\noindent\normalfont\color{black!92} \textbf{7} (\textit{Lépin}) cf. supra f\textsuperscript{o}~11~v\textsuperscript{o} (p.~\pageref{fol:23}), f\textsuperscript{o}~42~v\textsuperscript{o} (p.~\pageref{fol:85}). — \textbf{8} (\textit{Pizzi}) cf. supra f\textsuperscript{o}~24~r\textsuperscript{o} (p.~\pageref{fol:48}), f\textsuperscript{o}~42~v\textsuperscript{o} (p.~\pageref{fol:85}), infra f\textsuperscript{o}~46~r\textsuperscript{o} (p.~\pageref{fol:92}). — \textbf{10} (\textit{Major Hughes}) cf. infra f\textsuperscript{o}~48~r\textsuperscript{o} (p.~\pageref{fol:96}).\par\vspace{0.55mm}
\noindent\rule{38mm}{0.4pt}\par\vspace{0.5mm}
\noindent\hspace*{5mm}\parbox{\dimexpr\textwidth-5mm\relax}{\fontsize{6.6pt}{7.7pt}\selectfont\color{black!84}Le commandant du Train divisionnaire (\textit{18th Divisional Train Office}) qui signe l'ordre détachant Hippolyte Combes au service du fourrage à Acheux-en-Amiénois est le \textbf{lieutenant-colonel Daniel Charles Evans Grose}, A.S.C. (\texttt{Col. Grose}), commandant le \textit{18th Divisional Train} de 1915 à 1918 et décoré du \textit{Distinguished Service Order} (\textit{London Gazette}, 4 juin 1917, p.~5470 ; Nichols, 1922, p.~110).}\par}
\newpage
\phantomsection\label{fac:90}
% --- En-tête du fac-similé (Page impaire / droite : Photographie du manuscrit) ---
\noindent\rule{\textwidth}{0.5pt}\par\vspace{0.55mm}
\noindent\textbf{\fontsize{6.8pt}{7.8pt}\selectfont FAC-SIMILÉ EN VIS-À-VIS — FOLIO 45 RECTO} \hfill \textbf{\fontsize{6.8pt}{7.8pt}\selectfont NOVEMBRE 1916}\par\vspace{0.25mm}
\noindent{\fontsize{6.8pt}{8.0pt}\selectfont\textit{Livre-Journal d'Hippolyte Combes} — 2 novembre 1916 (Albert, Acheux-en-Amiénois (Somme))}\par\vspace{0.55mm}
\noindent\rule{\textwidth}{0.35pt}\par\vspace*{\fill}
{\centering
  \setlength{\fboxsep}{1.2mm}%
  \setlength{\fboxrule}{0.4pt}%
  \fcolorbox{gray!55}{white}{%
    \includegraphics[width=0.95\textwidth,height=0.81\textheight,keepaspectratio]{\facsimiledir/20251126_104601.jpg}%
  }\par}
\vspace*{\fill}
\newpage

% =============================================================================
% VIS-A-VIS SPREAD : FOLIO 45 VERSO (20251126_104623.jpg)
% =============================================================================
\cleardoubleevenpage
\renewcommand{\currentmonthyear}{Novembre 1916}
\renewcommand{\currentfoliolabel}{Folio 45 verso}
\phantomsection\label{fol:91}
% --- En-tête du folio (Page paire / gauche : Transcription, apparat & registres de bas de page) ---
\noindent\rule{\textwidth}{0.5pt}\par\vspace{0.55mm}
\noindent\textbf{\footnotesize FOLIO 45 VERSO} \hfill \textbf{\footnotesize NOVEMBRE 1916}\par\vspace{0.25mm}
\noindent{\fontsize{6.9pt}{8.0pt}\selectfont Course nocturne à Raincheval et dernier aller-retour au camp d'Albert}\par\vspace{0.2mm}
\noindent{\fontsize{6.9pt}{8.0pt}\selectfont\textit{2 (soir) – 3 novembre 1916} \hfill \textit{Acheux, Raincheval, Albert (Somme)}}\par\vspace{0.5mm}
\noindent\rule{\textwidth}{0.35pt}\par\vspace{0.8mm}
% --- Corps de la transcription diplomatique et apparat marginal en prolongement ---
{\renewcommand{\arraystretch}{0.87}%
\noindent\hspace*{3.5mm}\begin{tabular}{@{}p{87mm}@{\hspace{3.5mm}}p{39mm}@{}}
      \leavevmode\fontsize{6.9pt}{7.9pt}\selectfont\strut Après le dîner, course en auto à Raincheval.\strut &  \\
      \leavevmode\fontsize{6.9pt}{7.9pt}\selectfont\strut Vu Major Levy, of the D.O.S.T.\strut & \smash{\parbox[t]{39mm}{\fontsize{6.2pt}{7.0pt}\selectfont\raggedright\color{black!88}\textbf{2 D.O.S.T.} \textit{sigle militaire, développement non établi}. — \textbf{4} \textit{in marg. cum unco} (comptes).}} \\
      \leavevmode\fontsize{6.9pt}{7.9pt}\selectfont\strut Nuit très noire. — Couché assez tard.\strut &  \\
      \leavevmode\fontsize{6.9pt}{7.9pt}\selectfont\strut [ Donné à Victor Duchâteau : 5 f —\strut &  \\
      \leavevmode\llap{\fontsize{6.4pt}{7.3pt}\selectfont\color{darkgray}5\hspace{1.8mm}}\fontsize{6.9pt}{7.9pt}\selectfont\strut Mess, octobre : 47 f 70\strut &  \\
      \leavevmode\fontsize{6.9pt}{7.9pt}\selectfont\strut Acheux : ctes ples : 15 c ]\strut & \smash{\parbox[t]{39mm}{\fontsize{6.2pt}{7.0pt}\selectfont\raggedright\color{black!88}\textbf{6 ctes ples} \textit{abbr. pro} cartes postales. — \textbf{7 Ve.} \textit{scil.} Vendredi.}} \\
      \cline{1-1}
      \leavevmode\fontsize{6.9pt}{7.9pt}\selectfont\strut \underline{3 Ve.} Matin, ciel couvert, 1 peu de soleil\strut &  \\
      \leavevmode\fontsize{6.9pt}{7.9pt}\selectfont\strut l'après-midi.\strut &  \\
      \leavevmode\fontsize{6.9pt}{7.9pt}\selectfont\strut Partis en auto avec Lt Curtis. Resté à Albert\strut & \smash{\parbox[t]{39mm}{\fontsize{6.2pt}{7.0pt}\selectfont\raggedright\color{black!88}\textbf{9 Lt} \textit{abbr. pro} Lieutenant.}} \\
      \leavevmode\llap{\fontsize{6.4pt}{7.3pt}\selectfont\color{darkgray}10\hspace{1.8mm}}\fontsize{6.9pt}{7.9pt}\selectfont\strut et revenu seul s/ 1 lorry. Routes\strut & \smash{\parbox[t]{39mm}{\fontsize{6.2pt}{7.0pt}\selectfont\raggedright\color{black!88}\textbf{10 s/} \textit{abbr. pro} sur. \textbf{10 lorry} \textit{angl. in ms.}}} \\
      \leavevmode\fontsize{6.9pt}{7.9pt}\selectfont\strut mauvaises et sillonnées de véhicules divers.\strut &  \\
      \leavevmode\fontsize{6.9pt}{7.9pt}\selectfont\strut Revenu au camp, 152 Coy et au\strut & \smash{\parbox[t]{39mm}{\fontsize{6.2pt}{7.0pt}\selectfont\raggedright\color{black!88}\textbf{12 Coy} \textit{angl.} Company.}} \\
      \leavevmode\fontsize{6.9pt}{7.9pt}\selectfont\strut bureau mission.\strut &  \\
      \leavevmode\fontsize{6.9pt}{7.9pt}\selectfont\strut Après lunch, 2h 30, Acheux.\strut & \smash{\parbox[t]{39mm}{\fontsize{6.2pt}{7.0pt}\selectfont\raggedright\color{black!88}\textbf{14 lunch} \textit{angl. in ms.}}} \\
      \leavevmode\llap{\fontsize{6.4pt}{7.3pt}\selectfont\color{darkgray}15\hspace{1.8mm}}\fontsize{6.9pt}{7.9pt}\selectfont\strut Passé journée ds ma chambre — au bureau.\strut & \smash{\parbox[t]{39mm}{\fontsize{6.2pt}{7.0pt}\selectfont\raggedright\color{black!88}\textbf{15 ds} \textit{abbr. pro} dans.}} \\
      \cline{1-1}
      \leavevmode\fontsize{6.9pt}{7.9pt}\selectfont\strut [ — Mission, Albert : 58 f 75 (30 f retenus pour\strut & \smash{\parbox[t]{39mm}{\fontsize{6.2pt}{7.0pt}\selectfont\raggedright\color{black!88}\textbf{16} \textit{in marg. cum unco} (comptes).}} \\
      \leavevmode\fontsize{6.9pt}{7.9pt}\selectfont\strut 6 jours de permission) ]\strut &  \\
\end{tabular}}%
\par\vspace*{\fill}
% --- Registres de bas de page (1. Critique interne [si pertinent] ; 2. Loci paralleli en latin ; 3. Commentaire historique en retrait) ---
\noindent\rule{\textwidth}{0.5pt}\par\vspace{0.65mm}
{\fontsize{6.6pt}{7.7pt}\selectfont%
\noindent\normalfont \textbf{Doublon photographique :} Ce feuillet (fin de la journée du 2 novembre et journée du 3 novembre 1916) a fait l'objet de deux prises de vue lors de la numérisation du carnet (voir f\textsuperscript{o}~46~v\textsuperscript{o}) ; la transcription diplomatique est maintenue à l'identique sur les deux doubles pages.\par\vspace{0.5mm}
\noindent\rule{22mm}{0.3pt}\par\vspace{0.45mm}
\noindent\normalfont\color{black!92} \textbf{1--17} (\textit{eadem pagina bis photographata}) cf. infra f\textsuperscript{o}~46~v\textsuperscript{o} (p.~\pageref{fol:93}). — \textbf{1} (\textit{Raincheval}) cf. supra f\textsuperscript{o}~20~v\textsuperscript{o} (p.~\pageref{fol:41}), f\textsuperscript{o}~24~r\textsuperscript{o} (p.~\pageref{fol:48}), infra f\textsuperscript{o}~46~v\textsuperscript{o} (p.~\pageref{fol:93}), f\textsuperscript{o}~48~v\textsuperscript{o} (p.~\pageref{fol:97}), \textit{et passim}. — \textbf{2} (\textit{Major Levy}) cf. infra f\textsuperscript{o}~46~v\textsuperscript{o} (p.~\pageref{fol:93}), f\textsuperscript{o}~55~r\textsuperscript{o} (p.~\pageref{fol:110}).\par\vspace{0.55mm}
\noindent\rule{38mm}{0.4pt}\par\vspace{0.5mm}
\noindent\hspace*{5mm}\parbox{\dimexpr\textwidth-5mm\relax}{\fontsize{6.6pt}{7.7pt}\selectfont\color{black!84}Dès son arrivée à \textbf{Acheux-en-Amiénois} le 2 novembre 1916, Hippolyte Combes est présenté à Raincheval au \textbf{major Levy} (officier supérieur britannique dirigeant le service des achats et réquisitions d'armée). Le 3 novembre, il revient faire ses adieux au camp de la \textbf{152\textsuperscript{e} Compagnie A.S.C.} et perçoit à Albert le reliquat de sa solde.}\par}
\newpage
\phantomsection\label{fac:91}
% --- En-tête du fac-similé (Page impaire / droite : Photographie du manuscrit) ---
\noindent\rule{\textwidth}{0.5pt}\par\vspace{0.55mm}
\noindent\textbf{\fontsize{6.8pt}{7.8pt}\selectfont FAC-SIMILÉ EN VIS-À-VIS — FOLIO 45 VERSO} \hfill \textbf{\fontsize{6.8pt}{7.8pt}\selectfont NOVEMBRE 1916}\par\vspace{0.25mm}
\noindent{\fontsize{6.8pt}{8.0pt}\selectfont\textit{Livre-Journal d'Hippolyte Combes} — 2 (soir) – 3 novembre 1916 (Acheux, Raincheval, Albert (Somme))}\par\vspace{0.55mm}
\noindent\rule{\textwidth}{0.35pt}\par\vspace*{\fill}
{\centering
  \setlength{\fboxsep}{1.2mm}%
  \setlength{\fboxrule}{0.4pt}%
  \fcolorbox{gray!55}{white}{%
    \includegraphics[width=0.95\textwidth,height=0.81\textheight,keepaspectratio]{\facsimiledir/20251126_104623.jpg}%
  }\par}
\vspace*{\fill}
\newpage

% =============================================================================
% VIS-A-VIS SPREAD : FOLIO 46 RECTO (20251126_104636.jpg)
% =============================================================================
\cleardoubleevenpage
\renewcommand{\currentmonthyear}{Novembre 1916}
\renewcommand{\currentfoliolabel}{Folio 46 recto}
\phantomsection\label{fol:92}
% --- En-tête du folio (Page paire / gauche : Transcription, apparat & registres de bas de page) ---
\noindent\rule{\textwidth}{0.5pt}\par\vspace{0.55mm}
\noindent\textbf{\footnotesize FOLIO 46 RECTO} \hfill \textbf{\footnotesize NOVEMBRE 1916}\par\vspace{0.25mm}
\noindent{\fontsize{6.9pt}{8.0pt}\selectfont Course à pied à Bertrancourt et retrouvailles avec l'interprète Cordonnier}\par\vspace{0.2mm}
\noindent{\fontsize{6.9pt}{8.0pt}\selectfont\textit{4 novembre 1916} \hfill \textit{Acheux, Bertrancourt, Bus-lès-Artois, Albert (Somme)}}\par\vspace{0.5mm}
\noindent\rule{\textwidth}{0.35pt}\par\vspace{0.8mm}
% --- Corps de la transcription diplomatique et apparat marginal en prolongement ---
{\renewcommand{\arraystretch}{0.93}%
\noindent\hspace*{3.5mm}\begin{tabular}{@{}p{87mm}@{\hspace{3.5mm}}p{39mm}@{}}
      \leavevmode\fontsize{7.3pt}{8.4pt}\selectfont\strut \underline{Sam. 4} \hfill{} \underline{Novbre /16}\strut & \smash{\parbox[t]{39mm}{\fontsize{6.2pt}{7.0pt}\selectfont\raggedright\color{black!88}\textbf{1 Sam.} \textit{scil.} Samedi.}} \\
      \leavevmode\fontsize{7.3pt}{8.4pt}\selectfont\strut Matin, dû me rendre à Bertrancourt.\strut &  \\
      \leavevmode\fontsize{7.3pt}{8.4pt}\selectfont\strut Venu à pied. Retour, s/ 1 lorry, 3 kil.\strut & \smash{\parbox[t]{39mm}{\fontsize{6.2pt}{7.0pt}\selectfont\raggedright\color{black!88}\textbf{3 s/} \textit{abbr. pro} sur. \textbf{3 kil.} \textit{abbr. pro} kilomètres.}} \\
      \leavevmode\fontsize{7.3pt}{8.4pt}\selectfont\strut Routes entt boueuses. Région plate.\strut &  \\
      \leavevmode\llap{\fontsize{6.4pt}{7.3pt}\selectfont\color{darkgray}5\hspace{1.8mm}}\fontsize{7.3pt}{8.4pt}\selectfont\strut Terres dévastées. Partout des tentes, des huttes,\strut &  \\
      \leavevmode\fontsize{7.3pt}{8.4pt}\selectfont\strut des voitures, des caissons.\strut &  \\
      \leavevmode\fontsize{7.3pt}{8.4pt}\selectfont\strut Recherché interprète Jamet, à propos question\strut &  \\
      \leavevmode\fontsize{7.3pt}{8.4pt}\selectfont\strut de pétrole. Venu aux Divon Headquarters,\strut & \smash{\parbox[t]{39mm}{\fontsize{6.2pt}{7.0pt}\selectfont\raggedright\color{black!88}\textbf{8 Divon} \textit{abbr. pro} Division.}} \\
      \leavevmode\fontsize{7.3pt}{8.4pt}\selectfont\strut tentes et huttes, à gauche route de Bus.\strut &  \\
      \leavevmode\llap{\fontsize{6.4pt}{7.3pt}\selectfont\color{darkgray}10\hspace{1.8mm}}\fontsize{7.3pt}{8.4pt}\selectfont\strut Au bureau mission fr. 2e division,\strut &  \\
      \leavevmode\fontsize{7.3pt}{8.4pt}\selectfont\strut retrouvé interprète Cordonnier, R.F.A., que\strut & \smash{\parbox[t]{39mm}{\fontsize{6.2pt}{7.0pt}\selectfont\raggedright\color{black!88}\textbf{11 R.F.A.} \textit{angl.} Royal Field Artillery.}} \\
      \leavevmode\fontsize{7.3pt}{8.4pt}\selectfont\strut connu à Béthune.\strut &  \\
      \leavevmode\fontsize{7.3pt}{8.4pt}\selectfont\strut Après lunch, sauté s/ 1 lorry et\strut & \smash{\parbox[t]{39mm}{\fontsize{6.2pt}{7.0pt}\selectfont\raggedright\color{black!88}\textbf{13 s/} \textit{abbr. pro} sur. \textbf{13 lunch} \textit{angl. in ms.} — \textbf{14 Coy} \textit{angl.} Company. — \textbf{15 Lt} \textit{abbr. pro} Lieutenant.}} \\
      \leavevmode\fontsize{7.3pt}{8.4pt}\selectfont\strut retourné à Albert, à 152 Coy, puis au\strut &  \\
      \leavevmode\llap{\fontsize{6.4pt}{7.3pt}\selectfont\color{darkgray}15\hspace{1.8mm}}\fontsize{7.3pt}{8.4pt}\selectfont\strut bureau mission fr. Vu Lt Pizzi.\strut &  \\
      \leavevmode\fontsize{7.3pt}{8.4pt}\selectfont\strut Ciel noir. Pluie menaçante. Bois\strut &  \\
      \leavevmode\fontsize{7.3pt}{8.4pt}\selectfont\strut aux feuillages roussis et mordorés par les\strut &  \\
      \leavevmode\fontsize{7.3pt}{8.4pt}\selectfont\strut gelées d'automne. Routes, affrt boueuses.\strut & \smash{\parbox[t]{39mm}{\fontsize{6.2pt}{7.0pt}\selectfont\raggedright\color{black!88}\textbf{18 affrt} \textit{abbr. pro} affreusement.}} \\
\end{tabular}}%
\par\vspace*{\fill}
% --- Registres de bas de page (1. Critique interne [si pertinent] ; 2. Loci paralleli en latin ; 3. Commentaire historique en retrait) ---
\noindent\rule{\textwidth}{0.5pt}\par\vspace{0.65mm}
{\fontsize{6.6pt}{7.7pt}\selectfont%
\noindent\normalfont\color{black!92} \textbf{2} (\textit{Bertrancourt}) cf. infra f\textsuperscript{o}~60~v\textsuperscript{o} (p.~\pageref{fol:121}), f\textsuperscript{o}~84~r\textsuperscript{o} (p.~\pageref{fol:168}). — \textbf{9} (\textit{Bus}) cf. infra f\textsuperscript{o}~77~r\textsuperscript{o} (p.~\pageref{fol:154}). — \textbf{11} (\textit{R.F}) cf. supra f\textsuperscript{o}~1~r\textsuperscript{o} (p.~\pageref{fol:2}), f\textsuperscript{o}~37~v\textsuperscript{o} (p.~\pageref{fol:75}).\par\vspace{0.55mm}
\noindent\rule{38mm}{0.4pt}\par\vspace{0.5mm}
\noindent\hspace*{5mm}\parbox{\dimexpr\textwidth-5mm\relax}{\fontsize{6.6pt}{7.7pt}\selectfont\color{black!84}Le samedi 4 novembre 1916, une course de service pour l'approvisionnement en pétrole (nécessaire aux moteurs des batteuses mécaniques) conduit Hippolyte Combes à \textbf{Bertrancourt} (Somme, à 4~km à l'est d'Acheux), au quartier général de la \textbf{2\textsuperscript{e} Division britannique} (\textit{2nd Division}, V\textsuperscript{e} Corps, qui s'apprête à attaquer à Beaumont-Hamel et dans le vallon de l'Ancre le 13 novembre).}\par}
\newpage
\phantomsection\label{fac:92}
% --- En-tête du fac-similé (Page impaire / droite : Photographie du manuscrit) ---
\noindent\rule{\textwidth}{0.5pt}\par\vspace{0.55mm}
\noindent\textbf{\fontsize{6.8pt}{7.8pt}\selectfont FAC-SIMILÉ EN VIS-À-VIS — FOLIO 46 RECTO} \hfill \textbf{\fontsize{6.8pt}{7.8pt}\selectfont NOVEMBRE 1916}\par\vspace{0.25mm}
\noindent{\fontsize{6.8pt}{8.0pt}\selectfont\textit{Livre-Journal d'Hippolyte Combes} — 4 novembre 1916 (Acheux, Bertrancourt, Bus-lès-Artois, Albert (Somme))}\par\vspace{0.55mm}
\noindent\rule{\textwidth}{0.35pt}\par\vspace*{\fill}
{\centering
  \setlength{\fboxsep}{1.2mm}%
  \setlength{\fboxrule}{0.4pt}%
  \fcolorbox{gray!55}{white}{%
    \includegraphics[width=0.95\textwidth,height=0.81\textheight,keepaspectratio]{\facsimiledir/20251126_104636.jpg}%
  }\par}
\vspace*{\fill}
\newpage

% =============================================================================
% VIS-A-VIS SPREAD : FOLIO 46 VERSO (20251126_104647.jpg)
% =============================================================================
\cleardoubleevenpage
\renewcommand{\currentmonthyear}{Novembre 1916}
\renewcommand{\currentfoliolabel}{Folio 46 verso}
\phantomsection\label{fol:93}
% --- En-tête du folio (Page paire / gauche : Transcription, apparat & registres de bas de page) ---
\noindent\rule{\textwidth}{0.5pt}\par\vspace{0.55mm}
\noindent\textbf{\footnotesize FOLIO 46 VERSO} \hfill \textbf{\footnotesize NOVEMBRE 1916}\par\vspace{0.25mm}
\noindent{\fontsize{6.9pt}{8.0pt}\selectfont Second cliché archivistique du feuillet des 2–3 novembre 1916 (doublon de numérisation)}\par\vspace{0.2mm}
\noindent{\fontsize{6.9pt}{8.0pt}\selectfont\textit{2 (soir) – 3 novembre 1916 [cliché bis]} \hfill \textit{Acheux, Raincheval, Albert (Somme)}}\par\vspace{0.5mm}
\noindent\rule{\textwidth}{0.35pt}\par\vspace{0.8mm}
% --- Corps de la transcription diplomatique et apparat marginal en prolongement ---
{\renewcommand{\arraystretch}{0.87}%
\noindent\hspace*{3.5mm}\begin{tabular}{@{}p{87mm}@{\hspace{3.5mm}}p{39mm}@{}}
      \leavevmode\fontsize{6.9pt}{7.9pt}\selectfont\strut Après le dîner, course en auto à Raincheval.\strut &  \\
      \leavevmode\fontsize{6.9pt}{7.9pt}\selectfont\strut Vu Major Levy, of the D.O.S.T.\strut & \smash{\parbox[t]{39mm}{\fontsize{6.2pt}{7.0pt}\selectfont\raggedright\color{black!88}\textbf{2 D.O.S.T.} \textit{sigle militaire, développement non établi}. — \textbf{4} \textit{in marg. cum unco} (comptes).}} \\
      \leavevmode\fontsize{6.9pt}{7.9pt}\selectfont\strut Nuit très noire. — Couché assez tard.\strut &  \\
      \leavevmode\fontsize{6.9pt}{7.9pt}\selectfont\strut [ Donné à Victor Duchâteau : 5 f —\strut &  \\
      \leavevmode\llap{\fontsize{6.4pt}{7.3pt}\selectfont\color{darkgray}5\hspace{1.8mm}}\fontsize{6.9pt}{7.9pt}\selectfont\strut Mess, octobre : 47 f 70\strut &  \\
      \leavevmode\fontsize{6.9pt}{7.9pt}\selectfont\strut Acheux : ctes ples : 15 c ]\strut & \smash{\parbox[t]{39mm}{\fontsize{6.2pt}{7.0pt}\selectfont\raggedright\color{black!88}\textbf{6 ctes ples} \textit{abbr. pro} cartes postales. — \textbf{7 Ve.} \textit{scil.} Vendredi.}} \\
      \cline{1-1}
      \leavevmode\fontsize{6.9pt}{7.9pt}\selectfont\strut \underline{3 Ve.} Matin, ciel couvert, 1 peu de soleil\strut &  \\
      \leavevmode\fontsize{6.9pt}{7.9pt}\selectfont\strut l'après-midi.\strut &  \\
      \leavevmode\fontsize{6.9pt}{7.9pt}\selectfont\strut Partis en auto avec Lt Curtis. Resté à Albert\strut & \smash{\parbox[t]{39mm}{\fontsize{6.2pt}{7.0pt}\selectfont\raggedright\color{black!88}\textbf{9 Lt} \textit{abbr. pro} Lieutenant.}} \\
      \leavevmode\llap{\fontsize{6.4pt}{7.3pt}\selectfont\color{darkgray}10\hspace{1.8mm}}\fontsize{6.9pt}{7.9pt}\selectfont\strut et revenu seul s/ 1 lorry. Routes\strut & \smash{\parbox[t]{39mm}{\fontsize{6.2pt}{7.0pt}\selectfont\raggedright\color{black!88}\textbf{10 s/} \textit{abbr. pro} sur. \textbf{10 lorry} \textit{angl. in ms.}}} \\
      \leavevmode\fontsize{6.9pt}{7.9pt}\selectfont\strut mauvaises et sillonnées de véhicules divers.\strut &  \\
      \leavevmode\fontsize{6.9pt}{7.9pt}\selectfont\strut Revenu au camp, 152 Coy et au\strut & \smash{\parbox[t]{39mm}{\fontsize{6.2pt}{7.0pt}\selectfont\raggedright\color{black!88}\textbf{12 Coy} \textit{angl.} Company.}} \\
      \leavevmode\fontsize{6.9pt}{7.9pt}\selectfont\strut bureau mission.\strut &  \\
      \leavevmode\fontsize{6.9pt}{7.9pt}\selectfont\strut Après lunch, 2h 30, Acheux.\strut & \smash{\parbox[t]{39mm}{\fontsize{6.2pt}{7.0pt}\selectfont\raggedright\color{black!88}\textbf{14 lunch} \textit{angl. in ms.}}} \\
      \leavevmode\llap{\fontsize{6.4pt}{7.3pt}\selectfont\color{darkgray}15\hspace{1.8mm}}\fontsize{6.9pt}{7.9pt}\selectfont\strut Passé journée ds ma chambre — au bureau.\strut & \smash{\parbox[t]{39mm}{\fontsize{6.2pt}{7.0pt}\selectfont\raggedright\color{black!88}\textbf{15 ds} \textit{abbr. pro} dans.}} \\
      \cline{1-1}
      \leavevmode\fontsize{6.9pt}{7.9pt}\selectfont\strut [ — Mission, Albert : 58 f 75 (30 f retenus pour\strut & \smash{\parbox[t]{39mm}{\fontsize{6.2pt}{7.0pt}\selectfont\raggedright\color{black!88}\textbf{16} \textit{in marg. cum unco} (comptes).}} \\
      \leavevmode\fontsize{6.9pt}{7.9pt}\selectfont\strut 6 jours de permission) ]\strut &  \\
\end{tabular}}%
\par\vspace*{\fill}
% --- Registres de bas de page (1. Critique interne [si pertinent] ; 2. Loci paralleli en latin ; 3. Commentaire historique en retrait) ---
\noindent\rule{\textwidth}{0.5pt}\par\vspace{0.65mm}
{\fontsize{6.6pt}{7.7pt}\selectfont%
\noindent\normalfont \textbf{Seconde prise de vue du même feuillet :} Ce cliché reproduit la même page manuscrite que le f\textsuperscript{o}~45~v\textsuperscript{o} (journées des 2 et 3 novembre 1916).\par\vspace{0.5mm}
\noindent\rule{22mm}{0.3pt}\par\vspace{0.45mm}
\noindent\normalfont\color{black!92} \textbf{1--17} (\textit{eadem pagina bis photographata}) cf. supra f\textsuperscript{o}~45~v\textsuperscript{o} (p.~\pageref{fol:91}). — \textbf{1} (\textit{Raincheval}) cf. supra f\textsuperscript{o}~20~v\textsuperscript{o} (p.~\pageref{fol:41}), f\textsuperscript{o}~45~v\textsuperscript{o} (p.~\pageref{fol:91}), infra f\textsuperscript{o}~48~v\textsuperscript{o} (p.~\pageref{fol:97}), f\textsuperscript{o}~49~r\textsuperscript{o} (p.~\pageref{fol:98}), \textit{et passim}. — \textbf{2} (\textit{Major Levy}) cf. supra f\textsuperscript{o}~45~v\textsuperscript{o} (p.~\pageref{fol:91}), infra f\textsuperscript{o}~55~r\textsuperscript{o} (p.~\pageref{fol:110}), f\textsuperscript{o}~58~r\textsuperscript{o} (p.~\pageref{fol:116}).\par\vspace{0.55mm}
\noindent\rule{38mm}{0.4pt}\par\vspace{0.5mm}
\noindent\hspace*{5mm}\parbox{\dimexpr\textwidth-5mm\relax}{\fontsize{6.6pt}{7.7pt}\selectfont\color{black!84}Seconde prise de vue du feuillet des 2 et 3 novembre 1916 : présentation d'Hippolyte Combes au \textbf{major Levy} à Raincheval et règlement de sa solde au bureau de la Mission française à \textbf{Albert} après six jours de retenue de permission.}\par}
\newpage
\phantomsection\label{fac:93}
% --- En-tête du fac-similé (Page impaire / droite : Photographie du manuscrit) ---
\noindent\rule{\textwidth}{0.5pt}\par\vspace{0.55mm}
\noindent\textbf{\fontsize{6.8pt}{7.8pt}\selectfont FAC-SIMILÉ EN VIS-À-VIS — FOLIO 46 VERSO} \hfill \textbf{\fontsize{6.8pt}{7.8pt}\selectfont NOVEMBRE 1916}\par\vspace{0.25mm}
\noindent{\fontsize{6.8pt}{8.0pt}\selectfont\textit{Livre-Journal d'Hippolyte Combes} — 2 (soir) – 3 novembre 1916 [cliché bis] (Acheux, Raincheval, Albert (Somme))}\par\vspace{0.55mm}
\noindent\rule{\textwidth}{0.35pt}\par\vspace*{\fill}
{\centering
  \setlength{\fboxsep}{1.2mm}%
  \setlength{\fboxrule}{0.4pt}%
  \fcolorbox{gray!55}{white}{%
    \includegraphics[width=0.95\textwidth,height=0.81\textheight,keepaspectratio]{\facsimiledir/20251126_104647.jpg}%
  }\par}
\vspace*{\fill}
\newpage

% =============================================================================
% VIS-A-VIS SPREAD : FOLIO 47 RECTO (20251126_104702.jpg)
% =============================================================================
\cleardoubleevenpage
\renewcommand{\currentmonthyear}{Novembre 1916}
\renewcommand{\currentfoliolabel}{Folio 47 recto}
\phantomsection\label{fol:94}
% --- En-tête du folio (Page paire / gauche : Transcription, apparat & registres de bas de page) ---
\noindent\rule{\textwidth}{0.5pt}\par\vspace{0.55mm}
\noindent\textbf{\footnotesize FOLIO 47 RECTO} \hfill \textbf{\footnotesize NOVEMBRE 1916}\par\vspace{0.25mm}
\noindent{\fontsize{6.9pt}{8.0pt}\selectfont Emménagement au n° 1 route de Léalvillers et convocation chez le capitaine Huguet}\par\vspace{0.2mm}
\noindent{\fontsize{6.9pt}{8.0pt}\selectfont\textit{5 novembre 1916} \hfill \textit{Acheux-en-Amiénois (Somme)}}\par\vspace{0.5mm}
\noindent\rule{\textwidth}{0.35pt}\par\vspace{0.8mm}
% --- Corps de la transcription diplomatique et apparat marginal en prolongement ---
{\renewcommand{\arraystretch}{0.93}%
\noindent\hspace*{3.5mm}\begin{tabular}{@{}p{87mm}@{\hspace{3.5mm}}p{39mm}@{}}
      \leavevmode\fontsize{7.3pt}{8.4pt}\selectfont\strut \underline{Dim. 5} \hfill{} \underline{Novbre /16}\strut & \smash{\parbox[t]{39mm}{\fontsize{6.2pt}{7.0pt}\selectfont\raggedright\color{black!88}\textbf{1 Dim.} \textit{scil.} Dimanche.}} \\
      \leavevmode\fontsize{7.3pt}{8.4pt}\selectfont\strut Matinée, soleil brille 2 ou 3 h. Puis de\strut &  \\
      \leavevmode\fontsize{7.3pt}{8.4pt}\selectfont\strut nouveau, gros nuages assombri ciel bas.\strut &  \\
      \leavevmode\fontsize{7.3pt}{8.4pt}\selectfont\strut Délogé de mon billet (59) par le\strut &  \\
      \leavevmode\llap{\fontsize{6.4pt}{7.3pt}\selectfont\color{darkgray}5\hspace{1.8mm}}\fontsize{7.3pt}{8.4pt}\selectfont\strut Town Major. Grâce à l'1 de ses clerks,\strut & \smash{\parbox[t]{39mm}{\fontsize{6.2pt}{7.0pt}\selectfont\raggedright\color{black!88}\textbf{5 Town Major} \textit{angl. in ms.}}} \\
      \leavevmode\fontsize{7.3pt}{8.4pt}\selectfont\strut trouvé bientôt 1 autre chambre — et\strut &  \\
      \leavevmode\fontsize{7.3pt}{8.4pt}\selectfont\strut préférable, no 1, route de Léalvillers.\strut &  \\
      \leavevmode\fontsize{7.3pt}{8.4pt}\selectfont\strut Avant midi, terminé déménagement et\strut &  \\
      \leavevmode\fontsize{7.3pt}{8.4pt}\selectfont\strut installation ds mon nouveau local.\strut & \smash{\parbox[t]{39mm}{\fontsize{6.2pt}{7.0pt}\selectfont\raggedright\color{black!88}\textbf{9 ds} \textit{abbr. pro} dans.}} \\
      \leavevmode\llap{\fontsize{6.4pt}{7.3pt}\selectfont\color{darkgray}10\hspace{1.8mm}}\fontsize{7.3pt}{8.4pt}\selectfont\strut Après-midi, liberté entière.\strut &  \\
      \leavevmode\fontsize{7.3pt}{8.4pt}\selectfont\strut Achevé de m'installer maison. Fait\strut &  \\
      \leavevmode\fontsize{7.3pt}{8.4pt}\selectfont\strut connaissance avec proprre, très sympathique.\strut & \smash{\parbox[t]{39mm}{\fontsize{6.2pt}{7.0pt}\selectfont\raggedright\color{black!88}\textbf{12 proprre} \textit{abbr. pro} propriétaire.}} \\
      \leavevmode\fontsize{7.3pt}{8.4pt}\selectfont\strut A 6h soir, l'interprète de la Vth\strut &  \\
      \leavevmode\fontsize{7.3pt}{8.4pt}\selectfont\strut Army troops, venu me chercher : Cap.\strut & \smash{\parbox[t]{39mm}{\fontsize{6.2pt}{7.0pt}\selectfont\raggedright\color{black!88}\textbf{14 Cap.} \textit{abbr. pro} Capitaine. — \textbf{15} \textit{lect. incert. in ms.}}} \\
      \leavevmode\llap{\fontsize{6.4pt}{7.3pt}\selectfont\color{darkgray}15\hspace{1.8mm}}\fontsize{7.3pt}{8.4pt}\selectfont\strut Huguet, du [ajout interl. illis.] corps, désire me voir.\strut &  \\
      \leavevmode\fontsize{7.3pt}{8.4pt}\selectfont\strut Venu aussitôt bureau mission du corps.\strut &  \\
      \leavevmode\fontsize{7.3pt}{8.4pt}\selectfont\strut Très bien reçu. — qq froid,\strut & \smash{\parbox[t]{39mm}{\fontsize{6.2pt}{7.0pt}\selectfont\raggedright\color{black!88}\textbf{17 qq} \textit{abbr. pro} quelque(s).}} \\
      \leavevmode\fontsize{7.3pt}{8.4pt}\selectfont\strut — demande de renseignements divers.\strut &  \\
      \leavevmode\fontsize{7.3pt}{8.4pt}\selectfont\strut Dès 8h 1/2, ds ma chambre.\strut & \smash{\parbox[t]{39mm}{\fontsize{6.2pt}{7.0pt}\selectfont\raggedright\color{black!88}\textbf{19 ds} \textit{abbr. pro} dans. \textbf{19 1/2} \textit{scil.} demi / moitié.}} \\
\end{tabular}}%
\par\vspace*{\fill}
% --- Registres de bas de page (1. Critique interne [si pertinent] ; 2. Loci paralleli en latin ; 3. Commentaire historique en retrait) ---
\noindent\rule{\textwidth}{0.5pt}\par\vspace{0.65mm}
{\fontsize{6.6pt}{7.7pt}\selectfont%
\noindent\normalfont\color{black!92} \textbf{7} (\textit{Léalvillers}) cf. supra f\textsuperscript{o}~24~r\textsuperscript{o} (p.~\pageref{fol:48}), infra f\textsuperscript{o}~47~v\textsuperscript{o} (p.~\pageref{fol:95}), f\textsuperscript{o}~48~r\textsuperscript{o} (p.~\pageref{fol:96}), \textit{et passim}. — \textbf{15} (\textit{Huguet}) cf. infra f\textsuperscript{o}~61~v\textsuperscript{o} (p.~\pageref{fol:123}).\par\vspace{0.55mm}
\noindent\rule{38mm}{0.4pt}\par\vspace{0.5mm}
\noindent\hspace*{5mm}\parbox{\dimexpr\textwidth-5mm\relax}{\fontsize{6.6pt}{7.7pt}\selectfont\color{black!84}Le dimanche 5 novembre 1916, Hippolyte Combes emménage au \textbf{n\textsuperscript{o}~1 de la route de Léalvillers} à \textbf{Acheux-en-Amiénois}, chez le tailleur \textbf{M.~Carluvant} (nommé aux f\textsuperscript{o}~51~r\textsuperscript{o}, f\textsuperscript{o}~59~r\textsuperscript{o}, f\textsuperscript{o}~73~r\textsuperscript{o}--f\textsuperscript{o}~74~r\textsuperscript{o}), où il logera jusqu'au 7 février 1917.}\par}
\newpage
\phantomsection\label{fac:94}
% --- En-tête du fac-similé (Page impaire / droite : Photographie du manuscrit) ---
\noindent\rule{\textwidth}{0.5pt}\par\vspace{0.55mm}
\noindent\textbf{\fontsize{6.8pt}{7.8pt}\selectfont FAC-SIMILÉ EN VIS-À-VIS — FOLIO 47 RECTO} \hfill \textbf{\fontsize{6.8pt}{7.8pt}\selectfont NOVEMBRE 1916}\par\vspace{0.25mm}
\noindent{\fontsize{6.8pt}{8.0pt}\selectfont\textit{Livre-Journal d'Hippolyte Combes} — 5 novembre 1916 (Acheux-en-Amiénois (Somme))}\par\vspace{0.55mm}
\noindent\rule{\textwidth}{0.35pt}\par\vspace*{\fill}
{\centering
  \setlength{\fboxsep}{1.2mm}%
  \setlength{\fboxrule}{0.4pt}%
  \fcolorbox{gray!55}{white}{%
    \includegraphics[width=0.95\textwidth,height=0.81\textheight,keepaspectratio]{\facsimiledir/20251126_104702.jpg}%
  }\par}
\vspace*{\fill}
\newpage

% =============================================================================
% VIS-A-VIS SPREAD : FOLIO 47 VERSO (20251126_104714.jpg)
% =============================================================================
\cleardoubleevenpage
\renewcommand{\currentmonthyear}{Novembre 1916}
\renewcommand{\currentfoliolabel}{Folio 47 verso}
\phantomsection\label{fol:95}
% --- En-tête du folio (Page paire / gauche : Transcription, apparat & registres de bas de page) ---
\noindent\rule{\textwidth}{0.5pt}\par\vspace{0.55mm}
\noindent\textbf{\footnotesize FOLIO 47 VERSO} \hfill \textbf{\footnotesize NOVEMBRE 1916}\par\vspace{0.25mm}
\noindent{\fontsize{6.9pt}{8.0pt}\selectfont Courses à Belle-Église et Coisy, nostalgie de la 152e Compagnie et convois vers la ligne de feu}\par\vspace{0.2mm}
\noindent{\fontsize{6.9pt}{8.0pt}\selectfont\textit{6 – 7 novembre 1916} \hfill \textit{Acheux, Arquèves (Belle-Église), Hérissart, Rubempré, Coisy (Somme)}}\par\vspace{0.5mm}
\noindent\rule{\textwidth}{0.35pt}\par\vspace{0.8mm}
% --- Corps de la transcription diplomatique et apparat marginal en prolongement ---
{\renewcommand{\arraystretch}{0.83}%
\noindent\hspace*{3.5mm}\begin{tabular}{@{}p{87mm}@{\hspace{3.5mm}}p{39mm}@{}}
      \leavevmode\fontsize{6.6pt}{7.5pt}\selectfont\strut \underline{Lundi 6} \hfill{} \underline{Nov. /16}\strut &  \\
      \leavevmode\fontsize{6.6pt}{7.5pt}\selectfont\strut Tt le jour, rafales \& averses.\strut & \smash{\parbox[t]{39mm}{\fontsize{6.2pt}{7.0pt}\selectfont\raggedright\color{black!88}\textbf{2 Tt} \textit{abbr. pro} Tout.}} \\
      \leavevmode\fontsize{6.6pt}{7.5pt}\selectfont\strut Matinée, venu à pied \& s/ 1 lorry j'à Arquèves,\strut & \smash{\parbox[t]{39mm}{\fontsize{6.2pt}{7.0pt}\selectfont\raggedright\color{black!88}\textbf{3 j'à} \textit{abbr. pro} jusqu'à. \textbf{3 s/} \textit{abbr. pro} sur.}} \\
      \leavevmode\fontsize{6.6pt}{7.5pt}\selectfont\strut ferme de Belle-Église, pour 1 question de pétrole.\strut &  \\
      \leavevmode\llap{\fontsize{6.4pt}{7.3pt}\selectfont\color{darkgray}5\hspace{1.8mm}}\fontsize{6.6pt}{7.5pt}\selectfont\strut Après-midi, longue randonnée en auto\strut &  \\
      \leavevmode\fontsize{6.6pt}{7.5pt}\selectfont\strut à Coisy, par Hérissart, Rubempré etc.\strut &  \\
      \leavevmode\fontsize{6.6pt}{7.5pt}\selectfont\strut Achat de paille et foin. Partis vers 1h 1/2,\strut & \smash{\parbox[t]{39mm}{\fontsize{6.2pt}{7.0pt}\selectfont\raggedright\color{black!88}\textbf{7 1/2} \textit{scil.} demi / moitié.}} \\
      \leavevmode\fontsize{6.6pt}{7.5pt}\selectfont\strut rentrés seulement sur 6 h. Giboulées\strut &  \\
      \leavevmode\fontsize{6.6pt}{7.5pt}\selectfont\strut tour à tour, région de plaines et de petites\strut &  \\
      \leavevmode\llap{\fontsize{6.4pt}{7.3pt}\selectfont\color{darkgray}10\hspace{1.8mm}}\fontsize{6.6pt}{7.5pt}\selectfont\strut collines, avec bois dont les arbres, teintés\strut &  \\
      \leavevmode\fontsize{6.6pt}{7.5pt}\selectfont\strut de rouille, orange et or — et allées\strut &  \\
      \leavevmode\fontsize{6.6pt}{7.5pt}\selectfont\strut jaunies par jonchées de feuilles mortes.\strut &  \\
      \leavevmode\fontsize{6.6pt}{7.5pt}\selectfont\strut Regrette d'avoir quitté 152 Coy. Cherche\strut &  \\
      \leavevmode\fontsize{6.6pt}{7.5pt}\selectfont\strut moyen de ne pas quitter 18e Divon \& ma Coy.\strut & \smash{\parbox[t]{39mm}{\fontsize{6.2pt}{7.0pt}\selectfont\raggedright\color{black!88}\textbf{14 Divon} \textit{abbr. pro} Division.}} \\
      \cline{1-1}
      \leavevmode\llap{\fontsize{6.4pt}{7.3pt}\selectfont\color{darkgray}15\hspace{1.8mm}}\fontsize{6.6pt}{7.5pt}\selectfont\strut \underline{Mardi 7}\strut &  \\
      \leavevmode\fontsize{6.6pt}{7.5pt}\selectfont\strut Pluie durant partie journée.\strut &  \\
      \leavevmode\fontsize{6.6pt}{7.5pt}\selectfont\strut Routes boueuses pleines de trous \& ornières profondes.\strut &  \\
      \leavevmode\fontsize{6.6pt}{7.5pt}\selectfont\strut Ni matin, ni soir, rien à faire.\strut &  \\
      \leavevmode\fontsize{6.6pt}{7.5pt}\selectfont\strut Après-midi, essayé d'aller Albert : fatigué\strut &  \\
      \leavevmode\llap{\fontsize{6.4pt}{7.3pt}\selectfont\color{darkgray}20\hspace{1.8mm}}\fontsize{6.6pt}{7.5pt}\selectfont\strut d'attendre lorry, renoncé partir et rentré billet.\strut & \smash{\parbox[t]{39mm}{\fontsize{6.2pt}{7.0pt}\selectfont\raggedright\color{black!88}\textbf{20 lorry} \textit{angl. in ms.}}} \\
      \leavevmode\fontsize{6.6pt}{7.5pt}\selectfont\strut De la fenêtre de ma chambre, mouvement\strut &  \\
      \leavevmode\fontsize{6.6pt}{7.5pt}\selectfont\strut rue (route de Léalvillers). Flot presque ininterrompu\strut &  \\
      \leavevmode\fontsize{6.6pt}{7.5pt}\selectfont\strut de véhicules de ttes sortes : lorries, wagons, caissons\strut & \smash{\parbox[t]{39mm}{\fontsize{6.2pt}{7.0pt}\selectfont\raggedright\color{black!88}\textbf{23 ttes} \textit{abbr. pro} toutes. \textbf{23 lorries} \textit{angl. in ms.}}} \\
      \leavevmode\fontsize{6.6pt}{7.5pt}\selectfont\strut de munitions d'artrie. Et défilé incessant\strut &  \\
      \leavevmode\llap{\fontsize{6.4pt}{7.3pt}\selectfont\color{darkgray}25\hspace{1.8mm}}\fontsize{6.6pt}{7.5pt}\selectfont\strut de troupes allant au repos ou se rendant firing line.\strut &  \\
\end{tabular}}%
\par\vspace*{\fill}
% --- Registres de bas de page (1. Critique interne [si pertinent] ; 2. Loci paralleli en latin ; 3. Commentaire historique en retrait) ---
\noindent\rule{\textwidth}{0.5pt}\par\vspace{0.65mm}
{\fontsize{6.6pt}{7.7pt}\selectfont%
\noindent\normalfont\color{black!92} \textbf{3} (\textit{Arquèves}) cf. infra f\textsuperscript{o}~53~r\textsuperscript{o} (p.~\pageref{fol:106}), f\textsuperscript{o}~54~r\textsuperscript{o} (p.~\pageref{fol:108}), \textit{et passim}. — \textbf{4} (\textit{Belle-Église}) cf. infra f\textsuperscript{o}~53~r\textsuperscript{o} (p.~\pageref{fol:106}), f\textsuperscript{o}~54~r\textsuperscript{o} (p.~\pageref{fol:108}), \textit{et passim}. — \textbf{6} (\textit{Hérissart}) cf. infra f\textsuperscript{o}~49~v\textsuperscript{o} (p.~\pageref{fol:99}), f\textsuperscript{o}~51~v\textsuperscript{o} (p.~\pageref{fol:103}), \textit{et passim}.\par\vspace{0.55mm}
\noindent\rule{38mm}{0.4pt}\par\vspace{0.5mm}
\noindent\hspace*{5mm}\parbox{\dimexpr\textwidth-5mm\relax}{\fontsize{6.6pt}{7.7pt}\selectfont\color{black!84}Première mention le 6 novembre 1916 de la \textbf{ferme de Belle-Église}, grand domaine agricole isolé de la commune d'\textbf{Arquèves} (Somme, ancienne commanderie templière puis hospitalière de Belle-Église), exploitée par la famille belge \textbf{Houthouan} (f\textsuperscript{o}~54~r\textsuperscript{o} et f\textsuperscript{o}~85~v\textsuperscript{o}), où fonctionnent les batteuses à pétrole du service des fourrages.}\par}
\newpage
\phantomsection\label{fac:95}
% --- En-tête du fac-similé (Page impaire / droite : Photographie du manuscrit) ---
\noindent\rule{\textwidth}{0.5pt}\par\vspace{0.55mm}
\noindent\textbf{\fontsize{6.8pt}{7.8pt}\selectfont FAC-SIMILÉ EN VIS-À-VIS — FOLIO 47 VERSO} \hfill \textbf{\fontsize{6.8pt}{7.8pt}\selectfont NOVEMBRE 1916}\par\vspace{0.25mm}
\noindent{\fontsize{6.8pt}{8.0pt}\selectfont\textit{Livre-Journal d'Hippolyte Combes} — 6 – 7 novembre 1916 (Acheux, Arquèves (Belle-Église), Hérissart, Rubempré, Coisy (Somme))}\par\vspace{0.55mm}
\noindent\rule{\textwidth}{0.35pt}\par\vspace*{\fill}
{\centering
  \setlength{\fboxsep}{1.2mm}%
  \setlength{\fboxrule}{0.4pt}%
  \fcolorbox{gray!55}{white}{%
    \includegraphics[width=0.95\textwidth,height=0.81\textheight,keepaspectratio]{\facsimiledir/20251126_104714.jpg}%
  }\par}
\vspace*{\fill}
\newpage

% =============================================================================
% VIS-A-VIS SPREAD : FOLIO 48 RECTO (20251126_104728.jpg)
% =============================================================================
\cleardoubleevenpage
\renewcommand{\currentmonthyear}{Novembre 1916}
\renewcommand{\currentfoliolabel}{Folio 48 recto}
\phantomsection\label{fol:96}
% --- En-tête du folio (Page paire / gauche : Transcription, apparat & registres de bas de page) ---
\noindent\rule{\textwidth}{0.5pt}\par\vspace{0.55mm}
\noindent\textbf{\footnotesize FOLIO 48 RECTO} \hfill \textbf{\footnotesize NOVEMBRE 1916}\par\vspace{0.25mm}
\noindent{\fontsize{6.9pt}{8.0pt}\selectfont Promenades en auto à Forceville, Varennes et thé au mess de la 152e Compagnie à Albert}\par\vspace{0.2mm}
\noindent{\fontsize{6.9pt}{8.0pt}\selectfont\textit{8 – 9 novembre 1916} \hfill \textit{Acheux, Forceville, Varennes, Warloy, Senlis-le-Sec, Albert (Somme)}}\par\vspace{0.5mm}
\noindent\rule{\textwidth}{0.35pt}\par\vspace{0.8mm}
% --- Corps de la transcription diplomatique et apparat marginal en prolongement ---
{\renewcommand{\arraystretch}{0.87}%
\noindent\hspace*{3.5mm}\begin{tabular}{@{}p{87mm}@{\hspace{3.5mm}}p{39mm}@{}}
      \leavevmode\fontsize{6.9pt}{7.9pt}\selectfont\strut \underline{8 Me.} \hfill{} \underline{Novbre 1916}\strut &  \\
      \leavevmode\fontsize{6.9pt}{7.9pt}\selectfont\strut Matin, course en\strut &  \\
      \leavevmode\fontsize{6.9pt}{7.9pt}\selectfont\strut auto à Forceville — avec Major Hughes \& Curtis.\strut &  \\
      \leavevmode\fontsize{6.9pt}{7.9pt}\selectfont\strut Après lunch, venu en auto à Albert — en\strut & \smash{\parbox[t]{39mm}{\fontsize{6.2pt}{7.0pt}\selectfont\raggedright\color{black!88}\textbf{4 lunch} \textit{angl. in ms.}}} \\
      \leavevmode\llap{\fontsize{6.4pt}{7.3pt}\selectfont\color{darkgray}5\hspace{1.8mm}}\fontsize{6.9pt}{7.9pt}\selectfont\strut passant par Varennes – Warloy – Senlis.\strut &  \\
      \leavevmode\fontsize{6.9pt}{7.9pt}\selectfont\strut Soleil, brillé 1 peu. Délicieuse promenade.\strut &  \\
      \leavevmode\fontsize{6.9pt}{7.9pt}\selectfont\strut Venu bureau mission frse et demandé\strut & \smash{\parbox[t]{39mm}{\fontsize{6.2pt}{7.0pt}\selectfont\raggedright\color{black!88}\textbf{7 frse} \textit{abbr. pro} française.}} \\
      \leavevmode\fontsize{6.9pt}{7.9pt}\selectfont\strut renseignements.\strut &  \\
      \leavevmode\fontsize{6.9pt}{7.9pt}\selectfont\strut Retourné s/ 1 lorry.\strut & \smash{\parbox[t]{39mm}{\fontsize{6.2pt}{7.0pt}\selectfont\raggedright\color{black!88}\textbf{9 s/} \textit{abbr. pro} sur. \textbf{9 lorry} \textit{angl. in ms.} — \textbf{10 J.} \textit{scil.} Jeudi. — \textbf{11 1/2} \textit{scil.} demi / moitié. — \textbf{12 qq} \textit{abbr. pro} quelque(s).}} \\
      \cline{1-1}
      \leavevmode\llap{\fontsize{6.4pt}{7.3pt}\selectfont\color{darkgray}10\hspace{1.8mm}}\fontsize{6.9pt}{7.9pt}\selectfont\strut \underline{9 J.} Matin \& soir, temps couvert.\strut &  \\
      \leavevmode\fontsize{6.9pt}{7.9pt}\selectfont\strut 1/2 jour, enfin, soleil brillé — durant\strut &  \\
      \leavevmode\fontsize{6.9pt}{7.9pt}\selectfont\strut qq heures.\strut &  \\
      \leavevmode\fontsize{6.9pt}{7.9pt}\selectfont\strut Après-midi, course en auto à Varennes,\strut &  \\
      \leavevmode\fontsize{6.9pt}{7.9pt}\selectfont\strut par Léalvillers — et dû revenir par Acheux,\strut &  \\
      \leavevmode\llap{\fontsize{6.4pt}{7.3pt}\selectfont\color{darkgray}15\hspace{1.8mm}}\fontsize{6.9pt}{7.9pt}\selectfont\strut puis continuer j. à Albert.\strut & \smash{\parbox[t]{39mm}{\fontsize{6.2pt}{7.0pt}\selectfont\raggedright\color{black!88}\textbf{15 j. à} \textit{abbr. pro} jusqu'à.}} \\
      \leavevmode\fontsize{6.9pt}{7.9pt}\selectfont\strut Campagne ensoleillée — mais routes\strut &  \\
      \leavevmode\fontsize{6.9pt}{7.9pt}\selectfont\strut affreust boueuses.\strut &  \\
      \leavevmode\fontsize{6.9pt}{7.9pt}\selectfont\strut Resté pour le thé au mess-hut\strut & \smash{\parbox[t]{39mm}{\fontsize{6.2pt}{7.0pt}\selectfont\raggedright\color{black!88}\textbf{18 mess} \textit{angl. in ms.}}} \\
      \leavevmode\fontsize{6.9pt}{7.9pt}\selectfont\strut de la 152 Coy, s/ plateau au sol\strut & \smash{\parbox[t]{39mm}{\fontsize{6.2pt}{7.0pt}\selectfont\raggedright\color{black!88}\textbf{19 s/} \textit{abbr. pro} sur. \textbf{19 Coy} \textit{angl.} Company.}} \\
      \leavevmode\llap{\fontsize{6.4pt}{7.3pt}\selectfont\color{darkgray}20\hspace{1.8mm}}\fontsize{6.9pt}{7.9pt}\selectfont\strut transformé en marécage. Alors, ciel\strut &  \\
      \leavevmode\fontsize{6.9pt}{7.9pt}\selectfont\strut éclairci, qq étoiles, pleine lune magnifique.\strut & \smash{\parbox[t]{39mm}{\fontsize{6.2pt}{7.0pt}\selectfont\raggedright\color{black!88}\textbf{21 qq} \textit{abbr. pro} quelque(s).}} \\
      \leavevmode\fontsize{6.9pt}{7.9pt}\selectfont\strut Bombardement commence sur Albert.\strut &  \\
\end{tabular}}%
\par\vspace*{\fill}
% --- Registres de bas de page (1. Critique interne [si pertinent] ; 2. Loci paralleli en latin ; 3. Commentaire historique en retrait) ---
\noindent\rule{\textwidth}{0.5pt}\par\vspace{0.65mm}
{\fontsize{6.6pt}{7.7pt}\selectfont%
\noindent\normalfont\color{black!92} \textbf{3} (\textit{Curtis}) cf. supra f\textsuperscript{o}~1~r\textsuperscript{o} (p.~\pageref{fol:2}), f\textsuperscript{o}~46~v\textsuperscript{o} (p.~\pageref{fol:93}), infra f\textsuperscript{o}~48~v\textsuperscript{o} (p.~\pageref{fol:97}), f\textsuperscript{o}~50~v\textsuperscript{o} (p.~\pageref{fol:101}), \textit{et passim}. — \textbf{3} (\textit{Forceville}) cf. supra f\textsuperscript{o}~22~v\textsuperscript{o} (p.~\pageref{fol:45}), f\textsuperscript{o}~25~r\textsuperscript{o} (p.~\pageref{fol:50}). — \textbf{3} (\textit{Major Hughes}) cf. supra f\textsuperscript{o}~45~r\textsuperscript{o} (p.~\pageref{fol:90}).\par\vspace{0.55mm}
\noindent\rule{38mm}{0.4pt}\par\vspace{0.5mm}
\noindent\hspace*{5mm}\parbox{\dimexpr\textwidth-5mm\relax}{\fontsize{6.6pt}{7.7pt}\selectfont\color{black!84}Après une tournée à Forceville le 8 novembre avec le \textbf{major Hughes} (du \textit{18th Divisional Train}, f\textsuperscript{o}~45~r\textsuperscript{o}) et le lieutenant Curtis, Hippolyte Combes profite de ses courses des 8 et 9 novembre vers Varennes et \textbf{Albert} pour prendre le thé à la baraque-mess (\textit{mess-hut}) de son ancienne unité, la \textbf{152\textsuperscript{e} Compagnie A.S.C.}, toujours campée dans la boue du plateau d'Albert, tandis que le retour du temps clair sous la pleine lune du 9 novembre s'accompagne d'un bombardement sur Albert.}\par}
\newpage
\phantomsection\label{fac:96}
% --- En-tête du fac-similé (Page impaire / droite : Photographie du manuscrit) ---
\noindent\rule{\textwidth}{0.5pt}\par\vspace{0.55mm}
\noindent\textbf{\fontsize{6.8pt}{7.8pt}\selectfont FAC-SIMILÉ EN VIS-À-VIS — FOLIO 48 RECTO} \hfill \textbf{\fontsize{6.8pt}{7.8pt}\selectfont NOVEMBRE 1916}\par\vspace{0.25mm}
\noindent{\fontsize{6.8pt}{8.0pt}\selectfont\textit{Livre-Journal d'Hippolyte Combes} — 8 – 9 novembre 1916 (Acheux, Forceville, Varennes, Warloy, Senlis-le-Sec, Albert (Somme))}\par\vspace{0.55mm}
\noindent\rule{\textwidth}{0.35pt}\par\vspace*{\fill}
{\centering
  \setlength{\fboxsep}{1.2mm}%
  \setlength{\fboxrule}{0.4pt}%
  \fcolorbox{gray!55}{white}{%
    \includegraphics[width=0.95\textwidth,height=0.81\textheight,keepaspectratio]{\facsimiledir/20251126_104728.jpg}%
  }\par}
\vspace*{\fill}
\newpage

% =============================================================================
% VIS-A-VIS SPREAD : FOLIO 48 VERSO (20251126_104739.jpg)
% =============================================================================
\cleardoubleevenpage
\renewcommand{\currentmonthyear}{Novembre 1916}
\renewcommand{\currentfoliolabel}{Folio 48 verso}
\phantomsection\label{fol:97}
% --- En-tête du folio (Page paire / gauche : Transcription, apparat & registres de bas de page) ---
\noindent\rule{\textwidth}{0.5pt}\par\vspace{0.55mm}
\noindent\textbf{\footnotesize FOLIO 48 VERSO} \hfill \textbf{\footnotesize NOVEMBRE 1916}\par\vspace{0.25mm}
\noindent{\fontsize{6.9pt}{8.0pt}\selectfont Courses à Raincheval et Varennes, rencontre de l'interprète Sorgues et bombardement d'Acheux}\par\vspace{0.2mm}
\noindent{\fontsize{6.9pt}{8.0pt}\selectfont\textit{10 – 11 novembre 1916} \hfill \textit{Acheux, Raincheval, Léalvillers, Varennes (Somme)}}\par\vspace{0.5mm}
\noindent\rule{\textwidth}{0.35pt}\par\vspace{0.8mm}
% --- Corps de la transcription diplomatique et apparat marginal en prolongement ---
{\renewcommand{\arraystretch}{0.87}%
\noindent\hspace*{3.5mm}\begin{tabular}{@{}p{87mm}@{\hspace{3.5mm}}p{39mm}@{}}
      \leavevmode\fontsize{6.9pt}{7.9pt}\selectfont\strut \underline{Ven. 10} \hfill{} \underline{Novembre 1916}\strut & \smash{\parbox[t]{39mm}{\fontsize{6.2pt}{7.0pt}\selectfont\raggedright\color{black!88}\textbf{1 Ven.} \textit{scil.} Vendredi.}} \\
      \leavevmode\fontsize{6.9pt}{7.9pt}\selectfont\strut Très belle journée. Chaud\strut &  \\
      \leavevmode\fontsize{6.9pt}{7.9pt}\selectfont\strut soleil. Enfin!\strut &  \\
      \leavevmode\fontsize{6.9pt}{7.9pt}\selectfont\strut Matinée, au bureau, en l'absence du Lt Curtis\strut & \smash{\parbox[t]{39mm}{\fontsize{6.2pt}{7.0pt}\selectfont\raggedright\color{black!88}\textbf{4 Lt} \textit{abbr. pro} Lieutenant.}} \\
      \leavevmode\llap{\fontsize{6.4pt}{7.3pt}\selectfont\color{darkgray}5\hspace{1.8mm}}\fontsize{6.9pt}{7.9pt}\selectfont\strut \& du clerk. Lu, écrit, réfléchi, devant bon feu.\strut &  \\
      \leavevmode\fontsize{6.9pt}{7.9pt}\selectfont\strut Impression toute nouvelle.\strut &  \\
      \leavevmode\fontsize{6.9pt}{7.9pt}\selectfont\strut Tout l'après-midi, courses en auto : Raincheval\strut &  \\
      \leavevmode\fontsize{6.9pt}{7.9pt}\selectfont\strut et Léalvillers. Pu regarder à loisir campagne\strut &  \\
      \leavevmode\fontsize{6.9pt}{7.9pt}\selectfont\strut ensoleillée — et remarqué les détails du paysage.\strut &  \\
      \leavevmode\llap{\fontsize{6.4pt}{7.3pt}\selectfont\color{darkgray}10\hspace{1.8mm}}\fontsize{6.9pt}{7.9pt}\selectfont\strut Nuit extt belle, ciel éclairé par pleine lune.\strut & \smash{\parbox[t]{39mm}{\fontsize{6.2pt}{7.0pt}\selectfont\raggedright\color{black!88}\textbf{10 extt} \textit{abbr. pro} extrêmement. — \textbf{11 Sam.} \textit{scil.} Samedi.}} \\
      \cline{1-1}
      \leavevmode\fontsize{6.9pt}{7.9pt}\selectfont\strut \underline{Sam. 11} — Matin, dû me rendre à Varennes pour\strut &  \\
      \leavevmode\fontsize{6.9pt}{7.9pt}\selectfont\strut 1 question de paille à acheter. Aller \& retour à pied —\strut &  \\
      \leavevmode\fontsize{6.9pt}{7.9pt}\selectfont\strut Pour trouver, [addition interlinéaire illis.] Batt[suscrit illis.] F.A., B Battery, marché\strut & \smash{\parbox[t]{39mm}{\fontsize{6.2pt}{7.0pt}\selectfont\raggedright\color{black!88}\textbf{13} \textit{lect. incert. in ms.}}} \\
      \leavevmode\fontsize{6.9pt}{7.9pt}\selectfont\strut terrain mou, vrai cloaque où les pieds\strut &  \\
      \leavevmode\llap{\fontsize{6.4pt}{7.3pt}\selectfont\color{darkgray}15\hspace{1.8mm}}\fontsize{6.9pt}{7.9pt}\selectfont\strut enfonçaient et se collaient. Affreux.\strut &  \\
      \leavevmode\fontsize{6.9pt}{7.9pt}\selectfont\strut Interprète de cette unité, M. Sorgues [?], prof.\strut & \smash{\parbox[t]{39mm}{\fontsize{6.2pt}{7.0pt}\selectfont\raggedright\color{black!88}\textbf{16 Sorgues} \textit{lect. incert. in ms.}}} \\
      \leavevmode\fontsize{6.9pt}{7.9pt}\selectfont\strut au lycée de Toulon.\strut &  \\
      \leavevmode\fontsize{6.9pt}{7.9pt}\selectfont\strut Soir, à repas, sifflement obus. Acheux\strut &  \\
      \leavevmode\fontsize{6.9pt}{7.9pt}\selectfont\strut bombardé à intervalles rapprochés. Réfugiés cave.\strut &  \\
      \cline{1-1}
      \leavevmode\llap{\fontsize{6.4pt}{7.3pt}\selectfont\color{darkgray}20\hspace{1.8mm}}\fontsize{6.9pt}{7.9pt}\selectfont\strut [ — Ctes ples, à Varennes : 75 c — ]\strut & \smash{\parbox[t]{39mm}{\fontsize{6.2pt}{7.0pt}\selectfont\raggedright\color{black!88}\textbf{20} \textit{in marg. cum unco} (comptes). \textbf{20 Ctes ples} \textit{abbr. pro} Cartes postales.}} \\
\end{tabular}}%
\par\vspace*{\fill}
% --- Registres de bas de page (1. Critique interne [si pertinent] ; 2. Loci paralleli en latin ; 3. Commentaire historique en retrait) ---
\noindent\rule{\textwidth}{0.5pt}\par\vspace{0.65mm}
{\fontsize{6.6pt}{7.7pt}\selectfont%
\noindent\normalfont\color{black!92} \textbf{4} (\textit{Curtis}) cf. supra f\textsuperscript{o}~1~r\textsuperscript{o} (p.~\pageref{fol:2}), f\textsuperscript{o}~48~r\textsuperscript{o} (p.~\pageref{fol:96}), infra f\textsuperscript{o}~50~v\textsuperscript{o} (p.~\pageref{fol:101}), f\textsuperscript{o}~52~r\textsuperscript{o} (p.~\pageref{fol:104}), \textit{et passim}. — \textbf{7} (\textit{Raincheval}) cf. supra f\textsuperscript{o}~20~v\textsuperscript{o} (p.~\pageref{fol:41}), f\textsuperscript{o}~46~v\textsuperscript{o} (p.~\pageref{fol:93}), infra f\textsuperscript{o}~49~r\textsuperscript{o} (p.~\pageref{fol:98}), f\textsuperscript{o}~49~v\textsuperscript{o} (p.~\pageref{fol:99}), \textit{et passim}. — \textbf{8} (\textit{Léalvillers}) cf. supra f\textsuperscript{o}~24~r\textsuperscript{o} (p.~\pageref{fol:48}), f\textsuperscript{o}~48~r\textsuperscript{o} (p.~\pageref{fol:96}), infra f\textsuperscript{o}~50~v\textsuperscript{o} (p.~\pageref{fol:101}), f\textsuperscript{o}~56~r\textsuperscript{o} (p.~\pageref{fol:112}), \textit{et passim}.\par\vspace{0.55mm}
\noindent\rule{38mm}{0.4pt}\par\vspace{0.5mm}
\noindent\hspace*{5mm}\parbox{\dimexpr\textwidth-5mm\relax}{\fontsize{6.6pt}{7.7pt}\selectfont\color{black!84}Le samedi 11 novembre 1916, en allant régler à pied un achat de paille auprès de la \textit{B~Battery} d'une brigade de la \textit{Royal Field Artillery} (\textit{R.F.A.}) à \textbf{Varennes-en-Croix}, Hippolyte Combes rencontre un collègue universitaire mobilisé comme interprète, professeur au lycée de Toulon (ville où Hippolyte Combes avait lui-même été incorporé au 10\textsuperscript{e}~R.A.P. en mars--avril 1915).}\par}
\newpage
\phantomsection\label{fac:97}
% --- En-tête du fac-similé (Page impaire / droite : Photographie du manuscrit) ---
\noindent\rule{\textwidth}{0.5pt}\par\vspace{0.55mm}
\noindent\textbf{\fontsize{6.8pt}{7.8pt}\selectfont FAC-SIMILÉ EN VIS-À-VIS — FOLIO 48 VERSO} \hfill \textbf{\fontsize{6.8pt}{7.8pt}\selectfont NOVEMBRE 1916}\par\vspace{0.25mm}
\noindent{\fontsize{6.8pt}{8.0pt}\selectfont\textit{Livre-Journal d'Hippolyte Combes} — 10 – 11 novembre 1916 (Acheux, Raincheval, Léalvillers, Varennes (Somme))}\par\vspace{0.55mm}
\noindent\rule{\textwidth}{0.35pt}\par\vspace*{\fill}
{\centering
  \setlength{\fboxsep}{1.2mm}%
  \setlength{\fboxrule}{0.4pt}%
  \fcolorbox{gray!55}{white}{%
    \includegraphics[width=0.95\textwidth,height=0.81\textheight,keepaspectratio]{\facsimiledir/20251126_104739.jpg}%
  }\par}
\vspace*{\fill}
\newpage

% =============================================================================
% VIS-A-VIS SPREAD : FOLIO 49 RECTO (20251126_104801.jpg)
% =============================================================================
\cleardoubleevenpage
\renewcommand{\currentmonthyear}{Novembre 1916}
\renewcommand{\currentfoliolabel}{Folio 49 recto}
\phantomsection\label{fol:98}
% --- En-tête du folio (Page paire / gauche : Transcription, apparat & registres de bas de page) ---
\noindent\rule{\textwidth}{0.5pt}\par\vspace{0.55mm}
\noindent\textbf{\footnotesize FOLIO 49 RECTO} \hfill \textbf{\footnotesize NOVEMBRE 1916}\par\vspace{0.25mm}
\noindent{\fontsize{6.9pt}{8.0pt}\selectfont Bataille de Beaumont-Hamel (13–14 novembre), convois de blessés et foire de Doullens}\par\vspace{0.2mm}
\noindent{\fontsize{6.9pt}{8.0pt}\selectfont\textit{12 – 14 novembre 1916} \hfill \textit{Acheux, Varennes, Contay, Beauquesne, Raincheval, Doullens, Frévent (Somme / Pas-de-Calais)}}\par\vspace{0.5mm}
\noindent\rule{\textwidth}{0.35pt}\par\vspace{0.8mm}
% --- Corps de la transcription diplomatique et apparat marginal en prolongement ---
{\renewcommand{\arraystretch}{0.83}%
\noindent\hspace*{3.5mm}\begin{tabular}{@{}p{87mm}@{\hspace{3.5mm}}p{39mm}@{}}
      \leavevmode\fontsize{6.6pt}{7.5pt}\selectfont\strut \underline{Dim. 12.} \hfill{} \underline{Novembre 1916}\strut & \smash{\parbox[t]{39mm}{\fontsize{6.2pt}{7.0pt}\selectfont\raggedright\color{black!88}\textbf{1 Dim.} \textit{scil.} Dimanche.}} \\
      \leavevmode\fontsize{6.6pt}{7.5pt}\selectfont\strut Assez belle journée. Matin \& soir, brouillard.\strut &  \\
      \leavevmode\fontsize{6.6pt}{7.5pt}\selectfont\strut Routes mauvaises. Sol mou, pétri de sillons\strut &  \\
      \leavevmode\fontsize{6.6pt}{7.5pt}\selectfont\strut profonds par les innombrables véhicules qui\strut &  \\
      \leavevmode\llap{\fontsize{6.4pt}{7.3pt}\selectfont\color{darkgray}5\hspace{1.8mm}}\fontsize{6.6pt}{7.5pt}\selectfont\strut le labourent.\strut &  \\
      \leavevmode\fontsize{6.6pt}{7.5pt}\selectfont\strut Matin, liberté complète.\strut &  \\
      \leavevmode\fontsize{6.6pt}{7.5pt}\selectfont\strut Après-midi, course en auto à Beauquesne\strut &  \\
      \leavevmode\fontsize{6.6pt}{7.5pt}\selectfont\strut par Varennes \& Contay. Brume s'est levée\strut &  \\
      \leavevmode\fontsize{6.6pt}{7.5pt}\selectfont\strut peu à peu. Paysage anglais : arbres étriqués,\strut &  \\
      \leavevmode\llap{\fontsize{6.4pt}{7.3pt}\selectfont\color{darkgray}10\hspace{1.8mm}}\fontsize{6.6pt}{7.5pt}\selectfont\strut feuillages jaunes, marron semant la route.\strut &  \\
      \leavevmode\fontsize{6.6pt}{7.5pt}\selectfont\strut [ Journal semaine : 40 c\strut & \smash{\parbox[t]{39mm}{\fontsize{6.2pt}{7.0pt}\selectfont\raggedright\color{black!88}\textbf{11} \textit{in marg. cum unco} (comptes).}} \\
      \leavevmode\fontsize{6.6pt}{7.5pt}\selectfont\strut Casque peint khaki : 1 f — ]\strut &  \\
      \cline{1-1}
      \leavevmode\fontsize{6.6pt}{7.5pt}\selectfont\strut \underline{14 Ma.} — Tout hier, et principt matin, ainsi\strut & \smash{\parbox[t]{39mm}{\fontsize{6.2pt}{7.0pt}\selectfont\raggedright\color{black!88}\textbf{13 Ma.} \textit{scil.} Mardi. \textbf{13 principt} \textit{abbr. pro} principalement. — \textbf{14 auj'd'hui} \textit{abbr. pro} aujourd'hui. \textbf{14 bombardemt} \textit{abbr. pro} bombardement. — \textbf{16 s/} \textit{abbr. pro} sur. — \textbf{18 1/2} \textit{scil.} demi / moitié. \textbf{18 m.} \textit{abbr. pro} matin. — \textbf{19 tt} \textit{abbr. pro} tout. — \textbf{21 passim} \textit{lat.} çà et là. — \textbf{22} \textit{lect. incert. in ms.} — \textbf{23 gds} \textit{abbr. pro} grands. — \textbf{24 Cartes ples} \textit{abbr. pro} Cartes postales. — \textbf{25 fr.} \textit{abbr. pro} franc(s).}} \\
      \leavevmode\fontsize{6.6pt}{7.5pt}\selectfont\strut qu'auj'd'hui, très fort bombardemt — du\strut &  \\
      \leavevmode\llap{\fontsize{6.4pt}{7.3pt}\selectfont\color{darkgray}15\hspace{1.8mm}}\fontsize{6.6pt}{7.5pt}\selectfont\strut côté de Beaumont. Vu passer nbx blessés\strut &  \\
      \leavevmode\fontsize{6.6pt}{7.5pt}\selectfont\strut transportés s/ fourgons, autos \& cars —\strut &  \\
      \leavevmode\fontsize{6.6pt}{7.5pt}\selectfont\strut Courses à Raincheval, Doullens et Frévent.\strut &  \\
      \leavevmode\fontsize{6.6pt}{7.5pt}\selectfont\strut Partis 8h 1/2 m., rentrés 5h soir.\strut &  \\
      \leavevmode\fontsize{6.6pt}{7.5pt}\selectfont\strut Ciel nuageux — pluie menacé tt le jour.\strut &  \\
      \cline{1-1}
      \leavevmode\llap{\fontsize{6.4pt}{7.3pt}\selectfont\color{darkgray}20\hspace{1.8mm}}\fontsize{6.6pt}{7.5pt}\selectfont\strut \underline{A Doullens}, grande foire.\strut &  \\
      \leavevmode\fontsize{6.6pt}{7.5pt}\selectfont\strut A Frévent, promené passim. — Petite\strut &  \\
      \leavevmode\fontsize{6.6pt}{7.5pt}\selectfont\strut ville assez coquette : belle église, [carillonne ?] 1/4 d'h, usine,\strut &  \\
      \leavevmode\fontsize{6.6pt}{7.5pt}\selectfont\strut — Boutons simili-cuir : 4 f 75 (10 gds 0.40, 3 petits 0.25 —\strut &  \\
      \cline{1-1}
      \leavevmode\fontsize{6.6pt}{7.5pt}\selectfont\strut Cartes ples : 40 c, 1 f 50 et 10 c\strut &  \\
      \leavevmode\llap{\fontsize{6.4pt}{7.3pt}\selectfont\color{darkgray}25\hspace{1.8mm}}\fontsize{6.6pt}{7.5pt}\selectfont\strut Déjeuner à l'hôtel d'Amiens, ma part : 7 fr. ]\strut &  \\
\end{tabular}}%
\par\vspace*{\fill}
% --- Registres de bas de page (1. Critique interne [si pertinent] ; 2. Loci paralleli en latin ; 3. Commentaire historique en retrait) ---
\noindent\rule{\textwidth}{0.5pt}\par\vspace{0.65mm}
{\fontsize{6.6pt}{7.7pt}\selectfont%
\noindent\normalfont\color{black!92} \textbf{7} (\textit{Beauquesne}) cf. supra f\textsuperscript{o}~37~v\textsuperscript{o} (p.~\pageref{fol:75}), infra f\textsuperscript{o}~56~r\textsuperscript{o} (p.~\pageref{fol:112}), f\textsuperscript{o}~56~v\textsuperscript{o} (p.~\pageref{fol:113}), \textit{et passim}. — \textbf{8} (\textit{Varennes}) cf. supra f\textsuperscript{o}~25~r\textsuperscript{o} (p.~\pageref{fol:50}), f\textsuperscript{o}~48~v\textsuperscript{o} (p.~\pageref{fol:97}), infra f\textsuperscript{o}~60~v\textsuperscript{o} (p.~\pageref{fol:121}), f\textsuperscript{o}~64~r\textsuperscript{o} (p.~\pageref{fol:128}), \textit{et passim}. — \textbf{8} (\textit{Contay}) cf. infra f\textsuperscript{o}~50~v\textsuperscript{o} (p.~\pageref{fol:101}), f\textsuperscript{o}~79~v\textsuperscript{o} (p.~\pageref{fol:159}), \textit{et passim}.\par\vspace{0.55mm}
\noindent\rule{38mm}{0.4pt}\par\vspace{0.5mm}
\noindent\hspace*{5mm}\parbox{\dimexpr\textwidth-5mm\relax}{\fontsize{6.6pt}{7.7pt}\selectfont\color{black!84}Témoignage oculaire et auditif direct sur le déclenchement de la \textbf{bataille de l'Ancre} (\textit{Battle of the Ancre}, 13--18 novembre 1916): le lundi \textbf{13 novembre 1916} à 5~h~45 du matin et le mardi \textbf{14 novembre}, le \textbf{V\textsuperscript{e} Corps britannique} (notamment la 51\textsuperscript{e} \textit{Highland Division} et la 63\textsuperscript{e} \textit{Royal Naval Division}) prend d'assaut la position fortifiée de \textbf{Beaumont-Hamel} et \textbf{Beaucourt-sur-l'Ancre} (à 9~km à l'est d'Acheux), tandis que les convois d'ambulances évacuant les blessés vers les hôpitaux et trains sanitaires traversent Acheux.}\par}
\newpage
\phantomsection\label{fac:98}
% --- En-tête du fac-similé (Page impaire / droite : Photographie du manuscrit) ---
\noindent\rule{\textwidth}{0.5pt}\par\vspace{0.55mm}
\noindent\textbf{\fontsize{6.8pt}{7.8pt}\selectfont FAC-SIMILÉ EN VIS-À-VIS — FOLIO 49 RECTO} \hfill \textbf{\fontsize{6.8pt}{7.8pt}\selectfont NOVEMBRE 1916}\par\vspace{0.25mm}
\noindent{\fontsize{6.8pt}{8.0pt}\selectfont\textit{Livre-Journal d'Hippolyte Combes} — 12 – 14 novembre 1916 (Acheux, Varennes, Contay, Beauquesne, Raincheval, Doullens, Frévent (Somme / Pas-de-Calais))}\par\vspace{0.55mm}
\noindent\rule{\textwidth}{0.35pt}\par\vspace*{\fill}
{\centering
  \setlength{\fboxsep}{1.2mm}%
  \setlength{\fboxrule}{0.4pt}%
  \fcolorbox{gray!55}{white}{%
    \includegraphics[width=0.95\textwidth,height=0.81\textheight,keepaspectratio]{\facsimiledir/20251126_104801.jpg}%
  }\par}
\vspace*{\fill}
\newpage

% =============================================================================
% VIS-A-VIS SPREAD : FOLIO 49 VERSO (20251126_104814.jpg)
% =============================================================================
\cleardoubleevenpage
\renewcommand{\currentmonthyear}{Novembre 1916}
\renewcommand{\currentfoliolabel}{Folio 49 verso}
\phantomsection\label{fol:99}
% --- En-tête du folio (Page paire / gauche : Transcription, apparat & registres de bas de page) ---
\noindent\rule{\textwidth}{0.5pt}\par\vspace{0.55mm}
\noindent\textbf{\footnotesize FOLIO 49 VERSO} \hfill \textbf{\footnotesize NOVEMBRE 1916}\par\vspace{0.25mm}
\noindent{\fontsize{6.9pt}{8.0pt}\selectfont Canonnade sur Beaucourt-sur-l'Ancre et déjeuner à Hérissart chez le lieutenant Walsh}\par\vspace{0.2mm}
\noindent{\fontsize{6.9pt}{8.0pt}\selectfont\textit{15 – 16 novembre 1916} \hfill \textit{Acheux, Albert, Raincheval, Louvencourt, Toutencourt, Hérissart (Somme)}}\par\vspace{0.5mm}
\noindent\rule{\textwidth}{0.35pt}\par\vspace{0.8mm}
% --- Corps de la transcription diplomatique et apparat marginal en prolongement ---
{\renewcommand{\arraystretch}{0.87}%
\noindent\hspace*{3.5mm}\begin{tabular}{@{}p{87mm}@{\hspace{3.5mm}}p{39mm}@{}}
      \leavevmode\fontsize{6.9pt}{7.9pt}\selectfont\strut \underline{15 Me.} \hfill{} \underline{Novbre 1916}\strut &  \\
      \leavevmode\fontsize{6.9pt}{7.9pt}\selectfont\strut Toute la nuit, encore,\strut &  \\
      \leavevmode\fontsize{6.9pt}{7.9pt}\selectfont\strut et durant journée, principt à l'aube, très\strut & \smash{\parbox[t]{39mm}{\fontsize{6.2pt}{7.0pt}\selectfont\raggedright\color{black!88}\textbf{3 principt} \textit{abbr. pro} principalement. — \textbf{4 s/} \textit{abbr. pro} sur.}} \\
      \leavevmode\fontsize{6.9pt}{7.9pt}\selectfont\strut violente canonnade vers Beaumont s/ Ancre.\strut &  \\
      \leavevmode\llap{\fontsize{6.4pt}{7.3pt}\selectfont\color{darkgray}5\hspace{1.8mm}}\fontsize{6.9pt}{7.9pt}\selectfont\strut Très belle journée ensoleillée.\strut &  \\
      \leavevmode\fontsize{6.9pt}{7.9pt}\selectfont\strut Matin, venu en auto à Albert. Rentré de =\strut & \smash{\parbox[t]{39mm}{\fontsize{6.2pt}{7.0pt}\selectfont\raggedright\color{black!88}\textbf{6 de =} \textit{signum aequalitatis pro} de même. — \textbf{7 frse} \textit{abbr. pro} française.}} \\
      \leavevmode\fontsize{6.9pt}{7.9pt}\selectfont\strut Passé au bureau mission frse. Rien de\strut &  \\
      \leavevmode\fontsize{6.9pt}{7.9pt}\selectfont\strut nouveau me concernant.\strut &  \\
      \leavevmode\fontsize{6.9pt}{7.9pt}\selectfont\strut Après-midi, liberté entière. Promené\strut &  \\
      \leavevmode\llap{\fontsize{6.4pt}{7.3pt}\selectfont\color{darkgray}10\hspace{1.8mm}}\fontsize{6.9pt}{7.9pt}\selectfont\strut village —\strut &  \\
      \leavevmode\fontsize{6.9pt}{7.9pt}\selectfont\strut [ Cartes ples : Acheux : 35 c. — Journaux : 40 c\strut & \smash{\parbox[t]{39mm}{\fontsize{6.2pt}{7.0pt}\selectfont\raggedright\color{black!88}\textbf{11} \textit{in marg. cum unco} (comptes). \textbf{11 Cartes ples} \textit{abbr. pro} Cartes postales. — \textbf{12 Washing} \textit{angl. in ms.} — \textbf{13 J.} \textit{scil.} Jeudi. \textbf{13 bombardemt} \textit{abbr. pro} bombardement. — \textbf{17 fortt} \textit{abbr. pro} fortement.}} \\
      \leavevmode\fontsize{6.9pt}{7.9pt}\selectfont\strut Washing : 65 c — ]\strut &  \\
      \cline{1-1}
      \leavevmode\fontsize{6.9pt}{7.9pt}\selectfont\strut \underline{16 J.} — Matinée, toujours très violent bombardemt.\strut &  \\
      \leavevmode\fontsize{6.9pt}{7.9pt}\selectfont\strut Vitres et fenêtres ne cessent pas de\strut &  \\
      \leavevmode\llap{\fontsize{6.4pt}{7.3pt}\selectfont\color{darkgray}15\hspace{1.8mm}}\fontsize{6.9pt}{7.9pt}\selectfont\strut trembler.\strut &  \\
      \leavevmode\fontsize{6.9pt}{7.9pt}\selectfont\strut Soleil brillé journée entière. Gelé\strut &  \\
      \leavevmode\fontsize{6.9pt}{7.9pt}\selectfont\strut fortt durant nuit. Bise froide.\strut &  \\
      \leavevmode\fontsize{6.9pt}{7.9pt}\selectfont\strut Matin \& après-midi, en auto :\strut &  \\
      \leavevmode\fontsize{6.9pt}{7.9pt}\selectfont\strut Raincheval, Louvencourt, Toutencourt etc.\strut &  \\
      \leavevmode\llap{\fontsize{6.4pt}{7.3pt}\selectfont\color{darkgray}20\hspace{1.8mm}}\fontsize{6.9pt}{7.9pt}\selectfont\strut Déjeuné à Hérissart ch. Lt Walsh et\strut & \smash{\parbox[t]{39mm}{\fontsize{6.2pt}{7.0pt}\selectfont\raggedright\color{black!88}\textbf{20 ch.} \textit{abbr. pro} chez. \textbf{20 Lt} \textit{abbr. pro} Lieutenant.}} \\
      \leavevmode\fontsize{6.9pt}{7.9pt}\selectfont\strut avec son interprète. Fort apprécié ce dernier\strut &  \\
      \leavevmode\fontsize{6.9pt}{7.9pt}\selectfont\strut M. [blanc], directeur d'1 usine\strut & \smash{\parbox[t]{39mm}{\fontsize{6.2pt}{7.0pt}\selectfont\raggedright\color{black!88}\textbf{22} \textit{spatium relictum in ms.}}} \\
      \leavevmode\fontsize{6.9pt}{7.9pt}\selectfont\strut (tissus) près Paris. Intelligent \& distingué.\strut &  \\
\end{tabular}}%
\par\vspace*{\fill}
% --- Registres de bas de page (1. Critique interne [si pertinent] ; 2. Loci paralleli en latin ; 3. Commentaire historique en retrait) ---
\noindent\rule{\textwidth}{0.5pt}\par\vspace{0.65mm}
{\fontsize{6.6pt}{7.7pt}\selectfont%
\noindent\normalfont\color{black!92} \textbf{6} (\textit{Albert}) cf. supra f\textsuperscript{o}~39~r\textsuperscript{o} (p.~\pageref{fol:78}), f\textsuperscript{o}~48~r\textsuperscript{o} (p.~\pageref{fol:96}), infra f\textsuperscript{o}~50~r\textsuperscript{o} (p.~\pageref{fol:100}), f\textsuperscript{o}~82~r\textsuperscript{o} (p.~\pageref{fol:164}), \textit{et passim}. — \textbf{11} (\textit{Acheux}) cf. supra f\textsuperscript{o}~21~v\textsuperscript{o} (p.~\pageref{fol:43}), f\textsuperscript{o}~49~r\textsuperscript{o} (p.~\pageref{fol:98}), infra f\textsuperscript{o}~50~r\textsuperscript{o} (p.~\pageref{fol:100}), f\textsuperscript{o}~50~v\textsuperscript{o} (p.~\pageref{fol:101}), \textit{et passim}. — \textbf{19} (\textit{Raincheval}) cf. supra f\textsuperscript{o}~20~v\textsuperscript{o} (p.~\pageref{fol:41}), f\textsuperscript{o}~49~r\textsuperscript{o} (p.~\pageref{fol:98}), infra f\textsuperscript{o}~52~r\textsuperscript{o} (p.~\pageref{fol:104}), f\textsuperscript{o}~53~v\textsuperscript{o} (p.~\pageref{fol:107}), \textit{et passim}.\par\vspace{0.55mm}
\noindent\rule{38mm}{0.4pt}\par\vspace{0.5mm}
\noindent\hspace*{5mm}\parbox{\dimexpr\textwidth-5mm\relax}{\fontsize{6.6pt}{7.7pt}\selectfont\color{black!84}Les mercredi 15 et jeudi 16 novembre 1916, le roulement des tirs d'artillerie qui fait trembler les vitres d'Acheux répond à la poursuite des combats de la \textbf{bataille de l'Ancre} dans le secteur de \textbf{Beaumont-Hamel} et de \textbf{Beaucourt-sur-l'Ancre} (qu'Hippolyte Combes condense sous l'appellation \textit{Beaumont s/ Ancre}).}\par}
\newpage
\phantomsection\label{fac:99}
% --- En-tête du fac-similé (Page impaire / droite : Photographie du manuscrit) ---
\noindent\rule{\textwidth}{0.5pt}\par\vspace{0.55mm}
\noindent\textbf{\fontsize{6.8pt}{7.8pt}\selectfont FAC-SIMILÉ EN VIS-À-VIS — FOLIO 49 VERSO} \hfill \textbf{\fontsize{6.8pt}{7.8pt}\selectfont NOVEMBRE 1916}\par\vspace{0.25mm}
\noindent{\fontsize{6.8pt}{8.0pt}\selectfont\textit{Livre-Journal d'Hippolyte Combes} — 15 – 16 novembre 1916 (Acheux, Albert, Raincheval, Louvencourt, Toutencourt, Hérissart (Somme))}\par\vspace{0.55mm}
\noindent\rule{\textwidth}{0.35pt}\par\vspace*{\fill}
{\centering
  \setlength{\fboxsep}{1.2mm}%
  \setlength{\fboxrule}{0.4pt}%
  \fcolorbox{gray!55}{white}{%
    \includegraphics[width=0.95\textwidth,height=0.81\textheight,keepaspectratio]{\facsimiledir/20251126_104814.jpg}%
  }\par}
\vspace*{\fill}
\newpage

% =============================================================================
% VIS-A-VIS SPREAD : FOLIO 50 RECTO (20251126_104831.jpg)
% =============================================================================
\cleardoubleevenpage
\renewcommand{\currentmonthyear}{Novembre 1916}
\renewcommand{\currentfoliolabel}{Folio 50 recto}
\phantomsection\label{fol:100}
% --- En-tête du folio (Page paire / gauche : Transcription, apparat & registres de bas de page) ---
\noindent\rule{\textwidth}{0.5pt}\par\vspace{0.55mm}
\noindent\textbf{\footnotesize FOLIO 50 RECTO} \hfill \textbf{\footnotesize NOVEMBRE 1916}\par\vspace{0.25mm}
\noindent{\fontsize{6.9pt}{8.0pt}\selectfont Premières fortes gelées, conférence sur la Russie au Y.M.C.A. et première neige}\par\vspace{0.2mm}
\noindent{\fontsize{6.9pt}{8.0pt}\selectfont\textit{17 – 18 novembre 1916} \hfill \textit{Acheux, Authie, Orville, Amplier, Albert (Somme / Pas-de-Calais)}}\par\vspace{0.5mm}
\noindent\rule{\textwidth}{0.35pt}\par\vspace{0.8mm}
% --- Corps de la transcription diplomatique et apparat marginal en prolongement ---
{\renewcommand{\arraystretch}{0.83}%
\noindent\hspace*{3.5mm}\begin{tabular}{@{}p{87mm}@{\hspace{3.5mm}}p{39mm}@{}}
      \leavevmode\fontsize{6.6pt}{7.5pt}\selectfont\strut \underline{17 Ven.} \hfill{} \underline{Novbre -16}\strut & \smash{\parbox[t]{39mm}{\fontsize{6.2pt}{7.0pt}\selectfont\raggedright\color{black!88}\textbf{1 Ven.} \textit{scil.} Vendredi.}} \\
      \leavevmode\fontsize{6.6pt}{7.5pt}\selectfont\strut Nuit gelé très\strut &  \\
      \leavevmode\fontsize{6.6pt}{7.5pt}\selectfont\strut fortt. Sol, durci. Vitres, couvertes d'arabesques.\strut & \smash{\parbox[t]{39mm}{\fontsize{6.2pt}{7.0pt}\selectfont\raggedright\color{black!88}\textbf{3 fortt} \textit{abbr. pro} fortement. — \textbf{4 1/2} \textit{scil.} demi / moitié.}} \\
      \leavevmode\fontsize{6.6pt}{7.5pt}\selectfont\strut 1/2 jour, soleil brillé, mais bise glaciale\strut &  \\
      \leavevmode\llap{\fontsize{6.4pt}{7.3pt}\selectfont\color{darkgray}5\hspace{1.8mm}}\fontsize{6.6pt}{7.5pt}\selectfont\strut soufflé et soulevé routes, nuages de poussière.\strut &  \\
      \leavevmode\fontsize{6.6pt}{7.5pt}\selectfont\strut Toute matinée, roulé en auto. Authie,\strut &  \\
      \leavevmode\fontsize{6.6pt}{7.5pt}\selectfont\strut coquet bourg au fond d'1 vallon —\strut &  \\
      \leavevmode\fontsize{6.6pt}{7.5pt}\selectfont\strut Orville et Amplier. — Sites assez jolis.\strut &  \\
      \leavevmode\fontsize{6.6pt}{7.5pt}\selectfont\strut Après-midi, liberté complète.\strut &  \\
      \leavevmode\llap{\fontsize{6.4pt}{7.3pt}\selectfont\color{darkgray}10\hspace{1.8mm}}\fontsize{6.6pt}{7.5pt}\selectfont\strut A 6h s., venu hut de Y.M.C.A. et entendu\strut & \smash{\parbox[t]{39mm}{\fontsize{6.2pt}{7.0pt}\selectfont\raggedright\color{black!88}\textbf{10 Y.M.C.A.} \textit{angl.} Young Men's Christian Assoc. \textbf{10 s.} \textit{abbr. pro} soir. — \textbf{11 s/} \textit{abbr. pro} sur. — \textbf{12 ts} \textit{abbr. pro} tous. — \textbf{13 ds} \textit{abbr. pro} dans. — \textbf{14 Sam.} \textit{scil.} Samedi. \textbf{14 tte} \textit{abbr. pro} toute. — \textbf{16 ds} \textit{abbr. pro} dans.}} \\
      \leavevmode\fontsize{6.6pt}{7.5pt}\selectfont\strut conférence s/ Russia — —\strut &  \\
      \leavevmode\fontsize{6.6pt}{7.5pt}\selectfont\strut Comme ts les soirs précédents, nuit noire.\strut &  \\
      \leavevmode\fontsize{6.6pt}{7.5pt}\selectfont\strut Bombardement très violent ds le lointain.\strut &  \\
      \cline{1-1}
      \leavevmode\fontsize{6.6pt}{7.5pt}\selectfont\strut \underline{18 Sam.} Durant tte nuit \& + fort encore\strut &  \\
      \leavevmode\llap{\fontsize{6.4pt}{7.3pt}\selectfont\color{darkgray}15\hspace{1.8mm}}\fontsize{6.6pt}{7.5pt}\selectfont\strut aux 1ères heures jour, très violente\strut &  \\
      \leavevmode\fontsize{6.6pt}{7.5pt}\selectfont\strut canonnade, ds le lointain. Fenêtres \& portes,\strut &  \\
      \leavevmode\fontsize{6.6pt}{7.5pt}\selectfont\strut vibrant sans arrêt. Obus éclatent à coups\strut &  \\
      \leavevmode\fontsize{6.6pt}{7.5pt}\selectfont\strut de + en + précipités : on dirait énormes vagues\strut &  \\
      \leavevmode\fontsize{6.6pt}{7.5pt}\selectfont\strut déferlant s/ rochers avec 1 bruit de tonnerre.\strut & \smash{\parbox[t]{39mm}{\fontsize{6.2pt}{7.0pt}\selectfont\raggedright\color{black!88}\textbf{19 s/} \textit{abbr. pro} sur.}} \\
      \leavevmode\llap{\fontsize{6.4pt}{7.3pt}\selectfont\color{darkgray}20\hspace{1.8mm}}\fontsize{6.6pt}{7.5pt}\selectfont\strut Matinée, venu à Albert, aller \&\strut &  \\
      \leavevmode\fontsize{6.6pt}{7.5pt}\selectfont\strut retour, sur 1 lorry — —\strut & \smash{\parbox[t]{39mm}{\fontsize{6.2pt}{7.0pt}\selectfont\raggedright\color{black!88}\textbf{21 lorry} \textit{angl. in ms.}}} \\
      \leavevmode\fontsize{6.6pt}{7.5pt}\selectfont\strut Neige tombé matin — 1ère neige,\strut & \smash{\parbox[t]{39mm}{\fontsize{6.2pt}{7.0pt}\selectfont\raggedright\color{black!88}\textbf{22 1ère} \textit{scil.} première.}} \\
      \leavevmode\fontsize{6.6pt}{7.5pt}\selectfont\strut mais peu à peu, fondue sous pluie fine.\strut &  \\
      \leavevmode\fontsize{6.6pt}{7.5pt}\selectfont\strut Routes détrempées, vraies fondrières.\strut &  \\
      \cline{1-1}
      \leavevmode\llap{\fontsize{6.4pt}{7.3pt}\selectfont\color{darkgray}25\hspace{1.8mm}}\fontsize{6.6pt}{7.5pt}\selectfont\strut [ — Albert : solde : 88 f — 3 ctes ples : 30 c ]\strut & \smash{\parbox[t]{39mm}{\fontsize{6.2pt}{7.0pt}\selectfont\raggedright\color{black!88}\textbf{25} \textit{in marg. cum unco} (comptes). \textbf{25 ctes ples} \textit{abbr. pro} cartes postales.}} \\
\end{tabular}}%
\par\vspace*{\fill}
% --- Registres de bas de page (1. Critique interne [si pertinent] ; 2. Loci paralleli en latin ; 3. Commentaire historique en retrait) ---
\noindent\rule{\textwidth}{0.5pt}\par\vspace{0.65mm}
{\fontsize{6.6pt}{7.7pt}\selectfont%
\noindent\normalfont\color{black!92} \textbf{6} (\textit{Authie}) cf. infra f\textsuperscript{o}~51~r\textsuperscript{o} (p.~\pageref{fol:102}), f\textsuperscript{o}~53~v\textsuperscript{o} (p.~\pageref{fol:107}), \textit{et passim}. — \textbf{20} (\textit{Albert}) cf. supra f\textsuperscript{o}~39~r\textsuperscript{o} (p.~\pageref{fol:78}), f\textsuperscript{o}~49~v\textsuperscript{o} (p.~\pageref{fol:99}), infra f\textsuperscript{o}~82~r\textsuperscript{o} (p.~\pageref{fol:164}), f\textsuperscript{o}~84~v\textsuperscript{o} (p.~\pageref{fol:169}), \textit{et passim}.\par\vspace{0.55mm}
\noindent\rule{38mm}{0.4pt}\par\vspace{0.5mm}
\noindent\hspace*{5mm}\parbox{\dimexpr\textwidth-5mm\relax}{\fontsize{6.6pt}{7.7pt}\selectfont\color{black!84}Après une tournée dans la vallée de l'Authie (\textbf{Authie}, \textbf{Orville}, \textbf{Amplier}) et une conférence sur la Russie au foyer de la \textit{Y.M.C.A.} d'Acheux le 17 novembre, l'entrée du samedi \textbf{18 novembre 1916} consigne avec une exactitude remarquable le dernier grand assaut de l'\textbf{offensive de la Somme}: à l'aube du 18 novembre (6~h~10), au milieu de la première chute de neige de l'hiver, les \textbf{II\textsuperscript{e} et V\textsuperscript{e} Corps britanniques} (dont la 18\textsuperscript{e} Division devant \textit{Desire Trench}) attaquent vers Grandcourt et l'Ancre sous un barrage roulant d'artillerie qu'Hippolyte Combes compare à d'« énormes vagues déferlant sur rochers avec un bruit de tonnerre » (Nichols, \textit{op.~cit.}, p.~136--145).}\par}
\newpage
\phantomsection\label{fac:100}
% --- En-tête du fac-similé (Page impaire / droite : Photographie du manuscrit) ---
\noindent\rule{\textwidth}{0.5pt}\par\vspace{0.55mm}
\noindent\textbf{\fontsize{6.8pt}{7.8pt}\selectfont FAC-SIMILÉ EN VIS-À-VIS — FOLIO 50 RECTO} \hfill \textbf{\fontsize{6.8pt}{7.8pt}\selectfont NOVEMBRE 1916}\par\vspace{0.25mm}
\noindent{\fontsize{6.8pt}{8.0pt}\selectfont\textit{Livre-Journal d'Hippolyte Combes} — 17 – 18 novembre 1916 (Acheux, Authie, Orville, Amplier, Albert (Somme / Pas-de-Calais))}\par\vspace{0.55mm}
\noindent\rule{\textwidth}{0.35pt}\par\vspace*{\fill}
{\centering
  \setlength{\fboxsep}{1.2mm}%
  \setlength{\fboxrule}{0.4pt}%
  \fcolorbox{gray!55}{white}{%
    \includegraphics[width=0.95\textwidth,height=0.81\textheight,keepaspectratio]{\facsimiledir/20251126_104831.jpg}%
  }\par}
\vspace*{\fill}
\newpage

% =============================================================================
% VIS-A-VIS SPREAD : FOLIO 50 VERSO (20251126_104849.jpg)
% =============================================================================
\cleardoubleevenpage
\renewcommand{\currentmonthyear}{Novembre 1916}
\renewcommand{\currentfoliolabel}{Folio 50 verso}
\phantomsection\label{fol:101}
% --- En-tête du folio (Page paire / gauche : Transcription, apparat & registres de bas de page) ---
\noindent\rule{\textwidth}{0.5pt}\par\vspace{0.55mm}
\noindent\textbf{\footnotesize FOLIO 50 VERSO} \hfill \textbf{\footnotesize NOVEMBRE 1916}\par\vspace{0.25mm}
\noindent{\fontsize{6.9pt}{8.0pt}\selectfont Défilé de troupes et de prisonniers allemands à Acheux, nuit à Englebelmer et friction avec Curtis}\par\vspace{0.2mm}
\noindent{\fontsize{6.9pt}{8.0pt}\selectfont\textit{20 – 21 novembre 1916} \hfill \textit{Acheux, Bouzincourt, Contay, Englebelmer, Léalvillers (Somme)}}\par\vspace{0.5mm}
\noindent\rule{\textwidth}{0.35pt}\par\vspace{0.8mm}
% --- Corps de la transcription diplomatique et apparat marginal en prolongement ---
{\renewcommand{\arraystretch}{0.83}%
\noindent\hspace*{3.5mm}\begin{tabular}{@{}p{87mm}@{\hspace{3.5mm}}p{39mm}@{}}
      \leavevmode\fontsize{6.6pt}{7.5pt}\selectfont\strut \underline{20 Lundi} — \hfill{} \underline{Nov. /16}\strut &  \\
      \leavevmode\fontsize{6.6pt}{7.5pt}\selectfont\strut Auj'd'hui, comme hier, ciel nuageux. Fines\strut & \smash{\parbox[t]{39mm}{\fontsize{6.2pt}{7.0pt}\selectfont\raggedright\color{black!88}\textbf{2 Auj'd'hui} \textit{abbr. pro} Aujourd'hui.}} \\
      \leavevmode\fontsize{6.6pt}{7.5pt}\selectfont\strut ondées, matin \& soir.\strut &  \\
      \leavevmode\fontsize{6.6pt}{7.5pt}\selectfont\strut Rues \& routes, état lamentable, couvertes\strut &  \\
      \leavevmode\llap{\fontsize{6.4pt}{7.3pt}\selectfont\color{darkgray}5\hspace{1.8mm}}\fontsize{6.6pt}{7.5pt}\selectfont\strut de boue glissante \& épaisse.\strut &  \\
      \leavevmode\fontsize{6.6pt}{7.5pt}\selectfont\strut Ts les jours, passe des heures à regarder rue —\strut & \smash{\parbox[t]{39mm}{\fontsize{6.2pt}{7.0pt}\selectfont\raggedright\color{black!88}\textbf{6 Ts} \textit{abbr. pro} Tous.}} \\
      \leavevmode\fontsize{6.6pt}{7.5pt}\selectfont\strut debout devant fenêtre chambre ou assis s/ lit à côté.\strut & \smash{\parbox[t]{39mm}{\fontsize{6.2pt}{7.0pt}\selectfont\raggedright\color{black!88}\textbf{7 s/} \textit{abbr. pro} sur.}} \\
      \leavevmode\fontsize{6.6pt}{7.5pt}\selectfont\strut Mouvement inouï : lorries, wagons, cars etc. etc.\strut & \smash{\parbox[t]{39mm}{\fontsize{6.2pt}{7.0pt}\selectfont\raggedright\color{black!88}\textbf{8 lorries} \textit{angl. in ms.}}} \\
      \leavevmode\fontsize{6.6pt}{7.5pt}\selectfont\strut Et défilé incessant de troupes revenant de l'avant\strut &  \\
      \leavevmode\llap{\fontsize{6.4pt}{7.3pt}\selectfont\color{darkgray}10\hspace{1.8mm}}\fontsize{6.6pt}{7.5pt}\selectfont\strut ou des tranchées —\strut &  \\
      \leavevmode\fontsize{6.6pt}{7.5pt}\selectfont\strut Souvent aussi, cortèges de prisonniers allemands.\strut &  \\
      \leavevmode\fontsize{6.6pt}{7.5pt}\selectfont\strut Après-midi, longue course en auto : Bouzincourt.\strut &  \\
      \leavevmode\fontsize{6.6pt}{7.5pt}\selectfont\strut Au retour, arrêt à Contay [?]. Nuit noire. Cherché refuge\strut & \smash{\parbox[t]{39mm}{\fontsize{6.2pt}{7.0pt}\selectfont\raggedright\color{black!88}\textbf{13 Contay} \textit{lect. incert. in ms.} — \textbf{14 ds} \textit{abbr. pro} dans.}} \\
      \leavevmode\fontsize{6.6pt}{7.5pt}\selectfont\strut à Englebelmer. Dû attendre ds obscurité profonde, le long\strut &  \\
      \leavevmode\llap{\fontsize{6.4pt}{7.3pt}\selectfont\color{darkgray}15\hspace{1.8mm}}\fontsize{6.6pt}{7.5pt}\selectfont\strut hauts murs du château. Boue épaisse — —\strut &  \\
      \leavevmode\fontsize{6.6pt}{7.5pt}\selectfont\strut Au retour, ciel mervt étoilé... Nuit noire.\strut & \smash{\parbox[t]{39mm}{\fontsize{6.2pt}{7.0pt}\selectfont\raggedright\color{black!88}\textbf{16 mervt} \textit{abbr. pro} merveilleusement. — \textbf{17} \textit{in marg. cum unco} (comptes). \textbf{17 Cartes ples} \textit{abbr. pro} Cartes postales. — \textbf{19 Ma.} \textit{scil.} Mardi. — \textbf{20 1/2} \textit{scil.} demi / moitié.}} \\
      \leavevmode\fontsize{6.6pt}{7.5pt}\selectfont\strut [ — Cartes ples d'Acheux : 45 c. — Timbres : 2 f\strut &  \\
      \leavevmode\fontsize{6.6pt}{7.5pt}\selectfont\strut Carnet : 60 c ]\strut &  \\
      \cline{1-1}
      \leavevmode\fontsize{6.6pt}{7.5pt}\selectfont\strut \underline{21 Ma.} — Épais brouillard, matin \& soir.\strut &  \\
      \leavevmode\llap{\fontsize{6.4pt}{7.3pt}\selectfont\color{darkgray}20\hspace{1.8mm}}\fontsize{6.6pt}{7.5pt}\selectfont\strut 1/2 jour, soleil brillé 1 peu\strut &  \\
      \leavevmode\fontsize{6.6pt}{7.5pt}\selectfont\strut Routes état lamentable. Couche profonde\strut &  \\
      \leavevmode\fontsize{6.6pt}{7.5pt}\selectfont\strut de boue liquide. Humidité pénétrante.\strut &  \\
      \leavevmode\fontsize{6.6pt}{7.5pt}\selectfont\strut Matin, libre. Après-midi, s/ lorry\strut & \smash{\parbox[t]{39mm}{\fontsize{6.2pt}{7.0pt}\selectfont\raggedright\color{black!88}\textbf{23 s/} \textit{abbr. pro} sur. \textbf{23 lorry} \textit{angl. in ms.}}} \\
      \leavevmode\fontsize{6.6pt}{7.5pt}\selectfont\strut à Léalvillers. Au retour, office, remarque\strut &  \\
      \leavevmode\llap{\fontsize{6.4pt}{7.3pt}\selectfont\color{darkgray}25\hspace{1.8mm}}\fontsize{6.6pt}{7.5pt}\selectfont\strut +tôt désobligeante de Curtis : fâché\strut &  \\
      \leavevmode\fontsize{6.6pt}{7.5pt}\selectfont\strut contre lui.\strut &  \\
\end{tabular}}%
\par\vspace*{\fill}
% --- Registres de bas de page (1. Critique interne [si pertinent] ; 2. Loci paralleli en latin ; 3. Commentaire historique en retrait) ---
\noindent\rule{\textwidth}{0.5pt}\par\vspace{0.65mm}
{\fontsize{6.6pt}{7.7pt}\selectfont%
\noindent\normalfont\color{black!92} \textbf{12} (\textit{Bouzincourt}) cf. supra f\textsuperscript{o}~38~v\textsuperscript{o} (p.~\pageref{fol:77}), f\textsuperscript{o}~40~r\textsuperscript{o} (p.~\pageref{fol:80}), infra f\textsuperscript{o}~62~r\textsuperscript{o} (p.~\pageref{fol:124}), \textit{et passim}. — \textbf{13} (\textit{Contay}) cf. supra f\textsuperscript{o}~49~r\textsuperscript{o} (p.~\pageref{fol:98}), infra f\textsuperscript{o}~79~v\textsuperscript{o} (p.~\pageref{fol:159}), f\textsuperscript{o}~80~v\textsuperscript{o} (p.~\pageref{fol:161}), \textit{et passim}. — \textbf{17} (\textit{Acheux}) cf. supra f\textsuperscript{o}~21~v\textsuperscript{o} (p.~\pageref{fol:43}), f\textsuperscript{o}~50~r\textsuperscript{o} (p.~\pageref{fol:100}), infra f\textsuperscript{o}~51~r\textsuperscript{o} (p.~\pageref{fol:102}), f\textsuperscript{o}~51~v\textsuperscript{o} (p.~\pageref{fol:103}), \textit{et passim}.\par\vspace{0.55mm}
\noindent\rule{38mm}{0.4pt}\par\vspace{0.5mm}
\noindent\hspace*{5mm}\parbox{\dimexpr\textwidth-5mm\relax}{\fontsize{6.6pt}{7.7pt}\selectfont\color{black!84}Au lendemain de la bataille de l'Ancre (13--18 novembre 1916, au cours de laquelle la V\textsuperscript{e} Armée britannique captura plus de 7~000 prisonniers allemands à Beaumont-Hamel, Saint-Pierre-Divion et Beaucourt), Hippolyte Combes observe depuis sa fenêtre d'Acheux la relève des bataillons descendant des tranchées et le passage des colonnes de prisonniers allemands évacués vers l'arrière.}\par}
\newpage
\phantomsection\label{fac:101}
% --- En-tête du fac-similé (Page impaire / droite : Photographie du manuscrit) ---
\noindent\rule{\textwidth}{0.5pt}\par\vspace{0.55mm}
\noindent\textbf{\fontsize{6.8pt}{7.8pt}\selectfont FAC-SIMILÉ EN VIS-À-VIS — FOLIO 50 VERSO} \hfill \textbf{\fontsize{6.8pt}{7.8pt}\selectfont NOVEMBRE 1916}\par\vspace{0.25mm}
\noindent{\fontsize{6.8pt}{8.0pt}\selectfont\textit{Livre-Journal d'Hippolyte Combes} — 20 – 21 novembre 1916 (Acheux, Bouzincourt, Contay, Englebelmer, Léalvillers (Somme))}\par\vspace{0.55mm}
\noindent\rule{\textwidth}{0.35pt}\par\vspace*{\fill}
{\centering
  \setlength{\fboxsep}{1.2mm}%
  \setlength{\fboxrule}{0.4pt}%
  \fcolorbox{gray!55}{white}{%
    \includegraphics[width=0.95\textwidth,height=0.81\textheight,keepaspectratio]{\facsimiledir/20251126_104849.jpg}%
  }\par}
\vspace*{\fill}
\newpage

% =============================================================================
% VIS-A-VIS SPREAD : FOLIO 51 RECTO (20251126_104903.jpg)
% =============================================================================
\cleardoubleevenpage
\renewcommand{\currentmonthyear}{Novembre 1916}
\renewcommand{\currentfoliolabel}{Folio 51 recto}
\phantomsection\label{fol:102}
% --- En-tête du folio (Page paire / gauche : Transcription, apparat & registres de bas de page) ---
\noindent\rule{\textwidth}{0.5pt}\par\vspace{0.55mm}
\noindent\textbf{\footnotesize FOLIO 51 RECTO} \hfill \textbf{\footnotesize NOVEMBRE 1916}\par\vspace{0.25mm}
\noindent{\fontsize{6.9pt}{8.0pt}\selectfont Indisposition nocturne, couture de M. Carluvant et tournée à Amplier et Authie}\par\vspace{0.2mm}
\noindent{\fontsize{6.9pt}{8.0pt}\selectfont\textit{22 – 23 novembre 1916} \hfill \textit{Acheux, Amplier, Authieule, Authie (Somme / Pas-de-Calais)}}\par\vspace{0.5mm}
\noindent\rule{\textwidth}{0.35pt}\par\vspace{0.8mm}
% --- Corps de la transcription diplomatique et apparat marginal en prolongement ---
{\renewcommand{\arraystretch}{0.83}%
\noindent\hspace*{3.5mm}\begin{tabular}{@{}p{87mm}@{\hspace{3.5mm}}p{39mm}@{}}
      \leavevmode\fontsize{6.6pt}{7.5pt}\selectfont\strut \underline{22 Mer.} \hfill{} \underline{Novembre 1916 —}\strut & \smash{\parbox[t]{39mm}{\fontsize{6.2pt}{7.0pt}\selectfont\raggedright\color{black!88}\textbf{1 Mer.} \textit{scil.} Mercredi.}} \\
      \leavevmode\fontsize{6.6pt}{7.5pt}\selectfont\strut Hier soir, indigestion \& dérangement\strut &  \\
      \leavevmode\fontsize{6.6pt}{7.5pt}\selectfont\strut d'entrailles. Ni thé, ni souper. 7 selles\strut &  \\
      \leavevmode\fontsize{6.6pt}{7.5pt}\selectfont\strut liquides, dans la nuit. Dû me lever\strut &  \\
      \leavevmode\llap{\fontsize{6.4pt}{7.3pt}\selectfont\color{darkgray}5\hspace{1.8mm}}\fontsize{6.6pt}{7.5pt}\selectfont\strut chaque fois\strut &  \\
      \leavevmode\fontsize{6.6pt}{7.5pt}\selectfont\strut Auj'd'hui, léger mieux. [Craint ?] dysenterie\strut & \smash{\parbox[t]{39mm}{\fontsize{6.2pt}{7.0pt}\selectfont\raggedright\color{black!88}\textbf{6} \textit{lect. incert. in ms.} \textbf{6 Auj'd'hui} \textit{abbr. pro} Aujourd'hui. — \textbf{7} \textit{lect. incert. in ms.}}} \\
      \leavevmode\fontsize{6.6pt}{7.5pt}\selectfont\strut [illis.] selles [mot repassé illis.] clair[s ?] — 1 peu de\strut &  \\
      \leavevmode\fontsize{6.6pt}{7.5pt}\selectfont\strut fièvre. Me suis ménagé. Mangé très\strut &  \\
      \leavevmode\fontsize{6.6pt}{7.5pt}\selectfont\strut peu. Fin jour, amélioration très sensible\strut &  \\
      \leavevmode\llap{\fontsize{6.4pt}{7.3pt}\selectfont\color{darkgray}10\hspace{1.8mm}}\fontsize{6.6pt}{7.5pt}\selectfont\strut — presque guéri —\strut &  \\
      \leavevmode\fontsize{6.6pt}{7.5pt}\selectfont\strut [ A proprre M. Carluvant [?], tailleur : 75 c\strut & \smash{\parbox[t]{39mm}{\fontsize{6.2pt}{7.0pt}\selectfont\raggedright\color{black!88}\textbf{11 Carluvant} \textit{lect. incert. in ms.} \textbf{11} \textit{in marg. cum unco} (comptes). — \textbf{13 1/2} \textit{scil.} demi / moitié. \textbf{13 s.} \textit{abbr. pro} soir.}} \\
      \leavevmode\fontsize{6.6pt}{7.5pt}\selectfont\strut (élargi culottes neuves par derrière) ]\strut &  \\
      \leavevmode\fontsize{6.6pt}{7.5pt}\selectfont\strut Nuit très noire — dès 5h 1/2 s. — Rue\strut &  \\
      \leavevmode\fontsize{6.6pt}{7.5pt}\selectfont\strut boueuse \& impraticable.\strut &  \\
      \cline{1-1}
      \leavevmode\llap{\fontsize{6.4pt}{7.3pt}\selectfont\color{darkgray}15\hspace{1.8mm}}\fontsize{6.6pt}{7.5pt}\selectfont\strut \underline{23 Jeudi.} Enfin, 1 très belle journée.\strut &  \\
      \leavevmode\fontsize{6.6pt}{7.5pt}\selectfont\strut Mais temps 1 peu froid.\strut &  \\
      \leavevmode\fontsize{6.6pt}{7.5pt}\selectfont\strut Sol tjrs spongieux \& mou : sèche\strut & \smash{\parbox[t]{39mm}{\fontsize{6.2pt}{7.0pt}\selectfont\raggedright\color{black!88}\textbf{17 tjrs} \textit{abbr. pro} toujours.}} \\
      \leavevmode\fontsize{6.6pt}{7.5pt}\selectfont\strut très lentement.\strut &  \\
      \leavevmode\fontsize{6.6pt}{7.5pt}\selectfont\strut Matinée, reposé s/ lit du clerk à\strut & \smash{\parbox[t]{39mm}{\fontsize{6.2pt}{7.0pt}\selectfont\raggedright\color{black!88}\textbf{19 s/} \textit{abbr. pro} sur.}} \\
      \leavevmode\llap{\fontsize{6.4pt}{7.3pt}\selectfont\color{darkgray}20\hspace{1.8mm}}\fontsize{6.6pt}{7.5pt}\selectfont\strut l'office. Après-midi, course en auto :\strut &  \\
      \leavevmode\fontsize{6.6pt}{7.5pt}\selectfont\strut Amplier – Authieule et Authie. Rentré pour\strut &  \\
      \leavevmode\fontsize{6.6pt}{7.5pt}\selectfont\strut le thé.\strut &  \\
      \leavevmode\fontsize{6.6pt}{7.5pt}\selectfont\strut Nuit, extt obscure. Difficile marche rue.\strut & \smash{\parbox[t]{39mm}{\fontsize{6.2pt}{7.0pt}\selectfont\raggedright\color{black!88}\textbf{23 extt} \textit{abbr. pro} extrêmement.}} \\
      \leavevmode\fontsize{6.6pt}{7.5pt}\selectfont\strut Tout à fait rétabli légère indisposition.\strut &  \\
      \cline{1-1}
      \leavevmode\llap{\fontsize{6.4pt}{7.3pt}\selectfont\color{darkgray}25\hspace{1.8mm}}\fontsize{6.6pt}{7.5pt}\selectfont\strut [ Washing : 1 f 25 — ]\strut & \smash{\parbox[t]{39mm}{\fontsize{6.2pt}{7.0pt}\selectfont\raggedright\color{black!88}\textbf{25} \textit{in marg. cum unco} (comptes). \textbf{25 Washing} \textit{angl. in ms.}}} \\
\end{tabular}}%
\par\vspace*{\fill}
% --- Registres de bas de page (1. Critique interne [si pertinent] ; 2. Loci paralleli en latin ; 3. Commentaire historique en retrait) ---
\noindent\rule{\textwidth}{0.5pt}\par\vspace{0.65mm}
{\fontsize{6.6pt}{7.7pt}\selectfont%
\noindent\normalfont\color{black!92} \textbf{21} (\textit{Authie}) cf. supra f\textsuperscript{o}~50~r\textsuperscript{o} (p.~\pageref{fol:100}), infra f\textsuperscript{o}~53~v\textsuperscript{o} (p.~\pageref{fol:107}), f\textsuperscript{o}~54~r\textsuperscript{o} (p.~\pageref{fol:108}), \textit{et passim}.\par\vspace{0.55mm}
\noindent\rule{38mm}{0.4pt}\par\vspace{0.5mm}
\noindent\hspace*{5mm}\parbox{\dimexpr\textwidth-5mm\relax}{\fontsize{6.6pt}{7.7pt}\selectfont\color{black!84}Retenu le 22 novembre par une crise aiguë d'entérite, Hippolyte Combes mentionne pour la première fois par son nom et sa profession son logeur du n\textsuperscript{o}~1 route de Léalvillers à Acheux, le maître tailleur \textbf{M.~Carluvant} (qui retouche sa culotte militaire pour 75~c. et chez qui il loge jusqu'au 7 février 1917, f\textsuperscript{o}~59~r\textsuperscript{o}, f\textsuperscript{o}~73~r\textsuperscript{o}--f\textsuperscript{o}~74~r\textsuperscript{o}).}\par}
\newpage
\phantomsection\label{fac:102}
% --- En-tête du fac-similé (Page impaire / droite : Photographie du manuscrit) ---
\noindent\rule{\textwidth}{0.5pt}\par\vspace{0.55mm}
\noindent\textbf{\fontsize{6.8pt}{7.8pt}\selectfont FAC-SIMILÉ EN VIS-À-VIS — FOLIO 51 RECTO} \hfill \textbf{\fontsize{6.8pt}{7.8pt}\selectfont NOVEMBRE 1916}\par\vspace{0.25mm}
\noindent{\fontsize{6.8pt}{8.0pt}\selectfont\textit{Livre-Journal d'Hippolyte Combes} — 22 – 23 novembre 1916 (Acheux, Amplier, Authieule, Authie (Somme / Pas-de-Calais))}\par\vspace{0.55mm}
\noindent\rule{\textwidth}{0.35pt}\par\vspace*{\fill}
{\centering
  \setlength{\fboxsep}{1.2mm}%
  \setlength{\fboxrule}{0.4pt}%
  \fcolorbox{gray!55}{white}{%
    \includegraphics[width=0.95\textwidth,height=0.81\textheight,keepaspectratio]{\facsimiledir/20251126_104903.jpg}%
  }\par}
\vspace*{\fill}
\newpage

% =============================================================================
% VIS-A-VIS SPREAD : FOLIO 51 VERSO (20251126_104916.jpg)
% =============================================================================
\cleardoubleevenpage
\renewcommand{\currentmonthyear}{Novembre 1916}
\renewcommand{\currentfoliolabel}{Folio 51 verso}
\phantomsection\label{fol:103}
% --- En-tête du folio (Page paire / gauche : Transcription, apparat & registres de bas de page) ---
\noindent\rule{\textwidth}{0.5pt}\par\vspace{0.55mm}
\noindent\textbf{\footnotesize FOLIO 51 VERSO} \hfill \textbf{\footnotesize NOVEMBRE 1916}\par\vspace{0.25mm}
\noindent{\fontsize{6.9pt}{8.0pt}\selectfont Bombardement meurtrier de la sucrerie d'Acheux et visite de M. Deloney}\par\vspace{0.2mm}
\noindent{\fontsize{6.9pt}{8.0pt}\selectfont\textit{25 – 26 novembre 1916} \hfill \textit{Acheux, Hérissart (Somme)}}\par\vspace{0.5mm}
\noindent\rule{\textwidth}{0.35pt}\par\vspace{0.8mm}
% --- Corps de la transcription diplomatique et apparat marginal en prolongement ---
{\renewcommand{\arraystretch}{0.87}%
\noindent\hspace*{3.5mm}\begin{tabular}{@{}p{87mm}@{\hspace{3.5mm}}p{39mm}@{}}
      \leavevmode\fontsize{6.9pt}{7.9pt}\selectfont\strut \underline{Sam. 25} \hfill{} \underline{Nov. /16 —}\strut & \smash{\parbox[t]{39mm}{\fontsize{6.2pt}{7.0pt}\selectfont\raggedright\color{black!88}\textbf{1 Sam.} \textit{scil.} Samedi.}} \\
      \leavevmode\fontsize{6.9pt}{7.9pt}\selectfont\strut Auj'd'hui, + encore que hier, pluie averse.\strut & \smash{\parbox[t]{39mm}{\fontsize{6.2pt}{7.0pt}\selectfont\raggedright\color{black!88}\textbf{2 Auj'd'hui} \textit{abbr. pro} Aujourd'hui.}} \\
      \leavevmode\fontsize{6.9pt}{7.9pt}\selectfont\strut Terrains détrempés, vrais bourbiers. Routes \&\strut &  \\
      \leavevmode\fontsize{6.9pt}{7.9pt}\selectfont\strut rues, ornières profondes où l'on enfonce j'à\strut & \smash{\parbox[t]{39mm}{\fontsize{6.2pt}{7.0pt}\selectfont\raggedright\color{black!88}\textbf{4 j'à} \textit{abbr. pro} jusqu'à.}} \\
      \leavevmode\llap{\fontsize{6.4pt}{7.3pt}\selectfont\color{darkgray}5\hspace{1.8mm}}\fontsize{6.9pt}{7.9pt}\selectfont\strut 1/2 jambe. C'est désespérant.\strut & \smash{\parbox[t]{39mm}{\fontsize{6.2pt}{7.0pt}\selectfont\raggedright\color{black!88}\textbf{5 1/2} \textit{scil.} demi / moitié.}} \\
      \leavevmode\fontsize{6.9pt}{7.9pt}\selectfont\strut Durant de longues heures, langui bureau,\strut &  \\
      \leavevmode\fontsize{6.9pt}{7.9pt}\selectfont\strut assis caisse bidons de pétrole ou au pied du\strut &  \\
      \leavevmode\fontsize{6.9pt}{7.9pt}\selectfont\strut lit du clerk. Lis ou rumine. Triste vie.\strut &  \\
      \leavevmode\fontsize{6.9pt}{7.9pt}\selectfont\strut Matinée, + de 25 obus, tombés s/ Acheux,\strut & \smash{\parbox[t]{39mm}{\fontsize{6.2pt}{7.0pt}\selectfont\raggedright\color{black!88}\textbf{9 s/} \textit{abbr. pro} sur.}} \\
      \leavevmode\llap{\fontsize{6.4pt}{7.3pt}\selectfont\color{darkgray}10\hspace{1.8mm}}\fontsize{6.9pt}{7.9pt}\selectfont\strut du côté sucrerie et Fme Royer. Plusrs tués et blessés.\strut &  \\
      \leavevmode\fontsize{6.9pt}{7.9pt}\selectfont\strut Séjour ds ma chambre, impossible, à\strut & \smash{\parbox[t]{39mm}{\fontsize{6.2pt}{7.0pt}\selectfont\raggedright\color{black!88}\textbf{11 ds} \textit{abbr. pro} dans.}} \\
      \leavevmode\fontsize{6.9pt}{7.9pt}\selectfont\strut cause froid et surtout bruits énervants de la rue.\strut &  \\
      \cline{1-1}
      \leavevmode\fontsize{6.9pt}{7.9pt}\selectfont\strut \underline{Dim. 26}\strut & \smash{\parbox[t]{39mm}{\fontsize{6.2pt}{7.0pt}\selectfont\raggedright\color{black!88}\textbf{13 Dim.} \textit{scil.} Dimanche.}} \\
      \leavevmode\fontsize{6.9pt}{7.9pt}\selectfont\strut Toute la nuit et matinée entière,\strut &  \\
      \leavevmode\llap{\fontsize{6.4pt}{7.3pt}\selectfont\color{darkgray}15\hspace{1.8mm}}\fontsize{6.9pt}{7.9pt}\selectfont\strut pluie abondante. Habits ruissellent d'eau.\strut &  \\
      \leavevmode\fontsize{6.9pt}{7.9pt}\selectfont\strut Rues \& routes lavées : 1/2 des chaussées,\strut & \smash{\parbox[t]{39mm}{\fontsize{6.2pt}{7.0pt}\selectfont\raggedright\color{black!88}\textbf{16 1/2} \textit{scil.} demi / moitié.}} \\
      \leavevmode\fontsize{6.9pt}{7.9pt}\selectfont\strut côtés d., profondes couches de boue.\strut &  \\
      \leavevmode\fontsize{6.9pt}{7.9pt}\selectfont\strut Après-midi, pluie cesse peu à peu\strut &  \\
      \leavevmode\fontsize{6.9pt}{7.9pt}\selectfont\strut et fin journée, sans ondée.\strut &  \\
      \leavevmode\llap{\fontsize{6.4pt}{7.3pt}\selectfont\color{darkgray}20\hspace{1.8mm}}\fontsize{6.9pt}{7.9pt}\selectfont\strut Lt Walsh \& son interprète, M. Deloney, venus\strut & \smash{\parbox[t]{39mm}{\fontsize{6.2pt}{7.0pt}\selectfont\raggedright\color{black!88}\textbf{20 Lt} \textit{abbr. pro} Lieutenant. — \textbf{21 ds} \textit{abbr. pro} dans.}} \\
      \leavevmode\fontsize{6.9pt}{7.9pt}\selectfont\strut à Acheux, thé ds notre cabane. Les ai\strut &  \\
      \leavevmode\fontsize{6.9pt}{7.9pt}\selectfont\strut raccompagnés à Hérissart. Ainsi pu causer\strut &  \\
      \leavevmode\fontsize{6.9pt}{7.9pt}\selectfont\strut longt en fr. Nuit très belle,\strut & \smash{\parbox[t]{39mm}{\fontsize{6.2pt}{7.0pt}\selectfont\raggedright\color{black!88}\textbf{23 longt} \textit{abbr. pro} longtemps.}} \\
      \leavevmode\fontsize{6.9pt}{7.9pt}\selectfont\strut (magnift étoilée)\strut &  \\
\end{tabular}}%
\par\vspace*{\fill}
% --- Registres de bas de page (1. Critique interne [si pertinent] ; 2. Loci paralleli en latin ; 3. Commentaire historique en retrait) ---
\noindent\rule{\textwidth}{0.5pt}\par\vspace{0.65mm}
{\fontsize{6.6pt}{7.7pt}\selectfont%
\noindent\normalfont\color{black!92} \textbf{9} (\textit{Acheux}) cf. supra f\textsuperscript{o}~21~v\textsuperscript{o} (p.~\pageref{fol:43}), f\textsuperscript{o}~51~r\textsuperscript{o} (p.~\pageref{fol:102}), infra f\textsuperscript{o}~52~r\textsuperscript{o} (p.~\pageref{fol:104}), f\textsuperscript{o}~52~v\textsuperscript{o} (p.~\pageref{fol:105}), \textit{et passim}. — \textbf{20} (\textit{Lt Walsh}) cf. supra f\textsuperscript{o}~49~v\textsuperscript{o} (p.~\pageref{fol:99}), infra f\textsuperscript{o}~53~r\textsuperscript{o} (p.~\pageref{fol:106}), f\textsuperscript{o}~60~r\textsuperscript{o} (p.~\pageref{fol:120}). — \textbf{22} (\textit{Hérissart}) cf. supra f\textsuperscript{o}~47~v\textsuperscript{o} (p.~\pageref{fol:95}), f\textsuperscript{o}~49~v\textsuperscript{o} (p.~\pageref{fol:99}), infra f\textsuperscript{o}~52~r\textsuperscript{o} (p.~\pageref{fol:104}), f\textsuperscript{o}~53~r\textsuperscript{o} (p.~\pageref{fol:106}), \textit{et passim}.\par\vspace{0.55mm}
\noindent\rule{38mm}{0.4pt}\par\vspace{0.5mm}
\noindent\hspace*{5mm}\parbox{\dimexpr\textwidth-5mm\relax}{\fontsize{6.6pt}{7.7pt}\selectfont\color{black!84}Le samedi 25 novembre 1916 au matin, un tir d'artillerie allemande à longue portée de plus de vingt-cinq obus frappe \textbf{Acheux-en-Amiénois} aux abords de la râperie/sucrerie et de la ferme Royer, faisant plusieurs tués et blessés.}\par}
\newpage
\phantomsection\label{fac:103}
% --- En-tête du fac-similé (Page impaire / droite : Photographie du manuscrit) ---
\noindent\rule{\textwidth}{0.5pt}\par\vspace{0.55mm}
\noindent\textbf{\fontsize{6.8pt}{7.8pt}\selectfont FAC-SIMILÉ EN VIS-À-VIS — FOLIO 51 VERSO} \hfill \textbf{\fontsize{6.8pt}{7.8pt}\selectfont NOVEMBRE 1916}\par\vspace{0.25mm}
\noindent{\fontsize{6.8pt}{8.0pt}\selectfont\textit{Livre-Journal d'Hippolyte Combes} — 25 – 26 novembre 1916 (Acheux, Hérissart (Somme))}\par\vspace{0.55mm}
\noindent\rule{\textwidth}{0.35pt}\par\vspace*{\fill}
{\centering
  \setlength{\fboxsep}{1.2mm}%
  \setlength{\fboxrule}{0.4pt}%
  \fcolorbox{gray!55}{white}{%
    \includegraphics[width=0.95\textwidth,height=0.81\textheight,keepaspectratio]{\facsimiledir/20251126_104916.jpg}%
  }\par}
\vspace*{\fill}
\newpage

% =============================================================================
% VIS-A-VIS SPREAD : FOLIO 52 RECTO (20251126_104932.jpg)
% =============================================================================
\cleardoubleevenpage
\renewcommand{\currentmonthyear}{Novembre 1916}
\renewcommand{\currentfoliolabel}{Folio 52 recto}
\phantomsection\label{fol:104}
% --- En-tête du folio (Page paire / gauche : Transcription, apparat & registres de bas de page) ---
\noindent\rule{\textwidth}{0.5pt}\par\vspace{0.55mm}
\noindent\textbf{\footnotesize FOLIO 52 RECTO} \hfill \textbf{\footnotesize NOVEMBRE 1916}\par\vspace{0.25mm}
\noindent{\fontsize{6.9pt}{8.0pt}\selectfont Thé au château de Martinsart chez les Bertin et brouillard fantomatique sur la Somme}\par\vspace{0.2mm}
\noindent{\fontsize{6.9pt}{8.0pt}\selectfont\textit{28 – 30 novembre 1916} \hfill \textit{Acheux, Martinsart, Raincheval, Puchevillers, Rubempré, Septenville, Pierregot, Hérissart (Somme)}}\par\vspace{0.5mm}
\noindent\rule{\textwidth}{0.35pt}\par\vspace{0.8mm}
% --- Corps de la transcription diplomatique et apparat marginal en prolongement ---
{\renewcommand{\arraystretch}{0.83}%
\noindent\hspace*{3.5mm}\begin{tabular}{@{}p{87mm}@{\hspace{3.5mm}}p{39mm}@{}}
      \leavevmode\fontsize{6.6pt}{7.5pt}\selectfont\strut \underline{28 Ma.} \hfill{} \underline{Novbre 1916}\strut & \smash{\parbox[t]{39mm}{\fontsize{6.2pt}{7.0pt}\selectfont\raggedright\color{black!88}\textbf{1 Ma.} \textit{scil.} Mardi.}} \\
      \leavevmode\fontsize{6.6pt}{7.5pt}\selectfont\strut Tjrs même temps : brouillard \& pluie.\strut & \smash{\parbox[t]{39mm}{\fontsize{6.2pt}{7.0pt}\selectfont\raggedright\color{black!88}\textbf{2 Tjrs} \textit{abbr. pro} Toujours.}} \\
      \leavevmode\fontsize{6.6pt}{7.5pt}\selectfont\strut Routes, état lamentable. Matin, rien à faire.\strut &  \\
      \leavevmode\fontsize{6.6pt}{7.5pt}\selectfont\strut Après-midi, course en auto à Martinsart.\strut &  \\
      \leavevmode\llap{\fontsize{6.4pt}{7.3pt}\selectfont\color{darkgray}5\hspace{1.8mm}}\fontsize{6.6pt}{7.5pt}\selectfont\strut Château, reçu aimablt par M. \& Mme Bertin [?],\strut & \smash{\parbox[t]{39mm}{\fontsize{6.2pt}{7.0pt}\selectfont\raggedright\color{black!88}\textbf{5 Bertin} \textit{lect. incert. in ms.} \textbf{5 Mme} \textit{abbr. pro} Madame. — \textbf{6 tt} \textit{abbr. pro} tout.}} \\
      \leavevmode\fontsize{6.6pt}{7.5pt}\selectfont\strut et pris thé av. eux — Gros canons, tt\strut &  \\
      \leavevmode\fontsize{6.6pt}{7.5pt}\selectfont\strut proche, en action.\strut &  \\
      \leavevmode\fontsize{6.6pt}{7.5pt}\selectfont\strut Traversé régions plates, terrains défoncés.\strut &  \\
      \leavevmode\fontsize{6.6pt}{7.5pt}\selectfont\strut Convois de ttes sortes.\strut & \smash{\parbox[t]{39mm}{\fontsize{6.2pt}{7.0pt}\selectfont\raggedright\color{black!88}\textbf{9 ttes} \textit{abbr. pro} toutes.}} \\
      \leavevmode\llap{\fontsize{6.4pt}{7.3pt}\selectfont\color{darkgray}10\hspace{1.8mm}}\fontsize{6.6pt}{7.5pt}\selectfont\strut Au retour, nuit noire, éclairs bleuâtres\strut &  \\
      \leavevmode\fontsize{6.6pt}{7.5pt}\selectfont\strut à l'horizon — canonnade.\strut &  \\
      \cline{1-1}
      \leavevmode\fontsize{6.6pt}{7.5pt}\selectfont\strut \underline{29 Me.} Matin, très épais brouillard.\strut &  \\
      \leavevmode\fontsize{6.6pt}{7.5pt}\selectfont\strut Après-midi, longue course en auto : Raincheval,\strut &  \\
      \leavevmode\fontsize{6.6pt}{7.5pt}\selectfont\strut Puchevillers, Rubempré et Septenville \& Pierregot.\strut &  \\
      \leavevmode\llap{\fontsize{6.4pt}{7.3pt}\selectfont\color{darkgray}15\hspace{1.8mm}}\fontsize{6.6pt}{7.5pt}\selectfont\strut Brouillard se lève, de + en + dense et ensevelit\strut &  \\
      \leavevmode\fontsize{6.6pt}{7.5pt}\selectfont\strut paysage. Arbres, fantomatiques au 1/2 des champs.\strut & \smash{\parbox[t]{39mm}{\fontsize{6.2pt}{7.0pt}\selectfont\raggedright\color{black!88}\textbf{16 1/2} \textit{scil.} demi / moitié.}} \\
      \leavevmode\fontsize{6.6pt}{7.5pt}\selectfont\strut Moulins à (vent)\strut &  \\
      \leavevmode\fontsize{6.6pt}{7.5pt}\selectfont\strut [ — Washing : 70 c — ]\strut & \smash{\parbox[t]{39mm}{\fontsize{6.2pt}{7.0pt}\selectfont\raggedright\color{black!88}\textbf{18} \textit{in marg. cum unco} (comptes). \textbf{18 Washing} \textit{angl. in ms.} — \textbf{19 J.} \textit{scil.} Jeudi.}} \\
      \cline{1-1}
      \leavevmode\fontsize{6.6pt}{7.5pt}\selectfont\strut \underline{30 J.} — Matinée, à Acheux. Temps couvert\strut &  \\
      \leavevmode\llap{\fontsize{6.4pt}{7.3pt}\selectfont\color{darkgray}20\hspace{1.8mm}}\fontsize{6.6pt}{7.5pt}\selectfont\strut \& brouillard.\strut &  \\
      \leavevmode\fontsize{6.6pt}{7.5pt}\selectfont\strut Après-midi, course en auto : Septenville\strut &  \\
      \leavevmode\fontsize{6.6pt}{7.5pt}\selectfont\strut aller. Au retour, arrêt à Hérissart. Y pris thé.\strut &  \\
      \leavevmode\fontsize{6.6pt}{7.5pt}\selectfont\strut Rentré Acheux vers 7h s.\strut & \smash{\parbox[t]{39mm}{\fontsize{6.2pt}{7.0pt}\selectfont\raggedright\color{black!88}\textbf{23 s.} \textit{abbr. pro} soir.}} \\
      \leavevmode\fontsize{6.6pt}{7.5pt}\selectfont\strut Soir, dîné 1ère fois au nouveau mess\strut & \smash{\parbox[t]{39mm}{\fontsize{6.2pt}{7.0pt}\selectfont\raggedright\color{black!88}\textbf{24 1ère} \textit{scil.} première. \textbf{24 mess} \textit{angl. in ms.} — \textbf{25 Lt} \textit{abbr. pro} Lieutenant.}} \\
      \leavevmode\llap{\fontsize{6.4pt}{7.3pt}\selectfont\color{darkgray}25\hspace{1.8mm}}\fontsize{6.6pt}{7.5pt}\selectfont\strut 79 C, chambre du Lt Curtis.\strut &  \\
\end{tabular}}%
\par\vspace*{\fill}
% --- Registres de bas de page (1. Critique interne [si pertinent] ; 2. Loci paralleli en latin ; 3. Commentaire historique en retrait) ---
\noindent\rule{\textwidth}{0.5pt}\par\vspace{0.65mm}
{\fontsize{6.6pt}{7.7pt}\selectfont%
\noindent\normalfont\color{black!92} \textbf{4} (\textit{Martinsart}) cf. infra f\textsuperscript{o}~52~v\textsuperscript{o} (p.~\pageref{fol:105}), f\textsuperscript{o}~58~v\textsuperscript{o} (p.~\pageref{fol:117}), \textit{et passim}. — \textbf{13} (\textit{Raincheval}) cf. supra f\textsuperscript{o}~20~v\textsuperscript{o} (p.~\pageref{fol:41}), f\textsuperscript{o}~49~v\textsuperscript{o} (p.~\pageref{fol:99}), infra f\textsuperscript{o}~53~v\textsuperscript{o} (p.~\pageref{fol:107}), f\textsuperscript{o}~54~r\textsuperscript{o} (p.~\pageref{fol:108}), \textit{et passim}. — \textbf{14} (\textit{Puchevillers}) cf. supra f\textsuperscript{o}~37~v\textsuperscript{o} (p.~\pageref{fol:75}), infra f\textsuperscript{o}~55~r\textsuperscript{o} (p.~\pageref{fol:110}), f\textsuperscript{o}~56~r\textsuperscript{o} (p.~\pageref{fol:112}), \textit{et passim}.\par\vspace{0.55mm}
\noindent\rule{38mm}{0.4pt}\par\vspace{0.5mm}
\noindent\hspace*{5mm}\parbox{\dimexpr\textwidth-5mm\relax}{\fontsize{6.6pt}{7.7pt}\selectfont\color{black!84}Le mardi 28 novembre 1916, une tournée conduit Hippolyte Combes au château de \textbf{Martinsart} (Somme, à 3,5~km au nord d'Albert, en arrière immédiat des batteries lourdes britanniques tirant vers l'Ancre), dont les propriétaires âgés seront peu après évacués à Franvillers (où Hippolyte Combes les retrouvera le 30 décembre, f\textsuperscript{o}~58~v\textsuperscript{o}).}\par}
\newpage
\phantomsection\label{fac:104}
% --- En-tête du fac-similé (Page impaire / droite : Photographie du manuscrit) ---
\noindent\rule{\textwidth}{0.5pt}\par\vspace{0.55mm}
\noindent\textbf{\fontsize{6.8pt}{7.8pt}\selectfont FAC-SIMILÉ EN VIS-À-VIS — FOLIO 52 RECTO} \hfill \textbf{\fontsize{6.8pt}{7.8pt}\selectfont NOVEMBRE 1916}\par\vspace{0.25mm}
\noindent{\fontsize{6.8pt}{8.0pt}\selectfont\textit{Livre-Journal d'Hippolyte Combes} — 28 – 30 novembre 1916 (Acheux, Martinsart, Raincheval, Puchevillers, Rubempré, Septenville, Pierregot, Hérissart (Somme))}\par\vspace{0.55mm}
\noindent\rule{\textwidth}{0.35pt}\par\vspace*{\fill}
{\centering
  \setlength{\fboxsep}{1.2mm}%
  \setlength{\fboxrule}{0.4pt}%
  \fcolorbox{gray!55}{white}{%
    \includegraphics[width=0.95\textwidth,height=0.81\textheight,keepaspectratio]{\facsimiledir/20251126_104932.jpg}%
  }\par}
\vspace*{\fill}
\newpage

% =============================================================================
% CHAPITRE V (FOLIO 52 VERSO A FOLIO 63 VERSO)
% =============================================================================
\cleardoubleevenpage
\renewcommand{\currentmonthyear}{Décembre 1916 – Janvier 1917}
\renewcommand{\currentfoliolabel}{Chapitre V}
\chapter{L’hiver 1916–1917 à Acheux-en-Amiénois}

\noindent{\fontsize{7.8pt}{9.0pt}\selectfont\textit{Folios 52 verso à 63 verso — Le capitaine Williams, les tournées de fourrage, Noël et le grand gel de janvier (1er décembre 1916 – 25 janvier 1917)}}\par\vspace{1.5mm}
\noindent\rule{\textwidth}{0.5pt}\par\vspace{0.9mm}
\noindent\textbf{\scriptsize\textsc{Résumé analytique du Chapitre V (Folio 52 verso à Folio 63 verso)}}\par\vspace{0.5mm}
\noindent{\fontsize{6.9pt}{8.1pt}\selectfont Le cinquième chapitre (f\textsuperscript{o}~52~v\textsuperscript{o} à f\textsuperscript{o}~63~v\textsuperscript{o}) relate le rude hiver 1916–1917 à Acheux-en-Amiénois, du 1\textsuperscript{er} décembre 1916 au 25 janvier 1917. Dès le 2 décembre, le lieutenant Curtis est remplacé au service des fourrages par le très jeune capitaine Williams (\textit{Foraging Officer}), dont le caractère brusque et l'inexpérience provoquent de vives tensions morales avec l'universitaire français de quarante ans. Malgré la boue liquide de décembre puis la neige et le gel intense de janvier 1917, Hippolyte Combes accomplit chaque jour de longues tournées en automobile ou à cheval à travers toute la région entre Doullens, Albert et Amiens : Franvillers, Buire-sur-l'Ancre, Ribemont et Heilly (où il revoit avec émotion ses cantonnements de l'hiver 1915–1916), Toutencourt (où il assiste le major Levy lors d'une réunion des réfugiés d'Englebelmer), Authie, Saint-Léger, Beauquesne, Raincheval, Le Quesnoy, Contay, Vadencourt, Puchevillers, Hérissart et Rubempré. Après un réveillon de Noël au château d'Acheux (siège de l'école de mortiers de la V\textsuperscript{e} Armée) et les concerts des \textit{13th Corps Archie Singers} à la \textit{Y.M.C.A.}, il rend visite le 9 janvier près d'Hédauville à ses amis de la 150\textsuperscript{e} Compagnie A.S.C. (dont le lieutenant Morris) avant d'obtenir fin janvier son titre de permission pour le Midi.}\par\vspace{1.0mm}
\noindent\rule{\textwidth}{0.35pt}\par\vspace{1.4mm}

\noindent\textbf{\scriptsize\textsc{Introduction historique et philologique}}\par\vspace{0.6mm}
{\fontsize{7.1pt}{8.3pt}\selectfont Sur le plan militaire, cette période hivernale correspond à la réorganisation du dispositif britannique au nord de la Somme et de l'Ancre, avec l'extension du secteur de la \textbf{V\textsuperscript{e} Armée britannique} (\textit{Fifth Army}) et la relève progressive du XIII\textsuperscript{e} Corps par le \textbf{V\textsuperscript{e} Corps britannique} autour d'Acheux et de Toutencourt. La visite du 9 janvier 1917 au bivouac d'\textbf{Hédauville} (f\textsuperscript{o}~60~v\textsuperscript{o}) éclaire l'articulation interne du Train divisionnaire de la 18\textsuperscript{e} Division (\textit{18th Divisional Train}), qui réunissait les \textbf{150\textsuperscript{e}, 151\textsuperscript{e}, 152\textsuperscript{e} et 153\textsuperscript{e} Compagnies A.S.C.} (Maj. A.~F. Becke, \textit{Order of Battle of Divisions}, Part~3A, p.~81). Sur le plan humain, le portrait sans concession du jeune \textbf{capitaine Williams} met en lumière la condition délicate des universitaires français détachés comme interprètes auprès d'officiers subalternes britanniques beaucoup plus jeunes qu'eux.\par}
\vfill
\newpage
% --- Page de droite d'ouverture du chapitre : Carte vectorielle TikZ & Photographies d'archives ---
\noindent\rule{\textwidth}{0.5pt}\par\vspace{0.6mm}
\noindent\textbf{\fontsize{6.8pt}{7.8pt}\selectfont CARTE D'ITINÉRAIRE DU CHAPITRE V} \hfill \textbf{\fontsize{6.8pt}{7.8pt}\selectfont DÉCEMBRE 1916 – JANVIER 1917}\par\vspace{0.25mm}
\noindent{\fontsize{6.7pt}{7.8pt}\selectfont\textit{Déplacements, cantonnements et tournées d'Hippolyte Combes (Folio 52 verso à Folio 63 verso)}}\par\vspace{0.55mm}
\noindent\rule{\textwidth}{0.35pt}\par\vspace{1.2mm}
\begin{center}

\begin{tikzpicture}[
  scale=0.96,
  every node/.style={font=\fontsize{6.0pt}{7.0pt}\selectfont},
  loc/.style={circle, fill=black, inner sep=1.1pt},
  hub/.style={circle, draw=bordeaux, fill=white, line width=0.85pt, inner sep=1.8pt},
  tour/.style={draw=bordeaux!85, line width=0.55pt, ->, >=stealth},
  river/.style={draw=black!35, line width=1.1pt}
]
\draw[black!50, line width=0.4pt] (-60mm,-42mm) rectangle (60mm,42mm);
\fill[black!3] (-60mm,-42mm) rectangle (60mm,42mm);

% Rivière l'Ancre
\draw[river] (36mm,28mm) .. controls (32mm,8mm) and (28mm,-12mm) .. (18mm,-36mm)
  node[pos=0.82, right, font=\fontsize{5.6pt}{6.5pt}\selectfont\itshape]{L'Ancre};

% Centre d'Acheux-en-Amiénois
\node[hub] (acheux) at (4mm, 14mm) {};
\node[above right] at (acheux) {\textbf{ACHEUX-EN-AMIÉNOIS} (centre des tournées)};

% Nord et Nord-Ouest
\node[loc] (authie) at (-4mm, 32mm) {};
\node[above] at (authie) {\textbf{Authie} \& St-Léger (9 déc., 6 janv.)};
\node[loc] (louvencourt) at (2mm, 24mm) {};
\node[left] at (louvencourt) {\textbf{Louvencourt}};
\node[loc] (bertrancourt) at (18mm, 24mm) {};
\node[right] at (bertrancourt) {\textbf{Bertrancourt} \& Bus-lès-Artois};
\node[loc] (beauquesne) at (-26mm, 24mm) {};
\node[left] at (beauquesne) {\textbf{Beauquesne} (20 déc., 5 janv.)};
\node[loc] (quesnoy) at (-18mm, 16mm) {};
\node[left] at (quesnoy) {Ferme du Quesnoy};
\node[loc] (raincheval) at (-14mm, 10mm) {};
\node[below left] at (raincheval) {\textbf{Raincheval} (\textit{Army P. Board})};
\node[loc] (arqueves) at (-6mm, 16mm) {};
\node[above left] at (arqueves) {Arquèves (\textbf{Belle-Église})};

% Est et Sud-Est (vers l'Ancre)
\node[loc] (lealvillers) at (8mm, 8mm) {};
\node[right] at (lealvillers) {Léalvillers};
\node[loc] (varennes) at (14mm, 4mm) {};
\node[right] at (varennes) {Varennes};
\node[loc] (forceville) at (20mm, 12mm) {};
\node[right] at (forceville) {Forceville \& \textbf{Englebelmer}};
\node[loc] (hedauville) at (18mm, -1mm) {};
\node[right] at (hedauville) {\textbf{Hédauville} (\textit{150th Co. A.S.C.}, 9 janv.)};
\node[loc] (albert) at (30mm, -12mm) {};
\node[right] at (albert) {\textbf{Albert} \& Bouzincourt};
\node[loc] (buire) at (24mm, -24mm) {};
\node[right] at (buire) {\textbf{Buire-sur-l'Ancre}, Ribemont \& Heilly (1\textsuperscript{er} déc.)};
\node[loc] (franvillers) at (10mm, -30mm) {};
\node[below] at (franvillers) {\textbf{Franvillers}};

% Sud et Sud-Ouest
\node[loc] (toutencourt) at (-4mm, -2mm) {};
\node[left] at (toutencourt) {\textbf{Toutencourt} (Mission V\textsuperscript{e} Armée)};
\node[loc] (harponville) at (6mm, -4mm) {};
\node[right] at (harponville) {Harponville};
\node[loc] (contay) at (6mm, -12mm) {};
\node[left] at (contay) {\textbf{Contay} \& Vadencourt};
\node[loc] (puchevillers) at (-18mm, -2mm) {};
\node[left] at (puchevillers) {\textbf{Puchevillers}};
\node[loc] (herissart) at (-20mm, -10mm) {};
\node[left] at (herissart) {\textbf{Hérissart}};
\node[loc] (rubempre) at (-16mm, -18mm) {};
\node[left] at (rubempre) {\textbf{Rubempré} \& Septenville};

% Rayons des tournées quotidiennes
\foreach \dest in {authie,louvencourt,bertrancourt,beauquesne,raincheval,arqueves,forceville,hedauville,albert,toutencourt,contay,puchevillers,herissart,rubempre} {
  \draw[tour] (acheux) -- (\dest);
}
\draw[tour] (albert) -- (buire) -- (franvillers);

\draw[line width=0.6pt] (-54mm,-36mm) -- (-34mm,-36mm);
\node[above] at (-44mm,-36mm) {0 \hfill 4 \hfill 8 km};
\draw[->, line width=0.6pt] (52mm, 28mm) -- (52mm, 37mm) node[above]{\textbf{N}};
\end{tikzpicture}

\end{center}
\vfill
\noindent\rule{\textwidth}{0.3pt}\par\vspace{1.2mm}
\noindent\begin{center}
  \begin{minipage}[b]{0.50\textwidth}
    \centering
    \setlength{\fboxsep}{0.7mm}\setlength{\fboxrule}{0.4pt}%
    \fcolorbox{gray!55}{white}{\includegraphics[width=0.94\linewidth,height=28mm,keepaspectratio]{illustrations/camarades_ailly_1919.jpg}}
    \par\vspace{0.6mm}{\fontsize{5.8pt}{6.7pt}\selectfont Souvenir de camaraderie militaire sur le front de Picardie (carte-photo dédicacée à Hippolyte Combes, Ailly-sur-Noye, mai 1919). Collection familiale.}
  \end{minipage}
\end{center}
\vspace*{1mm}
\newpage

% =============================================================================
% VIS-A-VIS SPREAD : FOLIO 52 VERSO (20251126_104946.jpg)
% =============================================================================
\cleardoubleevenpage
\renewcommand{\currentmonthyear}{Décembre 1916}
\renewcommand{\currentfoliolabel}{Folio 52 verso}
\phantomsection\label{fol:105}
% --- En-tête du folio (Page paire / gauche : Transcription, apparat & registres de bas de page) ---
\noindent\rule{\textwidth}{0.5pt}\par\vspace{0.55mm}
\noindent\textbf{\footnotesize FOLIO 52 VERSO} \hfill \textbf{\footnotesize DÉCEMBRE 1916}\par\vspace{0.25mm}
\noindent{\fontsize{6.9pt}{8.0pt}\selectfont Tournée à Franvillers et Buire-sur-l'Ancre, arrivée imprévue du très jeune capitaine Williams}\par\vspace{0.2mm}
\noindent{\fontsize{6.9pt}{8.0pt}\selectfont\textit{1er – 2 décembre 1916} \hfill \textit{Acheux, Franvillers, Buire-sur-l'Ancre, Ribemont, Heilly, Toutencourt (Somme)}}\par\vspace{0.5mm}
\noindent\rule{\textwidth}{0.35pt}\par\vspace{0.8mm}
% --- Corps de la transcription diplomatique et apparat marginal en prolongement ---
{\renewcommand{\arraystretch}{0.83}%
\noindent\hspace*{3.5mm}\begin{tabular}{@{}p{87mm}@{\hspace{3.5mm}}p{39mm}@{}}
      \leavevmode\fontsize{6.6pt}{7.5pt}\selectfont\strut \underline{1er Ven.} \hfill{} \underline{Décembre 1916}\strut & \smash{\parbox[t]{39mm}{\fontsize{6.2pt}{7.0pt}\selectfont\raggedright\color{black!88}\textbf{1 Ven.} \textit{scil.} Vendredi. \textbf{1 1er} \textit{scil.} premier.}} \\
      \leavevmode\fontsize{6.6pt}{7.5pt}\selectfont\strut Matinée perdue à ne rien faire : office \& billet.\strut &  \\
      \leavevmode\fontsize{6.6pt}{7.5pt}\selectfont\strut Après early lunch, partis en auto. Venus\strut & \smash{\parbox[t]{39mm}{\fontsize{6.2pt}{7.0pt}\selectfont\raggedright\color{black!88}\textbf{3 lunch} \textit{angl. in ms.}}} \\
      \leavevmode\fontsize{6.6pt}{7.5pt}\selectfont\strut à Franvillers, puis à Buire s/ Ancre pour\strut & \smash{\parbox[t]{39mm}{\fontsize{6.2pt}{7.0pt}\selectfont\raggedright\color{black!88}\textbf{4 s/} \textit{abbr. pro} sur.}} \\
      \leavevmode\llap{\fontsize{6.4pt}{7.3pt}\selectfont\color{darkgray}5\hspace{1.8mm}}\fontsize{6.6pt}{7.5pt}\selectfont\strut affaires à traiter.\strut &  \\
      \leavevmode\fontsize{6.6pt}{7.5pt}\selectfont\strut Plaisir à revoir région de Ribemont – Heilly.\strut &  \\
      \leavevmode\fontsize{6.6pt}{7.5pt}\selectfont\strut Au retour, arrêté à Toutencourt \& monté\strut &  \\
      \leavevmode\fontsize{6.6pt}{7.5pt}\selectfont\strut bureau mission frse, Ve armée, Cap. Bresson.\strut & \smash{\parbox[t]{39mm}{\fontsize{6.2pt}{7.0pt}\selectfont\raggedright\color{black!88}\textbf{8 frse} \textit{abbr. pro} française. \textbf{8 Cap.} \textit{abbr. pro} Capitaine.}} \\
      \leavevmode\fontsize{6.6pt}{7.5pt}\selectfont\strut Trouvé message à Acheux — présenté au\strut &  \\
      \leavevmode\llap{\fontsize{6.4pt}{7.3pt}\selectfont\color{darkgray}10\hspace{1.8mm}}\fontsize{6.6pt}{7.5pt}\selectfont\strut bureau M.F., XIIIe corps et pris en charge.\strut &  \\
      \leavevmode\fontsize{6.6pt}{7.5pt}\selectfont\strut Perspective réintégration prochaine ds\strut & \smash{\parbox[t]{39mm}{\fontsize{6.2pt}{7.0pt}\selectfont\raggedright\color{black!88}\textbf{11 ds} \textit{abbr. pro} dans.}} \\
      \leavevmode\fontsize{6.6pt}{7.5pt}\selectfont\strut régiment d'origine! —\strut &  \\
      \leavevmode\fontsize{6.6pt}{7.5pt}\selectfont\strut [ — A Curtis, pour ma part de\strut & \smash{\parbox[t]{39mm}{\fontsize{6.2pt}{7.0pt}\selectfont\raggedright\color{black!88}\textbf{13} \textit{in marg. cum unco} (comptes). — \textbf{14 mess} \textit{angl. in ms.}}} \\
      \leavevmode\fontsize{6.6pt}{7.5pt}\selectfont\strut dépenses du mess, à Acheux : 40 f\strut &  \\
      \leavevmode\llap{\fontsize{6.4pt}{7.3pt}\selectfont\color{darkgray}15\hspace{1.8mm}}\fontsize{6.6pt}{7.5pt}\selectfont\strut (novembre) | Tin of black\strut &  \\
      \leavevmode\fontsize{6.6pt}{7.5pt}\selectfont\strut (Nugget) boot polish : 30 c ]\strut &  \\
      \cline{1-1}
      \leavevmode\fontsize{6.6pt}{7.5pt}\selectfont\strut \underline{2 Samedi.}\strut &  \\
      \leavevmode\fontsize{6.6pt}{7.5pt}\selectfont\strut Nuit, gelé très fortt. Temps sec \& froid.\strut & \smash{\parbox[t]{39mm}{\fontsize{6.2pt}{7.0pt}\selectfont\raggedright\color{black!88}\textbf{18 fortt} \textit{abbr. pro} fortement.}} \\
      \leavevmode\fontsize{6.6pt}{7.5pt}\selectfont\strut Pas de bise.\strut &  \\
      \leavevmode\llap{\fontsize{6.4pt}{7.3pt}\selectfont\color{darkgray}20\hspace{1.8mm}}\fontsize{6.6pt}{7.5pt}\selectfont\strut Tout le jour, resté à Acheux.\strut &  \\
      \leavevmode\fontsize{6.6pt}{7.5pt}\selectfont\strut Vers 4h après-midi, Cap. Williams, arrivé tout\strut & \smash{\parbox[t]{39mm}{\fontsize{6.2pt}{7.0pt}\selectfont\raggedright\color{black!88}\textbf{21 Cap.} \textit{abbr. pro} Capitaine. — \textbf{22 Lt} \textit{abbr. pro} Lieutenant.}} \\
      \leavevmode\fontsize{6.6pt}{7.5pt}\selectfont\strut jeune, mince. Il vient remplacer Lt Curtis. Étonnement profond.\strut &  \\
      \leavevmode\fontsize{6.6pt}{7.5pt}\selectfont\strut Curtis, à Martinsart. A son retour, tant\strut &  \\
      \leavevmode\fontsize{6.6pt}{7.5pt}\selectfont\strut surpris : visibt découragé — que faire?\strut &  \\
      \leavevmode\llap{\fontsize{6.4pt}{7.3pt}\selectfont\color{darkgray}25\hspace{1.8mm}}\fontsize{6.6pt}{7.5pt}\selectfont\strut Soir, venu au bureau mission \&\strut &  \\
      \leavevmode\fontsize{6.6pt}{7.5pt}\selectfont\strut vu Cap. H.\strut & \smash{\parbox[t]{39mm}{\fontsize{6.2pt}{7.0pt}\selectfont\raggedright\color{black!88}\textbf{26 Cap.} \textit{abbr. pro} Capitaine.}} \\
\end{tabular}}%
\par\vspace*{\fill}
% --- Registres de bas de page (1. Critique interne [si pertinent] ; 2. Loci paralleli en latin ; 3. Commentaire historique en retrait) ---
\noindent\rule{\textwidth}{0.5pt}\par\vspace{0.65mm}
{\fontsize{6.6pt}{7.7pt}\selectfont%
\noindent\normalfont \textbf{Renvoi interne :} L'appel de note \texttt{(1)} placé à la fin du récit du 1\textsuperscript{er} décembre (l.~14) renvoie à la mention d'achat de tabac pour le capitaine Huguet inscrite au bas du feuillet (l.~26).\par\vspace{0.5mm}
\noindent\rule{22mm}{0.3pt}\par\vspace{0.45mm}
\noindent\normalfont\color{black!92} \textbf{4} (\textit{Buire}) cf. infra f\textsuperscript{o}~90~v\textsuperscript{o} (p.~\pageref{fol:181}), f\textsuperscript{o}~91~r\textsuperscript{o} (p.~\pageref{fol:182}). — \textbf{4} (\textit{Franvillers}) cf. infra f\textsuperscript{o}~58~v\textsuperscript{o} (p.~\pageref{fol:117}). — \textbf{6} (\textit{Ribemont}) cf. infra f\textsuperscript{o}~76~v\textsuperscript{o} (p.~\pageref{fol:153}), f\textsuperscript{o}~84~v\textsuperscript{o} (p.~\pageref{fol:169}).\par\vspace{0.55mm}
\noindent\rule{38mm}{0.4pt}\par\vspace{0.5mm}
\noindent\hspace*{5mm}\parbox{\dimexpr\textwidth-5mm\relax}{\fontsize{6.6pt}{7.7pt}\selectfont\color{black!84}En se rendant à \textbf{Franvillers} et \textbf{Buire-sur-l'Ancre} le 1\textsuperscript{er} décembre 1916, Hippolyte Combes revoit la vallée de \textbf{Ribemont-sur-Ancre} et \textbf{Heilly}, où il avait cantonné avec la 18\textsuperscript{e} Division durant l'hiver 1915–1916. Le samedi 2 décembre survient à Acheux l'arrivée du jeune \textbf{capitaine Williams} (\textit{Foraging Officer}), appelé à remplacer le lieutenant Curtis.}\par}
\newpage
\phantomsection\label{fac:105}
% --- En-tête du fac-similé (Page impaire / droite : Photographie du manuscrit) ---
\noindent\rule{\textwidth}{0.5pt}\par\vspace{0.55mm}
\noindent\textbf{\fontsize{6.8pt}{7.8pt}\selectfont FAC-SIMILÉ EN VIS-À-VIS — FOLIO 52 VERSO} \hfill \textbf{\fontsize{6.8pt}{7.8pt}\selectfont DÉCEMBRE 1916}\par\vspace{0.25mm}
\noindent{\fontsize{6.8pt}{8.0pt}\selectfont\textit{Livre-Journal d'Hippolyte Combes} — 1er – 2 décembre 1916 (Acheux, Franvillers, Buire-sur-l'Ancre, Ribemont, Heilly, Toutencourt (Somme))}\par\vspace{0.55mm}
\noindent\rule{\textwidth}{0.35pt}\par\vspace*{\fill}
{\centering
  \setlength{\fboxsep}{1.2mm}%
  \setlength{\fboxrule}{0.4pt}%
  \fcolorbox{gray!55}{white}{%
    \includegraphics[width=0.95\textwidth,height=0.81\textheight,keepaspectratio]{\facsimiledir/20251126_104946.jpg}%
  }\par}
\vspace*{\fill}
\newpage

% =============================================================================
% VIS-A-VIS SPREAD : FOLIO 53 RECTO (20251126_105005.jpg)
% =============================================================================
\cleardoubleevenpage
\renewcommand{\currentmonthyear}{Décembre 1916}
\renewcommand{\currentfoliolabel}{Folio 53 recto}
\phantomsection\label{fol:106}
% --- En-tête du folio (Page paire / gauche : Transcription, apparat & registres de bas de page) ---
\noindent\rule{\textwidth}{0.5pt}\par\vspace{0.55mm}
\noindent\textbf{\footnotesize FOLIO 53 RECTO} \hfill \textbf{\footnotesize DÉCEMBRE 1916}\par\vspace{0.25mm}
\noindent{\fontsize{6.9pt}{8.0pt}\selectfont Longue randonnée jusqu'à Corbie, Villers-Bretonneux et Marcelcave dans la brume glaciale}\par\vspace{0.2mm}
\noindent{\fontsize{6.9pt}{8.0pt}\selectfont\textit{3 – 4 décembre 1916} \hfill \textit{Acheux, Septenville, Hérissart, Corbie, Villers-Bretonneux, Marcelcave, Warfusée-Abancourt, Arquèves, Louvencourt (Somme)}}\par\vspace{0.5mm}
\noindent\rule{\textwidth}{0.35pt}\par\vspace{0.8mm}
% --- Corps de la transcription diplomatique et apparat marginal en prolongement ---
{\renewcommand{\arraystretch}{0.83}%
\noindent\hspace*{3.5mm}\begin{tabular}{@{}p{87mm}@{\hspace{3.5mm}}p{39mm}@{}}
      \leavevmode\fontsize{6.6pt}{7.5pt}\selectfont\strut \underline{3 Dim.} — \hfill{} \underline{Décbre 1916}\strut & \smash{\parbox[t]{39mm}{\fontsize{6.2pt}{7.0pt}\selectfont\raggedright\color{black!88}\textbf{1 Dim.} \textit{scil.} Dimanche.}} \\
      \leavevmode\fontsize{6.6pt}{7.5pt}\selectfont\strut Nuit très froide. Gelé 1 peu moins que nuit\strut &  \\
      \leavevmode\fontsize{6.6pt}{7.5pt}\selectfont\strut précédente.\strut &  \\
      \leavevmode\fontsize{6.6pt}{7.5pt}\selectfont\strut Journée bien remplie \& fatigante. Partis vers\strut &  \\
      \leavevmode\llap{\fontsize{6.4pt}{7.3pt}\selectfont\color{darkgray}5\hspace{1.8mm}}\fontsize{6.6pt}{7.5pt}\selectfont\strut 8h 1/2 m. et rentrés seulement à 6h 1/2 s.\strut & \smash{\parbox[t]{39mm}{\fontsize{6.2pt}{7.0pt}\selectfont\raggedright\color{black!88}\textbf{5 1/2} \textit{scil.} demi / moitié. \textbf{5 m.} \textit{abbr. pro} matin.}} \\
      \leavevmode\fontsize{6.6pt}{7.5pt}\selectfont\strut Longues courses en auto, matin \& soir :\strut &  \\
      \leavevmode\fontsize{6.6pt}{7.5pt}\selectfont\strut Septenville et Hérissart. Déjeuné ch. Lieutt Walsh.\strut & \smash{\parbox[t]{39mm}{\fontsize{6.2pt}{7.0pt}\selectfont\raggedright\color{black!88}\textbf{7 ch.} \textit{abbr. pro} chez. \textbf{7 Lieutt} \textit{abbr. pro} Lieutenant. — \textbf{8 j'à} \textit{abbr. pro} jusqu'à. \textbf{8 lunch} \textit{angl. in ms.}}} \\
      \leavevmode\fontsize{6.6pt}{7.5pt}\selectfont\strut Après lunch, venus j'à Marcelcave. Passé\strut &  \\
      \leavevmode\fontsize{6.6pt}{7.5pt}\selectfont\strut à Corbie et Villers-Bretonneux.\strut &  \\
      \leavevmode\llap{\fontsize{6.4pt}{7.3pt}\selectfont\color{darkgray}10\hspace{1.8mm}}\fontsize{6.6pt}{7.5pt}\selectfont\strut Retour par Warfusée-Abancourt.\strut &  \\
      \leavevmode\fontsize{6.6pt}{7.5pt}\selectfont\strut Tt jour, brume glaciale enveloppé choses.\strut & \smash{\parbox[t]{39mm}{\fontsize{6.2pt}{7.0pt}\selectfont\raggedright\color{black!88}\textbf{11 Tt} \textit{abbr. pro} Tout.}} \\
      \leavevmode\fontsize{6.6pt}{7.5pt}\selectfont\strut Moulins squelettiques au 1/2 des champs...\strut & \smash{\parbox[t]{39mm}{\fontsize{6.2pt}{7.0pt}\selectfont\raggedright\color{black!88}\textbf{12 1/2} \textit{scil.} demi / moitié.}} \\
      \leavevmode\fontsize{6.6pt}{7.5pt}\selectfont\strut Heures mornes en auto : pu ruminer \&\strut &  \\
      \leavevmode\fontsize{6.6pt}{7.5pt}\selectfont\strut observer à loisir...\strut &  \\
      \leavevmode\llap{\fontsize{6.4pt}{7.3pt}\selectfont\color{darkgray}15\hspace{1.8mm}}\fontsize{6.6pt}{7.5pt}\selectfont\strut [ — A Hérissart, 2 bocks : 30 c\strut & \smash{\parbox[t]{39mm}{\fontsize{6.2pt}{7.0pt}\selectfont\raggedright\color{black!88}\textbf{15} \textit{in marg. cum unco} (comptes).}} \\
      \leavevmode\fontsize{6.6pt}{7.5pt}\selectfont\strut Au bureau mission XIIIe corps, Acheux,\strut &  \\
      \leavevmode\fontsize{6.6pt}{7.5pt}\selectfont\strut touché solde : 78 f 75 c (prochaine, le 17) ]\strut &  \\
      \cline{1-1}
      \leavevmode\fontsize{6.6pt}{7.5pt}\selectfont\strut \underline{4 L.} Matinée, dégelé rapidt et\strut & \smash{\parbox[t]{39mm}{\fontsize{6.2pt}{7.0pt}\selectfont\raggedright\color{black!88}\textbf{18 L.} \textit{scil.} Lundi. \textbf{18 rapidt} \textit{abbr. pro} rapidement.}} \\
      \leavevmode\fontsize{6.6pt}{7.5pt}\selectfont\strut routes redevenues boueuses.\strut &  \\
      \leavevmode\llap{\fontsize{6.4pt}{7.3pt}\selectfont\color{darkgray}20\hspace{1.8mm}}\fontsize{6.6pt}{7.5pt}\selectfont\strut Tout le matin, ds la grange de Mme Guerle,\strut & \smash{\parbox[t]{39mm}{\fontsize{6.2pt}{7.0pt}\selectfont\raggedright\color{black!88}\textbf{20 ds} \textit{abbr. pro} dans. \textbf{20 Mme} \textit{abbr. pro} Madame.}} \\
      \leavevmode\fontsize{6.6pt}{7.5pt}\selectfont\strut soleil brillé 1 peu —\strut &  \\
      \leavevmode\fontsize{6.6pt}{7.5pt}\selectfont\strut Après-midi, courses en auto à\strut &  \\
      \leavevmode\fontsize{6.6pt}{7.5pt}\selectfont\strut ferme de Belle-Église (Arquèves) et à Louvencourt.\strut &  \\
      \leavevmode\fontsize{6.6pt}{7.5pt}\selectfont\strut Rentré pour le thé et dès lors liberté.\strut &  \\
      \leavevmode\llap{\fontsize{6.4pt}{7.3pt}\selectfont\color{darkgray}25\hspace{1.8mm}}\fontsize{6.6pt}{7.5pt}\selectfont\strut Lt Walsh \& son interprète, invités à dîner.\strut & \smash{\parbox[t]{39mm}{\fontsize{6.2pt}{7.0pt}\selectfont\raggedright\color{black!88}\textbf{25 Lt} \textit{abbr. pro} Lieutenant. — \textbf{26 tt} \textit{abbr. pro} tout. \textbf{26 Capt.} \textit{angl.} Captain.}} \\
      \leavevmode\fontsize{6.6pt}{7.5pt}\selectfont\strut Capt. Williams, charmant et de tt à tt sympathique.\strut &  \\
\end{tabular}}%
\par\vspace*{\fill}
% --- Registres de bas de page (1. Critique interne [si pertinent] ; 2. Loci paralleli en latin ; 3. Commentaire historique en retrait) ---
\noindent\rule{\textwidth}{0.5pt}\par\vspace{0.65mm}
{\fontsize{6.6pt}{7.7pt}\selectfont%
\noindent\normalfont\color{black!92} \textbf{7} (\textit{Walsh}) cf. supra f\textsuperscript{o}~49~v\textsuperscript{o} (p.~\pageref{fol:99}), f\textsuperscript{o}~51~v\textsuperscript{o} (p.~\pageref{fol:103}), infra f\textsuperscript{o}~60~r\textsuperscript{o} (p.~\pageref{fol:120}). — \textbf{7} (\textit{Hérissart}) cf. supra f\textsuperscript{o}~47~v\textsuperscript{o} (p.~\pageref{fol:95}), f\textsuperscript{o}~52~r\textsuperscript{o} (p.~\pageref{fol:104}), infra f\textsuperscript{o}~54~r\textsuperscript{o} (p.~\pageref{fol:108}), f\textsuperscript{o}~62~r\textsuperscript{o} (p.~\pageref{fol:124}), \textit{et passim}. — \textbf{16} (\textit{Acheux}) cf. supra f\textsuperscript{o}~21~v\textsuperscript{o} (p.~\pageref{fol:43}), f\textsuperscript{o}~52~v\textsuperscript{o} (p.~\pageref{fol:105}), infra f\textsuperscript{o}~53~v\textsuperscript{o} (p.~\pageref{fol:107}), f\textsuperscript{o}~54~r\textsuperscript{o} (p.~\pageref{fol:108}), \textit{et passim}.\par\vspace{0.55mm}
\noindent\rule{38mm}{0.4pt}\par\vspace{0.5mm}
\noindent\hspace*{5mm}\parbox{\dimexpr\textwidth-5mm\relax}{\fontsize{6.6pt}{7.7pt}\selectfont\color{black!84}Le dimanche 3 décembre 1916, après un déjeuner à Hérissart chez le \textbf{lieutenant Walsh}, une longue tournée automobile entraîne Hippolyte Combes au sud de la Somme par \textbf{Corbie} et \textbf{Villers-Bretonneux} jusqu'à \textbf{Marcelcave} et \textbf{Warfusée-Abancourt} (commune de Lamotte-Warfusée, sur l'axe Amiens--Péronne), avant le paiement de sa quinzaine (78~fr.~75) au bureau de la Mission du XIII\textsuperscript{e} Corps à Acheux.}\par}
\newpage
\phantomsection\label{fac:106}
% --- En-tête du fac-similé (Page impaire / droite : Photographie du manuscrit) ---
\noindent\rule{\textwidth}{0.5pt}\par\vspace{0.55mm}
\noindent\textbf{\fontsize{6.8pt}{7.8pt}\selectfont FAC-SIMILÉ EN VIS-À-VIS — FOLIO 53 RECTO} \hfill \textbf{\fontsize{6.8pt}{7.8pt}\selectfont DÉCEMBRE 1916}\par\vspace{0.25mm}
\noindent{\fontsize{6.8pt}{8.0pt}\selectfont\textit{Livre-Journal d'Hippolyte Combes} — 3 – 4 décembre 1916 (Acheux, Septenville, Hérissart, Corbie, Villers-Bretonneux, Marcelcave, Warfusée-Abancourt, Arquèves, Louvencourt (Somme))}\par\vspace{0.55mm}
\noindent\rule{\textwidth}{0.35pt}\par\vspace*{\fill}
{\centering
  \setlength{\fboxsep}{1.2mm}%
  \setlength{\fboxrule}{0.4pt}%
  \fcolorbox{gray!55}{white}{%
    \includegraphics[width=0.95\textwidth,height=0.81\textheight,keepaspectratio]{\facsimiledir/20251126_105005.jpg}%
  }\par}
\vspace*{\fill}
\newpage

% =============================================================================
% VIS-A-VIS SPREAD : FOLIO 53 VERSO (20251126_105017.jpg)
% =============================================================================
\cleardoubleevenpage
\renewcommand{\currentmonthyear}{Décembre 1916}
\renewcommand{\currentfoliolabel}{Folio 53 verso}
\phantomsection\label{fol:107}
% --- En-tête du folio (Page paire / gauche : Transcription, apparat & registres de bas de page) ---
\noindent\rule{\textwidth}{0.5pt}\par\vspace{0.55mm}
\noindent\textbf{\footnotesize FOLIO 53 VERSO} \hfill \textbf{\footnotesize DÉCEMBRE 1916}\par\vspace{0.25mm}
\noindent{\fontsize{6.9pt}{8.0pt}\selectfont Course jusqu'aux ruines de Bayencourt et Coigneux, départ de Curtis et confidences de Williams}\par\vspace{0.2mm}
\noindent{\fontsize{6.9pt}{8.0pt}\selectfont\textit{5 – 6 décembre 1916} \hfill \textit{Acheux, Authie, Bayencourt, Coigneux, Louvencourt, Raincheval (Somme / Pas-de-Calais)}}\par\vspace{0.5mm}
\noindent\rule{\textwidth}{0.35pt}\par\vspace{0.8mm}
% --- Corps de la transcription diplomatique et apparat marginal en prolongement ---
{\renewcommand{\arraystretch}{0.87}%
\noindent\hspace*{3.5mm}\begin{tabular}{@{}p{87mm}@{\hspace{3.5mm}}p{39mm}@{}}
      \leavevmode\fontsize{6.9pt}{7.9pt}\selectfont\strut \underline{5 Mardi.} \hfill{} \underline{Décbre 1916}\strut &  \\
      \leavevmode\fontsize{6.9pt}{7.9pt}\selectfont\strut Matinée, liberté complète.\strut &  \\
      \leavevmode\fontsize{6.9pt}{7.9pt}\selectfont\strut Après-midi, course à Authie, puis à\strut &  \\
      \leavevmode\fontsize{6.9pt}{7.9pt}\selectfont\strut Bayencourt, près d'Hébuterne, à 3 kilomètres\strut &  \\
      \leavevmode\llap{\fontsize{6.4pt}{7.3pt}\selectfont\color{darkgray}5\hspace{1.8mm}}\fontsize{6.9pt}{7.9pt}\selectfont\strut de Gommecourt (Boches). Dépassé le village. Champs\strut &  \\
      \leavevmode\fontsize{6.9pt}{7.9pt}\selectfont\strut dévastés. Obus non éclatés... Village de Coigneux,\strut &  \\
      \leavevmode\fontsize{6.9pt}{7.9pt}\selectfont\strut bp de toits squelettiques, + de tuiles, à la\strut & \smash{\parbox[t]{39mm}{\fontsize{6.2pt}{7.0pt}\selectfont\raggedright\color{black!88}\textbf{7 bp} \textit{abbr. pro} beaucoup.}} \\
      \leavevmode\fontsize{6.9pt}{7.9pt}\selectfont\strut suite de bombardement.\strut &  \\
      \leavevmode\fontsize{6.9pt}{7.9pt}\selectfont\strut Brume épaisse et, par moments, pluie\strut &  \\
      \leavevmode\llap{\fontsize{6.4pt}{7.3pt}\selectfont\color{darkgray}10\hspace{1.8mm}}\fontsize{6.9pt}{7.9pt}\selectfont\strut fine. Après dîner, adieux avec Lieutt\strut & \smash{\parbox[t]{39mm}{\fontsize{6.2pt}{7.0pt}\selectfont\raggedright\color{black!88}\textbf{10 Lieutt} \textit{abbr. pro} Lieutenant.}} \\
      \leavevmode\fontsize{6.9pt}{7.9pt}\selectfont\strut Curtis qui doit repartir, demain matin, pour\strut &  \\
      \leavevmode\fontsize{6.9pt}{7.9pt}\selectfont\strut rejoindre la 152 Coy A.S.C.\strut & \smash{\parbox[t]{39mm}{\fontsize{6.2pt}{7.0pt}\selectfont\raggedright\color{black!88}\textbf{12 Coy} \textit{angl.} Company. \textbf{12 A.S.C.} \textit{angl.} Army Service Corps. — \textbf{13 Merc.} \textit{scil.} Mercredi. \textbf{13 Lt} \textit{abbr. pro} Lieutenant. — \textbf{14 1/2} \textit{scil.} demi / moitié. \textbf{14 m.} \textit{abbr. pro} matin. — \textbf{17 off.} \textit{abbr. pro} officier(s). \textbf{17 Cap.} \textit{abbr. pro} Capitaine.}} \\
      \cline{1-1}
      \leavevmode\fontsize{6.9pt}{7.9pt}\selectfont\strut \underline{6 Merc.} Levé 6h m. — Lt Curtis, parti vers\strut &  \\
      \leavevmode\fontsize{6.9pt}{7.9pt}\selectfont\strut 7h 1/2 m. — \& sans laisser de regret — A brainless\strut &  \\
      \leavevmode\llap{\fontsize{6.4pt}{7.3pt}\selectfont\color{darkgray}15\hspace{1.8mm}}\fontsize{6.9pt}{7.9pt}\selectfont\strut ass. Matin, course à Louvencourt, Raincheval\strut &  \\
      \leavevmode\fontsize{6.9pt}{7.9pt}\selectfont\strut etc.\strut &  \\
      \leavevmode\fontsize{6.9pt}{7.9pt}\selectfont\strut Nouvel off., Cap. Williams, 20 ans! Après\strut &  \\
      \leavevmode\fontsize{6.9pt}{7.9pt}\selectfont\strut repas soir, s'est 1 peu raconté : avant\strut &  \\
      \leavevmode\fontsize{6.9pt}{7.9pt}\selectfont\strut la guerre, était office boy à 50 £ a year.\strut &  \\
\end{tabular}}%
\par\vspace*{\fill}
% --- Registres de bas de page (1. Critique interne [si pertinent] ; 2. Loci paralleli en latin ; 3. Commentaire historique en retrait) ---
\noindent\rule{\textwidth}{0.5pt}\par\vspace{0.65mm}
{\fontsize{6.6pt}{7.7pt}\selectfont%
\noindent\normalfont\color{black!92} \textbf{3} (\textit{Authie}) cf. supra f\textsuperscript{o}~50~r\textsuperscript{o} (p.~\pageref{fol:100}), f\textsuperscript{o}~51~r\textsuperscript{o} (p.~\pageref{fol:102}), infra f\textsuperscript{o}~54~r\textsuperscript{o} (p.~\pageref{fol:108}), f\textsuperscript{o}~54~v\textsuperscript{o} (p.~\pageref{fol:109}), \textit{et passim}. — \textbf{11} (\textit{Curtis}) cf. supra f\textsuperscript{o}~1~r\textsuperscript{o} (p.~\pageref{fol:2}), f\textsuperscript{o}~52~v\textsuperscript{o} (p.~\pageref{fol:105}), \textit{et passim}. — \textbf{12} (\textit{152 Coy}) cf. supra f\textsuperscript{o}~33~r\textsuperscript{o} (p.~\pageref{fol:66}), f\textsuperscript{o}~48~r\textsuperscript{o} (p.~\pageref{fol:96}), \textit{et passim}.\par\vspace{0.55mm}
\noindent\rule{38mm}{0.4pt}\par\vspace{0.5mm}
\noindent\hspace*{5mm}\parbox{\dimexpr\textwidth-5mm\relax}{\fontsize{6.6pt}{7.7pt}\selectfont\color{black!84}Le mardi 5 décembre 1916, la tournée d'inspection s'avance au nord d'Authie jusqu'aux villages dévastés de \textbf{Bayencourt} et \textbf{Coigneux} (Somme, canton d'Acheux), situés immédiatement en arrière des tranchées d'\textbf{Hébuterne} et à 3~km du saillant allemand de \textbf{Gommecourt} (Pas-de-Calais).}\par}
\newpage
\phantomsection\label{fac:107}
% --- En-tête du fac-similé (Page impaire / droite : Photographie du manuscrit) ---
\noindent\rule{\textwidth}{0.5pt}\par\vspace{0.55mm}
\noindent\textbf{\fontsize{6.8pt}{7.8pt}\selectfont FAC-SIMILÉ EN VIS-À-VIS — FOLIO 53 VERSO} \hfill \textbf{\fontsize{6.8pt}{7.8pt}\selectfont DÉCEMBRE 1916}\par\vspace{0.25mm}
\noindent{\fontsize{6.8pt}{8.0pt}\selectfont\textit{Livre-Journal d'Hippolyte Combes} — 5 – 6 décembre 1916 (Acheux, Authie, Bayencourt, Coigneux, Louvencourt, Raincheval (Somme / Pas-de-Calais))}\par\vspace{0.55mm}
\noindent\rule{\textwidth}{0.35pt}\par\vspace*{\fill}
{\centering
  \setlength{\fboxsep}{1.2mm}%
  \setlength{\fboxrule}{0.4pt}%
  \fcolorbox{gray!55}{white}{%
    \includegraphics[width=0.95\textwidth,height=0.81\textheight,keepaspectratio]{\facsimiledir/20251126_105017.jpg}%
  }\par}
\vspace*{\fill}
\newpage

% =============================================================================
% VIS-A-VIS SPREAD : FOLIO 54 RECTO (20251126_105046.jpg)
% =============================================================================
\cleardoubleevenpage
\renewcommand{\currentmonthyear}{Décembre 1916}
\renewcommand{\currentfoliolabel}{Folio 54 recto}
\phantomsection\label{fol:108}
% --- En-tête du folio (Page paire / gauche : Transcription, apparat & registres de bas de page) ---
\noindent\rule{\textwidth}{0.5pt}\par\vspace{0.55mm}
\noindent\textbf{\footnotesize FOLIO 54 RECTO} \hfill \textbf{\footnotesize DÉCEMBRE 1916}\par\vspace{0.25mm}
\noindent{\fontsize{6.9pt}{8.0pt}\selectfont Les enfants de la fermière belge de Belle-Église et le village désert d'Englebelmer}\par\vspace{0.2mm}
\noindent{\fontsize{6.9pt}{8.0pt}\selectfont\textit{7 – 8 décembre 1916} \hfill \textit{Acheux, Arquèves (Belle-Église), Hédauville, Englebelmer, Saint-Léger-lès-Authie (Somme)}}\par\vspace{0.5mm}
\noindent\rule{\textwidth}{0.35pt}\par\vspace{0.8mm}
% --- Corps de la transcription diplomatique et apparat marginal en prolongement ---
{\renewcommand{\arraystretch}{0.83}%
\noindent\hspace*{3.5mm}\begin{tabular}{@{}p{87mm}@{\hspace{3.5mm}}p{39mm}@{}}
      \leavevmode\fontsize{6.6pt}{7.5pt}\selectfont\strut \underline{7 Jeudi} — — Décbre 1916 —\strut &  \\
      \leavevmode\fontsize{6.6pt}{7.5pt}\selectfont\strut Matin \& après-midi, en\strut &  \\
      \leavevmode\fontsize{6.6pt}{7.5pt}\selectfont\strut auto. — Matin, à Belle-Église, Arquèves,\strut &  \\
      \leavevmode\fontsize{6.6pt}{7.5pt}\selectfont\strut où dû attendre lgts retour du Cap. W.,\strut & \smash{\parbox[t]{39mm}{\fontsize{6.2pt}{7.0pt}\selectfont\raggedright\color{black!88}\textbf{4 lgts} \textit{abbr. pro} longtemps. \textbf{4 Cap. W.} \textit{scil.} Capitaine Williams. — \textbf{5 j'à} \textit{abbr. pro} jusqu'à. — \textbf{6 Houtteman} \textit{lect. incert. in ms.} \textbf{6 Mme} \textit{abbr. pro} Madame. — \textbf{8 ds} \textit{abbr. pro} dans. — \textbf{10 ds} \textit{abbr. pro} dans. — \textbf{11 lgts} \textit{abbr. pro} longtemps. — \textbf{12 Town Major} \textit{angl. in ms.}}} \\
      \leavevmode\llap{\fontsize{6.4pt}{7.3pt}\selectfont\color{darkgray}5\hspace{1.8mm}}\fontsize{6.6pt}{7.5pt}\selectfont\strut qui allé j'à Raincheval. Amusé par les\strut &  \\
      \leavevmode\fontsize{6.6pt}{7.5pt}\selectfont\strut enfants de la fermière belge, Mme Houtteman [?].\strut &  \\
      \leavevmode\fontsize{6.6pt}{7.5pt}\selectfont\strut Après-midi, venus à Englebelmer, par\strut &  \\
      \leavevmode\fontsize{6.6pt}{7.5pt}\selectfont\strut Hédauville. Resté près d'1 h. ds ce village\strut &  \\
      \leavevmode\fontsize{6.6pt}{7.5pt}\selectfont\strut à peu près désert : plus d'habitants. Rien\strut &  \\
      \leavevmode\llap{\fontsize{6.4pt}{7.3pt}\selectfont\color{darkgray}10\hspace{1.8mm}}\fontsize{6.6pt}{7.5pt}\selectfont\strut que des Anglais, qui restent ds leurs cantonnements.\strut &  \\
      \leavevmode\fontsize{6.6pt}{7.5pt}\selectfont\strut Attendu lgts en auto, devant ferme où le\strut &  \\
      \leavevmode\fontsize{6.6pt}{7.5pt}\selectfont\strut Town Major et le Daily Mail. A peu de\strut &  \\
      \leavevmode\fontsize{6.6pt}{7.5pt}\selectfont\strut distance, canons anglais ne cessaient point\strut &  \\
      \leavevmode\fontsize{6.6pt}{7.5pt}\selectfont\strut de tonner : après coup, vibrations métalliques\strut &  \\
      \leavevmode\llap{\fontsize{6.4pt}{7.3pt}\selectfont\color{darkgray}15\hspace{1.8mm}}\fontsize{6.6pt}{7.5pt}\selectfont\strut répercutées à travers l'espace — comme gémissement.\strut &  \\
      \leavevmode\fontsize{6.6pt}{7.5pt}\selectfont\strut Dîné seul : off. allé à Hérissart. Rentré\strut & \smash{\parbox[t]{39mm}{\fontsize{6.2pt}{7.0pt}\selectfont\raggedright\color{black!88}\textbf{16 off.} \textit{abbr. pro} officier(s). — \textbf{17 ds} \textit{abbr. pro} dans.}} \\
      \leavevmode\fontsize{6.6pt}{7.5pt}\selectfont\strut tôt ds ma chambre. — Rue silencieuse.\strut &  \\
      \leavevmode\fontsize{6.6pt}{7.5pt}\selectfont\strut Canonnade, intermittente.\strut &  \\
      \leavevmode\fontsize{6.6pt}{7.5pt}\selectfont\strut Brume épaisse s'étend de + en +.\strut &  \\
      \cline{1-1}
      \leavevmode\llap{\fontsize{6.4pt}{7.3pt}\selectfont\color{darkgray}20\hspace{1.8mm}}\fontsize{6.6pt}{7.5pt}\selectfont\strut \underline{8 Ve.} Tjrs même temps. Ciel caché par\strut & \smash{\parbox[t]{39mm}{\fontsize{6.2pt}{7.0pt}\selectfont\raggedright\color{black!88}\textbf{20 Ve.} \textit{scil.} Vendredi. \textbf{20 Tjrs} \textit{abbr. pro} Toujours.}} \\
      \leavevmode\fontsize{6.6pt}{7.5pt}\selectfont\strut brouillard. Bruine intermittente. Routes, état\strut &  \\
      \leavevmode\fontsize{6.6pt}{7.5pt}\selectfont\strut affreux.\strut &  \\
      \leavevmode\fontsize{6.6pt}{7.5pt}\selectfont\strut Matinée, course à St-Léger-lès-Authie,\strut & \smash{\parbox[t]{39mm}{\fontsize{6.2pt}{7.0pt}\selectfont\raggedright\color{black!88}\textbf{23 St} \textit{abbr. pro} Saint.}} \\
      \leavevmode\fontsize{6.6pt}{7.5pt}\selectfont\strut sous la pluie.\strut &  \\
      \leavevmode\llap{\fontsize{6.4pt}{7.3pt}\selectfont\color{darkgray}25\hspace{1.8mm}}\fontsize{6.6pt}{7.5pt}\selectfont\strut Après-midi, ferme de Belle-Église, puis\strut &  \\
      \leavevmode\fontsize{6.6pt}{7.5pt}\selectfont\strut Arquèves : on a commencé à battre paille.\strut &  \\
\end{tabular}}%
\par\vspace*{\fill}
% --- Registres de bas de page (1. Critique interne [si pertinent] ; 2. Loci paralleli en latin ; 3. Commentaire historique en retrait) ---
\noindent\rule{\textwidth}{0.5pt}\par\vspace{0.65mm}
{\fontsize{6.6pt}{7.7pt}\selectfont%
\noindent\normalfont\color{black!92} \textbf{3} (\textit{Arquèves}) cf. supra f\textsuperscript{o}~47~v\textsuperscript{o} (p.~\pageref{fol:95}), f\textsuperscript{o}~53~r\textsuperscript{o} (p.~\pageref{fol:106}), infra f\textsuperscript{o}~55~v\textsuperscript{o} (p.~\pageref{fol:111}), f\textsuperscript{o}~57~v\textsuperscript{o} (p.~\pageref{fol:115}), \textit{et passim}. — \textbf{3} (\textit{Belle-Église}) cf. supra f\textsuperscript{o}~47~v\textsuperscript{o} (p.~\pageref{fol:95}), f\textsuperscript{o}~53~r\textsuperscript{o} (p.~\pageref{fol:106}), infra f\textsuperscript{o}~54~v\textsuperscript{o} (p.~\pageref{fol:109}), f\textsuperscript{o}~55~r\textsuperscript{o} (p.~\pageref{fol:110}), \textit{et passim}. — \textbf{5} (\textit{Raincheval}) cf. supra f\textsuperscript{o}~20~v\textsuperscript{o} (p.~\pageref{fol:41}), f\textsuperscript{o}~53~v\textsuperscript{o} (p.~\pageref{fol:107}), infra f\textsuperscript{o}~55~v\textsuperscript{o} (p.~\pageref{fol:111}), f\textsuperscript{o}~56~v\textsuperscript{o} (p.~\pageref{fol:113}), \textit{et passim}.\par\vspace{0.55mm}
\noindent\rule{38mm}{0.4pt}\par\vspace{0.5mm}
\noindent\hspace*{5mm}\parbox{\dimexpr\textwidth-5mm\relax}{\fontsize{6.6pt}{7.7pt}\selectfont\color{black!84}Les 7 et 8 décembre 1916, Hippolyte Combes supervise la mise en route du battage mécanique de la paille à la ferme de \textbf{Belle-Église} (commune d'\textbf{Arquèves}), exploitée par une famille d'agriculteurs belges, \textbf{M. et Mme Houthouan}.}\par}
\newpage
\phantomsection\label{fac:108}
% --- En-tête du fac-similé (Page impaire / droite : Photographie du manuscrit) ---
\noindent\rule{\textwidth}{0.5pt}\par\vspace{0.55mm}
\noindent\textbf{\fontsize{6.8pt}{7.8pt}\selectfont FAC-SIMILÉ EN VIS-À-VIS — FOLIO 54 RECTO} \hfill \textbf{\fontsize{6.8pt}{7.8pt}\selectfont DÉCEMBRE 1916}\par\vspace{0.25mm}
\noindent{\fontsize{6.8pt}{8.0pt}\selectfont\textit{Livre-Journal d'Hippolyte Combes} — 7 – 8 décembre 1916 (Acheux, Arquèves (Belle-Église), Hédauville, Englebelmer, Saint-Léger-lès-Authie (Somme))}\par\vspace{0.55mm}
\noindent\rule{\textwidth}{0.35pt}\par\vspace*{\fill}
{\centering
  \setlength{\fboxsep}{1.2mm}%
  \setlength{\fboxrule}{0.4pt}%
  \fcolorbox{gray!55}{white}{%
    \includegraphics[width=0.95\textwidth,height=0.81\textheight,keepaspectratio]{\facsimiledir/20251126_105046.jpg}%
  }\par}
\vspace*{\fill}
\newpage

% =============================================================================
% VIS-A-VIS SPREAD : FOLIO 54 VERSO (20251126_105057.jpg)
% =============================================================================
\cleardoubleevenpage
\renewcommand{\currentmonthyear}{Décembre 1916}
\renewcommand{\currentfoliolabel}{Folio 54 verso}
\phantomsection\label{fol:109}
% --- En-tête du folio (Page paire / gauche : Transcription, apparat & registres de bas de page) ---
\noindent\rule{\textwidth}{0.5pt}\par\vspace{0.55mm}
\noindent\textbf{\footnotesize FOLIO 54 VERSO} \hfill \textbf{\footnotesize DÉCEMBRE 1916}\par\vspace{0.25mm}
\noindent{\fontsize{6.9pt}{8.0pt}\selectfont Randonnée pluvieuse jusqu'à Saulty et L'Arbret, autoritarisme de Williams et crise de dysenterie}\par\vspace{0.2mm}
\noindent{\fontsize{6.9pt}{8.0pt}\selectfont\textit{9 – 11 décembre 1916} \hfill \textit{Acheux, Belle-Église, Authie, Orville, Saulty, L'Arbret (Somme / Pas-de-Calais)}}\par\vspace{0.5mm}
\noindent\rule{\textwidth}{0.35pt}\par\vspace{0.8mm}
% --- Corps de la transcription diplomatique et apparat marginal en prolongement ---
{\renewcommand{\arraystretch}{0.87}%
\noindent\hspace*{3.5mm}\begin{tabular}{@{}p{87mm}@{\hspace{3.5mm}}p{39mm}@{}}
      \leavevmode\fontsize{6.9pt}{7.9pt}\selectfont\strut \underline{9 Samedi} \hfill{} \underline{Décbre 1916}\strut &  \\
      \leavevmode\fontsize{6.9pt}{7.9pt}\selectfont\strut Tout le jour, pluie, et\strut &  \\
      \leavevmode\fontsize{6.9pt}{7.9pt}\selectfont\strut parfois averse.\strut &  \\
      \leavevmode\fontsize{6.9pt}{7.9pt}\selectfont\strut Partis en auto, vers 9h m. Venus à $[\![$Belle-Église$]\!]$ Orville.\strut & \smash{\parbox[t]{39mm}{\fontsize{6.2pt}{7.0pt}\selectfont\raggedright\color{black!88}\textbf{4} $[\![$Belle-Église$]\!]\ $ \textit{del. in ms.} \textbf{4 m.} \textit{abbr. pro} matin. — \textbf{5 Divon} \textit{abbr. pro} Division. \textbf{5 Coy} \textit{angl.} Company. — \textbf{6 Croughton} \textit{lect. incert. in ms.} \textbf{6 Lt} \textit{abbr. pro} Lieutenant. — \textbf{8 j'à} \textit{abbr. pro} jusqu'à. — \textbf{11} \textit{lect. incert. in ms.}}} \\
      \leavevmode\llap{\fontsize{6.4pt}{7.3pt}\selectfont\color{darkgray}5\hspace{1.8mm}}\fontsize{6.9pt}{7.9pt}\selectfont\strut Déjeuné à Authie, au mess de no 3 Coy 31st Divon\strut &  \\
      \leavevmode\fontsize{6.9pt}{7.9pt}\selectfont\strut av. Lt Croughton [?].\strut &  \\
      \leavevmode\fontsize{6.9pt}{7.9pt}\selectfont\strut Après déjeuner, venus à Orville, puis\strut &  \\
      \leavevmode\fontsize{6.9pt}{7.9pt}\selectfont\strut continué route j'à Saulty – L'Arbret —\strut &  \\
      \leavevmode\fontsize{6.9pt}{7.9pt}\selectfont\strut Rentrés nuit noire. Temps affreux.\strut &  \\
      \leavevmode\llap{\fontsize{6.4pt}{7.3pt}\selectfont\color{darkgray}10\hspace{1.8mm}}\fontsize{6.9pt}{7.9pt}\selectfont\strut Très fatigué et tout mouillé. Mangé de\strut &  \\
      \leavevmode\fontsize{6.9pt}{7.9pt}\selectfont\strut suite et venu chambre. Extt [autoritaire ?]\strut &  \\
      \leavevmode\fontsize{6.9pt}{7.9pt}\selectfont\strut ce gosse de Williams... —\strut &  \\
      \cline{1-1}
      \leavevmode\fontsize{6.9pt}{7.9pt}\selectfont\strut \underline{11 L.} — Hier \& auj'd'hui, inoccupé, mais\strut & \smash{\parbox[t]{39mm}{\fontsize{6.2pt}{7.0pt}\selectfont\raggedright\color{black!88}\textbf{13 L.} \textit{scil.} Lundi. \textbf{13 auj'd'hui} \textit{abbr. pro} aujourd'hui. — \textbf{15 Tjrs} \textit{abbr. pro} Toujours.}} \\
      \leavevmode\fontsize{6.9pt}{7.9pt}\selectfont\strut resté office matin \& soir.\strut &  \\
      \leavevmode\llap{\fontsize{6.4pt}{7.3pt}\selectfont\color{darkgray}15\hspace{1.8mm}}\fontsize{6.9pt}{7.9pt}\selectfont\strut Tjrs temps affreux. Pluie \& boue.\strut &  \\
      \leavevmode\fontsize{6.9pt}{7.9pt}\selectfont\strut Nuit de dimanche à lundi, mauvaise\strut &  \\
      \leavevmode\fontsize{6.9pt}{7.9pt}\selectfont\strut digestion. Dysenterie durant tte journée\strut & \smash{\parbox[t]{39mm}{\fontsize{6.2pt}{7.0pt}\selectfont\raggedright\color{black!88}\textbf{17 tte} \textit{abbr. pro} toute.}} \\
      \leavevmode\fontsize{6.9pt}{7.9pt}\selectfont\strut de lundi. N'ai pris aucun aliment. —\strut &  \\
      \leavevmode\fontsize{6.9pt}{7.9pt}\selectfont\strut Tt le jour, langui ds ma chambre\strut & \smash{\parbox[t]{39mm}{\fontsize{6.2pt}{7.0pt}\selectfont\raggedright\color{black!88}\textbf{19 ds} \textit{abbr. pro} dans. \textbf{19 Tt} \textit{abbr. pro} Tout.}} \\
      \leavevmode\llap{\fontsize{6.4pt}{7.3pt}\selectfont\color{darkgray}20\hspace{1.8mm}}\fontsize{6.9pt}{7.9pt}\selectfont\strut ou au bureau...\strut &  \\
      \cline{1-1}
      \leavevmode\fontsize{6.9pt}{7.9pt}\selectfont\strut [ Tin of Wren's Dubbin : 30 c ]\strut & \smash{\parbox[t]{39mm}{\fontsize{6.2pt}{7.0pt}\selectfont\raggedright\color{black!88}\textbf{21} \textit{in marg. cum unco} (comptes). \textbf{21 Dubbin} \textit{angl. in ms.}}} \\
\end{tabular}}%
\par\vspace*{\fill}
% --- Registres de bas de page (1. Critique interne [si pertinent] ; 2. Loci paralleli en latin ; 3. Commentaire historique en retrait) ---
\noindent\rule{\textwidth}{0.5pt}\par\vspace{0.65mm}
{\fontsize{6.6pt}{7.7pt}\selectfont%
\noindent\normalfont\color{black!92} \textbf{4} (\textit{Belle-Église}) cf. supra f\textsuperscript{o}~47~v\textsuperscript{o} (p.~\pageref{fol:95}), f\textsuperscript{o}~54~r\textsuperscript{o} (p.~\pageref{fol:108}), infra f\textsuperscript{o}~55~r\textsuperscript{o} (p.~\pageref{fol:110}), f\textsuperscript{o}~55~v\textsuperscript{o} (p.~\pageref{fol:111}), \textit{et passim}. — \textbf{5} (\textit{Authie}) cf. supra f\textsuperscript{o}~50~r\textsuperscript{o} (p.~\pageref{fol:100}), f\textsuperscript{o}~54~r\textsuperscript{o} (p.~\pageref{fol:108}), infra f\textsuperscript{o}~56~r\textsuperscript{o} (p.~\pageref{fol:112}), f\textsuperscript{o}~59~r\textsuperscript{o} (p.~\pageref{fol:118}), \textit{et passim}. — \textbf{6} (\textit{Lt Croughton}) cf. infra f\textsuperscript{o}~75~r\textsuperscript{o} (p.~\pageref{fol:150}).\par\vspace{0.55mm}
\noindent\rule{38mm}{0.4pt}\par\vspace{0.5mm}
\noindent\hspace*{5mm}\parbox{\dimexpr\textwidth-5mm\relax}{\fontsize{6.6pt}{7.7pt}\selectfont\color{black!84}Le samedi 9 décembre 1916, au cours d'une longue tournée sous la pluie jusqu'en Artois (\textbf{Orville}, \textbf{Saulty} et \textbf{L'Arbret}, Pas-de-Calais, sur la route de Doullens à Arras), Hippolyte Combes déjeune à \textbf{Authie} au mess de la 3\textsuperscript{e} Compagnie du Train divisionnaire de la \textbf{31\textsuperscript{e} Division britannique} (\textit{No.~3 Company, 31st Divisional Train, A.S.C.}, c'est-à-dire la \textit{280th Company A.S.C.}; la 31\textsuperscript{e} Division, rattachée au XIII\textsuperscript{e} Corps, tenait alors le secteur de Serre--Hébuterne) avec le \textbf{lieutenant Croughton}.}\par}
\newpage
\phantomsection\label{fac:109}
% --- En-tête du fac-similé (Page impaire / droite : Photographie du manuscrit) ---
\noindent\rule{\textwidth}{0.5pt}\par\vspace{0.55mm}
\noindent\textbf{\fontsize{6.8pt}{7.8pt}\selectfont FAC-SIMILÉ EN VIS-À-VIS — FOLIO 54 VERSO} \hfill \textbf{\fontsize{6.8pt}{7.8pt}\selectfont DÉCEMBRE 1916}\par\vspace{0.25mm}
\noindent{\fontsize{6.8pt}{8.0pt}\selectfont\textit{Livre-Journal d'Hippolyte Combes} — 9 – 11 décembre 1916 (Acheux, Belle-Église, Authie, Orville, Saulty, L'Arbret (Somme / Pas-de-Calais))}\par\vspace{0.55mm}
\noindent\rule{\textwidth}{0.35pt}\par\vspace*{\fill}
{\centering
  \setlength{\fboxsep}{1.2mm}%
  \setlength{\fboxrule}{0.4pt}%
  \fcolorbox{gray!55}{white}{%
    \includegraphics[width=0.95\textwidth,height=0.81\textheight,keepaspectratio]{\facsimiledir/20251126_105057.jpg}%
  }\par}
\vspace*{\fill}
\newpage

% =============================================================================
% VIS-A-VIS SPREAD : FOLIO 55 RECTO (20251126_105112.jpg)
% =============================================================================
\cleardoubleevenpage
\renewcommand{\currentmonthyear}{Décembre 1916}
\renewcommand{\currentfoliolabel}{Folio 55 recto}
\phantomsection\label{fol:110}
% --- En-tête du folio (Page paire / gauche : Transcription, apparat & registres de bas de page) ---
\noindent\rule{\textwidth}{0.5pt}\par\vspace{0.55mm}
\noindent\textbf{\footnotesize FOLIO 55 RECTO} \hfill \textbf{\footnotesize DÉCEMBRE 1916}\par\vspace{0.25mm}
\noindent{\fontsize{6.9pt}{8.0pt}\selectfont Paysage de neige aux fenêtres d'Acheux et réunion des réfugiés d'Englebelmer à Toutencourt}\par\vspace{0.2mm}
\noindent{\fontsize{6.9pt}{8.0pt}\selectfont\textit{12 – 14 décembre 1916} \hfill \textit{Acheux, Senlis-le-Sec, Toutencourt, Belle-Église, Puchevillers, Harponville (Somme)}}\par\vspace{0.5mm}
\noindent\rule{\textwidth}{0.35pt}\par\vspace{0.8mm}
% --- Corps de la transcription diplomatique et apparat marginal en prolongement ---
{\renewcommand{\arraystretch}{0.87}%
\noindent\hspace*{3.5mm}\begin{tabular}{@{}p{87mm}@{\hspace{3.5mm}}p{39mm}@{}}
      \leavevmode\fontsize{6.9pt}{7.9pt}\selectfont\strut \underline{12 Mardi} \hfill{} \underline{Décbre 1916}\strut &  \\
      \leavevmode\fontsize{6.9pt}{7.9pt}\selectfont\strut Mieux, très sensible. Nuit, levé 2 fois. Mais\strut &  \\
      \leavevmode\fontsize{6.9pt}{7.9pt}\selectfont\strut constipation, durant journée.\strut &  \\
      \leavevmode\fontsize{6.9pt}{7.9pt}\selectfont\strut Au réveil, trouvé rues, toits \& arbres couverts\strut &  \\
      \leavevmode\llap{\fontsize{6.4pt}{7.3pt}\selectfont\color{darkgray}5\hspace{1.8mm}}\fontsize{6.9pt}{7.9pt}\selectfont\strut de neige. Resté lgts à ma fenêtre, s/ devant\strut & \smash{\parbox[t]{39mm}{\fontsize{6.2pt}{7.0pt}\selectfont\raggedright\color{black!88}\textbf{5 s/} \textit{abbr. pro} sur. \textbf{5 lgts} \textit{abbr. pro} longtemps. — \textbf{6 s/} \textit{abbr. pro} sur.}} \\
      \leavevmode\fontsize{6.9pt}{7.9pt}\selectfont\strut \& s/ derrière, à admirer paysage — —\strut &  \\
      \leavevmode\fontsize{6.9pt}{7.9pt}\selectfont\strut Matinée, course à Senlis. Après-midi, liberté.\strut &  \\
      \leavevmode\fontsize{6.9pt}{7.9pt}\selectfont\strut Me sens presque tt à fait guéri.\strut & \smash{\parbox[t]{39mm}{\fontsize{6.2pt}{7.0pt}\selectfont\raggedright\color{black!88}\textbf{8 tt} \textit{abbr. pro} tout.}} \\
      \cline{1-1}
      \leavevmode\fontsize{6.9pt}{7.9pt}\selectfont\strut \underline{13 Me.} Brouillard \& pluie. Routes, lamentables.\strut &  \\
      \leavevmode\llap{\fontsize{6.4pt}{7.3pt}\selectfont\color{darkgray}10\hspace{1.8mm}}\fontsize{6.9pt}{7.9pt}\selectfont\strut Matin, à Toutencourt, mairie, meeting de\strut &  \\
      \leavevmode\fontsize{6.9pt}{7.9pt}\selectfont\strut réfugiés d'Englebelmer, Major Levy...\strut &  \\
      \leavevmode\fontsize{6.9pt}{7.9pt}\selectfont\strut Après-midi, à Belle-Église \& Puchevillers, et\strut &  \\
      \leavevmode\fontsize{6.9pt}{7.9pt}\selectfont\strut Harponville. Secoué auto. Routes, très mauvaises.\strut &  \\
      \cline{1-1}
      \leavevmode\fontsize{6.9pt}{7.9pt}\selectfont\strut \underline{14 J.} Journée entière de liberté. Resté\strut & \smash{\parbox[t]{39mm}{\fontsize{6.2pt}{7.0pt}\selectfont\raggedright\color{black!88}\textbf{14 J.} \textit{scil.} Jeudi.}} \\
      \leavevmode\llap{\fontsize{6.4pt}{7.3pt}\selectfont\color{darkgray}15\hspace{1.8mm}}\fontsize{6.9pt}{7.9pt}\selectfont\strut lgts ds ma chambre : écrit, lu, ou rangé mes\strut & \smash{\parbox[t]{39mm}{\fontsize{6.2pt}{7.0pt}\selectfont\raggedright\color{black!88}\textbf{15 ds} \textit{abbr. pro} dans. \textbf{15 lgts} \textit{abbr. pro} longtemps.}} \\
      \leavevmode\fontsize{6.9pt}{7.9pt}\selectfont\strut affaires.\strut &  \\
      \leavevmode\fontsize{6.9pt}{7.9pt}\selectfont\strut Tt le jour, ciel brumeux et pluie fine.\strut & \smash{\parbox[t]{39mm}{\fontsize{6.2pt}{7.0pt}\selectfont\raggedright\color{black!88}\textbf{17 Tt} \textit{abbr. pro} Tout.}} \\
      \leavevmode\fontsize{6.9pt}{7.9pt}\selectfont\strut Température, pas froide. Pu écrire ds ma\strut & \smash{\parbox[t]{39mm}{\fontsize{6.2pt}{7.0pt}\selectfont\raggedright\color{black!88}\textbf{18 ds} \textit{abbr. pro} dans.}} \\
      \leavevmode\fontsize{6.9pt}{7.9pt}\selectfont\strut chambre, sans feu.\strut &  \\
      \leavevmode\llap{\fontsize{6.4pt}{7.3pt}\selectfont\color{darkgray}20\hspace{1.8mm}}\fontsize{6.9pt}{7.9pt}\selectfont\strut Nuit, extt noire. Routes, lamentablt sales.\strut & \smash{\parbox[t]{39mm}{\fontsize{6.2pt}{7.0pt}\selectfont\raggedright\color{black!88}\textbf{20 extt} \textit{abbr. pro} extrêmement. — \textbf{21} \textit{in marg. cum unco} (comptes).}} \\
      \cline{1-1}
      \leavevmode\fontsize{6.9pt}{7.9pt}\selectfont\strut [ — Essuie-mains : 90 c — ]\strut &  \\
\end{tabular}}%
\par\vspace*{\fill}
% --- Registres de bas de page (1. Critique interne [si pertinent] ; 2. Loci paralleli en latin ; 3. Commentaire historique en retrait) ---
\noindent\rule{\textwidth}{0.5pt}\par\vspace{0.65mm}
{\fontsize{6.6pt}{7.7pt}\selectfont%
\noindent\normalfont\color{black!92} \textbf{7} (\textit{Senlis}) cf. supra f\textsuperscript{o}~37~v\textsuperscript{o} (p.~\pageref{fol:75}), f\textsuperscript{o}~48~r\textsuperscript{o} (p.~\pageref{fol:96}), infra f\textsuperscript{o}~56~r\textsuperscript{o} (p.~\pageref{fol:112}), f\textsuperscript{o}~58~v\textsuperscript{o} (p.~\pageref{fol:117}), \textit{et passim}. — \textbf{10} (\textit{Toutencourt}) cf. supra f\textsuperscript{o}~49~v\textsuperscript{o} (p.~\pageref{fol:99}), f\textsuperscript{o}~52~v\textsuperscript{o} (p.~\pageref{fol:105}), infra f\textsuperscript{o}~55~v\textsuperscript{o} (p.~\pageref{fol:111}), f\textsuperscript{o}~56~r\textsuperscript{o} (p.~\pageref{fol:112}), \textit{et passim}. — \textbf{11} (\textit{Major Levy}) cf. supra f\textsuperscript{o}~45~v\textsuperscript{o} (p.~\pageref{fol:91}), f\textsuperscript{o}~46~v\textsuperscript{o} (p.~\pageref{fol:93}), infra f\textsuperscript{o}~58~r\textsuperscript{o} (p.~\pageref{fol:116}), f\textsuperscript{o}~83~r\textsuperscript{o} (p.~\pageref{fol:166}).\par\vspace{0.55mm}
\noindent\rule{38mm}{0.4pt}\par\vspace{0.5mm}
\noindent\hspace*{5mm}\parbox{\dimexpr\textwidth-5mm\relax}{\fontsize{6.6pt}{7.7pt}\selectfont\color{black!84}Le mercredi 13 décembre 1916 au matin, Hippolyte Combes assiste le \textbf{major Levy} (chef du service d'achats et réquisitions du secteur, voir f\textsuperscript{o}~45~v\textsuperscript{o}, f\textsuperscript{o}~58~r\textsuperscript{o} et f\textsuperscript{o}~83~r\textsuperscript{o}) lors d'une réunion officielle organisée à la mairie de \textbf{Toutencourt} avec les habitants évacués du village d'\textbf{Englebelmer} (dont Hippolyte Combes avait constaté l'abandon total le 7 décembre, f\textsuperscript{o}~54~r\textsuperscript{o}) pour le règlement des récoltes et fourrages laissés dans leurs exploitations.}\par}
\newpage
\phantomsection\label{fac:110}
% --- En-tête du fac-similé (Page impaire / droite : Photographie du manuscrit) ---
\noindent\rule{\textwidth}{0.5pt}\par\vspace{0.55mm}
\noindent\textbf{\fontsize{6.8pt}{7.8pt}\selectfont FAC-SIMILÉ EN VIS-À-VIS — FOLIO 55 RECTO} \hfill \textbf{\fontsize{6.8pt}{7.8pt}\selectfont DÉCEMBRE 1916}\par\vspace{0.25mm}
\noindent{\fontsize{6.8pt}{8.0pt}\selectfont\textit{Livre-Journal d'Hippolyte Combes} — 12 – 14 décembre 1916 (Acheux, Senlis-le-Sec, Toutencourt, Belle-Église, Puchevillers, Harponville (Somme))}\par\vspace{0.55mm}
\noindent\rule{\textwidth}{0.35pt}\par\vspace*{\fill}
{\centering
  \setlength{\fboxsep}{1.2mm}%
  \setlength{\fboxrule}{0.4pt}%
  \fcolorbox{gray!55}{white}{%
    \includegraphics[width=0.95\textwidth,height=0.81\textheight,keepaspectratio]{\facsimiledir/20251126_105112.jpg}%
  }\par}
\vspace*{\fill}
\newpage

% =============================================================================
% VIS-A-VIS SPREAD : FOLIO 55 VERSO (20251126_105125.jpg)
% =============================================================================
\cleardoubleevenpage
\renewcommand{\currentmonthyear}{Décembre 1916}
\renewcommand{\currentfoliolabel}{Folio 55 verso}
\phantomsection\label{fol:111}
% --- En-tête du folio (Page paire / gauche : Transcription, apparat & registres de bas de page) ---
\noindent\rule{\textwidth}{0.5pt}\par\vspace{0.55mm}
\noindent\textbf{\footnotesize FOLIO 55 VERSO} \hfill \textbf{\footnotesize DÉCEMBRE 1916}\par\vspace{0.25mm}
\noindent{\fontsize{6.9pt}{8.0pt}\selectfont Cahots sur les routes défoncées et portrait critique du jeune capitaine Williams}\par\vspace{0.2mm}
\noindent{\fontsize{6.9pt}{8.0pt}\selectfont\textit{15 – 18 décembre 1916} \hfill \textit{Acheux, Belle-Église, Arquèves, Raincheval, Toutencourt, Ferme du Quesnoy (Somme)}}\par\vspace{0.5mm}
\noindent\rule{\textwidth}{0.35pt}\par\vspace{0.8mm}
% --- Corps de la transcription diplomatique et apparat marginal en prolongement ---
{\renewcommand{\arraystretch}{0.87}%
\noindent\hspace*{3.5mm}\begin{tabular}{@{}p{87mm}@{\hspace{3.5mm}}p{39mm}@{}}
      \leavevmode\fontsize{6.9pt}{7.9pt}\selectfont\strut \underline{15 Vendredi} — \hfill{} \underline{Décbre 1916}\strut &  \\
      \leavevmode\fontsize{6.9pt}{7.9pt}\selectfont\strut Neigé 1 peu durant nuit. Mais pluie fine\strut &  \\
      \leavevmode\fontsize{6.9pt}{7.9pt}\selectfont\strut tt fondu matinée et chaussées, recouvertes [addition interlinéaire illis.]\strut & \smash{\parbox[t]{39mm}{\fontsize{6.2pt}{7.0pt}\selectfont\raggedright\color{black!88}\textbf{3} \textit{lect. incert. in ms.} \textbf{3 tt} \textit{abbr. pro} tout.}} \\
      \leavevmode\fontsize{6.9pt}{7.9pt}\selectfont\strut couche de boue liquide.\strut &  \\
      \leavevmode\llap{\fontsize{6.4pt}{7.3pt}\selectfont\color{darkgray}5\hspace{1.8mm}}\fontsize{6.9pt}{7.9pt}\selectfont\strut Pas quitté village de toute la journée.\strut &  \\
      \leavevmode\fontsize{6.9pt}{7.9pt}\selectfont\strut Resté lgts à l'office, devant brasero, ou ds ma\strut & \smash{\parbox[t]{39mm}{\fontsize{6.2pt}{7.0pt}\selectfont\raggedright\color{black!88}\textbf{6 ds} \textit{abbr. pro} dans. \textbf{6 lgts} \textit{abbr. pro} longtemps. — \textbf{7 bcp} \textit{abbr. pro} beaucoup. — \textbf{8 L.} \textit{scil.} Lundi. \textbf{8 Auj'd'hui} \textit{abbr. pro} Aujourd'hui.}} \\
      \leavevmode\fontsize{6.9pt}{7.9pt}\selectfont\strut chambre et bcp lu, pour me distraire.\strut &  \\
      \cline{1-1}
      \leavevmode\fontsize{6.9pt}{7.9pt}\selectfont\strut \underline{18 L.} — Auj'd'hui, comme hier, courses\strut &  \\
      \leavevmode\fontsize{6.9pt}{7.9pt}\selectfont\strut en auto à Belle-Église, Arquèves,\strut &  \\
      \leavevmode\llap{\fontsize{6.4pt}{7.3pt}\selectfont\color{darkgray}10\hspace{1.8mm}}\fontsize{6.9pt}{7.9pt}\selectfont\strut Raincheval, Toutencourt et ferme du\strut &  \\
      \leavevmode\fontsize{6.9pt}{7.9pt}\selectfont\strut Quesnoy.\strut &  \\
      \leavevmode\fontsize{6.9pt}{7.9pt}\selectfont\strut Sur les routes, très mauvais état, voiture\strut &  \\
      \leavevmode\fontsize{6.9pt}{7.9pt}\selectfont\strut ressaute : on est cahoté, secoué, jeté\strut &  \\
      \leavevmode\fontsize{6.9pt}{7.9pt}\selectfont\strut d'1 côté et d'autre. Fatigue extrême.\strut &  \\
      \leavevmode\llap{\fontsize{6.4pt}{7.3pt}\selectfont\color{darkgray}15\hspace{1.8mm}}\fontsize{6.9pt}{7.9pt}\selectfont\strut Tjrs même temps, ciel brumeux,\strut & \smash{\parbox[t]{39mm}{\fontsize{6.2pt}{7.0pt}\selectfont\raggedright\color{black!88}\textbf{15 Tjrs} \textit{abbr. pro} Toujours.}} \\
      \leavevmode\fontsize{6.9pt}{7.9pt}\selectfont\strut chaussées luisantes de boue liquide.\strut &  \\
      \leavevmode\fontsize{6.9pt}{7.9pt}\selectfont\strut A c. endroits, on y patauge à pleine\strut &  \\
      \leavevmode\fontsize{6.9pt}{7.9pt}\selectfont\strut cheville.\strut &  \\
      \leavevmode\fontsize{6.9pt}{7.9pt}\selectfont\strut Jours se suivent, monotones \&\strut &  \\
      \leavevmode\llap{\fontsize{6.4pt}{7.3pt}\selectfont\color{darkgray}20\hspace{1.8mm}}\fontsize{6.9pt}{7.9pt}\selectfont\strut ternes. Occupations uniformes. Société de mon\strut &  \\
      \leavevmode\fontsize{6.9pt}{7.9pt}\selectfont\strut off. — 1 cap. de 20 ans, présomptueux \& vain —\strut & \smash{\parbox[t]{39mm}{\fontsize{6.2pt}{7.0pt}\selectfont\raggedright\color{black!88}\textbf{21 off.} \textit{abbr. pro} officier(s).}} \\
      \leavevmode\fontsize{6.9pt}{7.9pt}\selectfont\strut peu agréable.\strut &  \\
      \cline{1-1}
      \leavevmode\fontsize{6.9pt}{7.9pt}\selectfont\strut [ — Cartes ples : 25 c ]\strut & \smash{\parbox[t]{39mm}{\fontsize{6.2pt}{7.0pt}\selectfont\raggedright\color{black!88}\textbf{23} \textit{in marg. cum unco} (comptes). \textbf{23 Cartes ples} \textit{abbr. pro} Cartes postales.}} \\
\end{tabular}}%
\par\vspace*{\fill}
% --- Registres de bas de page (1. Critique interne [si pertinent] ; 2. Loci paralleli en latin ; 3. Commentaire historique en retrait) ---
\noindent\rule{\textwidth}{0.5pt}\par\vspace{0.65mm}
{\fontsize{6.6pt}{7.7pt}\selectfont%
\noindent\normalfont\color{black!92} \textbf{9} (\textit{Arquèves}) cf. supra f\textsuperscript{o}~47~v\textsuperscript{o} (p.~\pageref{fol:95}), f\textsuperscript{o}~54~r\textsuperscript{o} (p.~\pageref{fol:108}), infra f\textsuperscript{o}~57~v\textsuperscript{o} (p.~\pageref{fol:115}), f\textsuperscript{o}~60~v\textsuperscript{o} (p.~\pageref{fol:121}), \textit{et passim}. — \textbf{9} (\textit{Belle-Église}) cf. supra f\textsuperscript{o}~47~v\textsuperscript{o} (p.~\pageref{fol:95}), f\textsuperscript{o}~55~r\textsuperscript{o} (p.~\pageref{fol:110}), infra f\textsuperscript{o}~56~v\textsuperscript{o} (p.~\pageref{fol:113}), f\textsuperscript{o}~58~v\textsuperscript{o} (p.~\pageref{fol:117}), \textit{et passim}. — \textbf{10} (\textit{Raincheval}) cf. supra f\textsuperscript{o}~20~v\textsuperscript{o} (p.~\pageref{fol:41}), f\textsuperscript{o}~54~r\textsuperscript{o} (p.~\pageref{fol:108}), infra f\textsuperscript{o}~56~v\textsuperscript{o} (p.~\pageref{fol:113}), f\textsuperscript{o}~57~v\textsuperscript{o} (p.~\pageref{fol:115}), \textit{et passim}.\par\vspace{0.55mm}
\noindent\rule{38mm}{0.4pt}\par\vspace{0.5mm}
\noindent\hspace*{5mm}\parbox{\dimexpr\textwidth-5mm\relax}{\fontsize{6.6pt}{7.7pt}\selectfont\color{black!84}Les tournées des 17 et 18 décembre 1916 sur les routes défoncées relient les centres de battage d'\textbf{Arquèves} (\textbf{Belle-Église}), \textbf{Raincheval}, \textbf{Toutencourt} et la \textbf{ferme du Quesnoy} (écart agricole près de Raincheval / Puchevillers, retrouvé aux f\textsuperscript{o}~56~r\textsuperscript{o}, f\textsuperscript{o}~56~v\textsuperscript{o}, f\textsuperscript{o}~57~v\textsuperscript{o}, f\textsuperscript{o}~58~r\textsuperscript{o} et f\textsuperscript{o}~59~r\textsuperscript{o}), tandis que s'accuse la mésentente morale entre l'universitaire de quarante ans et le jeune \textbf{capitaine Williams} (« un capitaine de 20 ans, présomptueux et vain »).}\par}
\newpage
\phantomsection\label{fac:111}
% --- En-tête du fac-similé (Page impaire / droite : Photographie du manuscrit) ---
\noindent\rule{\textwidth}{0.5pt}\par\vspace{0.55mm}
\noindent\textbf{\fontsize{6.8pt}{7.8pt}\selectfont FAC-SIMILÉ EN VIS-À-VIS — FOLIO 55 VERSO} \hfill \textbf{\fontsize{6.8pt}{7.8pt}\selectfont DÉCEMBRE 1916}\par\vspace{0.25mm}
\noindent{\fontsize{6.8pt}{8.0pt}\selectfont\textit{Livre-Journal d'Hippolyte Combes} — 15 – 18 décembre 1916 (Acheux, Belle-Église, Arquèves, Raincheval, Toutencourt, Ferme du Quesnoy (Somme))}\par\vspace{0.55mm}
\noindent\rule{\textwidth}{0.35pt}\par\vspace*{\fill}
{\centering
  \setlength{\fboxsep}{1.2mm}%
  \setlength{\fboxrule}{0.4pt}%
  \fcolorbox{gray!55}{white}{%
    \includegraphics[width=0.95\textwidth,height=0.81\textheight,keepaspectratio]{\facsimiledir/20251126_105125.jpg}%
  }\par}
\vspace*{\fill}
\newpage

% =============================================================================
% VIS-A-VIS SPREAD : FOLIO 56 RECTO (20251126_105142.jpg)
% =============================================================================
\cleardoubleevenpage
\renewcommand{\currentmonthyear}{Décembre 1916}
\renewcommand{\currentfoliolabel}{Folio 56 recto}
\phantomsection\label{fol:112}
% --- En-tête du folio (Page paire / gauche : Transcription, apparat & registres de bas de page) ---
\noindent\rule{\textwidth}{0.5pt}\par\vspace{0.55mm}
\noindent\textbf{\footnotesize FOLIO 56 RECTO} \hfill \textbf{\footnotesize DÉCEMBRE 1916}\par\vspace{0.25mm}
\noindent{\fontsize{6.9pt}{8.0pt}\selectfont Crise épistolaire avec Marguerite (« velléités de rupture ») et l'incident du « Bloody France ! »}\par\vspace{0.2mm}
\noindent{\fontsize{6.9pt}{8.0pt}\selectfont\textit{22 – 24 décembre 1916} \hfill \textit{Acheux, Senlis-le-Sec, Rubempré, Le Quesnoy, Léalvillers, Saint-Léger, Orville, Authie, Toutencourt, Puchevillers, Louvencourt (Somme)}}\par\vspace{0.5mm}
\noindent\rule{\textwidth}{0.35pt}\par\vspace{0.8mm}
% --- Corps de la transcription diplomatique et apparat marginal en prolongement ---
{\renewcommand{\arraystretch}{0.87}%
\noindent\hspace*{3.5mm}\begin{tabular}{@{}p{87mm}@{\hspace{3.5mm}}p{39mm}@{}}
      \leavevmode\fontsize{6.9pt}{7.9pt}\selectfont\strut \underline{22 Ven.} \hfill{} \underline{Déc. 1916}\strut & \smash{\parbox[t]{39mm}{\fontsize{6.2pt}{7.0pt}\selectfont\raggedright\color{black!88}\textbf{1 Ven.} \textit{scil.} Vendredi.}} \\
      \leavevmode\fontsize{6.9pt}{7.9pt}\selectfont\strut Dégel rapide. Routes, très\strut &  \\
      \leavevmode\fontsize{6.9pt}{7.9pt}\selectfont\strut mauvaises. Pluie fine. Matinée, à Senlis.\strut &  \\
      \leavevmode\fontsize{6.9pt}{7.9pt}\selectfont\strut Après-midi, à Rubempré \& Quesnoy —\strut &  \\
      \leavevmode\llap{\fontsize{6.4pt}{7.3pt}\selectfont\color{darkgray}5\hspace{1.8mm}}\fontsize{6.9pt}{7.9pt}\selectfont\strut Fait réparer lorgnon par horloger de\strut &  \\
      \leavevmode\fontsize{6.9pt}{7.9pt}\selectfont\strut Léalvillers.\strut &  \\
      \leavevmode\fontsize{6.9pt}{7.9pt}\selectfont\strut Soir, Cap. allé dîner à Beauquesne. Rentré\strut & \smash{\parbox[t]{39mm}{\fontsize{6.2pt}{7.0pt}\selectfont\raggedright\color{black!88}\textbf{7 Cap.} \textit{abbr. pro} Capitaine.}} \\
      \leavevmode\fontsize{6.9pt}{7.9pt}\selectfont\strut + tôt billet.\strut &  \\
      \leavevmode\fontsize{6.9pt}{7.9pt}\selectfont\strut [ — Lorgnon : 20 c. ]\strut & \smash{\parbox[t]{39mm}{\fontsize{6.2pt}{7.0pt}\selectfont\raggedright\color{black!88}\textbf{9} \textit{in marg. cum unco} (comptes). \textbf{9 c.} \textit{abbr. pro} centimes. — \textbf{10 bcp} \textit{abbr. pro} beaucoup. \textbf{10 Mte} \textit{scil.} Marguerite (Calmels). — \textbf{12 D.} \textit{scil.} Dimanche. \textbf{12 Tjrs} \textit{abbr. pro} Toujours. — \textbf{14 tte} \textit{abbr. pro} toute. — \textbf{15 St} \textit{abbr. pro} Saint. — \textbf{17 ch.} \textit{abbr. pro} chez. \textbf{17 Mme} \textit{abbr. pro} Madame. — \textbf{18 Cap. W.} \textit{scil.} Capitaine Williams. \textbf{18 Cap.} \textit{abbr. pro} Capitaine. — \textbf{22 Mte} \textit{scil.} Marguerite (Calmels). \textbf{22 ct} \textit{abbr. pro} courant. — \textbf{23 gdes} \textit{abbr. pro} grandes. — \textbf{24} \textit{in marg. cum unco} (comptes). \textbf{24 s/} \textit{abbr. pro} sur.}} \\
      \leavevmode\llap{\fontsize{6.4pt}{7.3pt}\selectfont\color{darkgray}10\hspace{1.8mm}}\fontsize{6.9pt}{7.9pt}\selectfont\strut Lettre de Mte, m'a bcp déplu. Très affecté.\strut &  \\
      \leavevmode\fontsize{6.9pt}{7.9pt}\selectfont\strut Regrets pour le passé \& appréhensions pour l'avenir.\strut &  \\
      \cline{1-1}
      \leavevmode\fontsize{6.9pt}{7.9pt}\selectfont\strut \underline{24 D.} — Tjrs même temps. Ciel couvert, averses\strut &  \\
      \leavevmode\fontsize{6.9pt}{7.9pt}\selectfont\strut intermittentes et aussi très violentes rafales de vent.\strut &  \\
      \leavevmode\fontsize{6.9pt}{7.9pt}\selectfont\strut Dehors tte la journée, en course. Matinée\strut &  \\
      \leavevmode\llap{\fontsize{6.4pt}{7.3pt}\selectfont\color{darkgray}15\hspace{1.8mm}}\fontsize{6.9pt}{7.9pt}\selectfont\strut à St-Léger \& Orville \& Authie. Après-midi\strut &  \\
      \leavevmode\fontsize{6.9pt}{7.9pt}\selectfont\strut à Louvencourt, Puchevillers etc. — Toutencourt etc.\strut &  \\
      \leavevmode\fontsize{6.9pt}{7.9pt}\selectfont\strut A Louvencourt, ch. Mme Cuisinier, gaffe\strut &  \\
      \leavevmode\fontsize{6.9pt}{7.9pt}\selectfont\strut du Cap. W. : « Bloody France etc.! »\strut &  \\
      \leavevmode\fontsize{6.9pt}{7.9pt}\selectfont\strut Très affecté par ce propos \& hésité à le reporter\strut &  \\
      \leavevmode\llap{\fontsize{6.4pt}{7.3pt}\selectfont\color{darkgray}20\hspace{1.8mm}}\fontsize{6.9pt}{7.9pt}\selectfont\strut à French Mission! — —\strut &  \\
      \leavevmode\fontsize{6.9pt}{7.9pt}\selectfont\strut Depuis 2 jours, déprimé et très triste, à\strut &  \\
      \leavevmode\fontsize{6.9pt}{7.9pt}\selectfont\strut la suite lettre de Mte du 17 ct. — Velléités de\strut &  \\
      \leavevmode\fontsize{6.9pt}{7.9pt}\selectfont\strut rupture — gdes appréhensions pour l'avenir.\strut &  \\
      \cline{1-1}
      \leavevmode\fontsize{6.9pt}{7.9pt}\selectfont\strut [ — Authie : 2 œufs s/ plat : 70 c. — Cartes ples : 30 c ]\strut &  \\
\end{tabular}}%
\par\vspace*{\fill}
% --- Registres de bas de page (1. Critique interne [si pertinent] ; 2. Loci paralleli en latin ; 3. Commentaire historique en retrait) ---
\noindent\rule{\textwidth}{0.5pt}\par\vspace{0.65mm}
{\fontsize{6.6pt}{7.7pt}\selectfont%
\noindent\normalfont \textbf{Inversion photographique :} Le f\textsuperscript{o}~56~r\textsuperscript{o} (journées des 22, 23 et vigile de Noël 24 décembre 1916) et le f\textsuperscript{o}~56~v\textsuperscript{o} (journées des 19, 20 et 21 décembre 1916) ont été photographiés dans l'ordre inverse de la chronologie.\par\vspace{0.5mm}
\noindent\rule{22mm}{0.3pt}\par\vspace{0.45mm}
\noindent\normalfont\color{black!92} \textbf{1} (\textit{22 Ven.}) de ordine chronologico cf. infra f\textsuperscript{o}~56~v\textsuperscript{o} (p.~\pageref{fol:113}). — \textbf{3} (\textit{Senlis}) cf. supra f\textsuperscript{o}~37~v\textsuperscript{o} (p.~\pageref{fol:75}), f\textsuperscript{o}~55~r\textsuperscript{o} (p.~\pageref{fol:110}), infra f\textsuperscript{o}~58~v\textsuperscript{o} (p.~\pageref{fol:117}), f\textsuperscript{o}~62~r\textsuperscript{o} (p.~\pageref{fol:124}), \textit{et passim}. — \textbf{4} (\textit{Rubempré}) cf. supra f\textsuperscript{o}~47~v\textsuperscript{o} (p.~\pageref{fol:95}), f\textsuperscript{o}~52~r\textsuperscript{o} (p.~\pageref{fol:104}), infra f\textsuperscript{o}~61~r\textsuperscript{o} (p.~\pageref{fol:122}), f\textsuperscript{o}~62~r\textsuperscript{o} (p.~\pageref{fol:124}), \textit{et passim}.\par\vspace{0.55mm}
\noindent\rule{38mm}{0.4pt}\par\vspace{0.5mm}
\noindent\hspace*{5mm}\parbox{\dimexpr\textwidth-5mm\relax}{\fontsize{6.6pt}{7.7pt}\selectfont\color{black!84}Réveillon et nuit de Noël 1916 au château d'\textbf{Acheux-en-Amiénois}, où est installée l'école de mortiers de tranchée de la V\textsuperscript{e} Armée britannique (\textit{Fifth Army School of Mortars}).}\par}
\newpage
\phantomsection\label{fac:112}
% --- En-tête du fac-similé (Page impaire / droite : Photographie du manuscrit) ---
\noindent\rule{\textwidth}{0.5pt}\par\vspace{0.55mm}
\noindent\textbf{\fontsize{6.8pt}{7.8pt}\selectfont FAC-SIMILÉ EN VIS-À-VIS — FOLIO 56 RECTO} \hfill \textbf{\fontsize{6.8pt}{7.8pt}\selectfont DÉCEMBRE 1916}\par\vspace{0.25mm}
\noindent{\fontsize{6.8pt}{8.0pt}\selectfont\textit{Livre-Journal d'Hippolyte Combes} — 22 – 24 décembre 1916 (Acheux, Senlis-le-Sec, Rubempré, Le Quesnoy, Léalvillers, Saint-Léger, Orville, Authie, Toutencourt, Puchevillers, Louvencourt (Somme))}\par\vspace{0.55mm}
\noindent\rule{\textwidth}{0.35pt}\par\vspace*{\fill}
{\centering
  \setlength{\fboxsep}{1.2mm}%
  \setlength{\fboxrule}{0.4pt}%
  \fcolorbox{gray!55}{white}{%
    \includegraphics[width=0.95\textwidth,height=0.81\textheight,keepaspectratio]{\facsimiledir/20251126_105142.jpg}%
  }\par}
\vspace*{\fill}
\newpage

% =============================================================================
% VIS-A-VIS SPREAD : FOLIO 56 VERSO (20251126_105155.jpg)
% =============================================================================
\cleardoubleevenpage
\renewcommand{\currentmonthyear}{Décembre 1916}
\renewcommand{\currentfoliolabel}{Folio 56 verso}
\phantomsection\label{fol:113}
% --- En-tête du folio (Page paire / gauche : Transcription, apparat & registres de bas de page) ---
\noindent\rule{\textwidth}{0.5pt}\par\vspace{0.55mm}
\noindent\textbf{\footnotesize FOLIO 56 VERSO} \hfill \textbf{\footnotesize DÉCEMBRE 1916}\par\vspace{0.25mm}
\noindent{\fontsize{6.9pt}{8.0pt}\selectfont Attente devant l'Army Purchase Board à Raincheval et grandes gelées d'avant Noël}\par\vspace{0.2mm}
\noindent{\fontsize{6.9pt}{8.0pt}\selectfont\textit{19 – 20 décembre 1916} \hfill \textit{Acheux, Belle-Église, Raincheval, Le Quesnoy, Beauquesne (Somme)}}\par\vspace{0.5mm}
\noindent\rule{\textwidth}{0.35pt}\par\vspace{0.8mm}
% --- Corps de la transcription diplomatique et apparat marginal en prolongement ---
{\renewcommand{\arraystretch}{0.87}%
\noindent\hspace*{3.5mm}\begin{tabular}{@{}p{87mm}@{\hspace{3.5mm}}p{39mm}@{}}
      \leavevmode\fontsize{6.9pt}{7.9pt}\selectfont\strut \underline{19 Ma.} — \hfill{} \underline{Décbre 1916 —}\strut & \smash{\parbox[t]{39mm}{\fontsize{6.2pt}{7.0pt}\selectfont\raggedright\color{black!88}\textbf{1 Ma.} \textit{scil.} Mardi.}} \\
      \leavevmode\fontsize{6.9pt}{7.9pt}\selectfont\strut Matinée, course à Belle-Église,\strut &  \\
      \leavevmode\fontsize{6.9pt}{7.9pt}\selectfont\strut et Raincheval, où dû attendre lgts devant\strut & \smash{\parbox[t]{39mm}{\fontsize{6.2pt}{7.0pt}\selectfont\raggedright\color{black!88}\textbf{3 lgts} \textit{abbr. pro} longtemps. — \textbf{4 ds} \textit{abbr. pro} dans.}} \\
      \leavevmode\fontsize{6.9pt}{7.9pt}\selectfont\strut Army P. Board. Village tr. calme \& désert. Froid ds auto.\strut &  \\
      \leavevmode\llap{\fontsize{6.4pt}{7.3pt}\selectfont\color{darkgray}5\hspace{1.8mm}}\fontsize{6.9pt}{7.9pt}\selectfont\strut Après-midi, resté à Acheux, office\strut &  \\
      \leavevmode\fontsize{6.9pt}{7.9pt}\selectfont\strut \& mon billet. Neige tombée à gros flocons.\strut &  \\
      \leavevmode\fontsize{6.9pt}{7.9pt}\selectfont\strut Comme d'habitude, passé soirée au\strut &  \\
      \leavevmode\fontsize{6.9pt}{7.9pt}\selectfont\strut rez-de-ch., avec mes propriétaires.\strut & \smash{\parbox[t]{39mm}{\fontsize{6.2pt}{7.0pt}\selectfont\raggedright\color{black!88}\textbf{8 rez-de-ch.} \textit{abbr. pro} rez-de-chaussée. — \textbf{9} \textit{in marg. cum unco} (comptes). — \textbf{10 fr.} \textit{abbr. pro} franc(s). — \textbf{11 c.} \textit{abbr. pro} centimes. \textbf{11 Christmas} \textit{angl. in ms.} — \textbf{12 Merc.} \textit{scil.} Mercredi. \textbf{12 fortt} \textit{abbr. pro} fortement.}} \\
      \leavevmode\fontsize{6.9pt}{7.9pt}\selectfont\strut [ — Bureau mission, XIIIe c., touché\strut &  \\
      \leavevmode\llap{\fontsize{6.4pt}{7.3pt}\selectfont\color{darkgray}10\hspace{1.8mm}}\fontsize{6.9pt}{7.9pt}\selectfont\strut solde etc. : 84 fr.\strut &  \\
      \leavevmode\fontsize{6.9pt}{7.9pt}\selectfont\strut Cartes Christmas : 70 c. ]\strut &  \\
      \cline{1-1}
      \leavevmode\fontsize{6.9pt}{7.9pt}\selectfont\strut \underline{20 Merc.} Nuit, gelé très fortt. Sol durci.\strut &  \\
      \leavevmode\fontsize{6.9pt}{7.9pt}\selectfont\strut Froid sec et très pénétrant.\strut &  \\
      \leavevmode\fontsize{6.9pt}{7.9pt}\selectfont\strut Après-midi, en auto. Soleil pâle. Air\strut &  \\
      \leavevmode\llap{\fontsize{6.4pt}{7.3pt}\selectfont\color{darkgray}15\hspace{1.8mm}}\fontsize{6.9pt}{7.9pt}\selectfont\strut glacial. Assis in front, moins jeté et\strut &  \\
      \leavevmode\fontsize{6.9pt}{7.9pt}\selectfont\strut moins secoué que à l'arrière —\strut &  \\
      \leavevmode\fontsize{6.9pt}{7.9pt}\selectfont\strut Venu à Raincheval, Quesnoy \& Beauquesne.\strut &  \\
      \leavevmode\fontsize{6.9pt}{7.9pt}\selectfont\strut Rentré nuit tombée : au ciel, étoiles très brillantes.\strut &  \\
      \leavevmode\fontsize{6.9pt}{7.9pt}\selectfont\strut Gelé à pierre fendre.\strut &  \\
      \leavevmode\llap{\fontsize{6.4pt}{7.3pt}\selectfont\color{darkgray}20\hspace{1.8mm}}\fontsize{6.9pt}{7.9pt}\selectfont\strut Mais préfère mille fois ce gd froid au temps\strut & \smash{\parbox[t]{39mm}{\fontsize{6.2pt}{7.0pt}\selectfont\raggedright\color{black!88}\textbf{20 gd} \textit{abbr. pro} grand.}} \\
      \leavevmode\fontsize{6.9pt}{7.9pt}\selectfont\strut brouillard et pluie des semaines précédentes.\strut &  \\
\end{tabular}}%
\par\vspace*{\fill}
% --- Registres de bas de page (1. Critique interne [si pertinent] ; 2. Loci paralleli en latin ; 3. Commentaire historique en retrait) ---
\noindent\rule{\textwidth}{0.5pt}\par\vspace{0.65mm}
{\fontsize{6.6pt}{7.7pt}\selectfont%
\noindent\normalfont \textbf{Chronologie interne :} Ce feuillet, qui réunit les entrées des mardi 19, mercredi 20 et jeudi 21 décembre 1916 séparées par des traits horizontaux, précède chronologiquement le f\textsuperscript{o}~56~r\textsuperscript{o}.\par\vspace{0.5mm}
\noindent\rule{22mm}{0.3pt}\par\vspace{0.45mm}
\noindent\normalfont\color{black!92} \textbf{1} (\textit{19 Ma.}) de ordine chronologico cf. supra f\textsuperscript{o}~56~r\textsuperscript{o} (p.~\pageref{fol:112}). — \textbf{2} (\textit{Belle-Église}) cf. supra f\textsuperscript{o}~47~v\textsuperscript{o} (p.~\pageref{fol:95}), f\textsuperscript{o}~55~v\textsuperscript{o} (p.~\pageref{fol:111}), infra f\textsuperscript{o}~58~v\textsuperscript{o} (p.~\pageref{fol:117}), f\textsuperscript{o}~59~v\textsuperscript{o} (p.~\pageref{fol:119}), \textit{et passim}. — \textbf{3} (\textit{Raincheval}) cf. supra f\textsuperscript{o}~20~v\textsuperscript{o} (p.~\pageref{fol:41}), f\textsuperscript{o}~55~v\textsuperscript{o} (p.~\pageref{fol:111}), infra f\textsuperscript{o}~57~v\textsuperscript{o} (p.~\pageref{fol:115}), f\textsuperscript{o}~58~r\textsuperscript{o} (p.~\pageref{fol:116}), \textit{et passim}.\par\vspace{0.55mm}
\noindent\rule{38mm}{0.4pt}\par\vspace{0.5mm}
\noindent\hspace*{5mm}\parbox{\dimexpr\textwidth-5mm\relax}{\fontsize{6.6pt}{7.7pt}\selectfont\color{black!84}À \textbf{Raincheval}, Hippolyte Combes attend devant le bureau de la commission d'achats de l'armée (\textbf{\textit{Army Purchase Board}}) et perçoit sa solde bimensuelle (84~francs) au bureau de la Mission française du \textbf{XIII\textsuperscript{e} Corps britannique}.}\par}
\newpage
\phantomsection\label{fac:113}
% --- En-tête du fac-similé (Page impaire / droite : Photographie du manuscrit) ---
\noindent\rule{\textwidth}{0.5pt}\par\vspace{0.55mm}
\noindent\textbf{\fontsize{6.8pt}{7.8pt}\selectfont FAC-SIMILÉ EN VIS-À-VIS — FOLIO 56 VERSO} \hfill \textbf{\fontsize{6.8pt}{7.8pt}\selectfont DÉCEMBRE 1916}\par\vspace{0.25mm}
\noindent{\fontsize{6.8pt}{8.0pt}\selectfont\textit{Livre-Journal d'Hippolyte Combes} — 19 – 20 décembre 1916 (Acheux, Belle-Église, Raincheval, Le Quesnoy, Beauquesne (Somme))}\par\vspace{0.55mm}
\noindent\rule{\textwidth}{0.35pt}\par\vspace*{\fill}
{\centering
  \setlength{\fboxsep}{1.2mm}%
  \setlength{\fboxrule}{0.4pt}%
  \fcolorbox{gray!55}{white}{%
    \includegraphics[width=0.95\textwidth,height=0.81\textheight,keepaspectratio]{\facsimiledir/20251126_105155.jpg}%
  }\par}
\vspace*{\fill}
\newpage

% =============================================================================
% VIS-A-VIS SPREAD : FOLIO 57 RECTO (20251126_105239.jpg)
% =============================================================================
\cleardoubleevenpage
\renewcommand{\currentmonthyear}{Décembre 1916}
\renewcommand{\currentfoliolabel}{Folio 57 recto}
\phantomsection\label{fol:114}
% --- En-tête du folio (Page paire / gauche : Transcription, apparat & registres de bas de page) ---
\noindent\rule{\textwidth}{0.5pt}\par\vspace{0.55mm}
\noindent\textbf{\footnotesize FOLIO 57 RECTO} \hfill \textbf{\footnotesize DÉCEMBRE 1916}\par\vspace{0.25mm}
\noindent{\fontsize{6.9pt}{8.0pt}\selectfont Noël 1916 à Acheux : insomnie, perplexité morale et concert au Y.M.C.A.}\par\vspace{0.2mm}
\noindent{\fontsize{6.9pt}{8.0pt}\selectfont\textit{25 décembre 1916} \hfill \textit{Acheux-en-Amiénois (Somme)}}\par\vspace{0.5mm}
\noindent\rule{\textwidth}{0.35pt}\par\vspace{0.8mm}
% --- Corps de la transcription diplomatique et apparat marginal en prolongement ---
{\renewcommand{\arraystretch}{0.87}%
\noindent\hspace*{3.5mm}\begin{tabular}{@{}p{87mm}@{\hspace{3.5mm}}p{39mm}@{}}
      \leavevmode\fontsize{6.9pt}{7.9pt}\selectfont\strut \underline{25 Lundi.} \hfill{} \underline{Décbre 1916}\strut &  \\
      \leavevmode\fontsize{6.9pt}{7.9pt}\selectfont\strut Nuit mauvaise. Insomnie.\strut &  \\
      \leavevmode\fontsize{6.9pt}{7.9pt}\selectfont\strut Tjrs sous l'impression propos fâcheux de\strut & \smash{\parbox[t]{39mm}{\fontsize{6.2pt}{7.0pt}\selectfont\raggedright\color{black!88}\textbf{3 Tjrs} \textit{abbr. pro} Toujours.}} \\
      \leavevmode\fontsize{6.9pt}{7.9pt}\selectfont\strut mon cap. à Louvencourt, hier après-midi.\strut &  \\
      \leavevmode\llap{\fontsize{6.4pt}{7.3pt}\selectfont\color{darkgray}5\hspace{1.8mm}}\fontsize{6.9pt}{7.9pt}\selectfont\strut Dois-je faire 1 rapport à Mission frse?\strut & \smash{\parbox[t]{39mm}{\fontsize{6.2pt}{7.0pt}\selectfont\raggedright\color{black!88}\textbf{5 frse} \textit{abbr. pro} française. — \textbf{6 Margte} \textit{scil.} Marguerite (Calmels). — \textbf{7 ct} \textit{abbr. pro} courant.}} \\
      \leavevmode\fontsize{6.9pt}{7.9pt}\selectfont\strut Hanté aussi par lettre de Margte du\strut &  \\
      \leavevmode\fontsize{6.9pt}{7.9pt}\selectfont\strut 17 ct. Pires pensées traversent mon esprit. Ne\strut &  \\
      \leavevmode\fontsize{6.9pt}{7.9pt}\selectfont\strut puis oublier. Cruelle perplexité. —\strut &  \\
      \leavevmode\fontsize{6.9pt}{7.9pt}\selectfont\strut Temps affreux. Averses. Vent souffle en\strut &  \\
      \leavevmode\llap{\fontsize{6.4pt}{7.3pt}\selectfont\color{darkgray}10\hspace{1.8mm}}\fontsize{6.9pt}{7.9pt}\selectfont\strut tempête.\strut &  \\
      \leavevmode\fontsize{6.9pt}{7.9pt}\selectfont\strut Journée de repos. Chambre...\strut &  \\
      \leavevmode\fontsize{6.9pt}{7.9pt}\selectfont\strut Peu à peu, rasséréné \& pris dessus.\strut &  \\
      \leavevmode\fontsize{6.9pt}{7.9pt}\selectfont\strut Après-midi, ds la Y.M.C.A. hut, concert : chansons\strut & \smash{\parbox[t]{39mm}{\fontsize{6.2pt}{7.0pt}\selectfont\raggedright\color{black!88}\textbf{13 ds} \textit{abbr. pro} dans. \textbf{13 Y.M.C.A.} \textit{angl.} Young Men's Christian Assoc. — \textbf{15} \textit{in marg. cum unco} (comptes). \textbf{15 c.} \textit{abbr. pro} centimes.}} \\
      \leavevmode\fontsize{6.9pt}{7.9pt}\selectfont\strut comiques \& saynètes — —\strut &  \\
      \leavevmode\llap{\fontsize{6.4pt}{7.3pt}\selectfont\color{darkgray}15\hspace{1.8mm}}\fontsize{6.9pt}{7.9pt}\selectfont\strut [ — Washing : 80 c. ]\strut &  \\
      \leavevmode\fontsize{6.9pt}{7.9pt}\selectfont\strut Fin journée, retrouvé compl. calme d'esprit\strut &  \\
      \leavevmode\fontsize{6.9pt}{7.9pt}\selectfont\strut \& parfait équilibre.\strut &  \\
      \leavevmode\fontsize{6.9pt}{7.9pt}\selectfont\strut Vers 9h s., quitté Cap. Williams \& regagné\strut & \smash{\parbox[t]{39mm}{\fontsize{6.2pt}{7.0pt}\selectfont\raggedright\color{black!88}\textbf{18 Cap.} \textit{abbr. pro} Capitaine. \textbf{18 s.} \textit{abbr. pro} soir. — \textbf{19 qq} \textit{abbr. pro} quelque(s). \textbf{19 Y.M.C.A.} \textit{angl.} Young Men's Christian Assoc.}} \\
      \leavevmode\fontsize{6.9pt}{7.9pt}\selectfont\strut chambre. Entré qq minutes Y.M.C.A. hut\strut &  \\
      \leavevmode\llap{\fontsize{6.4pt}{7.3pt}\selectfont\color{darkgray}20\hspace{1.8mm}}\fontsize{6.9pt}{7.9pt}\selectfont\strut où concert recommencé.\strut &  \\
      \leavevmode\fontsize{6.9pt}{7.9pt}\selectfont\strut Ciel merveillt étoilé.\strut &  \\
      \leavevmode\fontsize{6.9pt}{7.9pt}\selectfont\strut Temps doux, pas froid —\strut &  \\
      \leavevmode\fontsize{6.9pt}{7.9pt}\selectfont\strut Gd silence, nuit. Nul écho de canonnade.\strut & \smash{\parbox[t]{39mm}{\fontsize{6.2pt}{7.0pt}\selectfont\raggedright\color{black!88}\textbf{23 Gd} \textit{abbr. pro} Grand(e).}} \\
\end{tabular}}%
\par\vspace*{\fill}
% --- Registres de bas de page (1. Critique interne [si pertinent] ; 2. Loci paralleli en latin ; 3. Commentaire historique en retrait) ---
\noindent\rule{\textwidth}{0.5pt}\par\vspace{0.65mm}
{\fontsize{6.6pt}{7.7pt}\selectfont%
\noindent\normalfont\color{black!92} \textbf{4} (\textit{Louvencourt}) cf. supra f\textsuperscript{o}~49~v\textsuperscript{o} (p.~\pageref{fol:99}), f\textsuperscript{o}~56~r\textsuperscript{o} (p.~\pageref{fol:112}), infra f\textsuperscript{o}~58~r\textsuperscript{o} (p.~\pageref{fol:116}), f\textsuperscript{o}~59~r\textsuperscript{o} (p.~\pageref{fol:118}), \textit{et passim}. — \textbf{6} (\textit{Margte}) cf. supra f\textsuperscript{o}~1~r\textsuperscript{o} (p.~\pageref{fol:2}), f\textsuperscript{o}~56~r\textsuperscript{o} (p.~\pageref{fol:112}), infra f\textsuperscript{o}~66~v\textsuperscript{o} (p.~\pageref{fol:133}), f\textsuperscript{o}~67~r\textsuperscript{o} (p.~\pageref{fol:134}), \textit{et passim}. — \textbf{18} (\textit{Williams}) cf. supra f\textsuperscript{o}~52~v\textsuperscript{o} (p.~\pageref{fol:105}), f\textsuperscript{o}~55~v\textsuperscript{o} (p.~\pageref{fol:111}), infra f\textsuperscript{o}~58~v\textsuperscript{o} (p.~\pageref{fol:117}), f\textsuperscript{o}~73~v\textsuperscript{o} (p.~\pageref{fol:147}), \textit{et passim}.\par\vspace{0.55mm}
\noindent\rule{38mm}{0.4pt}\par\vspace{0.5mm}
\noindent\hspace*{5mm}\parbox{\dimexpr\textwidth-5mm\relax}{\fontsize{6.6pt}{7.7pt}\selectfont\color{black!84}Journée de Noël 1916 (\textbf{25 décembre 1916}) à \textbf{Acheux-en-Amiénois}: après une nuit d'insomnie partagée entre l'indignation causée par l'incident de Louvencourt avec le capitaine Williams et le trouble provoqué par la lettre de Marguerite du 17 décembre, Hippolyte Combes retrouve son équilibre moral dans sa chambre, assiste l'après-midi et le soir au spectacle de music-hall (\textit{chansons comiques et saynètes}) donné au foyer de la \textbf{\textit{Y.M.C.A.}} d'Acheux, et note le silence complet de l'artillerie dans la nuit étoilée.}\par}
\newpage
\phantomsection\label{fac:114}
% --- En-tête du fac-similé (Page impaire / droite : Photographie du manuscrit) ---
\noindent\rule{\textwidth}{0.5pt}\par\vspace{0.55mm}
\noindent\textbf{\fontsize{6.8pt}{7.8pt}\selectfont FAC-SIMILÉ EN VIS-À-VIS — FOLIO 57 RECTO} \hfill \textbf{\fontsize{6.8pt}{7.8pt}\selectfont DÉCEMBRE 1916}\par\vspace{0.25mm}
\noindent{\fontsize{6.8pt}{8.0pt}\selectfont\textit{Livre-Journal d'Hippolyte Combes} — 25 décembre 1916 (Acheux-en-Amiénois (Somme))}\par\vspace{0.55mm}
\noindent\rule{\textwidth}{0.35pt}\par\vspace*{\fill}
{\centering
  \setlength{\fboxsep}{1.2mm}%
  \setlength{\fboxrule}{0.4pt}%
  \fcolorbox{gray!55}{white}{%
    \includegraphics[width=0.95\textwidth,height=0.81\textheight,keepaspectratio]{\facsimiledir/20251126_105239.jpg}%
  }\par}
\vspace*{\fill}
\newpage

% =============================================================================
% VIS-A-VIS SPREAD : FOLIO 57 VERSO (20251126_105252.jpg)
% =============================================================================
\cleardoubleevenpage
\renewcommand{\currentmonthyear}{Décembre 1916}
\renewcommand{\currentfoliolabel}{Folio 57 verso}
\phantomsection\label{fol:115}
% --- En-tête du folio (Page paire / gauche : Transcription, apparat & registres de bas de page) ---
\noindent\rule{\textwidth}{0.5pt}\par\vspace{0.55mm}
\noindent\textbf{\footnotesize FOLIO 57 VERSO} \hfill \textbf{\footnotesize DÉCEMBRE 1916}\par\vspace{0.25mm}
\noindent{\fontsize{6.9pt}{8.0pt}\selectfont Location d'une machine à battre chez Mme Capron à Beauquesne et course rapide à Doullens}\par\vspace{0.2mm}
\noindent{\fontsize{6.9pt}{8.0pt}\selectfont\textit{26 – 27 décembre 1916} \hfill \textit{Acheux, Raincheval, Arquèves, Beauquesne, Le Quesnoy, Doullens (Somme)}}\par\vspace{0.5mm}
\noindent\rule{\textwidth}{0.35pt}\par\vspace{0.8mm}
% --- Corps de la transcription diplomatique et apparat marginal en prolongement ---
{\renewcommand{\arraystretch}{0.87}%
\noindent\hspace*{3.5mm}\begin{tabular}{@{}p{87mm}@{\hspace{3.5mm}}p{39mm}@{}}
      \leavevmode\fontsize{6.9pt}{7.9pt}\selectfont\strut \underline{26 Mardi} \hfill{} \underline{Décbre 1916.}\strut &  \\
      \leavevmode\fontsize{6.9pt}{7.9pt}\selectfont\strut Temps de pluie \& brouillard.\strut &  \\
      \leavevmode\fontsize{6.9pt}{7.9pt}\selectfont\strut Matin, à Raincheval \& Arquèves, vu\strut &  \\
      \leavevmode\fontsize{6.9pt}{7.9pt}\selectfont\strut M. Guibet...\strut &  \\
      \leavevmode\llap{\fontsize{6.4pt}{7.3pt}\selectfont\color{darkgray}5\hspace{1.8mm}}\fontsize{6.9pt}{7.9pt}\selectfont\strut Après-midi, à Beauquesne, Quesnoy etc.\strut &  \\
      \leavevmode\fontsize{6.9pt}{7.9pt}\selectfont\strut A Beauquesne, charmante Mme Capron,\strut & \smash{\parbox[t]{39mm}{\fontsize{6.2pt}{7.0pt}\selectfont\raggedright\color{black!88}\textbf{6 Mme} \textit{abbr. pro} Madame.}} \\
      \leavevmode\fontsize{6.9pt}{7.9pt}\selectfont\strut à propos location machine à battre.\strut &  \\
      \leavevmode\fontsize{6.9pt}{7.9pt}\selectfont\strut Rentrés à 6h. Fatigue extrême — Couché tôt.\strut &  \\
      \cline{1-1}
      \leavevmode\fontsize{6.9pt}{7.9pt}\selectfont\strut \underline{27 Merc.}\strut & \smash{\parbox[t]{39mm}{\fontsize{6.2pt}{7.0pt}\selectfont\raggedright\color{black!88}\textbf{9 Merc.} \textit{scil.} Mercredi.}} \\
      \leavevmode\llap{\fontsize{6.4pt}{7.3pt}\selectfont\color{darkgray}10\hspace{1.8mm}}\fontsize{6.9pt}{7.9pt}\selectfont\strut Enfin, très belle journée. Revu soleil.\strut &  \\
      \leavevmode\fontsize{6.9pt}{7.9pt}\selectfont\strut Matinée, course à Beauquesne. Puis\strut &  \\
      \leavevmode\fontsize{6.9pt}{7.9pt}\selectfont\strut continué j'à Doullens, mais resté\strut & \smash{\parbox[t]{39mm}{\fontsize{6.2pt}{7.0pt}\selectfont\raggedright\color{black!88}\textbf{12 j'à} \textit{abbr. pro} jusqu'à.}} \\
      \leavevmode\fontsize{6.9pt}{7.9pt}\selectfont\strut 20 m. à peine ds cette ville (hélas!).\strut & \smash{\parbox[t]{39mm}{\fontsize{6.2pt}{7.0pt}\selectfont\raggedright\color{black!88}\textbf{13 ds} \textit{abbr. pro} dans.}} \\
      \cline{1-1}
      \leavevmode\fontsize{6.9pt}{7.9pt}\selectfont\strut [ \underline{A Doullens} : 1 ramette papier à l. : 1 f 25 c\strut & \smash{\parbox[t]{39mm}{\fontsize{6.2pt}{7.0pt}\selectfont\raggedright\color{black!88}\textbf{14} \textit{in marg. cum unco} (comptes). — \textbf{15 gdes} \textit{abbr. pro} grandes. — \textbf{16 Cartes ples} \textit{abbr. pro} Cartes postales.}} \\
      \leavevmode\llap{\fontsize{6.4pt}{7.3pt}\selectfont\color{darkgray}15\hspace{1.8mm}}\fontsize{6.9pt}{7.9pt}\selectfont\strut 2 paquets gdes enveloppes : 45 + 45 c\strut &  \\
      \cline{1-1}
      \leavevmode\fontsize{6.9pt}{7.9pt}\selectfont\strut Cartes ples : 55 c ]\strut &  \\
\end{tabular}}%
\par\vspace*{\fill}
% --- Registres de bas de page (1. Critique interne [si pertinent] ; 2. Loci paralleli en latin ; 3. Commentaire historique en retrait) ---
\noindent\rule{\textwidth}{0.5pt}\par\vspace{0.65mm}
{\fontsize{6.6pt}{7.7pt}\selectfont%
\noindent\normalfont\color{black!92} \textbf{3} (\textit{Raincheval}) cf. supra f\textsuperscript{o}~20~v\textsuperscript{o} (p.~\pageref{fol:41}), f\textsuperscript{o}~56~v\textsuperscript{o} (p.~\pageref{fol:113}), infra f\textsuperscript{o}~58~r\textsuperscript{o} (p.~\pageref{fol:116}), f\textsuperscript{o}~58~v\textsuperscript{o} (p.~\pageref{fol:117}), \textit{et passim}. — \textbf{3} (\textit{Arquèves}) cf. supra f\textsuperscript{o}~47~v\textsuperscript{o} (p.~\pageref{fol:95}), f\textsuperscript{o}~55~v\textsuperscript{o} (p.~\pageref{fol:111}), infra f\textsuperscript{o}~60~v\textsuperscript{o} (p.~\pageref{fol:121}), f\textsuperscript{o}~61~r\textsuperscript{o} (p.~\pageref{fol:122}), \textit{et passim}. — \textbf{5} (\textit{Quesnoy}) cf. supra f\textsuperscript{o}~55~v\textsuperscript{o} (p.~\pageref{fol:111}), f\textsuperscript{o}~56~v\textsuperscript{o} (p.~\pageref{fol:113}), infra f\textsuperscript{o}~58~r\textsuperscript{o} (p.~\pageref{fol:116}), f\textsuperscript{o}~59~r\textsuperscript{o} (p.~\pageref{fol:118}), \textit{et passim}.\par\vspace{0.55mm}
\noindent\rule{38mm}{0.4pt}\par\vspace{0.5mm}
\noindent\hspace*{5mm}\parbox{\dimexpr\textwidth-5mm\relax}{\fontsize{6.6pt}{7.7pt}\selectfont\color{black!84}Reprise des tournées dès le lendemain de Noël: entretien avec \textbf{M.~Guibet} à Raincheval et négociation à \textbf{Beauquesne} avec \textbf{Mme Capron} pour la location d'une batteuse mécanique (\textit{machine à battre}) destinée au battage des pailles réquisitionnées par l'armée britannique (le contrat sera signé le 5 janvier 1917, f\textsuperscript{o}~60~r\textsuperscript{o}).}\par}
\newpage
\phantomsection\label{fac:115}
% --- En-tête du fac-similé (Page impaire / droite : Photographie du manuscrit) ---
\noindent\rule{\textwidth}{0.5pt}\par\vspace{0.55mm}
\noindent\textbf{\fontsize{6.8pt}{7.8pt}\selectfont FAC-SIMILÉ EN VIS-À-VIS — FOLIO 57 VERSO} \hfill \textbf{\fontsize{6.8pt}{7.8pt}\selectfont DÉCEMBRE 1916}\par\vspace{0.25mm}
\noindent{\fontsize{6.8pt}{8.0pt}\selectfont\textit{Livre-Journal d'Hippolyte Combes} — 26 – 27 décembre 1916 (Acheux, Raincheval, Arquèves, Beauquesne, Le Quesnoy, Doullens (Somme))}\par\vspace{0.55mm}
\noindent\rule{\textwidth}{0.35pt}\par\vspace*{\fill}
{\centering
  \setlength{\fboxsep}{1.2mm}%
  \setlength{\fboxrule}{0.4pt}%
  \fcolorbox{gray!55}{white}{%
    \includegraphics[width=0.95\textwidth,height=0.81\textheight,keepaspectratio]{\facsimiledir/20251126_105252.jpg}%
  }\par}
\vspace*{\fill}
\newpage

% =============================================================================
% VIS-A-VIS SPREAD : FOLIO 58 RECTO (20251126_105400.jpg)
% =============================================================================
\cleardoubleevenpage
\renewcommand{\currentmonthyear}{Décembre 1916}
\renewcommand{\currentfoliolabel}{Folio 58 recto}
\phantomsection\label{fol:116}
% --- En-tête du folio (Page paire / gauche : Transcription, apparat & registres de bas de page) ---
\noindent\rule{\textwidth}{0.5pt}\par\vspace{0.55mm}
\noindent\textbf{\footnotesize FOLIO 58 RECTO} \hfill \textbf{\footnotesize DÉCEMBRE 1916}\par\vspace{0.25mm}
\noindent{\fontsize{6.9pt}{8.0pt}\selectfont Concert des « 13th Corps Archie Singers » et tournée dans la voiture Sunbeam}\par\vspace{0.2mm}
\noindent{\fontsize{6.9pt}{8.0pt}\selectfont\textit{28 – 29 décembre 1916} \hfill \textit{Acheux, Léalvillers, Raincheval, Le Quesnoy, Louvencourt (Somme)}}\par\vspace{0.5mm}
\noindent\rule{\textwidth}{0.35pt}\par\vspace{0.8mm}
% --- Corps de la transcription diplomatique et apparat marginal en prolongement ---
{\renewcommand{\arraystretch}{0.83}%
\noindent\hspace*{3.5mm}\begin{tabular}{@{}p{87mm}@{\hspace{3.5mm}}p{39mm}@{}}
      \leavevmode\fontsize{6.6pt}{7.5pt}\selectfont\strut \underline{28 Jeudi} \hfill{} \underline{Décembre 1916}\strut &  \\
      \leavevmode\fontsize{6.6pt}{7.5pt}\selectfont\strut Très fortt gelé nuit.\strut & \smash{\parbox[t]{39mm}{\fontsize{6.2pt}{7.0pt}\selectfont\raggedright\color{black!88}\textbf{2 fortt} \textit{abbr. pro} fortement.}} \\
      \leavevmode\fontsize{6.6pt}{7.5pt}\selectfont\strut Toute nuit \& matinée, très violente canonnade\strut &  \\
      \leavevmode\fontsize{6.6pt}{7.5pt}\selectfont\strut du côté d'Hébuterne.\strut &  \\
      \leavevmode\llap{\fontsize{6.4pt}{7.3pt}\selectfont\color{darkgray}5\hspace{1.8mm}}\fontsize{6.6pt}{7.5pt}\selectfont\strut Au matin, dégelé rapidt et bientôt\strut & \smash{\parbox[t]{39mm}{\fontsize{6.2pt}{7.0pt}\selectfont\raggedright\color{black!88}\textbf{5 rapidt} \textit{abbr. pro} rapidement.}} \\
      \leavevmode\fontsize{6.6pt}{7.5pt}\selectfont\strut routes, redevenues affreust boueuses.\strut &  \\
      \leavevmode\fontsize{6.6pt}{7.5pt}\selectfont\strut Après breakfast, venu à pied à Léalvillers,\strut & \smash{\parbox[t]{39mm}{\fontsize{6.2pt}{7.0pt}\selectfont\raggedright\color{black!88}\textbf{7 breakfast} \textit{angl. in ms.}}} \\
      \leavevmode\fontsize{6.6pt}{7.5pt}\selectfont\strut retour à pied. Promenade hygiénique.\strut &  \\
      \leavevmode\fontsize{6.6pt}{7.5pt}\selectfont\strut Tout jour, liberté : car in workshop.\strut &  \\
      \leavevmode\llap{\fontsize{6.4pt}{7.3pt}\selectfont\color{darkgray}10\hspace{1.8mm}}\fontsize{6.6pt}{7.5pt}\selectfont\strut A 6h s., concert au Y.M.C.A. hut. 13th\strut & \smash{\parbox[t]{39mm}{\fontsize{6.2pt}{7.0pt}\selectfont\raggedright\color{black!88}\textbf{10 Y.M.C.A.} \textit{angl.} Young Men's Christian Assoc. \textbf{10 s.} \textit{abbr. pro} soir.}} \\
      \leavevmode\fontsize{6.6pt}{7.5pt}\selectfont\strut Corps Archie Singers. Chansons comiques\strut &  \\
      \leavevmode\fontsize{6.6pt}{7.5pt}\selectfont\strut \& pantomimes.\strut &  \\
      \leavevmode\fontsize{6.6pt}{7.5pt}\selectfont\strut [ — Brown Boot Polish : 20 c\strut & \smash{\parbox[t]{39mm}{\fontsize{6.2pt}{7.0pt}\selectfont\raggedright\color{black!88}\textbf{13} \textit{in marg. cum unco} (comptes).}} \\
      \leavevmode\fontsize{6.6pt}{7.5pt}\selectfont\strut — Boîte recommandée à Imbert, opticien\strut &  \\
      \leavevmode\llap{\fontsize{6.4pt}{7.3pt}\selectfont\color{darkgray}15\hspace{1.8mm}}\fontsize{6.6pt}{7.5pt}\selectfont\strut Paris : 30 c\strut &  \\
      \leavevmode\fontsize{6.6pt}{7.5pt}\selectfont\strut Concert Y.M.C.A. : 30 c ]\strut & \smash{\parbox[t]{39mm}{\fontsize{6.2pt}{7.0pt}\selectfont\raggedright\color{black!88}\textbf{16 Y.M.C.A.} \textit{angl.} Young Men's Christian Assoc. — \textbf{17 Ve.} \textit{scil.} Vendredi.}} \\
      \cline{1-1}
      \leavevmode\fontsize{6.6pt}{7.5pt}\selectfont\strut \underline{29 Ve.} Temps doux. Ciel couvert.\strut &  \\
      \leavevmode\fontsize{6.6pt}{7.5pt}\selectfont\strut Pluie fine —\strut &  \\
      \leavevmode\fontsize{6.6pt}{7.5pt}\selectfont\strut Matinée, avec Sunbeam Car, du Major\strut &  \\
      \leavevmode\llap{\fontsize{6.4pt}{7.3pt}\selectfont\color{darkgray}20\hspace{1.8mm}}\fontsize{6.6pt}{7.5pt}\selectfont\strut Levy, course à Raincheval, Quesnoy \&\strut &  \\
      \leavevmode\fontsize{6.6pt}{7.5pt}\selectfont\strut Louvencourt.\strut &  \\
      \leavevmode\fontsize{6.6pt}{7.5pt}\selectfont\strut Rentré à 3h s.\strut & \smash{\parbox[t]{39mm}{\fontsize{6.2pt}{7.0pt}\selectfont\raggedright\color{black!88}\textbf{22 s.} \textit{abbr. pro} soir.}} \\
      \leavevmode\fontsize{6.6pt}{7.5pt}\selectfont\strut Liberté entière fin après-midi.\strut &  \\
      \leavevmode\fontsize{6.6pt}{7.5pt}\selectfont\strut Tempérre printanière. On pt travailler\strut & \smash{\parbox[t]{39mm}{\fontsize{6.2pt}{7.0pt}\selectfont\raggedright\color{black!88}\textbf{24 pt} \textit{abbr. pro} peut.}} \\
      \leavevmode\llap{\fontsize{6.4pt}{7.3pt}\selectfont\color{darkgray}25\hspace{1.8mm}}\fontsize{6.6pt}{7.5pt}\selectfont\strut chambre sans feu, et fenêtres ouvertes.\strut &  \\
\end{tabular}}%
\par\vspace*{\fill}
% --- Registres de bas de page (1. Critique interne [si pertinent] ; 2. Loci paralleli en latin ; 3. Commentaire historique en retrait) ---
\noindent\rule{\textwidth}{0.5pt}\par\vspace{0.65mm}
{\fontsize{6.6pt}{7.7pt}\selectfont%
\noindent\normalfont\color{black!92} \textbf{7} (\textit{Léalvillers}) cf. supra f\textsuperscript{o}~24~r\textsuperscript{o} (p.~\pageref{fol:48}), f\textsuperscript{o}~56~r\textsuperscript{o} (p.~\pageref{fol:112}), infra f\textsuperscript{o}~88~v\textsuperscript{o} (p.~\pageref{fol:177}), \textit{et passim}. — \textbf{15} (\textit{Paris}) cf. supra in tegumento (p.~\pageref{fol:1}), f\textsuperscript{o}~49~v\textsuperscript{o} (p.~\pageref{fol:99}), infra f\textsuperscript{o}~61~r\textsuperscript{o} (p.~\pageref{fol:122}), f\textsuperscript{o}~65~r\textsuperscript{o} (p.~\pageref{fol:130}), \textit{et passim}. — \textbf{20} (\textit{Raincheval}) cf. supra f\textsuperscript{o}~20~v\textsuperscript{o} (p.~\pageref{fol:41}), f\textsuperscript{o}~57~v\textsuperscript{o} (p.~\pageref{fol:115}), infra f\textsuperscript{o}~58~v\textsuperscript{o} (p.~\pageref{fol:117}), f\textsuperscript{o}~59~r\textsuperscript{o} (p.~\pageref{fol:118}), \textit{et passim}.\par\vspace{0.55mm}
\noindent\rule{38mm}{0.4pt}\par\vspace{0.5mm}
\noindent\hspace*{5mm}\parbox{\dimexpr\textwidth-5mm\relax}{\fontsize{6.6pt}{7.7pt}\selectfont\color{black!84}Soirée de Noël prolongée à Acheux le 28 décembre 1916 au foyer de la \textit{Y.M.C.A.} par le concert de la troupe divisionnaire des \textbf{\textit{13th Corps Archie Singers}} (groupe artistique formé par les batteries de défense contre aéronefs, familièrement appelées \textit{Archies} dans l'argot militaire britannique). Le lendemain, tournée à Raincheval, Le Quesnoy et Louvencourt dans une automobile d'état-major \textit{Sunbeam}.}\par}
\newpage
\phantomsection\label{fac:116}
% --- En-tête du fac-similé (Page impaire / droite : Photographie du manuscrit) ---
\noindent\rule{\textwidth}{0.5pt}\par\vspace{0.55mm}
\noindent\textbf{\fontsize{6.8pt}{7.8pt}\selectfont FAC-SIMILÉ EN VIS-À-VIS — FOLIO 58 RECTO} \hfill \textbf{\fontsize{6.8pt}{7.8pt}\selectfont DÉCEMBRE 1916}\par\vspace{0.25mm}
\noindent{\fontsize{6.8pt}{8.0pt}\selectfont\textit{Livre-Journal d'Hippolyte Combes} — 28 – 29 décembre 1916 (Acheux, Léalvillers, Raincheval, Le Quesnoy, Louvencourt (Somme))}\par\vspace{0.55mm}
\noindent\rule{\textwidth}{0.35pt}\par\vspace*{\fill}
{\centering
  \setlength{\fboxsep}{1.2mm}%
  \setlength{\fboxrule}{0.4pt}%
  \fcolorbox{gray!55}{white}{%
    \includegraphics[width=0.95\textwidth,height=0.81\textheight,keepaspectratio]{\facsimiledir/20251126_105400.jpg}%
  }\par}
\vspace*{\fill}
\newpage

% =============================================================================
% VIS-A-VIS SPREAD : FOLIO 58 VERSO (20251126_105411.jpg)
% =============================================================================
\cleardoubleevenpage
\renewcommand{\currentmonthyear}{Décembre 1916}
\renewcommand{\currentfoliolabel}{Folio 58 verso}
\phantomsection\label{fol:117}
% --- En-tête du folio (Page paire / gauche : Transcription, apparat & registres de bas de page) ---
\noindent\rule{\textwidth}{0.5pt}\par\vspace{0.55mm}
\noindent\textbf{\footnotesize FOLIO 58 VERSO} \hfill \textbf{\footnotesize DÉCEMBRE 1916}\par\vspace{0.25mm}
\noindent{\fontsize{6.9pt}{8.0pt}\selectfont Nouvel éclat et excuses contrites du capitaine Williams à la Saint-Sylvestre}\par\vspace{0.2mm}
\noindent{\fontsize{6.9pt}{8.0pt}\selectfont\textit{30 – 31 décembre 1916} \hfill \textit{Acheux, Franvillers, Senlis-le-Sec, Belle-Église, Raincheval, Beauquesne (Somme)}}\par\vspace{0.5mm}
\noindent\rule{\textwidth}{0.35pt}\par\vspace{0.8mm}
% --- Corps de la transcription diplomatique et apparat marginal en prolongement ---
{\renewcommand{\arraystretch}{0.87}%
\noindent\hspace*{3.5mm}\begin{tabular}{@{}p{87mm}@{\hspace{3.5mm}}p{39mm}@{}}
      \leavevmode\fontsize{6.9pt}{7.9pt}\selectfont\strut \underline{30 Sam.} \hfill{} \underline{Décbre /16}\strut & \smash{\parbox[t]{39mm}{\fontsize{6.2pt}{7.0pt}\selectfont\raggedright\color{black!88}\textbf{1 Sam.} \textit{scil.} Samedi.}} \\
      \leavevmode\fontsize{6.9pt}{7.9pt}\selectfont\strut Toute la nuit \& à l'aube, pluie torrentielle.\strut &  \\
      \leavevmode\fontsize{6.9pt}{7.9pt}\selectfont\strut Température d'1 douceur exceptionnelle. Air tiède.\strut &  \\
      \leavevmode\fontsize{6.9pt}{7.9pt}\selectfont\strut Matinée, course en auto (Sunbeam) à Franvillers,\strut &  \\
      \leavevmode\llap{\fontsize{6.4pt}{7.3pt}\selectfont\color{darkgray}5\hspace{1.8mm}}\fontsize{6.9pt}{7.9pt}\selectfont\strut où revu trio de vieux, débonnaires, originaux\strut &  \\
      \leavevmode\fontsize{6.9pt}{7.9pt}\selectfont\strut extt sympathiques, réfugiés de Martinsart.\strut & \smash{\parbox[t]{39mm}{\fontsize{6.2pt}{7.0pt}\selectfont\raggedright\color{black!88}\textbf{6 extt} \textit{abbr. pro} extrêmement.}} \\
      \leavevmode\fontsize{6.9pt}{7.9pt}\selectfont\strut A Senlis, devant l'ordonnance, mot\strut &  \\
      \leavevmode\fontsize{6.9pt}{7.9pt}\selectfont\strut très regrettable du Cap. Williams. A l'arrivée à Acheux\strut & \smash{\parbox[t]{39mm}{\fontsize{6.2pt}{7.0pt}\selectfont\raggedright\color{black!88}\textbf{8 Cap.} \textit{abbr. pro} Capitaine.}} \\
      \leavevmode\fontsize{6.9pt}{7.9pt}\selectfont\strut explication très calme, mais très ferme. M'a fait des\strut &  \\
      \leavevmode\llap{\fontsize{6.4pt}{7.3pt}\selectfont\color{darkgray}10\hspace{1.8mm}}\fontsize{6.9pt}{7.9pt}\selectfont\strut excuses et m'a supplié de ne pas le quitter etc.\strut &  \\
      \leavevmode\fontsize{6.9pt}{7.9pt}\selectfont\strut Après-midi, entt libre —\strut &  \\
      \leavevmode\fontsize{6.9pt}{7.9pt}\selectfont\strut [ — Clous aux semelles de mes bottes : 1 f 60 ]\strut & \smash{\parbox[t]{39mm}{\fontsize{6.2pt}{7.0pt}\selectfont\raggedright\color{black!88}\textbf{12} \textit{in marg. cum unco} (comptes). — \textbf{13 Dim.} \textit{scil.} Dimanche.}} \\
      \cline{1-1}
      \leavevmode\fontsize{6.9pt}{7.9pt}\selectfont\strut \underline{31 Dim.} Tempérre +tôt chaude, et pluie\strut &  \\
      \leavevmode\fontsize{6.9pt}{7.9pt}\selectfont\strut fine, intermittente.\strut &  \\
      \leavevmode\llap{\fontsize{6.4pt}{7.3pt}\selectfont\color{darkgray}15\hspace{1.8mm}}\fontsize{6.9pt}{7.9pt}\selectfont\strut Matinée, ds ma chambre —\strut & \smash{\parbox[t]{39mm}{\fontsize{6.2pt}{7.0pt}\selectfont\raggedright\color{black!88}\textbf{15 ds} \textit{abbr. pro} dans.}} \\
      \leavevmode\fontsize{6.9pt}{7.9pt}\selectfont\strut Aussitôt après lunch, partis en auto : Belle-\strut & \smash{\parbox[t]{39mm}{\fontsize{6.2pt}{7.0pt}\selectfont\raggedright\color{black!88}\textbf{16 lunch} \textit{angl. in ms.}}} \\
      \leavevmode\fontsize{6.9pt}{7.9pt}\selectfont\strut Église, Raincheval \& Beauquesne.\strut &  \\
      \leavevmode\fontsize{6.9pt}{7.9pt}\selectfont\strut Légère brume estompe paysages lointains.\strut &  \\
      \leavevmode\fontsize{6.9pt}{7.9pt}\selectfont\strut Mon jeune capne, attitude gênée et\strut &  \\
      \leavevmode\llap{\fontsize{6.4pt}{7.3pt}\selectfont\color{darkgray}20\hspace{1.8mm}}\fontsize{6.9pt}{7.9pt}\selectfont\strut presque humiliée, contrite. Parle à peine (paraît\strut &  \\
      \leavevmode\fontsize{6.9pt}{7.9pt}\selectfont\strut et n'ose plus me regarder en face. (déprimé)\strut &  \\
      \leavevmode\fontsize{6.9pt}{7.9pt}\selectfont\strut Soir, Y.M.C.A. hut, assisté 1/2 séance —\strut & \smash{\parbox[t]{39mm}{\fontsize{6.2pt}{7.0pt}\selectfont\raggedright\color{black!88}\textbf{22 Y.M.C.A.} \textit{angl.} Young Men's Christian Assoc. \textbf{22 1/2} \textit{scil.} demi / moitié. — \textbf{23} \textit{in marg. cum unco} (comptes). \textbf{23 fr.} \textit{abbr. pro} franc(s). — \textbf{24 fr.} \textit{abbr. pro} franc(s).}} \\
      \cline{1-1}
      \leavevmode\fontsize{6.9pt}{7.9pt}\selectfont\strut [ — Dépenses mess : 45 f + 10 f + 1 f + 22 fr. —\strut &  \\
      \leavevmode\fontsize{6.9pt}{7.9pt}\selectfont\strut = 88 fr. — ]\strut &  \\
\end{tabular}}%
\par\vspace*{\fill}
% --- Registres de bas de page (1. Critique interne [si pertinent] ; 2. Loci paralleli en latin ; 3. Commentaire historique en retrait) ---
\noindent\rule{\textwidth}{0.5pt}\par\vspace{0.65mm}
{\fontsize{6.6pt}{7.7pt}\selectfont%
\noindent\normalfont\color{black!92} \textbf{4} (\textit{Franvillers}) cf. supra f\textsuperscript{o}~52~v\textsuperscript{o} (p.~\pageref{fol:105}). — \textbf{6} (\textit{Martinsart}) cf. supra f\textsuperscript{o}~52~r\textsuperscript{o} (p.~\pageref{fol:104}), f\textsuperscript{o}~52~v\textsuperscript{o} (p.~\pageref{fol:105}), infra f\textsuperscript{o}~75~r\textsuperscript{o} (p.~\pageref{fol:150}), f\textsuperscript{o}~79~v\textsuperscript{o} (p.~\pageref{fol:159}), \textit{et passim}. — \textbf{7} (\textit{Senlis}) cf. supra f\textsuperscript{o}~37~v\textsuperscript{o} (p.~\pageref{fol:75}), f\textsuperscript{o}~56~r\textsuperscript{o} (p.~\pageref{fol:112}), infra f\textsuperscript{o}~62~r\textsuperscript{o} (p.~\pageref{fol:124}), f\textsuperscript{o}~75~v\textsuperscript{o} (p.~\pageref{fol:151}), \textit{et passim}.\par\vspace{0.55mm}
\noindent\rule{38mm}{0.4pt}\par\vspace{0.5mm}
\noindent\hspace*{5mm}\parbox{\dimexpr\textwidth-5mm\relax}{\fontsize{6.6pt}{7.7pt}\selectfont\color{black!84}Après avoir revu le 30 décembre à \textbf{Franvillers} les propriétaires âgés évacués de \textbf{Martinsart} (rencontrés le 28 novembre au château de Martinsart, f\textsuperscript{o}~52~r\textsuperscript{o}), un nouvel éclat verbal du \textbf{capitaine Williams} à \textbf{Senlis-le-Sec} devant l'ordonnance provoque au retour à Acheux une mise au point ferme de l'interprète: redoutant une plainte officielle auprès de la Mission française, le jeune officier lui présente ses excuses et garde tout le jour de la Saint-Sylvestre (31 décembre 1916) une attitude contrite.}\par}
\newpage
\phantomsection\label{fac:117}
% --- En-tête du fac-similé (Page impaire / droite : Photographie du manuscrit) ---
\noindent\rule{\textwidth}{0.5pt}\par\vspace{0.55mm}
\noindent\textbf{\fontsize{6.8pt}{7.8pt}\selectfont FAC-SIMILÉ EN VIS-À-VIS — FOLIO 58 VERSO} \hfill \textbf{\fontsize{6.8pt}{7.8pt}\selectfont DÉCEMBRE 1916}\par\vspace{0.25mm}
\noindent{\fontsize{6.8pt}{8.0pt}\selectfont\textit{Livre-Journal d'Hippolyte Combes} — 30 – 31 décembre 1916 (Acheux, Franvillers, Senlis-le-Sec, Belle-Église, Raincheval, Beauquesne (Somme))}\par\vspace{0.55mm}
\noindent\rule{\textwidth}{0.35pt}\par\vspace*{\fill}
{\centering
  \setlength{\fboxsep}{1.2mm}%
  \setlength{\fboxrule}{0.4pt}%
  \fcolorbox{gray!55}{white}{%
    \includegraphics[width=0.95\textwidth,height=0.81\textheight,keepaspectratio]{\facsimiledir/20251126_105411.jpg}%
  }\par}
\vspace*{\fill}
\newpage

% =============================================================================
% VIS-A-VIS SPREAD : FOLIO 59 RECTO (20251126_105428.jpg)
% =============================================================================
\cleardoubleevenpage
\renewcommand{\currentmonthyear}{Janvier 1917}
\renewcommand{\currentfoliolabel}{Folio 59 recto}
\phantomsection\label{fol:118}
% --- En-tête du folio (Page paire / gauche : Transcription, apparat & registres de bas de page) ---
\noindent\rule{\textwidth}{0.5pt}\par\vspace{0.55mm}
\noindent\textbf{\footnotesize FOLIO 59 RECTO} \hfill \textbf{\footnotesize JANVIER 1917}\par\vspace{0.25mm}
\noindent{\fontsize{6.9pt}{8.0pt}\selectfont Jour de l'An 1917 à Acheux, tournées à Toutencourt, Saulty et Louvencourt}\par\vspace{0.2mm}
\noindent{\fontsize{6.9pt}{8.0pt}\selectfont\textit{1er – 2 janvier 1917} \hfill \textit{Acheux, Toutencourt, Saulty, Louvencourt, Raincheval (Somme / Pas-de-Calais)}}\par\vspace{0.5mm}
\noindent\rule{\textwidth}{0.35pt}\par\vspace{0.8mm}
% --- Corps de la transcription diplomatique et apparat marginal en prolongement ---
{\renewcommand{\arraystretch}{0.87}%
\noindent\hspace*{3.5mm}\begin{tabular}{@{}p{87mm}@{\hspace{3.5mm}}p{39mm}@{}}
      \leavevmode\fontsize{6.9pt}{7.9pt}\selectfont\strut \underline{1er Lundi — 1917} — \hfill{} \raisebox{0.5ex}{\underline{Janvier.}}\strut & \smash{\parbox[t]{39mm}{\fontsize{6.2pt}{7.0pt}\selectfont\raggedright\color{black!88}\textbf{1 1er} \textit{scil.} premier.}} \\
      \leavevmode\fontsize{6.9pt}{7.9pt}\selectfont\strut Temps pluvieux, brumeux, mais très doux.\strut &  \\
      \leavevmode\fontsize{6.9pt}{7.9pt}\selectfont\strut On ne croirait jamais être au milieu de l'hiver.\strut &  \\
      \leavevmode\fontsize{6.9pt}{7.9pt}\selectfont\strut Matinée, course en auto à Toutencourt \& Quesnoy.\strut &  \\
      \leavevmode\llap{\fontsize{6.4pt}{7.3pt}\selectfont\color{darkgray}5\hspace{1.8mm}}\fontsize{6.9pt}{7.9pt}\selectfont\strut Après-midi, venus à Saulty, par Authie,\strut &  \\
      \leavevmode\fontsize{6.9pt}{7.9pt}\selectfont\strut \& Pas. gde route d'Arras. Ondée très fine \&\strut & \smash{\parbox[t]{39mm}{\fontsize{6.2pt}{7.0pt}\selectfont\raggedright\color{black!88}\textbf{6 gde} \textit{scil.} grande.}} \\
      \leavevmode\fontsize{6.9pt}{7.9pt}\selectfont\strut intermittente.\strut &  \\
      \leavevmode\fontsize{6.9pt}{7.9pt}\selectfont\strut M. Carluvant m'a annoncé qu'il devra quitter\strut &  \\
      \leavevmode\fontsize{6.9pt}{7.9pt}\selectfont\strut maison ds qq jours : $[\![$nouveau$]\!]$ proprire tient à s'y\strut & \smash{\parbox[t]{39mm}{\fontsize{6.2pt}{7.0pt}\selectfont\raggedright\color{black!88}\textbf{9} $[\![$nouveau$]\!]\ $ \textit{del. in ms.} \textbf{9 ds} \textit{abbr. pro} dans.}} \\
      \leavevmode\llap{\fontsize{6.4pt}{7.3pt}\selectfont\color{darkgray}10\hspace{1.8mm}}\fontsize{6.9pt}{7.9pt}\selectfont\strut installer. Regrette de partir d'ici. Chambre chaude.\strut &  \\
      \cline{1-1}
      \leavevmode\fontsize{6.9pt}{7.9pt}\selectfont\strut \underline{2 Mardi}\strut &  \\
      \leavevmode\fontsize{6.9pt}{7.9pt}\selectfont\strut Ciel nuageux. Pluie menaçante.\strut &  \\
      \leavevmode\fontsize{6.9pt}{7.9pt}\selectfont\strut Temps doux, presque chaud.\strut &  \\
      \leavevmode\fontsize{6.9pt}{7.9pt}\selectfont\strut Matinée, à Louvencourt, ferme Cuisinier\strut &  \\
      \leavevmode\llap{\fontsize{6.4pt}{7.3pt}\selectfont\color{darkgray}15\hspace{1.8mm}}\fontsize{6.9pt}{7.9pt}\selectfont\strut expertise machine endommagée etc. \&\strut &  \\
      \leavevmode\fontsize{6.9pt}{7.9pt}\selectfont\strut puis à Raincheval.\strut &  \\
      \leavevmode\fontsize{6.9pt}{7.9pt}\selectfont\strut Après-midi, course en auto à\strut &  \\
      \leavevmode\fontsize{6.9pt}{7.9pt}\selectfont\strut Belle-Eglise, Ferme du Quesnoy\strut &  \\
      \leavevmode\fontsize{6.9pt}{7.9pt}\selectfont\strut \& Toutencourt.\strut &  \\
      \leavevmode\llap{\fontsize{6.4pt}{7.3pt}\selectfont\color{darkgray}20\hspace{1.8mm}}\fontsize{6.9pt}{7.9pt}\selectfont\strut Rentré 6h s., très fatigué pr plusieurs\strut & \smash{\parbox[t]{39mm}{\fontsize{6.2pt}{7.0pt}\selectfont\raggedright\color{black!88}\textbf{20 pr} \textit{scil.} pour. \textbf{20 s.} \textit{abbr. pro} soir.}} \\
      \leavevmode\fontsize{6.9pt}{7.9pt}\selectfont\strut heures d'auto.\strut &  \\
      \cline{1-1}
      \leavevmode\fontsize{6.9pt}{7.9pt}\selectfont\strut [ A Louvencourt canteen : Calvert Tooth Powder : 60 c.\strut & \smash{\parbox[t]{39mm}{\fontsize{6.2pt}{7.0pt}\selectfont\raggedright\color{black!88}\textbf{22} \textit{in marg. cum unco} (comptes). \textbf{22 c.} \textit{abbr. pro} centimes. — \textbf{23 boots} \textit{angl. in ms.}}} \\
      \leavevmode\fontsize{6.9pt}{7.9pt}\selectfont\strut Acheux : new boots, nails on : 1 f 50 ]\strut &  \\
\end{tabular}}%
\par\vspace*{\fill}
% --- Registres de bas de page (1. Critique interne [si pertinent] ; 2. Loci paralleli en latin ; 3. Commentaire historique en retrait) ---
\noindent\rule{\textwidth}{0.5pt}\par\vspace{0.65mm}
{\fontsize{6.6pt}{7.7pt}\selectfont%
\noindent\normalfont \textbf{Disposition :} L'année nouvelle (\textit{1917}) est mise en évidence dès l'en-tête souligné de la l.~1 (\textit{1\textsuperscript{er} Lundi — 1917 — \hfill Janvier.}) ; le feuillet regroupe les 1\textsuperscript{er} et 2 janvier 1917.\par\vspace{0.5mm}
\noindent\rule{22mm}{0.3pt}\par\vspace{0.45mm}
\noindent\normalfont\color{black!92} \textbf{4} (\textit{Toutencourt}) cf. supra f\textsuperscript{o}~49~v\textsuperscript{o} (p.~\pageref{fol:99}), f\textsuperscript{o}~56~r\textsuperscript{o} (p.~\pageref{fol:112}), infra f\textsuperscript{o}~62~r\textsuperscript{o} (p.~\pageref{fol:124}), f\textsuperscript{o}~75~r\textsuperscript{o} (p.~\pageref{fol:150}), \textit{et passim}. — \textbf{4} (\textit{Quesnoy}) cf. supra f\textsuperscript{o}~55~v\textsuperscript{o} (p.~\pageref{fol:111}), f\textsuperscript{o}~58~r\textsuperscript{o} (p.~\pageref{fol:116}), \textit{et passim}. — \textbf{5} (\textit{Authie}) cf. supra f\textsuperscript{o}~50~r\textsuperscript{o} (p.~\pageref{fol:100}), f\textsuperscript{o}~56~r\textsuperscript{o} (p.~\pageref{fol:112}), infra f\textsuperscript{o}~60~v\textsuperscript{o} (p.~\pageref{fol:121}), f\textsuperscript{o}~61~r\textsuperscript{o} (p.~\pageref{fol:122}), \textit{et passim}.\par\vspace{0.55mm}
\noindent\rule{38mm}{0.4pt}\par\vspace{0.5mm}
\noindent\hspace*{5mm}\parbox{\dimexpr\textwidth-5mm\relax}{\fontsize{6.6pt}{7.7pt}\selectfont\color{black!84}Les 1\textsuperscript{er} et 2 janvier 1917, le service des fourrages se poursuit sans trêve du Jour de l'An: tournée le lundi par \textbf{Authie} et \textbf{Pas-en-Artois} jusqu'à \textbf{Saulty} (Pas-de-Calais, sur la route de Doullens à Arras), et expertise contradictoire le mardi 2 janvier d'une machine agricole endommagée à la \textbf{ferme Cuisinier} à \textbf{Louvencourt} (où avait eu lieu l'incident du 24 décembre, f\textsuperscript{o}~56~r\textsuperscript{o}).}\par}
\newpage
\phantomsection\label{fac:118}
% --- En-tête du fac-similé (Page impaire / droite : Photographie du manuscrit) ---
\noindent\rule{\textwidth}{0.5pt}\par\vspace{0.55mm}
\noindent\textbf{\fontsize{6.8pt}{7.8pt}\selectfont FAC-SIMILÉ EN VIS-À-VIS — FOLIO 59 RECTO} \hfill \textbf{\fontsize{6.8pt}{7.8pt}\selectfont JANVIER 1917}\par\vspace{0.25mm}
\noindent{\fontsize{6.8pt}{8.0pt}\selectfont\textit{Livre-Journal d'Hippolyte Combes} — 1er – 2 janvier 1917 (Acheux, Toutencourt, Saulty, Louvencourt, Raincheval (Somme / Pas-de-Calais))}\par\vspace{0.55mm}
\noindent\rule{\textwidth}{0.35pt}\par\vspace*{\fill}
{\centering
  \setlength{\fboxsep}{1.2mm}%
  \setlength{\fboxrule}{0.4pt}%
  \fcolorbox{gray!55}{white}{%
    \includegraphics[width=0.95\textwidth,height=0.81\textheight,keepaspectratio]{\facsimiledir/20251126_105428.jpg}%
  }\par}
\vspace*{\fill}
\newpage

% =============================================================================
% VIS-A-VIS SPREAD : FOLIO 59 VERSO (20251126_105446.jpg)
% =============================================================================
\cleardoubleevenpage
\renewcommand{\currentmonthyear}{Janvier 1917}
\renewcommand{\currentfoliolabel}{Folio 59 verso}
\phantomsection\label{fol:119}
% --- En-tête du folio (Page paire / gauche : Transcription, apparat & registres de bas de page) ---
\noindent\rule{\textwidth}{0.5pt}\par\vspace{0.55mm}
\noindent\textbf{\footnotesize FOLIO 59 VERSO} \hfill \textbf{\footnotesize JANVIER 1917}\par\vspace{0.25mm}
\noindent{\fontsize{6.9pt}{8.0pt}\selectfont Course à Amiens : l'Hôtel de la Paix, le capitaine Trefcon et la cathédrale}\par\vspace{0.2mm}
\noindent{\fontsize{6.9pt}{8.0pt}\selectfont\textit{3 janvier 1917} \hfill \textit{Acheux, Belle-Église, Montigny, Amiens (Somme)}}\par\vspace{0.5mm}
\noindent\rule{\textwidth}{0.35pt}\par\vspace{0.8mm}
% --- Corps de la transcription diplomatique et apparat marginal en prolongement ---
{\renewcommand{\arraystretch}{0.87}%
\noindent\hspace*{3.5mm}\begin{tabular}{@{}p{87mm}@{\hspace{3.5mm}}p{39mm}@{}}
      \leavevmode\fontsize{6.9pt}{7.9pt}\selectfont\strut \underline{3 Mercredi} \hfill{} \underline{— Janvier /17 —}\strut &  \\
      \leavevmode\fontsize{6.9pt}{7.9pt}\selectfont\strut Course à Amiens\strut &  \\
      \leavevmode\fontsize{6.9pt}{7.9pt}\selectfont\strut Ciel couvert, pluie menacé tout le jour —\strut &  \\
      \leavevmode\fontsize{6.9pt}{7.9pt}\selectfont\strut Temps +tôt tiède \& doux —\strut &  \\
      \leavevmode\llap{\fontsize{6.4pt}{7.3pt}\selectfont\color{darkgray}5\hspace{1.8mm}}\fontsize{6.9pt}{7.9pt}\selectfont\strut Départ 9h m. 2 off., venus avec nous.\strut & \smash{\parbox[t]{39mm}{\fontsize{6.2pt}{7.0pt}\selectfont\raggedright\color{black!88}\textbf{5 off.} \textit{abbr. pro} officier(s). \textbf{5 m.} \textit{abbr. pro} matin.}} \\
      \leavevmode\fontsize{6.9pt}{7.9pt}\selectfont\strut Arrêt à Belle-Église et + long à Montigny,\strut &  \\
      \leavevmode\fontsize{6.9pt}{7.9pt}\selectfont\strut ch. M. Dubois Paul, maire — av. qui eu\strut & \smash{\parbox[t]{39mm}{\fontsize{6.2pt}{7.0pt}\selectfont\raggedright\color{black!88}\textbf{7 ch.} \textit{abbr. pro} chez.}} \\
      \leavevmode\fontsize{6.9pt}{7.9pt}\selectfont\strut conversation fort intéressante.\strut &  \\
      \leavevmode\fontsize{6.9pt}{7.9pt}\selectfont\strut Vers 11h, à Amiens.\strut &  \\
      \leavevmode\llap{\fontsize{6.4pt}{7.3pt}\selectfont\color{darkgray}10\hspace{1.8mm}}\fontsize{6.9pt}{7.9pt}\selectfont\strut Remisé, garé notre auto ds la cour de\strut & \smash{\parbox[t]{39mm}{\fontsize{6.2pt}{7.0pt}\selectfont\raggedright\color{black!88}\textbf{10 ds} \textit{abbr. pro} dans.}} \\
      \leavevmode\fontsize{6.9pt}{7.9pt}\selectfont\strut l'Hôtel de la Paix, rue — et dès lors,\strut &  \\
      \leavevmode\fontsize{6.9pt}{7.9pt}\selectfont\strut pour moi, liberté complète j'à 3 h 1/2 s.\strut & \smash{\parbox[t]{39mm}{\fontsize{6.2pt}{7.0pt}\selectfont\raggedright\color{black!88}\textbf{12 j'à} \textit{abbr. pro} jusqu'à. \textbf{12 1/2} \textit{scil.} demi / moitié. — \textbf{13 Cap.} \textit{abbr. pro} Capitaine. — \textbf{15 cathle} \textit{abbr. pro} cathédrale.}} \\
      \leavevmode\fontsize{6.9pt}{7.9pt}\selectfont\strut Rencontré Cap. Trefcon rue Sire Firmin Leroux\strut &  \\
      \leavevmode\fontsize{6.9pt}{7.9pt}\selectfont\strut (La Semeuse : instruments agricoles etc.)\strut &  \\
      \leavevmode\llap{\fontsize{6.4pt}{7.3pt}\selectfont\color{darkgray}15\hspace{1.8mm}}\fontsize{6.9pt}{7.9pt}\selectfont\strut Après déjeuner, visite à la cathle, y\strut &  \\
      \leavevmode\fontsize{6.9pt}{7.9pt}\selectfont\strut passé plus d'1 heure... Admiré à\strut &  \\
      \leavevmode\fontsize{6.9pt}{7.9pt}\selectfont\strut loisir choses \& détails déjà vus. —\strut &  \\
      \leavevmode\fontsize{6.9pt}{7.9pt}\selectfont\strut Promené çà et là, le long rues centre.\strut &  \\
      \leavevmode\fontsize{6.9pt}{7.9pt}\selectfont\strut Repartis vers 4h s.\strut & \smash{\parbox[t]{39mm}{\fontsize{6.2pt}{7.0pt}\selectfont\raggedright\color{black!88}\textbf{19 s.} \textit{abbr. pro} soir.}} \\
      \cline{1-1}
      \leavevmode\llap{\fontsize{6.4pt}{7.3pt}\selectfont\color{darkgray}20\hspace{1.8mm}}\fontsize{6.9pt}{7.9pt}\selectfont\strut [ — Amiens : 12 ctes ples : 40 c + 10 c. Ctes ples : 20 c\strut & \smash{\parbox[t]{39mm}{\fontsize{6.2pt}{7.0pt}\selectfont\raggedright\color{black!88}\textbf{20} \textit{in marg. cum unco} (comptes). \textbf{20 ctes ples} \textit{abbr. pro} cartes postales. — \textbf{21 c.} \textit{abbr. pro} centimes. — \textbf{22} \textit{lect. incert. in ms.}}} \\
      \leavevmode\fontsize{6.9pt}{7.9pt}\selectfont\strut 3 boutons khaki : 40 c. — Déjeuner : 1 f 95 c —\strut &  \\
      \leavevmode\fontsize{6.9pt}{7.9pt}\selectfont\strut 5 gravures [ocre ?] cathédrale : 2 f 50 —\strut &  \\
      \leavevmode\fontsize{6.9pt}{7.9pt}\selectfont\strut 1 bonnet police drap bleu : 3 f 75 c ]\strut &  \\
\end{tabular}}%
\par\vspace*{\fill}
% --- Registres de bas de page (1. Critique interne [si pertinent] ; 2. Loci paralleli en latin ; 3. Commentaire historique en retrait) ---
\noindent\rule{\textwidth}{0.5pt}\par\vspace{0.65mm}
{\fontsize{6.6pt}{7.7pt}\selectfont%
\noindent\normalfont\color{black!92} \textbf{2} (\textit{Amiens}) cf. supra f\textsuperscript{o}~25~v\textsuperscript{o} (p.~\pageref{fol:51}), f\textsuperscript{o}~49~r\textsuperscript{o} (p.~\pageref{fol:98}), infra f\textsuperscript{o}~63~r\textsuperscript{o} (p.~\pageref{fol:126}), f\textsuperscript{o}~64~v\textsuperscript{o} (p.~\pageref{fol:129}), \textit{et passim}. — \textbf{6} (\textit{Belle-Église}) cf. supra f\textsuperscript{o}~47~v\textsuperscript{o} (p.~\pageref{fol:95}), f\textsuperscript{o}~58~v\textsuperscript{o} (p.~\pageref{fol:117}), infra f\textsuperscript{o}~61~r\textsuperscript{o} (p.~\pageref{fol:122}), f\textsuperscript{o}~62~v\textsuperscript{o} (p.~\pageref{fol:125}), \textit{et passim}. — \textbf{6} (\textit{Montigny}) cf. infra f\textsuperscript{o}~61~v\textsuperscript{o} (p.~\pageref{fol:123}), f\textsuperscript{o}~79~r\textsuperscript{o} (p.~\pageref{fol:158}), \textit{et passim}.\par\vspace{0.55mm}
\noindent\rule{38mm}{0.4pt}\par\vspace{0.5mm}
\noindent\hspace*{5mm}\parbox{\dimexpr\textwidth-5mm\relax}{\fontsize{6.6pt}{7.7pt}\selectfont\color{black!84}Le mercredi 3 janvier 1917, en route vers \textbf{Amiens}, Hippolyte Combes s'arrête à \textbf{Montigny-sur-l'Hallue} (Somme, canton de Villers-Bocage) chez le maire \textbf{Paul Dubois} (et sa famille, qu'il retrouvera les 1\textsuperscript{er} et 2 mars 1917, f\textsuperscript{o}~80~v\textsuperscript{o}).}\par}
\newpage
\phantomsection\label{fac:119}
% --- En-tête du fac-similé (Page impaire / droite : Photographie du manuscrit) ---
\noindent\rule{\textwidth}{0.5pt}\par\vspace{0.55mm}
\noindent\textbf{\fontsize{6.8pt}{7.8pt}\selectfont FAC-SIMILÉ EN VIS-À-VIS — FOLIO 59 VERSO} \hfill \textbf{\fontsize{6.8pt}{7.8pt}\selectfont JANVIER 1917}\par\vspace{0.25mm}
\noindent{\fontsize{6.8pt}{8.0pt}\selectfont\textit{Livre-Journal d'Hippolyte Combes} — 3 janvier 1917 (Acheux, Belle-Église, Montigny, Amiens (Somme))}\par\vspace{0.55mm}
\noindent\rule{\textwidth}{0.35pt}\par\vspace*{\fill}
{\centering
  \setlength{\fboxsep}{1.2mm}%
  \setlength{\fboxrule}{0.4pt}%
  \fcolorbox{gray!55}{white}{%
    \includegraphics[width=0.95\textwidth,height=0.81\textheight,keepaspectratio]{\facsimiledir/20251126_105446.jpg}%
  }\par}
\vspace*{\fill}
\newpage

% =============================================================================
% VIS-A-VIS SPREAD : FOLIO 60 RECTO (20251126_105848.jpg)
% =============================================================================
\cleardoubleevenpage
\renewcommand{\currentmonthyear}{Janvier 1917}
\renewcommand{\currentfoliolabel}{Folio 60 recto}
\phantomsection\label{fol:120}
% --- En-tête du folio (Page paire / gauche : Transcription, apparat & registres de bas de page) ---
\noindent\rule{\textwidth}{0.5pt}\par\vspace{0.55mm}
\noindent\textbf{\footnotesize FOLIO 60 RECTO} \hfill \textbf{\footnotesize JANVIER 1917}\par\vspace{0.25mm}
\noindent{\fontsize{6.9pt}{8.0pt}\selectfont Concert des « XIII Corps Archie » et contrat de battage chez Mme Capron à Beauquesne}\par\vspace{0.2mm}
\noindent{\fontsize{6.9pt}{8.0pt}\selectfont\textit{4 – 5 janvier 1917} \hfill \textit{Acheux, Septenville, Beauquesne (Somme)}}\par\vspace{0.5mm}
\noindent\rule{\textwidth}{0.35pt}\par\vspace{0.8mm}
% --- Corps de la transcription diplomatique et apparat marginal en prolongement ---
{\renewcommand{\arraystretch}{0.87}%
\noindent\hspace*{3.5mm}\begin{tabular}{@{}p{87mm}@{\hspace{3.5mm}}p{39mm}@{}}
      \leavevmode\fontsize{6.9pt}{7.9pt}\selectfont\strut \underline{4 Jeudi} \hfill{} \underline{— Janvier /17}\strut &  \\
      \leavevmode\fontsize{6.9pt}{7.9pt}\selectfont\strut Tout le jour, repos.\strut &  \\
      \leavevmode\fontsize{6.9pt}{7.9pt}\selectfont\strut Matinée, fortes ondées. Après-midi, ciel s'est\strut &  \\
      \leavevmode\fontsize{6.9pt}{7.9pt}\selectfont\strut éclairci et tout ensoleillé.\strut &  \\
      \leavevmode\llap{\fontsize{6.4pt}{7.3pt}\selectfont\color{darkgray}5\hspace{1.8mm}}\fontsize{6.9pt}{7.9pt}\selectfont\strut Nuit, mervx clair de lune (entière) et\strut & \smash{\parbox[t]{39mm}{\fontsize{6.2pt}{7.0pt}\selectfont\raggedright\color{black!88}\textbf{5 mervx} \textit{abbr. pro} merveilleux. — \textbf{6 qq} \textit{abbr. pro} quelque(s). — \textbf{7 ds} \textit{abbr. pro} dans.}} \\
      \leavevmode\fontsize{6.9pt}{7.9pt}\selectfont\strut qq étoiles —\strut &  \\
      \leavevmode\fontsize{6.9pt}{7.9pt}\selectfont\strut Resté ds chambre, matin \& surtout après-\strut &  \\
      \leavevmode\fontsize{6.9pt}{7.9pt}\selectfont\strut midi, à ranger affaires, ou regarder spectacle\strut &  \\
      \leavevmode\fontsize{6.9pt}{7.9pt}\selectfont\strut rue, très animé.\strut &  \\
      \leavevmode\llap{\fontsize{6.4pt}{7.3pt}\selectfont\color{darkgray}10\hspace{1.8mm}}\fontsize{6.9pt}{7.9pt}\selectfont\strut Soir, concert des XIII Corps Archie, à la\strut &  \\
      \leavevmode\fontsize{6.9pt}{7.9pt}\selectfont\strut Y.M.C.A. hut. 2e partie programme, nos excellents\strut & \smash{\parbox[t]{39mm}{\fontsize{6.2pt}{7.0pt}\selectfont\raggedright\color{black!88}\textbf{11 Y.M.C.A.} \textit{angl.} Young Men's Christian Assoc.}} \\
      \leavevmode\fontsize{6.9pt}{7.9pt}\selectfont\strut chanteurs — duos très g.\strut &  \\
      \leavevmode\fontsize{6.9pt}{7.9pt}\selectfont\strut [ — Hair cut : 50 c\strut & \smash{\parbox[t]{39mm}{\fontsize{6.2pt}{7.0pt}\selectfont\raggedright\color{black!88}\textbf{13} \textit{in marg. cum unco} (comptes). \textbf{13 Hair cut} \textit{angl. in ms.} — \textbf{14 Y.M.C.A.} \textit{angl.} Young Men's Christian Assoc. — \textbf{15 Ven.} \textit{scil.} Vendredi. — \textbf{17 Lt} \textit{abbr. pro} Lieutenant.}} \\
      \leavevmode\fontsize{6.9pt}{7.9pt}\selectfont\strut — Concert Y.M.C.A. : 30 c ]\strut &  \\
      \cline{1-1}
      \leavevmode\llap{\fontsize{6.4pt}{7.3pt}\selectfont\color{darkgray}15\hspace{1.8mm}}\fontsize{6.9pt}{7.9pt}\selectfont\strut \underline{5 Ven.} — Tout le jour, en course.\strut &  \\
      \leavevmode\fontsize{6.9pt}{7.9pt}\selectfont\strut Septenville, ferme Moulard. —\strut &  \\
      \leavevmode\fontsize{6.9pt}{7.9pt}\selectfont\strut Beauquesne : déjeuné av. Lt Walsh, Lt Morris\strut &  \\
      \leavevmode\fontsize{6.9pt}{7.9pt}\selectfont\strut \& son interprète au Café Continental.\strut &  \\
      \leavevmode\fontsize{6.9pt}{7.9pt}\selectfont\strut Puis passé contrat location battage av.\strut &  \\
      \leavevmode\llap{\fontsize{6.4pt}{7.3pt}\selectfont\color{darkgray}20\hspace{1.8mm}}\fontsize{6.9pt}{7.9pt}\selectfont\strut Mme Capron.\strut & \smash{\parbox[t]{39mm}{\fontsize{6.2pt}{7.0pt}\selectfont\raggedright\color{black!88}\textbf{20 Mme} \textit{abbr. pro} Madame.}} \\
      \leavevmode\fontsize{6.9pt}{7.9pt}\selectfont\strut Rentré, nuit tombante. Mervx clair\strut &  \\
      \leavevmode\fontsize{6.9pt}{7.9pt}\selectfont\strut de lune. Routes boueuses comme ruisselantes\strut &  \\
      \leavevmode\fontsize{6.9pt}{7.9pt}\selectfont\strut de lumière —\strut &  \\
      \cline{1-1}
      \leavevmode\fontsize{6.9pt}{7.9pt}\selectfont\strut [ Bureau Mission fse XIIIe Corps Acheux : 78 f 75 ]\strut & \smash{\parbox[t]{39mm}{\fontsize{6.2pt}{7.0pt}\selectfont\raggedright\color{black!88}\textbf{24} \textit{in marg. cum unco} (comptes). \textbf{24 fse} \textit{abbr. pro} française.}} \\
\end{tabular}}%
\par\vspace*{\fill}
% --- Registres de bas de page (1. Critique interne [si pertinent] ; 2. Loci paralleli en latin ; 3. Commentaire historique en retrait) ---
\noindent\rule{\textwidth}{0.5pt}\par\vspace{0.65mm}
{\fontsize{6.6pt}{7.7pt}\selectfont%
\noindent\normalfont\color{black!92} \textbf{17} (\textit{Lt Walsh}) cf. supra f\textsuperscript{o}~49~v\textsuperscript{o} (p.~\pageref{fol:99}), f\textsuperscript{o}~53~r\textsuperscript{o} (p.~\pageref{fol:106}). — \textbf{17} (\textit{Beauquesne}) cf. supra f\textsuperscript{o}~37~v\textsuperscript{o} (p.~\pageref{fol:75}), f\textsuperscript{o}~58~v\textsuperscript{o} (p.~\pageref{fol:117}), infra f\textsuperscript{o}~61~r\textsuperscript{o} (p.~\pageref{fol:122}), f\textsuperscript{o}~61~v\textsuperscript{o} (p.~\pageref{fol:123}), \textit{et passim}. — \textbf{24} (\textit{Acheux}) cf. supra f\textsuperscript{o}~21~v\textsuperscript{o} (p.~\pageref{fol:43}), f\textsuperscript{o}~59~v\textsuperscript{o} (p.~\pageref{fol:119}), infra f\textsuperscript{o}~60~v\textsuperscript{o} (p.~\pageref{fol:121}), f\textsuperscript{o}~61~r\textsuperscript{o} (p.~\pageref{fol:122}), \textit{et passim}.\par\vspace{0.55mm}
\noindent\rule{38mm}{0.4pt}\par\vspace{0.5mm}
\noindent\hspace*{5mm}\parbox{\dimexpr\textwidth-5mm\relax}{\fontsize{6.6pt}{7.7pt}\selectfont\color{black!84}Après un second concert des \textbf{\textit{13th Corps Archie Singers}} au foyer de la Y.M.C.A. d'Acheux le 4 janvier, la journée du vendredi 5 janvier 1917 mène Hippolyte Combes à la ferme Moulard à \textbf{Septenville} (Rubempré) puis à \textbf{Beauquesne}, où il déjeune au Café Continental avec le \textbf{lieutenant Walsh} et le \textbf{lieutenant Morris} (officier britannique qu'il retrouvera avec amitié les 9 et 10 janvier, f\textsuperscript{o}~60~v\textsuperscript{o}--f\textsuperscript{o}~61~r\textsuperscript{o}), conclut avec \textbf{Mme Capron} le contrat de location de sa batteuse négocié le 26 décembre (f\textsuperscript{o}~57~v\textsuperscript{o}), et touche au retour à Acheux sa quinzaine de 78~fr.~75 au bureau du XIII\textsuperscript{e} Corps.}\par}
\newpage
\phantomsection\label{fac:120}
% --- En-tête du fac-similé (Page impaire / droite : Photographie du manuscrit) ---
\noindent\rule{\textwidth}{0.5pt}\par\vspace{0.55mm}
\noindent\textbf{\fontsize{6.8pt}{7.8pt}\selectfont FAC-SIMILÉ EN VIS-À-VIS — FOLIO 60 RECTO} \hfill \textbf{\fontsize{6.8pt}{7.8pt}\selectfont JANVIER 1917}\par\vspace{0.25mm}
\noindent{\fontsize{6.8pt}{8.0pt}\selectfont\textit{Livre-Journal d'Hippolyte Combes} — 4 – 5 janvier 1917 (Acheux, Septenville, Beauquesne (Somme))}\par\vspace{0.55mm}
\noindent\rule{\textwidth}{0.35pt}\par\vspace*{\fill}
{\centering
  \setlength{\fboxsep}{1.2mm}%
  \setlength{\fboxrule}{0.4pt}%
  \fcolorbox{gray!55}{white}{%
    \includegraphics[width=0.95\textwidth,height=0.81\textheight,keepaspectratio]{\facsimiledir/20251126_105848.jpg}%
  }\par}
\vspace*{\fill}
\newpage

% =============================================================================
% VIS-A-VIS SPREAD : FOLIO 60 VERSO (20251126_105908.jpg)
% =============================================================================
\cleardoubleevenpage
\renewcommand{\currentmonthyear}{Janvier 1917}
\renewcommand{\currentfoliolabel}{Folio 60 verso}
\phantomsection\label{fol:121}
% --- En-tête du folio (Page paire / gauche : Transcription, apparat & registres de bas de page) ---
\noindent\rule{\textwidth}{0.5pt}\par\vspace{0.55mm}
\noindent\textbf{\footnotesize FOLIO 60 VERSO} \hfill \textbf{\footnotesize JANVIER 1917}\par\vspace{0.25mm}
\noindent{\fontsize{6.9pt}{8.0pt}\selectfont Mouvements de troupes autour d'Acheux et visite à la 150th Company A.S.C. près d'Hédauville}\par\vspace{0.2mm}
\noindent{\fontsize{6.9pt}{8.0pt}\selectfont\textit{6 – 9 janvier 1917} \hfill \textit{Acheux, Louvencourt, Bertrancourt, Authie, Arquèves, Raincheval, Saint-Léger, Hédauville (Somme)}}\par\vspace{0.5mm}
\noindent\rule{\textwidth}{0.35pt}\par\vspace{0.8mm}
% --- Corps de la transcription diplomatique et apparat marginal en prolongement ---
{\renewcommand{\arraystretch}{0.87}%
\noindent\hspace*{3.5mm}\begin{tabular}{@{}p{87mm}@{\hspace{3.5mm}}p{39mm}@{}}
      \leavevmode\fontsize{6.9pt}{7.9pt}\selectfont\strut \underline{6 Sam.} \hfill{} \underline{Janv. 1917}\strut & \smash{\parbox[t]{39mm}{\fontsize{6.2pt}{7.0pt}\selectfont\raggedright\color{black!88}\textbf{1 Sam.} \textit{scil.} Samedi.}} \\
      \leavevmode\fontsize{6.9pt}{7.9pt}\selectfont\strut Très belle journée, ensoleillée\strut &  \\
      \leavevmode\fontsize{6.9pt}{7.9pt}\selectfont\strut mais 1 peu froide.\strut &  \\
      \leavevmode\fontsize{6.9pt}{7.9pt}\selectfont\strut Matinée, course en auto : Louvencourt,\strut &  \\
      \leavevmode\llap{\fontsize{6.4pt}{7.3pt}\selectfont\color{darkgray}5\hspace{1.8mm}}\fontsize{6.9pt}{7.9pt}\selectfont\strut Bertrancourt et Authie.\strut &  \\
      \leavevmode\fontsize{6.9pt}{7.9pt}\selectfont\strut Après-midi, liberté entière —\strut &  \\
      \leavevmode\fontsize{6.9pt}{7.9pt}\selectfont\strut Soir, au concert Y.M.C.A. Programme\strut & \smash{\parbox[t]{39mm}{\fontsize{6.2pt}{7.0pt}\selectfont\raggedright\color{black!88}\textbf{7 Y.M.C.A.} \textit{angl.} Young Men's Christian Assoc.}} \\
      \leavevmode\fontsize{6.9pt}{7.9pt}\selectfont\strut des 2 fois précédentes. Chants du jour. —\strut &  \\
      \leavevmode\fontsize{6.9pt}{7.9pt}\selectfont\strut Matinée, routes encombrées : déplacement\strut &  \\
      \leavevmode\llap{\fontsize{6.4pt}{7.3pt}\selectfont\color{darkgray}10\hspace{1.8mm}}\fontsize{6.9pt}{7.9pt}\selectfont\strut de troupes : convois d'artillerie, de A.S.C. etc.\strut & \smash{\parbox[t]{39mm}{\fontsize{6.2pt}{7.0pt}\selectfont\raggedright\color{black!88}\textbf{10 A.S.C.} \textit{angl.} Army Service Corps. — \textbf{11} \textit{in marg. cum unco} (comptes). \textbf{11 Y.M.C.A.} \textit{angl.} Young Men's Christian Assoc. — \textbf{12 Ma.} \textit{scil.} Mardi. — \textbf{14 St} \textit{abbr. pro} Saint.}} \\
      \leavevmode\fontsize{6.9pt}{7.9pt}\selectfont\strut [ — Washing : 70 c. — Concert Y.M.C.A. : 30 c ]\strut &  \\
      \cline{1-1}
      \leavevmode\fontsize{6.9pt}{7.9pt}\selectfont\strut \underline{9 Ma.} —\strut &  \\
      \leavevmode\fontsize{6.9pt}{7.9pt}\selectfont\strut Hier \& avant-hier, courses en auto\strut &  \\
      \leavevmode\fontsize{6.9pt}{7.9pt}\selectfont\strut matin \& soir, Arquèves, Raincheval, St-Léger.\strut &  \\
      \leavevmode\llap{\fontsize{6.4pt}{7.3pt}\selectfont\color{darkgray}15\hspace{1.8mm}}\fontsize{6.9pt}{7.9pt}\selectfont\strut Aujourd'hui, pareillement. Matinée,\strut &  \\
      \leavevmode\fontsize{6.9pt}{7.9pt}\selectfont\strut vu M. Guibet, à Raincheval...\strut &  \\
      \leavevmode\fontsize{6.9pt}{7.9pt}\selectfont\strut Après-midi, à Hédauville, route de\strut &  \\
      \leavevmode\fontsize{6.9pt}{7.9pt}\selectfont\strut Varennes, où retrouvé 150 Coy, 18th Divon\strut & \smash{\parbox[t]{39mm}{\fontsize{6.2pt}{7.0pt}\selectfont\raggedright\color{black!88}\textbf{18 Divon} \textit{abbr. pro} Division. \textbf{18 Coy} \textit{angl.} Company. — \textbf{19 ds} \textit{abbr. pro} dans.}} \\
      \leavevmode\fontsize{6.9pt}{7.9pt}\selectfont\strut ds 1 terrain fangeux, à 50 m. de la route.\strut &  \\
      \leavevmode\llap{\fontsize{6.4pt}{7.3pt}\selectfont\color{darkgray}20\hspace{1.8mm}}\fontsize{6.9pt}{7.9pt}\selectfont\strut Morris, était out riding —\strut &  \\
      \leavevmode\fontsize{6.9pt}{7.9pt}\selectfont\strut Temps + froid, neigé 1 peu, puis pluie\strut &  \\
      \leavevmode\fontsize{6.9pt}{7.9pt}\selectfont\strut fine — Toute la nuit \& durant journée\strut &  \\
      \leavevmode\fontsize{6.9pt}{7.9pt}\selectfont\strut très violente canonnade.\strut &  \\
      \cline{1-1}
      \leavevmode\fontsize{6.9pt}{7.9pt}\selectfont\strut [ — Avis de réception pour Imbert : 15 c — ]\strut & \smash{\parbox[t]{39mm}{\fontsize{6.2pt}{7.0pt}\selectfont\raggedright\color{black!88}\textbf{24} \textit{in marg. cum unco} (comptes).}} \\
\end{tabular}}%
\par\vspace*{\fill}
% --- Registres de bas de page (1. Critique interne [si pertinent] ; 2. Loci paralleli en latin ; 3. Commentaire historique en retrait) ---
\noindent\rule{\textwidth}{0.5pt}\par\vspace{0.65mm}
{\fontsize{6.6pt}{7.7pt}\selectfont%
\noindent\normalfont\color{black!92} \textbf{4} (\textit{Louvencourt}) cf. supra f\textsuperscript{o}~49~v\textsuperscript{o} (p.~\pageref{fol:99}), f\textsuperscript{o}~59~r\textsuperscript{o} (p.~\pageref{fol:118}), infra f\textsuperscript{o}~61~r\textsuperscript{o} (p.~\pageref{fol:122}), f\textsuperscript{o}~61~v\textsuperscript{o} (p.~\pageref{fol:123}), \textit{et passim}. — \textbf{5} (\textit{Bertrancourt}) cf. supra f\textsuperscript{o}~46~r\textsuperscript{o} (p.~\pageref{fol:92}), infra f\textsuperscript{o}~84~r\textsuperscript{o} (p.~\pageref{fol:168}). — \textbf{5} (\textit{Authie}) cf. supra f\textsuperscript{o}~50~r\textsuperscript{o} (p.~\pageref{fol:100}), f\textsuperscript{o}~59~r\textsuperscript{o} (p.~\pageref{fol:118}), infra f\textsuperscript{o}~61~r\textsuperscript{o} (p.~\pageref{fol:122}), f\textsuperscript{o}~77~r\textsuperscript{o} (p.~\pageref{fol:154}), \textit{et passim}.\par\vspace{0.55mm}
\noindent\rule{38mm}{0.4pt}\par\vspace{0.5mm}
\noindent\hspace*{5mm}\parbox{\dimexpr\textwidth-5mm\relax}{\fontsize{6.6pt}{7.7pt}\selectfont\color{black!84}Les mouvements de convois autour d'Acheux accompagnent l'extension du secteur du \textbf{V\textsuperscript{e} Corps britannique} dans la vallée de l'Ancre. À \textbf{Hédauville}, Hippolyte Combes rend visite le 9 janvier 1917 à la \textbf{150\textsuperscript{e} Compagnie de l'A.S.C.} (\textit{150th Company A.S.C.}, commandée par le major H.~G. Peachy), compagnie sœur de la 152\textsuperscript{e} Compagnie au sein du \textit{18th Divisional Train} (Maj. A.~F. Becke, \textit{Order of Battle}, Part~3A, p.~81), où sert son ami le \textbf{lieutenant Morris}.}\par}
\newpage
\phantomsection\label{fac:121}
% --- En-tête du fac-similé (Page impaire / droite : Photographie du manuscrit) ---
\noindent\rule{\textwidth}{0.5pt}\par\vspace{0.55mm}
\noindent\textbf{\fontsize{6.8pt}{7.8pt}\selectfont FAC-SIMILÉ EN VIS-À-VIS — FOLIO 60 VERSO} \hfill \textbf{\fontsize{6.8pt}{7.8pt}\selectfont JANVIER 1917}\par\vspace{0.25mm}
\noindent{\fontsize{6.8pt}{8.0pt}\selectfont\textit{Livre-Journal d'Hippolyte Combes} — 6 – 9 janvier 1917 (Acheux, Louvencourt, Bertrancourt, Authie, Arquèves, Raincheval, Saint-Léger, Hédauville (Somme))}\par\vspace{0.55mm}
\noindent\rule{\textwidth}{0.35pt}\par\vspace*{\fill}
{\centering
  \setlength{\fboxsep}{1.2mm}%
  \setlength{\fboxrule}{0.4pt}%
  \fcolorbox{gray!55}{white}{%
    \includegraphics[width=0.95\textwidth,height=0.81\textheight,keepaspectratio]{\facsimiledir/20251126_105908.jpg}%
  }\par}
\vspace*{\fill}
\newpage

% =============================================================================
% VIS-A-VIS SPREAD : FOLIO 61 RECTO (20251126_105920.jpg)
% =============================================================================
\cleardoubleevenpage
\renewcommand{\currentmonthyear}{Janvier 1917}
\renewcommand{\currentfoliolabel}{Folio 61 recto}
\phantomsection\label{fol:122}
% --- En-tête du folio (Page paire / gauche : Transcription, apparat & registres de bas de page) ---
\noindent\rule{\textwidth}{0.5pt}\par\vspace{0.55mm}
\noindent\textbf{\footnotesize FOLIO 61 RECTO} \hfill \textbf{\footnotesize JANVIER 1917}\par\vspace{0.25mm}
\noindent{\fontsize{6.9pt}{8.0pt}\selectfont Canonnade sur Miraumont, retrouvailles avec Morris et tempête de grésil sur les plateaux}\par\vspace{0.2mm}
\noindent{\fontsize{6.9pt}{8.0pt}\selectfont\textit{10 – 12 janvier 1917} \hfill \textit{Acheux, Saint-Léger, Belle-Église, Authie, Arquèves, Louvencourt, Raincheval, Beauquesne, Septenville, Rubempré (Somme)}}\par\vspace{0.5mm}
\noindent\rule{\textwidth}{0.35pt}\par\vspace{0.8mm}
% --- Corps de la transcription diplomatique et apparat marginal en prolongement ---
{\renewcommand{\arraystretch}{0.83}%
\noindent\hspace*{3.5mm}\begin{tabular}{@{}p{87mm}@{\hspace{3.5mm}}p{39mm}@{}}
      \leavevmode\fontsize{6.6pt}{7.5pt}\selectfont\strut \underline{10 Me.} \hfill{} \underline{— Janvier 1917 —}\strut &  \\
      \leavevmode\fontsize{6.6pt}{7.5pt}\selectfont\strut Toute la nuit \& matinée,\strut &  \\
      \leavevmode\fontsize{6.6pt}{7.5pt}\selectfont\strut très violente canonnade s/ ligne de Miraumont.\strut & \smash{\parbox[t]{39mm}{\fontsize{6.2pt}{7.0pt}\selectfont\raggedright\color{black!88}\textbf{3 s/} \textit{abbr. pro} sur.}} \\
      \leavevmode\fontsize{6.6pt}{7.5pt}\selectfont\strut Course en auto à St-Léger \& Belle-Église,\strut & \smash{\parbox[t]{39mm}{\fontsize{6.2pt}{7.0pt}\selectfont\raggedright\color{black!88}\textbf{4 St} \textit{abbr. pro} Saint.}} \\
      \leavevmode\llap{\fontsize{6.4pt}{7.3pt}\selectfont\color{darkgray}5\hspace{1.8mm}}\fontsize{6.6pt}{7.5pt}\selectfont\strut où rencontré, enfin, mon vieux Morris!\strut &  \\
      \leavevmode\fontsize{6.6pt}{7.5pt}\selectfont\strut Après-midi, liberté entière. Visite de\strut &  \\
      \leavevmode\fontsize{6.6pt}{7.5pt}\selectfont\strut Morris à mon billet. Causé de tt 1 peu.\strut & \smash{\parbox[t]{39mm}{\fontsize{6.2pt}{7.0pt}\selectfont\raggedright\color{black!88}\textbf{7 tt} \textit{abbr. pro} tout.}} \\
      \leavevmode\fontsize{6.6pt}{7.5pt}\selectfont\strut Temps froid, neige \& pluie glaciale.\strut &  \\
      \leavevmode\fontsize{6.6pt}{7.5pt}\selectfont\strut [ — Mandat-poste à Imbert, opticien,\strut & \smash{\parbox[t]{39mm}{\fontsize{6.2pt}{7.0pt}\selectfont\raggedright\color{black!88}\textbf{9} \textit{in marg. cum unco} (comptes).}} \\
      \leavevmode\llap{\fontsize{6.4pt}{7.3pt}\selectfont\color{darkgray}10\hspace{1.8mm}}\fontsize{6.6pt}{7.5pt}\selectfont\strut Paris, pour réparation pince-nez et pour\strut &  \\
      \leavevmode\fontsize{6.6pt}{7.5pt}\selectfont\strut 1 neuf d'après ordonnance : 19 fr. ]\strut & \smash{\parbox[t]{39mm}{\fontsize{6.2pt}{7.0pt}\selectfont\raggedright\color{black!88}\textbf{11 fr.} \textit{abbr. pro} franc(s).}} \\
      \cline{1-1}
      \leavevmode\fontsize{6.6pt}{7.5pt}\selectfont\strut \underline{11 J.} — Neige \& pluie fine, presque tt le jour.\strut & \smash{\parbox[t]{39mm}{\fontsize{6.2pt}{7.0pt}\selectfont\raggedright\color{black!88}\textbf{12 J.} \textit{scil.} Jeudi. \textbf{12 tt} \textit{abbr. pro} tout. — \textbf{13 off} \textit{abbr. pro} officier(s).}} \\
      \leavevmode\fontsize{6.6pt}{7.5pt}\selectfont\strut A whole day off —\strut &  \\
      \leavevmode\fontsize{6.6pt}{7.5pt}\selectfont\strut Mis à jour correspondance en arriéré.\strut &  \\
      \leavevmode\llap{\fontsize{6.4pt}{7.3pt}\selectfont\color{darkgray}15\hspace{1.8mm}}\fontsize{6.6pt}{7.5pt}\selectfont\strut Soir, dîné seul. Rentré tôt chambre.\strut &  \\
      \cline{1-1}
      \leavevmode\fontsize{6.6pt}{7.5pt}\selectfont\strut \underline{12 Ve.} Durant nuit, neigé encore.\strut & \smash{\parbox[t]{39mm}{\fontsize{6.2pt}{7.0pt}\selectfont\raggedright\color{black!88}\textbf{16 Ve.} \textit{scil.} Vendredi.}} \\
      \leavevmode\fontsize{6.6pt}{7.5pt}\selectfont\strut Toute matinée, averses et giboulées :\strut &  \\
      \leavevmode\fontsize{6.6pt}{7.5pt}\selectfont\strut grésil \& neige. Surpris par tempête au\strut &  \\
      \leavevmode\fontsize{6.6pt}{7.5pt}\selectfont\strut retour d'Authie \& sur la route d'Arquèves,\strut &  \\
      \leavevmode\llap{\fontsize{6.4pt}{7.3pt}\selectfont\color{darkgray}20\hspace{1.8mm}}\fontsize{6.6pt}{7.5pt}\selectfont\strut plateaux de Louvencourt.\strut &  \\
      \leavevmode\fontsize{6.6pt}{7.5pt}\selectfont\strut Courses en auto : Louvencourt, St-Léger,\strut & \smash{\parbox[t]{39mm}{\fontsize{6.2pt}{7.0pt}\selectfont\raggedright\color{black!88}\textbf{21 St} \textit{abbr. pro} Saint.}} \\
      \leavevmode\fontsize{6.6pt}{7.5pt}\selectfont\strut Raincheval, Beauquesne, où déjeuné au\strut &  \\
      \leavevmode\fontsize{6.6pt}{7.5pt}\selectfont\strut Café Continental.\strut &  \\
      \leavevmode\fontsize{6.6pt}{7.5pt}\selectfont\strut Après-midi, à Septenville \& Rubempré.\strut &  \\
      \leavevmode\llap{\fontsize{6.4pt}{7.3pt}\selectfont\color{darkgray}25\hspace{1.8mm}}\fontsize{6.6pt}{7.5pt}\selectfont\strut Rentré vers 6h, nuit noire.\strut &  \\
\end{tabular}}%
\par\vspace*{\fill}
% --- Registres de bas de page (1. Critique interne [si pertinent] ; 2. Loci paralleli en latin ; 3. Commentaire historique en retrait) ---
\noindent\rule{\textwidth}{0.5pt}\par\vspace{0.65mm}
{\fontsize{6.6pt}{7.7pt}\selectfont%
\noindent\normalfont\color{black!92} \textbf{3} (\textit{Miraumont}) cf. infra f\textsuperscript{o}~85~r\textsuperscript{o} (p.~\pageref{fol:170}). — \textbf{4} (\textit{Belle-Église}) cf. supra f\textsuperscript{o}~47~v\textsuperscript{o} (p.~\pageref{fol:95}), f\textsuperscript{o}~59~v\textsuperscript{o} (p.~\pageref{fol:119}), infra f\textsuperscript{o}~62~v\textsuperscript{o} (p.~\pageref{fol:125}), f\textsuperscript{o}~85~v\textsuperscript{o} (p.~\pageref{fol:171}), \textit{et passim}. — \textbf{4} (\textit{St-Léger}) cf. supra f\textsuperscript{o}~54~r\textsuperscript{o} (p.~\pageref{fol:108}), f\textsuperscript{o}~60~v\textsuperscript{o} (p.~\pageref{fol:121}), infra f\textsuperscript{o}~61~v\textsuperscript{o} (p.~\pageref{fol:123}).\par\vspace{0.55mm}
\noindent\rule{38mm}{0.4pt}\par\vspace{0.5mm}
\noindent\hspace*{5mm}\parbox{\dimexpr\textwidth-5mm\relax}{\fontsize{6.6pt}{7.7pt}\selectfont\color{black!84}La « très violente canonnade sur la ligne de Miraumont » notée dans la nuit et la matinée du mercredi \textbf{10 janvier 1917} correspond aux opérations offensives menées sur l'Ancre les \textbf{10 et 11 janvier 1917} par le \textbf{V\textsuperscript{e} Corps britannique} (attaque de la 7\textsuperscript{e} Division sur les tranchées allemandes dominant Beaumont-Hamel et le vallon de \textbf{Miraumont}; \textit{Official History, Military Operations France and Belgium, 1917}, vol.~I, p.~48--50).}\par}
\newpage
\phantomsection\label{fac:122}
% --- En-tête du fac-similé (Page impaire / droite : Photographie du manuscrit) ---
\noindent\rule{\textwidth}{0.5pt}\par\vspace{0.55mm}
\noindent\textbf{\fontsize{6.8pt}{7.8pt}\selectfont FAC-SIMILÉ EN VIS-À-VIS — FOLIO 61 RECTO} \hfill \textbf{\fontsize{6.8pt}{7.8pt}\selectfont JANVIER 1917}\par\vspace{0.25mm}
\noindent{\fontsize{6.8pt}{8.0pt}\selectfont\textit{Livre-Journal d'Hippolyte Combes} — 10 – 12 janvier 1917 (Acheux, Saint-Léger, Belle-Église, Authie, Arquèves, Louvencourt, Raincheval, Beauquesne, Septenville, Rubempré (Somme))}\par\vspace{0.55mm}
\noindent\rule{\textwidth}{0.35pt}\par\vspace*{\fill}
{\centering
  \setlength{\fboxsep}{1.2mm}%
  \setlength{\fboxrule}{0.4pt}%
  \fcolorbox{gray!55}{white}{%
    \includegraphics[width=0.95\textwidth,height=0.81\textheight,keepaspectratio]{\facsimiledir/20251126_105920.jpg}%
  }\par}
\vspace*{\fill}
\newpage

% =============================================================================
% VIS-A-VIS SPREAD : FOLIO 61 VERSO (20251126_105931.jpg)
% =============================================================================
\cleardoubleevenpage
\renewcommand{\currentmonthyear}{Janvier 1917}
\renewcommand{\currentfoliolabel}{Folio 61 verso}
\phantomsection\label{fol:123}
% --- En-tête du folio (Page paire / gauche : Transcription, apparat & registres de bas de page) ---
\noindent\rule{\textwidth}{0.5pt}\par\vspace{0.55mm}
\noindent\textbf{\footnotesize FOLIO 61 VERSO} \hfill \textbf{\footnotesize JANVIER 1917}\par\vspace{0.25mm}
\noindent{\fontsize{6.9pt}{8.0pt}\selectfont Fixation de la permission au 26 janvier, arbres ouatés de neige et course à Doullens}\par\vspace{0.2mm}
\noindent{\fontsize{6.9pt}{8.0pt}\selectfont\textit{13 – 15 janvier 1917} \hfill \textit{Acheux, Beaucourt, Montigny, Louvencourt, Saint-Léger, Beauquesne, Raincheval, Doullens (Somme)}}\par\vspace{0.5mm}
\noindent\rule{\textwidth}{0.35pt}\par\vspace{0.8mm}
% --- Corps de la transcription diplomatique et apparat marginal en prolongement ---
{\renewcommand{\arraystretch}{0.87}%
\noindent\hspace*{3.5mm}\begin{tabular}{@{}p{87mm}@{\hspace{3.5mm}}p{39mm}@{}}
      \leavevmode\fontsize{6.9pt}{7.9pt}\selectfont\strut \underline{13 Sam.} \hfill{} \underline{Janv. /17}\strut & \smash{\parbox[t]{39mm}{\fontsize{6.2pt}{7.0pt}\selectfont\raggedright\color{black!88}\textbf{1 Sam.} \textit{scil.} Samedi.}} \\
      \leavevmode\fontsize{6.9pt}{7.9pt}\selectfont\strut Tout le jour, pluie fine et (neige fondue)\strut &  \\
      \leavevmode\fontsize{6.9pt}{7.9pt}\selectfont\strut brouillard froid —\strut &  \\
      \leavevmode\fontsize{6.9pt}{7.9pt}\selectfont\strut Après-midi, course en auto à Beaucourt \&\strut &  \\
      \leavevmode\llap{\fontsize{6.4pt}{7.3pt}\selectfont\color{darkgray}5\hspace{1.8mm}}\fontsize{6.9pt}{7.9pt}\selectfont\strut Montigny.\strut &  \\
      \leavevmode\fontsize{6.9pt}{7.9pt}\selectfont\strut Matin, au bureau mission Ve corps : vu\strut &  \\
      \leavevmode\fontsize{6.9pt}{7.9pt}\selectfont\strut Cap. Huguet; mon départ fixé défint au 26 ct —\strut & \smash{\parbox[t]{39mm}{\fontsize{6.2pt}{7.0pt}\selectfont\raggedright\color{black!88}\textbf{7 Cap.} \textit{abbr. pro} Capitaine. \textbf{7 ct} \textit{abbr. pro} courant. — \textbf{8 D.} \textit{scil.} Dimanche.}} \\
      \cline{1-1}
      \leavevmode\fontsize{6.9pt}{7.9pt}\selectfont\strut \underline{14 D.} Neige tombé à gros flocons durant\strut &  \\
      \leavevmode\fontsize{6.9pt}{7.9pt}\selectfont\strut matinée entière. Toits \& arbres recouverts couche\strut &  \\
      \leavevmode\llap{\fontsize{6.4pt}{7.3pt}\selectfont\color{darkgray}10\hspace{1.8mm}}\fontsize{6.9pt}{7.9pt}\selectfont\strut blanche. Branches extt jolies avec ce léger\strut & \smash{\parbox[t]{39mm}{\fontsize{6.2pt}{7.0pt}\selectfont\raggedright\color{black!88}\textbf{10 extt} \textit{abbr. pro} extrêmement.}} \\
      \leavevmode\fontsize{6.9pt}{7.9pt}\selectfont\strut duvet — ainsi ouatées.\strut &  \\
      \leavevmode\fontsize{6.9pt}{7.9pt}\selectfont\strut Après-midi, courses à Louvencourt et\strut &  \\
      \leavevmode\fontsize{6.9pt}{7.9pt}\selectfont\strut à St-Léger.\strut & \smash{\parbox[t]{39mm}{\fontsize{6.2pt}{7.0pt}\selectfont\raggedright\color{black!88}\textbf{13 St} \textit{abbr. pro} Saint.}} \\
      \cline{1-1}
      \leavevmode\fontsize{6.9pt}{7.9pt}\selectfont\strut \underline{15 L.} Pluie fine, par intermittence.\strut & \smash{\parbox[t]{39mm}{\fontsize{6.2pt}{7.0pt}\selectfont\raggedright\color{black!88}\textbf{14 L.} \textit{scil.} Lundi.}} \\
      \leavevmode\llap{\fontsize{6.4pt}{7.3pt}\selectfont\color{darkgray}15\hspace{1.8mm}}\fontsize{6.9pt}{7.9pt}\selectfont\strut Course à Doullens. Paysages, asz pittoresques\strut &  \\
      \leavevmode\fontsize{6.9pt}{7.9pt}\selectfont\strut au-delà de Beauquesne. Champs couverts de\strut &  \\
      \leavevmode\fontsize{6.9pt}{7.9pt}\selectfont\strut neige. Masses sombres des arbres. Arbres \&\strut &  \\
      \leavevmode\fontsize{6.9pt}{7.9pt}\selectfont\strut maisons de village se détachent davantage\strut &  \\
      \leavevmode\fontsize{6.9pt}{7.9pt}\selectfont\strut au 1/2 du blanc (Raincheval). Brume monte.\strut & \smash{\parbox[t]{39mm}{\fontsize{6.2pt}{7.0pt}\selectfont\raggedright\color{black!88}\textbf{19 1/2} \textit{scil.} demi / moitié.}} \\
      \leavevmode\llap{\fontsize{6.4pt}{7.3pt}\selectfont\color{darkgray}20\hspace{1.8mm}}\fontsize{6.9pt}{7.9pt}\selectfont\strut 3 h. de liberté à Doullens : visité ville,\strut &  \\
      \leavevmode\fontsize{6.9pt}{7.9pt}\selectfont\strut arrêt église etc.\strut &  \\
      \cline{1-1}
      \leavevmode\fontsize{6.9pt}{7.9pt}\selectfont\strut [ — 2 douz. ctes ples : 1 f. — Lunch : 1 f 90\strut & \smash{\parbox[t]{39mm}{\fontsize{6.2pt}{7.0pt}\selectfont\raggedright\color{black!88}\textbf{22} \textit{in marg. cum unco} (comptes). \textbf{22 ctes ples} \textit{abbr. pro} cartes postales.}} \\
      \leavevmode\fontsize{6.9pt}{7.9pt}\selectfont\strut — 1 plaque chocolat : 60 c — ]\strut &  \\
\end{tabular}}%
\par\vspace*{\fill}
% --- Registres de bas de page (1. Critique interne [si pertinent] ; 2. Loci paralleli en latin ; 3. Commentaire historique en retrait) ---
\noindent\rule{\textwidth}{0.5pt}\par\vspace{0.65mm}
{\fontsize{6.6pt}{7.7pt}\selectfont%
\noindent\normalfont\color{black!92} \textbf{4} (\textit{Beaucourt}) cf. supra f\textsuperscript{o}~49~v\textsuperscript{o} (p.~\pageref{fol:99}), infra f\textsuperscript{o}~77~v\textsuperscript{o} (p.~\pageref{fol:155}), f\textsuperscript{o}~91~v\textsuperscript{o} (p.~\pageref{fol:183}). — \textbf{5} (\textit{Montigny}) cf. supra f\textsuperscript{o}~59~v\textsuperscript{o} (p.~\pageref{fol:119}), infra f\textsuperscript{o}~79~r\textsuperscript{o} (p.~\pageref{fol:158}), f\textsuperscript{o}~80~v\textsuperscript{o} (p.~\pageref{fol:161}), \textit{et passim}. — \textbf{7} (\textit{Huguet}) cf. supra f\textsuperscript{o}~47~r\textsuperscript{o} (p.~\pageref{fol:94}).\par\vspace{0.55mm}
\noindent\rule{38mm}{0.4pt}\par\vspace{0.5mm}
\noindent\hspace*{5mm}\parbox{\dimexpr\textwidth-5mm\relax}{\fontsize{6.6pt}{7.7pt}\selectfont\color{black!84}Le samedi 13 janvier 1917, au bureau de la Mission française d'Acheux — désormais rattaché au \textbf{V\textsuperscript{e} Corps britannique} (commandé par le lieutenant-général E.~A. Fanshawe) qui a relevé le XIII\textsuperscript{e} Corps dans ce secteur —, le \textbf{capitaine Huguet} fixe définitivement au vendredi 26 janvier le départ d'Hippolyte Combes pour sa deuxième permission.}\par}
\newpage
\phantomsection\label{fac:123}
% --- En-tête du fac-similé (Page impaire / droite : Photographie du manuscrit) ---
\noindent\rule{\textwidth}{0.5pt}\par\vspace{0.55mm}
\noindent\textbf{\fontsize{6.8pt}{7.8pt}\selectfont FAC-SIMILÉ EN VIS-À-VIS — FOLIO 61 VERSO} \hfill \textbf{\fontsize{6.8pt}{7.8pt}\selectfont JANVIER 1917}\par\vspace{0.25mm}
\noindent{\fontsize{6.8pt}{8.0pt}\selectfont\textit{Livre-Journal d'Hippolyte Combes} — 13 – 15 janvier 1917 (Acheux, Beaucourt, Montigny, Louvencourt, Saint-Léger, Beauquesne, Raincheval, Doullens (Somme))}\par\vspace{0.55mm}
\noindent\rule{\textwidth}{0.35pt}\par\vspace*{\fill}
{\centering
  \setlength{\fboxsep}{1.2mm}%
  \setlength{\fboxrule}{0.4pt}%
  \fcolorbox{gray!55}{white}{%
    \includegraphics[width=0.95\textwidth,height=0.81\textheight,keepaspectratio]{\facsimiledir/20251126_105931.jpg}%
  }\par}
\vspace*{\fill}
\newpage

% =============================================================================
% VIS-A-VIS SPREAD : FOLIO 62 RECTO (20251126_105945.jpg)
% =============================================================================
\cleardoubleevenpage
\renewcommand{\currentmonthyear}{Janvier 1917}
\renewcommand{\currentfoliolabel}{Folio 62 recto}
\phantomsection\label{fol:124}
% --- En-tête du folio (Page paire / gauche : Transcription, apparat & registres de bas de page) ---
\noindent\rule{\textwidth}{0.5pt}\par\vspace{0.55mm}
\noindent\textbf{\footnotesize FOLIO 62 RECTO} \hfill \textbf{\footnotesize JANVIER 1917}\par\vspace{0.25mm}
\noindent{\fontsize{6.9pt}{8.0pt}\selectfont Étendues blanches d'Hérissart et visite à Auxi-le-Château chez le collectionneur M. Renard}\par\vspace{0.2mm}
\noindent{\fontsize{6.9pt}{8.0pt}\selectfont\textit{16 – 17 janvier 1917} \hfill \textit{Acheux, Bouzincourt, Senlis-le-Sec, Rubempré, Hérissart, Toutencourt, Doullens, Auxi-le-Château (Somme / Pas-de-Calais)}}\par\vspace{0.5mm}
\noindent\rule{\textwidth}{0.35pt}\par\vspace{0.8mm}
% --- Corps de la transcription diplomatique et apparat marginal en prolongement ---
{\renewcommand{\arraystretch}{0.83}%
\noindent\hspace*{3.5mm}\begin{tabular}{@{}p{87mm}@{\hspace{3.5mm}}p{39mm}@{}}
      \leavevmode\fontsize{6.6pt}{7.5pt}\selectfont\strut \underline{16 Mardi} \hfill{} \underline{Janvier 1917}\strut &  \\
      \leavevmode\fontsize{6.6pt}{7.5pt}\selectfont\strut Nuit, gelé très fort.\strut &  \\
      \leavevmode\fontsize{6.6pt}{7.5pt}\selectfont\strut Durant journée, ciel nuageux. Temps froid, mais sec.\strut &  \\
      \leavevmode\fontsize{6.6pt}{7.5pt}\selectfont\strut Après-midi, course en auto (Sunbeam) à\strut &  \\
      \leavevmode\llap{\fontsize{6.4pt}{7.3pt}\selectfont\color{darkgray}5\hspace{1.8mm}}\fontsize{6.6pt}{7.5pt}\selectfont\strut Bouzincourt, Senlis \& Rubempré. Ds ce dernier village\strut & \smash{\parbox[t]{39mm}{\fontsize{6.2pt}{7.0pt}\selectfont\raggedright\color{black!88}\textbf{5 Ds} \textit{abbr. pro} Dans.}} \\
      \leavevmode\fontsize{6.6pt}{7.5pt}\selectfont\strut allé de ferme en ferme, pour trouver paille.\strut &  \\
      \leavevmode\fontsize{6.6pt}{7.5pt}\selectfont\strut Au retour, dévalé en vitesse coteaux pente\strut &  \\
      \leavevmode\fontsize{6.6pt}{7.5pt}\selectfont\strut Hérissart à Toutencourt. Coteaux et champs\strut &  \\
      \leavevmode\fontsize{6.6pt}{7.5pt}\selectfont\strut encore couverts neige : gdes \& mornes étendues\strut & \smash{\parbox[t]{39mm}{\fontsize{6.2pt}{7.0pt}\selectfont\raggedright\color{black!88}\textbf{9 gdes} \textit{abbr. pro} grandes.}} \\
      \leavevmode\llap{\fontsize{6.4pt}{7.3pt}\selectfont\color{darkgray}10\hspace{1.8mm}}\fontsize{6.6pt}{7.5pt}\selectfont\strut blanches.\strut &  \\
      \leavevmode\fontsize{6.6pt}{7.5pt}\selectfont\strut [ — Ctes ples : 30 c. ]\strut & \smash{\parbox[t]{39mm}{\fontsize{6.2pt}{7.0pt}\selectfont\raggedright\color{black!88}\textbf{11} \textit{in marg. cum unco} (comptes). \textbf{11 Ctes ples} \textit{abbr. pro} Cartes postales. — \textbf{12 Merc.} \textit{scil.} Mercredi. \textbf{12 abondt} \textit{abbr. pro} abondamment. — \textbf{14 j'à} \textit{abbr. pro} jusqu'à. — \textbf{15 ch.} \textit{abbr. pro} chez. — \textbf{18 ns} \textit{abbr. pro} nous.}} \\
      \cline{1-1}
      \leavevmode\fontsize{6.6pt}{7.5pt}\selectfont\strut \underline{17 Merc.} Nuit neigé abondt. Au réveil\strut &  \\
      \leavevmode\fontsize{6.6pt}{7.5pt}\selectfont\strut trouvé toits, jardins, rues, tout blancs.\strut &  \\
      \leavevmode\fontsize{6.6pt}{7.5pt}\selectfont\strut Tout le jour, en course. Venu j'à\strut &  \\
      \leavevmode\llap{\fontsize{6.4pt}{7.3pt}\selectfont\color{darkgray}15\hspace{1.8mm}}\fontsize{6.6pt}{7.5pt}\selectfont\strut Auxi-le-Château, ch. M. Renard. Jusqu'à Doullens\strut &  \\
      \leavevmode\fontsize{6.6pt}{7.5pt}\selectfont\strut plaines mornes. Au-delà, région pittoresque,\strut &  \\
      \leavevmode\fontsize{6.6pt}{7.5pt}\selectfont\strut vergers, petits coteaux boisés.\strut &  \\
      \leavevmode\fontsize{6.6pt}{7.5pt}\selectfont\strut M. Renard ns a invités à déjeuner.\strut &  \\
      \leavevmode\fontsize{6.6pt}{7.5pt}\selectfont\strut Ns a montré ses objets d'art précieux :\strut &  \\
      \leavevmode\llap{\fontsize{6.4pt}{7.3pt}\selectfont\color{darkgray}20\hspace{1.8mm}}\fontsize{6.6pt}{7.5pt}\selectfont\strut notamt et surtout un coffret en paille\strut & \smash{\parbox[t]{39mm}{\fontsize{6.2pt}{7.0pt}\selectfont\raggedright\color{black!88}\textbf{20 notamt} \textit{abbr. pro} notamment.}} \\
      \leavevmode\fontsize{6.6pt}{7.5pt}\selectfont\strut acheté en 1858 au bagne de Toulon.\strut &  \\
      \leavevmode\fontsize{6.6pt}{7.5pt}\selectfont\strut Matinée entière, neige fine, tourmente\strut &  \\
      \leavevmode\fontsize{6.6pt}{7.5pt}\selectfont\strut oblique, sous les rafales.\strut &  \\
      \leavevmode\fontsize{6.6pt}{7.5pt}\selectfont\strut Soir, température s'est radoucie considt\strut &  \\
      \leavevmode\llap{\fontsize{6.4pt}{7.3pt}\selectfont\color{darkgray}25\hspace{1.8mm}}\fontsize{6.6pt}{7.5pt}\selectfont\strut et dégel commencé —\strut &  \\
\end{tabular}}%
\par\vspace*{\fill}
% --- Registres de bas de page (1. Critique interne [si pertinent] ; 2. Loci paralleli en latin ; 3. Commentaire historique en retrait) ---
\noindent\rule{\textwidth}{0.5pt}\par\vspace{0.65mm}
{\fontsize{6.6pt}{7.7pt}\selectfont%
\noindent\normalfont\color{black!92} \textbf{5} (\textit{Senlis}) cf. supra f\textsuperscript{o}~37~v\textsuperscript{o} (p.~\pageref{fol:75}), f\textsuperscript{o}~58~v\textsuperscript{o} (p.~\pageref{fol:117}), infra f\textsuperscript{o}~75~v\textsuperscript{o} (p.~\pageref{fol:151}), f\textsuperscript{o}~79~v\textsuperscript{o} (p.~\pageref{fol:159}), \textit{et passim}. — \textbf{5} (\textit{Bouzincourt}) cf. supra f\textsuperscript{o}~38~v\textsuperscript{o} (p.~\pageref{fol:77}), f\textsuperscript{o}~50~v\textsuperscript{o} (p.~\pageref{fol:101}), \textit{et passim}. — \textbf{5} (\textit{Rubempré}) cf. supra f\textsuperscript{o}~47~v\textsuperscript{o} (p.~\pageref{fol:95}), f\textsuperscript{o}~61~r\textsuperscript{o} (p.~\pageref{fol:122}), infra f\textsuperscript{o}~74~v\textsuperscript{o} (p.~\pageref{fol:149}), f\textsuperscript{o}~80~v\textsuperscript{o} (p.~\pageref{fol:161}), \textit{et passim}.\par\vspace{0.55mm}
\noindent\rule{38mm}{0.4pt}\par\vspace{0.5mm}
\noindent\hspace*{5mm}\parbox{\dimexpr\textwidth-5mm\relax}{\fontsize{6.6pt}{7.7pt}\selectfont\color{black!84}Après une tournée de recherche de paille dans les fermes enneigées de Bouzincourt, Senlis-le-Sec et Rubempré le 16 janvier, l'excursion du mercredi 17 janvier 1917 conduit Hippolyte Combes par Doullens jusqu'à \textbf{Auxi-le-Château} (Pas-de-Calais, dans la vallée de l'Authie, à 38~km au nord-ouest d'Acheux).}\par}
\newpage
\phantomsection\label{fac:124}
% --- En-tête du fac-similé (Page impaire / droite : Photographie du manuscrit) ---
\noindent\rule{\textwidth}{0.5pt}\par\vspace{0.55mm}
\noindent\textbf{\fontsize{6.8pt}{7.8pt}\selectfont FAC-SIMILÉ EN VIS-À-VIS — FOLIO 62 RECTO} \hfill \textbf{\fontsize{6.8pt}{7.8pt}\selectfont JANVIER 1917}\par\vspace{0.25mm}
\noindent{\fontsize{6.8pt}{8.0pt}\selectfont\textit{Livre-Journal d'Hippolyte Combes} — 16 – 17 janvier 1917 (Acheux, Bouzincourt, Senlis-le-Sec, Rubempré, Hérissart, Toutencourt, Doullens, Auxi-le-Château (Somme / Pas-de-Calais))}\par\vspace{0.55mm}
\noindent\rule{\textwidth}{0.35pt}\par\vspace*{\fill}
{\centering
  \setlength{\fboxsep}{1.2mm}%
  \setlength{\fboxrule}{0.4pt}%
  \fcolorbox{gray!55}{white}{%
    \includegraphics[width=0.95\textwidth,height=0.81\textheight,keepaspectratio]{\facsimiledir/20251126_105945.jpg}%
  }\par}
\vspace*{\fill}
\newpage

% =============================================================================
% VIS-A-VIS SPREAD : FOLIO 62 VERSO (20251126_105956.jpg)
% =============================================================================
\cleardoubleevenpage
\renewcommand{\currentmonthyear}{Janvier 1917}
\renewcommand{\currentfoliolabel}{Folio 62 verso}
\phantomsection\label{fol:125}
% --- En-tête du folio (Page paire / gauche : Transcription, apparat & registres de bas de page) ---
\noindent\rule{\textwidth}{0.5pt}\par\vspace{0.55mm}
\noindent\textbf{\footnotesize FOLIO 62 VERSO} \hfill \textbf{\footnotesize JANVIER 1917}\par\vspace{0.25mm}
\noindent{\fontsize{6.9pt}{8.0pt}\selectfont Plaines blanches confondues avec le ciel mat et neige cristallisée à Puchevillers}\par\vspace{0.2mm}
\noindent{\fontsize{6.9pt}{8.0pt}\selectfont\textit{18 – 19 janvier 1917} \hfill \textit{Acheux, Belle-Église, Raincheval, Puchevillers (Somme)}}\par\vspace{0.5mm}
\noindent\rule{\textwidth}{0.35pt}\par\vspace{0.8mm}
% --- Corps de la transcription diplomatique et apparat marginal en prolongement ---
{\renewcommand{\arraystretch}{0.87}%
\noindent\hspace*{3.5mm}\begin{tabular}{@{}p{87mm}@{\hspace{3.5mm}}p{39mm}@{}}
      \leavevmode\fontsize{6.9pt}{7.9pt}\selectfont\strut \underline{18 Jeudi} Matin (canonnade) — \hfill{} \raisebox{0.5ex}{\underline{Janvier /17}}\strut &  \\
      \leavevmode\fontsize{6.9pt}{7.9pt}\selectfont\strut Temps froid. Neige\strut &  \\
      \leavevmode\fontsize{6.9pt}{7.9pt}\selectfont\strut restée s/ toits \& champs.\strut & \smash{\parbox[t]{39mm}{\fontsize{6.2pt}{7.0pt}\selectfont\raggedright\color{black!88}\textbf{3 s/} \textit{abbr. pro} sur.}} \\
      \leavevmode\fontsize{6.9pt}{7.9pt}\selectfont\strut Matin, course en auto à Belle-Église,\strut &  \\
      \leavevmode\llap{\fontsize{6.4pt}{7.3pt}\selectfont\color{darkgray}5\hspace{1.8mm}}\fontsize{6.9pt}{7.9pt}\selectfont\strut Raincheval \& Puchevillers. Plaines blanches\strut &  \\
      \leavevmode\fontsize{6.9pt}{7.9pt}\selectfont\strut se confondaient à l'horizon rapproché avec ciel\strut &  \\
      \leavevmode\fontsize{6.9pt}{7.9pt}\selectfont\strut mat. Légère buée —\strut &  \\
      \leavevmode\fontsize{6.9pt}{7.9pt}\selectfont\strut Après-midi, liberté entière.\strut &  \\
      \cline{1-1}
      \leavevmode\fontsize{6.9pt}{7.9pt}\selectfont\strut \underline{19 Vendredi}\strut &  \\
      \leavevmode\llap{\fontsize{6.4pt}{7.3pt}\selectfont\color{darkgray}10\hspace{1.8mm}}\fontsize{6.9pt}{7.9pt}\selectfont\strut Ciel mat \& bas. Tempérre, pas froide.\strut &  \\
      \leavevmode\fontsize{6.9pt}{7.9pt}\selectfont\strut Neige pas encore fondue. Toits, arbres et\strut &  \\
      \leavevmode\fontsize{6.9pt}{7.9pt}\selectfont\strut champs tout blancs.\strut &  \\
      \leavevmode\fontsize{6.9pt}{7.9pt}\selectfont\strut Après-midi, course à Puchevillers, par\strut &  \\
      \leavevmode\fontsize{6.9pt}{7.9pt}\selectfont\strut Raincheval. Neige comme cristallisée, durcie\strut &  \\
      \leavevmode\llap{\fontsize{6.4pt}{7.3pt}\selectfont\color{darkgray}15\hspace{1.8mm}}\fontsize{6.9pt}{7.9pt}\selectfont\strut ds les champs. Buée grise s'élève à l'horizon.\strut & \smash{\parbox[t]{39mm}{\fontsize{6.2pt}{7.0pt}\selectfont\raggedright\color{black!88}\textbf{15 ds} \textit{abbr. pro} dans.}} \\
      \leavevmode\fontsize{6.9pt}{7.9pt}\selectfont\strut Froid pénétrant. Il gèle déjà très fort.\strut &  \\
      \cline{1-1}
      \leavevmode\fontsize{6.9pt}{7.9pt}\selectfont\strut [ — Matin, au bureau mission\strut & \smash{\parbox[t]{39mm}{\fontsize{6.2pt}{7.0pt}\selectfont\raggedright\color{black!88}\textbf{17} \textit{in marg. cum unco} (comptes). — \textbf{18 fr.} \textit{abbr. pro} franc(s).}} \\
      \leavevmode\fontsize{6.9pt}{7.9pt}\selectfont\strut Ve corps, solde : 84 fr. — ]\strut &  \\
\end{tabular}}%
\par\vspace*{\fill}
% --- Registres de bas de page (1. Critique interne [si pertinent] ; 2. Loci paralleli en latin ; 3. Commentaire historique en retrait) ---
\noindent\rule{\textwidth}{0.5pt}\par\vspace{0.65mm}
{\fontsize{6.6pt}{7.7pt}\selectfont%
\noindent\normalfont\color{black!92} \textbf{4} (\textit{Belle-Église}) cf. supra f\textsuperscript{o}~47~v\textsuperscript{o} (p.~\pageref{fol:95}), f\textsuperscript{o}~61~r\textsuperscript{o} (p.~\pageref{fol:122}), infra f\textsuperscript{o}~85~v\textsuperscript{o} (p.~\pageref{fol:171}), f\textsuperscript{o}~87~v\textsuperscript{o} (p.~\pageref{fol:175}), \textit{et passim}. — \textbf{5} (\textit{Raincheval}) cf. supra f\textsuperscript{o}~20~v\textsuperscript{o} (p.~\pageref{fol:41}), f\textsuperscript{o}~61~v\textsuperscript{o} (p.~\pageref{fol:123}), infra f\textsuperscript{o}~73~v\textsuperscript{o} (p.~\pageref{fol:147}), f\textsuperscript{o}~77~v\textsuperscript{o} (p.~\pageref{fol:155}), \textit{et passim}. — \textbf{5} (\textit{Puchevillers}) cf. supra f\textsuperscript{o}~37~v\textsuperscript{o} (p.~\pageref{fol:75}), f\textsuperscript{o}~56~r\textsuperscript{o} (p.~\pageref{fol:112}), infra f\textsuperscript{o}~73~v\textsuperscript{o} (p.~\pageref{fol:147}), f\textsuperscript{o}~77~r\textsuperscript{o} (p.~\pageref{fol:154}), \textit{et passim}.\par\vspace{0.55mm}
\noindent\rule{38mm}{0.4pt}\par\vspace{0.5mm}
\noindent\hspace*{5mm}\parbox{\dimexpr\textwidth-5mm\relax}{\fontsize{6.6pt}{7.7pt}\selectfont\color{black!84}Sous la vague de froid polaire de la mi-janvier 1917 en Picardie, Hippolyte Combes poursuit les tournées quotidiennes de réquisition à Belle-Église, Raincheval et Puchevillers et perçoit sa solde bimensuelle (84~francs) au bureau de la Mission française du \textbf{V\textsuperscript{e} Corps britannique}.}\par}
\newpage
\phantomsection\label{fac:125}
% --- En-tête du fac-similé (Page impaire / droite : Photographie du manuscrit) ---
\noindent\rule{\textwidth}{0.5pt}\par\vspace{0.55mm}
\noindent\textbf{\fontsize{6.8pt}{7.8pt}\selectfont FAC-SIMILÉ EN VIS-À-VIS — FOLIO 62 VERSO} \hfill \textbf{\fontsize{6.8pt}{7.8pt}\selectfont JANVIER 1917}\par\vspace{0.25mm}
\noindent{\fontsize{6.8pt}{8.0pt}\selectfont\textit{Livre-Journal d'Hippolyte Combes} — 18 – 19 janvier 1917 (Acheux, Belle-Église, Raincheval, Puchevillers (Somme))}\par\vspace{0.55mm}
\noindent\rule{\textwidth}{0.35pt}\par\vspace*{\fill}
{\centering
  \setlength{\fboxsep}{1.2mm}%
  \setlength{\fboxrule}{0.4pt}%
  \fcolorbox{gray!55}{white}{%
    \includegraphics[width=0.95\textwidth,height=0.81\textheight,keepaspectratio]{\facsimiledir/20251126_105956.jpg}%
  }\par}
\vspace*{\fill}
\newpage

% =============================================================================
% VIS-A-VIS SPREAD : FOLIO 63 RECTO (20251126_110011.jpg)
% =============================================================================
\cleardoubleevenpage
\renewcommand{\currentmonthyear}{Janvier 1917}
\renewcommand{\currentfoliolabel}{Folio 63 recto}
\phantomsection\label{fol:126}
% --- En-tête du folio (Page paire / gauche : Transcription, apparat & registres de bas de page) ---
\noindent\rule{\textwidth}{0.5pt}\par\vspace{0.55mm}
\noindent\textbf{\footnotesize FOLIO 63 RECTO} \hfill \textbf{\footnotesize JANVIER 1917}\par\vspace{0.25mm}
\noindent{\fontsize{6.9pt}{8.0pt}\selectfont Grand gel sec, retour à pied de Louvencourt et préparatifs de permission}\par\vspace{0.2mm}
\noindent{\fontsize{6.9pt}{8.0pt}\selectfont\textit{20 – 23 janvier 1917} \hfill \textit{Acheux, Louvencourt (Somme)}}\par\vspace{0.5mm}
\noindent\rule{\textwidth}{0.35pt}\par\vspace{0.8mm}
% --- Corps de la transcription diplomatique et apparat marginal en prolongement ---
{\renewcommand{\arraystretch}{0.87}%
\noindent\hspace*{3.5mm}\begin{tabular}{@{}p{87mm}@{\hspace{3.5mm}}p{39mm}@{}}
      \leavevmode\fontsize{6.9pt}{7.9pt}\selectfont\strut \underline{20 Sam.} \hfill{} \underline{— Janvier 1917}\strut & \smash{\parbox[t]{39mm}{\fontsize{6.2pt}{7.0pt}\selectfont\raggedright\color{black!88}\textbf{1 Sam.} \textit{scil.} Samedi.}} \\
      \leavevmode\fontsize{6.9pt}{7.9pt}\selectfont\strut Repos tout le jour.\strut &  \\
      \leavevmode\fontsize{6.9pt}{7.9pt}\selectfont\strut Chambre ou office : lu — et toilette.\strut &  \\
      \leavevmode\fontsize{6.9pt}{7.9pt}\selectfont\strut Neige ne fond pas. Champs \& toits\strut &  \\
      \leavevmode\llap{\fontsize{6.4pt}{7.3pt}\selectfont\color{darkgray}5\hspace{1.8mm}}\fontsize{6.9pt}{7.9pt}\selectfont\strut encore tout blancs.\strut &  \\
      \leavevmode\fontsize{6.9pt}{7.9pt}\selectfont\strut Froid sec. Gelé très fortement.\strut &  \\
      \leavevmode\fontsize{6.9pt}{7.9pt}\selectfont\strut Sol resté durci. —\strut &  \\
      \cline{1-1}
      \leavevmode\fontsize{6.9pt}{7.9pt}\selectfont\strut \underline{21 Dim.} Nuit, gelé très fort. Froid sec.\strut & \smash{\parbox[t]{39mm}{\fontsize{6.2pt}{7.0pt}\selectfont\raggedright\color{black!88}\textbf{8 Dim.} \textit{scil.} Dimanche.}} \\
      \leavevmode\fontsize{6.9pt}{7.9pt}\selectfont\strut Matinée, en auto j'à Louvencourt, où\strut & \smash{\parbox[t]{39mm}{\fontsize{6.2pt}{7.0pt}\selectfont\raggedright\color{black!88}\textbf{9 j'à} \textit{abbr. pro} jusqu'à.}} \\
      \leavevmode\llap{\fontsize{6.4pt}{7.3pt}\selectfont\color{darkgray}10\hspace{1.8mm}}\fontsize{6.9pt}{7.9pt}\selectfont\strut vu Mme Cuis[?]. Sur ma demande, retour\strut & \smash{\parbox[t]{39mm}{\fontsize{6.2pt}{7.0pt}\selectfont\raggedright\color{black!88}\textbf{10 Cuis} \textit{lect. incert. in ms.} \textbf{10 Mme} \textit{abbr. pro} Madame.}} \\
      \leavevmode\fontsize{6.9pt}{7.9pt}\selectfont\strut à pied. Sol durci. Route excellente. 2 kilog.\strut &  \\
      \leavevmode\fontsize{6.9pt}{7.9pt}\selectfont\strut Champs tout blancs. Neige durcie ne\strut &  \\
      \leavevmode\fontsize{6.9pt}{7.9pt}\selectfont\strut fond pas.\strut &  \\
      \leavevmode\fontsize{6.9pt}{7.9pt}\selectfont\strut [ — Canteen : Wren's Dubbin : 3 tins : 30 c ]\strut & \smash{\parbox[t]{39mm}{\fontsize{6.2pt}{7.0pt}\selectfont\raggedright\color{black!88}\textbf{14} \textit{in marg. cum unco} (comptes). \textbf{14 Canteen} \textit{angl. in ms.}}} \\
      \leavevmode\llap{\fontsize{6.4pt}{7.3pt}\selectfont\color{darkgray}15\hspace{1.8mm}}\fontsize{6.9pt}{7.9pt}\selectfont\strut Après-midi, liberté. Commencé préparatifs\strut &  \\
      \leavevmode\fontsize{6.9pt}{7.9pt}\selectfont\strut départ.\strut &  \\
      \cline{1-1}
      \leavevmode\fontsize{6.9pt}{7.9pt}\selectfont\strut \underline{23 Mardi} | Hier \& auj'd'hui, liberté entière.\strut & \smash{\parbox[t]{39mm}{\fontsize{6.2pt}{7.0pt}\selectfont\raggedright\color{black!88}\textbf{17 auj'd'hui} \textit{abbr. pro} aujourd'hui.}} \\
      \leavevmode\fontsize{6.9pt}{7.9pt}\selectfont\strut Hier, Capne W., à Amiens — seul.\strut &  \\
      \leavevmode\fontsize{6.9pt}{7.9pt}\selectfont\strut Tempérre très basse, nuit \& matinée, gelé très\strut &  \\
      \leavevmode\llap{\fontsize{6.4pt}{7.3pt}\selectfont\color{darkgray}20\hspace{1.8mm}}\fontsize{6.9pt}{7.9pt}\selectfont\strut fort. Vitres, complèt mates de glace.\strut &  \\
      \leavevmode\fontsize{6.9pt}{7.9pt}\selectfont\strut Neige ne fond pas. 1/2 jour, très beau\strut & \smash{\parbox[t]{39mm}{\fontsize{6.2pt}{7.0pt}\selectfont\raggedright\color{black!88}\textbf{21 1/2} \textit{scil.} demi / moitié.}} \\
      \leavevmode\fontsize{6.9pt}{7.9pt}\selectfont\strut soleil. Neige étincelle sous ses rayons.\strut &  \\
\end{tabular}}%
\par\vspace*{\fill}
% --- Registres de bas de page (1. Critique interne [si pertinent] ; 2. Loci paralleli en latin ; 3. Commentaire historique en retrait) ---
\noindent\rule{\textwidth}{0.5pt}\par\vspace{0.65mm}
{\fontsize{6.6pt}{7.7pt}\selectfont%
\noindent\normalfont\color{black!92} \textbf{9} (\textit{Louvencourt}) cf. supra f\textsuperscript{o}~49~v\textsuperscript{o} (p.~\pageref{fol:99}), f\textsuperscript{o}~61~v\textsuperscript{o} (p.~\pageref{fol:123}), infra f\textsuperscript{o}~73~r\textsuperscript{o} (p.~\pageref{fol:146}), \textit{et passim}. — \textbf{14} (\textit{Canteen}) cf. supra f\textsuperscript{o}~21~v\textsuperscript{o} (p.~\pageref{fol:43}), f\textsuperscript{o}~59~r\textsuperscript{o} (p.~\pageref{fol:118}), infra f\textsuperscript{o}~83~v\textsuperscript{o} (p.~\pageref{fol:167}), f\textsuperscript{o}~89~v\textsuperscript{o} (p.~\pageref{fol:179}), \textit{et passim}. — \textbf{18} (\textit{Amiens}) cf. supra f\textsuperscript{o}~25~v\textsuperscript{o} (p.~\pageref{fol:51}), f\textsuperscript{o}~59~v\textsuperscript{o} (p.~\pageref{fol:119}), infra f\textsuperscript{o}~64~v\textsuperscript{o} (p.~\pageref{fol:129}), f\textsuperscript{o}~65~r\textsuperscript{o} (p.~\pageref{fol:130}), \textit{et passim}.\par\vspace{0.55mm}
\noindent\rule{38mm}{0.4pt}\par\vspace{0.5mm}
\noindent\hspace*{5mm}\parbox{\dimexpr\textwidth-5mm\relax}{\fontsize{6.6pt}{7.7pt}\selectfont\color{black!84}La vague de froid exceptionnelle de fin janvier--février 1917 (l'un des hivers les plus rigoureux du XX\textsuperscript{e}~s. sur le front occidental) durcit la neige et couvre les vitres de givre.}\par}
\newpage
\phantomsection\label{fac:126}
% --- En-tête du fac-similé (Page impaire / droite : Photographie du manuscrit) ---
\noindent\rule{\textwidth}{0.5pt}\par\vspace{0.55mm}
\noindent\textbf{\fontsize{6.8pt}{7.8pt}\selectfont FAC-SIMILÉ EN VIS-À-VIS — FOLIO 63 RECTO} \hfill \textbf{\fontsize{6.8pt}{7.8pt}\selectfont JANVIER 1917}\par\vspace{0.25mm}
\noindent{\fontsize{6.8pt}{8.0pt}\selectfont\textit{Livre-Journal d'Hippolyte Combes} — 20 – 23 janvier 1917 (Acheux, Louvencourt (Somme))}\par\vspace{0.55mm}
\noindent\rule{\textwidth}{0.35pt}\par\vspace*{\fill}
{\centering
  \setlength{\fboxsep}{1.2mm}%
  \setlength{\fboxrule}{0.4pt}%
  \fcolorbox{gray!55}{white}{%
    \includegraphics[width=0.95\textwidth,height=0.81\textheight,keepaspectratio]{\facsimiledir/20251126_110011.jpg}%
  }\par}
\vspace*{\fill}
\newpage

% =============================================================================
% VIS-A-VIS SPREAD : FOLIO 63 VERSO (20251126_110024.jpg)
% =============================================================================
\cleardoubleevenpage
\renewcommand{\currentmonthyear}{Janvier 1917}
\renewcommand{\currentfoliolabel}{Folio 63 verso}
\phantomsection\label{fol:127}
% --- En-tête du folio (Page paire / gauche : Transcription, apparat & registres de bas de page) ---
\noindent\rule{\textwidth}{0.5pt}\par\vspace{0.55mm}
\noindent\textbf{\footnotesize FOLIO 63 VERSO} \hfill \textbf{\footnotesize JANVIER 1917}\par\vspace{0.25mm}
\noindent{\fontsize{6.9pt}{8.0pt}\selectfont Conférence sur le Canada au Y.M.C.A., eau gelée dans la chambre et fin des préparatifs}\par\vspace{0.2mm}
\noindent{\fontsize{6.9pt}{8.0pt}\selectfont\textit{23 (soir) – 25 janvier 1917} \hfill \textit{Acheux-en-Amiénois (Somme)}}\par\vspace{0.5mm}
\noindent\rule{\textwidth}{0.35pt}\par\vspace{0.8mm}
% --- Corps de la transcription diplomatique et apparat marginal en prolongement ---
{\renewcommand{\arraystretch}{0.93}%
\noindent\hspace*{3.5mm}\begin{tabular}{@{}p{87mm}@{\hspace{3.5mm}}p{39mm}@{}}
      \leavevmode\fontsize{7.3pt}{8.4pt}\selectfont\strut \underline{23 Ma.} \hfill{} \underline{Janvier /17}\strut & \smash{\parbox[t]{39mm}{\fontsize{6.2pt}{7.0pt}\selectfont\raggedright\color{black!88}\textbf{1 Ma.} \textit{scil.} Mardi.}} \\
      \leavevmode\fontsize{7.3pt}{8.4pt}\selectfont\strut Soir, Y.M.C.A. hut,\strut & \smash{\parbox[t]{39mm}{\fontsize{6.2pt}{7.0pt}\selectfont\raggedright\color{black!88}\textbf{2 Y.M.C.A.} \textit{angl.} Young Men's Christian Assoc. — \textbf{3 very} \textit{angl. in ms.} — \textbf{4 Mer.} \textit{scil.} Mercredi. \textbf{4 extt} \textit{abbr. pro} extrêmement. — \textbf{6 tte} \textit{abbr. pro} toute.}} \\
      \leavevmode\fontsize{7.3pt}{8.4pt}\selectfont\strut a very interesting lecture on Canada.\strut &  \\
      \cline{1-1}
      \leavevmode\fontsize{7.3pt}{8.4pt}\selectfont\strut \underline{24 Mer.} Nuit extt froide. Gelé très\strut &  \\
      \leavevmode\llap{\fontsize{6.4pt}{7.3pt}\selectfont\color{darkgray}5\hspace{1.8mm}}\fontsize{7.3pt}{8.4pt}\selectfont\strut fort. Vitres couvertes d'1 épaisse couche\strut &  \\
      \leavevmode\fontsize{7.3pt}{8.4pt}\selectfont\strut de glace tte mate.\strut &  \\
      \leavevmode\fontsize{7.3pt}{8.4pt}\selectfont\strut Au réveil, eau pot \& seau gelée.\strut &  \\
      \leavevmode\fontsize{7.3pt}{8.4pt}\selectfont\strut 1/2 jour, soleil ardent.\strut & \smash{\parbox[t]{39mm}{\fontsize{6.2pt}{7.0pt}\selectfont\raggedright\color{black!88}\textbf{8 1/2} \textit{scil.} demi / moitié.}} \\
      \leavevmode\fontsize{7.3pt}{8.4pt}\selectfont\strut Aujourd'hui encore, liberté complète.\strut &  \\
      \leavevmode\llap{\fontsize{6.4pt}{7.3pt}\selectfont\color{darkgray}10\hspace{1.8mm}}\fontsize{7.3pt}{8.4pt}\selectfont\strut Terminé ts préparatifs départ.\strut & \smash{\parbox[t]{39mm}{\fontsize{6.2pt}{7.0pt}\selectfont\raggedright\color{black!88}\textbf{10 ts} \textit{abbr. pro} tous.}} \\
      \leavevmode\fontsize{7.3pt}{8.4pt}\selectfont\strut [ — Washing : 70 c ]\strut & \smash{\parbox[t]{39mm}{\fontsize{6.2pt}{7.0pt}\selectfont\raggedright\color{black!88}\textbf{11} \textit{in marg. cum unco} (comptes). \textbf{11 Washing} \textit{angl. in ms.}}} \\
      \cline{1-1}
      \leavevmode\fontsize{7.3pt}{8.4pt}\selectfont\strut \underline{25 Jeudi} — Gelé extr. + fort. Vitres\strut &  \\
      \leavevmode\fontsize{7.3pt}{8.4pt}\selectfont\strut couvertes épaisse couche mate pointillée.\strut &  \\
      \leavevmode\fontsize{7.3pt}{8.4pt}\selectfont\strut Milieu jour, très beau soleil, qui\strut &  \\
      \leavevmode\llap{\fontsize{6.4pt}{7.3pt}\selectfont\color{darkgray}15\hspace{1.8mm}}\fontsize{7.3pt}{8.4pt}\selectfont\strut fait étinceler neige durcie champs \& cours.\strut &  \\
      \leavevmode\fontsize{7.3pt}{8.4pt}\selectfont\strut Nuit \& matinée, violent bombardemt.\strut & \smash{\parbox[t]{39mm}{\fontsize{6.2pt}{7.0pt}\selectfont\raggedright\color{black!88}\textbf{16 bombardemt} \textit{abbr. pro} bombardement. — \textbf{17 Ts} \textit{abbr. pro} Tous. — \textbf{18 s.} \textit{abbr. pro} soir. — \textbf{19} \textit{in marg. cum unco} (comptes). \textbf{19 mess} \textit{angl. in ms.}}} \\
      \leavevmode\fontsize{7.3pt}{8.4pt}\selectfont\strut Ts préparatifs terminés et arrangements\strut &  \\
      \leavevmode\fontsize{7.3pt}{8.4pt}\selectfont\strut faits pour partir demain 3h s.\strut &  \\
      \cline{1-1}
      \leavevmode\fontsize{7.3pt}{8.4pt}\selectfont\strut [ — Dépenses mess pour janvier = 79 f 10 ]\strut &  \\
\end{tabular}}%
\par\vspace*{\fill}
% --- Registres de bas de page (1. Critique interne [si pertinent] ; 2. Loci paralleli en latin ; 3. Commentaire historique en retrait) ---
\noindent\rule{\textwidth}{0.5pt}\par\vspace{0.65mm}
{\fontsize{6.6pt}{7.7pt}\selectfont%
\noindent\normalfont\color{black!92} \textbf{1} (\textit{de eodem commorationis loco}) cf. supra f\textsuperscript{o}~63~r\textsuperscript{o} (p.~\pageref{fol:126}).\par\vspace{0.55mm}
\noindent\rule{38mm}{0.4pt}\par\vspace{0.5mm}
\noindent\hspace*{5mm}\parbox{\dimexpr\textwidth-5mm\relax}{\fontsize{6.6pt}{7.7pt}\selectfont\color{black!84}Après une conférence en anglais sur le Canada au foyer de la \textbf{\textit{Y.M.C.A.}} d'Acheux le 23 janvier au soir, le gel nocturne des 24 et 25 janvier 1917 atteint une telle intensité que l'eau gèle dans le pot à eau et le seau de sa chambre. À la veille de partir en permission, Hippolyte Combes solde intégralement sa part de frais de \textit{mess} pour les vingt-cinq premiers jours de janvier (79~fr.~10, soit 3~fr.~16 par jour).}\par}
\newpage
\phantomsection\label{fac:127}
% --- En-tête du fac-similé (Page impaire / droite : Photographie du manuscrit) ---
\noindent\rule{\textwidth}{0.5pt}\par\vspace{0.55mm}
\noindent\textbf{\fontsize{6.8pt}{7.8pt}\selectfont FAC-SIMILÉ EN VIS-À-VIS — FOLIO 63 VERSO} \hfill \textbf{\fontsize{6.8pt}{7.8pt}\selectfont JANVIER 1917}\par\vspace{0.25mm}
\noindent{\fontsize{6.8pt}{8.0pt}\selectfont\textit{Livre-Journal d'Hippolyte Combes} — 23 (soir) – 25 janvier 1917 (Acheux-en-Amiénois (Somme))}\par\vspace{0.55mm}
\noindent\rule{\textwidth}{0.35pt}\par\vspace*{\fill}
{\centering
  \setlength{\fboxsep}{1.2mm}%
  \setlength{\fboxrule}{0.4pt}%
  \fcolorbox{gray!55}{white}{%
    \includegraphics[width=0.95\textwidth,height=0.81\textheight,keepaspectratio]{\facsimiledir/20251126_110024.jpg}%
  }\par}
\vspace*{\fill}
\newpage

% =============================================================================
% CHAPITRE VI (FOLIO 64 RECTO A FOLIO 73 RECTO)
% =============================================================================
\cleardoubleevenpage
\renewcommand{\currentmonthyear}{Janvier 1917 – Février 1917}
\renewcommand{\currentfoliolabel}{Chapitre VI}
\chapter{La deuxième permission : Neuilly, Toulouse et Millau}

\noindent{\fontsize{7.8pt}{9.0pt}\selectfont\textit{Folios 64 recto à 73 recto — Étapes par Amiens et Neuilly-sur-Seine, séjour à Toulouse et pèlerinage familial à Millau (26 janvier – 7 février 1917)}}\par\vspace{1.5mm}
\noindent\rule{\textwidth}{0.5pt}\par\vspace{0.9mm}
\noindent\textbf{\scriptsize\textsc{Résumé analytique du Chapitre VI (Folio 64 recto à Folio 73 recto)}}\par\vspace{0.5mm}
\noindent{\fontsize{6.9pt}{8.1pt}\selectfont Le sixième chapitre (f\textsuperscript{o}~64~r\textsuperscript{o} à f\textsuperscript{o}~73~r\textsuperscript{o}) retrace la permission d'hiver d'Hippolyte Combes, du 26 janvier au 7 février 1917, par un froid exceptionnel. Parti d'Acheux le 26 janvier par Amiens, il fait escale le samedi 27 janvier à Neuilly-sur-Seine chez Mme Faure (où il s'entretient avec ses cousines Germaine et Marguerite et avec Marcel Faure en permission) avant de gagner Toulouse le dimanche 28 janvier. Durant six jours à Toulouse (28 janvier – 3 février), il retrouve Marguerite Calmels et sa sœur aînée Madeleine Calmels-Vincens, achète chez l'horloger Pfister la montre-bracelet en or et les bagues de fiançailles, visite avec elles la basilique Saint-Sernin et ses cryptes, les grands magasins du Capitole et la Maison Universelle d'Antoine Labit, assiste au cours public d'Ernest Zyromski sur Maurice de Guérin à la Faculté des lettres, parcourt l'Hôtel d'Assézat, la Daurade et la Dalbade et marche jusqu'à Lardenne. Le samedi 3 février, Hippolyte Combes et Marguerite Calmels voyagent ensemble en train vers Millau en saluant l'abbé Jean Crouzet à Carcassonne et en visitant la cathédrale Saint-Nazaire de Béziers. Accueilli au 10, rue Alfred-Guibert par la famille Calmels, Hippolyte Combes rend visite aux familles Vincens et Layrals (l'oncle Charles Layrals et la tante Thérèse), revoit la sœur de Marguerite, « Mimi », et se recueille au cimetière de Millau sur les tombes des siens avant un long voyage de retour à travers la Lozère, le Cantal et l'Auvergne enneigés (Saint-Chély-d'Apcher, Saint-Flour, Brioude, Clermont-Ferrand, Paris).}\par\vspace{1.0mm}
\noindent\rule{\textwidth}{0.35pt}\par\vspace{1.4mm}

\noindent\textbf{\scriptsize\textsc{Introduction historique et philologique}}\par\vspace{0.6mm}
{\fontsize{7.1pt}{8.3pt}\selectfont Ce deuxième séjour dans le Midi unit pour la première fois dans le carnet les deux villes maîtresses de la vie d'Hippolyte Combes : \textbf{Toulouse}, ville de ses études universitaires et de la famille Crouzet, et \textbf{Millau}, berceau natal des familles Combes, Alric, Calmels et Cartairade, où sera célébré son mariage avec Marguerite Calmels le 25 octobre 1917. La collation directe du manuscrit avec les archives familiales permet d'identifier avec certitude tous les parents visités à Millau les 4 et 5 février 1917 : au 10, rue Alfred-Guibert, les parents de Marguerite (le \textbf{D\textsuperscript{r} Gabriel Calmels} et \textbf{Valentine Cartairade}), son frère \textbf{Pierre Calmels}, ses sœurs \textbf{Madeleine Calmels} (épouse du \textbf{D\textsuperscript{r} Fernand Vincens}, mère de \textbf{Germaine} et \textbf{Robert Vincens}) et \textbf{Anne-Marie dite « Mimi » Calmels} (1888–1975, future épouse de l'avoué \textbf{Joseph Delcamps}), ainsi que son oncle et sa tante maternels \textbf{M\textsuperscript{e} Charles Layrals}, notaire à Millau, et \textbf{Thérèse Cartairade} (f\textsuperscript{o}~71~r\textsuperscript{o}, l.~6 : \texttt{visite à l'oncle \& tante Thérèse (Layral)}).\par}
\vfill
\newpage
% --- Page de droite d'ouverture du chapitre : Carte vectorielle TikZ & Photographies d'archives ---
\noindent\rule{\textwidth}{0.5pt}\par\vspace{0.6mm}
\noindent\textbf{\fontsize{6.8pt}{7.8pt}\selectfont CARTE D'ITINÉRAIRE DU CHAPITRE VI} \hfill \textbf{\fontsize{6.8pt}{7.8pt}\selectfont JANVIER 1917 – FÉVRIER 1917}\par\vspace{0.25mm}
\noindent{\fontsize{6.7pt}{7.8pt}\selectfont\textit{Déplacements, cantonnements et tournées d'Hippolyte Combes (Folio 64 recto à Folio 73 recto)}}\par\vspace{0.55mm}
\noindent\rule{\textwidth}{0.35pt}\par\vspace{1.2mm}
\begin{center}

\begin{tikzpicture}[
  scale=0.96,
  every node/.style={font=\fontsize{6.0pt}{7.0pt}\selectfont},
  loc/.style={circle, fill=black, inner sep=1.1pt},
  step/.style={circle, draw=bordeaux, fill=white, line width=0.75pt, inner sep=1.5pt},
  aller/.style={draw=bordeaux, line width=0.75pt, ->, >=stealth},
  retour/.style={draw=black!75, line width=0.65pt, dashed, ->, >=stealth}
]
\draw[black!50, line width=0.4pt] (-60mm,-42mm) rectangle (60mm,42mm);
\fill[black!3] (-60mm,-42mm) rectangle (60mm,42mm);

% Partie gauche : Boucle ferroviaire Amiens - Paris - Toulouse - Béziers - Millau - Clermont - Paris
\node[font=\fontsize{6.2pt}{7.2pt}\selectfont\bfseries] at (-29mm, 37mm) {A. Boucle de la 2\textsuperscript{e} permission (26 janv. – 7 févr. 1917)};
\node[step] (amiens) at (-30mm, 31mm) {};
\node[right] at (amiens) {\textbf{1. Acheux / Amiens} (26 janv.)};
\node[step] (paris) at (-30mm, 21mm) {};
\node[right] at (paris) {\textbf{2. Paris \& Neuilly} (famille Faure, 27 janv.)};
\node[step] (toulouse) at (-44mm, -22mm) {};
\node[above left] at (toulouse) {\textbf{3. Toulouse} (28 janv. – 3 févr.)};
\node[loc] (carca) at (-28mm, -28mm) {};
\node[below] at (carca) {Carcassonne (Jean Crouzet)};
\node[step] (beziers) at (-10mm, -28mm) {};
\node[below] at (beziers) {\textbf{4. Béziers} (cath. St-Nazaire, 3 févr.)};
\node[step] (millau) at (-12mm, -14mm) {};
\node[right] at (millau) {\textbf{5. Millau} (3–5 févr.)};
\node[loc] (stflour) at (-14mm, -3mm) {};
\node[right] at (stflour) {St-Chély \& St-Flour};
\node[step] (clermont) at (-16mm, 7mm) {};
\node[right] at (clermont) {\textbf{6. Brioude \& Clermont-F\textsuperscript{d}} (6 févr.)};

\draw[aller] (amiens) -- (paris);
\draw[aller] (paris) to[bend right=18] node[pos=0.5, left]{\textit{Aller (27–28 janv.)}} (toulouse);
\draw[aller] (toulouse) -- (carca) -- (beziers) -- (millau);
\draw[retour] (millau) -- (stflour) -- (clermont) -- (paris) -- (amiens);

% Filet vertical
\draw[black!40, line width=0.4pt] (3mm,-38mm) -- (3mm,38mm);

% Partie droite : Plan vectoriel de Millau (4mm à 58mm)
\node[font=\fontsize{6.2pt}{7.2pt}\selectfont\bfseries] at (31mm, 37mm) {B. Lieux familiaux à Millau (3–5 févr. 1917)};
% Le Tarn et la Dourbie
\draw[black!35, line width=1.6pt] (8mm,-32mm) .. controls (22mm,-18mm) and (36mm,-10mm) .. (54mm,8mm)
  node[pos=0.75, below right, font=\fontsize{5.6pt}{6.5pt}\selectfont\itshape]{Le Tarn};
\draw[black!35, line width=1.1pt] (36mm,-10mm) -- (55mm,-24mm)
  node[pos=0.7, below left, font=\fontsize{5.5pt}{6.5pt}\selectfont\itshape]{La Dourbie};

\node[loc] (gare_m) at (16mm, 18mm) {};
\node[right] at (gare_m) {Gare de Millau (arrivée 3 févr., 22 h 30)};
\node[step] (guibert) at (22mm, 6mm) {};
\node[right] at (guibert) {\textbf{10, rue Alfred-Guibert} (maison Calmels)};
\node[loc] (barris) at (20mm, -2mm) {};
\node[right] at (barris) {\textbf{Rue des Barris} (maison natale Combes)};
\node[loc] (mandarous) at (24mm, 12mm) {};
\node[right] at (mandarous) {Place du Mandarous \& Bd de l'Ayrolle};
\node[step] (layral_m) at (26mm, -9mm) {};
\node[right] at (layral_m) {\textbf{Étude \& maison Layrals} (oncle \& tante Thérèse)};
\node[loc] (cimetiere) at (18mm, 26mm) {};
\node[right] at (cimetiere) {\textbf{Cimetière de Millau} (tombeaux familiaux, 5 févr.)};
\end{tikzpicture}

\end{center}
\vfill
\noindent\rule{\textwidth}{0.3pt}\par\vspace{1.2mm}
\noindent\begin{center}
  \begin{minipage}[b]{0.50\textwidth}
    \centering
    \setlength{\fboxsep}{0.7mm}\setlength{\fboxrule}{0.4pt}%
    \fcolorbox{gray!55}{white}{\includegraphics[width=0.94\linewidth,height=28mm,keepaspectratio]{illustrations/mariage_millau_1917.jpg}}
    \par\vspace{0.6mm}{\fontsize{5.8pt}{6.7pt}\selectfont Photographie officielle du mariage d'Hippolyte Combes (en uniforme d'interprète militaire) et de Marguerite Calmels à Millau (25 octobre 1917), entourés des familles Combes et Calmels.}
  \end{minipage}\hfill
  \begin{minipage}[b]{0.36\textwidth}
    \centering
    \setlength{\fboxsep}{0.7mm}\setlength{\fboxrule}{0.4pt}%
    \fcolorbox{gray!55}{white}{\includegraphics[width=0.94\linewidth,height=28mm,keepaspectratio]{illustrations/menu_mariage_1917.jpg}}
    \par\vspace{0.6mm}{\fontsize{5.8pt}{6.7pt}\selectfont Menu du repas de noces d'Hippolyte Combes et Marguerite Calmels au 10, rue Alfred-Guibert à Millau (25 octobre 1917). Archives familiales.}
  \end{minipage}
\end{center}
\vspace*{1mm}
\newpage

% =============================================================================
% VIS-A-VIS SPREAD : FOLIO 64 RECTO (20251126_110048.jpg)
% =============================================================================
\cleardoubleevenpage
\renewcommand{\currentmonthyear}{Janvier 1917}
\renewcommand{\currentfoliolabel}{Folio 64 recto}
\phantomsection\label{fol:128}
% --- En-tête du folio (Page paire / gauche : Transcription, apparat & registres de bas de page) ---
\noindent\rule{\textwidth}{0.5pt}\par\vspace{0.55mm}
\noindent\textbf{\footnotesize FOLIO 64 RECTO} \hfill \textbf{\footnotesize JANVIER 1917}\par\vspace{0.25mm}
\noindent{\fontsize{6.9pt}{8.0pt}\selectfont Départ en permission de la gare de Varennes et longue attente nocturne à Gézaincourt}\par\vspace{0.2mm}
\noindent{\fontsize{6.9pt}{8.0pt}\selectfont\textit{26 janvier 1917} \hfill \textit{Acheux, Gare de Varennes, Gézaincourt (Somme)}}\par\vspace{0.5mm}
\noindent\rule{\textwidth}{0.35pt}\par\vspace{0.8mm}
% --- Corps de la transcription diplomatique et apparat marginal en prolongement ---
{\renewcommand{\arraystretch}{0.87}%
\noindent\hspace*{3.5mm}\begin{tabular}{@{}p{87mm}@{\hspace{3.5mm}}p{39mm}@{}}
      \leavevmode\fontsize{6.9pt}{7.9pt}\selectfont\strut \underline{26 Vendredi} — \hfill{} \underline{Janvier /17}\strut &  \\
      \leavevmode\fontsize{6.9pt}{7.9pt}\selectfont\strut Levé 8h 30 m.\strut & \smash{\parbox[t]{39mm}{\fontsize{6.2pt}{7.0pt}\selectfont\raggedright\color{black!88}\textbf{2 m.} \textit{abbr. pro} matin.}} \\
      \leavevmode\fontsize{6.9pt}{7.9pt}\selectfont\strut Nuit extt froide, gelé très fort. 1/2 jour, soleil\strut & \smash{\parbox[t]{39mm}{\fontsize{6.2pt}{7.0pt}\selectfont\raggedright\color{black!88}\textbf{3 extt} \textit{abbr. pro} extrêmement. \textbf{3 1/2} \textit{scil.} demi / moitié.}} \\
      \leavevmode\fontsize{6.9pt}{7.9pt}\selectfont\strut brillé et très vif éclat.\strut &  \\
      \leavevmode\llap{\fontsize{6.4pt}{7.3pt}\selectfont\color{darkgray}5\hspace{1.8mm}}\fontsize{6.9pt}{7.9pt}\selectfont\strut Matinée, achevé ranger mes affaires.\strut &  \\
      \leavevmode\fontsize{6.9pt}{7.9pt}\selectfont\strut A midi, fait prendre à gare 2 sacs par voiture\strut &  \\
      \leavevmode\fontsize{6.9pt}{7.9pt}\selectfont\strut qui fait service poste.\strut &  \\
      \leavevmode\fontsize{6.9pt}{7.9pt}\selectfont\strut Après déjeuner, venu à pied station, vers 2h.\strut &  \\
      \leavevmode\fontsize{6.9pt}{7.9pt}\selectfont\strut Attendu patiemment départ train — installé 2e cl. [addition interlinéaire : compartiment]\strut & \smash{\parbox[t]{39mm}{\fontsize{6.2pt}{7.0pt}\selectfont\raggedright\color{black!88}\textbf{9} \textit{lect. incert. in ms.}}} \\
      \leavevmode\llap{\fontsize{6.4pt}{7.3pt}\selectfont\color{darkgray}10\hspace{1.8mm}}\fontsize{6.9pt}{7.9pt}\selectfont\strut Soleil brûlant s/ les vitres, du côté de Varennes.\strut & \smash{\parbox[t]{39mm}{\fontsize{6.2pt}{7.0pt}\selectfont\raggedright\color{black!88}\textbf{10 s/} \textit{abbr. pro} sur.}} \\
      \leavevmode\fontsize{6.9pt}{7.9pt}\selectfont\strut Champs tt couverts couche neige durcie\strut & \smash{\parbox[t]{39mm}{\fontsize{6.2pt}{7.0pt}\selectfont\raggedright\color{black!88}\textbf{11 tt} \textit{abbr. pro} tout.}} \\
      \leavevmode\fontsize{6.9pt}{7.9pt}\selectfont\strut dont les particules, comme des grains de diamant,\strut &  \\
      \leavevmode\fontsize{6.9pt}{7.9pt}\selectfont\strut étincellent sous les rayons solaires.\strut &  \\
      \leavevmode\fontsize{6.9pt}{7.9pt}\selectfont\strut A la ligne de l'horizon, légère buée\strut &  \\
      \leavevmode\llap{\fontsize{6.4pt}{7.3pt}\selectfont\color{darkgray}15\hspace{1.8mm}}\fontsize{6.9pt}{7.9pt}\selectfont\strut blanchâtre, confondant terre \& ciel.\strut &  \\
      \leavevmode\fontsize{6.9pt}{7.9pt}\selectfont\strut 1 interprète, monté compartt. Entré en\strut & \smash{\parbox[t]{39mm}{\fontsize{6.2pt}{7.0pt}\selectfont\raggedright\color{black!88}\textbf{16 compartt} \textit{abbr. pro} compartiment.}} \\
      \leavevmode\fontsize{6.9pt}{7.9pt}\selectfont\strut conversn \& échangé nos impressions \& jugements\strut &  \\
      \leavevmode\fontsize{6.9pt}{7.9pt}\selectfont\strut sur les English.\strut &  \\
      \leavevmode\fontsize{6.9pt}{7.9pt}\selectfont\strut Enfin, train démarré après 3h s.\strut & \smash{\parbox[t]{39mm}{\fontsize{6.2pt}{7.0pt}\selectfont\raggedright\color{black!88}\textbf{19 s.} \textit{abbr. pro} soir.}} \\
      \leavevmode\llap{\fontsize{6.4pt}{7.3pt}\selectfont\color{darkgray}20\hspace{1.8mm}}\fontsize{6.9pt}{7.9pt}\selectfont\strut A Gézaincourt, à 4h 30.\strut &  \\
      \leavevmode\fontsize{6.9pt}{7.9pt}\selectfont\strut Là, dû attendre correspce, ds salle surchauffée,\strut & \smash{\parbox[t]{39mm}{\fontsize{6.2pt}{7.0pt}\selectfont\raggedright\color{black!88}\textbf{21 ds} \textit{abbr. pro} dans. \textbf{21 correspce} \textit{abbr. pro} correspondance. — \textbf{22 1/2} \textit{scil.} demi / moitié. — \textbf{24 s.} \textit{abbr. pro} soir.}} \\
      \leavevmode\fontsize{6.9pt}{7.9pt}\selectfont\strut au 1/2 de voyrs, soldats permissionnaires ou paysans, foule\strut &  \\
      \leavevmode\fontsize{6.9pt}{7.9pt}\selectfont\strut très bruyante. Train de Doullens, 3 h. de\strut &  \\
      \leavevmode\fontsize{6.9pt}{7.9pt}\selectfont\strut retard. Quitté Gézaincourt, à 10h 30 s.\strut &  \\
\end{tabular}}%
\par\vspace*{\fill}
% --- Registres de bas de page (1. Critique interne [si pertinent] ; 2. Loci paralleli en latin ; 3. Commentaire historique en retrait) ---
\noindent\rule{\textwidth}{0.5pt}\par\vspace{0.65mm}
{\fontsize{6.6pt}{7.7pt}\selectfont%
\noindent\normalfont\color{black!92} \textbf{10} (\textit{Varennes}) cf. supra f\textsuperscript{o}~25~r\textsuperscript{o} (p.~\pageref{fol:50}), f\textsuperscript{o}~60~v\textsuperscript{o} (p.~\pageref{fol:121}), infra f\textsuperscript{o}~88~r\textsuperscript{o} (p.~\pageref{fol:176}), f\textsuperscript{o}~92~r\textsuperscript{o} (p.~\pageref{fol:184}), \textit{et passim}. — \textbf{23} (\textit{Doullens}) cf. supra f\textsuperscript{o}~25~v\textsuperscript{o} (p.~\pageref{fol:51}), f\textsuperscript{o}~62~r\textsuperscript{o} (p.~\pageref{fol:124}), infra f\textsuperscript{o}~73~r\textsuperscript{o} (p.~\pageref{fol:146}), f\textsuperscript{o}~76~r\textsuperscript{o} (p.~\pageref{fol:152}), \textit{et passim}.\par\vspace{0.55mm}
\noindent\rule{38mm}{0.4pt}\par\vspace{0.5mm}
\noindent\hspace*{5mm}\parbox{\dimexpr\textwidth-5mm\relax}{\fontsize{6.6pt}{7.7pt}\selectfont\color{black!84}Départ pour sa deuxième permission le vendredi 26 janvier 1917: embarqué après 15~h à la gare d'\textbf{Acheux-Varennes} sur le chemin de fer économique départemental de la Somme (ligne à voie métrique d'Albert à Doullens), Hippolyte Combes atteint à 16~h~30 la gare de correspondance de \textbf{Gézaincourt} (au sud de Doullens, sur la ligne de Saint-Pol à Amiens), où le retard de trois heures du train venant de Doullens l'oblige à attendre six heures dans une salle surchauffée bondée de permissionnaires avant de repartir à 22~h~30.}\par}
\newpage
\phantomsection\label{fac:128}
% --- En-tête du fac-similé (Page impaire / droite : Photographie du manuscrit) ---
\noindent\rule{\textwidth}{0.5pt}\par\vspace{0.55mm}
\noindent\textbf{\fontsize{6.8pt}{7.8pt}\selectfont FAC-SIMILÉ EN VIS-À-VIS — FOLIO 64 RECTO} \hfill \textbf{\fontsize{6.8pt}{7.8pt}\selectfont JANVIER 1917}\par\vspace{0.25mm}
\noindent{\fontsize{6.8pt}{8.0pt}\selectfont\textit{Livre-Journal d'Hippolyte Combes} — 26 janvier 1917 (Acheux, Gare de Varennes, Gézaincourt (Somme))}\par\vspace{0.55mm}
\noindent\rule{\textwidth}{0.35pt}\par\vspace*{\fill}
{\centering
  \setlength{\fboxsep}{1.2mm}%
  \setlength{\fboxrule}{0.4pt}%
  \fcolorbox{gray!55}{white}{%
    \includegraphics[width=0.95\textwidth,height=0.81\textheight,keepaspectratio]{\facsimiledir/20251126_110048.jpg}%
  }\par}
\vspace*{\fill}
\newpage

% =============================================================================
% VIS-A-VIS SPREAD : FOLIO 64 VERSO (20251126_110059.jpg)
% =============================================================================
\cleardoubleevenpage
\renewcommand{\currentmonthyear}{Janvier 1917}
\renewcommand{\currentfoliolabel}{Folio 64 verso}
\phantomsection\label{fol:129}
% --- En-tête du folio (Page paire / gauche : Transcription, apparat & registres de bas de page) ---
\noindent\rule{\textwidth}{0.5pt}\par\vspace{0.55mm}
\noindent\textbf{\footnotesize FOLIO 64 VERSO} \hfill \textbf{\footnotesize JANVIER 1917}\par\vspace{0.25mm}
\noindent{\fontsize{6.9pt}{8.0pt}\selectfont Arrivée à Amiens à deux heures du matin et état de la caisse de voyage}\par\vspace{0.2mm}
\noindent{\fontsize{6.9pt}{8.0pt}\selectfont\textit{26 janvier 1917 (suite, nuit du 26 au 27)} \hfill \textit{Acheux, Gézaincourt, Amiens (Somme)}}\par\vspace{0.5mm}
\noindent\rule{\textwidth}{0.35pt}\par\vspace{0.8mm}
% --- Corps de la transcription diplomatique et apparat marginal en prolongement ---
{\renewcommand{\arraystretch}{0.87}%
\noindent\hspace*{3.5mm}\begin{tabular}{@{}p{87mm}@{\hspace{3.5mm}}p{39mm}@{}}
      \leavevmode\fontsize{6.9pt}{7.9pt}\selectfont\strut \underline{Ven. 26 (suite)} \hfill{} \raisebox{0.5ex}{\underline{Janvier 1917}}\strut & \smash{\parbox[t]{39mm}{\fontsize{6.2pt}{7.0pt}\selectfont\raggedright\color{black!88}\textbf{1 Ven.} \textit{scil.} Vendredi.}} \\
      \leavevmode\fontsize{6.9pt}{7.9pt}\selectfont\strut Arrivé Amiens vers 2h m.\strut & \smash{\parbox[t]{39mm}{\fontsize{6.2pt}{7.0pt}\selectfont\raggedright\color{black!88}\textbf{2 m.} \textit{abbr. pro} matin.}} \\
      \leavevmode\fontsize{6.9pt}{7.9pt}\selectfont\strut Froid pénétrant — gelé très dur.\strut &  \\
      \leavevmode\fontsize{6.9pt}{7.9pt}\selectfont\strut Au ciel sombre, qq étoiles très brillantes.\strut & \smash{\parbox[t]{39mm}{\fontsize{6.2pt}{7.0pt}\selectfont\raggedright\color{black!88}\textbf{4 qq} \textit{abbr. pro} quelque(s).}} \\
      \cline{1-1}
      \leavevmode\llap{\fontsize{6.4pt}{7.3pt}\selectfont\color{darkgray}5\hspace{1.8mm}}\fontsize{6.9pt}{7.9pt}\selectfont\strut [ — Voiturier, à Acheux : 1 f\strut & \smash{\parbox[t]{39mm}{\fontsize{6.2pt}{7.0pt}\selectfont\raggedright\color{black!88}\textbf{5} \textit{in marg. cum unco} (comptes).}} \\
      \leavevmode\fontsize{6.9pt}{7.9pt}\selectfont\strut (2 sacs à la gare)\strut &  \\
      \leavevmode\fontsize{6.9pt}{7.9pt}\selectfont\strut — A Gézaincourt : enregt bag. : 10 c\strut &  \\
      \leavevmode\fontsize{6.9pt}{7.9pt}\selectfont\strut (27 kg. — 2 sacs)\strut &  \\
      \cline{1-1}
      \leavevmode\fontsize{6.9pt}{7.9pt}\selectfont\strut \underline{En caisse}, ce matin : 319 f 50 c ]\strut &  \\
\end{tabular}}%
\par\vspace*{\fill}
% --- Registres de bas de page (1. Critique interne [si pertinent] ; 2. Loci paralleli en latin ; 3. Commentaire historique en retrait) ---
\noindent\rule{\textwidth}{0.5pt}\par\vspace{0.65mm}
{\fontsize{6.6pt}{7.7pt}\selectfont%
\noindent\normalfont\color{black!92} \textbf{2} (\textit{Amiens}) cf. supra f\textsuperscript{o}~25~v\textsuperscript{o} (p.~\pageref{fol:51}), f\textsuperscript{o}~63~r\textsuperscript{o} (p.~\pageref{fol:126}), infra f\textsuperscript{o}~65~r\textsuperscript{o} (p.~\pageref{fol:130}), f\textsuperscript{o}~72~v\textsuperscript{o} (p.~\pageref{fol:145}), \textit{et passim}. — \textbf{5} (\textit{Acheux}) cf. supra f\textsuperscript{o}~21~v\textsuperscript{o} (p.~\pageref{fol:43}), f\textsuperscript{o}~64~r\textsuperscript{o} (p.~\pageref{fol:128}), infra f\textsuperscript{o}~73~r\textsuperscript{o} (p.~\pageref{fol:146}), f\textsuperscript{o}~73~v\textsuperscript{o} (p.~\pageref{fol:147}), \textit{et passim}.\par\vspace{0.55mm}
\noindent\rule{38mm}{0.4pt}\par\vspace{0.5mm}
\noindent\hspace*{5mm}\parbox{\dimexpr\textwidth-5mm\relax}{\fontsize{6.6pt}{7.7pt}\selectfont\color{black!84}Suite de la nuit du 26 au 27 janvier 1917: parti d'Acheux-Varennes après 15~h, Hippolyte Combes n'atteint \textbf{Amiens} qu'à 2~h du matin. L'encadré financier consigne le poids de ses deux sacs enregistrés à Gézaincourt (27~kg) et son encaisse exacte au début de sa permission (\textbf{319~fr.~50}, qui lui servira à régler ses achats de voyage et de fiançailles à Paris, Toulouse et Millau).}\par}
\newpage
\phantomsection\label{fac:129}
% --- En-tête du fac-similé (Page impaire / droite : Photographie du manuscrit) ---
\noindent\rule{\textwidth}{0.5pt}\par\vspace{0.55mm}
\noindent\textbf{\fontsize{6.8pt}{7.8pt}\selectfont FAC-SIMILÉ EN VIS-À-VIS — FOLIO 64 VERSO} \hfill \textbf{\fontsize{6.8pt}{7.8pt}\selectfont JANVIER 1917}\par\vspace{0.25mm}
\noindent{\fontsize{6.8pt}{8.0pt}\selectfont\textit{Livre-Journal d'Hippolyte Combes} — 26 janvier 1917 (suite, nuit du 26 au 27) (Acheux, Gézaincourt, Amiens (Somme))}\par\vspace{0.55mm}
\noindent\rule{\textwidth}{0.35pt}\par\vspace*{\fill}
{\centering
  \setlength{\fboxsep}{1.2mm}%
  \setlength{\fboxrule}{0.4pt}%
  \fcolorbox{gray!55}{white}{%
    \includegraphics[width=0.95\textwidth,height=0.81\textheight,keepaspectratio]{\facsimiledir/20251126_110059.jpg}%
  }\par}
\vspace*{\fill}
\newpage

% =============================================================================
% VIS-A-VIS SPREAD : FOLIO 65 RECTO (20251126_110124.jpg)
% =============================================================================
\cleardoubleevenpage
\renewcommand{\currentmonthyear}{Janvier 1917}
\renewcommand{\currentfoliolabel}{Folio 65 recto}
\phantomsection\label{fol:130}
% --- En-tête du folio (Page paire / gauche : Transcription, apparat & registres de bas de page) ---
\noindent\rule{\textwidth}{0.5pt}\par\vspace{0.55mm}
\noindent\textbf{\footnotesize FOLIO 65 RECTO} \hfill \textbf{\footnotesize JANVIER 1917}\par\vspace{0.25mm}
\noindent{\fontsize{6.9pt}{8.0pt}\selectfont Nuit glaciale à la gare d'Amiens avec l'interprète Boilly et arrivée à Paris et Neuilly}\par\vspace{0.2mm}
\noindent{\fontsize{6.9pt}{8.0pt}\selectfont\textit{27 janvier 1917 (matin et midi)} \hfill \textit{Amiens, Paris (Gare du Nord, Gare d'Orsay), Neuilly-sur-Seine}}\par\vspace{0.5mm}
\noindent\rule{\textwidth}{0.35pt}\par\vspace{0.8mm}
% --- Corps de la transcription diplomatique et apparat marginal en prolongement ---
{\renewcommand{\arraystretch}{0.87}%
\noindent\hspace*{3.5mm}\begin{tabular}{@{}p{87mm}@{\hspace{3.5mm}}p{39mm}@{}}
      \leavevmode\fontsize{6.9pt}{7.9pt}\selectfont\strut \underline{Sam. 27} — \hfill{} \underline{Janv. 1917}\strut & \smash{\parbox[t]{39mm}{\fontsize{6.2pt}{7.0pt}\selectfont\raggedright\color{black!88}\textbf{1 Sam.} \textit{scil.} Samedi.}} \\
      \leavevmode\fontsize{6.9pt}{7.9pt}\selectfont\strut Train arrivé enfin à\strut &  \\
      \leavevmode\fontsize{6.9pt}{7.9pt}\selectfont\strut Amiens, vers 2h matin.\strut &  \\
      \leavevmode\fontsize{6.9pt}{7.9pt}\selectfont\strut En compagnie de 2 interprètes (dont l'1,\strut &  \\
      \leavevmode\llap{\fontsize{6.4pt}{7.3pt}\selectfont\color{darkgray}5\hspace{1.8mm}}\fontsize{6.9pt}{7.9pt}\selectfont\strut Boilly, est prof. à Pont-à-Mousson)\strut &  \\
      \leavevmode\fontsize{6.9pt}{7.9pt}\selectfont\strut sorti gare, et allé à l'Abri du Soldat,\strut &  \\
      \leavevmode\fontsize{6.9pt}{7.9pt}\selectfont\strut Bvd Alsace-Lorraine. Aucune place.\strut & \smash{\parbox[t]{39mm}{\fontsize{6.2pt}{7.0pt}\selectfont\raggedright\color{black!88}\textbf{7 Bvd} \textit{abbr. pro} Boulevard.}} \\
      \leavevmode\fontsize{6.9pt}{7.9pt}\selectfont\strut Entré à la Y.M.C.A. Hut, à côté gare.\strut & \smash{\parbox[t]{39mm}{\fontsize{6.2pt}{7.0pt}\selectfont\raggedright\color{black!88}\textbf{8 Y.M.C.A.} \textit{angl.} Young Men's Christian Assoc.}} \\
      \leavevmode\fontsize{6.9pt}{7.9pt}\selectfont\strut Puis durant des heures, déambulé\strut &  \\
      \leavevmode\llap{\fontsize{6.4pt}{7.3pt}\selectfont\color{darkgray}10\hspace{1.8mm}}\fontsize{6.9pt}{7.9pt}\selectfont\strut le long quai gare, en attendant train\strut &  \\
      \leavevmode\fontsize{6.9pt}{7.9pt}\selectfont\strut de Boulogne-Paris. Froid très vif.\strut &  \\
      \leavevmode\fontsize{6.9pt}{7.9pt}\selectfont\strut Enfin vers 5h 1/2 (2 h. de retard) train arrivé. Monté\strut & \smash{\parbox[t]{39mm}{\fontsize{6.2pt}{7.0pt}\selectfont\raggedright\color{black!88}\textbf{12 1/2} \textit{scil.} demi / moitié.}} \\
      \leavevmode\fontsize{6.9pt}{7.9pt}\selectfont\strut compartiment 2e cl., très confortable, mais\strut &  \\
      \leavevmode\fontsize{6.9pt}{7.9pt}\selectfont\strut y air brûlant, irrespirable.\strut &  \\
      \leavevmode\llap{\fontsize{6.4pt}{7.3pt}\selectfont\color{darkgray}15\hspace{1.8mm}}\fontsize{6.9pt}{7.9pt}\selectfont\strut Pu m'assoupir 1 peu.\strut &  \\
      \leavevmode\fontsize{6.9pt}{7.9pt}\selectfont\strut Paris, gare du Nord, vers 11h 15 m.\strut & \smash{\parbox[t]{39mm}{\fontsize{6.2pt}{7.0pt}\selectfont\raggedright\color{black!88}\textbf{16 m.} \textit{abbr. pro} matin.}} \\
      \leavevmode\fontsize{6.9pt}{7.9pt}\selectfont\strut Pris 1 fiacre et venu gare d'Orsay\strut &  \\
      \leavevmode\fontsize{6.9pt}{7.9pt}\selectfont\strut où billet et fait, en 10 minutes, enregistrer mes bagages.\strut &  \\
      \leavevmode\fontsize{6.9pt}{7.9pt}\selectfont\strut Me suis rendu par le métro à\strut &  \\
      \leavevmode\llap{\fontsize{6.4pt}{7.3pt}\selectfont\color{darkgray}20\hspace{1.8mm}}\fontsize{6.9pt}{7.9pt}\selectfont\strut Neuilly : Tuileries – Porte-Maillot.\strut &  \\
      \leavevmode\fontsize{6.9pt}{7.9pt}\selectfont\strut Après 5 m. de marche, arrivé\strut &  \\
      \leavevmode\fontsize{6.9pt}{7.9pt}\selectfont\strut rue d'Armenonville, no [?]6.\strut & \smash{\parbox[t]{39mm}{\fontsize{6.2pt}{7.0pt}\selectfont\raggedright\color{black!88}\textbf{22 no} \textit{lect. incert. in ms.}}} \\
      \leavevmode\fontsize{6.9pt}{7.9pt}\selectfont\strut Maison luxueuse. Au 4e étage, à g.\strut &  \\
\end{tabular}}%
\par\vspace*{\fill}
% --- Registres de bas de page (1. Critique interne [si pertinent] ; 2. Loci paralleli en latin ; 3. Commentaire historique en retrait) ---
\noindent\rule{\textwidth}{0.5pt}\par\vspace{0.65mm}
{\fontsize{6.6pt}{7.7pt}\selectfont%
\noindent\normalfont\color{black!92} \textbf{3} (\textit{Amiens}) cf. supra f\textsuperscript{o}~25~v\textsuperscript{o} (p.~\pageref{fol:51}), f\textsuperscript{o}~64~v\textsuperscript{o} (p.~\pageref{fol:129}), infra f\textsuperscript{o}~72~v\textsuperscript{o} (p.~\pageref{fol:145}), f\textsuperscript{o}~73~r\textsuperscript{o} (p.~\pageref{fol:146}), \textit{et passim}. — \textbf{11} (\textit{Paris}) cf. supra in tegumento (p.~\pageref{fol:1}), f\textsuperscript{o}~61~r\textsuperscript{o} (p.~\pageref{fol:122}), infra f\textsuperscript{o}~65~v\textsuperscript{o} (p.~\pageref{fol:131}), f\textsuperscript{o}~66~r\textsuperscript{o} (p.~\pageref{fol:132}), \textit{et passim}. — \textbf{20} (\textit{Neuilly}) cf. infra f\textsuperscript{o}~65~v\textsuperscript{o} (p.~\pageref{fol:131}).\par\vspace{0.55mm}
\noindent\rule{38mm}{0.4pt}\par\vspace{0.5mm}
\noindent\hspace*{5mm}\parbox{\dimexpr\textwidth-5mm\relax}{\fontsize{6.6pt}{7.7pt}\selectfont\color{black!84}En escale à Paris le samedi 27 janvier 1917 au matin, Hippolyte Combes se rend à \textbf{Neuilly-sur-Seine} pour rendre visite à la famille de \textbf{Mme Faure} (cousines de Marguerite Calmels).}\par}
\newpage
\phantomsection\label{fac:130}
% --- En-tête du fac-similé (Page impaire / droite : Photographie du manuscrit) ---
\noindent\rule{\textwidth}{0.5pt}\par\vspace{0.55mm}
\noindent\textbf{\fontsize{6.8pt}{7.8pt}\selectfont FAC-SIMILÉ EN VIS-À-VIS — FOLIO 65 RECTO} \hfill \textbf{\fontsize{6.8pt}{7.8pt}\selectfont JANVIER 1917}\par\vspace{0.25mm}
\noindent{\fontsize{6.8pt}{8.0pt}\selectfont\textit{Livre-Journal d'Hippolyte Combes} — 27 janvier 1917 (matin et midi) (Amiens, Paris (Gare du Nord, Gare d'Orsay), Neuilly-sur-Seine)}\par\vspace{0.55mm}
\noindent\rule{\textwidth}{0.35pt}\par\vspace*{\fill}
{\centering
  \setlength{\fboxsep}{1.2mm}%
  \setlength{\fboxrule}{0.4pt}%
  \fcolorbox{gray!55}{white}{%
    \includegraphics[width=0.95\textwidth,height=0.81\textheight,keepaspectratio]{\facsimiledir/20251126_110124.jpg}%
  }\par}
\vspace*{\fill}
\newpage

% =============================================================================
% VIS-A-VIS SPREAD : FOLIO 65 VERSO (20251126_110135.jpg)
% =============================================================================
\cleardoubleevenpage
\renewcommand{\currentmonthyear}{Janvier 1917}
\renewcommand{\currentfoliolabel}{Folio 65 verso}
\phantomsection\label{fol:131}
% --- En-tête du folio (Page paire / gauche : Transcription, apparat & registres de bas de page) ---
\noindent\rule{\textwidth}{0.5pt}\par\vspace{0.55mm}
\noindent\textbf{\footnotesize FOLIO 65 VERSO} \hfill \textbf{\footnotesize JANVIER 1917}\par\vspace{0.25mm}
\noindent{\fontsize{6.9pt}{8.0pt}\selectfont Déjeuner à Neuilly chez Mme Faure, visite chez l'opticien Imbert et départ pour Toulouse}\par\vspace{0.2mm}
\noindent{\fontsize{6.9pt}{8.0pt}\selectfont\textit{27 janvier 1917 (suite)} \hfill \textit{Neuilly-sur-Seine, Paris (Boulevard Saint-Michel, Hôtel de Ville, Gare d'Orsay)}}\par\vspace{0.5mm}
\noindent\rule{\textwidth}{0.35pt}\par\vspace{0.8mm}
% --- Corps de la transcription diplomatique et apparat marginal en prolongement ---
{\renewcommand{\arraystretch}{0.93}%
\noindent\hspace*{3.5mm}\begin{tabular}{@{}p{87mm}@{\hspace{3.5mm}}p{39mm}@{}}
      \leavevmode\fontsize{7.3pt}{8.4pt}\selectfont\strut \underline{Sam. 27 (suite)} \hfill{} \raisebox{0.5ex}{\underline{Janv. 1917}}\strut & \smash{\parbox[t]{39mm}{\fontsize{6.2pt}{7.0pt}\selectfont\raggedright\color{black!88}\textbf{1 Sam.} \textit{scil.} Samedi.}} \\
      \leavevmode\fontsize{7.3pt}{8.4pt}\selectfont\strut au 4e ét., à g., appartt de Mme Faure.\strut & \smash{\parbox[t]{39mm}{\fontsize{6.2pt}{7.0pt}\selectfont\raggedright\color{black!88}\textbf{2 appartt} \textit{abbr. pro} appartement. \textbf{2 ét.} \textit{abbr. pro} étage.}} \\
      \leavevmode\fontsize{7.3pt}{8.4pt}\selectfont\strut Y arrivé, vers midi 30 —\strut &  \\
      \leavevmode\fontsize{7.3pt}{8.4pt}\selectfont\strut Accueil très affectueux.\strut &  \\
      \leavevmode\llap{\fontsize{6.4pt}{7.3pt}\selectfont\color{darkgray}5\hspace{1.8mm}}\fontsize{7.3pt}{8.4pt}\selectfont\strut Déjeuné en famille. Mlle Germaine \& Marcel\strut & \smash{\parbox[t]{39mm}{\fontsize{6.2pt}{7.0pt}\selectfont\raggedright\color{black!88}\textbf{5 Mlle} \textit{abbr. pro} Mademoiselle.}} \\
      \leavevmode\fontsize{7.3pt}{8.4pt}\selectfont\strut présents. Les 2 autres fils, Henri, aspirant, et André,\strut &  \\
      \leavevmode\fontsize{7.3pt}{8.4pt}\selectfont\strut préparant Polyte, absents.\strut &  \\
      \leavevmode\fontsize{7.3pt}{8.4pt}\selectfont\strut Causé longuement, après repas.\strut &  \\
      \leavevmode\fontsize{7.3pt}{8.4pt}\selectfont\strut Vers 4h s., pris congé.\strut & \smash{\parbox[t]{39mm}{\fontsize{6.2pt}{7.0pt}\selectfont\raggedright\color{black!88}\textbf{9 s.} \textit{abbr. pro} soir.}} \\
      \leavevmode\llap{\fontsize{6.4pt}{7.3pt}\selectfont\color{darkgray}10\hspace{1.8mm}}\fontsize{7.3pt}{8.4pt}\selectfont\strut Par métro, venu sur Bvd St-Michel, ch. Imbert,\strut & \smash{\parbox[t]{39mm}{\fontsize{6.2pt}{7.0pt}\selectfont\raggedright\color{black!88}\textbf{10 ch.} \textit{abbr. pro} chez. \textbf{10 Bvd} \textit{abbr. pro} Boulevard.}} \\
      \leavevmode\fontsize{7.3pt}{8.4pt}\selectfont\strut mobilisé. Vu son remplaçant.\strut &  \\
      \leavevmode\fontsize{7.3pt}{8.4pt}\selectfont\strut Magasins de l'Hôtel de Ville, plusrs achats.\strut & \smash{\parbox[t]{39mm}{\fontsize{6.2pt}{7.0pt}\selectfont\raggedright\color{black!88}\textbf{12 plusrs} \textit{abbr. pro} plusieurs.}} \\
      \leavevmode\fontsize{7.3pt}{8.4pt}\selectfont\strut Remonté rue de Rivoli. Dîné restaurant \&\strut &  \\
      \leavevmode\fontsize{7.3pt}{8.4pt}\selectfont\strut à 7h s., gare d'Orsay.\strut & \smash{\parbox[t]{39mm}{\fontsize{6.2pt}{7.0pt}\selectfont\raggedright\color{black!88}\textbf{14 s.} \textit{abbr. pro} soir.}} \\
      \leavevmode\llap{\fontsize{6.4pt}{7.3pt}\selectfont\color{darkgray}15\hspace{1.8mm}}\fontsize{7.3pt}{8.4pt}\selectfont\strut Descendu quais. Installé compartt bien\strut & \smash{\parbox[t]{39mm}{\fontsize{6.2pt}{7.0pt}\selectfont\raggedright\color{black!88}\textbf{15 compartt} \textit{abbr. pro} compartiment.}} \\
      \leavevmode\fontsize{7.3pt}{8.4pt}\selectfont\strut chauffé.\strut &  \\
      \leavevmode\fontsize{7.3pt}{8.4pt}\selectfont\strut Départ train vers 7h 50 s. Qq soldats —\strut & \smash{\parbox[t]{39mm}{\fontsize{6.2pt}{7.0pt}\selectfont\raggedright\color{black!88}\textbf{17 Qq} \textit{abbr. pro} Quelque(s). \textbf{17 s.} \textit{abbr. pro} soir.}} \\
      \leavevmode\fontsize{7.3pt}{8.4pt}\selectfont\strut Air très chaud.\strut &  \\
      \leavevmode\fontsize{7.3pt}{8.4pt}\selectfont\strut Peu à peu, me suis assoupi.\strut &  \\
\end{tabular}}%
\par\vspace*{\fill}
% --- Registres de bas de page (1. Critique interne [si pertinent] ; 2. Loci paralleli en latin ; 3. Commentaire historique en retrait) ---
\noindent\rule{\textwidth}{0.5pt}\par\vspace{0.65mm}
{\fontsize{6.6pt}{7.7pt}\selectfont%
\noindent\normalfont\color{black!92} \textbf{2} (\textit{Faure}) cf. Appendicem I, doc.~5 (p.~\pageref{annexe:doc5}). — \textbf{5} (\textit{Germaine}) cf. infra f\textsuperscript{o}~70~v\textsuperscript{o} (p.~\pageref{fol:141}), f\textsuperscript{o}~71~v\textsuperscript{o} (p.~\pageref{fol:143}). — \textbf{10} (\textit{St-Michel}) cf. infra f\textsuperscript{o}~66~r\textsuperscript{o} (p.~\pageref{fol:132}).\par\vspace{0.55mm}
\noindent\rule{38mm}{0.4pt}\par\vspace{0.5mm}
\noindent\hspace*{5mm}\parbox{\dimexpr\textwidth-5mm\relax}{\fontsize{6.6pt}{7.7pt}\selectfont\color{black!84}Suite de la visite familiale du 27 janvier 1917 à \textbf{Neuilly-sur-Seine} chez \textbf{Mme Faure}, entourée de sa fille \textbf{Germaine} et de son fils \textbf{Marcel Faure} alors en permission, tandis que ses deux autres fils, \textbf{Henri} (aspirant d'infanterie) et \textbf{André} (candidat à l'École polytechnique), sont évoqués dans la conversation.}\par}
\newpage
\phantomsection\label{fac:131}
% --- En-tête du fac-similé (Page impaire / droite : Photographie du manuscrit) ---
\noindent\rule{\textwidth}{0.5pt}\par\vspace{0.55mm}
\noindent\textbf{\fontsize{6.8pt}{7.8pt}\selectfont FAC-SIMILÉ EN VIS-À-VIS — FOLIO 65 VERSO} \hfill \textbf{\fontsize{6.8pt}{7.8pt}\selectfont JANVIER 1917}\par\vspace{0.25mm}
\noindent{\fontsize{6.8pt}{8.0pt}\selectfont\textit{Livre-Journal d'Hippolyte Combes} — 27 janvier 1917 (suite) (Neuilly-sur-Seine, Paris (Boulevard Saint-Michel, Hôtel de Ville, Gare d'Orsay))}\par\vspace{0.55mm}
\noindent\rule{\textwidth}{0.35pt}\par\vspace*{\fill}
{\centering
  \setlength{\fboxsep}{1.2mm}%
  \setlength{\fboxrule}{0.4pt}%
  \fcolorbox{gray!55}{white}{%
    \includegraphics[width=0.95\textwidth,height=0.81\textheight,keepaspectratio]{\facsimiledir/20251126_110135.jpg}%
  }\par}
\vspace*{\fill}
\newpage

% =============================================================================
% VIS-A-VIS SPREAD : FOLIO 66 RECTO (20251126_110149.jpg)
% =============================================================================
\cleardoubleevenpage
\renewcommand{\currentmonthyear}{Janvier 1917}
\renewcommand{\currentfoliolabel}{Folio 66 recto}
\phantomsection\label{fol:132}
% --- En-tête du folio (Page paire / gauche : Transcription, apparat & registres de bas de page) ---
\noindent\rule{\textwidth}{0.5pt}\par\vspace{0.55mm}
\noindent\textbf{\footnotesize FOLIO 66 RECTO} \hfill \textbf{\footnotesize JANVIER 1917}\par\vspace{0.25mm}
\noindent{\fontsize{6.9pt}{8.0pt}\selectfont Dépenses parisiennes et arrivée matinale à Toulouse : souvenirs des disparus et visite de l'oncle Émile}\par\vspace{0.2mm}
\noindent{\fontsize{6.9pt}{8.0pt}\selectfont\textit{27 (comptes) – 28 janvier 1917 (matin)} \hfill \textit{Paris, Toulouse (Gare Matabiau, Saint-Michel, Sainte-Marie)}}\par\vspace{0.5mm}
\noindent\rule{\textwidth}{0.35pt}\par\vspace{0.8mm}
% --- Corps de la transcription diplomatique et apparat marginal en prolongement ---
{\renewcommand{\arraystretch}{0.87}%
\noindent\hspace*{3.5mm}\begin{tabular}{@{}p{87mm}@{\hspace{3.5mm}}p{39mm}@{}}
      \leavevmode\fontsize{6.9pt}{7.9pt}\selectfont\strut \underline{Sam. 27} \hfill{} \underline{Janv. /17}\strut & \smash{\parbox[t]{39mm}{\fontsize{6.2pt}{7.0pt}\selectfont\raggedright\color{black!88}\textbf{1 Sam.} \textit{scil.} Samedi.}} \\
      \leavevmode\fontsize{6.9pt}{7.9pt}\selectfont\strut Paris. Gare du Nord. Facteur : 20 c\strut &  \\
      \leavevmode\fontsize{6.9pt}{7.9pt}\selectfont\strut Gare d'Orsay : 10 c. | Cartes ples : 10 c\strut & \smash{\parbox[t]{39mm}{\fontsize{6.2pt}{7.0pt}\selectfont\raggedright\color{black!88}\textbf{3 Cartes ples} \textit{abbr. pro} Cartes postales. \textbf{3 c.} \textit{abbr. pro} centimes. — \textbf{4 Enregt} \textit{abbr. pro} Enregistrement.}} \\
      \leavevmode\fontsize{6.9pt}{7.9pt}\selectfont\strut Enregt 2 colis, 2 f 40 : 10 c\strut &  \\
      \leavevmode\llap{\fontsize{6.4pt}{7.3pt}\selectfont\color{darkgray}5\hspace{1.8mm}}\fontsize{6.9pt}{7.9pt}\selectfont\strut Fiacre taximètre : 1 f 45 + 30 c = 1 f 75 c\strut &  \\
      \leavevmode\fontsize{6.9pt}{7.9pt}\selectfont\strut Métro : 15, 15 c | Dîner : 1 f 95 + 10 c\strut &  \\
      \leavevmode\fontsize{6.9pt}{7.9pt}\selectfont\strut Billet pour Toulouse : 11 f 15 c\strut &  \\
      \leavevmode\fontsize{6.9pt}{7.9pt}\selectfont\strut Magasins Hôtel de Ville : 1 brosse à cheveux : 3 f 95\strut &  \\
      \leavevmode\fontsize{6.9pt}{7.9pt}\selectfont\strut Brosse à souliers : 65 c. — Peigne fin : 1 f 25 c\strut & \smash{\parbox[t]{39mm}{\fontsize{6.2pt}{7.0pt}\selectfont\raggedright\color{black!88}\textbf{9 c.} \textit{abbr. pro} centimes.}} \\
      \leavevmode\llap{\fontsize{6.4pt}{7.3pt}\selectfont\color{darkgray}10\hspace{1.8mm}}\fontsize{6.9pt}{7.9pt}\selectfont\strut Peigne ordre : 95 c. — Cirage noir : 50 c\strut & \smash{\parbox[t]{39mm}{\fontsize{6.2pt}{7.0pt}\selectfont\raggedright\color{black!88}\textbf{10 c.} \textit{abbr. pro} centimes.}} \\
      \cline{1-1}
      \leavevmode\fontsize{6.9pt}{7.9pt}\selectfont\strut \underline{Dimanche 28}\strut &  \\
      \leavevmode\fontsize{6.9pt}{7.9pt}\selectfont\strut Voyage, terminé sans incidents.\strut &  \\
      \leavevmode\fontsize{6.9pt}{7.9pt}\selectfont\strut En approchant de Toulouse, souvenirs jeunesse\strut &  \\
      \leavevmode\fontsize{6.9pt}{7.9pt}\selectfont\strut revenus en foule esprit. Cru revoir mère et\strut &  \\
      \leavevmode\llap{\fontsize{6.4pt}{7.3pt}\selectfont\color{darkgray}15\hspace{1.8mm}}\fontsize{6.9pt}{7.9pt}\selectfont\strut sœurs, m'attendant, cme autrefois, gare.\strut & \smash{\parbox[t]{39mm}{\fontsize{6.2pt}{7.0pt}\selectfont\raggedright\color{black!88}\textbf{15 cme} \textit{abbr. pro} comme.}} \\
      \leavevmode\fontsize{6.9pt}{7.9pt}\selectfont\strut Vers 8h 15, arrivée.\strut &  \\
      \leavevmode\fontsize{6.9pt}{7.9pt}\selectfont\strut Retiré 2 bagages. Pris tram St-Michel.\strut & \smash{\parbox[t]{39mm}{\fontsize{6.2pt}{7.0pt}\selectfont\raggedright\color{black!88}\textbf{17 St} \textit{abbr. pro} Saint.}} \\
      \leavevmode\fontsize{6.9pt}{7.9pt}\selectfont\strut Matinée, ds mon appartt, à ranger\strut & \smash{\parbox[t]{39mm}{\fontsize{6.2pt}{7.0pt}\selectfont\raggedright\color{black!88}\textbf{18 ds} \textit{abbr. pro} dans. \textbf{18 appartt} \textit{abbr. pro} appartement. — \textbf{20 1/2} \textit{scil.} demi / moitié. — \textbf{21 Ste} \textit{abbr. pro} Sainte.}} \\
      \leavevmode\fontsize{6.9pt}{7.9pt}\selectfont\strut affaires.\strut &  \\
      \leavevmode\llap{\fontsize{6.4pt}{7.3pt}\selectfont\color{darkgray}20\hspace{1.8mm}}\fontsize{6.9pt}{7.9pt}\selectfont\strut Oncle Émile, venu me surprendre à 10h 1/2.\strut &  \\
      \leavevmode\fontsize{6.9pt}{7.9pt}\selectfont\strut A midi, enfin, à Ste-Marie. Dû\strut &  \\
      \leavevmode\fontsize{6.9pt}{7.9pt}\selectfont\strut attendre tram 30 m. angle musée.\strut &  \\
\end{tabular}}%
\par\vspace*{\fill}
% --- Registres de bas de page (1. Critique interne [si pertinent] ; 2. Loci paralleli en latin ; 3. Commentaire historique en retrait) ---
\noindent\rule{\textwidth}{0.5pt}\par\vspace{0.65mm}
{\fontsize{6.6pt}{7.7pt}\selectfont%
\noindent\normalfont \textbf{Feuillet comptable :} Page entièrement dédiée au relevé des dépenses de voyage des 27 et 28 janvier 1917 (Paris, train d'Orsay et arrivée à Toulouse), réparties par localités soulignées.\par\vspace{0.5mm}
\noindent\rule{22mm}{0.3pt}\par\vspace{0.45mm}
\noindent\normalfont\color{black!92} \textbf{2} (\textit{Paris}) cf. supra in tegumento (p.~\pageref{fol:1}), f\textsuperscript{o}~65~v\textsuperscript{o} (p.~\pageref{fol:131}), infra f\textsuperscript{o}~72~v\textsuperscript{o} (p.~\pageref{fol:145}), f\textsuperscript{o}~74~r\textsuperscript{o} (p.~\pageref{fol:148}), \textit{et passim}, Appendicem I, doc.~5 (p.~\pageref{annexe:doc5}). — \textbf{7} (\textit{Toulouse}) cf. supra in tegumento (p.~\pageref{fol:1}), f\textsuperscript{o}~65~v\textsuperscript{o} (p.~\pageref{fol:131}), infra f\textsuperscript{o}~66~v\textsuperscript{o} (p.~\pageref{fol:133}), f\textsuperscript{o}~67~r\textsuperscript{o} (p.~\pageref{fol:134}), \textit{et passim}. — \textbf{17} (\textit{St-Michel}) cf. supra f\textsuperscript{o}~65~v\textsuperscript{o} (p.~\pageref{fol:131}).\par\vspace{0.55mm}
\noindent\rule{38mm}{0.4pt}\par\vspace{0.5mm}
\noindent\hspace*{5mm}\parbox{\dimexpr\textwidth-5mm\relax}{\fontsize{6.6pt}{7.7pt}\selectfont\color{black!84}En arrivant à la gare Matabiau de \textbf{Toulouse} le dimanche 28 janvier 1917 à 8~h~15, Hippolyte Combes évoque avec émotion sa mère \textbf{Marie Anaïs Alric} et ses sœurs disparues (\textbf{Maria Hippolyte Combes}, 1875--1903, et \textbf{Joséphine Hippolyte Combes}, 1873--1914) qui l'attendaient autrefois à la gare.}\par}
\newpage
\phantomsection\label{fac:132}
% --- En-tête du fac-similé (Page impaire / droite : Photographie du manuscrit) ---
\noindent\rule{\textwidth}{0.5pt}\par\vspace{0.55mm}
\noindent\textbf{\fontsize{6.8pt}{7.8pt}\selectfont FAC-SIMILÉ EN VIS-À-VIS — FOLIO 66 RECTO} \hfill \textbf{\fontsize{6.8pt}{7.8pt}\selectfont JANVIER 1917}\par\vspace{0.25mm}
\noindent{\fontsize{6.8pt}{8.0pt}\selectfont\textit{Livre-Journal d'Hippolyte Combes} — 27 (comptes) – 28 janvier 1917 (matin) (Paris, Toulouse (Gare Matabiau, Saint-Michel, Sainte-Marie))}\par\vspace{0.55mm}
\noindent\rule{\textwidth}{0.35pt}\par\vspace*{\fill}
{\centering
  \setlength{\fboxsep}{1.2mm}%
  \setlength{\fboxrule}{0.4pt}%
  \fcolorbox{gray!55}{white}{%
    \includegraphics[width=0.95\textwidth,height=0.81\textheight,keepaspectratio]{\facsimiledir/20251126_110149.jpg}%
  }\par}
\vspace*{\fill}
\newpage

% =============================================================================
% VIS-A-VIS SPREAD : FOLIO 66 VERSO (20251126_110200.jpg)
% =============================================================================
\cleardoubleevenpage
\renewcommand{\currentmonthyear}{Janvier 1917}
\renewcommand{\currentfoliolabel}{Folio 66 verso}
\phantomsection\label{fol:133}
% --- En-tête du folio (Page paire / gauche : Transcription, apparat & registres de bas de page) ---
\noindent\rule{\textwidth}{0.5pt}\par\vspace{0.55mm}
\noindent\textbf{\footnotesize FOLIO 66 VERSO} \hfill \textbf{\footnotesize JANVIER 1917}\par\vspace{0.25mm}
\noindent{\fontsize{6.9pt}{8.0pt}\selectfont Retrouvailles infiniment tendres avec Marguerite à Sainte-Marie (Toulouse)}\par\vspace{0.2mm}
\noindent{\fontsize{6.9pt}{8.0pt}\selectfont\textit{28 janvier 1917 (suite)} \hfill \textit{Toulouse (Quartier Sainte-Marie / Côte Pavée)}}\par\vspace{0.5mm}
\noindent\rule{\textwidth}{0.35pt}\par\vspace{0.8mm}
% --- Corps de la transcription diplomatique et apparat marginal en prolongement ---
{\renewcommand{\arraystretch}{0.87}%
\noindent\hspace*{3.5mm}\begin{tabular}{@{}p{87mm}@{\hspace{3.5mm}}p{39mm}@{}}
      \leavevmode\fontsize{6.9pt}{7.9pt}\selectfont\strut \underline{Dim. 28. (suite)} — Janv. /17 —\strut & \smash{\parbox[t]{39mm}{\fontsize{6.2pt}{7.0pt}\selectfont\raggedright\color{black!88}\textbf{1 Dim.} \textit{scil.} Dimanche.}} \\
      \leavevmode\fontsize{6.9pt}{7.9pt}\selectfont\strut Margte, sur la porte, m'attendait. Accueil\strut & \smash{\parbox[t]{39mm}{\fontsize{6.2pt}{7.0pt}\selectfont\raggedright\color{black!88}\textbf{2 Margte} \textit{scil.} Marguerite (Calmels).}} \\
      \leavevmode\fontsize{6.9pt}{7.9pt}\selectfont\strut infiniment tendre. S'est montrée charmante.\strut &  \\
      \leavevmode\fontsize{6.9pt}{7.9pt}\selectfont\strut Vu Tante \& Madeleine.\strut &  \\
      \leavevmode\llap{\fontsize{6.4pt}{7.3pt}\selectfont\color{darkgray}5\hspace{1.8mm}}\fontsize{6.9pt}{7.9pt}\selectfont\strut Passé journée entière à Ste-Marie. Longues\strut & \smash{\parbox[t]{39mm}{\fontsize{6.2pt}{7.0pt}\selectfont\raggedright\color{black!88}\textbf{5 Ste} \textit{abbr. pro} Sainte.}} \\
      \leavevmode\fontsize{6.9pt}{7.9pt}\selectfont\strut heures d'intime tête-à-tête av. Margte, dans\strut & \smash{\parbox[t]{39mm}{\fontsize{6.2pt}{7.0pt}\selectfont\raggedright\color{black!88}\textbf{6 Margte} \textit{scil.} Marguerite (Calmels).}} \\
      \leavevmode\fontsize{6.9pt}{7.9pt}\selectfont\strut le petit salon.\strut &  \\
      \leavevmode\fontsize{6.9pt}{7.9pt}\selectfont\strut Jean, monté à 6h 1/2 et dîné avec nous.\strut & \smash{\parbox[t]{39mm}{\fontsize{6.2pt}{7.0pt}\selectfont\raggedright\color{black!88}\textbf{8 1/2} \textit{scil.} demi / moitié.}} \\
      \leavevmode\fontsize{6.9pt}{7.9pt}\selectfont\strut Redescendu seul, à pied, à 9h s. Aucun tram.\strut & \smash{\parbox[t]{39mm}{\fontsize{6.2pt}{7.0pt}\selectfont\raggedright\color{black!88}\textbf{9 s.} \textit{abbr. pro} soir.}} \\
      \leavevmode\llap{\fontsize{6.4pt}{7.3pt}\selectfont\color{darkgray}10\hspace{1.8mm}}\fontsize{6.9pt}{7.9pt}\selectfont\strut Margte, vraiment exquise de simplicité \& de charme,\strut & \smash{\parbox[t]{39mm}{\fontsize{6.2pt}{7.0pt}\selectfont\raggedright\color{black!88}\textbf{10 Margte} \textit{scil.} Marguerite (Calmels).}} \\
      \leavevmode\fontsize{6.9pt}{7.9pt}\selectfont\strut vive, spontanée, aimante, affectueuse au possible.\strut &  \\
      \leavevmode\fontsize{6.9pt}{7.9pt}\selectfont\strut Mes sentiments à son égard, douceur infinie.\strut &  \\
      \leavevmode\fontsize{6.9pt}{7.9pt}\selectfont\strut Causeries délicieuses sur les sujets les plus\strut &  \\
      \leavevmode\fontsize{6.9pt}{7.9pt}\selectfont\strut divers.\strut &  \\
      \cline{1-1}
      \leavevmode\llap{\fontsize{6.4pt}{7.3pt}\selectfont\color{darkgray}15\hspace{1.8mm}}\fontsize{6.9pt}{7.9pt}\selectfont\strut [ — Trams : 10 c, 10 c — ]\strut & \smash{\parbox[t]{39mm}{\fontsize{6.2pt}{7.0pt}\selectfont\raggedright\color{black!88}\textbf{15} \textit{in marg. cum unco} (comptes).}} \\
\end{tabular}}%
\par\vspace*{\fill}
% --- Registres de bas de page (1. Critique interne [si pertinent] ; 2. Loci paralleli en latin ; 3. Commentaire historique en retrait) ---
\noindent\rule{\textwidth}{0.5pt}\par\vspace{0.65mm}
{\fontsize{6.6pt}{7.7pt}\selectfont%
\noindent\normalfont \textbf{Bilan financier :} Second feuillet comptable récapitulant les grands achats de fiançailles à Toulouse le 28 janvier 1917 (montre en or, bagues, chaîne) et l'état de la caisse (\textit{En caisse}).\par\vspace{0.5mm}
\noindent\rule{22mm}{0.3pt}\par\vspace{0.45mm}
\noindent\normalfont\color{black!92} \textbf{2} (\textit{Margte}) cf. supra f\textsuperscript{o}~1~r\textsuperscript{o} (p.~\pageref{fol:2}), f\textsuperscript{o}~57~r\textsuperscript{o} (p.~\pageref{fol:114}), infra f\textsuperscript{o}~67~r\textsuperscript{o} (p.~\pageref{fol:134}), f\textsuperscript{o}~67~v\textsuperscript{o} (p.~\pageref{fol:135}), \textit{et passim}, Appendicem I, doc.~5 (p.~\pageref{annexe:doc5}). — \textbf{4} (\textit{Tante}) cf. supra f\textsuperscript{o}~27~r\textsuperscript{o} (p.~\pageref{fol:54}), f\textsuperscript{o}~66~r\textsuperscript{o} (p.~\pageref{fol:132}), infra f\textsuperscript{o}~69~v\textsuperscript{o} (p.~\pageref{fol:139}), \textit{et passim}. — \textbf{4} (\textit{Madeleine}) cf. infra f\textsuperscript{o}~67~r\textsuperscript{o} (p.~\pageref{fol:134}).\par\vspace{0.55mm}
\noindent\rule{38mm}{0.4pt}\par\vspace{0.5mm}
\noindent\hspace*{5mm}\parbox{\dimexpr\textwidth-5mm\relax}{\fontsize{6.6pt}{7.7pt}\selectfont\color{black!84}Les retrouvailles chaleureuses du dimanche 28 janvier 1917 à Sainte-Marie dissipent entièrement le nuage épistolaire de la Noël (f\textsuperscript{o}~56~r\textsuperscript{o} et f\textsuperscript{o}~57~r\textsuperscript{o}).}\par}
\newpage
\phantomsection\label{fac:133}
% --- En-tête du fac-similé (Page impaire / droite : Photographie du manuscrit) ---
\noindent\rule{\textwidth}{0.5pt}\par\vspace{0.55mm}
\noindent\textbf{\fontsize{6.8pt}{7.8pt}\selectfont FAC-SIMILÉ EN VIS-À-VIS — FOLIO 66 VERSO} \hfill \textbf{\fontsize{6.8pt}{7.8pt}\selectfont JANVIER 1917}\par\vspace{0.25mm}
\noindent{\fontsize{6.8pt}{8.0pt}\selectfont\textit{Livre-Journal d'Hippolyte Combes} — 28 janvier 1917 (suite) (Toulouse (Quartier Sainte-Marie / Côte Pavée))}\par\vspace{0.55mm}
\noindent\rule{\textwidth}{0.35pt}\par\vspace*{\fill}
{\centering
  \setlength{\fboxsep}{1.2mm}%
  \setlength{\fboxrule}{0.4pt}%
  \fcolorbox{gray!55}{white}{%
    \includegraphics[width=0.95\textwidth,height=0.81\textheight,keepaspectratio]{\facsimiledir/20251126_110200.jpg}%
  }\par}
\vspace*{\fill}
\newpage

% =============================================================================
% VIS-A-VIS SPREAD : FOLIO 67 RECTO (20251126_110216.jpg)
% =============================================================================
\cleardoubleevenpage
\renewcommand{\currentmonthyear}{Janvier 1917}
\renewcommand{\currentfoliolabel}{Folio 67 recto}
\phantomsection\label{fol:134}
% --- En-tête du folio (Page paire / gauche : Transcription, apparat & registres de bas de page) ---
\noindent\rule{\textwidth}{0.5pt}\par\vspace{0.55mm}
\noindent\textbf{\footnotesize FOLIO 67 RECTO} \hfill \textbf{\footnotesize JANVIER 1917}\par\vspace{0.25mm}
\noindent{\fontsize{6.9pt}{8.0pt}\selectfont Le panier de fleurs de Nice et le pèlerinage de jeunesse aux cryptes de Saint-Sernin}\par\vspace{0.2mm}
\noindent{\fontsize{6.9pt}{8.0pt}\selectfont\textit{29 janvier 1917} \hfill \textit{Toulouse (Gare Matabiau, Sainte-Marie, Basilique Saint-Sernin, Place du Capitole, Rue d'Alsace, Côte Pavée)}}\par\vspace{0.5mm}
\noindent\rule{\textwidth}{0.35pt}\par\vspace{0.8mm}
% --- Corps de la transcription diplomatique et apparat marginal en prolongement ---
{\renewcommand{\arraystretch}{0.87}%
\noindent\hspace*{3.5mm}\begin{tabular}{@{}p{87mm}@{\hspace{3.5mm}}p{39mm}@{}}
      \leavevmode\fontsize{6.9pt}{7.9pt}\selectfont\strut \underline{Lundi 29.} \hfill{} \underline{— Janvier 1917 —}\strut &  \\
      \leavevmode\fontsize{6.9pt}{7.9pt}\selectfont\strut Levé 7h 1/2.\strut & \smash{\parbox[t]{39mm}{\fontsize{6.2pt}{7.0pt}\selectfont\raggedright\color{black!88}\textbf{2 1/2} \textit{scil.} demi / moitié.}} \\
      \leavevmode\fontsize{6.9pt}{7.9pt}\selectfont\strut Matin, temps 1 peu froid. Journée, soleil pâle.\strut &  \\
      \leavevmode\fontsize{6.9pt}{7.9pt}\selectfont\strut Allumé gd feu cheminée chambre rue. Goûté\strut & \smash{\parbox[t]{39mm}{\fontsize{6.2pt}{7.0pt}\selectfont\raggedright\color{black!88}\textbf{4 gd} \textit{abbr. pro} grand.}} \\
      \leavevmode\llap{\fontsize{6.4pt}{7.3pt}\selectfont\color{darkgray}5\hspace{1.8mm}}\fontsize{6.9pt}{7.9pt}\selectfont\strut confort chambre retrouvée.\strut &  \\
      \leavevmode\fontsize{6.9pt}{7.9pt}\selectfont\strut Venu gare. Me suis fait livrer, non\strut &  \\
      \leavevmode\fontsize{6.9pt}{7.9pt}\selectfont\strut sans peine, panier fleurs Nice. Colis sans\strut &  \\
      \leavevmode\fontsize{6.9pt}{7.9pt}\selectfont\strut étiquette, ouverture et feuille postale disparues.\strut &  \\
      \leavevmode\fontsize{6.9pt}{7.9pt}\selectfont\strut Vers midi, enfin, de nouveau à Ste-Marie.\strut & \smash{\parbox[t]{39mm}{\fontsize{6.2pt}{7.0pt}\selectfont\raggedright\color{black!88}\textbf{9 Ste} \textit{abbr. pro} Sainte.}} \\
      \leavevmode\llap{\fontsize{6.4pt}{7.3pt}\selectfont\color{darkgray}10\hspace{1.8mm}}\fontsize{6.9pt}{7.9pt}\selectfont\strut Après déjeuner, Marg. \& Madne, descendues\strut & \smash{\parbox[t]{39mm}{\fontsize{6.2pt}{7.0pt}\selectfont\raggedright\color{black!88}\textbf{10 Marg.} \textit{scil.} Marguerite (Calmels).}} \\
      \leavevmode\fontsize{6.9pt}{7.9pt}\selectfont\strut en ville avec moi.\strut &  \\
      \leavevmode\fontsize{6.9pt}{7.9pt}\selectfont\strut Visité longuement basilique St-Sernin\strut & \smash{\parbox[t]{39mm}{\fontsize{6.2pt}{7.0pt}\selectfont\raggedright\color{black!88}\textbf{12 St} \textit{abbr. pro} Saint.}} \\
      \leavevmode\fontsize{6.9pt}{7.9pt}\selectfont\strut et surtout les cryptes. Toute ma jeunesse\strut &  \\
      \leavevmode\fontsize{6.9pt}{7.9pt}\selectfont\strut réveillée, dans \& autour de cette admirable\strut &  \\
      \leavevmode\llap{\fontsize{6.4pt}{7.3pt}\selectfont\color{darkgray}15\hspace{1.8mm}}\fontsize{6.9pt}{7.9pt}\selectfont\strut église romane.\strut &  \\
      \leavevmode\fontsize{6.9pt}{7.9pt}\selectfont\strut Magasins du Capitole \& Labit.\strut &  \\
      \leavevmode\fontsize{6.9pt}{7.9pt}\selectfont\strut Remonté 1 peu Rue Alsace.\strut &  \\
      \leavevmode\fontsize{6.9pt}{7.9pt}\selectfont\strut Enfin, allé prendre au Capitole\strut &  \\
      \leavevmode\fontsize{6.9pt}{7.9pt}\selectfont\strut notre tram p. Côte Pavée.\strut &  \\
      \leavevmode\llap{\fontsize{6.4pt}{7.3pt}\selectfont\color{darkgray}20\hspace{1.8mm}}\fontsize{6.9pt}{7.9pt}\selectfont\strut A 9h s., quitté, hélas! Ste-Marie.\strut & \smash{\parbox[t]{39mm}{\fontsize{6.2pt}{7.0pt}\selectfont\raggedright\color{black!88}\textbf{20 Ste} \textit{abbr. pro} Sainte. \textbf{20 s.} \textit{abbr. pro} soir. — \textbf{21} \textit{in marg. cum unco} (comptes). — \textbf{23 Margte} \textit{scil.} Marguerite (Calmels). \textbf{23 Ctes ples} \textit{abbr. pro} Cartes postales. — \textbf{24 St} \textit{abbr. pro} Saint.}} \\
      \cline{1-1}
      \leavevmode\fontsize{6.9pt}{7.9pt}\selectfont\strut [ — Trams : 10, 10, 10 c, 60, 10 c —\strut &  \\
      \leavevmode\fontsize{6.9pt}{7.9pt}\selectfont\strut Mag. du Capitole : boîte de 3 savonnettes : 2 f 90\strut &  \\
      \leavevmode\fontsize{6.9pt}{7.9pt}\selectfont\strut Brosse à dents p. Margte : 2 f 25 c — Ctes ples : 30 c\strut &  \\
      \leavevmode\fontsize{6.9pt}{7.9pt}\selectfont\strut Cartes pour visiter cryptes de St-Sernin : 1 f — ]\strut &  \\
\end{tabular}}%
\par\vspace*{\fill}
% --- Registres de bas de page (1. Critique interne [si pertinent] ; 2. Loci paralleli en latin ; 3. Commentaire historique en retrait) ---
\noindent\rule{\textwidth}{0.5pt}\par\vspace{0.65mm}
{\fontsize{6.6pt}{7.7pt}\selectfont%
\noindent\normalfont\color{black!92} \textbf{9} (\textit{Ste-Marie}) cf. supra f\textsuperscript{o}~27~v\textsuperscript{o} (p.~\pageref{fol:55}), f\textsuperscript{o}~66~v\textsuperscript{o} (p.~\pageref{fol:133}), infra f\textsuperscript{o}~67~v\textsuperscript{o} (p.~\pageref{fol:135}), f\textsuperscript{o}~68~r\textsuperscript{o} (p.~\pageref{fol:136}), \textit{et passim}. — \textbf{10} (\textit{Madne}) cf. supra f\textsuperscript{o}~66~v\textsuperscript{o} (p.~\pageref{fol:133}). — \textbf{23} (\textit{Margte}) cf. supra f\textsuperscript{o}~1~r\textsuperscript{o} (p.~\pageref{fol:2}), f\textsuperscript{o}~66~v\textsuperscript{o} (p.~\pageref{fol:133}), infra f\textsuperscript{o}~67~v\textsuperscript{o} (p.~\pageref{fol:135}), f\textsuperscript{o}~68~r\textsuperscript{o} (p.~\pageref{fol:136}), \textit{et passim}.\par\vspace{0.55mm}
\noindent\rule{38mm}{0.4pt}\par\vspace{0.5mm}
\noindent\hspace*{5mm}\parbox{\dimexpr\textwidth-5mm\relax}{\fontsize{6.6pt}{7.7pt}\selectfont\color{black!84}À \textbf{Toulouse} le 29 janvier 1917, Hippolyte Combes retire en gare Matabiau un panier de fleurs expédié de Nice pour Marguerite Calmels, visite avec elle et sa sœur aînée \textbf{Madeleine Calmels-Vincens} la basilique romane \textbf{Saint-Sernin} et ses cryptes (voisines de l'ancien domicile familial des Combes rue du Taur), puis parcourt les grands magasins du Capitole et la \textbf{Maison Universelle} fondée par \textbf{Antoine Labit} (1832–1912, père du voyageur et collectionneur Georges Labit).}\par}
\newpage
\phantomsection\label{fac:134}
% --- En-tête du fac-similé (Page impaire / droite : Photographie du manuscrit) ---
\noindent\rule{\textwidth}{0.5pt}\par\vspace{0.55mm}
\noindent\textbf{\fontsize{6.8pt}{7.8pt}\selectfont FAC-SIMILÉ EN VIS-À-VIS — FOLIO 67 RECTO} \hfill \textbf{\fontsize{6.8pt}{7.8pt}\selectfont JANVIER 1917}\par\vspace{0.25mm}
\noindent{\fontsize{6.8pt}{8.0pt}\selectfont\textit{Livre-Journal d'Hippolyte Combes} — 29 janvier 1917 (Toulouse (Gare Matabiau, Sainte-Marie, Basilique Saint-Sernin, Place du Capitole, Rue d'Alsace, Côte Pavée))}\par\vspace{0.55mm}
\noindent\rule{\textwidth}{0.35pt}\par\vspace*{\fill}
{\centering
  \setlength{\fboxsep}{1.2mm}%
  \setlength{\fboxrule}{0.4pt}%
  \fcolorbox{gray!55}{white}{%
    \includegraphics[width=0.95\textwidth,height=0.81\textheight,keepaspectratio]{\facsimiledir/20251126_110216.jpg}%
  }\par}
\vspace*{\fill}
\newpage

% =============================================================================
% VIS-A-VIS SPREAD : FOLIO 67 VERSO (20251126_110233.jpg)
% =============================================================================
\cleardoubleevenpage
\renewcommand{\currentmonthyear}{Janvier 1917}
\renewcommand{\currentfoliolabel}{Folio 67 verso}
\phantomsection\label{fol:135}
% --- En-tête du folio (Page paire / gauche : Transcription, apparat & registres de bas de page) ---
\noindent\rule{\textwidth}{0.5pt}\par\vspace{0.55mm}
\noindent\textbf{\footnotesize FOLIO 67 VERSO} \hfill \textbf{\footnotesize JANVIER 1917}\par\vspace{0.25mm}
\noindent{\fontsize{6.9pt}{8.0pt}\selectfont Visite à Mme de Labarthe et conférence d'Ernest Zyromski sur Maurice de Guérin}\par\vspace{0.2mm}
\noindent{\fontsize{6.9pt}{8.0pt}\selectfont\textit{30 janvier 1917} \hfill \textit{Toulouse (Faculté des Lettres, Sainte-Marie)}}\par\vspace{0.5mm}
\noindent\rule{\textwidth}{0.35pt}\par\vspace{0.8mm}
% --- Corps de la transcription diplomatique et apparat marginal en prolongement ---
{\renewcommand{\arraystretch}{0.87}%
\noindent\hspace*{3.5mm}\begin{tabular}{@{}p{87mm}@{\hspace{3.5mm}}p{39mm}@{}}
      \leavevmode\fontsize{6.9pt}{7.9pt}\selectfont\strut \underline{Mardi 30.} \hfill{} \underline{Janvier 1917 —}\strut &  \\
      \leavevmode\fontsize{6.9pt}{7.9pt}\selectfont\strut Levé 7h m. Allumé gd feu\strut & \smash{\parbox[t]{39mm}{\fontsize{6.2pt}{7.0pt}\selectfont\raggedright\color{black!88}\textbf{2 gd} \textit{abbr. pro} grand. \textbf{2 m.} \textit{abbr. pro} matin. — \textbf{3 s/} \textit{abbr. pro} sur.}} \\
      \leavevmode\fontsize{6.9pt}{7.9pt}\selectfont\strut chambre. Me suis prélassé s/ fauteuil. Joui\strut &  \\
      \leavevmode\fontsize{6.9pt}{7.9pt}\selectfont\strut de mon bien-être \& repos de l'heure présente.\strut &  \\
      \leavevmode\llap{\fontsize{6.4pt}{7.3pt}\selectfont\color{darkgray}5\hspace{1.8mm}}\fontsize{6.9pt}{7.9pt}\selectfont\strut Évoqué passé sans mélancolie.\strut &  \\
      \leavevmode\fontsize{6.9pt}{7.9pt}\selectfont\strut Matin, visite à Mme de Labarthe. L'ai\strut & \smash{\parbox[t]{39mm}{\fontsize{6.2pt}{7.0pt}\selectfont\raggedright\color{black!88}\textbf{6 Mme} \textit{abbr. pro} Madame.}} \\
      \leavevmode\fontsize{6.9pt}{7.9pt}\selectfont\strut revue av. un plaisir extrême.\strut &  \\
      \leavevmode\fontsize{6.9pt}{7.9pt}\selectfont\strut Monté, cme d'habitude, à Ste-Marie pour\strut & \smash{\parbox[t]{39mm}{\fontsize{6.2pt}{7.0pt}\selectfont\raggedright\color{black!88}\textbf{8 cme} \textit{abbr. pro} comme. \textbf{8 Ste} \textit{abbr. pro} Sainte.}} \\
      \leavevmode\fontsize{6.9pt}{7.9pt}\selectfont\strut déjeuner.\strut &  \\
      \leavevmode\llap{\fontsize{6.4pt}{7.3pt}\selectfont\color{darkgray}10\hspace{1.8mm}}\fontsize{6.9pt}{7.9pt}\selectfont\strut Après-midi, assisté conférce de M. Zyromski\strut &  \\
      \leavevmode\fontsize{6.9pt}{7.9pt}\selectfont\strut sur Maurice de Guérin (Influence de la Bretagne\strut &  \\
      \leavevmode\fontsize{6.9pt}{7.9pt}\selectfont\strut \& de Lamennais). Admirablt écrite. Passages\strut &  \\
      \leavevmode\fontsize{6.9pt}{7.9pt}\selectfont\strut remarquables fond et forme. Vrai\strut &  \\
      \leavevmode\fontsize{6.9pt}{7.9pt}\selectfont\strut régal littéraire.\strut &  \\
      \leavevmode\llap{\fontsize{6.4pt}{7.3pt}\selectfont\color{darkgray}15\hspace{1.8mm}}\fontsize{6.9pt}{7.9pt}\selectfont\strut Fini journée en cie de Margte —\strut & \smash{\parbox[t]{39mm}{\fontsize{6.2pt}{7.0pt}\selectfont\raggedright\color{black!88}\textbf{15 Margte} \textit{scil.} Marguerite (Calmels). \textbf{15 cie} \textit{abbr. pro} compagnie. — \textbf{16} \textit{in marg. cum unco} (comptes). \textbf{16 Cartes ples} \textit{abbr. pro} Cartes postales. — \textbf{19 c.} \textit{abbr. pro} centimes. — \textbf{20 c.} \textit{abbr. pro} centimes.}} \\
      \cline{1-1}
      \leavevmode\fontsize{6.9pt}{7.9pt}\selectfont\strut [ — Cartes ples : 10 c.\strut &  \\
      \leavevmode\fontsize{6.9pt}{7.9pt}\selectfont\strut — 1 ramette papier vélin : 1 f 25\strut &  \\
      \leavevmode\fontsize{6.9pt}{7.9pt}\selectfont\strut (moyen format)\strut &  \\
      \leavevmode\fontsize{6.9pt}{7.9pt}\selectfont\strut 2 paquets d'enveloppes : 80 c.\strut &  \\
      \leavevmode\llap{\fontsize{6.4pt}{7.3pt}\selectfont\color{darkgray}20\hspace{1.8mm}}\fontsize{6.9pt}{7.9pt}\selectfont\strut Trams : 10 c, 60 c, 10 c.\strut &  \\
      \leavevmode\fontsize{6.9pt}{7.9pt}\selectfont\strut Le Miroir des Modes (no de janvier) : 1 f — ]\strut &  \\
\end{tabular}}%
\par\vspace*{\fill}
% --- Registres de bas de page (1. Critique interne [si pertinent] ; 2. Loci paralleli en latin ; 3. Commentaire historique en retrait) ---
\noindent\rule{\textwidth}{0.5pt}\par\vspace{0.65mm}
{\fontsize{6.6pt}{7.7pt}\selectfont%
\noindent\normalfont\color{black!92} \textbf{8} (\textit{Ste-Marie}) cf. supra f\textsuperscript{o}~27~v\textsuperscript{o} (p.~\pageref{fol:55}), f\textsuperscript{o}~67~r\textsuperscript{o} (p.~\pageref{fol:134}), infra f\textsuperscript{o}~68~r\textsuperscript{o} (p.~\pageref{fol:136}), f\textsuperscript{o}~68~v\textsuperscript{o} (p.~\pageref{fol:137}), \textit{et passim}. — \textbf{10} (\textit{Zyromski}) cf. supra f\textsuperscript{o}~26~v\textsuperscript{o} (p.~\pageref{fol:53}). — \textbf{15} (\textit{Margte}) cf. supra f\textsuperscript{o}~1~r\textsuperscript{o} (p.~\pageref{fol:2}), f\textsuperscript{o}~67~r\textsuperscript{o} (p.~\pageref{fol:134}), infra f\textsuperscript{o}~68~r\textsuperscript{o} (p.~\pageref{fol:136}), f\textsuperscript{o}~68~v\textsuperscript{o} (p.~\pageref{fol:137}), \textit{et passim}.\par\vspace{0.55mm}
\noindent\rule{38mm}{0.4pt}\par\vspace{0.5mm}
\noindent\hspace*{5mm}\parbox{\dimexpr\textwidth-5mm\relax}{\fontsize{6.6pt}{7.7pt}\selectfont\color{black!84}Journée toulousaine du 30 janvier 1917 : après l'achat d'une montre-bracelet en or (200~fr.) et de deux bagues chez l'horloger-bijoutier \textbf{Pfister} (rue d'Alsace-Lorraine / place Esquirol), Hippolyte Combes assiste à 17~h avec Marguerite Calmels et sa sœur Madeleine au cours public du professeur \textbf{Ernest Zyromski} (1862–1933) à la Faculté des lettres de Toulouse, consacré au poète tarnais \textbf{Maurice de Guérin} (1810–1839).}\par}
\newpage
\phantomsection\label{fac:135}
% --- En-tête du fac-similé (Page impaire / droite : Photographie du manuscrit) ---
\noindent\rule{\textwidth}{0.5pt}\par\vspace{0.55mm}
\noindent\textbf{\fontsize{6.8pt}{7.8pt}\selectfont FAC-SIMILÉ EN VIS-À-VIS — FOLIO 67 VERSO} \hfill \textbf{\fontsize{6.8pt}{7.8pt}\selectfont JANVIER 1917}\par\vspace{0.25mm}
\noindent{\fontsize{6.8pt}{8.0pt}\selectfont\textit{Livre-Journal d'Hippolyte Combes} — 30 janvier 1917 (Toulouse (Faculté des Lettres, Sainte-Marie))}\par\vspace{0.55mm}
\noindent\rule{\textwidth}{0.35pt}\par\vspace*{\fill}
{\centering
  \setlength{\fboxsep}{1.2mm}%
  \setlength{\fboxrule}{0.4pt}%
  \fcolorbox{gray!55}{white}{%
    \includegraphics[width=0.95\textwidth,height=0.81\textheight,keepaspectratio]{\facsimiledir/20251126_110233.jpg}%
  }\par}
\vspace*{\fill}
\newpage

% =============================================================================
% VIS-A-VIS SPREAD : FOLIO 68 RECTO (20251126_110247.jpg)
% =============================================================================
\cleardoubleevenpage
\renewcommand{\currentmonthyear}{Janvier 1917}
\renewcommand{\currentfoliolabel}{Folio 68 recto}
\phantomsection\label{fol:136}
% --- En-tête du folio (Page paire / gauche : Transcription, apparat & registres de bas de page) ---
\noindent\rule{\textwidth}{0.5pt}\par\vspace{0.55mm}
\noindent\textbf{\footnotesize FOLIO 68 RECTO} \hfill \textbf{\footnotesize JANVIER 1917}\par\vspace{0.25mm}
\noindent{\fontsize{6.9pt}{8.0pt}\selectfont Marché aux Halles Victor-Hugo avec Marguerite, visite de l'Hôtel d'Assézat et intimité à l'appartement}\par\vspace{0.2mm}
\noindent{\fontsize{6.9pt}{8.0pt}\selectfont\textit{31 janvier 1917} \hfill \textit{Toulouse (Saint-Sernin, Halles Victor-Hugo, Hôtel d'Assézat, Daurade, Dalbade, Sainte-Marie)}}\par\vspace{0.5mm}
\noindent\rule{\textwidth}{0.35pt}\par\vspace{0.8mm}
% --- Corps de la transcription diplomatique et apparat marginal en prolongement ---
{\renewcommand{\arraystretch}{0.87}%
\noindent\hspace*{3.5mm}\begin{tabular}{@{}p{87mm}@{\hspace{3.5mm}}p{39mm}@{}}
      \leavevmode\fontsize{6.9pt}{7.9pt}\selectfont\strut \underline{Mercredi 31.} — \hfill{} \underline{Janvier 1917}\strut &  \\
      \leavevmode\fontsize{6.9pt}{7.9pt}\selectfont\strut Nuit, gelé très fort. Ds la journée, neigé 1 peu,\strut & \smash{\parbox[t]{39mm}{\fontsize{6.2pt}{7.0pt}\selectfont\raggedright\color{black!88}\textbf{2 Ds} \textit{abbr. pro} Dans.}} \\
      \leavevmode\fontsize{6.9pt}{7.9pt}\selectfont\strut et soufflé bise glaciale.\strut &  \\
      \leavevmode\fontsize{6.9pt}{7.9pt}\selectfont\strut Matinée, venu rejoindre Mte p. son tram à St-Ser.\strut & \smash{\parbox[t]{39mm}{\fontsize{6.2pt}{7.0pt}\selectfont\raggedright\color{black!88}\textbf{4 Mte} \textit{scil.} Marguerite (Calmels). \textbf{4 St} \textit{abbr. pro} Saint.}} \\
      \leavevmode\llap{\fontsize{6.4pt}{7.3pt}\selectfont\color{darkgray}5\hspace{1.8mm}}\fontsize{6.9pt}{7.9pt}\selectfont\strut vers 10h. Ensemble, fait marché à V. Hugo.\strut &  \\
      \leavevmode\fontsize{6.9pt}{7.9pt}\selectfont\strut Après-midi, descendus en ville. Promené\strut &  \\
      \leavevmode\fontsize{6.9pt}{7.9pt}\selectfont\strut rues bien abritées. Visité Hôtel d'Assézat,\strut &  \\
      \leavevmode\fontsize{6.9pt}{7.9pt}\selectfont\strut Daurade \& Dalbade.\strut &  \\
      \leavevmode\fontsize{6.9pt}{7.9pt}\selectfont\strut Puis montés mon appartt et y reposé 2 longues\strut & \smash{\parbox[t]{39mm}{\fontsize{6.2pt}{7.0pt}\selectfont\raggedright\color{black!88}\textbf{9 appartt} \textit{abbr. pro} appartement. — \textbf{10} \textit{lect. incert. in ms.} \textbf{10 qq} \textit{abbr. pro} quelque(s). — \textbf{11 Ste} \textit{abbr. pro} Sainte. \textbf{11 1/2} \textit{scil.} demi / moitié. — \textbf{14} \textit{in marg. cum unco} (comptes).}} \\
      \leavevmode\llap{\fontsize{6.4pt}{7.3pt}\selectfont\color{darkgray}10\hspace{1.8mm}}\fontsize{6.9pt}{7.9pt}\selectfont\strut heures. Évité troubler qq [mot non établi] précieuses.\strut &  \\
      \leavevmode\fontsize{6.9pt}{7.9pt}\selectfont\strut A Ste-Marie, vers 6h 1/2. Retour à pied, vers\strut &  \\
      \leavevmode\fontsize{6.9pt}{7.9pt}\selectfont\strut 9h. Légers flocons neige tombaient sur le sol\strut &  \\
      \leavevmode\fontsize{6.9pt}{7.9pt}\selectfont\strut durci.\strut &  \\
      \leavevmode\fontsize{6.9pt}{7.9pt}\selectfont\strut [ — Halles Victor-Hugo : gigot d'agneau : 5 f\strut &  \\
      \leavevmode\llap{\fontsize{6.4pt}{7.3pt}\selectfont\color{darkgray}15\hspace{1.8mm}}\fontsize{6.9pt}{7.9pt}\selectfont\strut Petits fromages (gervais \& blancs) : 90 c\strut &  \\
      \leavevmode\fontsize{6.9pt}{7.9pt}\selectfont\strut 1/2 beurre : 60 c. — Boîte sardines : 2 f\strut & \smash{\parbox[t]{39mm}{\fontsize{6.2pt}{7.0pt}\selectfont\raggedright\color{black!88}\textbf{16 1/2} \textit{scil.} demi / moitié. \textbf{16 c.} \textit{abbr. pro} centimes. — \textbf{17 c.} \textit{abbr. pro} centimes. — \textbf{18 f.} \textit{abbr. pro} franc(s). — \textbf{19 Testet} \textit{lect. incert. in ms.} — \textbf{20 c.} \textit{abbr. pro} centimes. — \textbf{21 Mme} \textit{abbr. pro} Madame.}} \\
      \leavevmode\fontsize{6.9pt}{7.9pt}\selectfont\strut Gâteaux : 20 c. — 1 livre de dattes : 1 f — id. de\strut &  \\
      \leavevmode\fontsize{6.9pt}{7.9pt}\selectfont\strut figues : 1 f. — Gâteau aux noix : 1 f 25 c\strut &  \\
      \leavevmode\fontsize{6.9pt}{7.9pt}\selectfont\strut Gâteaux des allées : 1 f 10 chez Henri Testet [?]\strut &  \\
      \leavevmode\llap{\fontsize{6.4pt}{7.3pt}\selectfont\color{darkgray}20\hspace{1.8mm}}\fontsize{6.9pt}{7.9pt}\selectfont\strut Trams : 30 c., 30 c.\strut &  \\
      \cline{1-1}
      \leavevmode\fontsize{6.9pt}{7.9pt}\selectfont\strut Télégramme à Mme Spagnol, Nice : 2 f — ]\strut &  \\
\end{tabular}}%
\par\vspace*{\fill}
% --- Registres de bas de page (1. Critique interne [si pertinent] ; 2. Loci paralleli en latin ; 3. Commentaire historique en retrait) ---
\noindent\rule{\textwidth}{0.5pt}\par\vspace{0.65mm}
{\fontsize{6.6pt}{7.7pt}\selectfont%
\noindent\normalfont \textbf{Rature autographe (l.~15) :} À la fin du récit de la visite à son ancien appartement toulousain, Hippolyte Combes a biffé une proposition commençant par \textit{Évité troubler} sous une surcharge d'encre dense.\par\vspace{0.5mm}
\noindent\rule{22mm}{0.3pt}\par\vspace{0.45mm}
\noindent\normalfont\color{black!92} \textbf{4} (\textit{Mte}) cf. supra f\textsuperscript{o}~1~r\textsuperscript{o} (p.~\pageref{fol:2}), f\textsuperscript{o}~67~v\textsuperscript{o} (p.~\pageref{fol:135}), infra f\textsuperscript{o}~68~v\textsuperscript{o} (p.~\pageref{fol:137}), f\textsuperscript{o}~69~r\textsuperscript{o} (p.~\pageref{fol:138}), \textit{et passim}, Appendicem I, doc.~5 (p.~\pageref{annexe:doc5}). — \textbf{11} (\textit{Ste-Marie}) cf. supra f\textsuperscript{o}~27~v\textsuperscript{o} (p.~\pageref{fol:55}), f\textsuperscript{o}~67~v\textsuperscript{o} (p.~\pageref{fol:135}), infra f\textsuperscript{o}~68~v\textsuperscript{o} (p.~\pageref{fol:137}), f\textsuperscript{o}~69~r\textsuperscript{o} (p.~\pageref{fol:138}), \textit{et passim}. — \textbf{21} (\textit{Spagnol}) cf. infra f\textsuperscript{o}~79~r\textsuperscript{o} (p.~\pageref{fol:158}).\par\vspace{0.55mm}
\noindent\rule{38mm}{0.4pt}\par\vspace{0.5mm}
\noindent\hspace*{5mm}\parbox{\dimexpr\textwidth-5mm\relax}{\fontsize{6.6pt}{7.7pt}\selectfont\color{black!84}Promenade à \textbf{Toulouse} avec Marguerite Calmels le 31 janvier 1917 par le marché Victor-Hugo, la cour Renaissance de l'\textbf{Hôtel d'Assézat} (édifié à partir de 1555 sur les plans de Nicolas Bachelier), et les églises \textbf{Notre-Dame de la Daurade} et \textbf{Notre-Dame de la Dalbade} (dont le cousin Jean Crouzet deviendra curé), suivie au matin du 1\textsuperscript{er} février d'une marche vers \textbf{Lardenne} et les ruines de l'amphithéâtre gallo-romain de \textbf{Purpan-Ancely} (I\textsuperscript{er}~s. apr. J.-C.).}\par}
\newpage
\phantomsection\label{fac:136}
% --- En-tête du fac-similé (Page impaire / droite : Photographie du manuscrit) ---
\noindent\rule{\textwidth}{0.5pt}\par\vspace{0.55mm}
\noindent\textbf{\fontsize{6.8pt}{7.8pt}\selectfont FAC-SIMILÉ EN VIS-À-VIS — FOLIO 68 RECTO} \hfill \textbf{\fontsize{6.8pt}{7.8pt}\selectfont JANVIER 1917}\par\vspace{0.25mm}
\noindent{\fontsize{6.8pt}{8.0pt}\selectfont\textit{Livre-Journal d'Hippolyte Combes} — 31 janvier 1917 (Toulouse (Saint-Sernin, Halles Victor-Hugo, Hôtel d'Assézat, Daurade, Dalbade, Sainte-Marie))}\par\vspace{0.55mm}
\noindent\rule{\textwidth}{0.35pt}\par\vspace*{\fill}
{\centering
  \setlength{\fboxsep}{1.2mm}%
  \setlength{\fboxrule}{0.4pt}%
  \fcolorbox{gray!55}{white}{%
    \includegraphics[width=0.95\textwidth,height=0.81\textheight,keepaspectratio]{\facsimiledir/20251126_110247.jpg}%
  }\par}
\vspace*{\fill}
\newpage

% =============================================================================
% VIS-A-VIS SPREAD : FOLIO 68 VERSO (20251126_110302.jpg)
% =============================================================================
\cleardoubleevenpage
\renewcommand{\currentmonthyear}{Février 1917}
\renewcommand{\currentfoliolabel}{Folio 68 verso}
\phantomsection\label{fol:137}
% --- En-tête du folio (Page paire / gauche : Transcription, apparat & registres de bas de page) ---
\noindent\rule{\textwidth}{0.5pt}\par\vspace{0.55mm}
\noindent\textbf{\footnotesize FOLIO 68 VERSO} \hfill \textbf{\footnotesize FÉVRIER 1917}\par\vspace{0.25mm}
\noindent{\fontsize{6.9pt}{8.0pt}\selectfont Achat d'une montre en or chez Pfister et détachement intérieur vis-à-vis de Toulouse}\par\vspace{0.2mm}
\noindent{\fontsize{6.9pt}{8.0pt}\selectfont\textit{1er février 1917} \hfill \textit{Toulouse (Capitole, Sainte-Marie, Gare Matabiau, Rue de Languedoc)}}\par\vspace{0.5mm}
\noindent\rule{\textwidth}{0.35pt}\par\vspace{0.8mm}
% --- Corps de la transcription diplomatique et apparat marginal en prolongement ---
{\renewcommand{\arraystretch}{0.87}%
\noindent\hspace*{3.5mm}\begin{tabular}{@{}p{87mm}@{\hspace{3.5mm}}p{39mm}@{}}
      \leavevmode\fontsize{6.9pt}{7.9pt}\selectfont\strut \underline{1er Jeudi} \hfill{} \underline{Février 1917 —}\strut & \smash{\parbox[t]{39mm}{\fontsize{6.2pt}{7.0pt}\selectfont\raggedright\color{black!88}\textbf{1 1er} \textit{scil.} premier.}} \\
      \leavevmode\fontsize{6.9pt}{7.9pt}\selectfont\strut Temps très froid. Gelé dur. Dans la journée, soleil\strut &  \\
      \leavevmode\fontsize{6.9pt}{7.9pt}\selectfont\strut brillant.\strut &  \\
      \leavevmode\fontsize{6.9pt}{7.9pt}\selectfont\strut Matin, longue visite à Mme de Labarthe\strut & \smash{\parbox[t]{39mm}{\fontsize{6.2pt}{7.0pt}\selectfont\raggedright\color{black!88}\textbf{4 Mme} \textit{abbr. pro} Madame.}} \\
      \leavevmode\llap{\fontsize{6.4pt}{7.3pt}\selectfont\color{darkgray}5\hspace{1.8mm}}\fontsize{6.9pt}{7.9pt}\selectfont\strut \& aux dames Delsol, mes voisines.\strut &  \\
      \leavevmode\fontsize{6.9pt}{7.9pt}\selectfont\strut Au Capitole, retrouvé Mte, vers 11h 1/2 et sa\strut & \smash{\parbox[t]{39mm}{\fontsize{6.2pt}{7.0pt}\selectfont\raggedright\color{black!88}\textbf{6 Mte} \textit{scil.} Marguerite (Calmels). \textbf{6 1/2} \textit{scil.} demi / moitié. — \textbf{7 Ste} \textit{abbr. pro} Sainte. — \textbf{9 1/2} \textit{scil.} demi / moitié. — \textbf{10 tte} \textit{abbr. pro} toute. — \textbf{11 ch.} \textit{abbr. pro} chez.}} \\
      \leavevmode\fontsize{6.9pt}{7.9pt}\selectfont\strut sœur! Remontés ensemble tram Ste-Marie.\strut &  \\
      \leavevmode\fontsize{6.9pt}{7.9pt}\selectfont\strut Après déjeuner, descendu ville av. Henriette.\strut &  \\
      \leavevmode\fontsize{6.9pt}{7.9pt}\selectfont\strut Vers 4h 1/2, à gare, à propos panier fleurs\strut &  \\
      \leavevmode\llap{\fontsize{6.4pt}{7.3pt}\selectfont\color{darkgray}10\hspace{1.8mm}}\fontsize{6.9pt}{7.9pt}\selectfont\strut égaré. A Nazareth, revu tte ma famille.\strut &  \\
      \leavevmode\fontsize{6.9pt}{7.9pt}\selectfont\strut Choisi montre or ch. Pfister.\strut &  \\
      \leavevmode\fontsize{6.9pt}{7.9pt}\selectfont\strut Chapelle archevêché : conférence, pas eu lieu.\strut &  \\
      \leavevmode\fontsize{6.9pt}{7.9pt}\selectfont\strut Rentré dîner à Ste-Marie —\strut & \smash{\parbox[t]{39mm}{\fontsize{6.2pt}{7.0pt}\selectfont\raggedright\color{black!88}\textbf{13 Ste} \textit{abbr. pro} Sainte.}} \\
      \leavevmode\fontsize{6.9pt}{7.9pt}\selectfont\strut Vers 9h 1/2, à pied j'à rue de Languedoc.\strut & \smash{\parbox[t]{39mm}{\fontsize{6.2pt}{7.0pt}\selectfont\raggedright\color{black!88}\textbf{14 j'à} \textit{abbr. pro} jusqu'à. \textbf{14 1/2} \textit{scil.} demi / moitié. — \textbf{16 complètt} \textit{abbr. pro} complètement.}} \\
      \leavevmode\llap{\fontsize{6.4pt}{7.3pt}\selectfont\color{darkgray}15\hspace{1.8mm}}\fontsize{6.9pt}{7.9pt}\selectfont\strut Ville silencieuse. Maisons fermées. Je me\strut &  \\
      \leavevmode\fontsize{6.9pt}{7.9pt}\selectfont\strut sens complètt détaché de Toulouse. Partirai sans\strut &  \\
      \leavevmode\fontsize{6.9pt}{7.9pt}\selectfont\strut regret, et pour n'y retourner jamais plus.\strut &  \\
      \cline{1-1}
      \leavevmode\fontsize{6.9pt}{7.9pt}\selectfont\strut [ — Trams : 30 c., 10 c., 20 c., 30 c.\strut & \smash{\parbox[t]{39mm}{\fontsize{6.2pt}{7.0pt}\selectfont\raggedright\color{black!88}\textbf{18} \textit{in marg. cum unco} (comptes). \textbf{18 c.} \textit{abbr. pro} centimes. — \textbf{19 Food} \textit{angl. in ms.} — \textbf{21 1/2} \textit{scil.} demi / moitié.}} \\
      \leavevmode\fontsize{6.9pt}{7.9pt}\selectfont\strut Food (boulanger) : 25 c\strut &  \\
      \leavevmode\llap{\fontsize{6.4pt}{7.3pt}\selectfont\color{darkgray}20\hspace{1.8mm}}\fontsize{6.9pt}{7.9pt}\selectfont\strut Remboursé à mes voisines : 3 f 25\strut &  \\
      \leavevmode\fontsize{6.9pt}{7.9pt}\selectfont\strut (12 œufs, 1/2 beurre, 1 boulanger) ]\strut &  \\
\end{tabular}}%
\par\vspace*{\fill}
% --- Registres de bas de page (1. Critique interne [si pertinent] ; 2. Loci paralleli en latin ; 3. Commentaire historique en retrait) ---
\noindent\rule{\textwidth}{0.5pt}\par\vspace{0.65mm}
{\fontsize{6.6pt}{7.7pt}\selectfont%
\noindent\normalfont\color{black!92} \textbf{6} (\textit{Mte}) cf. supra f\textsuperscript{o}~1~r\textsuperscript{o} (p.~\pageref{fol:2}), f\textsuperscript{o}~68~r\textsuperscript{o} (p.~\pageref{fol:136}), infra f\textsuperscript{o}~69~r\textsuperscript{o} (p.~\pageref{fol:138}), f\textsuperscript{o}~69~v\textsuperscript{o} (p.~\pageref{fol:139}), \textit{et passim}, Appendicem I, doc.~5 (p.~\pageref{annexe:doc5}). — \textbf{7} (\textit{Ste-Marie}) cf. supra f\textsuperscript{o}~27~v\textsuperscript{o} (p.~\pageref{fol:55}), f\textsuperscript{o}~68~r\textsuperscript{o} (p.~\pageref{fol:136}), infra f\textsuperscript{o}~69~r\textsuperscript{o} (p.~\pageref{fol:138}), f\textsuperscript{o}~69~v\textsuperscript{o} (p.~\pageref{fol:139}), \textit{et passim}. — \textbf{11} (\textit{Pfister}) cf. supra f\textsuperscript{o}~28~r\textsuperscript{o} (p.~\pageref{fol:56}), infra f\textsuperscript{o}~69~r\textsuperscript{o} (p.~\pageref{fol:138}).\par\vspace{0.55mm}
\noindent\rule{38mm}{0.4pt}\par\vspace{0.5mm}
\noindent\hspace*{5mm}\parbox{\dimexpr\textwidth-5mm\relax}{\fontsize{6.6pt}{7.7pt}\selectfont\color{black!84}Le jeudi 1\textsuperscript{er} février 1917, après avoir visité Mme de Labarthe et ses voisines d'appartement les \textbf{dames Delsol}, Hippolyte Combes descend en ville avec \textbf{Henriette}, réclame en gare Matabiau au sujet d'un panier de fleurs égaré, se recueille à la chapelle \textbf{Notre-Dame-de-Nazareth} (évoquant le souvenir de ses parents et sœurs) et choisit chez l'horloger-bijoutier \textbf{Pfister} une montre en or pour Marguerite (réglée le lendemain, f\textsuperscript{o}~69~r\textsuperscript{o}), avant de noter le soir, rue de Languedoc, son détachement définitif à l'égard de Toulouse.}\par}
\newpage
\phantomsection\label{fac:137}
% --- En-tête du fac-similé (Page impaire / droite : Photographie du manuscrit) ---
\noindent\rule{\textwidth}{0.5pt}\par\vspace{0.55mm}
\noindent\textbf{\fontsize{6.8pt}{7.8pt}\selectfont FAC-SIMILÉ EN VIS-À-VIS — FOLIO 68 VERSO} \hfill \textbf{\fontsize{6.8pt}{7.8pt}\selectfont FÉVRIER 1917}\par\vspace{0.25mm}
\noindent{\fontsize{6.8pt}{8.0pt}\selectfont\textit{Livre-Journal d'Hippolyte Combes} — 1er février 1917 (Toulouse (Capitole, Sainte-Marie, Gare Matabiau, Rue de Languedoc))}\par\vspace{0.55mm}
\noindent\rule{\textwidth}{0.35pt}\par\vspace*{\fill}
{\centering
  \setlength{\fboxsep}{1.2mm}%
  \setlength{\fboxrule}{0.4pt}%
  \fcolorbox{gray!55}{white}{%
    \includegraphics[width=0.95\textwidth,height=0.81\textheight,keepaspectratio]{\facsimiledir/20251126_110302.jpg}%
  }\par}
\vspace*{\fill}
\newpage

% =============================================================================
% VIS-A-VIS SPREAD : FOLIO 69 RECTO (20251126_110324.jpg)
% =============================================================================
\cleardoubleevenpage
\renewcommand{\currentmonthyear}{Février 1917}
\renewcommand{\currentfoliolabel}{Folio 69 recto}
\phantomsection\label{fol:138}
% --- En-tête du folio (Page paire / gauche : Transcription, apparat & registres de bas de page) ---
\noindent\rule{\textwidth}{0.5pt}\par\vspace{0.55mm}
\noindent\textbf{\footnotesize FOLIO 69 RECTO} \hfill \textbf{\footnotesize FÉVRIER 1917}\par\vspace{0.25mm}
\noindent{\fontsize{6.9pt}{8.0pt}\selectfont Soins chez le dentiste M. de Lersay, achat du remontoir Louis XVI et dernière soirée à Toulouse}\par\vspace{0.2mm}
\noindent{\fontsize{6.9pt}{8.0pt}\selectfont\textit{2 février 1917} \hfill \textit{Toulouse (Rue Lafayette, Place Belfort, Sainte-Marie)}}\par\vspace{0.5mm}
\noindent\rule{\textwidth}{0.35pt}\par\vspace{0.8mm}
% --- Corps de la transcription diplomatique et apparat marginal en prolongement ---
{\renewcommand{\arraystretch}{0.87}%
\noindent\hspace*{3.5mm}\begin{tabular}{@{}p{87mm}@{\hspace{3.5mm}}p{39mm}@{}}
      \leavevmode\fontsize{6.9pt}{7.9pt}\selectfont\strut \underline{2 Vendredi.} \hfill{} \underline{Février 1917 —}\strut &  \\
      \leavevmode\fontsize{6.9pt}{7.9pt}\selectfont\strut Tjrs même temps, très froid. —\strut & \smash{\parbox[t]{39mm}{\fontsize{6.2pt}{7.0pt}\selectfont\raggedright\color{black!88}\textbf{2 Tjrs} \textit{abbr. pro} Toujours.}} \\
      \leavevmode\fontsize{6.9pt}{7.9pt}\selectfont\strut Neigé abondamment durant matinée.\strut &  \\
      \leavevmode\fontsize{6.9pt}{7.9pt}\selectfont\strut Matin \& soir, venu ch. dentiste, M. de Lersay,\strut & \smash{\parbox[t]{39mm}{\fontsize{6.2pt}{7.0pt}\selectfont\raggedright\color{black!88}\textbf{4 ch.} \textit{abbr. pro} chez.}} \\
      \leavevmode\llap{\fontsize{6.4pt}{7.3pt}\selectfont\color{darkgray}5\hspace{1.8mm}}\fontsize{6.9pt}{7.9pt}\selectfont\strut rue La Fayette. Très bien soigné par j. dame,\strut &  \\
      \leavevmode\fontsize{6.9pt}{7.9pt}\selectfont\strut opératrice.\strut &  \\
      \leavevmode\fontsize{6.9pt}{7.9pt}\selectfont\strut Après-midi, visite à Mme de Labarthe. Puis\strut & \smash{\parbox[t]{39mm}{\fontsize{6.2pt}{7.0pt}\selectfont\raggedright\color{black!88}\textbf{7 Mme} \textit{abbr. pro} Madame.}} \\
      \leavevmode\fontsize{6.9pt}{7.9pt}\selectfont\strut monté ch. Henriette : y retrouvé Margte etc.\strut & \smash{\parbox[t]{39mm}{\fontsize{6.2pt}{7.0pt}\selectfont\raggedright\color{black!88}\textbf{8 ch.} \textit{abbr. pro} chez. \textbf{8 Margte} \textit{scil.} Marguerite (Calmels). — \textbf{9 Ste} \textit{abbr. pro} Sainte. \textbf{9 s.} \textit{abbr. pro} soir. — \textbf{11} \textit{in marg. cum unco} (comptes). \textbf{11 Ch.} \textit{abbr. pro} Chez. — \textbf{12 fr.} \textit{abbr. pro} franc(s). — \textbf{14 c.} \textit{abbr. pro} centimes.}} \\
      \leavevmode\fontsize{6.9pt}{7.9pt}\selectfont\strut Descendu à pied de Ste-Marie, vers 10h s.\strut &  \\
      \leavevmode\llap{\fontsize{6.4pt}{7.3pt}\selectfont\color{darkgray}10\hspace{1.8mm}}\fontsize{6.9pt}{7.9pt}\selectfont\strut Froid piquant. Sol durci.\strut &  \\
      \leavevmode\fontsize{6.9pt}{7.9pt}\selectfont\strut [ — Ch. Pfister, remontoir or mat Louis XVI\strut &  \\
      \leavevmode\fontsize{6.9pt}{7.9pt}\selectfont\strut 105 fr. — Petit anneau pour attacher pendentif\strut &  \\
      \leavevmode\fontsize{6.9pt}{7.9pt}\selectfont\strut à sautoir : free. — 1 boîte pastilles de\strut &  \\
      \leavevmode\fontsize{6.9pt}{7.9pt}\selectfont\strut Chlorate de Potasse : 1 f 30 c. — Trams : 20 c., 10 c\strut &  \\
      \cline{1-1}
      \leavevmode\llap{\fontsize{6.4pt}{7.3pt}\selectfont\color{darkgray}15\hspace{1.8mm}}\fontsize{6.9pt}{7.9pt}\selectfont\strut Bain, Place Belfort : 1 f — ]\strut &  \\
      \leavevmode\fontsize{6.9pt}{7.9pt}\selectfont\strut Quitte à regret Tls, sacrifiant 3 jours entiers\strut &  \\
      \leavevmode\fontsize{6.9pt}{7.9pt}\selectfont\strut pour aller à Millau... Aurais préféré rester\strut &  \\
      \leavevmode\fontsize{6.9pt}{7.9pt}\selectfont\strut encore ici, afin de jouir + lgts de mon\strut & \smash{\parbox[t]{39mm}{\fontsize{6.2pt}{7.0pt}\selectfont\raggedright\color{black!88}\textbf{18 lgts} \textit{abbr. pro} longtemps. — \textbf{19 appartt} \textit{abbr. pro} appartement.}} \\
      \leavevmode\fontsize{6.9pt}{7.9pt}\selectfont\strut appartt, rempli de souvenirs.\strut &  \\
\end{tabular}}%
\par\vspace*{\fill}
% --- Registres de bas de page (1. Critique interne [si pertinent] ; 2. Loci paralleli en latin ; 3. Commentaire historique en retrait) ---
\noindent\rule{\textwidth}{0.5pt}\par\vspace{0.65mm}
{\fontsize{6.6pt}{7.7pt}\selectfont%
\noindent\normalfont\color{black!92} \textbf{8} (\textit{Margte}) cf. supra f\textsuperscript{o}~1~r\textsuperscript{o} (p.~\pageref{fol:2}), f\textsuperscript{o}~68~v\textsuperscript{o} (p.~\pageref{fol:137}), infra f\textsuperscript{o}~69~v\textsuperscript{o} (p.~\pageref{fol:139}), f\textsuperscript{o}~70~r\textsuperscript{o} (p.~\pageref{fol:140}), \textit{et passim}, Appendicem I, doc.~5 (p.~\pageref{annexe:doc5}). — \textbf{9} (\textit{Ste-Marie}) cf. supra f\textsuperscript{o}~27~v\textsuperscript{o} (p.~\pageref{fol:55}), f\textsuperscript{o}~68~v\textsuperscript{o} (p.~\pageref{fol:137}), infra f\textsuperscript{o}~69~v\textsuperscript{o} (p.~\pageref{fol:139}), \textit{et passim}. — \textbf{11} (\textit{Pfister}) cf. supra f\textsuperscript{o}~28~r\textsuperscript{o} (p.~\pageref{fol:56}), f\textsuperscript{o}~68~v\textsuperscript{o} (p.~\pageref{fol:137}).\par\vspace{0.55mm}
\noindent\rule{38mm}{0.4pt}\par\vspace{0.5mm}
\noindent\hspace*{5mm}\parbox{\dimexpr\textwidth-5mm\relax}{\fontsize{6.6pt}{7.7pt}\selectfont\color{black!84}Dernière journée à Toulouse (vendredi 2 février 1917): soins dentaires rue La Fayette chez \textbf{M.~de Lersay} (réglés 25~fr. le lendemain, f\textsuperscript{o}~70~r\textsuperscript{o}), visite à Mme de Labarthe et à Henriette, et règlement chez l'horloger-bijoutier \textbf{Pfister} de la montre à remontoir en or mat de style Louis~XVI (105~fr.) destinée à être suspendue au sautoir en or offert le 29 septembre 1916 (f\textsuperscript{o}~28~r\textsuperscript{o}).}\par}
\newpage
\phantomsection\label{fac:138}
% --- En-tête du fac-similé (Page impaire / droite : Photographie du manuscrit) ---
\noindent\rule{\textwidth}{0.5pt}\par\vspace{0.55mm}
\noindent\textbf{\fontsize{6.8pt}{7.8pt}\selectfont FAC-SIMILÉ EN VIS-À-VIS — FOLIO 69 RECTO} \hfill \textbf{\fontsize{6.8pt}{7.8pt}\selectfont FÉVRIER 1917}\par\vspace{0.25mm}
\noindent{\fontsize{6.8pt}{8.0pt}\selectfont\textit{Livre-Journal d'Hippolyte Combes} — 2 février 1917 (Toulouse (Rue Lafayette, Place Belfort, Sainte-Marie))}\par\vspace{0.55mm}
\noindent\rule{\textwidth}{0.35pt}\par\vspace*{\fill}
{\centering
  \setlength{\fboxsep}{1.2mm}%
  \setlength{\fboxrule}{0.4pt}%
  \fcolorbox{gray!55}{white}{%
    \includegraphics[width=0.95\textwidth,height=0.81\textheight,keepaspectratio]{\facsimiledir/20251126_110324.jpg}%
  }\par}
\vspace*{\fill}
\newpage

% =============================================================================
% VIS-A-VIS SPREAD : FOLIO 69 VERSO (20251126_110338.jpg)
% =============================================================================
\cleardoubleevenpage
\renewcommand{\currentmonthyear}{Février 1917}
\renewcommand{\currentfoliolabel}{Folio 69 verso}
\phantomsection\label{fol:139}
% --- En-tête du folio (Page paire / gauche : Transcription, apparat & registres de bas de page) ---
\noindent\rule{\textwidth}{0.5pt}\par\vspace{0.55mm}
\noindent\textbf{\footnotesize FOLIO 69 VERSO} \hfill \textbf{\footnotesize FÉVRIER 1917}\par\vspace{0.25mm}
\noindent{\fontsize{6.9pt}{8.0pt}\selectfont Voyage de Toulouse à Millau avec Marguerite par Carcassonne et Béziers}\par\vspace{0.2mm}
\noindent{\fontsize{6.9pt}{8.0pt}\selectfont\textit{3 février 1917} \hfill \textit{Toulouse, Carcassonne, Béziers, Bédarieux, Millau (Aveyron)}}\par\vspace{0.5mm}
\noindent\rule{\textwidth}{0.35pt}\par\vspace{0.8mm}
% --- Corps de la transcription diplomatique et apparat marginal en prolongement ---
{\renewcommand{\arraystretch}{0.87}%
\noindent\hspace*{3.5mm}\begin{tabular}{@{}p{87mm}@{\hspace{3.5mm}}p{39mm}@{}}
      \leavevmode\fontsize{6.9pt}{7.9pt}\selectfont\strut \underline{Sam. 3} \hfill{} \underline{Févr. /17}\strut & \smash{\parbox[t]{39mm}{\fontsize{6.2pt}{7.0pt}\selectfont\raggedright\color{black!88}\textbf{1 Sam.} \textit{scil.} Samedi.}} \\
      \leavevmode\fontsize{6.9pt}{7.9pt}\selectfont\strut Levé 5h —\strut &  \\
      \leavevmode\fontsize{6.9pt}{7.9pt}\selectfont\strut Mis ordre ds mes affaires \& caisses.\strut & \smash{\parbox[t]{39mm}{\fontsize{6.2pt}{7.0pt}\selectfont\raggedright\color{black!88}\textbf{3 ds} \textit{abbr. pro} dans.}} \\
      \leavevmode\fontsize{6.9pt}{7.9pt}\selectfont\strut A 8h 1/2, ch. Auriol, joaillier, puis ch. dentiste.\strut & \smash{\parbox[t]{39mm}{\fontsize{6.2pt}{7.0pt}\selectfont\raggedright\color{black!88}\textbf{4 ch.} \textit{abbr. pro} chez. \textbf{4 1/2} \textit{scil.} demi / moitié. — \textbf{5 Margte} \textit{scil.} Marguerite (Calmels). — \textbf{6 Ste} \textit{abbr. pro} Sainte.}} \\
      \leavevmode\llap{\fontsize{6.4pt}{7.3pt}\selectfont\color{darkgray}5\hspace{1.8mm}}\fontsize{6.9pt}{7.9pt}\selectfont\strut A 10h, quitté maison et rejoint Margte.\strut &  \\
      \leavevmode\fontsize{6.9pt}{7.9pt}\selectfont\strut Partis de Ste-Marie, vers 11h 40.\strut &  \\
      \leavevmode\fontsize{6.9pt}{7.9pt}\selectfont\strut Oncle m'a accompagné gare.\strut &  \\
      \leavevmode\fontsize{6.9pt}{7.9pt}\selectfont\strut Pris express midi 30. 15 m. de retard.\strut &  \\
      \leavevmode\fontsize{6.9pt}{7.9pt}\selectfont\strut A Carcassonne, pu embrasser Jean.\strut &  \\
      \leavevmode\llap{\fontsize{6.4pt}{7.3pt}\selectfont\color{darkgray}10\hspace{1.8mm}}\fontsize{6.9pt}{7.9pt}\selectfont\strut A Béziers, à 4h 1/2. Montés à pied ville\strut & \smash{\parbox[t]{39mm}{\fontsize{6.2pt}{7.0pt}\selectfont\raggedright\color{black!88}\textbf{10 1/2} \textit{scil.} demi / moitié. — \textbf{11 rapidt} \textit{abbr. pro} rapidement. \textbf{11 cathle} \textit{abbr. pro} cathédrale.}} \\
      \leavevmode\fontsize{6.9pt}{7.9pt}\selectfont\strut et visité rapidt cathle. Admiré panorama\strut &  \\
      \leavevmode\fontsize{6.9pt}{7.9pt}\selectfont\strut terrasse. Plaine, rivière et soleil couchant.\strut &  \\
      \leavevmode\fontsize{6.9pt}{7.9pt}\selectfont\strut Repartis de Béziers à 5h 23, petit train,\strut &  \\
      \leavevmode\fontsize{6.9pt}{7.9pt}\selectfont\strut vieux wagons, pas de 2e —\strut &  \\
      \leavevmode\llap{\fontsize{6.4pt}{7.3pt}\selectfont\color{darkgray}15\hspace{1.8mm}}\fontsize{6.9pt}{7.9pt}\selectfont\strut Délicieuse intimité s/ l'épaule de Mte.\strut & \smash{\parbox[t]{39mm}{\fontsize{6.2pt}{7.0pt}\selectfont\raggedright\color{black!88}\textbf{15 s/} \textit{abbr. pro} sur. \textbf{15 Mte} \textit{scil.} Marguerite (Calmels).}} \\
      \leavevmode\fontsize{6.9pt}{7.9pt}\selectfont\strut A 10h 30, à Millau.\strut &  \\
      \leavevmode\fontsize{6.9pt}{7.9pt}\selectfont\strut Pierre à la gare nous attendait —\strut &  \\
      \leavevmode\fontsize{6.9pt}{7.9pt}\selectfont\strut Impression silence profond en débarquant\strut &  \\
      \leavevmode\fontsize{6.9pt}{7.9pt}\selectfont\strut ville natale.\strut &  \\
      \leavevmode\llap{\fontsize{6.4pt}{7.3pt}\selectfont\color{darkgray}20\hspace{1.8mm}}\fontsize{6.9pt}{7.9pt}\selectfont\strut A la maison de Mte, parents m'ont\strut & \smash{\parbox[t]{39mm}{\fontsize{6.2pt}{7.0pt}\selectfont\raggedright\color{black!88}\textbf{20 Mte} \textit{scil.} Marguerite (Calmels).}} \\
      \leavevmode\fontsize{6.9pt}{7.9pt}\selectfont\strut fait accueil très affectueux.\strut &  \\
      \leavevmode\fontsize{6.9pt}{7.9pt}\selectfont\strut On m'a donné 1 des + belles chambres,\strut &  \\
      \leavevmode\fontsize{6.9pt}{7.9pt}\selectfont\strut au 1er étage — —\strut & \smash{\parbox[t]{39mm}{\fontsize{6.2pt}{7.0pt}\selectfont\raggedright\color{black!88}\textbf{23 1er} \textit{scil.} premier.}} \\
\end{tabular}}%
\par\vspace*{\fill}
% --- Registres de bas de page (1. Critique interne [si pertinent] ; 2. Loci paralleli en latin ; 3. Commentaire historique en retrait) ---
\noindent\rule{\textwidth}{0.5pt}\par\vspace{0.65mm}
{\fontsize{6.6pt}{7.7pt}\selectfont%
\noindent\normalfont\color{black!92} \textbf{5} (\textit{Margte}) cf. supra f\textsuperscript{o}~1~r\textsuperscript{o} (p.~\pageref{fol:2}), f\textsuperscript{o}~69~r\textsuperscript{o} (p.~\pageref{fol:138}), infra f\textsuperscript{o}~70~r\textsuperscript{o} (p.~\pageref{fol:140}), f\textsuperscript{o}~70~v\textsuperscript{o} (p.~\pageref{fol:141}), \textit{et passim}, Appendicem I, doc.~5 (p.~\pageref{annexe:doc5}). — \textbf{6} (\textit{Ste-Marie}) cf. supra f\textsuperscript{o}~27~v\textsuperscript{o} (p.~\pageref{fol:55}), f\textsuperscript{o}~69~r\textsuperscript{o} (p.~\pageref{fol:138}), \textit{et passim}. — \textbf{7} (\textit{Oncle}) cf. supra f\textsuperscript{o}~27~r\textsuperscript{o} (p.~\pageref{fol:54}), f\textsuperscript{o}~66~v\textsuperscript{o} (p.~\pageref{fol:133}), \textit{et passim}.\par\vspace{0.55mm}
\noindent\rule{38mm}{0.4pt}\par\vspace{0.5mm}
\noindent\hspace*{5mm}\parbox{\dimexpr\textwidth-5mm\relax}{\fontsize{6.6pt}{7.7pt}\selectfont\color{black!84}Voyage ferroviaire du samedi 3 février 1917 de \textbf{Toulouse} à \textbf{Millau} en compagnie de Marguerite Calmels : salut en gare de Carcassonne au cousin l'\textbf{abbé Jean Crouzet}, halte à \textbf{Béziers} avec visite de l'ancienne cathédrale gothique \textbf{Saint-Nazaire} et de sa terrasse dominant l'Orb, puis arrivée nocturne en gare de Millau où les attend \textbf{Pierre Calmels}, frère de Marguerite.}\par}
\newpage
\phantomsection\label{fac:139}
% --- En-tête du fac-similé (Page impaire / droite : Photographie du manuscrit) ---
\noindent\rule{\textwidth}{0.5pt}\par\vspace{0.55mm}
\noindent\textbf{\fontsize{6.8pt}{7.8pt}\selectfont FAC-SIMILÉ EN VIS-À-VIS — FOLIO 69 VERSO} \hfill \textbf{\fontsize{6.8pt}{7.8pt}\selectfont FÉVRIER 1917}\par\vspace{0.25mm}
\noindent{\fontsize{6.8pt}{8.0pt}\selectfont\textit{Livre-Journal d'Hippolyte Combes} — 3 février 1917 (Toulouse, Carcassonne, Béziers, Bédarieux, Millau (Aveyron))}\par\vspace{0.55mm}
\noindent\rule{\textwidth}{0.35pt}\par\vspace*{\fill}
{\centering
  \setlength{\fboxsep}{1.2mm}%
  \setlength{\fboxrule}{0.4pt}%
  \fcolorbox{gray!55}{white}{%
    \includegraphics[width=0.95\textwidth,height=0.81\textheight,keepaspectratio]{\facsimiledir/20251126_110338.jpg}%
  }\par}
\vspace*{\fill}
\newpage

% =============================================================================
% VIS-A-VIS SPREAD : FOLIO 70 RECTO (20251126_110353.jpg)
% =============================================================================
\cleardoubleevenpage
\renewcommand{\currentmonthyear}{Février 1917}
\renewcommand{\currentfoliolabel}{Folio 70 recto}
\phantomsection\label{fol:140}
% --- En-tête du folio (Page paire / gauche : Transcription, apparat & registres de bas de page) ---
\noindent\rule{\textwidth}{0.5pt}\par\vspace{0.55mm}
\noindent\textbf{\footnotesize FOLIO 70 RECTO} \hfill \textbf{\footnotesize FÉVRIER 1917}\par\vspace{0.25mm}
\noindent{\fontsize{6.9pt}{8.0pt}\selectfont Relevé des dépenses du départ de Toulouse et du voyage par Bédarieux}\par\vspace{0.2mm}
\noindent{\fontsize{6.9pt}{8.0pt}\selectfont\textit{3 février 1917 (suite, comptes)} \hfill \textit{Toulouse, Béziers, Bédarieux, Millau}}\par\vspace{0.5mm}
\noindent\rule{\textwidth}{0.35pt}\par\vspace{0.8mm}
% --- Corps de la transcription diplomatique et apparat marginal en prolongement ---
{\renewcommand{\arraystretch}{0.87}%
\noindent\hspace*{3.5mm}\begin{tabular}{@{}p{87mm}@{\hspace{3.5mm}}p{39mm}@{}}
      \leavevmode\fontsize{6.9pt}{7.9pt}\selectfont\strut \underline{Sam. 3} — \hfill{} \underline{Févr. /17}\strut & \smash{\parbox[t]{39mm}{\fontsize{6.2pt}{7.0pt}\selectfont\raggedright\color{black!88}\textbf{1 Sam.} \textit{scil.} Samedi.}} \\
      \leavevmode\fontsize{6.9pt}{7.9pt}\selectfont\strut — Dentiste, M. de Lersay, 1 grand nettoyage\strut &  \\
      \leavevmode\fontsize{6.9pt}{7.9pt}\selectfont\strut et 3 plombages : 40 f = net : 25 f —\strut &  \\
      \leavevmode\fontsize{6.9pt}{7.9pt}\selectfont\strut — Broche A.P.C. : 3 f\strut &  \\
      \leavevmode\llap{\fontsize{6.4pt}{7.3pt}\selectfont\color{darkgray}5\hspace{1.8mm}}\fontsize{6.9pt}{7.9pt}\selectfont\strut — Trams à Toulouse : 20 c., + 20 c\strut & \smash{\parbox[t]{39mm}{\fontsize{6.2pt}{7.0pt}\selectfont\raggedright\color{black!88}\textbf{5 c.} \textit{abbr. pro} centimes.}} \\
      \leavevmode\fontsize{6.9pt}{7.9pt}\selectfont\strut — Enregt bagages de Margte : 10 c\strut & \smash{\parbox[t]{39mm}{\fontsize{6.2pt}{7.0pt}\selectfont\raggedright\color{black!88}\textbf{6 Enregt} \textit{abbr. pro} Enregistrement. \textbf{6 Margte} \textit{scil.} Marguerite (Calmels).}} \\
      \leavevmode\fontsize{6.9pt}{7.9pt}\selectfont\strut — Étrennes à Marie : 5 f —\strut &  \\
      \leavevmode\fontsize{6.9pt}{7.9pt}\selectfont\strut — Supplément pour m'être déclassé en 2e : 1 f\strut &  \\
      \leavevmode\fontsize{6.9pt}{7.9pt}\selectfont\strut — Billet 1/2 pl. Béziers – Millau : 1 f 80\strut & \smash{\parbox[t]{39mm}{\fontsize{6.2pt}{7.0pt}\selectfont\raggedright\color{black!88}\textbf{9 1/2} \textit{scil.} demi / moitié.}} \\
      \leavevmode\llap{\fontsize{6.4pt}{7.3pt}\selectfont\color{darkgray}10\hspace{1.8mm}}\fontsize{6.9pt}{7.9pt}\selectfont\strut — Buffet de Bédarieux : 6 f 75 + 25 c\strut &  \\
\end{tabular}}%
\par\vspace*{\fill}
% --- Registres de bas de page (1. Critique interne [si pertinent] ; 2. Loci paralleli en latin ; 3. Commentaire historique en retrait) ---
\noindent\rule{\textwidth}{0.5pt}\par\vspace{0.65mm}
{\fontsize{6.6pt}{7.7pt}\selectfont%
\noindent\normalfont \textbf{Feuillet comptable :} Page récapitulative des dépenses toulousaines du samedi 3 février 1917 (dentiste M.~de Lersay, billets de chemin de fer pour Millau) et du solde en caisse (\textit{En caisse : 244 f. 85}).\par\vspace{0.5mm}
\noindent\rule{22mm}{0.3pt}\par\vspace{0.45mm}
\noindent\normalfont\color{black!92} \textbf{5} (\textit{Toulouse}) cf. supra in tegumento (p.~\pageref{fol:1}), f\textsuperscript{o}~69~v\textsuperscript{o} (p.~\pageref{fol:139}), infra f\textsuperscript{o}~78~r\textsuperscript{o} (p.~\pageref{fol:156}), \textit{et passim}, Appendicem I, doc.~1 et 3 (p.~\pageref{annexe:doc1}). — \textbf{6} (\textit{Margte}) cf. supra f\textsuperscript{o}~1~r\textsuperscript{o} (p.~\pageref{fol:2}), f\textsuperscript{o}~69~v\textsuperscript{o} (p.~\pageref{fol:139}), infra f\textsuperscript{o}~70~v\textsuperscript{o} (p.~\pageref{fol:141}), f\textsuperscript{o}~71~r\textsuperscript{o} (p.~\pageref{fol:142}), \textit{et passim}. — \textbf{9} (\textit{Millau}) cf. supra in tegumento (p.~\pageref{fol:1}), f\textsuperscript{o}~69~v\textsuperscript{o} (p.~\pageref{fol:139}), infra f\textsuperscript{o}~70~v\textsuperscript{o} (p.~\pageref{fol:141}), f\textsuperscript{o}~71~r\textsuperscript{o} (p.~\pageref{fol:142}), \textit{et passim}.\par\vspace{0.55mm}
\noindent\rule{38mm}{0.4pt}\par\vspace{0.5mm}
\noindent\hspace*{5mm}\parbox{\dimexpr\textwidth-5mm\relax}{\fontsize{6.6pt}{7.7pt}\selectfont\color{black!84}Relevé comptable de la journée de voyage du samedi 3 février 1917 (complétant le f\textsuperscript{o}~69~v\textsuperscript{o}): honoraires réduits du chirurgien-dentiste toulousain \textbf{M.~de Lersay} (25~fr. au lieu de 40~fr.), étrennes de 5~fr. à la domestique \textbf{Marie} à Toulouse, billet militaire au quart de tarif de Béziers à Millau (1~fr.~80 sur le réseau du Midi) et dîner des voyageurs au buffet de la gare de bifurcation de \textbf{Bédarieux} (Hérault, 7~fr.).}\par}
\newpage
\phantomsection\label{fac:140}
% --- En-tête du fac-similé (Page impaire / droite : Photographie du manuscrit) ---
\noindent\rule{\textwidth}{0.5pt}\par\vspace{0.55mm}
\noindent\textbf{\fontsize{6.8pt}{7.8pt}\selectfont FAC-SIMILÉ EN VIS-À-VIS — FOLIO 70 RECTO} \hfill \textbf{\fontsize{6.8pt}{7.8pt}\selectfont FÉVRIER 1917}\par\vspace{0.25mm}
\noindent{\fontsize{6.8pt}{8.0pt}\selectfont\textit{Livre-Journal d'Hippolyte Combes} — 3 février 1917 (suite, comptes) (Toulouse, Béziers, Bédarieux, Millau)}\par\vspace{0.55mm}
\noindent\rule{\textwidth}{0.35pt}\par\vspace*{\fill}
{\centering
  \setlength{\fboxsep}{1.2mm}%
  \setlength{\fboxrule}{0.4pt}%
  \fcolorbox{gray!55}{white}{%
    \includegraphics[width=0.95\textwidth,height=0.81\textheight,keepaspectratio]{\facsimiledir/20251126_110353.jpg}%
  }\par}
\vspace*{\fill}
\newpage

% =============================================================================
% VIS-A-VIS SPREAD : FOLIO 70 VERSO (20251126_110406.jpg)
% =============================================================================
\cleardoubleevenpage
\renewcommand{\currentmonthyear}{Février 1917}
\renewcommand{\currentfoliolabel}{Folio 70 verso}
\phantomsection\label{fol:141}
% --- En-tête du folio (Page paire / gauche : Transcription, apparat & registres de bas de page) ---
\noindent\rule{\textwidth}{0.5pt}\par\vspace{0.55mm}
\noindent\textbf{\footnotesize FOLIO 70 VERSO} \hfill \textbf{\footnotesize FÉVRIER 1917}\par\vspace{0.25mm}
\noindent{\fontsize{6.9pt}{8.0pt}\selectfont Nuit d'émotion et de larmes à Millau au souvenir de ses parents et de ses sœurs}\par\vspace{0.2mm}
\noindent{\fontsize{6.9pt}{8.0pt}\selectfont\textit{4 février 1917 (nuit et matin)} \hfill \textit{Millau (10 rue Alfred-Guibert, Aveyron)}}\par\vspace{0.5mm}
\noindent\rule{\textwidth}{0.35pt}\par\vspace{0.8mm}
% --- Corps de la transcription diplomatique et apparat marginal en prolongement ---
{\renewcommand{\arraystretch}{0.93}%
\noindent\hspace*{3.5mm}\begin{tabular}{@{}p{87mm}@{\hspace{3.5mm}}p{39mm}@{}}
      \leavevmode\fontsize{7.3pt}{8.4pt}\selectfont\strut \underline{Dim. 4} \hfill{} \underline{Févr. /17}\strut & \smash{\parbox[t]{39mm}{\fontsize{6.2pt}{7.0pt}\selectfont\raggedright\color{black!88}\textbf{1 Dim.} \textit{scil.} Dimanche.}} \\
      \leavevmode\fontsize{7.3pt}{8.4pt}\selectfont\strut [Monté ?] chambre hier soir vers 11h 1/4.\strut & \smash{\parbox[t]{39mm}{\fontsize{6.2pt}{7.0pt}\selectfont\raggedright\color{black!88}\textbf{2 Monté} \textit{lect. incert. in ms.} \textbf{2 chambre} \textit{supra lin. in ms.}}} \\
      \leavevmode\fontsize{7.3pt}{8.4pt}\selectfont\strut Avant d'éteindre bougie, ouvert volet de fer et\strut &  \\
      \leavevmode\fontsize{7.3pt}{8.4pt}\selectfont\strut regardé paysage. 1 peu à droite, et tout en\strut &  \\
      \leavevmode\llap{\fontsize{6.4pt}{7.3pt}\selectfont\color{darkgray}5\hspace{1.8mm}}\fontsize{7.3pt}{8.4pt}\selectfont\strut haut, pleine lune vif éclat\strut &  \\
      \leavevmode\fontsize{7.3pt}{8.4pt}\selectfont\strut au ciel, qq étoiles.\strut & \smash{\parbox[t]{39mm}{\fontsize{6.2pt}{7.0pt}\selectfont\raggedright\color{black!88}\textbf{6 qq} \textit{abbr. pro} quelque(s).}} \\
      \leavevmode\fontsize{7.3pt}{8.4pt}\selectfont\strut Silence impressionnant vallon natal.\strut &  \\
      \leavevmode\fontsize{7.3pt}{8.4pt}\selectfont\strut Couché depuis 1/2 h., quand souvenirs venus\strut & \smash{\parbox[t]{39mm}{\fontsize{6.2pt}{7.0pt}\selectfont\raggedright\color{black!88}\textbf{8 1/2} \textit{scil.} demi / moitié.}} \\
      \leavevmode\fontsize{7.3pt}{8.4pt}\selectfont\strut en foule et obsédé. Pensé à Papa,\strut &  \\
      \leavevmode\llap{\fontsize{6.4pt}{7.3pt}\selectfont\color{darkgray}10\hspace{1.8mm}}\fontsize{7.3pt}{8.4pt}\selectfont\strut Maman, Sœurs — que j'ai tant aimés.\strut &  \\
      \leavevmode\fontsize{7.3pt}{8.4pt}\selectfont\strut Peu à peu, gagné par l'émotion, et\strut &  \\
      \leavevmode\fontsize{7.3pt}{8.4pt}\selectfont\strut sangloté.\strut &  \\
      \leavevmode\fontsize{7.3pt}{8.4pt}\selectfont\strut Insomnie. Écouté sonner les heures au\strut &  \\
      \leavevmode\fontsize{7.3pt}{8.4pt}\selectfont\strut petit cadran cuivre posé sur la cheminée\strut &  \\
      \leavevmode\llap{\fontsize{6.4pt}{7.3pt}\selectfont\color{darkgray}15\hspace{1.8mm}}\fontsize{7.3pt}{8.4pt}\selectfont\strut et, chambre voisine, Germaine tousser.\strut &  \\
      \leavevmode\fontsize{7.3pt}{8.4pt}\selectfont\strut Pu reposer un peu sur le matin.\strut &  \\
      \leavevmode\fontsize{7.3pt}{8.4pt}\selectfont\strut Levé vers 7h.\strut &  \\
      \leavevmode\fontsize{7.3pt}{8.4pt}\selectfont\strut Impressions confort, bien-être, calme.\strut &  \\
      \leavevmode\fontsize{7.3pt}{8.4pt}\selectfont\strut Déjeuné av. Margte.\strut & \smash{\parbox[t]{39mm}{\fontsize{6.2pt}{7.0pt}\selectfont\raggedright\color{black!88}\textbf{19 Margte} \textit{scil.} Marguerite (Calmels). — \textbf{20 Mte} \textit{scil.} Marguerite (Calmels).}} \\
      \leavevmode\llap{\fontsize{6.4pt}{7.3pt}\selectfont\color{darkgray}20\hspace{1.8mm}}\fontsize{7.3pt}{8.4pt}\selectfont\strut Mte m'a laissé vers 11h pour aller à\strut &  \\
      \leavevmode\fontsize{7.3pt}{8.4pt}\selectfont\strut la messe. Alors monté chambre. Rêvé\strut &  \\
      \leavevmode\fontsize{7.3pt}{8.4pt}\selectfont\strut à mon aise. Soleil rendait appartt\strut & \smash{\parbox[t]{39mm}{\fontsize{6.2pt}{7.0pt}\selectfont\raggedright\color{black!88}\textbf{22 appartt} \textit{abbr. pro} appartement.}} \\
\end{tabular}}%
\par\vspace*{\fill}
% --- Registres de bas de page (1. Critique interne [si pertinent] ; 2. Loci paralleli en latin ; 3. Commentaire historique en retrait) ---
\noindent\rule{\textwidth}{0.5pt}\par\vspace{0.65mm}
{\fontsize{6.6pt}{7.7pt}\selectfont%
\noindent\normalfont\color{black!92} \textbf{15} (\textit{Germaine}) cf. supra f\textsuperscript{o}~65~v\textsuperscript{o} (p.~\pageref{fol:131}), infra f\textsuperscript{o}~71~v\textsuperscript{o} (p.~\pageref{fol:143}). — \textbf{19} (\textit{Margte}) cf. supra f\textsuperscript{o}~1~r\textsuperscript{o} (p.~\pageref{fol:2}), f\textsuperscript{o}~70~r\textsuperscript{o} (p.~\pageref{fol:140}), infra f\textsuperscript{o}~71~r\textsuperscript{o} (p.~\pageref{fol:142}), f\textsuperscript{o}~71~v\textsuperscript{o} (p.~\pageref{fol:143}), \textit{et passim}.\par\vspace{0.55mm}
\noindent\rule{38mm}{0.4pt}\par\vspace{0.5mm}
\noindent\hspace*{5mm}\parbox{\dimexpr\textwidth-5mm\relax}{\fontsize{6.6pt}{7.7pt}\selectfont\color{black!84}Première nuit d'Hippolyte Combes dans la demeure des Calmels au \textbf{10, rue Alfred-Guibert à Millau} (nuit du 3 au 4 février 1917).}\par}
\newpage
\phantomsection\label{fac:141}
% --- En-tête du fac-similé (Page impaire / droite : Photographie du manuscrit) ---
\noindent\rule{\textwidth}{0.5pt}\par\vspace{0.55mm}
\noindent\textbf{\fontsize{6.8pt}{7.8pt}\selectfont FAC-SIMILÉ EN VIS-À-VIS — FOLIO 70 VERSO} \hfill \textbf{\fontsize{6.8pt}{7.8pt}\selectfont FÉVRIER 1917}\par\vspace{0.25mm}
\noindent{\fontsize{6.8pt}{8.0pt}\selectfont\textit{Livre-Journal d'Hippolyte Combes} — 4 février 1917 (nuit et matin) (Millau (10 rue Alfred-Guibert, Aveyron))}\par\vspace{0.55mm}
\noindent\rule{\textwidth}{0.35pt}\par\vspace*{\fill}
{\centering
  \setlength{\fboxsep}{1.2mm}%
  \setlength{\fboxrule}{0.4pt}%
  \fcolorbox{gray!55}{white}{%
    \includegraphics[width=0.95\textwidth,height=0.81\textheight,keepaspectratio]{\facsimiledir/20251126_110406.jpg}%
  }\par}
\vspace*{\fill}
\newpage

% =============================================================================
% VIS-A-VIS SPREAD : FOLIO 71 RECTO (20251126_110422.jpg)
% =============================================================================
\cleardoubleevenpage
\renewcommand{\currentmonthyear}{Février 1917}
\renewcommand{\currentfoliolabel}{Folio 71 recto}
\phantomsection\label{fol:142}
% --- En-tête du folio (Page paire / gauche : Transcription, apparat & registres de bas de page) ---
\noindent\rule{\textwidth}{0.5pt}\par\vspace{0.55mm}
\noindent\textbf{\footnotesize FOLIO 71 RECTO} \hfill \textbf{\footnotesize FÉVRIER 1917}\par\vspace{0.25mm}
\noindent{\fontsize{6.9pt}{8.0pt}\selectfont Après-midi musical avec Marguerite Calmels, visite à l'oncle et la tante Thérèse (Layral) et arrivée de Mimi à Millau}\par\vspace{0.2mm}
\noindent{\fontsize{6.9pt}{8.0pt}\selectfont\textit{4 février 1917 (suite)} \hfill \textit{Millau (Aveyron)}}\par\vspace{0.5mm}
\noindent\rule{\textwidth}{0.35pt}\par\vspace{0.8mm}
% --- Corps de la transcription diplomatique et apparat marginal en prolongement ---
{\renewcommand{\arraystretch}{0.93}%
\noindent\hspace*{3.5mm}\begin{tabular}{@{}p{87mm}@{\hspace{3.5mm}}p{39mm}@{}}
      \leavevmode\fontsize{7.3pt}{8.4pt}\selectfont\strut \underline{Dim. 4 (suite)} \hfill{} \underline{Fév. /17}\strut & \smash{\parbox[t]{39mm}{\fontsize{6.2pt}{7.0pt}\selectfont\raggedright\color{black!88}\textbf{1 Dim.} \textit{scil.} Dimanche.}} \\
      \leavevmode\fontsize{7.3pt}{8.4pt}\selectfont\strut plus clair \& plus gai.\strut &  \\
      \leavevmode\fontsize{7.3pt}{8.4pt}\selectfont\strut Tout l'après-midi, av. Mte. Elle m'a\strut & \smash{\parbox[t]{39mm}{\fontsize{6.2pt}{7.0pt}\selectfont\raggedright\color{black!88}\textbf{3 Mte} \textit{scil.} Marguerite (Calmels). — \textbf{4 qq} \textit{abbr. pro} quelque(s).}} \\
      \leavevmode\fontsize{7.3pt}{8.4pt}\selectfont\strut joué qq morceaux et chanté 1 peu.\strut &  \\
      \leavevmode\llap{\fontsize{6.4pt}{7.3pt}\selectfont\color{darkgray}5\hspace{1.8mm}}\fontsize{7.3pt}{8.4pt}\selectfont\strut A 6h 15, cédant aux instances de tous,\strut &  \\
      \leavevmode\fontsize{7.3pt}{8.4pt}\selectfont\strut visite à l'oncle \& tante Thérèse (Layral).\strut & \smash{\parbox[t]{39mm}{\fontsize{6.2pt}{7.0pt}\selectfont\raggedright\color{black!88}\textbf{6 (Layral)} \textit{scil.} Charles et Thérèse Layrals.}} \\
      \leavevmode\fontsize{7.3pt}{8.4pt}\selectfont\strut Temps très sec. Rue déserte.\strut &  \\
      \leavevmode\fontsize{7.3pt}{8.4pt}\selectfont\strut Après dîner, prolongé causerie j'à\strut & \smash{\parbox[t]{39mm}{\fontsize{6.2pt}{7.0pt}\selectfont\raggedright\color{black!88}\textbf{8 j'à} \textit{abbr. pro} jusqu'à.}} \\
      \leavevmode\fontsize{7.3pt}{8.4pt}\selectfont\strut 11h 30 — av. Mimi, arrivée train 6h.\strut &  \\
      \leavevmode\llap{\fontsize{6.4pt}{7.3pt}\selectfont\color{darkgray}10\hspace{1.8mm}}\fontsize{7.3pt}{8.4pt}\selectfont\strut Ds ma chambre, à 11h 40.\strut & \smash{\parbox[t]{39mm}{\fontsize{6.2pt}{7.0pt}\selectfont\raggedright\color{black!88}\textbf{10 Ds} \textit{abbr. pro} Dans.}} \\
      \leavevmode\fontsize{7.3pt}{8.4pt}\selectfont\strut Ciel nuageux. Nuit très froide.\strut &  \\
      \leavevmode\fontsize{7.3pt}{8.4pt}\selectfont\strut Ce matin, en m'éveillant, charmé\strut &  \\
      \leavevmode\fontsize{7.3pt}{8.4pt}\selectfont\strut par sonneries cloches, églises et\strut &  \\
      \leavevmode\fontsize{7.3pt}{8.4pt}\selectfont\strut chapelles ville —\strut &  \\
      \leavevmode\llap{\fontsize{6.4pt}{7.3pt}\selectfont\color{darkgray}15\hspace{1.8mm}}\fontsize{7.3pt}{8.4pt}\selectfont\strut (Cf. Journal, feuille spéciale)\strut &  \\
\end{tabular}}%
\par\vspace*{\fill}
% --- Registres de bas de page (1. Critique interne [si pertinent] ; 2. Loci paralleli en latin ; 3. Commentaire historique en retrait) ---
\noindent\rule{\textwidth}{0.5pt}\par\vspace{0.65mm}
{\fontsize{6.6pt}{7.7pt}\selectfont%
\noindent\normalfont\color{black!92} \textbf{3} (\textit{Mte}) cf. supra f\textsuperscript{o}~1~r\textsuperscript{o} (p.~\pageref{fol:2}), f\textsuperscript{o}~70~v\textsuperscript{o} (p.~\pageref{fol:141}), infra f\textsuperscript{o}~71~v\textsuperscript{o} (p.~\pageref{fol:143}), \textit{et passim}, Appendicem I, doc.~5 (p.~\pageref{annexe:doc5}).\par\vspace{0.55mm}
\noindent\rule{38mm}{0.4pt}\par\vspace{0.5mm}
\noindent\hspace*{5mm}\parbox{\dimexpr\textwidth-5mm\relax}{\fontsize{6.6pt}{7.7pt}\selectfont\color{black!84}\textbf{Thérèse Cartairade} et son époux \textbf{M\textsuperscript{e}~Charles Layrals}, notaire à Millau (parents du notaire Louis Layrals, \dag~1961, et du magistrat Jean Layrals, \dag~1947), auxquels Hippolyte Combes et Marguerite Calmels rendent visite le dimanche 4 février 1917 à 18~h~15, sont la tante et l'oncle maternels de Marguerite. \textbf{« Mimi »}, arrivée par le train de 18~h, est le surnom familial d'\textbf{Anne-Marie Calmels} (1888–1975), sœur de Marguerite et future épouse de l'avoué millavois \textbf{Joseph Delcamps}.}\par}
\newpage
\phantomsection\label{fac:142}
% --- En-tête du fac-similé (Page impaire / droite : Photographie du manuscrit) ---
\noindent\rule{\textwidth}{0.5pt}\par\vspace{0.55mm}
\noindent\textbf{\fontsize{6.8pt}{7.8pt}\selectfont FAC-SIMILÉ EN VIS-À-VIS — FOLIO 71 RECTO} \hfill \textbf{\fontsize{6.8pt}{7.8pt}\selectfont FÉVRIER 1917}\par\vspace{0.25mm}
\noindent{\fontsize{6.8pt}{8.0pt}\selectfont\textit{Livre-Journal d'Hippolyte Combes} — 4 février 1917 (suite) (Millau (Aveyron))}\par\vspace{0.55mm}
\noindent\rule{\textwidth}{0.35pt}\par\vspace*{\fill}
{\centering
  \setlength{\fboxsep}{1.2mm}%
  \setlength{\fboxrule}{0.4pt}%
  \fcolorbox{gray!55}{white}{%
    \includegraphics[width=0.95\textwidth,height=0.81\textheight,keepaspectratio]{\facsimiledir/20251126_110422.jpg}%
  }\par}
\vspace*{\fill}
\newpage

% =============================================================================
% VIS-A-VIS SPREAD : FOLIO 71 VERSO (20251126_110441.jpg)
% =============================================================================
\cleardoubleevenpage
\renewcommand{\currentmonthyear}{Février 1917}
\renewcommand{\currentfoliolabel}{Folio 71 verso}
\phantomsection\label{fol:143}
% --- En-tête du folio (Page paire / gauche : Transcription, apparat & registres de bas de page) ---
\noindent\rule{\textwidth}{0.5pt}\par\vspace{0.55mm}
\noindent\textbf{\footnotesize FOLIO 71 VERSO} \hfill \textbf{\footnotesize FÉVRIER 1917}\par\vspace{0.25mm}
\noindent{\fontsize{6.9pt}{8.0pt}\selectfont Recueillement au cimetière de Millau avec Marguerite, adieux déchirants et arrivée à Saint-Flour}\par\vspace{0.2mm}
\noindent{\fontsize{6.9pt}{8.0pt}\selectfont\textit{5 février 1917} \hfill \textit{Millau, Séverac, Marvejols, Saint-Flour (Aveyron / Lozère / Cantal)}}\par\vspace{0.5mm}
\noindent\rule{\textwidth}{0.35pt}\par\vspace{0.8mm}
% --- Corps de la transcription diplomatique et apparat marginal en prolongement ---
{\renewcommand{\arraystretch}{0.87}%
\noindent\hspace*{3.5mm}\begin{tabular}{@{}p{87mm}@{\hspace{3.5mm}}p{39mm}@{}}
      \leavevmode\fontsize{6.9pt}{7.9pt}\selectfont\strut \underline{Lundi 5} \hfill{} \underline{Févr. /17}\strut &  \\
      \leavevmode\fontsize{6.9pt}{7.9pt}\selectfont\strut Sur pied 6h —\strut &  \\
      \leavevmode\fontsize{6.9pt}{7.9pt}\selectfont\strut Me suis complu rêver gd chambre.\strut & \smash{\parbox[t]{39mm}{\fontsize{6.2pt}{7.0pt}\selectfont\raggedright\color{black!88}\textbf{3 gd} \textit{abbr. pro} grande (devant « chambre »). — \textbf{4 ts} \textit{abbr. pro} tous. \textbf{4 Mte} \textit{scil.} Marguerite (Calmels).}} \\
      \leavevmode\fontsize{6.9pt}{7.9pt}\selectfont\strut Après déjeuner, sorti av. Mte et ts 2\strut &  \\
      \leavevmode\llap{\fontsize{6.4pt}{7.3pt}\selectfont\color{darkgray}5\hspace{1.8mm}}\fontsize{6.9pt}{7.9pt}\selectfont\strut seuls venus cimetière — —\strut &  \\
      \leavevmode\fontsize{6.9pt}{7.9pt}\selectfont\strut Pris renseignements gare. Décidé ne partir\strut &  \\
      \leavevmode\fontsize{6.9pt}{7.9pt}\selectfont\strut qu'à 5h 33 soir.\strut &  \\
      \leavevmode\fontsize{6.9pt}{7.9pt}\selectfont\strut Terminé journée + douce intimité.\strut &  \\
      \leavevmode\fontsize{6.9pt}{7.9pt}\selectfont\strut Quitté maison vers 5h s. Adieux, déchirants.\strut & \smash{\parbox[t]{39mm}{\fontsize{6.2pt}{7.0pt}\selectfont\raggedright\color{black!88}\textbf{9 s.} \textit{abbr. pro} soir.}} \\
      \leavevmode\llap{\fontsize{6.4pt}{7.3pt}\selectfont\color{darkgray}10\hspace{1.8mm}}\fontsize{6.9pt}{7.9pt}\selectfont\strut Pierre m'attendait devant gare.\strut &  \\
      \leavevmode\fontsize{6.9pt}{7.9pt}\selectfont\strut Train venant Béziers, long retard.\strut &  \\
      \leavevmode\fontsize{6.9pt}{7.9pt}\selectfont\strut Enfin, installé compartiment 2e cl.\strut &  \\
      \leavevmode\fontsize{6.9pt}{7.9pt}\selectfont\strut et départ 6h 15. —\strut &  \\
      \leavevmode\fontsize{6.9pt}{7.9pt}\selectfont\strut Arrêt à ttes les stations.\strut & \smash{\parbox[t]{39mm}{\fontsize{6.2pt}{7.0pt}\selectfont\raggedright\color{black!88}\textbf{14 ttes} \textit{abbr. pro} toutes.}} \\
      \leavevmode\llap{\fontsize{6.4pt}{7.3pt}\selectfont\color{darkgray}15\hspace{1.8mm}}\fontsize{6.9pt}{7.9pt}\selectfont\strut Neige partout.\strut &  \\
      \leavevmode\fontsize{6.9pt}{7.9pt}\selectfont\strut Arrivé Saint-Flour à minuit 30 s.\strut &  \\
      \cline{1-1}
      \leavevmode\fontsize{6.9pt}{7.9pt}\selectfont\strut [ — Remboursé à Mte cadeaux enfants Vincens : 6 f\strut & \smash{\parbox[t]{39mm}{\fontsize{6.2pt}{7.0pt}\selectfont\raggedright\color{black!88}\textbf{17} \textit{in marg. cum unco} (comptes). \textbf{17 Mte} \textit{scil.} Marguerite (Calmels). — \textbf{18 Mte} \textit{scil.} Marguerite (Calmels). \textbf{18 Gift} \textit{angl. in ms.}}} \\
      \leavevmode\fontsize{6.9pt}{7.9pt}\selectfont\strut Gift à la bonne de Mte : 3 f — ]\strut &  \\
\end{tabular}}%
\par\vspace*{\fill}
% --- Registres de bas de page (1. Critique interne [si pertinent] ; 2. Loci paralleli en latin ; 3. Commentaire historique en retrait) ---
\noindent\rule{\textwidth}{0.5pt}\par\vspace{0.65mm}
{\fontsize{6.6pt}{7.7pt}\selectfont%
\noindent\normalfont\color{black!92} \textbf{4} (\textit{Mte}) cf. supra f\textsuperscript{o}~1~r\textsuperscript{o} (p.~\pageref{fol:2}), f\textsuperscript{o}~71~r\textsuperscript{o} (p.~\pageref{fol:142}), \textit{et passim}, Appendicem I, doc.~5 (p.~\pageref{annexe:doc5}). — \textbf{10} (\textit{Pierre}) cf. supra f\textsuperscript{o}~69~v\textsuperscript{o} (p.~\pageref{fol:139}). — \textbf{17} (\textit{Vincens}) cf. supra f\textsuperscript{o}~65~v\textsuperscript{o} (p.~\pageref{fol:131}), f\textsuperscript{o}~70~v\textsuperscript{o} (p.~\pageref{fol:141}).\par\vspace{0.55mm}
\noindent\rule{38mm}{0.4pt}\par\vspace{0.5mm}
\noindent\hspace*{5mm}\parbox{\dimexpr\textwidth-5mm\relax}{\fontsize{6.6pt}{7.7pt}\selectfont\color{black!84}Le lundi 5 février 1917 après-midi, Hippolyte Combes et Marguerite se rendent seuls au \textbf{cimetière de Millau} (où reposait notamment le père d'Hippolyte Combes, François-Hippolyte Combes, mort à Millau le 11 août 1899). Avant de quitter le 10 rue Alfred-Guibert et de retrouver son futur beau-frère \textbf{Pierre Calmels} devant la gare de Millau, Hippolyte Combes rembourse 6~fr. de cadeaux pour les \textbf{enfants Vincens}.}\par}
\newpage
\phantomsection\label{fac:143}
% --- En-tête du fac-similé (Page impaire / droite : Photographie du manuscrit) ---
\noindent\rule{\textwidth}{0.5pt}\par\vspace{0.55mm}
\noindent\textbf{\fontsize{6.8pt}{7.8pt}\selectfont FAC-SIMILÉ EN VIS-À-VIS — FOLIO 71 VERSO} \hfill \textbf{\fontsize{6.8pt}{7.8pt}\selectfont FÉVRIER 1917}\par\vspace{0.25mm}
\noindent{\fontsize{6.8pt}{8.0pt}\selectfont\textit{Livre-Journal d'Hippolyte Combes} — 5 février 1917 (Millau, Séverac, Marvejols, Saint-Flour (Aveyron / Lozère / Cantal))}\par\vspace{0.55mm}
\noindent\rule{\textwidth}{0.35pt}\par\vspace*{\fill}
{\centering
  \setlength{\fboxsep}{1.2mm}%
  \setlength{\fboxrule}{0.4pt}%
  \fcolorbox{gray!55}{white}{%
    \includegraphics[width=0.95\textwidth,height=0.81\textheight,keepaspectratio]{\facsimiledir/20251126_110441.jpg}%
  }\par}
\vspace*{\fill}
\newpage

% =============================================================================
% VIS-A-VIS SPREAD : FOLIO 72 RECTO (20251126_110549.jpg)
% =============================================================================
\cleardoubleevenpage
\renewcommand{\currentmonthyear}{Février 1917}
\renewcommand{\currentfoliolabel}{Folio 72 recto}
\phantomsection\label{fol:144}
% --- En-tête du folio (Page paire / gauche : Transcription, apparat & registres de bas de page) ---
\noindent\rule{\textwidth}{0.5pt}\par\vspace{0.55mm}
\noindent\textbf{\footnotesize FOLIO 72 RECTO} \hfill \textbf{\footnotesize FÉVRIER 1917}\par\vspace{0.25mm}
\noindent{\fontsize{6.9pt}{8.0pt}\selectfont Nuit glaciale en gare de Saint-Flour et traversée de l'Auvergne sous les flocons}\par\vspace{0.2mm}
\noindent{\fontsize{6.9pt}{8.0pt}\selectfont\textit{6 février 1917 (matin et après-midi)} \hfill \textit{Saint-Flour, Brioude, Clermont-Ferrand, Gannat, Gien, Montargis}}\par\vspace{0.5mm}
\noindent\rule{\textwidth}{0.35pt}\par\vspace{0.8mm}
% --- Corps de la transcription diplomatique et apparat marginal en prolongement ---
{\renewcommand{\arraystretch}{0.93}%
\noindent\hspace*{3.5mm}\begin{tabular}{@{}p{87mm}@{\hspace{3.5mm}}p{39mm}@{}}
      \leavevmode\fontsize{7.3pt}{8.4pt}\selectfont\strut \underline{Mardi 6.} \hfill{} \underline{Février 1917}\strut &  \\
      \leavevmode\fontsize{7.3pt}{8.4pt}\selectfont\strut Nuit extt froide. Temps très dur.\strut & \smash{\parbox[t]{39mm}{\fontsize{6.2pt}{7.0pt}\selectfont\raggedright\color{black!88}\textbf{2 extt} \textit{abbr. pro} extrêmement.}} \\
      \leavevmode\fontsize{7.3pt}{8.4pt}\selectfont\strut D'abord salle d'attente, allongé banquette, puis\strut &  \\
      \leavevmode\fontsize{7.3pt}{8.4pt}\selectfont\strut compartiment 1re cl. train sur voie. Pieds gelés.\strut &  \\
      \leavevmode\llap{\fontsize{6.4pt}{7.3pt}\selectfont\color{darkgray}5\hspace{1.8mm}}\fontsize{7.3pt}{8.4pt}\selectfont\strut Partis enfin 5h 20. 15 m. de retard. En 2e cl.\strut &  \\
      \leavevmode\fontsize{7.3pt}{8.4pt}\selectfont\strut Voyagé av. qq permissionnaires et dames.\strut & \smash{\parbox[t]{39mm}{\fontsize{6.2pt}{7.0pt}\selectfont\raggedright\color{black!88}\textbf{6 qq} \textit{abbr. pro} quelque(s).}} \\
      \leavevmode\fontsize{7.3pt}{8.4pt}\selectfont\strut Champs couverts neige.\strut &  \\
      \leavevmode\fontsize{7.3pt}{8.4pt}\selectfont\strut A 7h 30, à Brioude. Correspce, gd retard.\strut & \smash{\parbox[t]{39mm}{\fontsize{6.2pt}{7.0pt}\selectfont\raggedright\color{black!88}\textbf{8 gd} \textit{abbr. pro} grand.}} \\
      \leavevmode\fontsize{7.3pt}{8.4pt}\selectfont\strut Voies, épaisse couche de neige.\strut &  \\
      \leavevmode\llap{\fontsize{6.4pt}{7.3pt}\selectfont\color{darkgray}10\hspace{1.8mm}}\fontsize{7.3pt}{8.4pt}\selectfont\strut Quitté Brioude à 8h 30. Enfin dans un\strut &  \\
      \leavevmode\fontsize{7.3pt}{8.4pt}\selectfont\strut grand train P.L.M., express.\strut & \smash{\parbox[t]{39mm}{\fontsize{6.2pt}{7.0pt}\selectfont\raggedright\color{black!88}\textbf{11 P.L.M.} \textit{scil.} Paris-Lyon-Méditerranée.}} \\
      \leavevmode\fontsize{7.3pt}{8.4pt}\selectfont\strut Clermont, à 10h.\strut &  \\
      \leavevmode\fontsize{7.3pt}{8.4pt}\selectfont\strut Gannat, 11h.\strut &  \\
      \leavevmode\fontsize{7.3pt}{8.4pt}\selectfont\strut Ciel bas \& mat, se confondant à\strut &  \\
      \leavevmode\llap{\fontsize{6.4pt}{7.3pt}\selectfont\color{darkgray}15\hspace{1.8mm}}\fontsize{7.3pt}{8.4pt}\selectfont\strut l'horizon proche avec campagne feutrée de\strut &  \\
      \leavevmode\fontsize{7.3pt}{8.4pt}\selectfont\strut blanc. Flocons serrés et menus tombent sans\strut &  \\
      \leavevmode\fontsize{7.3pt}{8.4pt}\selectfont\strut discontinuer, voilant l'atmosphère.\strut &  \\
      \leavevmode\fontsize{7.3pt}{8.4pt}\selectfont\strut Roulé tout le jour.\strut &  \\
      \leavevmode\fontsize{7.3pt}{8.4pt}\selectfont\strut A Gien, 3h 45. Montargis, 4h 20.\strut &  \\
\end{tabular}}%
\par\vspace*{\fill}
% --- Registres de bas de page (1. Critique interne [si pertinent] ; 2. Loci paralleli en latin ; 3. Commentaire historique en retrait) ---
\noindent\rule{\textwidth}{0.5pt}\par\vspace{0.65mm}
{\fontsize{6.6pt}{7.7pt}\selectfont%
\noindent\normalfont\color{black!92} \textbf{1} (\textit{de eodem commorationis loco}) cf. supra f\textsuperscript{o}~71~v\textsuperscript{o} (p.~\pageref{fol:143}).\par\vspace{0.55mm}
\noindent\rule{38mm}{0.4pt}\par\vspace{0.5mm}
\noindent\hspace*{5mm}\parbox{\dimexpr\textwidth-5mm\relax}{\fontsize{6.6pt}{7.7pt}\selectfont\color{black!84}Voyage de retour de Millau vers Paris à travers le Massif central enneigé le mardi 6 février 1917: après une nuit glaciale d'attente en gare de \textbf{Saint-Flour}, Hippolyte Combes gagne par Neussargues et Arvant la gare de \textbf{Brioude} (Haute-Loire, 7~h~30), où il rejoint à 8~h~30 la grande ligne du \textbf{P.-L.-M.} (\textit{Paris--Lyon--Méditerranée} par le Bourbonnais) et remonte vers Paris par \textbf{Clermont-Ferrand} (10~h), \textbf{Gannat} (11~h), \textbf{Gien} (15~h~45) et \textbf{Montargis} (16~h~20).}\par}
\newpage
\phantomsection\label{fac:144}
% --- En-tête du fac-similé (Page impaire / droite : Photographie du manuscrit) ---
\noindent\rule{\textwidth}{0.5pt}\par\vspace{0.55mm}
\noindent\textbf{\fontsize{6.8pt}{7.8pt}\selectfont FAC-SIMILÉ EN VIS-À-VIS — FOLIO 72 RECTO} \hfill \textbf{\fontsize{6.8pt}{7.8pt}\selectfont FÉVRIER 1917}\par\vspace{0.25mm}
\noindent{\fontsize{6.8pt}{8.0pt}\selectfont\textit{Livre-Journal d'Hippolyte Combes} — 6 février 1917 (matin et après-midi) (Saint-Flour, Brioude, Clermont-Ferrand, Gannat, Gien, Montargis)}\par\vspace{0.55mm}
\noindent\rule{\textwidth}{0.35pt}\par\vspace*{\fill}
{\centering
  \setlength{\fboxsep}{1.2mm}%
  \setlength{\fboxrule}{0.4pt}%
  \fcolorbox{gray!55}{white}{%
    \includegraphics[width=0.95\textwidth,height=0.81\textheight,keepaspectratio]{\facsimiledir/20251126_110549.jpg}%
  }\par}
\vspace*{\fill}
\newpage

% =============================================================================
% VIS-A-VIS SPREAD : FOLIO 72 VERSO (20251126_110601.jpg)
% =============================================================================
\cleardoubleevenpage
\renewcommand{\currentmonthyear}{Février 1917}
\renewcommand{\currentfoliolabel}{Folio 72 verso}
\phantomsection\label{fol:145}
% --- En-tête du folio (Page paire / gauche : Transcription, apparat & registres de bas de page) ---
\noindent\rule{\textwidth}{0.5pt}\par\vspace{0.55mm}
\noindent\textbf{\footnotesize FOLIO 72 VERSO} \hfill \textbf{\footnotesize FÉVRIER 1917}\par\vspace{0.25mm}
\noindent{\fontsize{6.9pt}{8.0pt}\selectfont Traversée de Paris en métro de la gare de Lyon à la gare du Nord et train de nuit vers Amiens}\par\vspace{0.2mm}
\noindent{\fontsize{6.9pt}{8.0pt}\selectfont\textit{6 février 1917 (suite, soir)} \hfill \textit{Corbeil-Essonnes, Paris (Gare de Lyon, Bastille, Gare du Nord)}}\par\vspace{0.5mm}
\noindent\rule{\textwidth}{0.35pt}\par\vspace{0.8mm}
% --- Corps de la transcription diplomatique et apparat marginal en prolongement ---
{\renewcommand{\arraystretch}{1.00}%
\noindent\hspace*{3.5mm}\begin{tabular}{@{}p{87mm}@{\hspace{3.5mm}}p{39mm}@{}}
      \leavevmode\fontsize{7.7pt}{9.0pt}\selectfont\strut \underline{Mardi 6 (suite)} \hfill{} \raisebox{0.5ex}{\underline{Févr. /17}}\strut &  \\
      \leavevmode\fontsize{7.7pt}{9.0pt}\selectfont\strut A mesure que ns remontons vers Paris, bcp\strut & \smash{\parbox[t]{39mm}{\fontsize{6.2pt}{7.0pt}\selectfont\raggedright\color{black!88}\textbf{2 ns} \textit{abbr. pro} nous. \textbf{2 bcp} \textit{abbr. pro} beaucoup.}} \\
      \leavevmode\fontsize{7.7pt}{9.0pt}\selectfont\strut moins de neige. Sol durci par fortes gelées.\strut &  \\
      \leavevmode\fontsize{7.7pt}{9.0pt}\selectfont\strut Après-midi, soleil brillé un peu. Rayons de +\strut &  \\
      \leavevmode\llap{\fontsize{6.4pt}{7.3pt}\selectfont\color{darkgray}5\hspace{1.8mm}}\fontsize{7.7pt}{9.0pt}\selectfont\strut en + faibles.\strut &  \\
      \leavevmode\fontsize{7.7pt}{9.0pt}\selectfont\strut Nuit s'annonce très froide.\strut &  \\
      \leavevmode\fontsize{7.7pt}{9.0pt}\selectfont\strut Corbeil-Essonnes, 5h 50 s. L'obscurité va\strut & \smash{\parbox[t]{39mm}{\fontsize{6.2pt}{7.0pt}\selectfont\raggedright\color{black!88}\textbf{7 s.} \textit{abbr. pro} soir.}} \\
      \leavevmode\fontsize{7.7pt}{9.0pt}\selectfont\strut augmentant très vite.\strut &  \\
      \leavevmode\fontsize{7.7pt}{9.0pt}\selectfont\strut \underline{A Paris}, vers 6h 30 s. —\strut & \smash{\parbox[t]{39mm}{\fontsize{6.2pt}{7.0pt}\selectfont\raggedright\color{black!88}\textbf{9 s.} \textit{abbr. pro} soir.}} \\
      \leavevmode\llap{\fontsize{6.4pt}{7.3pt}\selectfont\color{darkgray}10\hspace{1.8mm}}\fontsize{7.7pt}{9.0pt}\selectfont\strut Pris métro gare Lyon, changé à Bastille,\strut &  \\
      \leavevmode\fontsize{7.7pt}{9.0pt}\selectfont\strut et arrivé gare Nord assez à temps pour\strut &  \\
      \leavevmode\fontsize{7.7pt}{9.0pt}\selectfont\strut prendre train 7h 19. En 1re cl., avec type de\strut &  \\
      \leavevmode\fontsize{7.7pt}{9.0pt}\selectfont\strut soldats, fort curieux à observer et\strut &  \\
      \leavevmode\fontsize{7.7pt}{9.0pt}\selectfont\strut écouter.\strut &  \\
      \leavevmode\llap{\fontsize{6.4pt}{7.3pt}\selectfont\color{darkgray}15\hspace{1.8mm}}\fontsize{7.7pt}{9.0pt}\selectfont\strut Roulé toute la nuit.\strut &  \\
      \cline{1-1}
      \leavevmode\fontsize{7.7pt}{9.0pt}\selectfont\strut [ — Brioude, buffet, varia : 60 c ]\strut & \smash{\parbox[t]{39mm}{\fontsize{6.2pt}{7.0pt}\selectfont\raggedright\color{black!88}\textbf{16} \textit{in marg. cum unco} (comptes).}} \\
\end{tabular}}%
\par\vspace*{\fill}
% --- Registres de bas de page (1. Critique interne [si pertinent] ; 2. Loci paralleli en latin ; 3. Commentaire historique en retrait) ---
\noindent\rule{\textwidth}{0.5pt}\par\vspace{0.65mm}
{\fontsize{6.6pt}{7.7pt}\selectfont%
\noindent\normalfont\color{black!92} \textbf{2} (\textit{Paris}) cf. supra in tegumento (p.~\pageref{fol:1}), f\textsuperscript{o}~66~r\textsuperscript{o} (p.~\pageref{fol:132}), infra f\textsuperscript{o}~74~r\textsuperscript{o} (p.~\pageref{fol:148}), \textit{et passim}, Appendicem I, doc.~5 (p.~\pageref{annexe:doc5}).\par\vspace{0.55mm}
\noindent\rule{38mm}{0.4pt}\par\vspace{0.5mm}
\noindent\hspace*{5mm}\parbox{\dimexpr\textwidth-5mm\relax}{\fontsize{6.6pt}{7.7pt}\selectfont\color{black!84}Suite de la journée du mardi 6 février 1917: arrivé par Corbeil-Essonnes à \textbf{Paris-Gare de Lyon} vers 18~h~30, Hippolyte Combes effectue en moins de trois quarts d'heure la correspondance par le métropolitain (ligne~1 jusqu'à \textbf{Bastille}, puis ligne~5 jusqu'à la \textbf{Gare du Nord}) pour monter dans le train militaire de 19~h~19 à destination d'Amiens.}\par}
\newpage
\phantomsection\label{fac:145}
% --- En-tête du fac-similé (Page impaire / droite : Photographie du manuscrit) ---
\noindent\rule{\textwidth}{0.5pt}\par\vspace{0.55mm}
\noindent\textbf{\fontsize{6.8pt}{7.8pt}\selectfont FAC-SIMILÉ EN VIS-À-VIS — FOLIO 72 VERSO} \hfill \textbf{\fontsize{6.8pt}{7.8pt}\selectfont FÉVRIER 1917}\par\vspace{0.25mm}
\noindent{\fontsize{6.8pt}{8.0pt}\selectfont\textit{Livre-Journal d'Hippolyte Combes} — 6 février 1917 (suite, soir) (Corbeil-Essonnes, Paris (Gare de Lyon, Bastille, Gare du Nord))}\par\vspace{0.55mm}
\noindent\rule{\textwidth}{0.35pt}\par\vspace*{\fill}
{\centering
  \setlength{\fboxsep}{1.2mm}%
  \setlength{\fboxrule}{0.4pt}%
  \fcolorbox{gray!55}{white}{%
    \includegraphics[width=0.95\textwidth,height=0.81\textheight,keepaspectratio]{\facsimiledir/20251126_110601.jpg}%
  }\par}
\vspace*{\fill}
\newpage

% =============================================================================
% VIS-A-VIS SPREAD : FOLIO 73 RECTO (20251126_110616.jpg)
% =============================================================================
\cleardoubleevenpage
\renewcommand{\currentmonthyear}{Février 1917}
\renewcommand{\currentfoliolabel}{Folio 73 recto}
\phantomsection\label{fol:146}
% --- En-tête du folio (Page paire / gauche : Transcription, apparat & registres de bas de page) ---
\noindent\rule{\textwidth}{0.5pt}\par\vspace{0.55mm}
\noindent\textbf{\footnotesize FOLIO 73 RECTO} \hfill \textbf{\footnotesize FÉVRIER 1917}\par\vspace{0.25mm}
\noindent{\fontsize{6.9pt}{8.0pt}\selectfont Nuit glaciale à Amiens, combat aérien contre deux Taubes à Louvencourt et retour chez les Carluvant}\par\vspace{0.2mm}
\noindent{\fontsize{6.9pt}{8.0pt}\selectfont\textit{7 février 1917 (nuit et matin)} \hfill \textit{Amiens, Doullens, Louvencourt, Acheux-en-Amiénois (Somme)}}\par\vspace{0.5mm}
\noindent\rule{\textwidth}{0.35pt}\par\vspace{0.8mm}
% --- Corps de la transcription diplomatique et apparat marginal en prolongement ---
{\renewcommand{\arraystretch}{0.87}%
\noindent\hspace*{3.5mm}\begin{tabular}{@{}p{87mm}@{\hspace{3.5mm}}p{39mm}@{}}
      \leavevmode\fontsize{6.9pt}{7.9pt}\selectfont\strut \underline{Merc. 7} — \hfill{} \underline{Février 1917}\strut & \smash{\parbox[t]{39mm}{\fontsize{6.2pt}{7.0pt}\selectfont\raggedright\color{black!88}\textbf{1 Merc.} \textit{scil.} Mercredi.}} \\
      \leavevmode\fontsize{6.9pt}{7.9pt}\selectfont\strut Voyagé compartiment surchauffé. Très serré.\strut &  \\
      \leavevmode\fontsize{6.9pt}{7.9pt}\selectfont\strut Impossible de fermer yeux. Froid de +\strut &  \\
      \leavevmode\fontsize{6.9pt}{7.9pt}\selectfont\strut en + vif. Gelé très dur.\strut &  \\
      \leavevmode\llap{\fontsize{6.4pt}{7.3pt}\selectfont\color{darkgray}5\hspace{1.8mm}}\fontsize{6.9pt}{7.9pt}\selectfont\strut A minuit 30, à Amiens.\strut &  \\
      \leavevmode\fontsize{6.9pt}{7.9pt}\selectfont\strut Venu à Y.M.C.A. hut, où me suis reposé\strut & \smash{\parbox[t]{39mm}{\fontsize{6.2pt}{7.0pt}\selectfont\raggedright\color{black!88}\textbf{6 Y.M.C.A.} \textit{angl.} Young Men's Christian Assoc.}} \\
      \leavevmode\fontsize{6.9pt}{7.9pt}\selectfont\strut 1 h. et mangé un peu.\strut &  \\
      \leavevmode\fontsize{6.9pt}{7.9pt}\selectfont\strut Resté 2 h. salle d'attente gare, absolt bondée.\strut & \smash{\parbox[t]{39mm}{\fontsize{6.2pt}{7.0pt}\selectfont\raggedright\color{black!88}\textbf{8 absolt} \textit{abbr. pro} absolument. — \textbf{9 m.} \textit{abbr. pro} matin.}} \\
      \leavevmode\fontsize{6.9pt}{7.9pt}\selectfont\strut Enfin parti train de Doullens 6h 7 m.\strut &  \\
      \leavevmode\llap{\fontsize{6.4pt}{7.3pt}\selectfont\color{darkgray}10\hspace{1.8mm}}\fontsize{6.9pt}{7.9pt}\selectfont\strut Train, lenteur désespérante. Arrêts\strut &  \\
      \leavevmode\fontsize{6.9pt}{7.9pt}\selectfont\strut prolongés à toutes les stations.\strut &  \\
      \leavevmode\fontsize{6.9pt}{7.9pt}\selectfont\strut Peu à peu soleil se lève, très brillant.\strut &  \\
      \leavevmode\fontsize{6.9pt}{7.9pt}\selectfont\strut Ciel d'1 bleu foncé, pâle à l'horizon.\strut &  \\
      \leavevmode\fontsize{6.9pt}{7.9pt}\selectfont\strut En gare de Louvencourt, 2 taubes,\strut &  \\
      \leavevmode\llap{\fontsize{6.4pt}{7.3pt}\selectfont\color{darkgray}15\hspace{1.8mm}}\fontsize{6.9pt}{7.9pt}\selectfont\strut plané, évolué ciel bleu, juste au-\strut &  \\
      \leavevmode\fontsize{6.9pt}{7.9pt}\selectfont\strut dessus voie. Obus éclataient tout autour\strut &  \\
      \leavevmode\fontsize{6.9pt}{7.9pt}\selectfont\strut de ces avions et taches blanches.\strut &  \\
      \leavevmode\fontsize{6.9pt}{7.9pt}\selectfont\strut Arrivé à Acheux, vers 1h 1/2 s.\strut & \smash{\parbox[t]{39mm}{\fontsize{6.2pt}{7.0pt}\selectfont\raggedright\color{black!88}\textbf{18 1/2} \textit{scil.} demi / moitié. \textbf{18 s.} \textit{abbr. pro} soir.}} \\
      \leavevmode\fontsize{6.9pt}{7.9pt}\selectfont\strut Me suis rendu à pied ancien billet.\strut &  \\
      \leavevmode\llap{\fontsize{6.4pt}{7.3pt}\selectfont\color{darkgray}20\hspace{1.8mm}}\fontsize{6.9pt}{7.9pt}\selectfont\strut M. \& Mme Carluvant m'ont offert déjeuner\strut & \smash{\parbox[t]{39mm}{\fontsize{6.2pt}{7.0pt}\selectfont\raggedright\color{black!88}\textbf{20 Mme} \textit{abbr. pro} Madame.}} \\
      \leavevmode\fontsize{6.9pt}{7.9pt}\selectfont\strut sommaire.\strut &  \\
      \leavevmode\fontsize{6.9pt}{7.9pt}\selectfont\strut Après m'être rasé, lavé, brossé, sorti.\strut &  \\
\end{tabular}}%
\par\vspace*{\fill}
% --- Registres de bas de page (1. Critique interne [si pertinent] ; 2. Loci paralleli en latin ; 3. Commentaire historique en retrait) ---
\noindent\rule{\textwidth}{0.5pt}\par\vspace{0.65mm}
{\fontsize{6.6pt}{7.7pt}\selectfont%
\noindent\normalfont\color{black!92} \textbf{5} (\textit{Amiens}) cf. supra f\textsuperscript{o}~25~v\textsuperscript{o} (p.~\pageref{fol:51}), f\textsuperscript{o}~72~v\textsuperscript{o} (p.~\pageref{fol:145}), infra f\textsuperscript{o}~74~r\textsuperscript{o} (p.~\pageref{fol:148}), f\textsuperscript{o}~76~v\textsuperscript{o} (p.~\pageref{fol:153}), \textit{et passim}. — \textbf{9} (\textit{Doullens}) cf. supra f\textsuperscript{o}~25~v\textsuperscript{o} (p.~\pageref{fol:51}), f\textsuperscript{o}~64~r\textsuperscript{o} (p.~\pageref{fol:128}), infra f\textsuperscript{o}~76~r\textsuperscript{o} (p.~\pageref{fol:152}), f\textsuperscript{o}~79~r\textsuperscript{o} (p.~\pageref{fol:158}), \textit{et passim}. — \textbf{14} (\textit{Louvencourt}) cf. supra f\textsuperscript{o}~49~v\textsuperscript{o} (p.~\pageref{fol:99}), f\textsuperscript{o}~63~r\textsuperscript{o} (p.~\pageref{fol:126}), \textit{et passim}.\par\vspace{0.55mm}
\noindent\rule{38mm}{0.4pt}\par\vspace{0.5mm}
\noindent\hspace*{5mm}\parbox{\dimexpr\textwidth-5mm\relax}{\fontsize{6.6pt}{7.7pt}\selectfont\color{black!84}Après une halte nocturne en gare d'\textbf{Amiens} (foyer de la Y.M.C.A. et salle d'attente de minuit 30 à 6~h~07), le train omnibus remontant vers Doullens et Acheux est survolé en gare de \textbf{Louvencourt} (Somme) par deux avions allemands de reconnaissance (\textit{taubes}) encadrés par les flocons blancs des obus fusants de la défense contre avions britannique.}\par}
\newpage
\phantomsection\label{fac:146}
% --- En-tête du fac-similé (Page impaire / droite : Photographie du manuscrit) ---
\noindent\rule{\textwidth}{0.5pt}\par\vspace{0.55mm}
\noindent\textbf{\fontsize{6.8pt}{7.8pt}\selectfont FAC-SIMILÉ EN VIS-À-VIS — FOLIO 73 RECTO} \hfill \textbf{\fontsize{6.8pt}{7.8pt}\selectfont FÉVRIER 1917}\par\vspace{0.25mm}
\noindent{\fontsize{6.8pt}{8.0pt}\selectfont\textit{Livre-Journal d'Hippolyte Combes} — 7 février 1917 (nuit et matin) (Amiens, Doullens, Louvencourt, Acheux-en-Amiénois (Somme))}\par\vspace{0.55mm}
\noindent\rule{\textwidth}{0.35pt}\par\vspace*{\fill}
{\centering
  \setlength{\fboxsep}{1.2mm}%
  \setlength{\fboxrule}{0.4pt}%
  \fcolorbox{gray!55}{white}{%
    \includegraphics[width=0.95\textwidth,height=0.81\textheight,keepaspectratio]{\facsimiledir/20251126_110616.jpg}%
  }\par}
\vspace*{\fill}
\newpage

% =============================================================================
% CHAPITRE VII (FOLIO 73 VERSO A FOLIO 93 RECTO)
% =============================================================================
\cleardoubleevenpage
\renewcommand{\currentmonthyear}{Février 1917 – Avril 1917}
\renewcommand{\currentfoliolabel}{Chapitre VII}
\chapter{De février à avril 1917 : Hérissart et le printemps sur la Somme}

\noindent{\fontsize{7.8pt}{9.0pt}\selectfont\textit{Folios 73 verso à 93 recto — L'Army Purchase Board à Hérissart, le repli allemand sur l'Ancre, la famille Flutre et le sous-lieutenant Gilbert (7 février – 30 avril 1917)}}\par\vspace{1.5mm}
\noindent\rule{\textwidth}{0.5pt}\par\vspace{0.9mm}
\noindent\textbf{\scriptsize\textsc{Résumé analytique du Chapitre VII (Folio 73 verso à Folio 93 recto)}}\par\vspace{0.5mm}
\noindent{\fontsize{6.9pt}{8.1pt}\selectfont Le septième et dernier chapitre (f\textsuperscript{o}~73~v\textsuperscript{o} à f\textsuperscript{o}~93~r\textsuperscript{o}) couvre près de trois mois, du retour de permission le 7 février jusqu'à la clôture du carnet le lundi 30 avril 1917. Le service des achats et le logement du capitaine Williams ayant été transférés d'Acheux à Hérissart pendant son absence, Hippolyte Combes s'installe à Hérissart et reprend les tournées quotidiennes de fourrage vers Toutencourt, Raincheval, Puchevillers, Rubempré, Beaucourt-sur-l'Hallue, Molliens-au-Bois, Doullens et Auxi-le-Château, tout en visitant le 10 février son ancienne 152\textsuperscript{e} Compagnie A.S.C. au bois d'Aveluy. Le 10 mars, le départ définitif du capitaine Williams le délivre d'une cohabitation pesante ; dès le lendemain 11 mars, à la demande expresse du major Levy, Hippolyte Combes est maintenu à la commission d'achats de l'armée (\textit{Army Purchase Board}) à Hérissart et emménage dans la chambre paisible du \textit{billet} n\textsuperscript{o}~220 laissée par l'interprète Valet (de la 2\textsuperscript{e} Division australienne). Il rend visite le 13 mars au Mesnil-Martinsart au couple d'instituteurs Fortuné et Mme Flutre (parents du normalien Louis-Fernand Flutre), héberge le 19 mars un officier de la 63\textsuperscript{e} Division navale revenant des tranchées de Miraumont et d'Achiet-le-Petit lors du repli allemand sur la ligne Hindenburg, et admire le 21 mars à Senlis-le-Sec le défilé de la cavalerie indienne et des \textit{17th Lancers (Death or Glory Boys)}. Le mois d'avril 1917 est illuminé par l'arrivée de Pâques, par l'hospitalité chaleureuse de la famille Denis à Hérissart, par la canonnade de l'offensive d'Arras et surtout par les longues conversations littéraires et promenades printanières avec le sous-lieutenant poète Gilbert, originaire de Jersey, jusqu'à la dernière page radieuse du 30 avril 1917.}\par\vspace{1.0mm}
\noindent\rule{\textwidth}{0.35pt}\par\vspace{1.4mm}

\noindent\textbf{\scriptsize\textsc{Introduction historique et philologique}}\par\vspace{0.6mm}
{\fontsize{7.1pt}{8.3pt}\selectfont Ce dernier trimestre du \textit{Livre-Journal} offre un témoignage privilégié sur l'arrière-front de la Somme au moment du repli stratégique allemand vers la ligne Hindenburg (opération \textit{Alberich}, février–mars 1917) et des offensives alliées d'avril 1917 (batailles d'Arras, de la Scarpe et de Bullecourt). Trois rencontres majeures illustrent la richesse prosopographique et historique de ces pages : le 13 mars 1917 au \textbf{Mesnil-Martinsart} (f\textsuperscript{o}~84~r\textsuperscript{o}), la visite à l'instituteur \textbf{Fortuné Flutre} et à son épouse, parents du philologue romaniste \textbf{Louis-Fernand Flutre} (1892–1978, admis à l'École normale supérieure en 1914, agrégé en 1921 et auteur en 1955 des deux monographies de référence sur l'histoire et le parler picard du Mesnil-Martinsart ; R.~Debrie, \textit{Revue de linguistique romane}, 42, 1978, p.~238–240) ; les 19–21 mars (f\textsuperscript{o}~85~r\textsuperscript{o}–v\textsuperscript{o}), le passage de la \textbf{63\textsuperscript{e} Division (\textit{Royal Naval Division})} et de la \textbf{4\textsuperscript{e} Division de cavalerie} (\textit{Sialkot Cavalry Brigade}, unissant les \textbf{17\textsuperscript{e} Lanciers britanniques} \textit{Death or Glory Boys} et les régiments de lanciers indiens) ; enfin, à partir du 13 avril (f\textsuperscript{o}~89~v\textsuperscript{o}), l'amitié intellectuelle nouée avec le \textbf{sous-lieutenant Gilbert}, originaire de l'île anglo-normande de \textbf{Jersey}, qui rend à Hippolyte Combes le goût des lettres et de la sérénité au seuil du printemps 1917.\par}
\vfill
\newpage
% --- Page de droite d'ouverture du chapitre : Carte vectorielle TikZ & Photographies d'archives ---
\noindent\rule{\textwidth}{0.5pt}\par\vspace{0.6mm}
\noindent\textbf{\fontsize{6.8pt}{7.8pt}\selectfont CARTE D'ITINÉRAIRE DU CHAPITRE VII} \hfill \textbf{\fontsize{6.8pt}{7.8pt}\selectfont FÉVRIER 1917 – AVRIL 1917}\par\vspace{0.25mm}
\noindent{\fontsize{6.7pt}{7.8pt}\selectfont\textit{Déplacements, cantonnements et tournées d'Hippolyte Combes (Folio 73 verso à Folio 93 recto)}}\par\vspace{0.55mm}
\noindent\rule{\textwidth}{0.35pt}\par\vspace{1.2mm}
\begin{center}

\begin{tikzpicture}[
  scale=0.96,
  every node/.style={font=\fontsize{6.0pt}{7.0pt}\selectfont},
  loc/.style={circle, fill=black, inner sep=1.1pt},
  hub/.style={circle, draw=bordeaux, fill=white, line width=0.85pt, inner sep=1.8pt},
  tour/.style={draw=bordeaux!85, line width=0.55pt, ->, >=stealth},
  advance/.style={draw=black!75, line width=0.65pt, dashed, ->, >=stealth},
  river/.style={draw=black!35, line width=1.1pt}
]
\draw[black!50, line width=0.4pt] (-60mm,-42mm) rectangle (60mm,42mm);
\fill[black!3] (-60mm,-42mm) rectangle (60mm,42mm);

% Rivières : l'Authie au nord, l'Ancre à l'est, l'Hallue au sud
\draw[river] (-52mm,30mm) -- (12mm,26mm) node[right]{\textit{L'Authie}};
\draw[river] (42mm,24mm) .. controls (34mm,6mm) and (30mm,-10mm) .. (22mm,-34mm) node[right]{\textit{L'Ancre}};

% Centre d'Hérissart (Army Purchase Board)
\node[hub] (herissart) at (-8mm, 0mm) {};
\node[above left] at (herissart) {\textbf{HÉRISSART (\textit{Purchase Board})}};

% Nord-Ouest : Doullens, Auxi-le-Château, Candas
\node[loc] (auxi) at (-48mm, 28mm) {};
\node[below right] at (auxi) {\textbf{Auxi-le-Château} (21 févr.)};
\node[loc] (doullens) at (-20mm, 26mm) {};
\node[right] at (doullens) {\textbf{Doullens} (achats \& cantine)};
\node[loc] (candas) at (-30mm, 14mm) {};
\node[left] at (candas) {Candas};
\node[loc] (raincheval) at (-10mm, 14mm) {};
\node[left] at (raincheval) {Raincheval \& Puchevillers};
\node[loc] (acheux) at (6mm, 16mm) {};
\node[above] at (acheux) {Acheux \& Bertrancourt};

% Ouest et Sud : Villers-Bocage, Talmas, Rubempré, Molliens, Beaucourt
\node[loc] (talmas) at (-36mm, -6mm) {};
\node[left] at (talmas) {Talmas \& \textbf{Villers-Bocage}};
\node[loc] (pierregot) at (-16mm, -8mm) {};
\node[left] at (pierregot) {Pierregot \& Mirvaux};
\node[loc] (rubempre) at (-8mm, -10mm) {};
\node[below] at (rubempre) {\textbf{Rubempré}};
\node[loc] (molliens) at (-14mm, -20mm) {};
\node[left] at (molliens) {Molliens-au-Bois};
\node[loc] (beaucourt) at (-4mm, -24mm) {};
\node[below] at (beaucourt) {\textbf{Beaucourt-sur-l'Hallue} \& Montigny (16 févr.)};

% Est : Toutencourt, Contay, Warloy, Senlis, Mesnil-Martinsart, Albert
\node[loc] (toutencourt) at (2mm, 6mm) {};
\node[right] at (toutencourt) {\textbf{Toutencourt} (Mission V\textsuperscript{e} Armée)};
\node[loc] (contay) at (6mm, -6mm) {};
\node[right] at (contay) {\textbf{Contay} \& Warloy-Baillon};
\node[loc] (senlis) at (18mm, -2mm) {};
\node[right] at (senlis) {\textbf{Senlis-le-Sec} (\textit{17th Lancers}, 21 mars)};
\node[loc] (mesnil) at (28mm, 4mm) {};
\node[right] at (mesnil) {\textbf{Le Mesnil-Martinsart} (famille Flutre, 13 mars)};
\node[loc] (albert) at (28mm, -12mm) {};
\node[right] at (albert) {\textbf{Albert}, Aveluy \& Buire-sur-l'Ancre};

% Nord-Est : Secteur du repli allemand (Miraumont, Achiet-le-Petit, Bapaume, Arras)
\node[loc] (miraumont) at (38mm, 18mm) {};
\node[right] at (miraumont) {\textbf{Miraumont}};
\node[loc] (achiet) at (40mm, 26mm) {};
\node[right] at (achiet) {\textbf{Achiet-le-Petit} (\textit{63rd Naval Div.}, 19 mars)};
\node[loc] (arras) at (36mm, 36mm) {};
\node[above] at (arras) {\textbf{Vers Bapaume \& Arras} (canonnade, avr. 1917)};

% Rayons et mouvements
\foreach \dest in {doullens,candas,raincheval,acheux,toutencourt,contay,senlis,mesnil,albert,rubempre,pierregot,molliens,beaucourt,talmas} {
  \draw[tour] (herissart) -- (\dest);
}
\draw[tour] (doullens) -- (auxi);
\draw[advance] (albert) -- (mesnil) -- (miraumont) -- (achiet) -- (arras);

\draw[line width=0.6pt] (-54mm,-36mm) -- (-34mm,-36mm);
\node[above] at (-44mm,-36mm) {0 \hfill 5 \hfill 10 km};
\draw[->, line width=0.6pt] (52mm, 28mm) -- (52mm, 37mm) node[above]{\textbf{N}};
\end{tikzpicture}

\end{center}
\vfill
\noindent\rule{\textwidth}{0.3pt}\par\vspace{1.2mm}
\noindent\begin{center}
  \begin{minipage}[b]{0.31\textwidth}
    \centering
    \setlength{\fboxsep}{0.7mm}\setlength{\fboxrule}{0.4pt}%
    \fcolorbox{gray!55}{white}{\includegraphics[width=0.94\linewidth,height=28mm,keepaspectratio]{illustrations/andelys_ecole_militaire.jpg}}
    \par\vspace{0.6mm}{\fontsize{5.8pt}{6.7pt}\selectfont Photographie de groupe devant l'École militaire préparatoire des Andelys (où Hippolyte Combes est nommé professeur d'anglais le 24 novembre 1918).}
  \end{minipage}\hfill
  \begin{minipage}[b]{0.31\textwidth}
    \centering
    \setlength{\fboxsep}{0.7mm}\setlength{\fboxrule}{0.4pt}%
    \fcolorbox{gray!55}{white}{\includegraphics[width=0.94\linewidth,height=28mm,keepaspectratio]{illustrations/famille_combes_1925.jpg}}
    \par\vspace{0.6mm}{\fontsize{5.8pt}{6.7pt}\selectfont Hippolyte et Marguerite Combes entourés de leurs enfants à Montauban / Millau vers 1925.}
  \end{minipage}\hfill
  \begin{minipage}[b]{0.31\textwidth}
    \centering
    \setlength{\fboxsep}{0.7mm}\setlength{\fboxrule}{0.4pt}%
    \fcolorbox{gray!55}{white}{\includegraphics[width=0.94\linewidth,height=28mm,keepaspectratio]{illustrations/hippolyte_cambridge_1930.jpg}}
    \par\vspace{0.6mm}{\fontsize{5.8pt}{6.7pt}\selectfont Hippolyte Combes en Angleterre (Cambridge, été 1930).}
  \end{minipage}
\end{center}
\vspace*{1mm}
\newpage

% =============================================================================
% VIS-A-VIS SPREAD : FOLIO 73 VERSO (20251126_110628.jpg)
% =============================================================================
\cleardoubleevenpage
\renewcommand{\currentmonthyear}{Février 1917}
\renewcommand{\currentfoliolabel}{Folio 73 verso}
\phantomsection\label{fol:147}
% --- En-tête du folio (Page paire / gauche : Transcription, apparat & registres de bas de page) ---
\noindent\rule{\textwidth}{0.5pt}\par\vspace{0.55mm}
\noindent\textbf{\footnotesize FOLIO 73 VERSO} \hfill \textbf{\footnotesize FÉVRIER 1917}\par\vspace{0.25mm}
\noindent{\fontsize{6.9pt}{8.0pt}\selectfont Déménagement imprévu d'Acheux à Hérissart sous un gel intense}\par\vspace{0.2mm}
\noindent{\fontsize{6.9pt}{8.0pt}\selectfont\textit{7 février 1917 (après-midi)} \hfill \textit{Acheux, Raincheval, Puchevillers, Hérissart (Somme)}}\par\vspace{0.5mm}
\noindent\rule{\textwidth}{0.35pt}\par\vspace{0.8mm}
% --- Corps de la transcription diplomatique et apparat marginal en prolongement ---
{\renewcommand{\arraystretch}{0.87}%
\noindent\hspace*{3.5mm}\begin{tabular}{@{}p{87mm}@{\hspace{3.5mm}}p{39mm}@{}}
      \leavevmode\fontsize{6.9pt}{7.9pt}\selectfont\strut \underline{Merc. 7 (suite)} \hfill{} \raisebox{0.5ex}{\underline{Février 1917}}\strut & \smash{\parbox[t]{39mm}{\fontsize{6.2pt}{7.0pt}\selectfont\raggedright\color{black!88}\textbf{1 Merc.} \textit{scil.} Mercredi.}} \\
      \leavevmode\fontsize{6.9pt}{7.9pt}\selectfont\strut En me rendant à notre office,\strut &  \\
      \leavevmode\fontsize{6.9pt}{7.9pt}\selectfont\strut Mme Delille, qui au guichet, m'a appelé et\strut & \smash{\parbox[t]{39mm}{\fontsize{6.2pt}{7.0pt}\selectfont\raggedright\color{black!88}\textbf{3 Mme} \textit{abbr. pro} Madame.}} \\
      \leavevmode\fontsize{6.9pt}{7.9pt}\selectfont\strut m'a appris que Cap. Williams, quitté Acheux\strut & \smash{\parbox[t]{39mm}{\fontsize{6.2pt}{7.0pt}\selectfont\raggedright\color{black!88}\textbf{4 Cap.} \textit{abbr. pro} Capitaine. — \textbf{5} $[\![$...$]\!]\ $ \textit{del. in ms.} — \textbf{6 Cap. W.} \textit{scil.} Capitaine Williams. \textbf{6 Cap.} \textit{abbr. pro} Capitaine.}} \\
      \leavevmode\llap{\fontsize{6.4pt}{7.3pt}\selectfont\color{darkgray}5\hspace{1.8mm}}\fontsize{6.9pt}{7.9pt}\selectfont\strut pour $[\![$...$]\!]$ Hérissart. Venu aussitôt Mission\strut &  \\
      \leavevmode\fontsize{6.9pt}{7.9pt}\selectfont\strut française, où trouvé un mot du Cap. W.\strut &  \\
      \leavevmode\fontsize{6.9pt}{7.9pt}\selectfont\strut Lui ai téléphoné du bureau Mission pour lui\strut &  \\
      \leavevmode\fontsize{6.9pt}{7.9pt}\selectfont\strut annoncer mon arrivée.\strut &  \\
      \leavevmode\fontsize{6.9pt}{7.9pt}\selectfont\strut Rentré chez les Carluvant, auxquels annoncé mon\strut &  \\
      \leavevmode\llap{\fontsize{6.4pt}{7.3pt}\selectfont\color{darkgray}10\hspace{1.8mm}}\fontsize{6.9pt}{7.9pt}\selectfont\strut départ. Atterrés. Fait en hâte mes paquets. Attendu\strut &  \\
      \leavevmode\fontsize{6.9pt}{7.9pt}\selectfont\strut l'auto, qui arrivée vers 4h 1/2.\strut & \smash{\parbox[t]{39mm}{\fontsize{6.2pt}{7.0pt}\selectfont\raggedright\color{black!88}\textbf{11 1/2} \textit{scil.} demi / moitié.}} \\
      \leavevmode\fontsize{6.9pt}{7.9pt}\selectfont\strut Quitté village d'Acheux et surtout maison\strut &  \\
      \leavevmode\fontsize{6.9pt}{7.9pt}\selectfont\strut avec 1 serrement de cœur — —\strut &  \\
      \leavevmode\fontsize{6.9pt}{7.9pt}\selectfont\strut En auto à Hérissart, par Raincheval\strut &  \\
      \leavevmode\llap{\fontsize{6.4pt}{7.3pt}\selectfont\color{darkgray}15\hspace{1.8mm}}\fontsize{6.9pt}{7.9pt}\selectfont\strut \& Puchevillers.\strut &  \\
      \leavevmode\fontsize{6.9pt}{7.9pt}\selectfont\strut Bise glaciale. Froid intense.\strut &  \\
      \leavevmode\fontsize{6.9pt}{7.9pt}\selectfont\strut Town Major m'a donné petite chambre, n° 7b,\strut & \smash{\parbox[t]{39mm}{\fontsize{6.2pt}{7.0pt}\selectfont\raggedright\color{black!88}\textbf{17 Town Major} \textit{angl. in ms.}}} \\
      \leavevmode\fontsize{6.9pt}{7.9pt}\selectfont\strut 1er étage café Charlet, angle rue. De ma\strut & \smash{\parbox[t]{39mm}{\fontsize{6.2pt}{7.0pt}\selectfont\raggedright\color{black!88}\textbf{18 1er} \textit{scil.} premier.}} \\
      \leavevmode\fontsize{6.9pt}{7.9pt}\selectfont\strut fenêtre, as. belle vue sur la place et église.\strut & \smash{\parbox[t]{39mm}{\fontsize{6.2pt}{7.0pt}\selectfont\raggedright\color{black!88}\textbf{19 as.} \textit{abbr. pro} assez.}} \\
      \leavevmode\llap{\fontsize{6.4pt}{7.3pt}\selectfont\color{darkgray}20\hspace{1.8mm}}\fontsize{6.9pt}{7.9pt}\selectfont\strut Pleine lune brillait entre gds arbres.\strut & \smash{\parbox[t]{39mm}{\fontsize{6.2pt}{7.0pt}\selectfont\raggedright\color{black!88}\textbf{20 gds} \textit{abbr. pro} grands.}} \\
      \leavevmode\fontsize{6.9pt}{7.9pt}\selectfont\strut Changement plutôt pénible d'habitudes.\strut &  \\
\end{tabular}}%
\par\vspace*{\fill}
% --- Registres de bas de page (1. Critique interne [si pertinent] ; 2. Loci paralleli en latin ; 3. Commentaire historique en retrait) ---
\noindent\rule{\textwidth}{0.5pt}\par\vspace{0.65mm}
{\fontsize{6.6pt}{7.7pt}\selectfont%
\noindent\normalfont \textbf{Continuité :} Deuxième des trois feuillets consacrés à la seule journée du mercredi 7 février 1917 (\textit{Merc. 7 (suite)}), relatant le retour de permission entre Paris, Amiens et Hérissart.\par\vspace{0.5mm}
\noindent\rule{22mm}{0.3pt}\par\vspace{0.45mm}
\noindent\normalfont\color{black!92} \textbf{4} (\textit{Acheux}) cf. supra f\textsuperscript{o}~21~v\textsuperscript{o} (p.~\pageref{fol:43}), f\textsuperscript{o}~73~r\textsuperscript{o} (p.~\pageref{fol:146}), infra f\textsuperscript{o}~74~r\textsuperscript{o} (p.~\pageref{fol:148}), f\textsuperscript{o}~74~v\textsuperscript{o} (p.~\pageref{fol:149}), \textit{et passim}. — \textbf{4} (\textit{Williams}) cf. supra f\textsuperscript{o}~52~v\textsuperscript{o} (p.~\pageref{fol:105}), f\textsuperscript{o}~58~v\textsuperscript{o} (p.~\pageref{fol:117}), infra f\textsuperscript{o}~78~v\textsuperscript{o} (p.~\pageref{fol:157}), f\textsuperscript{o}~80~r\textsuperscript{o} (p.~\pageref{fol:160}), \textit{et passim}. — \textbf{5} (\textit{Hérissart}) cf. supra f\textsuperscript{o}~47~v\textsuperscript{o} (p.~\pageref{fol:95}), f\textsuperscript{o}~62~r\textsuperscript{o} (p.~\pageref{fol:124}), infra f\textsuperscript{o}~74~r\textsuperscript{o} (p.~\pageref{fol:148}), f\textsuperscript{o}~74~v\textsuperscript{o} (p.~\pageref{fol:149}), \textit{et passim}.\par\vspace{0.55mm}
\noindent\rule{38mm}{0.4pt}\par\vspace{0.5mm}
\noindent\hspace*{5mm}\parbox{\dimexpr\textwidth-5mm\relax}{\fontsize{6.6pt}{7.7pt}\selectfont\color{black!84}Le mercredi 7 février 1917 après-midi, à peine rentré de permission à Acheux, Hippolyte Combes apprend par la postière \textbf{Mme Delille} et par un pli laissé à la Mission française que le bureau des fourrages du \textbf{capitaine Williams} a été transféré à \textbf{Hérissart} (Somme, canton d'Acheux, à 11~km au sud-ouest d'Acheux).}\par}
\newpage
\phantomsection\label{fac:147}
% --- En-tête du fac-similé (Page impaire / droite : Photographie du manuscrit) ---
\noindent\rule{\textwidth}{0.5pt}\par\vspace{0.55mm}
\noindent\textbf{\fontsize{6.8pt}{7.8pt}\selectfont FAC-SIMILÉ EN VIS-À-VIS — FOLIO 73 VERSO} \hfill \textbf{\fontsize{6.8pt}{7.8pt}\selectfont FÉVRIER 1917}\par\vspace{0.25mm}
\noindent{\fontsize{6.8pt}{8.0pt}\selectfont\textit{Livre-Journal d'Hippolyte Combes} — 7 février 1917 (après-midi) (Acheux, Raincheval, Puchevillers, Hérissart (Somme))}\par\vspace{0.55mm}
\noindent\rule{\textwidth}{0.35pt}\par\vspace*{\fill}
{\centering
  \setlength{\fboxsep}{1.2mm}%
  \setlength{\fboxrule}{0.4pt}%
  \fcolorbox{gray!55}{white}{%
    \includegraphics[width=0.95\textwidth,height=0.81\textheight,keepaspectratio]{\facsimiledir/20251126_110628.jpg}%
  }\par}
\vspace*{\fill}
\newpage

% =============================================================================
% VIS-A-VIS SPREAD : FOLIO 74 RECTO (20251126_110646.jpg)
% =============================================================================
\cleardoubleevenpage
\renewcommand{\currentmonthyear}{Février 1917}
\renewcommand{\currentfoliolabel}{Folio 74 recto}
\phantomsection\label{fol:148}
% --- En-tête du folio (Page paire / gauche : Transcription, apparat & registres de bas de page) ---
\noindent\rule{\textwidth}{0.5pt}\par\vspace{0.55mm}
\noindent\textbf{\footnotesize FOLIO 74 RECTO} \hfill \textbf{\footnotesize FÉVRIER 1917}\par\vspace{0.25mm}
\noindent{\fontsize{6.9pt}{8.0pt}\selectfont Nuit glaciale à Hérissart, comptes de retour et premier jeudi de repos}\par\vspace{0.2mm}
\noindent{\fontsize{6.9pt}{8.0pt}\selectfont\textit{7 (fin) – 8 février 1917} \hfill \textit{Hérissart (Somme)}}\par\vspace{0.5mm}
\noindent\rule{\textwidth}{0.35pt}\par\vspace{0.8mm}
% --- Corps de la transcription diplomatique et apparat marginal en prolongement ---
{\renewcommand{\arraystretch}{0.87}%
\noindent\hspace*{3.5mm}\begin{tabular}{@{}p{87mm}@{\hspace{3.5mm}}p{39mm}@{}}
      \leavevmode\fontsize{6.9pt}{7.9pt}\selectfont\strut \underline{Merc. 7 (fin)} \hfill{} \underline{Février 1917}\strut & \smash{\parbox[t]{39mm}{\fontsize{6.2pt}{7.0pt}\selectfont\raggedright\color{black!88}\textbf{1 Merc.} \textit{scil.} Mercredi.}} \\
      \leavevmode\fontsize{6.9pt}{7.9pt}\selectfont\strut Notre bureau \& mess, au n° 104, mauvaise\strut & \smash{\parbox[t]{39mm}{\fontsize{6.2pt}{7.0pt}\selectfont\raggedright\color{black!88}\textbf{2 mess} \textit{angl. in ms.}}} \\
      \leavevmode\fontsize{6.9pt}{7.9pt}\selectfont\strut installation.\strut &  \\
      \leavevmode\fontsize{6.9pt}{7.9pt}\selectfont\strut Après dîner, retourné chambre vers 9h.\strut &  \\
      \leavevmode\llap{\fontsize{6.4pt}{7.3pt}\selectfont\color{darkgray}5\hspace{1.8mm}}\fontsize{6.9pt}{7.9pt}\selectfont\strut Vitres recouvertes très épaisse gelée glacée.\strut &  \\
      \leavevmode\fontsize{6.9pt}{7.9pt}\selectfont\strut Froid extrt rigoureux.\strut & \smash{\parbox[t]{39mm}{\fontsize{6.2pt}{7.0pt}\selectfont\raggedright\color{black!88}\textbf{6 extrt} \textit{abbr. pro} extrêmement. — \textbf{7} \textit{in marg. cum unco} (comptes). \textbf{7 f.} \textit{abbr. pro} franc(s). — \textbf{8 Y.M.C.A.} \textit{angl.} Young Men's Christian Assoc. \textbf{8 f.} \textit{abbr. pro} franc(s). — \textbf{9 f.} \textit{abbr. pro} franc(s). — \textbf{10 Gift} \textit{angl. in ms.} — \textbf{11 f.} \textit{abbr. pro} franc(s). \textbf{11 lunch, room} \textit{angl. in ms. (\emph{et room} infra lin. add.).} — \textbf{17 mess} \textit{angl. in ms.} — \textbf{18 breakfast} \textit{angl. in ms.}}} \\
      \leavevmode\fontsize{6.9pt}{7.9pt}\selectfont\strut [ Dîner à Paris, métro : 1 f.\strut &  \\
      \leavevmode\fontsize{6.9pt}{7.9pt}\selectfont\strut Amiens, Y.M.C.A. : 1 f. 15 c.\strut &  \\
      \leavevmode\fontsize{6.9pt}{7.9pt}\selectfont\strut Touché solde Mission fr., Acheux : 78 f. 75\strut &  \\
      \leavevmode\llap{\fontsize{6.4pt}{7.3pt}\selectfont\color{darkgray}10\hspace{1.8mm}}\fontsize{6.9pt}{7.9pt}\selectfont\strut Gift aux Carluvant, pour\strut &  \\
      \leavevmode\fontsize{6.9pt}{7.9pt}\selectfont\strut mon lunch et room : 5 f. ]\strut &  \\
      \cline{1-1}
      \leavevmode\fontsize{6.9pt}{7.9pt}\selectfont\strut \underline{Jeudi 8.}\strut &  \\
      \leavevmode\fontsize{6.9pt}{7.9pt}\selectfont\strut Très bonne nuit. Bien reposé. Lit\strut &  \\
      \leavevmode\fontsize{6.9pt}{7.9pt}\selectfont\strut très confortable.\strut &  \\
      \leavevmode\llap{\fontsize{6.4pt}{7.3pt}\selectfont\color{darkgray}15\hspace{1.8mm}}\fontsize{6.9pt}{7.9pt}\selectfont\strut Levé vers 8h —\strut &  \\
      \leavevmode\fontsize{6.9pt}{7.9pt}\selectfont\strut Soleil sur mes fenêtres.\strut &  \\
      \leavevmode\fontsize{6.9pt}{7.9pt}\selectfont\strut Descendu à l'office \& mess. Excellent\strut &  \\
      \leavevmode\fontsize{6.9pt}{7.9pt}\selectfont\strut breakfast.\strut &  \\
      \leavevmode\fontsize{6.9pt}{7.9pt}\selectfont\strut Tout le jour, liberté entière —\strut &  \\
\end{tabular}}%
\par\vspace*{\fill}
% --- Registres de bas de page (1. Critique interne [si pertinent] ; 2. Loci paralleli en latin ; 3. Commentaire historique en retrait) ---
\noindent\rule{\textwidth}{0.5pt}\par\vspace{0.65mm}
{\fontsize{6.6pt}{7.7pt}\selectfont%
\noindent\normalfont \textbf{Continuité :} Troisième et dernier feuillet de la journée du 7 février 1917 (\textit{Merc. 7 (fin)}), complété par le début de l'entrée du jeudi 8 février (l.~12--18).\par\vspace{0.5mm}
\noindent\rule{22mm}{0.3pt}\par\vspace{0.45mm}
\noindent\normalfont\color{black!92} \textbf{7} (\textit{Paris}) cf. supra in tegumento (p.~\pageref{fol:1}), f\textsuperscript{o}~72~v\textsuperscript{o} (p.~\pageref{fol:145}), \textit{et passim}, Appendicem I, doc.~5 (p.~\pageref{annexe:doc5}). — \textbf{8} (\textit{Amiens}) cf. supra f\textsuperscript{o}~25~v\textsuperscript{o} (p.~\pageref{fol:51}), f\textsuperscript{o}~73~r\textsuperscript{o} (p.~\pageref{fol:146}), infra f\textsuperscript{o}~76~v\textsuperscript{o} (p.~\pageref{fol:153}), f\textsuperscript{o}~85~v\textsuperscript{o} (p.~\pageref{fol:171}), \textit{et passim}. — \textbf{9} (\textit{Acheux}) cf. supra f\textsuperscript{o}~21~v\textsuperscript{o} (p.~\pageref{fol:43}), f\textsuperscript{o}~73~v\textsuperscript{o} (p.~\pageref{fol:147}), infra f\textsuperscript{o}~74~v\textsuperscript{o} (p.~\pageref{fol:149}), f\textsuperscript{o}~75~v\textsuperscript{o} (p.~\pageref{fol:151}), \textit{et passim}.\par\vspace{0.55mm}
\noindent\rule{38mm}{0.4pt}\par\vspace{0.5mm}
\noindent\hspace*{5mm}\parbox{\dimexpr\textwidth-5mm\relax}{\fontsize{6.6pt}{7.7pt}\selectfont\color{black!84}Fin des notes et comptes du mercredi 7 février 1917 et journée de repos du jeudi 8 février à \textbf{Hérissart}: le bureau et le \textit{mess} occupent le billet \textbf{n\textsuperscript{o}~104} (où Hippolyte Combes devra lui-même loger quelques jours fin février, f\textsuperscript{o}~78~v\textsuperscript{o}). Avant de quitter Acheux, il y a perçu sa quinzaine de 78~fr.~75 à la Mission française et laissé 5~fr. de cadeau à ses logeurs \textbf{M. et Mme Carluvant}.}\par}
\newpage
\phantomsection\label{fac:148}
% --- En-tête du fac-similé (Page impaire / droite : Photographie du manuscrit) ---
\noindent\rule{\textwidth}{0.5pt}\par\vspace{0.55mm}
\noindent\textbf{\fontsize{6.8pt}{7.8pt}\selectfont FAC-SIMILÉ EN VIS-À-VIS — FOLIO 74 RECTO} \hfill \textbf{\fontsize{6.8pt}{7.8pt}\selectfont FÉVRIER 1917}\par\vspace{0.25mm}
\noindent{\fontsize{6.8pt}{8.0pt}\selectfont\textit{Livre-Journal d'Hippolyte Combes} — 7 (fin) – 8 février 1917 (Hérissart (Somme))}\par\vspace{0.55mm}
\noindent\rule{\textwidth}{0.35pt}\par\vspace*{\fill}
{\centering
  \setlength{\fboxsep}{1.2mm}%
  \setlength{\fboxrule}{0.4pt}%
  \fcolorbox{gray!55}{white}{%
    \includegraphics[width=0.95\textwidth,height=0.81\textheight,keepaspectratio]{\facsimiledir/20251126_110646.jpg}%
  }\par}
\vspace*{\fill}
\newpage

% =============================================================================
% VIS-A-VIS SPREAD : FOLIO 74 VERSO (20251126_110659.jpg)
% =============================================================================
\cleardoubleevenpage
\renewcommand{\currentmonthyear}{Février 1917}
\renewcommand{\currentfoliolabel}{Folio 74 verso}
\phantomsection\label{fol:149}
% --- En-tête du folio (Page paire / gauche : Transcription, apparat & registres de bas de page) ---
\noindent\rule{\textwidth}{0.5pt}\par\vspace{0.55mm}
\noindent\textbf{\footnotesize FOLIO 74 VERSO} \hfill \textbf{\footnotesize FÉVRIER 1917}\par\vspace{0.25mm}
\noindent{\fontsize{6.9pt}{8.0pt}\selectfont Promenade à pied sur la neige durcie jusqu'à Rubempré}\par\vspace{0.2mm}
\noindent{\fontsize{6.9pt}{8.0pt}\selectfont\textit{8 (suite) – 9 février 1917 (matin)} \hfill \textit{Hérissart, Rubempré (Somme)}}\par\vspace{0.5mm}
\noindent\rule{\textwidth}{0.35pt}\par\vspace{0.8mm}
% --- Corps de la transcription diplomatique et apparat marginal en prolongement ---
{\renewcommand{\arraystretch}{0.87}%
\noindent\hspace*{3.5mm}\begin{tabular}{@{}p{87mm}@{\hspace{3.5mm}}p{39mm}@{}}
      \leavevmode\fontsize{6.9pt}{7.9pt}\selectfont\strut \underline{Jeudi 8 (suite)} \hfill{} \raisebox{0.5ex}{\underline{Février 1917}}\strut &  \\
      \leavevmode\fontsize{6.9pt}{7.9pt}\selectfont\strut Temps mervt beau. Ciel bleu, soleil ardent.\strut & \smash{\parbox[t]{39mm}{\fontsize{6.2pt}{7.0pt}\selectfont\raggedright\color{black!88}\textbf{2 mervt} \textit{abbr. pro} merveilleusement. — \textbf{3 j'à} \textit{abbr. pro} jusqu'à. \textbf{3 gde} \textit{abbr. pro} grande.}} \\
      \leavevmode\fontsize{6.9pt}{7.9pt}\selectfont\strut Gelé très fort. Terre durcie, j'à gde profondeur.\strut &  \\
      \leavevmode\fontsize{6.9pt}{7.9pt}\selectfont\strut Après-midi, promenade à pied à Rubempré.\strut &  \\
      \leavevmode\llap{\fontsize{6.4pt}{7.3pt}\selectfont\color{darkgray}5\hspace{1.8mm}}\fontsize{6.9pt}{7.9pt}\selectfont\strut Rayons ardents faisaient étinceler neige durcie,\strut &  \\
      \leavevmode\fontsize{6.9pt}{7.9pt}\selectfont\strut qui recouvrait champs.\strut &  \\
      \leavevmode\fontsize{6.9pt}{7.9pt}\selectfont\strut Région, moins banale qu'environs d'Acheux :\strut &  \\
      \leavevmode\fontsize{6.9pt}{7.9pt}\selectfont\strut vallonnements, petites hauteurs, couronnées de\strut &  \\
      \leavevmode\fontsize{6.9pt}{7.9pt}\selectfont\strut moulins à vent.\strut &  \\
      \leavevmode\llap{\fontsize{6.4pt}{7.3pt}\selectfont\color{darkgray}10\hspace{1.8mm}}\fontsize{6.9pt}{7.9pt}\selectfont\strut [durée incertaine] de marche.\strut & \smash{\parbox[t]{39mm}{\fontsize{6.2pt}{7.0pt}\selectfont\raggedright\color{black!88}\textbf{10} \textit{lect. incert. in ms.}}} \\
      \leavevmode\fontsize{6.9pt}{7.9pt}\selectfont\strut 1 h. ds ma chambre, où tt rangé et\strut & \smash{\parbox[t]{39mm}{\fontsize{6.2pt}{7.0pt}\selectfont\raggedright\color{black!88}\textbf{11 ds} \textit{abbr. pro} dans. \textbf{11 tt} \textit{abbr. pro} tout.}} \\
      \leavevmode\fontsize{6.9pt}{7.9pt}\selectfont\strut terminé installation.\strut &  \\
      \leavevmode\fontsize{6.9pt}{7.9pt}\selectfont\strut Froid toujours très vif.\strut &  \\
      \leavevmode\fontsize{6.9pt}{7.9pt}\selectfont\strut Couché vers 9h s. — 1 peu de lecture, lit\strut & \smash{\parbox[t]{39mm}{\fontsize{6.2pt}{7.0pt}\selectfont\raggedright\color{black!88}\textbf{14 s.} \textit{abbr. pro} soir.}} \\
      \cline{1-1}
      \leavevmode\llap{\fontsize{6.4pt}{7.3pt}\selectfont\color{darkgray}15\hspace{1.8mm}}\fontsize{6.9pt}{7.9pt}\selectfont\strut \underline{Ven. 9} — Nuit toujours excessvt froide.\strut & \smash{\parbox[t]{39mm}{\fontsize{6.2pt}{7.0pt}\selectfont\raggedright\color{black!88}\textbf{15 Ven.} \textit{scil.} Vendredi. \textbf{15 excessvt} \textit{abbr. pro} excessivement.}} \\
      \leavevmode\fontsize{6.9pt}{7.9pt}\selectfont\strut Gelé très fort. Ciel clair, pleine lune,\strut &  \\
      \leavevmode\fontsize{6.9pt}{7.9pt}\selectfont\strut vif éclat.\strut &  \\
      \leavevmode\fontsize{6.9pt}{7.9pt}\selectfont\strut Matinée, liberté complète. Mais trop froid\strut &  \\
      \leavevmode\fontsize{6.9pt}{7.9pt}\selectfont\strut pour pouvoir écrire chambre.\strut &  \\
\end{tabular}}%
\par\vspace*{\fill}
% --- Registres de bas de page (1. Critique interne [si pertinent] ; 2. Loci paralleli en latin ; 3. Commentaire historique en retrait) ---
\noindent\rule{\textwidth}{0.5pt}\par\vspace{0.65mm}
{\fontsize{6.6pt}{7.7pt}\selectfont%
\noindent\normalfont\color{black!92} \textbf{4} (\textit{Rubempré}) cf. supra f\textsuperscript{o}~47~v\textsuperscript{o} (p.~\pageref{fol:95}), f\textsuperscript{o}~62~r\textsuperscript{o} (p.~\pageref{fol:124}), infra f\textsuperscript{o}~80~v\textsuperscript{o} (p.~\pageref{fol:161}), f\textsuperscript{o}~81~r\textsuperscript{o} (p.~\pageref{fol:162}), \textit{et passim}. — \textbf{7} (\textit{Acheux}) cf. supra f\textsuperscript{o}~21~v\textsuperscript{o} (p.~\pageref{fol:43}), f\textsuperscript{o}~74~r\textsuperscript{o} (p.~\pageref{fol:148}), infra f\textsuperscript{o}~75~v\textsuperscript{o} (p.~\pageref{fol:151}), f\textsuperscript{o}~77~r\textsuperscript{o} (p.~\pageref{fol:154}), \textit{et passim}.\par\vspace{0.55mm}
\noindent\rule{38mm}{0.4pt}\par\vspace{0.5mm}
\noindent\hspace*{5mm}\parbox{\dimexpr\textwidth-5mm\relax}{\fontsize{6.6pt}{7.7pt}\selectfont\color{black!84}Première reconnaissance pédestre le jeudi 8 février 1917 autour de son nouveau cantonnement d'\textbf{Hérissart} jusqu'au village voisin de \textbf{Rubempré} (Somme, à 3~km au sud), à travers les croupes crayeuses couronnées de moulins à vent de l'Amiénois, par un gel intense qui interdit d'écrire dans sa chambre sans feu le 9 février au matin.}\par}
\newpage
\phantomsection\label{fac:149}
% --- En-tête du fac-similé (Page impaire / droite : Photographie du manuscrit) ---
\noindent\rule{\textwidth}{0.5pt}\par\vspace{0.55mm}
\noindent\textbf{\fontsize{6.8pt}{7.8pt}\selectfont FAC-SIMILÉ EN VIS-À-VIS — FOLIO 74 VERSO} \hfill \textbf{\fontsize{6.8pt}{7.8pt}\selectfont FÉVRIER 1917}\par\vspace{0.25mm}
\noindent{\fontsize{6.8pt}{8.0pt}\selectfont\textit{Livre-Journal d'Hippolyte Combes} — 8 (suite) – 9 février 1917 (matin) (Hérissart, Rubempré (Somme))}\par\vspace{0.55mm}
\noindent\rule{\textwidth}{0.35pt}\par\vspace*{\fill}
{\centering
  \setlength{\fboxsep}{1.2mm}%
  \setlength{\fboxrule}{0.4pt}%
  \fcolorbox{gray!55}{white}{%
    \includegraphics[width=0.95\textwidth,height=0.81\textheight,keepaspectratio]{\facsimiledir/20251126_110659.jpg}%
  }\par}
\vspace*{\fill}
\newpage

% =============================================================================
% VIS-A-VIS SPREAD : FOLIO 75 RECTO (20251126_110726.jpg)
% =============================================================================
\cleardoubleevenpage
\renewcommand{\currentmonthyear}{Février 1917}
\renewcommand{\currentfoliolabel}{Folio 75 recto}
\phantomsection\label{fol:150}
% --- En-tête du folio (Page paire / gauche : Transcription, apparat & registres de bas de page) ---
\noindent\rule{\textwidth}{0.5pt}\par\vspace{0.55mm}
\noindent\textbf{\footnotesize FOLIO 75 RECTO} \hfill \textbf{\footnotesize FÉVRIER 1917}\par\vspace{0.25mm}
\noindent{\fontsize{6.9pt}{8.0pt}\selectfont Visite à la Mission de Toutencourt et course au bois d'Aveluy et Martinsart}\par\vspace{0.2mm}
\noindent{\fontsize{6.9pt}{8.0pt}\selectfont\textit{9 (suite) – 10 février 1917 (matin)} \hfill \textit{Hérissart, Toutencourt, Martinsart (Somme)}}\par\vspace{0.5mm}
\noindent\rule{\textwidth}{0.35pt}\par\vspace{0.8mm}
% --- Corps de la transcription diplomatique et apparat marginal en prolongement ---
{\renewcommand{\arraystretch}{0.87}%
\noindent\hspace*{3.5mm}\begin{tabular}{@{}p{87mm}@{\hspace{3.5mm}}p{39mm}@{}}
      \leavevmode\fontsize{6.9pt}{7.9pt}\selectfont\strut \underline{— Vendredi 9 (suite)} — \hfill{} \raisebox{0.5ex}{\underline{Février 1917}}\strut &  \\
      \leavevmode\fontsize{6.9pt}{7.9pt}\selectfont\strut Après-midi, ds l'auto du Lieutt Croughton,\strut & \smash{\parbox[t]{39mm}{\fontsize{6.2pt}{7.0pt}\selectfont\raggedright\color{black!88}\textbf{2 ds} \textit{abbr. pro} dans. \textbf{2 Lieutt} \textit{abbr. pro} Lieutenant.}} \\
      \leavevmode\fontsize{6.9pt}{7.9pt}\selectfont\strut venu à Toutencourt. Bureau Mission fr. : appris\strut &  \\
      \leavevmode\fontsize{6.9pt}{7.9pt}\selectfont\strut que désormais serai administré par l'armée et\strut &  \\
      \leavevmode\llap{\fontsize{6.4pt}{7.3pt}\selectfont\color{darkgray}5\hspace{1.8mm}}\fontsize{6.9pt}{7.9pt}\selectfont\strut payé les 3 et 18 de chaque mois.\strut &  \\
      \leavevmode\fontsize{6.9pt}{7.9pt}\selectfont\strut Retour sur 1 lorry.\strut & \smash{\parbox[t]{39mm}{\fontsize{6.2pt}{7.0pt}\selectfont\raggedright\color{black!88}\textbf{6 lorry} \textit{angl. in ms.}}} \\
      \leavevmode\fontsize{6.9pt}{7.9pt}\selectfont\strut Écrit salle café maison où chambre.\strut &  \\
      \leavevmode\fontsize{6.9pt}{7.9pt}\selectfont\strut Vers 6h s., avec Cap. W., ds les fermes\strut & \smash{\parbox[t]{39mm}{\fontsize{6.2pt}{7.0pt}\selectfont\raggedright\color{black!88}\textbf{8 ds} \textit{abbr. pro} dans. \textbf{8 Cap. W.} \textit{scil.} Capitaine Williams. — \textbf{10 Sam.} \textit{scil.} Samedi.}} \\
      \leavevmode\fontsize{6.9pt}{7.9pt}\selectfont\strut pour acheter paille \& foin.\strut &  \\
      \cline{1-1}
      \leavevmode\llap{\fontsize{6.4pt}{7.3pt}\selectfont\color{darkgray}10\hspace{1.8mm}}\fontsize{6.9pt}{7.9pt}\selectfont\strut \underline{Sam. 10.}\strut &  \\
      \leavevmode\fontsize{6.9pt}{7.9pt}\selectfont\strut Sur pied dès 7h.\strut &  \\
      \leavevmode\fontsize{6.9pt}{7.9pt}\selectfont\strut Dû sortir aussitôt pour aller aux Lignes de M.T.,\strut & \smash{\parbox[t]{39mm}{\fontsize{6.2pt}{7.0pt}\selectfont\raggedright\color{black!88}\textbf{12 M.T.} \textit{angl.} Motor Transport.}} \\
      \leavevmode\fontsize{6.9pt}{7.9pt}\selectfont\strut at the end of the road.\strut &  \\
      \leavevmode\fontsize{6.9pt}{7.9pt}\selectfont\strut Nuit \& matinée, froid vif, très pénétrant.\strut &  \\
      \leavevmode\llap{\fontsize{6.4pt}{7.3pt}\selectfont\color{darkgray}15\hspace{1.8mm}}\fontsize{6.9pt}{7.9pt}\selectfont\strut Gelé toujours très fort.\strut &  \\
      \leavevmode\fontsize{6.9pt}{7.9pt}\selectfont\strut Après le breakfast, partis en auto à Martinsart.\strut & \smash{\parbox[t]{39mm}{\fontsize{6.2pt}{7.0pt}\selectfont\raggedright\color{black!88}\textbf{16 breakfast} \textit{angl. in ms.}}} \\
      \leavevmode\fontsize{6.9pt}{7.9pt}\selectfont\strut Dévalé routes en pente. Des 2 côtés, champs couverts\strut &  \\
      \leavevmode\fontsize{6.9pt}{7.9pt}\selectfont\strut de neige durcie. Soleil se lève, de + en + chaud.\strut &  \\
      \leavevmode\fontsize{6.9pt}{7.9pt}\selectfont\strut Au château de Martinsart, reçus par\strut &  \\
      \leavevmode\llap{\fontsize{6.4pt}{7.3pt}\selectfont\color{darkgray}20\hspace{1.8mm}}\fontsize{6.9pt}{7.9pt}\selectfont\strut un génal [Volké ?], très aimable.\strut & \smash{\parbox[t]{39mm}{\fontsize{6.2pt}{7.0pt}\selectfont\raggedright\color{black!88}\textbf{20} \textit{lect. incert. in ms.}}} \\
      \leavevmode\fontsize{6.9pt}{7.9pt}\selectfont\strut Village, rempli de troupes anglaises.\strut &  \\
\end{tabular}}%
\par\vspace*{\fill}
% --- Registres de bas de page (1. Critique interne [si pertinent] ; 2. Loci paralleli en latin ; 3. Commentaire historique en retrait) ---
\noindent\rule{\textwidth}{0.5pt}\par\vspace{0.65mm}
{\fontsize{6.6pt}{7.7pt}\selectfont%
\noindent\normalfont\color{black!92} \textbf{2} (\textit{Lt Croughton}) cf. supra f\textsuperscript{o}~54~v\textsuperscript{o} (p.~\pageref{fol:109}). — \textbf{3} (\textit{Toutencourt}) cf. supra f\textsuperscript{o}~49~v\textsuperscript{o} (p.~\pageref{fol:99}), f\textsuperscript{o}~62~r\textsuperscript{o} (p.~\pageref{fol:124}), infra f\textsuperscript{o}~76~r\textsuperscript{o} (p.~\pageref{fol:152}), f\textsuperscript{o}~77~r\textsuperscript{o} (p.~\pageref{fol:154}), \textit{et passim}. — \textbf{9} (\textit{paille \& foin}) cf. supra f\textsuperscript{o}~3~r\textsuperscript{o} (p.~\pageref{fol:6}), f\textsuperscript{o}~35~v\textsuperscript{o} (p.~\pageref{fol:71}), infra f\textsuperscript{o}~77~v\textsuperscript{o} (p.~\pageref{fol:155}), f\textsuperscript{o}~78~v\textsuperscript{o} (p.~\pageref{fol:157}), \textit{et passim}.\par\vspace{0.55mm}
\noindent\rule{38mm}{0.4pt}\par\vspace{0.5mm}
\noindent\hspace*{5mm}\parbox{\dimexpr\textwidth-5mm\relax}{\fontsize{6.6pt}{7.7pt}\selectfont\color{black!84}Conduit le 9 février 1917 par le \textbf{lieutenant Croughton} au bureau de la Mission militaire française de la \textbf{V\textsuperscript{e} Armée britannique} (\textit{Fifth Army}) à \textbf{Toutencourt}, Hippolyte Combes apprend qu'il relève désormais administrativement de la Mission d'armée, avec paiement de la solde à dates fixes les \textbf{3 et 18 de chaque mois}. Le lendemain, après un passage aux lignes du parc automobile (\textit{Motor Transport}, \textit{M.T. A.S.C.}), il se rend en tournée au château de \textbf{Martinsart}.}\par}
\newpage
\phantomsection\label{fac:150}
% --- En-tête du fac-similé (Page impaire / droite : Photographie du manuscrit) ---
\noindent\rule{\textwidth}{0.5pt}\par\vspace{0.55mm}
\noindent\textbf{\fontsize{6.8pt}{7.8pt}\selectfont FAC-SIMILÉ EN VIS-À-VIS — FOLIO 75 RECTO} \hfill \textbf{\fontsize{6.8pt}{7.8pt}\selectfont FÉVRIER 1917}\par\vspace{0.25mm}
\noindent{\fontsize{6.8pt}{8.0pt}\selectfont\textit{Livre-Journal d'Hippolyte Combes} — 9 (suite) – 10 février 1917 (matin) (Hérissart, Toutencourt, Martinsart (Somme))}\par\vspace{0.55mm}
\noindent\rule{\textwidth}{0.35pt}\par\vspace*{\fill}
{\centering
  \setlength{\fboxsep}{1.2mm}%
  \setlength{\fboxrule}{0.4pt}%
  \fcolorbox{gray!55}{white}{%
    \includegraphics[width=0.95\textwidth,height=0.81\textheight,keepaspectratio]{\facsimiledir/20251126_110726.jpg}%
  }\par}
\vspace*{\fill}
\newpage

% =============================================================================
% VIS-A-VIS SPREAD : FOLIO 75 VERSO (20251126_110737.jpg)
% =============================================================================
\cleardoubleevenpage
\renewcommand{\currentmonthyear}{Février 1917}
\renewcommand{\currentfoliolabel}{Folio 75 verso}
\phantomsection\label{fol:151}
% --- En-tête du folio (Page paire / gauche : Transcription, apparat & registres de bas de page) ---
\noindent\rule{\textwidth}{0.5pt}\par\vspace{0.55mm}
\noindent\textbf{\footnotesize FOLIO 75 VERSO} \hfill \textbf{\footnotesize FÉVRIER 1917}\par\vspace{0.25mm}
\noindent{\fontsize{6.9pt}{8.0pt}\selectfont Visite à son ancienne compagnie au bois d'Aveluy et combat aérien}\par\vspace{0.2mm}
\noindent{\fontsize{6.9pt}{8.0pt}\selectfont\textit{10 (suite) – 11 février 1917 (matin)} \hfill \textit{Bois d'Aveluy, Senlis-le-Sec, Acheux, Hérissart (Somme)}}\par\vspace{0.5mm}
\noindent\rule{\textwidth}{0.35pt}\par\vspace{0.8mm}
% --- Corps de la transcription diplomatique et apparat marginal en prolongement ---
{\renewcommand{\arraystretch}{0.87}%
\noindent\hspace*{3.5mm}\begin{tabular}{@{}p{87mm}@{\hspace{3.5mm}}p{39mm}@{}}
      \leavevmode\fontsize{6.9pt}{7.9pt}\selectfont\strut \underline{Sam. 10 (suite)} \hfill{} \raisebox{0.5ex}{\underline{Février 1917}}\strut & \smash{\parbox[t]{39mm}{\fontsize{6.2pt}{7.0pt}\selectfont\raggedright\color{black!88}\textbf{1 Sam.} \textit{scil.} Samedi.}} \\
      \leavevmode\fontsize{6.9pt}{7.9pt}\selectfont\strut Continué en auto route d'Aveluy, mon ancienne\strut &  \\
      \leavevmode\fontsize{6.9pt}{7.9pt}\selectfont\strut Cy, campée champ dévasté, sous des tentes.\strut &  \\
      \leavevmode\fontsize{6.9pt}{7.9pt}\selectfont\strut Lunché au mess, où vu Lt Bradfield (très aimable).\strut & \smash{\parbox[t]{39mm}{\fontsize{6.2pt}{7.0pt}\selectfont\raggedright\color{black!88}\textbf{4 Lt} \textit{abbr. pro} Lieutenant. \textbf{4 Lunché} \textit{angl. in ms.} — \textbf{5 qq} \textit{abbr. pro} quelque(s). \textbf{5 cantonnemt} \textit{abbr. pro} cantonnement. — \textbf{6 ts} \textit{abbr. pro} tous. — \textbf{9 ds} \textit{abbr. pro} dans. — \textbf{10 refilling point} \textit{angl. in ms.} — \textbf{11 Cap.} \textit{abbr. pro} Capitaine. — \textbf{12 1/2} \textit{scil.} demi / moitié. — \textbf{13 Cap.} \textit{abbr. pro} Capitaine.}} \\
      \leavevmode\llap{\fontsize{6.4pt}{7.3pt}\selectfont\color{darkgray}5\hspace{1.8mm}}\fontsize{6.9pt}{7.9pt}\selectfont\strut Puis visité cantonnemt et causé avec qq N.C.O's \& men.\strut &  \\
      \leavevmode\fontsize{6.9pt}{7.9pt}\selectfont\strut Ciel très bleu. Soleil brûlant. De ts côtés,\strut &  \\
      \leavevmode\fontsize{6.9pt}{7.9pt}\selectfont\strut vastes horizons. Canons de marine, en batterie\strut &  \\
      \leavevmode\fontsize{6.9pt}{7.9pt}\selectfont\strut \& cachés sous des branchages desséchés.\strut &  \\
      \leavevmode\fontsize{6.9pt}{7.9pt}\selectfont\strut 2 avions boches, canonnés ds l'azur —\strut &  \\
      \leavevmode\llap{\fontsize{6.4pt}{7.3pt}\selectfont\color{darkgray}10\hspace{1.8mm}}\fontsize{6.9pt}{7.9pt}\selectfont\strut Au retour, arrêté au refilling point, où revu\strut &  \\
      \leavevmode\fontsize{6.9pt}{7.9pt}\selectfont\strut Cap. Lovet. Charmant.\strut &  \\
      \leavevmode\fontsize{6.9pt}{7.9pt}\selectfont\strut Rentrés vers 3h 1/2, par Senlis. A la miss. fr.\strut &  \\
      \leavevmode\fontsize{6.9pt}{7.9pt}\selectfont\strut 2e corps, vu Cap. de Casteljau, énorme \& bête.\strut &  \\
      \leavevmode\fontsize{6.9pt}{7.9pt}\selectfont\strut Passé par Acheux.\strut &  \\
      \leavevmode\llap{\fontsize{6.4pt}{7.3pt}\selectfont\color{darkgray}15\hspace{1.8mm}}\fontsize{6.9pt}{7.9pt}\selectfont\strut Soleil rouge disparaissait à l'horizon.\strut &  \\
      \leavevmode\fontsize{6.9pt}{7.9pt}\selectfont\strut Nuit étoilée et sans lune.\strut &  \\
      \cline{1-1}
      \leavevmode\fontsize{6.9pt}{7.9pt}\selectfont\strut \underline{Dim. 11} — Nuit toujours très froide. Gelé\strut & \smash{\parbox[t]{39mm}{\fontsize{6.2pt}{7.0pt}\selectfont\raggedright\color{black!88}\textbf{17 Dim.} \textit{scil.} Dimanche.}} \\
      \leavevmode\fontsize{6.9pt}{7.9pt}\selectfont\strut extrt fort.\strut & \smash{\parbox[t]{39mm}{\fontsize{6.2pt}{7.0pt}\selectfont\raggedright\color{black!88}\textbf{18 extrt} \textit{abbr. pro} extrêmement. — \textbf{19 Ds} \textit{abbr. pro} Dans. \textbf{19 1/2 du jour} \textit{scil.} milieu du jour.}} \\
      \leavevmode\fontsize{6.9pt}{7.9pt}\selectfont\strut Ds le 1/2 du jour, ardent soleil réchauffe\strut &  \\
      \leavevmode\llap{\fontsize{6.4pt}{7.3pt}\selectfont\color{darkgray}20\hspace{1.8mm}}\fontsize{6.9pt}{7.9pt}\selectfont\strut sensiblement température. Bientôt ce sera\strut &  \\
      \leavevmode\fontsize{6.9pt}{7.9pt}\selectfont\strut le dégel, et fin de l'hiver.\strut &  \\
\end{tabular}}%
\par\vspace*{\fill}
% --- Registres de bas de page (1. Critique interne [si pertinent] ; 2. Loci paralleli en latin ; 3. Commentaire historique en retrait) ---
\noindent\rule{\textwidth}{0.5pt}\par\vspace{0.65mm}
{\fontsize{6.6pt}{7.7pt}\selectfont%
\noindent\normalfont\color{black!92} \textbf{2} (\textit{Aveluy}) cf. supra f\textsuperscript{o}~41~r\textsuperscript{o} (p.~\pageref{fol:82}), f\textsuperscript{o}~75~r\textsuperscript{o} (p.~\pageref{fol:150}). — \textbf{12} (\textit{Senlis}) cf. supra f\textsuperscript{o}~37~v\textsuperscript{o} (p.~\pageref{fol:75}), f\textsuperscript{o}~62~r\textsuperscript{o} (p.~\pageref{fol:124}), infra f\textsuperscript{o}~79~v\textsuperscript{o} (p.~\pageref{fol:159}), f\textsuperscript{o}~84~v\textsuperscript{o} (p.~\pageref{fol:169}), \textit{et passim}. — \textbf{14} (\textit{Acheux}) cf. supra f\textsuperscript{o}~21~v\textsuperscript{o} (p.~\pageref{fol:43}), f\textsuperscript{o}~74~v\textsuperscript{o} (p.~\pageref{fol:149}), infra f\textsuperscript{o}~77~r\textsuperscript{o} (p.~\pageref{fol:154}), f\textsuperscript{o}~84~r\textsuperscript{o} (p.~\pageref{fol:168}), \textit{et passim}.\par\vspace{0.55mm}
\noindent\rule{38mm}{0.4pt}\par\vspace{0.5mm}
\noindent\hspace*{5mm}\parbox{\dimexpr\textwidth-5mm\relax}{\fontsize{6.6pt}{7.7pt}\selectfont\color{black!84}Suite de la tournée du samedi 10 février 1917: sur la route d'\textbf{Aveluy} (au nord d'Albert), Hippolyte Combes rend visite à son ancienne compagnie (\textbf{\textit{152nd Company A.S.C.}} du \textit{18th Divisional Train}), campée sous tentes dans un terrain dévasté à proximité des pièces britanniques de marine à longue portée (\textit{naval guns}) camouflées sous des branchages.}\par}
\newpage
\phantomsection\label{fac:151}
% --- En-tête du fac-similé (Page impaire / droite : Photographie du manuscrit) ---
\noindent\rule{\textwidth}{0.5pt}\par\vspace{0.55mm}
\noindent\textbf{\fontsize{6.8pt}{7.8pt}\selectfont FAC-SIMILÉ EN VIS-À-VIS — FOLIO 75 VERSO} \hfill \textbf{\fontsize{6.8pt}{7.8pt}\selectfont FÉVRIER 1917}\par\vspace{0.25mm}
\noindent{\fontsize{6.8pt}{8.0pt}\selectfont\textit{Livre-Journal d'Hippolyte Combes} — 10 (suite) – 11 février 1917 (matin) (Bois d'Aveluy, Senlis-le-Sec, Acheux, Hérissart (Somme))}\par\vspace{0.55mm}
\noindent\rule{\textwidth}{0.35pt}\par\vspace*{\fill}
{\centering
  \setlength{\fboxsep}{1.2mm}%
  \setlength{\fboxrule}{0.4pt}%
  \fcolorbox{gray!55}{white}{%
    \includegraphics[width=0.95\textwidth,height=0.81\textheight,keepaspectratio]{\facsimiledir/20251126_110737.jpg}%
  }\par}
\vspace*{\fill}
\newpage

% =============================================================================
% VIS-A-VIS SPREAD : FOLIO 76 RECTO (20251126_110753.jpg)
% =============================================================================
\cleardoubleevenpage
\renewcommand{\currentmonthyear}{Février 1917}
\renewcommand{\currentfoliolabel}{Folio 76 recto}
\phantomsection\label{fol:152}
% --- En-tête du folio (Page paire / gauche : Transcription, apparat & registres de bas de page) ---
\noindent\rule{\textwidth}{0.5pt}\par\vspace{0.55mm}
\noindent\textbf{\footnotesize FOLIO 76 RECTO} \hfill \textbf{\footnotesize FÉVRIER 1917}\par\vspace{0.25mm}
\noindent{\fontsize{6.9pt}{8.0pt}\selectfont Excursion dominicale à Doullens et Auxi-le-Château chez M. Romert}\par\vspace{0.2mm}
\noindent{\fontsize{6.9pt}{8.0pt}\selectfont\textit{11 février 1917 (suite)} \hfill \textit{Toutencourt, Doullens, Auxi-le-Château (Somme / Pas-de-Calais)}}\par\vspace{0.5mm}
\noindent\rule{\textwidth}{0.35pt}\par\vspace{0.8mm}
% --- Corps de la transcription diplomatique et apparat marginal en prolongement ---
{\renewcommand{\arraystretch}{0.87}%
\noindent\hspace*{3.5mm}\begin{tabular}{@{}p{87mm}@{\hspace{3.5mm}}p{39mm}@{}}
      \leavevmode\fontsize{6.9pt}{7.9pt}\selectfont\strut \underline{Dim. 11} — \hfill{} \underline{— Février 1917 —}\strut & \smash{\parbox[t]{39mm}{\fontsize{6.2pt}{7.0pt}\selectfont\raggedright\color{black!88}\textbf{1 Dim.} \textit{scil.} Dimanche.}} \\
      \leavevmode\fontsize{6.9pt}{7.9pt}\selectfont\strut Tout le jour, en course.\strut &  \\
      \leavevmode\fontsize{6.9pt}{7.9pt}\selectfont\strut Matinée, à Toutencourt.\strut &  \\
      \leavevmode\fontsize{6.9pt}{7.9pt}\selectfont\strut Déjeuné à Doullens, où arrêt de 3 h.\strut &  \\
      \leavevmode\llap{\fontsize{6.4pt}{7.3pt}\selectfont\color{darkgray}5\hspace{1.8mm}}\fontsize{6.9pt}{7.9pt}\selectfont\strut Promenade solitaire à travers campagne, autour\strut &  \\
      \leavevmode\fontsize{6.9pt}{7.9pt}\selectfont\strut de la ville, le long d'1 petit ruisseau.\strut &  \\
      \leavevmode\fontsize{6.9pt}{7.9pt}\selectfont\strut Temps très beau.\strut &  \\
      \leavevmode\fontsize{6.9pt}{7.9pt}\selectfont\strut Continué, 3h s., jusqu'à Auxi-le-Château.\strut & \smash{\parbox[t]{39mm}{\fontsize{6.2pt}{7.0pt}\selectfont\raggedright\color{black!88}\textbf{8 s.} \textit{abbr. pro} soir.}} \\
      \leavevmode\fontsize{6.9pt}{7.9pt}\selectfont\strut Vu M. Romert. M'a offert café et excellent\strut &  \\
      \leavevmode\llap{\fontsize{6.4pt}{7.3pt}\selectfont\color{darkgray}10\hspace{1.8mm}}\fontsize{6.9pt}{7.9pt}\selectfont\strut rhum (1870) et kirsch Forêt Noire. Causé de\strut &  \\
      \leavevmode\fontsize{6.9pt}{7.9pt}\selectfont\strut tout un peu. Vieux célibataire, héritier de 3\strut &  \\
      \leavevmode\fontsize{6.9pt}{7.9pt}\selectfont\strut générations, laissera fortune à des étrangers.\strut &  \\
      \leavevmode\fontsize{6.9pt}{7.9pt}\selectfont\strut Au retour, soleil disparu, air de + en +\strut &  \\
      \leavevmode\fontsize{6.9pt}{7.9pt}\selectfont\strut glacé. Commence à geler très durement.\strut &  \\
      \leavevmode\llap{\fontsize{6.4pt}{7.3pt}\selectfont\color{darkgray}15\hspace{1.8mm}}\fontsize{6.9pt}{7.9pt}\selectfont\strut Nuit obscure, ciel étoilé mais sans lune.\strut &  \\
      \leavevmode\fontsize{6.9pt}{7.9pt}\selectfont\strut Il semble pourtant qu'il fasse un peu\strut &  \\
      \leavevmode\fontsize{6.9pt}{7.9pt}\selectfont\strut moins froid.\strut &  \\
      \cline{1-1}
      \leavevmode\fontsize{6.9pt}{7.9pt}\selectfont\strut [ Doullens : Papeterie : 2 f. 10. —\strut & \smash{\parbox[t]{39mm}{\fontsize{6.2pt}{7.0pt}\selectfont\raggedright\color{black!88}\textbf{18} \textit{in marg. cum unco} (comptes). \textbf{18 f.} \textit{abbr. pro} franc(s). — \textbf{19 Cartes ples} \textit{abbr. pro} Cartes postales. \textbf{19 c.} \textit{abbr. pro} centimes. — \textbf{20 c.} \textit{abbr. pro} centimes.}} \\
      \leavevmode\fontsize{6.9pt}{7.9pt}\selectfont\strut Cartes ples : 75 c.\strut &  \\
      \leavevmode\llap{\fontsize{6.4pt}{7.3pt}\selectfont\color{darkgray}20\hspace{1.8mm}}\fontsize{6.9pt}{7.9pt}\selectfont\strut Chocolat : 60 c. — Biscuits Petit Beurre : 50 c. ]\strut &  \\
\end{tabular}}%
\par\vspace*{\fill}
% --- Registres de bas de page (1. Critique interne [si pertinent] ; 2. Loci paralleli en latin ; 3. Commentaire historique en retrait) ---
\noindent\rule{\textwidth}{0.5pt}\par\vspace{0.65mm}
{\fontsize{6.6pt}{7.7pt}\selectfont%
\noindent\normalfont\color{black!92} \textbf{3} (\textit{Toutencourt}) cf. supra f\textsuperscript{o}~49~v\textsuperscript{o} (p.~\pageref{fol:99}), f\textsuperscript{o}~75~r\textsuperscript{o} (p.~\pageref{fol:150}), infra f\textsuperscript{o}~77~r\textsuperscript{o} (p.~\pageref{fol:154}), f\textsuperscript{o}~77~v\textsuperscript{o} (p.~\pageref{fol:155}), \textit{et passim}. — \textbf{4} (\textit{Doullens}) cf. supra f\textsuperscript{o}~25~v\textsuperscript{o} (p.~\pageref{fol:51}), f\textsuperscript{o}~73~r\textsuperscript{o} (p.~\pageref{fol:146}), infra f\textsuperscript{o}~79~r\textsuperscript{o} (p.~\pageref{fol:158}), f\textsuperscript{o}~81~v\textsuperscript{o} (p.~\pageref{fol:163}), \textit{et passim}. — \textbf{8} (\textit{Auxi-le-Château}) cf. supra f\textsuperscript{o}~62~r\textsuperscript{o} (p.~\pageref{fol:124}), infra f\textsuperscript{o}~79~r\textsuperscript{o} (p.~\pageref{fol:158}).\par\vspace{0.55mm}
\noindent\rule{38mm}{0.4pt}\par\vspace{0.5mm}
\noindent\hspace*{5mm}\parbox{\dimexpr\textwidth-5mm\relax}{\fontsize{6.6pt}{7.7pt}\selectfont\color{black!84}Longue tournée dominicale du 11 février 1917 par \textbf{Toutencourt} et \textbf{Doullens} jusqu'à \textbf{Auxi-le-Château} (Pas-de-Calais, déjà visité le 17 janvier, f\textsuperscript{o}~62~r\textsuperscript{o}), où Hippolyte Combes est reçu par un propriétaire terrien aisé, \textbf{M.~Romert}, qui lui sert un rhum millésimé de 1870 et du kirsch de la Forêt-Noire.}\par}
\newpage
\phantomsection\label{fac:152}
% --- En-tête du fac-similé (Page impaire / droite : Photographie du manuscrit) ---
\noindent\rule{\textwidth}{0.5pt}\par\vspace{0.55mm}
\noindent\textbf{\fontsize{6.8pt}{7.8pt}\selectfont FAC-SIMILÉ EN VIS-À-VIS — FOLIO 76 RECTO} \hfill \textbf{\fontsize{6.8pt}{7.8pt}\selectfont FÉVRIER 1917}\par\vspace{0.25mm}
\noindent{\fontsize{6.8pt}{8.0pt}\selectfont\textit{Livre-Journal d'Hippolyte Combes} — 11 février 1917 (suite) (Toutencourt, Doullens, Auxi-le-Château (Somme / Pas-de-Calais))}\par\vspace{0.55mm}
\noindent\rule{\textwidth}{0.35pt}\par\vspace*{\fill}
{\centering
  \setlength{\fboxsep}{1.2mm}%
  \setlength{\fboxrule}{0.4pt}%
  \fcolorbox{gray!55}{white}{%
    \includegraphics[width=0.95\textwidth,height=0.81\textheight,keepaspectratio]{\facsimiledir/20251126_110753.jpg}%
  }\par}
\vspace*{\fill}
\newpage

% =============================================================================
% VIS-A-VIS SPREAD : FOLIO 76 VERSO (20251126_110804.jpg)
% =============================================================================
\cleardoubleevenpage
\renewcommand{\currentmonthyear}{Février 1917}
\renewcommand{\currentfoliolabel}{Folio 76 verso}
\phantomsection\label{fol:153}
% --- En-tête du folio (Page paire / gauche : Transcription, apparat & registres de bas de page) ---
\noindent\rule{\textwidth}{0.5pt}\par\vspace{0.55mm}
\noindent\textbf{\footnotesize FOLIO 76 VERSO} \hfill \textbf{\footnotesize FÉVRIER 1917}\par\vspace{0.25mm}
\noindent{\fontsize{6.9pt}{8.0pt}\selectfont Visite de la minoterie de Méaulte et silhouettes des moulins d'Hérissart}\par\vspace{0.2mm}
\noindent{\fontsize{6.9pt}{8.0pt}\selectfont\textit{12 février 1917} \hfill \textit{Hérissart, Méaulte, Ville-sur-Ancre, Morlancourt, Ribemont (Somme)}}\par\vspace{0.5mm}
\noindent\rule{\textwidth}{0.35pt}\par\vspace{0.8mm}
% --- Corps de la transcription diplomatique et apparat marginal en prolongement ---
{\renewcommand{\arraystretch}{0.87}%
\noindent\hspace*{3.5mm}\begin{tabular}{@{}p{87mm}@{\hspace{3.5mm}}p{39mm}@{}}
      \leavevmode\fontsize{6.9pt}{7.9pt}\selectfont\strut \underline{Lundi 12} — \hfill{} \underline{Février 1917}\strut &  \\
      \leavevmode\fontsize{6.9pt}{7.9pt}\selectfont\strut Gelé nuit, mais bcp moins fort\strut & \smash{\parbox[t]{39mm}{\fontsize{6.2pt}{7.0pt}\selectfont\raggedright\color{black!88}\textbf{2 bcp} \textit{abbr. pro} beaucoup.}} \\
      \leavevmode\fontsize{6.9pt}{7.9pt}\selectfont\strut que jours précédents.\strut &  \\
      \leavevmode\fontsize{6.9pt}{7.9pt}\selectfont\strut Course en auto. Venus à Méaulte. Visité en détail\strut &  \\
      \leavevmode\llap{\fontsize{6.4pt}{7.3pt}\selectfont\color{darkgray}5\hspace{1.8mm}}\fontsize{6.9pt}{7.9pt}\selectfont\strut magnifique minoterie.\strut &  \\
      \leavevmode\fontsize{6.9pt}{7.9pt}\selectfont\strut Retour par Ville, Treux \& Méricourt. Ribemont.\strut &  \\
      \leavevmode\fontsize{6.9pt}{7.9pt}\selectfont\strut Repassé avec plaisir ds région qui m'est si familière\strut & \smash{\parbox[t]{39mm}{\fontsize{6.2pt}{7.0pt}\selectfont\raggedright\color{black!88}\textbf{7 ds} \textit{abbr. pro} dans.}} \\
      \leavevmode\fontsize{6.9pt}{7.9pt}\selectfont\strut et ces villages où vécu heures bien vides \& tristes.\strut &  \\
      \leavevmode\fontsize{6.9pt}{7.9pt}\selectfont\strut Ciel mort. Des champs couverts de neige, s'élevait\strut &  \\
      \leavevmode\llap{\fontsize{6.4pt}{7.3pt}\selectfont\color{darkgray}10\hspace{1.8mm}}\fontsize{6.9pt}{7.9pt}\selectfont\strut buée grise.\strut &  \\
      \leavevmode\fontsize{6.9pt}{7.9pt}\selectfont\strut En approchant de Hérissart, vu nombreux\strut &  \\
      \leavevmode\fontsize{6.9pt}{7.9pt}\selectfont\strut moulins à vent, détachant leurs ailes squelettiques\strut &  \\
      \leavevmode\fontsize{6.9pt}{7.9pt}\selectfont\strut sur ciel gris.\strut &  \\
      \leavevmode\fontsize{6.9pt}{7.9pt}\selectfont\strut Après-midi, liberté.\strut &  \\
      \leavevmode\llap{\fontsize{6.4pt}{7.3pt}\selectfont\color{darkgray}15\hspace{1.8mm}}\fontsize{6.9pt}{7.9pt}\selectfont\strut 2 h. dans ma chambrette. Achevé de m'installer\strut &  \\
      \leavevmode\fontsize{6.9pt}{7.9pt}\selectfont\strut \& rangé affaires ds mon sac.\strut & \smash{\parbox[t]{39mm}{\fontsize{6.2pt}{7.0pt}\selectfont\raggedright\color{black!88}\textbf{16 ds} \textit{abbr. pro} dans.}} \\
      \leavevmode\fontsize{6.9pt}{7.9pt}\selectfont\strut Dégel très rapide. Croûtes de glace tombant\strut &  \\
      \leavevmode\fontsize{6.9pt}{7.9pt}\selectfont\strut des toits.\strut &  \\
      \leavevmode\fontsize{6.9pt}{7.9pt}\selectfont\strut Parcouru 1 peu village d'Hérissart. Semble\strut &  \\
      \leavevmode\llap{\fontsize{6.4pt}{7.3pt}\selectfont\color{darkgray}20\hspace{1.8mm}}\fontsize{6.9pt}{7.9pt}\selectfont\strut inhabité. Rues en pente. Silence funèbre.\strut &  \\
      \cline{1-1}
      \leavevmode\fontsize{6.9pt}{7.9pt}\selectfont\strut [ 5 vues d'Amiens, cathédrale, achetées pour moi :\strut & \smash{\parbox[t]{39mm}{\fontsize{6.2pt}{7.0pt}\selectfont\raggedright\color{black!88}\textbf{21} \textit{in marg. cum unco} (comptes). — \textbf{22 f.} \textit{abbr. pro} franc(s).}} \\
      \leavevmode\fontsize{6.9pt}{7.9pt}\selectfont\strut 2 f. 50 ]\strut &  \\
\end{tabular}}%
\par\vspace*{\fill}
% --- Registres de bas de page (1. Critique interne [si pertinent] ; 2. Loci paralleli en latin ; 3. Commentaire historique en retrait) ---
\noindent\rule{\textwidth}{0.5pt}\par\vspace{0.65mm}
{\fontsize{6.6pt}{7.7pt}\selectfont%
\noindent\normalfont\color{black!92} \textbf{6} (\textit{Ribemont}) cf. supra f\textsuperscript{o}~52~v\textsuperscript{o} (p.~\pageref{fol:105}), infra f\textsuperscript{o}~84~v\textsuperscript{o} (p.~\pageref{fol:169}), f\textsuperscript{o}~90~v\textsuperscript{o} (p.~\pageref{fol:181}). — \textbf{11} (\textit{Hérissart}) cf. supra f\textsuperscript{o}~47~v\textsuperscript{o} (p.~\pageref{fol:95}), f\textsuperscript{o}~75~v\textsuperscript{o} (p.~\pageref{fol:151}), infra f\textsuperscript{o}~77~r\textsuperscript{o} (p.~\pageref{fol:154}), f\textsuperscript{o}~77~v\textsuperscript{o} (p.~\pageref{fol:155}), \textit{et passim}. — \textbf{21} (\textit{Amiens}) cf. supra f\textsuperscript{o}~25~v\textsuperscript{o} (p.~\pageref{fol:51}), f\textsuperscript{o}~74~r\textsuperscript{o} (p.~\pageref{fol:148}), infra f\textsuperscript{o}~85~v\textsuperscript{o} (p.~\pageref{fol:171}), f\textsuperscript{o}~91~r\textsuperscript{o} (p.~\pageref{fol:182}), \textit{et passim}.\par\vspace{0.55mm}
\noindent\rule{38mm}{0.4pt}\par\vspace{0.5mm}
\noindent\hspace*{5mm}\parbox{\dimexpr\textwidth-5mm\relax}{\fontsize{6.6pt}{7.7pt}\selectfont\color{black!84}Le lundi 12 février 1917, après l'inspection de la grande minoterie industrielle de \textbf{Méaulte} (au sud d'Albert sur l'Ancre), l'itinéraire de retour traverse \textbf{Ville-sur-Ancre}, \textbf{Treux}, \textbf{Méricourt-l'Abbé} et \textbf{Ribemont-sur-Ancre}: c'est dans ce secteur de la basse vallée de l'Ancre (entre Flesselles, Heilly, Ribemont et Morlancourt) qu'Hippolyte Combes avait vécu ses premiers mois de campagne avec la \textbf{18\textsuperscript{e} Division britannique} d'août 1915 au printemps 1916.}\par}
\newpage
\phantomsection\label{fac:153}
% --- En-tête du fac-similé (Page impaire / droite : Photographie du manuscrit) ---
\noindent\rule{\textwidth}{0.5pt}\par\vspace{0.55mm}
\noindent\textbf{\fontsize{6.8pt}{7.8pt}\selectfont FAC-SIMILÉ EN VIS-À-VIS — FOLIO 76 VERSO} \hfill \textbf{\fontsize{6.8pt}{7.8pt}\selectfont FÉVRIER 1917}\par\vspace{0.25mm}
\noindent{\fontsize{6.8pt}{8.0pt}\selectfont\textit{Livre-Journal d'Hippolyte Combes} — 12 février 1917 (Hérissart, Méaulte, Ville-sur-Ancre, Morlancourt, Ribemont (Somme))}\par\vspace{0.55mm}
\noindent\rule{\textwidth}{0.35pt}\par\vspace*{\fill}
{\centering
  \setlength{\fboxsep}{1.2mm}%
  \setlength{\fboxrule}{0.4pt}%
  \fcolorbox{gray!55}{white}{%
    \includegraphics[width=0.95\textwidth,height=0.81\textheight,keepaspectratio]{\facsimiledir/20251126_110804.jpg}%
  }\par}
\vspace*{\fill}
\newpage

% =============================================================================
% VIS-A-VIS SPREAD : FOLIO 77 RECTO (20251126_110831.jpg)
% =============================================================================
\cleardoubleevenpage
\renewcommand{\currentmonthyear}{Février 1917}
\renewcommand{\currentfoliolabel}{Folio 77 recto}
\phantomsection\label{fol:154}
% --- En-tête du folio (Page paire / gauche : Transcription, apparat & registres de bas de page) ---
\noindent\rule{\textwidth}{0.5pt}\par\vspace{0.55mm}
\noindent\textbf{\footnotesize FOLIO 77 RECTO} \hfill \textbf{\footnotesize FÉVRIER 1917}\par\vspace{0.25mm}
\noindent{\fontsize{6.9pt}{8.0pt}\selectfont Courses à Authie, Bus et Acheux : arabesques de givre et repas à l'Hôtel de Ville}\par\vspace{0.2mm}
\noindent{\fontsize{6.9pt}{8.0pt}\selectfont\textit{13 – 14 février 1917} \hfill \textit{Hérissart, Authie, Bus-lès-Artois, Gézaincourt, Acheux, Toutencourt (Somme)}}\par\vspace{0.5mm}
\noindent\rule{\textwidth}{0.35pt}\par\vspace{0.8mm}
% --- Corps de la transcription diplomatique et apparat marginal en prolongement ---
{\renewcommand{\arraystretch}{0.87}%
\noindent\hspace*{3.5mm}\begin{tabular}{@{}p{87mm}@{\hspace{3.5mm}}p{39mm}@{}}
      \leavevmode\fontsize{6.9pt}{7.9pt}\selectfont\strut \underline{Mardi 13} — \hfill{} \underline{Février 1917}\strut &  \\
      \leavevmode\fontsize{6.9pt}{7.9pt}\selectfont\strut Nuit un peu moins froide. Très peu gelé.\strut &  \\
      \leavevmode\fontsize{6.9pt}{7.9pt}\selectfont\strut Tout le jour, ciel bas, nuageux.\strut &  \\
      \leavevmode\fontsize{6.9pt}{7.9pt}\selectfont\strut Matinée, dégel commencé, assez rapide.\strut &  \\
      \leavevmode\llap{\fontsize{6.4pt}{7.3pt}\selectfont\color{darkgray}5\hspace{1.8mm}}\fontsize{6.9pt}{7.9pt}\selectfont\strut Course à Authie, par Puchevillers. Sur les\strut &  \\
      \leavevmode\fontsize{6.9pt}{7.9pt}\selectfont\strut hauteurs de P., bise glaciale, paupières,\strut &  \\
      \leavevmode\fontsize{6.9pt}{7.9pt}\selectfont\strut ruisselantes.\strut &  \\
      \leavevmode\fontsize{6.9pt}{7.9pt}\selectfont\strut Après-midi, liberté entière.\strut &  \\
      \leavevmode\fontsize{6.9pt}{7.9pt}\selectfont\strut Nuit extrt noire. Moins froid.\strut & \smash{\parbox[t]{39mm}{\fontsize{6.2pt}{7.0pt}\selectfont\raggedright\color{black!88}\textbf{9 extrt} \textit{abbr. pro} extrêmement. — \textbf{10 Merc.} \textit{scil.} Mercredi.}} \\
      \cline{1-1}
      \leavevmode\llap{\fontsize{6.4pt}{7.3pt}\selectfont\color{darkgray}10\hspace{1.8mm}}\fontsize{6.9pt}{7.9pt}\selectfont\strut \underline{Merc. 14.}\strut &  \\
      \leavevmode\fontsize{6.9pt}{7.9pt}\selectfont\strut Nuit, gelé encore. Au réveil, trouvé mes\strut &  \\
      \leavevmode\fontsize{6.9pt}{7.9pt}\selectfont\strut vitres embellies d'arabesques magnifiques.\strut &  \\
      \leavevmode\fontsize{6.9pt}{7.9pt}\selectfont\strut Journée, beau soleil fait en plusieurs endroits\strut &  \\
      \leavevmode\fontsize{6.9pt}{7.9pt}\selectfont\strut fondre glace \& neige.\strut &  \\
      \leavevmode\llap{\fontsize{6.4pt}{7.3pt}\selectfont\color{darkgray}15\hspace{1.8mm}}\fontsize{6.9pt}{7.9pt}\selectfont\strut Après breakfast, partis en auto : Authie,\strut & \smash{\parbox[t]{39mm}{\fontsize{6.2pt}{7.0pt}\selectfont\raggedright\color{black!88}\textbf{15 breakfast} \textit{angl. in ms.}}} \\
      \leavevmode\fontsize{6.9pt}{7.9pt}\selectfont\strut Gouin[?] et Acheux, où pris repas au café\strut & \smash{\parbox[t]{39mm}{\fontsize{6.2pt}{7.0pt}\selectfont\raggedright\color{black!88}\textbf{16 Gouin} \textit{lect. incert. in ms.}}} \\
      \leavevmode\fontsize{6.9pt}{7.9pt}\selectfont\strut de l'Hôtel de Ville (prix exorbitants).\strut &  \\
      \leavevmode\fontsize{6.9pt}{7.9pt}\selectfont\strut Au retour, arrêt à Toutencourt. Y revenus\strut &  \\
      \leavevmode\fontsize{6.9pt}{7.9pt}\selectfont\strut après tea, vers 6h s., pour rencontrer\strut & \smash{\parbox[t]{39mm}{\fontsize{6.2pt}{7.0pt}\selectfont\raggedright\color{black!88}\textbf{19 s.} \textit{abbr. pro} soir. \textbf{19 tea} \textit{angl. in ms.}}} \\
      \leavevmode\llap{\fontsize{6.4pt}{7.3pt}\selectfont\color{darkgray}20\hspace{1.8mm}}\fontsize{6.9pt}{7.9pt}\selectfont\strut M. Lefèvre.\strut &  \\
      \leavevmode\fontsize{6.9pt}{7.9pt}\selectfont\strut Nuit mervt étoilée, mais sans lune.\strut & \smash{\parbox[t]{39mm}{\fontsize{6.2pt}{7.0pt}\selectfont\raggedright\color{black!88}\textbf{21 mervt} \textit{abbr. pro} merveilleusement.}} \\
      \leavevmode\fontsize{6.9pt}{7.9pt}\selectfont\strut Obscurité presque complète. Froid de nouveau très vif.\strut &  \\
\end{tabular}}%
\par\vspace*{\fill}
% --- Registres de bas de page (1. Critique interne [si pertinent] ; 2. Loci paralleli en latin ; 3. Commentaire historique en retrait) ---
\noindent\rule{\textwidth}{0.5pt}\par\vspace{0.65mm}
{\fontsize{6.6pt}{7.7pt}\selectfont%
\noindent\normalfont\color{black!92} \textbf{5} (\textit{Puchevillers}) cf. supra f\textsuperscript{o}~37~v\textsuperscript{o} (p.~\pageref{fol:75}), f\textsuperscript{o}~73~v\textsuperscript{o} (p.~\pageref{fol:147}), infra f\textsuperscript{o}~77~v\textsuperscript{o} (p.~\pageref{fol:155}), f\textsuperscript{o}~79~r\textsuperscript{o} (p.~\pageref{fol:158}), \textit{et passim}. — \textbf{5} (\textit{Authie}) cf. supra f\textsuperscript{o}~50~r\textsuperscript{o} (p.~\pageref{fol:100}), f\textsuperscript{o}~61~r\textsuperscript{o} (p.~\pageref{fol:122}), \textit{et passim}. — \textbf{16} (\textit{Acheux}) cf. supra f\textsuperscript{o}~21~v\textsuperscript{o} (p.~\pageref{fol:43}), f\textsuperscript{o}~75~v\textsuperscript{o} (p.~\pageref{fol:151}), infra f\textsuperscript{o}~84~r\textsuperscript{o} (p.~\pageref{fol:168}), f\textsuperscript{o}~88~r\textsuperscript{o} (p.~\pageref{fol:176}), \textit{et passim}.\par\vspace{0.55mm}
\noindent\rule{38mm}{0.4pt}\par\vspace{0.5mm}
\noindent\hspace*{5mm}\parbox{\dimexpr\textwidth-5mm\relax}{\fontsize{6.6pt}{7.7pt}\selectfont\color{black!84}Courses des 13 et 14 février 1917 vers \textbf{Puchevillers}, \textbf{Authie}, \textbf{Acheux-en-Amiénois} (où Hippolyte Combes déjeune comme client de passage au café de l'Hôtel de Ville) et \textbf{Toutencourt}, où il retourne à 18~h pour un rendez-vous de service avec \textbf{M.~Lefèvre}.}\par}
\newpage
\phantomsection\label{fac:154}
% --- En-tête du fac-similé (Page impaire / droite : Photographie du manuscrit) ---
\noindent\rule{\textwidth}{0.5pt}\par\vspace{0.55mm}
\noindent\textbf{\fontsize{6.8pt}{7.8pt}\selectfont FAC-SIMILÉ EN VIS-À-VIS — FOLIO 77 RECTO} \hfill \textbf{\fontsize{6.8pt}{7.8pt}\selectfont FÉVRIER 1917}\par\vspace{0.25mm}
\noindent{\fontsize{6.8pt}{8.0pt}\selectfont\textit{Livre-Journal d'Hippolyte Combes} — 13 – 14 février 1917 (Hérissart, Authie, Bus-lès-Artois, Gézaincourt, Acheux, Toutencourt (Somme))}\par\vspace{0.55mm}
\noindent\rule{\textwidth}{0.35pt}\par\vspace*{\fill}
{\centering
  \setlength{\fboxsep}{1.2mm}%
  \setlength{\fboxrule}{0.4pt}%
  \fcolorbox{gray!55}{white}{%
    \includegraphics[width=0.95\textwidth,height=0.81\textheight,keepaspectratio]{\facsimiledir/20251126_110831.jpg}%
  }\par}
\vspace*{\fill}
\newpage

% =============================================================================
% VIS-A-VIS SPREAD : FOLIO 77 VERSO (20251126_110914.jpg)
% =============================================================================
\cleardoubleevenpage
\renewcommand{\currentmonthyear}{Février 1917}
\renewcommand{\currentfoliolabel}{Folio 77 verso}
\phantomsection\label{fol:155}
% --- En-tête du folio (Page paire / gauche : Transcription, apparat & registres de bas de page) ---
\noindent\rule{\textwidth}{0.5pt}\par\vspace{0.55mm}
\noindent\textbf{\footnotesize FOLIO 77 VERSO} \hfill \textbf{\footnotesize FÉVRIER 1917}\par\vspace{0.25mm}
\noindent{\fontsize{6.9pt}{8.0pt}\selectfont Visite à la mairie d'Hérissart et tournée au château de Beaucourt}\par\vspace{0.2mm}
\noindent{\fontsize{6.9pt}{8.0pt}\selectfont\textit{15 – 16 février 1917} \hfill \textit{Hérissart, Toutencourt, Raincheval, Puchevillers, Beaucourt-en-Santerre / sur-l'Hallue (Somme)}}\par\vspace{0.5mm}
\noindent\rule{\textwidth}{0.35pt}\par\vspace{0.8mm}
% --- Corps de la transcription diplomatique et apparat marginal en prolongement ---
{\renewcommand{\arraystretch}{0.87}%
\noindent\hspace*{3.5mm}\begin{tabular}{@{}p{87mm}@{\hspace{3.5mm}}p{39mm}@{}}
      \leavevmode\fontsize{6.9pt}{7.9pt}\selectfont\strut \underline{Jeudi 15} — \hfill{} \underline{Février 1917 —}\strut &  \\
      \leavevmode\fontsize{6.9pt}{7.9pt}\selectfont\strut Gelé nuit.\strut &  \\
      \leavevmode\fontsize{6.9pt}{7.9pt}\selectfont\strut Ds la journée, très beau soleil.\strut & \smash{\parbox[t]{39mm}{\fontsize{6.2pt}{7.0pt}\selectfont\raggedright\color{black!88}\textbf{3 Ds} \textit{abbr. pro} Dans.}} \\
      \leavevmode\fontsize{6.9pt}{7.9pt}\selectfont\strut Matinée, courses à Toutencourt, Raincheval\strut &  \\
      \leavevmode\llap{\fontsize{6.4pt}{7.3pt}\selectfont\color{darkgray}5\hspace{1.8mm}}\fontsize{6.9pt}{7.9pt}\selectfont\strut \& Puchevillers.\strut &  \\
      \leavevmode\fontsize{6.9pt}{7.9pt}\selectfont\strut Après-midi, liberté.\strut &  \\
      \leavevmode\fontsize{6.9pt}{7.9pt}\selectfont\strut Dans ma chambre : écrit un peu. Puis\strut &  \\
      \leavevmode\fontsize{6.9pt}{7.9pt}\selectfont\strut promené village, grand costume. Vu le\strut &  \\
      \leavevmode\fontsize{6.9pt}{7.9pt}\selectfont\strut maire. Population, sympathique.\strut &  \\
      \leavevmode\llap{\fontsize{6.4pt}{7.3pt}\selectfont\color{darkgray}10\hspace{1.8mm}}\fontsize{6.9pt}{7.9pt}\selectfont\strut Nuit extrt belle, mervt étoilée.\strut & \smash{\parbox[t]{39mm}{\fontsize{6.2pt}{7.0pt}\selectfont\raggedright\color{black!88}\textbf{10 extrt} \textit{abbr. pro} extrêmement. \textbf{10 mervt} \textit{abbr. pro} merveilleusement.}} \\
      \leavevmode\fontsize{6.9pt}{7.9pt}\selectfont\strut Toujours gelé 1 peu. Très difficile monter\strut &  \\
      \leavevmode\fontsize{6.9pt}{7.9pt}\selectfont\strut rues en pente, car couvertes de glace.\strut &  \\
      \cline{1-1}
      \leavevmode\fontsize{6.9pt}{7.9pt}\selectfont\strut \underline{Vendredi 16} — Au réveil, trouvé encore\strut &  \\
      \leavevmode\fontsize{6.9pt}{7.9pt}\selectfont\strut vitres ornées d'admirables dessins de glace.\strut &  \\
      \leavevmode\llap{\fontsize{6.4pt}{7.3pt}\selectfont\color{darkgray}15\hspace{1.8mm}}\fontsize{6.9pt}{7.9pt}\selectfont\strut Mais, ds la journée, température, de\strut & \smash{\parbox[t]{39mm}{\fontsize{6.2pt}{7.0pt}\selectfont\raggedright\color{black!88}\textbf{15 ds} \textit{abbr. pro} dans.}} \\
      \leavevmode\fontsize{6.9pt}{7.9pt}\selectfont\strut + en + douce et dégel très rapide.\strut &  \\
      \leavevmode\fontsize{6.9pt}{7.9pt}\selectfont\strut Toute matinée, en auto. [Roulé ?] 2 ou\strut & \smash{\parbox[t]{39mm}{\fontsize{6.2pt}{7.0pt}\selectfont\raggedright\color{black!88}\textbf{17 Roulé} \textit{vel} Rôdé, \textit{lect. incert. in ms.} — \textbf{18 ds} \textit{abbr. pro} dans.}} \\
      \leavevmode\fontsize{6.9pt}{7.9pt}\selectfont\strut 3 h. ds le village de Beaucourt, fermes\strut &  \\
      \leavevmode\fontsize{6.9pt}{7.9pt}\selectfont\strut et château, pour trouver paille \& foin.\strut &  \\
      \leavevmode\llap{\fontsize{6.4pt}{7.3pt}\selectfont\color{darkgray}20\hspace{1.8mm}}\fontsize{6.9pt}{7.9pt}\selectfont\strut Ciel nuageux. Pluie menace.\strut &  \\
      \leavevmode\fontsize{6.9pt}{7.9pt}\selectfont\strut Après-midi, liberté entière.\strut &  \\
\end{tabular}}%
\par\vspace*{\fill}
% --- Registres de bas de page (1. Critique interne [si pertinent] ; 2. Loci paralleli en latin ; 3. Commentaire historique en retrait) ---
\noindent\rule{\textwidth}{0.5pt}\par\vspace{0.65mm}
{\fontsize{6.6pt}{7.7pt}\selectfont%
\noindent\normalfont\color{black!92} \textbf{4} (\textit{Raincheval}) cf. supra f\textsuperscript{o}~20~v\textsuperscript{o} (p.~\pageref{fol:41}), f\textsuperscript{o}~73~v\textsuperscript{o} (p.~\pageref{fol:147}), infra f\textsuperscript{o}~83~r\textsuperscript{o} (p.~\pageref{fol:166}), f\textsuperscript{o}~83~v\textsuperscript{o} (p.~\pageref{fol:167}), \textit{et passim}. — \textbf{4} (\textit{Toutencourt}) cf. supra f\textsuperscript{o}~49~v\textsuperscript{o} (p.~\pageref{fol:99}), f\textsuperscript{o}~77~r\textsuperscript{o} (p.~\pageref{fol:154}), infra f\textsuperscript{o}~78~r\textsuperscript{o} (p.~\pageref{fol:156}), f\textsuperscript{o}~79~v\textsuperscript{o} (p.~\pageref{fol:159}), \textit{et passim}. — \textbf{5} (\textit{Puchevillers}) cf. supra f\textsuperscript{o}~37~v\textsuperscript{o} (p.~\pageref{fol:75}), f\textsuperscript{o}~77~r\textsuperscript{o} (p.~\pageref{fol:154}), infra f\textsuperscript{o}~79~r\textsuperscript{o} (p.~\pageref{fol:158}), f\textsuperscript{o}~81~r\textsuperscript{o} (p.~\pageref{fol:162}), \textit{et passim}.\par\vspace{0.55mm}
\noindent\rule{38mm}{0.4pt}\par\vspace{0.5mm}
\noindent\hspace*{5mm}\parbox{\dimexpr\textwidth-5mm\relax}{\fontsize{6.6pt}{7.7pt}\selectfont\color{black!84}Après une visite officielle en grande tenue au maire d'\textbf{Hérissart} le 15 février 1917, Hippolyte Combes parcourt le lendemain les fermes et le château de \textbf{Beaucourt-sur-l'Hallue} (Somme, canton de Villers-Bocage, à 7~km au sud d'Hérissart) pour l'achat de paille et de foin.}\par}
\newpage
\phantomsection\label{fac:155}
% --- En-tête du fac-similé (Page impaire / droite : Photographie du manuscrit) ---
\noindent\rule{\textwidth}{0.5pt}\par\vspace{0.55mm}
\noindent\textbf{\fontsize{6.8pt}{7.8pt}\selectfont FAC-SIMILÉ EN VIS-À-VIS — FOLIO 77 VERSO} \hfill \textbf{\fontsize{6.8pt}{7.8pt}\selectfont FÉVRIER 1917}\par\vspace{0.25mm}
\noindent{\fontsize{6.8pt}{8.0pt}\selectfont\textit{Livre-Journal d'Hippolyte Combes} — 15 – 16 février 1917 (Hérissart, Toutencourt, Raincheval, Puchevillers, Beaucourt-en-Santerre / sur-l'Hallue (Somme))}\par\vspace{0.55mm}
\noindent\rule{\textwidth}{0.35pt}\par\vspace*{\fill}
{\centering
  \setlength{\fboxsep}{1.2mm}%
  \setlength{\fboxrule}{0.4pt}%
  \fcolorbox{gray!55}{white}{%
    \includegraphics[width=0.95\textwidth,height=0.81\textheight,keepaspectratio]{\facsimiledir/20251126_110914.jpg}%
  }\par}
\vspace*{\fill}
\newpage

% =============================================================================
% VIS-A-VIS SPREAD : FOLIO 78 RECTO (20251126_110930.jpg)
% =============================================================================
\cleardoubleevenpage
\renewcommand{\currentmonthyear}{Février 1917}
\renewcommand{\currentfoliolabel}{Folio 78 recto}
\phantomsection\label{fol:156}
% --- En-tête du folio (Page paire / gauche : Transcription, apparat & registres de bas de page) ---
\noindent\rule{\textwidth}{0.5pt}\par\vspace{0.55mm}
\noindent\textbf{\footnotesize FOLIO 78 RECTO} \hfill \textbf{\footnotesize FÉVRIER 1917}\par\vspace{0.25mm}
\noindent{\fontsize{6.9pt}{8.0pt}\selectfont Grand dégel de février, fondrières de boue et prône du dimanche à Toutencourt}\par\vspace{0.2mm}
\noindent{\fontsize{6.9pt}{8.0pt}\selectfont\textit{17 – 18 février 1917} \hfill \textit{Hérissart, Toutencourt (Somme)}}\par\vspace{0.5mm}
\noindent\rule{\textwidth}{0.35pt}\par\vspace{0.8mm}
% --- Corps de la transcription diplomatique et apparat marginal en prolongement ---
{\renewcommand{\arraystretch}{0.83}%
\noindent\hspace*{3.5mm}\begin{tabular}{@{}p{87mm}@{\hspace{3.5mm}}p{39mm}@{}}
      \leavevmode\fontsize{6.6pt}{7.5pt}\selectfont\strut \underline{Sam. 17} — \hfill{} \underline{Février 1917}\strut & \smash{\parbox[t]{39mm}{\fontsize{6.2pt}{7.0pt}\selectfont\raggedright\color{black!88}\textbf{1 Sam.} \textit{scil.} Samedi.}} \\
      \leavevmode\fontsize{6.6pt}{7.5pt}\selectfont\strut Dégel continué nuit et jour. Bruine, matinée.\strut &  \\
      \leavevmode\fontsize{6.6pt}{7.5pt}\selectfont\strut Temps redevient de + en + doux. Pluvieux.\strut &  \\
      \leavevmode\fontsize{6.6pt}{7.5pt}\selectfont\strut Journée entière de liberté. Liquidé\strut &  \\
      \leavevmode\llap{\fontsize{6.4pt}{7.3pt}\selectfont\color{darkgray}5\hspace{1.8mm}}\fontsize{6.6pt}{7.5pt}\selectfont\strut correspce en retard.\strut & \smash{\parbox[t]{39mm}{\fontsize{6.2pt}{7.0pt}\selectfont\raggedright\color{black!88}\textbf{5 correspce} \textit{abbr. pro} correspondance.}} \\
      \leavevmode\fontsize{6.6pt}{7.5pt}\selectfont\strut Ce repos m'a fait grand bien,\strut &  \\
      \leavevmode\fontsize{6.6pt}{7.5pt}\selectfont\strut au physique \& encore plus au moral.\strut &  \\
      \leavevmode\fontsize{6.6pt}{7.5pt}\selectfont\strut Soirée, épais brouillard —\strut &  \\
      \cline{1-1}
      \leavevmode\fontsize{6.6pt}{7.5pt}\selectfont\strut \underline{Dim. 18.} — Ciel resté brouillé, nuageux.\strut & \smash{\parbox[t]{39mm}{\fontsize{6.2pt}{7.0pt}\selectfont\raggedright\color{black!88}\textbf{9 Dim.} \textit{scil.} Dimanche.}} \\
      \leavevmode\llap{\fontsize{6.4pt}{7.3pt}\selectfont\color{darkgray}10\hspace{1.8mm}}\fontsize{6.6pt}{7.5pt}\selectfont\strut Fines ondées de courte durée.\strut &  \\
      \leavevmode\fontsize{6.6pt}{7.5pt}\selectfont\strut Dégel continue, très rapide.\strut &  \\
      \leavevmode\fontsize{6.6pt}{7.5pt}\selectfont\strut Routes, de nouveau transformées en\strut &  \\
      \leavevmode\fontsize{6.6pt}{7.5pt}\selectfont\strut fondrières. On patauge ds boue liquide.\strut & \smash{\parbox[t]{39mm}{\fontsize{6.2pt}{7.0pt}\selectfont\raggedright\color{black!88}\textbf{13 ds} \textit{abbr. pro} dans.}} \\
      \leavevmode\fontsize{6.6pt}{7.5pt}\selectfont\strut Tempér., bien adoucie. Puis enfin rester ds\strut & \smash{\parbox[t]{39mm}{\fontsize{6.2pt}{7.0pt}\selectfont\raggedright\color{black!88}\textbf{14 ds} \textit{abbr. pro} dans.}} \\
      \leavevmode\llap{\fontsize{6.4pt}{7.3pt}\selectfont\color{darkgray}15\hspace{1.8mm}}\fontsize{6.6pt}{7.5pt}\selectfont\strut ma chambre, quoique sans feu.\strut &  \\
      \leavevmode\fontsize{6.6pt}{7.5pt}\selectfont\strut Jour entier de liberté.\strut &  \\
      \leavevmode\fontsize{6.6pt}{7.5pt}\selectfont\strut Matin, à Toutencourt, Mission fr. Vth Armée.\strut &  \\
      \leavevmode\fontsize{6.6pt}{7.5pt}\selectfont\strut Demandé renseignements au sujet cantonnement\strut &  \\
      \leavevmode\fontsize{6.6pt}{7.5pt}\selectfont\strut d'anglais, à Toulouse.\strut &  \\
      \leavevmode\llap{\fontsize{6.4pt}{7.3pt}\selectfont\color{darkgray}20\hspace{1.8mm}}\fontsize{6.6pt}{7.5pt}\selectfont\strut Église du village de Toutencourt. Assisté\strut &  \\
      \leavevmode\fontsize{6.6pt}{7.5pt}\selectfont\strut partie cérémonie, entendu prône —\strut &  \\
      \leavevmode\fontsize{6.6pt}{7.5pt}\selectfont\strut Retour à Hérissart sur un lorry.\strut & \smash{\parbox[t]{39mm}{\fontsize{6.2pt}{7.0pt}\selectfont\raggedright\color{black!88}\textbf{22 lorry} \textit{angl. in ms.}}} \\
      \leavevmode\fontsize{6.6pt}{7.5pt}\selectfont\strut Aujourd'hui, me suis tt à fait ressaisi. Retrouvé\strut & \smash{\parbox[t]{39mm}{\fontsize{6.2pt}{7.0pt}\selectfont\raggedright\color{black!88}\textbf{23 tt} \textit{abbr. pro} tout.}} \\
      \leavevmode\fontsize{6.6pt}{7.5pt}\selectfont\strut parfaite santé morale. — [ A Toutencourt,\strut & \smash{\parbox[t]{39mm}{\fontsize{6.2pt}{7.0pt}\selectfont\raggedright\color{black!88}\textbf{24} \textit{in marg. cum unco} (comptes). — \textbf{25 f.} \textit{abbr. pro} franc(s). \textbf{25 c.} \textit{abbr. pro} centimes.}} \\
      \leavevmode\llap{\fontsize{6.4pt}{7.3pt}\selectfont\color{darkgray}25\hspace{1.8mm}}\fontsize{6.6pt}{7.5pt}\selectfont\strut solde : 68 f. 25 c. ]\strut &  \\
\end{tabular}}%
\par\vspace*{\fill}
% --- Registres de bas de page (1. Critique interne [si pertinent] ; 2. Loci paralleli en latin ; 3. Commentaire historique en retrait) ---
\noindent\rule{\textwidth}{0.5pt}\par\vspace{0.65mm}
{\fontsize{6.6pt}{7.7pt}\selectfont%
\noindent\normalfont\color{black!92} \textbf{17} (\textit{Toutencourt}) cf. supra f\textsuperscript{o}~49~v\textsuperscript{o} (p.~\pageref{fol:99}), f\textsuperscript{o}~77~v\textsuperscript{o} (p.~\pageref{fol:155}), infra f\textsuperscript{o}~79~v\textsuperscript{o} (p.~\pageref{fol:159}), f\textsuperscript{o}~80~r\textsuperscript{o} (p.~\pageref{fol:160}), \textit{et passim}. — \textbf{19} (\textit{Toulouse}) cf. supra in tegumento (p.~\pageref{fol:1}), f\textsuperscript{o}~70~r\textsuperscript{o} (p.~\pageref{fol:140}), \textit{et passim}. — \textbf{22} (\textit{Hérissart}) cf. supra f\textsuperscript{o}~47~v\textsuperscript{o} (p.~\pageref{fol:95}), f\textsuperscript{o}~77~v\textsuperscript{o} (p.~\pageref{fol:155}), infra f\textsuperscript{o}~78~v\textsuperscript{o} (p.~\pageref{fol:157}), f\textsuperscript{o}~79~r\textsuperscript{o} (p.~\pageref{fol:158}), \textit{et passim}.\par\vspace{0.55mm}
\noindent\rule{38mm}{0.4pt}\par\vspace{0.5mm}
\noindent\hspace*{5mm}\parbox{\dimexpr\textwidth-5mm\relax}{\fontsize{6.6pt}{7.7pt}\selectfont\color{black!84}Le dimanche 18 février 1917 — jour d'échéance bimensuelle fixé au f\textsuperscript{o}~75~r\textsuperscript{o} —, Hippolyte Combes se rend au bureau de la Mission française de la \textbf{V\textsuperscript{e} Armée britannique} (\textit{Fifth Army}) à \textbf{Toutencourt}, où il perçoit 68~fr.~25 de solde (13 jours à 5~fr.~25), s'informe de l'existence éventuelle d'un détachement ou cantonnement britannique à \textbf{Toulouse} et assiste au prône dominical à l'église \textbf{Saint-Léger de Toutencourt}.}\par}
\newpage
\phantomsection\label{fac:156}
% --- En-tête du fac-similé (Page impaire / droite : Photographie du manuscrit) ---
\noindent\rule{\textwidth}{0.5pt}\par\vspace{0.55mm}
\noindent\textbf{\fontsize{6.8pt}{7.8pt}\selectfont FAC-SIMILÉ EN VIS-À-VIS — FOLIO 78 RECTO} \hfill \textbf{\fontsize{6.8pt}{7.8pt}\selectfont FÉVRIER 1917}\par\vspace{0.25mm}
\noindent{\fontsize{6.8pt}{8.0pt}\selectfont\textit{Livre-Journal d'Hippolyte Combes} — 17 – 18 février 1917 (Hérissart, Toutencourt (Somme))}\par\vspace{0.55mm}
\noindent\rule{\textwidth}{0.35pt}\par\vspace*{\fill}
{\centering
  \setlength{\fboxsep}{1.2mm}%
  \setlength{\fboxrule}{0.4pt}%
  \fcolorbox{gray!55}{white}{%
    \includegraphics[width=0.95\textwidth,height=0.81\textheight,keepaspectratio]{\facsimiledir/20251126_110930.jpg}%
  }\par}
\vspace*{\fill}
\newpage

% =============================================================================
% VIS-A-VIS SPREAD : FOLIO 78 VERSO (20251126_110943.jpg)
% =============================================================================
\cleardoubleevenpage
\renewcommand{\currentmonthyear}{Février 1917}
\renewcommand{\currentfoliolabel}{Folio 78 verso}
\phantomsection\label{fol:157}
% --- En-tête du folio (Page paire / gauche : Transcription, apparat & registres de bas de page) ---
\noindent\rule{\textwidth}{0.5pt}\par\vspace{0.55mm}
\noindent\textbf{\footnotesize FOLIO 78 VERSO} \hfill \textbf{\footnotesize FÉVRIER 1917}\par\vspace{0.25mm}
\noindent{\fontsize{6.9pt}{8.0pt}\selectfont Déménagement dans l'ancienne chambre humide du capitaine Williams}\par\vspace{0.2mm}
\noindent{\fontsize{6.9pt}{8.0pt}\selectfont\textit{19 – 20 février 1917} \hfill \textit{Hérissart, Fréchencourt, Bavelincourt, Pierregot (Somme)}}\par\vspace{0.5mm}
\noindent\rule{\textwidth}{0.35pt}\par\vspace{0.8mm}
% --- Corps de la transcription diplomatique et apparat marginal en prolongement ---
{\renewcommand{\arraystretch}{0.83}%
\noindent\hspace*{3.5mm}\begin{tabular}{@{}p{87mm}@{\hspace{3.5mm}}p{39mm}@{}}
      \leavevmode\fontsize{6.6pt}{7.5pt}\selectfont\strut \underline{19 Lundi} — (En caisse : 225 f. 95) \hfill{} \raisebox{0.5ex}{\underline{Février 1917}}\strut & \smash{\parbox[t]{39mm}{\fontsize{6.2pt}{7.0pt}\selectfont\raggedright\color{black!88}\textbf{1 f.} \textit{abbr. pro} franc(s).}} \\
      \leavevmode\fontsize{6.6pt}{7.5pt}\selectfont\strut Tout le jour, temps couvert.\strut &  \\
      \leavevmode\fontsize{6.6pt}{7.5pt}\selectfont\strut Petites ondées intermittentes.\strut &  \\
      \leavevmode\fontsize{6.6pt}{7.5pt}\selectfont\strut Matinée, malgré pluie \& routes boueuses,\strut &  \\
      \leavevmode\llap{\fontsize{6.4pt}{7.3pt}\selectfont\color{darkgray}5\hspace{1.8mm}}\fontsize{6.6pt}{7.5pt}\selectfont\strut exploré villages de Fréchencourt \& Bavelincourt,\strut &  \\
      \leavevmode\fontsize{6.6pt}{7.5pt}\selectfont\strut pour trouver paille et foin.\strut &  \\
      \leavevmode\fontsize{6.6pt}{7.5pt}\selectfont\strut Après-midi, repos. Mais au thé, appris par\strut &  \\
      \leavevmode\fontsize{6.6pt}{7.5pt}\selectfont\strut Cap. Leonard nouveaux arrangements pour billets.\strut & \smash{\parbox[t]{39mm}{\fontsize{6.2pt}{7.0pt}\selectfont\raggedright\color{black!88}\textbf{8 Cap.} \textit{abbr. pro} Capitaine. \textbf{8 billets} \textit{angl. in ms.}}} \\
      \leavevmode\fontsize{6.6pt}{7.5pt}\selectfont\strut Le soir même, vers 7h, déménagé du 7b\strut &  \\
      \leavevmode\llap{\fontsize{6.4pt}{7.3pt}\selectfont\color{darkgray}10\hspace{1.8mm}}\fontsize{6.6pt}{7.5pt}\selectfont\strut à la chambre qu'occupait Cap. Williams, à\strut & \smash{\parbox[t]{39mm}{\fontsize{6.2pt}{7.0pt}\selectfont\raggedright\color{black!88}\textbf{10 Cap.} \textit{abbr. pro} Capitaine.}} \\
      \leavevmode\fontsize{6.6pt}{7.5pt}\selectfont\strut côté office. Pièce très froide et surtout humide,\strut &  \\
      \leavevmode\fontsize{6.6pt}{7.5pt}\selectfont\strut murs \& briques du floor ruissellent d'eau.\strut &  \\
      \leavevmode\fontsize{6.6pt}{7.5pt}\selectfont\strut Vers 9h 1/2, prop.re du n° 104, malade : me\strut & \smash{\parbox[t]{39mm}{\fontsize{6.2pt}{7.0pt}\selectfont\raggedright\color{black!88}\textbf{13 1/2} \textit{scil.} demi / moitié. — \textbf{14 M.O.} \textit{angl.} Medical Officer. — \textbf{15} \textit{in marg. cum unco} (comptes). \textbf{15 c.} \textit{abbr. pro} centimes.}} \\
      \leavevmode\fontsize{6.6pt}{7.5pt}\selectfont\strut suis dévoué pour elle. Allé chercher curé et M.O.\strut &  \\
      \cline{1-1}
      \leavevmode\llap{\fontsize{6.4pt}{7.3pt}\selectfont\color{darkgray}15\hspace{1.8mm}}\fontsize{6.6pt}{7.5pt}\selectfont\strut \underline{20 Mardi.} — [ 1 hair cut : net : 30 c. ]\strut &  \\
      \leavevmode\fontsize{6.6pt}{7.5pt}\selectfont\strut Assez bien dormi nouvelle chambre. Mais,\strut &  \\
      \leavevmode\fontsize{6.6pt}{7.5pt}\selectfont\strut constaté avec regret combien elle est humide et\strut &  \\
      \leavevmode\fontsize{6.6pt}{7.5pt}\selectfont\strut froide. Impossible de m'y tenir, durant la\strut &  \\
      \leavevmode\fontsize{6.6pt}{7.5pt}\selectfont\strut journée. Obligé rester du matin au soir\strut &  \\
      \leavevmode\llap{\fontsize{6.4pt}{7.3pt}\selectfont\color{darkgray}20\hspace{1.8mm}}\fontsize{6.6pt}{7.5pt}\selectfont\strut dans notre office.\strut &  \\
      \leavevmode\fontsize{6.6pt}{7.5pt}\selectfont\strut Matin, à Pierregot, de ferme en ferme\strut &  \\
      \leavevmode\fontsize{6.6pt}{7.5pt}\selectfont\strut pour trouver paille \& foin. Route abst\strut & \smash{\parbox[t]{39mm}{\fontsize{6.2pt}{7.0pt}\selectfont\raggedright\color{black!88}\textbf{22 abst} \textit{abbr. pro} absolument.}} \\
      \leavevmode\fontsize{6.6pt}{7.5pt}\selectfont\strut boueuse \& impraticable.\strut &  \\
      \leavevmode\fontsize{6.6pt}{7.5pt}\selectfont\strut Après-midi, liberté.\strut &  \\
      \cline{1-1}
      \leavevmode\llap{\fontsize{6.4pt}{7.3pt}\selectfont\color{darkgray}25\hspace{1.8mm}}\fontsize{6.6pt}{7.5pt}\selectfont\strut [ A Albertine (Hérissart) : 1 f. 70 blanchissage,\strut & \smash{\parbox[t]{39mm}{\fontsize{6.2pt}{7.0pt}\selectfont\raggedright\color{black!88}\textbf{25} \textit{in marg. cum unco} (comptes). \textbf{25 f.} \textit{abbr. pro} franc(s). — \textbf{26 f.} \textit{abbr. pro} franc(s).}} \\
      \leavevmode\fontsize{6.6pt}{7.5pt}\selectfont\strut et pour service chambre : en tout : 2 f. ]\strut &  \\
\end{tabular}}%
\par\vspace*{\fill}
% --- Registres de bas de page (1. Critique interne [si pertinent] ; 2. Loci paralleli en latin ; 3. Commentaire historique en retrait) ---
\noindent\rule{\textwidth}{0.5pt}\par\vspace{0.65mm}
{\fontsize{6.6pt}{7.7pt}\selectfont%
\noindent\normalfont \textbf{Mention financière en en-tête :} Exceptionnellement, Hippolyte Combes inscrit son solde disponible (\textit{En caisse : 225 f. 95}) directement sur la première ligne entre la date soulignée et le mois.\par\vspace{0.5mm}
\noindent\rule{22mm}{0.3pt}\par\vspace{0.45mm}
\noindent\normalfont\color{black!92} \textbf{10} (\textit{Williams}) cf. supra f\textsuperscript{o}~52~v\textsuperscript{o} (p.~\pageref{fol:105}), f\textsuperscript{o}~73~v\textsuperscript{o} (p.~\pageref{fol:147}), infra f\textsuperscript{o}~80~r\textsuperscript{o} (p.~\pageref{fol:160}), f\textsuperscript{o}~81~v\textsuperscript{o} (p.~\pageref{fol:163}), \textit{et passim}. — \textbf{21} (\textit{Pierregot}) cf. supra f\textsuperscript{o}~52~r\textsuperscript{o} (p.~\pageref{fol:104}), infra f\textsuperscript{o}~84~v\textsuperscript{o} (p.~\pageref{fol:169}), f\textsuperscript{o}~85~v\textsuperscript{o} (p.~\pageref{fol:171}), \textit{et passim}. — \textbf{22} (\textit{paille \& foin}) cf. supra f\textsuperscript{o}~3~r\textsuperscript{o} (p.~\pageref{fol:6}), f\textsuperscript{o}~77~v\textsuperscript{o} (p.~\pageref{fol:155}), infra f\textsuperscript{o}~79~r\textsuperscript{o} (p.~\pageref{fol:158}), f\textsuperscript{o}~80~r\textsuperscript{o} (p.~\pageref{fol:160}), \textit{et passim}.\par\vspace{0.55mm}
\noindent\rule{38mm}{0.4pt}\par\vspace{0.5mm}
\noindent\hspace*{5mm}\parbox{\dimexpr\textwidth-5mm\relax}{\fontsize{6.6pt}{7.7pt}\selectfont\color{black!84}Après une tournée de recherche de fourrage dans la vallée de l'Hallue (\textbf{Fréchencourt} et \textbf{Bavelincourt}) le lundi 19 février 1917, le \textbf{capitaine Leonard} notifie une réorganisation des cantonnements qui oblige Hippolyte Combes à quitter sa chambre du n\textsuperscript{o}~7b (en réglant 2~fr. à la domestique \textbf{Albertine}, l.~25) pour la pièce humide et glaciale du \textbf{n\textsuperscript{o}~104} attenante au bureau.}\par}
\newpage
\phantomsection\label{fac:157}
% --- En-tête du fac-similé (Page impaire / droite : Photographie du manuscrit) ---
\noindent\rule{\textwidth}{0.5pt}\par\vspace{0.55mm}
\noindent\textbf{\fontsize{6.8pt}{7.8pt}\selectfont FAC-SIMILÉ EN VIS-À-VIS — FOLIO 78 VERSO} \hfill \textbf{\fontsize{6.8pt}{7.8pt}\selectfont FÉVRIER 1917}\par\vspace{0.25mm}
\noindent{\fontsize{6.8pt}{8.0pt}\selectfont\textit{Livre-Journal d'Hippolyte Combes} — 19 – 20 février 1917 (Hérissart, Fréchencourt, Bavelincourt, Pierregot (Somme))}\par\vspace{0.55mm}
\noindent\rule{\textwidth}{0.35pt}\par\vspace*{\fill}
{\centering
  \setlength{\fboxsep}{1.2mm}%
  \setlength{\fboxrule}{0.4pt}%
  \fcolorbox{gray!55}{white}{%
    \includegraphics[width=0.95\textwidth,height=0.81\textheight,keepaspectratio]{\facsimiledir/20251126_110943.jpg}%
  }\par}
\vspace*{\fill}
\newpage

% =============================================================================
% VIS-A-VIS SPREAD : FOLIO 79 RECTO (20251126_111008.jpg)
% =============================================================================
\cleardoubleevenpage
\renewcommand{\currentmonthyear}{Février 1917}
\renewcommand{\currentfoliolabel}{Folio 79 recto}
\phantomsection\label{fol:158}
% --- En-tête du folio (Page paire / gauche : Transcription, apparat & registres de bas de page) ---
\noindent\rule{\textwidth}{0.5pt}\par\vspace{0.55mm}
\noindent\textbf{\footnotesize FOLIO 79 RECTO} \hfill \textbf{\footnotesize FÉVRIER 1917}\par\vspace{0.25mm}
\noindent{\fontsize{6.9pt}{8.0pt}\selectfont Course à Doullens et Auxi-le-Château sous le brouillard : mandat pour Mme Spagnol}\par\vspace{0.2mm}
\noindent{\fontsize{6.9pt}{8.0pt}\selectfont\textit{21 – 23 février 1917} \hfill \textit{Hérissart, Doullens, Auxi-le-Château, Puchevillers, Molliens-au-Bois (Somme)}}\par\vspace{0.5mm}
\noindent\rule{\textwidth}{0.35pt}\par\vspace{0.8mm}
% --- Corps de la transcription diplomatique et apparat marginal en prolongement ---
{\renewcommand{\arraystretch}{0.83}%
\noindent\hspace*{3.5mm}\begin{tabular}{@{}p{87mm}@{\hspace{3.5mm}}p{39mm}@{}}
      \leavevmode\fontsize{6.6pt}{7.5pt}\selectfont\strut \underline{21 Merc.} — \hfill{} \underline{Février 1917}\strut & \smash{\parbox[t]{39mm}{\fontsize{6.2pt}{7.0pt}\selectfont\raggedright\color{black!88}\textbf{1 Merc.} \textit{scil.} Mercredi.}} \\
      \leavevmode\fontsize{6.6pt}{7.5pt}\selectfont\strut Temps nuageux. Brume pénétrante. Bruine\strut &  \\
      \leavevmode\fontsize{6.6pt}{7.5pt}\selectfont\strut intermittente. Dégel, transformé\strut &  \\
      \leavevmode\fontsize{6.6pt}{7.5pt}\selectfont\strut routes en fondrières véritables. Très profondes ornières.\strut &  \\
      \leavevmode\llap{\fontsize{6.4pt}{7.3pt}\selectfont\color{darkgray}5\hspace{1.8mm}}\fontsize{6.6pt}{7.5pt}\selectfont\strut Matin, course à Auxi-le-Château. Au-delà\strut &  \\
      \leavevmode\fontsize{6.6pt}{7.5pt}\selectfont\strut de Doullens, campagne as. pittoresque, mais\strut & \smash{\parbox[t]{39mm}{\fontsize{6.2pt}{7.0pt}\selectfont\raggedright\color{black!88}\textbf{6 as.} \textit{abbr. pro} assez.}} \\
      \leavevmode\fontsize{6.6pt}{7.5pt}\selectfont\strut brouillard et pluie rendent paysages banals.\strut &  \\
      \leavevmode\fontsize{6.6pt}{7.5pt}\selectfont\strut Arrêt à Doullens, 1h 1/2, pour lunch.\strut & \smash{\parbox[t]{39mm}{\fontsize{6.2pt}{7.0pt}\selectfont\raggedright\color{black!88}\textbf{8 1/2} \textit{scil.} demi / moitié. \textbf{8 lunch} \textit{angl. in ms.} — \textbf{9 abomles} \textit{abbr. pro} abominables. — \textbf{11} \textit{in marg. cum unco} (comptes). \textbf{11 Mme} \textit{abbr. pro} Madame. — \textbf{13 f.} \textit{abbr. pro} franc(s). — \textbf{14 f.} \textit{abbr. pro} franc(s).}} \\
      \leavevmode\fontsize{6.6pt}{7.5pt}\selectfont\strut Retour par Puchevillers, routes abomles.\strut &  \\
      \leavevmode\llap{\fontsize{6.4pt}{7.3pt}\selectfont\color{darkgray}10\hspace{1.8mm}}\fontsize{6.6pt}{7.5pt}\selectfont\strut Impressions de solitude \& détresse morale.\strut &  \\
      \leavevmode\fontsize{6.6pt}{7.5pt}\selectfont\strut [ \underline{A Doullens}, mandat de 42 f. pour Mme\strut &  \\
      \leavevmode\fontsize{6.6pt}{7.5pt}\selectfont\strut Spagnol, Nice — —\strut &  \\
      \leavevmode\fontsize{6.6pt}{7.5pt}\selectfont\strut \underline{A Doullens} : 2 œufs \& café : 1 f. —\strut &  \\
      \leavevmode\fontsize{6.6pt}{7.5pt}\selectfont\strut de semelles liège : 0 f. 80 ]\strut &  \\
      \leavevmode\llap{\fontsize{6.4pt}{7.3pt}\selectfont\color{darkgray}15\hspace{1.8mm}}\fontsize{6.6pt}{7.5pt}\selectfont\strut Après-midi, à Hérissart, village 1 peu\strut &  \\
      \leavevmode\fontsize{6.6pt}{7.5pt}\selectfont\strut plus désert, presque funèbre.\strut &  \\
      \cline{1-1}
      \leavevmode\fontsize{6.6pt}{7.5pt}\selectfont\strut \underline{22 Jeudi} — Journée d'1 monotonie extra.\strut &  \\
      \leavevmode\fontsize{6.6pt}{7.5pt}\selectfont\strut Temps couvert. Matin \& soir, épais\strut &  \\
      \leavevmode\fontsize{6.6pt}{7.5pt}\selectfont\strut brouillard.\strut &  \\
      \leavevmode\llap{\fontsize{6.4pt}{7.3pt}\selectfont\color{darkgray}20\hspace{1.8mm}}\fontsize{6.6pt}{7.5pt}\selectfont\strut Liberté entière tt le jour. Resté office, ou\strut & \smash{\parbox[t]{39mm}{\fontsize{6.2pt}{7.0pt}\selectfont\raggedright\color{black!88}\textbf{20 tt} \textit{abbr. pro} tout.}} \\
      \leavevmode\fontsize{6.6pt}{7.5pt}\selectfont\strut avec proprios, à lire, causer etc. —\strut &  \\
      \cline{1-1}
      \leavevmode\fontsize{6.6pt}{7.5pt}\selectfont\strut \underline{23 Vendredi} — Encore = temps, brouillard \& boue.\strut &  \\
      \leavevmode\fontsize{6.6pt}{7.5pt}\selectfont\strut Matinée, à Molliens-au-Bois, pour trouver\strut &  \\
      \leavevmode\fontsize{6.6pt}{7.5pt}\selectfont\strut fourrage. Petit village au fond d'1 vallon,\strut &  \\
      \leavevmode\llap{\fontsize{6.4pt}{7.3pt}\selectfont\color{darkgray}25\hspace{1.8mm}}\fontsize{6.6pt}{7.5pt}\selectfont\strut à côté de Montigny.\strut &  \\
      \leavevmode\fontsize{6.6pt}{7.5pt}\selectfont\strut Après-midi, liberté. Resté office et\strut &  \\
      \leavevmode\fontsize{6.6pt}{7.5pt}\selectfont\strut rangé paquets —\strut &  \\
\end{tabular}}%
\par\vspace*{\fill}
% --- Registres de bas de page (1. Critique interne [si pertinent] ; 2. Loci paralleli en latin ; 3. Commentaire historique en retrait) ---
\noindent\rule{\textwidth}{0.5pt}\par\vspace{0.65mm}
{\fontsize{6.6pt}{7.7pt}\selectfont%
\noindent\normalfont\color{black!92} \textbf{5} (\textit{Auxi-le-Château}) cf. supra f\textsuperscript{o}~62~r\textsuperscript{o} (p.~\pageref{fol:124}), f\textsuperscript{o}~76~r\textsuperscript{o} (p.~\pageref{fol:152}). — \textbf{6} (\textit{Doullens}) cf. supra f\textsuperscript{o}~25~v\textsuperscript{o} (p.~\pageref{fol:51}), f\textsuperscript{o}~76~r\textsuperscript{o} (p.~\pageref{fol:152}), infra f\textsuperscript{o}~81~v\textsuperscript{o} (p.~\pageref{fol:163}), f\textsuperscript{o}~89~v\textsuperscript{o} (p.~\pageref{fol:179}), \textit{et passim}. — \textbf{9} (\textit{Puchevillers}) cf. supra f\textsuperscript{o}~37~v\textsuperscript{o} (p.~\pageref{fol:75}), f\textsuperscript{o}~77~v\textsuperscript{o} (p.~\pageref{fol:155}), infra f\textsuperscript{o}~81~r\textsuperscript{o} (p.~\pageref{fol:162}), f\textsuperscript{o}~83~r\textsuperscript{o} (p.~\pageref{fol:166}), \textit{et passim}.\par\vspace{0.55mm}
\noindent\rule{38mm}{0.4pt}\par\vspace{0.5mm}
\noindent\hspace*{5mm}\parbox{\dimexpr\textwidth-5mm\relax}{\fontsize{6.6pt}{7.7pt}\selectfont\color{black!84}Lors de sa halte du 21 février 1917 à \textbf{Doullens} au retour d'\textbf{Auxi-le-Château}, Hippolyte Combes expédie un mandat-poste de 42~francs à \textbf{Mme Spagnol} à \textbf{Nice} en règlement du panier de fleurs reçu en gare de Toulouse le 29 janvier pour Marguerite Calmels.}\par}
\newpage
\phantomsection\label{fac:158}
% --- En-tête du fac-similé (Page impaire / droite : Photographie du manuscrit) ---
\noindent\rule{\textwidth}{0.5pt}\par\vspace{0.55mm}
\noindent\textbf{\fontsize{6.8pt}{7.8pt}\selectfont FAC-SIMILÉ EN VIS-À-VIS — FOLIO 79 RECTO} \hfill \textbf{\fontsize{6.8pt}{7.8pt}\selectfont FÉVRIER 1917}\par\vspace{0.25mm}
\noindent{\fontsize{6.8pt}{8.0pt}\selectfont\textit{Livre-Journal d'Hippolyte Combes} — 21 – 23 février 1917 (Hérissart, Doullens, Auxi-le-Château, Puchevillers, Molliens-au-Bois (Somme))}\par\vspace{0.55mm}
\noindent\rule{\textwidth}{0.35pt}\par\vspace*{\fill}
{\centering
  \setlength{\fboxsep}{1.2mm}%
  \setlength{\fboxrule}{0.4pt}%
  \fcolorbox{gray!55}{white}{%
    \includegraphics[width=0.95\textwidth,height=0.81\textheight,keepaspectratio]{\facsimiledir/20251126_111008.jpg}%
  }\par}
\vspace*{\fill}
\newpage

% =============================================================================
% VIS-A-VIS SPREAD : FOLIO 79 VERSO (20251126_111019.jpg)
% =============================================================================
\cleardoubleevenpage
\renewcommand{\currentmonthyear}{Février 1917}
\renewcommand{\currentfoliolabel}{Folio 79 verso}
\phantomsection\label{fol:159}
% --- En-tête du folio (Page paire / gauche : Transcription, apparat & registres de bas de page) ---
\noindent\rule{\textwidth}{0.5pt}\par\vspace{0.55mm}
\noindent\textbf{\footnotesize FOLIO 79 VERSO} \hfill \textbf{\footnotesize FÉVRIER 1917}\par\vspace{0.25mm}
\noindent{\fontsize{6.9pt}{8.0pt}\selectfont Promenades à Contay et Mirvaux : premières sensations de printemps}\par\vspace{0.2mm}
\noindent{\fontsize{6.9pt}{8.0pt}\selectfont\textit{24 – 26 février 1917} \hfill \textit{Hérissart, Contay, Mirvaux, Toutencourt (Somme)}}\par\vspace{0.5mm}
\noindent\rule{\textwidth}{0.35pt}\par\vspace{0.8mm}
% --- Corps de la transcription diplomatique et apparat marginal en prolongement ---
{\renewcommand{\arraystretch}{0.83}%
\noindent\hspace*{3.5mm}\begin{tabular}{@{}p{87mm}@{\hspace{3.5mm}}p{39mm}@{}}
      \leavevmode\fontsize{6.6pt}{7.5pt}\selectfont\strut \underline{24 Samedi} — \hfill{} \underline{Fév. 1917}\strut &  \\
      \leavevmode\fontsize{6.6pt}{7.5pt}\selectfont\strut Journée entière, inoccupé.\strut &  \\
      \leavevmode\fontsize{6.6pt}{7.5pt}\selectfont\strut Après-midi, courte promenade s/ route\strut & \smash{\parbox[t]{39mm}{\fontsize{6.2pt}{7.0pt}\selectfont\raggedright\color{black!88}\textbf{3 s/} \textit{abbr. pro} sur.}} \\
      \leavevmode\fontsize{6.6pt}{7.5pt}\selectfont\strut de Contay. Temps brumeux. Route très\strut &  \\
      \leavevmode\llap{\fontsize{6.4pt}{7.3pt}\selectfont\color{darkgray}5\hspace{1.8mm}}\fontsize{6.6pt}{7.5pt}\selectfont\strut boueuse. Arrivé jusqu'au moulin —\strut &  \\
      \cline{1-1}
      \leavevmode\fontsize{6.6pt}{7.5pt}\selectfont\strut \underline{25 Dim.}\strut & \smash{\parbox[t]{39mm}{\fontsize{6.2pt}{7.0pt}\selectfont\raggedright\color{black!88}\textbf{6 Dim.} \textit{scil.} Dimanche.}} \\
      \leavevmode\fontsize{6.6pt}{7.5pt}\selectfont\strut Matin, course à Mirvaux — hameau d'1\strut &  \\
      \leavevmode\fontsize{6.6pt}{7.5pt}\selectfont\strut 100ne d'habitants. Silence \& solitude funèbres.\strut &  \\
      \leavevmode\fontsize{6.6pt}{7.5pt}\selectfont\strut Enfin, durant qq heures, soleil brille, ciel\strut & \smash{\parbox[t]{39mm}{\fontsize{6.2pt}{7.0pt}\selectfont\raggedright\color{black!88}\textbf{9 qq} \textit{abbr. pro} quelque(s).}} \\
      \leavevmode\llap{\fontsize{6.4pt}{7.3pt}\selectfont\color{darkgray}10\hspace{1.8mm}}\fontsize{6.6pt}{7.5pt}\selectfont\strut bleu. 1ère sensation de printemps proche.\strut & \smash{\parbox[t]{39mm}{\fontsize{6.2pt}{7.0pt}\selectfont\raggedright\color{black!88}\textbf{10 1ère} \textit{scil.} première.}} \\
      \leavevmode\fontsize{6.6pt}{7.5pt}\selectfont\strut Après lunch, venu promener seul sur la\strut & \smash{\parbox[t]{39mm}{\fontsize{6.2pt}{7.0pt}\selectfont\raggedright\color{black!88}\textbf{11 lunch} \textit{angl. in ms.}}} \\
      \leavevmode\fontsize{6.6pt}{7.5pt}\selectfont\strut route de Toutencourt. Beau soleil, ciel bleu.\strut &  \\
      \leavevmode\fontsize{6.6pt}{7.5pt}\selectfont\strut Champs ensemencés. Masses sombres de petits\strut &  \\
      \leavevmode\fontsize{6.6pt}{7.5pt}\selectfont\strut bois à l'horizon. Air encore froid.\strut &  \\
      \leavevmode\llap{\fontsize{6.4pt}{7.3pt}\selectfont\color{darkgray}15\hspace{1.8mm}}\fontsize{6.6pt}{7.5pt}\selectfont\strut Village où ns sommes exilés, absolt\strut & \smash{\parbox[t]{39mm}{\fontsize{6.2pt}{7.0pt}\selectfont\raggedright\color{black!88}\textbf{15 ns} \textit{abbr. pro} nous. \textbf{15 absolt} \textit{abbr. pro} absolument. — \textbf{17 L.} \textit{scil.} Lundi. \textbf{17 Cap. W.} \textit{scil.} Capitaine Williams. — \textbf{19 s.} \textit{abbr. pro} soir. — \textbf{20 1/2} \textit{scil.} demi / moitié. — \textbf{21 bcp} \textit{abbr. pro} beaucoup. \textbf{21 English novel} \textit{angl. in ms.} — \textbf{22 ds} \textit{abbr. pro} dans. \textbf{22 tea} \textit{angl. in ms.}}} \\
      \leavevmode\fontsize{6.6pt}{7.5pt}\selectfont\strut perdu sur hauteurs — —\strut &  \\
      \cline{1-1}
      \leavevmode\fontsize{6.6pt}{7.5pt}\selectfont\strut \underline{26 L.} — Jour entier de liberté. Cap. W., seul\strut &  \\
      \leavevmode\fontsize{6.6pt}{7.5pt}\selectfont\strut parti vers 9h, allé à Senlis \& Martinsart,\strut &  \\
      \leavevmode\fontsize{6.6pt}{7.5pt}\selectfont\strut ne rentré qu'à 6h s.\strut &  \\
      \leavevmode\llap{\fontsize{6.4pt}{7.3pt}\selectfont\color{darkgray}20\hspace{1.8mm}}\fontsize{6.6pt}{7.5pt}\selectfont\strut Temps beau, soleil très agréable 1/2 journée.\strut &  \\
      \leavevmode\fontsize{6.6pt}{7.5pt}\selectfont\strut Écrit bcp et lu English novel.\strut &  \\
      \leavevmode\fontsize{6.6pt}{7.5pt}\selectfont\strut Après tea, promenade autour \& ds les\strut &  \\
      \leavevmode\fontsize{6.6pt}{7.5pt}\selectfont\strut environs du village.\strut &  \\
      \leavevmode\fontsize{6.6pt}{7.5pt}\selectfont\strut De nouveau, après dîner, nuit étoilée, et\strut &  \\
      \leavevmode\llap{\fontsize{6.4pt}{7.3pt}\selectfont\color{darkgray}25\hspace{1.8mm}}\fontsize{6.6pt}{7.5pt}\selectfont\strut pleine lune.\strut &  \\
\end{tabular}}%
\par\vspace*{\fill}
% --- Registres de bas de page (1. Critique interne [si pertinent] ; 2. Loci paralleli en latin ; 3. Commentaire historique en retrait) ---
\noindent\rule{\textwidth}{0.5pt}\par\vspace{0.65mm}
{\fontsize{6.6pt}{7.7pt}\selectfont%
\noindent\normalfont\color{black!92} \textbf{4} (\textit{Contay}) cf. supra f\textsuperscript{o}~49~r\textsuperscript{o} (p.~\pageref{fol:98}), f\textsuperscript{o}~50~v\textsuperscript{o} (p.~\pageref{fol:101}), infra f\textsuperscript{o}~80~v\textsuperscript{o} (p.~\pageref{fol:161}), f\textsuperscript{o}~83~v\textsuperscript{o} (p.~\pageref{fol:167}), \textit{et passim}. — \textbf{7} (\textit{Mirvaux}) cf. infra f\textsuperscript{o}~84~v\textsuperscript{o} (p.~\pageref{fol:169}), f\textsuperscript{o}~85~v\textsuperscript{o} (p.~\pageref{fol:171}). — \textbf{12} (\textit{Toutencourt}) cf. supra f\textsuperscript{o}~49~v\textsuperscript{o} (p.~\pageref{fol:99}), f\textsuperscript{o}~78~r\textsuperscript{o} (p.~\pageref{fol:156}), infra f\textsuperscript{o}~80~r\textsuperscript{o} (p.~\pageref{fol:160}), f\textsuperscript{o}~80~v\textsuperscript{o} (p.~\pageref{fol:161}), \textit{et passim}.\par\vspace{0.55mm}
\noindent\rule{38mm}{0.4pt}\par\vspace{0.5mm}
\noindent\hspace*{5mm}\parbox{\dimexpr\textwidth-5mm\relax}{\fontsize{6.6pt}{7.7pt}\selectfont\color{black!84}Fin février 1917 autour d'\textbf{Hérissart}: promenades vers le moulin de la route de \textbf{Contay}, le petit village voisin de \textbf{Mirvaux} (Somme, à 2,5~km au sud-est d'Hérissart) et la route de \textbf{Toutencourt}.}\par}
\newpage
\phantomsection\label{fac:159}
% --- En-tête du fac-similé (Page impaire / droite : Photographie du manuscrit) ---
\noindent\rule{\textwidth}{0.5pt}\par\vspace{0.55mm}
\noindent\textbf{\fontsize{6.8pt}{7.8pt}\selectfont FAC-SIMILÉ EN VIS-À-VIS — FOLIO 79 VERSO} \hfill \textbf{\fontsize{6.8pt}{7.8pt}\selectfont FÉVRIER 1917}\par\vspace{0.25mm}
\noindent{\fontsize{6.8pt}{8.0pt}\selectfont\textit{Livre-Journal d'Hippolyte Combes} — 24 – 26 février 1917 (Hérissart, Contay, Mirvaux, Toutencourt (Somme))}\par\vspace{0.55mm}
\noindent\rule{\textwidth}{0.35pt}\par\vspace*{\fill}
{\centering
  \setlength{\fboxsep}{1.2mm}%
  \setlength{\fboxrule}{0.4pt}%
  \fcolorbox{gray!55}{white}{%
    \includegraphics[width=0.95\textwidth,height=0.81\textheight,keepaspectratio]{\facsimiledir/20251126_111019.jpg}%
  }\par}
\vspace*{\fill}
\newpage

% =============================================================================
% VIS-A-VIS SPREAD : FOLIO 80 RECTO (20251126_111035.jpg)
% =============================================================================
\cleardoubleevenpage
\renewcommand{\currentmonthyear}{Février 1917}
\renewcommand{\currentfoliolabel}{Folio 80 recto}
\phantomsection\label{fol:160}
% --- En-tête du folio (Page paire / gauche : Transcription, apparat & registres de bas de page) ---
\noindent\rule{\textwidth}{0.5pt}\par\vspace{0.55mm}
\noindent\textbf{\footnotesize FOLIO 80 RECTO} \hfill \textbf{\footnotesize FÉVRIER 1917}\par\vspace{0.25mm}
\noindent{\fontsize{6.9pt}{8.0pt}\selectfont Rage de dents du capitaine Williams et passage de la cavalerie indienne à Hérissart}\par\vspace{0.2mm}
\noindent{\fontsize{6.9pt}{8.0pt}\selectfont\textit{27 – 28 février 1917} \hfill \textit{Hérissart, Toutencourt (Somme)}}\par\vspace{0.5mm}
\noindent\rule{\textwidth}{0.35pt}\par\vspace{0.8mm}
% --- Corps de la transcription diplomatique et apparat marginal en prolongement ---
{\renewcommand{\arraystretch}{0.87}%
\noindent\hspace*{3.5mm}\begin{tabular}{@{}p{87mm}@{\hspace{3.5mm}}p{39mm}@{}}
      \leavevmode\fontsize{6.9pt}{7.9pt}\selectfont\strut \underline{27 Ma.} — \hfill{} \underline{— Février 1917}\strut & \smash{\parbox[t]{39mm}{\fontsize{6.2pt}{7.0pt}\selectfont\raggedright\color{black!88}\textbf{1 Ma.} \textit{scil.} Mardi.}} \\
      \leavevmode\fontsize{6.9pt}{7.9pt}\selectfont\strut Même temps, très couvert. Tout le jour, ciel nuageux.\strut &  \\
      \leavevmode\fontsize{6.9pt}{7.9pt}\selectfont\strut Air assez doux.\strut &  \\
      \leavevmode\fontsize{6.9pt}{7.9pt}\selectfont\strut Capt. W., tooth ache, au lit j'à 1h après-midi.\strut & \smash{\parbox[t]{39mm}{\fontsize{6.2pt}{7.0pt}\selectfont\raggedright\color{black!88}\textbf{4 j'à} \textit{abbr. pro} jusqu'à. \textbf{4 Capt. W.} \textit{scil.} Captain Williams.}} \\
      \leavevmode\llap{\fontsize{6.4pt}{7.3pt}\selectfont\color{darkgray}5\hspace{1.8mm}}\fontsize{6.9pt}{7.9pt}\selectfont\strut D'où journée entière de liberté.\strut &  \\
      \leavevmode\fontsize{6.9pt}{7.9pt}\selectfont\strut Après-midi, courte promenade petite route\strut &  \\
      \leavevmode\fontsize{6.9pt}{7.9pt}\selectfont\strut de Toutencourt. Champs couverts de réseaux de\strut &  \\
      \leavevmode\fontsize{6.9pt}{7.9pt}\selectfont\strut barbed wire \& creusés de tranchées —\strut & \smash{\parbox[t]{39mm}{\fontsize{6.2pt}{7.0pt}\selectfont\raggedright\color{black!88}\textbf{8 barbed wire} \textit{angl. in ms.}}} \\
      \cline{1-1}
      \leavevmode\fontsize{6.9pt}{7.9pt}\selectfont\strut \underline{28 Merc.} — Temps 1 peu froid, mais très\strut & \smash{\parbox[t]{39mm}{\fontsize{6.2pt}{7.0pt}\selectfont\raggedright\color{black!88}\textbf{9 Merc.} \textit{scil.} Mercredi.}} \\
      \leavevmode\llap{\fontsize{6.4pt}{7.3pt}\selectfont\color{darkgray}10\hspace{1.8mm}}\fontsize{6.9pt}{7.9pt}\selectfont\strut beau. C'est le commt du printemps.\strut & \smash{\parbox[t]{39mm}{\fontsize{6.2pt}{7.0pt}\selectfont\raggedright\color{black!88}\textbf{10 commt} \textit{abbr. pro} commencement.}} \\
      \leavevmode\fontsize{6.9pt}{7.9pt}\selectfont\strut Matinée, liberté.\strut &  \\
      \leavevmode\fontsize{6.9pt}{7.9pt}\selectfont\strut Après lunch, à Toutencourt, pour trouver\strut & \smash{\parbox[t]{39mm}{\fontsize{6.2pt}{7.0pt}\selectfont\raggedright\color{black!88}\textbf{12 lunch} \textit{angl. in ms.}}} \\
      \leavevmode\fontsize{6.9pt}{7.9pt}\selectfont\strut fourrage \& paille —\strut &  \\
      \leavevmode\fontsize{6.9pt}{7.9pt}\selectfont\strut Après tea, promenade environs.\strut & \smash{\parbox[t]{39mm}{\fontsize{6.2pt}{7.0pt}\selectfont\raggedright\color{black!88}\textbf{14 tea} \textit{angl. in ms.}}} \\
      \leavevmode\llap{\fontsize{6.4pt}{7.3pt}\selectfont\color{darkgray}15\hspace{1.8mm}}\fontsize{6.9pt}{7.9pt}\selectfont\strut Matinée, plus. régiments de cavalerie\strut &  \\
      \leavevmode\fontsize{6.9pt}{7.9pt}\selectfont\strut indienne, traversé village d'Hérissart. Ces\strut &  \\
      \leavevmode\fontsize{6.9pt}{7.9pt}\selectfont\strut hindous, très curieux à regarder. Beaux\strut &  \\
      \leavevmode\fontsize{6.9pt}{7.9pt}\selectfont\strut chevaux, lances, à l'étrier, turbans etc.\strut &  \\
      \leavevmode\fontsize{6.9pt}{7.9pt}\selectfont\strut — Mess, pour février (depuis le 8) : 33 f. —\strut & \smash{\parbox[t]{39mm}{\fontsize{6.2pt}{7.0pt}\selectfont\raggedright\color{black!88}\textbf{19 f.} \textit{abbr. pro} franc(s).}} \\
\end{tabular}}%
\par\vspace*{\fill}
% --- Registres de bas de page (1. Critique interne [si pertinent] ; 2. Loci paralleli en latin ; 3. Commentaire historique en retrait) ---
\noindent\rule{\textwidth}{0.5pt}\par\vspace{0.65mm}
{\fontsize{6.6pt}{7.7pt}\selectfont%
\noindent\normalfont\color{black!92} \textbf{7} (\textit{Toutencourt}) cf. supra f\textsuperscript{o}~49~v\textsuperscript{o} (p.~\pageref{fol:99}), f\textsuperscript{o}~79~v\textsuperscript{o} (p.~\pageref{fol:159}), infra f\textsuperscript{o}~80~v\textsuperscript{o} (p.~\pageref{fol:161}), f\textsuperscript{o}~81~r\textsuperscript{o} (p.~\pageref{fol:162}), \textit{et passim}. — \textbf{13} (\textit{fourrage}) cf. supra f\textsuperscript{o}~3~r\textsuperscript{o} (p.~\pageref{fol:6}), f\textsuperscript{o}~79~r\textsuperscript{o} (p.~\pageref{fol:158}), infra f\textsuperscript{o}~84~v\textsuperscript{o} (p.~\pageref{fol:169}), f\textsuperscript{o}~88~r\textsuperscript{o} (p.~\pageref{fol:176}), \textit{et passim}. — \textbf{16} (\textit{Hérissart}) cf. supra f\textsuperscript{o}~47~v\textsuperscript{o} (p.~\pageref{fol:95}), f\textsuperscript{o}~79~v\textsuperscript{o} (p.~\pageref{fol:159}), infra f\textsuperscript{o}~80~v\textsuperscript{o} (p.~\pageref{fol:161}), f\textsuperscript{o}~81~r\textsuperscript{o} (p.~\pageref{fol:162}), \textit{et passim}.\par\vspace{0.55mm}
\noindent\rule{38mm}{0.4pt}\par\vspace{0.5mm}
\noindent\hspace*{5mm}\parbox{\dimexpr\textwidth-5mm\relax}{\fontsize{6.6pt}{7.7pt}\selectfont\color{black!84}Les régiments de \textbf{cavalerie indienne} traversant \textbf{Hérissart} le 28 février 1917 appartiennent aux \textbf{4\textsuperscript{e} et 5\textsuperscript{e} Divisions de cavalerie} (anciennes \textit{1st} et \textit{2nd Indian Cavalry Divisions} renommées en novembre 1916 au sein du \textit{Cavalry Corps} de la B.E.F. ; Maj. A.~F. Becke, \textit{Order of Battle}, Part~2A, p.~35–48).}\par}
\newpage
\phantomsection\label{fac:160}
% --- En-tête du fac-similé (Page impaire / droite : Photographie du manuscrit) ---
\noindent\rule{\textwidth}{0.5pt}\par\vspace{0.55mm}
\noindent\textbf{\fontsize{6.8pt}{7.8pt}\selectfont FAC-SIMILÉ EN VIS-À-VIS — FOLIO 80 RECTO} \hfill \textbf{\fontsize{6.8pt}{7.8pt}\selectfont FÉVRIER 1917}\par\vspace{0.25mm}
\noindent{\fontsize{6.8pt}{8.0pt}\selectfont\textit{Livre-Journal d'Hippolyte Combes} — 27 – 28 février 1917 (Hérissart, Toutencourt (Somme))}\par\vspace{0.55mm}
\noindent\rule{\textwidth}{0.35pt}\par\vspace*{\fill}
{\centering
  \setlength{\fboxsep}{1.2mm}%
  \setlength{\fboxrule}{0.4pt}%
  \fcolorbox{gray!55}{white}{%
    \includegraphics[width=0.95\textwidth,height=0.81\textheight,keepaspectratio]{\facsimiledir/20251126_111035.jpg}%
  }\par}
\vspace*{\fill}
\newpage

% =============================================================================
% VIS-A-VIS SPREAD : FOLIO 80 VERSO (20251126_111048.jpg)
% =============================================================================
\cleardoubleevenpage
\renewcommand{\currentmonthyear}{Mars 1917}
\renewcommand{\currentfoliolabel}{Folio 80 verso}
\phantomsection\label{fol:161}
% --- En-tête du folio (Page paire / gauche : Transcription, apparat & registres de bas de page) ---
\noindent\rule{\textwidth}{0.5pt}\par\vspace{0.55mm}
\noindent\textbf{\footnotesize FOLIO 80 VERSO} \hfill \textbf{\footnotesize MARS 1917}\par\vspace{0.25mm}
\noindent{\fontsize{6.9pt}{8.0pt}\selectfont Entrée en mars : château de Rubempré, halte au clair de lune et entretien avec le Dr Dubois}\par\vspace{0.2mm}
\noindent{\fontsize{6.9pt}{8.0pt}\selectfont\textit{1er – 2 mars 1917} \hfill \textit{Hérissart, Rubempré, Montigny, Fréchencourt, Contay, Toutencourt (Somme)}}\par\vspace{0.5mm}
\noindent\rule{\textwidth}{0.35pt}\par\vspace{0.8mm}
% --- Corps de la transcription diplomatique et apparat marginal en prolongement ---
{\renewcommand{\arraystretch}{0.87}%
\noindent\hspace*{3.5mm}\begin{tabular}{@{}p{87mm}@{\hspace{3.5mm}}p{39mm}@{}}
      \leavevmode\fontsize{6.9pt}{7.9pt}\selectfont\strut \underline{1er Jeudi} — \hfill{} \underline{— Mars 1917 —}\strut & \smash{\parbox[t]{39mm}{\fontsize{6.2pt}{7.0pt}\selectfont\raggedright\color{black!88}\textbf{1 1er} \textit{scil.} premier.}} \\
      \leavevmode\fontsize{6.9pt}{7.9pt}\selectfont\strut Très belle journée.\strut &  \\
      \leavevmode\fontsize{6.9pt}{7.9pt}\selectfont\strut Après-midi, courses en auto : Rubempré, vu château,\strut &  \\
      \leavevmode\fontsize{6.9pt}{7.9pt}\selectfont\strut puis Montigny \& Fréchencourt. Causé lgts avec le fils\strut & \smash{\parbox[t]{39mm}{\fontsize{6.2pt}{7.0pt}\selectfont\raggedright\color{black!88}\textbf{4 lgts} \textit{abbr. pro} longtemps. — \textbf{5 ds} \textit{abbr. pro} dans.}} \\
      \leavevmode\llap{\fontsize{6.4pt}{7.3pt}\selectfont\color{darkgray}5\hspace{1.8mm}}\fontsize{6.9pt}{7.9pt}\selectfont\strut du maire, ds ses champs ensoleillés, vers Pont-Noyelles.\strut &  \\
      \leavevmode\fontsize{6.9pt}{7.9pt}\selectfont\strut Après le thé, promenade sur route de Contay.\strut &  \\
      \leavevmode\fontsize{6.9pt}{7.9pt}\selectfont\strut Brume estompait horizons. Plateau, d'où belle\strut &  \\
      \leavevmode\fontsize{6.9pt}{7.9pt}\selectfont\strut vue panoramique. Moulins à vent. Aux\strut &  \\
      \leavevmode\fontsize{6.9pt}{7.9pt}\selectfont\strut carrefours, grands crucifix, cuivre et fer,\strut &  \\
      \leavevmode\llap{\fontsize{6.4pt}{7.3pt}\selectfont\color{darkgray}10\hspace{1.8mm}}\fontsize{6.9pt}{7.9pt}\selectfont\strut élancés.\strut &  \\
      \leavevmode\fontsize{6.9pt}{7.9pt}\selectfont\strut Soir, dîner pris, ressorti encore. Cheminé le\strut &  \\
      \leavevmode\fontsize{6.9pt}{7.9pt}\selectfont\strut long route de Toutencourt. Mervx clair de\strut &  \\
      \leavevmode\fontsize{6.9pt}{7.9pt}\selectfont\strut lune illuminant paysage. Ombre des arbres\strut &  \\
      \leavevmode\fontsize{6.9pt}{7.9pt}\selectfont\strut projetée sur le sol. Halo lunaire. Impression\strut &  \\
      \leavevmode\llap{\fontsize{6.4pt}{7.3pt}\selectfont\color{darkgray}15\hspace{1.8mm}}\fontsize{6.9pt}{7.9pt}\selectfont\strut de calme reposant \& silence absolu.\strut &  \\
      \cline{1-1}
      \leavevmode\fontsize{6.9pt}{7.9pt}\selectfont\strut \underline{2 Vendredi.} — Matinée, brouillard. Vers milieu\strut &  \\
      \leavevmode\fontsize{6.9pt}{7.9pt}\selectfont\strut jour, soleil, mais pâle. Temps redevenu\strut &  \\
      \leavevmode\fontsize{6.9pt}{7.9pt}\selectfont\strut froid.\strut &  \\
      \leavevmode\fontsize{6.9pt}{7.9pt}\selectfont\strut Matin, venus à Montigny. Très longue\strut &  \\
      \leavevmode\llap{\fontsize{6.4pt}{7.3pt}\selectfont\color{darkgray}20\hspace{1.8mm}}\fontsize{6.9pt}{7.9pt}\selectfont\strut \& intéressante conversation avec Dr Dubois —\strut &  \\
      \leavevmode\fontsize{6.9pt}{7.9pt}\selectfont\strut Après-midi, liberté entière.\strut &  \\
\end{tabular}}%
\par\vspace*{\fill}
% --- Registres de bas de page (1. Critique interne [si pertinent] ; 2. Loci paralleli en latin ; 3. Commentaire historique en retrait) ---
\noindent\rule{\textwidth}{0.5pt}\par\vspace{0.65mm}
{\fontsize{6.6pt}{7.7pt}\selectfont%
\noindent\normalfont\color{black!92} \textbf{3} (\textit{Rubempré}) cf. supra f\textsuperscript{o}~47~v\textsuperscript{o} (p.~\pageref{fol:95}), f\textsuperscript{o}~74~v\textsuperscript{o} (p.~\pageref{fol:149}), infra f\textsuperscript{o}~81~r\textsuperscript{o} (p.~\pageref{fol:162}), f\textsuperscript{o}~86~r\textsuperscript{o} (p.~\pageref{fol:172}), \textit{et passim}. — \textbf{4} (\textit{Montigny}) cf. supra f\textsuperscript{o}~59~v\textsuperscript{o} (p.~\pageref{fol:119}), f\textsuperscript{o}~79~r\textsuperscript{o} (p.~\pageref{fol:158}), infra f\textsuperscript{o}~81~r\textsuperscript{o} (p.~\pageref{fol:162}), f\textsuperscript{o}~91~v\textsuperscript{o} (p.~\pageref{fol:183}), \textit{et passim}. — \textbf{6} (\textit{Contay}) cf. supra f\textsuperscript{o}~49~r\textsuperscript{o} (p.~\pageref{fol:98}), f\textsuperscript{o}~79~v\textsuperscript{o} (p.~\pageref{fol:159}), infra f\textsuperscript{o}~83~v\textsuperscript{o} (p.~\pageref{fol:167}), f\textsuperscript{o}~84~v\textsuperscript{o} (p.~\pageref{fol:169}), \textit{et passim}.\par\vspace{0.55mm}
\noindent\rule{38mm}{0.4pt}\par\vspace{0.5mm}
\noindent\hspace*{5mm}\parbox{\dimexpr\textwidth-5mm\relax}{\fontsize{6.6pt}{7.7pt}\selectfont\color{black!84}Les jeudi 1\textsuperscript{er} et vendredi 2 mars 1917, les tournées d'achats mènent Hippolyte Combes au château de \textbf{Rubempré} puis dans la vallée de l'Hallue à \textbf{Fréchencourt} et \textbf{Montigny-sur-l'Hallue} (direction de \textbf{Pont-Noyelles}), où il s'entretient longuement avec le fils du maire Paul Dubois (déjà visité le 3 janvier, f\textsuperscript{o}~59~v\textsuperscript{o}) puis le lendemain avec le \textbf{docteur Dubois}.}\par}
\newpage
\phantomsection\label{fac:161}
% --- En-tête du fac-similé (Page impaire / droite : Photographie du manuscrit) ---
\noindent\rule{\textwidth}{0.5pt}\par\vspace{0.55mm}
\noindent\textbf{\fontsize{6.8pt}{7.8pt}\selectfont FAC-SIMILÉ EN VIS-À-VIS — FOLIO 80 VERSO} \hfill \textbf{\fontsize{6.8pt}{7.8pt}\selectfont MARS 1917}\par\vspace{0.25mm}
\noindent{\fontsize{6.8pt}{8.0pt}\selectfont\textit{Livre-Journal d'Hippolyte Combes} — 1er – 2 mars 1917 (Hérissart, Rubempré, Montigny, Fréchencourt, Contay, Toutencourt (Somme))}\par\vspace{0.55mm}
\noindent\rule{\textwidth}{0.35pt}\par\vspace*{\fill}
{\centering
  \setlength{\fboxsep}{1.2mm}%
  \setlength{\fboxrule}{0.4pt}%
  \fcolorbox{gray!55}{white}{%
    \includegraphics[width=0.95\textwidth,height=0.81\textheight,keepaspectratio]{\facsimiledir/20251126_111048.jpg}%
  }\par}
\vspace*{\fill}
\newpage

% =============================================================================
% VIS-A-VIS SPREAD : FOLIO 81 RECTO (20251126_111105.jpg)
% =============================================================================
\cleardoubleevenpage
\renewcommand{\currentmonthyear}{Mars 1917}
\renewcommand{\currentfoliolabel}{Folio 81 recto}
\phantomsection\label{fol:162}
% --- En-tête du folio (Page paire / gauche : Transcription, apparat & registres de bas de page) ---
\noindent\rule{\textwidth}{0.5pt}\par\vspace{0.55mm}
\noindent\textbf{\footnotesize FOLIO 81 RECTO} \hfill \textbf{\footnotesize MARS 1917}\par\vspace{0.25mm}
\noindent{\fontsize{6.9pt}{8.0pt}\selectfont Retenue de solde pour permission à Toutencourt et chute de neige sur les arbres}\par\vspace{0.2mm}
\noindent{\fontsize{6.9pt}{8.0pt}\selectfont\textit{3 – 5 mars 1917} \hfill \textit{Hérissart, Rubempré, Puchevillers, Toutencourt, Montigny, Fréchencourt (Somme)}}\par\vspace{0.5mm}
\noindent\rule{\textwidth}{0.35pt}\par\vspace{0.8mm}
% --- Corps de la transcription diplomatique et apparat marginal en prolongement ---
{\renewcommand{\arraystretch}{0.83}%
\noindent\hspace*{3.5mm}\begin{tabular}{@{}p{87mm}@{\hspace{3.5mm}}p{39mm}@{}}
      \leavevmode\fontsize{6.6pt}{7.5pt}\selectfont\strut \underline{3 Sam.} — \hfill{} \underline{— Mars 1917 —}\strut & \smash{\parbox[t]{39mm}{\fontsize{6.2pt}{7.0pt}\selectfont\raggedright\color{black!88}\textbf{1 Sam.} \textit{scil.} Samedi.}} \\
      \leavevmode\fontsize{6.6pt}{7.5pt}\selectfont\strut Matin \& soir, brouillard glacial. Ciel resté\strut &  \\
      \leavevmode\fontsize{6.6pt}{7.5pt}\selectfont\strut voilé de brume tout le jour.\strut &  \\
      \leavevmode\fontsize{6.6pt}{7.5pt}\selectfont\strut Matin, de 9h 30 à 1h, à Rubempré et à\strut &  \\
      \leavevmode\llap{\fontsize{6.4pt}{7.3pt}\selectfont\color{darkgray}5\hspace{1.8mm}}\fontsize{6.6pt}{7.5pt}\selectfont\strut Puchevillers (paille etc.), \& Town Major's office.\strut & \smash{\parbox[t]{39mm}{\fontsize{6.2pt}{7.0pt}\selectfont\raggedright\color{black!88}\textbf{5 Town Major} \textit{angl. in ms.}}} \\
      \leavevmode\fontsize{6.6pt}{7.5pt}\selectfont\strut Après-midi, liberté entière. Venu à pied\strut &  \\
      \leavevmode\fontsize{6.6pt}{7.5pt}\selectfont\strut à Toutencourt, Mission fr. \& qq fermes.\strut & \smash{\parbox[t]{39mm}{\fontsize{6.2pt}{7.0pt}\selectfont\raggedright\color{black!88}\textbf{7 qq} \textit{abbr. pro} quelque(s).}} \\
      \leavevmode\fontsize{6.6pt}{7.5pt}\selectfont\strut Retour sur 1 lorry. Temps froid. Brouillard.\strut & \smash{\parbox[t]{39mm}{\fontsize{6.2pt}{7.0pt}\selectfont\raggedright\color{black!88}\textbf{8 lorry} \textit{angl. in ms.}}} \\
      \leavevmode\fontsize{6.6pt}{7.5pt}\selectfont\strut Fini journée ds notre office, Hérissart, à lire\strut & \smash{\parbox[t]{39mm}{\fontsize{6.2pt}{7.0pt}\selectfont\raggedright\color{black!88}\textbf{9 ds} \textit{abbr. pro} dans.}} \\
      \leavevmode\llap{\fontsize{6.4pt}{7.3pt}\selectfont\color{darkgray}10\hspace{1.8mm}}\fontsize{6.6pt}{7.5pt}\selectfont\strut ou à écrire.\strut &  \\
      \leavevmode\fontsize{6.6pt}{7.5pt}\selectfont\strut [ Toutencourt, Mission fr. : 43 f. 75 (retenue\strut & \smash{\parbox[t]{39mm}{\fontsize{6.2pt}{7.0pt}\selectfont\raggedright\color{black!88}\textbf{11} \textit{in marg. cum unco} (comptes). \textbf{11 f.} \textit{abbr. pro} franc(s). — \textbf{12 f.} \textit{abbr. pro} franc(s). — \textbf{13 Dim.} \textit{scil.} Dimanche. — \textbf{14 ds} \textit{abbr. pro} dans. — \textbf{16 Cap.} \textit{abbr. pro} Capitaine.}} \\
      \leavevmode\fontsize{6.6pt}{7.5pt}\selectfont\strut de 35 f. pour 7 jours de permission) ]\strut &  \\
      \cline{1-1}
      \leavevmode\fontsize{6.6pt}{7.5pt}\selectfont\strut \underline{4 Dim.} — Nuit, gelé un peu.\strut &  \\
      \leavevmode\fontsize{6.6pt}{7.5pt}\selectfont\strut Matin \& soir, froid un peu vif. Mais ds la\strut &  \\
      \leavevmode\llap{\fontsize{6.4pt}{7.3pt}\selectfont\color{darkgray}15\hspace{1.8mm}}\fontsize{6.6pt}{7.5pt}\selectfont\strut journée, très beau temps, soleil chaud.\strut &  \\
      \leavevmode\fontsize{6.6pt}{7.5pt}\selectfont\strut Cap. \& son ami, en ballade à\strut &  \\
      \leavevmode\fontsize{6.6pt}{7.5pt}\selectfont\strut Flixecourt. Aussi, liberté entière,\strut &  \\
      \leavevmode\fontsize{6.6pt}{7.5pt}\selectfont\strut matin \& soir — \hfill{} [ Washing : 60 c. ]\strut & \smash{\parbox[t]{39mm}{\fontsize{6.2pt}{7.0pt}\selectfont\raggedright\color{black!88}\textbf{18} \textit{in marg. cum unco} (comptes). \textbf{18 c.} \textit{abbr. pro} centimes.}} \\
      \cline{1-1}
      \leavevmode\fontsize{6.6pt}{7.5pt}\selectfont\strut \underline{5 Lundi.} — Au réveil, épaisse couche de\strut &  \\
      \leavevmode\llap{\fontsize{6.4pt}{7.3pt}\selectfont\color{darkgray}20\hspace{1.8mm}}\fontsize{6.6pt}{7.5pt}\selectfont\strut neige recouvrait routes \& toits. Branches\strut &  \\
      \leavevmode\fontsize{6.6pt}{7.5pt}\selectfont\strut d'arbres, très belles, sous leur ouate blanche.\strut &  \\
      \leavevmode\fontsize{6.6pt}{7.5pt}\selectfont\strut Matin, course à Montigny \& à Fréchencourt.\strut &  \\
      \leavevmode\fontsize{6.6pt}{7.5pt}\selectfont\strut Après-midi, à Rubempré \& Puchevillers. Au\strut &  \\
      \leavevmode\fontsize{6.6pt}{7.5pt}\selectfont\strut retour, fait 500 m. à pied sur plateau —\strut &  \\
      \leavevmode\llap{\fontsize{6.4pt}{7.3pt}\selectfont\color{darkgray}25\hspace{1.8mm}}\fontsize{6.6pt}{7.5pt}\selectfont\strut Appris arrivée très prochaine off. australien, en Cte [?]\strut & \smash{\parbox[t]{39mm}{\fontsize{6.2pt}{7.0pt}\selectfont\raggedright\color{black!88}\textbf{25 Cte} \textit{lect. incert. in ms.} \textbf{25 off.} \textit{abbr. pro} officier(s). — \textbf{26 ns} \textit{abbr. pro} nous. \textbf{26 Cap.} \textit{abbr. pro} Capitaine. — \textbf{27 mervt} \textit{abbr. pro} merveilleusement.}} \\
      \leavevmode\fontsize{6.6pt}{7.5pt}\selectfont\strut interprète, pour ns remplacer. Très content de quitter le Cap.\strut &  \\
      \leavevmode\fontsize{6.6pt}{7.5pt}\selectfont\strut Nuit, mervt ciel étoilé, pleine lune — superbe.\strut &  \\
\end{tabular}}%
\par\vspace*{\fill}
% --- Registres de bas de page (1. Critique interne [si pertinent] ; 2. Loci paralleli en latin ; 3. Commentaire historique en retrait) ---
\noindent\rule{\textwidth}{0.5pt}\par\vspace{0.65mm}
{\fontsize{6.6pt}{7.7pt}\selectfont%
\noindent\normalfont\color{black!92} \textbf{4} (\textit{Rubempré}) cf. supra f\textsuperscript{o}~47~v\textsuperscript{o} (p.~\pageref{fol:95}), f\textsuperscript{o}~80~v\textsuperscript{o} (p.~\pageref{fol:161}), infra f\textsuperscript{o}~86~r\textsuperscript{o} (p.~\pageref{fol:172}), f\textsuperscript{o}~86~v\textsuperscript{o} (p.~\pageref{fol:173}), \textit{et passim}. — \textbf{5} (\textit{Puchevillers}) cf. supra f\textsuperscript{o}~37~v\textsuperscript{o} (p.~\pageref{fol:75}), f\textsuperscript{o}~79~r\textsuperscript{o} (p.~\pageref{fol:158}), infra f\textsuperscript{o}~83~r\textsuperscript{o} (p.~\pageref{fol:166}), f\textsuperscript{o}~83~v\textsuperscript{o} (p.~\pageref{fol:167}), \textit{et passim}. — \textbf{7} (\textit{Toutencourt}) cf. supra f\textsuperscript{o}~49~v\textsuperscript{o} (p.~\pageref{fol:99}), f\textsuperscript{o}~80~v\textsuperscript{o} (p.~\pageref{fol:161}), infra f\textsuperscript{o}~81~v\textsuperscript{o} (p.~\pageref{fol:163}), f\textsuperscript{o}~82~r\textsuperscript{o} (p.~\pageref{fol:164}), \textit{et passim}.\par\vspace{0.55mm}
\noindent\rule{38mm}{0.4pt}\par\vspace{0.5mm}
\noindent\hspace*{5mm}\parbox{\dimexpr\textwidth-5mm\relax}{\fontsize{6.6pt}{7.7pt}\selectfont\color{black!84}Le samedi 3 mars 1917 — échéance bimensuelle de la Mission de la V\textsuperscript{e} Armée à \textbf{Toutencourt} —, Hippolyte Combes perçoit 43~fr.~75 net (soit sa quinzaine normale de 78~fr.~75 diminuée d'une retenue réglementaire de 35~fr., à raison de 5~fr. par jour, pour ses 7 jours de permission de fin janvier--février).}\par}
\newpage
\phantomsection\label{fac:162}
% --- En-tête du fac-similé (Page impaire / droite : Photographie du manuscrit) ---
\noindent\rule{\textwidth}{0.5pt}\par\vspace{0.55mm}
\noindent\textbf{\fontsize{6.8pt}{7.8pt}\selectfont FAC-SIMILÉ EN VIS-À-VIS — FOLIO 81 RECTO} \hfill \textbf{\fontsize{6.8pt}{7.8pt}\selectfont MARS 1917}\par\vspace{0.25mm}
\noindent{\fontsize{6.8pt}{8.0pt}\selectfont\textit{Livre-Journal d'Hippolyte Combes} — 3 – 5 mars 1917 (Hérissart, Rubempré, Puchevillers, Toutencourt, Montigny, Fréchencourt (Somme))}\par\vspace{0.55mm}
\noindent\rule{\textwidth}{0.35pt}\par\vspace*{\fill}
{\centering
  \setlength{\fboxsep}{1.2mm}%
  \setlength{\fboxrule}{0.4pt}%
  \fcolorbox{gray!55}{white}{%
    \includegraphics[width=0.95\textwidth,height=0.81\textheight,keepaspectratio]{\facsimiledir/20251126_111105.jpg}%
  }\par}
\vspace*{\fill}
\newpage

% =============================================================================
% VIS-A-VIS SPREAD : FOLIO 81 VERSO (20251126_111116.jpg)
% =============================================================================
\cleardoubleevenpage
\renewcommand{\currentmonthyear}{Mars 1917}
\renewcommand{\currentfoliolabel}{Folio 81 verso}
\phantomsection\label{fol:163}
% --- En-tête du folio (Page paire / gauche : Transcription, apparat & registres de bas de page) ---
\noindent\rule{\textwidth}{0.5pt}\par\vspace{0.55mm}
\noindent\textbf{\footnotesize FOLIO 81 VERSO} \hfill \textbf{\footnotesize MARS 1917}\par\vspace{0.25mm}
\noindent{\fontsize{6.9pt}{8.0pt}\selectfont Courses à Doullens avec Williams et promesse du maréchal des logis Grauger à Toutencourt}\par\vspace{0.2mm}
\noindent{\fontsize{6.9pt}{8.0pt}\selectfont\textit{6 – 7 mars 1917} \hfill \textit{Hérissart, Doullens, Toutencourt (Somme)}}\par\vspace{0.5mm}
\noindent\rule{\textwidth}{0.35pt}\par\vspace{0.8mm}
% --- Corps de la transcription diplomatique et apparat marginal en prolongement ---
{\renewcommand{\arraystretch}{0.83}%
\noindent\hspace*{3.5mm}\begin{tabular}{@{}p{87mm}@{\hspace{3.5mm}}p{39mm}@{}}
      \leavevmode\fontsize{6.6pt}{7.5pt}\selectfont\strut \underline{6 Mardi} — \hfill{} \underline{— Mars 1917 —}\strut &  \\
      \leavevmode\fontsize{6.6pt}{7.5pt}\selectfont\strut Tout le jour, soufflé bise glaciale.\strut &  \\
      \leavevmode\fontsize{6.6pt}{7.5pt}\selectfont\strut Ciel resté nuageux.\strut &  \\
      \leavevmode\fontsize{6.6pt}{7.5pt}\selectfont\strut Matinée, entière libre. Nouvel off. arrivé,\strut & \smash{\parbox[t]{39mm}{\fontsize{6.2pt}{7.0pt}\selectfont\raggedright\color{black!88}\textbf{4 off.} \textit{abbr. pro} officier(s).}} \\
      \leavevmode\llap{\fontsize{6.4pt}{7.3pt}\selectfont\color{darkgray}5\hspace{1.8mm}}\fontsize{6.6pt}{7.5pt}\selectfont\strut australien, plutôt vulgaire.\strut &  \\
      \leavevmode\fontsize{6.6pt}{7.5pt}\selectfont\strut Après-midi, venu à Doullens, avec Cap. W.\strut & \smash{\parbox[t]{39mm}{\fontsize{6.2pt}{7.0pt}\selectfont\raggedright\color{black!88}\textbf{6 Cap. W.} \textit{scil.} Capitaine Williams. \textbf{6 Cap.} \textit{abbr. pro} Capitaine. — \textbf{8 s.} \textit{abbr. pro} soir. — \textbf{9 Cap.} \textit{abbr. pro} Capitaine. \textbf{9 Maral} \textit{abbr. pro} Maréchal. — \textbf{10 renseigts} \textit{abbr. pro} renseignements. \textbf{10 Maral} \textit{abbr. pro} Maréchal. — \textbf{12 cie} \textit{abbr. pro} compagnie. \textbf{12 mess} \textit{angl. in ms.} — \textbf{13} \textit{in marg. cum unco} (comptes). \textbf{13 f.} \textit{abbr. pro} franc(s). — \textbf{14 Cartes ples} \textit{abbr. pro} Cartes postales. \textbf{14 f.} \textit{abbr. pro} franc(s). — \textbf{15 c.} \textit{abbr. pro} centimes. \textbf{15 wounded} \textit{angl. in ms.}}} \\
      \leavevmode\fontsize{6.6pt}{7.5pt}\selectfont\strut Nouvel interprète, arrivé, très sympathique.\strut &  \\
      \leavevmode\fontsize{6.6pt}{7.5pt}\selectfont\strut Vers 6h s., allé à pied Toutencourt, Mission fr.\strut &  \\
      \leavevmode\fontsize{6.6pt}{7.5pt}\selectfont\strut Maral logis Grauger m'a conduit chz le Cap.\strut &  \\
      \leavevmode\llap{\fontsize{6.4pt}{7.3pt}\selectfont\color{darkgray}10\hspace{1.8mm}}\fontsize{6.6pt}{7.5pt}\selectfont\strut Raucoule, l. a demandé divers renseigts. Maral logis Grauger m'a promis\strut &  \\
      \leavevmode\fontsize{6.6pt}{7.5pt}\selectfont\strut de me faire donner bon poste —\strut &  \\
      \leavevmode\fontsize{6.6pt}{7.5pt}\selectfont\strut Fini journée au mess, en cie nouvel interprète.\strut &  \\
      \leavevmode\fontsize{6.6pt}{7.5pt}\selectfont\strut [ \underline{A Doullens} : papeterie : 3 f. + 1 f. 50\strut &  \\
      \leavevmode\fontsize{6.6pt}{7.5pt}\selectfont\strut Cartes ples : 1 f. — 2 boîtes pastilles chlorate : 1 f.\strut &  \\
      \leavevmode\llap{\fontsize{6.4pt}{7.3pt}\selectfont\color{darkgray}15\hspace{1.8mm}}\fontsize{6.6pt}{7.5pt}\selectfont\strut Quête pour wounded : 50 c.\strut &  \\
      \leavevmode\fontsize{6.6pt}{7.5pt}\selectfont\strut (Hôpital auxiliaire) ]\strut &  \\
      \cline{1-1}
      \leavevmode\fontsize{6.6pt}{7.5pt}\selectfont\strut \underline{7 Mercredi.}\strut &  \\
      \leavevmode\fontsize{6.6pt}{7.5pt}\selectfont\strut Matin \& soir, inoccupé.\strut &  \\
      \leavevmode\fontsize{6.6pt}{7.5pt}\selectfont\strut Préparé lentement paquets \& sac.\strut &  \\
      \leavevmode\llap{\fontsize{6.4pt}{7.3pt}\selectfont\color{darkgray}20\hspace{1.8mm}}\fontsize{6.6pt}{7.5pt}\selectfont\strut Ciel resté couvert. Bise glaciale. Après-midi,\strut &  \\
      \leavevmode\fontsize{6.6pt}{7.5pt}\selectfont\strut vent soufflé en violentes rafales et fait\strut &  \\
      \leavevmode\fontsize{6.6pt}{7.5pt}\selectfont\strut bruit de tempête à travers branches des\strut &  \\
      \leavevmode\fontsize{6.6pt}{7.5pt}\selectfont\strut arbres.\strut &  \\
      \leavevmode\fontsize{6.6pt}{7.5pt}\selectfont\strut Courte promenade vers 4h, autour\strut &  \\
      \leavevmode\llap{\fontsize{6.4pt}{7.3pt}\selectfont\color{darkgray}25\hspace{1.8mm}}\fontsize{6.6pt}{7.5pt}\selectfont\strut et ds les environs village. Visité cimetière.\strut & \smash{\parbox[t]{39mm}{\fontsize{6.2pt}{7.0pt}\selectfont\raggedright\color{black!88}\textbf{25 ds} \textit{abbr. pro} dans.}} \\
      \leavevmode\fontsize{6.6pt}{7.5pt}\selectfont\strut Temps funèbre.\strut &  \\
\end{tabular}}%
\par\vspace*{\fill}
% --- Registres de bas de page (1. Critique interne [si pertinent] ; 2. Loci paralleli en latin ; 3. Commentaire historique en retrait) ---
\noindent\rule{\textwidth}{0.5pt}\par\vspace{0.65mm}
{\fontsize{6.6pt}{7.7pt}\selectfont%
\noindent\normalfont\color{black!92} \textbf{6} (\textit{Doullens}) cf. supra f\textsuperscript{o}~25~v\textsuperscript{o} (p.~\pageref{fol:51}), f\textsuperscript{o}~79~r\textsuperscript{o} (p.~\pageref{fol:158}), infra f\textsuperscript{o}~89~v\textsuperscript{o} (p.~\pageref{fol:179}), \textit{et passim}. — \textbf{8} (\textit{Toutencourt}) cf. supra f\textsuperscript{o}~49~v\textsuperscript{o} (p.~\pageref{fol:99}), f\textsuperscript{o}~81~r\textsuperscript{o} (p.~\pageref{fol:162}), infra f\textsuperscript{o}~82~r\textsuperscript{o} (p.~\pageref{fol:164}), f\textsuperscript{o}~83~r\textsuperscript{o} (p.~\pageref{fol:166}), \textit{et passim}.\par\vspace{0.55mm}
\noindent\rule{38mm}{0.4pt}\par\vspace{0.5mm}
\noindent\hspace*{5mm}\parbox{\dimexpr\textwidth-5mm\relax}{\fontsize{6.6pt}{7.7pt}\selectfont\color{black!84}À l'arrivée de l'officier australien et de son interprète (nommé \textbf{Valet}, de la \textit{2nd Australian Division}, au f\textsuperscript{o}~82~r\textsuperscript{o}) venus relever le poste d'Hérissart le mardi 6 mars 1917, Hippolyte Combes accompagne une dernière fois le capitaine Williams à \textbf{Doullens} (où il verse 50~c. à une quête pour les blessés d'un hôpital auxiliaire), puis se présente à 18~h au bureau de la Mission française de la V\textsuperscript{e} Armée à \textbf{Toutencourt}, auprès du \textbf{maréchal des logis Grauger} et du chef de mission, le \textbf{capitaine Raucoule}.}\par}
\newpage
\phantomsection\label{fac:163}
% --- En-tête du fac-similé (Page impaire / droite : Photographie du manuscrit) ---
\noindent\rule{\textwidth}{0.5pt}\par\vspace{0.55mm}
\noindent\textbf{\fontsize{6.8pt}{7.8pt}\selectfont FAC-SIMILÉ EN VIS-À-VIS — FOLIO 81 VERSO} \hfill \textbf{\fontsize{6.8pt}{7.8pt}\selectfont MARS 1917}\par\vspace{0.25mm}
\noindent{\fontsize{6.8pt}{8.0pt}\selectfont\textit{Livre-Journal d'Hippolyte Combes} — 6 – 7 mars 1917 (Hérissart, Doullens, Toutencourt (Somme))}\par\vspace{0.55mm}
\noindent\rule{\textwidth}{0.35pt}\par\vspace*{\fill}
{\centering
  \setlength{\fboxsep}{1.2mm}%
  \setlength{\fboxrule}{0.4pt}%
  \fcolorbox{gray!55}{white}{%
    \includegraphics[width=0.95\textwidth,height=0.81\textheight,keepaspectratio]{\facsimiledir/20251126_111116.jpg}%
  }\par}
\vspace*{\fill}
\newpage

% =============================================================================
% VIS-A-VIS SPREAD : FOLIO 82 RECTO (20251126_111137.jpg)
% =============================================================================
\cleardoubleevenpage
\renewcommand{\currentmonthyear}{Mars 1917}
\renewcommand{\currentfoliolabel}{Folio 82 recto}
\phantomsection\label{fol:164}
% --- En-tête du folio (Page paire / gauche : Transcription, apparat & registres de bas de page) ---
\noindent\rule{\textwidth}{0.5pt}\par\vspace{0.55mm}
\noindent\textbf{\footnotesize FOLIO 82 RECTO} \hfill \textbf{\footnotesize MARS 1917}\par\vspace{0.25mm}
\noindent{\fontsize{6.9pt}{8.0pt}\selectfont Tempête de neige de mars et départ de l'interprète Valet pour Albert}\par\vspace{0.2mm}
\noindent{\fontsize{6.9pt}{8.0pt}\selectfont\textit{8 – 9 mars 1917} \hfill \textit{Hérissart, Toutencourt (Somme)}}\par\vspace{0.5mm}
\noindent\rule{\textwidth}{0.35pt}\par\vspace{0.8mm}
% --- Corps de la transcription diplomatique et apparat marginal en prolongement ---
{\renewcommand{\arraystretch}{0.83}%
\noindent\hspace*{3.5mm}\begin{tabular}{@{}p{87mm}@{\hspace{3.5mm}}p{39mm}@{}}
      \leavevmode\fontsize{6.6pt}{7.5pt}\selectfont\strut \underline{8 Jeudi} — \hfill{} \underline{— Mars 1917 —}\strut &  \\
      \leavevmode\fontsize{6.6pt}{7.5pt}\selectfont\strut Nuit très froide. Gelé fort!\strut &  \\
      \leavevmode\fontsize{6.6pt}{7.5pt}\selectfont\strut Durant matinée, neige tombée à gros flocons.\strut &  \\
      \leavevmode\fontsize{6.6pt}{7.5pt}\selectfont\strut Ciel mat \& comme voilé de brume.\strut &  \\
      \leavevmode\llap{\fontsize{6.4pt}{7.3pt}\selectfont\color{darkgray}5\hspace{1.8mm}}\fontsize{6.6pt}{7.5pt}\selectfont\strut Liberté entière, matin \& soir.\strut &  \\
      \leavevmode\fontsize{6.6pt}{7.5pt}\selectfont\strut Vers 3h s., intte Evar[?], revenu en cie d'1 autre\strut & \smash{\parbox[t]{39mm}{\fontsize{6.2pt}{7.0pt}\selectfont\raggedright\color{black!88}\textbf{6 Evar} \textit{lect. incert. in ms.} \textbf{6 cie} \textit{abbr. pro} compagnie. — \textbf{7 intte} \textit{abbr. pro} interprète. — \textbf{8 Evar} \textit{lect. incert. in ms.} \textbf{8 lorry} \textit{angl. in ms.} — \textbf{9 extrt} \textit{abbr. pro} extrêmement. \textbf{9 kit} \textit{angl. in ms.} — \textbf{12 j'à} \textit{abbr. pro} jusqu'à. — \textbf{14 walk} \textit{angl. in ms.} — \textbf{15 qq} \textit{abbr. pro} quelque(s). \textbf{15 mervt} \textit{abbr. pro} merveilleusement.}} \\
      \leavevmode\fontsize{6.6pt}{7.5pt}\selectfont\strut intte Valet, R.2, de la 2nd australian division.\strut &  \\
      \leavevmode\fontsize{6.6pt}{7.5pt}\selectfont\strut Evar[?], reparti à Albert sur 1 lorry avec\strut &  \\
      \leavevmode\fontsize{6.6pt}{7.5pt}\selectfont\strut tout son kit. Valet, extrt sympathique.\strut &  \\
      \leavevmode\llap{\fontsize{6.4pt}{7.3pt}\selectfont\color{darkgray}10\hspace{1.8mm}}\fontsize{6.6pt}{7.5pt}\selectfont\strut Aussitôt thé pris, allé à pied Mission fr.\strut &  \\
      \leavevmode\fontsize{6.6pt}{7.5pt}\selectfont\strut Toutencourt : maréchal logis..., m'a dit\strut &  \\
      \leavevmode\fontsize{6.6pt}{7.5pt}\selectfont\strut de rester av. mon cap. Hérissart j'à son départ,\strut &  \\
      \leavevmode\fontsize{6.6pt}{7.5pt}\selectfont\strut puis me rendre au dépôt, Toutencourt.\strut &  \\
      \leavevmode\fontsize{6.6pt}{7.5pt}\selectfont\strut Très agréable walk sur route très sèche. Temps\strut &  \\
      \leavevmode\llap{\fontsize{6.4pt}{7.3pt}\selectfont\color{darkgray}15\hspace{1.8mm}}\fontsize{6.6pt}{7.5pt}\selectfont\strut mervt pleine lune, ciel nuageux \& qq étoiles.\strut &  \\
      \cline{1-1}
      \leavevmode\fontsize{6.6pt}{7.5pt}\selectfont\strut \underline{9 Vendredi.} — Nuit, moins froide. Gelé un peu.\strut &  \\
      \leavevmode\fontsize{6.6pt}{7.5pt}\selectfont\strut Éveillé chants d'oiseaux.\strut &  \\
      \leavevmode\fontsize{6.6pt}{7.5pt}\selectfont\strut Matinée, fait ts préparatifs départ.\strut & \smash{\parbox[t]{39mm}{\fontsize{6.2pt}{7.0pt}\selectfont\raggedright\color{black!88}\textbf{18 ts} \textit{abbr. pro} tous.}} \\
      \leavevmode\fontsize{6.6pt}{7.5pt}\selectfont\strut Neige tombée à flocons serrés durant\strut &  \\
      \leavevmode\llap{\fontsize{6.4pt}{7.3pt}\selectfont\color{darkgray}20\hspace{1.8mm}}\fontsize{6.6pt}{7.5pt}\selectfont\strut plusieurs heures, surtout l'après-midi.\strut &  \\
      \leavevmode\fontsize{6.6pt}{7.5pt}\selectfont\strut Liberté entière journée. Bien reposé, office.\strut &  \\
      \leavevmode\fontsize{6.6pt}{7.5pt}\selectfont\strut Fin soirée, tt pensif, tristesse, changement\strut & \smash{\parbox[t]{39mm}{\fontsize{6.2pt}{7.0pt}\selectfont\raggedright\color{black!88}\textbf{22 tt} \textit{abbr. pro} tout.}} \\
      \leavevmode\fontsize{6.6pt}{7.5pt}\selectfont\strut de siton, \& incertitude au sujet nouvelle affectation.\strut & \smash{\parbox[t]{39mm}{\fontsize{6.2pt}{7.0pt}\selectfont\raggedright\color{black!88}\textbf{23 siton} \textit{abbr. pro} situation. — \textbf{24} \textit{in marg. cum unco} (comptes). \textbf{24 dinner} \textit{angl. in ms.} — \textbf{25} \textit{lect. incert. in ms.}}} \\
      \cline{1-1}
      \leavevmode\fontsize{6.6pt}{7.5pt}\selectfont\strut [ Après dinner, réglé comptes mess avec cap. Wms\strut &  \\
      \leavevmode\llap{\fontsize{6.4pt}{7.3pt}\selectfont\color{darkgray}25\hspace{1.8mm}}\fontsize{6.6pt}{7.5pt}\selectfont\strut jusqu'à ce soir : [18 f. ?] — ]\strut &  \\
\end{tabular}}%
\par\vspace*{\fill}
% --- Registres de bas de page (1. Critique interne [si pertinent] ; 2. Loci paralleli en latin ; 3. Commentaire historique en retrait) ---
\noindent\rule{\textwidth}{0.5pt}\par\vspace{0.65mm}
{\fontsize{6.6pt}{7.7pt}\selectfont%
\noindent\normalfont\color{black!92} \textbf{7} (\textit{Valet}) cf. infra f\textsuperscript{o}~83~r\textsuperscript{o} (p.~\pageref{fol:166}). — \textbf{8} (\textit{Albert}) cf. supra f\textsuperscript{o}~39~r\textsuperscript{o} (p.~\pageref{fol:78}), f\textsuperscript{o}~50~r\textsuperscript{o} (p.~\pageref{fol:100}), infra f\textsuperscript{o}~84~v\textsuperscript{o} (p.~\pageref{fol:169}), f\textsuperscript{o}~85~v\textsuperscript{o} (p.~\pageref{fol:171}), \textit{et passim}. — \textbf{11} (\textit{Toutencourt}) cf. supra f\textsuperscript{o}~49~v\textsuperscript{o} (p.~\pageref{fol:99}), f\textsuperscript{o}~81~v\textsuperscript{o} (p.~\pageref{fol:163}), infra f\textsuperscript{o}~83~r\textsuperscript{o} (p.~\pageref{fol:166}), f\textsuperscript{o}~83~v\textsuperscript{o} (p.~\pageref{fol:167}), \textit{et passim}.\par\vspace{0.55mm}
\noindent\rule{38mm}{0.4pt}\par\vspace{0.5mm}
\noindent\hspace*{5mm}\parbox{\dimexpr\textwidth-5mm\relax}{\fontsize{6.6pt}{7.7pt}\selectfont\color{black!84}L'interprète \textbf{Valet}, rencontré le 8 mars 1917 à Hérissart, est rattaché à la \textbf{2\textsuperscript{e} Division australienne} (\textit{2nd Australian Division}, unité de l'\textit{I~Anzac Corps} alors engagée devant Bapaume). Le 9 mars au soir, Hippolyte Combes solde ses comptes de \textit{mess} avec le \textbf{capitaine Williams} à la veille du départ de ce dernier.}\par}
\newpage
\phantomsection\label{fac:164}
% --- En-tête du fac-similé (Page impaire / droite : Photographie du manuscrit) ---
\noindent\rule{\textwidth}{0.5pt}\par\vspace{0.55mm}
\noindent\textbf{\fontsize{6.8pt}{7.8pt}\selectfont FAC-SIMILÉ EN VIS-À-VIS — FOLIO 82 RECTO} \hfill \textbf{\fontsize{6.8pt}{7.8pt}\selectfont MARS 1917}\par\vspace{0.25mm}
\noindent{\fontsize{6.8pt}{8.0pt}\selectfont\textit{Livre-Journal d'Hippolyte Combes} — 8 – 9 mars 1917 (Hérissart, Toutencourt (Somme))}\par\vspace{0.55mm}
\noindent\rule{\textwidth}{0.35pt}\par\vspace*{\fill}
{\centering
  \setlength{\fboxsep}{1.2mm}%
  \setlength{\fboxrule}{0.4pt}%
  \fcolorbox{gray!55}{white}{%
    \includegraphics[width=0.95\textwidth,height=0.81\textheight,keepaspectratio]{\facsimiledir/20251126_111137.jpg}%
  }\par}
\vspace*{\fill}
\newpage

% =============================================================================
% VIS-A-VIS SPREAD : FOLIO 82 VERSO (20251126_111151.jpg)
% =============================================================================
\cleardoubleevenpage
\renewcommand{\currentmonthyear}{Mars 1917}
\renewcommand{\currentfoliolabel}{Folio 82 verso}
\phantomsection\label{fol:165}
% --- En-tête du folio (Page paire / gauche : Transcription, apparat & registres de bas de page) ---
\noindent\rule{\textwidth}{0.5pt}\par\vspace{0.55mm}
\noindent\textbf{\footnotesize FOLIO 82 VERSO} \hfill \textbf{\footnotesize MARS 1917}\par\vspace{0.25mm}
\noindent{\fontsize{6.9pt}{8.0pt}\selectfont Départ définitif du capitaine Williams et mélancolie de l'attente à Hérissart}\par\vspace{0.2mm}
\noindent{\fontsize{6.9pt}{8.0pt}\selectfont\textit{10 mars 1917} \hfill \textit{Hérissart (Somme)}}\par\vspace{0.5mm}
\noindent\rule{\textwidth}{0.35pt}\par\vspace{0.8mm}
% --- Corps de la transcription diplomatique et apparat marginal en prolongement ---
{\renewcommand{\arraystretch}{0.93}%
\noindent\hspace*{3.5mm}\begin{tabular}{@{}p{87mm}@{\hspace{3.5mm}}p{39mm}@{}}
      \leavevmode\fontsize{7.3pt}{8.4pt}\selectfont\strut \underline{10 Samedi} — \hfill{} \underline{Mars 17 —}\strut &  \\
      \leavevmode\fontsize{7.3pt}{8.4pt}\selectfont\strut Avant breakfast, terminé ts mes\strut & \smash{\parbox[t]{39mm}{\fontsize{6.2pt}{7.0pt}\selectfont\raggedright\color{black!88}\textbf{2 ts} \textit{abbr. pro} tous. \textbf{2 breakfast} \textit{angl. in ms.}}} \\
      \leavevmode\fontsize{7.3pt}{8.4pt}\selectfont\strut paquets, sanglé valise.\strut &  \\
      \leavevmode\fontsize{7.3pt}{8.4pt}\selectfont\strut Étrange sensation de mélancolie, au moment de\strut &  \\
      \leavevmode\llap{\fontsize{6.4pt}{7.3pt}\selectfont\color{darkgray}5\hspace{1.8mm}}\fontsize{7.3pt}{8.4pt}\selectfont\strut quitter ce village et de changer d'affectation.\strut &  \\
      \leavevmode\fontsize{7.3pt}{8.4pt}\selectfont\strut Malgré moi, souvenirs affluent de ces derniers\strut &  \\
      \leavevmode\fontsize{7.3pt}{8.4pt}\selectfont\strut mois. Incertitude lendemain, \& situation précaire.\strut &  \\
      \leavevmode\fontsize{7.3pt}{8.4pt}\selectfont\strut Cap. nouvl, Williams, partis défintt vers\strut & \smash{\parbox[t]{39mm}{\fontsize{6.2pt}{7.0pt}\selectfont\raggedright\color{black!88}\textbf{8 défintt} \textit{abbr. pro} définitivement. \textbf{8 Cap.} \textit{abbr. pro} Capitaine.}} \\
      \leavevmode\fontsize{7.3pt}{8.4pt}\selectfont\strut 10h matin. Dès lors, suis sans fonctions, archi-\strut &  \\
      \leavevmode\llap{\fontsize{6.4pt}{7.3pt}\selectfont\color{darkgray}10\hspace{1.8mm}}\fontsize{7.3pt}{8.4pt}\selectfont\strut libre.\strut &  \\
      \leavevmode\fontsize{7.3pt}{8.4pt}\selectfont\strut Voulu rester encore 1 jour village et ferme\strut &  \\
      \leavevmode\fontsize{7.3pt}{8.4pt}\selectfont\strut où belle chambre. Fraternisé avec English connus\strut &  \\
      \leavevmode\fontsize{7.3pt}{8.4pt}\selectfont\strut depuis lgt.\strut & \smash{\parbox[t]{39mm}{\fontsize{6.2pt}{7.0pt}\selectfont\raggedright\color{black!88}\textbf{13 lgt} \textit{abbr. pro} longtemps.}} \\
      \leavevmode\fontsize{7.3pt}{8.4pt}\selectfont\strut Ciel nuageux. Temps très doux —\strut &  \\
      \leavevmode\llap{\fontsize{6.4pt}{7.3pt}\selectfont\color{darkgray}15\hspace{1.8mm}}\fontsize{7.3pt}{8.4pt}\selectfont\strut Fin journée, très calme.\strut &  \\
      \leavevmode\fontsize{7.3pt}{8.4pt}\selectfont\strut Reposé, écrit, lu, dans l'office.\strut &  \\
      \leavevmode\fontsize{7.3pt}{8.4pt}\selectfont\strut Temps maussade, bruine \& boue.\strut &  \\
      \leavevmode\fontsize{7.3pt}{8.4pt}\selectfont\strut Matinée, tt pénétré mélancolie.\strut & \smash{\parbox[t]{39mm}{\fontsize{6.2pt}{7.0pt}\selectfont\raggedright\color{black!88}\textbf{18 tt} \textit{abbr. pro} tout.}} \\
      \leavevmode\fontsize{7.3pt}{8.4pt}\selectfont\strut Mais ce soir, pris mon parti, et résigné\strut &  \\
      \leavevmode\llap{\fontsize{6.4pt}{7.3pt}\selectfont\color{darkgray}20\hspace{1.8mm}}\fontsize{7.3pt}{8.4pt}\selectfont\strut à mon nouveau sort.\strut &  \\
\end{tabular}}%
\par\vspace*{\fill}
% --- Registres de bas de page (1. Critique interne [si pertinent] ; 2. Loci paralleli en latin ; 3. Commentaire historique en retrait) ---
\noindent\rule{\textwidth}{0.5pt}\par\vspace{0.65mm}
{\fontsize{6.6pt}{7.7pt}\selectfont%
\noindent\normalfont \textbf{Unité de page :} La page réunit le départ du capitaine Williams le samedi 10 mars 1917 (l.~1--18, avec citation anglaise au style direct à la l.~9) et un état récapitulatif de caisse au bas du feuillet (l.~19--22).\par\vspace{0.5mm}
\noindent\rule{22mm}{0.3pt}\par\vspace{0.45mm}
\noindent\normalfont\color{black!92} \textbf{8} (\textit{Williams}) cf. supra f\textsuperscript{o}~52~v\textsuperscript{o} (p.~\pageref{fol:105}), f\textsuperscript{o}~82~r\textsuperscript{o} (p.~\pageref{fol:164}), \textit{et passim}.\par\vspace{0.55mm}
\noindent\rule{38mm}{0.4pt}\par\vspace{0.5mm}
\noindent\hspace*{5mm}\parbox{\dimexpr\textwidth-5mm\relax}{\fontsize{6.6pt}{7.7pt}\selectfont\color{black!84}Le samedi 10 mars 1917 à 10~h du matin s'achève le service d'Hippolyte Combes auprès du \textbf{capitaine Williams} (qui aura duré exactement trois mois et huit jours, du 2 décembre 1916 au 10 mars 1917). Resté vingt-quatre heures sans affectation à Hérissart avant de rejoindre le dépôt de la Mission française à Toutencourt, l'interprète fraternise avec les sous-officiers et soldats britanniques (\textit{English}) du détachement.}\par}
\newpage
\phantomsection\label{fac:165}
% --- En-tête du fac-similé (Page impaire / droite : Photographie du manuscrit) ---
\noindent\rule{\textwidth}{0.5pt}\par\vspace{0.55mm}
\noindent\textbf{\fontsize{6.8pt}{7.8pt}\selectfont FAC-SIMILÉ EN VIS-À-VIS — FOLIO 82 VERSO} \hfill \textbf{\fontsize{6.8pt}{7.8pt}\selectfont MARS 1917}\par\vspace{0.25mm}
\noindent{\fontsize{6.8pt}{8.0pt}\selectfont\textit{Livre-Journal d'Hippolyte Combes} — 10 mars 1917 (Hérissart (Somme))}\par\vspace{0.55mm}
\noindent\rule{\textwidth}{0.35pt}\par\vspace*{\fill}
{\centering
  \setlength{\fboxsep}{1.2mm}%
  \setlength{\fboxrule}{0.4pt}%
  \fcolorbox{gray!55}{white}{%
    \includegraphics[width=0.95\textwidth,height=0.81\textheight,keepaspectratio]{\facsimiledir/20251126_111151.jpg}%
  }\par}
\vspace*{\fill}
\newpage

% =============================================================================
% VIS-A-VIS SPREAD : FOLIO 83 RECTO (20251126_111240.jpg)
% =============================================================================
\cleardoubleevenpage
\renewcommand{\currentmonthyear}{Mars 1917}
\renewcommand{\currentfoliolabel}{Folio 83 recto}
\phantomsection\label{fol:166}
% --- En-tête du folio (Page paire / gauche : Transcription, apparat & registres de bas de page) ---
\noindent\rule{\textwidth}{0.5pt}\par\vspace{0.55mm}
\noindent\textbf{\footnotesize FOLIO 83 RECTO} \hfill \textbf{\footnotesize MARS 1917}\par\vspace{0.25mm}
\noindent{\fontsize{6.9pt}{8.0pt}\selectfont Maintien au Purchase Board sur demande du Major Levy et installation au billet 220}\par\vspace{0.2mm}
\noindent{\fontsize{6.9pt}{8.0pt}\selectfont\textit{11 mars 1917} \hfill \textit{Hérissart, Toutencourt, Puchevillers (Somme)}}\par\vspace{0.5mm}
\noindent\rule{\textwidth}{0.35pt}\par\vspace{0.8mm}
% --- Corps de la transcription diplomatique et apparat marginal en prolongement ---
{\renewcommand{\arraystretch}{0.83}%
\noindent\hspace*{3.5mm}\begin{tabular}{@{}p{87mm}@{\hspace{3.5mm}}p{39mm}@{}}
      \leavevmode\fontsize{6.6pt}{7.5pt}\selectfont\strut \underline{Dim. 11.} \hfill{} \underline{— Mars 1917 —}\strut & \smash{\parbox[t]{39mm}{\fontsize{6.2pt}{7.0pt}\selectfont\raggedright\color{black!88}\textbf{1 Dim.} \textit{scil.} Dimanche.}} \\
      \leavevmode\fontsize{6.6pt}{7.5pt}\selectfont\strut Levé 6h 1/2 m.\strut & \smash{\parbox[t]{39mm}{\fontsize{6.2pt}{7.0pt}\selectfont\raggedright\color{black!88}\textbf{2 1/2} \textit{scil.} demi / moitié. \textbf{2 m.} \textit{abbr. pro} matin. — \textbf{3 ts} \textit{abbr. pro} tous. — \textbf{4 extrt} \textit{abbr. pro} extrêmement.}} \\
      \leavevmode\fontsize{6.6pt}{7.5pt}\selectfont\strut Bientôt terminé ts mes paquets.\strut &  \\
      \leavevmode\fontsize{6.6pt}{7.5pt}\selectfont\strut Temps extrt doux. Soleil pâle. Éclat clair.\strut &  \\
      \leavevmode\llap{\fontsize{6.4pt}{7.3pt}\selectfont\color{darkgray}5\hspace{1.8mm}}\fontsize{6.6pt}{7.5pt}\selectfont\strut Maintenant, quitte sans regrets ce petit\strut &  \\
      \leavevmode\fontsize{6.6pt}{7.5pt}\selectfont\strut village \& Bureau Board.\strut &  \\
      \leavevmode\fontsize{6.6pt}{7.5pt}\selectfont\strut Après breakfast, 9h m., parti en auto pour\strut & \smash{\parbox[t]{39mm}{\fontsize{6.2pt}{7.0pt}\selectfont\raggedright\color{black!88}\textbf{7 m.} \textit{abbr. pro} matin. \textbf{7 breakfast} \textit{angl. in ms.} — \textbf{8 Ds} \textit{abbr. pro} Dans. \textbf{8 ts} \textit{abbr. pro} tous. — \textbf{9 mess} \textit{angl. in ms.}}} \\
      \leavevmode\fontsize{6.6pt}{7.5pt}\selectfont\strut Toutencourt. Ds ce village, laissé tout, ts mes\strut &  \\
      \leavevmode\fontsize{6.6pt}{7.5pt}\selectfont\strut sacs au mess des interprètes,\strut &  \\
      \leavevmode\llap{\fontsize{6.4pt}{7.3pt}\selectfont\color{darkgray}10\hspace{1.8mm}}\fontsize{6.6pt}{7.5pt}\selectfont\strut de Raincheval, puis cherché gîte au camp,\strut &  \\
      \leavevmode\fontsize{6.6pt}{7.5pt}\selectfont\strut n°, rue\strut &  \\
      \leavevmode\fontsize{6.6pt}{7.5pt}\selectfont\strut Mais au bureau Mission, appris que,\strut &  \\
      \leavevmode\fontsize{6.6pt}{7.5pt}\selectfont\strut sur demande instante du Major Levy, je\strut &  \\
      \leavevmode\fontsize{6.6pt}{7.5pt}\selectfont\strut resterai au Purchase Board.\strut & \smash{\parbox[t]{39mm}{\fontsize{6.2pt}{7.0pt}\selectfont\raggedright\color{black!88}\textbf{14 Purchase Board} \textit{angl. in ms.} — \textbf{15 Lieutt} \textit{abbr. pro} Lieutenant. \textbf{15 interview} \textit{angl. in ms.} — \textbf{16 tt} \textit{abbr. pro} tout. \textbf{16 off.} \textit{abbr. pro} officier(s). — \textbf{18 lorry} \textit{angl. in ms.} — \textbf{19 ts} \textit{abbr. pro} tous.}} \\
      \leavevmode\llap{\fontsize{6.4pt}{7.3pt}\selectfont\color{darkgray}15\hspace{1.8mm}}\fontsize{6.6pt}{7.5pt}\selectfont\strut Long interview avec Lieutt Godard,\strut &  \\
      \leavevmode\fontsize{6.6pt}{7.5pt}\selectfont\strut tt jeune off., charmant.\strut &  \\
      \leavevmode\fontsize{6.6pt}{7.5pt}\selectfont\strut Reparti de Toutencourt vers midi sur\strut &  \\
      \leavevmode\fontsize{6.6pt}{7.5pt}\selectfont\strut 1 lorry. Descendu au tournant de Puchevillers\strut &  \\
      \leavevmode\fontsize{6.6pt}{7.5pt}\selectfont\strut et porté ts mes paquets au billet 220. M'y\strut &  \\
      \leavevmode\llap{\fontsize{6.4pt}{7.3pt}\selectfont\color{darkgray}20\hspace{1.8mm}}\fontsize{6.6pt}{7.5pt}\selectfont\strut suis installé définitivement, après départ\strut &  \\
      \leavevmode\fontsize{6.6pt}{7.5pt}\selectfont\strut interprète Valet.\strut &  \\
      \leavevmode\fontsize{6.6pt}{7.5pt}\selectfont\strut Liberté entière après-midi. Achevé de\strut &  \\
      \leavevmode\fontsize{6.6pt}{7.5pt}\selectfont\strut m'installer pour le mieux ds ma nouvelle\strut & \smash{\parbox[t]{39mm}{\fontsize{6.2pt}{7.0pt}\selectfont\raggedright\color{black!88}\textbf{23 ds} \textit{abbr. pro} dans.}} \\
      \leavevmode\fontsize{6.6pt}{7.5pt}\selectfont\strut chambre. Y suis rentré aussitôt après dinner.\strut & \smash{\parbox[t]{39mm}{\fontsize{6.2pt}{7.0pt}\selectfont\raggedright\color{black!88}\textbf{24 dinner} \textit{angl. in ms.}}} \\
      \leavevmode\llap{\fontsize{6.4pt}{7.3pt}\selectfont\color{darkgray}25\hspace{1.8mm}}\fontsize{6.6pt}{7.5pt}\selectfont\strut M'y trouve très bien. 2 Tables pour écrire —\strut &  \\
\end{tabular}}%
\par\vspace*{\fill}
% --- Registres de bas de page (1. Critique interne [si pertinent] ; 2. Loci paralleli en latin ; 3. Commentaire historique en retrait) ---
\noindent\rule{\textwidth}{0.5pt}\par\vspace{0.65mm}
{\fontsize{6.6pt}{7.7pt}\selectfont%
\noindent\normalfont\color{black!92} \textbf{8} (\textit{Toutencourt}) cf. supra f\textsuperscript{o}~49~v\textsuperscript{o} (p.~\pageref{fol:99}), f\textsuperscript{o}~82~r\textsuperscript{o} (p.~\pageref{fol:164}), infra f\textsuperscript{o}~83~v\textsuperscript{o} (p.~\pageref{fol:167}), f\textsuperscript{o}~84~v\textsuperscript{o} (p.~\pageref{fol:169}), \textit{et passim}. — \textbf{10} (\textit{Raincheval}) cf. supra f\textsuperscript{o}~20~v\textsuperscript{o} (p.~\pageref{fol:41}), f\textsuperscript{o}~77~v\textsuperscript{o} (p.~\pageref{fol:155}), infra f\textsuperscript{o}~83~v\textsuperscript{o} (p.~\pageref{fol:167}), f\textsuperscript{o}~87~v\textsuperscript{o} (p.~\pageref{fol:175}), \textit{et passim}. — \textbf{13} (\textit{Major Levy}) cf. supra f\textsuperscript{o}~45~v\textsuperscript{o} (p.~\pageref{fol:91}), f\textsuperscript{o}~58~r\textsuperscript{o} (p.~\pageref{fol:116}).\par\vspace{0.55mm}
\noindent\rule{38mm}{0.4pt}\par\vspace{0.5mm}
\noindent\hspace*{5mm}\parbox{\dimexpr\textwidth-5mm\relax}{\fontsize{6.6pt}{7.7pt}\selectfont\color{black!84}Le dimanche 11 mars 1917, au bureau de la Mission française de la V\textsuperscript{e} Armée à \textbf{Toutencourt}, Hippolyte Combes apprend que, sur la demande expresse du \textbf{major Levy}, il est maintenu à la commission d'achats de l'armée (\textbf{\textit{Army Purchase Board}}) à \textbf{Hérissart}, où il s'installe dans le logement (\textit{billet} n\textsuperscript{o}~220) laissé libre par l'interprète Valet.}\par}
\newpage
\phantomsection\label{fac:166}
% --- En-tête du fac-similé (Page impaire / droite : Photographie du manuscrit) ---
\noindent\rule{\textwidth}{0.5pt}\par\vspace{0.55mm}
\noindent\textbf{\fontsize{6.8pt}{7.8pt}\selectfont FAC-SIMILÉ EN VIS-À-VIS — FOLIO 83 RECTO} \hfill \textbf{\fontsize{6.8pt}{7.8pt}\selectfont MARS 1917}\par\vspace{0.25mm}
\noindent{\fontsize{6.8pt}{8.0pt}\selectfont\textit{Livre-Journal d'Hippolyte Combes} — 11 mars 1917 (Hérissart, Toutencourt, Puchevillers (Somme))}\par\vspace{0.55mm}
\noindent\rule{\textwidth}{0.35pt}\par\vspace*{\fill}
{\centering
  \setlength{\fboxsep}{1.2mm}%
  \setlength{\fboxrule}{0.4pt}%
  \fcolorbox{gray!55}{white}{%
    \includegraphics[width=0.95\textwidth,height=0.81\textheight,keepaspectratio]{\facsimiledir/20251126_111240.jpg}%
  }\par}
\vspace*{\fill}
\newpage

% =============================================================================
% VIS-A-VIS SPREAD : FOLIO 83 VERSO (20251126_111251.jpg)
% =============================================================================
\cleardoubleevenpage
\renewcommand{\currentmonthyear}{Mars 1917}
\renewcommand{\currentfoliolabel}{Folio 83 verso}
\phantomsection\label{fol:167}
% --- En-tête du folio (Page paire / gauche : Transcription, apparat & registres de bas de page) ---
\noindent\rule{\textwidth}{0.5pt}\par\vspace{0.55mm}
\noindent\textbf{\footnotesize FOLIO 83 VERSO} \hfill \textbf{\footnotesize MARS 1917}\par\vspace{0.25mm}
\noindent{\fontsize{6.9pt}{8.0pt}\selectfont Bonheur de la nouvelle chambre (billet 220) et tournée en auto Wolseley}\par\vspace{0.2mm}
\noindent{\fontsize{6.9pt}{8.0pt}\selectfont\textit{11 (fin) – 12 mars 1917} \hfill \textit{Hérissart, Toutencourt, Contay, Beaumont, Molliens-au-Bois, Puchevillers, Raincheval, Candas (Somme)}}\par\vspace{0.5mm}
\noindent\rule{\textwidth}{0.35pt}\par\vspace{0.8mm}
% --- Corps de la transcription diplomatique et apparat marginal en prolongement ---
{\renewcommand{\arraystretch}{0.87}%
\noindent\hspace*{3.5mm}\begin{tabular}{@{}p{87mm}@{\hspace{3.5mm}}p{39mm}@{}}
      \leavevmode\fontsize{6.9pt}{7.9pt}\selectfont\strut \underline{Dim. 11 (fin)} \hfill{} \underline{Mars 1917}\strut & \smash{\parbox[t]{39mm}{\fontsize{6.2pt}{7.0pt}\selectfont\raggedright\color{black!88}\textbf{1 Dim.} \textit{scil.} Dimanche.}} \\
      \leavevmode\fontsize{6.9pt}{7.9pt}\selectfont\strut 2 fenêtres. Impressions de calme \& silence, favorables\strut &  \\
      \leavevmode\fontsize{6.9pt}{7.9pt}\selectfont\strut à rêverie \& étude —\strut &  \\
      \leavevmode\fontsize{6.9pt}{7.9pt}\selectfont\strut Matinée, vrai printemps. Soleil très chaud.\strut &  \\
      \leavevmode\llap{\fontsize{6.4pt}{7.3pt}\selectfont\color{darkgray}5\hspace{1.8mm}}\fontsize{6.9pt}{7.9pt}\selectfont\strut Après-midi, averses \& brume.\strut &  \\
      \leavevmode\fontsize{6.9pt}{7.9pt}\selectfont\strut [ A Hérissart, tommy m'a porté sacs : 50 c. ]\strut & \smash{\parbox[t]{39mm}{\fontsize{6.2pt}{7.0pt}\selectfont\raggedright\color{black!88}\textbf{6} \textit{in marg. cum unco} (comptes). \textbf{6 c.} \textit{abbr. pro} centimes. — \textbf{7 m.} \textit{abbr. pro} matin.}} \\
      \cline{1-1}
      \leavevmode\fontsize{6.9pt}{7.9pt}\selectfont\strut \underline{Lundi 12} — Levé 7h m. — Très bonne nuit.\strut &  \\
      \leavevmode\fontsize{6.9pt}{7.9pt}\selectfont\strut Achevé de m'installer nouvelle maison. Y serai bien.\strut &  \\
      \leavevmode\fontsize{6.9pt}{7.9pt}\selectfont\strut Tout le jour, temps brumeux \& par moments,\strut &  \\
      \leavevmode\llap{\fontsize{6.4pt}{7.3pt}\selectfont\color{darkgray}10\hspace{1.8mm}}\fontsize{6.9pt}{7.9pt}\selectfont\strut pluie fine.\strut &  \\
      \leavevmode\fontsize{6.9pt}{7.9pt}\selectfont\strut Matinée, courses à Toutencourt, Contay, Beaumont \&\strut &  \\
      \leavevmode\fontsize{6.9pt}{7.9pt}\selectfont\strut Molliens-au-Bois. Auto Wolseley, conduite par\strut &  \\
      \leavevmode\fontsize{6.9pt}{7.9pt}\selectfont\strut 1 australien, moins confortable que la Vauxhall\strut &  \\
      \leavevmode\fontsize{6.9pt}{7.9pt}\selectfont\strut de Jimmie.\strut &  \\
      \leavevmode\llap{\fontsize{6.4pt}{7.3pt}\selectfont\color{darkgray}15\hspace{1.8mm}}\fontsize{6.9pt}{7.9pt}\selectfont\strut Tout l'après-midi, à Puchevillers, puis\strut &  \\
      \leavevmode\fontsize{6.9pt}{7.9pt}\selectfont\strut Raincheval, Candas.\strut &  \\
      \cline{1-1}
      \leavevmode\fontsize{6.9pt}{7.9pt}\selectfont\strut [ Toutencourt : Exp. F. Canteen : 2 tins of\strut & \smash{\parbox[t]{39mm}{\fontsize{6.2pt}{7.0pt}\selectfont\raggedright\color{black!88}\textbf{17} \textit{in marg. cum unco} (comptes). \textbf{17 Canteen} \textit{angl. in ms.} — \textbf{18 f.} \textit{abbr. pro} franc(s). \textbf{18 c.} \textit{abbr. pro} centimes.}} \\
      \leavevmode\fontsize{6.9pt}{7.9pt}\selectfont\strut Calvert's Tooth Powder : 1 f. 75 c. ]\strut &  \\
\end{tabular}}%
\par\vspace*{\fill}
% --- Registres de bas de page (1. Critique interne [si pertinent] ; 2. Loci paralleli en latin ; 3. Commentaire historique en retrait) ---
\noindent\rule{\textwidth}{0.5pt}\par\vspace{0.65mm}
{\fontsize{6.6pt}{7.7pt}\selectfont%
\noindent\normalfont\color{black!92} \textbf{6} (\textit{Hérissart}) cf. supra f\textsuperscript{o}~47~v\textsuperscript{o} (p.~\pageref{fol:95}), f\textsuperscript{o}~83~r\textsuperscript{o} (p.~\pageref{fol:166}), infra f\textsuperscript{o}~84~r\textsuperscript{o} (p.~\pageref{fol:168}), f\textsuperscript{o}~84~v\textsuperscript{o} (p.~\pageref{fol:169}), \textit{et passim}. — \textbf{11} (\textit{Toutencourt}) cf. supra f\textsuperscript{o}~49~v\textsuperscript{o} (p.~\pageref{fol:99}), f\textsuperscript{o}~83~r\textsuperscript{o} (p.~\pageref{fol:166}), infra f\textsuperscript{o}~84~v\textsuperscript{o} (p.~\pageref{fol:169}), f\textsuperscript{o}~85~r\textsuperscript{o} (p.~\pageref{fol:170}), \textit{et passim}. — \textbf{11} (\textit{Contay}) cf. supra f\textsuperscript{o}~49~r\textsuperscript{o} (p.~\pageref{fol:98}), f\textsuperscript{o}~80~v\textsuperscript{o} (p.~\pageref{fol:161}), infra f\textsuperscript{o}~84~v\textsuperscript{o} (p.~\pageref{fol:169}), f\textsuperscript{o}~85~v\textsuperscript{o} (p.~\pageref{fol:171}), \textit{et passim}.\par\vspace{0.55mm}
\noindent\rule{38mm}{0.4pt}\par\vspace{0.5mm}
\noindent\hspace*{5mm}\parbox{\dimexpr\textwidth-5mm\relax}{\fontsize{6.6pt}{7.7pt}\selectfont\color{black!84}Reprise des tournées du \textit{Purchase Board} dès le lundi 12 mars 1917 (Toutencourt, Contay, Molliens-au-Bois, Puchevillers, Raincheval et \textbf{Candas}, Somme, canton de Bernaville).}\par}
\newpage
\phantomsection\label{fac:167}
% --- En-tête du fac-similé (Page impaire / droite : Photographie du manuscrit) ---
\noindent\rule{\textwidth}{0.5pt}\par\vspace{0.55mm}
\noindent\textbf{\fontsize{6.8pt}{7.8pt}\selectfont FAC-SIMILÉ EN VIS-À-VIS — FOLIO 83 VERSO} \hfill \textbf{\fontsize{6.8pt}{7.8pt}\selectfont MARS 1917}\par\vspace{0.25mm}
\noindent{\fontsize{6.8pt}{8.0pt}\selectfont\textit{Livre-Journal d'Hippolyte Combes} — 11 (fin) – 12 mars 1917 (Hérissart, Toutencourt, Contay, Beaumont, Molliens-au-Bois, Puchevillers, Raincheval, Candas (Somme))}\par\vspace{0.55mm}
\noindent\rule{\textwidth}{0.35pt}\par\vspace*{\fill}
{\centering
  \setlength{\fboxsep}{1.2mm}%
  \setlength{\fboxrule}{0.4pt}%
  \fcolorbox{gray!55}{white}{%
    \includegraphics[width=0.95\textwidth,height=0.81\textheight,keepaspectratio]{\facsimiledir/20251126_111251.jpg}%
  }\par}
\vspace*{\fill}
\newpage

% =============================================================================
% VIS-A-VIS SPREAD : FOLIO 84 RECTO (20251126_111309.jpg)
% =============================================================================
\cleardoubleevenpage
\renewcommand{\currentmonthyear}{Mars 1917}
\renewcommand{\currentfoliolabel}{Folio 84 recto}
\phantomsection\label{fol:168}
% --- En-tête du folio (Page paire / gauche : Transcription, apparat & registres de bas de page) ---
\noindent\rule{\textwidth}{0.5pt}\par\vspace{0.55mm}
\noindent\textbf{\footnotesize FOLIO 84 RECTO} \hfill \textbf{\footnotesize MARS 1917}\par\vspace{0.25mm}
\noindent{\fontsize{6.9pt}{8.0pt}\selectfont Visite à l'instituteur M. Flutre et déménagement du mess au billet 216}\par\vspace{0.2mm}
\noindent{\fontsize{6.9pt}{8.0pt}\selectfont\textit{13 – 15 mars 1917} \hfill \textit{Hérissart, Le Mesnil-Martinsart, Puchevillers, Acheux, Bertrancourt (Somme)}}\par\vspace{0.5mm}
\noindent\rule{\textwidth}{0.35pt}\par\vspace{0.8mm}
% --- Corps de la transcription diplomatique et apparat marginal en prolongement ---
{\renewcommand{\arraystretch}{0.87}%
\noindent\hspace*{3.5mm}\begin{tabular}{@{}p{87mm}@{\hspace{3.5mm}}p{39mm}@{}}
      \leavevmode\fontsize{6.9pt}{7.9pt}\selectfont\strut \underline{Mardi 13} — \hfill{} \underline{Mars 1917 —}\strut &  \\
      \leavevmode\fontsize{6.9pt}{7.9pt}\selectfont\strut Température douce. Tout le jour, ciel nuageux.\strut &  \\
      \leavevmode\fontsize{6.9pt}{7.9pt}\selectfont\strut Soleil brille très peu. Routes, horriblement boueuses.\strut &  \\
      \leavevmode\fontsize{6.9pt}{7.9pt}\selectfont\strut On patauge en pleine fange —\strut &  \\
      \leavevmode\llap{\fontsize{6.4pt}{7.3pt}\selectfont\color{darkgray}5\hspace{1.8mm}}\fontsize{6.9pt}{7.9pt}\selectfont\strut Mon lt, dû se rendre à Martinsart. Resté seul à\strut &  \\
      \leavevmode\fontsize{6.9pt}{7.9pt}\selectfont\strut Hérissart. Liberté entière.\strut &  \\
      \leavevmode\fontsize{6.9pt}{7.9pt}\selectfont\strut Après tea, visite à M. \& Mme Flutre, instituteurs,\strut & \smash{\parbox[t]{39mm}{\fontsize{6.2pt}{7.0pt}\selectfont\raggedright\color{black!88}\textbf{7 Mme} \textit{abbr. pro} Madame. \textbf{7 tea} \textit{angl. in ms.} — \textbf{8 Normle Supre} \textit{abbr. pro} Normale Supérieure.}} \\
      \leavevmode\fontsize{6.9pt}{7.9pt}\selectfont\strut au Mesnil. Ont le fils, élève de Normle Supre.\strut &  \\
      \cline{1-1}
      \leavevmode\fontsize{6.9pt}{7.9pt}\selectfont\strut \underline{Mercredi 14.} —\strut &  \\
      \leavevmode\llap{\fontsize{6.4pt}{7.3pt}\selectfont\color{darkgray}10\hspace{1.8mm}}\fontsize{6.9pt}{7.9pt}\selectfont\strut Matinée, pluie fine persistante.\strut &  \\
      \leavevmode\fontsize{6.9pt}{7.9pt}\selectfont\strut Rues \& routes détrempées, boue épaisse.\strut &  \\
      \leavevmode\fontsize{6.9pt}{7.9pt}\selectfont\strut Déménagé office et mess au billet 216,\strut & \smash{\parbox[t]{39mm}{\fontsize{6.2pt}{7.0pt}\selectfont\raggedright\color{black!88}\textbf{12 mess} \textit{angl. in ms.}}} \\
      \leavevmode\fontsize{6.9pt}{7.9pt}\selectfont\strut rue, tout en face de ma maison\strut &  \\
      \leavevmode\fontsize{6.9pt}{7.9pt}\selectfont\strut où j'ai ma chambre. Ns sommes 1000 fois\strut &  \\
      \leavevmode\llap{\fontsize{6.4pt}{7.3pt}\selectfont\color{darkgray}15\hspace{1.8mm}}\fontsize{6.9pt}{7.9pt}\selectfont\strut mieux dans notre nouvelle installation.\strut &  \\
      \leavevmode\fontsize{6.9pt}{7.9pt}\selectfont\strut Après-midi, course à Puchevillers; allé\strut &  \\
      \leavevmode\fontsize{6.9pt}{7.9pt}\selectfont\strut de ferme en ferme, à propos de paille.\strut &  \\
      \leavevmode\fontsize{6.9pt}{7.9pt}\selectfont\strut Fin jour, soleil brille un peu et dissipe\strut &  \\
      \leavevmode\fontsize{6.9pt}{7.9pt}\selectfont\strut nuages.\strut &  \\
      \cline{1-1}
      \leavevmode\llap{\fontsize{6.4pt}{7.3pt}\selectfont\color{darkgray}20\hspace{1.8mm}}\fontsize{6.9pt}{7.9pt}\selectfont\strut \underline{— Jeudi 15.} — Froid piquant matin \& fin journée.\strut &  \\
      \leavevmode\fontsize{6.9pt}{7.9pt}\selectfont\strut 1/2 jour, soleil brille 1 peu, mais vent glacial.\strut & \smash{\parbox[t]{39mm}{\fontsize{6.2pt}{7.0pt}\selectfont\raggedright\color{black!88}\textbf{21 1/2} \textit{scil.} demi / moitié.}} \\
      \leavevmode\fontsize{6.9pt}{7.9pt}\selectfont\strut Matin \& après-midi, courses à Acheux,\strut &  \\
      \leavevmode\fontsize{6.9pt}{7.9pt}\selectfont\strut \& Bertrancourt, etc. etc. A Acheux Wood,\strut &  \\
      \leavevmode\fontsize{6.9pt}{7.9pt}\selectfont\strut pataugé boue épaisse.\strut &  \\
\end{tabular}}%
\par\vspace*{\fill}
% --- Registres de bas de page (1. Critique interne [si pertinent] ; 2. Loci paralleli en latin ; 3. Commentaire historique en retrait) ---
\noindent\rule{\textwidth}{0.5pt}\par\vspace{0.65mm}
{\fontsize{6.6pt}{7.7pt}\selectfont%
\noindent\normalfont\color{black!92} \textbf{5} (\textit{Martinsart}) cf. supra f\textsuperscript{o}~52~r\textsuperscript{o} (p.~\pageref{fol:104}), f\textsuperscript{o}~79~v\textsuperscript{o} (p.~\pageref{fol:159}), \textit{et passim}. — \textbf{6} (\textit{Hérissart}) cf. supra f\textsuperscript{o}~47~v\textsuperscript{o} (p.~\pageref{fol:95}), f\textsuperscript{o}~83~v\textsuperscript{o} (p.~\pageref{fol:167}), infra f\textsuperscript{o}~84~v\textsuperscript{o} (p.~\pageref{fol:169}), f\textsuperscript{o}~85~r\textsuperscript{o} (p.~\pageref{fol:170}), \textit{et passim}. — \textbf{16} (\textit{Puchevillers}) cf. supra f\textsuperscript{o}~37~v\textsuperscript{o} (p.~\pageref{fol:75}), f\textsuperscript{o}~83~v\textsuperscript{o} (p.~\pageref{fol:167}), infra f\textsuperscript{o}~84~v\textsuperscript{o} (p.~\pageref{fol:169}), f\textsuperscript{o}~85~r\textsuperscript{o} (p.~\pageref{fol:170}), \textit{et passim}.\par\vspace{0.55mm}
\noindent\rule{38mm}{0.4pt}\par\vspace{0.5mm}
\noindent\hspace*{5mm}\parbox{\dimexpr\textwidth-5mm\relax}{\fontsize{6.6pt}{7.7pt}\selectfont\color{black!84}Au \textbf{Mesnil-Martinsart} (Somme), Hippolyte Combes rend visite le 13 mars 1917 à son collègue de l'Instruction publique \textbf{Fortuné Flutre}, instituteur de la commune (auteur en 1899 d'une notice historique sur Le Mesnil-Martinsart), et à son épouse, parents du romaniste \textbf{Louis-Fernand Flutre} (1892–1978), admis à l'\textbf{École normale supérieure} en 1914 (entré en 1919 après la guerre, agrégé de grammaire en 1921, professeur à la Faculté des lettres de Lyon et auteur en 1955 de deux monographies historiques et linguistiques sur Le Mesnil-Martinsart ; R.~Debrie, \textit{Revue de linguistique romane}, 42, 1978, p.~238–240).}\par}
\newpage
\phantomsection\label{fac:168}
% --- En-tête du fac-similé (Page impaire / droite : Photographie du manuscrit) ---
\noindent\rule{\textwidth}{0.5pt}\par\vspace{0.55mm}
\noindent\textbf{\fontsize{6.8pt}{7.8pt}\selectfont FAC-SIMILÉ EN VIS-À-VIS — FOLIO 84 RECTO} \hfill \textbf{\fontsize{6.8pt}{7.8pt}\selectfont MARS 1917}\par\vspace{0.25mm}
\noindent{\fontsize{6.8pt}{8.0pt}\selectfont\textit{Livre-Journal d'Hippolyte Combes} — 13 – 15 mars 1917 (Hérissart, Le Mesnil-Martinsart, Puchevillers, Acheux, Bertrancourt (Somme))}\par\vspace{0.55mm}
\noindent\rule{\textwidth}{0.35pt}\par\vspace*{\fill}
{\centering
  \setlength{\fboxsep}{1.2mm}%
  \setlength{\fboxrule}{0.4pt}%
  \fcolorbox{gray!55}{white}{%
    \includegraphics[width=0.95\textwidth,height=0.81\textheight,keepaspectratio]{\facsimiledir/20251126_111309.jpg}%
  }\par}
\vspace*{\fill}
\newpage

% =============================================================================
% VIS-A-VIS SPREAD : FOLIO 84 VERSO (20251126_111321.jpg)
% =============================================================================
\cleardoubleevenpage
\renewcommand{\currentmonthyear}{Mars 1917}
\renewcommand{\currentfoliolabel}{Folio 84 verso}
\phantomsection\label{fol:169}
% --- En-tête du folio (Page paire / gauche : Transcription, apparat & registres de bas de page) ---
\noindent\rule{\textwidth}{0.5pt}\par\vspace{0.55mm}
\noindent\textbf{\footnotesize FOLIO 84 VERSO} \hfill \textbf{\footnotesize MARS 1917}\par\vspace{0.25mm}
\noindent{\fontsize{6.9pt}{8.0pt}\selectfont Retour à Albert et Ribemont et promenades printanières avec le lieutenant Wentworth}\par\vspace{0.2mm}
\noindent{\fontsize{6.9pt}{8.0pt}\selectfont\textit{16 – 17 mars 1917} \hfill \textit{Hérissart, Mirvaux, Pierregot, Contay, Warloy, Senlis, Albert, Ribemont, Toutencourt, Puchevillers (Somme)}}\par\vspace{0.5mm}
\noindent\rule{\textwidth}{0.35pt}\par\vspace{0.8mm}
% --- Corps de la transcription diplomatique et apparat marginal en prolongement ---
{\renewcommand{\arraystretch}{0.87}%
\noindent\hspace*{3.5mm}\begin{tabular}{@{}p{87mm}@{\hspace{3.5mm}}p{39mm}@{}}
      \leavevmode\fontsize{6.9pt}{7.9pt}\selectfont\strut \underline{Ven. 16} \hfill{} \underline{— Mars 1917 —}\strut & \smash{\parbox[t]{39mm}{\fontsize{6.2pt}{7.0pt}\selectfont\raggedright\color{black!88}\textbf{1 Ven.} \textit{scil.} Vendredi.}} \\
      \leavevmode\fontsize{6.9pt}{7.9pt}\selectfont\strut Nuit très froide. Gelé dur. Au réveil trouvé boue séchée.\strut &  \\
      \leavevmode\fontsize{6.9pt}{7.9pt}\selectfont\strut Matin, courses à Mirvaux \& Pierregot.\strut &  \\
      \leavevmode\fontsize{6.9pt}{7.9pt}\selectfont\strut Early lunch terminé, venus à Albert, par\strut & \smash{\parbox[t]{39mm}{\fontsize{6.2pt}{7.0pt}\selectfont\raggedright\color{black!88}\textbf{4 lunch} \textit{angl. in ms.}}} \\
      \leavevmode\llap{\fontsize{6.4pt}{7.3pt}\selectfont\color{darkgray}5\hspace{1.8mm}}\fontsize{6.9pt}{7.9pt}\selectfont\strut Contay-Warloy-Senlis. Revu avec intérêt ville\strut &  \\
      \leavevmode\fontsize{6.9pt}{7.9pt}\selectfont\strut d'Albert : maisons atteintes par obus, basilique\strut &  \\
      \leavevmode\fontsize{6.9pt}{7.9pt}\selectfont\strut etc. etc. \& troupes australiennes cantonnées.\strut &  \\
      \leavevmode\fontsize{6.9pt}{7.9pt}\selectfont\strut Au retour, long arrêt à Ribemont, pour\strut &  \\
      \leavevmode\fontsize{6.9pt}{7.9pt}\selectfont\strut acheter fourrage. Heureux de revoir ce village\strut &  \\
      \leavevmode\llap{\fontsize{6.4pt}{7.3pt}\selectfont\color{darkgray}10\hspace{1.8mm}}\fontsize{6.9pt}{7.9pt}\selectfont\strut où laissé connaissances très sympathiques.\strut &  \\
      \leavevmode\fontsize{6.9pt}{7.9pt}\selectfont\strut [ Blacksmith d'Hérissart : répar. à ferrure sac : 30 c. ]\strut & \smash{\parbox[t]{39mm}{\fontsize{6.2pt}{7.0pt}\selectfont\raggedright\color{black!88}\textbf{11} \textit{in marg. cum unco} (comptes). \textbf{11 c.} \textit{abbr. pro} centimes. — \textbf{12 Sam.} \textit{scil.} Samedi.}} \\
      \cline{1-1}
      \leavevmode\fontsize{6.9pt}{7.9pt}\selectfont\strut \underline{Sam. 17} —\strut &  \\
      \leavevmode\fontsize{6.9pt}{7.9pt}\selectfont\strut Très belle journée. Soleil printanier.\strut &  \\
      \leavevmode\fontsize{6.9pt}{7.9pt}\selectfont\strut Matinée, venu à pied \& retour id. Toutencourt.\strut &  \\
      \leavevmode\llap{\fontsize{6.4pt}{7.3pt}\selectfont\color{darkgray}15\hspace{1.8mm}}\fontsize{6.9pt}{7.9pt}\selectfont\strut Après-midi, promenade à pied à Puchevillers\strut &  \\
      \leavevmode\fontsize{6.9pt}{7.9pt}\selectfont\strut \& retour de =, avec Lt. Wentworth. Air si\strut & \smash{\parbox[t]{39mm}{\fontsize{6.2pt}{7.0pt}\selectfont\raggedright\color{black!88}\textbf{16 Lt} \textit{abbr. pro} Lieutenant. \textbf{16 de =} \textit{signum aequalitatis pro} de même. — \textbf{17 bcp} \textit{abbr. pro} beaucoup. — \textbf{18 workshops} \textit{angl. in ms.}}} \\
      \leavevmode\fontsize{6.9pt}{7.9pt}\selectfont\strut chaud que sué bcp. Marche m'a fait bien.\strut &  \\
      \leavevmode\fontsize{6.9pt}{7.9pt}\selectfont\strut Auto est aux workshops.\strut &  \\
      \leavevmode\fontsize{6.9pt}{7.9pt}\selectfont\strut Troupes cavalerie indigène \& anglaise\strut &  \\
      \leavevmode\llap{\fontsize{6.4pt}{7.3pt}\selectfont\color{darkgray}20\hspace{1.8mm}}\fontsize{6.9pt}{7.9pt}\selectfont\strut traversé village au trot — — —\strut &  \\
      \cline{1-1}
      \leavevmode\fontsize{6.9pt}{7.9pt}\selectfont\strut [ A Toutencourt : c. ples : 1 f. + 15 c.\strut & \smash{\parbox[t]{39mm}{\fontsize{6.2pt}{7.0pt}\selectfont\raggedright\color{black!88}\textbf{21} \textit{in marg. cum unco} (comptes). \textbf{21 c. ples} \textit{abbr. pro} cartes postales.}} \\
      \leavevmode\fontsize{6.9pt}{7.9pt}\selectfont\strut et à soie ]\strut &  \\
\end{tabular}}%
\par\vspace*{\fill}
% --- Registres de bas de page (1. Critique interne [si pertinent] ; 2. Loci paralleli en latin ; 3. Commentaire historique en retrait) ---
\noindent\rule{\textwidth}{0.5pt}\par\vspace{0.65mm}
{\fontsize{6.6pt}{7.7pt}\selectfont%
\noindent\normalfont\color{black!92} \textbf{3} (\textit{Pierregot}) cf. supra f\textsuperscript{o}~52~r\textsuperscript{o} (p.~\pageref{fol:104}), f\textsuperscript{o}~78~v\textsuperscript{o} (p.~\pageref{fol:157}), infra f\textsuperscript{o}~85~v\textsuperscript{o} (p.~\pageref{fol:171}), f\textsuperscript{o}~87~r\textsuperscript{o} (p.~\pageref{fol:174}), \textit{et passim}. — \textbf{3} (\textit{Mirvaux}) cf. supra f\textsuperscript{o}~79~v\textsuperscript{o} (p.~\pageref{fol:159}), infra f\textsuperscript{o}~85~v\textsuperscript{o} (p.~\pageref{fol:171}), f\textsuperscript{o}~91~v\textsuperscript{o} (p.~\pageref{fol:183}). — \textbf{4} (\textit{Albert}) cf. supra f\textsuperscript{o}~39~r\textsuperscript{o} (p.~\pageref{fol:78}), f\textsuperscript{o}~82~r\textsuperscript{o} (p.~\pageref{fol:164}), infra f\textsuperscript{o}~85~v\textsuperscript{o} (p.~\pageref{fol:171}), f\textsuperscript{o}~90~v\textsuperscript{o} (p.~\pageref{fol:181}), \textit{et passim}.\par\vspace{0.55mm}
\noindent\rule{38mm}{0.4pt}\par\vspace{0.5mm}
\noindent\hspace*{5mm}\parbox{\dimexpr\textwidth-5mm\relax}{\fontsize{6.6pt}{7.7pt}\selectfont\color{black!84}Le vendredi 16 mars 1917, Hippolyte Combes revoit \textbf{Albert} (où cantonnent les troupes australiennes de l'\textit{I~Anzac Corps}) et \textbf{Ribemont-sur-Ancre}.}\par}
\newpage
\phantomsection\label{fac:169}
% --- En-tête du fac-similé (Page impaire / droite : Photographie du manuscrit) ---
\noindent\rule{\textwidth}{0.5pt}\par\vspace{0.55mm}
\noindent\textbf{\fontsize{6.8pt}{7.8pt}\selectfont FAC-SIMILÉ EN VIS-À-VIS — FOLIO 84 VERSO} \hfill \textbf{\fontsize{6.8pt}{7.8pt}\selectfont MARS 1917}\par\vspace{0.25mm}
\noindent{\fontsize{6.8pt}{8.0pt}\selectfont\textit{Livre-Journal d'Hippolyte Combes} — 16 – 17 mars 1917 (Hérissart, Mirvaux, Pierregot, Contay, Warloy, Senlis, Albert, Ribemont, Toutencourt, Puchevillers (Somme))}\par\vspace{0.55mm}
\noindent\rule{\textwidth}{0.35pt}\par\vspace*{\fill}
{\centering
  \setlength{\fboxsep}{1.2mm}%
  \setlength{\fboxrule}{0.4pt}%
  \fcolorbox{gray!55}{white}{%
    \includegraphics[width=0.95\textwidth,height=0.81\textheight,keepaspectratio]{\facsimiledir/20251126_111321.jpg}%
  }\par}
\vspace*{\fill}
\newpage

% =============================================================================
% VIS-A-VIS SPREAD : FOLIO 85 RECTO (20251126_111359.jpg)
% =============================================================================
\cleardoubleevenpage
\renewcommand{\currentmonthyear}{Mars 1917}
\renewcommand{\currentfoliolabel}{Folio 85 recto}
\phantomsection\label{fol:170}
% --- En-tête du folio (Page paire / gauche : Transcription, apparat & registres de bas de page) ---
\noindent\rule{\textwidth}{0.5pt}\par\vspace{0.55mm}
\noindent\textbf{\footnotesize FOLIO 85 RECTO} \hfill \textbf{\footnotesize MARS 1917}\par\vspace{0.25mm}
\noindent{\fontsize{6.9pt}{8.0pt}\selectfont Prône du vieux curé de Toutencourt et passage de la 63e Royal Naval Division}\par\vspace{0.2mm}
\noindent{\fontsize{6.9pt}{8.0pt}\selectfont\textit{18 – 19 mars 1917} \hfill \textit{Hérissart, Toutencourt, Puchevillers (Somme)}}\par\vspace{0.5mm}
\noindent\rule{\textwidth}{0.35pt}\par\vspace{0.8mm}
% --- Corps de la transcription diplomatique et apparat marginal en prolongement ---
{\renewcommand{\arraystretch}{0.83}%
\noindent\hspace*{3.5mm}\begin{tabular}{@{}p{87mm}@{\hspace{3.5mm}}p{39mm}@{}}
      \leavevmode\fontsize{6.6pt}{7.5pt}\selectfont\strut \underline{Dim. 18.} — \hfill{} \underline{— Mars 1917 —}\strut & \smash{\parbox[t]{39mm}{\fontsize{6.2pt}{7.0pt}\selectfont\raggedright\color{black!88}\textbf{1 Dim.} \textit{scil.} Dimanche.}} \\
      \leavevmode\fontsize{6.6pt}{7.5pt}\selectfont\strut Matin, ciel nuageux, temps pluvieux. Après-\strut &  \\
      \leavevmode\fontsize{6.6pt}{7.5pt}\selectfont\strut midi, soleil ardent.\strut &  \\
      \leavevmode\fontsize{6.6pt}{7.5pt}\selectfont\strut Matinée, venu à pied à Toutencourt,\strut &  \\
      \leavevmode\llap{\fontsize{6.4pt}{7.3pt}\selectfont\color{darkgray}5\hspace{1.8mm}}\fontsize{6.6pt}{7.5pt}\selectfont\strut Mission fr. Reposé église du village, q. mnts\strut & \smash{\parbox[t]{39mm}{\fontsize{6.2pt}{7.0pt}\selectfont\raggedright\color{black!88}\textbf{5 mnts} \textit{abbr. pro} moments / minutes.}} \\
      \leavevmode\fontsize{6.6pt}{7.5pt}\selectfont\strut \& entendu prône vieux curé —\strut &  \\
      \leavevmode\fontsize{6.6pt}{7.5pt}\selectfont\strut Après-midi, accompagné Lt Wentworth\strut & \smash{\parbox[t]{39mm}{\fontsize{6.2pt}{7.0pt}\selectfont\raggedright\color{black!88}\textbf{7 Lt} \textit{abbr. pro} Lieutenant.}} \\
      \leavevmode\fontsize{6.6pt}{7.5pt}\selectfont\strut à Puchevillers, aller \& retour à pied, notre\strut &  \\
      \leavevmode\fontsize{6.6pt}{7.5pt}\selectfont\strut auto tjours aux workshops, en réparation.\strut & \smash{\parbox[t]{39mm}{\fontsize{6.2pt}{7.0pt}\selectfont\raggedright\color{black!88}\textbf{9 workshops} \textit{angl. in ms.}}} \\
      \leavevmode\llap{\fontsize{6.4pt}{7.3pt}\selectfont\color{darkgray}10\hspace{1.8mm}}\fontsize{6.6pt}{7.5pt}\selectfont\strut Longuement regardé vue panoramique,\strut &  \\
      \leavevmode\fontsize{6.6pt}{7.5pt}\selectfont\strut champs labourés ou ensemencés \&\strut &  \\
      \leavevmode\fontsize{6.6pt}{7.5pt}\selectfont\strut l'horizon lointain, légères ondulations de\strut &  \\
      \leavevmode\fontsize{6.6pt}{7.5pt}\selectfont\strut collines —\strut &  \\
      \leavevmode\fontsize{6.6pt}{7.5pt}\selectfont\strut Ces courses, matin \& soir, m'ont fatigué 1 peu,\strut &  \\
      \leavevmode\llap{\fontsize{6.4pt}{7.3pt}\selectfont\color{darkgray}15\hspace{1.8mm}}\fontsize{6.6pt}{7.5pt}\selectfont\strut mais fait grand bien.\strut &  \\
      \leavevmode\fontsize{6.6pt}{7.5pt}\selectfont\strut [ A Toutencourt : Solde : 82 f. —\strut & \smash{\parbox[t]{39mm}{\fontsize{6.2pt}{7.0pt}\selectfont\raggedright\color{black!88}\textbf{16} \textit{in marg. cum unco} (comptes). \textbf{16 f.} \textit{abbr. pro} franc(s). — \textbf{17 Cartes ples} \textit{abbr. pro} Cartes postales. \textbf{17 f.} \textit{abbr. pro} franc(s). — \textbf{20 off.} \textit{abbr. pro} officier(s). — \textbf{21} \textit{lect. incert. in ms.} — \textbf{23 off.} \textit{abbr. pro} officier(s). — \textbf{24 div} \textit{abbr. pro} division. — \textbf{25 Pt} \textit{abbr. pro} Petit. \textbf{25 Achiet-le-Pt.} \textit{scil.} Achiet-le-Petit. — \textbf{26 renseigts} \textit{abbr. pro} renseignements. \textbf{26 off.} \textit{abbr. pro} officier(s).}} \\
      \leavevmode\fontsize{6.6pt}{7.5pt}\selectfont\strut Cartes ples : 2 f. — Washing : 80 c. ]\strut &  \\
      \cline{1-1}
      \leavevmode\fontsize{6.6pt}{7.5pt}\selectfont\strut \underline{19 Lundi.} — Tout le jour, temps couvert, pluie\strut &  \\
      \leavevmode\fontsize{6.6pt}{7.5pt}\selectfont\strut menace. Gros nuages. — Matinée, entière libre.\strut &  \\
      \leavevmode\llap{\fontsize{6.4pt}{7.3pt}\selectfont\color{darkgray}20\hspace{1.8mm}}\fontsize{6.6pt}{7.5pt}\selectfont\strut Après-midi, venu à pied avec off. à Puchevillers.\strut &  \\
      \leavevmode\fontsize{6.6pt}{7.5pt}\selectfont\strut Suivi même route qu'à l'ord[inaire ?], à travers champs\strut &  \\
      \leavevmode\fontsize{6.6pt}{7.5pt}\selectfont\strut et sentiers —\strut &  \\
      \leavevmode\fontsize{6.6pt}{7.5pt}\selectfont\strut Dû laisser occuper 1 des 2 lits par off. 63rd Div.,\strut &  \\
      \leavevmode\fontsize{6.6pt}{7.5pt}\selectfont\strut naval div., de passage. Son régt revient des\strut &  \\
      \leavevmode\llap{\fontsize{6.4pt}{7.3pt}\selectfont\color{darkgray}25\hspace{1.8mm}}\fontsize{6.6pt}{7.5pt}\selectfont\strut tranchées, Miraumont, Achiet-le-Pt. Fait parler ces\strut &  \\
      \leavevmode\fontsize{6.6pt}{7.5pt}\selectfont\strut 2 off., qui m'ont donné renseigts très intéressants.\strut &  \\
\end{tabular}}%
\par\vspace*{\fill}
% --- Registres de bas de page (1. Critique interne [si pertinent] ; 2. Loci paralleli en latin ; 3. Commentaire historique en retrait) ---
\noindent\rule{\textwidth}{0.5pt}\par\vspace{0.65mm}
{\fontsize{6.6pt}{7.7pt}\selectfont%
\noindent\normalfont\color{black!92} \textbf{4} (\textit{Toutencourt}) cf. supra f\textsuperscript{o}~49~v\textsuperscript{o} (p.~\pageref{fol:99}), f\textsuperscript{o}~84~v\textsuperscript{o} (p.~\pageref{fol:169}), infra f\textsuperscript{o}~87~r\textsuperscript{o} (p.~\pageref{fol:174}), f\textsuperscript{o}~87~v\textsuperscript{o} (p.~\pageref{fol:175}), \textit{et passim}. — \textbf{8} (\textit{Puchevillers}) cf. supra f\textsuperscript{o}~37~v\textsuperscript{o} (p.~\pageref{fol:75}), f\textsuperscript{o}~84~v\textsuperscript{o} (p.~\pageref{fol:169}), infra f\textsuperscript{o}~85~v\textsuperscript{o} (p.~\pageref{fol:171}), f\textsuperscript{o}~86~r\textsuperscript{o} (p.~\pageref{fol:172}), \textit{et passim}. — \textbf{23} (\textit{63rd Div}) cf. infra f\textsuperscript{o}~85~v\textsuperscript{o} (p.~\pageref{fol:171}).\par\vspace{0.55mm}
\noindent\rule{38mm}{0.4pt}\par\vspace{0.5mm}
\noindent\hspace*{5mm}\parbox{\dimexpr\textwidth-5mm\relax}{\fontsize{6.6pt}{7.7pt}\selectfont\color{black!84}Dans la nuit du 19 au 20 mars 1917, Hippolyte Combes cède le lit de sa chambre d'Hérissart à un jeune lieutenant britannique de la \textbf{63\textsuperscript{e} Division (\textit{63rd Royal Naval Division})}, dont le bataillon vient d'être relevé des tranchées boueuses de \textbf{Miraumont} et d'\textbf{Achiet-le-Petit} au début du repli stratégique allemand sur la ligne Hindenburg (opération \textit{Alberich}).}\par}
\newpage
\phantomsection\label{fac:170}
% --- En-tête du fac-similé (Page impaire / droite : Photographie du manuscrit) ---
\noindent\rule{\textwidth}{0.5pt}\par\vspace{0.55mm}
\noindent\textbf{\fontsize{6.8pt}{7.8pt}\selectfont FAC-SIMILÉ EN VIS-À-VIS — FOLIO 85 RECTO} \hfill \textbf{\fontsize{6.8pt}{7.8pt}\selectfont MARS 1917}\par\vspace{0.25mm}
\noindent{\fontsize{6.8pt}{8.0pt}\selectfont\textit{Livre-Journal d'Hippolyte Combes} — 18 – 19 mars 1917 (Hérissart, Toutencourt, Puchevillers (Somme))}\par\vspace{0.55mm}
\noindent\rule{\textwidth}{0.35pt}\par\vspace*{\fill}
{\centering
  \setlength{\fboxsep}{1.2mm}%
  \setlength{\fboxrule}{0.4pt}%
  \fcolorbox{gray!55}{white}{%
    \includegraphics[width=0.95\textwidth,height=0.81\textheight,keepaspectratio]{\facsimiledir/20251126_111359.jpg}%
  }\par}
\vspace*{\fill}
\newpage

% =============================================================================
% VIS-A-VIS SPREAD : FOLIO 85 VERSO (20251126_111410.jpg)
% =============================================================================
\cleardoubleevenpage
\renewcommand{\currentmonthyear}{Mars 1917}
\renewcommand{\currentfoliolabel}{Folio 85 verso}
\phantomsection\label{fol:171}
% --- En-tête du folio (Page paire / gauche : Transcription, apparat & registres de bas de page) ---
\noindent\rule{\textwidth}{0.5pt}\par\vspace{0.55mm}
\noindent\textbf{\footnotesize FOLIO 85 VERSO} \hfill \textbf{\footnotesize MARS 1917}\par\vspace{0.25mm}
\noindent{\fontsize{6.9pt}{8.0pt}\selectfont Départ de la Naval Division, ferme de Belle-Église et défilé des 17th Lancers (Death or Glory Boys)}\par\vspace{0.2mm}
\noindent{\fontsize{6.9pt}{8.0pt}\selectfont\textit{20 – 22 mars 1917} \hfill \textit{Hérissart, Puchevillers, Arquèves (Belle-Église), Contay, Warloy, Senlis, Pierregot, Molliens, Albert (Somme)}}\par\vspace{0.5mm}
\noindent\rule{\textwidth}{0.35pt}\par\vspace{0.8mm}
% --- Corps de la transcription diplomatique et apparat marginal en prolongement ---
{\renewcommand{\arraystretch}{0.87}%
\noindent\hspace*{3.5mm}\begin{tabular}{@{}p{87mm}@{\hspace{3.5mm}}p{39mm}@{}}
      \leavevmode\fontsize{6.9pt}{7.9pt}\selectfont\strut \underline{20 Mardi} — \hfill{} \underline{Mars 1917 —}\strut &  \\
      \leavevmode\fontsize{6.9pt}{7.9pt}\selectfont\strut Lieut. 63rd Div., couché ds lit\strut & \smash{\parbox[t]{39mm}{\fontsize{6.2pt}{7.0pt}\selectfont\raggedright\color{black!88}\textbf{2 ds} \textit{abbr. pro} dans.}} \\
      \leavevmode\fontsize{6.9pt}{7.9pt}\selectfont\strut ma chambre. Parti avec son unité 8h matin.\strut &  \\
      \leavevmode\fontsize{6.9pt}{7.9pt}\selectfont\strut Toute la nuit, vent soufflé en rafales extrt violentes.\strut & \smash{\parbox[t]{39mm}{\fontsize{6.2pt}{7.0pt}\selectfont\raggedright\color{black!88}\textbf{4 extrt} \textit{abbr. pro} extrêmement.}} \\
      \leavevmode\llap{\fontsize{6.4pt}{7.3pt}\selectfont\color{darkgray}5\hspace{1.8mm}}\fontsize{6.9pt}{7.9pt}\selectfont\strut Durant matinée, neige, tombée à gros flocons.\strut &  \\
      \leavevmode\fontsize{6.9pt}{7.9pt}\selectfont\strut Liberté entière, journée.\strut &  \\
      \cline{1-1}
      \leavevmode\fontsize{6.9pt}{7.9pt}\selectfont\strut \underline{21 Merc.}\strut & \smash{\parbox[t]{39mm}{\fontsize{6.2pt}{7.0pt}\selectfont\raggedright\color{black!88}\textbf{7 Merc.} \textit{scil.} Mercredi.}} \\
      \leavevmode\fontsize{6.9pt}{7.9pt}\selectfont\strut Nuit très froide. Tout le jour, ciel couvert,\strut &  \\
      \leavevmode\fontsize{6.9pt}{7.9pt}\selectfont\strut pluie fine \& neige, par intermittence —\strut &  \\
      \leavevmode\llap{\fontsize{6.4pt}{7.3pt}\selectfont\color{darkgray}10\hspace{1.8mm}}\fontsize{6.9pt}{7.9pt}\selectfont\strut Matin, course en auto à Puchevillers et\strut &  \\
      \leavevmode\fontsize{6.9pt}{7.9pt}\selectfont\strut Arquèves. Belle Église Farm, M. Houtteman[?].\strut & \smash{\parbox[t]{39mm}{\fontsize{6.2pt}{7.0pt}\selectfont\raggedright\color{black!88}\textbf{11 Houtteman} \textit{lect. incert. in ms.}}} \\
      \leavevmode\fontsize{6.9pt}{7.9pt}\selectfont\strut Après-midi, par Contay \& Warloy,\strut &  \\
      \leavevmode\fontsize{6.9pt}{7.9pt}\selectfont\strut venus à Senlis. Routes encombrées par troupes\strut &  \\
      \leavevmode\fontsize{6.9pt}{7.9pt}\selectfont\strut moving. Admiré cavrie indienne \& anglaise\strut & \smash{\parbox[t]{39mm}{\fontsize{6.2pt}{7.0pt}\selectfont\raggedright\color{black!88}\textbf{14 moving} \textit{angl. in ms.}}} \\
      \leavevmode\llap{\fontsize{6.4pt}{7.3pt}\selectfont\color{darkgray}15\hspace{1.8mm}}\fontsize{6.9pt}{7.9pt}\selectfont\strut (Death or Glory Boys, 17th Lancers).\strut & \smash{\parbox[t]{39mm}{\fontsize{6.2pt}{7.0pt}\selectfont\raggedright\color{black!88}\textbf{15 Death or Glory Boys} \textit{angl. in ms.}}} \\
      \cline{1-1}
      \leavevmode\fontsize{6.9pt}{7.9pt}\selectfont\strut \underline{22 Jeudi.}\strut &  \\
      \leavevmode\fontsize{6.9pt}{7.9pt}\selectfont\strut Matin \& soir, tourmentes de neige.\strut &  \\
      \leavevmode\fontsize{6.9pt}{7.9pt}\selectfont\strut Matin, courses en auto à Puchevillers, Mirvaux,\strut &  \\
      \leavevmode\fontsize{6.9pt}{7.9pt}\selectfont\strut Pierregot \& Molliens-au-Bois.\strut &  \\
      \leavevmode\llap{\fontsize{6.4pt}{7.3pt}\selectfont\color{darkgray}20\hspace{1.8mm}}\fontsize{6.9pt}{7.9pt}\selectfont\strut Après-midi, Senlis \& Albert. Très longuement\strut &  \\
      \leavevmode\fontsize{6.9pt}{7.9pt}\selectfont\strut regardé à Albert maisons atteintes par obus.\strut &  \\
      \leavevmode\fontsize{6.9pt}{7.9pt}\selectfont\strut Retour par gde route d'Amiens : à droite \&\strut & \smash{\parbox[t]{39mm}{\fontsize{6.2pt}{7.0pt}\selectfont\raggedright\color{black!88}\textbf{22 gde} \textit{abbr. pro} grande.}} \\
      \leavevmode\fontsize{6.9pt}{7.9pt}\selectfont\strut à g., champs couverts de neige.\strut &  \\
\end{tabular}}%
\par\vspace*{\fill}
% --- Registres de bas de page (1. Critique interne [si pertinent] ; 2. Loci paralleli en latin ; 3. Commentaire historique en retrait) ---
\noindent\rule{\textwidth}{0.5pt}\par\vspace{0.65mm}
{\fontsize{6.6pt}{7.7pt}\selectfont%
\noindent\normalfont\color{black!92} \textbf{2} (\textit{63rd Div}) cf. supra f\textsuperscript{o}~85~r\textsuperscript{o} (p.~\pageref{fol:170}). — \textbf{10} (\textit{Puchevillers}) cf. supra f\textsuperscript{o}~37~v\textsuperscript{o} (p.~\pageref{fol:75}), f\textsuperscript{o}~85~r\textsuperscript{o} (p.~\pageref{fol:170}), infra f\textsuperscript{o}~86~r\textsuperscript{o} (p.~\pageref{fol:172}), f\textsuperscript{o}~86~v\textsuperscript{o} (p.~\pageref{fol:173}), \textit{et passim}. — \textbf{11} (\textit{Arquèves}) cf. supra f\textsuperscript{o}~47~v\textsuperscript{o} (p.~\pageref{fol:95}), f\textsuperscript{o}~61~r\textsuperscript{o} (p.~\pageref{fol:122}), \textit{et passim}.\par\vspace{0.55mm}
\noindent\rule{38mm}{0.4pt}\par\vspace{0.5mm}
\noindent\hspace*{5mm}\parbox{\dimexpr\textwidth-5mm\relax}{\fontsize{6.6pt}{7.7pt}\selectfont\color{black!84}En mars 1917, lors de la marche en avant suivant le repli allemand sur la ligne Hindenburg, le \textbf{17\textsuperscript{e} Lanciers britannique} (\textit{17th Duke of Cambridge's Own Lancers}, surnommé \textit{Death or Glory Boys} d'après la devise de son insigne à tête de mort) formait l'unique régiment britannique de la \textit{Sialkot Cavalry Brigade} au sein de la \textbf{4\textsuperscript{e} Division de cavalerie} (ex-\textit{1st Indian Cavalry Division}), aux côtés des lanciers indiens des \textit{6th King Edward's Own Cavalry} et \textit{19th Lancers (Fane's Horse)} (Maj. A.~F. Becke, \textit{Order of Battle}, Part~2A, p.~35–40).}\par}
\newpage
\phantomsection\label{fac:171}
% --- En-tête du fac-similé (Page impaire / droite : Photographie du manuscrit) ---
\noindent\rule{\textwidth}{0.5pt}\par\vspace{0.55mm}
\noindent\textbf{\fontsize{6.8pt}{7.8pt}\selectfont FAC-SIMILÉ EN VIS-À-VIS — FOLIO 85 VERSO} \hfill \textbf{\fontsize{6.8pt}{7.8pt}\selectfont MARS 1917}\par\vspace{0.25mm}
\noindent{\fontsize{6.8pt}{8.0pt}\selectfont\textit{Livre-Journal d'Hippolyte Combes} — 20 – 22 mars 1917 (Hérissart, Puchevillers, Arquèves (Belle-Église), Contay, Warloy, Senlis, Pierregot, Molliens, Albert (Somme))}\par\vspace{0.55mm}
\noindent\rule{\textwidth}{0.35pt}\par\vspace*{\fill}
{\centering
  \setlength{\fboxsep}{1.2mm}%
  \setlength{\fboxrule}{0.4pt}%
  \fcolorbox{gray!55}{white}{%
    \includegraphics[width=0.95\textwidth,height=0.81\textheight,keepaspectratio]{\facsimiledir/20251126_111410.jpg}%
  }\par}
\vspace*{\fill}
\newpage

% =============================================================================
% VIS-A-VIS SPREAD : FOLIO 86 RECTO (20251126_111427.jpg)
% =============================================================================
\cleardoubleevenpage
\renewcommand{\currentmonthyear}{Mars 1917}
\renewcommand{\currentfoliolabel}{Folio 86 recto}
\phantomsection\label{fol:172}
% --- En-tête du folio (Page paire / gauche : Transcription, apparat & registres de bas de page) ---
\noindent\rule{\textwidth}{0.5pt}\par\vspace{0.55mm}
\noindent\textbf{\footnotesize FOLIO 86 RECTO} \hfill \textbf{\footnotesize MARS 1917}\par\vspace{0.25mm}
\noindent{\fontsize{6.9pt}{8.0pt}\selectfont Chants d'oiseaux dans le gel de mars et démarches pour conserver le billet 220}\par\vspace{0.2mm}
\noindent{\fontsize{6.9pt}{8.0pt}\selectfont\textit{23 – 25 mars 1917} \hfill \textit{Hérissart, Rubempré, Puchevillers (Somme)}}\par\vspace{0.5mm}
\noindent\rule{\textwidth}{0.35pt}\par\vspace{0.8mm}
% --- Corps de la transcription diplomatique et apparat marginal en prolongement ---
{\renewcommand{\arraystretch}{0.87}%
\noindent\hspace*{3.5mm}\begin{tabular}{@{}p{87mm}@{\hspace{3.5mm}}p{39mm}@{}}
      \leavevmode\fontsize{6.9pt}{7.9pt}\selectfont\strut \underline{23 Ven.} — \hfill{} \underline{— Mars 1917 —}\strut & \smash{\parbox[t]{39mm}{\fontsize{6.2pt}{7.0pt}\selectfont\raggedright\color{black!88}\textbf{1 Ven.} \textit{scil.} Vendredi.}} \\
      \leavevmode\fontsize{6.9pt}{7.9pt}\selectfont\strut Nuit, gelé très fort.\strut &  \\
      \leavevmode\fontsize{6.9pt}{7.9pt}\selectfont\strut Vitres fenêtre, couvertes de glace.\strut &  \\
      \leavevmode\fontsize{6.9pt}{7.9pt}\selectfont\strut Sol durci. Froid très vif.\strut &  \\
      \leavevmode\llap{\fontsize{6.4pt}{7.3pt}\selectfont\color{darkgray}5\hspace{1.8mm}}\fontsize{6.9pt}{7.9pt}\selectfont\strut Très beau ciel bleu, tout ensoleillé. Chants d'oiseaux.\strut &  \\
      \leavevmode\fontsize{6.9pt}{7.9pt}\selectfont\strut Notre auto, aux workshops. Matin \& soir,\strut & \smash{\parbox[t]{39mm}{\fontsize{6.2pt}{7.0pt}\selectfont\raggedright\color{black!88}\textbf{6 workshops} \textit{angl. in ms.}}} \\
      \leavevmode\fontsize{6.9pt}{7.9pt}\selectfont\strut liberté complète.\strut &  \\
      \leavevmode\fontsize{6.9pt}{7.9pt}\selectfont\strut Lu et écrit, office \& chambre.\strut &  \\
      \leavevmode\fontsize{6.9pt}{7.9pt}\selectfont\strut Fait démarches pour qu'on me laisse au billet\strut &  \\
      \leavevmode\llap{\fontsize{6.4pt}{7.3pt}\selectfont\color{darkgray}10\hspace{1.8mm}}\fontsize{6.9pt}{7.9pt}\selectfont\strut (220).\strut &  \\
      \leavevmode\fontsize{6.9pt}{7.9pt}\selectfont\strut Ds la journée, et afternoon, promené\strut & \smash{\parbox[t]{39mm}{\fontsize{6.2pt}{7.0pt}\selectfont\raggedright\color{black!88}\textbf{11 Ds} \textit{abbr. pro} Dans. \textbf{11 afternoon} \textit{angl. in ms.}}} \\
      \leavevmode\fontsize{6.9pt}{7.9pt}\selectfont\strut un peu campagne.\strut &  \\
      \leavevmode\fontsize{6.9pt}{7.9pt}\selectfont\strut Nuit très étoilée, mais froide.\strut &  \\
      \cline{1-1}
      \leavevmode\fontsize{6.9pt}{7.9pt}\selectfont\strut \underline{24 Sam.} — Nuit froide. Gelé assez fort! Sol durci.\strut & \smash{\parbox[t]{39mm}{\fontsize{6.2pt}{7.0pt}\selectfont\raggedright\color{black!88}\textbf{14 Sam.} \textit{scil.} Samedi.}} \\
      \leavevmode\llap{\fontsize{6.4pt}{7.3pt}\selectfont\color{darkgray}15\hspace{1.8mm}}\fontsize{6.9pt}{7.9pt}\selectfont\strut Tout le jour, soleil brille, très vif. Ciel d'un\strut &  \\
      \leavevmode\fontsize{6.9pt}{7.9pt}\selectfont\strut bleu profond. Vent souffle avec violence.\strut &  \\
      \leavevmode\fontsize{6.9pt}{7.9pt}\selectfont\strut Après-midi, venus à pied à Rubempré\strut &  \\
      \leavevmode\fontsize{6.9pt}{7.9pt}\selectfont\strut et j'à chemin Puchevillers —\strut & \smash{\parbox[t]{39mm}{\fontsize{6.2pt}{7.0pt}\selectfont\raggedright\color{black!88}\textbf{18 j'à} \textit{abbr. pro} jusqu'à.}} \\
      \leavevmode\fontsize{6.9pt}{7.9pt}\selectfont\strut Fin de journée magnifique. Superbe sunset.\strut & \smash{\parbox[t]{39mm}{\fontsize{6.2pt}{7.0pt}\selectfont\raggedright\color{black!88}\textbf{19 sunset} \textit{angl. in ms.}}} \\
      \cline{1-1}
      \leavevmode\llap{\fontsize{6.4pt}{7.3pt}\selectfont\color{darkgray}20\hspace{1.8mm}}\fontsize{6.9pt}{7.9pt}\selectfont\strut \underline{25 D.} — Matin \& soir, temps nuageux.\strut & \smash{\parbox[t]{39mm}{\fontsize{6.2pt}{7.0pt}\selectfont\raggedright\color{black!88}\textbf{20 D.} \textit{scil.} Dimanche.}} \\
      \leavevmode\fontsize{6.9pt}{7.9pt}\selectfont\strut Ds la journée, soleil ardent, ciel bleu —\strut & \smash{\parbox[t]{39mm}{\fontsize{6.2pt}{7.0pt}\selectfont\raggedright\color{black!88}\textbf{21 Ds} \textit{abbr. pro} Dans.}} \\
      \leavevmode\fontsize{6.9pt}{7.9pt}\selectfont\strut Course à pied à Rubempré, l'après-midi.\strut &  \\
      \leavevmode\fontsize{6.9pt}{7.9pt}\selectfont\strut Pas pris pardessus. Promenade fort agréable.\strut &  \\
\end{tabular}}%
\par\vspace*{\fill}
% --- Registres de bas de page (1. Critique interne [si pertinent] ; 2. Loci paralleli en latin ; 3. Commentaire historique en retrait) ---
\noindent\rule{\textwidth}{0.5pt}\par\vspace{0.65mm}
{\fontsize{6.6pt}{7.7pt}\selectfont%
\noindent\normalfont\color{black!92} \textbf{17} (\textit{Rubempré}) cf. supra f\textsuperscript{o}~47~v\textsuperscript{o} (p.~\pageref{fol:95}), f\textsuperscript{o}~81~r\textsuperscript{o} (p.~\pageref{fol:162}), infra f\textsuperscript{o}~86~v\textsuperscript{o} (p.~\pageref{fol:173}), f\textsuperscript{o}~87~r\textsuperscript{o} (p.~\pageref{fol:174}), \textit{et passim}. — \textbf{18} (\textit{Puchevillers}) cf. supra f\textsuperscript{o}~37~v\textsuperscript{o} (p.~\pageref{fol:75}), f\textsuperscript{o}~85~v\textsuperscript{o} (p.~\pageref{fol:171}), infra f\textsuperscript{o}~86~v\textsuperscript{o} (p.~\pageref{fol:173}), f\textsuperscript{o}~87~r\textsuperscript{o} (p.~\pageref{fol:174}), \textit{et passim}.\par\vspace{0.55mm}
\noindent\rule{38mm}{0.4pt}\par\vspace{0.5mm}
\noindent\hspace*{5mm}\parbox{\dimexpr\textwidth-5mm\relax}{\fontsize{6.6pt}{7.7pt}\selectfont\color{black!84}Trois journées de fin mars 1917 (23--25 mars) à \textbf{Hérissart}: l'automobile du service étant en réparation aux ateliers (\textit{workshops}), Hippolyte Combes obtient de conserver sa chambre d'étude du \textbf{billet n\textsuperscript{o}~220} (où il s'était installé le 11 mars, f\textsuperscript{o}~83~r\textsuperscript{o}--f\textsuperscript{o}~83~v\textsuperscript{o}) et effectue à pied ses courses vers \textbf{Rubempré} et \textbf{Puchevillers}.}\par}
\newpage
\phantomsection\label{fac:172}
% --- En-tête du fac-similé (Page impaire / droite : Photographie du manuscrit) ---
\noindent\rule{\textwidth}{0.5pt}\par\vspace{0.55mm}
\noindent\textbf{\fontsize{6.8pt}{7.8pt}\selectfont FAC-SIMILÉ EN VIS-À-VIS — FOLIO 86 RECTO} \hfill \textbf{\fontsize{6.8pt}{7.8pt}\selectfont MARS 1917}\par\vspace{0.25mm}
\noindent{\fontsize{6.8pt}{8.0pt}\selectfont\textit{Livre-Journal d'Hippolyte Combes} — 23 – 25 mars 1917 (Hérissart, Rubempré, Puchevillers (Somme))}\par\vspace{0.55mm}
\noindent\rule{\textwidth}{0.35pt}\par\vspace*{\fill}
{\centering
  \setlength{\fboxsep}{1.2mm}%
  \setlength{\fboxrule}{0.4pt}%
  \fcolorbox{gray!55}{white}{%
    \includegraphics[width=0.95\textwidth,height=0.81\textheight,keepaspectratio]{\facsimiledir/20251126_111427.jpg}%
  }\par}
\vspace*{\fill}
\newpage

% =============================================================================
% VIS-A-VIS SPREAD : FOLIO 86 VERSO (20251126_111439.jpg)
% =============================================================================
\cleardoubleevenpage
\renewcommand{\currentmonthyear}{Mars 1917}
\renewcommand{\currentfoliolabel}{Folio 86 verso}
\phantomsection\label{fol:173}
% --- En-tête du folio (Page paire / gauche : Transcription, apparat & registres de bas de page) ---
\noindent\rule{\textwidth}{0.5pt}\par\vspace{0.55mm}
\noindent\textbf{\footnotesize FOLIO 86 VERSO} \hfill \textbf{\footnotesize MARS 1917}\par\vspace{0.25mm}
\noindent{\fontsize{6.9pt}{8.0pt}\selectfont Visite du camarade Monin et départ de toutes les troupes d'Hérissart}\par\vspace{0.2mm}
\noindent{\fontsize{6.9pt}{8.0pt}\selectfont\textit{26 – 28 mars 1917} \hfill \textit{Hérissart, Rubempré, Puchevillers (Somme)}}\par\vspace{0.5mm}
\noindent\rule{\textwidth}{0.35pt}\par\vspace{0.8mm}
% --- Corps de la transcription diplomatique et apparat marginal en prolongement ---
{\renewcommand{\arraystretch}{0.93}%
\noindent\hspace*{3.5mm}\begin{tabular}{@{}p{87mm}@{\hspace{3.5mm}}p{39mm}@{}}
      \leavevmode\fontsize{7.3pt}{8.4pt}\selectfont\strut \underline{26 Lundi} \hfill{} \underline{Mars /17}\strut &  \\
      \leavevmode\fontsize{7.3pt}{8.4pt}\selectfont\strut Tt jour, pluie, intermittente. Rues \& routes,\strut & \smash{\parbox[t]{39mm}{\fontsize{6.2pt}{7.0pt}\selectfont\raggedright\color{black!88}\textbf{2 Tt} \textit{abbr. pro} Tout.}} \\
      \leavevmode\fontsize{7.3pt}{8.4pt}\selectfont\strut état lamentable, transformées en cloaques.\strut &  \\
      \leavevmode\fontsize{7.3pt}{8.4pt}\selectfont\strut On patauge boue j. à la cheville.\strut & \smash{\parbox[t]{39mm}{\fontsize{6.2pt}{7.0pt}\selectfont\raggedright\color{black!88}\textbf{4 j. à} \textit{abbr. pro} jusqu'à.}} \\
      \leavevmode\llap{\fontsize{6.4pt}{7.3pt}\selectfont\color{darkgray}5\hspace{1.8mm}}\fontsize{7.3pt}{8.4pt}\selectfont\strut Monin traversé village avec sa Coy pr cantonner\strut & \smash{\parbox[t]{39mm}{\fontsize{6.2pt}{7.0pt}\selectfont\raggedright\color{black!88}\textbf{5 pr} \textit{abbr. pro} pour. \textbf{5 Coy} \textit{angl.} Company. — \textbf{6 gd} \textit{abbr. pro} grand.}} \\
      \leavevmode\fontsize{7.3pt}{8.4pt}\selectfont\strut une nuit. L'ai revu avec gd plaisir. Passé ensemble\strut &  \\
      \leavevmode\fontsize{7.3pt}{8.4pt}\selectfont\strut plus. heures, après-midi. L'ai trouvé vieilli, enlaidi, usé.\strut &  \\
      \cline{1-1}
      \leavevmode\fontsize{7.3pt}{8.4pt}\selectfont\strut \underline{27 Mardi} —\strut &  \\
      \leavevmode\fontsize{7.3pt}{8.4pt}\selectfont\strut Giboulées. Tour à tour averses et chutes\strut &  \\
      \leavevmode\llap{\fontsize{6.4pt}{7.3pt}\selectfont\color{darkgray}10\hspace{1.8mm}}\fontsize{7.3pt}{8.4pt}\selectfont\strut abondantes de neige. Fin après-midi,\strut &  \\
      \leavevmode\fontsize{7.3pt}{8.4pt}\selectfont\strut gros flocons durant 1 h.\strut &  \\
      \leavevmode\fontsize{7.3pt}{8.4pt}\selectfont\strut Matinée, course à pied, à Rubempré.\strut &  \\
      \leavevmode\fontsize{7.3pt}{8.4pt}\selectfont\strut Après-midi, venus à Puchevillers. Routes détrempées\strut &  \\
      \leavevmode\fontsize{7.3pt}{8.4pt}\selectfont\strut affrt boueuses.\strut & \smash{\parbox[t]{39mm}{\fontsize{6.2pt}{7.0pt}\selectfont\raggedright\color{black!88}\textbf{14 affrt} \textit{abbr. pro} affreusement.}} \\
      \leavevmode\llap{\fontsize{6.4pt}{7.3pt}\selectfont\color{darkgray}15\hspace{1.8mm}}\fontsize{7.3pt}{8.4pt}\selectfont\strut Troupes quitté village à l'aube. Tout le jour,\strut &  \\
      \leavevmode\fontsize{7.3pt}{8.4pt}\selectfont\strut restés presq sans troupes. Village redevenu très calme. —\strut &  \\
      \cline{1-1}
      \leavevmode\fontsize{7.3pt}{8.4pt}\selectfont\strut \underline{28 Merc.} —\strut & \smash{\parbox[t]{39mm}{\fontsize{6.2pt}{7.0pt}\selectfont\raggedright\color{black!88}\textbf{17 Merc.} \textit{scil.} Mercredi.}} \\
      \leavevmode\fontsize{7.3pt}{8.4pt}\selectfont\strut Liberté, tout le jour.\strut &  \\
      \leavevmode\fontsize{7.3pt}{8.4pt}\selectfont\strut Froid assez vif. Pluie \& neige, intermittentes.\strut &  \\
      \leavevmode\llap{\fontsize{6.4pt}{7.3pt}\selectfont\color{darkgray}20\hspace{1.8mm}}\fontsize{7.3pt}{8.4pt}\selectfont\strut Fin après-midi, promenade s/ route de Puchevillers.\strut & \smash{\parbox[t]{39mm}{\fontsize{6.2pt}{7.0pt}\selectfont\raggedright\color{black!88}\textbf{20 s/} \textit{abbr. pro} sur.}} \\
      \leavevmode\fontsize{7.3pt}{8.4pt}\selectfont\strut Temps funèbre. Ciel nuageux. Horizon, noyé ds\strut & \smash{\parbox[t]{39mm}{\fontsize{6.2pt}{7.0pt}\selectfont\raggedright\color{black!88}\textbf{21 ds} \textit{abbr. pro} dans.}} \\
      \leavevmode\fontsize{7.3pt}{8.4pt}\selectfont\strut la brume —\strut &  \\
\end{tabular}}%
\par\vspace*{\fill}
% --- Registres de bas de page (1. Critique interne [si pertinent] ; 2. Loci paralleli en latin ; 3. Commentaire historique en retrait) ---
\noindent\rule{\textwidth}{0.5pt}\par\vspace{0.65mm}
{\fontsize{6.6pt}{7.7pt}\selectfont%
\noindent\normalfont\color{black!92} \textbf{12} (\textit{Rubempré}) cf. supra f\textsuperscript{o}~47~v\textsuperscript{o} (p.~\pageref{fol:95}), f\textsuperscript{o}~86~r\textsuperscript{o} (p.~\pageref{fol:172}), infra f\textsuperscript{o}~87~r\textsuperscript{o} (p.~\pageref{fol:174}), f\textsuperscript{o}~87~v\textsuperscript{o} (p.~\pageref{fol:175}), \textit{et passim}. — \textbf{13} (\textit{Puchevillers}) cf. supra f\textsuperscript{o}~37~v\textsuperscript{o} (p.~\pageref{fol:75}), f\textsuperscript{o}~86~r\textsuperscript{o} (p.~\pageref{fol:172}), infra f\textsuperscript{o}~87~r\textsuperscript{o} (p.~\pageref{fol:174}), f\textsuperscript{o}~91~v\textsuperscript{o} (p.~\pageref{fol:183}), \textit{et passim}.\par\vspace{0.55mm}
\noindent\rule{38mm}{0.4pt}\par\vspace{0.5mm}
\noindent\hspace*{5mm}\parbox{\dimexpr\textwidth-5mm\relax}{\fontsize{6.6pt}{7.7pt}\selectfont\color{black!84}Le lundi 26 mars 1917, Hippolyte Combes retrouve à Hérissart son camarade interprète \textbf{Monin}, de passage pour une nuit avec sa compagnie.}\par}
\newpage
\phantomsection\label{fac:173}
% --- En-tête du fac-similé (Page impaire / droite : Photographie du manuscrit) ---
\noindent\rule{\textwidth}{0.5pt}\par\vspace{0.55mm}
\noindent\textbf{\fontsize{6.8pt}{7.8pt}\selectfont FAC-SIMILÉ EN VIS-À-VIS — FOLIO 86 VERSO} \hfill \textbf{\fontsize{6.8pt}{7.8pt}\selectfont MARS 1917}\par\vspace{0.25mm}
\noindent{\fontsize{6.8pt}{8.0pt}\selectfont\textit{Livre-Journal d'Hippolyte Combes} — 26 – 28 mars 1917 (Hérissart, Rubempré, Puchevillers (Somme))}\par\vspace{0.55mm}
\noindent\rule{\textwidth}{0.35pt}\par\vspace*{\fill}
{\centering
  \setlength{\fboxsep}{1.2mm}%
  \setlength{\fboxrule}{0.4pt}%
  \fcolorbox{gray!55}{white}{%
    \includegraphics[width=0.95\textwidth,height=0.81\textheight,keepaspectratio]{\facsimiledir/20251126_111439.jpg}%
  }\par}
\vspace*{\fill}
\newpage

% =============================================================================
% VIS-A-VIS SPREAD : FOLIO 87 RECTO (20251126_111452.jpg)
% =============================================================================
\cleardoubleevenpage
\renewcommand{\currentmonthyear}{Mars 1917}
\renewcommand{\currentfoliolabel}{Folio 87 recto}
\phantomsection\label{fol:174}
% --- En-tête du folio (Page paire / gauche : Transcription, apparat & registres de bas de page) ---
\noindent\rule{\textwidth}{0.5pt}\par\vspace{0.55mm}
\noindent\textbf{\footnotesize FOLIO 87 RECTO} \hfill \textbf{\footnotesize MARS 1917}\par\vspace{0.25mm}
\noindent{\fontsize{6.9pt}{8.0pt}\selectfont Le calme délicieux des villages désertés par les troupes à la fin mars}\par\vspace{0.2mm}
\noindent{\fontsize{6.9pt}{8.0pt}\selectfont\textit{29 – 31 mars 1917} \hfill \textit{Hérissart, Puchevillers, Rubempré, Pierregot, Molliens-au-Bois, Toutencourt, Rainneville (Somme)}}\par\vspace{0.5mm}
\noindent\rule{\textwidth}{0.35pt}\par\vspace{0.8mm}
% --- Corps de la transcription diplomatique et apparat marginal en prolongement ---
{\renewcommand{\arraystretch}{0.83}%
\noindent\hspace*{3.5mm}\begin{tabular}{@{}p{87mm}@{\hspace{3.5mm}}p{39mm}@{}}
      \leavevmode\fontsize{6.6pt}{7.5pt}\selectfont\strut \underline{29 Jeudi} \hfill{} \underline{— Mars 1917 —}\strut &  \\
      \leavevmode\fontsize{6.6pt}{7.5pt}\selectfont\strut Temps moins froid. Mais, presque tt le jour,\strut & \smash{\parbox[t]{39mm}{\fontsize{6.2pt}{7.0pt}\selectfont\raggedright\color{black!88}\textbf{2 tt} \textit{abbr. pro} tout.}} \\
      \leavevmode\fontsize{6.6pt}{7.5pt}\selectfont\strut ondées torrentielles.\strut &  \\
      \leavevmode\fontsize{6.6pt}{7.5pt}\selectfont\strut Matinée, courses en auto : Puchevillers \&\strut &  \\
      \leavevmode\llap{\fontsize{6.4pt}{7.3pt}\selectfont\color{darkgray}5\hspace{1.8mm}}\fontsize{6.6pt}{7.5pt}\selectfont\strut Rubempré. Averses très fortes.\strut &  \\
      \leavevmode\fontsize{6.6pt}{7.5pt}\selectfont\strut Après-midi entière, liberté.\strut &  \\
      \leavevmode\fontsize{6.6pt}{7.5pt}\selectfont\strut Village paraît abandonné. Plus de troupes.\strut &  \\
      \leavevmode\fontsize{6.6pt}{7.5pt}\selectfont\strut Rues désertes. Calme délicieux —\strut &  \\
      \cline{1-1}
      \leavevmode\fontsize{6.6pt}{7.5pt}\selectfont\strut \underline{30 Ven.} —\strut & \smash{\parbox[t]{39mm}{\fontsize{6.2pt}{7.0pt}\selectfont\raggedright\color{black!88}\textbf{9 Ven.} \textit{scil.} Vendredi.}} \\
      \leavevmode\llap{\fontsize{6.4pt}{7.3pt}\selectfont\color{darkgray}10\hspace{1.8mm}}\fontsize{6.6pt}{7.5pt}\selectfont\strut Encore aujourd'hui, giboulées, saison. Averses très fortes,\strut &  \\
      \leavevmode\fontsize{6.6pt}{7.5pt}\selectfont\strut un peu de neige, rafales violentes.\strut &  \\
      \leavevmode\fontsize{6.6pt}{7.5pt}\selectfont\strut Routes, état lamentable, creusées de sillons\strut &  \\
      \leavevmode\fontsize{6.6pt}{7.5pt}\selectfont\strut profonds par gros camions \& lourds véhicules.\strut &  \\
      \leavevmode\fontsize{6.6pt}{7.5pt}\selectfont\strut Matinée, courses à Pierregot \& Molliens-au-B.\strut &  \\
      \leavevmode\llap{\fontsize{6.4pt}{7.3pt}\selectfont\color{darkgray}15\hspace{1.8mm}}\fontsize{6.6pt}{7.5pt}\selectfont\strut Après-midi, liberté complète.\strut &  \\
      \leavevmode\fontsize{6.6pt}{7.5pt}\selectfont\strut Mervx ciel étoilé, croissant lune brillant,\strut &  \\
      \leavevmode\fontsize{6.6pt}{7.5pt}\selectfont\strut nuages très rapides. État lyrique — —\strut &  \\
      \cline{1-1}
      \leavevmode\fontsize{6.6pt}{7.5pt}\selectfont\strut \underline{31 Sam.} — En course, matin \& soir : Rubempré,\strut & \smash{\parbox[t]{39mm}{\fontsize{6.2pt}{7.0pt}\selectfont\raggedright\color{black!88}\textbf{18 Sam.} \textit{scil.} Samedi.}} \\
      \leavevmode\fontsize{6.6pt}{7.5pt}\selectfont\strut par Toutencourt, Puchevillers, Pierregot et\strut &  \\
      \leavevmode\llap{\fontsize{6.4pt}{7.3pt}\selectfont\color{darkgray}20\hspace{1.8mm}}\fontsize{6.6pt}{7.5pt}\selectfont\strut Rainneville — Acheté paille.\strut &  \\
      \leavevmode\fontsize{6.6pt}{7.5pt}\selectfont\strut Temps très mauvais. Surpris en route par\strut &  \\
      \leavevmode\fontsize{6.6pt}{7.5pt}\selectfont\strut fortes averses, grêle — grêlons, frappé nos\strut &  \\
      \leavevmode\fontsize{6.6pt}{7.5pt}\selectfont\strut visages — Routes extrt détrempées — —\strut & \smash{\parbox[t]{39mm}{\fontsize{6.2pt}{7.0pt}\selectfont\raggedright\color{black!88}\textbf{23 extrt} \textit{abbr. pro} extrêmement.}} \\
      \leavevmode\fontsize{6.6pt}{7.5pt}\selectfont\strut Village tjours sans troupes, paraît inhabité.\strut &  \\
      \leavevmode\llap{\fontsize{6.4pt}{7.3pt}\selectfont\color{darkgray}25\hspace{1.8mm}}\fontsize{6.6pt}{7.5pt}\selectfont\strut J'aime ce calme \& ce silence.\strut &  \\
\end{tabular}}%
\par\vspace*{\fill}
% --- Registres de bas de page (1. Critique interne [si pertinent] ; 2. Loci paralleli en latin ; 3. Commentaire historique en retrait) ---
\noindent\rule{\textwidth}{0.5pt}\par\vspace{0.65mm}
{\fontsize{6.6pt}{7.7pt}\selectfont%
\noindent\normalfont\color{black!92} \textbf{4} (\textit{Puchevillers}) cf. supra f\textsuperscript{o}~37~v\textsuperscript{o} (p.~\pageref{fol:75}), f\textsuperscript{o}~86~v\textsuperscript{o} (p.~\pageref{fol:173}), infra f\textsuperscript{o}~91~v\textsuperscript{o} (p.~\pageref{fol:183}), f\textsuperscript{o}~92~r\textsuperscript{o} (p.~\pageref{fol:184}), \textit{et passim}. — \textbf{5} (\textit{Rubempré}) cf. supra f\textsuperscript{o}~47~v\textsuperscript{o} (p.~\pageref{fol:95}), f\textsuperscript{o}~86~v\textsuperscript{o} (p.~\pageref{fol:173}), infra f\textsuperscript{o}~87~v\textsuperscript{o} (p.~\pageref{fol:175}), f\textsuperscript{o}~88~v\textsuperscript{o} (p.~\pageref{fol:177}), \textit{et passim}. — \textbf{14} (\textit{Pierregot}) cf. supra f\textsuperscript{o}~52~r\textsuperscript{o} (p.~\pageref{fol:104}), f\textsuperscript{o}~85~v\textsuperscript{o} (p.~\pageref{fol:171}), infra f\textsuperscript{o}~87~v\textsuperscript{o} (p.~\pageref{fol:175}), f\textsuperscript{o}~89~r\textsuperscript{o} (p.~\pageref{fol:178}), \textit{et passim}.\par\vspace{0.55mm}
\noindent\rule{38mm}{0.4pt}\par\vspace{0.5mm}
\noindent\hspace*{5mm}\parbox{\dimexpr\textwidth-5mm\relax}{\fontsize{6.6pt}{7.7pt}\selectfont\color{black!84}Du 29 au 31 mars 1917, tandis que les tournées d'achat de paille du \textit{Purchase Board} rayonnent autour d'Hérissart jusqu'à \textbf{Pierregot}, \textbf{Molliens-au-Bois} et \textbf{Rainneville} (Somme, canton de Villers-Bocage, à 8~km au sud-ouest d'Hérissart), Hippolyte Combes savoure le silence et l'apaisement du village d'\textbf{Hérissart} désormais déserté par les troupes britanniques.}\par}
\newpage
\phantomsection\label{fac:174}
% --- En-tête du fac-similé (Page impaire / droite : Photographie du manuscrit) ---
\noindent\rule{\textwidth}{0.5pt}\par\vspace{0.55mm}
\noindent\textbf{\fontsize{6.8pt}{7.8pt}\selectfont FAC-SIMILÉ EN VIS-À-VIS — FOLIO 87 RECTO} \hfill \textbf{\fontsize{6.8pt}{7.8pt}\selectfont MARS 1917}\par\vspace{0.25mm}
\noindent{\fontsize{6.8pt}{8.0pt}\selectfont\textit{Livre-Journal d'Hippolyte Combes} — 29 – 31 mars 1917 (Hérissart, Puchevillers, Rubempré, Pierregot, Molliens-au-Bois, Toutencourt, Rainneville (Somme))}\par\vspace{0.55mm}
\noindent\rule{\textwidth}{0.35pt}\par\vspace*{\fill}
{\centering
  \setlength{\fboxsep}{1.2mm}%
  \setlength{\fboxrule}{0.4pt}%
  \fcolorbox{gray!55}{white}{%
    \includegraphics[width=0.95\textwidth,height=0.81\textheight,keepaspectratio]{\facsimiledir/20251126_111452.jpg}%
  }\par}
\vspace*{\fill}
\newpage

% =============================================================================
% VIS-A-VIS SPREAD : FOLIO 87 VERSO (20251126_111503.jpg)
% =============================================================================
\cleardoubleevenpage
\renewcommand{\currentmonthyear}{Avril 1917}
\renewcommand{\currentfoliolabel}{Folio 87 verso}
\phantomsection\label{fol:175}
% --- En-tête du folio (Page paire / gauche : Transcription, apparat & registres de bas de page) ---
\noindent\rule{\textwidth}{0.5pt}\par\vspace{0.55mm}
\noindent\textbf{\footnotesize FOLIO 87 VERSO} \hfill \textbf{\footnotesize AVRIL 1917}\par\vspace{0.25mm}
\noindent{\fontsize{6.9pt}{8.0pt}\selectfont Dimanche des Rameaux sous la neige et comptes du mess au 31 mars}\par\vspace{0.2mm}
\noindent{\fontsize{6.9pt}{8.0pt}\selectfont\textit{1er – 3 avril 1917} \hfill \textit{Hérissart, Belle-Église, Pierregot, Rainneville, Toutencourt, Rubempré, Raincheval (Somme)}}\par\vspace{0.5mm}
\noindent\rule{\textwidth}{0.35pt}\par\vspace{0.8mm}
% --- Corps de la transcription diplomatique et apparat marginal en prolongement ---
{\renewcommand{\arraystretch}{0.87}%
\noindent\hspace*{3.5mm}\begin{tabular}{@{}p{87mm}@{\hspace{3.5mm}}p{39mm}@{}}
      \leavevmode\fontsize{6.9pt}{7.9pt}\selectfont\strut \underline{1er Dim. (Rameaux)} \hfill{} \raisebox{0.5ex}{\underline{Avril 1917}}\strut & \smash{\parbox[t]{39mm}{\fontsize{6.2pt}{7.0pt}\selectfont\raggedright\color{black!88}\textbf{1 Dim.} \textit{scil.} Dimanche. \textbf{1 1er} \textit{scil.} premier.}} \\
      \leavevmode\fontsize{6.9pt}{7.9pt}\selectfont\strut Temps mauvais \& froid.\strut &  \\
      \leavevmode\fontsize{6.9pt}{7.9pt}\selectfont\strut Averses torrentielles et même, soirée, neige à\strut &  \\
      \leavevmode\fontsize{6.9pt}{7.9pt}\selectfont\strut gros flocons. Routes, état affreux.\strut &  \\
      \leavevmode\llap{\fontsize{6.4pt}{7.3pt}\selectfont\color{darkgray}5\hspace{1.8mm}}\fontsize{6.9pt}{7.9pt}\selectfont\strut Matin \& soir, courses en auto : Belle Église,\strut &  \\
      \leavevmode\fontsize{6.9pt}{7.9pt}\selectfont\strut Pierregot, Rainneville.\strut &  \\
      \leavevmode\fontsize{6.9pt}{7.9pt}\selectfont\strut Gde fatigue nerveuse. [ Mess : dû au 31 mars : 27 f. 50 ]\strut & \smash{\parbox[t]{39mm}{\fontsize{6.2pt}{7.0pt}\selectfont\raggedright\color{black!88}\textbf{7} \textit{in marg. cum unco} (comptes). \textbf{7 Gde} \textit{abbr. pro} Grande.}} \\
      \cline{1-1}
      \leavevmode\fontsize{6.9pt}{7.9pt}\selectfont\strut \underline{2 Lundi.} — Nuit froide. Au réveil, trouvé\strut &  \\
      \leavevmode\fontsize{6.9pt}{7.9pt}\selectfont\strut sol durci. Fortement gelé.\strut &  \\
      \leavevmode\llap{\fontsize{6.4pt}{7.3pt}\selectfont\color{darkgray}10\hspace{1.8mm}}\fontsize{6.9pt}{7.9pt}\selectfont\strut Tout le jour, liberté complète. Resté matin\strut &  \\
      \leavevmode\fontsize{6.9pt}{7.9pt}\selectfont\strut et soir, office, devant bon feu, à lire.\strut &  \\
      \leavevmode\fontsize{6.9pt}{7.9pt}\selectfont\strut Après-midi, durant plus de 2 heures,\strut &  \\
      \leavevmode\fontsize{6.9pt}{7.9pt}\selectfont\strut neige à gros flocons.\strut &  \\
      \leavevmode\fontsize{6.9pt}{7.9pt}\selectfont\strut Fin soirée, dégel et fonte neige. Rues\strut &  \\
      \leavevmode\llap{\fontsize{6.4pt}{7.3pt}\selectfont\color{darkgray}15\hspace{1.8mm}}\fontsize{6.9pt}{7.9pt}\selectfont\strut noyées boue liquide.\strut &  \\
      \leavevmode\fontsize{6.9pt}{7.9pt}\selectfont\strut [ Washing : 50 c. — — Clous souliers : 60 c. ]\strut & \smash{\parbox[t]{39mm}{\fontsize{6.2pt}{7.0pt}\selectfont\raggedright\color{black!88}\textbf{16} \textit{in marg. cum unco} (comptes). \textbf{16 c.} \textit{abbr. pro} centimes.}} \\
      \cline{1-1}
      \leavevmode\fontsize{6.9pt}{7.9pt}\selectfont\strut \underline{3 Mardi.} — Matin \& soir, courses en auto :\strut &  \\
      \leavevmode\fontsize{6.9pt}{7.9pt}\selectfont\strut Toutencourt, (Me Bien aimé, notaire)\strut &  \\
      \leavevmode\fontsize{6.9pt}{7.9pt}\selectfont\strut et Rubempré et Raincheval.\strut &  \\
      \leavevmode\llap{\fontsize{6.4pt}{7.3pt}\selectfont\color{darkgray}20\hspace{1.8mm}}\fontsize{6.9pt}{7.9pt}\selectfont\strut Surtout matin, en descendant côte de Hérissart\strut &  \\
      \leavevmode\fontsize{6.9pt}{7.9pt}\selectfont\strut admiré champs tout blancs de neige.\strut &  \\
      \leavevmode\fontsize{6.9pt}{7.9pt}\selectfont\strut Après-midi, dégel rapide. Routes\strut &  \\
      \leavevmode\fontsize{6.9pt}{7.9pt}\selectfont\strut changées en marécages.\strut &  \\
      \leavevmode\fontsize{6.9pt}{7.9pt}\selectfont\strut Temps redevenu plus doux.\strut &  \\
\end{tabular}}%
\par\vspace*{\fill}
% --- Registres de bas de page (1. Critique interne [si pertinent] ; 2. Loci paralleli en latin ; 3. Commentaire historique en retrait) ---
\noindent\rule{\textwidth}{0.5pt}\par\vspace{0.65mm}
{\fontsize{6.6pt}{7.7pt}\selectfont%
\noindent\normalfont\color{black!92} \textbf{5} (\textit{Belle Église}) cf. supra f\textsuperscript{o}~47~v\textsuperscript{o} (p.~\pageref{fol:95}), f\textsuperscript{o}~85~v\textsuperscript{o} (p.~\pageref{fol:171}), infra f\textsuperscript{o}~89~r\textsuperscript{o} (p.~\pageref{fol:178}), f\textsuperscript{o}~89~v\textsuperscript{o} (p.~\pageref{fol:179}), \textit{et passim}. — \textbf{6} (\textit{Pierregot}) cf. supra f\textsuperscript{o}~52~r\textsuperscript{o} (p.~\pageref{fol:104}), f\textsuperscript{o}~87~r\textsuperscript{o} (p.~\pageref{fol:174}), infra f\textsuperscript{o}~89~r\textsuperscript{o} (p.~\pageref{fol:178}), f\textsuperscript{o}~90~r\textsuperscript{o} (p.~\pageref{fol:180}), \textit{et passim}. — \textbf{18} (\textit{Toutencourt}) cf. supra f\textsuperscript{o}~49~v\textsuperscript{o} (p.~\pageref{fol:99}), f\textsuperscript{o}~87~r\textsuperscript{o} (p.~\pageref{fol:174}), infra f\textsuperscript{o}~88~r\textsuperscript{o} (p.~\pageref{fol:176}), f\textsuperscript{o}~88~v\textsuperscript{o} (p.~\pageref{fol:177}), \textit{et passim}.\par\vspace{0.55mm}
\noindent\rule{38mm}{0.4pt}\par\vspace{0.5mm}
\noindent\hspace*{5mm}\parbox{\dimexpr\textwidth-5mm\relax}{\fontsize{6.6pt}{7.7pt}\selectfont\color{black!84}Le dimanche des Rameaux (\textbf{1\textsuperscript{er} avril 1917}) et les 2--3 avril sont marqués par le retour de fortes chutes de neige tardives sur les plateaux de la Somme.}\par}
\newpage
\phantomsection\label{fac:175}
% --- En-tête du fac-similé (Page impaire / droite : Photographie du manuscrit) ---
\noindent\rule{\textwidth}{0.5pt}\par\vspace{0.55mm}
\noindent\textbf{\fontsize{6.8pt}{7.8pt}\selectfont FAC-SIMILÉ EN VIS-À-VIS — FOLIO 87 VERSO} \hfill \textbf{\fontsize{6.8pt}{7.8pt}\selectfont AVRIL 1917}\par\vspace{0.25mm}
\noindent{\fontsize{6.8pt}{8.0pt}\selectfont\textit{Livre-Journal d'Hippolyte Combes} — 1er – 3 avril 1917 (Hérissart, Belle-Église, Pierregot, Rainneville, Toutencourt, Rubempré, Raincheval (Somme))}\par\vspace{0.55mm}
\noindent\rule{\textwidth}{0.35pt}\par\vspace*{\fill}
{\centering
  \setlength{\fboxsep}{1.2mm}%
  \setlength{\fboxrule}{0.4pt}%
  \fcolorbox{gray!55}{white}{%
    \includegraphics[width=0.95\textwidth,height=0.81\textheight,keepaspectratio]{\facsimiledir/20251126_111503.jpg}%
  }\par}
\vspace*{\fill}
\newpage

% =============================================================================
% VIS-A-VIS SPREAD : FOLIO 88 RECTO (20251126_111525.jpg)
% =============================================================================
\cleardoubleevenpage
\renewcommand{\currentmonthyear}{Avril 1917}
\renewcommand{\currentfoliolabel}{Folio 88 recto}
\phantomsection\label{fol:176}
% --- En-tête du folio (Page paire / gauche : Transcription, apparat & registres de bas de page) ---
\noindent\rule{\textwidth}{0.5pt}\par\vspace{0.55mm}
\noindent\textbf{\footnotesize FOLIO 88 RECTO} \hfill \textbf{\footnotesize AVRIL 1917}\par\vspace{0.25mm}
\noindent{\fontsize{6.9pt}{8.0pt}\selectfont Jeudi Saint ensoleillé : les villages déserts de l'ancien front et le cimetière d'Hérissart}\par\vspace{0.2mm}
\noindent{\fontsize{6.9pt}{8.0pt}\selectfont\textit{4 – 5 avril 1917} \hfill \textit{Hérissart, Toutencourt, Acheux, Contay, Warloy, Varennes (Somme)}}\par\vspace{0.5mm}
\noindent\rule{\textwidth}{0.35pt}\par\vspace{0.8mm}
% --- Corps de la transcription diplomatique et apparat marginal en prolongement ---
{\renewcommand{\arraystretch}{0.87}%
\noindent\hspace*{3.5mm}\begin{tabular}{@{}p{87mm}@{\hspace{3.5mm}}p{39mm}@{}}
      \leavevmode\fontsize{6.9pt}{7.9pt}\selectfont\strut — Avril 1917 —\strut &  \\
      \leavevmode\fontsize{6.9pt}{7.9pt}\selectfont\strut \underline{4 Merc.} — Temps affreux. Matinée, tourmente\strut & \smash{\parbox[t]{39mm}{\fontsize{6.2pt}{7.0pt}\selectfont\raggedright\color{black!88}\textbf{2 Merc.} \textit{scil.} Mercredi.}} \\
      \leavevmode\fontsize{6.9pt}{7.9pt}\selectfont\strut de neige. Course à Toutencourt. Sous les flocons\strut &  \\
      \leavevmode\fontsize{6.9pt}{7.9pt}\selectfont\strut drus et serrés, allés de ferme en ferme, pr\strut & \smash{\parbox[t]{39mm}{\fontsize{6.2pt}{7.0pt}\selectfont\raggedright\color{black!88}\textbf{4 pr} \textit{abbr. pro} pour.}} \\
      \leavevmode\llap{\fontsize{6.4pt}{7.3pt}\selectfont\color{darkgray}5\hspace{1.8mm}}\fontsize{6.9pt}{7.9pt}\selectfont\strut acheter fourrage.\strut &  \\
      \leavevmode\fontsize{6.9pt}{7.9pt}\selectfont\strut Après-midi, liberté entière. Écrit ds ma\strut & \smash{\parbox[t]{39mm}{\fontsize{6.2pt}{7.0pt}\selectfont\raggedright\color{black!88}\textbf{6 ds} \textit{abbr. pro} dans.}} \\
      \leavevmode\fontsize{6.9pt}{7.9pt}\selectfont\strut chambre et lu jx au bureau.\strut &  \\
      \leavevmode\fontsize{6.9pt}{7.9pt}\selectfont\strut [ Hair cut : 30 c. — — Carnets de timbres : 4 f. — ]\strut & \smash{\parbox[t]{39mm}{\fontsize{6.2pt}{7.0pt}\selectfont\raggedright\color{black!88}\textbf{8} \textit{in marg. cum unco} (comptes). \textbf{8 f.} \textit{abbr. pro} franc(s).}} \\
      \cline{1-1}
      \leavevmode\fontsize{6.9pt}{7.9pt}\selectfont\strut \underline{5 Jeudi}-Saint — Enfin, belle journée. Ciel bleu,\strut &  \\
      \leavevmode\llap{\fontsize{6.4pt}{7.3pt}\selectfont\color{darkgray}10\hspace{1.8mm}}\fontsize{6.9pt}{7.9pt}\selectfont\strut tout ensoleillé.\strut &  \\
      \leavevmode\fontsize{6.9pt}{7.9pt}\selectfont\strut Après-midi, course en auto : Acheux, par\strut &  \\
      \leavevmode\fontsize{6.9pt}{7.9pt}\selectfont\strut Contay-Warloy-Varennes. Ts ces villages, très calmes,\strut & \smash{\parbox[t]{39mm}{\fontsize{6.2pt}{7.0pt}\selectfont\raggedright\color{black!88}\textbf{12 Ts} \textit{abbr. pro} Tous.}} \\
      \leavevmode\fontsize{6.9pt}{7.9pt}\selectfont\strut paraissent déserts, maintenant sans troupes.\strut &  \\
      \cline{1-1}
      \leavevmode\fontsize{6.9pt}{7.9pt}\selectfont\strut Dîner terminé, promenade habituelle, à\strut &  \\
      \leavevmode\llap{\fontsize{6.4pt}{7.3pt}\selectfont\color{darkgray}15\hspace{1.8mm}}\fontsize{6.9pt}{7.9pt}\selectfont\strut proximité d'Hérissart. Beau sunset, derrière\strut & \smash{\parbox[t]{39mm}{\fontsize{6.2pt}{7.0pt}\selectfont\raggedright\color{black!88}\textbf{15 sunset} \textit{angl. in ms.}}} \\
      \leavevmode\fontsize{6.9pt}{7.9pt}\selectfont\strut ligne de cyprès du cimetière. Contemplé vue\strut &  \\
      \leavevmode\fontsize{6.9pt}{7.9pt}\selectfont\strut panoramique circulaire.\strut &  \\
      \leavevmode\fontsize{6.9pt}{7.9pt}\selectfont\strut Soir, église village, entendu prêche passion.\strut &  \\
      \leavevmode\fontsize{6.9pt}{7.9pt}\selectfont\strut Nuit merve. Pleine lune, magnifique.\strut & \smash{\parbox[t]{39mm}{\fontsize{6.2pt}{7.0pt}\selectfont\raggedright\color{black!88}\textbf{19 merve} \textit{abbr. pro} merveilleuse.}} \\
\end{tabular}}%
\par\vspace*{\fill}
% --- Registres de bas de page (1. Critique interne [si pertinent] ; 2. Loci paralleli en latin ; 3. Commentaire historique en retrait) ---
\noindent\rule{\textwidth}{0.5pt}\par\vspace{0.65mm}
{\fontsize{6.6pt}{7.7pt}\selectfont%
\noindent\normalfont\color{black!92} \textbf{3} (\textit{Toutencourt}) cf. supra f\textsuperscript{o}~49~v\textsuperscript{o} (p.~\pageref{fol:99}), f\textsuperscript{o}~87~v\textsuperscript{o} (p.~\pageref{fol:175}), infra f\textsuperscript{o}~88~v\textsuperscript{o} (p.~\pageref{fol:177}), f\textsuperscript{o}~89~r\textsuperscript{o} (p.~\pageref{fol:178}), \textit{et passim}. — \textbf{5} (\textit{fourrage}) cf. supra f\textsuperscript{o}~3~r\textsuperscript{o} (p.~\pageref{fol:6}), f\textsuperscript{o}~84~v\textsuperscript{o} (p.~\pageref{fol:169}), \textit{et passim}. — \textbf{11} (\textit{Acheux}) cf. supra f\textsuperscript{o}~21~v\textsuperscript{o} (p.~\pageref{fol:43}), f\textsuperscript{o}~84~r\textsuperscript{o} (p.~\pageref{fol:168}), infra f\textsuperscript{o}~88~v\textsuperscript{o} (p.~\pageref{fol:177}), f\textsuperscript{o}~91~v\textsuperscript{o} (p.~\pageref{fol:183}), \textit{et passim}.\par\vspace{0.55mm}
\noindent\rule{38mm}{0.4pt}\par\vspace{0.5mm}
\noindent\hspace*{5mm}\parbox{\dimexpr\textwidth-5mm\relax}{\fontsize{6.6pt}{7.7pt}\selectfont\color{black!84}Après une tourmente de neige le mercredi 4 avril, le \textbf{Jeudi saint 5 avril 1917} offre une éclaircie printanière: traversant \textbf{Contay}, \textbf{Warloy-Baillon}, \textbf{Varennes-en-Croix} et \textbf{Acheux-en-Amiénois}, Hippolyte Combes constate l'abandon militaire complet de ces bourgs naguère fourmillants de troupes, puis assiste le soir au sermon de la Passion dans l'église Saint-Martin d'Hérissart.}\par}
\newpage
\phantomsection\label{fac:176}
% --- En-tête du fac-similé (Page impaire / droite : Photographie du manuscrit) ---
\noindent\rule{\textwidth}{0.5pt}\par\vspace{0.55mm}
\noindent\textbf{\fontsize{6.8pt}{7.8pt}\selectfont FAC-SIMILÉ EN VIS-À-VIS — FOLIO 88 RECTO} \hfill \textbf{\fontsize{6.8pt}{7.8pt}\selectfont AVRIL 1917}\par\vspace{0.25mm}
\noindent{\fontsize{6.8pt}{8.0pt}\selectfont\textit{Livre-Journal d'Hippolyte Combes} — 4 – 5 avril 1917 (Hérissart, Toutencourt, Acheux, Contay, Warloy, Varennes (Somme))}\par\vspace{0.55mm}
\noindent\rule{\textwidth}{0.35pt}\par\vspace*{\fill}
{\centering
  \setlength{\fboxsep}{1.2mm}%
  \setlength{\fboxrule}{0.4pt}%
  \fcolorbox{gray!55}{white}{%
    \includegraphics[width=0.95\textwidth,height=0.81\textheight,keepaspectratio]{\facsimiledir/20251126_111525.jpg}%
  }\par}
\vspace*{\fill}
\newpage

% =============================================================================
% VIS-A-VIS SPREAD : FOLIO 88 VERSO (20251126_111537.jpg)
% =============================================================================
\cleardoubleevenpage
\renewcommand{\currentmonthyear}{Avril 1917}
\renewcommand{\currentfoliolabel}{Folio 88 verso}
\phantomsection\label{fol:177}
% --- En-tête du folio (Page paire / gauche : Transcription, apparat & registres de bas de page) ---
\noindent\rule{\textwidth}{0.5pt}\par\vspace{0.55mm}
\noindent\textbf{\footnotesize FOLIO 88 VERSO} \hfill \textbf{\footnotesize AVRIL 1917}\par\vspace{0.25mm}
\noindent{\fontsize{6.9pt}{8.0pt}\selectfont Vendredi Saint : solde au Q.G. de la Ve Armée à Toutencourt et causerie sur l'Australie}\par\vspace{0.2mm}
\noindent{\fontsize{6.9pt}{8.0pt}\selectfont\textit{6 – 7 avril 1917} \hfill \textit{Hérissart, Toutencourt, Rubempré, Contay, Acheux, Léalvillers (Somme)}}\par\vspace{0.5mm}
\noindent\rule{\textwidth}{0.35pt}\par\vspace{0.8mm}
% --- Corps de la transcription diplomatique et apparat marginal en prolongement ---
{\renewcommand{\arraystretch}{0.87}%
\noindent\hspace*{3.5mm}\begin{tabular}{@{}p{87mm}@{\hspace{3.5mm}}p{39mm}@{}}
      \leavevmode\fontsize{6.9pt}{7.9pt}\selectfont\strut \underline{6 Vendr.} — \hfill{} \underline{Avril 1917 —}\strut & \smash{\parbox[t]{39mm}{\fontsize{6.2pt}{7.0pt}\selectfont\raggedright\color{black!88}\textbf{1 Vendr.} \textit{scil.} Vendredi.}} \\
      \leavevmode\fontsize{6.9pt}{7.9pt}\selectfont\strut Levé tôt. Après breakfast,\strut & \smash{\parbox[t]{39mm}{\fontsize{6.2pt}{7.0pt}\selectfont\raggedright\color{black!88}\textbf{2 breakfast} \textit{angl. in ms.}}} \\
      \leavevmode\fontsize{6.9pt}{7.9pt}\selectfont\strut conduit en auto à Toutencourt, 1 staff du bureau\strut &  \\
      \leavevmode\fontsize{6.9pt}{7.9pt}\selectfont\strut de la Mission fr., Vth Army H.Q., venu nous payer\strut & \smash{\parbox[t]{39mm}{\fontsize{6.2pt}{7.0pt}\selectfont\raggedright\color{black!88}\textbf{4 H.Q.} \textit{angl.} Headquarters.}} \\
      \leavevmode\llap{\fontsize{6.4pt}{7.3pt}\selectfont\color{darkgray}5\hspace{1.8mm}}\fontsize{6.9pt}{7.9pt}\selectfont\strut solde.\strut &  \\
      \leavevmode\fontsize{6.9pt}{7.9pt}\selectfont\strut Temps très beau. Délicieux promener rues ensoleillées\strut &  \\
      \leavevmode\fontsize{6.9pt}{7.9pt}\selectfont\strut de ce petit village aujourd'hui si tranquille, depuis\strut &  \\
      \leavevmode\fontsize{6.9pt}{7.9pt}\selectfont\strut départ H.Q. et troupes.\strut & \smash{\parbox[t]{39mm}{\fontsize{6.2pt}{7.0pt}\selectfont\raggedright\color{black!88}\textbf{8 H.Q.} \textit{angl.} Headquarters.}} \\
      \leavevmode\fontsize{6.9pt}{7.9pt}\selectfont\strut Revenu en auto à Hérissart. En reparti\strut &  \\
      \leavevmode\llap{\fontsize{6.4pt}{7.3pt}\selectfont\color{darkgray}10\hspace{1.8mm}}\fontsize{6.9pt}{7.9pt}\selectfont\strut avec off. pour Toutencourt \& Rubempré. Au retour,\strut & \smash{\parbox[t]{39mm}{\fontsize{6.2pt}{7.0pt}\selectfont\raggedright\color{black!88}\textbf{10 off.} \textit{abbr. pro} officier(s).}} \\
      \leavevmode\fontsize{6.9pt}{7.9pt}\selectfont\strut ciel nuageux. Bientôt bruine.\strut &  \\
      \leavevmode\fontsize{6.9pt}{7.9pt}\selectfont\strut Tout l'après-midi, liberté entière. Averses\strut &  \\
      \leavevmode\fontsize{6.9pt}{7.9pt}\selectfont\strut courtes, discontinues.\strut &  \\
      \leavevmode\fontsize{6.9pt}{7.9pt}\selectfont\strut After dinner, devant poêle rouge, causé lgt\strut & \smash{\parbox[t]{39mm}{\fontsize{6.2pt}{7.0pt}\selectfont\raggedright\color{black!88}\textbf{14 lgt} \textit{abbr. pro} longtemps. \textbf{14 dinner} \textit{angl. in ms.}}} \\
      \leavevmode\llap{\fontsize{6.4pt}{7.3pt}\selectfont\color{darkgray}15\hspace{1.8mm}}\fontsize{6.9pt}{7.9pt}\selectfont\strut de l'Australie avec M. Wentworth et lu pamphlets.\strut &  \\
      \leavevmode\fontsize{6.9pt}{7.9pt}\selectfont\strut [ Solde : 78 f. 75 c. — ]\strut & \smash{\parbox[t]{39mm}{\fontsize{6.2pt}{7.0pt}\selectfont\raggedright\color{black!88}\textbf{16} \textit{in marg. cum unco} (comptes). \textbf{16 f.} \textit{abbr. pro} franc(s). — \textbf{17 Sam.} \textit{scil.} Samedi. \textbf{17 gr.} \textit{abbr. pro} grand.}} \\
      \cline{1-1}
      \leavevmode\fontsize{6.9pt}{7.9pt}\selectfont\strut \underline{7 Sam.} — Temps gr. froid. Ciel resté\strut &  \\
      \leavevmode\fontsize{6.9pt}{7.9pt}\selectfont\strut couvert, gros nuages. Vent violent.\strut &  \\
      \leavevmode\fontsize{6.9pt}{7.9pt}\selectfont\strut Matinée, courses à Contay, Acheux \&\strut &  \\
      \leavevmode\llap{\fontsize{6.4pt}{7.3pt}\selectfont\color{darkgray}20\hspace{1.8mm}}\fontsize{6.9pt}{7.9pt}\selectfont\strut retour par route Léalvillers-Toutencourt \& Rubempré.\strut &  \\
      \leavevmode\fontsize{6.9pt}{7.9pt}\selectfont\strut Après-midi, liberté.\strut &  \\
      \leavevmode\fontsize{6.9pt}{7.9pt}\selectfont\strut Promené seul route Contay. Gdes rafales\strut &  \\
      \leavevmode\fontsize{6.9pt}{7.9pt}\selectfont\strut vent froid m'ont fait rentrer tôt.\strut &  \\
      \leavevmode\fontsize{6.9pt}{7.9pt}\selectfont\strut Admiré paysage —\strut &  \\
\end{tabular}}%
\par\vspace*{\fill}
% --- Registres de bas de page (1. Critique interne [si pertinent] ; 2. Loci paralleli en latin ; 3. Commentaire historique en retrait) ---
\noindent\rule{\textwidth}{0.5pt}\par\vspace{0.65mm}
{\fontsize{6.6pt}{7.7pt}\selectfont%
\noindent\normalfont\color{black!92} \textbf{3} (\textit{Toutencourt}) cf. supra f\textsuperscript{o}~49~v\textsuperscript{o} (p.~\pageref{fol:99}), f\textsuperscript{o}~88~r\textsuperscript{o} (p.~\pageref{fol:176}), infra f\textsuperscript{o}~89~r\textsuperscript{o} (p.~\pageref{fol:178}), f\textsuperscript{o}~91~r\textsuperscript{o} (p.~\pageref{fol:182}), \textit{et passim}. — \textbf{9} (\textit{Hérissart}) cf. supra f\textsuperscript{o}~47~v\textsuperscript{o} (p.~\pageref{fol:95}), f\textsuperscript{o}~88~r\textsuperscript{o} (p.~\pageref{fol:176}), infra f\textsuperscript{o}~89~r\textsuperscript{o} (p.~\pageref{fol:178}), f\textsuperscript{o}~89~v\textsuperscript{o} (p.~\pageref{fol:179}), \textit{et passim}. — \textbf{10} (\textit{Rubempré}) cf. supra f\textsuperscript{o}~47~v\textsuperscript{o} (p.~\pageref{fol:95}), f\textsuperscript{o}~87~v\textsuperscript{o} (p.~\pageref{fol:175}), infra f\textsuperscript{o}~89~r\textsuperscript{o} (p.~\pageref{fol:178}), f\textsuperscript{o}~90~r\textsuperscript{o} (p.~\pageref{fol:180}), \textit{et passim}.\par\vspace{0.55mm}
\noindent\rule{38mm}{0.4pt}\par\vspace{0.5mm}
\noindent\hspace*{5mm}\parbox{\dimexpr\textwidth-5mm\relax}{\fontsize{6.6pt}{7.7pt}\selectfont\color{black!84}Le Vendredi saint \textbf{6 avril 1917}, Hippolyte Combes note que le quartier général de la \textbf{V\textsuperscript{e} Armée britannique} (\textit{Fifth Army H.Q.} du général Sir Hubert Gough) et ses troupes ont quitté \textbf{Toutencourt} (l'état-major de la V\textsuperscript{e} Armée s'était en effet déplacé vers l'avant, à Albert, pour diriger les opérations devant la ligne Hindenburg et Bullecourt): un membre du personnel (\textit{staff}) de la Mission française de la V\textsuperscript{e} Armée vient néanmoins à Toutencourt payer la quinzaine des interprètes (78~fr.~75).}\par}
\newpage
\phantomsection\label{fac:177}
% --- En-tête du fac-similé (Page impaire / droite : Photographie du manuscrit) ---
\noindent\rule{\textwidth}{0.5pt}\par\vspace{0.55mm}
\noindent\textbf{\fontsize{6.8pt}{7.8pt}\selectfont FAC-SIMILÉ EN VIS-À-VIS — FOLIO 88 VERSO} \hfill \textbf{\fontsize{6.8pt}{7.8pt}\selectfont AVRIL 1917}\par\vspace{0.25mm}
\noindent{\fontsize{6.8pt}{8.0pt}\selectfont\textit{Livre-Journal d'Hippolyte Combes} — 6 – 7 avril 1917 (Hérissart, Toutencourt, Rubempré, Contay, Acheux, Léalvillers (Somme))}\par\vspace{0.55mm}
\noindent\rule{\textwidth}{0.35pt}\par\vspace*{\fill}
{\centering
  \setlength{\fboxsep}{1.2mm}%
  \setlength{\fboxrule}{0.4pt}%
  \fcolorbox{gray!55}{white}{%
    \includegraphics[width=0.95\textwidth,height=0.81\textheight,keepaspectratio]{\facsimiledir/20251126_111537.jpg}%
  }\par}
\vspace*{\fill}
\newpage

% =============================================================================
% VIS-A-VIS SPREAD : FOLIO 89 RECTO (20251126_111558.jpg)
% =============================================================================
\cleardoubleevenpage
\renewcommand{\currentmonthyear}{Avril 1917}
\renewcommand{\currentfoliolabel}{Folio 89 recto}
\phantomsection\label{fol:178}
% --- En-tête du folio (Page paire / gauche : Transcription, apparat & registres de bas de page) ---
\noindent\rule{\textwidth}{0.5pt}\par\vspace{0.55mm}
\noindent\textbf{\footnotesize FOLIO 89 RECTO} \hfill \textbf{\footnotesize AVRIL 1917}\par\vspace{0.25mm}
\noindent{\fontsize{6.9pt}{8.0pt}\selectfont Pâques 1917 à Hérissart et déclenchement de l'offensive d'Arras sous la neige}\par\vspace{0.2mm}
\noindent{\fontsize{6.9pt}{8.0pt}\selectfont\textit{8 – 10 avril 1917} \hfill \textit{Hérissart, Toutencourt, Harponville, Rubempré, Pierregot, Belle-Église, Fréchencourt, Contay (Somme)}}\par\vspace{0.5mm}
\noindent\rule{\textwidth}{0.35pt}\par\vspace{0.8mm}
% --- Corps de la transcription diplomatique et apparat marginal en prolongement ---
{\renewcommand{\arraystretch}{0.87}%
\noindent\hspace*{3.5mm}\begin{tabular}{@{}p{87mm}@{\hspace{3.5mm}}p{39mm}@{}}
      \leavevmode\fontsize{6.9pt}{7.9pt}\selectfont\strut \underline{Pâques — 8 —} \hfill{} \underline{Avril 1917}\strut &  \\
      \leavevmode\fontsize{6.9pt}{7.9pt}\selectfont\strut Matinée et 1/2 après-midi, temps mervt beau.\strut & \smash{\parbox[t]{39mm}{\fontsize{6.2pt}{7.0pt}\selectfont\raggedright\color{black!88}\textbf{2 mervt} \textit{abbr. pro} merveilleusement. \textbf{2 1/2} \textit{scil.} demi / moitié. — \textbf{4 ds} \textit{abbr. pro} dans.}} \\
      \leavevmode\fontsize{6.9pt}{7.9pt}\selectfont\strut Ciel d'1 bleu foncé. Arbres détachaient branches\strut &  \\
      \leavevmode\fontsize{6.9pt}{7.9pt}\selectfont\strut encore sans feuilles ds l'air lumineux \& limpide.\strut &  \\
      \leavevmode\llap{\fontsize{6.4pt}{7.3pt}\selectfont\color{darkgray}5\hspace{1.8mm}}\fontsize{6.9pt}{7.9pt}\selectfont\strut Chaleur assez forte.\strut &  \\
      \leavevmode\fontsize{6.9pt}{7.9pt}\selectfont\strut Matinée, à Toutencourt, chez notaire Bienaimé.\strut &  \\
      \leavevmode\fontsize{6.9pt}{7.9pt}\selectfont\strut Après-midi, à Harponville \& Rubempré — Beauquesne.\strut &  \\
      \leavevmode\fontsize{6.9pt}{7.9pt}\selectfont\strut Rentrés 5h s. Alors ciel couvert, gros nuages.\strut & \smash{\parbox[t]{39mm}{\fontsize{6.2pt}{7.0pt}\selectfont\raggedright\color{black!88}\textbf{8 s.} \textit{abbr. pro} soir.}} \\
      \leavevmode\fontsize{6.9pt}{7.9pt}\selectfont\strut Fin journée maussade. Visité famille Denis.\strut &  \\
      \leavevmode\llap{\fontsize{6.4pt}{7.3pt}\selectfont\color{darkgray}10\hspace{1.8mm}}\fontsize{6.9pt}{7.9pt}\selectfont\strut Nuit très noire.\strut &  \\
      \leavevmode\fontsize{6.9pt}{7.9pt}\selectfont\strut Promené 1 peu rues village abst désertes \& mortes.\strut & \smash{\parbox[t]{39mm}{\fontsize{6.2pt}{7.0pt}\selectfont\raggedright\color{black!88}\textbf{11 abst} \textit{abbr. pro} absolument.}} \\
      \cline{1-1}
      \leavevmode\fontsize{6.9pt}{7.9pt}\selectfont\strut \underline{9 Lundi.} — Temps très mauvais \& froid vif.\strut &  \\
      \leavevmode\fontsize{6.9pt}{7.9pt}\selectfont\strut Gdes averses, grêle et grésil. Rafales violentes.\strut &  \\
      \leavevmode\fontsize{6.9pt}{7.9pt}\selectfont\strut Courses matin \& soir : Pierregot \& Belle-Église.\strut &  \\
      \leavevmode\llap{\fontsize{6.4pt}{7.3pt}\selectfont\color{darkgray}15\hspace{1.8mm}}\fontsize{6.9pt}{7.9pt}\selectfont\strut Routes état lamentable. Auto rebondissait à cause\strut &  \\
      \leavevmode\fontsize{6.9pt}{7.9pt}\selectfont\strut des trous route, commotions très fortes —\strut &  \\
      \cline{1-1}
      \leavevmode\fontsize{6.9pt}{7.9pt}\selectfont\strut \underline{10 Mardi.} — Matinée orage, neige fine\strut &  \\
      \leavevmode\fontsize{6.9pt}{7.9pt}\selectfont\strut tombée abondante. Giboulées et rafales.\strut &  \\
      \leavevmode\fontsize{6.9pt}{7.9pt}\selectfont\strut Éclaircies de soleil sans durée.\strut &  \\
      \leavevmode\llap{\fontsize{6.4pt}{7.3pt}\selectfont\color{darkgray}20\hspace{1.8mm}}\fontsize{6.9pt}{7.9pt}\selectfont\strut Courses après-midi à Fréchencourt \& Contay.\strut &  \\
      \leavevmode\fontsize{6.9pt}{7.9pt}\selectfont\strut Temps très mauvais. Froid vif. C'est de\strut &  \\
      \leavevmode\fontsize{6.9pt}{7.9pt}\selectfont\strut nouveau l'hiver, hélas!\strut &  \\
\end{tabular}}%
\par\vspace*{\fill}
% --- Registres de bas de page (1. Critique interne [si pertinent] ; 2. Loci paralleli en latin ; 3. Commentaire historique en retrait) ---
\noindent\rule{\textwidth}{0.5pt}\par\vspace{0.65mm}
{\fontsize{6.6pt}{7.7pt}\selectfont%
\noindent\normalfont\color{black!92} \textbf{6} (\textit{Toutencourt}) cf. supra f\textsuperscript{o}~49~v\textsuperscript{o} (p.~\pageref{fol:99}), f\textsuperscript{o}~88~v\textsuperscript{o} (p.~\pageref{fol:177}), infra f\textsuperscript{o}~91~r\textsuperscript{o} (p.~\pageref{fol:182}), f\textsuperscript{o}~91~v\textsuperscript{o} (p.~\pageref{fol:183}), \textit{et passim}. — \textbf{7} (\textit{Rubempré}) cf. supra f\textsuperscript{o}~47~v\textsuperscript{o} (p.~\pageref{fol:95}), f\textsuperscript{o}~88~v\textsuperscript{o} (p.~\pageref{fol:177}), infra f\textsuperscript{o}~90~r\textsuperscript{o} (p.~\pageref{fol:180}), f\textsuperscript{o}~90~v\textsuperscript{o} (p.~\pageref{fol:181}), \textit{et passim}. — \textbf{7} (\textit{Harponville}) cf. supra f\textsuperscript{o}~55~r\textsuperscript{o} (p.~\pageref{fol:110}), infra f\textsuperscript{o}~92~r\textsuperscript{o} (p.~\pageref{fol:184}).\par\vspace{0.55mm}
\noindent\rule{38mm}{0.4pt}\par\vspace{0.5mm}
\noindent\hspace*{5mm}\parbox{\dimexpr\textwidth-5mm\relax}{\fontsize{6.6pt}{7.7pt}\selectfont\color{black!84}Le dimanche de \textbf{Pâques 8 avril 1917}, après une visite matinale au notaire \textbf{M\textsuperscript{e}~Bienaimé} à Toutencourt et une tournée à Harponville, Rubempré et Beauquesne, Hippolyte Combes passe pour la première fois la soirée chez la \textbf{famille Denis} à Hérissart (qu'il fréquentera assidûment jusqu'à la fin d'avril, f\textsuperscript{o}~90~r\textsuperscript{o}, f\textsuperscript{o}~91~r\textsuperscript{o}, f\textsuperscript{o}~91~v\textsuperscript{o}, f\textsuperscript{o}~92~r\textsuperscript{o}).}\par}
\newpage
\phantomsection\label{fac:178}
% --- En-tête du fac-similé (Page impaire / droite : Photographie du manuscrit) ---
\noindent\rule{\textwidth}{0.5pt}\par\vspace{0.55mm}
\noindent\textbf{\fontsize{6.8pt}{7.8pt}\selectfont FAC-SIMILÉ EN VIS-À-VIS — FOLIO 89 RECTO} \hfill \textbf{\fontsize{6.8pt}{7.8pt}\selectfont AVRIL 1917}\par\vspace{0.25mm}
\noindent{\fontsize{6.8pt}{8.0pt}\selectfont\textit{Livre-Journal d'Hippolyte Combes} — 8 – 10 avril 1917 (Hérissart, Toutencourt, Harponville, Rubempré, Pierregot, Belle-Église, Fréchencourt, Contay (Somme))}\par\vspace{0.55mm}
\noindent\rule{\textwidth}{0.35pt}\par\vspace*{\fill}
{\centering
  \setlength{\fboxsep}{1.2mm}%
  \setlength{\fboxrule}{0.4pt}%
  \fcolorbox{gray!55}{white}{%
    \includegraphics[width=0.95\textwidth,height=0.81\textheight,keepaspectratio]{\facsimiledir/20251126_111558.jpg}%
  }\par}
\vspace*{\fill}
\newpage

% =============================================================================
% VIS-A-VIS SPREAD : FOLIO 89 VERSO (20251126_111614.jpg)
% =============================================================================
\cleardoubleevenpage
\renewcommand{\currentmonthyear}{Avril 1917}
\renewcommand{\currentfoliolabel}{Folio 89 verso}
\phantomsection\label{fol:179}
% --- En-tête du folio (Page paire / gauche : Transcription, apparat & registres de bas de page) ---
\noindent\rule{\textwidth}{0.5pt}\par\vspace{0.55mm}
\noindent\textbf{\footnotesize FOLIO 89 VERSO} \hfill \textbf{\footnotesize AVRIL 1917}\par\vspace{0.25mm}
\noindent{\fontsize{6.9pt}{8.0pt}\selectfont Tempête de neige du 11 avril, emplettes à Doullens et entretien avec le lieutenant Gilbert}\par\vspace{0.2mm}
\noindent{\fontsize{6.9pt}{8.0pt}\selectfont\textit{11 – 13 avril 1917} \hfill \textit{Hérissart, Doullens, Warloy, Belle-Église (Somme)}}\par\vspace{0.5mm}
\noindent\rule{\textwidth}{0.35pt}\par\vspace{0.8mm}
% --- Corps de la transcription diplomatique et apparat marginal en prolongement ---
{\renewcommand{\arraystretch}{0.83}%
\noindent\hspace*{3.5mm}\begin{tabular}{@{}p{87mm}@{\hspace{3.5mm}}p{39mm}@{}}
      \leavevmode\fontsize{6.6pt}{7.5pt}\selectfont\strut \underline{Merc. 11} — \hfill{} \underline{— Avril 1917 —}\strut & \smash{\parbox[t]{39mm}{\fontsize{6.2pt}{7.0pt}\selectfont\raggedright\color{black!88}\textbf{1 Merc.} \textit{scil.} Mercredi.}} \\
      \leavevmode\fontsize{6.6pt}{7.5pt}\selectfont\strut Nuit très froide.\strut &  \\
      \leavevmode\fontsize{6.6pt}{7.5pt}\selectfont\strut Matinée, course à Doullens.\strut &  \\
      \leavevmode\fontsize{6.6pt}{7.5pt}\selectfont\strut Après-midi, course à Warloy.\strut &  \\
      \leavevmode\llap{\fontsize{6.4pt}{7.3pt}\selectfont\color{darkgray}5\hspace{1.8mm}}\fontsize{6.6pt}{7.5pt}\selectfont\strut Tout le jour et le soir davantage, neigé lgt\strut & \smash{\parbox[t]{39mm}{\fontsize{6.2pt}{7.0pt}\selectfont\raggedright\color{black!88}\textbf{5 lgt} \textit{abbr. pro} longtemps.}} \\
      \leavevmode\fontsize{6.6pt}{7.5pt}\selectfont\strut et abondamment. Routes, champs, toits, arbres, étaient\strut &  \\
      \leavevmode\fontsize{6.6pt}{7.5pt}\selectfont\strut blancs.\strut &  \\
      \leavevmode\fontsize{6.6pt}{7.5pt}\selectfont\strut Nuit, rafales très violentes.\strut &  \\
      \leavevmode\fontsize{6.6pt}{7.5pt}\selectfont\strut [ \underline{A Doullens} : ceinture cuir : (10 f. 50 + 14 f. = 24 f. 50)\strut & \smash{\parbox[t]{39mm}{\fontsize{6.2pt}{7.0pt}\selectfont\raggedright\color{black!88}\textbf{9} \textit{in marg. cum unco} (comptes). \textbf{9 f.} \textit{abbr. pro} franc(s). — \textbf{10 f.} \textit{abbr. pro} franc(s). \textbf{10 canteen} \textit{angl. in ms.} — \textbf{11 Cartes ples} \textit{abbr. pro} Cartes postales. \textbf{11 c.} \textit{abbr. pro} centimes. — \textbf{12 Tre} \textit{abbr. pro} Teinture. \textbf{12 f.} \textit{abbr. pro} franc(s). — \textbf{13 f.} \textit{abbr. pro} franc(s). \textbf{13 c.} \textit{abbr. pro} centimes. — \textbf{19 Ven.} \textit{scil.} Vendredi.}} \\
      \leavevmode\llap{\fontsize{6.4pt}{7.3pt}\selectfont\color{darkgray}10\hspace{1.8mm}}\fontsize{6.6pt}{7.5pt}\selectfont\strut (canteen) Boîte de cacao : 1 f.\strut &  \\
      \leavevmode\fontsize{6.6pt}{7.5pt}\selectfont\strut 2 Cartes géog. (Somme) : 60 c. — Cartes ples : 60 c. + 10 c.\strut &  \\
      \leavevmode\fontsize{6.6pt}{7.5pt}\selectfont\strut Varia : 1 f. — — 2 semelles liège : 40 c. — Tre d'iode : 35 c.\strut &  \\
      \leavevmode\fontsize{6.6pt}{7.5pt}\selectfont\strut 2 boîtes de chlorate potasse : 1 f. — 1 boîte pastilles menthe : 75 c. ]\strut &  \\
      \cline{1-1}
      \leavevmode\fontsize{6.6pt}{7.5pt}\selectfont\strut \underline{12 Jeudi.} — Nuit gelé un peu. Neige durcie.\strut &  \\
      \leavevmode\llap{\fontsize{6.4pt}{7.3pt}\selectfont\color{darkgray}15\hspace{1.8mm}}\fontsize{6.6pt}{7.5pt}\selectfont\strut Dans la journée, dégel fondu neige. Routes\strut &  \\
      \leavevmode\fontsize{6.6pt}{7.5pt}\selectfont\strut transformées en marécages boueux.\strut &  \\
      \leavevmode\fontsize{6.6pt}{7.5pt}\selectfont\strut Matinée, course à Belle-Église.\strut &  \\
      \leavevmode\fontsize{6.6pt}{7.5pt}\selectfont\strut Après-midi, liberté entière —\strut &  \\
      \cline{1-1}
      \leavevmode\fontsize{6.6pt}{7.5pt}\selectfont\strut \underline{13 Ven.} —\strut &  \\
      \leavevmode\llap{\fontsize{6.4pt}{7.3pt}\selectfont\color{darkgray}20\hspace{1.8mm}}\fontsize{6.6pt}{7.5pt}\selectfont\strut Matinée, très beau temps, mais un peu\strut &  \\
      \leavevmode\fontsize{6.6pt}{7.5pt}\selectfont\strut froid. Après-midi, ciel de nouveau nuageux\strut &  \\
      \leavevmode\fontsize{6.6pt}{7.5pt}\selectfont\strut et pluie menaçante.\strut &  \\
      \leavevmode\fontsize{6.6pt}{7.5pt}\selectfont\strut Liberté entière tout le jour.\strut &  \\
      \leavevmode\fontsize{6.6pt}{7.5pt}\selectfont\strut Après-midi, long talk avec Secd Lt Gilbert,\strut & \smash{\parbox[t]{39mm}{\fontsize{6.2pt}{7.0pt}\selectfont\raggedright\color{black!88}\textbf{24 Lt} \textit{abbr. pro} Lieutenant. \textbf{24 long talk} \textit{angl. in ms.}}} \\
      \leavevmode\llap{\fontsize{6.4pt}{7.3pt}\selectfont\color{darkgray}25\hspace{1.8mm}}\fontsize{6.6pt}{7.5pt}\selectfont\strut de Jersey, il est très instruit et poète.\strut &  \\
\end{tabular}}%
\par\vspace*{\fill}
% --- Registres de bas de page (1. Critique interne [si pertinent] ; 2. Loci paralleli en latin ; 3. Commentaire historique en retrait) ---
\noindent\rule{\textwidth}{0.5pt}\par\vspace{0.65mm}
{\fontsize{6.6pt}{7.7pt}\selectfont%
\noindent\normalfont\color{black!92} \textbf{3} (\textit{Doullens}) cf. supra f\textsuperscript{o}~25~v\textsuperscript{o} (p.~\pageref{fol:51}), f\textsuperscript{o}~81~v\textsuperscript{o} (p.~\pageref{fol:163}), \textit{et passim}. — \textbf{4} (\textit{Warloy}) cf. supra f\textsuperscript{o}~37~v\textsuperscript{o} (p.~\pageref{fol:75}), f\textsuperscript{o}~88~r\textsuperscript{o} (p.~\pageref{fol:176}), infra f\textsuperscript{o}~90~v\textsuperscript{o} (p.~\pageref{fol:181}), f\textsuperscript{o}~93~r\textsuperscript{o} (p.~\pageref{fol:186}), \textit{et passim}. — \textbf{10} (\textit{canteen}) cf. supra f\textsuperscript{o}~21~v\textsuperscript{o} (p.~\pageref{fol:43}), f\textsuperscript{o}~83~v\textsuperscript{o} (p.~\pageref{fol:167}), \textit{et passim}.\par\vspace{0.55mm}
\noindent\rule{38mm}{0.4pt}\par\vspace{0.5mm}
\noindent\hspace*{5mm}\parbox{\dimexpr\textwidth-5mm\relax}{\fontsize{6.6pt}{7.7pt}\selectfont\color{black!84}Après la tempête de neige du 11 avril 1917 (jour de la première attaque anglo-australienne devant Bullecourt) et ses achats à \textbf{Doullens}, Hippolyte Combes noue le vendredi 13 avril une amitié intellectuelle avec le nouveau sous-lieutenant britannique du service des achats, le \textbf{\textit{Second Lieutenant} Gilbert}, originaire de l'île anglo-normande de \textbf{Jersey}, francophone, cultivé et poète.}\par}
\newpage
\phantomsection\label{fac:179}
% --- En-tête du fac-similé (Page impaire / droite : Photographie du manuscrit) ---
\noindent\rule{\textwidth}{0.5pt}\par\vspace{0.55mm}
\noindent\textbf{\fontsize{6.8pt}{7.8pt}\selectfont FAC-SIMILÉ EN VIS-À-VIS — FOLIO 89 VERSO} \hfill \textbf{\fontsize{6.8pt}{7.8pt}\selectfont AVRIL 1917}\par\vspace{0.25mm}
\noindent{\fontsize{6.8pt}{8.0pt}\selectfont\textit{Livre-Journal d'Hippolyte Combes} — 11 – 13 avril 1917 (Hérissart, Doullens, Warloy, Belle-Église (Somme))}\par\vspace{0.55mm}
\noindent\rule{\textwidth}{0.35pt}\par\vspace*{\fill}
{\centering
  \setlength{\fboxsep}{1.2mm}%
  \setlength{\fboxrule}{0.4pt}%
  \fcolorbox{gray!55}{white}{%
    \includegraphics[width=0.95\textwidth,height=0.81\textheight,keepaspectratio]{\facsimiledir/20251126_111614.jpg}%
  }\par}
\vspace*{\fill}
\newpage

% =============================================================================
% VIS-A-VIS SPREAD : FOLIO 90 RECTO (20251126_111630.jpg)
% =============================================================================
\cleardoubleevenpage
\renewcommand{\currentmonthyear}{Avril 1917}
\renewcommand{\currentfoliolabel}{Folio 90 recto}
\phantomsection\label{fol:180}
% --- En-tête du folio (Page paire / gauche : Transcription, apparat & registres de bas de page) ---
\noindent\rule{\textwidth}{0.5pt}\par\vspace{0.55mm}
\noindent\textbf{\footnotesize FOLIO 90 RECTO} \hfill \textbf{\footnotesize AVRIL 1917}\par\vspace{0.25mm}
\noindent{\fontsize{6.9pt}{8.0pt}\selectfont Soirée chez la famille Denis sous l'orage et retour des giboulées}\par\vspace{0.2mm}
\noindent{\fontsize{6.9pt}{8.0pt}\selectfont\textit{14 – 17 avril 1917} \hfill \textit{Hérissart, Pierregot, Fréchencourt, Rubempré (Somme)}}\par\vspace{0.5mm}
\noindent\rule{\textwidth}{0.35pt}\par\vspace{0.8mm}
% --- Corps de la transcription diplomatique et apparat marginal en prolongement ---
{\renewcommand{\arraystretch}{0.87}%
\noindent\hspace*{3.5mm}\begin{tabular}{@{}p{87mm}@{\hspace{3.5mm}}p{39mm}@{}}
      \leavevmode\fontsize{6.9pt}{7.9pt}\selectfont\strut \underline{14 Sam.} \hfill{} \underline{— Avril 1917 —}\strut & \smash{\parbox[t]{39mm}{\fontsize{6.2pt}{7.0pt}\selectfont\raggedright\color{black!88}\textbf{1 Sam.} \textit{scil.} Samedi.}} \\
      \leavevmode\fontsize{6.9pt}{7.9pt}\selectfont\strut Enfin, très belle journée. Température se réchauffe,\strut &  \\
      \leavevmode\fontsize{6.9pt}{7.9pt}\selectfont\strut soleil ardent.\strut &  \\
      \leavevmode\fontsize{6.9pt}{7.9pt}\selectfont\strut Après-midi seulement, course à Pierregot. Matin\strut &  \\
      \leavevmode\llap{\fontsize{6.4pt}{7.3pt}\selectfont\color{darkgray}5\hspace{1.8mm}}\fontsize{6.9pt}{7.9pt}\selectfont\strut \& soirée, liberté entière.\strut &  \\
      \cline{1-1}
      \leavevmode\fontsize{6.9pt}{7.9pt}\selectfont\strut \underline{15 Dim.} — Pluie fine et tiède.\strut & \smash{\parbox[t]{39mm}{\fontsize{6.2pt}{7.0pt}\selectfont\raggedright\color{black!88}\textbf{6 Dim.} \textit{scil.} Dimanche.}} \\
      \leavevmode\fontsize{6.9pt}{7.9pt}\selectfont\strut Tout le jour, averses, temps froid.\strut &  \\
      \leavevmode\fontsize{6.9pt}{7.9pt}\selectfont\strut Matin \& soir, repos complet.\strut &  \\
      \leavevmode\fontsize{6.9pt}{7.9pt}\selectfont\strut Resté plusieurs heures à office, à lire —\strut &  \\
      \leavevmode\llap{\fontsize{6.4pt}{7.3pt}\selectfont\color{darkgray}10\hspace{1.8mm}}\fontsize{6.9pt}{7.9pt}\selectfont\strut [ Washing : 60 c. — ]\strut & \smash{\parbox[t]{39mm}{\fontsize{6.2pt}{7.0pt}\selectfont\raggedright\color{black!88}\textbf{10} \textit{in marg. cum unco} (comptes). \textbf{10 c.} \textit{abbr. pro} centimes.}} \\
      \cline{1-1}
      \leavevmode\fontsize{6.9pt}{7.9pt}\selectfont\strut \underline{16 Lundi.}\strut &  \\
      \leavevmode\fontsize{6.9pt}{7.9pt}\selectfont\strut Après-midi, course à Pierregot.\strut &  \\
      \leavevmode\fontsize{6.9pt}{7.9pt}\selectfont\strut Tout le jour, et le soir surtout, pluies très fortes.\strut &  \\
      \leavevmode\fontsize{6.9pt}{7.9pt}\selectfont\strut Passé soirée chez la famille Denis. Rentré\strut &  \\
      \leavevmode\llap{\fontsize{6.4pt}{7.3pt}\selectfont\color{darkgray}15\hspace{1.8mm}}\fontsize{6.9pt}{7.9pt}\selectfont\strut home sous 1 pluie battante, torrentielle,\strut & \smash{\parbox[t]{39mm}{\fontsize{6.2pt}{7.0pt}\selectfont\raggedright\color{black!88}\textbf{15 home} \textit{angl. in ms.}}} \\
      \leavevmode\fontsize{6.9pt}{7.9pt}\selectfont\strut et ds nuit très noire.\strut & \smash{\parbox[t]{39mm}{\fontsize{6.2pt}{7.0pt}\selectfont\raggedright\color{black!88}\textbf{16 ds} \textit{abbr. pro} dans.}} \\
      \cline{1-1}
      \leavevmode\fontsize{6.9pt}{7.9pt}\selectfont\strut \underline{17 Mardi.} — Temps toujours affreux. Averses\strut &  \\
      \leavevmode\fontsize{6.9pt}{7.9pt}\selectfont\strut intermittentes. Après-midi, neige et grêle,\strut &  \\
      \leavevmode\fontsize{6.9pt}{7.9pt}\selectfont\strut tour à tour. Froid de plus en plus vif.\strut &  \\
      \leavevmode\llap{\fontsize{6.4pt}{7.3pt}\selectfont\color{darkgray}20\hspace{1.8mm}}\fontsize{6.9pt}{7.9pt}\selectfont\strut C'est de nouveau l'hiver —\strut &  \\
      \leavevmode\fontsize{6.9pt}{7.9pt}\selectfont\strut Matin \& après-midi, courses en auto, à\strut &  \\
      \leavevmode\fontsize{6.9pt}{7.9pt}\selectfont\strut Fréchencourt \& Rubempré.\strut &  \\
\end{tabular}}%
\par\vspace*{\fill}
% --- Registres de bas de page (1. Critique interne [si pertinent] ; 2. Loci paralleli en latin ; 3. Commentaire historique en retrait) ---
\noindent\rule{\textwidth}{0.5pt}\par\vspace{0.65mm}
{\fontsize{6.6pt}{7.7pt}\selectfont%
\noindent\normalfont\color{black!92} \textbf{4} (\textit{Pierregot}) cf. supra f\textsuperscript{o}~52~r\textsuperscript{o} (p.~\pageref{fol:104}), f\textsuperscript{o}~89~r\textsuperscript{o} (p.~\pageref{fol:178}), infra f\textsuperscript{o}~91~r\textsuperscript{o} (p.~\pageref{fol:182}), f\textsuperscript{o}~91~v\textsuperscript{o} (p.~\pageref{fol:183}), \textit{et passim}. — \textbf{22} (\textit{Rubempré}) cf. supra f\textsuperscript{o}~47~v\textsuperscript{o} (p.~\pageref{fol:95}), f\textsuperscript{o}~89~r\textsuperscript{o} (p.~\pageref{fol:178}), infra f\textsuperscript{o}~90~v\textsuperscript{o} (p.~\pageref{fol:181}), f\textsuperscript{o}~92~r\textsuperscript{o} (p.~\pageref{fol:184}), \textit{et passim}.\par\vspace{0.55mm}
\noindent\rule{38mm}{0.4pt}\par\vspace{0.5mm}
\noindent\hspace*{5mm}\parbox{\dimexpr\textwidth-5mm\relax}{\fontsize{6.6pt}{7.7pt}\selectfont\color{black!84}Les tournées des 14--17 avril 1917 se limitent aux villages proches d'Hérissart (\textbf{Pierregot}, \textbf{Fréchencourt}, \textbf{Rubempré}) sous les pluies torrentielles du lundi 16 avril (jour du déclenchement dans l'Aisne de l'offensive française du Chemin des Dames) et les giboulées de neige et de grêle du mardi 17 avril; le lundi soir, Hippolyte Combes veille chez ses amis d'Hérissart, la \textbf{famille Denis}.}\par}
\newpage
\phantomsection\label{fac:180}
% --- En-tête du fac-similé (Page impaire / droite : Photographie du manuscrit) ---
\noindent\rule{\textwidth}{0.5pt}\par\vspace{0.55mm}
\noindent\textbf{\fontsize{6.8pt}{7.8pt}\selectfont FAC-SIMILÉ EN VIS-À-VIS — FOLIO 90 RECTO} \hfill \textbf{\fontsize{6.8pt}{7.8pt}\selectfont AVRIL 1917}\par\vspace{0.25mm}
\noindent{\fontsize{6.8pt}{8.0pt}\selectfont\textit{Livre-Journal d'Hippolyte Combes} — 14 – 17 avril 1917 (Hérissart, Pierregot, Fréchencourt, Rubempré (Somme))}\par\vspace{0.55mm}
\noindent\rule{\textwidth}{0.35pt}\par\vspace*{\fill}
{\centering
  \setlength{\fboxsep}{1.2mm}%
  \setlength{\fboxrule}{0.4pt}%
  \fcolorbox{gray!55}{white}{%
    \includegraphics[width=0.95\textwidth,height=0.81\textheight,keepaspectratio]{\facsimiledir/20251126_111630.jpg}%
  }\par}
\vspace*{\fill}
\newpage

% =============================================================================
% VIS-A-VIS SPREAD : FOLIO 90 VERSO (20251126_111644.jpg)
% =============================================================================
\cleardoubleevenpage
\renewcommand{\currentmonthyear}{Avril 1917}
\renewcommand{\currentfoliolabel}{Folio 90 verso}
\phantomsection\label{fol:181}
% --- En-tête du folio (Page paire / gauche : Transcription, apparat & registres de bas de page) ---
\noindent\rule{\textwidth}{0.5pt}\par\vspace{0.55mm}
\noindent\textbf{\footnotesize FOLIO 90 VERSO} \hfill \textbf{\footnotesize AVRIL 1917}\par\vspace{0.25mm}
\noindent{\fontsize{6.9pt}{8.0pt}\selectfont Retour du lieutenant Gilbert d'Abbeville et voiture désembourbée par les Australiens à Buire}\par\vspace{0.2mm}
\noindent{\fontsize{6.9pt}{8.0pt}\selectfont\textit{18 – 20 avril 1917} \hfill \textit{Hérissart, Bresle, Lavieville, Hénencourt, Warloy, Contay, Rubempré, Senlis, Albert, Buire-sur-l'Ancre, Ribemont (Somme)}}\par\vspace{0.5mm}
\noindent\rule{\textwidth}{0.35pt}\par\vspace{0.8mm}
% --- Corps de la transcription diplomatique et apparat marginal en prolongement ---
{\renewcommand{\arraystretch}{0.87}%
\noindent\hspace*{3.5mm}\begin{tabular}{@{}p{87mm}@{\hspace{3.5mm}}p{39mm}@{}}
      \leavevmode\fontsize{6.9pt}{7.9pt}\selectfont\strut \underline{18 Merc.} — \hfill{} \underline{— Avril 1917 —}\strut & \smash{\parbox[t]{39mm}{\fontsize{6.2pt}{7.0pt}\selectfont\raggedright\color{black!88}\textbf{1 Merc.} \textit{scil.} Mercredi.}} \\
      \leavevmode\fontsize{6.9pt}{7.9pt}\selectfont\strut Tout le jour, pluie fine.\strut &  \\
      \leavevmode\fontsize{6.9pt}{7.9pt}\selectfont\strut Routes, état lamentable.\strut &  \\
      \leavevmode\fontsize{6.9pt}{7.9pt}\selectfont\strut Notre auto, aux workshops. Repos complet.\strut & \smash{\parbox[t]{39mm}{\fontsize{6.2pt}{7.0pt}\selectfont\raggedright\color{black!88}\textbf{4 workshops} \textit{angl. in ms.}}} \\
      \cline{1-1}
      \leavevmode\llap{\fontsize{6.4pt}{7.3pt}\selectfont\color{darkgray}5\hspace{1.8mm}}\fontsize{6.9pt}{7.9pt}\selectfont\strut \underline{19 Jeudi.}\strut &  \\
      \leavevmode\fontsize{6.9pt}{7.9pt}\selectfont\strut Tout le jour, ciel couvert. Qq averses.\strut & \smash{\parbox[t]{39mm}{\fontsize{6.2pt}{7.0pt}\selectfont\raggedright\color{black!88}\textbf{6 Qq} \textit{abbr. pro} Quelque(s).}} \\
      \leavevmode\fontsize{6.9pt}{7.9pt}\selectfont\strut Temps très doux, presque tiède.\strut &  \\
      \leavevmode\fontsize{6.9pt}{7.9pt}\selectfont\strut Après-midi, courses en auto : Bresle, Lavieville,\strut &  \\
      \leavevmode\fontsize{6.9pt}{7.9pt}\selectfont\strut Hénencourt, Warloy et Contay. Revu avec plaisir\strut &  \\
      \leavevmode\llap{\fontsize{6.4pt}{7.3pt}\selectfont\color{darkgray}10\hspace{1.8mm}}\fontsize{6.9pt}{7.9pt}\selectfont\strut les 3 1ers villages, où pas retourné depuis fin\strut &  \\
      \leavevmode\fontsize{6.9pt}{7.9pt}\selectfont\strut juillet 1915. Aujourd'hui, vides et déserts.\strut &  \\
      \leavevmode\fontsize{6.9pt}{7.9pt}\selectfont\strut Lieut. Gilbert, retour d'Abbeville, venu passer soirée\strut &  \\
      \leavevmode\fontsize{6.9pt}{7.9pt}\selectfont\strut avec nous. Charmant causeur, très intéressant.\strut &  \\
      \cline{1-1}
      \leavevmode\fontsize{6.9pt}{7.9pt}\selectfont\strut \underline{20 Vendredi.} — Enfin, assez belle journée. Ciel\strut &  \\
      \leavevmode\llap{\fontsize{6.4pt}{7.3pt}\selectfont\color{darkgray}15\hspace{1.8mm}}\fontsize{6.9pt}{7.9pt}\selectfont\strut resté couvert, mais temps très doux, même chaud.\strut &  \\
      \leavevmode\fontsize{6.9pt}{7.9pt}\selectfont\strut Matinée, course à Rubempré, Senlis et\strut &  \\
      \leavevmode\fontsize{6.9pt}{7.9pt}\selectfont\strut Warloy. Après-midi, venus à Hénencourt,\strut &  \\
      \leavevmode\fontsize{6.9pt}{7.9pt}\selectfont\strut Albert et Buire [ajout souscrit : Ribemont]. Au retour, stuck, à la\strut & \smash{\parbox[t]{39mm}{\fontsize{6.2pt}{7.0pt}\selectfont\raggedright\color{black!88}\textbf{18 stuck} \textit{angl. in ms.} Ribemont ajouté sous Buire, sans parenthèses autographes.}} \\
      \leavevmode\fontsize{6.9pt}{7.9pt}\selectfont\strut sortie du village de Buire. Aidés par des\strut &  \\
      \leavevmode\llap{\fontsize{6.4pt}{7.3pt}\selectfont\color{darkgray}20\hspace{1.8mm}}\fontsize{6.9pt}{7.9pt}\selectfont\strut australiens, à dégager l'auto. Quels types!\strut &  \\
      \leavevmode\fontsize{6.9pt}{7.9pt}\selectfont\strut A Ribemont, revu avec plaisir les Colart,\strut &  \\
      \leavevmode\fontsize{6.9pt}{7.9pt}\selectfont\strut qui m'ont fait meilleur accueil —\strut &  \\
\end{tabular}}%
\par\vspace*{\fill}
% --- Registres de bas de page (1. Critique interne [si pertinent] ; 2. Loci paralleli en latin ; 3. Commentaire historique en retrait) ---
\noindent\rule{\textwidth}{0.5pt}\par\vspace{0.65mm}
{\fontsize{6.6pt}{7.7pt}\selectfont%
\noindent\normalfont\color{black!92} \textbf{9} (\textit{Warloy}) cf. supra f\textsuperscript{o}~37~v\textsuperscript{o} (p.~\pageref{fol:75}), f\textsuperscript{o}~89~v\textsuperscript{o} (p.~\pageref{fol:179}), infra f\textsuperscript{o}~93~r\textsuperscript{o} (p.~\pageref{fol:186}), \textit{et passim}. — \textbf{9} (\textit{Contay}) cf. supra f\textsuperscript{o}~49~r\textsuperscript{o} (p.~\pageref{fol:98}), f\textsuperscript{o}~89~r\textsuperscript{o} (p.~\pageref{fol:178}), infra f\textsuperscript{o}~91~r\textsuperscript{o} (p.~\pageref{fol:182}), f\textsuperscript{o}~92~r\textsuperscript{o} (p.~\pageref{fol:184}), \textit{et passim}. — \textbf{12} (\textit{Lt Gilbert}) cf. supra f\textsuperscript{o}~89~v\textsuperscript{o} (p.~\pageref{fol:179}), infra f\textsuperscript{o}~92~v\textsuperscript{o} (p.~\pageref{fol:185}), f\textsuperscript{o}~93~r\textsuperscript{o} (p.~\pageref{fol:186}).\par\vspace{0.55mm}
\noindent\rule{38mm}{0.4pt}\par\vspace{0.5mm}
\noindent\hspace*{5mm}\parbox{\dimexpr\textwidth-5mm\relax}{\fontsize{6.6pt}{7.7pt}\selectfont\color{black!84}Le jeudi 19 avril 1917, la tournée à \textbf{Bresle}, \textbf{Laviéville} et \textbf{Hénencourt} (Somme, à l'ouest d'Albert) fournit un repère autobiographique précieux: Hippolyte Combes n'était pas retourné dans ces trois villages depuis la \textbf{fin juillet 1915}, époque où la \textbf{18\textsuperscript{e} Division britannique}, tout juste débarquée en France et rassemblée autour de Flesselles (voir sa lettre du 21 février 1916, Annexe~I, doc.~1), venait prendre en charge le secteur d'Albert.}\par}
\newpage
\phantomsection\label{fac:181}
% --- En-tête du fac-similé (Page impaire / droite : Photographie du manuscrit) ---
\noindent\rule{\textwidth}{0.5pt}\par\vspace{0.55mm}
\noindent\textbf{\fontsize{6.8pt}{7.8pt}\selectfont FAC-SIMILÉ EN VIS-À-VIS — FOLIO 90 VERSO} \hfill \textbf{\fontsize{6.8pt}{7.8pt}\selectfont AVRIL 1917}\par\vspace{0.25mm}
\noindent{\fontsize{6.8pt}{8.0pt}\selectfont\textit{Livre-Journal d'Hippolyte Combes} — 18 – 20 avril 1917 (Hérissart, Bresle, Lavieville, Hénencourt, Warloy, Contay, Rubempré, Senlis, Albert, Buire-sur-l'Ancre, Ribemont (Somme))}\par\vspace{0.55mm}
\noindent\rule{\textwidth}{0.35pt}\par\vspace*{\fill}
{\centering
  \setlength{\fboxsep}{1.2mm}%
  \setlength{\fboxrule}{0.4pt}%
  \fcolorbox{gray!55}{white}{%
    \includegraphics[width=0.95\textwidth,height=0.81\textheight,keepaspectratio]{\facsimiledir/20251126_111644.jpg}%
  }\par}
\vspace*{\fill}
\newpage

% =============================================================================
% VIS-A-VIS SPREAD : FOLIO 91 RECTO (20251126_111702.jpg)
% =============================================================================
\cleardoubleevenpage
\renewcommand{\currentmonthyear}{Avril 1917}
\renewcommand{\currentfoliolabel}{Folio 91 recto}
\phantomsection\label{fol:182}
% --- En-tête du folio (Page paire / gauche : Transcription, apparat & registres de bas de page) ---
\noindent\rule{\textwidth}{0.5pt}\par\vspace{0.55mm}
\noindent\textbf{\footnotesize FOLIO 91 RECTO} \hfill \textbf{\footnotesize AVRIL 1917}\par\vspace{0.25mm}
\noindent{\fontsize{6.9pt}{8.0pt}\selectfont Violente canonnade vers Arras, moulin de Pierregot et batterie antiaérienne}\par\vspace{0.2mm}
\noindent{\fontsize{6.9pt}{8.0pt}\selectfont\textit{21 – 22 avril 1917} \hfill \textit{Hérissart, Toutencourt, Buire, Albert, Millencourt, Contay, Molliens-au-Bois, Pierregot (Somme)}}\par\vspace{0.5mm}
\noindent\rule{\textwidth}{0.35pt}\par\vspace{0.8mm}
% --- Corps de la transcription diplomatique et apparat marginal en prolongement ---
{\renewcommand{\arraystretch}{0.83}%
\noindent\hspace*{3.5mm}\begin{tabular}{@{}p{87mm}@{\hspace{3.5mm}}p{39mm}@{}}
      \leavevmode\fontsize{6.6pt}{7.5pt}\selectfont\strut \underline{21 Samedi} — \hfill{} \underline{— Avril 1917 —}\strut &  \\
      \leavevmode\fontsize{6.6pt}{7.5pt}\selectfont\strut Toute nuit, violente canonnade\strut &  \\
      \leavevmode\fontsize{6.6pt}{7.5pt}\selectfont\strut sud d'Arras.\strut &  \\
      \leavevmode\fontsize{6.6pt}{7.5pt}\selectfont\strut Encore assez belle journée. Temps doux, ciel couvert.\strut &  \\
      \leavevmode\llap{\fontsize{6.4pt}{7.3pt}\selectfont\color{darkgray}5\hspace{1.8mm}}\fontsize{6.6pt}{7.5pt}\selectfont\strut Matin, avant breakfast, venu à pied, en 40 m.,\strut & \smash{\parbox[t]{39mm}{\fontsize{6.2pt}{7.0pt}\selectfont\raggedright\color{black!88}\textbf{5 breakfast} \textit{angl. in ms.}}} \\
      \leavevmode\fontsize{6.6pt}{7.5pt}\selectfont\strut à Toutencourt, pour toucher solde. Dû revenir à\strut &  \\
      \leavevmode\fontsize{6.6pt}{7.5pt}\selectfont\strut pied.\strut &  \\
      \leavevmode\fontsize{6.6pt}{7.5pt}\selectfont\strut Course en auto à Buire, Albert \& Millencourt,\strut &  \\
      \leavevmode\fontsize{6.6pt}{7.5pt}\selectfont\strut Contay. Cherché vainement paille à acheter.\strut &  \\
      \leavevmode\llap{\fontsize{6.4pt}{7.3pt}\selectfont\color{darkgray}10\hspace{1.8mm}}\fontsize{6.6pt}{7.5pt}\selectfont\strut Après-midi, liberté entière.\strut &  \\
      \leavevmode\fontsize{6.6pt}{7.5pt}\selectfont\strut [ A Toutencourt : solde : 78 f. 75 c.\strut & \smash{\parbox[t]{39mm}{\fontsize{6.2pt}{7.0pt}\selectfont\raggedright\color{black!88}\textbf{11} \textit{in marg. cum unco} (comptes). \textbf{11 f.} \textit{abbr. pro} franc(s). — \textbf{13 ds} \textit{abbr. pro} dans. \textbf{13 f.} \textit{abbr. pro} franc(s). — \textbf{14 gde} \textit{abbr. pro} grande. \textbf{14 cathle} \textit{abbr. pro} cathédrale. \textbf{14 avec} \textit{lect. incert. in ms.}}} \\
      \leavevmode\fontsize{6.6pt}{7.5pt}\selectfont\strut Fait porter d'Amiens, par Bte Denis :\strut &  \\
      \leavevmode\fontsize{6.6pt}{7.5pt}\selectfont\strut A. Martin : Promenade ds Amiens : 1 f. 75\strut &  \\
      \leavevmode\fontsize{6.6pt}{7.5pt}\selectfont\strut et gde carte, [avec ?] vue cathle : 2 f. 75 c. ]\strut &  \\
      \cline{1-1}
      \leavevmode\llap{\fontsize{6.4pt}{7.3pt}\selectfont\color{darkgray}15\hspace{1.8mm}}\fontsize{6.6pt}{7.5pt}\selectfont\strut \underline{22 Dimanche} —\strut &  \\
      \leavevmode\fontsize{6.6pt}{7.5pt}\selectfont\strut Enfin, très belle journée. Soleil,\strut &  \\
      \leavevmode\fontsize{6.6pt}{7.5pt}\selectfont\strut mais vent très froid.\strut &  \\
      \leavevmode\fontsize{6.6pt}{7.5pt}\selectfont\strut Matinée, course à Toutencourt, Molliens-au-Bois\strut &  \\
      \leavevmode\fontsize{6.6pt}{7.5pt}\selectfont\strut et Pierregot. — Visité gd moulin à vent route de P.,\strut & \smash{\parbox[t]{39mm}{\fontsize{6.2pt}{7.0pt}\selectfont\raggedright\color{black!88}\textbf{19 gd} \textit{abbr. pro} grand.}} \\
      \leavevmode\llap{\fontsize{6.4pt}{7.3pt}\selectfont\color{darkgray}20\hspace{1.8mm}}\fontsize{6.6pt}{7.5pt}\selectfont\strut très intéressant.\strut &  \\
      \leavevmode\fontsize{6.6pt}{7.5pt}\selectfont\strut Après-midi, liberté. Promenade agréable\strut &  \\
      \leavevmode\fontsize{6.6pt}{7.5pt}\selectfont\strut sur la route d'Hérissart, anti-aircraft battery.\strut & \smash{\parbox[t]{39mm}{\fontsize{6.2pt}{7.0pt}\selectfont\raggedright\color{black!88}\textbf{22 anti-aircraft battery} \textit{angl. in ms.}}} \\
      \leavevmode\fontsize{6.6pt}{7.5pt}\selectfont\strut Très pittoresques points de vue : lointains bleuâtres,\strut &  \\
      \leavevmode\fontsize{6.6pt}{7.5pt}\selectfont\strut masses des bois rouillées —\strut &  \\
      \leavevmode\llap{\fontsize{6.4pt}{7.3pt}\selectfont\color{darkgray}25\hspace{1.8mm}}\fontsize{6.6pt}{7.5pt}\selectfont\strut Durant jour, et surtout tte la nuit, roulement\strut & \smash{\parbox[t]{39mm}{\fontsize{6.2pt}{7.0pt}\selectfont\raggedright\color{black!88}\textbf{25 tte} \textit{abbr. pro} toute.}} \\
      \leavevmode\fontsize{6.6pt}{7.5pt}\selectfont\strut ininterrompu, violente canonnade vers le nord,\strut &  \\
      \leavevmode\fontsize{6.6pt}{7.5pt}\selectfont\strut du côté d'Arras —\strut &  \\
\end{tabular}}%
\par\vspace*{\fill}
% --- Registres de bas de page (1. Critique interne [si pertinent] ; 2. Loci paralleli en latin ; 3. Commentaire historique en retrait) ---
\noindent\rule{\textwidth}{0.5pt}\par\vspace{0.65mm}
{\fontsize{6.6pt}{7.7pt}\selectfont%
\noindent\normalfont\color{black!92} \textbf{3} (\textit{Arras}) cf. supra f\textsuperscript{o}~16~r\textsuperscript{o} (p.~\pageref{fol:32}), f\textsuperscript{o}~89~r\textsuperscript{o} (p.~\pageref{fol:178}). — \textbf{6} (\textit{Toutencourt}) cf. supra f\textsuperscript{o}~49~v\textsuperscript{o} (p.~\pageref{fol:99}), f\textsuperscript{o}~89~r\textsuperscript{o} (p.~\pageref{fol:178}), infra f\textsuperscript{o}~91~v\textsuperscript{o} (p.~\pageref{fol:183}), f\textsuperscript{o}~93~r\textsuperscript{o} (p.~\pageref{fol:186}), \textit{et passim}. — \textbf{8} (\textit{Albert}) cf. supra f\textsuperscript{o}~39~r\textsuperscript{o} (p.~\pageref{fol:78}), f\textsuperscript{o}~90~v\textsuperscript{o} (p.~\pageref{fol:181}), \textit{et passim}.\par\vspace{0.55mm}
\noindent\rule{38mm}{0.4pt}\par\vspace{0.5mm}
\noindent\hspace*{5mm}\parbox{\dimexpr\textwidth-5mm\relax}{\fontsize{6.6pt}{7.7pt}\selectfont\color{black!84}La violente canonnade entendue les 21 et 22 avril 1917 au nord et au sud d'Arras correspond à la préparation d'artillerie britannique de la \textbf{seconde bataille de la Scarpe} (\textit{Second Battle of the Scarpe}, déclenchée à l'aube du 23 avril 1917 sur la ligne Croisilles–Guémappe). Le guide rapporté d'Amiens par Baptiste Denis est l'ouvrage illustré d'\textbf{Alexis Martin}, \textit{Une visite à Amiens} (Paris, Hennuyer, 1896).}\par}
\newpage
\phantomsection\label{fac:182}
% --- En-tête du fac-similé (Page impaire / droite : Photographie du manuscrit) ---
\noindent\rule{\textwidth}{0.5pt}\par\vspace{0.55mm}
\noindent\textbf{\fontsize{6.8pt}{7.8pt}\selectfont FAC-SIMILÉ EN VIS-À-VIS — FOLIO 91 RECTO} \hfill \textbf{\fontsize{6.8pt}{7.8pt}\selectfont AVRIL 1917}\par\vspace{0.25mm}
\noindent{\fontsize{6.8pt}{8.0pt}\selectfont\textit{Livre-Journal d'Hippolyte Combes} — 21 – 22 avril 1917 (Hérissart, Toutencourt, Buire, Albert, Millencourt, Contay, Molliens-au-Bois, Pierregot (Somme))}\par\vspace{0.55mm}
\noindent\rule{\textwidth}{0.35pt}\par\vspace*{\fill}
{\centering
  \setlength{\fboxsep}{1.2mm}%
  \setlength{\fboxrule}{0.4pt}%
  \fcolorbox{gray!55}{white}{%
    \includegraphics[width=0.95\textwidth,height=0.81\textheight,keepaspectratio]{\facsimiledir/20251126_111702.jpg}%
  }\par}
\vspace*{\fill}
\newpage

% =============================================================================
% VIS-A-VIS SPREAD : FOLIO 91 VERSO (20251126_111715.jpg)
% =============================================================================
\cleardoubleevenpage
\renewcommand{\currentmonthyear}{Avril 1917}
\renewcommand{\currentfoliolabel}{Folio 91 verso}
\phantomsection\label{fol:183}
% --- En-tête du folio (Page paire / gauche : Transcription, apparat & registres de bas de page) ---
\noindent\rule{\textwidth}{0.5pt}\par\vspace{0.55mm}
\noindent\textbf{\footnotesize FOLIO 91 VERSO} \hfill \textbf{\footnotesize AVRIL 1917}\par\vspace{0.25mm}
\noindent{\fontsize{6.9pt}{8.0pt}\selectfont Première vraie journée d'été, bourgeons vernissés et soirée chez la famille Denis}\par\vspace{0.2mm}
\noindent{\fontsize{6.9pt}{8.0pt}\selectfont\textit{23 – 25 avril 1917} \hfill \textit{Hérissart, Toutencourt, Pierregot, Acheux, Puchevillers, Mirvaux (Somme)}}\par\vspace{0.5mm}
\noindent\rule{\textwidth}{0.35pt}\par\vspace{0.8mm}
% --- Corps de la transcription diplomatique et apparat marginal en prolongement ---
{\renewcommand{\arraystretch}{0.83}%
\noindent\hspace*{3.5mm}\begin{tabular}{@{}p{87mm}@{\hspace{3.5mm}}p{39mm}@{}}
      \leavevmode\fontsize{6.6pt}{7.5pt}\selectfont\strut \underline{23 Lundi} — \hfill{} \underline{Avril 1917}\strut &  \\
      \leavevmode\fontsize{6.6pt}{7.5pt}\selectfont\strut Temps beau journée entière.\strut &  \\
      \leavevmode\fontsize{6.6pt}{7.5pt}\selectfont\strut Matin, course à Toutencourt et Pierregot. Métier à tisser.\strut &  \\
      \leavevmode\fontsize{6.6pt}{7.5pt}\selectfont\strut Après-midi, course à Acheux et Puchevillers.\strut &  \\
      \leavevmode\llap{\fontsize{6.4pt}{7.3pt}\selectfont\color{darkgray}5\hspace{1.8mm}}\fontsize{6.6pt}{7.5pt}\selectfont\strut Dîner terminé, courte promenade s/ la route du\strut & \smash{\parbox[t]{39mm}{\fontsize{6.2pt}{7.0pt}\selectfont\raggedright\color{black!88}\textbf{5 s/} \textit{abbr. pro} sur.}} \\
      \leavevmode\fontsize{6.6pt}{7.5pt}\selectfont\strut cimetière. Admiré sunset magnifique.\strut & \smash{\parbox[t]{39mm}{\fontsize{6.2pt}{7.0pt}\selectfont\raggedright\color{black!88}\textbf{6 sunset} \textit{angl. in ms.}}} \\
      \leavevmode\fontsize{6.6pt}{7.5pt}\selectfont\strut Très belle nuit étoilée —\strut &  \\
      \cline{1-1}
      \leavevmode\fontsize{6.6pt}{7.5pt}\selectfont\strut \underline{24 Mardi} — Vraie journée d'été. Tout le jour,\strut &  \\
      \leavevmode\fontsize{6.6pt}{7.5pt}\selectfont\strut ciel d'1 azur très foncé. Ds l'air lumineux,\strut & \smash{\parbox[t]{39mm}{\fontsize{6.2pt}{7.0pt}\selectfont\raggedright\color{black!88}\textbf{9 Ds} \textit{abbr. pro} Dans.}} \\
      \leavevmode\llap{\fontsize{6.4pt}{7.3pt}\selectfont\color{darkgray}10\hspace{1.8mm}}\fontsize{6.6pt}{7.5pt}\selectfont\strut branches des arbres se détachent nettement\strut &  \\
      \leavevmode\fontsize{6.6pt}{7.5pt}\selectfont\strut et bourgeons vernissés luisent au soleil.\strut &  \\
      \leavevmode\fontsize{6.6pt}{7.5pt}\selectfont\strut Tempér. chaude.\strut &  \\
      \leavevmode\fontsize{6.6pt}{7.5pt}\selectfont\strut Après-midi, course à Toutencourt \& Pierregot.\strut &  \\
      \leavevmode\fontsize{6.6pt}{7.5pt}\selectfont\strut Retour à 5h —\strut &  \\
      \leavevmode\llap{\fontsize{6.4pt}{7.3pt}\selectfont\color{darkgray}15\hspace{1.8mm}}\fontsize{6.6pt}{7.5pt}\selectfont\strut Dîner pris, promenade route de Mirvaux,\strut &  \\
      \leavevmode\fontsize{6.6pt}{7.5pt}\selectfont\strut chemin vicinal. Magnifiques échappées sur Beaucourt\strut &  \\
      \leavevmode\fontsize{6.6pt}{7.5pt}\selectfont\strut \& Montigny. Lointains bleuâtres. Horizons violets.\strut &  \\
      \leavevmode\fontsize{6.6pt}{7.5pt}\selectfont\strut Très beau sunset —\strut & \smash{\parbox[t]{39mm}{\fontsize{6.2pt}{7.0pt}\selectfont\raggedright\color{black!88}\textbf{18 sunset} \textit{angl. in ms.}}} \\
      \cline{1-1}
      \leavevmode\fontsize{6.6pt}{7.5pt}\selectfont\strut \underline{25 Merc.} — Tout le jour, ciel resté couvert. Fines ondées\strut & \smash{\parbox[t]{39mm}{\fontsize{6.2pt}{7.0pt}\selectfont\raggedright\color{black!88}\textbf{19 Merc.} \textit{scil.} Mercredi.}} \\
      \leavevmode\llap{\fontsize{6.4pt}{7.3pt}\selectfont\color{darkgray}20\hspace{1.8mm}}\fontsize{6.6pt}{7.5pt}\selectfont\strut de courte durée. Temps redevenu 1 peu froid.\strut &  \\
      \leavevmode\fontsize{6.6pt}{7.5pt}\selectfont\strut Matinée, course à Acheux — village presque désert\strut &  \\
      \leavevmode\fontsize{6.6pt}{7.5pt}\selectfont\strut et mort, plus de troupes cantonnées.\strut &  \\
      \leavevmode\fontsize{6.6pt}{7.5pt}\selectfont\strut Après-midi, à Pierregot. Lu chz Mme Mignot\strut & \smash{\parbox[t]{39mm}{\fontsize{6.2pt}{7.0pt}\selectfont\raggedright\color{black!88}\textbf{23 Mme} \textit{abbr. pro} Madame.}} \\
      \leavevmode\fontsize{6.6pt}{7.5pt}\selectfont\strut (et Louise)\strut &  \\
      \leavevmode\llap{\fontsize{6.4pt}{7.3pt}\selectfont\color{darkgray}25\hspace{1.8mm}}\fontsize{6.6pt}{7.5pt}\selectfont\strut Soirée passée chez la famille Denis, à Hérissart —\strut &  \\
\end{tabular}}%
\par\vspace*{\fill}
% --- Registres de bas de page (1. Critique interne [si pertinent] ; 2. Loci paralleli en latin ; 3. Commentaire historique en retrait) ---
\noindent\rule{\textwidth}{0.5pt}\par\vspace{0.65mm}
{\fontsize{6.6pt}{7.7pt}\selectfont%
\noindent\normalfont\color{black!92} \textbf{3} (\textit{Toutencourt}) cf. supra f\textsuperscript{o}~49~v\textsuperscript{o} (p.~\pageref{fol:99}), f\textsuperscript{o}~91~r\textsuperscript{o} (p.~\pageref{fol:182}), infra f\textsuperscript{o}~93~r\textsuperscript{o} (p.~\pageref{fol:186}), \textit{et passim}. — \textbf{3} (\textit{Pierregot}) cf. supra f\textsuperscript{o}~52~r\textsuperscript{o} (p.~\pageref{fol:104}), f\textsuperscript{o}~91~r\textsuperscript{o} (p.~\pageref{fol:182}), infra f\textsuperscript{o}~92~r\textsuperscript{o} (p.~\pageref{fol:184}), f\textsuperscript{o}~92~v\textsuperscript{o} (p.~\pageref{fol:185}), \textit{et passim}. — \textbf{4} (\textit{Acheux}) cf. supra f\textsuperscript{o}~21~v\textsuperscript{o} (p.~\pageref{fol:43}), f\textsuperscript{o}~88~v\textsuperscript{o} (p.~\pageref{fol:177}), \textit{et passim}.\par\vspace{0.55mm}
\noindent\rule{38mm}{0.4pt}\par\vspace{0.5mm}
\noindent\hspace*{5mm}\parbox{\dimexpr\textwidth-5mm\relax}{\fontsize{6.6pt}{7.7pt}\selectfont\color{black!84}Journées des 23, 24 et 25 avril 1917 autour d'Hérissart: observation d'un métier à tisser domestique (survivance du tissage rural picard de la bonneterie et du velours d'Amiens), promenade vespérale sur le chemin de Mirvaux offrant une vue sur la vallée de l'Hallue (\textbf{Beaucourt-sur-l'Hallue} et \textbf{Montigny-sur-l'Hallue}), constat de l'abandon militaire d'\textbf{Acheux-en-Amiénois} le 25 avril, lecture chez \textbf{Mme Mignot} (et sa fille \textbf{Louise}) à \textbf{Pierregot} et veillée chez la \textbf{famille Denis}.}\par}
\newpage
\phantomsection\label{fac:183}
% --- En-tête du fac-similé (Page impaire / droite : Photographie du manuscrit) ---
\noindent\rule{\textwidth}{0.5pt}\par\vspace{0.55mm}
\noindent\textbf{\fontsize{6.8pt}{7.8pt}\selectfont FAC-SIMILÉ EN VIS-À-VIS — FOLIO 91 VERSO} \hfill \textbf{\fontsize{6.8pt}{7.8pt}\selectfont AVRIL 1917}\par\vspace{0.25mm}
\noindent{\fontsize{6.8pt}{8.0pt}\selectfont\textit{Livre-Journal d'Hippolyte Combes} — 23 – 25 avril 1917 (Hérissart, Toutencourt, Pierregot, Acheux, Puchevillers, Mirvaux (Somme))}\par\vspace{0.55mm}
\noindent\rule{\textwidth}{0.35pt}\par\vspace*{\fill}
{\centering
  \setlength{\fboxsep}{1.2mm}%
  \setlength{\fboxrule}{0.4pt}%
  \fcolorbox{gray!55}{white}{%
    \includegraphics[width=0.95\textwidth,height=0.81\textheight,keepaspectratio]{\facsimiledir/20251126_111715.jpg}%
  }\par}
\vspace*{\fill}
\newpage

% =============================================================================
% VIS-A-VIS SPREAD : FOLIO 92 RECTO (20251126_111737.jpg)
% =============================================================================
\cleardoubleevenpage
\renewcommand{\currentmonthyear}{Avril 1917}
\renewcommand{\currentfoliolabel}{Folio 92 recto}
\phantomsection\label{fol:184}
% --- En-tête du folio (Page paire / gauche : Transcription, apparat & registres de bas de page) ---
\noindent\rule{\textwidth}{0.5pt}\par\vspace{0.55mm}
\noindent\textbf{\footnotesize FOLIO 92 RECTO} \hfill \textbf{\footnotesize AVRIL 1917}\par\vspace{0.25mm}
\noindent{\fontsize{6.9pt}{8.0pt}\selectfont Promenade le long des anciennes tranchées de Puchevillers et visite à la famille Mortreux à Varennes}\par\vspace{0.2mm}
\noindent{\fontsize{6.9pt}{8.0pt}\selectfont\textit{26 – 27 avril 1917} \hfill \textit{Hérissart, Contay, Puchevillers, Pierregot, Harponville, Varennes, Rubempré (Somme)}}\par\vspace{0.5mm}
\noindent\rule{\textwidth}{0.35pt}\par\vspace{0.8mm}
% --- Corps de la transcription diplomatique et apparat marginal en prolongement ---
{\renewcommand{\arraystretch}{0.87}%
\noindent\hspace*{3.5mm}\begin{tabular}{@{}p{87mm}@{\hspace{3.5mm}}p{39mm}@{}}
      \leavevmode\fontsize{6.9pt}{7.9pt}\selectfont\strut \underline{26 Jeudi.} \hfill{} \underline{— Avril 1917 —}\strut &  \\
      \leavevmode\fontsize{6.9pt}{7.9pt}\selectfont\strut Matinée, temps nuageux.\strut &  \\
      \leavevmode\fontsize{6.9pt}{7.9pt}\selectfont\strut Fin après-midi, soleil brille. Fin de journée très\strut &  \\
      \leavevmode\fontsize{6.9pt}{7.9pt}\selectfont\strut belle. Courses à Contay, Puchevillers \& Pierregot.\strut &  \\
      \leavevmode\llap{\fontsize{6.4pt}{7.3pt}\selectfont\color{darkgray}5\hspace{1.8mm}}\fontsize{6.9pt}{7.9pt}\selectfont\strut Après dîner, très agréable promenade petit\strut &  \\
      \leavevmode\fontsize{6.9pt}{7.9pt}\selectfont\strut chemin de Puchevillers, le long de dug-outs et\strut & \smash{\parbox[t]{39mm}{\fontsize{6.2pt}{7.0pt}\selectfont\raggedright\color{black!88}\textbf{6 dug-outs} \textit{angl. in ms.}}} \\
      \leavevmode\fontsize{6.9pt}{7.9pt}\selectfont\strut de tranchées. Ciel bleu, chaud soleil.\strut &  \\
      \leavevmode\fontsize{6.9pt}{7.9pt}\selectfont\strut Nuit merveillet étoilée —\strut & \smash{\parbox[t]{39mm}{\fontsize{6.2pt}{7.0pt}\selectfont\raggedright\color{black!88}\textbf{8 merveillet} \textit{abbr. pro} merveilleusement. — \textbf{9 Ven.} \textit{scil.} Vendredi.}} \\
      \cline{1-1}
      \leavevmode\fontsize{6.9pt}{7.9pt}\selectfont\strut \underline{27 Ven.}\strut &  \\
      \leavevmode\llap{\fontsize{6.4pt}{7.3pt}\selectfont\color{darkgray}10\hspace{1.8mm}}\fontsize{6.9pt}{7.9pt}\selectfont\strut Matin, courses à Contay, Harponville \& Varennes,\strut &  \\
      \leavevmode\fontsize{6.9pt}{7.9pt}\selectfont\strut où revu av. plaisir famille Mortreux —\strut &  \\
      \leavevmode\fontsize{6.9pt}{7.9pt}\selectfont\strut Après-midi, Pierregot \& Rubempré. Rentré\strut &  \\
      \leavevmode\fontsize{6.9pt}{7.9pt}\selectfont\strut seul à pied de Rubempré; l'auto, continuant\strut &  \\
      \leavevmode\fontsize{6.9pt}{7.9pt}\selectfont\strut sur Varennes.\strut &  \\
      \leavevmode\llap{\fontsize{6.4pt}{7.3pt}\selectfont\color{darkgray}15\hspace{1.8mm}}\fontsize{6.9pt}{7.9pt}\selectfont\strut Tout le jour, temps nuageux, 1 peu froid.\strut &  \\
      \leavevmode\fontsize{6.9pt}{7.9pt}\selectfont\strut Soirée, chez les Denis, à Hérissart.\strut &  \\
\end{tabular}}%
\par\vspace*{\fill}
% --- Registres de bas de page (1. Critique interne [si pertinent] ; 2. Loci paralleli en latin ; 3. Commentaire historique en retrait) ---
\noindent\rule{\textwidth}{0.5pt}\par\vspace{0.65mm}
{\fontsize{6.6pt}{7.7pt}\selectfont%
\noindent\normalfont\color{black!92} \textbf{4} (\textit{Puchevillers}) cf. supra f\textsuperscript{o}~37~v\textsuperscript{o} (p.~\pageref{fol:75}), f\textsuperscript{o}~91~v\textsuperscript{o} (p.~\pageref{fol:183}), infra f\textsuperscript{o}~93~r\textsuperscript{o} (p.~\pageref{fol:186}), \textit{et passim}. — \textbf{4} (\textit{Contay}) cf. supra f\textsuperscript{o}~49~r\textsuperscript{o} (p.~\pageref{fol:98}), f\textsuperscript{o}~91~r\textsuperscript{o} (p.~\pageref{fol:182}), infra f\textsuperscript{o}~92~v\textsuperscript{o} (p.~\pageref{fol:185}), \textit{et passim}. — \textbf{4} (\textit{Pierregot}) cf. supra f\textsuperscript{o}~52~r\textsuperscript{o} (p.~\pageref{fol:104}), f\textsuperscript{o}~91~v\textsuperscript{o} (p.~\pageref{fol:183}), infra f\textsuperscript{o}~92~v\textsuperscript{o} (p.~\pageref{fol:185}), f\textsuperscript{o}~93~r\textsuperscript{o} (p.~\pageref{fol:186}), \textit{et passim}.\par\vspace{0.55mm}
\noindent\rule{38mm}{0.4pt}\par\vspace{0.5mm}
\noindent\hspace*{5mm}\parbox{\dimexpr\textwidth-5mm\relax}{\fontsize{6.6pt}{7.7pt}\selectfont\color{black!84}Les jeudi 26 et vendredi 27 avril 1917, Hippolyte Combes poursuit ses tournées à \textbf{Contay}, \textbf{Puchevillers} (où il se promène le soir le long d'anciens abris souterrains, \textit{dug-outs}, et tranchées d'exercice), \textbf{Pierregot}, \textbf{Harponville} et \textbf{Varennes-en-Croix}, où il rend visite à la \textbf{famille Mortreux}.}\par}
\newpage
\phantomsection\label{fac:184}
% --- En-tête du fac-similé (Page impaire / droite : Photographie du manuscrit) ---
\noindent\rule{\textwidth}{0.5pt}\par\vspace{0.55mm}
\noindent\textbf{\fontsize{6.8pt}{7.8pt}\selectfont FAC-SIMILÉ EN VIS-À-VIS — FOLIO 92 RECTO} \hfill \textbf{\fontsize{6.8pt}{7.8pt}\selectfont AVRIL 1917}\par\vspace{0.25mm}
\noindent{\fontsize{6.8pt}{8.0pt}\selectfont\textit{Livre-Journal d'Hippolyte Combes} — 26 – 27 avril 1917 (Hérissart, Contay, Puchevillers, Pierregot, Harponville, Varennes, Rubempré (Somme))}\par\vspace{0.55mm}
\noindent\rule{\textwidth}{0.35pt}\par\vspace*{\fill}
{\centering
  \setlength{\fboxsep}{1.2mm}%
  \setlength{\fboxrule}{0.4pt}%
  \fcolorbox{gray!55}{white}{%
    \includegraphics[width=0.95\textwidth,height=0.81\textheight,keepaspectratio]{\facsimiledir/20251126_111737.jpg}%
  }\par}
\vspace*{\fill}
\newpage

% =============================================================================
% VIS-A-VIS SPREAD : FOLIO 92 VERSO (20251126_111749.jpg)
% =============================================================================
\cleardoubleevenpage
\renewcommand{\currentmonthyear}{Avril 1917}
\renewcommand{\currentfoliolabel}{Folio 92 verso}
\phantomsection\label{fol:185}
% --- En-tête du folio (Page paire / gauche : Transcription, apparat & registres de bas de page) ---
\noindent\rule{\textwidth}{0.5pt}\par\vspace{0.55mm}
\noindent\textbf{\footnotesize FOLIO 92 VERSO} \hfill \textbf{\footnotesize AVRIL 1917}\par\vspace{0.25mm}
\noindent{\fontsize{6.9pt}{8.0pt}\selectfont Annonce du départ prochain pour Cambrai et flânerie avec le lieutenant Gilbert}\par\vspace{0.2mm}
\noindent{\fontsize{6.9pt}{8.0pt}\selectfont\textit{28 – 29 avril 1917} \hfill \textit{Hérissart, Pierregot, Contay, Mirvaux (Somme)}}\par\vspace{0.5mm}
\noindent\rule{\textwidth}{0.35pt}\par\vspace{0.8mm}
% --- Corps de la transcription diplomatique et apparat marginal en prolongement ---
{\renewcommand{\arraystretch}{0.87}%
\noindent\hspace*{3.5mm}\begin{tabular}{@{}p{87mm}@{\hspace{3.5mm}}p{39mm}@{}}
      \leavevmode\fontsize{6.9pt}{7.9pt}\selectfont\strut \underline{28 Samedi} — \hfill{} \underline{— Avril 1917 —}\strut &  \\
      \leavevmode\fontsize{6.9pt}{7.9pt}\selectfont\strut Commencemt jour, temps couvert.\strut &  \\
      \leavevmode\fontsize{6.9pt}{7.9pt}\selectfont\strut Fin matin et tout le jour, très beau soleil.\strut &  \\
      \leavevmode\fontsize{6.9pt}{7.9pt}\selectfont\strut Vers 5h 1/4, sortis en auto, course à Pierregot.\strut &  \\
      \leavevmode\llap{\fontsize{6.4pt}{7.3pt}\selectfont\color{darkgray}5\hspace{1.8mm}}\fontsize{6.9pt}{7.9pt}\selectfont\strut Off., retour de Bournemouth, m'a appris que\strut &  \\
      \leavevmode\fontsize{6.9pt}{7.9pt}\selectfont\strut bientôt ns quitterons le village et ns irons à\strut & \smash{\parbox[t]{39mm}{\fontsize{6.2pt}{7.0pt}\selectfont\raggedright\color{black!88}\textbf{6 ns} \textit{abbr. pro} nous.}} \\
      \leavevmode\fontsize{6.9pt}{7.9pt}\selectfont\strut Cambrai. Regretterai bcp les points de vue\strut & \smash{\parbox[t]{39mm}{\fontsize{6.2pt}{7.0pt}\selectfont\raggedright\color{black!88}\textbf{7 bcp} \textit{abbr. pro} beaucoup.}} \\
      \leavevmode\fontsize{6.9pt}{7.9pt}\selectfont\strut et panoramas circulaires des environs d'Hérissart :\strut &  \\
      \leavevmode\fontsize{6.9pt}{7.9pt}\selectfont\strut c'est la guerre.\strut &  \\
      \leavevmode\llap{\fontsize{6.4pt}{7.3pt}\selectfont\color{darkgray}10\hspace{1.8mm}}\fontsize{6.9pt}{7.9pt}\selectfont\strut Après lunch, très agréable promenade\strut & \smash{\parbox[t]{39mm}{\fontsize{6.2pt}{7.0pt}\selectfont\raggedright\color{black!88}\textbf{10 lunch} \textit{angl. in ms.}}} \\
      \leavevmode\fontsize{6.9pt}{7.9pt}\selectfont\strut en pleine campagne.\strut &  \\
      \cline{1-1}
      \leavevmode\fontsize{6.9pt}{7.9pt}\selectfont\strut \underline{29 Dim.}\strut & \smash{\parbox[t]{39mm}{\fontsize{6.2pt}{7.0pt}\selectfont\raggedright\color{black!88}\textbf{12 Dim.} \textit{scil.} Dimanche.}} \\
      \leavevmode\fontsize{6.9pt}{7.9pt}\selectfont\strut Journée très belle. C'est enfin le printemps!\strut &  \\
      \leavevmode\fontsize{6.9pt}{7.9pt}\selectfont\strut Matin, partis en auto. En arrivant à Contay,\strut &  \\
      \leavevmode\llap{\fontsize{6.4pt}{7.3pt}\selectfont\color{darkgray}15\hspace{1.8mm}}\fontsize{6.9pt}{7.9pt}\selectfont\strut avarie : 2 ressorts cassés. Dû retourner à\strut &  \\
      \leavevmode\fontsize{6.9pt}{7.9pt}\selectfont\strut Hérissart, d'abord à pied, puis lift sur 1\strut & \smash{\parbox[t]{39mm}{\fontsize{6.2pt}{7.0pt}\selectfont\raggedright\color{black!88}\textbf{16 lift} \textit{angl. in ms.}}} \\
      \leavevmode\fontsize{6.9pt}{7.9pt}\selectfont\strut water-cart.\strut & \smash{\parbox[t]{39mm}{\fontsize{6.2pt}{7.0pt}\selectfont\raggedright\color{black!88}\textbf{17 water-cart} \textit{angl. in ms.}}} \\
      \leavevmode\fontsize{6.9pt}{7.9pt}\selectfont\strut Liberté matin et soir.\strut &  \\
      \leavevmode\fontsize{6.9pt}{7.9pt}\selectfont\strut Après-midi, promenade en cie Lt Gilbert\strut & \smash{\parbox[t]{39mm}{\fontsize{6.2pt}{7.0pt}\selectfont\raggedright\color{black!88}\textbf{19 Lt} \textit{abbr. pro} Lieutenant. \textbf{19 cie} \textit{abbr. pro} compagnie.}} \\
      \leavevmode\llap{\fontsize{6.4pt}{7.3pt}\selectfont\color{darkgray}20\hspace{1.8mm}}\fontsize{6.9pt}{7.9pt}\selectfont\strut sur route Mirvaux. Assis sur l'herbe rebord\strut &  \\
      \leavevmode\fontsize{6.9pt}{7.9pt}\selectfont\strut talus en pleins champs et causé 1 peu de tout.\strut &  \\
      \leavevmode\fontsize{6.9pt}{7.9pt}\selectfont\strut Fin jour magnifique. Beau sunset.\strut & \smash{\parbox[t]{39mm}{\fontsize{6.2pt}{7.0pt}\selectfont\raggedright\color{black!88}\textbf{22 sunset} \textit{angl. in ms.}}} \\
      \leavevmode\fontsize{6.9pt}{7.9pt}\selectfont\strut Nuit superbe, mervt étoilée et claire.\strut & \smash{\parbox[t]{39mm}{\fontsize{6.2pt}{7.0pt}\selectfont\raggedright\color{black!88}\textbf{23 mervt} \textit{abbr. pro} merveilleusement.}} \\
\end{tabular}}%
\par\vspace*{\fill}
% --- Registres de bas de page (1. Critique interne [si pertinent] ; 2. Loci paralleli en latin ; 3. Commentaire historique en retrait) ---
\noindent\rule{\textwidth}{0.5pt}\par\vspace{0.65mm}
{\fontsize{6.6pt}{7.7pt}\selectfont%
\noindent\normalfont\color{black!92} \textbf{4} (\textit{Pierregot}) cf. supra f\textsuperscript{o}~52~r\textsuperscript{o} (p.~\pageref{fol:104}), f\textsuperscript{o}~92~r\textsuperscript{o} (p.~\pageref{fol:184}), infra f\textsuperscript{o}~93~r\textsuperscript{o} (p.~\pageref{fol:186}), \textit{et passim}. — \textbf{8} (\textit{Hérissart}) cf. supra f\textsuperscript{o}~47~v\textsuperscript{o} (p.~\pageref{fol:95}), f\textsuperscript{o}~92~r\textsuperscript{o} (p.~\pageref{fol:184}), infra f\textsuperscript{o}~93~r\textsuperscript{o} (p.~\pageref{fol:186}), \textit{et passim}. — \textbf{14} (\textit{Contay}) cf. supra f\textsuperscript{o}~49~r\textsuperscript{o} (p.~\pageref{fol:98}), f\textsuperscript{o}~92~r\textsuperscript{o} (p.~\pageref{fol:184}), \textit{et passim}.\par\vspace{0.55mm}
\noindent\rule{38mm}{0.4pt}\par\vspace{0.5mm}
\noindent\hspace*{5mm}\parbox{\dimexpr\textwidth-5mm\relax}{\fontsize{6.6pt}{7.7pt}\selectfont\color{black!84}Le samedi 28 avril 1917, un officier britannique rentrant de la station balnéaire anglaise de \textbf{Bournemouth} (Dorset / Hampshire) annonce le déplacement prochain du service des achats vers le secteur de \textbf{Cambrai} (vers lequel le front britannique s'est avancé depuis le repli allemand de mars 1917).}\par}
\newpage
\phantomsection\label{fac:185}
% --- En-tête du fac-similé (Page impaire / droite : Photographie du manuscrit) ---
\noindent\rule{\textwidth}{0.5pt}\par\vspace{0.55mm}
\noindent\textbf{\fontsize{6.8pt}{7.8pt}\selectfont FAC-SIMILÉ EN VIS-À-VIS — FOLIO 92 VERSO} \hfill \textbf{\fontsize{6.8pt}{7.8pt}\selectfont AVRIL 1917}\par\vspace{0.25mm}
\noindent{\fontsize{6.8pt}{8.0pt}\selectfont\textit{Livre-Journal d'Hippolyte Combes} — 28 – 29 avril 1917 (Hérissart, Pierregot, Contay, Mirvaux (Somme))}\par\vspace{0.55mm}
\noindent\rule{\textwidth}{0.35pt}\par\vspace*{\fill}
{\centering
  \setlength{\fboxsep}{1.2mm}%
  \setlength{\fboxrule}{0.4pt}%
  \fcolorbox{gray!55}{white}{%
    \includegraphics[width=0.95\textwidth,height=0.81\textheight,keepaspectratio]{\facsimiledir/20251126_111749.jpg}%
  }\par}
\vspace*{\fill}
\newpage

% =============================================================================
% VIS-A-VIS SPREAD : FOLIO 93 RECTO (20251126_111833.jpg)
% =============================================================================
\cleardoubleevenpage
\renewcommand{\currentmonthyear}{Avril 1917}
\renewcommand{\currentfoliolabel}{Folio 93 recto}
\phantomsection\label{fol:186}
% --- En-tête du folio (Page paire / gauche : Transcription, apparat & registres de bas de page) ---
\noindent\rule{\textwidth}{0.5pt}\par\vspace{0.55mm}
\noindent\textbf{\footnotesize FOLIO 93 RECTO} \hfill \textbf{\footnotesize AVRIL 1917}\par\vspace{0.25mm}
\noindent{\fontsize{6.9pt}{8.0pt}\selectfont Dernier folio du Livre-Journal : chaleur d'été, course à Villers-Bocage et nuit étoilée du 30 avril 1917}\par\vspace{0.2mm}
\noindent{\fontsize{6.9pt}{8.0pt}\selectfont\textit{30 avril 1917} \hfill \textit{Hérissart, Vadencourt, Warloy, Hénencourt, Millencourt, Lavieville, Puchevillers, Pierregot, Villers-Bocage, Toutencourt (Somme)}}\par\vspace{0.5mm}
\noindent\rule{\textwidth}{0.35pt}\par\vspace{0.8mm}
% --- Corps de la transcription diplomatique et apparat marginal en prolongement ---
{\renewcommand{\arraystretch}{0.87}%
\noindent\hspace*{3.5mm}\begin{tabular}{@{}p{87mm}@{\hspace{3.5mm}}p{39mm}@{}}
      \leavevmode\fontsize{6.9pt}{7.9pt}\selectfont\strut \underline{30 Lundi.} — \hfill{} \underline{Avril 1917 —}\strut &  \\
      \leavevmode\fontsize{6.9pt}{7.9pt}\selectfont\strut Aujourd'hui encore, temps magnifique.\strut &  \\
      \leavevmode\fontsize{6.9pt}{7.9pt}\selectfont\strut Soleil ardent. Chaleur se fait sentir. Courses\strut &  \\
      \leavevmode\fontsize{6.9pt}{7.9pt}\selectfont\strut en auto, sans pardessus.\strut &  \\
      \leavevmode\llap{\fontsize{6.4pt}{7.3pt}\selectfont\color{darkgray}5\hspace{1.8mm}}\fontsize{6.9pt}{7.9pt}\selectfont\strut Matinée, venus à Vadencourt, Warloy, Hénencourt,\strut &  \\
      \leavevmode\fontsize{6.9pt}{7.9pt}\selectfont\strut Millencourt \& Lavieville.\strut &  \\
      \leavevmode\fontsize{6.9pt}{7.9pt}\selectfont\strut Après-midi, en cie Lt Gilbert, allés en auto\strut & \smash{\parbox[t]{39mm}{\fontsize{6.2pt}{7.0pt}\selectfont\raggedright\color{black!88}\textbf{7 Lt} \textit{abbr. pro} Lieutenant. \textbf{7 cie} \textit{abbr. pro} compagnie.}} \\
      \leavevmode\fontsize{6.9pt}{7.9pt}\selectfont\strut à Puchevillers, Pierregot \& Villers-Bocage.\strut &  \\
      \leavevmode\fontsize{6.9pt}{7.9pt}\selectfont\strut Route, auto chassé lapin sauvage.\strut &  \\
      \leavevmode\llap{\fontsize{6.4pt}{7.3pt}\selectfont\color{darkgray}10\hspace{1.8mm}}\fontsize{6.9pt}{7.9pt}\selectfont\strut Ciel d'un bleu pur \& intense. Aux arbres,\strut &  \\
      \leavevmode\fontsize{6.9pt}{7.9pt}\selectfont\strut bourgeons grossissent à vue d'œil.\strut &  \\
      \leavevmode\fontsize{6.9pt}{7.9pt}\selectfont\strut Après tea-dinner, rentré chambre où me\strut & \smash{\parbox[t]{39mm}{\fontsize{6.2pt}{7.0pt}\selectfont\raggedright\color{black!88}\textbf{12 tea} \textit{angl. in ms.}}} \\
      \leavevmode\fontsize{6.9pt}{7.9pt}\selectfont\strut suis mis à l'aise.\strut &  \\
      \leavevmode\fontsize{6.9pt}{7.9pt}\selectfont\strut Fatigué, couché de bonne heure.\strut &  \\
      \leavevmode\llap{\fontsize{6.4pt}{7.3pt}\selectfont\color{darkgray}15\hspace{1.8mm}}\fontsize{6.9pt}{7.9pt}\selectfont\strut Nuit très belle, mervt éclairée, lune superbe.\strut & \smash{\parbox[t]{39mm}{\fontsize{6.2pt}{7.0pt}\selectfont\raggedright\color{black!88}\textbf{15 mervt} \textit{abbr. pro} merveilleusement. — \textbf{16} \textit{in marg. cum unco} (comptes). \textbf{16 c.} \textit{abbr. pro} centimes. — \textbf{17 c.} \textit{abbr. pro} centimes.}} \\
      \cline{1-1}
      \leavevmode\fontsize{6.9pt}{7.9pt}\selectfont\strut [ Washing : 50 c.\strut &  \\
      \leavevmode\fontsize{6.9pt}{7.9pt}\selectfont\strut A Hénencourt, miroir de poche : 50 c. ]\strut &  \\
\end{tabular}}%
\par\vspace*{\fill}
% --- Registres de bas de page (1. Critique interne [si pertinent] ; 2. Loci paralleli en latin ; 3. Commentaire historique en retrait) ---
\noindent\rule{\textwidth}{0.5pt}\par\vspace{0.65mm}
{\fontsize{6.6pt}{7.7pt}\selectfont%
\noindent\normalfont \textbf{Clôture du carnet :} Dernier feuillet du \textit{Livre-Journal}, daté du lundi 30 avril 1917 (date qui correspond au \textit{terminus ad quem} inscrit par Hippolyte Combes sur l'étiquette de la couverture).\par\vspace{0.5mm}
\noindent\rule{22mm}{0.3pt}\par\vspace{0.45mm}
\noindent\normalfont\color{black!92} \textbf{5} (\textit{Warloy}) cf. supra f\textsuperscript{o}~37~v\textsuperscript{o} (p.~\pageref{fol:75}), f\textsuperscript{o}~90~v\textsuperscript{o} (p.~\pageref{fol:181}), \textit{et passim}. — \textbf{7} (\textit{Lt Gilbert}) cf. supra f\textsuperscript{o}~89~v\textsuperscript{o} (p.~\pageref{fol:179}), f\textsuperscript{o}~92~v\textsuperscript{o} (p.~\pageref{fol:185}). — \textbf{8} (\textit{Puchevillers}) cf. supra f\textsuperscript{o}~37~v\textsuperscript{o} (p.~\pageref{fol:75}), f\textsuperscript{o}~92~r\textsuperscript{o} (p.~\pageref{fol:184}), \textit{et passim}.\par\vspace{0.55mm}
\noindent\rule{38mm}{0.4pt}\par\vspace{0.5mm}
\noindent\hspace*{5mm}\parbox{\dimexpr\textwidth-5mm\relax}{\fontsize{6.6pt}{7.7pt}\selectfont\color{black!84}Dernière entrée datée du \textit{Livre-Journal} (\textbf{lundi 30 avril 1917}, correspondant à la borne finale inscrite sur l'étiquette de couverture, Couverture): par une radieuse journée de printemps, Hippolyte Combes parcourt le matin \textbf{Vadencourt}, \textbf{Warloy-Baillon}, \textbf{Hénencourt}, \textbf{Millencourt} et \textbf{Laviéville}, puis l'après-midi, en compagnie du \textbf{sous-lieutenant Gilbert}, les villages de \textbf{Puchevillers}, \textbf{Pierregot} et le chef-lieu de canton de \textbf{Villers-Bocage} (Somme).}\par}
\newpage
\phantomsection\label{fac:186}
% --- En-tête du fac-similé (Page impaire / droite : Photographie du manuscrit) ---
\noindent\rule{\textwidth}{0.5pt}\par\vspace{0.55mm}
\noindent\textbf{\fontsize{6.8pt}{7.8pt}\selectfont FAC-SIMILÉ EN VIS-À-VIS — FOLIO 93 RECTO} \hfill \textbf{\fontsize{6.8pt}{7.8pt}\selectfont AVRIL 1917}\par\vspace{0.25mm}
\noindent{\fontsize{6.8pt}{8.0pt}\selectfont\textit{Livre-Journal d'Hippolyte Combes} — 30 avril 1917 (Hérissart, Vadencourt, Warloy, Hénencourt, Millencourt, Lavieville, Puchevillers, Pierregot, Villers-Bocage, Toutencourt (Somme))}\par\vspace{0.55mm}
\noindent\rule{\textwidth}{0.35pt}\par\vspace*{\fill}
{\centering
  \setlength{\fboxsep}{1.2mm}%
  \setlength{\fboxrule}{0.4pt}%
  \fcolorbox{gray!55}{white}{%
    \includegraphics[width=0.95\textwidth,height=0.81\textheight,keepaspectratio]{\facsimiledir/20251126_111833.jpg}%
  }\par}
\vspace*{\fill}
\newpage

% =============================================================================
% ANNEXES DOCUMENTAIRES ET ÉPISTOLAIRES
% =============================================================================
\cleardoublepage
\appendix
\renewcommand{\currentmonthyear}{Annexe I — Correspondance}
\renewcommand{\currentfoliolabel}{Annexe I}
\titleformat{\chapter}[display]
  {\normalfont\normalsize\bfseries}
  {\textsc{Annexe\ \Roman{chapter}}}
  {0.8mm}
  {\large\color{black}}
\chapter{Correspondance de guerre inédite et pièces justificatives (1916–1918)}
\phantomsection\label{annexe:lettres}
{\small\sloppy

Les archives familiales d'Hippolyte Combes renferment huit pièces autographes majeures de la Grande Guerre qui éclairent directement la genèse et le contexte du \textit{Livre-Journal}.

\section{Document 1 : Copie autographe d'une lettre d'Hippolyte Combes à un capitaine, ancien collègue de la Faculté de Toulouse (Au front, 21 février 1916)}
\phantomsection\label{annexe:doc1}

\begin{quote}\footnotesize
« \underline{\underline{Copie}} \hfill Au front, le 21 février 1916\\[1mm]
Mon Capitaine,\\
Vous vous êtes peut-être demandé ce qu'a \textsuperscript{bien} pu devenir l'un de vos collègues \textsuperscript{de l'enseignement}, ancien camarade à la Faculté de Toulouse, et retrouvé, simple interprète à St-Omer, l'été dernier. Vous en souvient-il ?\\
À cette époque, je fus attaché à la 18\textsuperscript{e} division qui se formait dans les parages de Flesselles. Je fus choisi $[\![$comme$]\!]$ \textsuperscript{pour être} secrétaire au $[\![$membre$]\!]$ bureau de la mission française et vous ne fûtes pas étranger à ma nomination. Au bout de quinze jours, je me fatiguai de cette vie trop sédentaire et monotone et sur ma demande réitérée, je reçus une autre affectation. \textit{[verso]}\\
Depuis le mois d'août, je suis avec une compagnie d'\textit{A.S.C.} Mes officiers anglais sont de parfaits \textit{gentlemen}. Je n'ai \textsuperscript{aussi} qu'à me louer de l'officier interprète de la 18\textsuperscript{e} division [M. Pizzi est d'une haute intelligence et d'une distinction rare].\\
Vous doutez-vous, \textsuperscript{mon Capitaine}, que les $[\![$professeurs$]\!]$ \textsuperscript{universitaires, pourtant en majorité,} sont assez mal vus de quelques interprètes ? Ainsi le sergent Schumacher, qui m'a remplacé, \textsuperscript{en mon absence}, comme secrétaire au bureau de la mission, a osé dire, sans rire et devant témoins, que les professeurs sont de « rudes ânes » (\textit{sic}). Vous pensez bien que cela ne nous touche guère. Nous qui avons donné des preuves de notre intelligence et de notre culture générale, nous pouvons mépriser les boutades d'un autodidacte jaloux [N'est-il pas $[\![$un ancien$]\!]$ contremaître d'usine métallurgique \textsuperscript{et d'instruction primaire} ?].\\
Y aura-t-il bientôt un concours pour être nommé stagiaire ?\\[1mm]
\textit{[En marge gauche verticale :]} Veuillez agréer, mon Capitaine, l'expression de mes meilleurs sentiments et me croire votre bien dévoué et reconnaissant collègue — (signé) \textbf{H. C.} — »
\end{quote}

\noindent Cette copie autographe du 21 février 1916 documente le passage d'Hippolyte Combes par Saint-Omer et Flesselles à l'été 1915, son rattachement à la 18\textsuperscript{e} Division britannique sous l'officier interprète divisionnaire M.~Pizzi et son affectation depuis août 1915 à une compagnie d'\textit{Army Service Corps}.

\section{Document 2 : Rapport d'Hippolyte Combes au sous-lieutenant Ferran sur l'incident de Morlancourt (mars 1916) et requête au sujet de l'indemnité des interprètes}
\phantomsection\label{annexe:doc2}

\begin{quote}\footnotesize
« \textbf{1.} [p.~26] \textbf{L'interprète Combes au sous-lieutenant Ferran, officier de liaison à la 18\textsuperscript{e} division}\\[1mm]
J'ai l'honneur de porter à votre connaissance un incident qui s'est produit lundi 2[?] mars, au village de Morlancourt.\\
La matinée \textsuperscript{de ce jour}, j'avais, dans ce village, avec un de mes officiers, le lieutenant May, fait un cantonnement et nous avions \textsuperscript{choisi} une ferme que nous avait indiquée le maire de la localité. Or, l'après-midi, en l'absence de mon officier, parti au devant de notre compagnie, un lieutenant du \textit{6th Liverpool Regiment} (5\textsuperscript{th} division), \textsuperscript{dont j'ignore le nom}, arriva dans la même ferme, et \textsuperscript{après avoir inspecté les lieux}, dit à la propriétaire qu'il comptait y cantonner ses hommes. Comme cette dernière \textsuperscript{lui} faisait comprendre qu'on avait déjà retenu cette ferme, le lieutenant s'écria \textsuperscript{qu'il était chez lui et que l'interprète n'était} à prendre pour un sous-officier \textbf{2.} [p.~27] et me donna l'ordre de chercher un autre cantonnement. Devant ma protestation ferme, mais respectueuse, l'officier reprit sur un ton dédaigneux :\\
\textit{« This billet is all right for me, it shall be mine — I don't want to discuss with an interpreter. »} Je n'insistai pas et m'éloignai aussitôt.\\
Deux jours plus tard, mon capitaine a reçu une plainte dans laquelle l'officier en question \textsuperscript{(sans doute)} prétend que j'ai été insolent et que je ne l'ai pas salué. Je proteste contre ces deux allégations avec la plus grande énergie.\\[1mm]
[Je ne vous avais pas mis au courant de ces faits, car je ne croyais pas que l'incident aurait de suites]\\
\hspace*{0pt}\hfill \textbf{H. Combes}, \textit{A.S.C.} 18\textsuperscript{e} D. »
\end{quote}

\noindent Au verso du second feuillet figure la fin d'une minute de lettre adressée à un parlementaire (« Monsieur le Président ») pour la défense de l'indemnité des interprètes militaires :
\begin{quote}\footnotesize
« \textbf{3.} [p.~37] Nous sommes très utiles. Les Anglais ne sauraient se passer de nous, car très peu \textsuperscript{d'entre eux} parlent assez \textsuperscript{bien} le français pour se faire comprendre. Ils sont les premiers à reconnaître les grands services que nous leur rendons.\\
Aussi, Monsieur le Président, si vous voulez bien vous intéresser à notre sort et faire valoir \textsuperscript{à qui de droit} les considérations qui militent en faveur du maintien intégral de notre \textsuperscript{faible} indemnité.\\
Daignez agréer, Monsieur le Président, avec toutes mes excuses pour mon importunité, l'hommage de mon très profond respect.\\
\noindent \textbf{Interprète H. C.} etc. \hfill \textbf{H. C.}, prof. au coll. de Nogent-le-R. »
\end{quote}

\section{Document 3 : Lettre confidentielle intégrale d'Hippolyte Combes à sa sœur Françoise au Carmel de Périgueux (Au front, vendredi 16 juin 1916)}
\phantomsection\label{annexe:doc3}

\begin{quote}\footnotesize
« \textit{[En marge supérieure gauche verticale :]} \textbf{Brûle ce feuillet — Je ne veux pas que tu le gardes dans tes papiers.}\\[1mm]
\hspace*{0pt}\hfill Au front, vendredi 16 juin 1916\\
Ma bien chère Sœur,\\
Ta dernière lettre m'est parvenue assez rapidement. D'ordinaire toutes les correspondances du midi mettent au moins 4 et 5 jours pour arriver à destination. Tu es toujours très contente et en parfaite santé. Que cela dure ainsi ! Amen !\\
Ici, c'est à peu près la même vie, monotone et banale. On végète, tout simplement. Je n'ai jamais aimé la campagne : maintenant je l'ai en horreur. Avec les beaux \textit{[p.~2]} jours, la vie dans de petits villages devrait être plus facilement supportable. On se consume de nostalgie et d'ennui. Voudrait-on s'amuser ? on ne trouve rien. MM. les offs, eux, vont chercher des distractions dans les grandes villes les plus rapprochées.\\
Ah ! quand enfin, démuselés, nous pourrons gueuler tout ce que nous avons vu, on ne voudra pas nous croire. J'ai entendu de mauvais Français exprimer déjà du regret à la pensée que la guerre finira tôt ou tard. C'est révoltant. À combien de gens aura-t-elle profité ?...\\
À cause de la censure, je dois me taire. \textit{[p.~3]}\\
À Toulouse, j'ai retrouvé des camarades en permission — nous avons causé, j'ai constaté avec plaisir qu'ils pensent absolument comme moi. Mon ami, le docteur de B—, dont la mère m'est tout dévouée, disait : « Jamais un officier n'entrera dans mon salon, après la guerre !.. » Sa mère et sa jeune femme, après l'avoir écouté, l'approuvaient en pensées, et je faisais chorus.\\
En tout, nous en avons tous plus qu'assez de cette vie qui se prolonge. Pour ma part, j'aimerais autant être sous terre que de rester encore une année en esclavage. On est aigri, découragé, \textit{[p.~4]} on ne prend goût à rien. Toutes ces abstractions : gloire, honneur etc. quel tas de blagues ! Vivre avant tout. Dans la mort, il n'y a plus d'alliés ni d'ennemis... Être gouverné par 1 empereur, un roi, ou n'avoir qu'un vulgaire président, c'est kif-kif — on paie les impôts — et puis qu'on nous f. la paix, tonnerre de — diable !\\
Je vais bien. Ne m'écris qu'une fois par mois, c'est plus qu'assez. Si tu savais avec quelle indifférence je parcours tes pages.\\
Je me fous de tout et de tous — Je ne demande que la liberté ou la mort.\\[1mm]
\textit{[En marge gauche verticale :]} Bien à toi, quand même — \textbf{Hippolyte} »
\end{quote}

\section{Document 4 : Poème anglais intégral de Beatrice Chilcott (\textit{Wesley Band of Hope})}
\phantomsection\label{annexe:doc4}

\begin{quote}\footnotesize\itshape
What, love a passion ! Heaven forbid,\\
That man should call true love by such a name,\\
Passion is but the creature of an hour,\\
A glowing ember, a bright, short lived flame ;\\
A fire, kindled from a tiny spark,\\
Glowing and live, but soon the brightness wanes ;\\
Colder and paler still, the embers grow\\
Till nought but ashes, of the fire remains.\\[2mm]
No passion lives for ever ; so\\
True love is not a passion I contend,\\
Love is, a mutual affinity,\\
A strong affection 'twixt two souls that blend,\\
Placid and pure, and deep as is the sea.\\
Love can make two souls one in heart, in thought,\\
And love comes not at will ; it is not hired,\\
Cannot be sold, neither can it be bought.\\[2mm]
A mystic power, unfathomed still,\\
Which baffles all the probing powers of man.\\
Potent, enthralling, it shapes men's destinies,\\
And silent ruled since first the world began ;\\
Heav'ns greatest gift, love is the source of life\\
Which sweetens all the bitterness and pain ;\\
And ever points toward the divine love\\
Beside which, human love is poor and vain.\\[1mm]
\normalfont\footnotesize\hspace*{0pt}\hfill — \textbf{Beatrice Chilcott}, \textit{Wesley Band of Hope}.
\end{quote}

\section{Document 5 : Correspondance de l'armistice à Marguerite Calmels-Combes (Saint-Pol, 12 novembre 1918 ; Paris, 25–26 novembre 1918 ; Les Andelys, 28 novembre – 2 décembre 1918)}
\phantomsection\label{annexe:doc5}

\begin{enumerate}
  \item \textbf{Carte de Saint-Pol-sur-Ternoise au lendemain de l'Armistice (mardi 12 novembre 1918) :} Adressée à \textit{« Madame Combes, 10 rue Alfred Guibert, Millau (Aveyron) »} (cachet postal : \textit{Trésor et Postes 115, 13-11-18} ; franchise militaire : \textit{Field Service Post Card}) :
  \begin{quote}\scriptsize
  « Saint-Pol, mardi 12 nov\textsuperscript{bre} 1918. Affectueuses pensées. Bons baisers. \textbf{Hippolyte}. »
  \end{quote}
  \item \textbf{Première lettre-carte de Paris, Bibliothèque Sainte-Geneviève (lundi 25 novembre 1918, 11 h matin) :}
  \begin{quote}\scriptsize
  « Paris, lundi 25 novembre 1918, 11 h. matin.\\
  Ma Chérie, ce matin à 6 h., enfin je me suis retrouvé dans cette ville que je préfère à toutes les autres capitales du monde. Je ne repartirai que demain. Je n’ai pas voulu descendre chez tes cousines, à Neuilly, afin de ne point les déranger, et aussi pour me sentir plus libre. Cet après-midi j’irai les voir. Toutefois, si l’on me prie, j’accepterai à déjeuner demain seulement. [Avec ?] plaisir, je suis comme un voyageur $[\![$[illis.]$]\!]$ [illis.] pays [neuf ?]. Je n’ai pu me loger [ajout suscrit : faute de place] dans mon hôtel de prédilection, University Hôtel, 160, rue St-Jacques, tout près de l’École de Droit. J’ai pris une jolie chambre à l’Hôtel du Panthéon, en face, sur le devant. Tout à l’heure, je me suis débarbouillé et bien reposé ! Je t’écris dans la grande salle de la Bibliothèque Geneviève. Que de fois, [illis.] m’y a-t-elle accompagné ! Elle [illis.] à côté de moi. Comme je serais heureux de te posséder, de te sentir tout près de moi ! Courage et confiance, notre séparation finira bientôt.\\
  Ma nouvelle affectation, à défaut d’avantages plus appréciables, aura du moins celui de m’avoir bien rapproché de Paris, 102 kilomètres environ. De jeudi au dimanche, sans doute j’obtiendrai facilement des permissions de 24 heures. Ainsi je pourrai me rendre soit à Rouen, soit dans la capitale et aussi — et surtout — à Nogent-le-Rotrou. J’éprouverai une satisfaction profonde quand je reverrai mon chez moi du Perche, quitté depuis près de 4 ans. La maison que j’habite à Nogent, 17, rue Saint-Hilaire, est autrement jolie et pittoresque que celle de Toulouse, rue de Languedoc, [24 ?]. Elle a un petit air flamand qui lui donne une certaine originalité. [Santé ?] toujours parfaite. Je ne me suis jamais mieux porté, me semble-t-il. [Enfin ?] je me ressemble tout à fait, malgré les [illis.], et attends sans fièvre les événements. Je ne tarderai pas à recouvrer ma liberté entière et définitive. Ne t’alarme point au sujet de notre fillette. Elle prendra le dessus et grossira à l’aise à partir du 3ème mois.\\
  \textit{[En marge :]} Mes plus doux baisers. Ton petitou qui t’aime si tendrement — \textbf{Hippolyte}. Tu peux m’écrire à l’adresse que je t’ai indiquée dimanche. »
  \end{quote}
  \item \textbf{Deuxième lettre-carte de Paris, Bibliothèque Sainte-Geneviève (mardi 26 novembre 1918, 11 h matin) :}
  \begin{quote}\scriptsize
  « Paris, mardi 26 novembre 1918, 11 h. matin.\\
  Ma Margueritou tant aimée, Me voici de nouveau installé dans la grande salle de la Bibliothèque Sainte-Geneviève. Comme je m’y plais ! Tant de souvenirs me la rendent chère ! Devant moi, sur la table, j’ai des revues françaises et un volume de Maeterlinck. Mais, au lieu de lire, je fais ma petite correspondance. J’ai passé une excellente nuit. Je me suis bien reposé. Ce soir, je continuerai et terminerai mon voyage. Je prendrai l’express du Havre, à [illis.] [groupe traversé de traits : gare St-Lazare, j’arriverai] à certainement après minuit. Peu m’importe ! Partant par le dernier train, j’aurai pu rester qq heures de plus dans la Capitale.\\
  Hier, j’ai fait une très longue visite aux Cousines Faure. M[me ?] [illis.] descendait l’escalier, comme j’arrivais sur le 2ème étage, elle sortait avec ses [voisines ?], les dames [illis.], auxquelles elle m’a présenté. Elle a renoncé à sa promenade, pour me recevoir. Je l’avais priée de n’en rien faire. Germaine était aussi là, elle préparait ses devoirs. J’ai pu causer durant plusieurs heures avec Marguerite. Elle m’a raconté son pèlerinage bien triste à [Moyenneville ?]. Nous avons longuement parlé du pauvre défunt. [Toi ?] et notre fillette, naturellement, [avez ?] été l’inépuisable sujet de nos conversations. Ta Cousine [veut ?] t’écrire elle-même ce qu’elle pense et qu’elle te conseillerait. Je suis entièrement de son avis. Je te résumerai en qq mots ce qu’elle m’a dit à propos de notre mignonne Valentine. Tu ne dois point t’inquiéter, ni t’énerver ; avec le temps, l’enfant prendra le dessus et deviendra magnifique, comme le petit Bonhomme, que tu as vu, paraît-il, à Paris, chez les locataires du rez-de-chaussée, 26, rue d’[Angoulême ?] ; il était alors bien chétif, bien menu. Sois calme, patiente et surtout écoute ton excellente mère : apprends à dormir. Le sommeil t’est aussi nécessaire que le lait à la fillette.\\
  \textit{[En marge :]} Mes plus tendres baisers pour toutes deux. Ton mari qui t’aime [tendrement ?]. »
  \end{quote}
  \item \textbf{Lettres-cartes de l'École militaire préparatoire des Andelys (jeudi 28 novembre et lundi 2 décembre 1918) :}
  \begin{quote}\scriptsize
  « \textit{[1. Carte postale du 28 novembre 1918 :]} Les Andelys, jeudi soir 28 nov\textsuperscript{bre} 1918, École Militaire.\\
  Ma jolie Perlounette, Aujourd'hui encore, liberté entière. Je suis allé ce matin à l'École. L'emploi du temps pour l'anglais n'est pas encore fixé. Jusqu'à présent, on n'enseignait pas cette langue. Quand mon horaire aura été établi, samedi ou lundi, je commencerai mon nouveau service. À ma honte, je t'avouerai que le plus tard sera le mieux, je ne suis nullement pressé. Pour le moment, je jouis en paix d'un repos complet. Oh ! le paresseux ! Tu sais, il fait délicieux dans ma grande chambre, rue de la Sous-Préfecture. Affalé sur mon fauteuil — un vieux pouf — les jambes sur une chaise, je rumine, je lis et — je fume ! Hier soir, après dîner, j'ai travaillé deux heures à peine et je me suis mis au lit de bonne heure. Réveillé ce matin à 7 h. ½, j'ai flâné sous les couvertures une grosse demi-heure. Ensuite j'ai procédé à ma toilette qui a été très minutieuse, donc très longue. Parti à 9 h. pour l'École (20 minutes de marche environ), j'étais de retour avant 11 h. Cet après-midi, le temps étant très couvert, je ne suis pas sorti pour visiter la ville en détail.\\
  \textit{[Recto :]} Non loin de ma rue, se trouve l'église [Notre-Dame du Grand-Andely], admirable ! J'y suis entré deux minutes. Il y a de merveilleux vitraux (XIII\textsuperscript{e} et XIV\textsuperscript{e}) qui reproduisent les principaux faits de l'existence de St\textsuperscript{e} Clotilde. Cette reine avait fondé au VI\textsuperscript{e} siècle un monastère important aux Andelys. C'est là l'origine de cette petite cité. — J'ai des collègues fort sympathiques. Ce sont tous des universitaires, mobilisés eux aussi, et ne désirant qu'une chose : être libérés. Comme nous sommes des privilégiés (la plupart ont près de quarante ans ou ont dépassé la quarantaine), nous ne reprendrons nos fonctions d'avant-guerre que des premiers. En ce qui me concerne, ce sera sûrement en janvier au plus tard. Amen ! »\\[1mm]
  « \textit{[2. Carte-lettre du 2 décembre 1918 :]} École Militaire des Andelys, lundi 2 décembre 1918.\\
  Chérie, cette carte, écrite à l'encre noire, vient directement de ma chambre militaire. Tu le sais, je ne me sers que d'encre violette, dans ma chambre civile. Je ne commencerai pas les cours d'anglais avant demain mardi. Samedi 21 sera le dernier jour de classe de l'année. À mon retour de permission, en janvier, je serai presque à la veille du plus beau, du plus heureux jour de ma vie, celui où je serai enfin démobilisé. Je n'ai jamais autant désiré une chose, depuis que j'existe. En réalité, je n'aurai pas trop à me plaindre de mon séjour ici, étant donné les avantages considérables attachés à ma fonction : liberté, tranquillité et loisirs nombreux. Par moments, je me figure que je suis déjà tout à fait libre, et alors la pensée de mon bonheur très prochain fait battre mon cœur avec force. Notre séparation n'est plus qu'une question de semaines. Oh ! Dieu, la liberté, je vais donc la recouvrer ! Cette certitude me transporte d'allégresse.\\
  \textit{[En marge gauche verticale :]} Mes plus doux baisers pour Toi et pour la fillette. Ton mari qui t'aime tant, \textbf{Hippolyte}. »
  \end{quote}
\end{enumerate}
\par}

\cleardoublepage
\renewcommand{\currentmonthyear}{Annexe II — État des services}
\renewcommand{\currentfoliolabel}{Annexe II}
\chapter{État signalétique des services et table synoptique des folios}
{\small\sloppy

\section{Résumé de la fiche matricule d'Hippolyte Combes (Classe 1896, Matricule 911, Bureau de recrutement de Toulouse)}

\begin{figure}[htbp]
  \centering
  \begin{minipage}[b]{0.46\textwidth}
    \centering
    \includegraphics[height=40mm,keepaspectratio]{illustrations/mariage_millau_1917.jpg}
    \par\vspace{1mm}{\scriptsize\textit{Fig. 3.} Mariage d'Hippolyte Combes et Marguerite Calmels à Millau (25 octobre 1917).}
  \end{minipage}\hfill
  \begin{minipage}[b]{0.46\textwidth}
    \centering
    \includegraphics[height=40mm,keepaspectratio]{illustrations/famille_combes_1925.jpg}
    \par\vspace{1mm}{\scriptsize\textit{Fig. 4.} Hippolyte et Marguerite Combes entourés de leurs enfants (vers 1925).}
  \end{minipage}
\end{figure}

{\raggedright\noindent D'après le registre matricule du recrutement de Toulouse et les actes d'état civil de Millau :\par}
\begin{itemize}
  \item \textbf{État civil :} \textbf{Marie Lucien Hippolyte Combes}, né le \textbf{14 août 1876} à \textbf{Millau} (Aveyron), rue des Barris (inscrit sous la date du 22 août 1876 sur la fiche matricule), et décédé le \textbf{2 septembre 1970} à l'\textbf{Hôtel du Centre} à \textbf{Saint-Chély-d'Apcher} (Lozère ; inhumé à Millau), fils de \textbf{François Augustin Hippolyte Combes} et de \textbf{Marie Anaïs Alric}, domiciliés à Toulouse (47, rue du Taur). Signalement : cheveux et sourcils bruns, yeux châtains, taille 1,66 m. Professeur d'anglais au collège de Nogent-le-Rotrou (Eure-et-Loir) puis au lycée Ingres de Montauban.
  \item \textbf{Parcours militaire et campagnes (1914–1919) :}
    \begin{itemize}
      \item Exempté après deux ans d'ajournement pour « faiblesse générale » (numéro de tirage 74, canton de Toulouse-Nord) ; classé \textbf{Service armé} par le conseil de révision d'Eure-et-Loir le \textbf{18 décembre 1914}.
      \item Appelé à l'activité le \textbf{20 mars 1915} ; arrivé au corps le \textbf{21 mars 1915} au \textbf{10\textsuperscript{e} Régiment d'Artillerie à pied} (39\textsuperscript{e} batterie) à Toulon (campagne simple à l'intérieur du 21 mars au 3 mai 1915).
      \item Détaché à la \textbf{Mission militaire française attachée à l'Armée britannique} aux armées (campagne double) du \textbf{4 mai 1915 au 11 novembre 1918} ; passé administrativement au \textbf{19\textsuperscript{e} Escadron du Train des Équipages Militaires} le \textbf{15 octobre 1915}.
      \item Affecté auprès de la Mission française de Saint-Omer, de la 18\textsuperscript{e} Division britannique (Flesselles, juillet 1915), de la \textbf{152\textsuperscript{e} Compagnie d'\textit{Army Service Corps}} de la 18\textsuperscript{e} Division (à partir d'août 1915), puis des services d'approvisionnement (\textit{Foraging} / \textit{Army Purchase Board}) dans la Somme (novembre 1916 – avril 1917).
      \item Affecté le \textbf{24 novembre 1918} à l'\textbf{École militaire préparatoire des Andelys} comme professeur d'anglais ; nommé \textbf{caporal} le \textbf{4 décembre 1918}.
      \item Mis en congé illimité de démobilisation le \textbf{7 février 1919} (se retire à Nogent-le-Rotrou, 17, rue Saint-Hilaire, puis à Montauban le 24 janvier 1920 comme professeur au lycée Ingres).
    \end{itemize}
\end{itemize}

\section{Table synoptique des folios du \textit{Livre-Journal}}

{\fontsize{6.2pt}{7.3pt}\selectfont
\renewcommand{\arraystretch}{1.12}
\setlength{\tabcolsep}{2.5pt}
\noindent\begin{longtable}{@{}>{\raggedright\arraybackslash}p{21mm}>{\raggedright\arraybackslash}p{12mm}>{\raggedright\arraybackslash}p{26mm}>{\raggedright\arraybackslash}p{67mm}@{}}
\toprule
\textbf{Folio} & \textbf{Page} & \textbf{Date} & \textbf{Lieu \& Intitulé analytique} \\
\midrule
\endfirsthead
\toprule
\textbf{Folio} & \textbf{Page} & \textbf{Date} & \textbf{Lieu \& Intitulé analytique (suite)} \\
\midrule
\endhead
\bottomrule
\endfoot
Couverture (plat supérieur) & p.~\pageref{fol:1} & 25 juillet 1916 – 30 avril 1917 & \textbf{Somme, Flandres, Artois, Paris, Toulouse, Millau} — Couverture rouge à motif losangé et étiquette bleue autographe : Livre-Journal, du 25 juillet 1916 au 30 avril 1917 \\
Folio 1 recto & p.~\pageref{fol:2} & Mardi 25 juillet 1916 & \textbf{Lynde et Sercus (Nord)} — Mardi 25 juillet 1916 — Lynde et Sercus : rassemblement avec le lieutenant Curtis devant l'église de Sercus et cantonnement de 36 officiers des Royal Fusiliers \\
Folio 1 verso & p.~\pageref{fol:3} & Mercredi 26 juillet 1916 & \textbf{Lynde (Nord)} — Mercredi 26 juillet 1916 — Lynde : préparation de 36 chambres d'officiers, château de la brasserie et annonce d'un nouveau départ \\
Folio 2 recto & p.~\pageref{fol:4} & Jeudi 27 juillet 1916 & \textbf{Lynde, Renescure et Hazebrouck (Nord)} — Jeudi 27 juillet 1916 — De Lynde à Renescure et Hazebrouck : châteaux du Moyen Âge, jardins potagers et achats de cartes postales d'Ypres et de Renescure \\
Folio 2 verso & p.~\pageref{fol:5} & Vendredi 28 juillet 1916 & \textbf{De Lynde vers Saint-Jans-Cappel par Hazebrouck (Nord)} — Vendredi 28 juillet 1916 — Départ de Lynde (Mme Donze) et lente chevauchée sous la chaleur par Hazebrouck jusqu'à Saint-Jans-Cappel \\
Folio 3 recto & p.~\pageref{fol:6} & Samedi 29 juillet 1916 & \textbf{Saint-Jans-Cappel et le Mont-Noir (Nord)} — Samedi 29 juillet 1916 — Tournée avec le lieutenant Curtis à Saint-Jans-Cappel et achat de fourrages au Mont-Noir \\
Folio 3 verso & p.~\pageref{fol:7} & Dimanche 30 et lundi 31 juillet 1916 & \textbf{Saint-Jans-Cappel et Bailleul (Nord)} — Messe dominicale à Saint-Jans-Cappel et corvée de fourrage vert vers Bailleul \\
Folio 4 recto & p.~\pageref{fol:8} & Mardi 1er août 1916 & \textbf{Saint-Jans-Cappel et Bailleul (Nord)} — Première promenade à Bailleul : la Grand-Place et le beffroi \\
Folio 4 verso & p.~\pageref{fol:9} & Mercredi 2 et jeudi 3 août 1916 & \textbf{Bailleul (Nord)} — Mercredi 2 et jeudi 3 août 1916 — Promenades dans Bailleul, maisons flamandes et dentellières au fuseau \\
Folio 5 recto & p.~\pageref{fol:10} & Vendredi 4 août 1916 (matin et après-midi) & \textbf{Saint-Jans-Cappel, Bailleul, Steenwerck, La Croix-du-Bac (Nord)} — Chevauchée à La Croix-du-Bac par Steenwerck et repos dans le « garret » \\
Folio 5 verso & p.~\pageref{fol:11} & Vendredi 4 août 1916 (soir) & \textbf{Saint-Jans-Cappel (Nord)} — Soirée solitaire à la ferme et prière du soir de la famille flamande \\
Folio 6 recto & p.~\pageref{fol:12} & Samedi 5 août 1916 & \textbf{De Saint-Jans-Cappel à La Croix-du-Bac / Steenwerck (Nord)} — Départ du convoi avec les A.S.C. men et étape à La Croix-du-Bac \\
Folio 6 verso & p.~\pageref{fol:13} & Dimanche 6 août 1916 & \textbf{La Croix-du-Bac et Bac-Saint-Maur (Nord)} — Dimanche 6 août 1916 — Au refilling point sur les bords de la Lys avec le lieutenant Curtis, rencontre de Maurin et concert militaire \\
Folio 7 recto & p.~\pageref{fol:14} & Lundi 7 août 1916 & \textbf{Steenwerck et route de Bac-Saint-Maur (Nord)} — Quête de cantonnements avec le lieutenant Curtis, départ inopiné de la compagnie et nouvelle chambre \\
Folio 7 verso & p.~\pageref{fol:15} & Mardi 8 et mercredi 9 août 1916 & \textbf{Route de Bac-Saint-Maur, Steenwerck (Nord)} — Deux journées de canicule, de poussière et de moustiques dans la vallée de la Lys \\
Folio 8 recto & p.~\pageref{fol:16} & Jeudi 10 août 1916 & \textbf{Bac-Saint-Maur et Sailly-sur-la-Lys (Nord / Pas-de-Calais)} — Promenade à Sailly-sur-la-Lys, l'église bombardée et conversation avec l'éclusier sur le général Maxse \\
Folio 8 verso & p.~\pageref{fol:17} & Vendredi 11 août 1916 & \textbf{Steenwerck / Erquinghem-Lys (Nord)} — Vendredi 11 août 1916 — Bombardement meurtrier autour de la grande blanchisserie Barbry à Steenwerck \\
Folio 9 recto & p.~\pageref{fol:18} & Samedi 12 août 1916 & \textbf{Rives de la Lys, Bac-Saint-Maur (Nord)} — Sieste au bord de la Lys dans un champ de luzerne face aux blés mûrs \\
Folio 9 verso & p.~\pageref{fol:19} & Dimanche 13 et lundi 14 août 1916 & \textbf{Bac-Saint-Maur et direction d'Armentières (Nord)} — Attente vaine d'une lettre de Marguerite, « pot-au-noir » et orage sur la Lys \\
Folio 10 recto & p.~\pageref{fol:20} & Mardi 15 août 1916 (Assomption) & \textbf{Bac-Saint-Maur / Steenwerck (Nord)} — Journée pluvieuse du 15 août dans la mansarde et comptes de blanchissage \\
Folio 10 verso & p.~\pageref{fol:21} & Vendredi 18 et samedi 19 août 1916 & \textbf{Steenwerck, Fleurbaix et secteur d'Armentières (Nord)} — Vendredi 18 et Samedi 19 août 1916 — Giboulées d'août, perception de la solde et chevauchée sous l'orage avec le lieutenant Curtis \\
Folio 11 recto & p.~\pageref{fol:22} & Samedi 19 août 1916 (suite) & \textbf{Erquinghem-Lys, Fleurbaix et Armentières (Pas-de-Calais / Nord)} — Traversée de Fleurbaix et d'Armentières dévastées par les obus \\
Folio 11 verso & p.~\pageref{fol:23} & Dimanche 20 et lundi 21 août 1916 & \textbf{Erquinghem-Lys et Steenwerck (Nord)} — Café chez M. Barbry avec les adjudants Lupin et Schumacher et préparatifs de départ \\
Folio 12 recto & p.~\pageref{fol:24} & Mardi 22 août 1916 & \textbf{Steenwerck, Bac-Saint-Maur et Erquinghem-Lys (Nord)} — Mardi 22 août 1916 — Bière chez un charcutier avec Maurin évoquant les beautés d'Antibes et envoi d’un sac de linge à Nogent \\
Folio 12 verso & p.~\pageref{fol:25} & Mercredi 23 août 1916 & \textbf{Estaires (Nord)} — Chevauchée jusqu'à Estaires : l'église Saint-Vaast, l'Hôtel de Ville et l'hygiène par l'exercice \\
Folio 13 recto & p.~\pageref{fol:26} & Jeudi 24 août 1916 & \textbf{Steenwerck (Nord)} — Jeudi 24 août 1916 — Église de Steenwerck, transport des sacs sous la tente de l'office par Fisher et adieux aux Flandres \\
Folio 13 verso & p.~\pageref{fol:27} & Vendredi 25 août 1916 & \textbf{De Steenwerck et Bailleul à Saint-Pol-sur-Ternoise (Pas-de-Calais)} — Embarquement ferroviaire à Bailleul et descente par Hazebrouck et Lillers vers Saint-Pol-sur-Ternoise \\
Folio 14 recto & p.~\pageref{fol:28} & Vendredi 25 août 1916 (suite) & \textbf{De Bailleul à Saint-Pol et Marquay (Pas-de-Calais)} — Mélancolie du voyage en train au son de l'harmonica et arrivée nocturne à Marquay \\
Folio 14 verso & p.~\pageref{fol:29} & Samedi 26 août 1916 & \textbf{Marquay et Bailleul-aux-Cornailles (Pas-de-Calais)} — Promenade sur le plateau de Bailleul-aux-Cornailles et portrait au vitriol de la logeuse \\
Folio 15 recto & p.~\pageref{fol:30} & Lundi 28 et mardi 29 août 1916 & \textbf{Marquay (Pas-de-Calais)} — Lundi 28 et Mardi 29 août 1916 — Bourbiers d'Artois, corvées de fourrage avec le lieutenant Curtis et entrée en guerre de la Roumanie \\
Folio 15 verso & p.~\pageref{fol:31} & Mercredi 30 août 1916 & \textbf{Marquay et Bailleul-aux-Cornailles (Pas-de-Calais)} — Tempête de vent dans les grands arbres du voisinage et fourrage sous les ondées \\
Folio 16 recto & p.~\pageref{fol:32} & Jeudi 31 août 1916 & \textbf{Marquay, Ramecourt, Saint-Michel, Saint-Pol-sur-Ternoise et Ostreville (Pas-de-Calais)} — Marche à pied par Ramecourt jusqu'à Saint-Pol-sur-Ternoise et souvenir de saint Benoît Labre \\
Folio 16 verso & p.~\pageref{fol:33} & Vendredi 1er septembre 1916 & \textbf{Marquay et Ligny-Saint-Flochel (Pas-de-Calais)} — L'église médiévale de Ligny-Saint-Flochel et les « manœuvres spéciales » britanniques \\
Folio 17 recto & p.~\pageref{fol:34} & Samedi 2 et dimanche 3 septembre 1916 & \textbf{Marquay et Saint-Pol-sur-Ternoise (Pas-de-Calais)} — Rosée matinale dans les terres argileuses, liste de blanchissage en anglais et marche à Saint-Pol \\
Folio 17 verso & p.~\pageref{fol:35} & Lundi 4 et mardi 5 septembre 1916 & \textbf{Marquay et Roëllecourt (Pas-de-Calais)} — Courses à Roëllecourt, préparatifs inutiles et déception du départ retardé \\
Folio 18 recto & p.~\pageref{fol:36} & Dimanche 3 (addition) / mardi 5 septembre 1916 & \textbf{Marquay et Saint-Pol-sur-Ternoise (Pas-de-Calais)} — Champs d'avoine sous la brise et achat de la carte du département de la Somme \\
Folio 18 verso & p.~\pageref{fol:37} & Mercredi 6 et jeudi 7 septembre 1916 & \textbf{Marquay, Bailleul-aux-Cornailles, Saint-Pol et Saint-Michel-sur-Ternoise (Pas-de-Calais)} — Conversation avec Mme Vauchel et halte au cimetière de Saint-Michel-sur-Ternoise \\
Folio 19 recto & p.~\pageref{fol:38} & Vendredi 8 septembre 1916 & \textbf{Marquay et ferme de Ryonval (Pas-de-Calais)} — Attente des fourgons des « Bedfords » et promenade au clair de lune à Ryonval \\
Folio 19 verso & p.~\pageref{fol:39} & Samedi 9 septembre 1916 & \textbf{De Marquay et Ligny-Saint-Flochel à Monts-en-Ternois (Pas-de-Calais)} — Départ de Marquay, l'école de mortiers de tranchée et l'ancien monastère de Monts-en-Ternois \\
Folio 20 recto & p.~\pageref{fol:40} & Dimanche 10 septembre 1916 & \textbf{De Monts-en-Ternois à Oppy et Beaudricourt (Pas-de-Calais)} — Nuit devant le fournil à Monts-en-Ternois et étape au hameau d'Oppy (Beaudricourt) \\
Folio 20 verso & p.~\pageref{fol:41} & Lundi 11 septembre 1916 & \textbf{D'Oppy (Beaudricourt) à Lucheux, Thièvres et Raincheval (Pas-de-Calais / Somme)} — Longue chevauchée d'Oppy à Raincheval par Lucheux (l'Édit des postes de Louis XI) et Thièvres \\
Folio 21 recto & p.~\pageref{fol:42} & Mardi 12 et mercredi 13 septembre 1916 & \textbf{Raincheval (Somme)} — Admission permanente au mess des officiers du capitaine West et promenade en forêt de Raincheval \\
Folio 21 verso & p.~\pageref{fol:43} & Jeudi 14 et vendredi 15 septembre 1916 & \textbf{Raincheval et Acheux-en-Amiénois (Somme)} — Visite à la Mission française au parc du château d'Acheux et exercices de grenades à Raincheval \\
Folio 22 recto & p.~\pageref{fol:44} & Dimanche 17 septembre 1916 & \textbf{Raincheval (Somme)} — Les grands chênes du bois de Raincheval, le vieux collectionneur de faïences et l’office méthodiste (Wesleyans) \\
Folio 22 verso & p.~\pageref{fol:45} & Lundi 18 septembre 1916 & \textbf{Raincheval, Hédauville, Forceville et Acheux-en-Amiénois (Somme)} — Chevauchée sous l'averse par Hédauville et Forceville jusqu'à Acheux et malaise gastrique \\
Folio 23 recto & p.~\pageref{fol:46} & Jeudi 21 septembre 1916 & \textbf{Raincheval et Acheux-en-Amiénois (Somme)} — Routes transformées en ornières et perception de la solde à la hutte du château d'Acheux \\
Folio 23 verso & p.~\pageref{fol:47} & Vendredi 22 septembre 1916 & \textbf{Raincheval (Somme)} — Match de football des « Bedfs », annonce du départ (« we are moving demain ») et nuit sombre \\
Folio 24 recto & p.~\pageref{fol:48} & Samedi 23 septembre 1916 & \textbf{Raincheval, Léalvillers, Acheux et Forceville (Somme)} — Samedi 23 septembre 1916 — Marche vers Léalvillers, message apporté en auto par le lieutenant Pizzi et bivouac sous la tente à Forceville \\
Folio 24 verso & p.~\pageref{fol:49} & Dimanche 24 septembre 1916 & \textbf{Forceville et Acheux-en-Amiénois (Somme)} — Dimanche 24 septembre 1916 — Construction d'une hutte derrière la haie d'aubépine à Forceville et préparatifs du départ en permission \\
Folio 25 recto & p.~\pageref{fol:50} & Lundi 25 septembre 1916 (matin et après-midi) & \textbf{Forceville et gare d'Acheux-Varennes (Somme)} — Lundi 25 septembre 1916 — Réveil à 7 h, corvée de green forage à Forceville avec le Lt. Ford et départ en permission à la gare d'Acheux-Varennes \\
Folio 25 verso & p.~\pageref{fol:51} & Lundi 25 septembre 1916 (soir et nuit) & \textbf{Gézaincourt, Amiens, Longueau et train vers Paris} — De Gézaincourt à Amiens et raid nocturne de « Taubes » sur la gare de Longueau \\
Folio 26 recto & p.~\pageref{fol:52} & Mardi 26 septembre 1916 (matin et après-midi) & \textbf{Le Bourget et Paris (Gare du Nord, Tuileries, Orsay, Invalides, rive gauche)} — Arrivée matinale au Bourget, traversée de Paris et démarches aux Invalides \\
Folio 26 verso & p.~\pageref{fol:53} & Mardi 26 septembre 1916 (après-midi et soir) & \textbf{Paris (Tour Saint-Jacques, rue de Rivoli, librairie Colin, gare d'Orsay)} — Mardi 26 septembre 1916 (suite) — Square de la Tour Saint-Jacques, rue de Rivoli, achats d'ouvrages d'Ernest Zyromski et de Maurice Maeterlinck chez Armand Colin et départ à la gare d'Orsay \\
Folio 27 recto & p.~\pageref{fol:54} & Mercredi 27 septembre 1916 & \textbf{Toulouse (Sainte-Marie-des-Champs)} — Mercredi 27 septembre 1916 — Arrivée à Toulouse, montée à Sainte-Marie-des-Champs et déjeuner intime chez les Crouzet \\
Folio 27 verso & p.~\pageref{fol:55} & Jeudi 28 septembre 1916 & \textbf{Toulouse (Sainte-Marie, barrière Saint-Agne, Pouvourville, Pech-David)} — Jeudi 28 septembre 1916 — Rangement de l'appartement toulousain, bouquet pour Marguerite et promenade par Pouvourville et la crête de Pech-David \\
Folio 28 recto & p.~\pageref{fol:56} & Vendredi 29 septembre 1916 & \textbf{Toulouse, Purpan, Blagnac et les Ponts-Jumeaux (Haute-Garonne)} — Vendredi 29 septembre 1916 — Visite à Mme de Labarthe, promenade avec sa fiancée Marguerite à Blagnac (ruines de l'amphithéâtre romain) et achat d'un sautoir en or chez Pfister (295 fr.) \\
Folio 28 verso & p.~\pageref{fol:57} & Samedi 30 septembre 1916 & \textbf{Toulouse (Sainte-Marie, Montaudran et cimetière de Terre-Cabade)} — Les cinq cloches du clocher de Montaudran et le serment au cimetière de Terre-Cabade \\
Folio 29 recto & p.~\pageref{fol:58} & Dimanche 1er octobre 1916 & \textbf{Toulouse (Cathédrale Saint-Étienne, Sainte-Marie et Lardenne)} — Messe matinale à la cathédrale Saint-Étienne et promenade au bord du Touch à Lardenne \\
Folio 29 verso & p.~\pageref{fol:59} & Lundi 2 octobre 1916 & \textbf{Toulouse (Appartement, Sainte-Marie, chapelle de Nazareth, rue d'Alsace-Lorraine)} — Méditation solitaire dans l'appartement toulousain et visite de Marguerite et de sa mère \\
Folio 30 recto & p.~\pageref{fol:60} & Mardi 3 octobre 1916 & \textbf{Toulouse (Cathédrale Saint-Étienne, rue d'Alsace-Lorraine, marché Victor-Hugo, route de Castanet)} — Mardi 3 octobre 1916 — Messe à la cathédrale Saint-Étienne, déjeuner à l'Automate, promenade vers la route de Castanet et présentation de sa fiancée à Mme de Labarthe \\
Folio 30 verso & p.~\pageref{fol:61} & Mardi 3 octobre 1916 (suite) & \textbf{Toulouse} — Règlement du tailleur Guiraud, du loyer semestriel à Mlle Guilhem et thé chez Mme de Labarthe \\
Folio 31 recto & p.~\pageref{fol:62} & Mercredi 4 octobre 1916 & \textbf{Toulouse (Sainte-Marie et gare Matabiau)} — Mercredi 4 octobre 1916 — Dernier après-midi à Sainte-Marie et séparation déchirante à la gare Matabiau \\
Folio 31 verso & p.~\pageref{fol:63} & Jeudi 5 octobre 1916 & \textbf{Paris (Gare d'Orsay, Gare du Nord, église Saint-Vincent-de-Paul, Grands Magasins du Louvre)} — Traversée de Paris sous la pluie : l'église Saint-Vincent-de-Paul et les salons du Louvre \\
Folio 32 recto & p.~\pageref{fol:64} & Mercredi 4 octobre 1916 (soir et nuit) & \textbf{Gare de Toulouse et train de nuit Toulouse – Paris-Orsay} — Fin des adieux en gare de Toulouse, nuit blanche dans l'express d'Orsay et comptes de départ \\
Folio 32 verso & p.~\pageref{fol:65} & Jeudi 5 octobre 1916 (après-midi et soir) & \textbf{Paris (Gare du Nord), Amiens et Doullens (Somme)} — Train de l’après-midi de Paris-Nord à Amiens et arrivée à minuit en gare de Doullens \\
Folio 33 recto & p.~\pageref{fol:66} & Vendredi 6 octobre 1916 (matin et midi) & \textbf{Doullens, Canaples et Barlette (Somme)} — Vendredi 6 octobre 1916 — Fin de la nuit à la gare de Doullens, train de 7 h 25 et retour en lorry de la 54e Brigade auprès de la 152nd Coy A.S.C. \\
Folio 33 verso & p.~\pageref{fol:67} & Vendredi 6 octobre 1916 (après-midi et soir) & \textbf{Hameau de Barlette (Ribeaucourt) et Bernaville (Somme)} — Installation au hameau de Barlette (Ribeaucourt) et visa de permission à Bernaville \\
Folio 34 recto & p.~\pageref{fol:68} & Samedi 7 octobre 1916 & \textbf{Hameau de Barlette, Ribeaucourt (Somme)} — Tempête nocturne sur les vergers de Barlette et règlement des frais de mess au capitaine Panton \\
Folio 34 verso & p.~\pageref{fol:69} & Dimanche 8 octobre 1916 & \textbf{Barlette (Ribeaucourt) et Bernaville (Somme)} — Dimanche de repos à Barlette et perception de la solde à pied à Bernaville sous la rafale \\
Folio 35 recto & p.~\pageref{fol:70} & Lundi 9 octobre 1916 & \textbf{Barlette et forêt de Ribeaucourt (Somme)} — Promenade sous les futaies centenaires de Ribeaucourt et clair de lune d'octobre \\
Folio 35 verso & p.~\pageref{fol:71} & Mardi 10 octobre 1916 & \textbf{Barlette et bois de Ribeaucourt (Somme)} — Corvée de fourrage vert en lisière des bois aux couleurs d'automne et nuit de pleine lune \\
Folio 36 recto & p.~\pageref{fol:72} & Mercredi 11 octobre 1916 & \textbf{Barlette et Ribeaucourt (Somme)} — Convoi de luzerne sous une pluie battante derrière le grand bois de Ribeaucourt \\
Folio 36 verso & p.~\pageref{fol:73} & Jeudi 12 et vendredi 13 octobre 1916 & \textbf{Barlette, Ribeaucourt et Prouville (Somme)} — Conduite d'un « limber G.S. », symphonie du vent dans les bois et chevauchée à Prouville avec le capitaine Panton \\
Folio 37 recto & p.~\pageref{fol:74} & Dimanche 15 octobre 1916 & \textbf{De Barlette (Ribeaucourt) à Beauval (Somme)} — Départ définitif de Barlette sous les averses et arrivée de la colonne à Beauval \\
Folio 37 verso & p.~\pageref{fol:75} & 16 octobre 1916 & \textbf{Beauval, Beauquesne, Puchevillers, Warloy-Baillon (Somme)} — Départ de Beauval sous un air glacial et marche vers Warloy-Baillon \\
Folio 38 recto & p.~\pageref{fol:76} & 16 (comptes) – 17 octobre 1916 (matin) & \textbf{Beauval, Warloy-Baillon (Somme)} — Nuit glaciale sous la tente à Warloy-Baillon et départ de la colonne \\
Folio 38 verso & p.~\pageref{fol:77} & 17 octobre 1916 (suite) & \textbf{Senlis-le-Sec, Bouzincourt, Hédauville (Somme)} — Bivouac dans un champ dévasté entre Bouzincourt et Hédauville \\
Folio 39 recto & p.~\pageref{fol:78} & 18 octobre 1916 & \textbf{Camp de Bouzincourt, Albert (Somme)} — Visite d'Albert bombardée et de sa basilique en ruines \\
Folio 39 verso & p.~\pageref{fol:79} & 18 octobre 1916 (suite) & \textbf{Camp de Bouzincourt, Albert (Somme)} — Le flot ininterrompu des convois vers Albert et la solde à la Mission \\
Folio 40 recto & p.~\pageref{fol:80} & 19 octobre 1916 & \textbf{Camp de Bouzincourt, Albert (Brickfield Area, Somme)} — Pluie torrentielle et installation au Brickfield Area d'Albert \\
Folio 40 verso & p.~\pageref{fol:81} & 20 octobre 1916 & \textbf{Albert (Brickfield Area, Somme)} — Duel d'artillerie formidable et horizon embrasé au-dessus d'Albert \\
Folio 41 recto & p.~\pageref{fol:82} & 21 octobre 1916 & \textbf{Albert, Aveluy (Somme)} — Gelée d'octobre, reconnaissance vers Aveluy et veillée près du poêle \\
Folio 41 verso & p.~\pageref{fol:83} & 22 octobre 1916 & \textbf{Albert (Somme)} — La « Vierge penchée » d'Albert et la rumeur de rappel des jeunes R.A.T. \\
Folio 42 recto & p.~\pageref{fol:84} & 23 octobre 1916 & \textbf{Albert, La Boisselle (Somme)} — Visite des abris souterrains allemands et du grand cratère de La Boisselle \\
Folio 42 verso & p.~\pageref{fol:85} & 24 – 25 octobre 1916 & \textbf{Albert (Somme)} — Maintien confirmé à la Mission française et averses dans les fondrières d'Albert \\
Folio 43 recto & p.~\pageref{fol:86} & 26 – 28 octobre 1916 & \textbf{Albert (Somme)} — Journées de pluie, d'inaction et d'accalmie dans Albert dévastée \\
Folio 43 verso & p.~\pageref{fol:87} & 29 – 30 octobre 1916 & \textbf{Albert, Acheux, Senlis-le-Sec (Somme)} — Enquête sur les pailles à Acheux, rencontre d'un artilleur belge et tempête \\
Folio 44 recto & p.~\pageref{fol:88} & 31 octobre 1916 & \textbf{Albert (Somme)} — Éclaircie de fin octobre et bombardement d'Albert \\
Folio 44 verso & p.~\pageref{fol:89} & 1er novembre 1916 & \textbf{Albert (Somme)} — Toussaint 1916 au cimetière dévasté d'Albert, prise d'armes écossaise et ordre de départ \\
Folio 45 recto & p.~\pageref{fol:90} & 2 novembre 1916 & \textbf{Albert, Acheux-en-Amiénois (Somme)} — Ordre de mission du colonel Grose et départ sous la pluie vers Acheux \\
Folio 45 verso & p.~\pageref{fol:91} & 2 (soir) – 3 novembre 1916 & \textbf{Acheux, Raincheval, Albert (Somme)} — Course nocturne à Raincheval et dernier aller-retour au camp d'Albert \\
Folio 46 recto & p.~\pageref{fol:92} & 4 novembre 1916 & \textbf{Acheux, Bertrancourt, Bus-lès-Artois, Albert (Somme)} — Course à pied à Bertrancourt et retrouvailles avec l'interprète Cordonnier \\
Folio 46 verso & p.~\pageref{fol:93} & 2 (soir) – 3 novembre 1916 [cliché bis] & \textbf{Acheux, Raincheval, Albert (Somme)} — Second cliché archivistique du feuillet des 2–3 novembre 1916 (doublon de numérisation) \\
Folio 47 recto & p.~\pageref{fol:94} & 5 novembre 1916 & \textbf{Acheux-en-Amiénois (Somme)} — Emménagement au n° 1 route de Léalvillers et convocation chez le capitaine Huguet \\
Folio 47 verso & p.~\pageref{fol:95} & 6 – 7 novembre 1916 & \textbf{Acheux, Arquèves (Belle-Église), Hérissart, Rubempré, Coisy (Somme)} — Courses à Belle-Église et Coisy, nostalgie de la 152e Compagnie et convois vers la ligne de feu \\
Folio 48 recto & p.~\pageref{fol:96} & 8 – 9 novembre 1916 & \textbf{Acheux, Forceville, Varennes, Warloy, Senlis-le-Sec, Albert (Somme)} — Promenades en auto à Forceville, Varennes et thé au mess de la 152e Compagnie à Albert \\
Folio 48 verso & p.~\pageref{fol:97} & 10 – 11 novembre 1916 & \textbf{Acheux, Raincheval, Léalvillers, Varennes (Somme)} — Courses à Raincheval et Varennes, rencontre de l'interprète Sorgues et bombardement d'Acheux \\
Folio 49 recto & p.~\pageref{fol:98} & 12 – 14 novembre 1916 & \textbf{Acheux, Varennes, Contay, Beauquesne, Raincheval, Doullens, Frévent (Somme / Pas-de-Calais)} — Bataille de Beaumont-Hamel (13–14 novembre), convois de blessés et foire de Doullens \\
Folio 49 verso & p.~\pageref{fol:99} & 15 – 16 novembre 1916 & \textbf{Acheux, Albert, Raincheval, Louvencourt, Toutencourt, Hérissart (Somme)} — Canonnade sur Beaucourt-sur-l'Ancre et déjeuner à Hérissart chez le lieutenant Walsh \\
Folio 50 recto & p.~\pageref{fol:100} & 17 – 18 novembre 1916 & \textbf{Acheux, Authie, Orville, Amplier, Albert (Somme / Pas-de-Calais)} — Premières fortes gelées, conférence sur la Russie au Y.M.C.A. et première neige \\
Folio 50 verso & p.~\pageref{fol:101} & 20 – 21 novembre 1916 & \textbf{Acheux, Bouzincourt, Contay, Englebelmer, Léalvillers (Somme)} — Défilé de troupes et de prisonniers allemands à Acheux, nuit à Englebelmer et friction avec Curtis \\
Folio 51 recto & p.~\pageref{fol:102} & 22 – 23 novembre 1916 & \textbf{Acheux, Amplier, Authieule, Authie (Somme / Pas-de-Calais)} — Indisposition nocturne, couture de M. Carluvant et tournée à Amplier et Authie \\
Folio 51 verso & p.~\pageref{fol:103} & 25 – 26 novembre 1916 & \textbf{Acheux, Hérissart (Somme)} — Bombardement meurtrier de la sucrerie d'Acheux et visite de M. Deloney \\
Folio 52 recto & p.~\pageref{fol:104} & 28 – 30 novembre 1916 & \textbf{Acheux, Martinsart, Raincheval, Puchevillers, Rubempré, Septenville, Pierregot, Hérissart (Somme)} — Thé au château de Martinsart chez les Bertin et brouillard fantomatique sur la Somme \\
Folio 52 verso & p.~\pageref{fol:105} & 1er – 2 décembre 1916 & \textbf{Acheux, Franvillers, Buire-sur-l'Ancre, Ribemont, Heilly, Toutencourt (Somme)} — Tournée à Franvillers et Buire-sur-l'Ancre, arrivée imprévue du très jeune capitaine Williams \\
Folio 53 recto & p.~\pageref{fol:106} & 3 – 4 décembre 1916 & \textbf{Acheux, Septenville, Hérissart, Corbie, Villers-Bretonneux, Marcelcave, Warfusée-Abancourt, Arquèves, Louvencourt (Somme)} — Longue randonnée jusqu'à Corbie, Villers-Bretonneux et Marcelcave dans la brume glaciale \\
Folio 53 verso & p.~\pageref{fol:107} & 5 – 6 décembre 1916 & \textbf{Acheux, Authie, Bayencourt, Coigneux, Louvencourt, Raincheval (Somme / Pas-de-Calais)} — Course jusqu'aux ruines de Bayencourt et Coigneux, départ de Curtis et confidences de Williams \\
Folio 54 recto & p.~\pageref{fol:108} & 7 – 8 décembre 1916 & \textbf{Acheux, Arquèves (Belle-Église), Hédauville, Englebelmer, Saint-Léger-lès-Authie (Somme)} — Les enfants de la fermière belge de Belle-Église et le village désert d'Englebelmer \\
Folio 54 verso & p.~\pageref{fol:109} & 9 – 11 décembre 1916 & \textbf{Acheux, Belle-Église, Authie, Orville, Saulty, L'Arbret (Somme / Pas-de-Calais)} — Randonnée pluvieuse jusqu'à Saulty et L'Arbret, autoritarisme de Williams et crise de dysenterie \\
Folio 55 recto & p.~\pageref{fol:110} & 12 – 14 décembre 1916 & \textbf{Acheux, Senlis-le-Sec, Toutencourt, Belle-Église, Puchevillers, Harponville (Somme)} — Paysage de neige aux fenêtres d'Acheux et réunion des réfugiés d'Englebelmer à Toutencourt \\
Folio 55 verso & p.~\pageref{fol:111} & 15 – 18 décembre 1916 & \textbf{Acheux, Belle-Église, Arquèves, Raincheval, Toutencourt, Ferme du Quesnoy (Somme)} — Cahots sur les routes défoncées et portrait critique du jeune capitaine Williams \\
Folio 56 recto & p.~\pageref{fol:112} & 22 – 24 décembre 1916 & \textbf{Acheux, Senlis-le-Sec, Rubempré, Le Quesnoy, Léalvillers, Saint-Léger, Orville, Authie, Toutencourt, Puchevillers, Louvencourt (Somme)} — Crise épistolaire avec Marguerite (« velléités de rupture ») et l'incident du « Bloody France ! » \\
Folio 56 verso & p.~\pageref{fol:113} & 19 – 20 décembre 1916 & \textbf{Acheux, Belle-Église, Raincheval, Le Quesnoy, Beauquesne (Somme)} — Attente devant l'Army Purchase Board à Raincheval et grandes gelées d'avant Noël \\
Folio 57 recto & p.~\pageref{fol:114} & 25 décembre 1916 & \textbf{Acheux-en-Amiénois (Somme)} — Noël 1916 à Acheux : insomnie, perplexité morale et concert au Y.M.C.A. \\
Folio 57 verso & p.~\pageref{fol:115} & 26 – 27 décembre 1916 & \textbf{Acheux, Raincheval, Arquèves, Beauquesne, Le Quesnoy, Doullens (Somme)} — Location d'une machine à battre chez Mme Capron à Beauquesne et course rapide à Doullens \\
Folio 58 recto & p.~\pageref{fol:116} & 28 – 29 décembre 1916 & \textbf{Acheux, Léalvillers, Raincheval, Le Quesnoy, Louvencourt (Somme)} — Concert des « 13th Corps Archie Singers » et tournée dans la voiture Sunbeam \\
Folio 58 verso & p.~\pageref{fol:117} & 30 – 31 décembre 1916 & \textbf{Acheux, Franvillers, Senlis-le-Sec, Belle-Église, Raincheval, Beauquesne (Somme)} — Nouvel éclat et excuses contrites du capitaine Williams à la Saint-Sylvestre \\
Folio 59 recto & p.~\pageref{fol:118} & 1er – 2 janvier 1917 & \textbf{Acheux, Toutencourt, Saulty, Louvencourt, Raincheval (Somme / Pas-de-Calais)} — Jour de l'An 1917 à Acheux, tournées à Toutencourt, Saulty et Louvencourt \\
Folio 59 verso & p.~\pageref{fol:119} & 3 janvier 1917 & \textbf{Acheux, Belle-Église, Montigny, Amiens (Somme)} — Course à Amiens : l'Hôtel de la Paix, le capitaine Trefcon et la cathédrale \\
Folio 60 recto & p.~\pageref{fol:120} & 4 – 5 janvier 1917 & \textbf{Acheux, Septenville, Beauquesne (Somme)} — Concert des « XIII Corps Archie » et contrat de battage chez Mme Capron à Beauquesne \\
Folio 60 verso & p.~\pageref{fol:121} & 6 – 9 janvier 1917 & \textbf{Acheux, Louvencourt, Bertrancourt, Authie, Arquèves, Raincheval, Saint-Léger, Hédauville (Somme)} — Mouvements de troupes autour d'Acheux et visite à la 150th Company A.S.C. près d'Hédauville \\
Folio 61 recto & p.~\pageref{fol:122} & 10 – 12 janvier 1917 & \textbf{Acheux, Saint-Léger, Belle-Église, Authie, Arquèves, Louvencourt, Raincheval, Beauquesne, Septenville, Rubempré (Somme)} — Canonnade sur Miraumont, retrouvailles avec Morris et tempête de grésil sur les plateaux \\
Folio 61 verso & p.~\pageref{fol:123} & 13 – 15 janvier 1917 & \textbf{Acheux, Beaucourt, Montigny, Louvencourt, Saint-Léger, Beauquesne, Raincheval, Doullens (Somme)} — Fixation de la permission au 26 janvier, arbres ouatés de neige et course à Doullens \\
Folio 62 recto & p.~\pageref{fol:124} & 16 – 17 janvier 1917 & \textbf{Acheux, Bouzincourt, Senlis-le-Sec, Rubempré, Hérissart, Toutencourt, Doullens, Auxi-le-Château (Somme / Pas-de-Calais)} — Étendues blanches d'Hérissart et visite à Auxi-le-Château chez le collectionneur M. Renard \\
Folio 62 verso & p.~\pageref{fol:125} & 18 – 19 janvier 1917 & \textbf{Acheux, Belle-Église, Raincheval, Puchevillers (Somme)} — Plaines blanches confondues avec le ciel mat et neige cristallisée à Puchevillers \\
Folio 63 recto & p.~\pageref{fol:126} & 20 – 23 janvier 1917 & \textbf{Acheux, Louvencourt (Somme)} — Grand gel sec, retour à pied de Louvencourt et préparatifs de permission \\
Folio 63 verso & p.~\pageref{fol:127} & 23 (soir) – 25 janvier 1917 & \textbf{Acheux-en-Amiénois (Somme)} — Conférence sur le Canada au Y.M.C.A., eau gelée dans la chambre et fin des préparatifs \\
Folio 64 recto & p.~\pageref{fol:128} & 26 janvier 1917 & \textbf{Acheux, Gare de Varennes, Gézaincourt (Somme)} — Départ en permission de la gare de Varennes et longue attente nocturne à Gézaincourt \\
Folio 64 verso & p.~\pageref{fol:129} & 26 janvier 1917 (suite, nuit du 26 au 27) & \textbf{Acheux, Gézaincourt, Amiens (Somme)} — Arrivée à Amiens à deux heures du matin et état de la caisse de voyage \\
Folio 65 recto & p.~\pageref{fol:130} & 27 janvier 1917 (matin et midi) & \textbf{Amiens, Paris (Gare du Nord, Gare d'Orsay), Neuilly-sur-Seine} — Nuit glaciale à la gare d'Amiens avec l'interprète Boilly et arrivée à Paris et Neuilly \\
Folio 65 verso & p.~\pageref{fol:131} & 27 janvier 1917 (suite) & \textbf{Neuilly-sur-Seine, Paris (Boulevard Saint-Michel, Hôtel de Ville, Gare d'Orsay)} — Déjeuner à Neuilly chez Mme Faure, visite chez l'opticien Imbert et départ pour Toulouse \\
Folio 66 recto & p.~\pageref{fol:132} & 27 (comptes) – 28 janvier 1917 (matin) & \textbf{Paris, Toulouse (Gare Matabiau, Saint-Michel, Sainte-Marie)} — Dépenses parisiennes et arrivée matinale à Toulouse : souvenirs des disparus et visite de l'oncle Émile \\
Folio 66 verso & p.~\pageref{fol:133} & 28 janvier 1917 (suite) & \textbf{Toulouse (Quartier Sainte-Marie / Côte Pavée)} — Retrouvailles infiniment tendres avec Marguerite à Sainte-Marie (Toulouse) \\
Folio 67 recto & p.~\pageref{fol:134} & 29 janvier 1917 & \textbf{Toulouse (Gare Matabiau, Sainte-Marie, Basilique Saint-Sernin, Place du Capitole, Rue d'Alsace, Côte Pavée)} — Le panier de fleurs de Nice et le pèlerinage de jeunesse aux cryptes de Saint-Sernin \\
Folio 67 verso & p.~\pageref{fol:135} & 30 janvier 1917 & \textbf{Toulouse (Faculté des Lettres, Sainte-Marie)} — Visite à Mme de Labarthe et conférence d'Ernest Zyromski sur Maurice de Guérin \\
Folio 68 recto & p.~\pageref{fol:136} & 31 janvier 1917 & \textbf{Toulouse (Saint-Sernin, Halles Victor-Hugo, Hôtel d'Assézat, Daurade, Dalbade, Sainte-Marie)} — Marché aux Halles Victor-Hugo avec Marguerite, visite de l'Hôtel d'Assézat et intimité à l'appartement \\
Folio 68 verso & p.~\pageref{fol:137} & 1er février 1917 & \textbf{Toulouse (Capitole, Sainte-Marie, Gare Matabiau, Rue de Languedoc)} — Achat d'une montre en or chez Pfister et détachement intérieur vis-à-vis de Toulouse \\
Folio 69 recto & p.~\pageref{fol:138} & 2 février 1917 & \textbf{Toulouse (Rue Lafayette, Place Belfort, Sainte-Marie)} — Soins chez le dentiste M. de Lersay, achat du remontoir Louis XVI et dernière soirée à Toulouse \\
Folio 69 verso & p.~\pageref{fol:139} & 3 février 1917 & \textbf{Toulouse, Carcassonne, Béziers, Bédarieux, Millau (Aveyron)} — Voyage de Toulouse à Millau avec Marguerite par Carcassonne et Béziers \\
Folio 70 recto & p.~\pageref{fol:140} & 3 février 1917 (suite, comptes) & \textbf{Toulouse, Béziers, Bédarieux, Millau} — Relevé des dépenses du départ de Toulouse et du voyage par Bédarieux \\
Folio 70 verso & p.~\pageref{fol:141} & 4 février 1917 (nuit et matin) & \textbf{Millau (10 rue Alfred-Guibert, Aveyron)} — Nuit d'émotion et de larmes à Millau au souvenir de ses parents et de ses sœurs \\
Folio 71 recto & p.~\pageref{fol:142} & 4 février 1917 (suite) & \textbf{Millau (Aveyron)} — Après-midi musical avec Marguerite Calmels, visite à l'oncle et la tante Thérèse (Layral) et arrivée de Mimi à Millau \\
Folio 71 verso & p.~\pageref{fol:143} & 5 février 1917 & \textbf{Millau, Séverac, Marvejols, Saint-Flour (Aveyron / Lozère / Cantal)} — Recueillement au cimetière de Millau avec Marguerite, adieux déchirants et arrivée à Saint-Flour \\
Folio 72 recto & p.~\pageref{fol:144} & 6 février 1917 (matin et après-midi) & \textbf{Saint-Flour, Brioude, Clermont-Ferrand, Gannat, Gien, Montargis} — Nuit glaciale en gare de Saint-Flour et traversée de l'Auvergne sous les flocons \\
Folio 72 verso & p.~\pageref{fol:145} & 6 février 1917 (suite, soir) & \textbf{Corbeil-Essonnes, Paris (Gare de Lyon, Bastille, Gare du Nord)} — Traversée de Paris en métro de la gare de Lyon à la gare du Nord et train de nuit vers Amiens \\
Folio 73 recto & p.~\pageref{fol:146} & 7 février 1917 (nuit et matin) & \textbf{Amiens, Doullens, Louvencourt, Acheux-en-Amiénois (Somme)} — Nuit glaciale à Amiens, combat aérien contre deux Taubes à Louvencourt et retour chez les Carluvant \\
Folio 73 verso & p.~\pageref{fol:147} & 7 février 1917 (après-midi) & \textbf{Acheux, Raincheval, Puchevillers, Hérissart (Somme)} — Déménagement imprévu d'Acheux à Hérissart sous un gel intense \\
Folio 74 recto & p.~\pageref{fol:148} & 7 (fin) – 8 février 1917 & \textbf{Hérissart (Somme)} — Nuit glaciale à Hérissart, comptes de retour et premier jeudi de repos \\
Folio 74 verso & p.~\pageref{fol:149} & 8 (suite) – 9 février 1917 (matin) & \textbf{Hérissart, Rubempré (Somme)} — Promenade à pied sur la neige durcie jusqu'à Rubempré \\
Folio 75 recto & p.~\pageref{fol:150} & 9 (suite) – 10 février 1917 (matin) & \textbf{Hérissart, Toutencourt, Martinsart (Somme)} — Visite à la Mission de Toutencourt et course au bois d'Aveluy et Martinsart \\
Folio 75 verso & p.~\pageref{fol:151} & 10 (suite) – 11 février 1917 (matin) & \textbf{Bois d'Aveluy, Senlis-le-Sec, Acheux, Hérissart (Somme)} — Visite à son ancienne compagnie au bois d'Aveluy et combat aérien \\
Folio 76 recto & p.~\pageref{fol:152} & 11 février 1917 (suite) & \textbf{Toutencourt, Doullens, Auxi-le-Château (Somme / Pas-de-Calais)} — Excursion dominicale à Doullens et Auxi-le-Château chez M. Romert \\
Folio 76 verso & p.~\pageref{fol:153} & 12 février 1917 & \textbf{Hérissart, Méaulte, Ville-sur-Ancre, Morlancourt, Ribemont (Somme)} — Visite de la minoterie de Méaulte et silhouettes des moulins d'Hérissart \\
Folio 77 recto & p.~\pageref{fol:154} & 13 – 14 février 1917 & \textbf{Hérissart, Authie, Bus-lès-Artois, Gézaincourt, Acheux, Toutencourt (Somme)} — Courses à Authie, Bus et Acheux : arabesques de givre et repas à l'Hôtel de Ville \\
Folio 77 verso & p.~\pageref{fol:155} & 15 – 16 février 1917 & \textbf{Hérissart, Toutencourt, Raincheval, Puchevillers, Beaucourt-en-Santerre / sur-l'Hallue (Somme)} — Visite à la mairie d'Hérissart et tournée au château de Beaucourt \\
Folio 78 recto & p.~\pageref{fol:156} & 17 – 18 février 1917 & \textbf{Hérissart, Toutencourt (Somme)} — Grand dégel de février, fondrières de boue et prône du dimanche à Toutencourt \\
Folio 78 verso & p.~\pageref{fol:157} & 19 – 20 février 1917 & \textbf{Hérissart, Fréchencourt, Bavelincourt, Pierregot (Somme)} — Déménagement dans l'ancienne chambre humide du capitaine Williams \\
Folio 79 recto & p.~\pageref{fol:158} & 21 – 23 février 1917 & \textbf{Hérissart, Doullens, Auxi-le-Château, Puchevillers, Molliens-au-Bois (Somme)} — Course à Doullens et Auxi-le-Château sous le brouillard : mandat pour Mme Spagnol \\
Folio 79 verso & p.~\pageref{fol:159} & 24 – 26 février 1917 & \textbf{Hérissart, Contay, Mirvaux, Toutencourt (Somme)} — Promenades à Contay et Mirvaux : premières sensations de printemps \\
Folio 80 recto & p.~\pageref{fol:160} & 27 – 28 février 1917 & \textbf{Hérissart, Toutencourt (Somme)} — Rage de dents du capitaine Williams et passage de la cavalerie indienne à Hérissart \\
Folio 80 verso & p.~\pageref{fol:161} & 1er – 2 mars 1917 & \textbf{Hérissart, Rubempré, Montigny, Fréchencourt, Contay, Toutencourt (Somme)} — Entrée en mars : château de Rubempré, halte au clair de lune et entretien avec le Dr Dubois \\
Folio 81 recto & p.~\pageref{fol:162} & 3 – 5 mars 1917 & \textbf{Hérissart, Rubempré, Puchevillers, Toutencourt, Montigny, Fréchencourt (Somme)} — Retenue de solde pour permission à Toutencourt et chute de neige sur les arbres \\
Folio 81 verso & p.~\pageref{fol:163} & 6 – 7 mars 1917 & \textbf{Hérissart, Doullens, Toutencourt (Somme)} — Courses à Doullens avec Williams et promesse du maréchal des logis Grauger à Toutencourt \\
Folio 82 recto & p.~\pageref{fol:164} & 8 – 9 mars 1917 & \textbf{Hérissart, Toutencourt (Somme)} — Tempête de neige de mars et départ de l'interprète Valet pour Albert \\
Folio 82 verso & p.~\pageref{fol:165} & 10 mars 1917 & \textbf{Hérissart (Somme)} — Départ définitif du capitaine Williams et mélancolie de l'attente à Hérissart \\
Folio 83 recto & p.~\pageref{fol:166} & 11 mars 1917 & \textbf{Hérissart, Toutencourt, Puchevillers (Somme)} — Maintien au Purchase Board sur demande du Major Levy et installation au billet 220 \\
Folio 83 verso & p.~\pageref{fol:167} & 11 (fin) – 12 mars 1917 & \textbf{Hérissart, Toutencourt, Contay, Beaumont, Molliens-au-Bois, Puchevillers, Raincheval, Candas (Somme)} — Bonheur de la nouvelle chambre (billet 220) et tournée en auto Wolseley \\
Folio 84 recto & p.~\pageref{fol:168} & 13 – 15 mars 1917 & \textbf{Hérissart, Le Mesnil-Martinsart, Puchevillers, Acheux, Bertrancourt (Somme)} — Visite à l'instituteur M. Flutre et déménagement du mess au billet 216 \\
Folio 84 verso & p.~\pageref{fol:169} & 16 – 17 mars 1917 & \textbf{Hérissart, Mirvaux, Pierregot, Contay, Warloy, Senlis, Albert, Ribemont, Toutencourt, Puchevillers (Somme)} — Retour à Albert et Ribemont et promenades printanières avec le lieutenant Wentworth \\
Folio 85 recto & p.~\pageref{fol:170} & 18 – 19 mars 1917 & \textbf{Hérissart, Toutencourt, Puchevillers (Somme)} — Prône du vieux curé de Toutencourt et passage de la 63e Royal Naval Division \\
Folio 85 verso & p.~\pageref{fol:171} & 20 – 22 mars 1917 & \textbf{Hérissart, Puchevillers, Arquèves (Belle-Église), Contay, Warloy, Senlis, Pierregot, Molliens, Albert (Somme)} — Départ de la Naval Division, ferme de Belle-Église et défilé des 17th Lancers (Death or Glory Boys) \\
Folio 86 recto & p.~\pageref{fol:172} & 23 – 25 mars 1917 & \textbf{Hérissart, Rubempré, Puchevillers (Somme)} — Chants d'oiseaux dans le gel de mars et démarches pour conserver le billet 220 \\
Folio 86 verso & p.~\pageref{fol:173} & 26 – 28 mars 1917 & \textbf{Hérissart, Rubempré, Puchevillers (Somme)} — Visite du camarade Monin et départ de toutes les troupes d'Hérissart \\
Folio 87 recto & p.~\pageref{fol:174} & 29 – 31 mars 1917 & \textbf{Hérissart, Puchevillers, Rubempré, Pierregot, Molliens-au-Bois, Toutencourt, Rainneville (Somme)} — Le calme délicieux des villages désertés par les troupes à la fin mars \\
Folio 87 verso & p.~\pageref{fol:175} & 1er – 3 avril 1917 & \textbf{Hérissart, Belle-Église, Pierregot, Rainneville, Toutencourt, Rubempré, Raincheval (Somme)} — Dimanche des Rameaux sous la neige et comptes du mess au 31 mars \\
Folio 88 recto & p.~\pageref{fol:176} & 4 – 5 avril 1917 & \textbf{Hérissart, Toutencourt, Acheux, Contay, Warloy, Varennes (Somme)} — Jeudi Saint ensoleillé : les villages déserts de l'ancien front et le cimetière d'Hérissart \\
Folio 88 verso & p.~\pageref{fol:177} & 6 – 7 avril 1917 & \textbf{Hérissart, Toutencourt, Rubempré, Contay, Acheux, Léalvillers (Somme)} — Vendredi Saint : solde au Q.G. de la Ve Armée à Toutencourt et causerie sur l'Australie \\
Folio 89 recto & p.~\pageref{fol:178} & 8 – 10 avril 1917 & \textbf{Hérissart, Toutencourt, Harponville, Rubempré, Pierregot, Belle-Église, Fréchencourt, Contay (Somme)} — Pâques 1917 à Hérissart et déclenchement de l'offensive d'Arras sous la neige \\
Folio 89 verso & p.~\pageref{fol:179} & 11 – 13 avril 1917 & \textbf{Hérissart, Doullens, Warloy, Belle-Église (Somme)} — Tempête de neige du 11 avril, emplettes à Doullens et entretien avec le lieutenant Gilbert \\
Folio 90 recto & p.~\pageref{fol:180} & 14 – 17 avril 1917 & \textbf{Hérissart, Pierregot, Fréchencourt, Rubempré (Somme)} — Soirée chez la famille Denis sous l'orage et retour des giboulées \\
Folio 90 verso & p.~\pageref{fol:181} & 18 – 20 avril 1917 & \textbf{Hérissart, Bresle, Lavieville, Hénencourt, Warloy, Contay, Rubempré, Senlis, Albert, Buire-sur-l'Ancre, Ribemont (Somme)} — Retour du lieutenant Gilbert d'Abbeville et voiture désembourbée par les Australiens à Buire \\
Folio 91 recto & p.~\pageref{fol:182} & 21 – 22 avril 1917 & \textbf{Hérissart, Toutencourt, Buire, Albert, Millencourt, Contay, Molliens-au-Bois, Pierregot (Somme)} — Violente canonnade vers Arras, moulin de Pierregot et batterie antiaérienne \\
Folio 91 verso & p.~\pageref{fol:183} & 23 – 25 avril 1917 & \textbf{Hérissart, Toutencourt, Pierregot, Acheux, Puchevillers, Mirvaux (Somme)} — Première vraie journée d'été, bourgeons vernissés et soirée chez la famille Denis \\
Folio 92 recto & p.~\pageref{fol:184} & 26 – 27 avril 1917 & \textbf{Hérissart, Contay, Puchevillers, Pierregot, Harponville, Varennes, Rubempré (Somme)} — Promenade le long des anciennes tranchées de Puchevillers et visite à la famille Mortreux à Varennes \\
Folio 92 verso & p.~\pageref{fol:185} & 28 – 29 avril 1917 & \textbf{Hérissart, Pierregot, Contay, Mirvaux (Somme)} — Annonce du départ prochain pour Cambrai et flânerie avec le lieutenant Gilbert \\
Folio 93 recto & p.~\pageref{fol:186} & 30 avril 1917 & \textbf{Hérissart, Vadencourt, Warloy, Hénencourt, Millencourt, Lavieville, Puchevillers, Pierregot, Villers-Bocage, Toutencourt (Somme)} — Dernier folio du Livre-Journal : chaleur d'été, course à Villers-Bocage et nuit étoilée du 30 avril 1917 \\
\end{longtable}
}
\par}

% =============================================================================
% INDEX DES NOMS DE PERSONNES (INDEX NOMINUM) ET INDEX DES NOMS DE LIEUX (INDEX LOCORUM)
% =============================================================================
\cleardoublepage
\chapter*{Index des noms de personnes (\textit{Index nominum})}
\addcontentsline{toc}{chapter}{Index des noms de personnes (\textit{Index nominum})}
\markboth{Index des noms de personnes}{Index des noms de personnes}
{\fontsize{7.0pt}{8.3pt}\selectfont
\noindent Chaque entrée renvoie simultanément à la foliotation du \textit{Livre-Journal} (\texttt{f\textsuperscript{o}~X~r\textsuperscript{o}/v\textsuperscript{o}}) et à la pagination exacte de la présente édition (\texttt{p.~Y}).\par\vspace{2.0mm}
\noindent\hangindent=4.5mm\textbf{Alric, Marie Anaïs (1838–1907)}, mère d'Hippolyte Combes, épouse de François-Hippolyte Combes : f\textsuperscript{o}~29~v\textsuperscript{o} (p.~\pageref{fol:59}), f\textsuperscript{o}~34~r\textsuperscript{o} (p.~\pageref{fol:68}), f\textsuperscript{o}~46~r\textsuperscript{o} (p.~\pageref{fol:92}), f\textsuperscript{o}~66~r\textsuperscript{o} (p.~\pageref{fol:132}), f\textsuperscript{o}~70~v\textsuperscript{o} (p.~\pageref{fol:141}), f\textsuperscript{o}~71~v\textsuperscript{o} (p.~\pageref{fol:143}).\par\vspace{0.55mm}
\noindent\hangindent=4.5mm\textbf{Auriol}, joaillier à Toulouse : f\textsuperscript{o}~69~v\textsuperscript{o} (p.~\pageref{fol:139}).\par\vspace{0.55mm}
\noindent\hangindent=4.5mm\textbf{Backer, M\textsuperscript{me}}, logeuse d'Hippolyte Combes à Lynde (Nord) : f\textsuperscript{o}~1~r\textsuperscript{o} (p.~\pageref{fol:2}).\par\vspace{0.55mm}
\noindent\hangindent=4.5mm\textbf{Barbry, Émile}, industriel blanchisseur sur la Lys près d'Armentières : f\textsuperscript{o}~8~v\textsuperscript{o} (p.~\pageref{fol:17}), f\textsuperscript{o}~11~r\textsuperscript{o} (p.~\pageref{fol:22}), f\textsuperscript{o}~11~v\textsuperscript{o} (p.~\pageref{fol:23}).\par\vspace{0.55mm}
\noindent\hangindent=4.5mm\textbf{Benoît-Joseph Labre, saint (1748–1783)}, pèlerin vénéré en l'église de Saint-Pol-sur-Ternoise : f\textsuperscript{o}~16~r\textsuperscript{o} (p.~\pageref{fol:32}).\par\vspace{0.55mm}
\noindent\hangindent=4.5mm\textbf{Bertin (famille)}, hôtes au château de Martinsart (Somme) : f\textsuperscript{o}~52~r\textsuperscript{o} (p.~\pageref{fol:104}).\par\vspace{0.55mm}
\noindent\hangindent=4.5mm\textbf{Bienaimé, M\textsuperscript{e}}, notaire à Toutencourt (Somme) : f\textsuperscript{o}~89~r\textsuperscript{o} (p.~\pageref{fol:178}).\par\vspace{0.55mm}
\noindent\hangindent=4.5mm\textbf{Bradfield, Lieutenant}, officier britannique de la 152\textsuperscript{e} Compagnie A.S.C. : f\textsuperscript{o}~75~v\textsuperscript{o} (p.~\pageref{fol:151}).\par\vspace{0.55mm}
\noindent\hangindent=4.5mm\textbf{Brammas}, collègue interprète militaire : f\textsuperscript{o}~22~v\textsuperscript{o} (p.~\pageref{fol:45}).\par\vspace{0.55mm}
\noindent\hangindent=4.5mm\textbf{Bresson, Capitaine}, officier de la Mission militaire française auprès de la V\textsuperscript{e} Armée à Toutencourt : f\textsuperscript{o}~52~v\textsuperscript{o} (p.~\pageref{fol:105}).\par\vspace{0.55mm}
\noindent\hangindent=4.5mm\textbf{Calmels, Anne-Marie dite « Mimi » (1888–1975)}, sœur de Marguerite Calmels, future épouse de M\textsuperscript{e} Joseph Delcamps : f\textsuperscript{o}~69~v\textsuperscript{o} (p.~\pageref{fol:139}), f\textsuperscript{o}~70~v\textsuperscript{o} (p.~\pageref{fol:141}), f\textsuperscript{o}~71~r\textsuperscript{o} (p.~\pageref{fol:142}), f\textsuperscript{o}~71~v\textsuperscript{o} (p.~\pageref{fol:143}).\par\vspace{0.55mm}
\noindent\hangindent=4.5mm\textbf{Calmels, D\textsuperscript{r} Gabriel (1852–1921) \& Valentine Cartairade (1858–1943)}, parents de Marguerite Calmels au 10, rue Alfred-Guibert à Millau : f\textsuperscript{o}~27~r\textsuperscript{o} (p.~\pageref{fol:54}), f\textsuperscript{o}~30~r\textsuperscript{o} (p.~\pageref{fol:60}), f\textsuperscript{o}~69~v\textsuperscript{o} (p.~\pageref{fol:139}), f\textsuperscript{o}~70~v\textsuperscript{o} (p.~\pageref{fol:141}), f\textsuperscript{o}~71~r\textsuperscript{o} (p.~\pageref{fol:142}), f\textsuperscript{o}~71~v\textsuperscript{o} (p.~\pageref{fol:143}).\par\vspace{0.55mm}
\noindent\hangindent=4.5mm\textbf{Calmels, Madeleine (1882–1961, \textit{Mad\textsuperscript{ne}})}, sœur aînée de Marguerite Calmels, épouse du D\textsuperscript{r} Fernand Vincens : f\textsuperscript{o}~66~v\textsuperscript{o} (p.~\pageref{fol:133}), f\textsuperscript{o}~67~r\textsuperscript{o} (p.~\pageref{fol:134}), f\textsuperscript{o}~67~v\textsuperscript{o} (p.~\pageref{fol:135}), f\textsuperscript{o}~70~v\textsuperscript{o} (p.~\pageref{fol:141}), f\textsuperscript{o}~71~v\textsuperscript{o} (p.~\pageref{fol:143}).\par\vspace{0.55mm}
\noindent\hangindent=4.5mm\textbf{Calmels, Marie Marguerite (1892–1974, \textit{Marg\textsuperscript{te}}, \textit{M\textsuperscript{te}}, \textit{Gueïte})}, fiancée puis épouse d'Hippolyte Combes : f\textsuperscript{o}~1~r\textsuperscript{o} (p.~\pageref{fol:2}), f\textsuperscript{o}~3~r\textsuperscript{o} (p.~\pageref{fol:6}), f\textsuperscript{o}~9~v\textsuperscript{o} (p.~\pageref{fol:19}), f\textsuperscript{o}~27~r\textsuperscript{o} (p.~\pageref{fol:54}), f\textsuperscript{o}~27~v\textsuperscript{o} (p.~\pageref{fol:55}), f\textsuperscript{o}~28~r\textsuperscript{o} (p.~\pageref{fol:56}), f\textsuperscript{o}~28~v\textsuperscript{o} (p.~\pageref{fol:57}), f\textsuperscript{o}~29~r\textsuperscript{o} (p.~\pageref{fol:58}), f\textsuperscript{o}~29~v\textsuperscript{o} (p.~\pageref{fol:59}), f\textsuperscript{o}~30~r\textsuperscript{o} (p.~\pageref{fol:60}), f\textsuperscript{o}~31~r\textsuperscript{o} (p.~\pageref{fol:62}), f\textsuperscript{o}~32~r\textsuperscript{o} (p.~\pageref{fol:64}), f\textsuperscript{o}~56~r\textsuperscript{o} (p.~\pageref{fol:112}), f\textsuperscript{o}~57~r\textsuperscript{o} (p.~\pageref{fol:114}), f\textsuperscript{o}~66~v\textsuperscript{o} (p.~\pageref{fol:133}), f\textsuperscript{o}~67~r\textsuperscript{o} (p.~\pageref{fol:134}), f\textsuperscript{o}~67~v\textsuperscript{o} (p.~\pageref{fol:135}), f\textsuperscript{o}~68~r\textsuperscript{o} (p.~\pageref{fol:136}), f\textsuperscript{o}~68~v\textsuperscript{o} (p.~\pageref{fol:137}), f\textsuperscript{o}~69~r\textsuperscript{o} (p.~\pageref{fol:138}), f\textsuperscript{o}~69~v\textsuperscript{o} (p.~\pageref{fol:139}), f\textsuperscript{o}~70~r\textsuperscript{o} (p.~\pageref{fol:140}), f\textsuperscript{o}~70~v\textsuperscript{o} (p.~\pageref{fol:141}), f\textsuperscript{o}~71~r\textsuperscript{o} (p.~\pageref{fol:142}), f\textsuperscript{o}~71~v\textsuperscript{o} (p.~\pageref{fol:143}).\par\vspace{0.55mm}
\noindent\hangindent=4.5mm\textbf{Calmels, Pierre}, frère de Marguerite Calmels à Millau : f\textsuperscript{o}~69~v\textsuperscript{o} (p.~\pageref{fol:139}), f\textsuperscript{o}~70~v\textsuperscript{o} (p.~\pageref{fol:141}), f\textsuperscript{o}~71~v\textsuperscript{o} (p.~\pageref{fol:143}).\par\vspace{0.55mm}
\noindent\hangindent=4.5mm\textbf{Capron, M\textsuperscript{me}}, propriétaire d'une batteuse agricole à Beauquesne (Somme) : f\textsuperscript{o}~57~v\textsuperscript{o} (p.~\pageref{fol:115}), f\textsuperscript{o}~60~r\textsuperscript{o} (p.~\pageref{fol:120}).\par\vspace{0.55mm}
\noindent\hangindent=4.5mm\textbf{Carluvant, M.}, maître tailleur et logeur d'Hippolyte Combes au n\textsuperscript{o}~1, route de Léalvillers à Acheux : f\textsuperscript{o}~47~r\textsuperscript{o} (p.~\pageref{fol:94}), f\textsuperscript{o}~51~r\textsuperscript{o} (p.~\pageref{fol:102}), f\textsuperscript{o}~59~r\textsuperscript{o} (p.~\pageref{fol:118}), f\textsuperscript{o}~73~r\textsuperscript{o} (p.~\pageref{fol:146}), f\textsuperscript{o}~73~v\textsuperscript{o} (p.~\pageref{fol:147}), f\textsuperscript{o}~74~r\textsuperscript{o} (p.~\pageref{fol:148}).\par\vspace{0.55mm}
\noindent\hangindent=4.5mm\textbf{Combes, François Augustin Hippolyte (1825–1899)}, avocat à Millau puis à Toulouse, père d'Hippolyte Combes : f\textsuperscript{o}~5~v\textsuperscript{o} (p.~\pageref{fol:11}), f\textsuperscript{o}~30~r\textsuperscript{o} (p.~\pageref{fol:60}), f\textsuperscript{o}~71~v\textsuperscript{o} (p.~\pageref{fol:143}).\par\vspace{0.55mm}
\noindent\hangindent=4.5mm\textbf{Combes, Françoise (1871–1942, \textit{Sœur Marie de la Croix})}, sœur aînée d'Hippolyte Combes, carmélite à Périgueux : Prolégomènes \& Annexe~I (p.~\pageref{annexe:doc1}--\pageref{annexe:doc5}).\par\vspace{0.55mm}
\noindent\hangindent=4.5mm\textbf{Combes, Joséphine (1873–1914, \textit{Joseph\textsuperscript{e}})}, sœur d'Hippolyte Combes, décédée à Nogent-le-Rotrou en octobre 1914 : f\textsuperscript{o}~28~v\textsuperscript{o} (p.~\pageref{fol:57}), f\textsuperscript{o}~30~r\textsuperscript{o} (p.~\pageref{fol:60}).\par\vspace{0.55mm}
\noindent\hangindent=4.5mm\textbf{Combes, Maria (1875–1903)}, sœur d'Hippolyte Combes, religieuse décédée à Carthage : f\textsuperscript{o}~32~r\textsuperscript{o} (p.~\pageref{fol:64}).\par\vspace{0.55mm}
\noindent\hangindent=4.5mm\textbf{Combes, Marie Lucien Hippolyte (1876–1970)}, auteur du \textit{Livre-Journal}, professeur d'anglais et interprète auprès de l'Armée britannique (mention nominale autographe \textit{Hipp\textsuperscript{te}}) : Couverture (p.~\pageref{fol:1}), f\textsuperscript{o}~13~r\textsuperscript{o} (p.~\pageref{fol:26}), f\textsuperscript{o}~67~v\textsuperscript{o} (p.~\pageref{fol:135}).\par\vspace{0.55mm}
\noindent\hangindent=4.5mm\textbf{Connelly}, soldat britannique, ordonnance (\textit{groom}) d'Hippolyte Combes à la 152\textsuperscript{e} Compagnie A.S.C. : f\textsuperscript{o}~6~r\textsuperscript{o} (p.~\pageref{fol:12}), f\textsuperscript{o}~19~v\textsuperscript{o} (p.~\pageref{fol:39}), f\textsuperscript{o}~20~v\textsuperscript{o} (p.~\pageref{fol:41}), f\textsuperscript{o}~22~v\textsuperscript{o} (p.~\pageref{fol:45}), f\textsuperscript{o}~37~r\textsuperscript{o} (p.~\pageref{fol:74}).\par\vspace{0.55mm}
\noindent\hangindent=4.5mm\textbf{Croughton, Lieutenant}, officier britannique du service automobile et des achats : f\textsuperscript{o}~54~v\textsuperscript{o} (p.~\pageref{fol:109}), f\textsuperscript{o}~75~r\textsuperscript{o} (p.~\pageref{fol:150}).\par\vspace{0.55mm}
\noindent\hangindent=4.5mm\textbf{Crouzet, Abbé Jean}, cousin germain de Marguerite Calmels, prêtre à Toulouse (futur curé de la Dalbade) : f\textsuperscript{o}~27~r\textsuperscript{o} (p.~\pageref{fol:54}), f\textsuperscript{o}~29~r\textsuperscript{o} (p.~\pageref{fol:58}), f\textsuperscript{o}~29~v\textsuperscript{o} (p.~\pageref{fol:59}), f\textsuperscript{o}~30~r\textsuperscript{o} (p.~\pageref{fol:60}), f\textsuperscript{o}~66~v\textsuperscript{o} (p.~\pageref{fol:133}), f\textsuperscript{o}~69~v\textsuperscript{o} (p.~\pageref{fol:139}).\par\vspace{0.55mm}
\noindent\hangindent=4.5mm\textbf{Crouzet, Émile (1842–1921, \textit{M.~C\textsuperscript{zt}}, \textit{Oncle}) \& Élisabeth Calmels (\textit{Tante})}, oncle et tante de Marguerite Calmels à Toulouse (rue du Taur et rue Pharaon) : f\textsuperscript{o}~27~r\textsuperscript{o} (p.~\pageref{fol:54}), f\textsuperscript{o}~27~v\textsuperscript{o} (p.~\pageref{fol:55}), f\textsuperscript{o}~28~r\textsuperscript{o} (p.~\pageref{fol:56}), f\textsuperscript{o}~29~r\textsuperscript{o} (p.~\pageref{fol:58}), f\textsuperscript{o}~30~r\textsuperscript{o} (p.~\pageref{fol:60}), f\textsuperscript{o}~31~r\textsuperscript{o} (p.~\pageref{fol:62}), f\textsuperscript{o}~32~r\textsuperscript{o} (p.~\pageref{fol:64}), f\textsuperscript{o}~66~r\textsuperscript{o} (p.~\pageref{fol:132}), f\textsuperscript{o}~66~v\textsuperscript{o} (p.~\pageref{fol:133}), f\textsuperscript{o}~68~v\textsuperscript{o} (p.~\pageref{fol:137}), f\textsuperscript{o}~69~v\textsuperscript{o} (p.~\pageref{fol:139}).\par\vspace{0.55mm}
\noindent\hangindent=4.5mm\textbf{Curtis, Lieutenant (\textit{C\textsuperscript{tis}})}, officier britannique d'approvisionnement et de fourrage (18\textsuperscript{e} Division puis Acheux) : f\textsuperscript{o}~1~r\textsuperscript{o} (p.~\pageref{fol:2}), f\textsuperscript{o}~3~r\textsuperscript{o} (p.~\pageref{fol:6}), f\textsuperscript{o}~6~v\textsuperscript{o} (p.~\pageref{fol:13}), f\textsuperscript{o}~7~r\textsuperscript{o} (p.~\pageref{fol:14}), f\textsuperscript{o}~9~v\textsuperscript{o} (p.~\pageref{fol:19}), f\textsuperscript{o}~10~v\textsuperscript{o} (p.~\pageref{fol:21}), f\textsuperscript{o}~11~r\textsuperscript{o} (p.~\pageref{fol:22}), f\textsuperscript{o}~11~v\textsuperscript{o} (p.~\pageref{fol:23}), f\textsuperscript{o}~15~r\textsuperscript{o} (p.~\pageref{fol:30}), f\textsuperscript{o}~18~v\textsuperscript{o} (p.~\pageref{fol:37}), f\textsuperscript{o}~22~v\textsuperscript{o} (p.~\pageref{fol:45}), f\textsuperscript{o}~37~r\textsuperscript{o} (p.~\pageref{fol:74}), f\textsuperscript{o}~43~v\textsuperscript{o} (p.~\pageref{fol:87}), f\textsuperscript{o}~44~v\textsuperscript{o} (p.~\pageref{fol:89}), f\textsuperscript{o}~45~r\textsuperscript{o} (p.~\pageref{fol:90}), f\textsuperscript{o}~45~v\textsuperscript{o} (p.~\pageref{fol:91}), f\textsuperscript{o}~46~v\textsuperscript{o} (p.~\pageref{fol:93}), f\textsuperscript{o}~48~r\textsuperscript{o} (p.~\pageref{fol:96}), f\textsuperscript{o}~48~v\textsuperscript{o} (p.~\pageref{fol:97}), f\textsuperscript{o}~50~v\textsuperscript{o} (p.~\pageref{fol:101}), f\textsuperscript{o}~52~r\textsuperscript{o} (p.~\pageref{fol:104}), f\textsuperscript{o}~52~v\textsuperscript{o} (p.~\pageref{fol:105}), f\textsuperscript{o}~53~v\textsuperscript{o} (p.~\pageref{fol:107}).\par\vspace{0.55mm}
\noindent\hangindent=4.5mm\textbf{Denis, Baptiste \& famille}, famille d'Hérissart hôte d'Hippolyte Combes au printemps 1917 : f\textsuperscript{o}~89~r\textsuperscript{o} (p.~\pageref{fol:178}), f\textsuperscript{o}~90~r\textsuperscript{o} (p.~\pageref{fol:180}), f\textsuperscript{o}~91~r\textsuperscript{o} (p.~\pageref{fol:182}), f\textsuperscript{o}~91~v\textsuperscript{o} (p.~\pageref{fol:183}), f\textsuperscript{o}~92~r\textsuperscript{o} (p.~\pageref{fol:184}).\par\vspace{0.55mm}
\noindent\hangindent=4.5mm\textbf{Donze, M\textsuperscript{me}}, logeuse d'Hippolyte Combes à Lynde (Nord) : f\textsuperscript{o}~2~v\textsuperscript{o} (p.~\pageref{fol:5}).\par\vspace{0.55mm}
\noindent\hangindent=4.5mm\textbf{Dubois, D\textsuperscript{r}}, médecin rencontré par Hippolyte Combes dans la Somme : f\textsuperscript{o}~59~v\textsuperscript{o} (p.~\pageref{fol:119}), f\textsuperscript{o}~80~v\textsuperscript{o} (p.~\pageref{fol:161}).\par\vspace{0.55mm}
\noindent\hangindent=4.5mm\textbf{Duchâteau, Victor \& famille}, famille logeant Hippolyte Combes près du camp des briqueteries à Albert : f\textsuperscript{o}~40~r\textsuperscript{o} (p.~\pageref{fol:80}), f\textsuperscript{o}~41~r\textsuperscript{o} (p.~\pageref{fol:82}), f\textsuperscript{o}~43~v\textsuperscript{o} (p.~\pageref{fol:87}), f\textsuperscript{o}~45~v\textsuperscript{o} (p.~\pageref{fol:91}), f\textsuperscript{o}~46~v\textsuperscript{o} (p.~\pageref{fol:93}).\par\vspace{0.55mm}
\noindent\hangindent=4.5mm\textbf{Faure, M\textsuperscript{me} \& ses enfants (Marguerite, Germaine, Marcel, Henri, André)}, cousines de Marguerite Calmels à Neuilly-sur-Seine (16, rue d'Armenonville) : f\textsuperscript{o}~26~r\textsuperscript{o} (p.~\pageref{fol:52}), f\textsuperscript{o}~65~r\textsuperscript{o} (p.~\pageref{fol:130}), f\textsuperscript{o}~65~v\textsuperscript{o} (p.~\pageref{fol:131}), f\textsuperscript{o}~68~r\textsuperscript{o} (p.~\pageref{fol:136}).\par\vspace{0.55mm}
\noindent\hangindent=4.5mm\textbf{First, Lieutenant}, officier britannique à Forceville : f\textsuperscript{o}~25~r\textsuperscript{o} (p.~\pageref{fol:50}).\par\vspace{0.55mm}
\noindent\hangindent=4.5mm\textbf{Flutre, Fortuné \& Louis-Fernand (1892–1978)}, instituteur du Mesnil-Martinsart et son fils élève de l'École normale supérieure : f\textsuperscript{o}~84~r\textsuperscript{o} (p.~\pageref{fol:168}).\par\vspace{0.55mm}
\noindent\hangindent=4.5mm\textbf{Ford, Capitaine E.~M.}, officier britannique du \textit{18th Divisional Train, A.S.C.}, décoré de la \textit{Military Cross} : f\textsuperscript{o}~25~r\textsuperscript{o} (p.~\pageref{fol:50}), f\textsuperscript{o}~38~r\textsuperscript{o} (p.~\pageref{fol:76}).\par\vspace{0.55mm}
\noindent\hangindent=4.5mm\textbf{Gilbert, Second Lieutenant}, sous-lieutenant britannique originaire de Jersey, poète, attaché au \textit{Purchase Board} à Hérissart : f\textsuperscript{o}~83~r\textsuperscript{o} (p.~\pageref{fol:166}), f\textsuperscript{o}~89~v\textsuperscript{o} (p.~\pageref{fol:179}), f\textsuperscript{o}~90~v\textsuperscript{o} (p.~\pageref{fol:181}), f\textsuperscript{o}~92~v\textsuperscript{o} (p.~\pageref{fol:185}), f\textsuperscript{o}~93~r\textsuperscript{o} (p.~\pageref{fol:186}).\par\vspace{0.55mm}
\noindent\hangindent=4.5mm\textbf{Godard, Lieutenant}, jeune officier britannique du \textit{Purchase Board} à Toutencourt et Hérissart : f\textsuperscript{o}~83~r\textsuperscript{o} (p.~\pageref{fol:166}).\par\vspace{0.55mm}
\noindent\hangindent=4.5mm\textbf{Grose, Lieutenant-Colonel Daniel Charles Evans}, commandant (\textit{Officer Commanding}) le \textit{18th Divisional Train, A.S.C.} : f\textsuperscript{o}~45~r\textsuperscript{o} (p.~\pageref{fol:90}).\par\vspace{0.55mm}
\noindent\hangindent=4.5mm\textbf{Guérin, Maurice de (1810–1839)}, poète tarnais étudié dans le cours public d'Ernest Zyromski à Toulouse : f\textsuperscript{o}~67~v\textsuperscript{o} (p.~\pageref{fol:135}).\par\vspace{0.55mm}
\noindent\hangindent=4.5mm\textbf{Guibet, M.}, habitant de Raincheval (Somme) : f\textsuperscript{o}~57~v\textsuperscript{o} (p.~\pageref{fol:115}), f\textsuperscript{o}~60~v\textsuperscript{o} (p.~\pageref{fol:121}).\par\vspace{0.55mm}
\noindent\hangindent=4.5mm\textbf{Guiraud}, tailleur à Toulouse : f\textsuperscript{o}~30~v\textsuperscript{o} (p.~\pageref{fol:61}).\par\vspace{0.55mm}
\noindent\hangindent=4.5mm\textbf{Houthouan, M. \& famille}, exploitants agricoles belges de la ferme de Belle-Église à Arquèves (Somme) : f\textsuperscript{o}~47~v\textsuperscript{o} (p.~\pageref{fol:95}), f\textsuperscript{o}~54~r\textsuperscript{o} (p.~\pageref{fol:108}), f\textsuperscript{o}~85~v\textsuperscript{o} (p.~\pageref{fol:171}).\par\vspace{0.55mm}
\noindent\hangindent=4.5mm\textbf{Hughes, Major}, officier supérieur britannique au \textit{18th Divisional Train} : f\textsuperscript{o}~45~r\textsuperscript{o} (p.~\pageref{fol:90}), f\textsuperscript{o}~48~r\textsuperscript{o} (p.~\pageref{fol:96}).\par\vspace{0.55mm}
\noindent\hangindent=4.5mm\textbf{Huguet, Capitaine (\textit{Cap.~H.})}, officier de la Mission militaire française auprès du XIII\textsuperscript{e} Corps à Acheux : f\textsuperscript{o}~47~r\textsuperscript{o} (p.~\pageref{fol:94}), f\textsuperscript{o}~52~v\textsuperscript{o} (p.~\pageref{fol:105}), f\textsuperscript{o}~61~v\textsuperscript{o} (p.~\pageref{fol:123}).\par\vspace{0.55mm}
\noindent\hangindent=4.5mm\textbf{Imbert}, opticien à Paris : f\textsuperscript{o}~58~r\textsuperscript{o} (p.~\pageref{fol:116}), f\textsuperscript{o}~60~v\textsuperscript{o} (p.~\pageref{fol:121}), f\textsuperscript{o}~61~r\textsuperscript{o} (p.~\pageref{fol:122}), f\textsuperscript{o}~65~v\textsuperscript{o} (p.~\pageref{fol:131}).\par\vspace{0.55mm}
\noindent\hangindent=4.5mm\textbf{Jimmie}, conducteur britannique de l'automobile Vauxhall du service des achats : f\textsuperscript{o}~83~v\textsuperscript{o} (p.~\pageref{fol:167}).\par\vspace{0.55mm}
\noindent\hangindent=4.5mm\textbf{Labit, Antoine (1832–1912) \& Georges (1862–1899)}, fondateur du grand magasin toulousain \textit{La Maison Universelle} et son fils : f\textsuperscript{o}~67~r\textsuperscript{o} (p.~\pageref{fol:134}).\par\vspace{0.55mm}
\noindent\hangindent=4.5mm\textbf{Layrals, M\textsuperscript{e} Charles \& Thérèse Cartairade (\textit{Layral})}, notaire à Millau et son épouse, oncle et tante maternels de Marguerite Calmels : f\textsuperscript{o}~71~r\textsuperscript{o} (p.~\pageref{fol:142}).\par\vspace{0.55mm}
\noindent\hangindent=4.5mm\textbf{Lépin (ou \textit{Liepin}, \textit{Lupin})}, adjudant interprète à la Mission française de la 18\textsuperscript{e} Division : f\textsuperscript{o}~11~v\textsuperscript{o} (p.~\pageref{fol:23}), f\textsuperscript{o}~42~v\textsuperscript{o} (p.~\pageref{fol:85}), f\textsuperscript{o}~45~r\textsuperscript{o} (p.~\pageref{fol:90}).\par\vspace{0.55mm}
\noindent\hangindent=4.5mm\textbf{Lersay, M. de}, chirurgien-dentiste à Toulouse : f\textsuperscript{o}~69~r\textsuperscript{o} (p.~\pageref{fol:138}), f\textsuperscript{o}~70~r\textsuperscript{o} (p.~\pageref{fol:140}).\par\vspace{0.55mm}
\noindent\hangindent=4.5mm\textbf{Levy, Major}, officier supérieur britannique dirigeant le service des achats (\textit{Purchase Board}) dans la Somme : f\textsuperscript{o}~45~v\textsuperscript{o} (p.~\pageref{fol:91}), f\textsuperscript{o}~46~v\textsuperscript{o} (p.~\pageref{fol:93}), f\textsuperscript{o}~55~r\textsuperscript{o} (p.~\pageref{fol:110}), f\textsuperscript{o}~58~r\textsuperscript{o} (p.~\pageref{fol:116}), f\textsuperscript{o}~83~r\textsuperscript{o} (p.~\pageref{fol:166}).\par\vspace{0.55mm}
\noindent\hangindent=4.5mm\textbf{Louis XI (1423–1483)}, roi de France, associé au chêne historique du beffroi de Lucheux (édit de 1464) : f\textsuperscript{o}~20~v\textsuperscript{o} (p.~\pageref{fol:41}).\par\vspace{0.55mm}
\noindent\hangindent=4.5mm\textbf{Lovet, Capitaine}, officier britannique revu au \textit{refilling point} près d'Albert : f\textsuperscript{o}~75~v\textsuperscript{o} (p.~\pageref{fol:151}).\par\vspace{0.55mm}
\noindent\hangindent=4.5mm\textbf{Maeterlinck, Maurice (1862–1949)}, écrivain belge, auteur du recueil \textit{Les Débris de la guerre} (1916) acquis par Hippolyte Combes : f\textsuperscript{o}~26~v\textsuperscript{o} (p.~\pageref{fol:53}).\par\vspace{0.55mm}
\noindent\hangindent=4.5mm\textbf{Martin, Alexis}, auteur du guide illustré \textit{Une visite à Amiens} rapporté à Hippolyte Combes : f\textsuperscript{o}~91~r\textsuperscript{o} (p.~\pageref{fol:182}).\par\vspace{0.55mm}
\noindent\hangindent=4.5mm\textbf{Maurin}, collègue interprète militaire dans la vallée de la Lys : f\textsuperscript{o}~6~v\textsuperscript{o} (p.~\pageref{fol:13}), f\textsuperscript{o}~7~r\textsuperscript{o} (p.~\pageref{fol:14}), f\textsuperscript{o}~9~v\textsuperscript{o} (p.~\pageref{fol:19}), f\textsuperscript{o}~10~v\textsuperscript{o} (p.~\pageref{fol:21}), f\textsuperscript{o}~12~r\textsuperscript{o} (p.~\pageref{fol:24}), f\textsuperscript{o}~13~r\textsuperscript{o} (p.~\pageref{fol:26}).\par\vspace{0.55mm}
\noindent\hangindent=4.5mm\textbf{Maxse, Major-General Sir Ivor (1862–1958)}, commandant la 18\textsuperscript{e} Division britannique (\textit{18th Eastern Division}) de 1914 à janvier 1917 : f\textsuperscript{o}~8~r\textsuperscript{o} (p.~\pageref{fol:16}).\par\vspace{0.55mm}
\noindent\hangindent=4.5mm\textbf{Michel, Lieutenant}, officier britannique de cantonnement (\textit{billeting officer}) à Lynde : f\textsuperscript{o}~1~v\textsuperscript{o} (p.~\pageref{fol:3}).\par\vspace{0.55mm}
\noindent\hangindent=4.5mm\textbf{Monin}, collègue interprète retrouvé de passage à Hérissart : f\textsuperscript{o}~86~v\textsuperscript{o} (p.~\pageref{fol:173}).\par\vspace{0.55mm}
\noindent\hangindent=4.5mm\textbf{Morris, Lieutenant}, officier britannique ami d'Hippolyte Combes à la 150\textsuperscript{e} Compagnie A.S.C. : f\textsuperscript{o}~60~r\textsuperscript{o} (p.~\pageref{fol:120}), f\textsuperscript{o}~60~v\textsuperscript{o} (p.~\pageref{fol:121}), f\textsuperscript{o}~61~r\textsuperscript{o} (p.~\pageref{fol:122}).\par\vspace{0.55mm}
\noindent\hangindent=4.5mm\textbf{Mortreux (famille)}, famille d'agriculteurs de Varennes (Somme) amie d'Hippolyte Combes : f\textsuperscript{o}~48~v\textsuperscript{o} (p.~\pageref{fol:97}), f\textsuperscript{o}~92~r\textsuperscript{o} (p.~\pageref{fol:184}).\par\vspace{0.55mm}
\noindent\hangindent=4.5mm\textbf{Moulard (ferme)}, exploitation agricole à Septenville (Rubempré, Somme) : f\textsuperscript{o}~60~r\textsuperscript{o} (p.~\pageref{fol:120}).\par\vspace{0.55mm}
\noindent\hangindent=4.5mm\textbf{Newton}, militaire britannique mentionné à Forceville : f\textsuperscript{o}~24~v\textsuperscript{o} (p.~\pageref{fol:49}).\par\vspace{0.55mm}
\noindent\hangindent=4.5mm\textbf{Panton, Capitaine}, officier britannique de la 152\textsuperscript{e} Compagnie A.S.C. : f\textsuperscript{o}~23~v\textsuperscript{o} (p.~\pageref{fol:47}), f\textsuperscript{o}~34~r\textsuperscript{o} (p.~\pageref{fol:68}), f\textsuperscript{o}~36~v\textsuperscript{o} (p.~\pageref{fol:73}), f\textsuperscript{o}~37~r\textsuperscript{o} (p.~\pageref{fol:74}), f\textsuperscript{o}~45~r\textsuperscript{o} (p.~\pageref{fol:90}).\par\vspace{0.55mm}
\noindent\hangindent=4.5mm\textbf{Pfister}, horloger-bijoutier à Toulouse (rue d'Alsace-Lorraine / place Esquirol) : f\textsuperscript{o}~28~r\textsuperscript{o} (p.~\pageref{fol:56}), f\textsuperscript{o}~66~v\textsuperscript{o} (p.~\pageref{fol:133}), f\textsuperscript{o}~67~r\textsuperscript{o} (p.~\pageref{fol:134}), f\textsuperscript{o}~67~v\textsuperscript{o} (p.~\pageref{fol:135}), f\textsuperscript{o}~68~v\textsuperscript{o} (p.~\pageref{fol:137}), f\textsuperscript{o}~69~r\textsuperscript{o} (p.~\pageref{fol:138}).\par\vspace{0.55mm}
\noindent\hangindent=4.5mm\textbf{Pizzi, Lieutenant}, officier interprète chef du détachement français auprès de la 18\textsuperscript{e} Division britannique : f\textsuperscript{o}~24~r\textsuperscript{o} (p.~\pageref{fol:48}), f\textsuperscript{o}~42~v\textsuperscript{o} (p.~\pageref{fol:85}), f\textsuperscript{o}~45~r\textsuperscript{o} (p.~\pageref{fol:90}), f\textsuperscript{o}~46~r\textsuperscript{o} (p.~\pageref{fol:92}).\par\vspace{0.55mm}
\noindent\hangindent=4.5mm\textbf{Raucoule}, collègue interprète militaire originaire du Midi : f\textsuperscript{o}~44~r\textsuperscript{o} (p.~\pageref{fol:88}), f\textsuperscript{o}~45~r\textsuperscript{o} (p.~\pageref{fol:90}), f\textsuperscript{o}~81~v\textsuperscript{o} (p.~\pageref{fol:163}).\par\vspace{0.55mm}
\noindent\hangindent=4.5mm\textbf{Schumacher}, adjudant interprète au bureau de la Mission française de la 18\textsuperscript{e} Division : f\textsuperscript{o}~11~v\textsuperscript{o} (p.~\pageref{fol:23}), f\textsuperscript{o}~41~v\textsuperscript{o} (p.~\pageref{fol:83}).\par\vspace{0.55mm}
\noindent\hangindent=4.5mm\textbf{Serpin}, sous-officier interprète à la Mission française : f\textsuperscript{o}~45~r\textsuperscript{o} (p.~\pageref{fol:90}).\par\vspace{0.55mm}
\noindent\hangindent=4.5mm\textbf{Spagnol, M\textsuperscript{me}}, horticultrice-fleuriste à Nice ayant expédié un panier de fleurs à Toulouse : f\textsuperscript{o}~67~r\textsuperscript{o} (p.~\pageref{fol:134}), f\textsuperscript{o}~68~r\textsuperscript{o} (p.~\pageref{fol:136}), f\textsuperscript{o}~79~r\textsuperscript{o} (p.~\pageref{fol:158}).\par\vspace{0.55mm}
\noindent\hangindent=4.5mm\textbf{Valet}, interprète militaire français rattaché à la 2\textsuperscript{e} Division australienne à Hérissart : f\textsuperscript{o}~82~r\textsuperscript{o} (p.~\pageref{fol:164}), f\textsuperscript{o}~83~r\textsuperscript{o} (p.~\pageref{fol:166}).\par\vspace{0.55mm}
\noindent\hangindent=4.5mm\textbf{Vauchel, M\textsuperscript{me}}, exploitante agricole à Marquay (Pas-de-Calais) : f\textsuperscript{o}~18~v\textsuperscript{o} (p.~\pageref{fol:37}).\par\vspace{0.55mm}
\noindent\hangindent=4.5mm\textbf{Vincens, D\textsuperscript{r} Fernand, Germaine \& Robert}, beau-frère, nièce et neveu de Marguerite Calmels à Millau : f\textsuperscript{o}~28~r\textsuperscript{o} (p.~\pageref{fol:56}), f\textsuperscript{o}~65~v\textsuperscript{o} (p.~\pageref{fol:131}), f\textsuperscript{o}~70~v\textsuperscript{o} (p.~\pageref{fol:141}), f\textsuperscript{o}~71~v\textsuperscript{o} (p.~\pageref{fol:143}).\par\vspace{0.55mm}
\noindent\hangindent=4.5mm\textbf{Walsh, Lieutenant}, officier britannique du \textit{18th Divisional Train, A.S.C.} : f\textsuperscript{o}~37~r\textsuperscript{o} (p.~\pageref{fol:74}), f\textsuperscript{o}~49~v\textsuperscript{o} (p.~\pageref{fol:99}), f\textsuperscript{o}~51~v\textsuperscript{o} (p.~\pageref{fol:103}), f\textsuperscript{o}~53~r\textsuperscript{o} (p.~\pageref{fol:106}), f\textsuperscript{o}~60~r\textsuperscript{o} (p.~\pageref{fol:120}).\par\vspace{0.55mm}
\noindent\hangindent=4.5mm\textbf{Wentworth, Lieutenant}, officier britannique du \textit{Purchase Board} à Hérissart : f\textsuperscript{o}~83~r\textsuperscript{o} (p.~\pageref{fol:166}), f\textsuperscript{o}~84~v\textsuperscript{o} (p.~\pageref{fol:169}), f\textsuperscript{o}~85~r\textsuperscript{o} (p.~\pageref{fol:170}), f\textsuperscript{o}~88~v\textsuperscript{o} (p.~\pageref{fol:177}).\par\vspace{0.55mm}
\noindent\hangindent=4.5mm\textbf{West, Capitaine E.~M.}, commandant la 152\textsuperscript{e} Compagnie A.S.C., décoré de la \textit{Military Cross} : f\textsuperscript{o}~19~v\textsuperscript{o} (p.~\pageref{fol:39}), f\textsuperscript{o}~21~r\textsuperscript{o} (p.~\pageref{fol:42}), f\textsuperscript{o}~24~r\textsuperscript{o} (p.~\pageref{fol:48}), f\textsuperscript{o}~37~r\textsuperscript{o} (p.~\pageref{fol:74}).\par\vspace{0.55mm}
\noindent\hangindent=4.5mm\textbf{Williams, Capitaine (\textit{Cap.~W.}, \textit{W\textsuperscript{ms}})}, jeune officier britannique (\textit{Foraging Officer}) à Acheux puis Hérissart de décembre 1916 à mars 1917 : f\textsuperscript{o}~52~v\textsuperscript{o} (p.~\pageref{fol:105}), f\textsuperscript{o}~53~r\textsuperscript{o} (p.~\pageref{fol:106}), f\textsuperscript{o}~53~v\textsuperscript{o} (p.~\pageref{fol:107}), f\textsuperscript{o}~54~r\textsuperscript{o} (p.~\pageref{fol:108}), f\textsuperscript{o}~54~v\textsuperscript{o} (p.~\pageref{fol:109}), f\textsuperscript{o}~55~v\textsuperscript{o} (p.~\pageref{fol:111}), f\textsuperscript{o}~57~r\textsuperscript{o} (p.~\pageref{fol:114}), f\textsuperscript{o}~57~v\textsuperscript{o} (p.~\pageref{fol:115}), f\textsuperscript{o}~58~v\textsuperscript{o} (p.~\pageref{fol:117}), f\textsuperscript{o}~63~r\textsuperscript{o} (p.~\pageref{fol:126}), f\textsuperscript{o}~73~v\textsuperscript{o} (p.~\pageref{fol:147}), f\textsuperscript{o}~74~r\textsuperscript{o} (p.~\pageref{fol:148}), f\textsuperscript{o}~78~v\textsuperscript{o} (p.~\pageref{fol:157}), f\textsuperscript{o}~80~r\textsuperscript{o} (p.~\pageref{fol:160}), f\textsuperscript{o}~81~v\textsuperscript{o} (p.~\pageref{fol:163}), f\textsuperscript{o}~82~r\textsuperscript{o} (p.~\pageref{fol:164}), f\textsuperscript{o}~82~v\textsuperscript{o} (p.~\pageref{fol:165}).\par\vspace{0.55mm}
\noindent\hangindent=4.5mm\textbf{Zyromski, Ernest (1862–1933)}, professeur de littérature française à la Faculté des lettres de Toulouse, auteur de \textit{L'Orgueil humain} : f\textsuperscript{o}~26~v\textsuperscript{o} (p.~\pageref{fol:53}), f\textsuperscript{o}~67~v\textsuperscript{o} (p.~\pageref{fol:135}).\par\vspace{0.55mm}
\par}

\cleardoublepage
\chapter*{Index des noms de lieux (\textit{Index locorum})}
\addcontentsline{toc}{chapter}{Index des noms de lieux (\textit{Index locorum})}
\markboth{Index des noms de lieux}{Index des noms de lieux}
{\fontsize{7.0pt}{8.3pt}\selectfont
\noindent Chaque toponyme renvoie simultanément au folio du \textit{Livre-Journal} (\texttt{f\textsuperscript{o}~X~r\textsuperscript{o}/v\textsuperscript{o}}) et au numéro de page correspondant du présent volume (\texttt{p.~Y}).\par\vspace{2.0mm}
\noindent\hangindent=4.5mm\textbf{Abbeville (Somme)}, chef-lieu d'arrondissement de la Somme : f\textsuperscript{o}~90~v\textsuperscript{o} (p.~\pageref{fol:181}).\par\vspace{0.55mm}
\noindent\hangindent=4.5mm\textbf{Acheux-en-Amiénois (Somme)}, chef-lieu de canton de la Somme, gare (\textit{railhead} d'Acheux-Varennes), château et cantonnement d'Hippolyte Combes (nov. 1916 – févr. 1917) : f\textsuperscript{o}~21~v\textsuperscript{o} (p.~\pageref{fol:43}), f\textsuperscript{o}~22~v\textsuperscript{o} (p.~\pageref{fol:45}), f\textsuperscript{o}~23~r\textsuperscript{o} (p.~\pageref{fol:46}), f\textsuperscript{o}~24~r\textsuperscript{o} (p.~\pageref{fol:48}), f\textsuperscript{o}~24~v\textsuperscript{o} (p.~\pageref{fol:49}), f\textsuperscript{o}~25~r\textsuperscript{o} (p.~\pageref{fol:50}), f\textsuperscript{o}~43~v\textsuperscript{o} (p.~\pageref{fol:87}), f\textsuperscript{o}~44~v\textsuperscript{o} (p.~\pageref{fol:89}), f\textsuperscript{o}~45~r\textsuperscript{o} (p.~\pageref{fol:90}), f\textsuperscript{o}~45~v\textsuperscript{o} (p.~\pageref{fol:91}), f\textsuperscript{o}~46~r\textsuperscript{o} (p.~\pageref{fol:92}), f\textsuperscript{o}~46~v\textsuperscript{o} (p.~\pageref{fol:93}), f\textsuperscript{o}~47~r\textsuperscript{o} (p.~\pageref{fol:94}), f\textsuperscript{o}~47~v\textsuperscript{o} (p.~\pageref{fol:95}), f\textsuperscript{o}~48~r\textsuperscript{o} (p.~\pageref{fol:96}), f\textsuperscript{o}~48~v\textsuperscript{o} (p.~\pageref{fol:97}), f\textsuperscript{o}~49~r\textsuperscript{o} (p.~\pageref{fol:98}), f\textsuperscript{o}~49~v\textsuperscript{o} (p.~\pageref{fol:99}), f\textsuperscript{o}~50~r\textsuperscript{o} (p.~\pageref{fol:100}), f\textsuperscript{o}~50~v\textsuperscript{o} (p.~\pageref{fol:101}), f\textsuperscript{o}~51~r\textsuperscript{o} (p.~\pageref{fol:102}), f\textsuperscript{o}~51~v\textsuperscript{o} (p.~\pageref{fol:103}), f\textsuperscript{o}~52~r\textsuperscript{o} (p.~\pageref{fol:104}), f\textsuperscript{o}~52~v\textsuperscript{o} (p.~\pageref{fol:105}), f\textsuperscript{o}~53~r\textsuperscript{o} (p.~\pageref{fol:106}), f\textsuperscript{o}~53~v\textsuperscript{o} (p.~\pageref{fol:107}), f\textsuperscript{o}~54~r\textsuperscript{o} (p.~\pageref{fol:108}), f\textsuperscript{o}~54~v\textsuperscript{o} (p.~\pageref{fol:109}), f\textsuperscript{o}~55~r\textsuperscript{o} (p.~\pageref{fol:110}), f\textsuperscript{o}~55~v\textsuperscript{o} (p.~\pageref{fol:111}), f\textsuperscript{o}~56~r\textsuperscript{o} (p.~\pageref{fol:112}), f\textsuperscript{o}~56~v\textsuperscript{o} (p.~\pageref{fol:113}), f\textsuperscript{o}~57~r\textsuperscript{o} (p.~\pageref{fol:114}), f\textsuperscript{o}~57~v\textsuperscript{o} (p.~\pageref{fol:115}), f\textsuperscript{o}~58~r\textsuperscript{o} (p.~\pageref{fol:116}), f\textsuperscript{o}~58~v\textsuperscript{o} (p.~\pageref{fol:117}), f\textsuperscript{o}~59~r\textsuperscript{o} (p.~\pageref{fol:118}), f\textsuperscript{o}~59~v\textsuperscript{o} (p.~\pageref{fol:119}), f\textsuperscript{o}~60~r\textsuperscript{o} (p.~\pageref{fol:120}), f\textsuperscript{o}~60~v\textsuperscript{o} (p.~\pageref{fol:121}), f\textsuperscript{o}~61~r\textsuperscript{o} (p.~\pageref{fol:122}), f\textsuperscript{o}~61~v\textsuperscript{o} (p.~\pageref{fol:123}), f\textsuperscript{o}~62~r\textsuperscript{o} (p.~\pageref{fol:124}), f\textsuperscript{o}~62~v\textsuperscript{o} (p.~\pageref{fol:125}), f\textsuperscript{o}~63~r\textsuperscript{o} (p.~\pageref{fol:126}), f\textsuperscript{o}~63~v\textsuperscript{o} (p.~\pageref{fol:127}), f\textsuperscript{o}~64~r\textsuperscript{o} (p.~\pageref{fol:128}), f\textsuperscript{o}~64~v\textsuperscript{o} (p.~\pageref{fol:129}), f\textsuperscript{o}~73~r\textsuperscript{o} (p.~\pageref{fol:146}), f\textsuperscript{o}~73~v\textsuperscript{o} (p.~\pageref{fol:147}), f\textsuperscript{o}~74~r\textsuperscript{o} (p.~\pageref{fol:148}), f\textsuperscript{o}~74~v\textsuperscript{o} (p.~\pageref{fol:149}), f\textsuperscript{o}~75~v\textsuperscript{o} (p.~\pageref{fol:151}), f\textsuperscript{o}~77~r\textsuperscript{o} (p.~\pageref{fol:154}), f\textsuperscript{o}~84~r\textsuperscript{o} (p.~\pageref{fol:168}), f\textsuperscript{o}~88~r\textsuperscript{o} (p.~\pageref{fol:176}), f\textsuperscript{o}~88~v\textsuperscript{o} (p.~\pageref{fol:177}), f\textsuperscript{o}~91~v\textsuperscript{o} (p.~\pageref{fol:183}).\par\vspace{0.55mm}
\noindent\hangindent=4.5mm\textbf{Achiet-le-Petit (Pas-de-Calais)}, village au nord-est de Miraumont repris lors du repli allemand de mars 1917 : f\textsuperscript{o}~85~r\textsuperscript{o} (p.~\pageref{fol:170}).\par\vspace{0.55mm}
\noindent\hangindent=4.5mm\textbf{Albert (Somme)}, ville de la vallée de l'Ancre, basilique Notre-Dame de Brebières (\textit{Vierge penchée}) et bivouac des briqueteries (\textit{Brickfields}) : f\textsuperscript{o}~39~r\textsuperscript{o} (p.~\pageref{fol:78}), f\textsuperscript{o}~39~v\textsuperscript{o} (p.~\pageref{fol:79}), f\textsuperscript{o}~40~r\textsuperscript{o} (p.~\pageref{fol:80}), f\textsuperscript{o}~40~v\textsuperscript{o} (p.~\pageref{fol:81}), f\textsuperscript{o}~41~r\textsuperscript{o} (p.~\pageref{fol:82}), f\textsuperscript{o}~41~v\textsuperscript{o} (p.~\pageref{fol:83}), f\textsuperscript{o}~42~r\textsuperscript{o} (p.~\pageref{fol:84}), f\textsuperscript{o}~42~v\textsuperscript{o} (p.~\pageref{fol:85}), f\textsuperscript{o}~43~r\textsuperscript{o} (p.~\pageref{fol:86}), f\textsuperscript{o}~43~v\textsuperscript{o} (p.~\pageref{fol:87}), f\textsuperscript{o}~44~r\textsuperscript{o} (p.~\pageref{fol:88}), f\textsuperscript{o}~44~v\textsuperscript{o} (p.~\pageref{fol:89}), f\textsuperscript{o}~45~r\textsuperscript{o} (p.~\pageref{fol:90}), f\textsuperscript{o}~45~v\textsuperscript{o} (p.~\pageref{fol:91}), f\textsuperscript{o}~46~r\textsuperscript{o} (p.~\pageref{fol:92}), f\textsuperscript{o}~46~v\textsuperscript{o} (p.~\pageref{fol:93}), f\textsuperscript{o}~47~v\textsuperscript{o} (p.~\pageref{fol:95}), f\textsuperscript{o}~48~r\textsuperscript{o} (p.~\pageref{fol:96}), f\textsuperscript{o}~49~v\textsuperscript{o} (p.~\pageref{fol:99}), f\textsuperscript{o}~50~r\textsuperscript{o} (p.~\pageref{fol:100}), f\textsuperscript{o}~75~v\textsuperscript{o} (p.~\pageref{fol:151}), f\textsuperscript{o}~82~r\textsuperscript{o} (p.~\pageref{fol:164}), f\textsuperscript{o}~84~v\textsuperscript{o} (p.~\pageref{fol:169}), f\textsuperscript{o}~85~v\textsuperscript{o} (p.~\pageref{fol:171}), f\textsuperscript{o}~90~v\textsuperscript{o} (p.~\pageref{fol:181}), f\textsuperscript{o}~91~r\textsuperscript{o} (p.~\pageref{fol:182}).\par\vspace{0.55mm}
\noindent\hangindent=4.5mm\textbf{Amiens (Somme)}, préfecture de la Somme, cathédrale Notre-Dame et nœud ferroviaire des permissions : f\textsuperscript{o}~25~v\textsuperscript{o} (p.~\pageref{fol:51}), f\textsuperscript{o}~32~v\textsuperscript{o} (p.~\pageref{fol:65}), f\textsuperscript{o}~44~v\textsuperscript{o} (p.~\pageref{fol:89}), f\textsuperscript{o}~49~r\textsuperscript{o} (p.~\pageref{fol:98}), f\textsuperscript{o}~59~v\textsuperscript{o} (p.~\pageref{fol:119}), f\textsuperscript{o}~63~r\textsuperscript{o} (p.~\pageref{fol:126}), f\textsuperscript{o}~64~r\textsuperscript{o} (p.~\pageref{fol:128}), f\textsuperscript{o}~64~v\textsuperscript{o} (p.~\pageref{fol:129}), f\textsuperscript{o}~65~r\textsuperscript{o} (p.~\pageref{fol:130}), f\textsuperscript{o}~72~v\textsuperscript{o} (p.~\pageref{fol:145}), f\textsuperscript{o}~73~r\textsuperscript{o} (p.~\pageref{fol:146}), f\textsuperscript{o}~73~v\textsuperscript{o} (p.~\pageref{fol:147}), f\textsuperscript{o}~74~r\textsuperscript{o} (p.~\pageref{fol:148}), f\textsuperscript{o}~76~v\textsuperscript{o} (p.~\pageref{fol:153}), f\textsuperscript{o}~85~v\textsuperscript{o} (p.~\pageref{fol:171}), f\textsuperscript{o}~91~r\textsuperscript{o} (p.~\pageref{fol:182}).\par\vspace{0.55mm}
\noindent\hangindent=4.5mm\textbf{Ancre (vallée de l')}, affluent de la Somme traversant Miraumont, Albert, Aveluy, Buire, Ribemont et Heilly : f\textsuperscript{o}~49~v\textsuperscript{o} (p.~\pageref{fol:99}), f\textsuperscript{o}~52~v\textsuperscript{o} (p.~\pageref{fol:105}), f\textsuperscript{o}~76~v\textsuperscript{o} (p.~\pageref{fol:153}), f\textsuperscript{o}~90~v\textsuperscript{o} (p.~\pageref{fol:181}).\par\vspace{0.55mm}
\noindent\hangindent=4.5mm\textbf{Antibes (Alpes-Maritimes)}, ville de la Côte d'Azur évoquée avec l'interprète Maurin : f\textsuperscript{o}~12~r\textsuperscript{o} (p.~\pageref{fol:24}).\par\vspace{0.55mm}
\noindent\hangindent=4.5mm\textbf{Armentières (Nord)}, ville industrielle linière sur la Lys, bombardée par l'artillerie allemande : f\textsuperscript{o}~8~v\textsuperscript{o} (p.~\pageref{fol:17}), f\textsuperscript{o}~9~v\textsuperscript{o} (p.~\pageref{fol:19}), f\textsuperscript{o}~10~v\textsuperscript{o} (p.~\pageref{fol:21}), f\textsuperscript{o}~11~r\textsuperscript{o} (p.~\pageref{fol:22}).\par\vspace{0.55mm}
\noindent\hangindent=4.5mm\textbf{Arquèves \& ferme de Belle-Église (Somme)}, commune du canton d'Acheux et ancienne commanderie exploitée par la famille Houthouan : f\textsuperscript{o}~47~v\textsuperscript{o} (p.~\pageref{fol:95}), f\textsuperscript{o}~53~r\textsuperscript{o} (p.~\pageref{fol:106}), f\textsuperscript{o}~54~r\textsuperscript{o} (p.~\pageref{fol:108}), f\textsuperscript{o}~54~v\textsuperscript{o} (p.~\pageref{fol:109}), f\textsuperscript{o}~55~r\textsuperscript{o} (p.~\pageref{fol:110}), f\textsuperscript{o}~55~v\textsuperscript{o} (p.~\pageref{fol:111}), f\textsuperscript{o}~56~v\textsuperscript{o} (p.~\pageref{fol:113}), f\textsuperscript{o}~57~v\textsuperscript{o} (p.~\pageref{fol:115}), f\textsuperscript{o}~58~v\textsuperscript{o} (p.~\pageref{fol:117}), f\textsuperscript{o}~59~v\textsuperscript{o} (p.~\pageref{fol:119}), f\textsuperscript{o}~60~v\textsuperscript{o} (p.~\pageref{fol:121}), f\textsuperscript{o}~61~r\textsuperscript{o} (p.~\pageref{fol:122}), f\textsuperscript{o}~62~v\textsuperscript{o} (p.~\pageref{fol:125}), f\textsuperscript{o}~85~v\textsuperscript{o} (p.~\pageref{fol:171}), f\textsuperscript{o}~87~v\textsuperscript{o} (p.~\pageref{fol:175}), f\textsuperscript{o}~89~r\textsuperscript{o} (p.~\pageref{fol:178}), f\textsuperscript{o}~89~v\textsuperscript{o} (p.~\pageref{fol:179}).\par\vspace{0.55mm}
\noindent\hangindent=4.5mm\textbf{Arras (Pas-de-Calais)}, préfecture du Pas-de-Calais et secteur de l'offensive britannique d'avril 1917 (bataille de la Scarpe) : f\textsuperscript{o}~16~r\textsuperscript{o} (p.~\pageref{fol:32}), f\textsuperscript{o}~59~r\textsuperscript{o} (p.~\pageref{fol:118}), f\textsuperscript{o}~89~r\textsuperscript{o} (p.~\pageref{fol:178}), f\textsuperscript{o}~91~r\textsuperscript{o} (p.~\pageref{fol:182}).\par\vspace{0.55mm}
\noindent\hangindent=4.5mm\textbf{Authie (Somme)}, village de la haute vallée de l'Authie au nord d'Acheux : f\textsuperscript{o}~50~r\textsuperscript{o} (p.~\pageref{fol:100}), f\textsuperscript{o}~51~r\textsuperscript{o} (p.~\pageref{fol:102}), f\textsuperscript{o}~53~v\textsuperscript{o} (p.~\pageref{fol:107}), f\textsuperscript{o}~54~r\textsuperscript{o} (p.~\pageref{fol:108}), f\textsuperscript{o}~54~v\textsuperscript{o} (p.~\pageref{fol:109}), f\textsuperscript{o}~56~r\textsuperscript{o} (p.~\pageref{fol:112}), f\textsuperscript{o}~59~r\textsuperscript{o} (p.~\pageref{fol:118}), f\textsuperscript{o}~60~v\textsuperscript{o} (p.~\pageref{fol:121}), f\textsuperscript{o}~61~r\textsuperscript{o} (p.~\pageref{fol:122}), f\textsuperscript{o}~77~r\textsuperscript{o} (p.~\pageref{fol:154}).\par\vspace{0.55mm}
\noindent\hangindent=4.5mm\textbf{Auxi-le-Château (Pas-de-Calais)}, chef-lieu de canton sur l'Authie visité en tournée : f\textsuperscript{o}~62~r\textsuperscript{o} (p.~\pageref{fol:124}), f\textsuperscript{o}~76~r\textsuperscript{o} (p.~\pageref{fol:152}), f\textsuperscript{o}~79~r\textsuperscript{o} (p.~\pageref{fol:158}).\par\vspace{0.55mm}
\noindent\hangindent=4.5mm\textbf{Aveluy (Somme)}, village et bois au nord d'Albert sur l'Ancre : f\textsuperscript{o}~41~r\textsuperscript{o} (p.~\pageref{fol:82}), f\textsuperscript{o}~75~r\textsuperscript{o} (p.~\pageref{fol:150}), f\textsuperscript{o}~75~v\textsuperscript{o} (p.~\pageref{fol:151}).\par\vspace{0.55mm}
\noindent\hangindent=4.5mm\textbf{Bac-Saint-Maur (Pas-de-Calais)}, écart de Sailly-sur-la-Lys et écluse sur la Lys (\textit{refilling point}) : f\textsuperscript{o}~6~v\textsuperscript{o} (p.~\pageref{fol:13}), f\textsuperscript{o}~7~r\textsuperscript{o} (p.~\pageref{fol:14}), f\textsuperscript{o}~7~v\textsuperscript{o} (p.~\pageref{fol:15}), f\textsuperscript{o}~8~r\textsuperscript{o} (p.~\pageref{fol:16}), f\textsuperscript{o}~9~r\textsuperscript{o} (p.~\pageref{fol:18}), f\textsuperscript{o}~9~v\textsuperscript{o} (p.~\pageref{fol:19}), f\textsuperscript{o}~10~r\textsuperscript{o} (p.~\pageref{fol:20}), f\textsuperscript{o}~12~r\textsuperscript{o} (p.~\pageref{fol:24}).\par\vspace{0.55mm}
\noindent\hangindent=4.5mm\textbf{Bailleul (Nord)}, ville de Flandre flamingante au pied des Monts, Grand-Place, beffroi et gare d'embarquement : f\textsuperscript{o}~3~v\textsuperscript{o} (p.~\pageref{fol:7}), f\textsuperscript{o}~4~r\textsuperscript{o} (p.~\pageref{fol:8}), f\textsuperscript{o}~4~v\textsuperscript{o} (p.~\pageref{fol:9}), f\textsuperscript{o}~5~r\textsuperscript{o} (p.~\pageref{fol:10}), f\textsuperscript{o}~6~r\textsuperscript{o} (p.~\pageref{fol:12}), f\textsuperscript{o}~13~v\textsuperscript{o} (p.~\pageref{fol:27}), f\textsuperscript{o}~14~r\textsuperscript{o} (p.~\pageref{fol:28}).\par\vspace{0.55mm}
\noindent\hangindent=4.5mm\textbf{Bailleul-aux-Cornailles (Pas-de-Calais)}, village du Ternois voisin de Marquay : f\textsuperscript{o}~14~v\textsuperscript{o} (p.~\pageref{fol:29}), f\textsuperscript{o}~15~v\textsuperscript{o} (p.~\pageref{fol:31}), f\textsuperscript{o}~18~v\textsuperscript{o} (p.~\pageref{fol:37}).\par\vspace{0.55mm}
\noindent\hangindent=4.5mm\textbf{Barlette (Somme)}, hameau de la commune de Ribeaucourt où cantonne la 152\textsuperscript{e} Compagnie A.S.C. : f\textsuperscript{o}~33~r\textsuperscript{o} (p.~\pageref{fol:66}), f\textsuperscript{o}~33~v\textsuperscript{o} (p.~\pageref{fol:67}), f\textsuperscript{o}~34~r\textsuperscript{o} (p.~\pageref{fol:68}), f\textsuperscript{o}~34~v\textsuperscript{o} (p.~\pageref{fol:69}), f\textsuperscript{o}~35~r\textsuperscript{o} (p.~\pageref{fol:70}), f\textsuperscript{o}~35~v\textsuperscript{o} (p.~\pageref{fol:71}), f\textsuperscript{o}~36~r\textsuperscript{o} (p.~\pageref{fol:72}), f\textsuperscript{o}~36~v\textsuperscript{o} (p.~\pageref{fol:73}), f\textsuperscript{o}~37~r\textsuperscript{o} (p.~\pageref{fol:74}).\par\vspace{0.55mm}
\noindent\hangindent=4.5mm\textbf{Beaucourt-sur-l'Hallue (Somme)}, village et château au sud d'Hérissart : f\textsuperscript{o}~49~v\textsuperscript{o} (p.~\pageref{fol:99}), f\textsuperscript{o}~61~v\textsuperscript{o} (p.~\pageref{fol:123}), f\textsuperscript{o}~77~v\textsuperscript{o} (p.~\pageref{fol:155}), f\textsuperscript{o}~91~v\textsuperscript{o} (p.~\pageref{fol:183}).\par\vspace{0.55mm}
\noindent\hangindent=4.5mm\textbf{Beaudricourt \& hameau d'Oppy (Pas-de-Calais)}, commune et hameau du canton d'Avesnes-le-Comte, étape du 9–10 septembre 1916 : f\textsuperscript{o}~19~v\textsuperscript{o} (p.~\pageref{fol:39}), f\textsuperscript{o}~20~r\textsuperscript{o} (p.~\pageref{fol:40}), f\textsuperscript{o}~20~v\textsuperscript{o} (p.~\pageref{fol:41}).\par\vspace{0.55mm}
\noindent\hangindent=4.5mm\textbf{Beaumont-Hamel (Somme)}, village du front de l'Ancre : f\textsuperscript{o}~49~r\textsuperscript{o} (p.~\pageref{fol:98}), f\textsuperscript{o}~49~v\textsuperscript{o} (p.~\pageref{fol:99}), f\textsuperscript{o}~83~v\textsuperscript{o} (p.~\pageref{fol:167}).\par\vspace{0.55mm}
\noindent\hangindent=4.5mm\textbf{Beauquesne (Somme)}, bourg du canton de Doullens : f\textsuperscript{o}~37~v\textsuperscript{o} (p.~\pageref{fol:75}), f\textsuperscript{o}~49~r\textsuperscript{o} (p.~\pageref{fol:98}), f\textsuperscript{o}~56~r\textsuperscript{o} (p.~\pageref{fol:112}), f\textsuperscript{o}~56~v\textsuperscript{o} (p.~\pageref{fol:113}), f\textsuperscript{o}~57~v\textsuperscript{o} (p.~\pageref{fol:115}), f\textsuperscript{o}~58~v\textsuperscript{o} (p.~\pageref{fol:117}), f\textsuperscript{o}~60~r\textsuperscript{o} (p.~\pageref{fol:120}), f\textsuperscript{o}~61~r\textsuperscript{o} (p.~\pageref{fol:122}), f\textsuperscript{o}~61~v\textsuperscript{o} (p.~\pageref{fol:123}), f\textsuperscript{o}~89~r\textsuperscript{o} (p.~\pageref{fol:178}).\par\vspace{0.55mm}
\noindent\hangindent=4.5mm\textbf{Beauval (Somme)}, bourg entre Doullens et Amiens, étape des 15–16 octobre 1916 : f\textsuperscript{o}~37~r\textsuperscript{o} (p.~\pageref{fol:74}), f\textsuperscript{o}~37~v\textsuperscript{o} (p.~\pageref{fol:75}), f\textsuperscript{o}~38~r\textsuperscript{o} (p.~\pageref{fol:76}).\par\vspace{0.55mm}
\noindent\hangindent=4.5mm\textbf{Bédarieux (Hérault)}, gare de bifurcation vers Millau sur la ligne des Causses : f\textsuperscript{o}~69~v\textsuperscript{o} (p.~\pageref{fol:139}), f\textsuperscript{o}~70~r\textsuperscript{o} (p.~\pageref{fol:140}).\par\vspace{0.55mm}
\noindent\hangindent=4.5mm\textbf{Bertrancourt (Somme)}, commune à l'est d'Acheux-en-Amiénois : f\textsuperscript{o}~46~r\textsuperscript{o} (p.~\pageref{fol:92}), f\textsuperscript{o}~52~v\textsuperscript{o} (p.~\pageref{fol:105}), f\textsuperscript{o}~60~v\textsuperscript{o} (p.~\pageref{fol:121}), f\textsuperscript{o}~84~r\textsuperscript{o} (p.~\pageref{fol:168}).\par\vspace{0.55mm}
\noindent\hangindent=4.5mm\textbf{Béziers (Hérault)}, ville du Languedoc visitée avec Marguerite Calmels (cathédrale Saint-Nazaire et terrasse sur l'Orb) : f\textsuperscript{o}~69~v\textsuperscript{o} (p.~\pageref{fol:139}), f\textsuperscript{o}~70~r\textsuperscript{o} (p.~\pageref{fol:140}), f\textsuperscript{o}~71~v\textsuperscript{o} (p.~\pageref{fol:143}).\par\vspace{0.55mm}
\noindent\hangindent=4.5mm\textbf{Bouzincourt (Somme)}, village à l'ouest d'Albert où bivouaque la compagnie en octobre 1916 : f\textsuperscript{o}~38~v\textsuperscript{o} (p.~\pageref{fol:77}), f\textsuperscript{o}~39~r\textsuperscript{o} (p.~\pageref{fol:78}), f\textsuperscript{o}~39~v\textsuperscript{o} (p.~\pageref{fol:79}), f\textsuperscript{o}~40~r\textsuperscript{o} (p.~\pageref{fol:80}), f\textsuperscript{o}~50~v\textsuperscript{o} (p.~\pageref{fol:101}), f\textsuperscript{o}~62~r\textsuperscript{o} (p.~\pageref{fol:124}).\par\vspace{0.55mm}
\noindent\hangindent=4.5mm\textbf{Bresle (Somme)}, village entre Warloy-Baillon et Albert : f\textsuperscript{o}~90~v\textsuperscript{o} (p.~\pageref{fol:181}).\par\vspace{0.55mm}
\noindent\hangindent=4.5mm\textbf{Brioude \& Arvant (Haute-Loire)}, gares de la ligne des Cévennes lors du retour de Millau en février 1917 : f\textsuperscript{o}~72~r\textsuperscript{o} (p.~\pageref{fol:144}), f\textsuperscript{o}~72~v\textsuperscript{o} (p.~\pageref{fol:145}).\par\vspace{0.55mm}
\noindent\hangindent=4.5mm\textbf{Buire-sur-l'Ancre (Somme)}, village de la vallée de l'Ancre en aval d'Albert : f\textsuperscript{o}~52~v\textsuperscript{o} (p.~\pageref{fol:105}), f\textsuperscript{o}~90~v\textsuperscript{o} (p.~\pageref{fol:181}), f\textsuperscript{o}~91~r\textsuperscript{o} (p.~\pageref{fol:182}).\par\vspace{0.55mm}
\noindent\hangindent=4.5mm\textbf{Bus-lès-Artois (Somme)}, commune au nord-est d'Acheux-en-Amiénois : f\textsuperscript{o}~46~r\textsuperscript{o} (p.~\pageref{fol:92}), f\textsuperscript{o}~50~r\textsuperscript{o} (p.~\pageref{fol:100}), f\textsuperscript{o}~77~r\textsuperscript{o} (p.~\pageref{fol:154}).\par\vspace{0.55mm}
\noindent\hangindent=4.5mm\textbf{Cambrai (Nord)}, ville du Nord évoquée comme destination d'offensive au printemps 1917 : f\textsuperscript{o}~92~v\textsuperscript{o} (p.~\pageref{fol:185}).\par\vspace{0.55mm}
\noindent\hangindent=4.5mm\textbf{Canaples (Somme)}, commune et gare de ravitaillement (\textit{railhead}) au nord d'Amiens : f\textsuperscript{o}~33~r\textsuperscript{o} (p.~\pageref{fol:66}).\par\vspace{0.55mm}
\noindent\hangindent=4.5mm\textbf{Candas (Somme)}, commune entre Doullens et Rubempré : f\textsuperscript{o}~83~v\textsuperscript{o} (p.~\pageref{fol:167}).\par\vspace{0.55mm}
\noindent\hangindent=4.5mm\textbf{Carcassonne (Aude)}, escale ferroviaire où Hippolyte Combes et Marguerite Calmels embrassent l'abbé Jean Crouzet : f\textsuperscript{o}~69~v\textsuperscript{o} (p.~\pageref{fol:139}).\par\vspace{0.55mm}
\noindent\hangindent=4.5mm\textbf{Clermont-Ferrand (Puy-de-Dôme)}, étape ferroviaire du retour de Millau le 6 février 1917 : f\textsuperscript{o}~72~r\textsuperscript{o} (p.~\pageref{fol:144}), f\textsuperscript{o}~72~v\textsuperscript{o} (p.~\pageref{fol:145}).\par\vspace{0.55mm}
\noindent\hangindent=4.5mm\textbf{Coisy (Somme)}, village au sud-ouest de Rubempré : f\textsuperscript{o}~47~v\textsuperscript{o} (p.~\pageref{fol:95}).\par\vspace{0.55mm}
\noindent\hangindent=4.5mm\textbf{Contay (Somme)}, village et château de la vallée de l'Hallue entre Toutencourt et Warloy : f\textsuperscript{o}~49~r\textsuperscript{o} (p.~\pageref{fol:98}), f\textsuperscript{o}~50~v\textsuperscript{o} (p.~\pageref{fol:101}), f\textsuperscript{o}~53~v\textsuperscript{o} (p.~\pageref{fol:107}), f\textsuperscript{o}~79~v\textsuperscript{o} (p.~\pageref{fol:159}), f\textsuperscript{o}~80~v\textsuperscript{o} (p.~\pageref{fol:161}), f\textsuperscript{o}~83~v\textsuperscript{o} (p.~\pageref{fol:167}), f\textsuperscript{o}~84~v\textsuperscript{o} (p.~\pageref{fol:169}), f\textsuperscript{o}~85~v\textsuperscript{o} (p.~\pageref{fol:171}), f\textsuperscript{o}~88~r\textsuperscript{o} (p.~\pageref{fol:176}), f\textsuperscript{o}~88~v\textsuperscript{o} (p.~\pageref{fol:177}), f\textsuperscript{o}~89~r\textsuperscript{o} (p.~\pageref{fol:178}), f\textsuperscript{o}~90~v\textsuperscript{o} (p.~\pageref{fol:181}), f\textsuperscript{o}~91~r\textsuperscript{o} (p.~\pageref{fol:182}), f\textsuperscript{o}~92~r\textsuperscript{o} (p.~\pageref{fol:184}), f\textsuperscript{o}~92~v\textsuperscript{o} (p.~\pageref{fol:185}).\par\vspace{0.55mm}
\noindent\hangindent=4.5mm\textbf{Creil (Oise)}, nœud ferroviaire entre Amiens et Paris : f\textsuperscript{o}~25~v\textsuperscript{o} (p.~\pageref{fol:51}), f\textsuperscript{o}~32~v\textsuperscript{o} (p.~\pageref{fol:65}).\par\vspace{0.55mm}
\noindent\hangindent=4.5mm\textbf{Domqueur (Somme)}, village du Ponthieu voisin de Ribeaucourt : f\textsuperscript{o}~35~r\textsuperscript{o} (p.~\pageref{fol:70}).\par\vspace{0.55mm}
\noindent\hangindent=4.5mm\textbf{Doullens (Somme)}, sous-préfecture de la Somme, église Notre-Dame, cantine militaire et commerces : f\textsuperscript{o}~25~v\textsuperscript{o} (p.~\pageref{fol:51}), f\textsuperscript{o}~32~v\textsuperscript{o} (p.~\pageref{fol:65}), f\textsuperscript{o}~33~r\textsuperscript{o} (p.~\pageref{fol:66}), f\textsuperscript{o}~33~v\textsuperscript{o} (p.~\pageref{fol:67}), f\textsuperscript{o}~49~r\textsuperscript{o} (p.~\pageref{fol:98}), f\textsuperscript{o}~57~v\textsuperscript{o} (p.~\pageref{fol:115}), f\textsuperscript{o}~61~v\textsuperscript{o} (p.~\pageref{fol:123}), f\textsuperscript{o}~62~r\textsuperscript{o} (p.~\pageref{fol:124}), f\textsuperscript{o}~64~r\textsuperscript{o} (p.~\pageref{fol:128}), f\textsuperscript{o}~73~r\textsuperscript{o} (p.~\pageref{fol:146}), f\textsuperscript{o}~76~r\textsuperscript{o} (p.~\pageref{fol:152}), f\textsuperscript{o}~79~r\textsuperscript{o} (p.~\pageref{fol:158}), f\textsuperscript{o}~81~v\textsuperscript{o} (p.~\pageref{fol:163}), f\textsuperscript{o}~89~v\textsuperscript{o} (p.~\pageref{fol:179}).\par\vspace{0.55mm}
\noindent\hangindent=4.5mm\textbf{Englebelmer (Somme)}, village du front de l'Ancre évacué par ses habitants : f\textsuperscript{o}~50~v\textsuperscript{o} (p.~\pageref{fol:101}), f\textsuperscript{o}~54~r\textsuperscript{o} (p.~\pageref{fol:108}), f\textsuperscript{o}~55~r\textsuperscript{o} (p.~\pageref{fol:110}).\par\vspace{0.55mm}
\noindent\hangindent=4.5mm\textbf{Erquinghem-Lys (Nord)}, commune de la vallée de la Lys en amont d'Armentières : f\textsuperscript{o}~8~v\textsuperscript{o} (p.~\pageref{fol:17}), f\textsuperscript{o}~11~r\textsuperscript{o} (p.~\pageref{fol:22}), f\textsuperscript{o}~11~v\textsuperscript{o} (p.~\pageref{fol:23}), f\textsuperscript{o}~12~r\textsuperscript{o} (p.~\pageref{fol:24}).\par\vspace{0.55mm}
\noindent\hangindent=4.5mm\textbf{Estaires (Nord)}, ville sur la Lys (église Saint-Vaast et hôtel de ville Renaissance détruits en 1918) : f\textsuperscript{o}~8~r\textsuperscript{o} (p.~\pageref{fol:16}), f\textsuperscript{o}~12~v\textsuperscript{o} (p.~\pageref{fol:25}).\par\vspace{0.55mm}
\noindent\hangindent=4.5mm\textbf{Flandre}, plaine flamande et Flandre intérieure traversées en juillet–août 1916 : f\textsuperscript{o}~1~r\textsuperscript{o} (p.~\pageref{fol:2}), f\textsuperscript{o}~26~r\textsuperscript{o} (p.~\pageref{fol:52}).\par\vspace{0.55mm}
\noindent\hangindent=4.5mm\textbf{Fleurbaix (Pas-de-Calais)}, village du front au sud d'Armentières : f\textsuperscript{o}~10~v\textsuperscript{o} (p.~\pageref{fol:21}), f\textsuperscript{o}~11~r\textsuperscript{o} (p.~\pageref{fol:22}).\par\vspace{0.55mm}
\noindent\hangindent=4.5mm\textbf{Forceville (Somme)}, village au sud-est d'Acheux-en-Amiénois, cantonnement de septembre 1916 : f\textsuperscript{o}~22~v\textsuperscript{o} (p.~\pageref{fol:45}), f\textsuperscript{o}~24~r\textsuperscript{o} (p.~\pageref{fol:48}), f\textsuperscript{o}~24~v\textsuperscript{o} (p.~\pageref{fol:49}), f\textsuperscript{o}~25~r\textsuperscript{o} (p.~\pageref{fol:50}), f\textsuperscript{o}~48~r\textsuperscript{o} (p.~\pageref{fol:96}).\par\vspace{0.55mm}
\noindent\hangindent=4.5mm\textbf{Franvillers (Somme)}, commune sur la route d'Albert à Amiens : f\textsuperscript{o}~52~v\textsuperscript{o} (p.~\pageref{fol:105}), f\textsuperscript{o}~58~v\textsuperscript{o} (p.~\pageref{fol:117}).\par\vspace{0.55mm}
\noindent\hangindent=4.5mm\textbf{Fréchencourt (Somme)}, village de la vallée de l'Hallue : f\textsuperscript{o}~78~v\textsuperscript{o} (p.~\pageref{fol:157}), f\textsuperscript{o}~80~v\textsuperscript{o} (p.~\pageref{fol:161}), f\textsuperscript{o}~81~r\textsuperscript{o} (p.~\pageref{fol:162}), f\textsuperscript{o}~89~r\textsuperscript{o} (p.~\pageref{fol:178}), f\textsuperscript{o}~90~r\textsuperscript{o} (p.~\pageref{fol:180}).\par\vspace{0.55mm}
\noindent\hangindent=4.5mm\textbf{Harponville (Somme)}, village entre Toutencourt et Warloy-Baillon : f\textsuperscript{o}~55~r\textsuperscript{o} (p.~\pageref{fol:110}), f\textsuperscript{o}~89~r\textsuperscript{o} (p.~\pageref{fol:178}), f\textsuperscript{o}~92~r\textsuperscript{o} (p.~\pageref{fol:184}).\par\vspace{0.55mm}
\noindent\hangindent=4.5mm\textbf{Hazebrouck (Nord)}, sous-préfecture du Nord, église Saint-Éloi et nœud ferroviaire : f\textsuperscript{o}~2~r\textsuperscript{o} (p.~\pageref{fol:4}), f\textsuperscript{o}~2~v\textsuperscript{o} (p.~\pageref{fol:5}), f\textsuperscript{o}~4~r\textsuperscript{o} (p.~\pageref{fol:8}), f\textsuperscript{o}~8~r\textsuperscript{o} (p.~\pageref{fol:16}), f\textsuperscript{o}~13~v\textsuperscript{o} (p.~\pageref{fol:27}).\par\vspace{0.55mm}
\noindent\hangindent=4.5mm\textbf{Hébuterne (Pas-de-Calais)}, secteur du front au nord d'Acheux et de Serre : f\textsuperscript{o}~53~v\textsuperscript{o} (p.~\pageref{fol:107}), f\textsuperscript{o}~58~r\textsuperscript{o} (p.~\pageref{fol:116}).\par\vspace{0.55mm}
\noindent\hangindent=4.5mm\textbf{Hédauville (Somme)}, village entre Forceville et Warloy-Baillon (bivouac de la 150\textsuperscript{e} Compagnie A.S.C.) : f\textsuperscript{o}~22~v\textsuperscript{o} (p.~\pageref{fol:45}), f\textsuperscript{o}~38~v\textsuperscript{o} (p.~\pageref{fol:77}), f\textsuperscript{o}~54~r\textsuperscript{o} (p.~\pageref{fol:108}), f\textsuperscript{o}~60~v\textsuperscript{o} (p.~\pageref{fol:121}).\par\vspace{0.55mm}
\noindent\hangindent=4.5mm\textbf{Heilly (Somme)}, village de la vallée de l'Ancre où Hippolyte Combes avait cantonné durant l'hiver 1915–1916 : f\textsuperscript{o}~52~v\textsuperscript{o} (p.~\pageref{fol:105}).\par\vspace{0.55mm}
\noindent\hangindent=4.5mm\textbf{Hénencourt (Somme)}, village et château entre Warloy-Baillon et Albert : f\textsuperscript{o}~90~v\textsuperscript{o} (p.~\pageref{fol:181}), f\textsuperscript{o}~93~r\textsuperscript{o} (p.~\pageref{fol:186}).\par\vspace{0.55mm}
\noindent\hangindent=4.5mm\textbf{Hérissart (Somme)}, commune du canton d'Acheux, siège du \textit{Purchase Board} et résidence d'Hippolyte Combes de février à avril 1917 (logements n\textsuperscript{os}~104, 220 et 216) : f\textsuperscript{o}~47~v\textsuperscript{o} (p.~\pageref{fol:95}), f\textsuperscript{o}~49~v\textsuperscript{o} (p.~\pageref{fol:99}), f\textsuperscript{o}~51~v\textsuperscript{o} (p.~\pageref{fol:103}), f\textsuperscript{o}~52~r\textsuperscript{o} (p.~\pageref{fol:104}), f\textsuperscript{o}~53~r\textsuperscript{o} (p.~\pageref{fol:106}), f\textsuperscript{o}~54~r\textsuperscript{o} (p.~\pageref{fol:108}), f\textsuperscript{o}~62~r\textsuperscript{o} (p.~\pageref{fol:124}), f\textsuperscript{o}~73~v\textsuperscript{o} (p.~\pageref{fol:147}), f\textsuperscript{o}~74~r\textsuperscript{o} (p.~\pageref{fol:148}), f\textsuperscript{o}~74~v\textsuperscript{o} (p.~\pageref{fol:149}), f\textsuperscript{o}~75~r\textsuperscript{o} (p.~\pageref{fol:150}), f\textsuperscript{o}~75~v\textsuperscript{o} (p.~\pageref{fol:151}), f\textsuperscript{o}~76~r\textsuperscript{o} (p.~\pageref{fol:152}), f\textsuperscript{o}~76~v\textsuperscript{o} (p.~\pageref{fol:153}), f\textsuperscript{o}~77~r\textsuperscript{o} (p.~\pageref{fol:154}), f\textsuperscript{o}~77~v\textsuperscript{o} (p.~\pageref{fol:155}), f\textsuperscript{o}~78~r\textsuperscript{o} (p.~\pageref{fol:156}), f\textsuperscript{o}~78~v\textsuperscript{o} (p.~\pageref{fol:157}), f\textsuperscript{o}~79~r\textsuperscript{o} (p.~\pageref{fol:158}), f\textsuperscript{o}~79~v\textsuperscript{o} (p.~\pageref{fol:159}), f\textsuperscript{o}~80~r\textsuperscript{o} (p.~\pageref{fol:160}), f\textsuperscript{o}~80~v\textsuperscript{o} (p.~\pageref{fol:161}), f\textsuperscript{o}~81~r\textsuperscript{o} (p.~\pageref{fol:162}), f\textsuperscript{o}~81~v\textsuperscript{o} (p.~\pageref{fol:163}), f\textsuperscript{o}~82~r\textsuperscript{o} (p.~\pageref{fol:164}), f\textsuperscript{o}~82~v\textsuperscript{o} (p.~\pageref{fol:165}), f\textsuperscript{o}~83~r\textsuperscript{o} (p.~\pageref{fol:166}), f\textsuperscript{o}~83~v\textsuperscript{o} (p.~\pageref{fol:167}), f\textsuperscript{o}~84~r\textsuperscript{o} (p.~\pageref{fol:168}), f\textsuperscript{o}~84~v\textsuperscript{o} (p.~\pageref{fol:169}), f\textsuperscript{o}~85~r\textsuperscript{o} (p.~\pageref{fol:170}), f\textsuperscript{o}~85~v\textsuperscript{o} (p.~\pageref{fol:171}), f\textsuperscript{o}~86~r\textsuperscript{o} (p.~\pageref{fol:172}), f\textsuperscript{o}~86~v\textsuperscript{o} (p.~\pageref{fol:173}), f\textsuperscript{o}~87~r\textsuperscript{o} (p.~\pageref{fol:174}), f\textsuperscript{o}~87~v\textsuperscript{o} (p.~\pageref{fol:175}), f\textsuperscript{o}~88~r\textsuperscript{o} (p.~\pageref{fol:176}), f\textsuperscript{o}~88~v\textsuperscript{o} (p.~\pageref{fol:177}), f\textsuperscript{o}~89~r\textsuperscript{o} (p.~\pageref{fol:178}), f\textsuperscript{o}~89~v\textsuperscript{o} (p.~\pageref{fol:179}), f\textsuperscript{o}~90~r\textsuperscript{o} (p.~\pageref{fol:180}), f\textsuperscript{o}~90~v\textsuperscript{o} (p.~\pageref{fol:181}), f\textsuperscript{o}~91~r\textsuperscript{o} (p.~\pageref{fol:182}), f\textsuperscript{o}~91~v\textsuperscript{o} (p.~\pageref{fol:183}), f\textsuperscript{o}~92~r\textsuperscript{o} (p.~\pageref{fol:184}), f\textsuperscript{o}~92~v\textsuperscript{o} (p.~\pageref{fol:185}), f\textsuperscript{o}~93~r\textsuperscript{o} (p.~\pageref{fol:186}).\par\vspace{0.55mm}
\noindent\hangindent=4.5mm\textbf{Jersey (île anglo-normande)}, patrie du sous-lieutenant poète Gilbert : f\textsuperscript{o}~89~v\textsuperscript{o} (p.~\pageref{fol:179}).\par\vspace{0.55mm}
\noindent\hangindent=4.5mm\textbf{La Croix-du-Bac (Steenwerck, Nord)}, hameau de la vallée de la Lys, église Notre-Dame-des-Sept-Douleurs et cantonnement d'août 1916 : f\textsuperscript{o}~5~r\textsuperscript{o} (p.~\pageref{fol:10}), f\textsuperscript{o}~6~r\textsuperscript{o} (p.~\pageref{fol:12}), f\textsuperscript{o}~6~v\textsuperscript{o} (p.~\pageref{fol:13}), f\textsuperscript{o}~7~r\textsuperscript{o} (p.~\pageref{fol:14}).\par\vspace{0.55mm}
\noindent\hangindent=4.5mm\textbf{Lardenne \& Purpan-Ancely (Toulouse, Haute-Garonne)}, faubourg occidental de Toulouse et ruines de l'amphithéâtre gallo-romain : f\textsuperscript{o}~29~r\textsuperscript{o} (p.~\pageref{fol:58}), f\textsuperscript{o}~68~r\textsuperscript{o} (p.~\pageref{fol:136}), f\textsuperscript{o}~68~v\textsuperscript{o} (p.~\pageref{fol:137}).\par\vspace{0.55mm}
\noindent\hangindent=4.5mm\textbf{Laviéville (Somme)}, village entre Warloy-Baillon et Albert : f\textsuperscript{o}~90~v\textsuperscript{o} (p.~\pageref{fol:181}), f\textsuperscript{o}~93~r\textsuperscript{o} (p.~\pageref{fol:186}).\par\vspace{0.55mm}
\noindent\hangindent=4.5mm\textbf{Le Mesnil-Martinsart (Somme)}, commune du canton d'Albert où enseignent M. et M\textsuperscript{me} Fortuné Flutre : f\textsuperscript{o}~52~r\textsuperscript{o} (p.~\pageref{fol:104}), f\textsuperscript{o}~84~r\textsuperscript{o} (p.~\pageref{fol:168}).\par\vspace{0.55mm}
\noindent\hangindent=4.5mm\textbf{Le Quesnoy (ferme, Somme)}, écart agricole entre Raincheval et Puchevillers : f\textsuperscript{o}~55~v\textsuperscript{o} (p.~\pageref{fol:111}), f\textsuperscript{o}~56~r\textsuperscript{o} (p.~\pageref{fol:112}), f\textsuperscript{o}~56~v\textsuperscript{o} (p.~\pageref{fol:113}), f\textsuperscript{o}~57~v\textsuperscript{o} (p.~\pageref{fol:115}), f\textsuperscript{o}~58~r\textsuperscript{o} (p.~\pageref{fol:116}), f\textsuperscript{o}~59~r\textsuperscript{o} (p.~\pageref{fol:118}).\par\vspace{0.55mm}
\noindent\hangindent=4.5mm\textbf{Léalvillers (Somme)}, village immédiatement voisin d'Acheux-en-Amiénois : f\textsuperscript{o}~24~r\textsuperscript{o} (p.~\pageref{fol:48}), f\textsuperscript{o}~47~r\textsuperscript{o} (p.~\pageref{fol:94}), f\textsuperscript{o}~47~v\textsuperscript{o} (p.~\pageref{fol:95}), f\textsuperscript{o}~48~r\textsuperscript{o} (p.~\pageref{fol:96}), f\textsuperscript{o}~48~v\textsuperscript{o} (p.~\pageref{fol:97}), f\textsuperscript{o}~50~v\textsuperscript{o} (p.~\pageref{fol:101}), f\textsuperscript{o}~51~r\textsuperscript{o} (p.~\pageref{fol:102}), f\textsuperscript{o}~56~r\textsuperscript{o} (p.~\pageref{fol:112}), f\textsuperscript{o}~58~r\textsuperscript{o} (p.~\pageref{fol:116}), f\textsuperscript{o}~88~v\textsuperscript{o} (p.~\pageref{fol:177}).\par\vspace{0.55mm}
\noindent\hangindent=4.5mm\textbf{Ligny-Saint-Flochel (Pas-de-Calais)}, village du Ternois à l'église gothique classée (XV\textsuperscript{e}–XVI\textsuperscript{e}~s.) : f\textsuperscript{o}~16~v\textsuperscript{o} (p.~\pageref{fol:33}), f\textsuperscript{o}~19~v\textsuperscript{o} (p.~\pageref{fol:39}).\par\vspace{0.55mm}
\noindent\hangindent=4.5mm\textbf{Louvencourt (Somme)}, commune au nord d'Acheux-en-Amiénois : f\textsuperscript{o}~46~r\textsuperscript{o} (p.~\pageref{fol:92}), f\textsuperscript{o}~49~v\textsuperscript{o} (p.~\pageref{fol:99}), f\textsuperscript{o}~53~r\textsuperscript{o} (p.~\pageref{fol:106}), f\textsuperscript{o}~53~v\textsuperscript{o} (p.~\pageref{fol:107}), f\textsuperscript{o}~56~r\textsuperscript{o} (p.~\pageref{fol:112}), f\textsuperscript{o}~57~r\textsuperscript{o} (p.~\pageref{fol:114}), f\textsuperscript{o}~58~r\textsuperscript{o} (p.~\pageref{fol:116}), f\textsuperscript{o}~59~r\textsuperscript{o} (p.~\pageref{fol:118}), f\textsuperscript{o}~60~v\textsuperscript{o} (p.~\pageref{fol:121}), f\textsuperscript{o}~61~r\textsuperscript{o} (p.~\pageref{fol:122}), f\textsuperscript{o}~61~v\textsuperscript{o} (p.~\pageref{fol:123}), f\textsuperscript{o}~63~r\textsuperscript{o} (p.~\pageref{fol:126}), f\textsuperscript{o}~73~r\textsuperscript{o} (p.~\pageref{fol:146}).\par\vspace{0.55mm}
\noindent\hangindent=4.5mm\textbf{Lucheux (Somme)}, bourg de la vallée de la Grouches (ruines du château des comtes de Saint-Pol et beffroi au chêne de Louis~XI) : f\textsuperscript{o}~20~v\textsuperscript{o} (p.~\pageref{fol:41}).\par\vspace{0.55mm}
\noindent\hangindent=4.5mm\textbf{Lynde (Nord)}, village de Flandre intérieure entre Renescure et Hazebrouck, premier cantonnement du carnet : f\textsuperscript{o}~1~r\textsuperscript{o} (p.~\pageref{fol:2}), f\textsuperscript{o}~1~v\textsuperscript{o} (p.~\pageref{fol:3}), f\textsuperscript{o}~2~r\textsuperscript{o} (p.~\pageref{fol:4}), f\textsuperscript{o}~2~v\textsuperscript{o} (p.~\pageref{fol:5}).\par\vspace{0.55mm}
\noindent\hangindent=4.5mm\textbf{Lys (vallée de la)}, rivière franco-belge canalisée entre Aire, Estaires, Sailly, Bac-Saint-Maur et Armentières : f\textsuperscript{o}~6~v\textsuperscript{o} (p.~\pageref{fol:13}), f\textsuperscript{o}~7~v\textsuperscript{o} (p.~\pageref{fol:15}), f\textsuperscript{o}~8~r\textsuperscript{o} (p.~\pageref{fol:16}), f\textsuperscript{o}~8~v\textsuperscript{o} (p.~\pageref{fol:17}), f\textsuperscript{o}~9~r\textsuperscript{o} (p.~\pageref{fol:18}), f\textsuperscript{o}~9~v\textsuperscript{o} (p.~\pageref{fol:19}), f\textsuperscript{o}~11~r\textsuperscript{o} (p.~\pageref{fol:22}), f\textsuperscript{o}~11~v\textsuperscript{o} (p.~\pageref{fol:23}), f\textsuperscript{o}~12~r\textsuperscript{o} (p.~\pageref{fol:24}), f\textsuperscript{o}~12~v\textsuperscript{o} (p.~\pageref{fol:25}).\par\vspace{0.55mm}
\noindent\hangindent=4.5mm\textbf{Marquay (Pas-de-Calais)}, commune du Ternois, cantonnement de la compagnie du 25 août au 9 septembre 1916 : f\textsuperscript{o}~14~r\textsuperscript{o} (p.~\pageref{fol:28}), f\textsuperscript{o}~14~v\textsuperscript{o} (p.~\pageref{fol:29}), f\textsuperscript{o}~15~r\textsuperscript{o} (p.~\pageref{fol:30}), f\textsuperscript{o}~15~v\textsuperscript{o} (p.~\pageref{fol:31}), f\textsuperscript{o}~16~r\textsuperscript{o} (p.~\pageref{fol:32}), f\textsuperscript{o}~16~v\textsuperscript{o} (p.~\pageref{fol:33}), f\textsuperscript{o}~17~r\textsuperscript{o} (p.~\pageref{fol:34}), f\textsuperscript{o}~17~v\textsuperscript{o} (p.~\pageref{fol:35}), f\textsuperscript{o}~18~r\textsuperscript{o} (p.~\pageref{fol:36}), f\textsuperscript{o}~18~v\textsuperscript{o} (p.~\pageref{fol:37}), f\textsuperscript{o}~19~r\textsuperscript{o} (p.~\pageref{fol:38}), f\textsuperscript{o}~19~v\textsuperscript{o} (p.~\pageref{fol:39}).\par\vspace{0.55mm}
\noindent\hangindent=4.5mm\textbf{Martinsart (Somme)}, hameau et château de la commune du Mesnil-Martinsart près d'Albert : f\textsuperscript{o}~52~r\textsuperscript{o} (p.~\pageref{fol:104}), f\textsuperscript{o}~52~v\textsuperscript{o} (p.~\pageref{fol:105}), f\textsuperscript{o}~58~v\textsuperscript{o} (p.~\pageref{fol:117}), f\textsuperscript{o}~75~r\textsuperscript{o} (p.~\pageref{fol:150}), f\textsuperscript{o}~79~v\textsuperscript{o} (p.~\pageref{fol:159}), f\textsuperscript{o}~84~r\textsuperscript{o} (p.~\pageref{fol:168}).\par\vspace{0.55mm}
\noindent\hangindent=4.5mm\textbf{Millau (Aveyron)}, ville natale d'Hippolyte Combes et de Marguerite Calmels (10, rue Alfred-Guibert ; rue des Barris ; cimetière) : Couverture (p.~\pageref{fol:1}), f\textsuperscript{o}~69~r\textsuperscript{o} (p.~\pageref{fol:138}), f\textsuperscript{o}~69~v\textsuperscript{o} (p.~\pageref{fol:139}), f\textsuperscript{o}~70~r\textsuperscript{o} (p.~\pageref{fol:140}), f\textsuperscript{o}~70~v\textsuperscript{o} (p.~\pageref{fol:141}), f\textsuperscript{o}~71~r\textsuperscript{o} (p.~\pageref{fol:142}), f\textsuperscript{o}~71~v\textsuperscript{o} (p.~\pageref{fol:143}), f\textsuperscript{o}~72~r\textsuperscript{o} (p.~\pageref{fol:144}).\par\vspace{0.55mm}
\noindent\hangindent=4.5mm\textbf{Miraumont (Somme)}, village aux sources de l'Ancre évacué par l'armée allemande en février–mars 1917 : f\textsuperscript{o}~61~r\textsuperscript{o} (p.~\pageref{fol:122}), f\textsuperscript{o}~85~r\textsuperscript{o} (p.~\pageref{fol:170}).\par\vspace{0.55mm}
\noindent\hangindent=4.5mm\textbf{Mirvaux (Somme)}, village voisin de Pierregot et Rubempré : f\textsuperscript{o}~79~v\textsuperscript{o} (p.~\pageref{fol:159}), f\textsuperscript{o}~84~v\textsuperscript{o} (p.~\pageref{fol:169}), f\textsuperscript{o}~85~v\textsuperscript{o} (p.~\pageref{fol:171}), f\textsuperscript{o}~91~v\textsuperscript{o} (p.~\pageref{fol:183}), f\textsuperscript{o}~92~v\textsuperscript{o} (p.~\pageref{fol:185}).\par\vspace{0.55mm}
\noindent\hangindent=4.5mm\textbf{Molliens-au-Bois (Somme)}, commune du canton de Villers-Bocage : f\textsuperscript{o}~79~r\textsuperscript{o} (p.~\pageref{fol:158}), f\textsuperscript{o}~83~v\textsuperscript{o} (p.~\pageref{fol:167}), f\textsuperscript{o}~85~v\textsuperscript{o} (p.~\pageref{fol:171}), f\textsuperscript{o}~87~r\textsuperscript{o} (p.~\pageref{fol:174}), f\textsuperscript{o}~91~r\textsuperscript{o} (p.~\pageref{fol:182}).\par\vspace{0.55mm}
\noindent\hangindent=4.5mm\textbf{Mont-Noir (Nord)}, colline des Monts des Flandres au nord de Saint-Jans-Cappel : f\textsuperscript{o}~3~r\textsuperscript{o} (p.~\pageref{fol:6}).\par\vspace{0.55mm}
\noindent\hangindent=4.5mm\textbf{Montauban (Tarn-et-Garonne)}, ville traversée en train lors des permissions et futur poste d'Hippolyte Combes au lycée Ingres : f\textsuperscript{o}~27~r\textsuperscript{o} (p.~\pageref{fol:54}).\par\vspace{0.55mm}
\noindent\hangindent=4.5mm\textbf{Montigny-sur-l'Hallue (Somme)}, village voisin de Beaucourt-sur-l'Hallue : f\textsuperscript{o}~59~v\textsuperscript{o} (p.~\pageref{fol:119}), f\textsuperscript{o}~61~v\textsuperscript{o} (p.~\pageref{fol:123}), f\textsuperscript{o}~79~r\textsuperscript{o} (p.~\pageref{fol:158}), f\textsuperscript{o}~80~v\textsuperscript{o} (p.~\pageref{fol:161}), f\textsuperscript{o}~81~r\textsuperscript{o} (p.~\pageref{fol:162}), f\textsuperscript{o}~91~v\textsuperscript{o} (p.~\pageref{fol:183}).\par\vspace{0.55mm}
\noindent\hangindent=4.5mm\textbf{Neuilly-sur-Seine (Hauts-de-Seine)}, résidence de la famille Faure (16, rue d'Armenonville) visitée lors des permissions : f\textsuperscript{o}~26~r\textsuperscript{o} (p.~\pageref{fol:52}), f\textsuperscript{o}~65~r\textsuperscript{o} (p.~\pageref{fol:130}), f\textsuperscript{o}~65~v\textsuperscript{o} (p.~\pageref{fol:131}).\par\vspace{0.55mm}
\noindent\hangindent=4.5mm\textbf{Nice (Alpes-Maritimes)}, ville d'où provient le panier de fleurs commandé chez M\textsuperscript{me} Spagnol : f\textsuperscript{o}~67~r\textsuperscript{o} (p.~\pageref{fol:134}), f\textsuperscript{o}~68~r\textsuperscript{o} (p.~\pageref{fol:136}), f\textsuperscript{o}~79~r\textsuperscript{o} (p.~\pageref{fol:158}).\par\vspace{0.55mm}
\noindent\hangindent=4.5mm\textbf{Nogent-le-Rotrou (Eure-et-Loir)}, ville du Perche où Hippolyte Combes enseignait au collège communal (17, rue Saint-Hilaire) : f\textsuperscript{o}~11~v\textsuperscript{o} (p.~\pageref{fol:23}), f\textsuperscript{o}~12~r\textsuperscript{o} (p.~\pageref{fol:24}).\par\vspace{0.55mm}
\noindent\hangindent=4.5mm\textbf{Ostreville (Pas-de-Calais)}, village du Ternois entre Marquay et Saint-Pol : f\textsuperscript{o}~16~r\textsuperscript{o} (p.~\pageref{fol:32}).\par\vspace{0.55mm}
\noindent\hangindent=4.5mm\textbf{Paris}, capitale traversée aux permissions de septembre–octobre 1916 et janvier–février 1917 (gare du Nord, gare d'Orsay, boulevard Saint-Germain, boulevard Saint-Michel, rue Saint-Jacques, Panthéon, église Saint-Vincent-de-Paul) : Couverture (p.~\pageref{fol:1}), f\textsuperscript{o}~25~v\textsuperscript{o} (p.~\pageref{fol:51}), f\textsuperscript{o}~26~r\textsuperscript{o} (p.~\pageref{fol:52}), f\textsuperscript{o}~26~v\textsuperscript{o} (p.~\pageref{fol:53}), f\textsuperscript{o}~31~v\textsuperscript{o} (p.~\pageref{fol:63}), f\textsuperscript{o}~32~r\textsuperscript{o} (p.~\pageref{fol:64}), f\textsuperscript{o}~32~v\textsuperscript{o} (p.~\pageref{fol:65}), f\textsuperscript{o}~49~v\textsuperscript{o} (p.~\pageref{fol:99}), f\textsuperscript{o}~58~r\textsuperscript{o} (p.~\pageref{fol:116}), f\textsuperscript{o}~61~r\textsuperscript{o} (p.~\pageref{fol:122}), f\textsuperscript{o}~65~r\textsuperscript{o} (p.~\pageref{fol:130}), f\textsuperscript{o}~65~v\textsuperscript{o} (p.~\pageref{fol:131}), f\textsuperscript{o}~66~r\textsuperscript{o} (p.~\pageref{fol:132}), f\textsuperscript{o}~72~v\textsuperscript{o} (p.~\pageref{fol:145}), f\textsuperscript{o}~73~r\textsuperscript{o} (p.~\pageref{fol:146}), f\textsuperscript{o}~74~r\textsuperscript{o} (p.~\pageref{fol:148}).\par\vspace{0.55mm}
\noindent\hangindent=4.5mm\textbf{Pech-David \& Pouvourville (Toulouse, Haute-Garonne)}, coteaux dominant le sud de Toulouse et la Garonne : f\textsuperscript{o}~27~v\textsuperscript{o} (p.~\pageref{fol:55}).\par\vspace{0.55mm}
\noindent\hangindent=4.5mm\textbf{Pierregot (Somme)}, village à l'ouest de Rubempré : f\textsuperscript{o}~52~r\textsuperscript{o} (p.~\pageref{fol:104}), f\textsuperscript{o}~78~v\textsuperscript{o} (p.~\pageref{fol:157}), f\textsuperscript{o}~84~v\textsuperscript{o} (p.~\pageref{fol:169}), f\textsuperscript{o}~85~v\textsuperscript{o} (p.~\pageref{fol:171}), f\textsuperscript{o}~87~r\textsuperscript{o} (p.~\pageref{fol:174}), f\textsuperscript{o}~87~v\textsuperscript{o} (p.~\pageref{fol:175}), f\textsuperscript{o}~89~r\textsuperscript{o} (p.~\pageref{fol:178}), f\textsuperscript{o}~90~r\textsuperscript{o} (p.~\pageref{fol:180}), f\textsuperscript{o}~91~r\textsuperscript{o} (p.~\pageref{fol:182}), f\textsuperscript{o}~91~v\textsuperscript{o} (p.~\pageref{fol:183}), f\textsuperscript{o}~92~r\textsuperscript{o} (p.~\pageref{fol:184}), f\textsuperscript{o}~92~v\textsuperscript{o} (p.~\pageref{fol:185}), f\textsuperscript{o}~93~r\textsuperscript{o} (p.~\pageref{fol:186}).\par\vspace{0.55mm}
\noindent\hangindent=4.5mm\textbf{Puchevillers (Somme)}, commune entre Raincheval et Hérissart, centre d'hôpitaux d'évacuation britanniques (\textit{Casualty Clearing Stations}) : f\textsuperscript{o}~37~v\textsuperscript{o} (p.~\pageref{fol:75}), f\textsuperscript{o}~52~r\textsuperscript{o} (p.~\pageref{fol:104}), f\textsuperscript{o}~55~r\textsuperscript{o} (p.~\pageref{fol:110}), f\textsuperscript{o}~56~r\textsuperscript{o} (p.~\pageref{fol:112}), f\textsuperscript{o}~62~v\textsuperscript{o} (p.~\pageref{fol:125}), f\textsuperscript{o}~73~v\textsuperscript{o} (p.~\pageref{fol:147}), f\textsuperscript{o}~77~r\textsuperscript{o} (p.~\pageref{fol:154}), f\textsuperscript{o}~77~v\textsuperscript{o} (p.~\pageref{fol:155}), f\textsuperscript{o}~79~r\textsuperscript{o} (p.~\pageref{fol:158}), f\textsuperscript{o}~81~r\textsuperscript{o} (p.~\pageref{fol:162}), f\textsuperscript{o}~83~r\textsuperscript{o} (p.~\pageref{fol:166}), f\textsuperscript{o}~83~v\textsuperscript{o} (p.~\pageref{fol:167}), f\textsuperscript{o}~84~r\textsuperscript{o} (p.~\pageref{fol:168}), f\textsuperscript{o}~84~v\textsuperscript{o} (p.~\pageref{fol:169}), f\textsuperscript{o}~85~r\textsuperscript{o} (p.~\pageref{fol:170}), f\textsuperscript{o}~85~v\textsuperscript{o} (p.~\pageref{fol:171}), f\textsuperscript{o}~86~r\textsuperscript{o} (p.~\pageref{fol:172}), f\textsuperscript{o}~86~v\textsuperscript{o} (p.~\pageref{fol:173}), f\textsuperscript{o}~87~r\textsuperscript{o} (p.~\pageref{fol:174}), f\textsuperscript{o}~91~v\textsuperscript{o} (p.~\pageref{fol:183}), f\textsuperscript{o}~92~r\textsuperscript{o} (p.~\pageref{fol:184}), f\textsuperscript{o}~93~r\textsuperscript{o} (p.~\pageref{fol:186}).\par\vspace{0.55mm}
\noindent\hangindent=4.5mm\textbf{Raincheval (Somme)}, commune et château au nord d'Hérissart, cantonnement de septembre 1916 et siège de l'\textit{Army Purchase Board} : f\textsuperscript{o}~20~v\textsuperscript{o} (p.~\pageref{fol:41}), f\textsuperscript{o}~21~r\textsuperscript{o} (p.~\pageref{fol:42}), f\textsuperscript{o}~21~v\textsuperscript{o} (p.~\pageref{fol:43}), f\textsuperscript{o}~22~r\textsuperscript{o} (p.~\pageref{fol:44}), f\textsuperscript{o}~22~v\textsuperscript{o} (p.~\pageref{fol:45}), f\textsuperscript{o}~23~r\textsuperscript{o} (p.~\pageref{fol:46}), f\textsuperscript{o}~23~v\textsuperscript{o} (p.~\pageref{fol:47}), f\textsuperscript{o}~24~r\textsuperscript{o} (p.~\pageref{fol:48}), f\textsuperscript{o}~45~v\textsuperscript{o} (p.~\pageref{fol:91}), f\textsuperscript{o}~46~v\textsuperscript{o} (p.~\pageref{fol:93}), f\textsuperscript{o}~48~v\textsuperscript{o} (p.~\pageref{fol:97}), f\textsuperscript{o}~49~r\textsuperscript{o} (p.~\pageref{fol:98}), f\textsuperscript{o}~49~v\textsuperscript{o} (p.~\pageref{fol:99}), f\textsuperscript{o}~52~r\textsuperscript{o} (p.~\pageref{fol:104}), f\textsuperscript{o}~53~v\textsuperscript{o} (p.~\pageref{fol:107}), f\textsuperscript{o}~54~r\textsuperscript{o} (p.~\pageref{fol:108}), f\textsuperscript{o}~55~v\textsuperscript{o} (p.~\pageref{fol:111}), f\textsuperscript{o}~56~v\textsuperscript{o} (p.~\pageref{fol:113}), f\textsuperscript{o}~57~v\textsuperscript{o} (p.~\pageref{fol:115}), f\textsuperscript{o}~58~r\textsuperscript{o} (p.~\pageref{fol:116}), f\textsuperscript{o}~58~v\textsuperscript{o} (p.~\pageref{fol:117}), f\textsuperscript{o}~59~r\textsuperscript{o} (p.~\pageref{fol:118}), f\textsuperscript{o}~60~v\textsuperscript{o} (p.~\pageref{fol:121}), f\textsuperscript{o}~61~r\textsuperscript{o} (p.~\pageref{fol:122}), f\textsuperscript{o}~61~v\textsuperscript{o} (p.~\pageref{fol:123}), f\textsuperscript{o}~62~v\textsuperscript{o} (p.~\pageref{fol:125}), f\textsuperscript{o}~73~v\textsuperscript{o} (p.~\pageref{fol:147}), f\textsuperscript{o}~77~v\textsuperscript{o} (p.~\pageref{fol:155}), f\textsuperscript{o}~83~r\textsuperscript{o} (p.~\pageref{fol:166}), f\textsuperscript{o}~83~v\textsuperscript{o} (p.~\pageref{fol:167}), f\textsuperscript{o}~87~v\textsuperscript{o} (p.~\pageref{fol:175}).\par\vspace{0.55mm}
\noindent\hangindent=4.5mm\textbf{Ramecourt (Pas-de-Calais)}, commune et plateau à l'ouest de Saint-Pol-sur-Ternoise : f\textsuperscript{o}~16~r\textsuperscript{o} (p.~\pageref{fol:32}).\par\vspace{0.55mm}
\noindent\hangindent=4.5mm\textbf{Renescure (Nord)}, bourg de Flandre intérieure (château de Zuthove, château dit de Philippe de Comines et église Notre-Dame-de-l'Assomption) : f\textsuperscript{o}~2~r\textsuperscript{o} (p.~\pageref{fol:4}).\par\vspace{0.55mm}
\noindent\hangindent=4.5mm\textbf{Ribeaucourt (Somme)}, commune et château du Ponthieu où la compagnie prend son repos du 6 au 15 octobre 1916 : f\textsuperscript{o}~33~v\textsuperscript{o} (p.~\pageref{fol:67}), f\textsuperscript{o}~34~r\textsuperscript{o} (p.~\pageref{fol:68}), f\textsuperscript{o}~34~v\textsuperscript{o} (p.~\pageref{fol:69}), f\textsuperscript{o}~35~r\textsuperscript{o} (p.~\pageref{fol:70}), f\textsuperscript{o}~35~v\textsuperscript{o} (p.~\pageref{fol:71}), f\textsuperscript{o}~36~r\textsuperscript{o} (p.~\pageref{fol:72}), f\textsuperscript{o}~36~v\textsuperscript{o} (p.~\pageref{fol:73}), f\textsuperscript{o}~37~r\textsuperscript{o} (p.~\pageref{fol:74}), f\textsuperscript{o}~37~v\textsuperscript{o} (p.~\pageref{fol:75}).\par\vspace{0.55mm}
\noindent\hangindent=4.5mm\textbf{Ribemont-sur-Ancre (Somme)}, village de la vallée de l'Ancre en aval d'Albert : f\textsuperscript{o}~52~v\textsuperscript{o} (p.~\pageref{fol:105}), f\textsuperscript{o}~76~v\textsuperscript{o} (p.~\pageref{fol:153}), f\textsuperscript{o}~84~v\textsuperscript{o} (p.~\pageref{fol:169}), f\textsuperscript{o}~90~v\textsuperscript{o} (p.~\pageref{fol:181}).\par\vspace{0.55mm}
\noindent\hangindent=4.5mm\textbf{Roëllecourt (Pas-de-Calais)}, village du Ternois entre Saint-Pol et Marquay : f\textsuperscript{o}~17~v\textsuperscript{o} (p.~\pageref{fol:35}).\par\vspace{0.55mm}
\noindent\hangindent=4.5mm\textbf{Roumanie}, royaume entré en guerre aux côtés des Alliés le 27 août 1916 : f\textsuperscript{o}~15~r\textsuperscript{o} (p.~\pageref{fol:30}), f\textsuperscript{o}~26~v\textsuperscript{o} (p.~\pageref{fol:53}).\par\vspace{0.55mm}
\noindent\hangindent=4.5mm\textbf{Rubempré \& Septenville (Somme)}, chef-lieu de commune et hameau voisins d'Hérissart : f\textsuperscript{o}~47~v\textsuperscript{o} (p.~\pageref{fol:95}), f\textsuperscript{o}~52~r\textsuperscript{o} (p.~\pageref{fol:104}), f\textsuperscript{o}~53~r\textsuperscript{o} (p.~\pageref{fol:106}), f\textsuperscript{o}~56~r\textsuperscript{o} (p.~\pageref{fol:112}), f\textsuperscript{o}~60~r\textsuperscript{o} (p.~\pageref{fol:120}), f\textsuperscript{o}~61~r\textsuperscript{o} (p.~\pageref{fol:122}), f\textsuperscript{o}~62~r\textsuperscript{o} (p.~\pageref{fol:124}), f\textsuperscript{o}~74~v\textsuperscript{o} (p.~\pageref{fol:149}), f\textsuperscript{o}~80~v\textsuperscript{o} (p.~\pageref{fol:161}), f\textsuperscript{o}~81~r\textsuperscript{o} (p.~\pageref{fol:162}), f\textsuperscript{o}~86~r\textsuperscript{o} (p.~\pageref{fol:172}), f\textsuperscript{o}~86~v\textsuperscript{o} (p.~\pageref{fol:173}), f\textsuperscript{o}~87~r\textsuperscript{o} (p.~\pageref{fol:174}), f\textsuperscript{o}~87~v\textsuperscript{o} (p.~\pageref{fol:175}), f\textsuperscript{o}~88~v\textsuperscript{o} (p.~\pageref{fol:177}), f\textsuperscript{o}~89~r\textsuperscript{o} (p.~\pageref{fol:178}), f\textsuperscript{o}~90~r\textsuperscript{o} (p.~\pageref{fol:180}), f\textsuperscript{o}~90~v\textsuperscript{o} (p.~\pageref{fol:181}), f\textsuperscript{o}~92~r\textsuperscript{o} (p.~\pageref{fol:184}).\par\vspace{0.55mm}
\noindent\hangindent=4.5mm\textbf{Sailly-sur-la-Lys (Pas-de-Calais)}, commune de la vallée de la Lys (pont, écluse et ruines de l'église Saint-Vaast) : f\textsuperscript{o}~8~r\textsuperscript{o} (p.~\pageref{fol:16}).\par\vspace{0.55mm}
\noindent\hangindent=4.5mm\textbf{Saint-Chély-d'Apcher (Lozère)}, ville de Lozère traversée sur la ligne des Causses (et où Hippolyte Combes décédera à l'Hôtel du Centre le 2 septembre 1970) : Couverture (p.~\pageref{fol:1}), f\textsuperscript{o}~72~v\textsuperscript{o} (p.~\pageref{fol:145}).\par\vspace{0.55mm}
\noindent\hangindent=4.5mm\textbf{Saint-Flour (Cantal)}, sous-préfecture du Cantal traversée sous la neige le 6 février 1917 : f\textsuperscript{o}~71~v\textsuperscript{o} (p.~\pageref{fol:143}), f\textsuperscript{o}~72~r\textsuperscript{o} (p.~\pageref{fol:144}), f\textsuperscript{o}~72~v\textsuperscript{o} (p.~\pageref{fol:145}).\par\vspace{0.55mm}
\noindent\hangindent=4.5mm\textbf{Saint-Jans-Cappel (Nord)}, village au pied du Mont-Noir où cantonne Hippolyte Combes du 28 juillet au 5 août 1916 : f\textsuperscript{o}~2~v\textsuperscript{o} (p.~\pageref{fol:5}), f\textsuperscript{o}~3~r\textsuperscript{o} (p.~\pageref{fol:6}), f\textsuperscript{o}~3~v\textsuperscript{o} (p.~\pageref{fol:7}), f\textsuperscript{o}~4~r\textsuperscript{o} (p.~\pageref{fol:8}), f\textsuperscript{o}~5~r\textsuperscript{o} (p.~\pageref{fol:10}), f\textsuperscript{o}~5~v\textsuperscript{o} (p.~\pageref{fol:11}), f\textsuperscript{o}~6~r\textsuperscript{o} (p.~\pageref{fol:12}).\par\vspace{0.55mm}
\noindent\hangindent=4.5mm\textbf{Saint-Léger-lès-Authie (Somme)}, village de la vallée de l'Authie : f\textsuperscript{o}~54~r\textsuperscript{o} (p.~\pageref{fol:108}), f\textsuperscript{o}~54~v\textsuperscript{o} (p.~\pageref{fol:109}), f\textsuperscript{o}~56~r\textsuperscript{o} (p.~\pageref{fol:112}), f\textsuperscript{o}~60~v\textsuperscript{o} (p.~\pageref{fol:121}), f\textsuperscript{o}~61~r\textsuperscript{o} (p.~\pageref{fol:122}), f\textsuperscript{o}~61~v\textsuperscript{o} (p.~\pageref{fol:123}).\par\vspace{0.55mm}
\noindent\hangindent=4.5mm\textbf{Saint-Michel-sur-Ternoise (Pas-de-Calais)}, village et église ancienne à l'est de Saint-Pol-sur-Ternoise : f\textsuperscript{o}~16~r\textsuperscript{o} (p.~\pageref{fol:32}), f\textsuperscript{o}~18~v\textsuperscript{o} (p.~\pageref{fol:37}), f\textsuperscript{o}~65~v\textsuperscript{o} (p.~\pageref{fol:131}), f\textsuperscript{o}~66~r\textsuperscript{o} (p.~\pageref{fol:132}).\par\vspace{0.55mm}
\noindent\hangindent=4.5mm\textbf{Saint-Pol-sur-Ternoise (Pas-de-Calais)}, sous-préfecture du Pas-de-Calais, gare de débarquement et église des Carmes (1697) : f\textsuperscript{o}~13~v\textsuperscript{o} (p.~\pageref{fol:27}), f\textsuperscript{o}~14~r\textsuperscript{o} (p.~\pageref{fol:28}), f\textsuperscript{o}~16~r\textsuperscript{o} (p.~\pageref{fol:32}), f\textsuperscript{o}~17~r\textsuperscript{o} (p.~\pageref{fol:34}), f\textsuperscript{o}~18~r\textsuperscript{o} (p.~\pageref{fol:36}), f\textsuperscript{o}~18~v\textsuperscript{o} (p.~\pageref{fol:37}).\par\vspace{0.55mm}
\noindent\hangindent=4.5mm\textbf{Senlis-le-Sec (Somme)}, village entre Warloy-Baillon et Bouzincourt : f\textsuperscript{o}~37~v\textsuperscript{o} (p.~\pageref{fol:75}), f\textsuperscript{o}~38~v\textsuperscript{o} (p.~\pageref{fol:77}), f\textsuperscript{o}~43~v\textsuperscript{o} (p.~\pageref{fol:87}), f\textsuperscript{o}~48~r\textsuperscript{o} (p.~\pageref{fol:96}), f\textsuperscript{o}~55~r\textsuperscript{o} (p.~\pageref{fol:110}), f\textsuperscript{o}~56~r\textsuperscript{o} (p.~\pageref{fol:112}), f\textsuperscript{o}~58~v\textsuperscript{o} (p.~\pageref{fol:117}), f\textsuperscript{o}~62~r\textsuperscript{o} (p.~\pageref{fol:124}), f\textsuperscript{o}~75~v\textsuperscript{o} (p.~\pageref{fol:151}), f\textsuperscript{o}~79~v\textsuperscript{o} (p.~\pageref{fol:159}), f\textsuperscript{o}~84~v\textsuperscript{o} (p.~\pageref{fol:169}), f\textsuperscript{o}~85~v\textsuperscript{o} (p.~\pageref{fol:171}), f\textsuperscript{o}~90~v\textsuperscript{o} (p.~\pageref{fol:181}).\par\vspace{0.55mm}
\noindent\hangindent=4.5mm\textbf{Sercus (Nord)}, village de Flandre intérieure et église Saint-Érasme visités le 25 juillet 1916 : f\textsuperscript{o}~1~r\textsuperscript{o} (p.~\pageref{fol:2}).\par\vspace{0.55mm}
\noindent\hangindent=4.5mm\textbf{Steenwerck (Nord)}, bourg de la plaine de la Lys, gare et église Saint-Jean-Baptiste : f\textsuperscript{o}~5~r\textsuperscript{o} (p.~\pageref{fol:10}), f\textsuperscript{o}~6~r\textsuperscript{o} (p.~\pageref{fol:12}), f\textsuperscript{o}~7~r\textsuperscript{o} (p.~\pageref{fol:14}), f\textsuperscript{o}~7~v\textsuperscript{o} (p.~\pageref{fol:15}), f\textsuperscript{o}~8~v\textsuperscript{o} (p.~\pageref{fol:17}), f\textsuperscript{o}~10~r\textsuperscript{o} (p.~\pageref{fol:20}), f\textsuperscript{o}~10~v\textsuperscript{o} (p.~\pageref{fol:21}), f\textsuperscript{o}~11~v\textsuperscript{o} (p.~\pageref{fol:23}), f\textsuperscript{o}~12~r\textsuperscript{o} (p.~\pageref{fol:24}), f\textsuperscript{o}~13~r\textsuperscript{o} (p.~\pageref{fol:26}), f\textsuperscript{o}~13~v\textsuperscript{o} (p.~\pageref{fol:27}).\par\vspace{0.55mm}
\noindent\hangindent=4.5mm\textbf{Talmas (Somme)}, commune sur la route d'Amiens à Doullens : f\textsuperscript{o}~81~r\textsuperscript{o} (p.~\pageref{fol:162}).\par\vspace{0.55mm}
\noindent\hangindent=4.5mm\textbf{Thièvres (Somme / Pas-de-Calais)}, village de la vallée de l'Authie entre Lucheux et Raincheval : f\textsuperscript{o}~20~v\textsuperscript{o} (p.~\pageref{fol:41}).\par\vspace{0.55mm}
\noindent\hangindent=4.5mm\textbf{Toulouse (Haute-Garonne)}, métropole languedocienne où séjournent Hippolyte Combes et Marguerite Calmels lors des deux permissions (rue du Taur, rue Pharaon, cathédrale Saint-Étienne, basilique Saint-Sernin, Daurade, Dalbade, Hôtel d'Assézat, Faculté des lettres, Sainte-Marie-des-Champs, Capitole, place Esquirol, marché Victor-Hugo, canal du Midi) : Couverture (p.~\pageref{fol:1}), f\textsuperscript{o}~23~v\textsuperscript{o} (p.~\pageref{fol:47}), f\textsuperscript{o}~26~r\textsuperscript{o} (p.~\pageref{fol:52}), f\textsuperscript{o}~26~v\textsuperscript{o} (p.~\pageref{fol:53}), f\textsuperscript{o}~27~r\textsuperscript{o} (p.~\pageref{fol:54}), f\textsuperscript{o}~27~v\textsuperscript{o} (p.~\pageref{fol:55}), f\textsuperscript{o}~28~r\textsuperscript{o} (p.~\pageref{fol:56}), f\textsuperscript{o}~28~v\textsuperscript{o} (p.~\pageref{fol:57}), f\textsuperscript{o}~29~r\textsuperscript{o} (p.~\pageref{fol:58}), f\textsuperscript{o}~29~v\textsuperscript{o} (p.~\pageref{fol:59}), f\textsuperscript{o}~30~r\textsuperscript{o} (p.~\pageref{fol:60}), f\textsuperscript{o}~30~v\textsuperscript{o} (p.~\pageref{fol:61}), f\textsuperscript{o}~31~r\textsuperscript{o} (p.~\pageref{fol:62}), f\textsuperscript{o}~32~r\textsuperscript{o} (p.~\pageref{fol:64}), f\textsuperscript{o}~65~v\textsuperscript{o} (p.~\pageref{fol:131}), f\textsuperscript{o}~66~r\textsuperscript{o} (p.~\pageref{fol:132}), f\textsuperscript{o}~66~v\textsuperscript{o} (p.~\pageref{fol:133}), f\textsuperscript{o}~67~r\textsuperscript{o} (p.~\pageref{fol:134}), f\textsuperscript{o}~67~v\textsuperscript{o} (p.~\pageref{fol:135}), f\textsuperscript{o}~68~r\textsuperscript{o} (p.~\pageref{fol:136}), f\textsuperscript{o}~68~v\textsuperscript{o} (p.~\pageref{fol:137}), f\textsuperscript{o}~69~r\textsuperscript{o} (p.~\pageref{fol:138}), f\textsuperscript{o}~69~v\textsuperscript{o} (p.~\pageref{fol:139}), f\textsuperscript{o}~70~r\textsuperscript{o} (p.~\pageref{fol:140}), f\textsuperscript{o}~78~r\textsuperscript{o} (p.~\pageref{fol:156}).\par\vspace{0.55mm}
\noindent\hangindent=4.5mm\textbf{Toutencourt (Somme)}, commune entre Raincheval et Contay, siège de la Mission militaire française auprès de la V\textsuperscript{e} Armée britannique : f\textsuperscript{o}~49~v\textsuperscript{o} (p.~\pageref{fol:99}), f\textsuperscript{o}~52~v\textsuperscript{o} (p.~\pageref{fol:105}), f\textsuperscript{o}~55~r\textsuperscript{o} (p.~\pageref{fol:110}), f\textsuperscript{o}~55~v\textsuperscript{o} (p.~\pageref{fol:111}), f\textsuperscript{o}~56~r\textsuperscript{o} (p.~\pageref{fol:112}), f\textsuperscript{o}~59~r\textsuperscript{o} (p.~\pageref{fol:118}), f\textsuperscript{o}~62~r\textsuperscript{o} (p.~\pageref{fol:124}), f\textsuperscript{o}~75~r\textsuperscript{o} (p.~\pageref{fol:150}), f\textsuperscript{o}~76~r\textsuperscript{o} (p.~\pageref{fol:152}), f\textsuperscript{o}~77~r\textsuperscript{o} (p.~\pageref{fol:154}), f\textsuperscript{o}~77~v\textsuperscript{o} (p.~\pageref{fol:155}), f\textsuperscript{o}~78~r\textsuperscript{o} (p.~\pageref{fol:156}), f\textsuperscript{o}~79~v\textsuperscript{o} (p.~\pageref{fol:159}), f\textsuperscript{o}~80~r\textsuperscript{o} (p.~\pageref{fol:160}), f\textsuperscript{o}~80~v\textsuperscript{o} (p.~\pageref{fol:161}), f\textsuperscript{o}~81~r\textsuperscript{o} (p.~\pageref{fol:162}), f\textsuperscript{o}~81~v\textsuperscript{o} (p.~\pageref{fol:163}), f\textsuperscript{o}~82~r\textsuperscript{o} (p.~\pageref{fol:164}), f\textsuperscript{o}~83~r\textsuperscript{o} (p.~\pageref{fol:166}), f\textsuperscript{o}~83~v\textsuperscript{o} (p.~\pageref{fol:167}), f\textsuperscript{o}~84~v\textsuperscript{o} (p.~\pageref{fol:169}), f\textsuperscript{o}~85~r\textsuperscript{o} (p.~\pageref{fol:170}), f\textsuperscript{o}~87~r\textsuperscript{o} (p.~\pageref{fol:174}), f\textsuperscript{o}~87~v\textsuperscript{o} (p.~\pageref{fol:175}), f\textsuperscript{o}~88~r\textsuperscript{o} (p.~\pageref{fol:176}), f\textsuperscript{o}~88~v\textsuperscript{o} (p.~\pageref{fol:177}), f\textsuperscript{o}~89~r\textsuperscript{o} (p.~\pageref{fol:178}), f\textsuperscript{o}~91~r\textsuperscript{o} (p.~\pageref{fol:182}), f\textsuperscript{o}~91~v\textsuperscript{o} (p.~\pageref{fol:183}), f\textsuperscript{o}~93~r\textsuperscript{o} (p.~\pageref{fol:186}).\par\vspace{0.55mm}
\noindent\hangindent=4.5mm\textbf{Vadencourt (Somme)}, village et château voisins de Contay : f\textsuperscript{o}~53~v\textsuperscript{o} (p.~\pageref{fol:107}), f\textsuperscript{o}~93~r\textsuperscript{o} (p.~\pageref{fol:186}).\par\vspace{0.55mm}
\noindent\hangindent=4.5mm\textbf{Varennes-en-Croix (Somme)}, commune voisine d'Acheux-en-Amiénois (famille Mortreux) : f\textsuperscript{o}~25~r\textsuperscript{o} (p.~\pageref{fol:50}), f\textsuperscript{o}~48~r\textsuperscript{o} (p.~\pageref{fol:96}), f\textsuperscript{o}~48~v\textsuperscript{o} (p.~\pageref{fol:97}), f\textsuperscript{o}~49~r\textsuperscript{o} (p.~\pageref{fol:98}), f\textsuperscript{o}~60~v\textsuperscript{o} (p.~\pageref{fol:121}), f\textsuperscript{o}~64~r\textsuperscript{o} (p.~\pageref{fol:128}), f\textsuperscript{o}~88~r\textsuperscript{o} (p.~\pageref{fol:176}), f\textsuperscript{o}~92~r\textsuperscript{o} (p.~\pageref{fol:184}).\par\vspace{0.55mm}
\noindent\hangindent=4.5mm\textbf{Villers-Bocage (Somme)}, chef-lieu de canton entre Amiens et Doullens : f\textsuperscript{o}~81~r\textsuperscript{o} (p.~\pageref{fol:162}), f\textsuperscript{o}~93~r\textsuperscript{o} (p.~\pageref{fol:186}).\par\vspace{0.55mm}
\noindent\hangindent=4.5mm\textbf{Warloy-Baillon (Somme)}, bourg entre Contay et Senlis-le-Sec : f\textsuperscript{o}~37~v\textsuperscript{o} (p.~\pageref{fol:75}), f\textsuperscript{o}~38~r\textsuperscript{o} (p.~\pageref{fol:76}), f\textsuperscript{o}~48~r\textsuperscript{o} (p.~\pageref{fol:96}), f\textsuperscript{o}~84~v\textsuperscript{o} (p.~\pageref{fol:169}), f\textsuperscript{o}~85~v\textsuperscript{o} (p.~\pageref{fol:171}), f\textsuperscript{o}~88~r\textsuperscript{o} (p.~\pageref{fol:176}), f\textsuperscript{o}~89~v\textsuperscript{o} (p.~\pageref{fol:179}), f\textsuperscript{o}~90~v\textsuperscript{o} (p.~\pageref{fol:181}), f\textsuperscript{o}~93~r\textsuperscript{o} (p.~\pageref{fol:186}).\par\vspace{0.55mm}
\noindent\hangindent=4.5mm\textbf{Ypres (Belgique)}, ville martyre de Flandre occidentale (cartes postales acquises à Renescure) : f\textsuperscript{o}~2~r\textsuperscript{o} (p.~\pageref{fol:4}), f\textsuperscript{o}~4~v\textsuperscript{o} (p.~\pageref{fol:9}).\par\vspace{0.55mm}
\par}

\end{document}
